% Auto-generated from data/segments.json + layout.json (+ transforms) — do not edit by hand.
% volume: 30Abhi02
% mode: reading
% segments: 3307
% transform_rules: 0
% toc_marks: 444

% Volume layout defaults (layout.json ``layout``).
\csromanlayoutapply{1}{1.5}{21.6pt}{5pt}{6.3pt}{65pt}{2.5em}%

% Reading physical-page layout (layout.json ``page_layout_reading_mode``).
\csromanreadingpagelayoutvolume{1}{1.5}{21.6pt}{5pt}{6.3pt}{65pt}{2.5em}%
\csromanreadingpagelayoutenable

\pitaka{\textbf{อภิธมฺมปิฏก}}
\csromanmatikaanchor{mtk.0}
\csromantocmarknopage{chapter}{วิภงฺคปาฬิ}{mtk.0}
\bookmark[dest={mtk.0},level=0]{วิภงฺคปาฬิ}
\gambhira{\textbf{วิภงฺคปาฬิ}}
\namakkaram{นโม ตสฺส ภควโต อรหโต สมฺมาสมฺพุทฺธสฺส.}
\csromanheader{2}{1. สุตฺตนฺตภาชนีย}
\proseitem{1}{\csromanfolio{1}ปญฺจกฺขนฺธา รูปกฺขนฺโธ เวทนากฺขนฺโธ สญฺญากฺขนฺโธ สงฺขารกฺขนฺโธ วิญฺญาณกฺขนฺโธ.}

\csromanmatikaanchor{mtk.1}
\csromanmatikaanchor{mtk.2}
\csromanmatikaanchor{mtk.3}
\csromantocmarknopage{chapter}{1. ขนฺธวิภงฺค}{mtk.1}
\bookmark[dest={mtk.1},level=0]{1. ขนฺธวิภงฺค}
\csromantocmarknopage{section}{1. สุตฺตนฺตภาชนีย}{mtk.2}
\bookmark[dest={mtk.2},level=1]{1. สุตฺตนฺตภาชนีย}
\csromantocmarkref{subsection}{1. รูปกฺขนฺธ}{mtk.3}
\bookmark[dest={mtk.3},level=2]{1. รูปกฺขนฺธ}
\csromanheader{2}{1. รูปกฺขนฺธ}
\proseitem{2}{ตตฺถ กตโม รูปกฺขนฺโธ, ยํ กิญฺจิ รูปํ อตีตานาคตปจฺจุปฺปนฺนํ อชฺฌตฺตํ วา พหิทฺธา วา โอฬาริกํ วา สุขุมํ วา หีนํ วา ปณีตํ วา ยํ ทูเร สนฺติเก วา, ตเทกชฺฌํ อภิสญฺญูหิตฺวา อภิสงฺขิปิตฺวา อยํ วุจฺจติ รูปกฺขนฺโธ.}

\proseitem{3}{ตตฺถ กตมํ รูปํ อตีตํ, ยํ รูปํ อตีตํ นิรุทฺธํ วิคตํ วิปริณตํ อตฺถงฺคตํ อพฺภตฺถงฺคตํ อุปฺปชฺชิตฺวา วิคตํ อตีตํ อตีตํเสน สงฺคหิตํ, จตฺตาโร จ มหาภูตา จตุนฺนญฺจ มหาภูตานํ อุปาทายรูปํ, อิทํ วุจฺจติ รูปํ อตีตํ.}

\prose{ตตฺถ กตมํ รูปํ อนาคตํ, ยํ รูปํ อชาตํ อภูตํ อสญฺชาตํ อนิพฺพตฺตํ อนภินิพฺพตฺตํ อปาตุภูตํ อนุปฺปนฺนํ อสมุปฺปนฺนํ อนุฏฺฐิตํ อสมุฏฺฐิตํ อนาคตํ อนาคตํเสน สงฺคหิตํ, จตฺตาโร จ มหาภูตา จตุนฺนญฺจ มหาภูตานํ อุปาทายรูปํ, อิทํ วุจฺจติ รูปํ อนาคตํ.}

\prose{\csromanfolio{2}ตตฺถ กตมํ รูปํ ปจฺจุปฺปนฺนํ, ยํ รูปํ ชาตํ ภูตํ สญฺชาตํ นิพฺพตฺตํ อภินิพฺพตฺตํ ปาตุภูตํ อุปฺปนฺนํ สมุปฺปนฺนํ อุฏฺฐิตํ สมุฏฺฐิตํ ปจฺจุปฺปนฺนํ ปจฺจุปฺปนฺนํเสน สงฺคหิตํ, จตฺตาโร จ มหาภูตา จตุนฺนญฺจ มหาภูตานํ อุปาทายรูปํ, อิทํ วุจฺจติ รูปํ ปจฺจุปฺปนฺนํ.}

\proseitem{4}{ตตฺถ กตมํ รูปํ อชฺฌตฺตํ, ยํ รูปํ เตสํ เตสํ สตฺตานํ อชฺฌตฺตํ ปจฺจตฺตํ นิยกํ ปาฏิปุคฺคลิกํ อุปาทินฺนํ, จตฺตาโร จ มหาภูตา จตุนฺนญฺจ มหาภูตานํ อุปาทายรูปํ, อิทํ วุจฺจติ รูปํ อชฺฌตฺตํ.}

\prose{ตตฺถ กตมํ รูปํ พหิทฺธา, ยํ รูปํ เตสํ เตสํ ปรสตฺตานํ ปรปุคฺคลานํ อชฺฌตฺตํ ปจฺจตฺตํ นิยกํ ปาฏิปุคฺคลิกํ อุปาทินฺนํ, จตฺตาโร จ มหาภูตา จตุนฺนญฺจ มหาภูตานํ อุปาทายรูปํ, อิทํ วุจฺจติ รูปํ พหิทฺธา.}

\proseitem{5}{ตตฺถ กตมํ รูปํ โอฬาริกํ, จกฺขายตนํ ฯเปฯ โผฏฺฐพฺพายตนํ, อิทํ วุจฺจติ รูปํ โอฬาริกํ.}

\prose{ตตฺถ กตมํ รูปํ สุขุมํ, อิตฺถินฺทฺริยํ ฯเปฯ กพฬีกาโร\footnote{กพฬิํกาโร (สี, สฺยา)} อาหาโร, อิทํ วุจฺจติ รูปํ สุขุมํ.}

\proseitem{6}{ตตฺถ กตมํ รูปํ หีนํ, ยํ รูปํ เตสํ เตสํ สตฺตานํ อุญฺญาตํ อวญฺญาตํ หีฬิตํ ปริภูตํ อจิตฺตีกตํ หีนํ หีนมตํ หีนสมฺมตํ อนิฏฺฐํ อกนฺตํ อมนาปํ, รูปา สทฺทา คนฺธา รสา โผฏฺฐพฺพา, อิทํ วุจฺจติ รูปํ หีนํ.}

\prose{ตตฺถ กตมํ รูปํ ปณีตํ, ยํ รูปํ เตสํ เตสํ สตฺตานํ อนุญฺญาตํ อนวญฺญาตํ อหีฬิตํ อปริภูตํ จิตฺตีกตํ ปณีตํ ปณีตมตํ ปณีตสมฺมตํ อิฏฺฐํ กนฺตํ มนาปํ, รูปา สทฺทา คนฺธา รสา โผฏฺฐพฺพา, อิทํ วุจฺจติ รูปํ ปณีตํ.\csromanspacer{} ตํ ตํ วา ปน รูปํ อุปาทายุปาทาย รูปํ หีนํ ปณีตํ ทฏฺฐพฺพํ.}

\proseitem{7}{ตตฺถ กตมํ รูปํ ทูเร, อิตฺถินฺทฺริยํ ฯเปฯ กพฬีกาโร อาหาโร, ยํ วา ปนญฺญมฺปิ อตฺถิ รูปํ อนาสนฺเน อนุปกฏฺเฐ ทูเร อสนฺติเก, อิทํ วุจฺจติ รูปํ ทูเร.}

\prose{\csromanfolio{3}ตตฺถ กตมํ รูปํ สนฺติเก, จกฺขายตนํ ฯเปฯ โผฏฺฐพฺพายตนํ, ยํ วา ปนญฺญมฺปิ อตฺถิ รูปํ อาสนฺเน อุปกฏฺเฐ อวิทูเร สนฺติเก, อิทํ วุจฺจติ รูปํ สนฺติเก.\csromanspacer{} ตํ ตํ วา ปน รูปํ อุปาทายุปาทาย รูปํ ทูเร สนฺติเก ทฏฺฐพฺพํ.}

\csromanmatikaanchor{mtk.4}
\csromantocmarkref{subsection}{2. เวทนากฺขนฺธ}{mtk.4}
\bookmark[dest={mtk.4},level=2]{2. เวทนากฺขนฺธ}
\csromanheader{2}{2. เวทนากฺขนฺธ}
\proseitem{8}{ตตฺถ กตโม เวทนากฺขนฺโธ, ยา กาจิ เวทนา อตีตานาคตปจฺจุปฺปนฺนา อชฺฌตฺตา วา พหิทฺธา วา โอฬาริกา วา สุขุมา วา หีนา วา ปณีตา วา ยา ทูเร สนฺติเก วา, ตเทกชฺฌํ อภิสญฺญูหิตฺวา อภิสงฺขิปิตฺวา อยํ วุจฺจติ เวทนากฺขนฺโธ.}

\proseitem{9}{ตตฺถ กตมา เวทนา อตีตา, ยา เวทนา อตีตา นิรุทฺธา วิคตา วิปริณตา อตฺถงฺคตา อพฺภตฺถงฺคตา อุปฺปชฺชิตฺวา วิคตา อตีตา อตีตํเสน สงฺคหิตา สุขา เวทนา ทุกฺขา เวทนา อทุกฺขมสุขา เวทนา, อยํ วุจฺจติ เวทนา อตีตา.}

\prose{ตตฺถ กตมา เวทนา อนาคตา, ยา เวทนา อชาตา อภูตา อสญฺชาตา อนิพฺพตฺตา อนภินิพฺพตฺตา อปาตุภูตา อนุปฺปนฺนา อสมุปฺปนฺนา อนุฏฺฐิตา อสมุฏฺฐิตา อนาคตา อนาคตํเสน สงฺคหิตา สุขา เวทนา ทุกฺขา เวทนา อทุกฺขมสุขา เวทนา, อยํ วุจฺจติ เวทนา อนาคตา.}

\prose{ตตฺถ กตมา เวทนา ปจฺจุปฺปนฺนา, ยา เวทนา ชาตา ภูตา สญฺชาตา นิพฺพตฺตา อภินิพฺพตฺตา ปาตุภูตา อุปฺปนฺนา สมุปฺปนฺนา อุฏฺฐิตา สมุฏฺฐิตา ปจฺจุปฺปนฺนา ปจฺจุปฺปนฺนํเสน สงฺคหิตา สุขา เวทนา ทุกฺขา เวทนา อทุกฺขมสุขา เวทนา, อยํ วุจฺจติ เวทนา ปจฺจุปฺปนฺนา.}

\proseitem{10}{ตตฺถ กตมา เวทนา อชฺฌตฺตา, ยา เวทนา เตสํ เตสํ สตฺตานํ อชฺฌตฺตํ ปจฺจตฺตํ นิยกา ปาฏิปุคฺคลิกา อุปาทินฺนา สุขา เวทนา ทุกฺขา เวทนา อทุกฺขมสุขา เวทนา, อยํ วุจฺจติ เวทนา อชฺฌตฺตา.}

\prose{ตตฺถ กตมา เวทนา พหิทฺธา, ยา เวทนา เตสํ เตสํ ปรสตฺตานํ ปรปุคฺคลานํ อชฺฌตฺตํ ปจฺจตฺตํ นิยกา ปาฏิปุคฺคลิกา อุปาทินฺนา สุขา เวทนา ทุกฺขา เวทนา อทุกฺขมสุขา เวทนา, อยํ วุจฺจติ เวทนา พหิทฺธา.}

\proseitem{11}{\csromanfolio{4}ตตฺถ กตมา เวทนา โอฬาริกา สุขุมา, อกุสลา เวทนา โอฬาริกา, กุสลาพฺยากตา เวทนา สุขุมา.\csromanspacer{} กุสลากุสลา เวทนา โอฬาริกา, อพฺยากตา เวทนา สุขุมา.\csromanspacer{} ทุกฺขา เวทนา โอฬาริกา, สุขา จ อทุกฺขมสุขา จ เวทนา สุขุมา.\csromanspacer{} สุขทุกฺขา เวทนา โอฬาริกา, อทุกฺขมสุขา เวทนา สุขุมา.\csromanspacer{} อสมาปนฺนสฺส เวทนา โอฬาริกา, สมาปนฺนสฺส เวทนา สุขุมา.\csromanspacer{} สาสวา เวทนา โอฬาริกา, อนาสวา เวทนา สุขุมา.\csromanspacer{} ตํ ตํ วา ปน เวทนํ อุปาทายุปาทาย เวทนา โอฬาริกา สุขุมา ทฏฺฐพฺพา.}

\proseitem{12}{ตตฺถ กตมา เวทนา หีนา ปณีตา, อกุสลา เวทนา หีนา, กุสลาพฺยากตา เวทนา ปณีตา.\csromanspacer{} กุสลากุสลา เวทนา หีนา, อพฺยากตา เวทนา ปณีตา.\csromanspacer{} ทุกฺขา เวทนา หีนา, สุขา จ อทุกฺขมสุขา จ เวทนา ปณีตา.\csromanspacer{} สุขทุกฺขา เวทนา หีนา, อทุกฺขมสุขา เวทนา ปณีตา.\csromanspacer{} อสมาปนฺนสฺส เวทนา หีนา, สมาปนฺนสฺส เวทนา ปณีตา.\csromanspacer{} สาสวา เวทนา หีนา, อนาสวา เวทนา ปณีตา.\csromanspacer{} ตํ ตํ วา ปน เวทนํ อุปาทายุปาทาย เวทนา หีนา ปณีตา ทฏฺฐพฺพา.}

\proseitem{13}{ตตฺถ กตมา เวทนา ทูเร, อกุสลา เวทนา กุสลาพฺยากตาหิ เวทนาหิ ทูเร, กุสลาพฺยากตา เวทนา อกุสลาย เวทนาย ทูเร, กุสลา เวทนา อกุสลาพฺยากตาหิ เวทนาหิ ทูเร, อกุสลาพฺยากตา เวทนา กุสลาย เวทนาย ทูเร, อพฺยากตา เวทนา กุสลากุสลาหิ เวทนาหิ ทูเร, กุสลากุสลา เวทนา อพฺยากตาย เวทนาย ทูเร.\csromanspacer{} ทุกฺขา เวทนา สุขาย จ อทุกฺขมสุขาย จ เวทนาหิ ทูเร, สุขา จ อทุกฺขมสุขา จ เวทนา ทุกฺขาย เวทนาย ทูเร, สุขา เวทนา ทุกฺขาย จ อทุกฺขมสุขาย จ เวทนาหิ ทูเร, ทุกฺขา จ อทุกฺขมสุขา จ เวทนา สุขาย เวทนาย ทูเร, อทุกฺขมสุขา เวทนา สุขทุกฺขาหิ เวทนาหิ ทูเร, สุขทุกฺขา เวทนา อทุกฺขมสุขาย เวทนาย ทูเร.\csromanspacer{} อสมาปนฺนสฺส เวทนา สมาปนฺนสฺส เวทนาย ทูเร, สมาปนฺนสฺส เวทนา อสมาปนฺนสฺส เวทนาย ทูเร.\csromanspacer{} สาสวา เวทนา อนาสวาย เวทนาย ทูเร, อนาสวา เวทนา สาสวาย เวทนาย ทูเร.\csromanspacer{} อยํ วุจฺจติ เวทนา ทูเร.}

\prose{ตตฺถ กตมา เวทนา สนฺติเก, อกุสลา เวทนา อกุสลาย เวทนาย สนฺติเก, กุสลา เวทนา กุสลาย เวทนาย สนฺติเก, \csromanfolio{5}อพฺยากตา เวทนา อพฺยากตาย เวทนาย สนฺติเก.\csromanspacer{} ทุกฺขา เวทนา ทุกฺขาย เวทนาย สนฺติเก, สุขา เวทนา สุขาย เวทนาย สนฺติเก, อทุกฺขมสุขา เวทนา อทุกฺขมสุขาย เวทนาย สนฺติเก.\csromanspacer{} อสมาปนฺนสฺส เวทนา อสมาปนฺนสฺส เวทนาย สนฺติเก, สมาปนฺนสฺส เวทนา สมาปนฺนสฺส เวทนาย สนฺติเก.\csromanspacer{} สาสวา เวทนา สาสวาย เวทนาย สนฺติเก, อนาสวา เวทนา อนาสวาย เวทนาย สนฺติเก.\csromanspacer{} อยํ วุจฺจติ เวทนา สนฺติเก.\csromanspacer{} ตํ ตํ วา ปน เวทนํ อุปาทายุปาทาย เวทนา ทูเร สนฺติเก ทฏฺฐพฺพา.}

\csromanmatikaanchor{mtk.5}
\csromantocmarkref{subsection}{3. สญฺญากฺขนฺธ}{mtk.5}
\bookmark[dest={mtk.5},level=2]{3. สญฺญากฺขนฺธ}
\csromanheader{2}{3. สญฺญากฺขนฺธ}
\proseitem{14}{ตตฺถ กตโม สญฺญากฺขนฺโธ, ยา กาจิ สญฺญา อตีตานาคตปจฺจุปฺปนฺนา อชฺฌตฺตา วา พหิทฺธา วา โอฬาริกา วา สุขุมา วา หีนา วา ปณีตา วา ยา ทูเร สนฺติเก วา, ตเทกชฺฌํ อภิสญฺญูหิตฺวา อภิสงฺขิปิตฺวา อยํ วุจฺจติ สญฺญากฺขนฺโธ.}

\proseitem{15}{ตตฺถ กตมา สญฺญา อตีตา, ยา สญฺญา อตีตา นิรุทฺธา วิคตา วิปริณตา อตฺถงฺคตา อพฺภตฺถงฺคตา อุปฺปชฺชิตฺวา วิคตา อตีตา อตีตํเสน สงฺคหิตา จกฺขุสมฺผสฺสชา สญฺญา โสตสมฺผสฺสชา สญฺญา ฆานสมฺผสฺสชา สญฺญา ชิวฺหาสมฺผสฺสชา สญฺญา กายสมฺผสฺสชา สญฺญา มโนสมฺผสฺสชา สญฺญา, อยํ วุจฺจติ สญฺญา อตีตา.}

\prose{ตตฺถ กตมา สญฺญา อนาคตา, ยา สญฺญา อชาตา อภูตา อสญฺชาตา อนิพฺพตฺตา อนภินิพฺพตฺตา อปาตุภูตา อนุปฺปนฺนา อสมุปฺปนฺนา อนุฏฺฐิตา อสมุฏฺฐิตา อนาคตา อนาคตํเสน สงฺคหิตา จกฺขุสมฺผสฺสชา สญฺญา โสตสมฺผสฺสชา สญฺญา ฆานสมฺผสฺสชา สญฺญา ชิวฺหาสมฺผสฺสชา สญฺญา กายสมฺผสฺสชา สญฺญา มโนสมฺผสฺสชา สญฺญา, อยํ วุจฺจติ สญฺญา อนาคตา.}

\prose{ตตฺถ กตมา สญฺญา ปจฺจุปฺปนฺนา, ยา สญฺญา ชาตา ภูตา สญฺชาตา นิพฺพตฺตา อภินิพฺพตฺตา ปาตุภูตา อุปฺปนฺนา สมุปฺปนฺนา อุฏฺฐิตา สมุฏฺฐิตา ปจฺจุปฺปนฺนา ปจฺจุปฺปนฺนํเสน สงฺคหิตา จกฺขุสมฺผสฺสชา สญฺญา โสตสมฺผสฺสชา สญฺญา ฆานสมฺผสฺสชา สญฺญา ชิวฺหาสมฺผสฺสชา สญฺญา กายสมฺผสฺสชา สญฺญา มโนสมฺผสฺสชา สญฺญา, อยํ วุจฺจติ สญฺญา ปจฺจุปฺปนฺนา.}

\proseitem{16}{\csromanfolio{6}ตตฺถ กตมา สญฺญา อชฺฌตฺตา, ยา สญฺญา เตสํ เตสํ สตฺตานํ อชฺฌตฺตํ ปจฺจตฺตํ นิยกา ปาฏิปุคฺคลิกา อุปาทินฺนา จกฺขุสมฺผสฺสชา สญฺญา โสตสมฺผสฺสชา สญฺญา ฆานสมฺผสฺสชา สญฺญา ชิวฺหาสมฺผสฺสชา สญฺญา กายสมฺผสฺสชา สญฺญา มโนสมฺผสฺสชา สญฺญา, อยํ วุจฺจติ สญฺญา อชฺฌตฺตา.}

\prose{ตตฺถ กตมา สญฺญา พหิทฺธา, ยา สญฺญา เตสํ เตสํ ปรสตฺตานํ ปรปุคฺคลานํ อชฺฌตฺตํ ปจฺจตฺตํ นิยกา ปาฏิปุคฺคลิกา อุปาทินฺนา จกฺขุสมฺผสฺสชา สญฺญา โสตสมฺผสฺสชา สญฺญา ฆานสมฺผสฺสชา สญฺญา ชิวฺหาสมฺผสฺสชา สญฺญา กายสมฺผสฺสชา สญฺญา มโนสมฺผสฺสชา สญฺญา, อยํ วุจฺจติ สญฺญา พหิทฺธา.}

\proseitem{17}{ตตฺถ กตมา สญฺญา โอฬาริกา สุขุมา, ปฏิฆสมฺผสฺสชา สญฺญา โอฬาริกา, อธิวจนสมฺผสฺสชา สญฺญา สุขุมา.\csromanspacer{} อกุสลา สญฺญา โอฬาริกา, กุสลาพฺยากตา สญฺญา สุขุมา.\csromanspacer{} กุสลากุสลา สญฺญา โอฬาริกา, อพฺยากตา สญฺญา สุขุมา.\csromanspacer{} ทุกฺขาย เวทนาย สมฺปยุตฺตา สญฺญา โอฬาริกา, สุขาย จ อทุกฺขมสุขาย จ เวทนาหิ สมฺปยุตฺตา สญฺญา สุขุมา.\csromanspacer{} สุขทุกฺขาหิ เวทนาหิ สมฺปยุตฺตา สญฺญา โอฬาริกา, อทุกฺขมสุขาย เวทนาย สมฺปยุตฺตา สญฺญา สุขุมา.\csromanspacer{} อสมาปนฺนสฺส สญฺญา โอฬาริกา, สมาปนฺนสฺส สญฺญา สุขุมา.\csromanspacer{} สาสวา สญฺญา โอฬาริกา, อนาสวา สญฺญา สุขุมา.\csromanspacer{} ตํ ตํ วา ปน สญฺญํ อุปาทายุปาทาย สญฺญา โอฬาริกา สุขุมา ทฏฺฐพฺพา.}

\proseitem{18}{ตตฺถ กตมา สญฺญา หีนา ปณีตา, อกุสลา สญฺญา หีนา, กุสลาพฺยากตา สญฺญา ปณีตา.\csromanspacer{} กุสลากุสลา สญฺญา หีนา, อพฺยากตา สญฺญา ปณีตา.\csromanspacer{} ทุกฺขาย เวทนาย สมฺปยุตฺตา สญฺญา หีนา, สุขาย จ อทุกฺขมสุขาย จ เวทนาหิ สมฺปยุตฺตา สญฺญา ปณีตา.\csromanspacer{} สุขทุกฺขาหิ เวทนาหิ สมฺปยุตฺตา สญฺญา หีนา, อทุกฺขมสุขาย เวทนาย สมฺปยุตฺตา สญฺญา ปณีตา.\csromanspacer{} อสมาปนฺนสฺส สญฺญา หีนา, สมาปนฺนสฺส สญฺญา ปณีตา.\csromanspacer{} สาสวา สญฺญา หีนา, อนาสวา สญฺญา ปณีตา.\csromanspacer{} ตํ ตํ วา ปน สญฺญํ อุปาทายุปาทาย สญฺญา หีนา ปณีตา ทฏฺฐพฺพา.}

\proseitem{19}{ตตฺถ กตมา สญฺญา ทูเร, อกุสลา สญฺญา กุสลาพฺยากตาหิ สญฺญาหิ ทูเร, กุสลาพฺยากตา สญฺญา อกุสลาย \csromanfolio{7}สญฺญาย ทูเร, กุสลา สญฺญา อกุสลาพฺยากตาหิ สญฺญาหิ ทูเร, อกุสลาพฺยากตา สญฺญา กุสลาย สญฺญาย ทูเร.\csromanspacer{} อพฺยากตา สญฺญา กุสลากุสลาหิ สญฺญาหิ ทูเร, กุสลากุสลา สญฺญา อพฺยากตาย สญฺญาย ทูเร.\csromanspacer{} ทุกฺขาย เวทนาย สมฺปยุตฺตา สญฺญา สุขาย จ อทุกฺขมสุขาย จ เวทนาหิ สมฺปยุตฺตาหิ สญฺญาหิ ทูเร, สุขาย จ อทุกฺขมสุขาย จ เวทนาหิ สมฺปยุตฺตา สญฺญา ทุกฺขาย เวทนาย สมฺปยุตฺตาย สญฺญาย ทูเร, สุขาย เวทนาย สมฺปยุตฺตา สญฺญา ทุกฺขาย จ อทุกฺขมสุขาย จ เวทนาหิ สมฺปยุตฺตาหิ สญฺญาหิ ทูเร, ทุกฺขาย จ อทุกฺขมสุขาย จ เวทนาหิ สมฺปยุตฺตา สญฺญา สุขาย เวทนาย สมฺปยุตฺตาย สญฺญาย ทูเร, อทุกฺขมสุขาย เวทนาย สมฺปยุตฺตา สญฺญา สุขทุกฺขาหิ เวทนาหิ สมฺปยุตฺตาหิ สญฺญาหิ ทูเร, สุขทุกฺขาหิ เวทนาหิ สมฺปยุตฺตา สญฺญา อทุกฺขมสุขาย เวทนาย สมฺปยุตฺตาย สญฺญาย ทูเร.\csromanspacer{} อสมาปนฺนสฺส สญฺญา สมาปนฺนสฺส สญฺญาย ทูเร, สมาปนฺนสฺส สญฺญา อสมาปนฺนสฺส สญฺญาย ทูเร.\csromanspacer{} สาสวา สญฺญา อนาสวาย สญฺญาย ทูเร, อนาสวา สญฺญา สาสวาย สญฺญาย ทูเร.\csromanspacer{} อยํ วุจฺจติ สญฺญา ทูเร.}

\prose{ตตฺถ กตมา สญฺญา สนฺติเก, อกุสลา สญฺญา อกุสลาย สญฺญาย สนฺติเก, กุสลา สญฺญา กุสลาย สญฺญาย สนฺติเก, อพฺยากตา สญฺญา อพฺยากตาย สญฺญาย สนฺติเก, ทุกฺขาย เวทนาย สมฺปยุตฺตา สญฺญา ทุกฺขาย เวทนาย สมฺปยุตฺตาย สญฺญาย สนฺติเก, สุขาย เวทนาย สมฺปยุตฺตา สญฺญา สุขาย เวทนาย สมฺปยุตฺตาย สญฺญาย สนฺติเก, อทุกฺขมสุขาย เวทนาย สมฺปยุตฺตา สญฺญา อทุกฺขมสุขาย เวทนาย สมฺปยุตฺตาย สญฺญาย สนฺติเก.\csromanspacer{} อสมาปนฺนสฺส สญฺญา อสมาปนฺนสฺส สญฺญาย สนฺติเก, สมาปนฺนสฺส สญฺญา สมาปนฺนสฺส สญฺญาย สนฺติเก.\csromanspacer{} สาสวา สญฺญา สาสวาย สญฺญาย สนฺติเก, อนาสวา สญฺญา อนาสวาย สญฺญาย สนฺติเก.\csromanspacer{} อยํ วุจฺจติ สญฺญา สนฺติเก.\csromanspacer{} ตํ ตํ วา ปน สญฺญํ อุปาทายุปาทาย สญฺญา ทูเร สนฺติเก ทฏฺฐพฺพา.}

\csromanmatikaanchor{mtk.6}
\csromantocmarkref{subsection}{4. สงฺขารกฺขนฺธ}{mtk.6}
\bookmark[dest={mtk.6},level=2]{4. สงฺขารกฺขนฺธ}
\csromanheader{2}{4. สงฺขารกฺขนฺธ}
\proseitem{20}{ตตฺถ กตโม สงฺขารกฺขนฺโธ, เย เกจิ สงฺขารา อตีตานาคตปจฺจุปฺปนฺนา อชฺฌตฺตา วา พหิทฺธา วา โอฬาริกา วา สุขุมา วา \csromanfolio{8}หีนา วา ปณีตา วา เย ทูเร สนฺติเก วา, ตเทกชฺฌํ อภิสญฺญูหิตฺวา อภิสงฺขิปิตฺวา อยํ วุจฺจติ สงฺขารกฺขนฺโธ.}

\proseitem{21}{ตตฺถ กตเม สงฺขารา อตีตา, เย สงฺขารา อตีตา นิรุทฺธา วิคตา วิปริณตา อตฺถงฺคตา อพฺภตฺถงฺคตา อุปฺปชฺชิตฺวา วิคตา อตีตา อตีตํเสน สงฺคหิตา จกฺขุสมฺผสฺสชา เจตนา โสตสมฺผสฺสชา เจตนา ฆานสมฺผสฺสชา เจตนา ชิวฺหาสมฺผสฺสชา เจตนา กายสมฺผสฺสชา เจตนา มโนสมฺผสฺสชา เจตนา, อิเม วุจฺจนฺติ สงฺขารา อตีตา.}

\prose{ตตฺถ กตเม สงฺขารา อนาคตา, เย สงฺขารา อชาตา อภูตา อสญฺชาตา อนิพฺพตฺตา อนภินิพฺพตฺตา อปาตุภูตา อนุปฺปนฺนา อสมุปฺปนฺนา อนุฏฺฐิตา อสมุฏฺฐิตา อนาคตา อนาคตํเสน สงฺคหิตา จกฺขุสมฺผสฺสชา เจตนา โสตสมฺผสฺสชา เจตนา ฆานสมฺผสฺสชา เจตนา ชิวฺหาสมฺผสฺสชา เจตนา กายสมฺผสฺสชา เจตนา มโนสมฺผสฺสชา เจตนา, อิเม วุจฺจนฺติ สงฺขารา อนาคตา.}

\prose{ตตฺถ กตเม สงฺขารา ปจฺจุปฺปนฺนา, เย สงฺขารา ชาตา ภูตา สญฺชาตา นิพฺพตฺตา อภินิพฺพตฺตา ปาตุภูตา อุปฺปนฺนา สมุปฺปนฺนา อุฏฺฐิตา สมุฏฺฐิตา ปจฺจุปฺปนฺนา ปจฺจุปฺปนฺนํเสน สงฺคหิตา จกฺขุสมฺผสฺสชา เจตนา โสตสมฺผสฺสชา เจตนา ฆานสมฺผสฺสชา เจตนา ชิวฺหาสมฺผสฺสชา เจตนา กายสมฺผสฺสชา เจตนา มโนสมฺผสฺสชา เจตนา, อิเม วุจฺจนฺติ สงฺขารา ปจฺจุปฺปนฺนา.}

\proseitem{22}{ตตฺถ กตเม สงฺขารา อชฺฌตฺตา, เย สงฺขารา เตสํ เตสํ สตฺตานํ อชฺฌตฺตํ ปจฺจตฺตํ นิยกา ปาฏิปุคฺคลิกา อุปาทินฺนา จกฺขุสมฺผสฺสชา เจตนา โสตสมฺผสฺสชา เจตนา ฆานสมฺผสฺสชา เจตนา ชิวฺหาสมฺผสฺสชา เจตนา กายสมฺผสฺสชา เจตนา มโนสมฺผสฺสชา เจตนา อิเม วุจฺจนฺติ สงฺขารา อชฺฌตฺตา.}

\prose{ตตฺถ กตเม สงฺขารา พหิทฺธา, เย สงฺขารา เตสํ เตสํ ปรสตฺตานํ ปรปุคฺคลานํ อชฺฌตฺตํ ปจฺจตฺตํ นิยกา ปาฏิปุคฺคลิกา อุปาทินฺนา จกฺขุสมฺผสฺสชา เจตนา โสตสมฺผสฺสชา เจตนา ฆานสมฺผสฺสชา เจตนา ชิวฺหาสมฺผสฺสชา เจตนา กายสมฺผสฺสชา เจตนา มโนสมฺผสฺสชา เจตนา, อิเม วุจฺจนฺติ สงฺขารา พหิทฺธา.}

\proseitem{23}{\csromanfolio{9}ตตฺถ กตเม สงฺขารา โอฬาริกา สุขุมา, อกุสลา สงฺขารา โอฬาริกา, กุสลาพฺยากตา สงฺขารา สุขุมา.\csromanspacer{} กุสลากุสลา สงฺขารา โอฬาริกา, อพฺยากตา สงฺขารา สุขุมา.\csromanspacer{} ทุกฺขาย เวทนาย สมฺปยุตฺตา สงฺขารา โอฬาริกา, สุขาย จ อทุกฺขมสุขาย จ เวทนาหิ สมฺปยุตฺตา สงฺขารา สุขุมา.\csromanspacer{} สุขทุกฺขาหิ เวทนาหิ สมฺปยุตฺตา สงฺขารา โอฬาริกา, อทุกฺขมสุขาย เวทนาย สมฺปยุตฺตา สงฺขารา สุขุมา.\csromanspacer{} อสมาปนฺนสฺส สงฺขารา โอฬาริกา, สมาปนฺนสฺส สงฺขารา สุขุมา.\csromanspacer{} สาสวา สงฺขารา โอฬาริกา, อนาสวา สงฺขารา สุขุมา.\csromanspacer{} เต เต วา ปน สงฺขาเร อุปาทายุปาทาย สงฺขารา โอฬาริกา สุขุมา ทฏฺฐพฺพา.}

\proseitem{24}{ตตฺถ กตเม สงฺขารา หีนา ปณีตา, อกุสลา สงฺขารา หีนา, กุสลาพฺยากตา สงฺขารา ปณีตา.\csromanspacer{} กุสลากุสลา สงฺขารา หีนา, อพฺยากตา สงฺขารา ปณีตา.\csromanspacer{} ทุกฺขาย เวทนาย สมฺปยุตฺตา สงฺขารา หีนา, สุขาย จ อทุกฺขมสุขาย จ เวทนาหิ สมฺปยุตฺตา สงฺขารา ปณีตา.\csromanspacer{} สุขทุกฺขาหิ เวทนาหิ สมฺปยุตฺตา สงฺขารา หีนา, อทุกฺขมสุขาย เวทนาย สมฺปยุตฺตา สงฺขารา ปณีตา.\csromanspacer{} อสมาปนฺนสฺส สงฺขารา หีนา, สมาปนฺนสฺส สงฺขารา ปณีตา.\csromanspacer{} สาสวา สงฺขารา หีนา, อนาสวา สงฺขารา ปณีตา.\csromanspacer{} เต เต วา ปน สงฺขาเร อุปาทายุปาทาย สงฺขารา หีนา ปณีตา ทฏฺฐพฺพา.}

\proseitem{25}{ตตฺถ กตเม สงฺขารา ทูเร, อกุสลา สงฺขารา กุสลาพฺยากเตหิ สงฺขาเรหิ ทูเร, กุสลาพฺยากตา สงฺขารา อกุสเลหิ สงฺขาเรหิ ทูเร, กุสลา สงฺขารา อกุสลาพฺยากเตหิ สงฺขาเรหิ ทูเร, อกุสลาพฺยากตา สงฺขารา กุสเลหิ สงฺขาเรหิ ทูเร, อพฺยากตา สงฺขารา กุสลากุสเลหิ สงฺขาเรหิ ทูเร, กุสลากุสลา สงฺขารา อพฺยากเตหิ สงฺขาเรหิ ทูเร.\csromanspacer{} ทุกฺขาย เวทนาย สมฺปยุตฺตา สงฺขารา สุขาย จ อทุกฺขมสุขาย จ เวทนาหิ สมฺปยุตฺเตหิ สงฺขาเรหิ ทูเร, สุขาย จ อทุกฺขมสุขาย จ เวทนาหิ สมฺปยุตฺตา สงฺขารา ทุกฺขาย เวทนาย สมฺปยุตฺเตหิ สงฺขาเรหิ ทูเร, สุขาย เวทนาย สมฺปยุตฺตา สงฺขารา ทุกฺขาย จ อทุกฺขมสุขาย จ เวทนาหิ สมฺปยุตฺเตหิ สงฺขาเรหิ ทูเร, ทุกฺขาย จ อทุกฺขมสุขาย จ เวทนาหิ สมฺปยุตฺตา สงฺขารา สุขาย เวทนาย สมฺปยุตฺเตหิ สงฺขาเรหิ ทูเร, \csromanfolio{10}อทุกฺขมสุขาย เวทนาย สมฺปยุตฺตา สงฺขารา สุขทุกฺขาหิ เวทนาหิ สมฺปยุตฺเตหิ สงฺขาเรหิ ทูเร, สุขทุกฺขาหิ เวทนาหิ สมฺปยุตฺตา สงฺขารา อทุกฺขมสุขาย เวทนาย สมฺปยุตฺเตหิ สงฺขาเรหิ ทูเร.\csromanspacer{} อสมาปนฺนสฺส สงฺขารา สมาปนฺนสฺส สงฺขาเรหิ ทูเร, สมาปนฺนสฺส สงฺขารา อสมาปนฺนสฺส สงฺขาเรหิ ทูเร.\csromanspacer{} สาสวา สงฺขารา อนาสเวหิ สงฺขาเรหิ ทูเร, อนาสวา สงฺขารา สาสเวหิ สงฺขาเรหิ ทูเร.\csromanspacer{} อิเม วุจฺจนฺติ สงฺขารา ทูเร.}

\prose{ตตฺถ กตเม สงฺขารา สนฺติเก, อกุสลา สงฺขารา อกุสลานํ สงฺขารานํ สนฺติเก, กุสลา สงฺขารา กุสลานํ สงฺขารานํ สนฺติเก, อพฺยากตา สงฺขารา อพฺยากตานํ สงฺขารานํ สนฺติเก.\csromanspacer{} ทุกฺขาย เวทนาย สมฺปยุตฺตา สงฺขารา ทุกฺขาย เวทนาย สมฺปยุตฺตานํ สงฺขารานํ สนฺติเก, สุขาย เวทนาย สมฺปยุตฺตา สงฺขารา สุขาย เวทนาย สมฺปยุตฺตานํ สงฺขารานํ สนฺติเก, อทุกฺขมสุขาย เวทนาย สมฺปยุตฺตา สงฺขารา อทุกฺขมสุขาย เวทนาย สมฺปยุตฺตานํ สงฺขารานํ สนฺติเก.\csromanspacer{} อสมาปนฺนสฺส สงฺขารา อสมาปนฺนสฺส สงฺขารานํ สนฺติเก, สมาปนฺนสฺส สงฺขารา สมาปนฺนสฺส สงฺขารานํ สนฺติเก.\csromanspacer{} สาสวา สงฺขารา สาสวานํ สงฺขารานํ สนฺติเก, อนาสวา สงฺขารา อนาสวานํ สงฺขารานํ สนฺติเก.\csromanspacer{} อิเม วุจฺจนฺติ สงฺขารา สนฺติเก.\csromanspacer{} เต เต วา ปน สงฺขาเร อุปาทายุปาทาย สงฺขารา ทูเร สนฺติเก ทฏฺฐพฺพา.}

\csromanmatikaanchor{mtk.7}
\csromantocmarkref{subsection}{5. วิญฺญาณกฺขนฺธ}{mtk.7}
\bookmark[dest={mtk.7},level=2]{5. วิญฺญาณกฺขนฺธ}
\csromanheader{2}{5. วิญฺญาณกฺขนฺธ}
\proseitem{26}{ตตฺถ กตโม วิญฺญาณกฺขนฺโธ, ยํ กิญฺจิ วิญฺญาณํ อตีตานาคตปจฺจุปฺปนฺนํ อชฺฌตฺตํ วา พหิทฺธา วา โอฬาริกํ วา สุขุมํ วา หีนํ วา ปณีตํ วา ยํ ทูเร สนฺติเก วา, ตเทกชฺฌํ อภิสญฺญูหิตฺวา อภิสงฺขิปิตฺวา อยํ วุจฺจติ วิญฺญาณกฺขนฺโธ.}

\proseitem{27}{ตตฺถ กตมํ วิญฺญาณํ อตีตํ, ยํ วิญฺญาณํ อตีตํ นิรุทฺธํ วิคตํ วิปริณตํ อตฺถงฺคตํ อพฺภตฺถงฺคตํ อุปฺปชฺชิตฺวา วิคตํ อตีตํ อตีตํเสน สงฺคหิตํ จกฺขุวิญฺญาณํ โสตวิญฺญาณํ ฆานวิญฺญาณํ ชิวฺหาวิญฺญาณํ กายวิญฺญาณํ มโนวิญฺญาณํ, อิทํ วุจฺจติ วิญฺญาณํ อตีตํ.}

\prose{ตตฺถ กตมํ วิญฺญาณํ อนาคตํ, ยํ วิญฺญาณํ อชาตํ อภูตํ อสญฺชาตํ อนิพฺพตฺตํ อนภินิพฺพตฺตํ อปาตุภูตํ อนุปฺปนฺนํ อสมุปฺปนฺนํ อนุฏฺฐิตํ \csromanfolio{11}อสมุฏฺฐิตํ อนาคตํ อนาคตํเสน สงฺคหิตํ จกฺขุวิญฺญาณํ โสตวิญฺญาณํ ฆานวิญฺญาณํ ชิวฺหาวิญฺญาณํ กายวิญฺญาณํ มโนวิญฺญาณํ, อิทํ วุจฺจติ วิญฺญาณํ อนาคตํ.}

\prose{ตตฺถ กตมํ วิญฺญาณํ ปจฺจุปฺปนฺนํ, ยํ วิญฺญาณํ ชาตํ ภูตํ สญฺชาตํ นิพฺพตฺตํ อภินิพฺพตฺตํ ปาตุภูตํ อุปฺปนฺนํ สมุปฺปนฺนํ อุฏฺฐิตํ สมุฏฺฐิตํ ปจฺจุปฺปนฺนํ ปจฺจุปฺปนฺนํเสน สงฺคหิตํ จกฺขุวิญฺญาณํ โสตวิญฺญาณํ ฆานวิญฺญาณํ ชิวฺหาวิญฺญาณํ กายวิญฺญาณํ มโนวิญฺญาณํ, อิทํ วุจฺจติ วิญฺญาณํ ปจฺจุปฺปนฺนํ.}

\proseitem{28}{ตตฺถ กตมํ วิญฺญาณํ อชฺฌตฺตํ, ยํ วิญฺญาณํ เตสํ เตสํ สตฺตานํ อชฺฌตฺตํ ปจฺจตฺตํ นิยกํ ปาฏิปุคฺคลิกํ อุปาทินฺนํ จกฺขุวิญฺญาณํ โสตวิญฺญาณํ ฆานวิญฺญาณํ ชิวฺหาวิญฺญาณํ กายวิญฺญาณํ มโนวิญฺญาณํ, อิทํ วุจฺจติ วิญฺญาณํ อชฺฌตฺตํ.}

\prose{ตตฺถ กตมํ วิญฺญาณํ พหิทฺธา, ยํ วิญฺญาณํ เตสํ เตสํ ปรสตฺตานํ ปรปุคฺคลานํ อชฺฌตฺตํ ปจฺจตฺตํ นิยกํ ปาฏิคฺคลิกํ อุปาทินฺนํ จกฺขุวิญฺญาณํ โสตวิญฺญาณํ ฆานวิญฺญาณํ ชิวฺหาวิญฺญาณํ กายวิญฺญาณํ มโนวิญฺญาณํ, อิทํ วุจฺจติ วิญฺญาณํ พหิทฺธา.}

\proseitem{29}{ตตฺถ กตมํ วิญฺญาณํ โอฬาริกํ สุขุมํ, อกุสลํ วิญฺญาณํ โอฬาริกํ, กุสลาพฺยากตา วิญฺญาณา สุขุมา.\csromanspacer{} กุสลากุสลา วิญฺญาณา โอฬาริกา, อพฺยากตํ วิญฺญาณํ สุขุมํ.\csromanspacer{} ทุกฺขาย เวทนาย สมฺปยุตฺตํ วิญฺญาณํ โอฬาริกํ, สุขาย จ อทุกฺขมสุขาย จ เวทนาหิ สมฺปยุตฺตา วิญฺญาณา สุขุมา.\csromanspacer{} สุขทุกฺขาหิ เวทนาหิ สมฺปยุตฺตา วิญฺญาณา โอฬาริกา, อทุกฺขมสุขาย เวทนาย สมฺปยุตฺตํ วิญฺญาณํ สุขุมํ.\csromanspacer{} อสมาปนฺนสฺส วิญฺญาณํ โอฬาริกํ, สมาปนฺนสฺส วิญฺญาณํ สุขุมํ.\csromanspacer{} สาสวํ วิญฺญาณํ โอฬาริกํ, อนาสวํ วิญฺญาณํ สุขุมํ.\csromanspacer{} ตํ ตํ วา ปน วิญฺญาณํ อุปาทายุปาทาย วิญฺญาณํ โอฬาริกํ สุขุมํ ทฏฺฐพฺพํ.}

\proseitem{30}{ตตฺถ กตมํ วิญฺญาณํ หีนํ ปณีตํ, อกุสลํ วิญฺญาณํ หีนํ, กุสลาพฺยากตา วิญฺญาณา ปณีตา.\csromanspacer{} กุสลากุสลา วิญฺญาณา หีนา, อพฺยากตํ วิญฺญาณํ ปณีตํ.\csromanspacer{} ทุกฺขาย เวทนาย สมฺปยุตฺตํ วิญฺญาณํ หีนํ, สุขาย จ อทุกฺขมสุขาย จ เวทนาหิ สมฺปยุตฺตา วิญฺญาณา ปณีตา.\csromanspacer{} สุขทุกฺขาหิ เวทนาหิ สมฺปยุตฺตา วิญฺญาณา หีนา, อทุกฺขมสุขาย \csromanfolio{12}เวทนาย สมฺปยุตฺตํ วิญฺญาณํ ปณีตํ.\csromanspacer{} อสมาปนฺนสฺส วิญฺญาณํ หีนํ, สมาปนฺนสฺส วิญฺญาณํ ปณีตํ.\csromanspacer{} สาสวํ วิญฺญาณํ หีนํ, อนาสวํ วิญฺญาณํ ปณีตํ.\csromanspacer{} ตํ ตํ วา ปน วิญฺญาณํ อุปาทายุปาทาย วิญฺญาณํ หีนํ ปณีตํ ทฏฺฐพฺพํ.}

\proseitem{31}{ตตฺถ กตมํ วิญฺญาณํ ทูเร, อกุสลํ วิญฺญาณํ กุสลาพฺยากเตหิ วิญฺญาเณหิ ทูเร, กุสลาพฺยากตา วิญฺญาณา อกุสลา วิญฺญาณา ทูเร, กุสลํ วิญฺญาณํ อกุสลาพฺยากเตหิ วิญฺญาเณหิ ทูเร, อกุสลาพฺยากตา วิญฺญาณา กุสลา วิญฺญาณา ทูเร, อพฺยากตํ วิญฺญาณํ กุสลากุสเลหิ วิญฺญาเณหิ ทูเร, กุสลากุสลา วิญฺญาณา อพฺยากตา วิญฺญาณา ทูเร.\csromanspacer{} ทุกฺขาย เวทนาย สมฺปยุตฺตํ วิญฺญาณํ สุขาย จ อทุกฺขมสุขาย จ เวทนาหิ สมฺปยุตฺเตหิ วิญฺญาเณหิ ทูเร, สุขาย จ อทุกฺขมสุขาย จ เวทนาหิ สมฺปยุตฺตา วิญฺญาณา ทุกฺขาย เวทนาย สมฺปยุตฺตา วิญฺญาณา ทูเร, สุขาย เวทนาย สมฺปยุตฺตํ วิญฺญาณํ ทุกฺขาย จ อทุกฺขมสุขาย จ เวทนาหิ สมฺปยุตฺเตหิ วิญฺญาเณหิ ทูเร, ทุกฺขาย จ อทุกฺขมสุขาย จ เวทนาหิ สมฺปยุตฺตา วิญฺญาณา สุขาย เวทนาย สมฺปยุตฺตา วิญฺญาณา ทูเร.\csromanspacer{} อทุกฺขมสุขาย เวทนาย สมฺปยุตฺตํ วิญฺญาณํ สุขทุกฺขาหิ เวทนาหิ สมฺปยุตฺเตหิ วิญฺญาเณหิ ทูเร, สุขทุกฺขาหิ เวทนาหิ สมฺปยุตฺตา วิญฺญาณา อทุกฺขมสุขาย เวทนาย สมฺปยุตฺตา วิญฺญาณา ทูเร.\csromanspacer{} อสมาปนฺนสฺส วิญฺญาณํ สมาปนฺนสฺส วิญฺญาณา ทูเร, สมาปนฺนสฺส วิญฺญาณํ อสมาปนฺนสฺส วิญฺญาณา ทูเร.\csromanspacer{} สาสวํ วิญฺญาณํ อนาสวา วิญฺญาณา ทูเร, อนาสวํ วิญฺญาณํ สาสวา วิญฺญาณา ทูเร.\csromanspacer{} อิทํ วุจฺจติ วิญฺญาณํ ทูเร.}

\prose{ตตฺถ กตมํ วิญฺญาณํ สนฺติเก, อกุสลํ วิญฺญาณํ อกุสลสฺส วิญฺญาณสฺส สนฺติเก, กุสลํ วิญฺญาณํ กุสลสฺส วิญฺญาณสฺส สนฺติเก, อพฺยากตํ วิญฺญาณํ อพฺยากตสฺส วิญฺญาณสฺส สนฺติเก.\csromanspacer{} ทุกฺขาย เวทนาย สมฺปยุตฺตํ วิญฺญาณํ ทุกฺขาย เวทนาย สมฺปยุตฺตสฺส วิญฺญาณสฺส สนฺติเก, สุขาย เวทนาย สมฺปยุตฺตํ วิญฺญาณํ สุขาย เวทนาย สมฺปยุตฺตสฺส วิญฺญาณสฺส สนฺติเก, อทุกฺขมสุขาย เวทนาย สมฺปยุตฺตํ วิญฺญาณํ อทุกฺขมสุขาย เวทนาย สมฺปยุตฺตสฺส วิญฺญาณสฺส สนฺติเก.\csromanspacer{} อสมาปนฺนสฺส วิญฺญาณํ อสมาปนฺนสฺส วิญฺญาณสฺส สนฺติเก, สมาปนฺนสฺส วิญฺญาณํ สมาปนฺนสฺส วิญฺญาณสฺส สนฺติเก.\csromanspacer{} สาสวํ วิญฺญาณํ สาสวสฺส \csromanfolio{13}วิญฺญาณสฺส สนฺติเก, อนาสวํ วิญฺญาณํ อนาสวสฺส วิญฺญาณสฺส สนฺติเก.\csromanspacer{} อิทํ วุจฺจติ วิญฺญาณํ สนฺติเก.\csromanspacer{} ตํ ตํ วา ปน วิญฺญาณํ อุปาทายุปาทาย วิญฺญาณํ ทูเร สนฺติเก ทฏฺฐพฺพํ.}

\csromanheader{2}{สุตฺตนฺตภาชนียํ.}
\csromansectionrule
\csromanheader{2}{2. อภิธมฺมภาชนีย}
\proseitem{32}{ปญฺจกฺขนฺธา รูปกฺขนฺโธ เวทนากฺขนฺโธ สญฺญากฺขนฺโธ สงฺขารกฺขนฺโธ วิญฺญาณกฺขนฺโธ.}

\csromanmatikaanchor{mtk.8}
\csromanmatikaanchor{mtk.9}
\csromantocmarknopage{section}{2. อภิธมฺมภาชนีย}{mtk.8}
\bookmark[dest={mtk.8},level=1]{2. อภิธมฺมภาชนีย}
\csromantocmarkref{subsection}{1. รูปกฺขนฺธ}{mtk.9}
\bookmark[dest={mtk.9},level=2]{1. รูปกฺขนฺธ}
\csromanheader{2}{1. รูปกฺขนฺธ}
\proseitem{33}{ตตฺถ กตโม รูปกฺขนฺโธ, เอกวิเธน รูปกฺขนฺโธ\csromandash{}สพฺพํ รูปํ น เหตุ, อเหตุกํ, เหตุวิปฺปยุตฺตํ, สปฺปจฺจยํ, สงฺขตํ, รูปํ, โลกิยํ สาสวํ, สํโยชนิยํ, คนฺถนิยํ, โอฆนิยํ, โยคนิยํ, นีวรณิยํ, ปรามฏฺฐํ, อุปาทานิยํ, สํกิเลสิกํ, อพฺยากตํ, อนารมฺมณํ, อเจตสิกํ, จิตฺตวิปฺปยุตฺตํ, เนววิปากนวิปากธมฺมธมฺมํ, อสํกิลิฏฺฐสํกิเลสิกํ, น สวิตกฺกสวิจารํ, น อวิตกฺกวิจารมตฺตํ, อวิตกฺก-อวิจารํ, น ปีติสหคตํ, น สุขสหคตํ, น อุเปกฺขาสหคตํ, เนว ทสฺสเนน น ภาวนาย ปหาตพฺพํ, เนว ทสฺสเนน น ภาวนาย ปหาตพฺพเหตุกํ, เนวาจยคามินาปจยคามี, เนวเสกฺขนาเสกฺขํ, ปริตฺตํ, กามาวจรํ, น รูปาวจรํ, น อรูปาวจรํ, ปริยาปนฺนํ, โน อปริยาปนฺนํ, อนิยตํ, อนิยฺยานิกํ, อุปฺปนฺนํ, ฉหิ วิญฺญาเณหิ วิญฺเญยฺยํ, อนิจฺจํ, ชราภิภูตํ.\csromanspacer{} เอวํ เอกวิเธน รูปกฺขนฺโธ.}

\prose{ทุวิเธน รูปกฺขนฺโธ\csromandash{}อตฺถิ รูปํ อุปาทา, อตฺถิ รูปํ โน อุปาทา\footnote{นุปาทา (สี, ก)}.\csromanspacer{} อตฺถิ รูปํ อุปาทินฺนํ, อตฺถิ รูปํ อนุปาทินฺนํ.\csromanspacer{} อตฺถิ รูปํ อุปาทินฺนุปาทานิยํ, อตฺถิ รูปํ อนุปาทินฺนุปาทานิยํ.\csromanspacer{} อตฺถิ รูปํ สนิทสฺสนํ, อตฺถิ รูปํ อนิทสฺสนํ.\csromanspacer{} อตฺถิ รูปํ สปฺปฏิฆํ, อตฺถิ รูปํ อปฺปฏิฆํ.\csromanspacer{} อตฺถิ รูปํ อินฺทฺริยํ, อตฺถิ รูปํ น อินฺทฺริยํ.\csromanspacer{} อตฺถิ รูปํ มหาภูตํ, อตฺถิ รูปํ น มหาภูตํ.\csromanspacer{} อตฺถิ รูปํ วิญฺญตฺติ, อตฺถิ รูปํ น วิญฺญตฺติ.\csromanspacer{} อตฺถิ รูปํ จิตฺตสมุฏฺฐานํ, อตฺถิ รูปํ น จิตฺตสมุฏฺฐานํ.\csromanspacer{} อตฺถิ รูปํ จิตฺตสหภุ, อตฺถิ รูปํ น จิตฺตสหภุ.\csromanspacer{} อตฺถิ รูปํ จิตฺตานุปริวตฺติ, อตฺถิ รูปํ \csromanfolio{14}น จิตฺตานุปริวตฺติ.\csromanspacer{} อตฺถิ รูปํ อชฺฌตฺติกํ, อตฺถิ รูปํ พาหิรํ.\csromanspacer{} อตฺถิ รูปํ โอฬาริกํ, อตฺถิ รูปํ สุขุมํ.\csromanspacer{} อตฺถิ รูปํ ทูเร, อตฺถิ รูปํ สนฺติเก ฯเปฯ.\csromanspacer{} อตฺถิ รูปํ กพฬีกาโร อาหาโร, อตฺถิ รูปํ น กพฬีกาโร อาหาโร.\csromanspacer{} เอวํ ทุวิเธน รูปกฺขนฺโธ. (ยถา รูปกณฺเฑ วิภตฺตํ, ตถา อิธ วิภชิตพฺพํ.)}

\prose{ติวิเธน รูปกฺขนฺโธ\csromandash{}ยํ ตํ รูปํ อชฺฌตฺติกํ ตํ อุปาทา, ยํ ตํ รูปํ พาหิรํ ตํ อตฺถิ อุปาทา, อตฺถิ โน อุปาทา.\csromanspacer{} ยํ ตํ รูปํ อชฺฌตฺติกํ ตํ อุปาทินฺนํ, ยํ ตํ รูปํ พาหิรํ ตํ อตฺถิ อุปาทินฺนํ, อตฺถิ อนุปาทินฺนํ.\csromanspacer{} ยํ ตํ รูปํ อชฺฌตฺติกํ ตํ อุปาทินฺนุปาทานิยํ, ยํ ตํ รูปํ พาหิรํ ตํ อตฺถิ อุปาทินฺนุปาทานิยํ, อตฺถิ อนุปาทินฺนุปาทานิยํ ฯเปฯ.\csromanspacer{} ยํ ตํ รูปํ อชฺฌตฺติกํ ตํ น กพฬีกาโร อาหาโร, ยํ ตํ รูปํ พาหิรํ ตํ อตฺถิ กพฬีกาโร อาหาโร, อตฺถิ น กพฬีกาโร อาหาโร.\csromanspacer{} เอวํ ติวิเธน รูปกฺขนฺโธ.}

\prose{จตุพฺพิเธน รูปกฺขนฺโธ\csromandash{}ยํ ตํ รูปํ อุปาทา ตํ อตฺถิ อุปาทินฺนํ, อตฺถิ อนุปาทินฺนํ, ยํ ตํ รูปํ โน อุปาทา ตํ อตฺถิ อุปาทินฺนํ, อตฺถิ อนุปาทินฺนํ.\csromanspacer{} ยํ ตํ รูปํ อุปาทา ตํ อตฺถิ อุปาทินฺนุปาทานิยํ, อตฺถิ อนุปาทินฺนุปาทานิยํ, ยํ ตํ รูปํ โน อุปาทา ตํ อตฺถิ อุปาทินฺนุปาทานิยํ, อตฺถิ อนุปาทินฺนุปาทานิยํ.\csromanspacer{} ยํ ตํ รูปํ อุปาทา ตํ อตฺถิ สปฺปฏิฆํ, อตฺถิ อปฺปฏิฆํ, ยํ ตํ รูปํ โน อุปาทา ตํ อตฺถิ สปฺปฏิฆํ, อตฺถิ อปฺปฏิฆํ.\csromanspacer{} ยํ ตํ รูปํ อุปาทา ตํ อตฺถิ โอฬาริกํ, อตฺถิ สุขุมํ, ยํ ตํ รูปํ โน อุปาทา ตํ อตฺถิ โอฬาริกํ, อตฺถิ สุขุมํ.\csromanspacer{} ยํ ตํ รูปํ อุปาทา ตํ อตฺถิ ทูเร, อตฺถิ สนฺติเก, ยํ ตํ รูปํ โน อุปาทา ตํ อตฺถิ ทูเร, อตฺถิ สนฺติเก ฯเปฯ.\csromanspacer{} ทิฏฺฐํ สุตํ มุตํ วิญฺญาตํ รูปํ.\csromanspacer{} เอวํ จตุพฺพิเธน รูปกฺขนฺโธ.}

\prose{ปญฺจวิเธน รูปกฺขนฺโธ\csromandash{}ปถวีธาตุ, อาโปธาตุ, เตโชธาตุ, วาโยธาตุ, ยญฺจ รูปํ อุปาทา.\csromanspacer{} เอวํ ปญฺจวิเธน รูปกฺขนฺโธ.}

\prose{ฉพฺพิเธน รูปกฺขนฺโธ\csromandash{}จกฺขุวิญฺเญยฺยํ รูปํ, โสตวิญฺเญยฺยํ รูปํ, ฆานวิญฺเญยฺยํ รูปํ, ชิวฺหาวิญฺเญยฺยํ รูปํ, กายวิญฺเญยฺยํ รูปํ, มโนวิญฺเญยฺยํ รูปํ.\csromanspacer{} เอวํ ฉพฺพิเธน รูปกฺขนฺโธ.}

\prose{สตฺตวิเธน รูปกฺขนฺโธ\csromandash{}จกฺขุวิญฺเญยฺยํ รูปํ, โสตวิญฺเญยฺยํ รูปํ, ฆานวิญฺเญยฺยํ รูปํ, ชิวฺหาวิญฺเญยฺยํ รูปํ, กายวิญฺเญยฺยํ รูปํ, มโนธาตุวิญฺเญยฺยํ รูปํ, มโนวิญฺญาณธาตุวิญฺเญยฺยํ รูปํ.\csromanspacer{} เอวํ สตฺตวิเธน รูปกฺขนฺโธ.}

\prose{\csromanfolio{15}อฏฺฐวิเธน รูปกฺขนฺโธ\csromandash{}จกฺขุวิญฺเญยฺยํ รูปํ, โสตวิญฺเญยฺยํ รูปํ, ฆานวิญฺเญยฺยํ รูปํ, ชิวฺหาวิญฺเญยฺยํ รูปํ, กายวิญฺเญยฺยํ รูปํ, อตฺถิ สุขสมฺผสฺสํ, อตฺถิ ทุกฺขสมฺผสฺสํ มโนธาตุวิญฺเญยฺยํ รูปํ, มโนวิญฺญาณธาตุวิญฺเญยฺยํ รูปํ.\csromanspacer{} เอวํ อฏฺฐวิเธน รูปกฺขนฺโธ.}

\prose{นววิเธน รูปกฺขนฺโธ\csromandash{}จกฺขุนฺทฺริยํ, โสตินฺทฺริยํ, ฆานินฺทฺริยํ, ชิวฺหินฺทฺริยํ, กายินฺทฺริยํ, อิตฺถินฺทฺริยํ, ปุริสินฺทฺริยํ, ชีวิตินฺทฺริยํ, ยญฺจ รูปํ น อินฺทฺริยํ.\csromanspacer{} เอวํ นววิเธน รูปกฺขนฺโธ.}

\prose{ทสวิเธน รูปกฺขนฺโธ\csromandash{}จกฺขุนฺทฺริยํ, โสตินฺทฺริยํ, ฆานินฺทฺริยํ, ชิวฺหินฺทฺริยํ, กายินฺทฺริยํ, อิตฺถินฺทฺริยํ, ปุริสินฺทฺริยํ, ชีวิตินฺทฺริยํ, น อินฺทฺริยํ รูปํ อตฺถิ สปฺปฏิฆํ, อตฺถิ อปฺปฏิฆํ.\csromanspacer{} เอวํ ทสวิเธน รูปกฺขนฺโธ.}

\prose{เอกาทสวิเธน รูปกฺขนฺโธ\csromandash{}จกฺขายตนํ, โสตายตนํ, ฆานายตนํ, ชิวฺหายตนํ, กายายตนํ, รูปายตนํ, สทฺทายตนํ, คนฺธายตนํ, รสายตนํ, โผฏฺฐพฺพายตนํ, ยญฺจ รูปํ อนิทสฺสน- อปฺปฏิฆํ ธมฺมายตนปริยาปนฺนํ.\csromanspacer{} เอวํ เอกาทสวิเธน รูปกฺขนฺโธ.}

\csromanheader{2}{อยํ วุจฺจติ รูปกฺขนฺโธ.}
\csromansectionrule
\csromanmatikaanchor{mtk.10}
\csromantocmarkref{subsection}{2. เวทนากฺขนฺธ}{mtk.10}
\bookmark[dest={mtk.10},level=2]{2. เวทนากฺขนฺธ}
\csromanheader{2}{2. เวทนากฺขนฺธ}
\proseitem{34}{ตตฺถ กตโม เวทนากฺขนฺโธ, เอกวิเธน เวทนากฺขนฺโธ\csromandash{} ผสฺสสมฺปยุตฺโต.}

\prose{ทุวิเธน เวทนากฺขนฺโธ\csromandash{}อตฺถิ สเหตุโก, อตฺถิ อเหตุโก.}

\prose{ติวิเธน เวทนากฺขนฺโธ\csromandash{}อตฺถิ กุสโล, อตฺถิ อกุสโล, อตฺถิ อพฺยากโต.}

\prose{จตุพฺพิเธน เวทนากฺขนฺโธ\csromandash{}อตฺถิ กามาวจโร, อตฺถิ รูปาวจโร, อตฺถิ อรูปาวจโร, อตฺถิ อปริยาปนฺโน.}

\prose{ปญฺจวิเธน เวทนากฺขนฺโธ\csromandash{}อตฺถิ สุขินฺทฺริยํ, อตฺถิ ทุกฺขินฺทฺริยํ, อตฺถิ โสมนสฺสินฺทฺริยํ, อตฺถิ โทมนสฺสินฺทฺริยํ, อตฺถิ อุเปกฺขินฺทฺริยํ.\csromanspacer{} เอวํ ปญฺจวิเธน เวทนากฺขนฺโธ.}

\prose{\csromanfolio{16}ฉพฺพิเธน เวทนากฺขนฺโธ\csromandash{}จกฺขุสมฺผสฺสชา เวทนา, โสตสมฺผสฺสชา เวทนา, ฆานสมฺผสฺสชา เวทนา, ชิวฺหาสมฺผสฺสชา เวทนา, กายสมฺผสฺสชา เวทนา, มโนสมฺผสฺสชา เวทนา.\csromanspacer{} เอวํ ฉพฺพิเธน เวทนากฺขนฺโธ.}

\prose{สตฺตวิเธน เวทนากฺขนฺโธ\csromandash{}จกฺขุสมฺผสฺสชา เวทนา, โสตสมฺผสฺสชา เวทนา, ฆานสมฺผสฺสชา เวทนา, ชิวฺหาสมฺผสฺสชา เวทนา, กายสมฺผสฺสชา เวทนา, มโนธาตุสมฺผสฺสชา เวทนา, มโนวิญฺญาณธาตุสมฺผสฺสชา เวทนา.\csromanspacer{} เอวํ สตฺตวิเธน เวทนากฺขนฺโธ.}

\prose{อฏฺฐวิเธน เวทนากฺขนฺโธ\csromandash{}จกฺขุสมฺผสฺสชา เวทนา, โสตสมฺผสฺสชา เวทนา, ฆานสมฺผสฺสชา เวทนา, ชิวฺหาสมฺผสฺสชา เวทนา, กายสมฺผสฺสชา เวทนา อตฺถิ สุขา, อตฺถิ ทุกฺขา, มโนธาตุสมฺผสฺสชา เวทนา, มโนวิญฺญาณธาตุสมฺผสฺสชา เวทนา.\csromanspacer{} เอวํ อฏฺฐวิเธน เวทนากฺขนฺโธ.}

\prose{นววิเธน เวทนากฺขนฺโธ\csromandash{}จกฺขุสมฺผสฺสชา เวทนา, โสตสมฺผสฺสชา เวทนา, ฆานสมฺผสฺสชา เวทนา, ชิวฺหาสมฺผสฺสชา เวทนา, กายสมฺผสฺสชา เวทนา, มโนธาตุสมฺผสฺสชา เวทนา, มโนวิญฺญาณธาตุสมฺผสฺสชา เวทนา อตฺถิ กุสลา, อตฺถิ อกุสลา, อตฺถิ อพฺยากตา.\csromanspacer{} เอวํ นววิเธน เวทนากฺขนฺโธ.}

\prose{ทสวิเธน เวทนากฺขนฺโธ\csromandash{}จกฺขุสมฺผสฺสชา เวทนา, โสตสมฺผสฺสชา เวทนา, ฆานสมฺผสฺสชา เวทนา, ชิวฺหาสมฺผสฺสชา เวทนา, กายสมฺผสฺสชา เวทนา อตฺถิ สุขา, อตฺถิ ทุกฺขา, มโนธาตุสมฺผสฺสชา เวทนา, มโนวิญฺญาณธาตุสมฺผสฺสชา เวทนา อตฺถิ กุสลา, อตฺถิ อกุสลา, อตฺถิ อพฺยากตา.\csromanspacer{} เอวํ ทสวิเธน เวทนากฺขนฺโธ.}

\proseitem{35}{เอกวิเธน เวทนากฺขนฺโธ\csromandash{}ผสฺสสมฺปยุตฺโต.}

\prose{ทุวิเธน เวทนากฺขนฺโธ\csromandash{}อตฺถิ สเหตุโก, อตฺถิ อเหตุโก.}

\prose{ติวิเธน เวทนากฺขนฺโธ\csromandash{}อตฺถิ วิปาโก, อตฺถิ วิปากธมฺมธมฺโม, อตฺถิ เนววิปากนวิปากธมฺมธมฺโม.\csromanspacer{} อตฺถิ อุปาทินฺนุปาทานิโย, อตฺถิ อนุปาทินฺนุปาทานิโย, อตฺถิ อนุปาทินฺน- อนุปาทานิโย.\csromanspacer{} อตฺถิ สํกิลิฏฺฐสํกิเลสิโก, อตฺถิ อสํกิลิฏฺฐสํกิเลสิโก, อตฺถิ อสํกิลิฏฺฐ-อสํกิเลสิโก.\csromanspacer{} อตฺถิ สวิตกฺกสวิจาโร, อตฺถิ อวิตกฺกวิจารมตฺโต, \csromanfolio{17}อตฺถิ อวิตกฺก-อวิจาโร.\csromanspacer{} อตฺถิ ทสฺสเนน ปหาตพฺโพ, อตฺถิ ภาวนาย ปหาตพฺโพ, อตฺถิ เนว ทสฺสเนน น ภาวนาย ปหาตพฺโพ.\csromanspacer{} อตฺถิ ทสฺสเนน ปหาตพฺพเหตุโก, อตฺถิ ภาวนาย ปหาตพฺพเหตุโก, อตฺถิ เนว ทสฺสเนน น ภาวนาย ปหาตพฺพเหตุโก.\csromanspacer{} อตฺถิ อาจยคามี, อตฺถิ อปจยคามี, อตฺถิ เนวาจยคามินาปจยคามี.\csromanspacer{} อตฺถิ เสกฺโข, อตฺถิ อเสกฺโข, อตฺถิ เนวเสกฺขนาเสกฺโข.\csromanspacer{} อตฺถิ ปริตฺโต, อตฺถิ มหคฺคโต, อตฺถิ อปฺปมาโณ.\csromanspacer{} อตฺถิ ปริตฺตารมฺมโณ, อตฺถิ มหคฺคตารมฺมโณ, อตฺถิ อปฺปมาณารมฺมโณ.\csromanspacer{} อตฺถิ หีโน, อตฺถิ มชฺฌิโม, อตฺถิ ปณีโต.\csromanspacer{} อตฺถิ มิจฺฉตฺตนิยโต, อตฺถิ สมฺมตฺตนิยโต, อตฺถิ อนิยโต.\csromanspacer{} อตฺถิ มคฺคารมฺมโณ, อตฺถิ มคฺคเหตุโก, อตฺถิ มคฺคาธิปติ.\csromanspacer{} อตฺถิ อุปฺปนฺโน, อตฺถิ อนุปฺปนฺโน, อตฺถิ อุปฺปาที.\csromanspacer{} อตฺถิ อตีโต, อตฺถิ อนาคโต, อตฺถิ ปจฺจุปฺปนฺโน.\csromanspacer{} อตฺถิ อตีตารมฺมโณ, อตฺถิ อนาคตารมฺมโณ, อตฺถิ ปจฺจุปฺปนฺนารมฺมโณ.\csromanspacer{} อตฺถิ อชฺฌตฺโต, อตฺถิ พหิทฺโธ, อตฺถิ อชฺฌตฺตพหิทฺโธ.\csromanspacer{} อตฺถิ อชฺฌตฺตารมฺมโณ, อตฺถิ พหิทฺธารมฺมโณ, อตฺถิ อชฺฌตฺตพหิทฺธารมฺมโณ ฯเปฯ.\csromanspacer{} เอวํ ทสวิเธน เวทนากฺขนฺโธ.}

\proseitem{36}{เอกวิเธน เวทนากฺขนฺโธ\csromandash{}ผสฺสสมฺปยุตฺโต.}

\prose{ทุวิเธน เวทนากฺขนฺโธ\csromandash{}อตฺถิ เหตุสมฺปยุตฺโต, อตฺถิ เหตุวิปฺปยุตฺโต ฯเปฯ.\csromanspacer{} อตฺถิ น เหตุสเหตุโก, อตฺถิ น เหตุ-อเหตุโก.\csromanspacer{} อตฺถิ โลกิโย, อตฺถิ โลกุตฺตโร.\csromanspacer{} อตฺถิ เกนจิ วิญฺเญยฺโย, อตฺถิ เกนจิ น วิญฺเญยฺโย.\csromanspacer{} อตฺถิ สาสโว, อตฺถิ อนาสโว.\csromanspacer{} อตฺถิ อาสวสมฺปยุตฺโต, อตฺถิ อาสววิปฺปยุตฺโต.\csromanspacer{} อตฺถิ อาสววิปฺปยุตฺตสาสโว, อตฺถิ อาสววิปฺปยุตฺต-อนาสโว.\csromanspacer{} อตฺถิ สํโยชนิโย, อตฺถิ อสํโยชนิโย.\csromanspacer{} อตฺถิ สํโยชนสมฺปยุตฺโต, อตฺถิ สํโยชนวิปฺปยุตฺโต.\csromanspacer{} อตฺถิ สํโยชนวิปฺปยุตฺตสํโยชนิโย, อตฺถิ สํโยชนวิปฺปยุตฺต-อสํโยชนิโย.\csromanspacer{} อตฺถิ คนฺถนิโย, อตฺถิ อคนฺถนิโย.\csromanspacer{} อตฺถิ คนฺถสมฺปยุตฺโต, อตฺถิ คนฺถวิปฺปยุตฺโต.\csromanspacer{} อตฺถิ คนฺถวิปฺปยุตฺตคนฺถนิโย, อตฺถิ คนฺถวิปฺปยุตฺต-อคนฺถนิโย.\csromanspacer{} อตฺถิ โอฆนิโย, อตฺถิ อโนฆนิโย.\csromanspacer{} อตฺถิ โอฆสมฺปยุตฺโต, อตฺถิ โอฆวิปฺปยุตฺโต.\csromanspacer{} อตฺถิ โอฆวิปฺปยุตฺต-โอฆนิโย, อตฺถิ โอฆวิปฺปยุตฺต-อโนฆนิโย.\csromanspacer{} อตฺถิ โยคนิโย, อตฺถิ อโยคนิโย.\csromanspacer{} อตฺถิ โยคสมฺปยุตฺโต, อตฺถิ โยควิปฺปยุตฺโต. \csromanfolio{18}อตฺถิ โยควิปฺปยุตฺตโยคนิโย, อตฺถิ โยควิปฺปยุตฺต-อโยคนิโย.\csromanspacer{} อตฺถิ นีวรณิโย, อตฺถิ อนีวรณิโย.\csromanspacer{} อตฺถิ นีวรณสมฺปยุตฺโต, อตฺถิ นีวรณวิปฺปยุตฺโต.\csromanspacer{} อตฺถิ นีวรณวิปฺปยุตฺตนีวรณิโย, อตฺถิ นีวรณวิปฺปยุตฺต-อนีวรณิโย.\csromanspacer{} อตฺถิ ปรามฏฺโฐ, อตฺถิ อปรามฏฺโฐ.\csromanspacer{} อตฺถิ ปรามาสสมฺปยุตฺโต, อตฺถิ ปรามาสวิปฺปยุตฺโต.\csromanspacer{} อตฺถิ ปรามาสวิปฺปยุตฺตปรามฏฺโฐ, อตฺถิ ปรามาสวิปฺปยุตฺต-อปรามฏฺโฐ.\csromanspacer{} อตฺถิ อุปาทินฺโน, อตฺถิ อนุปาทินฺโน.\csromanspacer{} อตฺถิ อุปาทานิโย, อตฺถิ อนุปาทานิโย.\csromanspacer{} อตฺถิ อุปาทานสมฺปยุตฺโต, อตฺถิ อุปาทานวิปฺปยุตฺโต.\csromanspacer{} อตฺถิ อุปาทานวิปฺปยุตฺต- อุปาทานิโย, อตฺถิ อุปาทานวิปฺปยุตฺต-อนุปาทานิโย.\csromanspacer{} อตฺถิ สํกิเลสิโก, อตฺถิ อสํกิเลสิโก.\csromanspacer{} อตฺถิ สํกิลิฏฺโฐ, อตฺถิ อสํกิลิฏฺโฐ.\csromanspacer{} อตฺถิ กิเลสสมฺปยุตฺโต, อตฺถิ กิเลสวิปฺปยุตฺโต.\csromanspacer{} อตฺถิ กิเลสวิปฺปยุตฺตสํกิเลสิโก, อตฺถิ กิเลสวิปฺปยุตฺต-อสํกิเลสิโก.\csromanspacer{} อตฺถิ ทสฺสเนน ปหาตพฺโพ, อตฺถิ น ทสฺสเนน ปหาตพฺโพ.\csromanspacer{} อตฺถิ ภาวนาย ปหาตพฺโพ, อตฺถิ น ภาวนาย ปหาตพฺโพ.\csromanspacer{} อตฺถิ ทสฺสเนน ปหาตพฺพเหตุโก, อตฺถิ น ทสฺสเนน ปหาตพฺพเหตุโก.\csromanspacer{} อตฺถิ ภาวนาย ปหาตพฺพเหตุโก, อตฺถิ น ภาวนาย ปหาตพฺพเหตุโก.\csromanspacer{} อตฺถิ สวิตกฺโก, อตฺถิ อวิตกฺโก.\csromanspacer{} อตฺถิ สวิจาโร, อตฺถิ อวิจาโร.\csromanspacer{} อตฺถิ สปฺปีติโก, อตฺถิ อปฺปีติโก, อตฺถิ ปีติสหคโต, อตฺถิ น ปีติสหคโต.\csromanspacer{} อตฺถิ กามาวจโร, อตฺถิ น กามาวจโร.\csromanspacer{} อตฺถิ รูปาวจโร, อตฺถิ น รูปาวจโร.\csromanspacer{} อตฺถิ อรูปาวจโร, อตฺถิ น อรูปาวจโร.\csromanspacer{} อตฺถิ ปริยาปนฺโน, อตฺถิ อปริยาปนฺโน.\csromanspacer{} อตฺถิ นิยฺยานิโก, อตฺถิ อนิยฺยานิโก.\csromanspacer{} อตฺถิ นิยโต, อตฺถิ อนิยโต.\csromanspacer{} อตฺถิ ส-อุตฺตโร, อตฺถิ อนุตฺตโร.\csromanspacer{} อตฺถิ สรโณ, อตฺถิ อรโณ.}

\prose{ติวิเธน เวทนากฺขนฺโธ\csromandash{}อตฺถิ กุสโล, อตฺถิ อกุสโล, อตฺถิ อพฺยากโต ฯเปฯ.\csromanspacer{} เอวํ ทสวิเธน เวทนากฺขนฺโธ.}

\proseitem{37}{เอกวิเธน เวทนากฺขนฺโธ\csromandash{}ผสฺสสมฺปยุตฺโต.}

\prose{ทุวิเธน เวทนากฺขนฺโธ\csromandash{}อตฺถิ สรโณ, อตฺถิ อรโณ.}

\prose{ติวิเธน เวทนากฺขนฺโธ\csromandash{}อตฺถิ วิปาโก, อตฺถิ วิปากธมฺมธมฺโม, อตฺถิ เนววิปากนวิปากธมฺมธมฺโม.\csromanspacer{} อตฺถิ อุปาทินฺนุปาทานิโย ฯเปฯ.\csromanspacer{} อตฺถิ \csromanfolio{19}อชฺฌตฺตารมฺมโณ, อตฺถิ พหิทฺธารมฺมโณ, อตฺถิ อชฺฌตฺตพหิทฺธารมฺมโณ ฯเปฯ.\csromanspacer{} เอวํ ทสวิเธน เวทนากฺขนฺโธ.\csromanspacer{} ทุกมูลกํ.}

\proseitem{38}{เอกวิเธน เวทนากฺขนฺโธ\csromandash{}ผสฺสสมฺปยุตฺโต.}

\prose{ทุวิเธน เวทนากฺขนฺโธ\csromandash{}อตฺถิ สเหตุโก, อตฺถิ อเหตุโก.}

\prose{ติวิเธน เวทนากฺขนฺโธ\csromandash{}อตฺถิ กุสโล, อตฺถิ อกุสโล, อตฺถิ อพฺยากโต ฯเปฯ.\csromanspacer{} เอวํ ทสวิเธน เวทนากฺขนฺโธ.}

\proseitem{39}{เอกวิเธน เวทนฺทากฺขนฺโธ\csromandash{}ผสฺสสมฺปยุตฺโต.}

\prose{ทุวิเธน เวทนากฺขนฺโธ\csromandash{}อตฺถิ เหตุสมฺปยุตฺโต, อตฺถิ เหตุวิปฺปยุตฺโต.}

\prose{ติวิเธน เวทนากฺขนฺโธ\csromandash{}อตฺถิ กุสโล, อตฺถิ อกุสโล, อตฺถิ อพฺยากโต ฯเปฯ.\csromanspacer{} เอวํ ทสวิเธน เวทนากฺขนฺโธ.}

\proseitem{40}{เอกวิเธน เวทนากฺขนฺโธ\csromandash{}ผสฺสสมฺปยุตฺโต.}

\prose{ทุวิเธน เวทนากฺขนฺโธ\csromandash{}อตฺถิ น เหตุสเหตุโก, อตฺถิ น เหตุ- อเหตุโก.\csromanspacer{} อตฺถิ โลกิโย, อตฺถิ โลกุตฺตโร.\csromanspacer{} อตฺถิ เกนจิ วิญฺเญยฺโย, อตฺถิ เกนจิ น วิญฺเญยฺโย.\csromanspacer{} อตฺถิ สาสโว, อตฺถิ อนาสโว.\csromanspacer{} อตฺถิ อาสวสมฺปยุตฺโต, อตฺถิ อาสววิปฺปยุตฺโต.\csromanspacer{} อตฺถิ อาสววิปฺปยุตฺตสาสโว, อตฺถิ อาสววิปฺปยุตฺต- อนาสโว ฯเปฯ.\csromanspacer{} อตฺถิ สรโณ, อตฺถิ อรโณ.}

\prose{ติวิเธน เวทนากฺขนฺโธ\csromandash{}อตฺถิ กุสโล, อตฺถิ อกุสโล, อตฺถิ อพฺยากโต ฯเปฯ.\csromanspacer{} เอวํ ทสวิเธน เวทนากฺขนฺโธ.}

\proseitem{41}{เอกวิเธน เวทนากฺขนฺโธ\csromandash{}ผสฺสสมฺปยุตฺโต.}

\prose{ทุวิเธน เวทนากฺขนฺโธ\csromandash{}อตฺถิ สเหตุโก, อตฺถิ อเหตุโก.}

\prose{ติวิเธน เวทนากฺขนฺโธ\csromandash{}อตฺถิ วิปาโก, อตฺถิ วิปากธมฺมธมฺโม, อตฺถิ เนววิปากนวิปากธมฺมธมฺโม.\csromanspacer{} อตฺถิ อุปาทินฺนุปาทานิโย, อตฺถิ \csromanfolio{20}อนุปาทินฺนุปาทานิโย, อตฺถิ อนุปาทินฺน-อนุปาทานิโย.\csromanspacer{} อตฺถิ สํกิลิฏฺฐสํกิเลสิโก, อตฺถิ อสํกิลิฏฺฐสํกิเลสิโก, อตฺถิ อสํกิลิฏฺฐ- อสํกิเลสิโก.\csromanspacer{} อตฺถิ สวิตกฺกสวิจาโร, อตฺถิ อวิตกฺกวิจารมตฺโต, อตฺถิ อวิตกฺก-อวิจาโร.\csromanspacer{} อตฺถิ ทสฺสเนน ปหาตพฺโพ, อตฺถิ ภาวนาย ปหาตพฺโพ, อตฺถิ เนว ทสฺสเนน น ภาวนาย ปหาตพฺโพ.\csromanspacer{} อตฺถิ ทสฺสเนน ปหาตพฺพเหตุโก, อตฺถิ ภาวนาย ปหาตพฺพเหตุโก, อตฺถิ เนว ทสฺสเนน น ภาวนาย ปหาตพฺพเหตุโก.\csromanspacer{} อตฺถิ อาจยคามี, อตฺถิ อปจยคามี, อตฺถิ เนวาจยคามินาปจยคามี.\csromanspacer{} อตฺถิ เสกฺโข, อตฺถิ อเสกฺโข, อตฺถิ เนวเสกฺขนาเสกฺโข.\csromanspacer{} อตฺถิ ปริตฺโต, อตฺถิ มหคฺคโต, อตฺถิ อปฺปมาโณ.\csromanspacer{} อตฺถิ ปริตฺตารมฺมโณ, อตฺถิ มหคฺคตารมฺมโณ, อตฺถิ อปฺปมาณารมฺมโณ.\csromanspacer{} อตฺถิ หีโน, อตฺถิ มชฺฌิโม, อตฺถิ ปณีโต.\csromanspacer{} อตฺถิ มิจฺฉตฺตนิยโต, อตฺถิ สมฺมตฺตนิยโต, อตฺถิ อนิยโต.\csromanspacer{} อตฺถิ มคฺคารมฺมโณ, อตฺถิ มคฺคเหตุโก, อตฺถิ มคฺคาธิปติ.\csromanspacer{} อตฺถิ อุปฺปนฺโน, อตฺถิ อนุปฺปนฺโน, อตฺถิ อุปฺปาที.\csromanspacer{} อตฺถิ อตีโต, อตฺถิ อนาคโต, อตฺถิ ปจฺจุปฺปนฺโน.\csromanspacer{} อตฺถิ อตีตารมฺมโณ, อตฺถิ อนาคตารมฺมโณ, อตฺถิ ปจฺจุปฺปนฺนารมฺมโณ.\csromanspacer{} อตฺถิ อชฺฌตฺโต, อตฺถิ พหิทฺโธ, อตฺถิ อชฺฌตฺตพหิทฺโธ.\csromanspacer{} อตฺถิ อชฺฌตฺตารมฺมโณ, อตฺถิ พหิทฺธารมฺมโณ, อตฺถิ อชฺฌตฺตพหิทฺธารมฺมโณ ฯเปฯ.\csromanspacer{} เอวํ ทสวิเธน เวทนากฺขนฺโธ.}

\proseitem{42}{เอกวิเธน เวทนากฺขนฺโธ\csromandash{}ผสฺสสมฺปยุตฺโต.}

\prose{ทุวิเธน เวทนากฺขนฺโธ\csromandash{}อตฺถิ เหตุสมฺปยุตฺโต, อตฺถิ เหตุวิปฺปยุตฺโต.\csromanspacer{} อตฺถิ น เหตุสเหตุโก, อตฺถิ น เหตุ-อเหตุโก.\csromanspacer{} อตฺถิ โลกิโย, อตฺถิ โลกุตฺตโร.\csromanspacer{} อตฺถิ เกนจิ วิญฺเญยฺโย, อตฺถิ เกนจิ น วิญฺเญยฺโย.\csromanspacer{} อตฺถิ สาสโว, อตฺถิ อนาสโว.\csromanspacer{} อตฺถิ อาสวสมฺปยุตฺโต, อตฺถิ อาสววิปฺปยุตฺโต.\csromanspacer{} อตฺถิ อาสววิปฺปยุตฺตสาสโว, อตฺถิ อาสววิปฺปยุตฺต-อนาสโว.\csromanspacer{} อตฺถิ สํโยชนิโย อตฺถิ อสํโยชนิโย ป-.\csromanspacer{} อตฺถิ สรโณ, อตฺถิ อรโณ.}

\prose{ติวิเธน เวทนากฺขนฺโธ\csromandash{}อตฺถิ อชฺฌตฺตารมฺมโณ, อตฺถิ พหิทฺธารมฺมโณ, อตฺถิ อชฺฌตฺตพหิทฺธารมฺมโณ ฯเปฯ.\csromanspacer{} เอวํ ทสวิเธน เวทนากฺขนฺโธ.}

\vspace{\tipitakaparskip}
\csromancenter{ติกมูลกํ.}
\proseitem{43}{\csromanfolio{21}เอกวิเธน เวทนากฺขนฺโธ\csromandash{}ผสฺสสมฺปยุตฺโต.}

\prose{ทุวิเธน เวทนากฺขนฺโธ\csromandash{}อตฺถิ สเหตุโก, อตฺถิ อเหตุโก.}

\prose{ติวิเธน เวทนากฺขนฺโธ\csromandash{}อตฺถิ กุสโล, อตฺถิ อกุสโล, อตฺถิ อพฺยากโต ฯเปฯ.\csromanspacer{} เอวํ ทสวิเธน เวทนากฺขนฺโธ.}

\proseitem{44}{เอกวิเธน เวทนากฺขนฺโธ\csromandash{}ผสฺสสมฺปยุตฺโต.}

\prose{ทุวิเธน เวทนากฺขนฺโธ\csromandash{}อตฺถิ เหตุสมฺปยุตฺโต, อตฺถิ เหตุวิปฺปยุตฺโต.}

\prose{ติวิเธน เวทนากฺขนฺโธ\csromandash{}อตฺถิ วิปาโก, อตฺถิ วิปากธมฺมธมฺโม, อตฺถิ เนววิปากนวิปากธมฺมธมฺโม ฯเปฯ.\csromanspacer{} เอวํ ทสวิเธน เวทนากฺขนฺโธ.}

\proseitem{45}{เอกวิเธน เวทนากฺขนฺโธ\csromandash{}ผสฺสสมฺปยุตฺโต.}

\prose{ทุวิเธน เวทนากฺขนฺโธ\csromandash{}อตฺถิ น เหตุสเหตุโก, อตฺถิ น เหตุ- อเหตุโก.}

\prose{ติวิเธน เวทนากฺขนฺโธ\csromandash{}อตฺถิ อุปาทินฺนุปาทานิโย, อตฺถิ อนุปาทินฺนุปาทานิโย, อตฺถิ อนุปาทินฺน-อนุปาทานิโย ฯเปฯ.\csromanspacer{} เอวํ ทสวิเธน เวทนากฺขนฺโธ.}

\proseitem{46}{เอกวิเธน เวทนากฺขนฺโธ\csromandash{}ผสฺสสมฺปยุตฺโต.}

\prose{ทุวิเธน เวทนากฺขนฺโธ\csromandash{}อตฺถิ โลกิโย, อตฺถิ โลกุตฺตโร.}

\prose{ติวิเธน เวทนากฺขนฺโธ\csromandash{}อตฺถิ สํกิลิฏฺฐสํกิเลสิโก, อตฺถิ อสํกิลิฏฺฐสํกิเลสิโก, อตฺถิ อสํกิลิฏฺฐ-อสํกิเลสิโก ฯเปฯ.\csromanspacer{} เอวํ ทสวิเธน เวทนากฺขนฺโธ.}

\proseitem{47}{เอกวิเธน เวทนากฺขนฺโธ\csromandash{}ผสฺสสมฺปยุตฺโต.}

\prose{ทุวิเธน เวทนากฺขนฺโธ\csromandash{}อตฺถิ เกนจิ วิญฺเญยฺโย, อตฺถิ เกนจิ น วิญฺเญยฺโย.}

\prose{ติวิเธน เวทนากฺขนฺโธ\csromandash{}อตฺถิ สวิตกฺกสวิจาโร, อตฺถิ อวิตกฺกวิจารมตฺโต, อตฺถิ อวิตกฺก-อวิจาโร ฯเปฯ.\csromanspacer{} เอวํ ทสวิเธน เวทนากฺขนฺโธ.}

\proseitem{48}{\csromanfolio{22}เอกวิเธน เวทนากฺขนฺโธ\csromandash{}ผสฺสสมฺปยุตฺโต.}

\prose{ทุวิเธน เวทนากฺขนฺโธ\csromandash{}อตฺถิ สาสโว, อตฺถิ อนาสโว.}

\prose{ติวิเธน เวทนากฺขนฺโธ\csromandash{}อตฺถิ ทสฺสเนน ปหาตพฺโพ, อตฺถิ ภาวนาย ปหาตพฺโพ, อตฺถิ เนว ทสฺสเนน น ภาวนาย ปหาตพฺโพ ฯเปฯ.\csromanspacer{} เอวํ ทสวิเธน เวทนากฺขนฺโธ.}

\proseitem{49}{เอกวิเธน เวทนากฺขนฺโธ\csromandash{}ผสฺสสมฺปยุตฺโต.}

\prose{ทุวิเธน เวทนากฺขนฺโธ\csromandash{}อตฺถิ อาสวสมฺปยุตฺโต, อตฺถิ อาสววิปฺปยุตฺโต.}

\prose{ติวิเธน เวทนากฺขนฺโธ\csromandash{}อตฺถิ ทสฺสเนน ปหาตพฺพเหตุโก, อตฺถิ ภาวนาย ปหาตพฺพเหตุโก, อตฺถิ เนว ทสฺสเนน น ภาวนาย ปหาตพฺพเหตุโก ฯเปฯ.\csromanspacer{} เอวํ ทสวิเธน เวทนากฺขนฺโธ.}

\proseitem{50}{เอกวิเธน เวทนากฺขนฺโธ\csromandash{}ผสฺสสมฺปยุตฺโต.}

\prose{ทุวิเธน เวทนากฺขนฺโธ\csromandash{}อตฺถิ อาสววิปฺปยุตฺตสาสโว, อตฺถิ อาสววิปฺปยุตฺต-อนาสโว.}

\prose{ติวิเธน เวทนากฺขนฺโธ\csromandash{}อตฺถิ อาจยคามี, อตฺถิ อปจยคามี, อตฺถิ เนวาจยคามินาปจยคามี ฯเปฯ.\csromanspacer{} เอวํ ทสวิเธน เวทนากฺขนฺโธ.}

\proseitem{51}{เอกวิเธน เวทนากฺขนฺโธ\csromandash{}ผสฺสสมฺปยุตฺโต.}

\prose{ทุวิเธน เวทนากฺขนฺโธ\csromandash{}อตฺถิ สํโยชนิโย, อตฺถิ อสํโยชนิโย.}

\prose{ติวิเธน เวทนากฺขนฺโธ\csromandash{}อตฺถิ เสกฺโข, อตฺถิ อเสกฺโข, อตฺถิ เนวเสกฺขนาเสกฺโข ฯเปฯ.\csromanspacer{} เอวํ ทสวิเธน เวทนากฺขนฺโธ.}

\proseitem{52}{เอกวิเธน เวทนากฺขนฺโธ\csromandash{}ผสฺสสมฺปยุตฺโต.}

\prose{ทุวิเธน เวทนากฺขนฺโธ\csromandash{}อตฺถิ สํโยชนสมฺปยุตฺโต, อตฺถิ สํโยชนวิปฺปยุตฺโต.}

\prose{ติวิเธน เวทนากฺขนฺโธ\csromandash{}อตฺถิ ปริตฺโต, อตฺถิ มหคฺคโต, อตฺถิ อปฺปมาโณ ฯเปฯ.\csromanspacer{} เอวํ ทสวิเธน เวทนากฺขนฺโธ.}

\proseitem{53}{\csromanfolio{23}เอกวิเธน เวทนากฺขนฺโธ\csromandash{}ผสฺสสมฺปยุตฺโต.}

\prose{ทุวิเธน เวทนากฺขนฺโธ\csromandash{}อตฺถิ สํโยชนวิปฺปยุตฺตสํโยชนิโย, อตฺถิ สํโยชนวิปฺปยุตฺต-อสํโยชนิโย.}

\prose{ติวิเธน เวทนากฺขนฺโธ\csromandash{}อตฺถิ ปริตฺตารมฺมโณ, อตฺถิ มหคฺคตารมฺมโณ, อตฺถิ อปฺปมาณารมฺมโณ ฯเปฯ.\csromanspacer{} เอวํ ทสวิเธน เวทนากฺขนฺโธ.}

\proseitem{54}{เอกวิเธน เวทนากฺขนฺโธ\csromandash{}ผสฺสสมฺปยุตฺโต.}

\prose{ทุวิเธน เวทนากฺขนฺโธ\csromandash{}อตฺถิ คนฺถนิโย, อตฺถิ อคนฺถนิโย.}

\prose{ติวิเธน เวทนากฺขนฺโธ\csromandash{}อตฺถิ หีโน, อตฺถิ มชฺฌิโม, อตฺถิ ปณีโต ฯเปฯ.\csromanspacer{} เอวํ ทสวิเธน เวทนากฺขนฺโธ.}

\proseitem{55}{เอกวิเธน เวทนากฺขนฺโธ\csromandash{}ผสฺสสมฺปยุตฺโต.}

\prose{ทุวิเธน เวทนากฺขนฺโธ\csromandash{}อตฺถิ คนฺถสมฺปยุตฺโต, อตฺถิ คนฺถวิปฺปยุตฺโต.}

\prose{ติวิเธน เวทนากฺขนฺโธ\csromandash{}อตฺถิ มิจฺฉตฺตนิยโต, อตฺถิ สมฺมตฺตนิยโต, อตฺถิ อนิยโต ฯเปฯ.\csromanspacer{} เอวํ ทสวิเธน เวทนากฺขนฺโธ.}

\proseitem{56}{เอกวิเธน เวทนากฺขนฺโธ\csromandash{}ผสฺสสมฺปยุตฺโต.}

\prose{ทุวิเธน เวทนากฺขนฺโธ\csromandash{}อตฺถิ คนฺถวิปฺปยุตฺตคนฺถนิโย, อตฺถิ คนฺถวิปฺปยุตฺต-อคนฺถนิโย.}

\prose{ติวิเธน เวทนากฺขนฺโธ\csromandash{}อตฺถิ มคฺคารมฺมโณ, อตฺถิ มคฺคเหตุโก, อตฺถิ มคฺคาธิปติ ฯเปฯ.\csromanspacer{} เอวํ ทสวิเธน เวทนากฺขนฺโธ.}

\proseitem{57}{เอกวิเธน เวทนากฺขนฺโธ\csromandash{}ผสฺสสมฺปยุตฺโต.}

\prose{ทุวิเธน เวทนากฺขนฺโธ\csromandash{}อตฺถิ โอฆนิโย, อตฺถิ อโนฆนิโย.}

\prose{ติวิเธน เวทนากฺขนฺโธ\csromandash{}อตฺถิ อุปฺปนฺโน, อตฺถิ อนุปฺปนฺโน, อตฺถิ อุปฺปาที ฯเปฯ.\csromanspacer{} เอวํ ทสวิเธน เวทนากฺขนฺโธ.}

\proseitem{58}{เอกวิเธน เวทนากฺขนฺโธ\csromandash{}ผสฺสสมฺปยุตฺโต.}

\prose{ทุวิเธน เวทนากฺขนฺโธ\csromandash{}อตฺถิ โอฆสมฺปยุตฺโต, อตฺถิ โอฆวิปฺปยุตฺโต.}

\prose{\csromanfolio{24}ติวิเธน เวทนากฺขนฺโธ\csromandash{}อตฺถิ อตีโต, อตฺถิ อนาคโต, อตฺถิ ปจฺจุปฺปนฺโน ฯเปฯ.\csromanspacer{} เอวํ ทสวิเธน เวทนากฺขนฺโธ.}

\proseitem{59}{เอกวิเธน เวทนากฺขนฺโธ\csromandash{}ผสฺสสมฺปยุตฺโต.}

\prose{ทุวิเธน เวทนากฺขนฺโธ\csromandash{}อตฺถิ โอฆวิปฺปยุตฺต-โอฆนิโย, อตฺถิ โอฆวิปฺปยุตฺต-อโนฆนิโย.}

\prose{ติวิเธน เวทนากฺขนฺโธ\csromandash{}อตฺถิ อตีตารมฺมโณ, อตฺถิ อนาคตารมฺมโณ, อตฺถิ ปจฺจุปฺปนฺนารมฺมโณ ฯเปฯ.\csromanspacer{} เอวํ ทสวิเธน เวทนากฺขนฺโธ.}

\proseitem{60}{เอกวิเธน เวทนากฺขนฺโธ\csromandash{}ผสฺสสมฺปยุตฺโต.}

\prose{ทุวิเธน เวทนากฺขนฺโธ\csromandash{}อตฺถิ โยคนิโย, อตฺถิ อโยคนิโย.}

\prose{ติวิเธน เวทนากฺขนฺโธ\csromandash{}อตฺถิ อชฺฌตฺโต, อตฺถิ พหิทฺโธ, อตฺถิ อชฺฌตฺตพหิทฺโธ ฯเปฯ.\csromanspacer{} เอวํ ทสวิเธน เวทนากฺขนฺโธ.}

\proseitem{61}{เอกวิเธน เวทนากฺขนฺโธ\csromandash{}ผสฺสสมฺปยุตฺโต.}

\prose{ทุวิเธน เวทนากฺขนฺโธ\csromandash{}อตฺถิ โยคสมฺปยุตฺโต, อตฺถิ โยควิปฺปยุตฺโต.}

\prose{ติวิเธน เวทนากฺขนฺโธ\csromandash{}อตฺถิ อชฺฌตฺตารมฺมโณ, อตฺถิ พหิทฺธารมฺมโณ, อตฺถิ อชฺฌตฺตพหิทฺธารมฺมโณ ฯเปฯ.\csromanspacer{} เอวํ ทสวิเธน เวทนากฺขนฺโธ.\csromanspacer{} อุภโตวฑฺฒกํ.}

\prose{สตฺตวิเธน เวทนากฺขนฺโธ\csromandash{}อตฺถิ กุสโล, อตฺถิ อกุสโล, อตฺถิ อพฺยากโต, อตฺถิ กามาวจโร, อตฺถิ รูปาวจโร, อตฺถิ อรูปาวจโร, อตฺถิ อปริยาปนฺโน.\csromanspacer{} เอวํ สตฺตวิเธน เวทนากฺขนฺโธ.}

\prose{อปโรปิ สตฺตวิเธน เวทนากฺขนฺโธ\csromandash{}อตฺถิ วิปาโก ฯเปฯ อตฺถิ อชฺฌตฺตารมฺมโณ, อตฺถิ พหิทฺธารมฺมโณ, อตฺถิ อชฺฌตฺตพหิทฺธารมฺมโณ, อตฺถิ กามาวจโร, อตฺถิ รูปาวจโร, อตฺถิ อรูปาวจโร, อตฺถิ อปริยาปนฺโน.\csromanspacer{} เอวํ สตฺตวิเธน เวทนากฺขนฺโธ.}

\prose{\csromanfolio{25}จตุวีสติวิเธน เวทนากฺขนฺโธ\csromandash{}จกฺขุสมฺผสฺสปจฺจยา เวทนากฺขนฺโธ อตฺถิ กุสโล, อตฺถิ อกุสโล, อตฺถิ อพฺยากโต.\csromanspacer{} โสตสมฺผสฺสปจฺจยา เวทนากฺขนฺโธ ฯเปฯ.\csromanspacer{} ฆานสมฺผสฺสปจฺจยา เวทนากฺขนฺโธ ฯเปฯ.\csromanspacer{} ชิวฺหาสมฺผสฺสปจฺจยา เวทนากฺขนฺโธ ฯเปฯ.\csromanspacer{} กายสมฺผสฺสปจฺจยา เวทนากฺขนฺโธ ฯเปฯ.\csromanspacer{} มโนสมฺผสฺสปจฺจยา เวทนากฺขนฺโธ อตฺถิ กุสโล, อตฺถิ อกุสโล, อตฺถิ อพฺยากโต.\csromanspacer{} จกฺขุสมฺผสฺสชา เวทนา, โสตสมฺผสฺสชา เวทนา, ฆานสมฺผสฺสชา เวทนา, ชิวฺหาสมฺผสฺสชา เวทนา, กายสมฺผสฺสชา เวทนา, มโนสมฺผสฺสชา เวทนา.\csromanspacer{} เอวํ จตุวีสติวิเธน เวทนากฺขนฺโธ.}

\prose{อปโรปิ จตุวีสติวิเธน เวทนากฺขนฺโธ\csromandash{}จกฺขุสมฺผสฺสปจฺจยา เวทนากฺขนฺโธ อตฺถิ วิปาโก ฯเปฯ อตฺถิ อชฺฌตฺตารมฺมโณ, อตฺถิ พหิทฺธารมฺมโณ, อตฺถิ อชฺฌตฺตพหิทฺธารมฺมโณ.\csromanspacer{} จกฺขุสมฺผสฺสชา เวทนา ฯเปฯ มโนสมฺผสฺสชา เวทนา.\csromanspacer{} โสตสมฺผสฺสปจฺจยา เวทนากฺขนฺโธ ฯเปฯ.\csromanspacer{} ฆานสมฺผสฺสปจฺจยา เวทนากฺขนฺโธ ฯเปฯ.\csromanspacer{} ชิวฺหาสมฺผสฺสปจฺจยา เวทนากฺขนฺโธ ฯเปฯ.\csromanspacer{} กายสมฺผสฺสปจฺจยา เวทนากฺขนฺโธ ฯเปฯ.\csromanspacer{} มโนสมฺผสฺสปจฺจยา เวทนากฺขนฺโธ อตฺถิ วิปาโก ฯเปฯ อตฺถิ อชฺฌตฺตารมฺมโณ, อตฺถิ พหิทฺธารมฺมโณ, อตฺถิ อชฺฌตฺตพหิทฺธารมฺมโณ.\csromanspacer{} จกฺขุสมฺผสฺสชา เวทนา ฯเปฯ มโนสมฺผสฺสชา เวทนา.\csromanspacer{} เอวํ จตุวีสติวิเธน เวทนากฺขนฺโธ.}

\prose{ติํสวิเธน เวทนากฺขนฺโธ\csromandash{}จกฺขุสมฺผสฺสปจฺจยา เวทนากฺขนฺโธ อตฺถิ กามาวจโร, อตฺถิ รูปาวจโร, อตฺถิ อรูปาวจโร, อตฺถิ อปริยาปนฺโน.\csromanspacer{} โสตสมฺผสฺสปจฺจยา ฯเปฯ.\csromanspacer{} ฆานสมฺผสฺสปจฺจยา ฯเปฯ.\csromanspacer{} ชิวฺหาสมฺผสฺสปจฺจยา ฯเปฯ.\csromanspacer{} กายสมฺผสฺสปจฺจยา ฯเปฯ.\csromanspacer{} มโนสมฺผสฺสปจฺจยา เวทนากฺขนฺโธ อตฺถิ กามาวจโร, อตฺถิ รูปาวจโร, อตฺถิ อรูปาวจโร, อตฺถิ อปริยาปนฺโน.\csromanspacer{} จกฺขุสมฺผสฺสชา เวทนา, โสตสมฺผสฺสชา เวทนา, ฆานสมฺผสฺสชา เวทนา, ชิวฺหาสมฺผสฺสชา เวทนา, กายสมฺผสฺสชา เวทนา, มโนสมฺผสฺสชา เวทนา.\csromanspacer{} เอวํ ติํสวิเธน เวทนากฺขนฺโธ.}

\prose{พหุวิเธน เวทนากฺขนฺโธ\csromandash{}จกฺขุสมฺผสฺสปจฺจยา เวทนากฺขนฺโธ อตฺถิ กุสโล, อตฺถิ อกุสโล, อตฺถิ อพฺยากโต, อตฺถิ กามาวจโร, อตฺถิ รูปาวจโร, อตฺถิ อรูปาวจโร, อตฺถิ อปริยาปนฺโน.\csromanspacer{} โสตสมฺผสฺสปจฺจยา ฯเปฯ.\csromanspacer{} ฆานสมฺผสฺสปจฺจยา ฯเปฯ.\csromanspacer{} ชิวฺหาสมฺผสฺสปจฺจยา ฯเปฯ.\csromanspacer{} กายสมฺผสฺสปจฺจยา ฯเปฯ.\csromanspacer{} มโนสมฺผสฺสปจฺจยา เวทนากฺขนฺโธ อตฺถิ \csromanfolio{26}กุสโล, อตฺถิ อกุสโล, อตฺถิ อพฺยากโต, อตฺถิ กามาวจโร, อตฺถิ รูปาวจโร, อตฺถิ อรูปาวจโร, อตฺถิ อปริยาปนฺโน.\csromanspacer{} จกฺขุสมฺผสฺสชา เวทนา, โสตสมฺผสฺสชา เวทนา, ฆานสมฺผสฺสชา เวทนา, ชิวฺหาสมฺผสฺสชา เวทนา, กายสมฺผสฺสชา เวทนา, มโนสมฺผสฺสชา เวทนา.\csromanspacer{} เอวํ พหุวิเธน เวทนากฺขนฺโธ.}

\prose{อปโรปิ พหุวิเธน เวทนากฺขนฺโธ\csromandash{}จกฺขุสมฺผสฺสปจฺจยา เวทนากฺขนฺโธ อตฺถิ วิปาโก ฯเปฯ.\csromanspacer{} อตฺถิ อชฺฌตฺตารมฺมโณ, อตฺถิ พหิทฺธารมฺมโณ, อตฺถิ อชฺฌตฺตพหิทฺธารมฺมโณ, อตฺถิ กามาวจโร, อตฺถิ รูปาวจโร, อตฺถิ อรูปาวจโร, อตฺถิ อปริยาปนฺโน.\csromanspacer{} โสตสมฺผสฺสปจฺจยา เวทนากฺขนฺโธ ฯเปฯ.\csromanspacer{} ฆานสมฺผสฺสปจฺจยา เวทนากฺขนฺโธ ฯเปฯ.\csromanspacer{} ชิวฺหาสมฺผสฺสปจฺจยา เวทนากฺขนฺโธ ฯเปฯ.\csromanspacer{} กายสมฺผสฺสปจฺจยา เวทนากฺขนฺโธ ฯเปฯ.\csromanspacer{} มโนสมฺผสฺสปจฺจยา เวทนากฺขนฺโธ อตฺถิ วิปาโก ฯเปฯ อตฺถิ อชฺฌตฺตารมฺมโณ, อตฺถิ พหิทฺธารมฺมโณ, อตฺถิ อชฺฌตฺตพหิทฺธารมฺมโณ, อตฺถิ กามาวจโร, อตฺถิ รูปาวจโร, อตฺถิ อรูปาวจโร, อตฺถิ อปริยาปนฺโน.\csromanspacer{} จกฺขุสมฺผสฺสชา เวทนา, โสตสมฺผสฺสชา เวทนา, ฆานสมฺผสฺสชา เวทนา, ชิวฺหาสมฺผสฺสชา เวทนา, กายสมฺผสฺสชา เวทนา, มโนสมฺผสฺสชา เวทนา.\csromanspacer{} เอวํ พหุวิเทน เวทนากฺขนฺโธ.}

\csromanheader{2}{อยํ วุจฺจติ เวทนากฺขนฺโธ.}
\csromansectionrule
\csromanmatikaanchor{mtk.11}
\csromantocmarkref{subsection}{3. สญฺญากฺขนฺธ}{mtk.11}
\bookmark[dest={mtk.11},level=2]{3. สญฺญากฺขนฺธ}
\csromanheader{2}{3. สญฺญากฺขนฺธ}
\proseitem{62}{ตตฺถ กตโม สญฺญากฺขนฺโธ, เอกวิเธน สญฺญากฺขนฺโธ\csromandash{} ผสฺสสมฺปยุตฺโต.}

\prose{ทุวิเธน สญฺญากฺขนฺโธ\csromandash{}อตฺถิ สเหตุโก, อตฺถิ อเหตุโก.}

\prose{ติวิเธน สญฺญากฺขนฺโธ\csromandash{}อตฺถิ กุสโล, อตฺถิ อกุสโล, อตฺถิ อพฺยากโต.}

\prose{จตุพฺพิเธน สญฺญากฺขนฺโธ\csromandash{}อตฺถิ กามาวจโร, อตฺถิ รูปาวจโร, อตฺถิ อรูปาวจโร, อตฺถิ อปริยาปนฺโน.}

\prose{\csromanfolio{27}ปญฺจวิเธน สญฺญากฺขนฺโธ\csromandash{}อตฺถิ สุขินฺทฺริยสมฺปยุตฺโต, อตฺถิ ทุกฺขินฺทฺริยสมฺปยุตฺโต, อตฺถิ โสมนสฺสินฺทฺริยสมฺปยุตฺโต, อตฺถิ โทมนสฺสินฺทฺริยสมฺปยุตฺโต, อตฺถิ อุเปกฺขินฺทฺริยสมฺปยุตฺโต.\csromanspacer{} เอวํ ปญฺจวิเธน สญฺญากฺขนฺโธ.}

\prose{ฉพฺพิเธน สญฺญากฺขนฺโธ\csromandash{}จกฺขุสมฺผสฺสชา สญฺญา, โสตสมฺผสฺสชา สญฺญา, ฆานสมฺผสฺสชา สญฺญา, ชิวฺหาสมฺผสฺสชา สญฺญา, กายสมฺผสฺสชา สญฺญา, มโนสมฺผสฺสชา สญฺญา.\csromanspacer{} เอวํ ฉพฺพิเธน สญฺญากฺขนฺโธ.}

\prose{สตฺตวิเธน สญฺญากฺขนฺโธ\csromandash{}จกฺขุสมฺผสฺสชา สญฺญา, โสตสมฺผสฺสชา สญฺญา, ฆานสมฺผสฺสชา สญฺญา, ชิวฺหาสมฺผสฺสชา สญฺญา, กายสมฺผสฺสชา สญฺญา, มโนธาตุสมฺผสฺสชา สญฺญา, มโนวิญฺญาณธาตุสมฺผสฺสชา สญฺญา.\csromanspacer{} เอวํ สตฺตวิเธน สญฺญากฺขนฺโธ.}

\prose{อฏฺฐวิเธน สญฺญากฺขนฺโธ\csromandash{}จกฺขุสมฺผสฺสชา สญฺญา ฯเปฯ กายสมฺผสฺสชา สญฺญา อตฺถิ สุขสหคตา, อตฺถิ ทุกฺขสหคตา, มโนธาตุสมฺผสฺสชา สญฺญา, มโนวิญฺญาณธาตุสมฺผสฺสชา สญฺญา.\csromanspacer{} เอวํ อฏฺฐวิเธน สญฺญากฺขนฺโธ.}

\prose{นววิเธน สญฺญากฺขนฺโธ\csromandash{}จกฺขุสมฺผสฺสชา สญฺญา ฯเปฯ กายสมฺผสฺสชา สญฺญา, มโนธาตุสมฺผสฺสชา สญฺญา, มโนวิญฺญาณธาตุสมฺผสฺสชา สญฺญา อตฺถิ กุสลา, อตฺถิ อกุสลา, อตฺถิ อพฺยากตา.\csromanspacer{} เอวํ นววิเธน สญฺญากฺขนฺโธ.}

\prose{ทสวิเธน สญฺญากฺขนฺโธ\csromandash{}จกฺขุสมฺผสฺสชา สญฺญา ฯเปฯ กายสมฺผสฺสชา สญฺญา อตฺถิ สุขสหคตา, อตฺถิ ทุกฺขสหคตา, มโนธาตุสมฺผสฺสชา สญฺญา, มโนวิญฺญาณธาตุสมฺผสฺสชา สญฺญา อตฺถิ กุสลา, อตฺถิ อกุสลา, อตฺถิ อพฺยากตา.\csromanspacer{} เอวํ ทสวิเธนสญฺญากฺขนฺโธ.}

\proseitem{63}{เอกวิเธน สญฺญากฺขนฺโธ\csromandash{}ผสฺสสมฺปยุตฺโต.}

\prose{ทุวิเธน สญฺญากฺขนฺโธ\csromandash{}อตฺถิ สเหตุโก, อตฺถิ อเหตุโก.}

\prose{ติวิเธน สญฺญากฺขนฺโธ\csromandash{}อตฺถิ สุขาย เวทนาย สมฺปยุตฺโต, อตฺถิ ทุกฺขาย เวทนาย สมฺปยุตฺโต, อตฺถิ อทุกฺขมสุขาย เวทนาย สมฺปยุตฺโต ฯเปฯ.\csromanspacer{} เอวํ ทสวิเธน สญฺญากฺขนฺโธ.}

\proseitem{64}{\csromanfolio{28}เอกวิเธน สญฺญากฺขนฺโธ\csromandash{}ผสฺสสมฺปยุตฺโต.}

\prose{ทุวิเธน สญฺญากฺขนฺโธ\csromandash{}อตฺถิ สเหตุโก, อตฺถิ อเหตุโก.}

\prose{ติวิเธน สญฺญากฺขนฺโธ\csromandash{}อตฺถิ วิปาโก, อตฺถิ วิปากธมฺมธมฺโม, อตฺถิ เนววิปากนวิปากธมฺมธมฺโม.\csromanspacer{} อตฺถิ อุปาทินฺนุปาทานิโย, อตฺถิ อนุปาทินฺนุปาทานิโย, อตฺถิ อนุปาทินฺน-อนุปาทานิโย.\csromanspacer{} อตฺถิ สํกิลิฏฺฐสํกิเลสิโก, อตฺถิ อสํกิลิฏฺฐสํกิเลสิโก, อตฺถิ อสํกิลิฏฺฐ- อสํกิเลสิโก.\csromanspacer{} อตฺถิ สวิตกฺกสวิจาโร, อตฺถิ อวิตกฺกวิจารมตฺโต, อตฺถิ อวิตกฺก-อวิจาโร.\csromanspacer{} อตฺถิ ปีติสหคโต, อตฺถิ สุขสหคโต, อตฺถิ อุเปกฺขาสหคโต.\csromanspacer{} อตฺถิ ทสฺสเนน ปหาตพฺโพ, อตฺถิ ภาวนาย ปหาตพฺโพ, อตฺถิ เนว ทสฺสเนน น ภาวนาย ปหาตพฺโพ.\csromanspacer{} อตฺถิ ทสฺสเนน ปหาตพฺพเหตุโก, อตฺถิ ภาวนาย ปหาตพฺพเหตุโก, อตฺถิ เนว ทสฺสเนน น ภาวนาย ปหาตพฺพเหตุโก.\csromanspacer{} อตฺถิ อาจยคามี, อตฺถิ อปจยคามี, อตฺถิ เนวาจยคามินาปจยคามี.\csromanspacer{} อตฺถิ เสกฺโข, อตฺถิ อเสกฺโข, อตฺถิ เนวเสกฺขนาเสกฺโข.\csromanspacer{} อตฺถิ ปริตฺโต, อตฺถิ มหคฺคโต, อตฺถิ อปฺปมาโณ.\csromanspacer{} อตฺถิ ปริตฺตารมฺมโณ, อตฺถิ มหคฺคตารมฺมโณ, อตฺถิ อปฺปมาณารมฺมโณ.\csromanspacer{} อตฺถิ หีโน, อตฺถิ มชฺฌิโม, อตฺถิ ปณีโต.\csromanspacer{} อตฺถิ มิจฺฉตฺตนิยโต, อตฺถิ สมฺมตฺตนิยโต, อตฺถิ อนิยโต.\csromanspacer{} อตฺถิ มคฺคารมฺมโณ, อตฺถิ มคฺคเหตุโก, อตฺถิ มคฺคาธิปติ.\csromanspacer{} อตฺถิ อุปฺปนฺโน, อตฺถิ อนุปฺปนฺโน, อตฺถิ อุปฺปาที.\csromanspacer{} อตฺถิ อตีโต, อตฺถิ อนาคโต, อตฺถิ ปจฺจุปฺปนฺโน.\csromanspacer{} อตฺถิ อตีตารมฺมโณ, อตฺถิ อนาคตารมฺมโณ, อตฺถิ ปจฺจุปฺปนฺนารมฺมโณ.\csromanspacer{} อตฺถิ อชฺฌตฺโต, อตฺถิ พหิทฺโธ, อตฺถิ อชฺฌตฺตพหิทฺโธ.\csromanspacer{} อตฺถิ อชฺฌตฺตารมฺมโณ, อตฺถิ พหิทฺธารมฺมโณ, อตฺถิ อชฺฌตฺตพหิทฺธารมฺมโณ ฯเปฯ.\csromanspacer{} เอวํ ทสวิเธน สญฺญากฺขนฺโธ.}

\proseitem{65}{เอกวิเธน สญฺญากฺขนฺโธ\csromandash{}ผสฺสสมฺปยุตฺโต.}

\prose{ทุวิเธน สญฺญากฺขนฺโธ\csromandash{}อตฺถิ เหตุสมฺปยุตฺโต, อตฺถิ เหตุวิปฺปยุตฺโต.\csromanspacer{} อตฺถิ น เหตุสเหตุโก, อตฺถิ น เหตุ-อเหตุโก.\csromanspacer{} อตฺถิ โลกิโย, อตฺถิ โลกุตฺตโร.\csromanspacer{} อตฺถิ เกนจิ วิญฺเญยฺโย, อตฺถิ เกนจิ น วิญฺเญยฺโย.\csromanspacer{} อตฺถิ สาสโว, อตฺถิ อนาสโว.\csromanspacer{} อตฺถิ อาสวสมฺปยุตฺโต, อตฺถิ อาสววิปฺปยุตฺโต.\csromanspacer{} อตฺถิ \csromanfolio{29}อาสววิปฺปยุตฺตสาสโว, อตฺถิ อาสววิปฺปยุตฺต-อนาสโว.\csromanspacer{} อตฺถิสํโยชนิโย, อตฺถิ อสํโยชนิโย.\csromanspacer{} อตฺถิ สํโยชนสมฺปยุตฺโต, อตฺถิ สํโยชนวิปฺปยุตฺโต.\csromanspacer{} อตฺถิ สํโยชนวิปฺปยุตฺตสํโยชนิโย, อตฺถิ สํโยชนวิปฺปยุตฺต- อสํโยชนิโย.\csromanspacer{} อตฺถิ คนฺถนิโย, อตฺถิ อคนฺถนิโย.\csromanspacer{} อตฺถิ คนฺถสมฺปยุตฺโต, อตฺถิ คนฺถวิปฺปยุตฺโต.\csromanspacer{} อตฺถิ คนฺถวิปฺปยุตฺตคนฺถนิโย, อตฺถิ คนฺถวิปฺปยุตฺต-อคนฺถนิโย.\csromanspacer{} อตฺถิ โอฆนิโย, อตฺถิ อโนฆนิโย.\csromanspacer{} อตฺถิ โอฆสมฺปยุตฺโต, อตฺถิ โอฆวิปฺปยุตฺโต.\csromanspacer{} อตฺถิ โอฆวิปฺปยุตฺต-โอฆนิโย, อตฺถิ โอฆวิปฺปยุตฺต-อโนฆนิโย.\csromanspacer{} อตฺถิ โยคนิโย, อตฺถิ อโยคนิโย.\csromanspacer{} อตฺถิ โยคสมฺปยุตฺโต, อตฺถิ โยควิปฺปยุตฺโต.\csromanspacer{} อตฺถิ โยควิปฺปยุตฺตโยคนิโย, อตฺถิ โยควิปฺปยุตฺต-อโยคนิโย.\csromanspacer{} อตฺถิ นีวรณิโย, อตฺถิ อนีวรณิโย.\csromanspacer{} อตฺถิ นีวรณสมฺปยุตฺโต, อตฺถิ นีวรณวิปฺปยุตฺโต.\csromanspacer{} อตฺถิ นีวรณวิปฺปยุตฺตนีวรณิโย, อตฺถิ นีวรณวิปฺปยุตฺต-อนีวรณิโย.\csromanspacer{} อตฺถิ ปรามฏฺโฐ, อตฺถิ อปรามฏฺโฐ.\csromanspacer{} อตฺถิ ปรามาสสมฺปยุตฺโต, อตฺถิ ปรามาสวิปฺปยุตฺโต.\csromanspacer{} อตฺถิ ปรามาสวิปฺปยุตฺตปรามฏฺโฐ, อตฺถิ ปรามาสวิปฺปยุตฺต-อปรามฏฺโฐ.\csromanspacer{} อตฺถิ อุปาทินฺโน, อตฺถิ อนุปาทินฺโน.\csromanspacer{} อตฺถิ อุปาทานิโย, อตฺถิ อนุปาทานิโย.\csromanspacer{} อตฺถิ อุปาทานสมฺปยุตฺโต, อตฺถิ อุปาทานวิปฺปยุตฺโต.\csromanspacer{} อตฺถิ อุปาทานวิปฺปยุตฺต-อุปาทานิโย, อตฺถิ อุปาทานวิปฺปยุตฺต-อนุปาทานิโย.\csromanspacer{} อตฺถิ สํกิเลสิโก, อตฺถิ อสํกิเลสิโก.\csromanspacer{} อตฺถิ สํกิลิฏฺโฐ, อตฺถิ อสํกิลิฏฺโฐ.\csromanspacer{} อตฺถิ กิเลสสมฺปยุตฺโต, อตฺถิ กิเลสวิปฺปยุตฺโต.\csromanspacer{} อตฺถิ กิเลสวิปฺปยุตฺตสํกิเลสิโก, อตฺถิ กิเลสวิปฺปยุตฺต- อสํกิเลสิโก.\csromanspacer{} อตฺถิ ทสฺสเนน ปหาตพฺโพ, อตฺถิ น ทสฺสเนน ปหาตพฺโพ.\csromanspacer{} อตฺถิ ภาวนาย ปหาตพฺโพ, อตฺถิ น ภาวนาย ปหาตพฺโพ.\csromanspacer{} อตฺถิ ทสฺสเนน ปหาตพฺพเหตุโก, อตฺถิ น ทสฺสเนน ปหาตพฺพเหตุโก.\csromanspacer{} อตฺถิ ภาวนาย ปหาตพฺพเหตุโก, อตฺถิ น ภาวนาย ปหาตพฺพเหตุโก.\csromanspacer{} อตฺถิ สวิตกฺโก, อตฺถิ อวิตกฺโก.\csromanspacer{} อตฺถิ สวิจาโร, อตฺถิ อวิจาโร.\csromanspacer{} อตฺถิ สปฺปีติโก, อตฺถิ อปฺปีติโก.\csromanspacer{} อตฺถิ ปีติสหคโต, อตฺถิ น ปีติสหคโต.\csromanspacer{} อตฺถิ สุขสหคโต, อตฺถิ น สุขสหคโต.\csromanspacer{} อตฺถิ อุเปกฺขาสหคโต, อตฺถิ น อุเปกฺขาสหคโต.\csromanspacer{} อตฺถิ กามาวจโร, อตฺถิ น กามาวจโร.\csromanspacer{} อตฺถิ รูปาวจโร, อตฺถิ น รูปาวจโร.\csromanspacer{} อตฺถิ อรูปาวจโร, อตฺถิ น อรูปาวจโร.\csromanspacer{} อตฺถิ \csromanfolio{30}ปริยาปนฺโน, อตฺถิ อปริยาปนฺโน.\csromanspacer{} อตฺถิ นิยฺยานิโก, อตฺถิ อนิยฺยานิโก.\csromanspacer{} อตฺถิ นิยโต, อตฺถิ อนิยโต.\csromanspacer{} อตฺถิ ส-อุตฺตโร, อตฺถิ อนุตฺตโร.\csromanspacer{} อตฺถิ สรโณ, อตฺถิ อรโณ.}

\prose{ติวิเธน สญฺญากฺขนฺโธ\csromandash{}อตฺถิ กุสโล, อตฺถิ อกุสโล, อตฺถิ อพฺยากโต ฯเปฯ.\csromanspacer{} เอวํ ทสวิเธน สญฺญากฺขนฺโธ.}

\proseitem{66}{เอกวิเธน สญฺญากฺขนฺโธ\csromandash{}ผสฺสสมฺปยุตฺโต.}

\prose{ทุวิเธน สญฺญากฺขนฺโธ\csromandash{}อตฺถิ สรโณ, อตฺถิ อรโณ.}

\prose{ติวิเธน สญฺญากฺขนฺโธ\csromandash{}อตฺถิ สุขาย เวทนาย สมฺปยุตฺโต, อตฺถิ ทุกฺขาย เวทนาย สมฺปยุตฺโต, อตฺถิ อทุกฺขมสุขาย เวทนาย สมฺปยุตฺโต.\csromanspacer{} อตฺถิ วิปาโก ฯเปฯ อตฺถิ อชฺฌตฺตารมฺมโณ, อตฺถิ พหิทฺธารมฺมโณ, อตฺถิ อชฺฌตฺตพหิทฺธารมฺมโณ ฯเปฯ.\csromanspacer{} เอวํ ทสวิเธน สญฺญากฺขนฺโธ. (ยถา กุสลตฺติเก วิตฺถาโร, เอวํ สพฺเพปิ ติกา วิตฺถาเรตพฺพา.) ทุกมูลกํ.}

\proseitem{67}{เอกวิเธน สญฺญากฺขนฺโธ\csromandash{}ผสฺสสมฺปยุตฺโต.}

\prose{ทุวิเธน สญฺญากฺขนฺโธ\csromandash{}อตฺถิ สเหตุโก, อตฺถิ อเหตุโก.}

\prose{ติวิเธน สญฺญากฺขนฺโธ\csromandash{}อตฺถิ กุสโล, อตฺถิ อกุสโล, อตฺถิ อพฺยากโต ฯเปฯ.\csromanspacer{} เอวํ ทสวิเธน สญฺญากฺขนฺโธ.}

\proseitem{68}{เอกวิเธน สญฺญากฺขนฺโธ\csromandash{}ผสฺสสมฺปยุตฺโต.}

\prose{ทุวิเธน สญฺญากฺขนฺโธ\csromandash{}อตฺถิ เหตุสมฺปยุตฺโต, อตฺถิ เหตุวิปฺปยุตฺโต ฯเปฯ อตฺถิ สรโณ, อตฺถิ อรโณ.}

\prose{ติวิเธน สญฺญากฺขนฺโธ\csromandash{}อตฺถิ กุสโล, อตฺถิ อกุสโล, อตฺถิ อพฺยากโต ฯเปฯ.\csromanspacer{} เอวํ ทสวิเธน สญฺญากฺขนฺโธ.}

\proseitem{69}{เอกวิเธน สญฺญากฺขนฺโธ\csromandash{}ผสฺสสมฺปยุตฺโต.}

\prose{ทุวิเธน สญฺญากฺขนฺโธ\csromandash{}อตฺถิ สเหตุโก, อตฺถิ อเหตุโก.}

\prose{ติวิเธน สญฺญากฺขนฺโธ\csromandash{}อตฺถิ สุขาย เวทนาย สมฺปยุตฺโต, อตฺถิ ทุกฺขาย เวทนาย สมฺปยุตฺโต, อตฺถิ อทุกฺขมสุขาย เวทนาย \csromanfolio{31}สมฺปยุตฺโต.\csromanspacer{} อตฺถิ วิปาโก ฯเปฯ อตฺถิ อชฺฌตฺตารมฺมโณ, อตฺถิ พหิทฺธารมฺมโณ, อตฺถิ อชฺฌตฺตพหิทฺธารมฺมโณ ฯเปฯ.\csromanspacer{} เอวํ ทสวิเธน สญฺญากฺขนฺโธ.}

\proseitem{70}{เอกวิเธน สญฺญากฺขนฺโธ\csromandash{}ผสฺสสมฺปยุตฺโต.}

\prose{ทุวิเธน สญฺญากฺขนฺโธ\csromandash{}อตฺถิ เหตุสมฺปยุตฺโต, อตฺถิ เหตุวิปฺปยุตฺโต ฯเปฯ อตฺถิ สรโณ, อตฺถิ อรโณ.}

\prose{ติวิเธน สญฺญากฺขนฺโธ ฯเปฯ อตฺถิ อชฺฌตฺตารมฺมโณ, อตฺถิ พหิทฺธารมฺมโณ, อตฺถิ อชฺฌตฺตพหิทฺธารมฺมโณ ฯเปฯ.\csromanspacer{} เอวํ ทสวิเธน สญฺญากฺขนฺโธ.\csromanspacer{} ติกมูลกํ.}

\proseitem{71}{เอกวิเธน สญฺญากฺขนฺโธ\csromandash{}ผสฺสสมฺปยุตฺโต.}

\prose{ทุวิเธน สญฺญากฺขนฺโธ\csromandash{}อตฺถิ สเหตุโก, อตฺถิ อเหตุโก.}

\prose{ติวิเธน สญฺญากฺขนฺโธ\csromandash{}อตฺถิ กุสโล, อตฺถิ อกุสโล, อตฺถิ อพฺยากโต ฯเปฯ.\csromanspacer{} เอวํ ทสวิเธน สญฺญากฺขนฺโธ.}

\proseitem{72}{เอกวิเธน สญฺญากฺขนฺโธ\csromandash{}ผสฺสสมฺปยุตฺโต.}

\prose{ทุวิเธน สญฺญากฺขนฺโธ\csromandash{}อตฺถิ เหตุสมฺปยุตฺโต, อตฺถิ เหตุวิปฺปยุตฺโต.}

\prose{ติวิเธน สญฺญากฺขนฺโธ\csromandash{}อตฺถิ สุขาย เวทนาย สมฺปยุตฺโต, อตฺถิ ทุกฺขาย เวทนาย สมฺปยุตฺโต, อตฺถิ อทุกฺขมสุขาย เวทนาย สมฺปยุตฺโต ฯเปฯ.\csromanspacer{} เอวํ ทสวิเธน สญฺญากฺขนฺโธ.}

\proseitem{73}{เอกวิเธน สญฺญากฺขนฺโธ\csromandash{}ผสฺสสมฺปยุตฺโต.}

\prose{ทุวิเธน สญฺญากฺขนฺโธ\csromandash{}อตฺถิ น เหตุสเหตุโก, อตฺถิ น เหตุ- อเหตุโก.}

\prose{ติวิเธน สญฺญากฺขนฺโธ\csromandash{}อตฺถิ วิปาโก, อตฺถิ วิปากธมฺมธมฺโม, อตฺถิ เนววิปากนวิปากธมฺมธมฺโม ฯเปฯ.\csromanspacer{} เอวํ ทสวิเธน สญฺญากฺขนฺโธ.}

\proseitem{74}{\csromanfolio{32}เอกวิเธน สญฺญากฺขนฺโธ\csromandash{}ผสฺสสมฺปยุตฺโต.}

\prose{ทุวิเธน สญฺญากฺขนฺโธ\csromandash{}อตฺถิ โลกิโย, อตฺถิ โลกุตฺตโร.}

\prose{ติวิเธน สญฺญากฺขนฺโธ\csromandash{}อตฺถิ อุปาทินฺนุปาทานิโย, อตฺถิ อนุปาทินฺนุปาทานิโย, อตฺถิ อนุปาทินฺน-อนุปาทานิโย ฯเปฯ.\csromanspacer{} เอวํ ทสวิเธน สญฺญากฺขนฺโธ.}

\proseitem{75}{เอกวิเธน สญฺญากฺขนฺโธ\csromandash{}ผสฺสสมฺปยุตฺโต.}

\prose{ทุวิเธน สญฺญากฺขนฺโธ\csromandash{}อตฺถิ เกนจิ วิญฺเญยฺโย, อตฺถิ เกนจิ น วิญฺเญยฺโย.}

\prose{ติวิเธน สญฺญากฺขนฺโธ\csromandash{}อตฺถิ สํกิลิฏฺฐสํกิเลสิโก, อตฺถิ อสํกิลิฏฺฐสํกิเลสิโก, อตฺถิ อสํกิลิฏฺฐ-อสํกิเลสิโก ฯเปฯ.\csromanspacer{} เอวํ ทสวิเธน สญฺญากฺขนฺโธ.}

\proseitem{76}{เอกวิเธน สญฺญากฺขนฺโธ\csromandash{}ผสฺสสมฺปยุตฺโต.}

\prose{ทุวิเธน สญฺญากฺขนฺโธ\csromandash{}อตฺถิ สาสโว อตฺถิ อนาสโว.}

\prose{ติวิเธน สญฺญากฺขนฺโธ\csromandash{}อตฺถิ สวิตกฺกสวิจาโร, อตฺถิ อวิตกฺกวิจารมตฺโต, อตฺถิ อวิตกฺก-อวิจาโร ฯเปฯ.\csromanspacer{} เอวํ ทสวิเธน สญฺญากฺขนฺโธ.}

\proseitem{77}{เอกวิเธน สญฺญากฺขนฺโธ\csromandash{}ผสฺสสมฺปยุตฺโต.}

\prose{ทุวิเธน สญฺญากฺขนฺโธ\csromandash{}อตฺถิ อาสวสมฺปยุตฺโต, อตฺถิ อาสววิปฺปยุตฺโต.}

\prose{ติวิเธน สญฺญากฺขนฺโธ\csromandash{}อตฺถิ ปีติสหคโต, อตฺถิ สุขสหคโต, อตฺถิ อุเปกฺขาสหคโต ฯเปฯ.\csromanspacer{} เอวํ ทสวิเธน สญฺญากฺขนฺโธ.}

\proseitem{78}{เอกวิเธน สญฺญากฺขนฺโธ\csromandash{}ผสฺสสมฺปยุตฺโต.}

\prose{ทุวิเธน สญฺญากฺขนฺโธ\csromandash{}อตฺถิ อาสววิปฺปยุตฺตสาสโว, อตฺถิ อาสววิปฺปยุตฺต-อนาสโว.}

\prose{ติวิเธน สญฺญากฺขนฺโธ\csromandash{}อตฺถิ ทสฺสเนน ปหาตพฺโพ, อตฺถิ ภาวนาย ปหาตพฺโพ, อตฺถิ เนว ทสฺสเนน น ภาวนาย ปหาตพฺโพ ฯเปฯ.\csromanspacer{} เอวํ ทสวิเธน สญฺญากฺขนฺโธ.}

\proseitem{79}{\csromanfolio{33}เอกวิเธน เวทนากฺขนฺโธ\csromandash{}ผสฺสสมฺปยุตฺโต.}

\prose{ทุวิเธน สญฺญากฺขนฺโธ\csromandash{}อตฺถิ สํโยชนิโย, อตฺถิ อสํโยชนิโย.}

\prose{ติวิเธน สญฺญากฺขนฺโธ\csromandash{}อตฺถิ ทสฺสเนน ปหาตพฺพเหตุโก, อตฺถิ ภาวนาย ปหาตพฺพเหตุโก, อตฺถิ เนว ทสฺสเนน น ภาวนาย ปหาตพฺพเหตุโก ฯเปฯ.\csromanspacer{} เอวํ ทสวิเธน สญฺญากฺขนฺโธ.}

\proseitem{80}{เอกวิเธน สญฺญากฺขนฺโธ\csromandash{}ผสฺสสมฺปยุตฺโต.}

\prose{ทุวิเธน สญฺญากฺขนฺโธ\csromandash{}อตฺถิ สํโยชนสมฺปยุตฺโต, อตฺถิ สํโยชนวิปฺปยุตฺโต.}

\prose{ติวิเธน สญฺญากฺขนฺโธ\csromandash{}อตฺถิ อาจยคามี, อตฺถิ อปจยคามี, อตฺถิ เนวาจยคามินาปจยคามี ฯเปฯ.\csromanspacer{} เอวํ ทสวิเธน สญฺญากฺขนฺโธ.}

\proseitem{81}{เอกวิเธน สญฺญากฺขนฺโธ\csromandash{}ผสฺสสมฺปยุตฺโต.}

\prose{ทุวิเธน สญฺญากฺขนฺโธ\csromandash{}อตฺถิ สํโยชนวิปฺปยุตฺตสํโยชนิโย, อตฺถิ สํโยชนวิปฺปยุตฺต-อสํโยชนิโย.}

\prose{ติวิเธน สญฺญากฺขนฺโธ\csromandash{}อตฺถิ เสกฺโข, อตฺถิ อเสกฺโข, อตฺถิ เนวเสกฺขนาเสกฺโข ฯเปฯ.\csromanspacer{} เอวํ ทสวิเธน สญฺญากฺขนฺโธ.}

\proseitem{82}{เอกวิเธน สญฺญากฺขนฺโธ\csromandash{}ผสฺสสมฺปยุตฺโต.}

\prose{ทุวิเธน สญฺญากฺขนฺโธ\csromandash{}อตฺถิ คนฺถนิโย, อตฺถิ อคนฺถนิโย.}

\prose{ติวิเธน สญฺญากฺขนฺโธ\csromandash{}อตฺถิ ปริตฺโต, อตฺถิ มหคฺคโต, อตฺถิ อปฺปมาโณ ฯเปฯ.\csromanspacer{} เอวํ ทสวิเธน สญฺญากฺขนฺโธ.}

\proseitem{83}{เอกวิเธน สญฺญากฺขนฺโธ\csromandash{}ผสฺสสมฺปยุตฺโต.}

\prose{ทุวิเธน สญฺญากฺขนฺโธ\csromandash{}อตฺถิ คนฺถสมฺปยุตฺโต, อตฺถิ คนฺถวิปฺปยุตฺโต.}

\prose{ติวิเธน สญฺญากฺขนฺโธ\csromandash{}อตฺถิ ปริตฺตารมฺมโณ, อตฺถิ มหคฺคตารมฺมโณ, อตฺถิ อปฺปมาณารมฺมโณ ฯเปฯ.\csromanspacer{} เอวํ ทสวิเธน สญฺญากฺขนฺโธ.}

\proseitem{84}{\csromanfolio{34}เอกวิเธน สญฺญากฺขนฺโธ\csromandash{}ผสฺสสมฺปยุตฺโต.}

\prose{ทุวิเธน สญฺญากฺขนฺโธ\csromandash{}อตฺถิ คนฺถวิปฺปยุตฺตคนฺถนิโย, อตฺถิ คนฺถวิปฺปยุตฺต-อคนฺถนิโย.}

\prose{ติวิเธน สญฺญากฺขนฺโธ\csromandash{}อตฺถิ หีโน, อตฺถิ มชฺฌิโม, อตฺถิ ปณีโต ฯเปฯ.\csromanspacer{} เอวํ ทสวิเธน สญฺญากฺขนฺโธ.}

\proseitem{85}{เอกวิเธน สญฺญากฺขนฺโธ\csromandash{}ผสฺสสมฺปยุตฺโต.}

\prose{ทุวิเธน สญฺญากฺขนฺโธ\csromandash{}อตฺถิ โอฆนิโย, อตฺถิ อโนฆนิโย.}

\prose{ติวิเธน สญฺญากฺขนฺโธ\csromandash{}อตฺถิ มิจฺฉตฺตนิยโต, อตฺถิ สมฺมตฺตนิยโต, อตฺถิ อนิยโต ฯเปฯ.\csromanspacer{} เอวํ ทสวิเธน สญฺญากฺขนฺโธ.}

\proseitem{86}{เอกวิเธน สญฺญากฺขนฺโธ\csromandash{}ผสฺสสมฺปยุตฺโต.}

\prose{ทุวิเธน สญฺญากฺขนฺโธ\csromandash{}อตฺถิ โอฆสมฺปยุตฺโต, อตฺถิ โอฆวิปฺปยุตฺโต.}

\prose{ติวิเธน สญฺญากฺขนฺโธ\csromandash{}อตฺถิ มคฺคารมฺมโณ, อตฺถิ มคฺคเหตุโก, อตฺถิ มคฺคาธิปติ ฯเปฯ.\csromanspacer{} เอวํ ทสวิเธน สญฺญากฺขนฺโธ.}

\proseitem{87}{เอกวิเธน สญฺญากฺขนฺโธ\csromandash{}ผสฺสสมฺปยุตฺโต.}

\prose{ทุวิเธน สญฺญากฺขนฺโธ\csromandash{}อตฺถิ โอฆวิปฺปยุตฺต-โอฆนิโย, อตฺถิ โอฆวิปฺปยุตฺต-อโนฆนิโย.}

\prose{ติวิเธน สญฺญากฺขนฺโธ\csromandash{}อตฺถิ อุปฺปนฺโน, อตฺถิ อนุปฺปนฺโน, อตฺถิ อุปฺปาที ฯเปฯ.\csromanspacer{} เอวํ ทสวิเธน สญฺญากฺขนฺโธ.}

\proseitem{88}{เอกวิเธน สญฺญากฺขนฺโธ\csromandash{}ผสฺสสมฺปยุตฺโต.}

\prose{ทุวิเธน สญฺญากฺขนฺโธ\csromandash{}อตฺถิ โยคนิโย, อตฺถิ อโยคนิโย.}

\prose{ติวิเธน สญฺญากฺขนฺโธ\csromandash{}อตฺถิ อตีโต, อตฺถิ อนาคโต, อตฺถิ ปจฺจุปฺปนฺโน ฯเปฯ.\csromanspacer{} เอวํ ทสวิเธน สญฺญากฺขนฺโธ.}

\proseitem{89}{เอกวิเธน สญฺญากฺขนฺโธ\csromandash{}ผสฺสสมฺปยุตฺโต.}

\prose{ทุวิเธน สญฺญากฺขนฺโธ\csromandash{}อตฺถิ โยคสมฺปยุตฺโต, อตฺถิ โยควิปฺปยุตฺโต.}

\prose{\csromanfolio{35}ติวิเธน สญฺญากฺขนฺโธ\csromandash{}อตฺถิ อตีตารมฺมโณ, อตฺถิ อนาคตารมฺมโณ, อตฺถิ ปจฺจุปฺปนฺนารมฺมโณ ฯเปฯ.\csromanspacer{} เอวํ ทสวิเธน สญฺญากฺขนฺโธ.}

\proseitem{90}{เอกวิเธน สญฺญากฺขนฺโธ\csromandash{}ผสฺสสมฺปยุตฺโต.}

\prose{ทุวิเธน สญฺญากฺขนฺโธ\csromandash{}อตฺถิ โยควิปฺปยุตฺตโยคนิโย, อตฺถิ โยควิปฺปยุตฺต-อโยคนิโย.}

\prose{ติวิเธน สญฺญากฺขนฺโธ\csromandash{}อตฺถิ อชฺฌตฺโต, อตฺถิ พหิทฺโธ, อตฺถิ อชฺฌตฺตพหิทฺโธ ฯเปฯ.\csromanspacer{} เอวํ ทสวิเธน สญฺญากฺขนฺโธ.}

\proseitem{91}{เอกวิเธน สญฺญากฺขนฺโธ\csromandash{}ผสฺสสมฺปยุตฺโต.}

\prose{ทุวิเธน สญฺญากฺขนฺโธ\csromandash{}อตฺถิ นีวรณิโย, อตฺถิ อนีวรณิโย.}

\prose{ติวิเธน สญฺญากฺขนฺโธ\csromandash{}อตฺถิ อชฺฌตฺตารมฺมโณ, อตฺถิ พหิทฺธารมฺมโณ, อตฺถิ อชฺฌตฺตพหิทฺธารมฺมโณ ฯเปฯ.\csromanspacer{} เอวํ ทสวิเธน สญฺญากฺขนฺโธ.\csromanspacer{} อุภโตวฑฺฒกํ.}

\prose{สตฺตวิเธน สญฺญากฺขนฺโธ\csromandash{}อตฺถิ กุสโล, อตฺถิ อกุสโล, อตฺถิ อพฺยากโต, อตฺถิ กามาวจโร, อตฺถิ รูปาวจโร, อตฺถิ อรูปาวจโร, อตฺถิ อปริยาปนฺโน.\csromanspacer{} เอวํ สตฺตวิเธน สญฺญากฺขนฺโธ.}

\prose{อปโรปิ สตฺตวิเธน สญฺญากฺขนฺโธ\csromandash{}อตฺถิ สุขาย เวทนาย สมฺปยุตฺโต, อตฺถิ ทุกฺขาย เวทนาย สมฺปยุตฺโต, อตฺถิ อทุกฺขมสุขาย เวทนาย สมฺปยุตฺโต, อตฺถิ กามาวจโร, อตฺถิ รูปาวจโร, อตฺถิ อรูปาวจโร, อตฺถิ อปริยาปนฺโน ฯเปฯ อตฺถิ อชฺฌตฺตารมฺมโณ, อตฺถิ พหิทฺธารมฺมโณ, อตฺถิ อชฺฌตฺตพหิทฺธารมฺมโณ, อตฺถิ กามาวจโร, อตฺถิ รูปาวจโร, อตฺถิ อรูปาวจโร, อตฺถิ อปริยาปนฺโน.\csromanspacer{} เอวํ สตฺตวิเธน สญฺญากฺขนฺโธ.}

\prose{จตุวีสติวิเธน สญฺญากฺขนฺโธ\csromandash{}จกฺขุสมฺผสฺสปจฺจยา สญฺญากฺขนฺโธ อตฺถิ กุสโล, อตฺถิ อกุสโล, อตฺถิ อพฺยากโต.\csromanspacer{} โสตสมฺผสฺสปจฺจยา ฯเปฯ.\csromanspacer{} ฆานสมฺผสฺสปจฺจยา ฯเปฯ.\csromanspacer{} ชิวฺหาสมฺผสฺสปจฺจยา ฯเปฯ.\csromanspacer{} กายสมฺผสฺสปจฺจยา ฯเปฯ.\csromanspacer{} มโนสมฺผสฺสปจฺจยา สญฺญากฺขนฺโธ อตฺถิ กุสโล, \csromanfolio{36}อตฺถิ อกุสโล, อตฺถิ อพฺยากโต, จกฺขุสมฺผสฺสชา สญฺญา, โสตสมฺผสฺสชา สญฺญา, ฆานสมฺผสฺสชา สญฺญา, ชิวฺหาสมฺผสฺสชา สญฺญา, กายสมฺผสฺสชา สญฺญา, มโนสมฺผสฺสชา สญฺญา.\csromanspacer{} เอวํ จตุวีสติวิเธน สญฺญากฺขนฺโธ.}

\prose{อปโรปิ จตุวีสติวิเธน สญฺญากฺขนฺโธ\csromandash{}จกฺขุสมฺผสฺสปจฺจยา สญฺญากฺขนฺโธ อตฺถิ สุขาย เวทนาย สมฺปยุตฺโต ฯเปฯ อตฺถิ อชฺฌตฺตารมฺมโณ, อตฺถิ พหิทฺธารมฺมโณ, อตฺถิ อชฺฌตฺตพหิทฺธารมฺมโณ, จกฺขุสมฺผสฺสชา สญฺญา ฯเปฯ มโนสมฺผสฺสชา สญฺญา.\csromanspacer{} โสตสมฺผสฺสปจฺจยา ฯเปฯ.\csromanspacer{} ฆานสมฺผสฺสปจฺจยา ฯเปฯ.\csromanspacer{} ชิวฺหาสมฺผสฺสปจฺจยา ฯเปฯ.\csromanspacer{} กายสมฺผสฺสปจฺจยา ฯเปฯ.\csromanspacer{} มโนสมฺผสฺสปจฺจยา สญฺญากฺขนฺโธ อตฺถิ สุขาย เวทนาย สมฺปยุตฺโต ฯเปฯ อตฺถิ อชฺฌตฺตารมฺมโณ, อตฺถิ พหิทฺธารมฺมโณ, อตฺถิ อชฺฌตฺตพหิทฺธารมฺมโณ, จกฺขุสมฺผสฺสชา สญฺญา ฯเปฯ มโนสมฺผสฺสชา สญฺญา.\csromanspacer{} เอวํ จตุวีสติวิเธน สญฺญากฺขนฺโธ.}

\prose{ติํสติวิเธน สญฺญากฺขนฺโธ\csromandash{}จกฺขุสมฺผสฺสปจฺจยา สญฺญากฺขนฺโธ อตฺถิ กามาวจโร, อตฺถิ รูปาวจโร, อตฺถิ อรูปาวจโร, อตฺถิ อปริยาปนฺโน.\csromanspacer{} โสตสมฺผสฺสปจฺจยา ฯเปฯ.\csromanspacer{} ฆานสมฺผสฺสปจฺจยา ฯเปฯ.\csromanspacer{} ชิวฺหาสมฺผสฺสปจฺจยา ฯเปฯ กายสมฺผสฺสปจฺจยา ฯเปฯ.\csromanspacer{} มโนสมฺผสฺสปจฺจยา สญฺญากฺขนฺโธ อตฺถิ กามาวจโร, อตฺถิ รูปาวจโร, อตฺถิ อรูปาวจโร, อตฺถิ อปริยาปนฺโน.\csromanspacer{} จกฺขุสมฺผสฺสชา สญฺญา ฯเปฯ มโนสมฺผสฺสชา สญฺญา.\csromanspacer{} เอวํ ติํสติวิเธน สญฺญากฺขนฺโธ.}

\prose{พหุวิเธน สญฺญากฺขนฺโธ\csromandash{}จกฺขุสมฺผสฺสปจฺจยา สญฺญากฺขนฺโธ อตฺถิ กุสโล, อตฺถิ อกุสโล, อตฺถิ อพฺยากโต, อตฺถิ กามาวจโร, อตฺถิ รูปาวจโร, อตฺถิ อรูปาวจโร, อตฺถิ อปริยาปนฺโน.\csromanspacer{} โสตสมฺผสฺสปจฺจยา ฯเปฯ.\csromanspacer{} ฆานสมฺผสฺสปจฺจยา ฯเปฯ.\csromanspacer{} ชิวฺหาสมฺผสฺสปจฺจยา ฯเปฯ.\csromanspacer{} กายสมฺผสฺสปจฺจยา ฯเปฯ.\csromanspacer{} มโนสมฺผสฺสปจฺจยา สญฺญากฺขนฺโธ อตฺถิ กุสโล, อตฺถิ อกุสโล, อตฺถิ อพฺยากโต, อตฺถิ กามาวจโร, อตฺถิ รูปาวจโร, อตฺถิ อรูปาวจโร, อตฺถิ อปริยาปนฺโน, จกฺขุสมฺผสฺสชา สญฺญา ฯเปฯ มโนสมฺผสฺสชา สญฺญา.\csromanspacer{} เอวํ พหุวิเธน สญฺญากฺขนฺโธ.}

\prose{\csromanfolio{37}อปโรปิ พหุวิเธน สญฺญากฺขนฺโธ\csromandash{}จกฺขุสมฺผสฺสปจฺจยา สญฺญากฺขนฺโธ อตฺถิ สุขาย เวทนาย สมฺปยุตฺโต ฯเปฯ อตฺถิ อชฺฌตฺตารมฺมโณ, อตฺถิ พหิทฺธารมฺมโณ, อตฺถิ อชฺฌตฺตพหิทฺธารมฺมโณ, อตฺถิ กามาวจโร, อตฺถิ รูปาวจโร, อตฺถิ อรูปาวจโร, อตฺถิ อปริยาปนฺโน.\csromanspacer{} โสตสมฺผสฺสปจฺจยา ฯเปฯ.\csromanspacer{} ฆานสมฺผสฺสปจฺจยา ฯเปฯ.\csromanspacer{} ชิวฺหาสมฺผสฺสปจฺจยา ฯเปฯ.\csromanspacer{} กายสมฺผสฺสปจฺจยา ฯเปฯ.\csromanspacer{} มโนสมฺผสฺสปจฺจยา สญฺญากฺขนฺโธ อตฺถิ อชฺฌตฺตารมฺมโณ, อตฺถิ พหิทฺธารมฺมโณ, อตฺถิ อชฺฌตฺตพหิทฺธารมฺมโณ, อตฺถิ กามาวจโร, อตฺถิ รูปาวจโร, อตฺถิ อรูปาวจโร, อตฺถิ อปริยาปนฺโน.\csromanspacer{} จกฺขุสมฺผสฺสชา สญฺญา, โสตสมฺผสฺสชา สญฺญา, ฆานสมฺผสฺสชา สญฺญา, ชิวฺหาสมฺผสฺสชา สญฺญา, กายสมฺผสฺสชา สญฺญา, มโนสมฺผสฺสชา สญฺญา.\csromanspacer{} เอวํ พหุวิเธน สญฺญากฺขนฺโธ.}

\csromanheader{2}{อยํ วุจฺจติ สญฺญากฺขนฺโธ.}
\csromansectionrule
\csromanmatikaanchor{mtk.12}
\csromantocmarkref{subsection}{4. สงฺขารกฺขนฺธ}{mtk.12}
\bookmark[dest={mtk.12},level=2]{4. สงฺขารกฺขนฺธ}
\csromanheader{2}{4. สงฺขารกฺขนฺธ}
\proseitem{92}{ตตฺถ กตโม สงฺขารกฺขนฺโธ, เอกวิเธน สงฺขารกฺขนฺโธ\csromandash{} จิตฺตสมฺปยุตฺโต.}

\prose{ทุวิเธน สงฺขารกฺขนฺโธ\csromandash{}อตฺถิ เหตุ, อตฺถิ น เหตุ.}

\prose{ติวิเธน สงฺขารกฺขนฺโธ\csromandash{}อตฺถิ กุสโล, อตฺถิ อกุสโล, อตฺถิ อพฺยากโต.}

\prose{จตุพฺพิเธน สงฺขารกฺขนฺโธ\csromandash{}อตฺถิ กามาวจโร, อตฺถิ รูปาวจโร, อตฺถิ อรูปาวจโร, อตฺถิ อปริยาปนฺโน.}

\prose{ปญฺจวิเธน สงฺขารกฺขนฺโธ\csromandash{}อตฺถิ สุขินฺทฺริยสมฺปยุตฺโต, อตฺถิ ทุกฺขินฺทฺริยสมฺปยุตฺโต, อตฺถิ โสมนสฺสินฺทฺริยสมฺปยุตฺโต, อตฺถิ โทมนสฺสินฺทฺริยสมฺปยุตฺโต, อตฺถิ อุเปกฺขินฺทฺริยสมฺปยุตฺโต.\csromanspacer{} เอวํ ปญฺจวิเธน สงฺขารกฺขนฺโธ.}

\prose{ฉพฺพิเธน สงฺขารกฺขนฺโธ\csromandash{}จกฺขุสมฺผสฺสชา เจตนา, โสตสมฺผสฺสชา เจตนา, ฆานสมฺผสฺสชา เจตนา, ชิวฺหาสมฺผสฺสชา เจตนา, กายสมฺผสฺสชา เจตนา, มโนสมฺผสฺสชา เจตนา.\csromanspacer{} เอวํ ฉพฺพิเธน สงฺขารกฺขนฺโธ.}

\prose{\csromanfolio{38}สตฺตวิเธน สงฺขารกฺขนฺโธ\csromandash{}จกฺขุสมฺผสฺสชา เจตนา, โสตสมฺผสฺสชา เจตนา, ฆานสมฺผสฺสชา เจตนา, ชิวฺหาสมฺผสฺสชา เจตนา, กายสมฺผสฺสชา เจตนา, มโนธาตุสมฺผสฺสชา เจตนา, มโนวิญฺญาณธาตุสมฺผสฺสชา เจตนา.\csromanspacer{} เอวํ สตฺตวิเธน สงฺขารกฺขนฺโธ.}

\prose{อฏฺฐวิเธน สงฺขารกฺขนฺโธ\csromandash{}จกฺขุสมฺผสฺสชา เจตนา ฯเปฯ กายสมฺผสฺสชา เจตนา อตฺถิ สุขสหคตา, อตฺถิ ทุกฺขสหคตา, มโนธาตุสมฺผสฺสชา เจตนา, มโนวิญฺญาณธาตุสมฺผสฺสชา เจตนา.\csromanspacer{} เอวํ อฏฺฐวิเธน สงฺขารกฺขนฺโธ.}

\prose{นววิเธน สงฺขารกฺขนฺโธ\csromandash{}จกฺขุสมฺผสฺสชา เจตนา ฯเปฯ มโนธาตุสมฺผสฺสชา เจตนา, มโนวิญฺญาณธาตุสมฺผสฺสชา เจตนา อตฺถิ กุสลา, อตฺถิ อกุสลา, อตฺถิ อพฺยากตา.\csromanspacer{} เอวํ นววิเธน สงฺขารกฺขนฺโธ.}

\prose{ทสวิเธน สงฺขารกฺขนฺโธ\csromandash{}จกฺขุสมฺผสฺสชา เจตนา ฯเปฯ กายสมฺผสฺสชา เจตนา อตฺถิ สุขสหคตา, อตฺถิ ทุกฺขสหคตา, มโนธาตุสมฺผสฺสชา เจตนา, มโนวิญฺญาณธาตุสมฺผสฺสชา เจตนา อตฺถิ กุสลา, อตฺถิ อกุสลา, อตฺถิ อพฺยากตา.\csromanspacer{} เอวํ ทสวิเธน สงฺขารกฺขนฺโธ.}

\proseitem{93}{เอกวิเธน สงฺขารกฺขนฺโธ\csromandash{}จิตฺตสมฺปยุตฺโต.}

\prose{ทุวิเธน สงฺขารกฺขนฺโธ\csromandash{}อตฺถิ เหตุ, อตฺถิ น เหตุ.}

\prose{ติวิเธน สงฺขารกฺขนฺโธ\csromandash{}อตฺถิ สุขาย เวทนาย สมฺปยุตฺโต, อตฺถิ ทุกฺขาย เวทนาย สมฺปยุตฺโต, อตฺถิ อทุกฺขมสุขาย เวทนาย สมฺปยุตฺโต.\csromanspacer{} อตฺถิ วิปาโก, อตฺถิ วิปากธมฺมธมฺโม, อตฺถิ เนววิปากนวิปากธมฺมธมฺโม.\csromanspacer{} อตฺถิ อุปาทินฺนุปาทานิโย, อตฺถิ อนุปาทินฺนุปาทานิโย, อตฺถิ อนุปาทินฺน-อนุปาทานิโย.\csromanspacer{} อตฺถิ สํกิลิฏฺฐสํกิเลสิโก, อตฺถิ อสํกิลิฏฺฐสํกิเลสิโก, อตฺถิ อสํกิลิฏฺฐ- อสํกิเลสิโก.\csromanspacer{} อตฺถิ สวิตกฺกสวิจาโร, อตฺถิ อวิตกฺกวิจารมตฺโต, อตฺถิ อวิตกฺก-อวิจาโร.\csromanspacer{} อตฺถิ ปีติสหคโต, อตฺถิ สุขสหคโต, อตฺถิ อุเปกฺขาสหคโต.\csromanspacer{} อตฺถิ ทสฺสเนน ปหาตพฺโพ, อตฺถิ ภาวนาย ปหาตพฺโพ, อตฺถิ เนว ทสฺสเนน น ภาวนาย ปหาตพฺโพ.\csromanspacer{} อตฺถิ \csromanfolio{39}ทสฺสเนน ปหาตพฺพเหตุโก, อตฺถิ ภาวนาย ปหาตพฺพเหตุโก, อตฺถิ เนว ทสฺสเนน น ภาวนาย ปหาตพฺพเหตุโก.\csromanspacer{} อตฺถิ อาจยคามี, อตฺถิ อปจยคามี, อตฺถิ เนวาจยคามินาปจยคามี.\csromanspacer{} อตฺถิ เสกฺโข, อตฺถิ อเสกฺโข, อตฺถิ เนวเสกฺขนาเสกฺโข.\csromanspacer{} อตฺถิ ปริตฺโต, อตฺถิ มหคฺคโต, อตฺถิ อปฺปมาโณ.\csromanspacer{} อตฺถิ ปริตฺตารมฺมโณ, อตฺถิ มหคฺคตารมฺมโณ, อตฺถิ อปฺปมาณารมฺมโณ.\csromanspacer{} อตฺถิ หีโน, อตฺถิ มชฺฌิโม, อตฺถิ ปณีโต.\csromanspacer{} อตฺถิ มิจฺฉตฺตนิยโต, อตฺถิ สมฺมตฺตนิยโต, อตฺถิ อนิยโต.\csromanspacer{} อตฺถิ มคฺคารมฺมโณ, อตฺถิ มคฺคเหตุโก, อตฺถิ มคฺคาธิปติ.\csromanspacer{} อตฺถิ อุปฺปนฺโน, อตฺถิ อนุปฺปนฺโน, อตฺถิ อุปฺปาที.\csromanspacer{} อตฺถิ อตีโต, อตฺถิ อนาคโต, อตฺถิ ปจฺจุปฺปนฺโน.\csromanspacer{} อตฺถิ อตีตารมฺมโณ, อตฺถิ อนาคตารมฺมโณ, อตฺถิ ปจฺจุปฺปนฺนารมฺมโณ.\csromanspacer{} อตฺถิ อชฺฌตฺโต, อตฺถิ พหิทฺโธ, อตฺถิ อชฺฌตฺตพหิทฺโธ.\csromanspacer{} อตฺถิ อชฺฌตฺตารมฺมโณ, อตฺถิ พหิทฺธารมฺมโณ, อตฺถิ อชฺฌตฺตพหิทฺธารมฺมโณ ฯเปฯ.\csromanspacer{} เอวํ ทสวิเธน สงฺขารกฺขนฺโธ.}

\proseitem{94}{เอกวิเธน สงฺขารกฺขนฺโธ\csromandash{}จิตฺตสมฺปยุตฺโต.}

\prose{ทุวิเธน สงฺขารกฺขนฺโธ\csromandash{}อตฺถิ สเหตุโก, อตฺถิ อเหตุโก.\csromanspacer{} อตฺถิ เหตุสมฺปยุตฺโต, อตฺถิ เหตุวิปฺปยุตฺโต.\csromanspacer{} อตฺถิ เหตุ เจว สเหตุโก จ, อตฺถิ สเหตุโก เจว น จ เหตุ.\csromanspacer{} อตฺถิ เหตุ เจว เหตุสมฺปยุตฺโต จ, อตฺถิ เหตุสมฺปยุตฺโต เจว น จ เหตุ.\csromanspacer{} อตฺถิ น เหตุ สเหตุโก, อตฺถิ น เหตุ อเหตุโก.\csromanspacer{} อตฺถิ โลกิโย, อตฺถิ โลกุตฺตโร.\csromanspacer{} อตฺถิ เกนจิ วิญฺเญยฺโย, อตฺถิ เกนจิ น วิญฺเญยฺโย.\csromanspacer{} อตฺถิ อาสโว, อตฺถิ โน อาสโว.\csromanspacer{} อตฺถิ สาสโว, อตฺถิ อนาสโว.\csromanspacer{} อตฺถิ อาสวสมฺปยุตฺโต, อตฺถิ อาสววิปฺปยุตฺโต.\csromanspacer{} อตฺถิ อาสโว เจว สาสโว จ, อตฺถิ สาสโว เจว โน จ อาสโว.\csromanspacer{} อตฺถิ อาสโว เจว อาสวสมฺปยุตฺโต จ, อตฺถิ อาสวสมฺปยุตฺโต เจว โน จ อาสโว.\csromanspacer{} อตฺถิ อาสววิปฺปยุตฺตสาสโว, อตฺถิ อาสววิปฺปยุตฺต-อนาสโว.\csromanspacer{} อตฺถิ สํโยชนํ, อตฺถิ โน สํโยชนํ.\csromanspacer{} อตฺถิ สํโยชนิโย, อตฺถิ อสํโยชนิโย.\csromanspacer{} อตฺถิ สํโยชนสมฺปยุตฺโต, อตฺถิ สํโยชนวิปฺปยุตฺโต.\csromanspacer{} อตฺถิ สํโยชนญฺเจว สํโยชนิโย จ, อตฺถิ สํโยชนิโย เจว โน จ สํโยชนํ.\csromanspacer{} อตฺถิ สํโยชนญฺเจว สํโยชนสมฺปยุตฺโต จ, อตฺถิ \csromanfolio{40}สํโยชนสมฺปยุตฺโต เจว โน จ สํโยชนํ.\csromanspacer{} อตฺถิ สํโยชนวิปฺปยุตฺตสํโยชนิโย, อตฺถิ สํโยชนวิปฺปยุตฺต-อสํโยชนิโย.}

\prose{อตฺถิ คนฺโถ, อตฺถิ โน คนฺโถ.\csromanspacer{} อตฺถิ คนฺถนิโย, อตฺถิ อคนฺถนิโย.\csromanspacer{} อตฺถิ คนฺถสมฺปยุตฺโต, อตฺถิ คนฺถวิปฺปยุตฺโต.\csromanspacer{} อตฺถิ คนฺโถ เจว คนฺถนิโย จ, อตฺถิ คนฺถนิโย เจว โน จ คนฺโถ.\csromanspacer{} อตฺถิ คนฺโถ เจว คนฺถสมฺปยุตฺโต จ, อตฺถิ คนฺถสมฺปยุตฺโต เจว โน จ คนฺโถ.\csromanspacer{} อตฺถิ คนฺถวิปฺปยุตฺตคนฺถนิโย, อตฺถิ คนฺถวิปฺปยุตฺต-อคนฺถนิโย.\csromanspacer{} อตฺถิ โอโฆ, อตฺถิ โน โอโฆ.\csromanspacer{} อตฺถิ โอฆนิโย, อตฺถิ อโนฆนิโย.\csromanspacer{} อตฺถิ โอฆสมฺปยุตฺโต, อตฺถิ โอฆวิปฺปยุตฺโต.\csromanspacer{} อตฺถิ โอโฆ เจว โอฆนิโย จ, อตฺถิ โอฆนิโย เจว โน จ โอโฆ.\csromanspacer{} อตฺถิ โอโฆ เจว โอฆสมฺปยุตฺโต จ, อตฺถิ โอฆสมฺปยุตฺโต เจว โน จ โอโฆ.\csromanspacer{} อตฺถิ โอฆวิปฺปยุตฺต-โอฆนิโย, อตฺถิ โอฆวิปฺปยุตฺต-อโนฆนิโย.\csromanspacer{} อตฺถิ โยโค, อตฺถิ โน โยโค.\csromanspacer{} อตฺถิ โยคนิโย, อตฺถิ อโยคนิโย.\csromanspacer{} อตฺถิ โยคสมฺปยุตฺโต, อตฺถิ โยควิปฺปยุตฺโต.\csromanspacer{} อตฺถิ โยโค เจว โยคนิโย จ, อตฺถิ โยคนิโย เจว โน จ โยโค.\csromanspacer{} อตฺถิ โยโค เจว โยคสมฺปยุตฺโต จ, อตฺถิ โยคสมฺปยุตฺโต เจว โน จ โยโค.\csromanspacer{} อตฺถิ โยควิปฺปยุตฺตโยคนิโย, อตฺถิ โยควิปฺปยุตฺต-อโยคนิโย.\csromanspacer{} อตฺถิ นีวรณํ, อตฺถิ โน นีวรณํ.\csromanspacer{} อตฺถิ นีวรณิโย, อตฺถิ อนีวรณิโย.\csromanspacer{} อตฺถิ นีวรณสมฺปยุตฺโต, อตฺถิ นีวรณวิปฺปยุตฺโต.\csromanspacer{} อตฺถิ นีวรณญฺเจว นีวรณิโย จ, อตฺถิ นีวรณิโย เจว โน จ นีวรณํ.\csromanspacer{} อตฺถิ นีวรณญฺเจว นีวรณสมฺปยุตฺโต จ, อตฺถิ นีวรณสมฺปยุตฺโต เจว โน จ นีวรณํ.\csromanspacer{} อตฺถิ นีวรณวิปฺปยุตฺตนีวรณิโย, อตฺถิ นีวรณวิปฺปยุตฺต-อนีวรณิโย.}

\prose{อตฺถิ ปรามาโส, อตฺถิ โน ปรามาโส.\csromanspacer{} อตฺถิ ปรามฏฺโฐ, อตฺถิ อปรามฏฺโฐ.\csromanspacer{} อตฺถิ ปรามาสสมฺปยุตฺโต, อตฺถิ ปรามาสวิปฺปยุตฺโต.\csromanspacer{} อตฺถิ ปรามาโส เจว ปรามฏฺโฐ จ, อตฺถิ ปรามฏฺโฐ เจว โน จ ปรามาโส.\csromanspacer{} อตฺถิ ปรามาสวิปฺปยุตฺตปรามฏฺโฐ, อตฺถิ ปรามาสวิปฺปยุตฺต-อปรามฏฺโฐ.\csromanspacer{} อตฺถิ อุปาทินฺโน, อตฺถิ อนุปาทินฺโน.\csromanspacer{} อตฺถิ อุปาทานํ, อตฺถิ โน อุปาทานํ.\csromanspacer{} อตฺถิ อุปาทานิโย, อตฺถิ อนุปาทานิโย.\csromanspacer{} อตฺถิ อุปาทานสมฺปยุตฺโต, อตฺถิ อุปาทานวิปฺปยุตฺโต.\csromanspacer{} อตฺถิ อุปาทานญฺเจว อุปาทานิโย จ, อตฺถิ อุปาทานิโย เจว โน จ อุปาทานํ.\csromanspacer{} อตฺถิ อุปาทานญฺเจว อุปาทานสมฺปยุตฺโต \csromanfolio{41}จ, อตฺถิ อุปาทานสมฺปยุตฺโต เจว โน จ อุปาทานํ.\csromanspacer{} อตฺถิ อุปาทานวิปฺปยุตฺต-อุปาทานิโย, อตฺถิ อุปาทานวิปฺปยุตฺต-อนุปาทานิโย.}

\prose{อตฺถิ กิเลโส, อตฺถิ โน กิเลโส.\csromanspacer{} อตฺถิ สํกิเลสิโก, อตฺถิ อสํกิเลสิโก.\csromanspacer{} อตฺถิ สํกิลิฏฺโฐ, อตฺถิ อสํกิลิฏฺโฐ.\csromanspacer{} อตฺถิ กิเลสสมฺปยุตฺโต, อตฺถิ กิเลสวิปฺปยุตฺโต.\csromanspacer{} อตฺถิ กิเลโส เจว สํกิเลสิโก จ, อตฺถิ สํกิเลสิโก เจว โน จ กิเลโส.\csromanspacer{} อตฺถิ กิเลโส เจว สํกิลิฏฺโฐ จ, อตฺถิ สํกิลิฏฺโฐ เจว โน จ กิเลโส.\csromanspacer{} อตฺถิ กิเลโส เจว กิเลสสมฺปยุตฺโต จ, อตฺถิ กิเลสสมฺปยุตฺโต เจว โน จ กิเลโส.\csromanspacer{} อตฺถิ กิเลสวิปฺปยุตฺตสํกิเลสิโก, อตฺถิ กิเลสวิปฺปยุตฺต- อสํกิเลสิโก.\csromanspacer{} อตฺถิ ทสฺสเนน ปหาตพฺโพ, อตฺถิ น ทสฺสเนน ปหาตพฺโพ.\csromanspacer{} อตฺถิ ภาวนาย ปหาตพฺโพ, อตฺถิ น ภาวนาย ปหาตพฺโพ.\csromanspacer{} อตฺถิ ทสฺสเนน ปหาตพฺพเหตุโก, อตฺถิ น ทสฺสเนน ปหาตพฺพเหตุโก.\csromanspacer{} อตฺถิ ภาวนาย ปหาตพฺพเหตุโก, อตฺถิ น ภาวนาย ปหาตพฺพเหตุโก.}

\prose{อตฺถิ สวิตกฺโก, อตฺถิ อวิตกฺโก.\csromanspacer{} อตฺถิ สวิจาโร, อตฺถิ อวิจาโร.\csromanspacer{} อตฺถิ สปฺปีติโก, อตฺถิ อปฺปีติโก.\csromanspacer{} อตฺถิ ปีติสหคโต, อตฺถิ น ปีติสหคโต.\csromanspacer{} อตฺถิ สุขสหคโต, อตฺถิ น สุขสหคโต.\csromanspacer{} อตฺถิ อุเปกฺขาสหคโต, อตฺถิ น อุเปกฺขาสหคโต.\csromanspacer{} อตฺถิ กามาวจโร, อตฺถิ น กามาวจโร.\csromanspacer{} อตฺถิ รูปาวจโร, อตฺถิ น รูปาวจโร.\csromanspacer{} อตฺถิ อรูปาวจโร, อตฺถิ น อรูปาวจโร.\csromanspacer{} อตฺถิ ปริยาปนฺโน, อตฺถิ อปริยาปนฺโน.\csromanspacer{} อตฺถิ นิยฺยานิโก, อตฺถิ อนิยฺยานิโก.\csromanspacer{} อตฺถิ นิยโต, อตฺถิ อนิยโต.\csromanspacer{} อตฺถิ ส-อุตฺตโร, อตฺถิ อนุตฺตโร.\csromanspacer{} อตฺถิ สรโณ, อตฺถิ อรโณ.}

\prose{ติวิเธน สงฺขารกฺขนฺโธ\csromandash{}อตฺถิ กุสโล, อตฺถิ อกุสโล, อตฺถิ อพฺยากโต ฯเปฯ.\csromanspacer{} เอวํ ทสวิเธน สงฺขารกฺขนฺโธ.}

\proseitem{95}{เอกวิเธน สงฺขารกฺขนฺโธ\csromandash{}จิตฺตสมฺปยุตฺโต.}

\prose{ทุวิเธน สงฺขารกฺขนฺโธ\csromandash{}อตฺถิ สรโณ, อตฺถิ อรโณ.}

\prose{\csromanfolio{42}ติวิเธน สงฺขารกฺขนฺโธ\csromandash{}อตฺถิ สุขาย เวทนาย สมฺปยุตฺโต ฯเปฯ อตฺถิ อชฺฌตฺตารมฺมโณ, อตฺถิ พหิทฺธารมฺมโณ, อตฺถิ อชฺฌตฺตพหิทฺธารมฺมโณ ฯเปฯ.\csromanspacer{} เอวํ ทสวิเธน สงฺขารกฺขนฺโธ.\csromanspacer{} ทุกมูลกํ.}

\proseitem{96}{เอกวิเธน สงฺขารกฺขนฺโธ\csromandash{}จิตฺตสมฺปยุตฺโต.}

\prose{ทุวิเธน สงฺขารกฺขนฺโธ\csromandash{}อตฺถิ เหตุ, อตฺถิ น เหตุ.}

\prose{ติวิเธน สงฺขารกฺขนฺโธ\csromandash{}อตฺถิ กุสโล, อตฺถิ อกุสโล, อตฺถิ อพฺยากโต ฯเปฯ.\csromanspacer{} เอวํ ทสวิเธน สงฺขารกฺขนฺโธ.}

\proseitem{97}{เอกวิเธน สงฺขารกฺขนฺโธ\csromandash{}จิตฺตสมฺปยุตฺโต.}

\prose{ทุวิเธน สงฺขารกฺขนฺโธ\csromandash{}อตฺถิ สรโณ, อตฺถิ อรโณ.}

\prose{ติวิเธน สงฺขารกฺขนฺโธ\csromandash{}อตฺถิ กุสโล, อตฺถิ อกุสโล, อตฺถิ อพฺยากโต ฯเปฯ.\csromanspacer{} เอวํ ทสวิเธน สงฺขารกฺขนฺโธ.}

\proseitem{98}{เอกวิเธน สงฺขารกฺขนฺโธ\csromandash{}จิตฺตสมฺปยุตฺโต.}

\prose{ทุวิเธน สงฺขารกฺขนฺโธ\csromandash{}อตฺถิ เหตุ, อตฺถิ น เหตุ.}

\prose{ติวิเธน สงฺขารกฺขนฺโธ\csromandash{}อตฺถิ อชฺฌตฺตารมฺมโณ, อตฺถิ พหิทฺธารมฺมโณ, อตฺถิ อชฺฌตฺตพหิทฺธารมฺมโณ ฯเปฯ.\csromanspacer{} เอวํ ทสวิเธน สงฺขารกฺขนฺโธ.}

\proseitem{99}{เอกวิเธน สงฺขารกฺขนฺโธ\csromandash{}จิตฺตสมฺปยุตฺโต.}

\prose{ทุวิเธน สงฺขารกฺขนฺโธ\csromandash{}อตฺถิ สรโณ, อตฺถิ อรโณ.}

\prose{ติวิเธน สงฺขารกฺขนฺโธ\csromandash{}อตฺถิ อชฺฌตฺตารมฺมโณ, อตฺถิ พหิทฺธารมฺมโณ, อตฺถิ อชฺฌตฺตพหิทฺธารมฺมโณ ฯเปฯ.\csromanspacer{} เอวํ ทสวิเธน สงฺขารกฺขนฺโธ.\csromanspacer{} ติกมูลกํ.}

\proseitem{100}{เอกวิเธน สงฺขารกฺขนฺโธ\csromandash{}จิตฺตสมฺปยุตฺโต.}

\prose{ทุวิเธน สงฺขารกฺขนฺโธ\csromandash{}อตฺถิ เหตุ, อตฺถิ น เหตุ.}

\prose{\csromanfolio{43}ติวิเธน สงฺขารกฺขนฺโธ\csromandash{}อตฺถิ กุสโล, อตฺถิ อกุสโล, อตฺถิ อพฺยากโต ฯเปฯ.\csromanspacer{} เอวํ ทสวิเธน สงฺขารกฺขนฺโธ.}

\proseitem{101}{เอกวิเธน สงฺขารกฺขนฺโธ\csromandash{}จิตฺตสมฺปยุตฺโต.}

\prose{ทุวิเธน สงฺขารกฺขนฺโธ\csromandash{}อตฺถิ สเหตุโก, อตฺถิ อเหตุโก.}

\prose{ติวิเธน สงฺขารกฺขนฺโธ\csromandash{}อตฺถิ สุขาย เวทนาย สมฺปยุตฺโต, อตฺถิ ทุกฺขาย เวทนาย สมฺปยุตฺโต, อตฺถิ อทุกฺขมสุขาย เวทนาย สมฺปยุตฺโต ฯเปฯ.\csromanspacer{} เอวํ ทสวิเธน สงฺขารกฺขนฺโธ.}

\proseitem{102}{เอกวิเธน สงฺขารกฺขนฺโธ\csromandash{}จิตฺตสมฺปยุตฺโต.}

\prose{ทุวิเธน สงฺขารกฺขนฺโธ\csromandash{}อตฺถิ เหตุสมฺปยุตฺโต, อตฺถิ เหตุวิปฺปยุตฺโต.}

\prose{ติวิเธน สงฺขารกฺขนฺโธ\csromandash{}อตฺถิ วิปาโก, อตฺถิ วิปากธมฺมธมฺโม, อตฺถิ เนววิปากนวิปากธมฺมธมฺโม ฯเปฯ.\csromanspacer{} เอวํ ทสวิเธน สงฺขารกฺขนฺโธ.}

\proseitem{103}{เอกวิเธน สงฺขารกฺขนฺโธ\csromandash{}จิตฺตสมฺปยุตฺโต.}

\prose{ทุวิเธน สงฺขารกฺขนฺโธ\csromandash{}อตฺถิ เหตุ เจว สเหตุโก จ, อตฺถิ สเหตุโก เจว น จ เหตุ.}

\prose{ติวิเธน สงฺขารกฺขนฺโธ\csromandash{}อตฺถิ อุปาทินฺนุปาทานิโย, อตฺถิ อนุปาทินฺนุปาทานิโย, อตฺถิ อนุปาทินฺน-อนุปาทานิโย ฯเปฯ.\csromanspacer{} เอวํ ทสวิเธน สงฺขารกฺขนฺโธ.}

\proseitem{104}{เอกวิเธน สงฺขารกฺขนฺโธ\csromandash{}จิตฺตสมฺปยุตฺโต.}

\prose{ทุวิเธน สงฺขารกฺขนฺโธ\csromandash{}อตฺถิ เหตุ เจว เหตุสมฺปยุตฺโต จ, อตฺถิ เหตุสมฺปยุตฺโต เจว น จ เหตุ.}

\prose{ติวิเธน สงฺขารกฺขนฺโธ\csromandash{}อตฺถิ สํกิลิฏฺฐสํกิเลสิโก, อตฺถิ อสํกิลิฏฺฐสํกิเลสิโก, อตฺถิ อสํกิลิฏฺฐ-อสํกิเลสิโก ฯเปฯ.\csromanspacer{} เอวํ ทสวิเธน สงฺขารกฺขนฺโธ.}

\proseitem{105}{เอกวิเธน สงฺขารกฺขนฺโธ\csromandash{}จิตฺตสมฺปยุตฺโต.}

\prose{ทุวิเธน สงฺขารกฺขนฺโธ\csromandash{}อตฺถิ น เหตุ สเหตุโก, อตฺถิ น เหตุ อเหตุโก.}

\prose{\csromanfolio{44}ติวิเธน สงฺขารกฺขนฺโธ\csromandash{}อตฺถิ สวิตกฺกสวิจาโร, อตฺถิ อวิตกฺกวิจารมตฺโต, อตฺถิ อวิตกฺก-อวิจาโร ฯเปฯ.\csromanspacer{} เอวํ ทสวิเธน สงฺขารกฺขนฺโธ.}

\proseitem{106}{เอกวิเธน สงฺขารกฺขนฺโธ\csromandash{}จิตฺตสมฺปยุตฺโต.}

\prose{ทุวิเธน สงฺขารกฺขนฺโธ\csromandash{}อตฺถิ โลกิโย, อตฺถิ โลกุตฺตโร.}

\prose{ติวิเธน สงฺขารกฺขนฺโธ\csromandash{}อตฺถิ ปีติสหคโต, อตฺถิ สุขสหคโต, อตฺถิ อุเปกฺขาสหคโต ฯเปฯ.\csromanspacer{} เอวํ ทสวิเธน สงฺขารกฺขนฺโธ.}

\proseitem{107}{เอกวิเธน สงฺขารกฺขนฺโธ\csromandash{}จิตฺตสมฺปยุตฺโต.}

\prose{ทุวิเธน สงฺขารกฺขนฺโธ\csromandash{}อตฺถิ เกนจิ วิญฺเญยฺโย, อตฺถิ เกนจิ น วิญฺเญยฺโย.}

\prose{ติวิเธน สงฺขารกฺขนฺโธ\csromandash{}อตฺถิ ทสฺสเนน ปหาตพฺโพ, อตฺถิ ภาวนาย ปหาตพฺโพ, อตฺถิ เนว ทสฺสเนน น ภาวนาย ปหาตพฺโพ ฯเปฯ.\csromanspacer{} เอวํ ทสวิเธน สงฺขารกฺขนฺโธ.}

\proseitem{108}{เอกวิเธน สงฺขารกฺขนฺโธ\csromandash{}จิตฺตสมฺปยุตฺโต.}

\prose{ทุวิเธน สงฺขารกฺขนฺโธ\csromandash{}อตฺถิ อาสโว, อตฺถิ โน อาสโว.}

\prose{ติวิเธน สงฺขารกฺขนฺโธ\csromandash{}อตฺถิ ทสฺสเนน ปหาตพฺพเหตุโก, อตฺถิ ภาวนาย ปหาตพฺพเหตุโก, อตฺถิ เนว ทสฺสเนน น ภาวนาย ปหาตพฺพเหตุโก ฯเปฯ.\csromanspacer{} เอวํ ทสวิเธน สงฺขารกฺขนฺโธ.}

\proseitem{109}{เอกวิเธน สงฺขารกฺขนฺโธ\csromandash{}จิตฺตสมฺปยุตฺโต.}

\prose{ทุวิเธน สงฺขารกฺขนฺโธ\csromandash{}อตฺถิ สาสโว, อตฺถิ อนาสโว.}

\prose{ติวิเธน สงฺขารกฺขนฺโธ\csromandash{}อตฺถิ อาจยคามี, อตฺถิ อปจยคามี, อตฺถิ เนวาจยคามินาปจยคามี ฯเปฯ.\csromanspacer{} เอวํ ทสวิเธน สงฺขารกฺขนฺโธ.}

\proseitem{110}{เอกวิเธน สงฺขารกฺขนฺโธ\csromandash{}จิตฺตสมฺปยุตฺโต.}

\prose{ทุวิเธน สงฺขารกฺขนฺโธ\csromandash{}อตฺถิ อาสวสมฺปยุตฺโต, อตฺถิ อาสววิปฺปยุตฺโต.}

\prose{ติวิเธน สงฺขารกฺขนฺโธ\csromandash{}อตฺถิ เสกฺโข, อตฺถิ อเสกฺโข, อตฺถิ เนวเสกฺขนาเสกฺโข ฯเปฯ.\csromanspacer{} เอวํ ทสวิเธน สงฺขารกฺขนฺโธ.}

\proseitem{111}{\csromanfolio{45}เอกวิเธน สงฺขารกฺขนฺโธ\csromandash{}จิตฺตสมฺปยุตฺโต.}

\prose{ทุวิเธน สงฺขารกฺขนฺโธ\csromandash{}อตฺถิ อาสโว เจว สาสโว จ, อตฺถิ สาสโว เจว โน จ อาสโว.}

\prose{ติวิเธน สงฺขารกฺขนฺโธ\csromandash{}อตฺถิ ปริตฺโต, อตฺถิ มหคฺคโต, อตฺถิ อปฺปมาโณ ฯเปฯ.\csromanspacer{} เอวํ ทสวิเธน สงฺขารกฺขนฺโธ.}

\proseitem{112}{เอกวิเธน สงฺขารกฺขนฺโธ\csromandash{}จิตฺตสมฺปยุตฺโต.}

\prose{ทุวิเธน สงฺขารกฺขนฺโธ\csromandash{}อตฺถิ อาสโว เจว อาสวสมฺปยุตฺโต จ, อตฺถิ อาสวสมฺปยุตฺโต เจว โน จ อาสโว.}

\prose{ติวิเธน สงฺขารกฺขนฺโธ\csromandash{}อตฺถิ ปริตฺตารมฺมโณ, อตฺถิ มหคฺคตารมฺมโณ, อตฺถิ อปฺปมาณารมฺมโณ ฯเปฯ.\csromanspacer{} เอวํ ทสวิเธน สงฺขารกฺขนฺโธ.}

\proseitem{113}{เอกวิเธน สงฺขารกฺขนฺโธ\csromandash{}จิตฺตสมฺปยุตฺโต.}

\prose{ทุวิเธน สงฺขารกฺขนฺโธ\csromandash{}อตฺถิ อาสววิปฺปยุตฺตสาสโว, อตฺถิ อาสววิปฺปยุตฺต-อนาสโว.}

\prose{ติวิเธน สงฺขารกฺขนฺโธ\csromandash{}อตฺถิ หีโน, อตฺถิ มชฺฌิโม, อตฺถิ ปณีโต ฯเปฯ.\csromanspacer{} เอวํ ทสวิเธน สงฺขารกฺขนฺโธ.}

\proseitem{114}{เอกวิเธน สงฺขารกฺขนฺโธ\csromandash{}จิตฺตสมฺปยุตฺโต.}

\prose{ทุวิเธน สงฺขารกฺขนฺโธ\csromandash{}อตฺถิ สํโยชนํ, อตฺถิ โน สํโยชนํ.}

\prose{ติวิเธน สงฺขารกฺขนฺโธ\csromandash{}อตฺถิ มิจฺฉตฺตนิยโต, อตฺถิ สมฺมตฺตนิยโต, อตฺถิ อนิยโต ฯเปฯ.\csromanspacer{} เอวํ ทสวิเธน สงฺขารกฺขนฺโธ.}

\proseitem{115}{เอกวิเธน สงฺขารกฺขนฺโธ\csromandash{}จิตฺตสมฺปยุตฺโต.}

\prose{ทุวิเธน สงฺขารกฺขนฺโธ\csromandash{}อตฺถิ สํโยชนิโย, อตฺถิ อสํโยชนิโย.}

\prose{ติวิเธน สงฺขารกฺขนฺโธ\csromandash{}อตฺถิ มคฺคารมฺมโณ, อตฺถิ มคฺคเหตุโก, อตฺถิ มคฺคาธิปติ ฯเปฯ.\csromanspacer{} เอวํ ทสวิเธน สงฺขารกฺขนฺโธ.}

\proseitem{116}{\csromanfolio{46}เอกวิเธน สงฺขารกฺขนฺโธ\csromandash{}จิตฺตสมฺปยุตฺโต.}

\prose{ทุวิเธน สงฺขารกฺขนฺโธ\csromandash{}อตฺถิ สํโยชนสมฺปยุตฺโต, อตฺถิ สํโยชนวิปฺปยุตฺโต.}

\prose{ติวิเธน สงฺขารกฺขนฺโธ\csromandash{}อตฺถิ อุปฺปนฺโน, อตฺถิ อนุปฺปนฺโน, อตฺถิ อุปฺปาที ฯเปฯ.\csromanspacer{} เอวํ ทสวิเธน สงฺขารกฺขนฺโธ.}

\proseitem{117}{เอกวิเธน สงฺขารกฺขนฺโธ\csromandash{}จิตฺตสมฺปยุตฺโต.}

\prose{ทุวิเธน สงฺขารกฺขนฺโธ\csromandash{}อตฺถิ สํโยชนญฺเจว สํโยชนิโย จ, อตฺถิ สํโยชนิโย เจว โน จ สํโยชนํ.}

\prose{ติวิเธน สงฺขารกฺขนฺโธ\csromandash{}อตฺถิ อตีโต, อตฺถิ อนาคโต, อตฺถิ ปจฺจุปฺปนฺโน ฯเปฯ.\csromanspacer{} เอวํ ทสวิเธน สงฺขารกฺขนฺโธ.}

\proseitem{118}{เอกวิเธน สงฺขารกฺขนฺโธ\csromandash{}จิตฺตสมฺปยุตฺโต.}

\prose{ทุวิเธน สงฺขารกฺขนฺโธ\csromandash{}อตฺถิ สํโยชนญฺเจว สํโยชนสมฺปยุตฺโต จ, อตฺถิ สํโยชนสมฺปยุตฺโต เจว โน จ สํโยชนํ.}

\prose{ติวิเธน สงฺขารกฺขนฺโธ\csromandash{}อตฺถิ อตีตารมฺมโณ, อตฺถิ อนาคตารมฺมโณ, อตฺถิ ปจฺจุปฺปนฺนารมฺมโณ ฯเปฯ.\csromanspacer{} เอวํ ทสวิเธน สงฺขารกฺขนฺโธ.}

\proseitem{119}{เอกวิเธน สงฺขารกฺขนฺโธ\csromandash{}จิตฺตสมฺปยุตฺโต.}

\prose{ทุวิเธน สงฺขารกฺขนฺโธ\csromandash{}อตฺถิ สํโยชนวิปฺปยุตฺตสํโยชนิโย, อตฺถิ สํโยชนวิปฺปยุตฺต-อสํโยชนิโย.}

\prose{ติวิเธน สงฺขารกฺขนฺโธ\csromandash{}อตฺถิ อชฺฌตฺโต, อตฺถิ พหิทฺโธ, อตฺถิ อชฺฌตฺตพหิทฺโธ ฯเปฯ.\csromanspacer{} เอวํ ทสวิเธน สงฺขารกฺขนฺโธ.}

\proseitem{120}{เอกวิเธน สงฺขารกฺขนฺโธ\csromandash{}จิตฺตสมฺปยุตฺโต.}

\prose{ทุวิเธน สงฺขารกฺขนฺโธ\csromandash{}อตฺถิ คนฺโถ, อตฺถิ โน คนฺโถ.}

\prose{ติวิเธน สงฺขารกฺขนฺโธ\csromandash{}อตฺถิ อชฺฌตฺตารมฺมโณ, อตฺถิ พหิทฺธารมฺมโณ, อตฺถิ อชฺฌตฺตพหิทฺธารมฺมโณ ฯเปฯ.\csromanspacer{} เอวํ ทสวิเธน สงฺขารกฺขนฺโธ.}

\vspace{\tipitakaparskip}
\csromancenter{อุภโตวฑฺฒกํ.}
\prose{\csromanfolio{47}สตฺตวิเธน สงฺขารกฺขนฺโธ\csromandash{}อตฺถิ กุสโล, อตฺถิ อกุสโล, อตฺถิ อพฺยากโต, อตฺถิ กามาวจโร, อตฺถิ รูปาวจโร, อตฺถิ อรูปาวจโร, อตฺถิ อปริยาปนฺโน.\csromanspacer{} เอวํ สตฺตวิเธน สงฺขารกฺขนฺโธ.}

\prose{อปโรปิ สตฺตวิเธน สงฺขารกฺขนฺโธ\csromandash{}อตฺถิ สุขาย เวทนาย สมฺปยุตฺโต, อตฺถิ ทุกฺขาย เวทนาย สมฺปยุตฺโต, อตฺถิ อทุกฺขมสุขาย เวทนาย สมฺปยุตฺโต, อตฺถิ กามาวจโร, อตฺถิ รูปาวจโร, อตฺถิ อรูปาวจโร, อตฺถิ อปริยาปนฺโน ฯเปฯ อตฺถิ อชฺฌตฺตารมฺมโณ, อตฺถิ พหิทฺธารมฺมโณ, อตฺถิ อชฺฌตฺตพหิทฺธารมฺมโณ, อตฺถิ กามาวจโร, อตฺถิ รูปาวจโร, อตฺถิ อรูปาวจโร, อตฺถิ อปริยาปนฺโน.\csromanspacer{} เอวํ สตฺตวิเธน สงฺขารกฺขนฺโธ.}

\prose{จตุวีสติวิเธน สงฺขารกฺขนฺโธ\csromandash{}จกฺขุสมฺผสฺสปจฺจยา สงฺขารกฺขนฺโธ อตฺถิ กุสโล, อตฺถิ อกุสโล, อตฺถิ อพฺยากโต.\csromanspacer{} โสตสมฺผสฺสปจฺจยา ฯเปฯ.\csromanspacer{} ฆานสมฺผสฺสปจฺจยา ฯเปฯ.\csromanspacer{} ชิวฺหาสมฺผสฺสปจฺจยา ฯเปฯ.\csromanspacer{} กายสมฺผสฺสปจฺจยา ฯเปฯ.\csromanspacer{} มโนสมฺผสฺสปจฺจยา สงฺขารกฺขนฺโธ อตฺถิ กุสโล, อตฺถิ อกุสโล, อตฺถิ อพฺยากโต, จกฺขุสมฺผสฺสชา เจตนา ฯเปฯ มโนสมฺผสฺสชา เจตนา.\csromanspacer{} เอวํ จตุวีสติวิเธน สงฺขารกฺขนฺโธ.}

\prose{อปโรปิ จตุวีสติวิเธน สงฺขารกฺขนฺโธ\csromandash{} จกฺขุสมฺผสฺสปจฺจยา สงฺขารกฺขนฺโธ อตฺถิ สุขาย เวทนาย สมฺปยุตฺโต ฯเปฯ อตฺถิ อชฺฌตฺตารมฺมโณ, อตฺถิ พหิทฺธารมฺมโณ, อตฺถิ อชฺฌตฺตพหิทฺธารมฺมโณ.\csromanspacer{} จกฺขุสมฺผสฺสชา เจตนา ฯเปฯ มโนสมฺผสฺสชา เจตนา.\csromanspacer{} โสตสมฺผสฺสปจฺจยา ฯเปฯ.\csromanspacer{} ฆานสมฺผสฺสปจฺจยา ฯเปฯ.\csromanspacer{} ชิวฺหาสมฺผสฺสปจฺจยา ฯเปฯ.\csromanspacer{} กายสมฺผสฺสปจฺจยา ฯเปฯ.\csromanspacer{} มโนสมฺผสฺสปจฺจยา สงฺขารกฺขนฺโธ อตฺถิ สุขาย เวทนาย สมฺผยุตฺโต ฯเปฯ อตฺถิ อชฺฌตฺตารมฺมโณ, อตฺถิ พหิทฺธารมฺมโณ, อตฺถิ อชฺฌตฺตพหิทฺธารมฺมโณ.\csromanspacer{} จกฺขุสมฺผสฺสชา เจตนา, โสตสมฺผสฺสชา เจตนา, ฆานสมฺผสฺสชา เจตนา, ชิวฺหาสมฺผสฺสชา เจตนา, กายสมฺผสฺสชา เจตนา, มโนสมฺผสฺสชา เจตนา.\csromanspacer{} เอวํ จตุวีสติวิเธน สงฺขารกฺขนฺโธ.}

\prose{ติํสติวิเธน สงฺขารกฺขนฺโธ\csromandash{}จกฺขุสมฺผสฺสปจฺจยา สงฺขารกฺขนฺโธ อตฺถิ กามาวจโร, อตฺถิ รูปาวจโร, อตฺถิ อรูปาวจโร, อตฺถิ \csromanfolio{48}อปริยาปนฺโน.\csromanspacer{} โสตสมฺผสฺสปจฺจยา ฯเปฯ.\csromanspacer{} ฆานสมฺผสฺสปจฺจยา ฯเปฯ.\csromanspacer{} ชิวฺหาสมฺผสฺสปจฺจยา ฯเปฯ.\csromanspacer{} กายสมฺผสฺสปจฺจยา ฯเปฯ.\csromanspacer{} มโนสมฺผสฺสปจฺจยา สงฺขารกฺขนฺโธ อตฺถิ กามาวจโร, อตฺถิ รูปาวจโร, อตฺถิ อรูปาวจโร, อตฺถิ อปริยาปนฺโน.\csromanspacer{} จกฺขุสมฺผสฺสชา เจตนา ฯเปฯ มโนสมฺผสฺสชา เจตนา.\csromanspacer{} เอวํ ติํสติวิเธน สงฺขารกฺขนฺโธ.}

\prose{พหุวิเธน สงฺขารกฺขนฺโธ\csromandash{}จกฺขุสมฺผสฺสปจฺจยา สงฺขารกฺขนฺโธ อตฺถิ กุสโล, อตฺถิ อกุสโล, อตฺถิ อพฺยากโต, อตฺถิ กามาวจโร, อตฺถิ รูปาวจโร, อตฺถิ อรูปาวจโร, อตฺถิ อปริยาปนฺโน.\csromanspacer{} โสตสมฺผสฺสปจฺจยา ฯเปฯ.\csromanspacer{} ฆานสมฺผสฺสปจฺจยา ฯเปฯ.\csromanspacer{} ชิวฺหาสมฺผสฺสปจฺจยา ฯเปฯ.\csromanspacer{} กายสมฺผสฺสปจฺจยา ฯเปฯ.\csromanspacer{} มโนสมฺผสฺสปจฺจยา สงฺขารกฺขนฺโธ อตฺถิ กุสโล, อตฺถิ อกุสโล, อตฺถิ อพฺยากโต, อตฺถิ กามาวจโร, อตฺถิ รูปาวจโร, อตฺถิ อรูปาวจโร, อตฺถิ อปริยาปนฺโน.\csromanspacer{} จกฺขุสมฺผสฺสชา เจตนา, โสตสมฺผสฺสชา เจตนา, ฆานสมฺผสฺสชา เจตนา, ชิวฺหาสมฺผสฺสชา เจตนา, กายสมฺผสฺสชา เจตนา, มโนสมฺผสฺสชา เจตนา.\csromanspacer{} เอวํ พหุวิเธน สงฺขารกฺขนฺโธ.}

\prose{อปโรปิ พหุวิเธน สงฺขารกฺขนฺโธ\csromandash{}จกฺขุสมฺผสฺสปจฺจยา สงฺขารกฺขนฺโธ อตฺถิ สุขาย เวทนาย สมฺปยุตฺโต, อตฺถิ ทุกฺขาย เวทนาย สมฺปยุตฺโต, อตฺถิ อทุกฺขมสุขาย เวทนาย สมฺปยุตฺโต ฯเปฯ อตฺถิ อชฺฌตฺตารมฺมโณ, อตฺถิ พหิทฺธารมฺมโณ, อตฺถิ อชฺฌตฺตพหิทฺธารมฺมโณ, อตฺถิ กามาวจโร, อตฺถิ รูปาวจโร, อตฺถิ อรูปาวจโร อตฺถิ อปริยาปนฺโน.\csromanspacer{} โสตสมฺผสฺสปจฺจยา ฯเปฯ.\csromanspacer{} ฆานสมฺผสฺสปจฺจยา ฯเปฯ.\csromanspacer{} ชิวฺหาสมฺผสฺสปจฺจยา ฯเปฯ.\csromanspacer{} กายสมฺผสฺสปจฺจยา ฯเปฯ.\csromanspacer{} มโนสมฺผสฺสปจฺจยา สงฺขารกฺขนฺโธ อตฺถิ สุขาย เวทนาย สมฺปยุตฺโต ฯเปฯ อตฺถิ อชฺฌตฺตารมฺมโณ, อตฺถิ พหิทฺธารมฺมโณ, อตฺถิ อชฺฌตฺตพหิทฺธารมฺมโณ, อตฺถิ กามาวจโร, อตฺถิ รูปาวจโร, อตฺถิ อรูปาวจโร, อตฺถิ อปริยาปนฺโน.\csromanspacer{} จกฺขุสมฺผสฺสชา เจตนา, โสตสมฺผสฺสชา เจตนา, ฆานสมฺผสฺสชา เจตนา, ชิวฺหาสมฺผสฺสชา เจตนา, กายสมฺผสฺสชา เจตนา, มโนสมฺผสฺสชา เจตนา.\csromanspacer{} เอวํ พหุวิเธน สงฺขารกฺขนฺโธ.}

\vspace{\tipitakaparskip}
\csromancenter{อยํ วุจฺจติ สงฺขารกฺขนฺโธ.}
\csromanmatikaanchor{mtk.13}
\csromantocmarkref{subsection}{5. วิญฺญาณกฺขนฺธ}{mtk.13}
\bookmark[dest={mtk.13},level=2]{5. วิญฺญาณกฺขนฺธ}
\csromanheader{2}{5. วิญฺญาณกฺขนฺธ}
\proseitem{121}{\csromanfolio{49}ตตฺถ กตโม วิญฺญาณกฺขนฺโธ, เอกวิเธน วิญฺญาณกฺขนฺโธ\csromandash{} ผสฺสสมฺปยุตฺโต.}

\prose{ทุวิเธน วิญฺญาณกฺขนฺโธ\csromandash{}อตฺถิ สเหตุโก, อตฺถิ อเหตุโก.}

\prose{ติวิเธน วิญฺญาณกฺขนฺโธ\csromandash{}อตฺถิ กุสโล, อตฺถิ อกุสโล, อตฺถิ อพฺยากโต.}

\prose{จตุพฺพิเธน วิญฺญาณกฺขนฺโธ\csromandash{}อตฺถิ กามาวจโร, อตฺถิ รูปาวจโร, อตฺถิ อรูปาวจโร, อตฺถิ อปริยาปนฺโน.}

\prose{ปญฺจวิเธน วิญฺญาณกฺขนฺโธ\csromandash{}อตฺถิ สุขินฺทฺริยสมฺปยุตฺโต, อตฺถิ ทุกฺขินฺทฺริยสมฺปยุตฺโต, อตฺถิ โสมนสฺสินฺทฺริยสมฺปยุตฺโต, อตฺถิ โทมนสฺสินฺทฺริยสมฺปยุตฺโต, อตฺถิ อุเปกฺขินฺทฺริยสมฺปยุตฺโต.\csromanspacer{} เอวํ ปญฺจวิเธน วิญฺญาณกฺขนฺโธ.}

\prose{ฉพฺพิเธน วิญฺญาณกฺขนฺโธ\csromandash{}จกฺขุวิญฺญาณํ, โสตวิญฺญาณํ, ฆานวิญฺญาณํ, ชิวฺหาวิญฺญาณํ, กายวิญฺญาณํ, มโนวิญฺญาณํ.\csromanspacer{} เอวํ ฉพฺพิเธน วิญฺญาณกฺขนฺโธ.}

\prose{สตฺตวิเธน วิญฺญาณกฺขนฺโธ\csromandash{}จกฺขุวิญฺญาณํ, โสตวิญฺญาณํ, ฆานวิญฺญาณํ, ชิวฺหาวิญฺญาณํ, กายวิญฺญาณํ, มโนธาตุ, มโนวิญฺญาณธาตุ.\csromanspacer{} เอวํ สตฺตวิเธน วิญฺญาณกฺขนฺโธ.}

\prose{อฏฺฐวิเธน วิญฺญาณกฺขนฺโธ\csromandash{}จกฺขุวิญฺญาณํ, โสตวิญฺญาณํ, ฆานวิญฺญาณํ, ชิวฺหาวิญฺญาณํ, กายวิญฺญาณํ อตฺถิ สุขสหคตํ, อตฺถิ ทุกฺขสหคตํ, มโนธาตุ, มโนวิญฺญาณธาตุ.\csromanspacer{} เอวํ อฏฺฐวิเธน วิญฺญาณกฺขนฺโธ.}

\prose{นววิเธน วิญฺญาณกฺขนฺโธ\csromandash{}จกฺขุวิญฺญาณํ, โสตวิญฺญาณํ, ฆานวิญฺญาณํ, ชิวฺหาวิญฺญาณํ, กายวิญฺญาณํ, มโนธาตุ, มโนวิญฺญาณธาตุ อตฺถิ กุสโล, อตฺถิ อกุสโล, อตฺถิ อพฺยากโต.\csromanspacer{} เอวํ นววิเธน วิญฺญาณกฺขนฺโธ.}

\prose{ทสวิเธน วิญฺญาณกฺขนฺโธ\csromandash{}จกฺขุวิญฺญาณํ ฯเปฯ กายวิญฺญาณํ อตฺถิ สุขสหคตํ, อตฺถิ ทุกฺขสหคตํ, มโนธาตุ, มโนวิญฺญาณธาตุ อตฺถิ กุสลํ, อตฺถิ อกุสลํ, อตฺถิ อพฺยากตํ.\csromanspacer{} เอวํ ทสวิเธน วิญฺญาณกฺขนฺโธ.}

\proseitem{122}{\csromanfolio{50}เอกวิเธน วิญฺญาณกฺขนฺโธ\csromandash{}ผสฺสสมฺปยุตฺโต.}

\prose{ทุวิเธน วิญฺญาณกฺขนฺโธ\csromandash{}อตฺถิ สเหตุโก, อตฺถิ อเหตุโก.}

\prose{ติวิเธน วิญฺญาณกฺขนฺโธ\csromandash{}อตฺถิ สุขาย เวทนาย สมฺปยุตฺโต, อตฺถิ ทุกฺขาย เวทนาย สมฺปยุตฺโต, อตฺถิ อทุกฺขมสุขาย เวทนาย สมฺปยุตฺโต.\csromanspacer{} อตฺถิ วิปาโก, อตฺถิ วิปากธมฺมธมฺโม, อตฺถิ เนววิปากนวิปากธมฺมธมฺโม.\csromanspacer{} อตฺถิ อุปาทินฺนุปาทานิโย, อตฺถิ อนุปาทินฺนุปาทานิโย, อตฺถิ อนุปาทินฺน-อนุปาทานิโย.\csromanspacer{} อตฺถิ สํกิลิฏฺฐสํกิเลสิโก, อตฺถิ อสํกิลิฏฺฐสํกิเลสิโก, อตฺถิ อสํกิเลฏฺฐ- อสํกิเลสิโก.\csromanspacer{} อตฺถิ สวิตกฺกสวิจาโร, อตฺถิ อวิตกฺกวิจารมตฺโต, อตฺถิ อวิตกฺก-อวิจาโร.\csromanspacer{} อตฺถิ ปีติสหคโต, อตฺถิ สุขสหคโต, อตฺถิ อุเปกฺขาสหคโต.\csromanspacer{} อตฺถิ ทสฺสเนน ปหาตพฺโพ, อตฺถิ ภาวนาย ปหาตพฺโพ, อตฺถิ เนว ทสฺสเนน น ภาวนาย ปหาตพฺโพ.\csromanspacer{} อตฺถิ ทสฺสเนน ปหาตพฺพเหตุโก, อตฺถิ ภาวนาย ปหาตพฺพเหตุโก, อตฺถิ เนว ทสฺสเนน น ภาวนาย ปหาตพฺพเหตุโก.\csromanspacer{} อตฺถิ อาจยคามี, อตฺถิ อปจยคามี, อตฺถิ เนวาจยคามินาปจยคามี.\csromanspacer{} อตฺถิ เสกฺโข, อตฺถิ อเสกฺโข, อตฺถิ เนวเสกฺขนาเสกฺโข.\csromanspacer{} อตฺถิ ปริตฺโต, อตฺถิ มหคฺคโต, อตฺถิ อปฺปมาโณ.\csromanspacer{} อตฺถิ ปริตฺตารมฺมโณ, อตฺถิ มหคฺคตารมฺมโณ, อตฺถิ อปฺปมาณารมฺมโณ.\csromanspacer{} อตฺถิ หีโน, อตฺถิ มชฺฌิโม, อตฺถิ ปณีโต.\csromanspacer{} อตฺถิ มิจฺฉตฺตนิยโต, อตฺถิ สมฺมตฺตนิยโต, อตฺถิ อนิยโต.\csromanspacer{} อตฺถิ มคฺคารมฺมโณ, อตฺถิ มคฺคเหตุโก, อตฺถิ มคฺคาธิปติ.\csromanspacer{} อตฺถิ อุปฺปนฺโน, อตฺถิ อนุปฺปนฺโน, อตฺถิ อุปฺปาที.\csromanspacer{} อตฺถิ อตีโต, อตฺถิ อนาคโต, อตฺถิ ปจฺจุปฺปนฺโน.\csromanspacer{} อตฺถิ อตีตารมฺมโณ, อตฺถิ อนาคตารมฺมโณ, อตฺถิ ปจฺจุปฺปนฺนารมฺมโณ.\csromanspacer{} อตฺถิ อชฺฌตฺโต, อตฺถิ พหิทฺโธ, อตฺถิ อชฺฌตฺตพหิทฺโธ.\csromanspacer{} อตฺถิ อชฺฌตฺตารมฺมโณ, อตฺถิ พหิทฺธารมฺมโณ, อตฺถิ อชฺฌตฺตพหิทฺธารมฺมโณ ฯเปฯ.\csromanspacer{} เอวํ ทสวิเธน วิญฺญาณกฺขนฺโธ.}

\proseitem{123}{เอกวิเธน วิญฺญาณกฺขนฺโธ\csromandash{}ผสฺสสมฺปยุตฺโต.}

\prose{ทุวิเธน วิญฺญาณกฺขนฺโธ\csromandash{}อตฺถิ เหตุสมฺปยุตฺโต, อตฺถิ เหตุวิปฺปยุตฺโต.\csromanspacer{} อตฺถิ น เหตุ สเหตุโก, อตฺถิ น เหตุ อเหตุโก.\csromanspacer{} อตฺถิ โลกิโย, อตฺถิ โลกุตฺตโร.\csromanspacer{} อตฺถิ เกนจิ วิญฺเญยฺโย, อตฺถิ เกนจิ น วิญฺเญยฺโย.\csromanspacer{} อตฺถิ สาสโว, อตฺถิ อนาสโว.\csromanspacer{} อตฺถิ \csromanfolio{51}อาสวสมฺปยุตฺโต, อตฺถิ อาสววิปฺปยุตฺโต.\csromanspacer{} อตฺถิ อาสววิปฺปยุตฺตสาสโว, อตฺถิ อาสววิปฺปยุตฺต-อนาสโว.\csromanspacer{} อตฺถิ สํโยชนิโย, อตฺถิ อสํโยชนิโย.\csromanspacer{} อตฺถิ สํโยชนสมฺปยุตฺโต, อตฺถิ สํโยชนวิปฺปยุตฺโต.\csromanspacer{} อตฺถิ สํโยชนวิปฺปยุตฺตสํโยชนิโย, อตฺถิ สํโยชนวิปฺปยุตฺต-อสํโยชนิโย.}

\prose{อตฺถิ คนฺถนิโย, อตฺถิ อคนฺถนิโย.\csromanspacer{} อตฺถิ คนฺถสมฺปยุตฺโต, อตฺถิ คนฺถวิปฺปยุตฺโต.\csromanspacer{} อตฺถิ คนฺถวิปฺปยุตฺตคนฺถนิโย, อตฺถิ คนฺถวิปฺปยุตฺต- อคนฺถนิโย.\csromanspacer{} อตฺถิ โอฆนิโย, อตฺถิ อโนฆนิโย.\csromanspacer{} อตฺถิ โอฆสมฺปยุตฺโต, อตฺถิ โอฆวิปฺปยุตฺโต.\csromanspacer{} อตฺถิ โอฆวิปฺปยุตฺต-โอฆนิโย, อตฺถิ โอฆวิปฺปยุตฺต- อโนฆนิโย.\csromanspacer{} อตฺถิ โยคนิโย, อตฺถิ อโยคนิโย.\csromanspacer{} อตฺถิ โยคสมฺปยุตฺโต, อตฺถิ โยควิปฺปยุตฺโต.\csromanspacer{} อตฺถิ โยควิปฺปยุตฺตโยคนิโย, อตฺถิ โยควิปฺปยุตฺต- อโยคนิโย.\csromanspacer{} อตฺถิ นีวรณิโย, อตฺถิ อนีวรณิโย.\csromanspacer{} อตฺถิ นีวรณสมฺปยุตฺโต, อตฺถิ นีวรณวิปฺปยุตฺโต.\csromanspacer{} อตฺถิ นีวรณวิปฺปยุตฺตนีวรณิโย, อตฺถิ นีวรณวิปฺปยุตฺต-อนีวรณิโย.}

\prose{อตฺถิ ปรามฏฺโฐ, อตฺถิ อปรามฏฺโฐ.\csromanspacer{} อตฺถิ ปรามาสสมฺปยุตฺโต, อตฺถิ ปรามาสวิปฺปยุตฺโต.\csromanspacer{} อตฺถิ ปรามาสวิปฺปยุตฺตปรามฏฺโฐ, อตฺถิ ปรามาสวิปฺปยุตฺต-อปรามฏฺโฐ.\csromanspacer{} อตฺถิ อุปาทินฺโน, อตฺถิ อนุปาทินฺโน.\csromanspacer{} อตฺถิ อุปาทานิโย, อตฺถิ อนุปาทานิโย.\csromanspacer{} อตฺถิ อุปาทานสมฺปยุตฺโต, อตฺถิ อุปาทานวิปฺปยุตฺโต.\csromanspacer{} อตฺถิ อุปาทานวิปฺปยุตฺต-อุปาทานิโย, อตฺถิ อุปาทานวิปฺปยุตฺต-อนุปาทานิโย.\csromanspacer{} อตฺถิ สํกิเลสิโก, อตฺถิ อสํกิเลสิโก.\csromanspacer{} อตฺถิ สํกิลิฏฺโฐ, อตฺถิ อสํกิลิฏฺโฐ.\csromanspacer{} อตฺถิ กิเลสสมฺปยุตฺโต, อตฺถิ กิเลสวิปฺปยุตฺโต.\csromanspacer{} อตฺถิ กิเลสวิปฺปยุตฺตสํกิเลสิโก, อตฺถิ กิเลสวิปฺปยุตฺต- อสํกิเลสิโก.\csromanspacer{} อตฺถิ ทสฺสเนน ปหาตพฺโพ, อตฺถิ น ทสฺสเนน ปหาตพฺโพ.\csromanspacer{} อตฺถิ ภาวนาย ปหาตพฺโพ, อตฺถิ น ภาวนาย ปหาตพฺโพ.\csromanspacer{} อตฺถิ ทสฺสเนน ปหาตพฺพเหตุโก, อตฺถิ น ทสฺสเนน ปหาตพฺพเหตุโก.\csromanspacer{} อตฺถิ ภาวนาย ปหาตพฺพเหตุโก, อตฺถิ น ภาวนาย ปหาตพฺพเหตุโก.}

\prose{อตฺถิ สวิตกฺโก, อตฺถิ อวิตกฺโก.\csromanspacer{} อตฺถิ สวิจาโร, อตฺถิ อวิจาโร.\csromanspacer{} อตฺถิ สปฺปีติโก, อตฺถิ อปฺปีติโก.\csromanspacer{} อตฺถิ ปีติสหคโต, อตฺถิ น ปีติสหคโต.\csromanspacer{} อตฺถิ สุขสหคโต, อตฺถิ น สุขสหคโต.\csromanspacer{} อตฺถิ อุเปกฺขาสหคโต, อตฺถิ น อุเปกฺขาสหคโต.\csromanspacer{} อตฺถิ \csromanfolio{52}กามาวจโร, อตฺถิ น กามาวจโร.\csromanspacer{} อตฺถิ รูปาวจโร, อตฺถิ น รูปาวจโร.\csromanspacer{} อตฺถิ อรูปาวจโร, อตฺถิ น อรูปาวจโร.\csromanspacer{} อตฺถิ ปริยาปนฺโน, อตฺถิ อปริยาปนฺโน.\csromanspacer{} อตฺถิ นิยฺยานิโก, อตฺถิ อนิยฺยานิโก.\csromanspacer{} อตฺถิ นิยโต, อตฺถิ อนิยโต.\csromanspacer{} อตฺถิ ส-อุตฺตโร, อตฺถิ อนุตฺตโร.\csromanspacer{} อตฺถิ สรโณ, อตฺถิ อรโณ.}

\prose{ติวิเธน วิญฺญาณกฺขนฺโธ\csromandash{}อตฺถิ กุสโล, อตฺถิ อกุสโล อตฺถิ อพฺยากโต ฯเปฯ.\csromanspacer{} เอวํ ทสวิเธน วิญฺญาณกฺขนฺโธ.}

\proseitem{124}{เอกวิเธน วิญฺญาณกฺขนฺโธ\csromandash{}ผสฺสสมฺปยุตฺโต.}

\prose{ทุวิเธน วิญฺญาณกฺขนฺโธ\csromandash{}อตฺถิ สรโณ, อตฺถิ อรโณ.}

\prose{ติวิเธน วิญฺญาณกฺขนฺโธ\csromandash{}อตฺถิ สุขาย เวทนาย สมฺปยุตฺโต, อตฺถิ ทุกฺขาย เวทนาย สมฺปยุตฺโต, อตฺถิ อทุกฺขมสุขาย เวทนาย สมฺปยุตฺโต.\csromanspacer{} อตฺถิ วิปาโก ฯเปฯ อตฺถิ อชฺฌตฺตารมฺมโณ, อตฺถิ พหิทฺธารมฺมโณ, อตฺถิ อชฺฌตฺตพหิทฺธารมฺมโณ ฯเปฯ.\csromanspacer{} เอวํ ทสวิเธน วิญฺญาณกฺขนฺโธ.\csromanspacer{} ทุกมูลกํ.}

\proseitem{125}{เอกวิเธน วิญฺญาณกฺขนฺโธ\csromandash{}ผสฺสสมฺปยุตฺโต.}

\prose{ทุวิเธน วิญฺญาณกฺขนฺโธ\csromandash{}อตฺถิ สเหตุโก, อตฺถิ อเหตุโก.}

\prose{ติวิเธน วิญฺญาณกฺขนฺโธ\csromandash{}อตฺถิ กุสโล, อตฺถิ อกุสโล, อตฺถิ อพฺยากโต ฯเปฯ.\csromanspacer{} เอวํ ทสวิเธน วิญฺญาณกฺขนฺโธ.}

\proseitem{126}{เอกวิเธน วิญฺญาณกฺขนฺโธ\csromandash{}ผสฺสสมฺปยุตฺโต.}

\prose{ทุวิเธน วิญฺญาณกฺขนฺโธ\csromandash{}อตฺถิ เหตุสมฺปยุตฺโต, อตฺถิ เหตุวิปฺปยุตฺโต ฯเปฯ.\csromanspacer{} อตฺถิ สรโณ, อตฺถิ อรโณ.}

\prose{ติวิเธน วิญฺญาณกฺขนฺโธ\csromandash{}อตฺถิ กุสโล, อตฺถิ อกุสโล, อตฺถิ อพฺยากโต ฯเปฯ.\csromanspacer{} เอวํ ทสวิเธน วิญฺญาณกฺขนฺโธ.}

\proseitem{127}{เอกวิเธน วิญฺญาณกฺขนฺโธ\csromandash{}ผสฺสสมฺปยุตฺโต.}

\prose{ทุวิเธน วิญฺญาณกฺขนฺโธ\csromandash{}อตฺถิ สเหตุโก, อตฺถิ อเหตุโก.}

\prose{ติวิเธน วิญฺญาณกฺขนฺโธ\csromandash{}อตฺถิ สุขาย เวทนาย สมฺปยุตฺโต, อตฺถิ ทุกฺขาย เวทนาย สมฺปยุตฺโต, อตฺถิ อทุกฺขมสุขาย เวทนาย \csromanfolio{53}สมฺปยุตฺโต.\csromanspacer{} อตฺถิ วิปาโก ฯเปฯ.\csromanspacer{} อตฺถิ อชฺฌตฺตารมฺมโณ, อตฺถิ พหิทฺธารมฺมโณ, อตฺถิ อชฺฌตฺตพหิทฺธารมฺมโณ ฯเปฯ.\csromanspacer{} เอวํ ทสวิเธน วิญฺญาณกฺขนฺโธ.}

\proseitem{128}{เอกวิเธน วิญฺญาณกฺขนฺโธ\csromandash{}ผสฺสสมฺปยุตฺโต.}

\prose{ทุวิเธน วิญฺญาณกฺขนฺโธ\csromandash{}อตฺถิ เหตุสมฺปยุตฺโต, อตฺถิ เหตุวิปฺปยุตฺโต ฯเปฯ.\csromanspacer{} อตฺถิ สรโณ, อตฺถิ อรโณ.}

\prose{ติวิเธน วิญฺญาณกฺขนฺโธ\csromandash{}อตฺถิ อชฺฌตฺตารมฺมโณ, อตฺถิ พหิทฺธารมฺมโณ, อตฺถิ อชฺฌตฺตพหิทฺธารมฺมโณ ฯเปฯ.\csromanspacer{} เอวํ ทสวิเธน วิญฺญาณกฺขนฺโธ.\csromanspacer{} ติกมูลกํ.}

\proseitem{129}{เอกวิเธน วิญฺญาณกฺขนฺโธ\csromandash{}ผสฺสสมฺปยุตฺโต.}

\prose{ทุวิเธน วิญฺญาณกฺขนฺโธ\csromandash{}อตฺถิ สเหตุโก, อตฺถิ อเหตุโก.}

\prose{ติวิเธน วิญฺญาณกฺขนฺโธ\csromandash{}อตฺถิ กุสโล, อตฺถิ อกุสโล, อตฺถิ อพฺยากโต ฯเปฯ.\csromanspacer{} เอวํ ทสวิเธน วิญฺญาณกฺขนฺโธ.}

\proseitem{130}{เอกวิเธน วิญฺญาณกฺขนฺโธ\csromandash{}ผสฺสสมฺปยุตฺโต.}

\prose{ทุวิเธน วิญฺญาณกฺขนฺโธ\csromandash{}อตฺถิ เหตุสมฺปยุตฺโต, อตฺถิ เหตุวิปฺปยุตฺโต.}

\prose{ติวิเธน วิญฺญาณกฺขนฺโธ\csromandash{}อตฺถิ สุขาย เวทนาย สมฺปยุตฺโต, อตฺถิ ทุกฺขาย เวทนาย สมฺปยุตฺโต, อตฺถิ อทุกฺขมสุขาย เวทนาย สมฺปยุตฺโต ฯเปฯ.\csromanspacer{} เอวํ ทสวิเธน วิญฺญาณกฺขนฺโธ.}

\proseitem{131}{เอกวิเธน วิญฺญาณกฺขนฺโธ\csromandash{}ผสฺสสมฺปยุตฺโต.}

\prose{ทุวิเธน วิญฺญาณกฺขนฺโธ\csromandash{}อตฺถิ น เหตุ สเหตุโก, อตฺถิ น เหตุ- อเหตุโก.}

\prose{ติวิเธน วิญฺญาณกฺขนฺโธ\csromandash{}อตฺถิ วิปาโก, อตฺถิ วิปากธมฺมธมฺโม, อตฺถิ เนววิปากนวิปากธมฺมธมฺโม ฯเปฯ.\csromanspacer{} เอวํ ทสวิเธน วิญฺญาณกฺขนฺโธ.}

\proseitem{132}{เอกวิเธน วิญฺญาณกฺขนฺโธ\csromandash{}ผสฺสสมฺปยุตฺโต.}

\prose{ทุวิเธน วิญฺญาณกฺขนฺโธ\csromandash{}อตฺถิ โลกิโย, อตฺถิ โลกุตฺตโร.}

\prose{\csromanfolio{54}ติวิเธน วิญฺญาณกฺขนฺโธ\csromandash{}อตฺถิ อุปาทินฺนุปาทานิโย, อตฺถิ อนุปาทินฺนุปาทานิโย, อตฺถิ อนุปาทินฺน-อนุปาทานิโย ฯเปฯ.\csromanspacer{} เอวํ ทสวิเธน วิญฺญาณกฺขนฺโธ.}

\proseitem{133}{เอกวิเธน วิญฺญาณกฺขนฺโธ\csromandash{}ผสฺสสมฺปยุตฺโต.}

\prose{ทุวิเธน วิญฺญาณกฺขนฺโธ\csromandash{}อตฺถิ เกนจิ วิญฺเญยฺโย, อตฺถิ เกนจิ น วิญฺเญยฺโย.}

\prose{ติวิเธน วิญฺญาณกฺขนฺโธ\csromandash{}อตฺถิ สํกิลิฏฺฐสํกิเลสิโก, อตฺถิ อสํกิลิฏฺฐสํกิเลสิโก, อตฺถิ อสํกิลิฏฺฐ-อสํกิเลสิโก ฯเปฯ.\csromanspacer{} เอวํ ทสวิเธน วิญฺญาณกฺขนฺโธ.}

\proseitem{134}{เอกวิเธน วิญฺญาณกฺขนฺโธ\csromandash{}ผสฺสสมฺปยุตฺโต.}

\prose{ทุวิเธน วิญฺญาณกฺขนฺโธ\csromandash{}อตฺถิ สาสโว, อตฺถิ อนาสโว.}

\prose{ติวิเธน วิญฺญาณกฺขนฺโธ\csromandash{}อตฺถิ สวิตกฺกสวิจาโร, อตฺถิ อวิตกฺกวิจารมตฺโต, อตฺถิ อวิตกฺก-อวิจาโร ฯเปฯ.\csromanspacer{} เอวํ ทสวิเธน วิญฺญาณกฺขนฺโธ.}

\proseitem{135}{เอกวิเธน วิญฺญาณกฺขนฺโธ\csromandash{}ผสฺสสมฺปยุตฺโต.}

\prose{ทุวิเธน วิญฺญาณกฺขนฺโธ\csromandash{}อตฺถิ อาสวสมฺปยุตฺโต, อตฺถิ อาสววิปฺปยุตฺโต.}

\prose{ติวิเธน วิญฺญาณกฺขนฺโธ\csromandash{}อตฺถิ ปีติสหคโต, อตฺถิ สุขสหคโต, อตฺถิ อุเปกฺขาสหคโต ฯเปฯ.\csromanspacer{} เอวํ ทสวิเธน วิญฺญาณกฺขนฺโธ.}

\proseitem{136}{เอกวิเธน วิญฺญาณกฺขนฺโธ\csromandash{}ผสฺสสมฺปยุตฺโต.}

\prose{ทุวิเธน วิญฺญาณกฺขนฺโธ\csromandash{}อตฺถิ อาสววิปฺปยุตฺตสาสโว, อตฺถิ อาสววิปฺปยุตฺต-อนาสโว.}

\prose{ติวิเธน วิญฺญาณกฺขนฺโธ\csromandash{}อตฺถิ ทสฺสเนน ปหาตพฺโพ, อตฺถิ ภาวนาย ปหาตพฺโพ, อตฺถิ เนว ทสฺสเนน น ภาวนาย ปหาตพฺโพ ฯเปฯ.\csromanspacer{} เอวํ ทสวิเธน วิญฺญาณกฺขนฺโธ.}

\proseitem{137}{\csromanfolio{55}เอกวิเธน วิญฺญาณกฺขนฺโธ\csromandash{}ผสฺสสมฺปยุตฺโต.}

\prose{ทุวิเธน วิญฺญาณกฺขนฺโธ\csromandash{}อตฺถิ สํโยชนิโย, อตฺถิ อสํโยชนิโย.}

\prose{ติวิเธน วิญฺญาณกฺขนฺโธ\csromandash{}อตฺถิ ทสฺสเนน ปหาตพฺพเหตุโก, อตฺถิ ภาวนาย ปหาตพฺพเหตุโก, อตฺถิ เนว ทสฺสเนน น ภาวนาย ปหาตพฺพเหตุโก ฯเปฯ.\csromanspacer{} เอวํ ทสวิเธน วิญฺญาณกฺขนฺโธ.}

\proseitem{138}{เอกวิเธน วิญฺญาณกฺขนฺโธ\csromandash{}ผสฺสสมฺปยุตฺโต.}

\prose{ทุวิเธน วิญฺญาณกฺขนฺโธ\csromandash{}อตฺถิ สํโยชนสมฺปยุตฺโต, อตฺถิ สํโยชนวิปฺปยุตฺโต.}

\prose{ติวิเธน วิญฺญาณกฺขนฺโธ\csromandash{}อตฺถิ อาจยคามี, อตฺถิ อปจยคามี, อตฺถิ เนวาจยคามินาปจยคามี ฯเปฯ.\csromanspacer{} เอวํ ทสวิเธน วิญฺญาณกฺขนฺโธ.}

\proseitem{139}{เอกวิเธน วิญฺญาณกฺขนฺโธ\csromandash{}ผสฺสสมฺปยุตฺโต.}

\prose{ทุวิเธน วิญฺญาณกฺขนฺโธ\csromandash{}อตฺถิ สํโยชนวิปฺปยุตฺตสํโยชนิโย, อตฺถิ สํโยชนวิปฺปยุตฺต-อสํโยชนิโย.}

\prose{ติวิเธน วิญฺญาณกฺขนฺโธ\csromandash{}อตฺถิ เสกฺโข, อตฺถิ อเสกฺโข, อตฺถิ เนวเสกฺขนาเสกฺโข ฯเปฯ.\csromanspacer{} เอวํ ทสวิเธน วิญฺญาณกฺขนฺโธ.}

\proseitem{140}{เอกวิเธน วิญฺญาณกฺขนฺโธ\csromandash{}ผสฺสสมฺปยุตฺโต.}

\prose{ทุวิเธน วิญฺญาณกฺขนฺโธ\csromandash{}อตฺถิ คนฺถนิโย, อตฺถิ อคนฺถนิโย.}

\prose{ติวิเธน วิญฺญาณกฺขนฺโธ\csromandash{}อตฺถิ ปริตฺโต, อตฺถิ มหคฺคโต, อตฺถิ อปฺปมาโณ ฯเปฯ.\csromanspacer{} เอวํ ทสวิเธน วิญฺญาณกฺขนฺโธ.}

\proseitem{141}{เอกวิเธน วิญฺญาณกฺขนฺโธ\csromandash{}ผสฺสสมฺปยุตฺโต.}

\prose{ทุวิเธน วิญฺญาณกฺขนฺโธ\csromandash{}อตฺถิ คนฺถสมฺปยุตฺโต, อตฺถิ คนฺถวิปฺปยุตฺโต.}

\prose{ติวิเธน วิญฺญาณกฺขนฺโธ\csromandash{}อตฺถิ ปริตฺตารมฺมโณ, อตฺถิ มหคฺคตารมฺมโณ, อตฺถิ อปฺปมาณารมฺมโณ ฯเปฯ.\csromanspacer{} เอวํ ทสวิเธน วิญฺญาณกฺขนฺโธ.}

\proseitem{142}{\csromanfolio{56}เอกวิเธน วิญฺญาณกฺขนฺโธ\csromandash{}ผสฺสสมฺปยุตฺโต.}

\prose{ทุวิเธน วิญฺญาณกฺขนฺโธ\csromandash{}อตฺถิ คนฺถวิปฺปยุตฺตคนฺถนิโย, อตฺถิ คนฺถวิปฺปยุตฺต-อคนฺถนิโย.}

\prose{ติวิเธน วิญฺญาณกฺขนฺโธ\csromandash{}อตฺถิ หีโน, อตฺถิ มชฺฌิโม, อตฺถิ ปณีโต ฯเปฯ.\csromanspacer{} เอวํ ทสวิเธน วิญฺญาณกฺขนฺโธ.}

\proseitem{143}{เอกวิเธน วิญฺญาณกฺขนฺโธ\csromandash{}ผสฺสสมฺปยุตฺโต.}

\prose{ทุวิเธน วิญฺญาณกฺขนฺโธ\csromandash{}อตฺถิ โอฆนิโย, อตฺถิ อโนฆนิโย.}

\prose{ติวิเธน วิญฺญาณกฺขนฺโธ\csromandash{}อตฺถิ มิจฺฉตฺตนิยโต, อตฺถิ สมฺมตฺตนิยโต, อตฺถิ อนิยโต ฯเปฯ.\csromanspacer{} เอวํ ทสวิเธน วิญฺญาณกฺขนฺโธ.}

\proseitem{144}{เอกวิเธน วิญฺญาณกฺขนฺโธ\csromandash{}ผสฺสสมฺปยุตฺโต.}

\prose{ทุวิเธน วิญฺญาณกฺขนฺโธ\csromandash{}อตฺถิ โอฆสมฺปยุตฺโต, อตฺถิ โอฆวิปฺปยุตฺโต.}

\prose{ติวิเธน วิญฺญาณกฺขนฺโธ\csromandash{}อตฺถิ มคฺคารมฺมโณ, อตฺถิ มคฺคเหตุโก, อตฺถิ มคฺคาธิปติ ฯเปฯ.\csromanspacer{} เอวํ ทสวิเธน วิญฺญาณกฺขนฺโธ.}

\proseitem{145}{เอกวิเธน วิญฺญาณกฺขนฺโธ\csromandash{}ผสฺสสมฺปยุตฺโต.}

\prose{ทุวิเธน วิญฺญาณกฺขนฺโธ\csromandash{}อตฺถิ โอฆวิปฺปยุตฺต-โอฆนิโย, อตฺถิ โอฆวิปฺปยุตฺต-อโนฆนิโย.}

\prose{ติวิเธน วิญฺญาณกฺขนฺโธ\csromandash{}อตฺถิ อุปฺปนฺโน, อตฺถิ อนุปฺปนฺโน, อตฺถิ อุปฺปาที ฯเปฯ.\csromanspacer{} เอวํ ทสวิเธน วิญฺญาณกฺขนฺโธ.}

\proseitem{146}{เอกวิเธน วิญฺญาณกฺขนฺโธ\csromandash{}ผสฺสสมฺปยุตฺโต.}

\prose{ทุวิเธน วิญฺญาณกฺขนฺโธ\csromandash{}อตฺถิ โยคนิโย, อตฺถิ อโยคนิโย.}

\prose{ติวิเธน วิญฺญาณกฺขนฺโธ\csromandash{}อตฺถิ อตีโต, อตฺถิ อนาคโต, อตฺถิ ปจฺจุปฺปนฺโน ฯเปฯ.\csromanspacer{} เอวํ ทสวิเธน วิญฺญาณกฺขนฺโธ.}

\proseitem{147}{เอกวิเธน วิญฺญาณกฺขนฺโธ\csromandash{}ผสฺสสมฺปยุตฺโต.}

\prose{ทุวิเธน วิญฺญาณกฺขนฺโธ\csromandash{}อตฺถิ โยคสมฺปยุตฺโต, อตฺถิ โยควิปฺปยุตฺโต.}

\prose{\csromanfolio{57}ติวิเธน วิญฺญาณกฺขนฺโธ\csromandash{}อตฺถิ อตีตารมฺมโณ, อตฺถิ อนาคตารมฺมโณ, อตฺถิ ปจฺจุปฺปนฺนารมฺมโณ ฯเปฯ.\csromanspacer{} เอวํ ทสวิเธน วิญฺญาณกฺขนฺโธ.}

\proseitem{148}{เอกวิเธน วิญฺญาณกฺขนฺโธ\csromandash{}ผสฺสสมฺปยุตฺโต.}

\prose{ทุวิเธน วิญฺญาณกฺขนฺโธ\csromandash{}อตฺถิ โยควิปฺปยุตฺตโยคนิโย, อตฺถิ โยควิปฺปยุตฺต-อโยคนิโย.}

\prose{ติวิเธน วิญฺญาณกฺขนฺโธ\csromandash{}อตฺถิ อชฺฌตฺโต, อตฺถิ พหิทฺโธ, อตฺถิ อชฺฌตฺตพหิทฺโธ ฯเปฯ.\csromanspacer{} เอวํ ทสวิเธน วิญฺญาณกฺขนฺโธ.}

\proseitem{149}{เอกวิเธน วิญฺญาณกฺขนฺโธ\csromandash{}ผสฺสสมฺปยุตฺโต.}

\prose{ทุวิเธน วิญฺญาณกฺขนฺโธ\csromandash{}อตฺถิ นีวรณิโย, อตฺถิ อนีวรณิโย.}

\prose{ติวิเธน วิญฺญาณกฺขนฺโธ\csromandash{}อตฺถิ อชฺฌตฺตารมฺมโณ, อตฺถิ พหิทฺธารมฺมโณ, อตฺถิ อชฺฌตฺตพหิทฺธารมฺมโณ ฯเปฯ.\csromanspacer{} เอวํ ทสวิเธน วิญฺญาณกฺขนฺโธ.\csromanspacer{} อุภโตวฑฺฒกํ.}

\prose{สตฺตวิเธน วิญฺญาณกฺขนฺโธ\csromandash{}อตฺถิ กุสโล, อตฺถิ อกุสโล, อตฺถิ อพฺยากโต, อตฺถิ กามาวจโร, อตฺถิ รูปาวจโร, อตฺถิ อรูปาวจโร, อตฺถิ อปริยาปนฺโน.\csromanspacer{} เอวํ สตฺตวิเธน วิญฺญาณกฺขนฺโน.}

\prose{อปโรปิ สตฺตวิเธน วิญฺญาณกฺขนฺโธ\csromandash{}อตฺถิ สุขาย เวทนาย สมฺปยุตฺโต, อตฺถิ ทุกฺขาย เวทนาย สมฺปยุตฺโต, อตฺถิ อทุกฺขมสุขาย เวทนาย สมฺปยุตฺโต, อตฺถิ กามาวจโร, อตฺถิ รูปาวจโร, อตฺถิ อรูปาวจโร, อตฺถิ อปริยาปนฺโน ฯเปฯ อตฺถิ อชฺฌตฺตารมฺมโณ, อตฺถิ พหิทฺธารมฺมโณ, อตฺถิ อชฺฌตฺตพหิทฺธารมฺมโณ, อตฺถิ กามาวจโร, อตฺถิ รูปาวจโร, อตฺถิ อรูปาวจโร, อตฺถิ อปริยาปนฺโน.\csromanspacer{} เอวํ สตฺตวิเธน วิญฺญาณกฺขนฺโธ.}

\prose{จตุวีสติวิเธน วิญฺญาณกฺขนฺโธ\csromandash{}จกฺขุสมฺผสฺสปจฺจยา วิญฺญาณกฺขนฺโธ อตฺถิ กุสโล, อตฺถิ อกุสโล, อตฺถิ อพฺยากโต.\csromanspacer{} โสตสมฺผสฺสปจฺจยา ฯเปฯ.\csromanspacer{} ฆานสมฺผสฺสปจฺจยา ฯเปฯ.\csromanspacer{} ชิวฺหาสมฺผสฺสปจฺจยา ฯเปฯ.\csromanspacer{} กายสมฺผสฺสปจฺจยา ฯเปฯ.\csromanspacer{} มโนสมฺผสฺสปจฺจยา วิญฺญาณกฺขนฺโธ อตฺถิ กุสโล, อตฺถิ อกุสโล, อตฺถิ อพฺยากโต.\csromanspacer{} จกฺขุวิญฺญาณํ, โสตวิญฺญาณํ, \csromanfolio{58}ฆานวิญฺญาณํ, ชิวฺหาวิญฺญาณํ, กายวิญฺญาณํ, มโนวิญฺญาณํ.\csromanspacer{} เอวํ จตุวีสติวิเธน วิญฺญาณกฺขนฺโธ.}

\prose{อปโรปิ จตุวีสติวิเธน วิญฺญาณกฺขนฺโธ\csromandash{} จกฺขุสมฺผสฺสปจฺจยา วิญฺญาณกฺขนฺโธ อตฺถิ สุขาย เวทนาย สมฺปยุตฺโต, อตฺถิ ทุกฺขาย เวทนาย สมฺปยุตฺโต, อตฺถิ อทุกฺขมสุขาย เวทนาย สมฺปยุตฺโต ฯเปฯ อตฺถิ อชฺฌตฺตารมฺมโณ, อตฺถิ พหิทฺธารมฺมโณ, อตฺถิ อชฺฌตฺตพหิทฺธารมฺมโณ.\csromanspacer{} จกฺขุวิญฺญาณํ ฯเปฯ.\csromanspacer{} กายวิญฺญาณํ, มโนวิญฺญาณํ.\csromanspacer{} โสตสมฺผสฺสปจฺจยา ฯเปฯ.\csromanspacer{} ฆานสมฺผสฺสปจฺจยา ฯเปฯ.\csromanspacer{} ชิวฺหาสมฺผสฺสปจฺจยา ฯเปฯ.\csromanspacer{} กายสมฺผสฺสปจฺจยา ฯเปฯ.\csromanspacer{} มโนสมฺผสฺสปจฺจยา วิญฺญาณกฺขนฺโธ อตฺถิ อชฺฌตฺตารมฺมโณ, อตฺถิ พหิทฺธารมฺมโณ, อตฺถิ อชฺฌตฺตพหิทฺธารมฺมโณ.\csromanspacer{} จกฺขุวิญฺญาณํ, โสตวิญฺญาณํ, ฆานวิญฺญาณํ, ชิวฺหาวิญฺญาณํ, กายวิญฺญาณํ, มโนวิญฺญาณํ.\csromanspacer{} เอวํ จตุวีสติวิเธน วิญฺญาณกฺขนฺโธ.}

\prose{ติํสติวิเธน วิญฺญาณกฺขนฺโธ\csromandash{}จกฺขุสมฺผสฺสปจฺจยา วิญฺญาณกฺขนฺโธ อตฺถิ กามาวจโร, อตฺถิ รูปาวจโร, อตฺถิ อรูปาวจโร, อตฺถิ อปริยาปนฺโน.\csromanspacer{} โสตสมฺผสฺสปจฺจยา ฯเปฯ.\csromanspacer{} ฆานสมฺผสฺสปจฺจยา ฯเปฯ.\csromanspacer{} ชิวฺหาสมฺผสฺสปจฺจยา ฯเปฯ.\csromanspacer{} กายสมฺผสฺสปจฺจยา ฯเปฯ.\csromanspacer{} มโนสมฺผสฺสปจฺจยา วิญฺญาณกฺขนฺโธ อตฺถิ กามาวจโร, อตฺถิ รูปาวจโร, อตฺถิ อรูปาวจโร, อตฺถิ อปริยาปนฺโน.\csromanspacer{} จกฺขุวิญฺญาณํ, โสตวิญฺญาณํ, ฆานวิญฺญาณํ, ชิวฺหาวิญฺญาณํ, กายวิญฺญาณํ, มโนวิญฺญาณํ.\csromanspacer{} เอวํ ติํสติวิเธน วิญฺญาณกฺขนฺโธ.}

\prose{พหุวิเธน วิญฺญาณกฺขนฺโธ\csromandash{}จกฺขุสมฺผสฺสปจฺจยา วิญฺญาณกฺขนฺโธ อตฺถิ กุสโล, อตฺถิ อกุสโล, อตฺถิ อพฺยากโต, อตฺถิ กามาวจโร, อตฺถิ รูปาวจโร, อตฺถิ อรูปาวจโร, อตฺถิ อปริยาปนฺโน.\csromanspacer{} จกฺขุวิญฺญาณํ ฯเปฯ มโนวิญฺญาณํ, โสตสมฺผสฺสปจฺจยา ฯเปฯ.\csromanspacer{} ฆานสมฺผสฺสปจฺจยา ฯเปฯ.\csromanspacer{} ชิวฺหาสมฺผสฺสปจฺจยา ฯเปฯ.\csromanspacer{} กายสมฺผสฺสปจฺจยา ฯเปฯ.\csromanspacer{} มโนสมฺผสฺสปจฺจยา วิญฺญาณกฺขนฺโธ อตฺถิ กุสโล, อตฺถิ อกุสโล, อตฺถิ อพฺยากโต, อตฺถิ กามาวจโร, อตฺถิ รูปาวจโร, อตฺถิ อรูปาวจโร, อตฺถิ อปริยาปนฺโน.\csromanspacer{} จกฺขุวิญฺญาณํ ฯเปฯ มโนวิญฺญาณํ.\csromanspacer{} เอวํ พหุวิเธน วิญฺญาณกฺขนฺโธ.}

\prose{\csromanfolio{59}อปโรปิ พหุวิเธน วิญฺญาณกฺขนฺโธ\csromandash{}จกฺขุสมฺผสฺสปจฺจยา วิญฺญาณกฺขนฺโธ อตฺถิ สุขาย เวทนา สมฺปยุตฺโต ฯเปฯ อตฺถิ อชฺฌตฺตารมฺมโณ, อตฺถิ พหิทฺธารมฺมโณ, อตฺถิ อชฺฌตฺตพหิทฺธารมฺมโณ, อตฺถิ กามาวจโร, อตฺถิ รูปาวจโร, อตฺถิ อรูปาวจโร, อตฺถิ อปริยาปนฺโน.\csromanspacer{} โสตสมฺผสฺสปจฺจยา ฯเปฯ.\csromanspacer{} ฆานสมฺผสฺสปจฺจยา ฯเปฯ.\csromanspacer{} ชิวฺหาสมฺผสฺสปจฺจยา ฯเปฯ.\csromanspacer{} กายสมฺผสฺสปจฺจยา ฯเปฯ.\csromanspacer{} มโนสมฺปสฺสปจฺจยา วิญฺญาณกฺขนฺโธ อตฺถิ สุขาย เวทนาย สมฺปยุตฺโต ฯเปฯ อตฺถิ อชฺฌตฺตารมฺมโณ, อตฺถิ พหิทฺธารมฺมโณ, อตฺถิ อชฺฌตฺตพหิทฺธารมฺมโณ, อตฺถิ กามาวจโร, อตฺถิ รูปาวจโร, อตฺถิ อรูปาวจโร, อตฺถิ อปริยาปนฺโน.\csromanspacer{} จกฺขุวิญฺญาณํ, โสตวิญฺญาณํ, ฆานวิญฺญาณํ, ชิวฺหาวิญฺญาณํ, กายวิญฺญาณํ, มโนวิญฺญาณํ.\csromanspacer{} เอวํ พหุวิเธน วิญฺญาณกฺขนฺโธ.}

\vspace{\tipitakaparskip}
\csromancenter{อยํ วุจฺจติ วิญฺญาณกฺขนฺโธ.}
\csromanheader{2}{อภิธมฺมภาชนียํ.}
\csromansectionrule
\csromanheader{2}{3. ปญฺหาปุจฺฉก}
\proseitem{150}{ปญฺจกฺขนฺธา รูปกฺขนฺโธ เวทนากฺขนฺโธ สญฺญากฺขนฺโธ สงฺขารกฺขนฺโธ วิญฺญาณกฺขนฺโธ.}

\proseitem{151}{ปญฺจนฺนํ ขนฺธานํ กติ กุสลา, กติ อกุสลา, กติ อพฺยากตา ฯเปฯ กติ สรณา, กติ อรณา.}

\csromanmatikaanchor{mtk.14}
\csromanmatikaanchor{mtk.15}
\csromantocmarknopage{section}{3. ปญฺหาปุจฺฉก}{mtk.14}
\bookmark[dest={mtk.14},level=1]{3. ปญฺหาปุจฺฉก}
\csromantocmarkref{subsection}{1. ติก}{mtk.15}
\bookmark[dest={mtk.15},level=2]{1. ติก}
\csromanheader{2}{1. ติก}
\proseitem{152}{รูปกฺขนฺโธ อพฺยากโต.\csromanspacer{} จตฺตาโร ขนฺธา สิยา กุสลา, สิยา อกุสลา, สิยา อพฺยากตา.\csromanspacer{} เทฺว ขนฺธา น วตฺตพฺพา “สุขาย เวทนาย สมฺปยุตฺตา”ติปิ, “ทุกฺขาย เวทนาย สมฺปยุตฺตา”ติปิ, “อทุกฺขมสุขาย เวทนาย สมฺปยุตฺตา”ติปิ.\csromanspacer{} ตโย ขนฺธา สิยา สุขาย เวทนาย สมฺปยุตฺตา, สิยา ทุกฺขาย เวทนาย สมฺปยุตฺตา, สิยา อทุกฺขมสุขาย เวทนาย สมฺปยุตฺตา.\csromanspacer{} รูปกฺขนฺโธ เนววิปากนวิปากธมฺมธมฺโม.\csromanspacer{} จตฺตาโร ขนฺธา สิยา วิปากา, สิยา วิปากธมฺมธมฺมา, สิยา นิววิปากนวิปากธมฺมธมฺมา.\csromanspacer{} รูปกฺขนฺโธ สิยา อุปาทินฺนุปาทานิโย, สิย \csromanfolio{60}อนุปาทินฺนุปาทานิโย, จตฺตาโร ขนฺธา สิยา อุปาทินฺนุปาทานิยา, สิยา อนุปาทินฺนุปาทานิยา, สิยา อนุปาทินฺน-อนุปาทานิยา.}

\prose{รูปกฺขนฺโธ อสํกิลิฏฺฐสํกิเลสิโก.\csromanspacer{} จตฺตาโร ขนฺธา สิยา สํกิลิฏฺฐสํกิเลสิกา, สิยา อสํกิลิฏฺฐสํกิเลสิกา, สิยา อสํกิลิฏฺฐ- อสํกิเลสิกา.\csromanspacer{} รูปกฺขนฺโธ อวิตกฺก-อวิจาโร.\csromanspacer{} ตโย ขนฺธา สิยา สวิตกฺกสวิจารา, สิยา อวิตกฺกวิจารมตฺตา, สิยา อวิตกฺก-อวิจารา.\csromanspacer{} สงฺขารกฺขนฺโธ สิยา สวิตกฺกสวิจาโร, สิยา อวิตกฺกวิจารมตฺโต, สิยา อวิตกฺก-อวิจาโร, สิยา น วตฺตพฺโพ “สวิตกฺกสวิจาโร”ติปิ, “อวิตกฺกวิจารมตฺโต”ติปิ, “อวิตกฺก-อวิจาโร”ติปิ.\csromanspacer{} รูปกฺขนฺโธ น วตฺตพฺโพ “ปีติสหคโต”ติปิ, “สุขสหคโต”ติปิ, “อุเปกฺขาสหคโต”ติปิ.\csromanspacer{} เวทนากฺขนฺโธ สิยา ปีติสหคโต, น สุขสหคโต, น อุเปกฺขาสหคโต, สิยา น วตฺตพฺโพ “ปีติสหคโต”ติ.\csromanspacer{} ตโย ขนฺธา สิยา ปีติสหคตา, สิยา สุขสหคตา, สิยา อุเปกฺขาสหคตา, สิยา น วตฺตพฺพา “ปีติสหคตา”ติปิ, “สุขสหคตา”ติปิ, “อุเปกฺขาสหคตา”ติปิ.}

\prose{รูปกฺขนฺโธ เนว ทสฺสเนน น ภาวนาย ปหาตพฺโพ.\csromanspacer{} จตฺตาโร ขนฺธา สิยา ทสฺสเนน ปหาตพฺพา, สิยา ภาวนาย ปหาตพฺพา, สิยา เนว ทสฺสเนน น ภาวนาย ปหาตพฺพา.\csromanspacer{} รูปกฺขนฺโธ เนว ทสฺสเนน น ภาวนาย ปหาตพฺพเหตุโก.\csromanspacer{} จตฺตาโร ขนฺธา สิยา ทสฺสเนน ปหาตพฺพเหตุกา, สิยา ภาวนาย ปหาตพฺพเหตุกา, สิยา เนว ทสฺสเนน น ภาวนาย ปหาตพฺพเหตุกา.\csromanspacer{} รูปกฺขนฺโธ เนวาจยคามินาปจยคามี.\csromanspacer{} จตฺตาโร ขนฺธา สิยา อาจยคามิโน, สิยา อปจยคามิโน, สิยา เนวาจยคามินาปจยคามิโน.\csromanspacer{} รูปกฺขนฺโธ เนวเสกฺขนาเสกฺโข.\csromanspacer{} จตฺตาโร ขนฺธา สิยา เสกฺขา, สิยา อเสกฺขา, สิยา เนวเสกฺขนาเสกฺขา.\csromanspacer{} รูปกฺขนฺโธ ปริตฺโต.\csromanspacer{} จตฺตาโร ขนฺธา สิยา ปริตฺตา, สิยา มหคฺคตา, สิยา อปฺปมาณา.\csromanspacer{} รูปกฺขนฺโธ อนารมฺมโณ.\csromanspacer{} จตฺตาโร ขนฺธา สิยา ปริตฺตารมฺมณา, สิยา มหคฺคตารมฺมณา, สิยา อปฺปมาณารมฺมณา, สิยา น วตฺตพฺพา “ปริตฺตารมฺมณา”ติปิ, “มหคฺคตารมฺมณา”ติปิ, “อปฺปมาณารมฺมณา”ติปิ.\csromanspacer{} รูปกฺขนฺโธ มชฺฌิโม.\csromanspacer{} จตฺตาโร ขนฺธา สิยา หีนา, สิยา มชฺฌิมา, สิยา ปณีตา.\csromanspacer{} รูปกฺขนฺโธ \csromanfolio{61}อนิยโต.\csromanspacer{} จตฺตาโร ขนฺธา สิยา มิจฺฉตฺตนิยตา, สิยา สมฺมตฺตนิยตา, สิยา อนิยตา.}

\prose{รูปกฺขนฺโธ อนารมฺมโณ.\csromanspacer{} จตฺตาโร ขนฺธา สิยา มคฺคารมฺมณา, สิยา มคฺคเหตุกา, สิยา มคฺคาธิปติโน, สิยา น วตฺตพฺพา “มคฺคารมฺมณา”ติปิ, “มคฺคเหตุกา”ติปิ, “มคฺคาธิปติโน”ติปิ, สิยา อุปฺปนฺนา, สิยา อนุปฺปนฺนา, สิยา อุปฺปาทิโน, สิยา อตีตา, สิยา อนาคตา, สิยา ปจฺจุปฺปนฺนา.\csromanspacer{} รูปกฺขนฺโธ อนารมฺมโณ.\csromanspacer{} จตฺตาโร ขนฺธา สิยา อตีตารมฺมณา, สิยา อนาคตารมฺมณา, สิยา ปจฺจุปฺปนฺนารมฺมณา, สิยา น วตฺตพฺพา “อตีตารมฺมณา”ติปิ, “อนาคตารมฺมณา”ติปิ, “ปจฺจุปฺปนฺนารมฺมณา”ติปิ, สิยา อชฺฌตฺตา, สิยา พหิทฺธา, สิยา อชฺฌตฺตพหิทฺธา.\csromanspacer{} รูปกฺขนฺโธ อนารมฺมโณ.\csromanspacer{} จตฺตาโร ขนฺธา สิยา อชฺฌตฺตารมฺมณา, สิยา พหิทฺธารมฺมณา, สิยา อชฺฌตฺตพหิทฺธารมฺมณา, สิยา น วตฺตพฺพา “อชฺฌตฺตารมฺมณา”ติปิ, “พหิทฺธารมฺมณา”ติปิ, “อชฺฌตฺตพหิทฺธารมฺมณา”ติปิ.\csromanspacer{} จตฺตาโร ขนฺธา อนิทสฺสน-อปฺปฏิฆา.\csromanspacer{} รูปกฺขนฺโธ สิยา สนิทสฺสนสปฺปฏิโฆ, สิยา อนิทสฺสนสปฺปฏิโฆ, สิยา อนิทสฺสน-อปฺปฏิโฆ.}

\csromanmatikaanchor{mtk.16}
\csromantocmarkref{subsection}{2. ทุก}{mtk.16}
\bookmark[dest={mtk.16},level=2]{2. ทุก}
\csromanheader{2}{2. ทุก}
\proseitem{153}{จตฺตาโร ขนฺธา น เหตู.\csromanspacer{} สงฺขารกฺขนฺโธ สิยา เหตุ, สิยา น เหตุ.\csromanspacer{} รูปกฺขนฺโธ อเหตุโก.\csromanspacer{} จตฺตาโร ขนฺธา สิยา สเหตุกา, สิยา อเหตุกา.\csromanspacer{} รูปกฺขนฺโธ เหตุวิปฺปยุตฺโต.\csromanspacer{} จตฺตาโร ขนฺธา สิยา เหตุสมฺปยุตฺตา, สิยา เหตุวิปฺปยุตฺตา.\csromanspacer{} รูปกฺขนฺโธ น วตฺตพฺโพ “เหตุ เจว สเหตุโก จา”ติปิ, “สเหตุโก เจว น จ เหตู”ติปิ.\csromanspacer{} ตโย ขนฺธา น วตฺตพฺพา “เหตู เจว สเหตุกา จา”ติ, สิยา สเหตุกา เจว น จ เหตู, สิยา น วตฺตพฺพา “สเหตุกา เจว น จ เหตู”ติ.\csromanspacer{} สงฺขารกฺขนฺโธ สิยา เหตุ เจว สเหตุโก จ, สิยา สเหตุโก เจว น จ เหตุ, สิยา น วตฺตพฺโพ “เหตุ เจว สเหตุโก จา”ติปิ, “สเหตุโก เจว น ว เหตู”ติปิ.\csromanspacer{} รูปกฺขนฺโธ น วตฺตพฺโพ “เหตุ เจว เหตุสมฺปยุตฺโต จา”ติปิ, “เหตุสมฺปยุตฺโต เจว น จ เหตู”ติปิ.\csromanspacer{} ตโย ขนฺธา น วตฺตพฺพา “เหตู เจว เหตุสมฺปยุตฺตา จา”ติ.\csromanspacer{} สิยา เหตุสมฺปยุตฺตา เจว น จ เหตู, สิยา น วตฺตพฺพา “เหตุสมฺปยุตฺตา เจว น จ เหตู”ติ.\csromanspacer{} สงฺขารกฺขนฺโธ สิยา เหตุ เจว เหตุสมฺปยุตฺโต จ, สิยา เหตุสมฺปยุตฺโต เจว \csromanfolio{62}น จ เหตุ, สิยา น วตฺตพฺโพ “เหตุ เจว เหตุสมฺปยุตฺโต จา”ติปิ, เหตุสมฺปยุตฺโต เจว น จ เหตู”ติปิ.\csromanspacer{} รูปกฺขนฺโธ น เหตุ อเหตุโก.\csromanspacer{} ตโย ขนฺธา สิยา น เหตุ สเหตุกา, สิยา น เหตุ อเหตุกา.\csromanspacer{} สงฺขารกฺขนฺโธ สิยา น เหตุ สเหตุโก, สิยา น เหตุ อเหตุโก, สิยา น วตฺตพฺโพ “น เหตุ สเหตุโก”ติปิ, “น เหตุ อเหตุโก”ติปิ. (1)}

\prose{สปฺปจฺจยา, สงฺขตา.}

\prose{จตฺตาโร ขนฺธา อนิทสฺสนา.\csromanspacer{} รูปกฺขนฺโธ สิยา สนิทสฺสโน สิยา อนิทสฺสโน, จตฺตาโร ขนฺธา อปฺปฏิฆา.\csromanspacer{} รูปกฺขนฺโธ สิยา สปฺปฏิโฆ, สิยา อปฺปฏิโฆ.\csromanspacer{} รูปกฺขนฺโธ รูปํ.\csromanspacer{} จตฺตาโร ขนฺธา อรูปา.\csromanspacer{} รูปกฺขนฺโธ โลกิโย.\csromanspacer{} จตฺตาโร ขนฺธา สิยา โลกิยา, สิยา โลกุตฺตรา.\csromanspacer{} เกนจิ วิญฺเญยฺยา, เกนจิ น วิญฺเญยฺยา. (2)}

\prose{จตฺตาโร ขนฺธา โน อาสวา.\csromanspacer{} สงฺขารกฺขนฺโธ สิยา อาสโว, สิยา โน อาสโว.\csromanspacer{} รูปกฺขนฺโธ สาสโว.\csromanspacer{} จตฺตาโร ขนฺธา สิยา สาสวา, สิยา อนาสวา.\csromanspacer{} รูปกฺขนฺโธ อาสววิปฺปยุตฺโต.\csromanspacer{} จตฺตาโร ขนฺธา สิยา อาสวสมฺปยุตฺตา, สิยา อาสววิปฺปยุตฺตา.\csromanspacer{} รูปกฺขนฺโธ น วตฺตพฺโพ “อาสโว เจว สาสโว จา”ติ, สาสโว เจว โน จ อาสโว.\csromanspacer{} ตโย ขนฺธา น วตฺตพฺพา “อาสวา เจว สาสวา จา”ติ, สิยา สาสวา เจว โน จ อาสวา, สิยา น วตฺตพฺพา “สาสวา เจว โน จ อาสวา”ติ.\csromanspacer{} สงฺขารกฺขนฺโธ สิยา อาสโว เจว สาสโว จ, สิยา สาสโว เจว โน จ อาสโว, สิยา น วตฺตพฺโพ “อาสโว เจว สาสโว จา”ติปิ, “สาสโว เจว โน จ อาสโว”ติปิ.\csromanspacer{} รูปกฺขนฺโธ น วตฺตพฺโพ “อาสโว เจว อาสวสมฺปยุตฺโต จา”ติปิ, “อาสวสมฺปยุตฺโต เจว โน จ อาสโว”ติปิ.\csromanspacer{} ตโย ขนฺธา น วตฺตพฺพา “อาสวา เจว อาสวสมฺปยุตฺตา จา”ติ, สิยา อาสวสมฺปยุตฺตา เจว โน จ อาสวา, สิยา น วตฺตพฺพา “อาสวสมฺปยุตฺตา เจว โน จ อาสวา”ติ.\csromanspacer{} สงฺขารกฺขนฺโธ สิยา อาสโว เจว อาสวสมฺปยุตฺโต จ, สิยา อาสวสมฺปยุตฺโต เจว โน จ อาสโว, สิยา น วตฺตพฺโพ “อาสโว เจว อาสวสมฺปยุตฺโต จา”ติปิ, “อาสวสมฺปยุตฺโต เจว โน จ อาสโว”ติปิ.\csromanspacer{} รูปกฺขนฺโธ อาสววิปฺปยุตฺตสาสโว.\csromanspacer{} จตฺตาโร ขนฺธา สิยา อาสววิปฺปยุตฺตสาสวา, สิยา อาสววิปฺปยุตฺต-อนาสวา, \csromanfolio{63}สิยา น วตฺตพฺพา “อาสววิปฺปยุตฺตสาสวา”ติปิ, “อาสววิปฺปยุตฺต-อนาสวา”ติปิ. (3)}

\prose{จตฺตาโร ขนฺธา โน สํโยชนา.\csromanspacer{} สงฺขารกฺขนฺโธ สิยา สํโยชนํ, สิยา โน สํโยชนํ.\csromanspacer{} รูปกฺขนฺโธ สํโยชนิโย.\csromanspacer{} จตฺตาโร ขนฺธา สิยา สํโยชนิยา, สิยา อสํโยชนิยา.\csromanspacer{} รูปกฺขนฺโธ สํโยชนวิปฺปยุตฺโต.\csromanspacer{} จตฺตาโร ขนฺธา สิยา สํโยชนสมฺปยุตฺตา, สิยา สํโยชนวิปฺปยุตฺตา.\csromanspacer{} รูปกฺขนฺโธ น วตฺตพฺโพ “สํโยชนญฺเจว สํโยชนิโย จา”ติ, สํโยชนิโย เจว โน จ สํโยชนํ.\csromanspacer{} ตโย ขนฺธา น วตฺตพฺพา “สํโยชนา เจว สํโยชนิยา จา”ติ, สิยา สํโยชนิยา เจว โน จ สํโยชนา, สิยา น วตฺตพฺพา “สํโยชนิยา เจว โน จ สํโยชนา”ติ.\csromanspacer{} สงฺขารกฺขนฺโธ สิยา สํโยชนญฺเจว สํโยชนิโย จ, สิยา สํโยชนิโย เจว โน จ สํโยชนํ, สิยา น วตฺตพฺโพ “สํโยชนญฺเจว สํโยชนิโย จา”ติปิ, “สํโยชนิโย เจว โน จ สํโยชนนฺ”ติปิ.\csromanspacer{} รูปกฺขนฺโธ น วตฺตพฺโพ “สํโยชนญฺเจว สํโยชนสมฺปยุตฺโต จา”ติปิ, “สํโยชนสมฺปยุตฺโต เจว โน จ สํโยชนนฺ”ติปิ.\csromanspacer{} ตโย ขนฺธา น วตฺตพฺพา “สํโยชนา เจว สํโยชนสมฺปยุตฺตา จา”ติ, สิยา สํโยชนสมฺปยุตฺตา เจว โน จ สํโยชนา, สิยา น วตฺตพฺพา “สํโยชนสมฺปยุตฺตา เจว โน จ สํโยชนา”ติ.\csromanspacer{} สงฺขารกฺขนฺโธ สิยา สํโยชนญฺเจว สํโยชนสมฺปยุตฺโต จ, สิยา สํโยชนสมฺปยุตฺโต เจว โน จ สํโยชนํ, สิยา น วตฺตพฺโพ “สํโยชนญฺเจว สํโยชนสมฺปยุตฺโต จา”ติปิ, “สํโยชนสมฺปยุตฺโต เจว โน จ สํโยชนนฺ”ติปิ.\csromanspacer{} รูปกฺขนฺโธ สํโยชนวิปฺปยุตฺตสํโยชนิโย.\csromanspacer{} จตฺตาโร ขนฺธา สิยา สํโยชนวิปฺปยุตฺตสํโยชนิยา, สิยา สํโยชนวิปฺปยุตฺต- อสํโยชนิยา, สิยา น วตฺตพฺพา “สํโยชนวิปฺปยุตฺตสํโยชนิยา “ติปิ, “สํโยชนวิปฺปยุตฺต-อสํโยชนิยา”ติปิ. (4)}

\prose{จตฺตาโร ขนฺธา โน คนฺถา.\csromanspacer{} สงฺขารกฺขนฺโธ สิยา คนฺโถ, สิยา โน คนฺโถ.\csromanspacer{} รูปกฺขนฺโธ คนฺถนิโย.\csromanspacer{} จตฺตาโร ขนฺธา สิยา คนฺถนิยา, สิยา อคนฺถนิยา.\csromanspacer{} รูปกฺขนฺโธ คนฺถวิปฺปยุตฺโต.\csromanspacer{} จตฺตาโร ขนฺธา สิยา คนฺถสมฺปยุตฺตา, สิยา คนฺถวิปฺปยุตฺตา.\csromanspacer{} รูปกฺขนฺโธ น วตฺตพฺโพ “คนฺโถ เจว คนฺถนิโย จา”ติ, คนฺถนิโย เจว โน จ คนฺโถ.\csromanspacer{} ตโย ขนฺธา น \csromanfolio{64}วตฺตพฺพา “คนฺถา เจว คนฺถนิยา จา”ติ, สิยา คนฺถนิยา เจว โน จ คนฺถา, สิยา น วตฺตพฺพา “คนฺถนิยา เจว โน จ คนฺถา”ติ.\csromanspacer{} สงฺขารกฺขนฺโธ สิยา คนฺโถ เจว คนฺถนิโย จ, สิยา คนฺถนิโย เจว โน จ คนฺโถ, สิยา น วตฺตพฺโพ “คนฺโถ เจว คนฺถนิโย จา”ติปิ, “คนฺถนิโย เจว โน จ คนฺโถ”ติปิ.\csromanspacer{} รูปกฺขนฺโธ น วตฺตพฺโพ “คนฺโถ เจว คนฺถสมฺปยุตฺโต จา”ติปิ, “คนฺถสมฺปยุตฺโต เจว โน จ คนฺโถ”ติปิ.\csromanspacer{} ตโย ขนฺธา น วตฺตพฺพา “คนฺถา เจว คนฺถสมฺปยุตฺตา จา”ติ, สิยา คนฺถสมฺปยุตฺตา เจว โน จ คนฺถา, สิยา น วตฺตพฺพา “คนฺถสมฺปยุตฺตา เจว โน จ คนฺถา”ติ.\csromanspacer{} สงฺขารกฺขนฺโธ สิยา คนฺโถ เจว คนฺถสมฺปยุตฺโต จ, สิยา คนฺถสมฺปยุตฺโต เจว โน จ คนฺโถ, สิยา น วตฺตพฺโพ “คนฺโถ เจว คนฺถสมฺปยุตฺโต จา”ติปิ, “คนฺถสมฺปยุตฺโต เจว โน จ คนฺโถ”ติปิ.\csromanspacer{} รูปกฺขนฺโธ คนฺถวิปฺปยุตฺตคนฺถนิโย.\csromanspacer{} จตฺตาโร ขนฺธา สิยา คนฺถวิปฺปยุตฺตคนฺถนิยา, สิยา คนฺถวิปฺปยุตฺต-อคนฺถนิยา, สิยา น วตฺตพฺพา “คนฺถวิปฺปยุตฺตคนฺถนิยา”ติปิ, “คนฺถวิปฺปยุตฺต- อคนฺถนิยา”ติปิ. (5)}

\prose{จตฺตาโร ขนฺธา โน โอฆา ฯเปฯ โน โยคา ฯเปฯ โน นีวรณา.\csromanspacer{} สงฺขารกฺขนฺโธ สิยา นีวรณํ, สิยา โน นีวรณํ.\csromanspacer{} รูปกฺขนฺโธ นีวรณิโย.\csromanspacer{} จตฺตาโร ขนฺธา สิยา นีวรณิยา, สิยา อนีวรณิยา.\csromanspacer{} รูปกฺขนฺโธ นีวรณวิปฺปยุตฺโต.\csromanspacer{} จตฺตาโร ขนฺธา สิยา นีวรณสมฺปยุตฺตา, สิยา นีวรณวิปฺปยุตฺตา.\csromanspacer{} รูปกฺขนฺโธ น วตฺตพฺโพ “นีวรณญฺเจว นีวรณิโย จา”ติ, นีวรณิโย เจว โน จ นีวรณํ.\csromanspacer{} ตโย ขนฺธา น วตฺตพฺพา “นีวรณา เจว นีวรณิยา จา”ติ, สิยา นีวรณิยา เจว โน จ นีวรณา, สิยา น วตฺตพฺพา “นีวรณิยา เจว โน จ นีวรณา”ติ.\csromanspacer{} สงฺขารกฺขนฺโธ สิยา นีวรณญฺเจว นีวรณิโย จ, สิยา นีวรณิโย เจว โน จ นีวรณํ, สิยา น วตฺตพฺโพ “นีวรณญฺเจว นีวรณิโย จา”ติปิ, “นีวรณิโย เจว โน จ นีวรณนฺ”ติปิ, รูปกฺขนฺโธ น วตฺตพฺโพ “นีวรณญฺเจว นีวรณสมฺปยุตฺโต จา”ติปิ, “นีวรณสมฺปยุตฺโต เจว โน จ นีวรณนฺ”ติปิ.\csromanspacer{} ตโย ขนฺธา น วตฺตพฺพา “นีวรณา เจว นีวรณสมฺปยุตฺตา จา”ติ, สิยา นีวรณสมฺปยุตฺตา เจว โน จ นีวรณา, สิยา น วตฺตพฺพา “นีวรณสมฺปยุตฺตา เจว โน จ นีวรณา”ติ.\csromanspacer{} สงฺขารกฺขนฺโธ สิยา นีวรณญฺเจว นีวรณสมฺปยุตฺโต จ, สิยา นีวรณสมฺปยุตฺโต เจว โน จ นีวรณํ, สิยา น วตฺตพฺโพ “นีวรณญฺเจว นีวรณสมฺปยุตฺโต จา”ติปิ, “นีวรณสมฺปยุตฺโต เจว \csromanfolio{65}โน จ นีวรณนฺ”ติปิ.\csromanspacer{} รูปกฺขนฺโธ นีวรณวิปฺปยุตฺตนีวรณิโย.\csromanspacer{} จตฺตาโร ขนฺธา สิยา นีวรณวิปฺปยุตฺตนีวรณิยา, สิยา นีวรณวิปฺปยุตฺต- อนีวรณิยา, สิยา น วตฺตพฺพา “นีวรณวิปฺปยุตฺตนีวรณิยา”ติปิ, “นีวรณวิปฺปยุตฺต-อนีวรณิยา”ติปิ. (8)}

\prose{จตฺตาโร ขนฺธา โน ปรามาสา.\csromanspacer{} สงฺขารกฺขนฺโธ สิยา ปรามาโส, สิยา โน ปรามาโส.\csromanspacer{} รูปกฺขนฺโธ ปรามฏฺโฐ.\csromanspacer{} จตฺตาโร ขนฺธา สิยา ปรามฏฺฐา, สิยา อปรามฏฺฐา.\csromanspacer{} รูปกฺขนฺโธ ปรามาสวิปฺปยุตฺโต.\csromanspacer{} ตโย ขนฺธา สิยา ปรามาสสมฺปยุตฺตา, สิยา ปรามาสวิปฺปยุตฺตา.\csromanspacer{} สงฺขารกฺขนฺโธ สิยา ปรามาสสมฺปยุตฺโต, สิยา ปรามาสวิปฺปยุตฺโต, สิยา น วตฺตพฺโพ “ปรามาสสมฺปยุตฺโต”ติปิ, “ปรามาสวิปฺปยุตฺโต”ติปิ.\csromanspacer{} รูปกฺขนฺโธ น วตฺตพฺโพ “ปรามาโส เจว ปรามฏฺโฐ จา”ติ, ปรามฏฺโฐ เจว โน จ ปรามาโส.\csromanspacer{} ตโย ขนฺธา น วตฺตพฺพา “ปรามาสา เจว ปรามฏฺฐา จา”ติ, สิยา ปรามฏฺฐา เจว โน จ ปรามาสา, สิยา น วตฺตพฺพา “ปรามฏฺฐา เจว โน จ ปรามาสา”ติ.\csromanspacer{} สงฺขารกฺขนฺโธ สิยา ปรามาโส เจว ปรามฏฺโฐ จ, สิยา ปรามฏฺโฐ เจว โน จ ปรามาโส, สิยา น วตฺตพฺโพ “ปรามาโส เจว ปรามฏฺโฐ จา”ติปิ, “ปรามฏฺโฐ เจว โน จ ปรามาโส”ติปิ.\csromanspacer{} รูปกฺขนฺโธ ปรามาสวิปฺปยุตฺตปรามฏฺโฐ.\csromanspacer{} จตฺตาโร ขนฺธา สิยา ปรามาสวิปฺปยุตฺตปรามฏฺฐา, สิยา ปรามาสวิปฺปยุตฺต-อปรามฏฺฐา, สิยา น วตฺตพฺพา “ปรามาสวิปฺปยุตฺตปรามฏฺฐา”ติปิ, “ปรามาสวิปฺปยุตฺต- อปรามฏฺฐา”ติปิ. (9)}

\prose{รูปกฺขนฺโธ อนารมฺมโณ.\csromanspacer{} จตฺตาโร ขนฺธา สารมฺมณา.\csromanspacer{} วิญฺญาณกฺขนฺโธ จิตฺตํ.\csromanspacer{} จตฺตาโร ขนฺธา โน จิตฺตา.\csromanspacer{} ตโย ขนฺธา เจตสิกา.\csromanspacer{} เทฺว ขนฺธา อเจตสิกา.\csromanspacer{} ตโย ขนฺธา จิตฺตสมฺปยุตฺตา.\csromanspacer{} รูปกฺขนฺโธ จิตฺตวิปฺปยุตฺโต.\csromanspacer{} วิญฺญาณกฺขนฺโธ น วตฺตพฺโพ “จิตฺเตน สมฺปยุตฺโต”ติปิ, “จิตฺเตน วิปฺปยุตฺโต”ติปิ.\csromanspacer{} ตโย ขนฺธา จิตฺตสํสฏฺฐา.\csromanspacer{} รูปกฺขนฺโธ จิตฺตวิสํสฏฺโฐ.\csromanspacer{} วิญฺญาณกฺขนฺโธ น วตฺตพฺโพ “จิตฺเตน สํสฏฺโฐ”ติปิ, “จิตฺเตน วิสํสฏฺโฐ”ติปิ.\csromanspacer{} ตโย ขนฺธา จิตฺตสมุฏฺฐานา.\csromanspacer{} วิญฺญาณกฺขนฺโธ โน จิตฺตสมุฏฺฐาโน.\csromanspacer{} รูปกฺขนฺโธ สิยา จิตฺตสมุฏฺฐาโน, สิยา โน จิตฺตสมุฏฺฐาโน.\csromanspacer{} ตโย ขนฺธา จิตฺตสหภุโน.\csromanspacer{} วิญฺญาณกฺขนฺโธ โน จิตฺตสหภู.\csromanspacer{} รูปกฺขนฺโธ สิยา จิตฺตสหภู, สิยา โน จิตฺตสหภู.\csromanspacer{} ตโย ขนฺธา จิตฺตานุปริวตฺติโน.\csromanspacer{} วิญฺญาณกฺขนฺโธ โน จิตฺตานุปริวตฺติ. \csromanfolio{66}รูปกฺขนฺโธ สิยา จิตฺตานุปริวตฺติ, สิยา โน จิตฺตานุปริวตฺติ.\csromanspacer{} ตโย ขนฺธา จิตฺตสํสฏฺฐสมุฏฺฐานา.\csromanspacer{} เทฺว ขนฺธา โน จิตฺตสํสฏฺฐสมุฏฺฐานา.\csromanspacer{} ตโย ขนฺธา จิตฺตสํสฏฺฐสมุฏฺฐานสหภุโน.\csromanspacer{} เทฺว ขนฺธา โน จิตฺตสํสฏฺฐสมุฏฺฐานสหภุโน.\csromanspacer{} ตโย ขนฺธา จิตฺตสํสฏฺฐสมุฏฺฐานานุปริวตฺติโน.\csromanspacer{} เทฺว ขนฺธา โน จิตฺตสํสฏฺฐสมุฏฺฐานานุปริวตฺติโน. (10)}

\prose{วิญฺญาณกฺขนฺโธ อชฺฌตฺติโก.\csromanspacer{} ตโย ขนฺธา พาหิรา.\csromanspacer{} รูปกฺขนฺโธ สิยา อชฺฌตฺติโก, สิยา พาหิโร.\csromanspacer{} จตฺตาโร ขนฺธา โน อุปาทา.\csromanspacer{} รูปกฺขนฺโธ สิยา อุปาทา, สิยา โน อุปาทา, สิยา อุปาทินฺนา, สิยา อนุปาทินฺนา.\csromanspacer{} จตฺตาโร ขนฺธา โน อุปาทานา.\csromanspacer{} สงฺขารกฺขนฺโธ สิยา อุปาทานํ, สิยา โน อุปาทานํ.\csromanspacer{} รูปกฺขนฺโธ อุปาทานิโย.\csromanspacer{} จตฺตาโร ขนฺธา สิยา อุปาทานิยา, สิยา อนุปาทานิยา.\csromanspacer{} รูปกฺขนฺโธ อุปาทานวิปฺปยุตฺโต.\csromanspacer{} จตฺตาโร ขนฺธา สิยา อุปาทานสมฺปยุตฺตา, สิยา อุปาทานวิปฺปยุตฺตา.\csromanspacer{} รูปกฺขนฺโธ น วตฺตพฺโพ “อุปาทานญฺเจว อุปาทานิโย จา”ติ, อุปาทานิโย เจว โน จ อุปาทานํ.\csromanspacer{} ตโย ขนฺธา น วตฺตพฺพา “อุปาทานญฺเจว อุปาทานิยา จา”ติ, สิยา อุปาทานิยา เจว โน จ อุปาทานา, สิยา น วตฺตพฺพา “อุปาทานิยา เจว โน จ อุปาทานา”ติ.\csromanspacer{} สงฺขารกฺขนฺโธ สิยา อุปาทานญฺเจว อุปาทานิโย จ, สิยา อุปาทานิโย เจว โน จ อุปาทานํ, สิยา น วตฺตพฺโพ “อุปาทานญฺเจว อุปาทานิโย จา”ติปิ, “อุปาทานิโย เจว โน จ อุปาทานนฺ”ติปิ.\csromanspacer{} รูปกฺขนฺโธ น วตฺตพฺโพ “อุปาทานญฺเจว อุปาทานสมฺปยุตฺโต จา”ติปิ, “อุปาทานสมฺปยุตฺโต เจว โน จ อุปาทานนฺ”ติปิ.\csromanspacer{} ตโย ขนฺธา น วตฺตพฺพา “อุปาทานา เจว อุปาทานสมฺปยุตฺตา จา”ติ, สิยา อุปาทานสมฺปยุตฺตา เจว โน จ อุปาทานา, สิยา น วตฺตพฺพา “อุปาทานสมฺปยุตฺตา เจว โน จ อุปาทานา”ติ.\csromanspacer{} สงฺขารกฺขนฺโธ สิยา อุปาทานญฺเจว อุปาทานสมฺปยุตฺโต จ, สิยา อุปาทานสมฺปยุตฺโต เจว โน จ อุปาทานํ, สิยา น วตฺตพฺโพ “อุปาทานญฺเจว อุปาทานสมฺปยุตฺโต จา”ติปิ, “อุปาทานสมฺปยุตฺโต เจว โน จ อุปาทานนฺ”ติปิ.\csromanspacer{} รูปกฺขนฺโธ อุปาทานวิปฺปยุตฺต-อุปาทานิโย.\csromanspacer{} จตฺตาโร ขนฺธา สิยา อุปาทานวิปฺปยุตฺต- อุปาทานิยา, สิยา อุปาทานวิปฺปยุตฺต-อนุปาทานิยา, สิยา น วตฺตพฺพา “อุปาทานวิปฺปยุตฺต-อุปาทานิยา”ติปิ, “อุปาทานวิปฺปยุตฺต-อนุปาทานิยา”ติปิ. (11)}

\prose{จตฺตาโร ขนฺธา โน กิเลสา.\csromanspacer{} สงฺขารกฺขนฺโธ สิยา กิเลโส, สิยา โน กิเลโส.\csromanspacer{} รูปกฺขนฺโธ สํกิเลสิโก.\csromanspacer{} จตฺตาโร ขนฺธา สิยา \csromanfolio{67}สํกิเลสิกา, สิยา อสํกิเลสิกา, รูปกฺขนฺโธ อสํกิลิฏฺโฐ.\csromanspacer{} จตฺตาโร ขนฺธา สิยา สํกิลิฏฺฐา, สิยา อสํกิลิฏฺฐา.\csromanspacer{} รูปกฺขนฺโธ กิเลสวิปฺปยุตฺโต.\csromanspacer{} จตฺตาโร ขนฺธา สิยา กิเลสสมฺปยุตฺตา, สิยา กิเลสวิปฺปยุตฺตา.\csromanspacer{} รูปกฺขนฺโธ น วตฺตพฺโพ “กิเลโส เจว สํกิเลสิโก จา”ติ, สํกิเลสิโก เจว โน จ กิเลโส.\csromanspacer{} ตโย ขนฺธา น วตฺตพฺพา “กิเลสา เจว สํกิเลสิกา จา”ติ, สิยา สํกิเลสิกา เจว โน จ กิเลสา, สิยา น วตฺตพฺพา “สํกิเลสิกา เจว โน จ กิเลสา”ติ.\csromanspacer{} สงฺขารกฺขนฺโธ สิยา กิเลโส เจว สํกิเลสิโก จ, สิยา สํกิเลสิโก เจว โน จ กิเลโส, สิยา น วตฺตพฺโพ “กิเลโส เจว สํกิเลสิโก จา”ติปิ, “สํกิเลสิโก เจว โน จ กิเลโส”ติปิ.\csromanspacer{} รูปกฺขนฺโธ น วตฺตพฺโพ “กิเลโส เจว สํกิลิฏฺโฐ จา”ติปิ, “สํกิลิฏฺโฐ เจว โน จ กิเลโส”ติปิ.\csromanspacer{} ตโย ขนฺธา น วตฺตพฺพา “กิเลโส เจว สํกิลิฏฺฐา จา”ติ, สิยา สํกิลิฏฺฐา เจว โน จ กิเลสา, สิยา น วตฺตพฺพา “สํกิลิฏฺฐา เจว โน จ กิเลสา”ติ.\csromanspacer{} สงฺขารกฺขนฺโธ สิยา กิเลโส เจว สํกิลิฏฺโฐ จ, สิยา สํกิลิฏฺโฐ เจว โน จ กิเลโส, สิยา น วตฺตพฺโพ “กิเลโส เจว สํกิลิฏฺโฐ จา”ติปิ, “สํกิลิฏฺโฐ เจว โน จ กิเลโส”ติปิ.}

\prose{รูปกฺขนฺโธ น วตฺตพฺโพ “กิเลโส เจว กิเลสสมฺปยุตฺโต จา”ติปิ, “กิเลสสมฺปยุตฺโต เจว โน จ กิเลโส”ติปิ.\csromanspacer{} ตโย ขนฺธา น วตฺตพฺพา “กิเลสา เจว กิเลสสมฺปยุตฺตา จา”ติ.\csromanspacer{} สิยา กิเลสสมฺปยุตฺตา เจว โน จ กิเลสา, สิยา น วตฺตพฺพา “กิเลสสมฺปยุตฺตา เจว โน จ กิเลสา”ติ, สงฺขารกฺขนฺโธ สิยา กิเลโส เจว กิเลสสมฺปยุตฺโต จ, สิยา กิเลสสมฺปยุตฺโต เจว โน จ กิเลโส, สิยา น วตฺตพฺโพ “กิเลโส เจว กิเลสสมฺปยุตฺโต จา”ติปิ, กิเลสสมฺปยุตฺโต เจว โน จ กิเลโส”ติปิ.\csromanspacer{} รูปกฺขนฺโธ กิเลสวิปฺปยุตฺตสํกิเลสิโก.\csromanspacer{} จตฺตาโร ขนฺธา สิยา กิเลสวิปฺปยุตฺตสํกิเลสิกา, สิยา กิเลสวิปฺปยุตฺต- อสํกิเลสิกา, สิยา น วตฺตพฺพา “กิเลสวิปฺปยุตฺตสํกิเลสิกา”ติปิ, “กิเลสวิปฺปยุตฺต-อสํกิเลสิกา”ติปิ. (12)}

\prose{รูปกฺขนฺโธ น ทสฺสเนน ปหาตพฺโพ.\csromanspacer{} จตฺตาโร ขนฺธา สิยา ทสฺสเนน ปหาตพฺพา, สิยา น ทสฺสเนน ปหาตพฺพา.\csromanspacer{} รูปกฺขนฺโธ น ภาวนาย \csromanfolio{68}ปหาตพฺโพ.\csromanspacer{} จตฺตาโร ขนฺธา สิยา ภาวนาย ปหาตพฺพา, สิยา น ภาวนาย ปหาตพฺพา.\csromanspacer{} รูปกฺขนฺโธ น ทสฺสเนน ปหาตพฺพเหตุโก.\csromanspacer{} จตฺตาโร ขนฺธา สิยา ทสฺสเนน ปหาตพฺพเหตุกา, สิยา น ทสฺสเนน ปหาตพฺพเหตุกา.\csromanspacer{} รูปกฺขนฺโธ น ภาวนาย ปหาตพฺพเหตุโก.\csromanspacer{} จตฺตาโร ขนฺธา สิยา ภาวนาย ปหาตพฺพเหตุกา, สิยา น ภาวนาย ปหาตพฺพเหตุกา.\csromanspacer{} รูปกฺขนฺโธ อวิตกฺโก.\csromanspacer{} จตฺตาโร ขนฺธา สิยา สวิตกฺกา, สิยา อวิตกฺกา.\csromanspacer{} รูปกฺขนฺโธ อวิจาโร.\csromanspacer{} จตฺตาโร ขนฺธา สิยา สวิจารา, สิยา อวิจารา.\csromanspacer{} รูปกฺขนฺโธ อปฺปีติโก.\csromanspacer{} จตฺตาโร ขนฺธา สิยา สปฺปีติกา, สิยา อปฺปีติกา.\csromanspacer{} รูปกฺขนฺโธ น ปีติสหคโต.\csromanspacer{} จตฺตาโร ขนฺธา สิยา ปีติสหคตา, สิยา น ปีติสหคตา.\csromanspacer{} เทฺว ขนฺธา น สุขสหคตา.\csromanspacer{} ตโย ขนฺธา สิยา สุขสหคตา, สิยา น สุขสหคตา.\csromanspacer{} เทฺว ขนฺธา น อุเปกฺขาสหคตา.\csromanspacer{} ตโย ขนฺธา สิยา อุเปกฺขาสหคตา, สิยา น อุเปกฺขาสหคตา.}

\prose{รูปกฺขนฺโธ กามาวจโร.\csromanspacer{} จตฺตาโร ขนฺธา สิยา กามาวจรา, สิยา น กามาวจรา.\csromanspacer{} รูปกฺขนฺโธ น รูปาวจโร.\csromanspacer{} จตฺตาโร ขนฺธา สิยา รูปาวจรา, สิยา น รูปาวจรา.\csromanspacer{} รูปกฺขนฺโธ น อรูปาวจโร.\csromanspacer{} จตฺตาโร ขนฺธา สิยา อรูปาวจรา, สิยา น อรูปาวจรา.\csromanspacer{} รูปกฺขนฺโธ ปริยาปนฺโน.\csromanspacer{} จตฺตาโร ขนฺธา สิยา ปริยาปนฺนา, สิยา อปริยาปนฺนา.\csromanspacer{} รูปกฺขนฺโธ อนิยฺยานิโก.\csromanspacer{} จตฺตาโร ขนฺธา สิยา นิยฺยานิกา, สิยา อนิยฺยานิกา.\csromanspacer{} รูปกฺขนฺโธ อนิยโต.\csromanspacer{} จตฺตาโร ขนฺธา สิยา นิยตา, สิยา อนิยตา.\csromanspacer{} รูปกฺขนฺโธ ส-อุตฺตโร.\csromanspacer{} จตฺตาโร ขนฺธา สิยา ส-อุตฺตรา, สิยา อนุตฺตรา.\csromanspacer{} รูปกฺขนฺโธ อรโณ.\csromanspacer{} จตฺตาโร ขนฺธา สิยา สรณา, สิยา อรณาติ. (13)}

\vspace{\tipitakaparskip}
\csromancenter{ปญฺหาปุจฺฉกํ.}
\nitthitam{\textbf{ขนฺธวิภงฺโค นิฏฺฐิโต.}}
\csromanmatikaanchor{mtk.17}
\csromantocmarknopage{chapter}{2. อายตนวิภงฺค}{mtk.17}
\bookmark[dest={mtk.17},level=0]{2. อายตนวิภงฺค}
\chapterheadpageread{2. อายตนวิภงฺค}
\csromanmatikaanchor{mtk.18}
\csromantocmarkref{section}{1. สุตฺตนฺตภาชนีย}{mtk.18}
\bookmark[dest={mtk.18},level=1]{1. สุตฺตนฺตภาชนีย}
\csromanheader{1}{1. สุตฺตนฺตภาชนีย}
\proseitem{154}{\csromanfolio{69}ทฺวาทสายตนานิ จกฺขายตนํ รูปายตนํ โสตายตนํ สทฺทายตนํ ฆานายตนํ คนฺธายตนํ ชิวฺหายตนํ รสายตนํ กายายตนํ โผฏฺฐพฺพายตนํ มนายตนํ ธมฺมายตนํ.}

\prose{จกฺขุํ อนิจฺจํ ทุกฺขํ อนตฺตา วิปริณามธมฺมํ, รูปา อนิจฺจา ทุกฺขา อนตฺตา วิปริณามธมฺมา.\csromanspacer{} โสตํ อนิจฺจํ ทุกฺขํ อนตฺตา วิปริณามธมฺมํ, สทฺทา อนิจฺจา ทุกฺขา อนตฺตา วิปริณามธมฺมา.\csromanspacer{} ฆานํ อนิจฺจํ ทุกฺขํ อนตฺตา วิปริณามธมฺมํ, คนฺธา อนิจฺจา ทุกฺขา อนตฺตา วิปริณามธมฺมา.\csromanspacer{} ชิวฺหา อนิจฺจา ทุกฺขา อนตฺตา วิปริณามธมฺมา, รสา อนิจฺจา ทุกฺขา อนตฺตา วิปริณามธมฺมา.\csromanspacer{} กาโย อนิจฺโจ ทุกฺโข อนตฺตา วิปริณามธมฺโม, โผฏฺฐพฺพา อนิจฺจา ทุกฺขา อนตฺตา วิปริณามธมฺมา.\csromanspacer{} มโน อนิจฺโจ ทุกฺโข อนตฺตา วิปริณามธมฺโม, ธมฺมา อนิจฺจา ทุกฺขา อนตฺตา วิปริณามธมฺมา.}

\csromanheader{2}{สุตฺตนฺตภาชนียํ.}
\csromansectionrule
\csromanmatikaanchor{mtk.19}
\csromantocmarkref{section}{2. อภิธมฺมภาชนีย}{mtk.19}
\bookmark[dest={mtk.19},level=1]{2. อภิธมฺมภาชนีย}
\csromanheader{1}{2. อภิธมฺมภาชนีย}
\proseitem{155}{ทฺวาทสายตนานิ จกฺขายตนํ โสตายตนํ ฆานายตนํ ชิวฺหายตนํ กายายตนํ มนายตนํ รูปายตนํ สทฺทายตนํ คนฺธายตนํ รสายตนํ โผฏฺฐพฺพายตนํ ธมฺมายตนํ.}

\proseitem{156}{ตตฺถ กตมํ จกฺขายตนํ, ยํ จกฺขุ\footnote{จกฺขุํ (สี, สฺยา, ก) อภิ 1. 157 ปิฏฺเฐปิ.} จตุนฺนํ มหาภูตานํ อุปาทาย ปสาโท อตฺตภาวปริยาปนฺโน อนิทสฺสโน สปฺปฏิโฆ, เยน จกฺขุนา อนิทสฺสเนน สปฺปฏิเฆน รูปํ สนิทสฺสนํ สปฺปฏิฆํ ปสฺสิ วา ปสฺสติ วา ปสฺสิสฺสติ วา ปสฺเส วา, จกฺขุมฺเปตํ, จกฺขายตนมฺเปตํ, จกฺขุธาตุเปสา, จกฺขุนฺทฺริยมฺเปตํ, โลโกเปโส, ทฺวาราเปสา, สมุทฺโทเปโส, ปณฺฑรมฺเปตํ, เขตฺตมฺเปตํ, วตฺถุมฺเปตํ, เนตฺตมฺเปตํ, นยนมฺเปตํ, โอริมํ ตีรมฺเปตํ, สุญฺโญ คาโมเปโส.\csromanspacer{} อิทํ วุจฺจติ จกฺขายตนํ. (1)}

\proseitem{157}{\csromanfolio{70}ตตฺถ กตมํ โสตายตนํ, ยํ โสตํ จตุนฺนํ มหาภูตานํ อุปาทาย ปสาโท อตฺตภาวปริยาปนฺโน อนิทสฺสโน สปฺปฏิโฆ, เยน โสเตน อนิทสฺสเนน สปฺปฏิเฆน สทฺทํ อนิทสฺสนํ สปฺปฏิฆํ สุณิ วา สุณาติ วา สุณิสฺสติ วา สุเณ วา, โสตมฺเปตํ, โสตายตนมฺเปตํ, โสตธาตุเปสา, โสตินฺทฺริยมฺเปตํ, โลโกเปโส, ทฺวาราเปสา, สมุทฺโทเปโส, ปณฺฑรมฺเปตํ, เขตฺตมฺเปตํ, วตฺถุมฺเปตํ, โอริมํ ตีรมฺเปตํ, สุญฺโญ คาโมเปโส.\csromanspacer{} อิทํ วุจฺจติ โสตายตนํ. (2)}

\proseitem{158}{ตตฺถ กตมํ ฆานายตนํ, ยํ ฆานํ จตุนฺนํ มหาภูตานํ อุปาทาย ปสาโท อตฺตภาวปริยาปนฺโน อนิทสฺสโน สปฺปฏิโฆ, เยน ฆาเนน อนิทสฺสเนน สปฺปฏิเฆน คนฺธํ อนิทสฺสนํ สปฺปฏิฆํ ฆายิ วา ฆายติ วา ฆายิสฺสติ วา ฆาเย วา, ฆานมฺเปตํ, ฆานายตนมฺเปตํ, ฆานธาตุเปสา, ฆานินฺทฺริยมฺเปตํ, โลโกเปโส, ทฺวาราเปสา, สมุทฺโทเปโส, ปณฺฑรมฺเปตํ, เขตฺตมฺเปตํ, วตฺถุมฺเปตํ, โอริมํ ตีรมฺเปตํ, สุญฺโญ คาโมเปโส.\csromanspacer{} อิทํ วุจฺจติ ฆานายตนํ. (3)}

\proseitem{159}{ตตฺถ กตมํ ชิวฺหายตนํ, ยา ชิวฺหา จตุนฺนํ มหาภูตานํ อุปาทาย ปสาโท อตฺตภาวปริยาปนฺโน อนิทสฺสโน สปฺปฏิโฆ, ยาย ชิวฺหาย อนิทสฺสนาย สปฺปฏิฆาย รสํ อนิทสฺสนํ สปฺปฏิฆํ สายิ วา สายติ วา สายิสฺสติ วา สาเย วา, ชิวฺหาเปสา, ชิวฺหายตนมฺเปตํ, ชิวฺหาธาตุเปสา, ชิวฺหินฺทฺริยมฺเปตํ, โลโกเปโส, ทฺวาราเปสา, สมุทฺโทเปโส, ปณฺฑรมฺเปตํ, เขตฺตมฺเปตํ, วตฺถุมฺเปตํ, โอริมํ ตีรมฺเปตํ, สุญฺโญ คาโมเปโส.\csromanspacer{} อิทํ วุจฺจติ ชิวฺหายตนํ. (4)}

\proseitem{160}{ตตฺถ กตมํ กายายตนํ, โย กาโย จตุนฺนํ มหาภูตานํ อุปาทาย ปสาโท อตฺตภาวปริยาปนฺโน อนิทสฺสโน สปฺปฏิโฆ, เยน กาเยน อนิทสฺสเนน สปฺปฏิเฆน โผฏฺฐพฺพํ อนิทสฺสนํ สปฺปฏิฆํ ผุสิ วา ผุสติ วา ผุสิสฺสติ วา ผุเส วา, กาโยเปโส กายายตนมฺเปตํ, กายธาตุเปสา, กายินฺทฺริยมฺเปตํ, โลโกเปโส, ทฺวาราเปสา, สมุทฺโทเปโส, ปณฺฑรมฺเปตํ, เขตฺตมฺเปตํ, วตฺถุมฺเปตํ, โอริมํ ตีรมฺเปตํ, สุญฺโญ คาโมเปโส.\csromanspacer{} อิทํ วุจฺจติ กายายตนํ. (5)}

\csromanheader{2}{2. อายตนวิภงฺค}
\proseitem{161}{\csromanfolio{71}ตตฺถ กตมํ มนายตนํ, เอกวิเธน มนายตนํ\csromandash{} ผสฺสสมฺปยุตฺตํ.}

\prose{ทุวิเธน มนายตนํ\csromandash{}อตฺถิ สเหตุกํ, อตฺถิ อเหตุกํ.}

\prose{ติวิเธน มนายตนํ\csromandash{}อตฺถิ กุสลํ, อตฺถิ อกุสลํ, อตฺถิ อพฺยากตํ.}

\prose{จตุพฺพิเธน มนายตนํ\csromandash{}อตฺถิ กามาวจรํ, อตฺถิ รูปาวจรํ, อตฺถิ อรูปาวจรํ, อตฺถิ อปริยาปนฺนํ.}

\prose{ปญฺจวิเธน มนายตนํ\csromandash{}อตฺถิ สุขินฺทฺริยสมฺปยุตฺตํ, อตฺถิ ทุกฺขินฺทฺริยสมฺปยุตฺตํ, อตฺถิ โสมนสฺสินฺทฺริยสมฺปยุตฺตํ, อตฺถิ โทมนสฺสินฺทฺริยสมฺปยุตฺตํ, อตฺถิ อุเปกฺขินฺทฺริยสมฺปยุตฺตํ.}

\prose{ฉพฺพิเธน มนายตนํ\csromandash{}จกฺขุวิญฺญาณํ, โสตวิญฺญาณํ, ฆานวิญฺญาณํ, ชิวฺหาวิญฺญาณํ, กายวิญฺญาณํ, มโนวิญฺญาณํ.\csromanspacer{} เอวํ ฉพฺพิเธน มนายตนํ.}

\prose{สตฺตวิเธน มนายตนํ\csromandash{}จกฺขุวิญฺญาณํ, โสตวิญฺญาณํ, ฆานวิญฺญาณํ, ชิวฺหาวิญฺญาณํ, กายวิญฺญาณํ, มโนธาตุ, มโนวิญฺญาณธาตุ.\csromanspacer{} เอวํ สตฺตวิเธน มนายตนํ.}

\prose{อฏฺฐวิเธน มนายตนํ\csromandash{}จกฺขุวิญฺญาณํ, โสตวิญฺญาณํ, ฆานวิญฺญาณํ, ชิวฺหาวิญฺญาณํ, กายวิญฺญาณํ อตฺถิ สุขสหคตํ, อตฺถิ ทุกฺขสหคตํ, มโนธาตุ, มโนวิญฺญาณธาตุ.\csromanspacer{} เอวํ อฏฺฐวิเธน มนายตนํ.}

\prose{นววิเธน มนายตนํ\csromandash{}จกฺขุวิญฺญาณํ, โสตวิญฺญาณํ, ฆานวิญฺญาณํ, ชิวฺหาวิญฺญาณํ, กายวิญฺญาณํ, มโนธาตุ, มโนวิญฺญาณธาตุ อตฺถิ กุสลํ, อตฺถิ อกุสลํ, อตฺถิ อพฺยากตํ.\csromanspacer{} เอวํ นววิเธน มนายตนํ.}

\prose{ทสวิเธน มนายตนํ\csromandash{}จกฺขุวิญฺญาณํ, โสตวิญฺญาณํ, ฆานวิญฺญาณํ, ชิวฺหาวิญฺญาณํ, กายวิญฺญาณํ อตฺถิ สุขสหคตํ, อตฺถิ ทุกฺขสหคตํ, มโนธาตุ, มโนวิญฺญาณธาตุ อตฺถิ กุสลํ, อตฺถิ อกุสลํ, อตฺถิ อพฺยากตํ, เอวํ ทสวิเธน มนายตนํ.}

\prose{เอกวิเธน มนายตนํ\csromandash{}ผสฺสสมฺปยุตฺตํ.}

\prose{ทุวิเธน มนายตนํ\csromandash{}อตฺถิ สเหตุกํ, อตฺถิ อเหตุกํ.}

\prose{\csromanfolio{72}ติวิเธน มนายตนํ\csromandash{}อตฺถิ สุขาย เวทนาย สมฺปยุตฺตํ, อตฺถิ ทุกฺขาย เวทนาย สมฺปยุตฺตํ, อตฺถิ อทุกฺขมสุขาย เวทนาย สมฺปยุตฺตํ ฯเปฯ.\csromanspacer{} เอวํ พหุวิเธน มนายตนํ.\csromanspacer{} อิทํ วุจฺจติ มนายตนํ. (6)}

\proseitem{162}{ตตฺถ กตมํ รูปายตนํ, ยํ รูปํ จตุนฺนํ มหาภูตานํ อุปาทาย วณฺณนิภา สนิทสฺสนํ สปฺปฏิฆํ นีลํ ปีตกํ โลหิตกํ\footnote{ปีตํ โลหิตํ (สี)} โอทาตํ กาฬกํ มญฺชิฏฺฐกํ\footnote{มญฺเชฏฺฐกํ (สี, สฺยา)} หริ หริวณฺณํ อมฺพงฺกุรวณฺณํ ทีฆํ รสฺสํ อณุํ ถูลํ วฏฺฏํ ปริมณฺฑลํ จตุรสฺสํ ฉฬํสํ อฏฺฐํสํ โสฬสํสํ นินฺนํ ถลํ ฉายา อาตโป อาโลโก อนฺธกาโร อพฺภา มหิกา ธูโม รโช จนฺทมณฺฑลสฺส วณฺณนิภา สูริยมณฺฑลสฺส\footnote{สุริยมณฺฑลสฺส (สี, สฺยา, กํ)} วณฺณนิภา ตารกรูปานํ วณฺณนิภา อาทาสมณฺฑลสฺส วณฺณนิภา มณิสงฺขมุตฺตเวฬุริยสฺส วณฺณนิภา ชาตรูปรชตสฺส วณฺณนิภา, ยํ วา ปนญฺญมฺปิ อตฺถิ รูปํ จตุนฺนํ มหาภูตานํ อุปาทาย วณฺณนิภา สนิทสฺสนํ สปฺปฏิฆํ, ยํ รูปํ สนิทสฺสนํ สปฺปฏิฆํ จกฺขุนา อนิทสฺสเนน สปฺปฏิเฆน ปสฺสิ วา ปสฺสติ วา ปสฺสิสฺสติ วา ปสฺเส วา, รูปมฺเปตํ, รูปายตนมฺเปตํ, รูปธาตุเปสา.\csromanspacer{} อิทํ วุจฺจติ รูปายตนํ. (7)}

\proseitem{163}{ตตฺถ กตมํ สทฺทายตนํ, โย สทฺโท จตุนฺนํ มหาภูตานํ อุปาทาย อนิทสฺสโน สปฺปฏิโฆ เภริสทฺโท มุทิงฺคสทฺโท\footnote{มุติงฺคสทฺโท (สี)} สงฺขสทฺโท ปณวสทฺโท คีตสทฺโท วาทิตสทฺโท สมฺมสทฺโท ปาณิสทฺโท สตฺตานํ นิคฺโฆสสทฺโท ธาตูนํ สนฺนิฆาตสทฺโท วาตสทฺโท อุทกสทฺโท มนุสฺสสทฺโท อมนุสฺสสทฺโท, โย วา ปนญฺโญปิ อตฺถิ สทฺโท จตุนฺนํ มหาภูตานํ อุปาทาย อนิทสฺสโน สปฺปฏิโฆ, ยํ สทฺทํ อนิทสฺสนํ สปฺปฏิฆํ โสเตน อนิทสฺสเนน สปฺปฏิเฆน สุณิ วา สุณาติ วา สุณิสฺสติ วา สุเณ วา, สทฺโทเปโส, สทฺทายตนมฺเปตํ, สทฺทธาตุเปสา.\csromanspacer{} อิทํ วุจฺจติ สทฺทายตนํ. (8)}

\proseitem{164}{ตตฺถ กตมํ คนฺธายตนํ, โย คนฺโธ จตุนฺนํ มหาภูตานํ อุปาทาย อนิทสฺสโน สปฺปฏิโฆ มูลคนฺโธ สารคนฺโธ ตจคนฺโธ ปตฺตคนฺโธ ปุปฺผคนฺโธ ผลคนฺโธ อามคนฺโธ วิสฺสคนฺโธ สุคนฺโธ ทุคฺคนฺโธ, โย วา ปนญฺโญปิ อตฺถิ คนฺโธ จตุนฺนํ มหาภูตานํ อุปาทาย อนิทสฺสโน}

\vspace{\tipitakaparskip}
\csromancenter{อายตนวิภงฺค สปฺปฏิโฆ, ยํ คนฺธํ อนิทสฺสนํ สปฺปฏิฆํ ฆาเนน อนิทสฺสเนน สปฺปฏิเฆน ฆายิ วา ฆายติ วา ฆายิสฺสติ วา ฆาเย วา, คนฺโธเปโส, คนฺธายตนมฺเปตํ, คนฺธธาตุเปสา.\csromanspacer{} อิทํ วุจฺจติ คนฺธายตนํ. (9)}
\proseitem{165}{\csromanfolio{73}ตตฺถ กตมํ รสายตนํ, โย รโส จตุนฺนํ มหาภูตานํ อุปาทาย อนิทสฺสโน สปฺปฏิโฆ มูลรโส ขนฺธรโส ตจรโส ปตฺตรโส ปุปฺผรโส ผลรโส อมฺพิลํ มธุรํ ติตฺตกํ กฏุกํ โลณิกํ ขาริกํ ลมฺพิกํ\footnote{ลปิลกํ (สี), ลมฺปิกํ (ก-สี)} กสาโว สาทุ อสาทุ, โย วา ปนญฺโญปิ อตฺถิ รโส จตุนฺนํ มหาภูตานํ อุปาทาย อนิทสฺสโน สปฺปฏิโฆ, ยํ รสํ อนิทสฺสนํ สปฺปฏิฆํ ชิวฺหาย อนิทสฺสนาย สปฺปฏิฆาย สายิ วา สายติ วา สายิสฺสติ วา สาเย วา, รโสเปโส, รสายตนมฺเปตํ, รสธาตุเปสา.\csromanspacer{} อิทํ วุจฺจติ รสายตนํ. (10)}

\proseitem{166}{ตตฺถ กตมํ โผฏฺฐพฺพายตนํ, ปถวีธาตุ เตโชธาตุ วาโยธาตุ กกฺขฬํ มุทุกํ สณฺหํ ผรุสํ สุขสมฺผสฺสํ ทุกฺขสมฺผสฺสํ ครุกํ ลหุกํ, ยํ โผฏฺฐพฺพํ อนิทสฺสนํ สปฺปฏิฆํ กาเยน อนิทสฺสเนน สปฺปฏิเฆน ผุสิ วา ผุสติ วา ผุสิสฺสติ วา ผุเส วา, โผฏฺฐพฺโพเปโส, โผฏฺฐพฺพายตนมฺเปตํ, โผฏฺฐพฺพธาตุเปสา.\csromanspacer{} อิทํ วุจฺจติ โผฏฺฐพฺพายตนํ. (11)}

\proseitem{167}{ตตฺถ กตมํ ธมฺมายตนํ, เวทนากฺขนฺโธ สญฺญากฺขนฺโธ สงฺขารกฺขนฺโธ, ยญฺจ รูปํ อนิทสฺสน-อปฺปฏิฆํ ธมฺมายตนปริยาปนฺนํ อสงฺขตา จ ธาตุ.}

\prose{ตตฺถ กตโม เวทนากฺขนฺโธ, เอกวิเธน เวทนากฺขนฺโธ\csromandash{} ผสฺสสมฺปยุตฺโต.\csromanspacer{} ทุวิเธน เวทนากฺขนฺโธ\csromandash{}อตฺถิ สเหตุโก, อตฺถิ อเหตุโก.\csromanspacer{} ติวิเธน เวทนากฺขนฺโธ\csromandash{}อตฺถิ กุสโล, อตฺถิ อกุสโล, อตฺถิ อพฺยากโต ฯเปฯ.\csromanspacer{} เอวํ ทสวิเธน เวทนากฺขนฺโธ.\csromanspacer{} เอวํ พหุวิเธน เวทนากฺขนฺโธ.\csromanspacer{} อยํ วุจฺจติ เวทนากฺขนฺโธ. (1)}

\prose{ตตฺถ กตโม สญฺญากฺขนฺโธ, เอกวิเธน สญฺญากฺขนฺโธ\csromandash{} ผสฺสสมฺปยุตฺโต.\csromanspacer{} ทุวิเธน สญฺญากฺขนฺโธ\csromandash{}อตฺถิ สเหตุโก, อตฺถิ อเหตุโก.\csromanspacer{} ติวิเธน สญฺญากฺขนฺโธ\csromandash{}อตฺถิ กุสโล, อตฺถิ อกุสโล, \csromanfolio{74}อตฺถิ อพฺยากโต ฯเปฯ.\csromanspacer{} เอวํ ทสวิเธน สญฺญากฺขนฺโธ ฯเปฯ.\csromanspacer{} เอวํ พหุวิเธน สญฺญากฺขนฺโธ.\csromanspacer{} อยํ วุจฺจติ สญฺญากฺขนฺโธ. (2)}

\prose{ตตฺถ กตโม สงฺขารกฺขนฺโธ, เอกวิเธน สงฺขารกฺขนฺโธ\csromandash{} จิตฺตสมฺปยุตฺโต.\csromanspacer{} ทุวิเธน สงฺขารกฺขนฺโธ\csromandash{}อตฺถิ เหตุ, อตฺถิ น เหตุ.\csromanspacer{} ติวิเธน สงฺขารกฺขนฺโธ\csromandash{}อตฺถิ กุสโล, อตฺถิ อกุสโล, อตฺถิ อพฺยากโต ฯเปฯ.\csromanspacer{} เอวํ ทสวิเธน สงฺขารกฺขนฺโธ ฯเปฯ.\csromanspacer{} เอวํ พหุวิเธน สงฺขารกฺขนฺโธ.\csromanspacer{} อยํ วุจฺจติ สงฺขารกฺขนฺโธ. (3)}

\prose{ตตฺถ กตมํ รูปํ อนิทสฺสน-อปฺปฏิฆํ ธมฺมายตนปริยาปนฺนํ, อิตฺถินฺทฺริยํ ปุริสินฺทฺริยํ ฯเปฯ กพฬีกาโร อาหาโร.\csromanspacer{} อิทํ วุจฺจติ รูปํ อนิทสฺสน-อปฺปฏิฆํ ธมฺมายตนปริยาปนฺนํ. (4)}

\prose{ตตฺถ กตมา อสงฺขตา ธาตุ, ราคกฺขโย โทสกฺขโย, โมหกฺขโย.\csromanspacer{} อยํ วุจฺจติ อสงฺขตา ธาตุ. (5)}

\vspace{\tipitakaparskip}
\csromancenter{อิทํ วุจฺจติ ธมฺมายตนํ.}
\csromanheader{2}{อภิธมฺมภาชนียํ.}
\csromansectionrule
\csromanmatikaanchor{mtk.20}
\csromantocmarkref{section}{3. ปญฺหาปุจฺฉก}{mtk.20}
\bookmark[dest={mtk.20},level=1]{3. ปญฺหาปุจฺฉก}
\csromanheader{1}{3. ปญฺหาปุจฺฉก}
\proseitem{168}{ทฺวาทสายตนานิ จกฺขายตนํ รูปายตนํ โสตายตนํ สทฺทายตนํ ฆานายตนํ คนฺธายตนํ ชิวฺหายตนํ รสายตนํ กายายตนํ โผฏฺฐพฺพายตนํ มนายตนํ ธมฺมายตนํ.}

\proseitem{169}{ทฺวาทสนฺนํ อายตนานํ กติ กุสลา, กติ อกุสลา, กติ อพฺยากตา ฯเปฯ กติ สรณา, กติ อรณา.}

\csromanmatikaanchor{mtk.21}
\csromantocmarkref{subsection}{1. ติก}{mtk.21}
\bookmark[dest={mtk.21},level=2]{1. ติก}
\csromanheader{2}{1. ติก}
\proseitem{170}{ทสายตนา อพฺยากตา.\csromanspacer{} ทฺวายตนา สิยา กุสลา, สิยา อกุสลา, สิยา อพฺยากตา.\csromanspacer{} ทสายตนา น วตฺตพฺพา “สุขาย เวทนาย สมฺปยุตฺตา”ติปิ, “ทุกฺขาย เวทนาย สมฺปยุตฺตา”ติปิ, “อทุกฺขมสุขาย เวทนาย สมฺปยุตฺตา”ติปิ.\csromanspacer{} มนายตนํ สิยา สุขาย}

\vspace{\tipitakaparskip}
\csromancenter{อายตนวิภงฺค เวทนาย สมฺปยุตฺตํ, สิยา ทุกฺขาย เวทนาย สมฺปยุตฺตํ, สิยา อทุกฺขมสุขาย เวทนาย สมฺปยุตฺตํ.\csromanspacer{} ธมฺมายตนํ สิยา สุขาย เวทนาย สมฺปยุตฺตํ, สิยา ทุกฺขาย เวทนาย สมฺปยุตฺตํ, สิยา อทุกฺขมสุขาย เวทนาย สมฺปยุตฺตํ, สิยา น วตฺตพฺพํ “สุขาย เวทนาย สมฺปยุตฺตนฺ”ติปิ, “ทุกฺขาย เวทนาย สมฺปยุตฺตนฺ”ติปิ, “อทุกฺขมสุขาย เวทนาย สมฺปยุตฺตนฺ”ติปิ.\csromanspacer{} ทสายตนา เนววิปากนวิปากธมฺมธมฺมา, ทฺวายตนา สิยา วิปากา, สิยา วิปากธมฺมธมฺมา, สิยา เนววิปากนวิปากธมฺมธมฺมา.}
\prose{\csromanfolio{75}ปญฺจายตนา อุปาทินฺนุปาทานิยา.\csromanspacer{} สทฺทายตนํ อนุปาทินฺนุปาทานิยํ.\csromanspacer{} จตฺตาโร อายตนา สิยา อุปาทินฺนุปาทานิยา, สิยา อนุปาทินฺนุปาทานิยา, สิยา อนุปาทานิยา.\csromanspacer{} ทฺวายตนา สิยา อุปาทินฺนุปาทานิยา, สิยา อนุปาทินฺนุปาทานิยา, สิยา อนุปาทินฺน-อนุปาทานิยา.\csromanspacer{} ทสายตนา อสํกิลิฏฺฐสํกิเลสิกา, ทฺวายตนา สิยา สํกิลิฏฺฐสํกิเลสิกา, สิยา อสํกิลิฏฺฐสํกิเลสิกา, สิยา อสํกิลิฏฺฐ-อสํกิเลสิกา.\csromanspacer{} ทสายตนา อวิตกฺก-อวิจารา.\csromanspacer{} มนายตนํ สิยา สวิตกฺกสวิจารํ, สิยา อวิตกฺกวิจารมตฺตํ, สิยา อวิตกฺก-อวิจารํ.\csromanspacer{} ธมฺมายตนํ สิยา สวิตกฺกสวิจารํ, สิยา อวิตกฺกวิจารมตฺตํ, สิยา อวิตกฺก-อวิจารํ, สิยา น วตฺตพฺพํ “สวิตกฺกสวิจารนฺ”ติปิ, “อวิตกฺกวิจารมตฺตนฺ”ติปิ, “อวิตกฺก-อวิจารนฺ”ติปิ.\csromanspacer{} ทสายตนา น วตฺตพฺพา “ปีติสหคตา”ติปิ, “สุขสหคตา”ติปิ, “อุเปกฺขาสหคตา”ติปิ.\csromanspacer{} ทฺวายตนา สิยา ปีติสหคตา, สิยา สุขสหคตา, สิยา อุเปกฺขาสหคตา, สิยา น วตฺตพฺพา “ปีติสหคตา”ติปิ, “สุขสหคตา”ติปิ, “อุเปกฺขาสหคตา”ติปิ.}

\prose{ทสายตนา เนว ทสฺสเนน น ภาวนาย ปหาตพฺพา.\csromanspacer{} ทฺวายตนา สิยา ทสฺสเนน ปหาตพฺพา, สิยา ภาวนาย ปหาตพฺพา, สิยา เนว ทสฺสเนน น ภาวนาย ปหาตพฺพา.\csromanspacer{} ทสายตนา เนว ทสฺสเนน น ภาวนาย ปหาตพฺพเหตุกา.\csromanspacer{} ทฺวายตนา สิยา ทสฺสเนน น ปหาตพฺพเหตุกา, สิยา ภาวนาย ปหาตพฺพเหตุกา, สิยา เนว ทสฺสเนน น ภาวนาย ปหาตพฺพเหตุกา.\csromanspacer{} ทสายตนา เนวาจยคามินาปจยคามิโน.\csromanspacer{} ทฺวายตนา สิยา อาจยคามิโน, สิยา อปจยคามิโน, สิยา เนวาจยคามินาปจยคามิโน.\csromanspacer{} ทสายตนา เนวเสกฺขนาเสกฺขา.\csromanspacer{} ทฺวายตนา สิยา เสกฺขา, สิยา \csromanfolio{76}อเสกฺขา, สิยา เนวเสกฺขนาเสกฺขา.\csromanspacer{} ทสายตนา ปริตฺตา.\csromanspacer{} ทฺวายตนา สิยา ปริตฺตา, สิยา มหคฺคตา, สิยา อปฺปมาณา.\csromanspacer{} ทสายตนา อนารมฺมณา.\csromanspacer{} ทฺวายตนา สิยา ปริตฺตารมฺมณา, สิยา มหคฺคตารมฺมณา, สิยา อปฺปมาณารมฺมณา, สิยา น วตฺตพฺพา “ปริตฺตารมฺมณา”ติปิ, “มหคฺคตารมฺมณา”ติปิ, “อปฺปมาณารมฺมณา”ติปิ.\csromanspacer{} ทสายตนา มชฺฌิมา.\csromanspacer{} ทฺวายตนา สิยา หีนา, สิยา มชฺฌิมา, สิยา ปณีตา.\csromanspacer{} ทสายตนา อนิยตา.\csromanspacer{} ทฺวายตนา สิยา มิจฺฉตฺตนิยตา, สิยา สมฺมตฺตนิยตา, สิยา อนิยตา.}

\prose{ทสายตนา อนารมฺมณา.\csromanspacer{} ทฺวายตนา สิยา มคฺคารมฺมณา, สิยา มคฺคเหตุกา, สิยา มคฺคาธิปติโน, สิยา น วตฺตพฺพา “มคฺคารมฺมณา”ติปิ, “มคฺคเหตุกา”ติปิ, “มคฺคาธิปติโน”ติปิ.\csromanspacer{} ปญฺจายตนา สิยา อุปฺปนฺนา, สิยา อุปฺปาทิโน, น วตฺตพฺพา “อนุปฺปนฺนา”ติ.\csromanspacer{} สทฺทายตนํ สิยา อุปฺปนฺนํ, สิยา อนุปฺปนฺนํ, น วตฺตพฺพํ “อุปฺปาที”ติ.\csromanspacer{} ปญฺจายตนา สิยา อุปฺปนฺนา, สิยา อนุปฺปนฺนา, สิยา อุปฺปาทิโน.\csromanspacer{} ธมฺมายตนํ สิยา อุปฺปนฺนํ, สิยา อนุปฺปนฺนํ, สิยา อุปฺปาที, สิยา น วตฺตพฺพํ “อุปฺปนฺนนฺ”ติปิ, “อนุปฺปนฺนนฺ”ติปิ, “อุปฺปาที”ติปิ.\csromanspacer{} เอกาทสายตนา สิยา อตีตา, สิยา อนาคตา, สิยา ปจฺจุปฺปนฺนา.\csromanspacer{} ธมฺมายตนํ สิยา อตีตํ, สิยา อนาคตํ, สิยา ปจฺจุปฺปนฺนํ, สิยา น วตฺตพฺพํ “อตีตนฺ”ติปิ, “อนาคตนฺ”ติปิ, “ปจฺจุปฺปนฺนนฺ”ติปิ.\csromanspacer{} ทสายตนา อนารมฺมณา.\csromanspacer{} ทฺวายตนา สิยา อตีตารมฺมณา, สิยา อนาคตารมฺมณา, สิยา ปจฺจุปฺปนฺนารมฺมณา, สิยา น วตฺตพฺพา “อตีตารมฺมณา”ติปิ, “อนาคตารมฺมณา”ติปิ, “ปจฺจุปฺปนฺนารมฺมณา”ติปิ, สิยา อชฺฌตฺตา, สิยา พหิทฺธา, สิยา อชฺฌตฺตพหิทฺธา.\csromanspacer{} ทสายตนา อนารมฺมณา.\csromanspacer{} ทฺวายตนา สิยา อชฺฌตฺตารมฺมณา, สิยา พหิทฺธารมฺมณา, สิยา อชฺฌตฺตพหิทฺธารมฺมณา, สิยา น วตฺตพฺพา “อชฺฌตฺตารมฺมณา”ติปิ, “พหิทฺธารมฺมณา”ติปิ, “อชฺฌตฺตพหิทฺธารมฺมณา”ติปิ.\csromanspacer{} รูปายตนํ สนิทสฺสนสปฺปฏิฆํ.\csromanspacer{} นวายตนา อนิทสฺสนสปฺปฏิฆา.\csromanspacer{} ทฺวายตนา อนิทสฺสน-อปฺปฏิฆา.}

\csromanmatikaanchor{mtk.22}
\csromantocmarkref{subsection}{2. ทุก}{mtk.22}
\bookmark[dest={mtk.22},level=2]{2. ทุก}
\csromanheader{2}{2. ทุก}
\proseitem{171}{เอกาทสายตนา น เหตู.\csromanspacer{} ธมฺมายตนํ สิยา เหตุ, สิยา น เหตุ.\csromanspacer{} ทสายตนา อเหตุกา.\csromanspacer{} ทฺวายตนา สิยา สเหตุกา, สิยา อเหตุกา.\csromanspacer{} ทสายตนา เหตุวิปฺปยุตฺตา.\csromanspacer{} ทฺวายตนา สิยา เหตุสมฺปยุตฺตา, สิยา เหตุวิปฺปยุตฺตา.\csromanspacer{} ทสายตนา น วตฺตพฺพา}

\vspace{\tipitakaparskip}
\csromancenter{อายตนวิภงฺค “เหตู เจว สเหตุกา จา”ติปิ, “สเหตุกา เจว น จ เหตู”ติปิ.\csromanspacer{} มนายตนํ น วตฺตพฺพํ “เหตุ เจว สเหตุกญฺจา”ติ, สิยา สเหตุกญฺเจว น จ เหตุ, สิยา น วตฺตพฺพํ “สเหตุกญฺเจว น จ เหตู”ติ.\csromanspacer{} ธมฺมายตนํ สิยาเหตุ เจว สเหตุกญฺจ, สิยา สเหตุกญฺเจว น จ เหตุ, สิยา น วตฺตพฺพํ “เหตุ เจว สเหตุกญฺจา”ติปิ, “สเหตุกญฺเจว น จ เหตู”ติปิ.\csromanspacer{} ทสายตนา น วตฺตพฺพา “เหตู เจว เหตุสมฺปยุตฺตา จา”ติปิ, “เหตุสมฺปยุตฺตา เจว น จ เหตู”ติปิ.\csromanspacer{} มนายตนํ น วตฺตพฺพํ “เหตุ เจว เหตุสมฺปยุตฺตญฺจา”ติ, สิยา เหตุสมฺปยุตฺตญฺเจว น จ เหตุ, สิยา น วตฺตพฺพํ “เหตุสมฺปยุตฺตญฺเจว น จ เหตู”ติ.\csromanspacer{} ธมฺมายตนํ สิยา เหตุ เจว เหตุสมฺปยุตฺตญฺจ, สิยา เหตุสมฺปยุตฺตญฺเจว น จ เหตุ, สิยา น วตฺตพฺพํ “เหตุ เจว เหตุสมฺปยุตฺตญฺจา”ติปิ, “เหตุสมฺปยุตฺตญฺเจว น จ เหตู”ติปิ.\csromanspacer{} ทสายตนา น เหตุ-อเหตุกา.\csromanspacer{} มนายตนํ สิยา น เหตุสเหตุกํ, สิยา น เหตุ- อเหตุกํ.\csromanspacer{} ธมฺมายตนํ สิยา น เหตุสเหตุกํ, สิยา น เหตุ-อเหตุกํ, สิยา น วตฺตพฺพํ “น เหตุสเหตุกนฺ”ติปิ, “น เหตุ-อเหตุกนฺ”ติปิ. (1)}
\prose{\csromanfolio{77}เอกาทสายตนา สปฺปจฺจยา.\csromanspacer{} ธมฺมายตนํ สิยา สปฺปจฺจยํ, สิยา อปฺปจฺจยํ.\csromanspacer{} เอกาทสายตนา สงฺขตํ.\csromanspacer{} ธมฺมายตนํ สิยา สงฺขตํ, สิยา อสงฺขตํ.\csromanspacer{} รูปายตนํ สนิทสฺสนํ.\csromanspacer{} เอกาทสายตนา อนิทสฺสนา.\csromanspacer{} ทสายตนา สปฺปฏิฆา.\csromanspacer{} ทฺวายตนา อปฺปฏิฆา.\csromanspacer{} ทสายตนา รูปา.\csromanspacer{} มนายตนํ อรูปํ.\csromanspacer{} ธมฺมายตนํ สิยา รูปํ, สิยา อรูปํ.\csromanspacer{} ทสายตนา โลกิยา.\csromanspacer{} ทฺวายตนา สิยา โลกิยา, สิยา โลกุตฺตรา, เกนจิ วิญฺเญยฺยา, เกนจิ น วิญฺเญยฺยา. (2)}

\prose{เอกาทสายตนา โน อาสวา.\csromanspacer{} ธมฺมายตนํ สิยา อาสโว, สิยา โน อาสโว.\csromanspacer{} ทสายตนา สาสวา.\csromanspacer{} ทฺวายตนา สิยา สาสวา, สิยา อนาสวา.\csromanspacer{} ทสายตนา อาสววิปฺปยุตฺตา.\csromanspacer{} ทฺวายตนา สิยา อาสวสมฺปยุตฺตา, สิยา อาสววิปฺปยุตฺตา.\csromanspacer{} ทสายตนา น วตฺตพฺพา “อาสวา เจว สาสวา จา”ติ, สาสวา เจว โน จ อาสวา.\csromanspacer{} มนายตนํ น วตฺตพฺพํ “อาสโว เจว สาสวญฺจา”ติ, สิยา สาสวญฺเจว โน จ อาสโว, สิยา น วตฺตพฺพํ “สาสวญฺเจว โน จ อาสโว”ติ.\csromanspacer{} ธมฺมายตนํ สิยา อาสโว เจว สาสวญฺจ, สิยา สาสวญฺเจว โน \csromanfolio{78}จ อาสโว, สิยา น วตฺตพฺพํ “อาสโว เจว สาสวญฺจา”ติปิ, “สาสวญฺเจว โน จ อาสโว”ติปิ.\csromanspacer{} ทสายตนา น วตฺตพฺพา “อาสวา เจว อาสวสมฺปยุตฺตา จา”ติปิ, “อาสวสมฺปยุตฺตา เจว โน จ อาสวา”ติปิ.\csromanspacer{} มนายตนํ น วตฺตพฺพํ “อาสโว เจว อาสวสมฺปยุตฺตญฺจา”ติ, สิยา อาสวสมฺปยุตฺตญฺเจว โน จ อาสโว, สิยา น วตฺตพฺพํ “อาสวสมฺปยุตฺตญฺเจว โน จ อาสโว”ติ.\csromanspacer{} ธมฺมายตนํ สิยา อาสโว เจว อาสวสมฺปยุตฺตญฺจ, สิยา อาสวสมฺปยุตฺตญฺเจว โน จ อาสโว, สิยา น วตฺตพฺพํ “อาสโว เจว อาสวสมฺปยุตฺตญฺจา”ติปิ, “อาสวสมฺปยุตฺตญฺเจว โน จ อาสโว”ติปิ.\csromanspacer{} ทสายตนา อาสววิปฺปยุตฺตสาสวา, ทฺวายตนา สิยา อาสววิปฺปยุตฺตสาสวา, สิยา อาสววิปฺปยุตฺต-อนาสวา, สิยา น วตฺตพฺพา “อาสววิปฺปยุตฺตสาสวา”ติปิ, “อาสววิปฺปยุตฺต-อนาสวา”ติปิ. (3)}

\prose{เอกาทสายตนา โน สํโยชนา.\csromanspacer{} ธมฺมายตนํ สิยา สํโยชนํ, สิยา โน สํโยชนํ.\csromanspacer{} ทสายตนา สํโยชนิยา.\csromanspacer{} ทฺวายตนา สิยา สํโยชนิยา, สิยา อสํโยชนิยา.\csromanspacer{} ทสายตนา สํโยชนวิปฺปยุตฺตา.\csromanspacer{} ทฺวายตนา สิยา สํโยชนสมฺปยุตฺตา, สิยา สํโยชนวิปฺปยุตฺตา.\csromanspacer{} ทสายตนา น วตฺตพฺพา “สํโยชนา เจว สํโยชนิยา จา”ติ, สํโยชนิยา เจว โน จ สํโยชนา.\csromanspacer{} มนายตนํ น วตฺตพฺพํ “สํโยชนญฺเจว สํโยชนิยญฺจา”ติ, สิยา สํโยชนิยญฺเจว โน จ สํโยชนํ, สิยา น วตฺตพฺพํ “สํโยชนิยญฺเจว โน จ สํโยชนนฺ”ติ.\csromanspacer{} ธมฺมายตนํ สิยา สํโยชนญฺเจว สํโยชนิยญฺจ, สิยา สํโยชนิยญฺเจว โน จ สํโยชนํ, สิยา น วตฺตพฺพํ “สํโยชนญฺจว สํโยชนิยนฺ”ติปิ, “สํโยชนิยญฺเจว โน จ สํโยชนนฺ”ติปิ.\csromanspacer{} ทสายตนา น วตฺตพฺพา “สํโยชนา เจว สํโยชนสมฺปยุตฺตา จา”ติปิ, “สํโยชนสมฺปยุตฺตา เจว โน จ สํโยชนา”ติปิ.\csromanspacer{} มนายตนํ น วตฺตพฺพํ “สํโยชนญฺเจว สํโยชนสมฺปยุตฺตญฺจา”ติ, สิยา สํโยชนสมฺปยุตฺตญฺเจว โน จ สํโยชนํ, สิยา น วตฺตพฺพํ “สํโยชนสมฺปยุตฺตญฺเจว โน จ สํโยชนนฺ”ติ.\csromanspacer{} ธมฺมายตนํ สิยา สํโยชนญฺเจว สํโยชนสมฺปยุตฺตญฺจ, สิยา สํโยชนสมฺปยุตฺตญฺเจว โน จ สํโยชนํ, สิยา น วตฺตพฺพํ “สํโยชนญฺเจว สํโยชนสมฺปยุตฺตญฺจา”ติปิ, “สํโยชนสมฺปยุตฺตญฺเจว โน จ สํโยชนนฺ”ติปิ.\csromanspacer{} ทสายตนา}

\vspace{\tipitakaparskip}
\csromancenter{อายตนวิภงฺค สํโยชนวิปฺปยุตฺตสํโยชนิยา.\csromanspacer{} ทฺวายตนา สิยา สํโยชนวิปฺปยุตฺตสํโยชนิยา, สิยา สํโยชนวิปฺปยุตฺต-อสํโยชนิยา, สิยา น วตฺตพฺพา “สํโยชนวิปฺปยุตฺตสํโยชนิยา”ติปิ, “สํโยชนวิปฺปยุตฺต-อสํโยชนิยา”ติปิ. (4)}
\prose{\csromanfolio{79}เอกาทสายตนา โน คนฺถา.\csromanspacer{} ธมฺมายตนํ สิยา คนฺโถ, สิยา โน คนฺโถ.\csromanspacer{} ทสายตนา คนฺถนิยา.\csromanspacer{} ทฺวายตนา สิยา คนฺถนิยา, สิยา อคนฺถนิยา.\csromanspacer{} ทสายตนา คนฺถวิปฺปยุตฺตา.\csromanspacer{} ทฺวายตนา สิยา คนฺถสมฺปยุตฺตา, สิยา คนฺถวิปฺปยุตฺตา.\csromanspacer{} ทสายตนา น วตฺตพฺพา “คนฺถา เจว คนฺถนิยา จา”ติ, คนฺถนิยา เจว โน จ คนฺถา.\csromanspacer{} มนายตนํ น วตฺตพฺพํ “คนฺโถ เจว คนฺถนิยญฺจา”ติ, สิยา คนฺถนิยญฺเจว โน จ คนฺโถ, สิยา น วตฺตพฺพํ “คนฺถนิยญฺเจว โน จ คนฺโถ”ติ.\csromanspacer{} ธมฺมายตนํ สิยา คนฺโถ เจว คนฺถนิยญฺจ, สิยา คนฺถนิยญฺเจว โน จ คนฺโถ, สิยา น วตฺตพฺพํ “คนฺโถ เจว คนฺถนิยญฺจา”ติปิ, “คนฺถนิยญฺเจว โน จ คนฺโถ”ติปิ, ทสายตนา น วตฺตพฺพา “คนฺถา เจว คนฺถสมฺปยุตฺตา จา”ติปิ, “คนฺถสมฺปยุตฺตา เจว โน จ คนฺถา”ติปิ.\csromanspacer{} มนายตนํ น วตฺตพฺพํ “คนฺโถ เจว คนฺถสมฺปยุตฺตญฺจา”ติ, สิยา คนฺถสมฺปยุตฺตญฺเจว โน จ คนฺโถ, สิยา น วตฺตพฺพํ “คนฺถสมฺปยุตฺตญฺเจว โน จ คนฺโถ”ติ.\csromanspacer{} ธมฺมายตนํ สิยา คนฺโถ เจว คนฺถสมฺปยุตฺตญฺจ, สิยา คนฺถสมฺปยุตฺตญฺเจว โน จ คนฺโถ, สิยา น วตฺตพฺพํ “คนฺโถ เจว คนฺถสมฺปยุตฺตญฺจา”ติปิ, “คนฺถสมฺปยุตฺตญฺเจว โน จ คนฺโถ”ติปิ.\csromanspacer{} ทสายตนา คนฺถวิปฺปยุตฺตคนฺถนิยา.\csromanspacer{} ทฺวายตนา สิยา คนฺถวิปฺปยุตฺตคนฺถนิยา, สิยา คนฺถวิปฺปยุตฺต-อคนฺถนิยา, สิยา น วตฺตพฺพา “คนฺถวิปฺปยุตฺตคนฺถนิยา”ติปิ, “คนฺถวิปฺปยุตฺต-อคนฺถนิยา”ติปิ. (5)}

\prose{เอกาทสายตนา โน โอฆา ฯเปฯ โน โยคา ฯเปฯ โน นีวรณา.\csromanspacer{} ธมฺมายตนํ สิยา นีวรณํ, สิยา โน นีวรณํ.\csromanspacer{} ทสายตนา นีวรณิยา.\csromanspacer{} ทฺวายตนา สิยา นีวรณิยา, สิยา อนีวรณิยา.\csromanspacer{} ทสายตนา นีวรณวิปฺปยุตฺตา.\csromanspacer{} ทฺวายตนา สิยา นีวรณสมฺปยุตฺตา, สิยา นีวรณวิปฺปยุตฺตา.\csromanspacer{} ทสายตนา น วตฺตพฺพา “นีวรณา เจว นีวรณิยา จา”ติ.\csromanspacer{} นีวรณิยา เจว โน จ นีวรณา.\csromanspacer{} มนายตนํ น วตฺตพฺพํ “นีวรณญฺเจว นีวรณิยญฺจา”ติ, สิยา นีวรณิยญฺเจว โน จ นีวรณํ, สิยา น วตฺตพฺพํ “นีวรณิยญฺเจว โน จ นีวรณนฺ”ติ.\csromanspacer{} ธมฺมายตนํ สิยา นีวรณญฺเจว นีวรณิยญฺจ, สิยา นีวรณิยญฺเจว โน จ นีวรณํ, \csromanfolio{80}สิยา น วตฺตพฺพํ “นีวรณญฺเจว นีวรณิยญฺจา”ติปิ, “นีวรณิยญฺเจว โน จ นีวรณนฺ”ติปิ.\csromanspacer{} ทสายตนา น วตฺตพฺพา “นีวรณา เจว นีวรณสมฺปยุตฺตา จา”ติปิ, “นีวรณสมฺปยุตฺตา เจว โน จ นีวรณา”ติปิ.\csromanspacer{} มนายตนํ น วตฺตพฺพํ “นีวรณญฺเจว นีวรณสมฺปยุตฺตญฺจา”ติ, สิยา นีวรณสมฺปยุตฺตญฺเจว โน จ นีวรณํ, สิยา น วตฺตพฺพํ “นีวรณสมฺปยุตฺตญฺเจว โน จ นีวรณนฺ”ติ.\csromanspacer{} ธมฺมายตนํ สิยา นีวรณญฺเจว นีวรณสมฺปยุตฺตญฺจ, สิยา นีวรณสมฺปยุตฺตญฺเจว โน จ นีวรณํ, สิยา น วตฺตพฺพํ “นีวรณญฺเจว นีวรณสมฺปยุตฺตญฺจา”ติปิ, “นีวรณสมฺปยุตฺตญฺเจว โน จ นีวรณนฺ”ติปิ.\csromanspacer{} ทสายตนา นีวรณวิปฺปยุตฺตนีวรณิยา.\csromanspacer{} ทฺวายตนา สิยา นีวรณวิปฺปยุตฺตนีวรณิยา, สิยา นีวรณวิปฺปยุตฺต-อนีวรณิยา, สิยา น วตฺตพฺพา “นีวรณวิปฺปยุตฺตนีวรณิยา”ติปิ, “นีวรณวิปฺปยุตฺต- อนีวรณิยา”ติปิ. (8)}

\prose{เอกาทสายตนา โน ปรามาสา.\csromanspacer{} ธมฺมายตนํ สิยา ปรามาโส, สิยา โน ปรามาโส.\csromanspacer{} ทสายตนา ปรามฏฺฐา.\csromanspacer{} ทฺวายตนา สิยา ปรามฏฺฐา, สิยา อปรามฏฺฐา.\csromanspacer{} ทสายตนา ปรามาสวิปฺปยุตฺตา.\csromanspacer{} มนายตนํ สิยา ปรามาสสมฺปยุตฺตํ, สิยา ปรามาสวิปฺปยุตฺตํ.\csromanspacer{} ธมฺมายตนํ สิยา ปรามาสสมฺปยุตฺตํ, สิยา ปรามาสวิปฺปยุตฺตํ, สิยา น วตฺตพฺพํ “ปรามาสสมฺปยุตฺตนฺ”ติปิ, “ปรามาสวิปฺปยุตฺตนฺ”ติปิ.\csromanspacer{} ทสายตนา น วตฺตพฺพา “ปรามาสา เจว ปรามฏฺฐา จา”ติ, ปรามฏฺฐา เจว โน จ ปรามาสา.\csromanspacer{} มนายตนํ น วตฺตพฺพํ “ปรามาโส เจว ปรามฏฺฐญฺจา”ติ, สิยา ปรามฏฺฐญฺเจว โน จ ปรามาโส, สิยา น วตฺตพฺพํ “ปรามฏฺฐญฺเจว โน จ ปรามาโส”ติ.\csromanspacer{} ธมฺมายตนํ สิยา ปรามาโส เจว ปรามฏฺฐญฺจ, สิยา ปรามฏฺฐญฺเจว โน จ ปรามาโส, สิยา น วตฺตพฺพํ “ปรามาโส เจว ปรามฏฺฐญฺจา”ติปิ, “ปรามฏฺฐญฺเจว โน จ ปรามาโส”ติปิ.\csromanspacer{} ทสายตนา ปรามาสวิปฺปยุตฺตปรามฏฺฐา.\csromanspacer{} ทฺวายตนา สิยา ปรามาสวิปฺปยุตฺตปรามฏฺฐา, สิยา ปรามาสวิปฺปยุตฺต-อปรามฏฺฐา, สิยา น วตฺตพฺพา “ปรามาสวิปฺปยุตฺตปรามฏฺฐา”ติปิ, “ปรามาสวิปฺปยุตฺต- อปรามฏฺฐา”ติปิ. (9)}

\prose{ทสายตนา อนารมฺมณา.\csromanspacer{} มนายตนํ สารมฺมณํ.\csromanspacer{} ธมฺมายตนํ สิยา สารมฺมณํ, สิยา อนารมฺมณํ.\csromanspacer{} มนายตนํ จิตฺตํ.\csromanspacer{} เอกาทสายตนา โน จิตฺตา.\csromanspacer{} เอกาทสายตนา อเจตสิกา.\csromanspacer{} ธมฺมายตนํ สิยา เจตสิกํ, สิยา อเจตสิกํ, ทสายตนา จิตฺตวิปฺปยุตฺตา.\csromanspacer{} ธมฺมายตนํ สิยา จิตฺตสมฺปยุตฺตํ, สิยา จิตฺตวิปฺปยุตฺตํ.\csromanspacer{} มนายตนํ น วตฺตพฺพํ “จิตฺเตน}

\vspace{\tipitakaparskip}
\csromancenter{อายตนวิภงฺค สมฺปยุตฺตนฺ”ติปิ, “จิตฺเตน วิปฺปยุตฺตนฺ”ติปิ, ทสายตนํ จิตฺตวิสํสฏฺฐา.\csromanspacer{} ธมฺมายตนํ สิยา จิตฺตสํสฏฺฐํ, สิยา จิตฺตวิสํสฏฺฐํ.\csromanspacer{} มนายตนํ น วตฺตพฺพํ “จิตฺเตน สํสฏฺฐนฺ”ติปิ, “จิตฺเตน วิสํสฏฺฐนฺ”ติปิ.\csromanspacer{} ฉายตนา โน จิตฺตสมุฏฺฐานา.\csromanspacer{} ฉายตนา สิยา จิตฺตสมุฏฺฐานา, สิยา โน จิตฺตสมุฏฺฐานา.\csromanspacer{} เอกาทสายตนา โน จิตฺตสหภุโน.\csromanspacer{} ธมฺมายตนํ สิยา จิตฺตสหภู, สิยา โน จิตฺตสหภู.\csromanspacer{} เอกาทสายตนา โน จิตฺตานุปริวตฺติโน.\csromanspacer{} ธมฺมายตนํ สิยา จิตฺตานุปริวตฺติ, สิยา โน จิตฺตานุปริวตฺติ.\csromanspacer{} เอกาทสายตนา โน จิตฺตสํสฏฺฐสมุฏฺฐานา.\csromanspacer{} ธมฺมายตนํ สิยา จิตฺตสํสฏฺฐสมุฏฺฐานํ, สิยา โน จิตฺตสํสฏฺฐสมุฏฺฐานํ.\csromanspacer{} เอกาทสายตนา โน จิตฺตสํสฏฺฐสมุฏฺฐานสหภุโน.\csromanspacer{} ธมฺมายตนํ สิยา จิตฺตสํสฏฺฐสมุฏฺฐานสหภู, สิยา โน จิตฺตสํสฏฺฐสมุฏฺฐานสหภู.\csromanspacer{} เอกาทสายตนา โน จิตฺตสํสฏฺฐสมุฏฺฐานานุปริวตฺติโน.\csromanspacer{} ธมฺมายตนํ สิยา จิตฺตสํสฏฺฐสมุฏฺฐานานุปริวตฺติ, สิยา โน จิตฺตสํสฏฺฐสมุฏฺฐานานุปริวตฺติ. (10)}
\prose{\csromanfolio{81}ฉายตนา อชฺฌตฺติกา.\csromanspacer{} ฉายตนา พาหิรา.\csromanspacer{} นวายตนา อุปาทา.\csromanspacer{} ทฺวายตนา โน อุปาทา.\csromanspacer{} ธมฺมายตนํ สิยา อุปาทา, สิยา โน อุปาทา.\csromanspacer{} ปญฺจายตนา อุปาทินฺนา.\csromanspacer{} สทฺทายตนํ อนุปาทินฺนํ.\csromanspacer{} ฉายตนา สิยา อุปาทินฺนา, สิยา อนุปาทินฺนา.\csromanspacer{} เอกาทสายตนา โน อุปาทานา.\csromanspacer{} ธมฺมายตนํ สิยา อุปาทานํ, สิยา โน อุปาทานํ.\csromanspacer{} ทสายตนา อุปาทานิยา.\csromanspacer{} ทฺวายตนา สิยา อุปาทานิยา, สิยา อนุปาทานิยา.\csromanspacer{} ทสายตนา อุปาทานวิปฺปยุตฺตา.\csromanspacer{} ทฺวายตนา สิยา อุปาทานสมฺปยุตฺตา, สิยา อุปาทานวิปฺปยุตฺตา.\csromanspacer{} ทสายตนา น วตฺตพฺพา “อุปาทานา เจว อุปาทานิยา จา”ติ, อุปาทานิยา เจว โน จ อุปาทานา.\csromanspacer{} มนายตนํ น วตฺตพฺพํ “อุปาทานญฺเจว อุปาทานิยญฺจา”ติ, สิยา อุปาทานิยญฺเจว โน จ อุปาทานํ, สิยา น วตฺตพฺพํ “อุปาทานิยญฺเจว โน จ อุปาทานนฺ”ติ.\csromanspacer{} ธมฺมายตนํ สิยา อุปาทานญฺเจว อุปาทานิยญฺจ, สิยา อุปาทานิยญฺเจว โน จ อุปาทานํ, สิยา น วตฺตพฺพํ “อุปาทานญฺเจว อุปาทานิยญฺจา”ติปิ, “อุปาทานิยญฺเจว โน จ อุปาทานนฺ”ติปิ.\csromanspacer{} ทสายตนา น วตฺตพฺพา “อุปาทานา เจว อุปาทานสมฺปยุตฺตา จา”ติปิ, “อุปาทานสมฺปยุตฺตา เจว โน จ อุปาทานา”ติปิ.\csromanspacer{} มนายตนํ น วตฺตพฺพํ “อุปาทานิยญฺเจว อุปาทานสมฺปยุตฺตญฺจา”ติ, สิยา อุปาทานสมฺปยุตฺตญฺเจว โน จ อุปาทานํ, สิยา น วตฺตพฺพํ “อุปาทานสมฺปยุตฺตญฺเจว โน จ อุปาทานนฺ”ติ.\csromanspacer{} ธมฺมายตนํ สิยา อุปาทานญฺเจว อุปาทานสมฺปยุตฺตญฺจ, สิยา อุปาทานสมฺปยุตฺตญฺเจว โน จ \csromanfolio{82}อุปาทานํ, สิยา น วตฺตพฺพํ “อุปาทานญฺเจว อุปาทานสมฺปยุตฺตญฺจา”ติปิ, “อุปาทานสมฺปยุตฺตญฺเจว โน จ อุปาทานนฺ”ติปิ.\csromanspacer{} ทสายตนา อุปาทานวิปฺปยุตฺต อุปาทานิยา.\csromanspacer{} ทฺวายตนา สิยา อุปาทานสมฺปยุตฺต- อุปาทานิยา, สิยา อุปาทานวิปฺปยุตฺต-อนุปาทานิยา, สิยา น วตฺตพฺพา “อุปาทานวิปฺปยุตฺต อุปาทานิยา”ติปิ, “อุปาทานวิปฺปยุตฺต-อนุปาทานิยา”ติปิ. (11)}

\prose{เอกาทสายตนา โน กิเลสา.\csromanspacer{} ธมฺมายตนํ สิยา กิเลโส, สิยา โน กิเลโส.\csromanspacer{} ทสายตนา สํกิเลสิกา.\csromanspacer{} ทฺวายตนา สิยา สํกิเลสิกา, สิยา อสํกิเลสิกา.\csromanspacer{} ทสายตนา อสํกิลิฏฺฐา.\csromanspacer{} ทฺวายตนา สิยา สํกิลิฏฺฐา, สิยา อสํกิลิฏฺฐา.\csromanspacer{} ทสายตนา กิเลสวิปฺปยุตฺตา.\csromanspacer{} ทฺวายตนา สิยา กิเลสสมฺปยุตฺตา, สิยา กิเลสวิปฺปยุตฺตา.\csromanspacer{} ทสายตนา น วตฺตพฺพา “กิเลสา เจว สํกิเลสิกา จา”ติ, สํกิเลสิกา เจว โน จ กิเลสา.\csromanspacer{} มนายตนํ น วตฺตพฺพํ “กิเลโส เจว สํกิเลสิกญฺจา”ติ, สิยา สํกิเลสิกญฺเจว โน จ กิเลโส, สิยา น วตฺตพฺพํ “สํกิเลสิกญฺเจว โน จ กิเลโส”ติ.\csromanspacer{} ธมฺมายตนํ สิยา กิเลโส เจว สํกิเลสิกญฺจ, สิยา สํกิเลสิกญฺเจว โน จ กิเลโส, สิยา น วตฺตพฺพํ “กิเลโส เจว สํกิเลสิกญฺจา”ติปิ, “สํกิเลสิกญฺเจว โน จ กิเลโส”ติปิ.\csromanspacer{} ทสายตนา น วตฺตพฺพา “กิเลสา เจว สํกิลิฏฺฐา จา”ติปิ, “สํกิลิฏฺฐา เจว โน จ กิเลสา”ติปิ.\csromanspacer{} มนายตนํ น วตฺตพฺพํ “กิเลโส เจว สํกิลิฏฺฐญฺจา”ติ, สิยา สํกิลิฏฺฐญฺเจว โน จ กิเลโส, สิยา น วตฺตพฺพํ “สํกิลิฏฺฐญฺเจว โน จ กิเลโส”ติ.\csromanspacer{} ธมฺมายตนํ สิยา กิเลโส เจว สํกิลิฏฺฐญฺจ, สิยา สํกิลิฏฺฐญฺเจว โน จ กิเลโส, สิยา น วตฺตพฺพํ “กิเลโส เจว สํกิลิฏฺฐญฺจา”ติปิ, “สํกิลิฏฺฐญฺเจว โน จ กิเลโส”ติปิ. ทสายตนา น วตฺตพฺพา “กิเลสา เจว กิเลสสมฺปยุตฺตา จา”ติปิ, “กิเลสสมฺปยุตฺตา เจว โน จ กิเลสา”ติปิ.\csromanspacer{} มนายตนํ น วตฺตพฺพํ “กิเลโส เจว กิเลสสมฺปยุตฺตญฺจา”ติ, สิยา กิเลสสมฺปยุตฺตญฺเจว โน จ กิเลโส, สิยา น วตฺตพฺพํ “กิเลสสมฺปยุตฺตญฺเจว โน จ กิเลโส”ติ.\csromanspacer{} ธมฺมายตนํ สิยา กิเลโส เจว กิเลสสมฺปยุตฺตญฺจ, สิยา กิเลสสมฺปยุตฺตญฺเจว โน จ กิเลโส, สิยา น วตฺตพฺพํ “กิเลโส เจว กิเลสสมฺปยุตฺตญฺจา”ติปิ, “กิเลสสมฺปยุตฺตญฺเจว}

\vspace{\tipitakaparskip}
\csromancenter{อายตนวิภงฺค โน จ กิเลโส”ติปิ.\csromanspacer{} ทสายตนา กิเลสวิปฺปยุตฺตสํกิเลสิกา.\csromanspacer{} ทฺวายตนา สิยา กิเลสวิปฺปยุตฺตสํกิเลสิกา, สิยา กิเลสวิปฺปยุตฺต-อสํกิเลสิกา, สิยา น วตฺตพฺพา “กิเลสวิปฺปยุตฺตสํกิเลสิกา”ติปิ, “กิเลสวิปฺปยุตฺต- อสํกิเลสิกา”ติปิ. (12)}
\prose{\csromanfolio{83}ทสายตนา น ทสฺสเนน ปหาตพฺพา.\csromanspacer{} ทฺวายตนา สิยา ทสฺสเนน ปหาตพฺพา, สิยา น ทสฺสเนน ปหาตพฺพา.\csromanspacer{} ทสายตนา น ภาวนาย ปหาตพฺพา.\csromanspacer{} ทฺวายตนา สิยา ภาวนาย ปหาตพฺพา, สิยา น ภาวนาย ปหาตพฺพา.\csromanspacer{} ทสายตนา น ทสฺสเนน ปหาตพฺพเหตุกา.\csromanspacer{} ทฺวายตนา สิยา ทสฺสเนน ปหาตพฺพเหตุกา, สิยา น ทสฺสเนน ปหาตพฺพเหตุกา.\csromanspacer{} ทสายตนา น ภาวนาย ปหาตพฺพเหตุกา.\csromanspacer{} ทฺวายตนา สิยา ภาวนาย ปหาตพฺพเหตุกา, สิยา น ภาวนาย ปหาตพฺพเหตุกา.\csromanspacer{} ทสายตนา อวิตกฺกา.\csromanspacer{} ทฺวายตนา สิยา สวิตกฺกา, สิยา อวิตกฺกา.\csromanspacer{} ทสายตนา อวิจารา.\csromanspacer{} ทฺวายตนา สิยา สวิจารา, สิยา อวิจารา.\csromanspacer{} ทสายตนา อปฺปีติกา.\csromanspacer{} ทฺวายตนา สิยา สปฺปีติกา, สิยา อปฺปีติกา.\csromanspacer{} ทสายตนา น ปีติสหคตา.\csromanspacer{} ทฺวายตนา สิยา ปีติสหคตา, สิยา น ปีติสหคตา.\csromanspacer{} ทสายตนา น สุขสหคตา.\csromanspacer{} ทฺวายตนา สิยา สุขสหคตา, สิยา น สุขสหคตา.\csromanspacer{} ทสายตนา น อุเปกฺขาสหคตา.\csromanspacer{} ทฺวายตนา สิยา อุเปกฺขาสหคตา, สิยา น อุเปกฺขาสหคตา.}

\prose{ทสายตนา กามาวจรา.\csromanspacer{} ทฺวายตนา สิยา กามาวจรา, สิยา น กามาวจรา.\csromanspacer{} ทสายตนา น รูปาวจรา.\csromanspacer{} ทฺวายตนา สิยา รูปาวจรา, สิยา น รูปาวจรา.\csromanspacer{} ทสายตนา น อรูปาวจรา.\csromanspacer{} ทฺวายตนา สิยา อรูปาวจรา, สิยา น อรูปาวจรา.\csromanspacer{} ทสายตนา ปริยาปนฺนา.\csromanspacer{} ทฺวายตนา สิยา ปริยาปนฺนา, สิยา อปริยาปนฺนา.\csromanspacer{} ทสายตนา อนิยฺยานิกา.\csromanspacer{} ทฺวายตนา สิยา นิยฺยานิกา, สิยา อนิยฺยานิกา.\csromanspacer{} ทสายตนา อนิยตา.\csromanspacer{} ทฺวายตนา สิยา นิยตา, สิยา อนิยตา.\csromanspacer{} ทสายตนา ส-อุตฺตรา.\csromanspacer{} ทฺวายตนา สิยา ส-อุตฺตรา, สิยา อนุตฺตรา.\csromanspacer{} ทสายตนา อรณา.\csromanspacer{} ทฺวายตนา สิยา สรณา, สิยา อรณาติ. (13)}

\vspace{\tipitakaparskip}
\csromancenter{ปญฺหาปุจฺฉกํ.}
\nitthitam{\textbf{อายตนวิภงฺโค นิฏฺฐิโต.}}
\csromanmatikaanchor{mtk.23}
\csromantocmarknopage{chapter}{3. ธาตุวิภงฺค}{mtk.23}
\bookmark[dest={mtk.23},level=0]{3. ธาตุวิภงฺค}
\chapterheadpageread{3. ธาตุวิภงฺค}
\csromanmatikaanchor{mtk.24}
\csromantocmarkref{section}{1. สุตฺตนฺตภาชนีย}{mtk.24}
\bookmark[dest={mtk.24},level=1]{1. สุตฺตนฺตภาชนีย}
\csromanheader{1}{1. สุตฺตนฺตภาชนีย}
\proseitem{172}{\csromanfolio{84}ฉ ธาตุโย ปถวีธาตุ\footnote{ปฐวีธาตุ (สี, สฺยา) เอวมุปริปิ.} อาโปธาตุ เตโชธาตุ วาโยธาตุ อากาสธาตุ วิญฺญาณธาตุ.}

\proseitem{173}{ตตฺถ กตมา ปถวีธาตุ, ปถวีธาตุทฺวยํ อตฺถิ อชฺฌตฺติกา, อตฺถิ พาหิรา.\csromanspacer{} ตตฺถ กตมา อชฺฌตฺติกา ปถวีธาตุ, ยํ อชฺฌตฺตํ ปจฺจตฺตํ กกฺขฬํ ขริคตํ กกฺขฬตฺตํ กกฺขฬภาโว อชฺฌตฺตํ อุปาทินฺนํ.\csromanspacer{} เสยฺยถิทํ, เกสา โลมา นขา ทนฺตา ตโจ, มํสํ นฺหารุ\footnote{นหารุ (สี) 3. อฏฺฐิมิญฺชา (สี)} อฏฺฐิ อฏฺฐิมิญฺชํ3 วกฺกํ, หทยํ ยกนํ กิโลมกํ ปิหกํ ปปฺผาสํ, อนฺตํ อนฺตคุณํ อุทริยํ กรีสํ, ยํ วา ปนญฺญมฺปิ อตฺถิ อชฺฌตฺตํ ปจฺจตฺตํ กกฺขฬํ ขริคตํ กกฺขฬตฺตํ กกฺขฬภาโว อชฺฌตฺตํ อุปาทินฺนํ.\csromanspacer{} อยํ วุจฺจติ อชฺฌตฺติกา ปถวีธาตุ.}

\prose{ตตฺถ กตมา พาหิรา ปถวีธาตุ, ยํ พาหิรํ กกฺขฬํ ขริคตํ กกฺขฬตฺตํ กกฺขฬภาโว พหิทฺธา อนุปาทินฺนํ.\csromanspacer{} เสยฺยถิทํ, อโย โลหํ ติปุ สีสํ สชฺฌํ\footnote{สชฺฌุ (สฺยา) 5. โลหิตงฺโค (สฺยา, ก), โลหิตโก (?)} มุตฺตา มณิ เวฬุริโย สงฺโข สิลา ปวาฬํ รชตํ ชาตรูปํ โลหิตงฺโก5 มสารคลฺลํ ติณํ กฏฺฐํ สกฺขรา กฐลํ\footnote{กถลํ (ก)} ภูมิ ปาสาโณ ปพฺพโต, ยํ วา ปนญฺญมฺปิ อตฺถิ พาหิรํ กกฺขฬํ ขริคตํ กกฺขฬตฺตํ กกฺขฬภาโว พหิทฺธา อนุปาทินฺนํ.\csromanspacer{} อยํ วุจฺจติ พาหิรา ปถวีธาตุ.\csromanspacer{} ยา จ อชฺฌตฺติกา ปถวีธาตุ ยา จ พาหิรา ปถวีธาตุ, ตเทกชฺฌํ อภิสญฺญูหิตฺวา อภิสงฺขิปิตฺวา อยํ วุจฺจติ ปถวีธาตุ. (1)}

\proseitem{174}{ตตฺถ กตมา อาโปธาตุ, อาโปธาตุทฺวยํ อตฺถิ อชฺฌตฺติกา, อตฺถิ พาหิรา.\csromanspacer{} ตตฺถ กตมา อชฺฌตฺติกา อาโปธาตุ, ยํ อชฺฌตฺตํ ปจฺจตฺตํ อาโป อาโปคตํ สิเนโห สิเนหคตํ\footnote{สฺเนโหสฺเนหคตํ(สฺยา)} พนฺธนตฺตํ รูปสฺส อชฺฌตฺตํ อุปาทินฺนํ.\csromanspacer{} เสยฺยถิทํ, ปิตฺตํ เสมฺหํ ปุพฺโพ โลหิตํ เสโท เมโท อสฺสุ วสา เขโฬ สิงฺฆาณิกา ลสิกา มุตฺตํ, ยํ วา ปนญฺญมฺปิ อตฺถิ}

\vspace{\tipitakaparskip}
\csromancenter{ธาตุวิภงฺค อชฺฌตฺตํ ปจฺจตฺตํ อาโป อาโปคตํ สิเนโห สิเนหคตํ พนฺธนตฺตํ รูปสฺส อชฺฌตฺตํ อุปาทินฺนํ.\csromanspacer{} อยํ วุจฺจติ อชฺฌตฺติกา อาโปธาตุ.}
\prose{\csromanfolio{85}ตตฺถ กตมา พาหิรา อาโปธาตุ, ยํ พาหิรํ อาโป อาโปคตํ สิเนโห สิเนหคตํ พนฺธนตฺตํ รูปสฺส พหิทฺธา อนุปาทินฺนํ.\csromanspacer{} เสยฺยถิทํ, มูลรโส ขนฺธรโส ตจรโส ปตฺตรโส ปุปฺผรโส ผลรโส ขีรํ ทธิ สปฺปิ นวนีตํ เตลํ มธุ ผาณิตํ ภุมฺมานิ วา อุทกานิ อนฺตลิกฺขานิ วา, ยํ วา ปนญฺญมฺปิ อตฺถิ พาหิรํ อาโป อาโปคตํ สิเนโห สิเนหคตํ พนฺธนตฺตํ รูปสฺส พหิทฺธา อนุปาทินฺนํ.\csromanspacer{} อยํ วุจฺจติ พาหิรา อาโปธาตุ.\csromanspacer{} ยา จ อชฺฌตฺติกา อาโปธาตุ ยา จ พาหิรา อาโปธาตุ, ตเทกชฺฌํ อภิสญฺญูหิตฺวา อภิสงฺขิปิตฺวา อยํ วุจฺจติ อาโปธาตุ. (2)}

\proseitem{175}{ตตฺถ กตมา เตโชธาตุ, เตโชธาตุทฺวยํ อตฺถิ อชฺฌตฺติกา, อตฺถิ พาหิรา.\csromanspacer{} ตตฺถ กตมา อชฺฌตฺติกา เตโชธาตุ, ยํ อชฺฌตฺตํ ปจฺจตฺตํ เตโช เตโชคตํ อุสฺมา อุสฺมาคตํ อุสุมํ อุสุมคตํ อชฺฌตฺตํ อุปาทินฺนํ.\csromanspacer{} เสยฺยถิทํ, เยน จ สนฺตปฺปติ, เยน จ ชีรียติ, เยน จ ปริฑยฺหติ, เยน จ อสิตปีตขายิตสายิตํ สมฺมา ปริณามํ คจฺฉติ, ยํ วา ปนญฺญมฺปิ อตฺถิ อชฺฌตฺตํ ปจฺจตฺตํ เตโช เตโชคตํ อุสฺมา อุสฺมาคตํ อุสุมํ อุสุมคตํ อชฺฌตฺตํ อุปาทินฺนํ.\csromanspacer{} อยํ วุจฺจติ อชฺฌตฺติกา เตโชธาตุ.}

\prose{ตตฺถ กตมา พาหิรา เตโชธาตุ, ยํ พาหิรํ เตโช เตโชคตํ อุสฺมา อุสฺมาคตํ อุสุมํ อุสุมคตํ พหิทฺธา อนุปาทินฺนํ.\csromanspacer{} เสยฺยถิทํ, กฏฺฐคฺคิ ปลาลคฺคิ\footnote{สกลิกคฺคิ (สพฺพตฺถ)} ติณคฺคิ โคมยคฺคิ ถุสคฺคิ สงฺการคฺคิ อินฺทคฺคิ อคฺคิสนฺตาโป สูริยสนฺตาโป กฏฺฐสนฺนิจยสนฺตาโป ติณสนฺนิจยสนฺตาโป ธญฺญสนฺนิจยสนฺตาโป ภณฺฑสนฺนิจยสนฺตาโป, ยํ วา ปนญฺญมฺปิ อตฺถิ พาหิรํ เตโช เตโชคตํ อุสฺมา อุสฺมาคตํ อุสุมํ อุสุมคตํ พหิทฺธา อนุปาทินฺนํ.\csromanspacer{} อยํ วุจฺจติ พาหิรา เตโชธาตุ.\csromanspacer{} ยา จ อชฺฌตฺติกา เตโชธาตุ ยา จ พาหิรา เตโชธาตุ, ตเทกชฺฌํ อภิสญฺญูหิตฺวา อภิสงฺขิปิตฺวา อยํ วุจฺจติ เตโชธาตุ. (3)}

\proseitem{176}{\csromanfolio{86}ตตฺถ กตมา วาโยธาตุ, วาโยธาตุทฺวยํ อตฺถิ อชฺฌตฺติกา, อตฺถิ พาหิรา.\csromanspacer{} ตตฺถ กตมา อชฺฌตฺติกา วาโยธาตุ, ยํ อชฺฌตฺตํ ปจฺจตฺตํ วาโย วาโยคตํ ถมฺภิตตฺตํ รูปสฺส อชฺฌตฺตํ อุปาทินฺนํ.\csromanspacer{} เสยฺยถิทํ, อุทฺธงฺคมา วาตา, อโธคมา วาตา, กุจฺฉิสยา วาตา, โกฏฺฐาสยา\footnote{โกฏฺฐสยา (สี, สฺยา)} วาตา, องฺคมงฺคานุสาริโน วาตา, สตฺถกวาตา, ขุรกวาตา, อุปฺปลกวาตา, อสฺสาโส ปสฺสาโส อิติ วา, ยํ วา ปนญฺญมฺปิ อตฺถิ อชฺฌตฺตํ ปจฺจตฺตํ วาโย วาโยคตํ ถมฺภิตตฺตํ รูปสฺส อชฺฌตฺตํ อุปาทินฺนํ.\csromanspacer{} อยํ วุจฺจติ อชฺฌตฺติกา วาโยธาตุ.}

\prose{ตตฺถ กตมา พาหิรา วาโยธาตุ, ยํ พาหิรํ วาโย วาโยคตํ ถมฺภิตตฺตํ รูปสฺส พหิทฺธา อนุปาทินฺนํ.\csromanspacer{} เสยฺยถิทํ, ปุรตฺถิมา วาตา, ปจฺฉิมา วาตา, อุตฺตรา วาตา, ทกฺขิณา วาตา, สรชา วาตา, อรชา วาตา, สีตา วาตา, อุณฺหา วาตา, ปริตฺตา วาตา, อธิมตฺตา วาตา, กาฬวาตา, เวรมฺภวาตา, ปกฺขวาตา, สุปณฺณวาตา, ตาลวณฺฏวาตา, วิธูปนวาตา, ยํ วา ปนญฺญมฺปิ อตฺถิ พาหิรํ วาโย วาโยคตํ ถมฺภิตตฺตํ รูปสฺส พหิทฺธา อนุปาทินฺนํ.\csromanspacer{} อยํ วุจฺจติ พาหิรา วาโยธาตุ.\csromanspacer{} ยา จ อชฺฌตฺติกา วาโยธาตุ ยา จ พาหิรา วาโยธาตุ, ตเทกชฺฌํ อภิสญฺญูหิตฺวา อภิสงฺขิปิตฺวา อยํ วุจฺจติ วาโยธาตุ. (4)}

\proseitem{177}{ตตฺถ กตมา อากาสธาตุ, อากาสธาตุทฺวยํ อตฺถิ อชฺฌตฺติกา, อตฺถิ พาหิรา.\csromanspacer{} ตตฺถ กตมา อชฺฌตฺติกา อากาสธาตุ, ยํ อชฺฌตฺตํ ปจฺจตฺตํ อากาโส อากาสคตํ อฆํ อฆคตํ วิวโร วิวรคตํ อสมฺผุฏฺฐํ มํสโลหิเตหิ อชฺฌตฺตํ อุปาทินฺนํ.\csromanspacer{} เสยฺยถิทํ, กณฺณจฺฉิทฺทํ นาสจฺฉิทฺทํ มุขทฺวารํ, เยน จ อสิตปีตขายิตสายิตํ อชฺโฌหรติ, ยตฺถ จ อสิตปีตขายิตสายิตํ สนฺติฏฺฐติ, เยน จ อสิตปีตขายิตสายิตํ อโธภาคํ นิกฺขมติ, ยํ วา ปนญฺญมฺปิ อตฺถิ อชฺฌตฺตํ ปจฺจตฺตํ อากาโส อากาสคตํ อฆํ อฆคตํ วิวโร วิวรคตํ อสมฺผุฏฺฐํ มํสโลหิเตหิ อชฺฌตฺตํ อุปาทินฺนํ.\csromanspacer{} อยํ วุจฺจติ อชฺฌตฺติกา อากาสธาตุ.}

\prose{ตตฺถ กตมา พาหิรา อากาสธาตุ, ยํ พาหิรํ อากาโส อากาสคตํ อฆํ อฆคตํ วิวโร วิวรคตํ อสมฺผุฏฺฐํ จตูหิ}

\vspace{\tipitakaparskip}
\csromancenter{ธาตุวิภงฺค มหาภูเตหิ พหิทฺธา อนุปาทินฺนํ.\csromanspacer{} อยํ วุจฺจติ พาหิรา อากาสธาตุ.\csromanspacer{} ยา จ อชฺฌตฺติกา อากาสธาตุ ยา จ พาหิรา อากาสธาตุ, ตเทกชฺฌํ อภิสญฺญุหิตฺวา อภิสงฺขิปิตฺวา อยํ วุจฺจติ อากาสธาตุ. (5)}
\proseitem{178}{\csromanfolio{87}ตตฺถ กตมา วิญฺญาณธาตุ, จกฺขุวิญฺญาณธาตุ โสตวิญฺญาณธาตุ ฆานวิญฺญาณธาตุ ชิวฺหาวิญฺญาณธาตุ กายวิญฺญาณธาตุ มโนวิญฺญาณธาตุ.\csromanspacer{} อยํ วุจฺจติ วิญฺญาณธาตุ. (6) อิมา ฉ ธาตุโย.}

\proseitem{179}{อปราปิ ฉ ธาตุโย สุขธาตุ ทุกฺขธาตุ โสมนสฺสธาตุ โทมนสฺสธาตุ อุเปกฺขาธาตุ อวิชฺชาธาตุ.}

\proseitem{180}{ตตฺถ กตมา สุขธาตุ, ยํ กายิกํ สาตํ กายิกํ สุขํ กายสมฺผสฺสชํ สาตํ สุขํ เวทยิตํ กายสมฺผสฺสชา สาตา สุขา เวทนา.\csromanspacer{} อยํ วุจฺจติ สุขธาตุ. (1)}

\prose{ตตฺถ กตมา ทุกฺขธาตุ, ยํ กายิกํ อสาตํ กายิกํ ทุกฺขํ กายสมฺผสฺสชํ อสาตํ ทุกฺขํ เวทยิตํ กายสมฺผสฺสชา อสาตา ทุกฺขา เวทนา.\csromanspacer{} อยํ วุจฺจติ ทุกฺขธาตุ. (2)}

\prose{ตตฺถ กตมา โสมนสฺสธาตุ, ยํ เจตสิกํ สาตํ เจตสิกํ สุขํ เจโตสมฺผสฺสชํ สาตํ สุขํ เวทยิตํ เจโตสมฺผสฺสชา สาตา สุขา เวทนา.\csromanspacer{} อยํ วุจฺจติ โสมนสฺสธาตุ. (3)}

\prose{ตตฺถ กตมา โทมนสฺสธาตุ, ยํ เจตสิกํ อสาตํ เจตสิกํ ทุกฺขํ เจโตสมฺผสฺสชํ อสาตํ ทุกฺขํ เวทยิตํ เจโตสมฺผสฺสชา อสาตา ทุกฺขา เวทนา.\csromanspacer{} อยํ วุจฺจติ โทมนสฺสธาตุ. (4)}

\prose{ตตฺถ กตมา อุเปกฺขาธาตุ, ยํ เจตสิกํ เนว สาตํ นาสาตํ เจโตสมฺผสฺสชํ อทุกฺขมสุขํ เวทยิตํ เจโตสมฺผสฺสชา อทุกฺขมสุขา เวทนา.\csromanspacer{} อยํ วุจฺจติ อุเปกฺขาธาตุ. (5)}

\prose{ตตฺถ กตมา อวิชฺชาธาตุ, ยํ อญฺญาณํ อทสฺสนํ อนภิสมโย อนนุโพโธ อสมฺโพโธ อปฺปฏิเวโธ อสงฺคาหณา อปริโยคาหณา อสมเปกฺขนา อปจฺจเวกฺขณา อปจฺจกฺขกมฺมํ ทุมฺเมชฺฌํ พาลฺยํ \csromanfolio{88}อสมฺปชญฺญํ โมโห ปโมโห สมฺโมโห อวิชฺชา อวิชฺโชโฆ อวิชฺชาโยโค อวิชฺชานุสโย อวิชฺชาปริยุฏฺฐานํ, อวิชฺชาลงฺคี โมโห อกุสลมูลํ.\csromanspacer{} อยํ วุจฺจติ อวิชฺชาธาตุ. (6) อิมา ฉ ธาตุโย.}

\proseitem{181}{อปราปิ ฉ ธาตุโย กามธาตุ พฺยาปาทธาตุ วิหิํสาธาตุ เนกฺขมฺมธาตุ อพฺยาปาทธาตุ อวิหิํสาธาตุ.}

\proseitem{182}{ตตฺถ กตมา กามธาตุ, กามปฏิสํยุตฺโต ตกฺโก วิตกฺโก สงฺกปฺโป อปฺปนา พฺยปฺปนา เจตโส อภินิโรปนา มิจฺฉาสงฺกปฺโป.\csromanspacer{} อยํ วุจฺจติ กามธาตุ.\csromanspacer{} เหฏฺฐโต อวีจินิรยํ ปริยนฺตํ กริตฺวา อุปริโต ปรนิมฺมิตวสวตฺตี เทเว อนฺโต กริตฺวา ยํ เอตสฺมิํ อนฺตเร เอตฺถาวจรา เอตฺถ ปริยาปนฺนา ขนฺธธาตุ-อายตนา รูปา\footnote{รูปํ (สฺยา)} เวทนา สญฺญา สงฺขารา วิญฺญาณํ.\csromanspacer{} อยํ วุจฺจติ กามธาตุ. (1)}

\prose{ตตฺถ กตมา พฺยาปาทธาตุ, พฺยาปาทปฏิสํยุตฺโต ตกฺโก วิตกฺโก ฯเปฯ.\csromanspacer{} มิจฺฉาสงฺกปฺโป.\csromanspacer{} อยํ วุจฺจติ พฺยาปาทธาตุ.\csromanspacer{} ทสสุ วา อาฆาตวตฺถูสุ จิตฺตสฺส อาฆาโต ปฏิฆาโต ปฏิฆํ ปฏิวิโรโธ โกโป ปโกโป สมฺปโกโป โทโส ปโทโส สมฺปโทโส จิตฺตสฺส พฺยาปตฺติ มโนปโทโส โกโธ กุชฺฌนา กุชฺฌิตตฺตํ โทโส ทุสฺสนา ทุสฺสิตตฺตํ พฺยาปตฺติ พฺยาปชฺชนา วิโรโธ ปฏิวิโรโธ จณฺฑิกฺกํ อสุโรโป อนตฺตมนตา จิตฺตสฺส.\csromanspacer{} อยํ วุจฺจติ พฺยาปาทธาตุ. (2)}

\prose{ตตฺถ กตมา วิหิํสาธาตุ, วิหิํสาปฏิสํยุตฺโต ตกฺโก วิตกฺโก ฯเปฯ มิจฺฉาสงฺกปฺโป.\csromanspacer{} อยํ วุจฺจติ วิหิํสาธาตุ.\csromanspacer{} อิเธกจฺโจ ปาณินา วา เลฑฺฑุนา วา ทณฺเฑน วา สตฺเถน วา รชฺชุยา วา อญฺญตรญฺญตเรน สตฺเต วิเหเฐติ, ยา เอวรูปา เหฐนา วิเหฐนา หิํสนา วิหิํสนา โรสนา วิโรสนา ปรูปฆาโต.\csromanspacer{} อยํ วุจฺจติ วิหิํสาธาตุ. (3)}

\prose{ตตฺถ กตมา เนกฺขมฺมธาตุ, เนกฺขมฺมปฏิสํยุตฺโต ตกฺโก วิตกฺโก ฯเปฯ สมฺมาสงฺกปฺโป.\csromanspacer{} อยํ วุจฺจติ เนกฺขมฺมธาตุ.\csromanspacer{} สพฺเพปิ กุสลา ธมฺมา เนกฺขมฺมธาตุ. (4)}

\csromanheader{2}{3. ธาตุวิภงฺค}
\prose{\csromanfolio{89}ตตฺถ กตมา อพฺยาปาทธาตุ, อพฺยาปาทปฏิสํยุตฺโต ตกฺโก วิตกฺโก ฯเปฯ สมฺมาสงฺกปฺโป.\csromanspacer{} อยํ วุจฺจติ อพฺยาปาทธาตุ.\csromanspacer{} ยา สตฺเตสุ เมตฺติ เมตฺตายนา เมตฺตายิตตฺตํ เมตฺตาเจโตวิมุตฺติ.\csromanspacer{} อยํ วุจฺจติ อพฺยาปาทธาตุ. (5)}

\prose{ตตฺถ กตมา อวิหิํสาธาตุ, อวิหิํสาปฏิสํยุตฺโต ตกฺโก วิตกฺโก สงฺกปฺโป อปฺปนา พฺยปฺปนา เจตโส อภินิโรปนา สมฺมาสงฺกปฺโป.\csromanspacer{} อยํ วุจฺจติ อวิหิํสาธาตุ.\csromanspacer{} ยา สตฺเตสุ กรุณายนา กรุณายิตตฺตํ กรุณาเจโตวิมุตฺติ.\csromanspacer{} อยํ วุจฺจติ อวิหิํสาธาตุ. (6) อิมา ฉ ธาตุโย.}

\prose{อิติ อิมานิ ตีณิ ฉกฺกานิ ตเทกชฺฌํ อภิสญฺญูหิตฺวา อภิสงฺขิปิตฺวา อฏฺฐารส ธาตุโย โหนฺติ.}

\csromanheader{2}{สุตฺตนฺตภาชนียํ.}
\csromansectionrule
\csromanmatikaanchor{mtk.25}
\csromantocmarkref{section}{2. อภิธมฺมภาชนีย}{mtk.25}
\bookmark[dest={mtk.25},level=1]{2. อภิธมฺมภาชนีย}
\csromanheader{1}{2. อภิธมฺมภาชนีย}
\proseitem{183}{อฏฺฐารส ธาตุโย จกฺขุธาตุ รูปธาตุ จกฺขุวิญฺญาณธาตุ โสตธาตุ สทฺทธาตุ โสตวิญฺญาณธาตุ ฆานธาตุ คนฺธธาตุ ฆานวิญฺญาณธาตุ ชิวฺหาธาตุ รสธาตุ ชิวฺหาวิญฺญาณธาตุ กายธาตุ โผฏฺฐพฺพธาตุ กายวิญฺญาณธาตุ มโนธาตุ ธมฺมธาตุ มโนวิญฺญาณธาตุ.}

\proseitem{184}{ตตฺถ กตมา จกฺขุธาตุ, ยํ จกฺขุ จตุนฺนํ มหาภูตานํ อุปาทาย ปสาโท ฯเปฯ สุญฺโญ คาโมเปโส.\csromanspacer{} อยํ วุจฺจติ จกฺขุธาตุ. (1)}

\prose{ตตฺถ กตมา รูปธาตุ, ยํ รูปํ จตุนฺนํ มหาภูตานํ อุปาทาย วณฺณนิภา ฯเปฯ รูปธาตุเปสา.\csromanspacer{} อยํ วุจฺจติ รูปธาตุ. (2)}

\prose{ตตฺถ กตมา จกฺขุวิญฺญาณธาตุ, จกฺขุญฺจ ปฏิจฺจ รูเป จ อุปฺปชฺชติ จิตฺตํ มโน มานสํ หทยํ ปณฺฑรํ มโน มนายตนํ มนินฺทฺริยํ วิญฺญาณํ วิญฺญาณกฺขนฺโธ ตชฺชาจกฺขุวิญฺญาณธาตุ.\csromanspacer{} อยํ วุจฺจติ จกฺขุวิญฺญาณธาตุ. (3)}

\prose{\csromanfolio{90}ตตฺถ กตมา โสตธาตุ, ยํ โสตํ จตุนฺนํ มหาภูตานํ อุปาทาย ปสาโท ฯเปฯ สุญฺโญ คาโมเปโส.\csromanspacer{} อยํ วุจฺจติ โสตธาตุ. (4)}

\prose{ตตฺถ กตมา สทฺทธาตุ, โย สทฺโท จตุนฺนํ มหาภูตานํ อุปาทาย อนิทสฺสโน สปฺปฏิโฆ ฯเปฯ สทฺทธาตุเปสา.\csromanspacer{} อยํ วุจฺจติ สทฺทธาตุ. (5)}

\prose{ตตฺถ กตมา โสตวิญฺญาณธาตุ, โสตญฺจ ปฏิจฺจ สทฺเท จ อุปฺปชฺชติ จิตฺตํ มโน มานสํ หทยํ ปณฺฑรํ มโน มนายตนํ มนินฺทฺริยํ วิญฺญาณํ วิญฺญาณกฺขนฺโธ ตชฺชาโสตวิญฺญาณธาตุ.\csromanspacer{} อยํ วุจฺจติ โสตวิญฺญาณธาตุ. (6)}

\prose{ตตฺถ กตมา ฆานธาตุ, ยํ ฆานํ จตุนฺนํ มหาภูตานํ อุปาทาย ปสาโท ฯเปฯ สุญฺโญ คาโมเปโส.\csromanspacer{} อยํ วุจฺจติ ฆานธาตุ. (7)}

\prose{ตตฺถ กตมา คนฺธธาตุ, โย คนฺโธ จตุนฺนํ มหาภูตานํ อุปาทาย อนิทสฺสโน สปฺปฏิโฆ ฯเปฯ คนฺธธาตุเปสา.\csromanspacer{} อยํ วุจฺจติ คนฺธธาตุ. (8)}

\prose{ตตฺถ กตมา ฆานวิญฺญาณธาตุ, ฆานญฺจ ปฏิจฺจ คนฺเธ จ อุปฺปชฺชติ จิตฺตํ มโน มานสํ หทยํ ปณฺฑรํ มโน มนายตนํ มนินฺทฺริยํ วิญฺญาณํ วิญฺญาณกฺขนฺโธ ตชฺชาฆานวิญฺญาณธาตุ.\csromanspacer{} อยํ วุจฺจติ ฆานวิญฺญาณธาตุ. (9)}

\prose{ตตฺถ กตมา ชิวฺหาธาตุ, ยา ชิวฺหา จตุนฺนํ มหาภูตานํ อุปาทาย ปสาโท ฯเปฯ สุญฺโญ คาโมเปโส.\csromanspacer{} อยํ วุจฺจติ ชิวฺหาธาตุ. (10)}

\prose{ตตฺถ กตมา รสธาตุ, โย รโส จตุนฺนํ มหาภูตานํ อุปาทาย อนิทสฺสโน สปฺปฏิโฆ ฯเปฯ รสธาตุเปสา.\csromanspacer{} อยํ วุจฺจติ รสธาตุ. (11)}

\prose{ตตฺถ กตมา ชิวฺหาวิญฺญาณธาตุ, ชิวฺหญฺจ ปฏิจฺจ รเส จ อุปฺปชฺชติ จิตฺตํ มโน มานสํ หทยํ ปณฺฑรํ มโน มนายตนํ มนินฺทฺริยํ วิญฺญาณํ วิญฺญาณกฺขนฺโธ ตชฺชาชิวฺหาวิญฺญาณธาตุ.\csromanspacer{} อยํ วุจฺจติ ชิวฺหาวิญฺญาณธาตุ. (12)}

\prose{ตตฺถ กตมา กายธาตุ, โย กาโย จตุนฺนํ มหาภูตานํ อุปาทาย ปสาโท ฯเปฯ สุญฺโญ คาโมเปโส.\csromanspacer{} อยํ วุจฺจติ กายธาตุ. (13)}

\csromanheader{2}{3. ธาตุวิภงฺค}
\prose{\csromanfolio{91}ตตฺถ กตมา โผฏฺฐพฺพธาตุ, ปถวีธาตุ ฯเปฯ โผฏฺฐพฺพธาตุเปสา.\csromanspacer{} อยํ วุจฺจติ โผฏฺฐพฺพธาตุ. (14)}

\prose{ตตฺถ กตมา กายวิญฺญาณธาตุ, กายญฺจ ปฏิจฺจ โผฏฺฐพฺเพ จ อุปฺปชฺชติ จิตฺตํ มโน มานสํ หทยํ ปณฺฑรํ มโน มนายตนํ มนินฺทฺริยํ วิญฺญาณํ วิญฺญาณกฺขนฺโธ ตชฺชากายวิญฺญาณธาตุ.\csromanspacer{} อยํ วุจฺจติ กายวิญฺญาณธาตุ. (15)}

\prose{ตตฺถ กตมา มโนธาตุ, จกฺขุวิญฺญาณธาตุยา อุปฺปชฺชิตฺวา นิรุทฺธสมนนฺตรา อุปฺปชฺชติ จิตฺตํ มโน มานสํ หทยํ ปณฺฑรํ มโน มนายตนํ มนินฺทฺริยํ วิญฺญาณํ วิญฺญาณกฺขนฺโธ ตชฺชามโนธาตุ.\csromanspacer{} โสตวิญฺญาณธาตุยา ฯเปฯ.\csromanspacer{} ฆานวิญฺญาณธาตุยา ฯเปฯ.\csromanspacer{} ชิวฺหาวิญฺญาณธาตุยา ฯเปฯ.\csromanspacer{} กายวิญฺญาณธาตุยา อุปฺปชฺชิตฺวา นิรุทฺธสมนนฺตรา อุปฺปชฺชติ จิตฺตํ มโน มานสํ หทยํ ปณฺฑรํ มโน มนายตนํ มนินฺทฺริยํ วิญฺญาณํ วิญฺญาณกฺขนฺโธ ตชฺชามโนธาตุ, สพฺพธมฺเมสุ วา ปน ปฐมสมนฺนาหาโร อุปฺปชฺชติ จิตฺตํ มโน มานสํ หทยํ ปณฺฑรํ มโน มนายตนํ มนินฺทฺริยํ วิญฺญาณํ วิญฺญาณกฺขนฺโธ ตชฺชามโนธาตุ.\csromanspacer{} อยํ วุจฺจติ มโนธาตุ. (16)}

\prose{ตตฺถ กตมา ธมฺมธาตุ, เวทนากฺขนฺโธ สญฺญากฺขนฺโธ สงฺขารกฺขนฺโธ, ยญฺจ รูปํ อนิทสฺสน-อปฺปฏิฆํ ธมฺมายตนปริยาปนฺนํ อสงฺขตา จ ธาตุ.}

\prose{ตตฺถ กตโม เวทนากฺขนฺโธ, เอกวิเธน เวทนากฺขนฺโธ\csromandash{} ผสฺสสมฺปยุตฺโต.\csromanspacer{} ทุวิเธน เวทนากฺขนฺโธ\csromandash{}อตฺถิ สเหตุโก, อตฺถิ อเหตุโก.\csromanspacer{} ติวิเธน เวทนากฺขนฺโธ\csromandash{}อตฺถิ กุสโล, อตฺถิ อกุสโล, อตฺถิ อพฺยากโต ฯเปฯ.\csromanspacer{} เอวํ ทสวิเธน เวทนากฺขนฺโธ ฯเปฯ.\csromanspacer{} เอวํ พหุวิเธน เวทนากฺขนฺโธ.\csromanspacer{} อยํ วุจฺจติ เวทนากฺขนฺโธ. (1)}

\prose{ตตฺถ กตโม สญฺญากฺขนฺโธ, เอกวิเธน สญฺญากฺขนฺโธ\csromandash{} ผสฺสสมฺปยุตฺโต.\csromanspacer{} ทุวิเธน สญฺญากฺขนฺโธ\csromandash{}อตฺถิ สเหตุโก, อตฺถิ อเหตุโก.\csromanspacer{} ติวิเธน สญฺญากฺขนฺโธ\csromandash{}อตฺถิ กุสโล, อตฺถิ อกุสโล, อตฺถิ อพฺยากโต ฯเปฯ.\csromanspacer{} เอวํ ทสวิเธน สญฺญากฺขนฺโธ ฯเปฯ.\csromanspacer{} เอวํ พหุวิเธน สญฺญากฺขนฺโธ.\csromanspacer{} อยํ วุจฺจติ สญฺญากฺขนฺโธ. (2)}

\prose{ตตฺถ กตโม สงฺขารกฺขนฺโธ, เอกวิเธน สงฺขารกฺขนฺโธ\csromandash{} จิตฺตสมฺปยุตฺโต.\csromanspacer{} ทุวิเธน สงฺขารกฺขนฺโธ\csromandash{}อตฺถิ เหตุ, อตฺถิ อเหตุ.\csromanspacer{} ติวิเธน สงฺขารกฺขนฺโธ\csromandash{}อตฺถิ กุสโล, อตฺถิ อกุสโล, อตฺถิ อพฺยากโต ฯเปฯ.}

\prose{\csromanfolio{92}เอวํ ทสวิเธน สงฺขารกฺขนฺโธ ฯเปฯ.\csromanspacer{} เอวํ พหุวิเธน สงฺขารกฺขนฺโธ.\csromanspacer{} อยํ วุจฺจติ สงฺขารกฺขนฺโธ. (1)}

\prose{ตตฺถ กตมํ รูปํ อนิทสฺสน-อปฺปฏิฆํ ธมฺมายตนปริยาปนฺนํ, อิตฺถินฺทฺริยํ ฯเปฯ กพฬีกาโร อาหาโร.\csromanspacer{} อิทํ วุจฺจติ รูปํ อนิทสฺสน-อปฺปฏิฆํ ธมฺมายตนปริยาปนฺนํ. (4)}

\prose{ตตฺถ กตมา อสงฺขตา ธาตุ, ราคกฺขโย โทสกฺขโย โมหกฺขโย.\csromanspacer{} อยํ วุจฺจติ อสงฺขตา ธาตุ.\csromanspacer{} อยํ วุจฺจติ ธมฺมธาตุ. (5-17)}

\prose{ตตฺถ กตมา มโนวิญฺญาณธาตุ, จกฺขุวิญฺญาณธาตุยา อุปฺปชฺชิตฺวา นิรุทฺธสมนนฺตรา อุปฺปชฺชติ มโนธาตุ, มโนธาตุยา อุปฺปชฺชิตฺวา นิรุทฺธสมนนฺตรา อุปฺปชฺชติ จิตฺตํ มโน มานสํ ฯเปฯ ตชฺชามโนวิญฺญาณธาตุ, โสตวิญฺญาณธาตุยา ฯเปฯ.\csromanspacer{} ฆานวิญฺญาณธาตุยา ฯเปฯ.\csromanspacer{} ชิวฺหาวิญฺญาณธาตุยา ฯเปฯ.\csromanspacer{} กายวิญฺญาณธาตุยา อุปฺปชฺชิตฺวา นิรุทฺธสมนนฺตรา อุปฺปชฺชติ มโนธาตุ, มโนธาตุยาปิ อุปฺปชฺชิตฺวา นิรุทฺธสมนนฺตรา อุปฺปชฺชติ จิตฺตํ มโน มานสํ ฯเปฯ ตชฺชามโนวิญฺญาณธาตุ, มนญฺจ ปฏิจฺจ ธมฺเม จ อุปฺปชฺชติ จิตฺตํ มโน มานสํ หทยํ ปณฺฑรํ มโน มนายตนํ มนินฺทฺริยํ วิญฺญาณํ วิญฺญาณกฺขนฺโธ ตชฺชามโนวิญฺญาณธาตุ.\csromanspacer{} อยํ วุจฺจติ มโนวิญฺญาณธาตุ. (18)}

\csromanheader{2}{อภิธมฺมภาชนียํ.}
\csromansectionrule
\csromanmatikaanchor{mtk.26}
\csromantocmarkref{section}{3. ปญฺหาปุจฺฉก}{mtk.26}
\bookmark[dest={mtk.26},level=1]{3. ปญฺหาปุจฺฉก}
\csromanheader{1}{3. ปญฺหาปุจฺฉก}
\proseitem{185}{อฏฺฐารส ธาตุโย จกฺขุธาตุ รูปธาตุ จกฺขุวิญฺญาณธาตุ โสตธาตุ สทฺทธาตุ โสตวิญฺญาณธาตุ ฆานธาตุ คนฺธธาตุ ฆานวิญฺญาณธาตุ ชิวฺหาธาตุ รสธาตุ ชิวฺหาวิญฺญาณธาตุ กายธาตุ โผฏฺฐพฺพธาตุ กายวิญฺญาณธาตุ มโนธาตุ ธมฺมธาตุ มโนวิญฺญาณธาตุ.}

\proseitem{186}{อฏฺฐารสนฺนํ ธาตูนํ กติ กุสลา, กติ อกุสลา, กติ อพฺยากตา ฯเปฯ กติ สรณา, กติ อรณา.}

\csromanmatikaanchor{mtk.27}
\csromantocmarkref{subsection}{1. ติก}{mtk.27}
\bookmark[dest={mtk.27},level=2]{1. ติก}
\csromanheader{2}{3. ธาตุวิภงฺค \textbf{1. ติก}}
\proseitem{187}{\csromanfolio{93}โสฬส ธาตุโย อพฺยากตา.\csromanspacer{} เทฺว ธาตุโย สิยา กุสลา, สิยา อกุสลา, สิยา อพฺยากตา.\csromanspacer{} ทส ธาตุโย น วตฺตพฺพา “สุขาย เวทนาย สมฺปยุตฺตา”ติปิ, “ทุกฺขาย เวทนาย สมฺปยุตฺตา”ติปิ, “อทุกฺขมสุขาย เวทนาย สมฺปยุตฺตา”ติปิ.\csromanspacer{} ปญฺจ ธาตุโย อทุกฺขมสุขาย เวทนาย สมฺปยุตฺตา.\csromanspacer{} กายวิญฺญาณธาตุ สิยา สุขาย เวทนาย สมฺปยุตฺตา, สิยา ทุกฺขาย เวทนาย สมฺปยุตฺตา.\csromanspacer{} มโนวิญฺญาณธาตุ สิยา สุขาย เวทนาย สมฺปยุตฺตา, สิยา ทุกฺขาย เวทนาย สมฺปยุตฺตา, สิยา อทุกฺขมสุขาย เวทนาย สมฺปยุตฺตา.\csromanspacer{} ธมฺมธาตุ สิยา สุขาย เวทนาย สมฺปยุตฺตา, สิยา ทุกฺขาย เวทนาย สมฺปยุตฺตา, สิยา อทุกฺขมสุขาย เวทนาย สมฺปยุตฺตา, สิยา น วตฺตพฺพา “สุขาย เวทนาย สมฺปยุตฺตา”ติปิ, “ทุกฺขาย เวทนาย สมฺปยุตฺตา”ติปิ, “อทุกฺขมสุขาย เวทนาย สมฺปยุตฺตา”ติปิ.}

\prose{ทส ธาตุโย เนววิปากนวิปากธมฺมธมฺมา.\csromanspacer{} ปญฺจธาตุโย วิปากา.\csromanspacer{} มโนธาตุ สิยา วิปากา, สิยา เนววิปากนวิปากธมฺมธมฺมา.\csromanspacer{} เทฺว ธาตุโย สิยา วิปากา, สิยา วิปากธมฺมธมฺมา.\csromanspacer{} สิยา เนววิปากนวิปากธมฺมธมฺมา.}

\prose{ทส ธาตุโย อุปาทินฺนุปาทานิยา.\csromanspacer{} สทฺทธาตุ อนุปาทินฺนุปาทานิยา.\csromanspacer{} ปญฺจ ธาตุโย สิยา อุปาทินฺนุปาทานิยา, สิยา อนุปาทินฺนุปาทานิยา.\csromanspacer{} เทฺว ธาตุโย สิยา อุปาทินฺนุปาทานิยา, สิยา อนุปาทินฺนุปาทานิยา, สิยา อนุปาทินฺน- อนุปาทานิยา.}

\prose{โสฬส ธาตุโย อสํกิลิฏฺฐสํกิเลสิกา.\csromanspacer{} เทฺว ธาตุโย สิยา สํกิลิฏฺฐสํกิเลสิกา, สิยา อสํกิลิฏฺฐสํกิเลสิกา, สิยา อสํกิลิฏฺฐ- อสํกิเลสิกา.\csromanspacer{} ปนฺนรส ธาตุโย อวิตกฺก-อวิจารา.\csromanspacer{} มโนธาตุ สวิตกฺกสวิจารา.\csromanspacer{} มโนวิญฺญาณธาตุ สิยา สวิตกฺกสวิจารา, สิยา อวิตกฺกวิจารมตฺตา, สิยา อวิตกฺก-อวิจารา.\csromanspacer{} ธมฺมธาตุ สิยา สวิตกฺกสวิจารา, สิยา อวิตกฺกวิจารมตฺตา, สิยา อวิตกฺก-อวิจารา, สิยา น วตฺตพฺพา “สวิตกฺกสวิจารา”ติปิ, “อวิตกฺกวิจารมตฺตา”ติปิ, “อวิตกฺก- อวิจารา”ติปิ.\csromanspacer{} ทส ธาตุโย น วตฺตพฺพา “ปีติสหคตา”ติปิ, “สุขสหคตา”ติปิ, “อุเปกฺขาสหคตา”ติปิ.\csromanspacer{} ปญฺจ ธาตุโย อุเปกฺขาสหคตา.\csromanspacer{} กายวิญฺญาณธาตุ น ปีติสหคตา, \csromanfolio{94}สิยา สุขสหคตา, น อุเปกฺขาสหคตา, สิยา น วตฺตพฺพา “สุขสหคตา”ติ.\csromanspacer{} เทฺว ธาตุโย สิยา ปีติสหคตา, สิยา สุขสหคตา, สิยา อุเปกฺขาสหคตา, สิยา น วตฺตพฺพา “ปีติสหคตา”ติปิ, “สุขสหคตา”ติปิ, “อุเปกฺขาสหคตา”ติปิ.}

\prose{โสฬส ธาตุโย เนว ทสฺสเนน น ภาวนาย ปหาตพฺพา.\csromanspacer{} เทฺว ธาตุโย สิยา ทสฺสเนน ปหาตพฺพา, สิยา ภาวนาย ปหาตพฺพา, สิยา เนว ทสฺสเนน น ภาวนาย ปหาตพฺพา.\csromanspacer{} โสฬส ธาตุโย เนว ทสฺสเนน น ภาวนาย ปหาตพฺพเหตุกา.\csromanspacer{} เทฺว ธาตุโย สิยา ทสฺสเนน ปหาตพฺพเหตุกา, สิยา ภาวนาย ปหาตพฺพเหตุกา, สิยา เนว ทสฺสเนน น ภาวนาย ปหาตพฺพเหตุกา.\csromanspacer{} โสฬส ธาตุโย เนวาจยคามินาปจยคามิโน.\csromanspacer{} เทฺว ธาตุโย สิยา อาจยคามิโน, สิยา อปจยคามิโน, สิยา เนวาจยคามินาปจยคามิโน.\csromanspacer{} โสฬส ธาตุโย เนวเสกฺขนาเสกฺขา.\csromanspacer{} เทฺว ธาตุโย สิยา เสกฺขา, สิยา อเสกฺขา.\csromanspacer{} สิยา เนวเสกฺขนาเสกฺขา.}

\prose{โสฬส ธาตุโย ปริตฺตา, เทฺว ธาตุโย สิยา ปริตฺตา, สิยา มหคฺคตา, สิยา อปฺปมาณา.\csromanspacer{} ทส ธาตุโย อนารมฺมณา.\csromanspacer{} ฉ ธาตุโย ปริตฺตารมฺมณา.\csromanspacer{} เทฺว ธาตุโย สิยา ปริตฺตารมฺมณา, สิยา มหคฺคตารมฺมณา, สิยา อปฺปมาณารมฺมณา, สิยา น วตฺตพฺพา “ปริตฺตารมฺมณา”ติปิ, “มหคฺคตารมฺมณา”ติปิ, “อปฺปมาณารมฺมณา”ติปิ.\csromanspacer{} โสฬส ธาตุโย มชฺฌิมา.\csromanspacer{} เทฺว ธาตุโย สิยา หีนา, สิยา มชฺฌิมา, สิยา ปณีตา.\csromanspacer{} โสฬส ธาตุโย อนิยตา.\csromanspacer{} เทฺว ธาตุโย สิยา มิจฺฉตฺตนิยตา, สิยา สมฺมตฺตนิยตา, สิยา อนิยตา. ทส ธาตุโย อนารมฺมณา.\csromanspacer{} ฉ ธาตุโย น วตฺตพฺพา “มคฺคารมฺมณา”ติปิ, “มคฺคเหตุกา”ติปิ, “มคฺคาธิปติโน”ติปิ.\csromanspacer{} เทฺว ธาตุโย สิยา มคฺคารมฺมณา, สิยา มคฺคเหตุกา, สิยา มคฺคาธิปติโน, สิยา น วตฺตพฺพา “มคฺคารมฺมณา”ติปิ, “มคฺคเหตุกา”ติปิ, “มคฺคาธิปติโน”ติปิ.\csromanspacer{} ทส ธาตุโย สิยา อุปฺปนฺนา, สิยา อุปฺปาทิโน, น วตฺตพฺพา “อนุปฺปนฺนา”ติ.\csromanspacer{} สทฺทธาตุ สิยา อุปฺปนฺนา, สิยา อนุปฺปนฺนา, น วตฺตพฺพา “อุปฺปาทินี”ติ.\csromanspacer{} ฉ ธาตุโย สิยา อุปฺปนฺนา, สิยา อนุปฺปนฺนา, สิยา อุปฺปาทิโน.\csromanspacer{} ธมฺมธาตุ สิยา อุปฺปนฺนา, สิยา อนุปฺปนฺนา, สิยา อุปฺปาทินี, สิยา น วตฺตพฺพา “อุปฺปนฺนา”ติปิ, “อนุปฺปนฺนา”ติปิ, “อุปฺปาทินี”ติปิ.}

\csromanheader{2}{3. ธาตุวิภงฺค}
\prose{\csromanfolio{95}สตฺตรส ธาตุโย สิยา อตีตา, สิยา อนาคตา, สิยา ปจฺจุปฺปนฺนา.\csromanspacer{} ธมฺมธาตุ สิยา อตีตา, สิยา อนาคตา, สิยา ปจฺจุปฺปนฺนา, สิยา น วตฺตพฺพา “อตีตา”ติปิ, “อนาคตา”ติปิ, “ปจฺจุปฺปนฺนา”ติปิ.\csromanspacer{} ทส ธาตุโย อนารมฺมณา.\csromanspacer{} ฉ ธาตุโย ปจฺจุปฺปนฺนารมฺมณา.\csromanspacer{} เทฺว ธาตุโย สิยา อตีตารมฺมณา, สิยา อนาคตารมฺมณา, สิยา ปจฺจุปฺปนฺนารมฺมณา, สิยา น วตฺตพฺพา “อตีตารมฺมณา”ติปิ, “อนาคตารมฺมณา”ติปิ, “ปจฺจุปฺปนฺนารมฺมณา”ติปิ, สิยา อชฺฌตฺตา, สิยา พหิทฺธา, สิยา อชฺฌตฺตพหิทฺธา.}

\prose{ทส ธาตุโย อนารมฺมณา.\csromanspacer{} ฉ ธาตุโย สิยา อชฺฌตฺตารมฺมณา, สิยา พหิทฺธารมฺมณา, สิยา อชฺฌตฺตพหิทฺธารมฺมณา.\csromanspacer{} เทฺว ธาตุโย สิยา อชฺฌตฺตารมฺมณา, สิยา พหิทฺธารมฺมณา, สิยา อชฺฌตฺตพหิทฺธารมฺมณา, สิยา น วตฺตพฺพา “อชฺฌตฺตารมฺมณา”ติปิ, “พหิทฺธารมฺมณา”ติปิ, “อชฺฌตฺตพหิทฺธารมฺมณา”ติปิ.\csromanspacer{} รูปธาตุ สนิทสฺสนสปฺปฏิฆา.\csromanspacer{} นว ธาตุโย อนิทสฺสน-อปฺปฏิฆา.\csromanspacer{} อฏฺฐ ธาตุโย อนิทสฺสน-อปฺปฏิฆา.}

\csromanmatikaanchor{mtk.28}
\csromantocmarkref{subsection}{2. ทุก}{mtk.28}
\bookmark[dest={mtk.28},level=2]{2. ทุก}
\csromanheader{2}{2. ทุก}
\proseitem{188}{สตฺตรส ธาตุโย น เหตู.\csromanspacer{} ธมฺมธาตุ สิยา เหตุ, สิยา น เหตุ.\csromanspacer{} โสฬส ธาตุโย อเหตุกา.\csromanspacer{} เทฺว ธาตุโย สิยา สเหตุกา, สิยา อเหตุกา.\csromanspacer{} โสฬส ธาตุโย เหตุวิปฺปยุตฺตา.\csromanspacer{} เทฺว ธาตุโย สิยา เหตุสมฺปยุตฺตา, สิยา เหตุวิปฺปยุตฺตา.\csromanspacer{} โสฬส ธาตุโย น วตฺตพฺพา “เหตุ เจว สเหตุกา จา”ติปิ, “สเหตุกา เจว น จ เหตู”ติปิ.\csromanspacer{} มโนวิญฺญาณธาตุ น วตฺตพฺพา “เหตุ เจว สเหตุกา จา”ติ, สิยา สเหตุกา เจว น จ เหตุ, สิยา น วตฺตพฺพา “สเหตุกา เจว น จ เหตู”ติ.\csromanspacer{} ธมฺมธาตุ สิยา เหตุ เจว สเหตุกา จ, สิยา สเหตุกา เจว น จ เหตุ, สิยา น วตฺตพฺพา “เหตุ เจว สเหตุกา จา”ติปิ, “สเหตุกา เจว น จ เหตู”ติปิ.\csromanspacer{} โสฬส ธาตุโย น วตฺตพฺพา “เหตู เจว เหตุสมฺปยุตฺตา จา”ติปิ, “เหตุสมฺปยุตฺตา เจว น จ เหตู”ติปิ.\csromanspacer{} มโนวิญฺญาณธาตุ น วตฺตพฺพา “เหตุ เจว เหตุสมฺปยุตฺตา จา”ติ, สิยา เหตุสมฺปยุตฺตา เจว น จ เหตุ, สิยา น วตฺตพฺพา “เหตุสมฺปยุตฺตา เจว น จ เหตู”ติ.\csromanspacer{} ธมฺมธาตุ สิยา เหตุ เจว เหตุสมฺปยุตฺตา จ, สิยา เหตุสมฺปยุตฺตา เจว น จ เหตุ, สิยา น วตฺตพฺพา “เหตุ เจว \csromanfolio{96}เหตุสมฺปยุตฺตา จา”ติปิ, “เหตุสมฺปยุตฺตา เจว น จ เหตู”ติปิ.\csromanspacer{} โสฬส ธาตุโย น เหตุ-อเหตุกา.\csromanspacer{} มโนวิญฺญาณธาตุ สิยา น เหตุสเหตุกา, สิยา น เหตุ-อเหตุกา.\csromanspacer{} ธมฺมธาตุ สิยา น เหตุสเหตุกา, สิยา น เหตุ-อเหตุกา, สิยา น วตฺตพฺพา “น เหตุสเหตุกา”ติปิ, “น เหตุ-อเหตุกา”ติปิ. (1)}

\prose{สตฺตรส ธาตุโย สปฺปจฺจยา.\csromanspacer{} ธมฺมธาตุ สิยา สปฺปจฺจยา, สิยา อปฺปจฺจยา.\csromanspacer{} สตฺตรส ธาตุโย สงฺขตา.\csromanspacer{} ธมฺมธาตุ สิยา สงฺขตา, สิยา อสงฺขตา.\csromanspacer{} รูปธาตุ สนิทสฺสนา.\csromanspacer{} สตฺตรส ธาตุโย อนิทสฺสนา.\csromanspacer{} ทส ธาตุโย สปฺปฏิฆา.\csromanspacer{} อฏฺฐ ธาตุโย อปฺปฏิฆา.\csromanspacer{} ทส ธาตุโย รูปา.\csromanspacer{} สตฺต ธาตุโย อรูปา.\csromanspacer{} ธมฺมธาตุ สิยา รูปา, สิยา อรูปา.\csromanspacer{} โสฬส ธาตุโย โลกิยา.\csromanspacer{} เทฺว ธาตุโย สิยา โลกิยา, สิยา โลกุตฺตรา, เกนจิ วิญฺเญยฺยา, เกนจิ น วิญฺเญยฺยา. (2)}

\prose{สตฺตรส ธาตุโย โน อาสวา.\csromanspacer{} ธมฺมธาตุ สิยา อาสวา, สิยา โน อาสวา.\csromanspacer{} โสฬส ธาตุโย สาสวา.\csromanspacer{} เทฺว ธาตุโย สิยา สาสวา, สิยา อนาสวา.\csromanspacer{} โสฬส ธาตุโย อาสววิปฺปยุตฺตา.\csromanspacer{} เทฺว ธาตุโย สิยา อาสวสมฺปยุตฺตา, สิยา อาสววิปฺปยุตฺตา.\csromanspacer{} โสฬส ธาตุโย น วตฺตพฺพา “อาสวา เจว สาสวา จา”ติ, สาสวา เจว โน จ อาสวา.\csromanspacer{} มโนวิญฺญาณธาตุ น วตฺตพฺพา “อาสโว เจว สาสวา จา”ติ, สิยา สาสวา เจว โน จ อาสโว, สิยา น วตฺตพฺพา “สาสวา เจว โน จ อาสโว”ติ.\csromanspacer{} ธมฺมธาตุ สิยา อาสโว เจว สาสวา จ, สิยา สาสวา เจว โน จ อาสโว, สิยา น วตฺตพฺพา “อาสโว เจว สาสวา จา”ติปิ, “สาสวา เจว โน จ อาสโว”ติปิ. โสฬส ธาตุโย น วตฺตพฺพา “อาสวา เจว อาสวสมฺปยุตฺตา จา”ติปิ, “อาสวสมฺปยุตฺตา เจว โน จ อาสวา”ติปิ.\csromanspacer{} มโนวิญฺญาณธาตุ น วตฺตพฺพา “อาสโว เจว อาสวสมฺปยุตฺตา จา”ติ, สิยา อาสวสมฺปยุตฺตา เจว โน จ อาสโว, สิยา น วตฺตพฺพา “อาสวสมฺปยุตฺตา เจว โน จ อาสโว”ติ, ธมฺมธาตุ สิยา อาสโว เจว อาสวสมฺปยุตฺตา จ, สิยา อาสวสมฺปยุตฺตา เจว โน จ อาสโว, สิยา น วตฺตพฺพา “อาสโว เจว อาสวสมฺปยุตฺตา จา”ติปิ, “อาสวสมฺปยุตฺตา เจว โน จ อาสโว”ติปิ.\csromanspacer{} โสฬส ธาตุโย อาสววิปฺปยุตฺตสาสวา.\csromanspacer{} เทฺว ธาตุโย สิยา อาสววิปฺปยุตฺตสาสวา, สิยา}

\vspace{\tipitakaparskip}
\csromancenter{ธาตุวิภงฺค อาสววิปฺปยุตฺต-อนาสวา, สิยา น วตฺตพฺพา “อาสววิปฺปยุตฺตสาสวา”ติปิ, “อาสววิปฺปยุตฺต-อนาสวา”ติปิ. (3)}
\prose{\csromanfolio{97}สตฺตรส ธาตุโย โน สํโยชนา.\csromanspacer{} ธมฺมธาตุ สิยา สํโยชนํ, สิยา โน สํโยชนํ.\csromanspacer{} โสฬส ธาตุโย สํโยชนิยา.\csromanspacer{} เทฺว ธาตุโย สิยา สํโยชนิยา, สิยา อสํโยชนิยา.\csromanspacer{} โสฬส ธาตุโย สํโยชนวิปฺปยุตฺตา.\csromanspacer{} เทฺว ธาตุโย สิยา สํโยชนสมฺปยุตฺตา, สิยา สํโยชนวิปฺปยุตฺตา.\csromanspacer{} โสฬส ธาตุโย น วตฺตพฺพา “สํโยชนา เจว สํโยชนิยา จา”ติ, สํโยชนิยา เจว โน จ สํโยชนา.\csromanspacer{} มโนวิญฺญาณธาตุ น วตฺตพฺพา “สํโยชนญฺเจว สํโยชนิยา จา”ติ, สิยา สํโยชนิยา เจว โน จ สํโยชนํ, สิยา น วตฺตพฺพา “สํโยชนิยา เจว โน จ สํโยชนนฺ”ติ.\csromanspacer{} ธมฺมธาตุ สิยา สํโยชนญฺเจว สํโยชนิยา จ, สิยา สํโยชนิยา เจว โน จ สํโยชนํ, สิยา น วตฺตพฺพา “สํโยชนญฺเจว สํโยชนิยา จา”ติปิ, “สํโยชนิยา เจว โน จ สํโยชนนฺ”ติปิ.\csromanspacer{} โสฬส ธาตุโย น วตฺตพฺพา “สํโยชนา เจว สํโยชนสมฺปยุตฺตา จา”ติปิ, “สํโยชนสมฺปยุตฺตา เจว โน จ สํโยชนา”ติปิ.\csromanspacer{} มโนวิญฺญาณธาตุ น วตฺตพฺพา “สํโยชนญฺเจว สํโยชนสมฺปยุตฺตา จา”ติ, สิยา สํโยชนสมฺปยุตฺตา เจว โน จ สํโยชนํ, สิยา น วตฺตพฺพา “สํโยชนสมฺปยุตฺตา เจว โน จ สํโยชนนฺ”ติ.\csromanspacer{} ธมฺมธาตุ สิยา สํโยชนญฺเจว สํโยชนสมฺปยุตฺตา จ, สิยา สํโยชนสมฺปยุตฺตา เจว โน จ สํโยชนํ, สิยา น วตฺตพฺพา “สํโยชนญฺเจว สํโยชนสมฺปยุตฺตา จา”ติปิ, “สํโยชนสมฺปยุตฺตา เจว โน จ สํโยชนนฺ”ติปิ.\csromanspacer{} โสฬส ธาตุโย สํโยชนวิปฺปยุตฺตสํโยชนิยา.\csromanspacer{} เทฺว ธาตุโย สิยา สํโยชนวิปฺปยุตฺตสํโยชนิยา, สิยา สํโยชนวิปฺปยุตฺต-อสํโยชนิยา, สิยา น วตฺตพฺพา “สํโยชนวิปฺปยุตฺตสํโยชนิยา”ติปิ, “สํโยชนวิปฺปยุตฺต- อสํโยชนิยา”ติปิ. (4)}

\prose{สตฺตรส ธาตุโย โน คนฺถา.\csromanspacer{} ธมฺมธาตุ สิยา คนฺโถ, สิยา โน คนฺโถ.\csromanspacer{} โสฬส ธาตุโย คนฺถนิยา.\csromanspacer{} เทฺว ธาตุโย สิยา คนฺถนิยา, สิยา อคนฺถนิยา.\csromanspacer{} โสฬส ธาตุโย คนฺถวิปฺปยุตฺตา.\csromanspacer{} เทฺว ธาตุโย สิยา คนฺถสมฺปยุตฺตา, สิยา คนฺถวิปฺปยุตฺตา โสฬส \csromanfolio{98}ธาตุโย น วตฺตพฺพา “คนฺถา เจว คนฺถนิยา จา”ติ, คนฺถนิยา เจว โน จ คนฺถา, มโนวิญฺญาณธาตุ น วตฺตพฺพา “คนฺโถ เจว คนฺถนิยา จา”ติ, สิยา คนฺถนิยา เจว โน จ คนฺโถ, สิยา น วตฺตพฺพา “คนฺถนิยา เจว โน จ คนฺโถ”ติ.\csromanspacer{} ธมฺมธาตุ สิยา คนฺโถ เจว คนฺถนิยา จ, สิยา คนฺถนิยา เจว โน จ คนฺโถ, สิยา น วตฺตพฺพา “คนฺโถ เจว คนฺถนิยา จา”ติปิ, “คนฺถนิยา เจว โน จ คนฺโถ”ติปิ.\csromanspacer{} โสฬส ธาตุโย น วตฺตพฺพา “คนฺถา เจว คนฺถสมฺปยุตฺตา จา”ติปิ, “คนฺถสมฺปยุตฺตา เจว โน จ คนฺถา”ติปิ.\csromanspacer{} มโนวิญฺญาณธาตุ น วตฺตพฺพา “คนฺโถ เจว คนฺถสมฺปยุตฺตา จา”ติ, สิยา คนฺถสมฺปยุตฺตา เจว โน จ คนฺโถ, สิยา น วตฺตพฺพา “คนฺถสมฺปยุตฺตา เจว โน จ คนฺโถ”ติ.\csromanspacer{} ธมฺมธาตุ สิยา คนฺโถ เจว คนฺถสมฺปยุตฺตา จ, สิยา คนฺถสมฺปยุตฺตา เจว โน จ คนฺโถ, สิยา น วตฺตพฺพา “คนฺโถ เจว คนฺถสมฺปยุตฺตา จา”ติปิ, “คนฺถสมฺปยุตฺตา เจว โน จ คนฺโถ”ติปิ.\csromanspacer{} โสฬส ธาตุโย คนฺถวิปฺปยุตฺตคนฺถนิยา.\csromanspacer{} เทฺว ธาตุโย สิยา คนฺถวิปฺปยุตฺตคนฺถนิยา, สิยา คนฺถวิปฺปยุตฺต-อคนฺถนิยา, สิยา น วตฺตพฺพา “คนฺถวิปฺปยุตฺตคนฺถนิยา”ติปิ, “คนฺถวิปฺปยุตฺต- อคนฺถนิยา”ติปิ. (5)}

\prose{สตฺตรส ธาตุโย โน โอฆา ฯเปฯ โน โยคา ฯเปฯ โน นีวรณา.\csromanspacer{} ธมฺมธาตุ สิยา นีวรณํ, สิยา โน นีวรณํ.\csromanspacer{} โสฬส ธาตุโย นีวรณิยา.\csromanspacer{} เทฺว ธาตุโย สิยา นีวรณิยา, สิยา อนีวรณิยา.\csromanspacer{} โสฬส ธาตุโย นีวรณวิปฺปยุตฺตา.\csromanspacer{} เทฺว ธาตุโย สิยา นีวรณสมฺปยุตฺตา, สิยา นีวรณวิปฺปยุตฺตา.\csromanspacer{} โสฬส ธาตุโย น วตฺตพฺพา “นีวรณา เจว นีวรณิยา จา”ติ, นีวรณิยา เจว โน จ นีวรณา.\csromanspacer{} มโนวิญฺญาณธาตุ น วตฺตพฺพา “นีวรณญฺเจว นีวรณิยา จา”ติ, สิยา นีวรณิยา เจว โน จ นีวรณํ, สิยา น วตฺตพฺพา “นีวรณิยา เจว โน จ นีวรณนฺ”ติ.\csromanspacer{} ธมฺมธาตุ สิยา นีวรณญฺเจว นีวรณิยา จ, สิยา นีวรณิยา เจว โน จ นีวรณํ, สิยา น วตฺตพฺพา “นีวรณญฺเจว นีวรณิยา จา”ติปิ, “นีวรณิยา เจว โน จ นีวรณนฺ”ติปิ.\csromanspacer{} โสฬส ธาตุโย น วตฺตพฺพา “นีวรณา เจว นีวรณสมฺปยุตฺตา จา”ติปิ, “นีวรณสมฺปยุตฺตา เจว โน จ นีวรณา”ติปิ.\csromanspacer{} มโนวิญฺญาณธาตุ น วตฺตพฺพา “นีวรณญฺเจว นีวรณสมฺปยุตฺตา จา”ติ, สิยา นีวรณสมฺปยุตฺตา เจว โน จ นีวรณํ, สิยา น วตฺตพฺพา “นีวรณสมฺปยุตฺตา เจว}

\vspace{\tipitakaparskip}
\csromancenter{ธาตุวิภงฺค โน จ นีวรณนฺ”ติ.\csromanspacer{} ธมฺมธาตุ สิยา นีวรณญฺเจว นีวรณสมฺปยุตฺตา จ, สิยา นีวรณสมฺปยุตฺตา เจว โน จ นีวรณํ, สิยา น วตฺตพฺพา “นีวรณญฺเจว นีวรณสมฺปยุตฺตา จา”ติปิ, “นีวรณสมฺปยุตฺตา เจว โน จ นีวรณนฺ”ติปิ.\csromanspacer{} โสฬส ธาตุโย นีวรณวิปฺปยุตฺตนีวรณิยา.\csromanspacer{} เทฺว ธาตุโย สิยา นีวรณวิปฺปยุตฺตนีวรณิยา, สิยา นีวรณวิปฺปยุตฺต- อนีวรณิยา, สิยา น วตฺตพฺพา “นีวรณวิปฺปยุตฺตนีวรณิยา”ติปิ, “นีวรณวิปฺปยุตฺต-อนีวรณิยา”ติปิ. (8)}
\prose{\csromanfolio{99}สตฺตรส ธาตุโย โน ปรามาสา.\csromanspacer{} ธมฺมธาตุ สิยา ปรามาโส, สิยา โน ปรามาโส.\csromanspacer{} โสฬส ธาตุโย ปรามฏฺฐา.\csromanspacer{} เทฺว ธาตุโย สิยา ปรามฏฺฐา, สิยา อปรามฏฺฐา.\csromanspacer{} โสฬส ธาตุโย ปรามาสวิปฺปยุตฺตา.\csromanspacer{} มโนวิญฺญาณธาตุ สิยา ปรามาสสมฺปยุตฺตา, สิยา ปรามาสวิปฺปยุตฺตา.\csromanspacer{} ธมฺมธาตุ สิยา ปรามาสสมฺปยุตฺตา, สิยา ปรามาสวิปฺปยุตฺตา, สิยา น วตฺตพฺพา “ปรามาสสมฺปยุตฺตา”ติปิ, “ปรามาสวิปฺปยุตฺตา”ติปิ.\csromanspacer{} โสฬส ธาตุโย น วตฺตพฺพา “ปรามาสา เจว ปรามฏฺฐา จา”ติ, ปรามฏฺฐา เจว โน จ ปรามาสา.\csromanspacer{} มโนวิญฺญาณธาตุ น วตฺตพฺพา “ปรามาโส เจว ปรามฏฺฐา จา”ติ, สิยา ปรามฏฺฐา เจว โน จ ปรามาโส, สิยา น วตฺตพฺพา “ปรามฏฺฐา เจว โน จ ปรามาโส”ติ.\csromanspacer{} ธมฺมธาตุ สิยา ปรามาโส เจว ปรามฏฺฐา จ, สิยา ปรามฏฺฐา เจว โน จ ปรามาโส, สิยา น วตฺตพฺพา “ปรามาโส เจว ปรามฏฺฐา จา”ติปิ, “ปรามฏฺฐา เจว โน จ ปรามาโส”ติปิ.\csromanspacer{} โสฬส ธาตุโย ปรามาสวิปฺปยุตฺตปรามฏฺฐา.\csromanspacer{} เทฺว ธาตุโย สิยา ปรามาสวิปฺปยุตฺตปรามฏฺฐา, สิยา ปรามาสวิปฺปยุตฺต-อปรามฏฺฐา, สิยา น วตฺตพฺพา “ปรามาสวิปฺปยุตฺตปรามฏฺฐา”ติปิ, “ปรามาสวิปฺปยุตฺต- อปรามฏฺฐา”ติปิ. (9)}

\prose{ทส ธาตุโย อนารมฺมณา.\csromanspacer{} สตฺต ธาตุโย สารมฺมณา.\csromanspacer{} ธมฺมธาตุ สิยา สารมฺมณา, สิยา อนารมฺมณา.\csromanspacer{} สตฺต ธาตุโย จิตฺตา.\csromanspacer{} เอกาทส ธาตุโย โน จิตฺตา.\csromanspacer{} สตฺตรส ธาตุโย อเจตสิกา.\csromanspacer{} ธมฺมธาตุ สิยา เจตสิกา, สิยา อเจตสิกา.\csromanspacer{} ทส ธาตุโย จิตฺตวิปฺปยุตฺตา.\csromanspacer{} ธมฺมธาตุ สิยา จิตฺตสมฺปยุตฺตา, สิยา จิตฺตวิปฺปยุตฺตา.\csromanspacer{} สตฺต ธาตุโย น วตฺตพฺพา “จิตฺเตน สมฺปยุตฺตา”ติปิ, “จิตฺเตน วิปฺปยุตฺตา”ติปิ.\csromanspacer{} ทส ธาตุโย จิตฺตวิสํสฏฺฐา.\csromanspacer{} ธมฺมธาตุ สิยา \csromanfolio{100}จิตฺตสํสฏฺฐา, สิยา จิตฺตวิสํสฏฺฐา.\csromanspacer{} สตฺต ธาตุโย น วตฺตพฺพา “จิตฺเตน สํสฏฺฐา”ติปิ, “จิตฺเตน วิสํสฏฺฐา”ติปิ.}

\prose{ทฺวาทส ธาตุโย โน จิตฺตสมุฏฺฐานา.\csromanspacer{} ฉ ธาตุโย สิยา จิตฺตสมุฏฺฐานา, สิยา โน จิตฺตสมุฏฺฐานา, สตฺตรส ธาตุโย โน จิตฺตสหภุโน.\csromanspacer{} ธมฺมธาตุ สิยา จิตฺตสหภู, สิยา โน จิตฺตสหภู.\csromanspacer{} สตฺตรส ธาตุโย โน จิตฺตานุปริวตฺติโน.\csromanspacer{} ธมฺมธาตุ สิยา จิตฺตานุปริวตฺตี, สิยา โน จิตฺตานุปริวตฺตี.\csromanspacer{} สตฺตรส ธาตุโย โน จิตฺตสํสฏฺฐสมุฏฺฐานา.\csromanspacer{} ธมฺมธาตุ สิยา จิตฺตสํสฏฺฐสมุฏฺฐานา, สิยา โน จิตฺตสํสฏฺฐสมุฏฺฐานา.\csromanspacer{} สตฺตรส ธาตุโย โน จิตฺตสํสฏฺฐสมุฏฺฐานสหภุโน.\csromanspacer{} ธมฺมธาตุ สิยา จิตฺตสํสฏฺฐสมุฏฺฐานสหภู, สิยา โน จิตฺตสํสฏฺฐสมุฏฺฐานสหภู.\csromanspacer{} สตฺตรส ธาตุโย โน จิตฺตสํสฏฺฐสมุฏฺฐานานุปริวตฺติโน.\csromanspacer{} ธมฺมธาตุ สิยา จิตฺตสํสฏฺฐสมุฏฺฐานานุปริวตฺตี, สิยา โน จิตฺตสํสฏฺฐสมุฏฺฐานานุปริวตฺตี.\csromanspacer{} ทฺวาทส ธาตุโย อชฺฌตฺติกา.\csromanspacer{} ฉ ธาตุโย พาหิรา. (10)}

\prose{นว ธาตุโย อุปาทา.\csromanspacer{} อฏฺฐ ธาตุโย โน อุปาทา.\csromanspacer{} ธมฺมธาตุ สิยา อุปาทา, สิยา โน อุปาทา.\csromanspacer{} ทส ธาตุโย อุปาทินฺนา.\csromanspacer{} สทฺทธาตุ อนุปาทินฺนา.\csromanspacer{} สตฺต ธาตุโย สิยา อุปาทินฺนา, สิยา อนุปาทินฺนา.\csromanspacer{} สตฺตรส ธาตุโย โน อุปาทานา.\csromanspacer{} ธมฺมธาตุ สิยา อุปาทานํ, สิยา โน อุปาทานํ.\csromanspacer{} โสฬส ธาตุโย อุปาทานิยา.\csromanspacer{} เทฺว ธาตุโย สิยา อุปาทานิยา, สิยา อนุปาทานิยา.\csromanspacer{} โสฬส ธาตุโย อุปาทานวิปฺปยุตฺตา.\csromanspacer{} เทฺว ธาตุโย สิยา อุปาทานสมฺปยุตฺตา, สิยา อุปาทานวิปฺปยุตฺตา.\csromanspacer{} โสฬส ธาตุโย น วตฺตพฺพา “อุปาทานา เจว อุปาทานิยา จา”ติ, อุปาทานิยา เจว โน จ อุปาทานา.\csromanspacer{} มโนวิญฺญาณธาตุ น วตฺตพฺพา “อุปาทานญฺเจว อุปาทานิยา จา”ติ, สิยา อุปาทานิยา เจว โน จ อุปาทานํ, สิยา น วตฺตพฺพา “อุปาทานิยา เจว โน จ อุปาทานนฺ”ติ.\csromanspacer{} ธมฺมธาตุ สิยา อุปาทานญฺเจว อุปาทานิยา จ, สิยา อุปาทานิยา เจว โน จ อุปาทานํ, สิยา น วตฺตพฺพา “อุปาทานญฺเจว อุปาทานิยา จา”ติปิ, “อุปาทานิยา เจว โน จ อุปาทานนฺ”ติปิ. โสฬส ธาตุโย น วตฺตพฺพา “อุปาทานา เจว อุปาทานสมฺปยุตฺตา จา”ติปิ, “อุปาทานสมฺปยุตฺตา เจว โน จ อุปาทานา”ติปิ.\csromanspacer{} มโนวิญฺญาณธาตุ น วตฺตพฺพา “อุปาทานญฺเจว อุปาทานสมฺปยุตฺตา จา”ติ,}

\vspace{\tipitakaparskip}
\csromancenter{ธาตุวิภงฺค สิยา อุปาทานสมฺปยุตฺตา เจว โน จ อุปาทานํ, สิยา น วตฺตพฺพา “อุปาทานสมฺปยุตฺตา เจว โน จ อุปาทานนฺ”ติ.\csromanspacer{} ธมฺมธาตุ สิยา อุปาทานญฺเจว อุปาทานสมฺปยุตฺตา จ, สิยา อุปาทานสมฺปยุตฺตา เจว โน จ อุปาทานํ, สิยา น วตฺตพฺพา “อุปาทานญฺเจว อุปาทานสมฺปยุตฺตา จา”ติปิ, “อุปาทานสมฺปยุตฺตา เจว โน จ อุปาทานนฺ”ติปิ.\csromanspacer{} โสฬส ธาตุโย อุปาทานวิปฺปยุตฺต-อุปาทานิยา.\csromanspacer{} เทฺว ธาตุโย สิยา อุปาทานวิปฺปยุตฺต- อุปาทานิยา, สิยา อุปาทานวิปฺปยุตฺต-อนุปาทานิยา, สิยา น วตฺตพฺพา “อุปาทานวิปฺปยุตฺต-อุปาทานิยา”ติปิ, “อุปาทานวิปฺปยุตฺต-อนุปาทานิยา”ติปิ. (11)}
\prose{\csromanfolio{101}สตฺตรส ธาตุโย โน กิเลสา.\csromanspacer{} ธมฺมธาตุ สิยา กิเลสา, สิยา โน กิเลสา.\csromanspacer{} โสฬส ธาตุโย สํกิเลสิกา.\csromanspacer{} เทฺว ธาตุโย สิยา สํกิเลสิกา, สิยา อสํกิเลสิกา.\csromanspacer{} โสฬส ธาตุโย อสํกิลิฏฺฐา.\csromanspacer{} เทฺว ธาตุโย สิยา สํกิลิฏฺฐา, สิยา อสํกิลิฏฺฐา.\csromanspacer{} โสฬส ธาตุโย กิเลสวิปฺปยุตฺตา.\csromanspacer{} เทฺว ธาตุโย สิยา กิเลสสมฺปยุตฺตา, สิยา กิเลสวิปฺปยุตฺตา.\csromanspacer{} โสฬส ธาตุโย น วตฺตพฺพา “กิเลสา เจว สํกิเลสิกา จา”ติ, สํกิเลสิกา เจว โน จ กิเลสา.\csromanspacer{} มโนวิญฺญาณธาตุ น วตฺตพฺพา “กิเลโส เจว สํกิเลสิกา จา”ติ, สิยา สํกิเลสิกา เจว โน จ กิเลโส, สิยา น วตฺตพฺพา “สํกิเลสิกา เจว โน จ กิเลโส”ติ.\csromanspacer{} ธมฺมธาตุ สิยา กิเลโส เจว สํกิเลสิกา จ, สิยา สํกิเลสิกา เจว โน จ กิเลโส, สิยา น วตฺตพฺพา “กิเลโส เจว สํกิเลสิกา จา”ติปิ, “สํกิเลสิกา เจว โน จ กิเลโส”ติปิ.}

\prose{โสฬส ธาตุโย น วตฺตพฺพา “กิเลสา เจว สํกิลิฏฺฐา จา”ติปิ, “สํกิลิฏฺฐา เจว โน จ กิเลสา”ติปิ.\csromanspacer{} มโนวิญฺญาณธาตุ น วตฺตพฺพา “กิเลโส เจว สํกิลิฏฺฐา จา”ติ, สิยา สํกิลิฏฺฐา เจว โน จ กิเลโส, สิยา น วตฺตพฺพา “สํกิลิฏฺฐา เจว โน จ กิเลโส”ติ.\csromanspacer{} ธมฺมธาตุ สิยา กิเลโส เจว สํกิลิฏฺฐา จ, สิยา สํกิลิฏฺฐา เจว โน จ กิเลโส, สิยา น วตฺตพฺพา “กิเลโส เจว สํกิลิฏฺฐา จา”ติปิ, “สํกิลิฏฺฐา เจว โน จ กิเลโส”ติปิ.\csromanspacer{} โสฬส ธาตุโย น วตฺตพฺพา “กิเลโส เจว กิเลสสมฺปยุตฺตา จา”ติปิ, “กิเลสสมฺปยุตฺตา เจว โน จ กิเลสา”ติปิ.\csromanspacer{} มโนวิญฺญาณธาตุ \csromanfolio{102}น วตฺตพฺพา “กิเลโส เจว กิเลสสมฺปยุตฺตา จา”ติ, สิยา กิเลสสมฺปยุตฺตา เจว โน จ กิเลโส, สิยา น วตฺตพฺพา “กิเลสสมฺปยุตฺตา เจว โน จ กิเลโส”ติ.\csromanspacer{} ธมฺมธาตุ สิยา กิเลโส เจว กิเลสสมฺปยุตฺตา จ, สิยา กิเลสสมฺปยุตฺตา เจว โน จ กิเลโส, สิยา น วตฺตพฺพา “กิเลโส เจว กิเลสสมฺปยุตฺตา จา”ติปิ, “กิเลสสมฺปยุตฺตา เจว โน จ กิเลโส”ติปิ.\csromanspacer{} โสฬส ธาตุโย กิเลสวิปฺปยุตฺตสํกิเลสิกา.\csromanspacer{} เทฺว ธาตุโย สิยา กิเลสวิปฺปยุตฺตสํกิเลสิกา, สิยา กิเลสวิปฺปยุตฺต-อสํกิเลสิกา, สิยา น วตฺตพฺพา “กิเลสวิปฺปยุตฺตสํกิเลสิกา”ติปิ, “กิเลสวิปฺปยุตฺต-อสํกิเลสิกา”ติปิ. (12)}

\prose{โสฬส ธาตุโย น ทสฺสเนน ปหาตพฺพา.\csromanspacer{} เทฺว ธาตุโย สิยา ทสฺสเนน ปหาตพฺพา, สิยา น ทสฺสเนน ปหาตพฺพา.\csromanspacer{} โสฬส ธาตุโย น ภาวนาย ปหาตพฺพา.\csromanspacer{} เทฺว ธาตุโย สิยา ภาวนาย ปหาตพฺพา, สิยา น ภาวนาย ปหาตพฺพา.\csromanspacer{} โสฬส ธาตุโย น ทสฺสเนน ปหาตพฺพเหตุกา.\csromanspacer{} เทฺว ธาตุโย สิยา ทสฺสเนน ปหาตพฺพเหตุกา, สิยา น ทสฺสเนน ปหาตพฺพเหตุกา.\csromanspacer{} โสฬส ธาตุโย น ภาวนาย ปหาตพฺพเหตุกา.\csromanspacer{} เทฺว ธาตุโย สิยา ภาวนาย ปหาตพฺพเหตุกา, สิยา น ภาวนาย ปหาตพฺพเหตุกา.}

\prose{ปนฺนรส ธาตุโย อวิตกฺกา.\csromanspacer{} มโนธาตุ สวิตกฺกา.\csromanspacer{} เทฺว ธาตุโย สิยา สวิตกฺกา, สิยา อวิตกฺกา.\csromanspacer{} ปนฺนรส ธาตุโย อวิจารา.\csromanspacer{} มโนธาตุ สวิจารา.\csromanspacer{} เทฺว ธาตุโย สิยา สวิจารา, สิยา อวิจารา.\csromanspacer{} โสฬส ธาตุโย อปฺปีติกา, เทฺว ธาตุโย สิยา สปฺปีติกา, สิยา อปฺปีติกา.\csromanspacer{} โสฬส ธาตุโย น ปีติสหคตา.\csromanspacer{} เทฺว ธาตุโย สิยา ปีติสหคตา, สิยา น ปีติสหคตา.\csromanspacer{} ปนฺนรส ธาตุโย น สุขสหคตา.\csromanspacer{} ติสฺโส ธาตุโย สิยา สุขสหคตา, สิยา น สุขสหคตา.\csromanspacer{} เอกาทส ธาตุโย น อุเปกฺขาสหคตา.\csromanspacer{} ปญฺจ ธาตุโย อุเปกฺขาสหคตา.\csromanspacer{} เทฺว ธาตุโย สิยา อุเปกฺขาสหคตา, สิยา น อุเปกฺขาสหคตา.}

\prose{โสฬส ธาตุโย กามาวจรา.\csromanspacer{} เทฺว ธาตุโย สิยา กามาวจรา, สิยา น กามาวจรา.\csromanspacer{} โสฬส ธาตุโย น รูปาวจรา.\csromanspacer{} เทฺว ธาตุโย สิยา รูปาวจรา, สิยา น รูปาวจรา.\csromanspacer{} โสฬส ธาตุโย น อรูปาวจรา.}

\vspace{\tipitakaparskip}
\csromancenter{ธาตุวิภงฺค เทฺว ธาตุโย สิยา อรูปาวจรา, สิยา น อรูปาวจรา.\csromanspacer{} โสฬส ธาตุโย ปริยาปนฺนา.\csromanspacer{} เทฺว ธาตุโย สิยา ปริยาปนฺนา.\csromanspacer{} สิยา อปริยาปนฺนา.\csromanspacer{} โสฬส ธาตุโย อนิยฺยานิกา.\csromanspacer{} เทฺว ธาตุโย สิยา นิยฺยานิกา, สิยา อนิยฺยานิกา.\csromanspacer{} โสฬส ธาตุโย อนิยตา.\csromanspacer{} เทฺว ธาตุโย สิยา นิยตา, สิยา อนิยตา.\csromanspacer{} โสฬส ธาตุโย ส- อุตฺตรา.\csromanspacer{} เทฺว ธาตุโย สิยา ส-อุตฺตรา, สิยา อนุตฺตรา.\csromanspacer{} โสฬส ธาตุโย อรณา.\csromanspacer{} เทฺว ธาตุโย สิยา สรณา, สิยา อรณาติ. (13)}
\csromancenter{ปญฺหาปุจฺฉกํ.}
\nitthitam{\textbf{ธาตุวิภงฺโค นิฏฺฐิโต.}}
\csromanmatikaanchor{mtk.29}
\csromantocmarknopage{section}{4. สจฺจวิภงฺค}{mtk.29}
\bookmark[dest={mtk.29},level=1]{4. สจฺจวิภงฺค}
\csromanheader{1}{4. สจฺจวิภงฺค}
\csromanheader{2}{1. สุตฺตนฺตภาชนีย}
\proseitem{189}{\csromanfolio{104}จตฺตาริ อริยสจฺจานิ ทุกฺขํ อริยสจฺจํ, ทุกฺขสมุทยํ\footnote{ทุกฺขสมุทโย (สฺยา)} อริยสจฺจํ, ทุกฺขนิโรธํ\footnote{ทุกฺขนิโรโธ (สฺยา)} อริยสจฺจํ, ทุกฺขนิโรธคามินี ปฏิปทา อริยสจฺจํ.}

\csromanmatikaanchor{mtk.30}
\csromanmatikaanchor{mtk.31}
\csromantocmarknopage{subsection}{1. สุตฺตนฺตภาชนีย}{mtk.30}
\bookmark[dest={mtk.30},level=2]{1. สุตฺตนฺตภาชนีย}
\csromantocmarkref{subsubsection}{1. ทุกฺขสจฺจ}{mtk.31}
\bookmark[dest={mtk.31},level=3]{1. ทุกฺขสจฺจ}
\csromanheader{3}{1. ทุกฺขสจฺจ}
\proseitem{190}{ตตฺถ กตมํ ทุกฺขํ อริยสจฺจํ, ชาติปิ ทุกฺขา, ชราปิ ทุกฺขา, มรณมฺปิ ทุกฺขํ, โสกปริเทวทุกฺขโทมนสฺสุปายาสาปิ ทุกฺขา, อปฺปิเยหิ สมฺปโยโค ทุกฺโข, ปิเยหิ วิปฺปโยโค ทุกฺโข, ยมฺปิจฺฉํ น ลภติ ตมฺปิ ทุกฺขํ, สํขิตฺเตน ปญฺจุปาทานกฺขนฺธา ทุกฺขา.}

\proseitem{191}{ตตฺถ กตมา ชาติ, ยา เตสํ เตสํ สตฺตานํ ตมฺหิ ตมฺหิ สตฺตนิกาเย ชาติ สญฺชาติ โอกฺกนฺติ อภินิพฺพตฺติ ขนฺธานํ ปาตุภาโว อายตนานํ ปฏิลาโภ.\csromanspacer{} อยํ วุจฺจติ ชาติ.}

\proseitem{192}{ตตฺถ กตมา ชรา, ยา เตสํ เตสํ สตฺตานํ ตมฺหิ ตมฺหิ สตฺตนิกาเย ชรา ชีรณตา ขณฺฑิจฺจํ ปาลิจฺจํ วลิตฺตจตา อายุโน สํหานิ อินฺทฺริยานํ ปริปาโก.\csromanspacer{} อยํ วุจฺจติ ชรา.}

\proseitem{193}{ตตฺถ กตมํ มรณํ, ยา เตสํ เตสํ สตฺตานํ ตมฺหา ตมฺหา สตฺตนิกายา จุติ จวนตา เภโท อนฺตรธานํ มจฺจุ มรณํ กาลกิริยา ขนฺธานํ เภโท กเฬวรสฺส นิกฺเขโป ชีวิตินฺทฺริยสฺสุปจฺเฉโท.\csromanspacer{} อิทํ วุจฺจติ มรณํ.}

\proseitem{194}{ตตฺถ กตโม โสโก, ญาติพฺยสเนน วา ผุฏฺฐสฺส โภคพฺยสเนน วา ผุฏฺฐสฺส โรคพฺยสเนน วา ผุฏฺฐสฺส สีลพฺยสเนน วา ผุฏฺฐสฺส ทิฏฺฐิพฺยสเนน วา ผุฏฺฐสฺส อญฺญตรญฺญตเรน พฺยสเนน สมนฺนาคตสฺส อญฺญตรญฺญตเรน ทุกฺขธมฺเมน ผุฏฺฐสฺส โสโก โสจนา โสจิตตฺตํ อนฺโตโสโก อนฺโตปริโสโก เจตโส ปริชฺฌายนา โทมนสฺสํ โสกสลฺลํ.\csromanspacer{} อยํ วุจฺจติ โสโก.}

\csromanheader{2}{4. สจฺจวิภงฺค}
\proseitem{195}{\csromanfolio{105}ตตฺถ กตโม ปริเทโว, ญาติพฺยสเนน วา ผุฏฺฐสฺส โภคพฺยสเนน วา ผุฏฺฐสฺส โรคพฺยสเนน วา ผุฏฺฐสฺส สีลพฺยสเนน วา ผุฏฺฐสฺส ทิฏฺฐิพฺยสเนน วา ผุฏฺฐสฺส อญฺญตรญฺญตเรน พฺยสเนน สมนฺนาคตสฺส อญฺญตรญฺญตเรน ทุกฺขธมฺเมน ผุตฺถสฺส อาเทโว ปริเทโว อาเทวนา ปริเทวนา อาเทวิตตฺตํ ปริเทวิตตฺตํ วาจา ปลาโป วิปฺปลาโป ลาลปฺโป ลาลปฺปนา ลาลปฺปิตตฺตํ\footnote{ลาลโป ลาลปนา ลาลปิตตฺตํ (สฺยา)}.\csromanspacer{} อยํ วุจฺจติ ปริเทโว.}

\proseitem{196}{ตตฺถ กตมํ ทุกฺขํ, ยํ กายิกํ อสาตํ กายิกํ ทุกฺขํ กายสมฺผสฺสชํ อสาตํ ทุกฺขํ เวทยิตํ กายสมฺผสฺสชา อสาตา ทุกฺขา เวทนา.\csromanspacer{} อิทํ วุจฺจติ ทุกฺขํ.}

\proseitem{197}{ตตฺถ กตมํ โทมนสฺสํ, ยํ เจตสิกํ อสาตํ เจตสิกํ ทุกฺขํ เจโตสมฺผสฺสชํ อสาตํ ทุกฺขํ เวทยิตํ เจโตสมฺผสฺสชา อสาตา ทุกฺขา เวทนา.\csromanspacer{} อิทํ วุจฺจติ โทมนสฺสํ.}

\proseitem{198}{ตตฺถ กตโม อุปายาโส, ญาติพฺยสเนน วา ผุฏฺฐสฺส โภคพฺยสเนน วา ผุฏฺฐสฺส โรคพฺยสเนน วา ผุฏฺฐสฺส สีลพฺยสเนน วา ผุฏฺฐสฺส ทิฏฺฐิพฺยสเนน วา ผุฏฺฐสฺส อญฺญตรญฺญตเรน พฺยสเนน สมนฺนาคตสฺส อญฺญตรญฺญตเรน ทุกฺขธมฺเมน ผุฏฺฐสฺส อายาโส อุปายาโส อายาสิตตฺตํ อุปายาสิตตฺตํ.\csromanspacer{} อยํ วุจฺจติ อุปายาโส.}

\proseitem{199}{ตตฺถ กตโม อปฺปิเยหิ สมฺปโยโค ทุกฺโข, อิธ ยสฺส เต โหนฺติ อนิฏฺฐา อกนฺตา อมนาปา รูปา สทฺทา คนฺธา รสา โผฏฺฐพฺพา, เย วา ปนสฺส เต โหนฺติ อนตฺถกามา อหิตกามา อผาสุกกามา อโยคกฺเขมกามา, ยา เตหิ สงฺคติ สมาคโม สโมธานํ มิสฺสีภาโว.\csromanspacer{} อยํ วุจฺจติ อปฺปิเยหิ สมฺปโยโค ทุกฺโข.}

\proseitem{200}{ตตฺถ กตโม ปิเยหิ วิปฺปโยโค ทุกฺโข, อิธ ยสฺส เต โหนฺติ อิฏฺฐา กนฺตา มนาปา รูปา สทฺทา คนฺธา รสา โผฏฺฐพฺพา, เย วา ปนสฺส เต โหนฺติ อตฺถกามา หิตกามา ผาสุกกามา โยคกฺเขมกามา มาตา วา ปิตา วา ภาตา วา ภคินี วา มิตฺตา วา อมจฺจา วา ญาตี วา \csromanfolio{106}สาโลหิตา วา, ยา เตหิ อสงฺคติ อสมาคโม อสโมธานํ อมิสฺสีภาโว.\csromanspacer{} อยํ วุจฺจติ ปิเยหิ วิปฺปโยโค ทุกฺโข.}

\proseitem{201}{ตตฺถ กตมํ ยมฺปิจฺฉํ น ลภติ ตมฺปิ ทุกฺขํ, ชาติธมฺมานํ สตฺตานํ เอวํ อิจฺฉา อุปฺปชฺชติ “อโห วต มยํ น ชาติธมฺมา อสฺสาม, น จ วต โน ชาติ อาคจฺเฉยฺยา”ติ, น โข ปเนตํ อิจฺฉาย ปตฺตพฺพํ, อิทมฺปิ ยมฺปิจฺฉํ น ลภติ ตมฺปิ ทุกฺขํ.}

\prose{ชราธมฺมานํ สตฺตานํ ฯเปฯ.\csromanspacer{} พฺยาธิธมฺมานํ สตฺตานํ ฯเปฯ.\csromanspacer{} มรณธมฺมานํ สตฺตานํ ฯเปฯ.\csromanspacer{} โสกปริเทวทุกฺขโทมนสฺสุปายาสธมฺมานํ สตฺตานํ เอวํ อิจฺฉา อุปฺปชฺชติ “อโห วต มยํ น โสกปริเทวทุกฺขโทมนสฺสุปายาสธมฺมา อสฺสาม, น จ วต โน โสกปริเทวทุกฺขโทมนสฺสุปายาสา อาคจฺเฉยฺยุนฺ”ติ, น โข ปเนตํ อิจฺฉาย ปตฺตพฺพํ, อิทมฺปิ ยมฺปิจฺฉํ น ลภติ ตมฺปิ ทุกฺขํ.}

\proseitem{202}{ตตฺถ กตเม สํขิตฺเตน ปญฺจุปาทานกฺขนฺธา ทุกฺขา.\csromanspacer{} เสยฺยถิทํ, รูปุปาทานกฺขนฺโธ เวทนุปาทานกฺขนฺโธ สญฺญุปาทานกฺขนฺโธ สงฺขารุปาทานกฺขนฺโธ วิญฺญาณุปาทานกฺขนฺโธ, อิเม วุจฺจนฺติ สํขิตฺเตน ปญฺจุปาทานกฺขนฺธา ทุกฺขา.}

\csromanheader{2}{อิทํ วุจฺจติ ทุกฺขํ อริยสจฺจํ.}
\csromansectionrule
\csromanmatikaanchor{mtk.32}
\csromantocmarkref{subsubsection}{2. สมุทยสจฺจ}{mtk.32}
\bookmark[dest={mtk.32},level=3]{2. สมุทยสจฺจ}
\csromanheader{3}{2. สมุทยสจฺจ}
\proseitem{203}{ตตฺถ กตมํ ทุกฺขสมุทยํ อริยสจฺจํ, ยายํ ตณฺหา โปโนภวิกา\footnote{โปโนพฺภวิกา (สฺยา, ก)} นนฺทิราคสหคตา ตตฺรตตฺราภินนฺทินี.\csromanspacer{} เสยฺยถิทํ, กามตณฺหา ภวตณฺหา วิภวตณฺหา.}

\prose{สา โข ปเนสา ตณฺหา กตฺถ อุปฺปชฺชมานา อุปฺปชฺชติ, กตฺถ นิวิสมานา นิวิสติ, ยํ โลเก ปิยรูปํ สาตรูปํ, เอตฺเถสา ตณฺหา อุปฺปชฺชมานา อุปฺปชฺชติ, เอตฺถ นิวิสมานา นิวิสติ.}

\prose{กิญฺจ โลเก ปิยรูปํ สาตรูปํ, จกฺขุํ โลเก ปิยรูปํ สาตรูปํ, เอตฺเถสา ตณฺหา อุปฺปชฺชมานา อุปฺปชฺชติ, เอตฺถ นิวิสมานา นิวิสติ.\csromanspacer{} โสตํ}

\vspace{\tipitakaparskip}
\csromancenter{สจฺจวิภงฺค โลเก ฯเปฯ.\csromanspacer{} ฆานํ โลเก.\csromanspacer{} ชิวฺหา โลเก.\csromanspacer{} กาโย โลเก.\csromanspacer{} มโน โลเก ปิยรูปํ สาตรูปํ, เอตฺเถสา ตณฺหา อุปฺปชฺชมานา อุปฺปชฺชติ, เอตฺถ นิวิสมานา นิวิสติ.}
\prose{\csromanfolio{107}รูปา โลเก ปิยรูปํ สาตรูปํ, เอตฺเถสา ตณฺหา อุปฺปชฺชมานา อุปฺปชฺชติ, เอตฺถ นิวิสมานา นิวิสติ.\csromanspacer{} สทฺทา โลเก ฯเปฯ.\csromanspacer{} คนฺธา โลเก.\csromanspacer{} รสา โลเก.\csromanspacer{} โผฏฺฐพฺพา โลเก.\csromanspacer{} ธมฺมา โลเก ปิยรูปํ สาตรูปํ, เอตฺเถสา ตณฺหา อุปฺปชฺชมานา อุปฺปชฺชติ, เอตฺถ นิวิสมานา นิวิสติ.}

\prose{จกฺขุวิญฺญาณํ โลเก ปิยรูปํ สาตรูปํ, เอตฺเถสา ตณฺหา อุปฺปชฺชมานา อุปฺปชฺชติ, เอตฺถ นิวิสมานา นิวิสติ.\csromanspacer{} โสตวิญฺญาณํ โลเก ฯเปฯ.\csromanspacer{} ฆานวิญฺญาณํ โลเก.\csromanspacer{} ชิวฺหาวิญฺญาณํ โลเก.\csromanspacer{} กายวิญฺญาณํ โลเก.\csromanspacer{} มโนวิญฺญาณํ โลเก ปิยรูปํ สาตรูปํ, เอตฺเถสา ตณฺหา อุปฺปชฺชมานา อุปฺปชฺชติ, เอตฺถ นิวิสมานา นิวิสติ.}

\prose{จกฺขุสมฺผสฺโส โลเก ปิยรูปํ สาตรูปํ, เอตฺเถสา ตณฺหา อุปฺปชฺชมานา อุปฺปชฺชติ, เอตฺถ นิวิสมานา นิวิสติ.\csromanspacer{} โสตสมฺผสฺโส โลเก ฯเปฯ.\csromanspacer{} ฆานสมฺผสฺโส โลเก.\csromanspacer{} ชิวฺหาสมฺผสฺโส โลเก.\csromanspacer{} กายสมฺผสฺโส โลเก.\csromanspacer{} มโนสมฺผสฺโส โลเก ปิยรูปํ สาตรูปํ, เอตฺเถสา ตณฺหา อุปฺปชฺชมานา อุปฺปชฺชติ, เอตฺถ นิวิสมานา นิวิสติ.}

\prose{จกฺขุสมฺผสฺสชา เวทนา โลเก ปิยรูปํ สาตรูปํ, เอตฺเถสา ตณฺหา อุปฺปชฺชมานา อุปฺปชฺชติ, เอตฺถ นิวิสมานา นิวิสติ.\csromanspacer{} โสตสมฺผสฺสชา เวทนา โลเก ฯเปฯ.\csromanspacer{} ฆานสมฺผสฺสชา เวทนา โลเก.\csromanspacer{} ชิวฺหาสมฺผสฺสชา เวทนา โลเก.\csromanspacer{} กายสมฺผสฺสชา เวทนา โลเก.\csromanspacer{} มโนสมฺผสฺสชา เวทนา โลเก ปิยรูปํ สาตรูปํ, เอตฺเถสา ตณฺหา อุปฺปชฺชมานา อุปฺปชฺชติ, เอตฺถ นิวิสมานา นิวิสติ.}

\prose{รูปสญฺญา โลเก ปิยรูปํ สาตรูปํ, เอตฺเถสา ตณฺหา อุปฺปชฺชมานา อุปฺปชฺชติ, เอตฺถ นิวิสมานา นิวิสติ.\csromanspacer{} สทฺทสญฺญา โลเก ฯเปฯ.\csromanspacer{} คนฺธสญฺญา โลเก.\csromanspacer{} รสสญฺญา โลเก.\csromanspacer{} โผฏฺฐพฺพสญฺญา โลเก.\csromanspacer{} ธมฺมสญฺญา โลเก ปิยรูปํ สาตรูปํ, เอตฺเถสา ตณฺหา อุปฺปชฺชมานา อุปฺปชฺชติ, เอตฺถ นิวิสมานา นิวิสติ.}

\prose{รูปสญฺเจตนา โลเก ปิยรูปํ สาตรูปํ, เอตฺเถสา ตณฺหา อุปฺปชฺชมานา อุปฺปชฺชติ, เอตฺถ นิวิสมานา นิวิสติ.\csromanspacer{} สทฺทสญฺเจตนา โลเก ฯเปฯ. \csromanfolio{108}คนฺธสญฺเจตนา โลเก.\csromanspacer{} รสสญฺเจตนา โลเก.\csromanspacer{} โผฏฺฐพฺพสญฺเจตนา โลเก.\csromanspacer{} ธมฺมสญฺเจตนา โลเก ปิยรูปํ สาตรูปํ, เอตฺเถสา ตณฺหา อุปฺปชฺชมานา อุปฺปชฺชติ, เอตฺถ นิวิสมานา นิวิสติ.}

\prose{รูปตณฺหา โลเก ปิยรูปํ สาตรูปํ, เอตฺเถสา ตณฺหา อุปฺปชฺชมานา อุปฺปชฺชติ, เอตฺถ นิวิสมานา นิวิสติ.\csromanspacer{} สทฺทตณฺหา โลเก ฯเปฯ.\csromanspacer{} คนฺธตณฺหา โลเก.\csromanspacer{} รสตณฺหา โลเก.\csromanspacer{} โผฏฺฐพฺพตณฺหา โลเก.\csromanspacer{} ธมฺมตณฺหา โลเก ปิยรูปํ สาตรูปํ, เอตฺเถสา ตณฺหา อุปฺปชฺชมานา อุปฺปชฺชติ, เอตฺถ นิวิสมานา นิวิสติ.}

\prose{รูปวิตกฺโก โลเก ปิยรูปํ สาตรูปํ, เอตฺเถสา ตณฺหา อุปฺปชฺชมานา อุปฺปชฺชติ, เอตฺถ นิวิสมานา นิวิสติ.\csromanspacer{} สทฺทวิตกฺโก โลเก ฯเปฯ.\csromanspacer{} คนฺธวิตกฺโก โลเก.\csromanspacer{} รสวิตกฺโก โลเก.\csromanspacer{} โผฏฺฐพฺพวิตกฺโก โลเก.\csromanspacer{} ธมฺมวิตกฺโก โลเก ปิยรูปํ สาตรูปํ, เอตฺเถสา ตณฺหา อุปฺปชฺชมานา อุปฺปชฺชติ, เอตฺถ นิวิสมานา นิวิสติ.}

\prose{รูปวิจาโร โลเก ปิยรูปํ สาตรูปํ, เอตฺเถสา ตณฺหา อุปฺปชฺชมานา อุปฺปชฺชติ, เอตฺถ นิวิสมานา นิวิสติ.\csromanspacer{} สทฺทวิจาโร โลเก ปิยรูปํ สาตรูปํ, เอตฺเถสา ตณฺหา อุปฺปชฺชมานา อุปฺปชฺชติ, เอตฺถ นิวิสมานา นิวิสติ.\csromanspacer{} คนฺธวิจาโร โลเก ปิยรูปํ สาตรูปํ, เอตฺเถสา ตณฺหา อุปฺปชฺชมานา อุปฺปชฺชติ, เอตฺถ นิวิสมานา นิวิสติ.\csromanspacer{} รสวิจาโร โลเก ปิยรูปํ สาตรูปํ, เอตฺเถสา ตณฺหา อุปฺปชฺชมานา อุปฺปชฺชติ, เอตฺถ นิวิสมานา นิวิสติ.\csromanspacer{} โผฏฺฐพฺพวิจาโร โลเก ปิยรูปํ สาตรูปํ, เอตฺเถสา ตณฺหา อุปฺปชฺชมานา อุปฺปชฺชติ, เอตฺถ นิวิสมานา นิวิสติ.\csromanspacer{} ธมฺมวิจาโร โลเก ปิยรูปํ สาตรูปํ, เอตฺเถสา ตณฺหา อุปฺปชฺชมานา อุปฺปชฺชติ, เอตฺถ นิวิสมานา นิวิสติ.}

\csromanheader{2}{อิทํ วุจฺจติ ทุกฺขสมุทยํ อริยสจฺจํ.}
\csromansectionrule
\csromanmatikaanchor{mtk.33}
\csromantocmarkref{subsubsection}{3. นิโรธสจฺจ}{mtk.33}
\bookmark[dest={mtk.33},level=3]{3. นิโรธสจฺจ}
\csromanheader{3}{3. นิโรธสจฺจ}
\proseitem{204}{ตตฺถ กตมํ ทุกฺขนิโรธํ อริยสจฺจํ, โย ตสฺสาเยว ตณฺหาย อเสสวิราคนิโรโธ จาโค ปฏินิสฺสคฺโค มุตฺติ อนาลโย.}

\csromanheader{2}{4. สจฺจวิภงฺค}
\prose{\csromanfolio{109}สา โข ปเนสา ตณฺหา กตฺถ ปหียมานา ปหียติ, กตฺถ นิรุชฺฌมานา นิรุชฺฌติ.\csromanspacer{} ยํ โลเก ปิยรูปํ สาตรูปํ, เอตฺเถสา ตณฺหา ปหียมานา ปหียติ, เอตฺถ นิรุชฺฌมานา นิรุชฺฌติ.}

\prose{กิญฺจ โลเก ปิยรูปํ สาตรูปํ, จกฺขุํ โลเก ปิยรูปํ สาตรูปํ, เอตฺเถสา ตณฺหา ปหียมานา ปหียติ, เอตฺถ นิรุชฺฌมานา นิรุชฺฌติ.\csromanspacer{} โสตํ โลเก ฯเปฯ.\csromanspacer{} ฆานํ โลเก.\csromanspacer{} ชิวฺหา โลเก.\csromanspacer{} กาโย โลเก.\csromanspacer{} มโน โลเก ปิยรูปํ สาตรูปํ, เอตฺเถสา ตณฺหา ปหียมานา ปหียติ, เอตฺถ นิรุชฺฌมานา นิรุชฺฌติ.}

\prose{รูปา โลเก ปิยรูปํ สาตรูปํ, เอตฺเถสา ตณฺหา ปหียมานา ปหียติ, เอตฺถ นิรุชฺฌมานา นิรุชฺฌติ.\csromanspacer{} สทฺทา โลเก ฯเปฯ.\csromanspacer{} คนฺธา โลเก.\csromanspacer{} รสา โลเก.\csromanspacer{} โผฏฺฐพฺพา โลเก.\csromanspacer{} ธมฺมา โลเก ปิยรูปํ สาตรูปํ, เอตฺเถสา ตณฺหา ปหียมานา ปหียติ, เอตฺถ นิรุชฺฌมานา นิรุชฺฌติ.}

\prose{จกฺขุวิญฺญาณํ โลเก ปิยรูปํ สาตรูปํ, เอตฺเถสา ตณฺหา ปหียมานา ปหียติ, เอตฺถ นิรุชฺฌมานา นิรุชฺฌติ.\csromanspacer{} โสตวิญฺญาณํ โลเก ฯเปฯ.\csromanspacer{} ฆานวิญฺญาณํ โลเก.\csromanspacer{} ชิวฺหาวิญฺญาณํ โลเก.\csromanspacer{} กายวิญฺญาณํ โลเก.\csromanspacer{} มโนวิญฺญาณํ โลเก ปิยรูปํ สาตรูปํ, เอตฺเถสา ตณฺหา ปหียมานา ปหียติ, เอตฺถ นิรุชฺฌมานา นิรุชฺฌติ.}

\prose{จกฺขุสมฺผสฺโส โลเก ปิยรูปํ สาตรูปํ, เอตฺเถสา ตณฺหา ปหียมานา ปหียติ, เอตฺถ นิรุชฺฌมานา นิรุชฺฌติ.\csromanspacer{} โสตสมฺผสฺโส โลเก ฯเปฯ.\csromanspacer{} ฆานสมฺผสฺโส โลเก ฯเปฯ.\csromanspacer{} ชิวฺหาสมฺผสฺโส โลเก.\csromanspacer{} กายสมฺผสฺโส โลเก.\csromanspacer{} มโนสมฺผสฺโส โลเก ปิยรูปํ สาตรูปํ, เอตฺเถสา ตณฺหา ปหียมานา ปหียติ, เอตฺถ นิรุชฺฌมานา นิรุชฺฌติ.}

\prose{จกฺขุสมฺผสฺสชา เวทนา โลเก ปิยรูปํ สาตรูปํ, เอตฺเถสา ตณฺหา ปหียมานา ปหียติ, เอตฺถ นิรุชฺฌมานา นิรุชฺฌติ.\csromanspacer{} โสตสมฺผสฺสชา เวทนา โลเก ฯเปฯ.\csromanspacer{} ฆานสมฺผสฺสชา เวทนา โลเก.\csromanspacer{} ชิวฺหาสมฺผสฺสชา เวทนา โลเก.\csromanspacer{} กายสมฺผสฺสชา เวทนา โลเก.\csromanspacer{} มโนสมฺผสฺสชา เวทนา โลเก ปิยรูปํ สาตรูปํ, เอตฺเถสา ตณฺหา ปหียมานา ปหียติ, เอตฺถ นิรุชฺฌมานา นิรุชฺฌติ.}

\prose{รูปสญฺญา โลเก.\csromanspacer{} สทฺทสญฺญา โลเก.\csromanspacer{} คนฺธสญฺญา โลเก.\csromanspacer{} รสสญฺญา โลเก.\csromanspacer{} โผฏฺฐพฺพสญฺญา โลเก.\csromanspacer{} ธมฺมสญฺญา โลเก}

\prose{\csromanfolio{110}ปิยรูปํ สาตรูปํ, เอตฺเถสา ตณฺหา ปหียมานา ปหียติ, เอตฺถ นิรุชฺฌมานา นิรุชฺฌติ.}

\prose{รูปสญฺเจตนา โลเก.\csromanspacer{} สทฺทสญฺเจตนา โลเก.\csromanspacer{} คนฺธสญฺเจตนา โลเก.\csromanspacer{} รสสญฺเจตนา โลเก.\csromanspacer{} โผฏฺฐพฺพสญฺเจตนา โลเก.\csromanspacer{} ธมฺมสญฺเจตนา โลเก ปิยรูปํ สาตรูปํ, เอตฺเถสา ตณฺหา ปหียมานา ปหียติ, เอตฺถ นิรุชฺฌมานา นิรุชฺฌติ.}

\prose{รูปตณฺหา โลเก.\csromanspacer{} สทฺทตณฺหา โลเก.\csromanspacer{} คนฺธตณฺหา โลเก.\csromanspacer{} รสตณฺหา โลเก.\csromanspacer{} โผฏฺฐพฺพตณฺหา โลเก.\csromanspacer{} ธมฺมตณฺหา โลเก ปิยรูปํ สาตรูปํ, เอตฺเถสา ตณฺหา ปหียมานา ปหียติ, เอตฺถ นิรุชฺฌมานา นิรุชฺฌติ.}

\prose{รูปวิตกฺโก โลเก.\csromanspacer{} สทฺทวิตกฺโก โลเก.\csromanspacer{} คนฺธวิตกฺโก โลเก.\csromanspacer{} รสวิตกฺโก โลเก.\csromanspacer{} โผฏฺฐพฺพวิตกฺโก โลเก.\csromanspacer{} ธมฺมวิตกฺโก โลเก ปิยรูปํ สาตรูปํ, เอตฺเถสา ตณฺหา ปหียมานา ปหียติ, เอตฺถ นิรุชฺฌมานา นิรุชฺฌติ.}

\prose{รูปวิจาโร โลเก.\csromanspacer{} สทฺทวิจาโร โลเก.\csromanspacer{} คนฺธวิจาโร โลเก.\csromanspacer{} รสวิจาโร โลเก.\csromanspacer{} โผฏฺฐพฺพวิจาโร โลเก.\csromanspacer{} ธมฺมวิจาโร โลเก ปิยรูปํ สาตรูปํ, เอตฺเถสา ตณฺหา ปหียมานา ปหียติ, เอตฺถ นิรุชฺฌมานา นิรุชฺฌติ.}

\csromanheader{2}{อิทํ วุจฺจติ ทุกฺขนิโรธํ อริยสจฺจํ.}
\csromansectionrule
\csromanmatikaanchor{mtk.34}
\csromantocmarkref{subsubsection}{4. มคฺคสจฺจ}{mtk.34}
\bookmark[dest={mtk.34},level=3]{4. มคฺคสจฺจ}
\csromanheader{3}{4. มคฺคสจฺจ}
\proseitem{205}{ตตฺถ กตมํ ทุกฺขนิโรธคามินี ปฏิปทา อริยสจฺจํ, อยเมว อริโย อฏฺฐงฺคิโก มคฺโค.\csromanspacer{} เสยฺยถิทํ, สมฺมาทิฏฺฐิ สมฺมาสงฺกปฺโป สมฺมาวาจา สมฺมากมฺมนฺโต สมฺมา-อาชีโว สมฺมาวายาโม สมฺมาสติ สมฺมาสมาธิ.}

\prose{ตตฺถ กตมา สมฺมาทิฏฺฐิ, ทุกฺเข ญาณํ ทุกฺขสมุทเย ญาณํ ทุกฺขนิโรเธ ญาณํ ทุกฺขนิโรธคามินิยา ปฏิปทาย ญาณํ.\csromanspacer{} อยํ วุจฺจติ สมฺมาทิฏฺฐิ.}

\csromanheader{2}{4. สจฺจวิภงฺค}
\prose{\csromanfolio{111}ตตฺถ กตโม สมฺมาสงฺกปฺโป, เนกฺขมฺมสงฺกปฺโป อพฺยาปาทสงฺกปฺโป อวิหิํสาสงฺกปฺโป.\csromanspacer{} อยํ วุจฺจติ สมฺมาสงฺกปฺโป.}

\prose{ตตฺถ กตมา สมฺมาวาจา, มุสาวาทา เวรมณี ปิสุณาย วาจาย เวรมณี ผรุสาย วาจาย เวรมณี สมฺผปฺปลาปา เวรมณี.\csromanspacer{} อยํ วุจฺจติ สมฺมาวาจา.}

\prose{ตตฺถ กตโม สมฺมากมฺมนฺโต, ปาณาติปาตา เวรมณี อทินฺนาทานา เวรมณี กาเมสุมิจฺฉาจารา เวรมณี.\csromanspacer{} อยํ วุจฺจติ สมฺมากมฺมนฺโต.}

\prose{ตตฺถ กตโม สมฺมา-อาชีโว, อิธ อริยสาวโก มิจฺฉา-อาชีวํ ปหาย สมฺมา-อาชีเวน ชีวิกํ กปฺเปติ.\csromanspacer{} อยํ วุจฺจติ สมฺมา-อาชีโว.}

\prose{ตตฺถ กตโม สมฺมาวายาโม, อิธ ภิกฺขุ อนุปฺปนฺนานํ ปาปกานํ อกุสลานํ ธมฺมานํ อนุปฺปาทาย ฉนฺทํ ชเนติ วายมติ วีริยํ อารภติ จิตฺตํ ปคฺคณฺหาติ ปทหติ, อุปฺปนฺนานํ ปาปกานํ อกุสลานํ ธมฺมานํ ปหานาย ฉนฺทํ ชเนติ วายมติ วีริยํ อารภติ จิตฺตํ ปคฺคณฺหาติ ปทหติ, อนุปฺปนฺนานํ กุสลานํ ธมฺมานํ อุปฺปาทาย ฉนฺทํ ชเนติ วายมติ วิริย อารภติ จิตฺตํ ปคฺคณฺหาติ ปทหติ, อุปฺปนฺนานํ กุสลานํ ธมฺมานํ ฐิติยา อสมฺโมสาย ภิยฺโยภาวาย เวปุลฺลาย ภาวนาย ปาริปูริยา ฉนฺทํ ชเนติ วายมติ วีริยํ อารภติ จิตฺตํ ปคฺคณฺหาติ ปทหติ.\csromanspacer{} อยํ วุจฺจติ สมฺมาวายาโม.}

\prose{ตตฺถ กตมา สมฺมาสติ, อิธ ภิกฺขุ กาเย กายานุปสฺสี วิหรติ อาตาปี สมฺปชาโน สติมา วิเนยฺย โลเก อภิชฺฌาโทมนสฺสํ.\csromanspacer{} เวทนาสุ ฯเปฯ.\csromanspacer{} จิตฺเต ฯเปฯ.\csromanspacer{} ธมฺเมสุ ธมฺมานุปสฺสี วิหรติ อาตาปี สมฺปชาโน สติมา วิเนยฺย โลเก อภิชฺฌาโทมนสฺสํ.\csromanspacer{} อยํ วุจฺจติ สมฺมาสติ.}

\prose{ตตฺถ กตโม สมฺมาสมาธิ, อิธ ภิกฺขุ วิวิจฺเจย กาเมหิ วิวิจฺจ อกุสเลหิ ธมฺเมหิ สวิตกฺกํ สวิจารํ วิเวกชํ ปีติสุขํ ปฐมํ ฌานํ อุปสมฺปชฺช วิหรติ, วิตกฺกวิจารานํ วูปสมา อชฺฌตฺตํ สมฺปสาทนํ เจตโส เอโกทิภาวํ อวิตกฺกํ อวิจารํ สมาธิชํ ปีติสุขํ ทุติยํ ฌานํ อุปสมฺปชฺช วิหรติ, ปีติยา จ วิราคา อุเปกฺขโก จ วิหรติ สโต จ สมฺปชาโน สุขญฺจ กาเยน ปฏิสํเวเทติ, ยํ ตํ อริยา อาจิกฺขนฺติ “อุเปกฺขโก สติมา สุขวิหารี”ติ, ตติยํ ฌานํ อุปสมฺปชฺช วิหรติ,}

\prose{\csromanfolio{112}สุขสฺส จ ปหานา ทุกฺขสฺส จ ปหานา ปุพฺเพว โสมนสฺสโทมนสฺสานํ อตฺถงฺคมา อทุกฺขมสุขํ อุเปกฺขาสติปาริสุทฺธิํ จตุตฺถํ ฌานํ อุปสมฺปชฺช วิหรติ.\csromanspacer{} อยํ วุจฺจติ สมฺมาสมาธิ.}

\prose{อิทํ วุจฺจติ ทุกฺขนิโรธคามินี ปฏิปทา อริยสจฺจํ.}

\csromanheader{2}{สุตฺตนฺตภาชนียํ.}
\csromansectionrule
\csromanmatikaanchor{mtk.35}
\csromantocmarkref{subsection}{2. อภิธมฺมภาชนีย}{mtk.35}
\bookmark[dest={mtk.35},level=2]{2. อภิธมฺมภาชนีย}
\csromanheader{2}{2. อภิธมฺมภาชนีย}
\proseitem{206}{จตฺตาริ สจฺจานิ ทุกฺขํ ทุกฺขสมุทโย ทุกฺขนิโรโธ ทุกฺขนิโรธคามินี ปฏิปทา.}

\prose{ตตฺถ กตโม ทุกฺขสมุทโย, ตณฺหา.\csromanspacer{} อยํ วุจฺจติ ทุกฺขสมุทโย.}

\prose{ตตฺถ กตมํ ทุกฺขํ, อวเสสา จ กิเลสา อวเสสา จ อกุสลา ธมฺมา ตีณิ จ กุสลมูลานิ สาสวานิ อวเสสา จ สาสวา กุสลา ธมฺมา สาสวา จ กุสลากุสลานํ ธมฺมานํ วิปากา, เย จ ธมฺมา กิริยา เนว กุสลา นากุสลา น จ กมฺมวิปากา สพฺพญฺจ รูปํ.\csromanspacer{} อิทํ วุจฺจติ ทุกฺขํ.}

\prose{ตตฺถ กตโม ทุกฺขนิโรโธ, ตณฺหาย ปหานํ.\csromanspacer{} อยํ วุจฺจติ ทุกฺขนิโรโธ.}

\prose{ตตฺถ กตมา ทุกฺขนิโรธคามินี ปฏิปทา, อิธ ภิกฺขุ ยสฺมิํ สมเย โลกุตฺตรํ ฌานํ ภาเวติ นิยฺยานิกํ อปจยคามิํ ทิฏฺฐิคตานํ ปหานาย ปฐมาย ภูมิยา ปตฺติยา วิวิจฺเจว กาเมหิ ฯเปฯ ปฐมํ ฌานํ อุปสมฺปชฺช วิหรติ ทุกฺขปฏิปทํ ทนฺธาภิญฺญํ, ตสฺมิํ สมเย อฏฺฐงฺคิโก มคฺโค โหติ สมฺมาทิฏฺฐิ ฯเปฯ สมฺมาสมาธิ.}

\prose{ตตฺถ กตมา สมฺมาทิฏฺฐิ, ยา ปญฺญา ปชานนา ฯเปฯ อโมโห ธมฺมวิจโย สมฺมาทิฏฺฐิ ธมฺมวิจยสมฺโพชฺฌงฺโค มคฺคงฺคํ มคฺคปริยาปนฺนํ.\csromanspacer{} อยํ วุจฺจติ สมฺมาทิฏฺฐิ.}

\prose{ตตฺถ กตโม สมฺมาสงฺกปฺโป, โย ตกฺโก วิตกฺโก ฯเปฯ สมฺมาสงฺกปฺโป มคฺคงฺคํ มคฺคปริยาปนฺนํ.\csromanspacer{} อยํ วุจฺจติ สมฺมาสงฺกปฺโป.}

\csromanheader{2}{4. สจฺจวิภงฺค}
\prose{\csromanfolio{113}ตตฺถ กตมา สมฺมาวาจา, ยา จตูหิ วจีทุจฺจริเตหิ อารติ วิรติ ปฏิวิรติ เวรมณี\footnote{เวรมณิ (ก) เอวมุปริปิ.} อกิริยา อกรณํ อนชฺฌาปตฺติ เวลา-อนติกฺกโม เสตุฆาโต สมฺมาวาจา มคฺคงฺคํ มคฺคปริยาปนฺนํ.\csromanspacer{} อยํ วุจฺจติ สมฺมาวาจา.}

\prose{ตตฺถ กตโม สมฺมากมฺมนฺโต, ยา ตีหิ กายทุจฺจริเตหิ อารติ วิรติ ปฏิวิรติ เวรมณี อกิริยา อกรณํ อนชฺฌาปตฺติ เวลา-อนติกฺกโม เสตุฆาโต สมฺมากมฺมนฺโต มคฺคงฺคํ มคฺคปริยาปนฺนํ.\csromanspacer{} อยํ วุจฺจติ สมฺมากมฺมนฺโต.}

\prose{ตตฺถ กตโม สมฺมา-อาชีโว, ยา มิจฺฉา อาชีวา อารติ วิรติ ปฏิวิรติ เวรมณี อกิริยา อกรณํ อนชฺฌาปตฺติ เวลา-อนติกฺกโม เสตุฆาโต สมฺมา- อาชีโว มคฺคงฺคํ มคฺคปริยาปนฺนํ.\csromanspacer{} อยํ วุจฺจติ สมฺมา-อาชีโว.}

\prose{ตตฺถ กตโม สมฺมาวายาโม, โย เจตสิโก วีริยารมฺโภ\footnote{วิริยารมฺโภ (สี, สฺยา)} ฯเปฯ สมฺมาวายาโม วีริยสมฺโพชฺฌงฺโค มคฺคงฺคํ มคฺคปริยาปนฺนํ.\csromanspacer{} อยํ วุจฺจติ สมฺมาวายาโม.}

\prose{ตตฺถ กตมา สมฺมาสติ, ยา สติ อนุสฺสติ ฯเปฯ สมฺมาสติ สติสมฺโพชฺฌงฺโค มคฺคงฺคํ มคฺคปริยาปนฺนํ.\csromanspacer{} อยํ วุจฺจติ สมฺมาสติ.}

\prose{ตตฺถ กตโม สมฺมาสมาธิ, ยา จิตฺตสฺส ฐิติ สณฺฐิติ ฯเปฯ สมฺมาสมาธิ สมาธิสมฺโพชฺฌงฺโค มคฺคงฺคํ มคฺคปริยาปนฺนํ.\csromanspacer{} อยํ วุจฺจติ สมฺมาสมาธิ.\csromanspacer{} อยํ วุจฺจติ ทุกฺขนิโรธคามินี ปฏิปทา, อวเสสา ธมฺมา ทุกฺขนิโรธคามินิยา ปฏิปทาย สมฺปยุตฺตา.}

\proseitem{207}{ตตฺถ กตโม ทุกฺขสมุทโย, ตณฺหา จ อวเสสา จ กิเลสา.\csromanspacer{} อยํ วุจฺจติ ทุกฺขสมุทโย.}

\prose{ตตฺถ กตมํ ทุกฺขํ, อวเสสา จ อกุสลา ธมฺมา ตีณิ จ กุสลมูลานิ สาสวานิ อวเสสา จ สาสวา กุสลา ธมฺมา สาสวา จ กุสลากุสลานํ ธมฺมานํ วิปากา, เย จ ธมฺมา กิริยา เนว กุสลา นากุสลา น จ กมฺมวิปากา สพฺพญฺจ รูปํ.\csromanspacer{} อิทํ วุจฺจติ ทุกฺขํ.}

\prose{\csromanfolio{114}ตตฺถ กตโม ทุกฺขนิโรโธ, ตณฺหาย จ อวเสสานญฺจ กิเลสานํ ปหานํ.\csromanspacer{} อยํ วุจฺจติ ทุกฺขนิโรโธ.}

\prose{ตตฺถ กตมา ทุกฺขนิโรธคามินี ปฏิปทา, อิธ ภิกฺขุ ยสฺมิํ สมเย โลกุตฺตรํ ฌานํ ภาเวติ นิยฺยานิกํ อปจยคามิํ ทิฏฺฐิคตานํ ปหานาย ปฐมาย ภูมิยา ปตฺติยา วิวิจฺเจว กาเมหิ ฯเปฯ ปฐมํ ฌานํ อุปสมฺปชฺช วิหรติ ทุกฺขปฏิปทํ ทนฺธาภิญฺญํ, ตสฺมิํ สมเย อฏฺฐงฺคิโก มคฺโค โหติ สมฺมาทิฏฺฐิ ฯเปฯ สมฺมาสมาธิ.\csromanspacer{} อยํ วุจฺจติ ทุกฺขนิโรธคามินี ปฏิปทา, อวเสสา ธมฺมา ทุกฺขนิโรธคามินิยา ปฏิปทาย สมฺปยุตฺตา.}

\proseitem{208}{ตตฺถ กตโม ทุกฺขสมุทโย, ตณฺหา จ อวเสสา จ กิเลสา อวเสสา จ อกุสลา ธมฺมา.\csromanspacer{} อยํ วุจฺจติ ทุกฺขสมุทโย.}

\prose{ตตฺถ กตมํ ทุกฺขํ, ตีณิ จ กุสลมูลานิ สาสวานิ อวเสสา จ สาสวา กุสลา ธมฺมา สาสวา จ กุสลากุสลานํ ธมฺมานํ วิปากา เย จ ธมฺมา กิริยา เนว กุสลา นากุสลา น จ กมฺมวิปากา สพฺพญฺจ รูปํ.\csromanspacer{} อิทํ วุจฺจติ ทุกฺขํ.}

\prose{ตตฺถ กตโม ทุกฺขนิโรโธ, ตณฺหาย จ อวเสสานญฺจ กิเลสานํ อวเสสานญฺจ อกุสลานํ ธมฺมานํ ปหานํ.\csromanspacer{} อยํ วุจฺจติ ทุกฺขนิโรโธ.}

\prose{ตตฺถ กตมา ทุกฺขนิโรธคามินี ปฏิปทา, อิธ ภิกฺขุ ยสฺมิํ สมเย โลกุตฺตรํ ฌานํ ภาเวติ นิยฺยานิกํ อปจยคามิํ ทิฏฺฐิคตานํ ปหานาย ปฐมาย ภูมิยา ปตฺติยา วิวิจฺเจว กาเมหิ ฯเปฯ ปฐมํ ฌานํ อุปสมฺปชฺช วิหรติ ทุกฺขปฏิปทํ ทนฺธาภิญฺญํ, ตสฺมิํ สมเย อฏฺฐงฺคิโก มคฺโค โหติ, สมฺมาทิฏฺฐิ ฯเปฯ สมฺมาสมาธิ.\csromanspacer{} อยํ วุจฺจติ ทุกฺขนิโรทคามินี ปฏิปทา, อวเสสา ธมฺมา ทุกฺขนิโรธคามินิยา ปฏิปทาย สมฺปยุตฺตา.}

\proseitem{209}{ตตฺถ กตโม ทุกฺขสมุทโย, ตณฺหา จ อวเสสา จ กิเลสา อวเสสา จ อกุสลา ธมฺมา ตีณิ จ กุสลมูลานิ สาสวานิ.\csromanspacer{} อยํ วุจฺจติ ทุกฺขสมุทโย.}

\prose{ตตฺถ กตมํ ทุกฺขํ, อวเสสา จ สาสวา กุสลา ธมฺมา สาสวา จ กุสลากุสลานํ ธมฺมานํ วิปากา เย จ ธมฺมา กิริยา เนว กุสลา นากุสลา น จ กมฺมวิปากา สพฺพญฺจ รูปํ.\csromanspacer{} อิทํ วุจฺจติ ทุกฺขํ.}

\csromanheader{2}{4. สจฺจวิภงฺค}
\prose{\csromanfolio{115}ตตฺถ กตโม ทุกฺขนิโรโธ, ตณฺหาย จ อวเสสานญฺจ กิเลสานํ อวเสสานญฺจ อกุสลานํ ธมฺมานํ ติณฺณญฺจ กุสลมูลานํ สาสวานํ ปหานํ.\csromanspacer{} อยํ วุจฺจติ ทุกฺขนิโรโธ.}

\prose{ตตฺถ กตมา ทุกฺขนิโรธคามินี ปฏิปทา, อิธ ภิกฺขุ ยสฺมิํ สมเย โลกุตฺตรํ ฌานํ ภาเวติ นิยฺยานิกํ อปจยคามิํ ทิฏฺฐิคตานํ ปหานาย ปฐมาย ภูมิยา ปตฺติยา วิวิจฺเจว กาเมหิ ฯเปฯ ปฐมํ ฌานํ อุปสมฺปชฺช วิหรติ ทุกฺขปฏิปทํ ทนฺธาภิญฺญํ, ตสฺมิํ สมเย อฏฺฐงฺคิโก มคฺโค โหติ, สมฺมาทิฏฺฐิ ฯเปฯ สมฺมาสมาธิ.\csromanspacer{} อยํ วุจฺจติ ทุกฺขนิโรธคามินี ปฏิปทา, อวเสสา ธมฺมา ทุกฺขนิโรธคามินิยา ปฏิปทาย สมฺปยุตฺตา.}

\proseitem{210}{ตตฺถ กตโม ทุกฺขสมุทโย, ตณฺหา จ อวเสสา จ กิเลสา อวเสสา จ อกุสลา ธมฺมา ตีณิ จ กุสลมูลานิ สาสวานิ อวเสสา จ สาสวา กุสลา ธมฺมา.\csromanspacer{} อยํ วุจฺจติ ทุกฺขสมุทโย.}

\prose{ตตฺถ กตมํ ทุกฺขํ, สาสวา กุสลากุสลานํ ธมฺมานํ วิปากา เย จ ธมฺมา กิริยา เนว กุสลา นากุสลา น จ กมฺมวิปากา สพฺพญฺจ รูปํ.\csromanspacer{} อิทํ วุจฺจติ ทุกฺขํ.}

\prose{ตตฺถ กตโม ทุกฺขนิโรโธ, ตณฺหาย จ อวเสสานญฺจ กิเลสานํ อวเสสานญฺจ อกุสลานํ ธมฺมานํ ติณฺณญฺจ กุสลมูลานํ สาสวานํ อวเสสานญฺจ สาสวานํ กุสลานํ ธมฺมานํ ปหานํ.\csromanspacer{} อยํ วุจฺจติ ทุกฺขนิโรโธ.}

\prose{ตตฺถ กตมา ทุกฺขนิโรธคามินี ปฏิปทา, อิธ ภิกฺขุ ยสฺมิํ สมเย โลกุตฺตรํ ฌานํ ภาเวติ นิยฺยานิกํ อปจยคามิํ ทิฏฺฐิคตานํ ปหานาย ปฐมาย ภูมิยา ปตฺติยา วิวิจฺเจว กาเมหิ ฯเปฯ ปฐมํ ฌานํ อุปสมฺปชฺช วิหรติ ทุกฺขปฏิปทํ ทนฺธาภิญฺญํ, ตสฺมิํ สมเย อฏฺฐงฺคิโก มคฺโค โหติ, สมฺมาทิฏฺฐิ ฯเปฯ สมฺมาสมาธิ.\csromanspacer{} อยํ วุจฺจติ ทุกฺขนิโรธคามินี ปฏิปทา, อวเสสา ธมฺมา ทุกฺขนิโรธคามินิยา ปฏิปทาย สมฺปยุตฺตา.}

\proseitem{211}{จตฺตาริ สจฺจานิ ทุกฺขํ ทุกฺขสมุทโย ทุกฺขนิโรโธ ทุกฺขนิโรธคามินี ปฏิปทา.}

\prose{\csromanfolio{116}ตตฺถ กตโม ทุกฺขสมุทโย, ตณฺหา.\csromanspacer{} อยํ วุจฺจติ ทุกฺขสมุทโย.}

\prose{ตตฺถ กตมํ ทุกฺขํ, อวเสสา จ กิเลสา อวเสสา จ อกุสลา ธมฺมา ตีณิ จ กุสลมูลานิ สาสวานิ อวเสสา จ สาสวา กุสลา ธมฺมา สาสวา จ กุสลากุสลานํ ธมฺมานํ วิปากา เย จ ธมฺมา กิริยา เนว กุสลา นากุสลา น จ กมฺมวิปากา สพฺพญฺจ รูปํ, อิทํ วุจฺจติ ทุกฺขํ.}

\prose{ตตฺถ กตโม ทุกฺขนิโรโธ, ตณฺหาย ปหานํ.\csromanspacer{} อยํ วุจฺจติ ทุกฺขนิโรโธ.}

\prose{ตตฺถ กตมา ทุกฺขนิโรธคามินี ปฏิปทา, อิธ ภิกฺขุ ยสฺมิํ สมเย โลกุตฺตรํ ฌานํ ภาเวติ นิยฺยานิกํ อปจยคามิํ ทิฏฺฐิคตานํ ปหานาย ปฐมาย ภูมิยา ปตฺติยา วิวิจฺเจว กาเมหิ ฯเปฯ ปฐมํ ฌานํ อุปสมฺปชฺช วิหรติ ทุกฺขปฏิปทํ ทนฺธาภิญฺญํ, ตสฺมิํ สมเย ปญฺจงฺคิโก มคฺโค โหติ สมฺมาทิฏฺฐิ สมฺมาสงฺกปฺโป สมฺมาวายาโม สมฺมาสติ สมฺมาสมาธิ.}

\prose{ตตฺถ กตมา สมฺมาทิฏฺฐิ, ยา ปญฺญา ปชานนา ฯเปฯ อโมโห ธมฺมวิจโย สมฺมาทิฏฺฐิ ธมฺมวิจยสมฺโพชฺฌงฺโค มคฺคงฺคํ มคฺคปริยาปนฺนํ.\csromanspacer{} อยํ วุจฺจติ สมฺมาทิฏฺฐิ.}

\prose{ตตฺถ กตโม สมฺมาสงฺกปฺโป, โย ตกฺโก วิตกฺโก ฯเปฯ สมฺมาสงฺกปฺโป มคฺคงฺคํ มคฺคปริยาปนฺนํ.\csromanspacer{} อยํ วุจฺจติ สมฺมาสงฺกปฺโป.}

\prose{ตตฺถ กตโม สมฺมาวายาโม, โย เจตสิโก วีริยารมฺโภ ฯเปฯ สมฺมาวายาโม วีริยสมฺโพชฺฌงฺโค มคฺคงฺคํ มคฺคปริยาปนฺนํ.\csromanspacer{} อยํ วุจฺจติ สมฺมาวายาโม.}

\prose{ตตฺถ กตมา สมฺมาสติ, ยา สติ อนุสฺสติ ฯเปฯ สมฺมาสติ สติสมฺโพชฺฌงฺโค มคฺคงฺคํ มคฺคปริยาปนฺนํ.\csromanspacer{} อยํ วุจฺจติ สมฺมาสติ.}

\prose{ตตฺถ กตโม สมฺมาสมาธิ, ยา จิตฺตสฺส ฐิติ ฯเปฯ สมฺมาสมาธิ สมาธิสมฺโพชฺฌงฺโค มคฺคงฺคํ มคฺคปริยาปนฺนํ.\csromanspacer{} อยํ วุจฺจติ สมฺมาสมาธิ.\csromanspacer{} อยํ วุจฺจติ ทุกฺขนิโรธคามินี ปฏิปทา, อวเสสา ธมฺมา ทุกฺขนิโรธคามินิยา ปฏิปทาย สมฺปยุตฺตา.}

\proseitem{212}{ตตฺถ กตโม ทุกฺขสมุทโย, ตณฺหา จ อวเสสา จ กิเลสา อวเสสา จ อกุสลา ธมฺมา ตีณิ จ กุสลมูลานิ}

\vspace{\tipitakaparskip}
\csromancenter{สจฺจวิภงฺค สาสวานิ อวเสสา จ สาสวา กุสลา ธมฺมา.\csromanspacer{} อยํ วุจฺจติ ทุกฺขสมุทโย.}
\prose{\csromanfolio{117}ตตฺถ กตมํ ทุกฺขํ, สาสวา กุสลากุสลานํ ธมฺมานํ วิปากา เย จ ธมฺมา กิริยา เนว กุสลา นากุสลา น จ กมฺมวิปากา สพฺพญฺจ รูปํ.\csromanspacer{} อิทํ วุจฺจติ ทุกฺขํ.}

\prose{ตตฺถ กตโม ทุกฺขนิโรโธ, ตณฺหาย จ อวเสสานญฺจ กิเลสานํ อวเสสานญฺจ อกุสลานํ ธมฺมานํ ติณฺณญฺจ กุสลมูลานํ สาสวานํ อวเสสานญฺจ สาสวานํ กุสลานํ ธมฺมานํ ปหานํ.\csromanspacer{} อยํ วุจฺจติ ทุกฺขนิโรโธ.}

\prose{ตตฺถ กตมา ทุกฺขนิโรธคามินี ปฏิปทา, อิธ ภิกฺขุ ยสฺมิํ สมเย โลกุตฺตรํ ฌานํ ภาเวติ นิยฺยานิกํ อปจยคามิํ ทิฏฺฐิคตานํ ปหานาย ปฐมาย ภูมิยา ปตฺติยา วิวิจฺเจว กาเมหิ ฯเปฯ ปฐมํ ฌานํ อุปสมฺปชฺช วิหรติ ทุกฺขปฏิปทํ ทนฺธาภิญฺญํ, ตสฺมิํ สมเย ปญฺจงฺคิโก มคฺโค โหติ สมฺมาทิฏฺฐิ สมฺมาสงฺกปฺโป สมฺมาวายาโม สมฺมาสติ สมฺมาสมาธิ.\csromanspacer{} อยํ วุจฺจติ ทุกฺขนิโรธคามินี ปฏิปทา, อวเสสา ธมฺมา ทุกฺขนิโรธคามินิยา ปฏิปทาย สมฺปยุตฺตา.}

\proseitem{213}{จตฺตาริ สจฺจานิ ทุกฺขํ ทุกฺขสมุทโย ทุกฺขนิโรโธ ทุกฺขนิโรธคามินี ปฏิปทา.}

\prose{ตตฺถ กตโม ทุกฺขสมุทโย, ตณฺหา.\csromanspacer{} อยํ วุจฺจติ ทุกฺขสมุทโย.}

\prose{ตตฺถ กตมํ ทุกฺขํ, อวเสสา จ กิเลสา อวเสสา จ อกุสลา ธมฺมา ตีณิ จ กุสลมูลานิ สาสวานิ อวเสสา จ สาสวา กุสลา ธมฺมา สาสวา จ กุสลากุสลานํ ธมฺมานํ วิปากา เย จ ธมฺมา กิริยา เนว กุสลา นากุสลา น จ กมฺมวิปากา สพฺพญฺจ รูปํ.\csromanspacer{} อิทํ วุจฺจติ ทุกฺขํ.}

\prose{ตตฺถ กตโม ทุกฺขนิโรโธ, ตณฺหาย ปหานํ.\csromanspacer{} อยํ วุจฺจติ ทุกฺขนิโรโธ.}

\prose{ตตฺถ กตมา ทุกฺขนิโรธคามินี ปฏิปทา, อิธ ภิกฺขุ ยสฺมิํ สมเย โลกุตฺตรํ ฌานํ ภาเวติ นิยฺยานิกํ อปจยคามิํ ทิฏฺฐิคตานํ ปหานาย ปฐมาย ภูมิยา ปตฺติยา วิวิจฺเจว กาเมหิ ฯเปฯ ปฐมํ ฌานํ \csromanfolio{118}อุปสมฺปชฺช วิหรติ ทุกฺขปฏิปทํ ทนฺธาภิญฺญํ, ตสฺมิํ สมเย ผสฺโส โหติ ฯเปฯ อวิกฺเขโป โหติ.\csromanspacer{} อยํ วุจฺจติ ทุกฺขนิโรธคามินี ปฏิปทา.}

\proseitem{214}{ตตฺถ กตโม ทุกฺขสมุสโย, ตณฺหา จ อวเสสา จ กิเลสา อวเสสา จ อกุสลา ธมฺมา ตีณิ จ กุสลมูลานิ สาสวานิ อวเสสา จ สาสวา กุสลา ธมฺมา.\csromanspacer{} อยํ วุจฺจติ ทุกฺขสมุทโย.}

\prose{ตตฺถ กตมํ ทุกฺขํ, สาสวา กุสลากุสลานํ ธมฺมานํ วิปากา เย จ ธมฺมา กิริยา เนว กุสลา นากุสลา น จ กมฺมวิปากา สพฺพญฺจ รูปํ.\csromanspacer{} อิทํ วุจฺจติ ทุกฺขํ.}

\prose{ตตฺถ กตโม ทุกฺขนิโรโธ, ตณฺหาย จ อวเสสานญฺจ กิเลสานํ อวเสสานญฺจ อกุสลานํ ธมฺมานํ ติณฺณญฺจ กุสลมูลานํ สาสวานํ อวเสสานญฺจ สาสวานํ กุสลานํ ธมฺมานํ ปหานํ.\csromanspacer{} อยํ วุจฺจติ ทุกฺขนิโรโธ.}

\prose{ตตฺถ กตมา ทุกฺขนิโรธคามินี ปฏิปทา, อิธ ภิกฺขุ ยสฺมิํ สมเย โลกุตฺตรํ ฌานํ ภาเวติ นิยฺยานิกํ อปจยคามิํ ทิฏฺฐิคตานํ ปหานาย ปฐมาย ภูมิยา ปตฺติยา วิวิจฺเจว กาเมหิ ฯเปฯ ปฐมํ ฌานํ อุปสมฺปชฺช วิหรติ ทุกฺขปฏิปทํ ทนฺธาภิญฺญํ, ตสฺมิํ สมเย ผสฺโส โหติ ฯเปฯ อวิกฺเขโป โหติ.\csromanspacer{} อยํ วุจฺจติ ทุกฺขนิโรธคามินี ปฏิปทา.}

\csromanheader{2}{อภิธมฺมภาชนียํ.}
\csromansectionrule
\csromanmatikaanchor{mtk.36}
\csromantocmarkref{subsection}{3. ปญฺหาปุจฺฉก}{mtk.36}
\bookmark[dest={mtk.36},level=2]{3. ปญฺหาปุจฺฉก}
\csromanheader{2}{3. ปญฺหาปุจฺฉก}
\proseitem{215}{จตฺตาริ อริยสจฺจานิ ทุกฺขํ อริยสจฺจํ ทุกฺขสมุทยํ อริยสจฺจํ ทุกฺขนิโรธํ อริยสจฺจํ ทุกฺขนิโรธคามินี ปฏิปทา อริยสจฺจํ.}

\proseitem{216}{จตุนฺนํ อริยสจฺจานํ กติ กุสลา, กติ อกุสลา, กติ อพฺยากตา ฯเปฯ กติ สรณา, กติ อรณา.}

\csromanmatikaanchor{mtk.37}
\csromantocmarkref{subsubsection}{1. ติก}{mtk.37}
\bookmark[dest={mtk.37},level=3]{1. ติก}
\csromanheader{3}{4. สจฺจวิภงฺค \textbf{1. ติก}}
\proseitem{217}{\csromanfolio{119}สมุทยสจฺจํ อกุสลํ.\csromanspacer{} มคฺคสจฺจํ กุสลํ.\csromanspacer{} นิโรธสจฺจํ อพฺยากตํ.\csromanspacer{} ทุกฺขสจฺจํ สิยา กุสลํ, สิยา อกุสลํ, สิยา อพฺยากตํ.\csromanspacer{} เทฺว สจฺจา สิยา สุขาย เวทนาย สมฺปยุตฺตา, สิยา อทุกฺขมสุขาย เวทนาย สมฺปยุตฺตา.\csromanspacer{} นิโรธสจฺจํ น วตฺตพฺพํ “สุขาย เวทนาย สมฺปยุตฺตนฺ”ติปิ, “ทุกฺขาย เวทนาย สมฺปยุตฺตนฺ”ติปิ, “อทุกฺขมสุขาย เวทนาย สมฺปยุตฺตนฺ”ติปิ.\csromanspacer{} ทุกฺขสจฺจํ สิยา สุขาย เวทนาย สมฺปยุตฺตํ, สิยา ทุกฺขาย เวทนาย สมฺปยุตฺตํ, สิยา อทุกฺขมสุขาย เวทนาย สมฺปยุตฺตํ, สิยา น วตฺตพฺพํ “สุขาย เวทนาย สมฺปยุตฺตนฺ”ติปิ, “ทุกฺขาย เวทนาย สมฺปยุตฺตนฺ”ติปิ, “อทุกฺขมสุขาย เวทนาย สมฺปยุตฺตนฺ”ติปิ.\csromanspacer{} เทฺว สจฺจา วิปากธมฺมธมฺมา.\csromanspacer{} นิโรธสจฺจํ เนววิปากนวิปากธมฺมธมฺมํ.\csromanspacer{} ทุกฺขสจฺจํ สิยา วิปากํ, สิยา วิปากธมฺมธมฺมํ, สิยา เนววิปากนวิปากธมฺมธมฺมํ.\csromanspacer{} สมุทยสจฺจํ อนุปาทินฺนุปาทานิยํ.\csromanspacer{} เทฺว สจฺจา อนุปาทินฺน- อนุปาทานิยา.\csromanspacer{} ทุกฺขสจฺจํ สิยา อุปาทินฺนุปาทานิยํ, สิยา อนุปาทินฺนุปาทานิยํ.}

\prose{สมุทยสจฺจํ สํกิลิฏฺฐสํกิเลสิกํ.\csromanspacer{} เทฺว สจฺจา อสํกิลิฏฺฐ- อสํกิเลสิกา.\csromanspacer{} ทุกฺขสจฺจํ สิยา สํกิลิฏฺฐสํกิเลสิกํ, สิยา อสํกิลิฏฺฐสํกิเลสิกํ.\csromanspacer{} สมุทยสจฺจํ สวิตกฺกสวิจารํ.\csromanspacer{} นิโรธสจฺจํ อวิตกฺก-อวิจารํ.\csromanspacer{} มคฺคสจฺจํ สิยา สวิตกฺกสวิจารํ, สิยา อวิตกฺกวิจารมตฺตํ, สิยา อวิตกฺก-อวิจารํ.\csromanspacer{} ทุกฺขสจฺจํ สิยา สวิตกฺกสวิจารํ, สิยา อวิตกฺกวิจารมตฺตํ, สิยา อวิตกฺก-อวิจารํ, สิยา น วตฺตพฺพํ “สวิตกฺกสวิจารนฺ”ติปิ, “อวิตกฺกวิจารมตฺตนฺ”ติปิ, “อวิตกฺก-อวิจารนฺ”ติปิ.\csromanspacer{} เทฺว สจฺจา สิยา ปีติสหคตา, สิยา สุขสหคตา, สิยา อุเปกฺขาสหคตา.\csromanspacer{} นิโรธสจฺจํ น วตฺตพฺพํ “ปีติสหคตนฺ”ติปิ, “สุขสหคตนฺ”ติปิ, “อุเปกฺขาสหคตนฺ”ติปิ.\csromanspacer{} ทุกฺขสจฺจํ สิยา ปีติสหคตํ, สิยา สุขสหคตํ, สิยา อุเปกฺขาสหคตํ, สิยา น วตฺตพฺพํ “ปีติสหคตนฺ”ติปิ, “สุขสหคตนฺ”ติปิ, “อุเปกฺขาสหคตนฺ”ติปิ.}

\prose{เทฺว สจฺจา เนว ทสฺสเนน น ภาวนาย ปหาตพฺพา.\csromanspacer{} สมุทยสจฺจํ สิยา ทสฺสเนน ปหาตพฺพํ, สิยา ภาวนาย ปหาตพฺพํ.\csromanspacer{} ทุกฺขสจฺจํ สิยา ทสฺสเนน ปหาตพฺพํ, สิยา ภาวนาย ปหาตพฺพํ, สิยา เนว ทสฺสเนน น ภาวนาย ปหาตพฺพํ.\csromanspacer{} เทฺว สจฺจา เนว ทสฺสเนน น ภาวนาย ปหาตพฺพเหตุกา.\csromanspacer{} สมุทยสจฺจํ สิยา ทสฺสเนน ปหาตพฺพเหตุกํ, สิยา ภาวนาย \csromanfolio{120}ปหาตพฺพเหตุกํ.\csromanspacer{} ทุกฺขสจฺจํ สิยา ทสฺสเนน ปหาตพฺพเหตุกํ, สิยา ภาวนาย ปหาตพฺพเหตุกํ, สิยา เนว ทสฺสเนน น ภาวนาย ปหาตพฺพเหตุกํ.\csromanspacer{} สมุทยสจฺจํ อาจยคามิ.\csromanspacer{} มคฺคสจฺจํ อปจยคามิ.\csromanspacer{} นิโรธสจฺจํ เนวาจยคามินาปจยคามิ.\csromanspacer{} ทุกฺขสจฺจํ สิยา อาจยคามิ, สิยา เนวาจยคามินาปจยคามิ.\csromanspacer{} มคฺคสจฺจํ เสกฺขํ.\csromanspacer{} ตีณิ สจฺจานิ เนวเสกฺขนาเสกฺขา.\csromanspacer{} สมุทยสจฺจํ ปริตฺตํ.\csromanspacer{} เทฺว สจฺจา อปฺปมาณา.\csromanspacer{} ทุกฺขสจฺจํ สิยา ปริตฺตํ, สิยา มหคฺคตํ.\csromanspacer{} นิโรธสจฺจํ อนารมฺมณํ.\csromanspacer{} มคฺคสจฺจํ อปฺปมาณารมฺมณํ.\csromanspacer{} สมุทยสจฺจํ สิยา ปริตฺตารมฺมณํ, สิยา มหคฺคตารมฺมณํ, น อปฺปมาณารมฺมณํ, สิยา น วตฺตพฺพํ “ปริตฺตารมฺมณนฺ”ติปิ, “มหคฺคตารมฺมณนฺ”ติปิ.\csromanspacer{} ทุกฺขสจฺจํ สิยา ปริตฺตารมฺมณํ, สิยา มหคฺคตารมฺมณํ, สิยา อปฺปมาณารมฺมณํ, สิยา น วตฺตพฺพํ “ปริตฺตารมฺมณนฺ”ติปิ, “มหคฺคตารมฺมณนฺ”ติปิ, “อปฺปมาณารมฺมณนฺ”ติปิ.}

\prose{สมุทยสจฺจํ หีนํ.\csromanspacer{} เทฺว สจฺจา ปณีตา.\csromanspacer{} ทุกฺขสจฺจํ สิยา หีนํ, สิยา มชฺฌิมํ.\csromanspacer{} นิโรธสจฺจํ อนิยตํ.\csromanspacer{} มคฺคสจฺจํ สมฺมตฺตนิยตํ.\csromanspacer{} เทฺว สจฺจา สิยา มิจฺฉตฺตนิยตา สิยา อนิยตา.\csromanspacer{} นิโรธสจฺจํ อนารมฺมณํ.\csromanspacer{} สมุทยสจฺจํ น วตฺตพฺพํ “มคฺคารมฺมณนฺ”ติปิ, “มคฺคเหตุกนฺ”ติปิ, “มคฺคาธิปตี”ติปิ.\csromanspacer{} มคฺคสจฺจํ น มคฺคารมฺมณํ มคฺคเหตุกํ, สิยา มคฺคาธิปติ, สิยา น วตฺตพฺพํ “มคฺคาธิปตี”ติ.\csromanspacer{} ทุกฺขสจฺจํ สิยา มคฺคารมฺมณํ, น มคฺคเหตุกํ.\csromanspacer{} สิยา มคฺคาธิปติ, สิยา น วตฺตพฺพํ “มคฺคารมฺมณนฺ”ติปิ, “มคฺคาธิปตี”ติปิ.\csromanspacer{} เทฺว สจฺจา สิยา อุปฺปนฺนา, สิยา อนุปฺปนฺนา, น วตฺตพฺพา “อุปฺปาทิโน”ติ.\csromanspacer{} นิโรธสจฺจํ น วตฺตพฺพํ “อุปฺปนฺนนฺ”ติปิ, “อนุปฺปนฺนนฺ”ติปิ, “อุปฺปาที”ติปิ.\csromanspacer{} ทุกฺขสจฺจํ สิยา อุปฺปนฺนํ, สิยา อนุปฺปนฺนํ, สิยา อุปฺปาทิ.\csromanspacer{} ตีณิ สจฺจานิ สิยา อตีตา, สิยา อนาคตา, สิยา ปจฺจุปฺปนฺนา.\csromanspacer{} นิโรธสจฺจํ น วตฺตพฺพํ “อตีตนฺ”ติปิ, “อนาคตนฺ”ติปิ, “ปจฺจุปฺปนฺนนฺ”ติปิ.\csromanspacer{} นิโรธสจฺจํ อนารมฺมณํ.\csromanspacer{} มคฺคสจฺจํ น วตฺตพฺพํ “อตีตารมฺมณนฺ”ติปิ, “อนาคตารมฺมณนฺ”ติปิ, “ปจฺจุปฺปนฺนารมฺมณนฺ”ติปิ.\csromanspacer{} เทฺว สจฺจา สิยา อตีตารมฺมณา, สิยา อนาคตารมฺมณา, สิยา ปจฺจุปฺปนฺนารมฺมณา, สิยา น วตฺตพฺพา “อตีตารมฺมณา”ติปิ, “อนาคตารมฺมณา”ติปิ, “ปจฺจุปฺปนฺนารมฺมณา”ติปิ.\csromanspacer{} นิโรธสจฺจํ พหิทฺธา.\csromanspacer{} ตีณิ สจฺจานิ สิยา อชฺฌตฺตา, สิยา พหิทฺธา, สิยา อชฺฌตฺตพหิทฺธา.\csromanspacer{} นิโรธสจฺจํ อนารมฺมณํ.\csromanspacer{} มคฺคสจฺจํ พหิทฺธารมฺมณํ.\csromanspacer{} สมุทยสจฺจํ สิยา อชฺฌตฺตารมฺมณํ, สิยา พหิทฺธารมฺมณํ, สิยา อชฺฌตฺตพหิทฺธารมฺมณํ.\csromanspacer{} ทุกฺขสจฺจํ สิยา อชฺฌตฺตารมฺมณํ, สิยา พหิทฺธารมฺมณํ, สิยา}

\vspace{\tipitakaparskip}
\csromancenter{สจฺจวิภงฺค อชฺฌตฺตพหิทฺธารมฺมณํ, สิยา น วตฺตพฺพํ “อชฺฌตฺตารมฺมณนฺ”ติปิ, “พหิทฺธารมฺมณนฺ”ติปิ, “อชฺฌตฺตพหิทฺธารมฺมณนฺ”ติปิ.\csromanspacer{} ตีณิ สจฺจานิ อนิทสฺสน-อปฺปฏิฆา.\csromanspacer{} ทุกฺขสจฺจํ สิยา สนิทสฺสนสปฺปฏิฆํ, สิยา อนิทสฺสนสปฺปฏิฆํ, สิยา อนิทสฺสน-อปฺปฏิฆํ.}
\csromanmatikaanchor{mtk.38}
\csromantocmarkref{subsubsection}{2. ทุก}{mtk.38}
\bookmark[dest={mtk.38},level=3]{2. ทุก}
\csromanheader{3}{2. ทุก}
\proseitem{218}{\csromanfolio{121}สมุทยสจฺจํ เหตุ.\csromanspacer{} นิโรธสจฺจํ น เหตุ.\csromanspacer{} เทฺว สจฺจา สิยา เหตู, สิยา น เหตู.\csromanspacer{} เทฺว สจฺจา สเหตุกา.\csromanspacer{} นิโรธสจฺจํ อเหตุกํ.\csromanspacer{} ทุกฺขสจฺจํ สิยา สเหตุกํ, สิยา อเหตุกํ.\csromanspacer{} เทฺว สจฺจา เหตุสมฺปยุตฺตา.\csromanspacer{} นิโรธสจฺจํ เหตุวิปฺปยุตฺตํ.\csromanspacer{} ทุกฺขสจฺจํ สิยา เหตุสมฺปยุตฺตํ, สิยา เหตุวิปฺปยุตฺตํ.\csromanspacer{} สมุทยสจฺจํ เหตุ เจว สเหตุกญฺจ.\csromanspacer{} นิโรธสจฺจํ น วตฺตพฺพํ “เหตุ เจว สเหตุกญฺจา”ติปิ, “สเหตุกญฺเจว น จ เหตู”ติปิ.\csromanspacer{} มคฺคสจฺจํ สิยา เหตุ เจว สเหตุกญฺจ, สิยา สเหตุกญฺเจว น จ เหตุ.\csromanspacer{} ทุกฺขสจฺจํ สิยา เหตุ เจว สเหตุกญฺจ, สิยา สเหตุกญฺเจว น จ เหตุ, สิยา น วตฺตพฺพํ “เหตุ เจว สเหตุกญฺจา”ติปิ, “สเหตุกญฺเจว น จ เหตู”ติปิ.\csromanspacer{} สมุทยสจฺจํ เหตุ เจว เหตุสมฺปยุตฺตญฺจ.\csromanspacer{} นิโรธสจฺจํ น วตฺตพฺพํ “เหตุ เจว เหตุสมฺปยุตฺตญฺจา”ติปิ, “เหตุสมฺปยุตฺตญฺเจว น จ เหตู”ติปิ.\csromanspacer{} มคฺคสจฺจํ สิยา เหตุ เจว เหตุสมฺปยุตฺตญฺจ, สิยา เหตุสมฺปยุตฺตญฺเจว น จ เหตุ.\csromanspacer{} ทุกฺขสจฺจํ สิยา เหตุ เจว เหตุสมฺปยุตฺตญฺจ, สิยา เหตุสมฺปยุตฺตญฺเจว น จ เหตุ, สิยา น วตฺตพฺพํ “เหตุ เจว เหตุสมฺปยุตฺตญฺจา”ติปิ, “เหตุสมฺปยุตฺตญฺเจว น จ เหตู”ติปิ.\csromanspacer{} นิโรธสจฺจํ น เหตุ-อเหตุกํ.\csromanspacer{} สมุทยสจฺจํ น วตฺตพฺพํ “น เหตุสเหตุกนฺ”ติปิ, “น เหตุ-อเหตุกนฺ”ติปิ.\csromanspacer{} มคฺคสจฺจํ สิยา น เหตุสเหตุกํ, สิยา น วตฺตพฺพํ “น เหตุสเหตุกนฺ”ติปิ, “น เหตุ- อเหตุกนฺ”ติปิ.\csromanspacer{} ทุกฺขสจฺจํ สิยา น เหตุสเหตุกํ, สิยา น เหตุ- อเหตุกํ, สิยา น วตฺตพฺพํ “น เหตุสเหตุกนฺ”ติปิ, “น เหตุ-อเหตุกนฺ”ติปิ. (1)}

\prose{ตีณิ สจฺจานิ สปฺปจฺจยา.\csromanspacer{} นิโรธสจฺจํ อปฺปจฺจยํ.\csromanspacer{} ตีณิ สจฺจานิ สงฺขตา.\csromanspacer{} นิโรธสจฺจํ อสงฺขตํ.\csromanspacer{} ตีณิ สจฺจานิ อนิทสฺสนา.\csromanspacer{} ทุกฺขสจฺจํ สิยา สนิทสฺสนํ, สิยา อนิทสฺสนํ.\csromanspacer{} ตีณิ สจฺจานิ อปฺปฏิฆา.\csromanspacer{} ทุกฺขสจฺจํ สิยา สปฺปฏิฆํ, สิยา อปฺปฏิฆํ.\csromanspacer{} ตีณิ สจฺจานิ อรูปานิ.\csromanspacer{} ทุกฺขสจฺจํ สิยา รูปํ, สิยา อรูปํ.\csromanspacer{} เทฺว \csromanfolio{122}สจฺจา โลกิยา.\csromanspacer{} เทฺว สจฺจา โลกุตฺตรา, เกนจิ วิญฺเญยฺยา, เกนจิ น วิญฺเญยฺยา. (2)}

\prose{สมุทยสจฺจํ อาสโว.\csromanspacer{} เทฺว สจฺจา โน อาสวา.\csromanspacer{} ทุกฺขสจฺจํ สิยา อาสโว, สิยา โน อาสโว.\csromanspacer{} เทฺว สจฺจา สาสวา.\csromanspacer{} เทฺว สจฺจา อนาสวา.\csromanspacer{} สมุทยสจฺจํ อาสวสมฺปยุตฺตา.\csromanspacer{} เทฺว สจฺจา อาสววิปฺปยุตฺตา.\csromanspacer{} ทุกฺขสจฺจํ สิยา อาสวสมฺปยุตฺตํ, สิยา อาสววิปฺปยุตฺตํ.\csromanspacer{} สมุทยสจฺจํ อาสโว เจว สาสวญฺจ.\csromanspacer{} เทฺว สจฺจา น วตฺตพฺพา “อาสวา เจว สาสวา จา”ติปิ, “สาสวา เจว โน จ อาสวา”ติปิ.\csromanspacer{} ทุกฺขสจฺจํ สิยา อาสโว เจว สาสวญฺจ, สิยา สาสวญฺเจว โน จ อาสโว.\csromanspacer{} สมุทยสจฺจํ อาสโว เจว อาสวสมฺปยุตฺตญฺจ.\csromanspacer{} เทฺว สจฺจา น วตฺตพฺพา “อาสวา เจว อาสวสมฺปยุตฺตา จา”ติปิ, “อาสวสมฺปยุตฺตา เจว โน จ อาสวา”ติปิ.\csromanspacer{} ทุกฺขสจฺจํ สิยา อาสโว เจว อาสวสมฺปยุตฺตญฺจ, สิยา อาสวสมฺปยุตฺตญฺเจว โน จ อาสโว, สิยา น วตฺตพฺพํ “อาสโว เจว อาสวสมฺปยุตฺตญฺจา”ติปิ, “อาสวสมฺปยุตฺตญฺเจว โน จ อาสโว”ติปิ.\csromanspacer{} เทฺว สจฺจา อาสววิปฺปยุตฺต-อนาสวา.\csromanspacer{} สมุทยสจฺจํ น วตฺตพฺพํ “อาสววิปฺปยุตฺตสาสวนฺ”ติปิ, “อาสววิปฺปยุตฺต-อนาสวนฺ”ติปิ.\csromanspacer{} ทุกฺขสจฺจํ สิยา อาสววิปฺปยุตฺตสาสวํ, สิยา น วตฺตพฺพํ “อาสววิปฺปยุตฺตสาสวนฺ”ติปิ, “อาสววิปฺปยุตฺต-อนาสวนฺ”ติปิ. (3)}

\prose{สมุทยสจฺจํ สํโยชนํ.\csromanspacer{} เทฺว สจฺจา โน สํโยชนา.\csromanspacer{} ทุกฺขสจฺจํ สิยา สํโยชนํ, สิยา โน สํโยชนํ.\csromanspacer{} เทฺว สจฺจา สํโยชนิยา.\csromanspacer{} เทฺว สจฺจา อสํโยชนิยา.\csromanspacer{} สมุทยสจฺจํ สํโยชนสมฺปยุตฺตํ.\csromanspacer{} เทฺว สจฺจา สํโยชนวิปฺปยุตฺตา.\csromanspacer{} ทุกฺขสจฺจํ สิยา สํโยชนสมฺปยุตฺตํ, สิยา สํโยชนวิปฺปยุตฺตํ.\csromanspacer{} สมุทยสจฺจํ สํโยชนญฺเจว สํโยชนิยญฺจ.\csromanspacer{} เทฺว สจฺจา น วตฺตพฺพา “สํโยชนา เจว สํโยชนิยา จา”ติปิ, “สํโยชนิยา เจว โน จ สํโยชนา”ติปิ.\csromanspacer{} ทุกฺขสจฺจํ สิยา สํโยชนญฺเจว สํโยชนิยญฺจ, สิยา สํโยชนิยญฺเจว โน จ สํโยชนํ.\csromanspacer{} สมุทยสจฺจํ สํโยชนญฺเจว สํโยชนสมฺปยุตฺตญฺจ.\csromanspacer{} เทฺว สจฺจา น วตฺตพฺพา “สํโยชนา เจว สํโยชนสมฺปยุตฺตา จา”ติปิ, “สํโยชนสมฺปยุตฺตา เจว โน จ สํโยชนา”ติปิ.\csromanspacer{} ทุกฺขสจฺจํ สิยา สํโยชนญฺเจว สํโยชนสมฺปยุตฺตญฺจ, สิยา สํโยชนสมฺปยุตฺตญฺเจว โน จ สํโยชนํ, สิยา น วตฺตพฺพํ “สํโยชนญฺเจว สํโยชนสมฺปยุตฺตญฺจา”ติปิ, สํโยชนสมฺปยุตฺตญฺเจว โน จ}

\vspace{\tipitakaparskip}
\csromancenter{สจฺจวิภงฺค สํโยชนนฺ”ติปิ.\csromanspacer{} เทฺว สจฺจา สํโยชนวิปฺปยุตฺต-อสํโยชนิยา.\csromanspacer{} สมุทยสจฺจํ น วตฺตพฺพํ “สํโยชนวิปฺปยุตฺตสํโยชนิยนฺ”ติปิ, “สํโยชนวิปฺปยุตฺต-อสํโยชนิยนฺ”ติปิ.\csromanspacer{} ทุกฺขสจฺจํ สิยา สํโยชนวิปฺปยุตฺตสํโยชนิยํ, สิยา น วตฺตพฺพํ “สํโยชนวิปฺปยุตฺตสํโยชนิยนฺ”ติปิ, “สํโยชนวิปฺปยุตฺต- อสํโยชนิยนฺ”ติปิ. (4)}
\prose{\csromanfolio{123}สมุทยสจฺจํ คนฺโถ.\csromanspacer{} เทฺว สจฺจา โน คนฺถา.\csromanspacer{} ทุกฺขสจฺจํ สิยา คนฺโถ, สิยา โน คนฺโถ.\csromanspacer{} เทฺว สจฺจา คนฺถนิยา.\csromanspacer{} เทฺว สจฺจา อคนฺถนิยา.\csromanspacer{} เทฺว สจฺจา คนฺถวิปฺปยุตฺตา.\csromanspacer{} เทฺว สจฺจา สิยา คนฺถสมฺปยุตฺตา, สิยา คนฺถวิปฺปยุตฺตา.\csromanspacer{} สมุทยสจฺจํ คนฺโถ เจว คนฺถนิยญฺจ.\csromanspacer{} เทฺว สจฺจา น วตฺตพฺพา “คนฺถา เจว คนฺถนิยา จา”ติปิ, “คนฺถนิยา เจว โน จ คนฺถา”ติปิ.\csromanspacer{} ทุกฺขสจฺจํ สิยา คนฺโถ เจว คนฺถนิยญฺจ, สิยา คนฺถนิยญฺเจว โน จ คนฺโถ.\csromanspacer{} สมุทยสจฺจํ คนฺโถ เจว คนฺถสมฺปยุตฺตญฺจ, สิยา น วตฺตพฺพํ “คนฺโถ เจว คนฺถสมฺปยุตฺตญฺจา”ติปิ, “คนฺถสมฺปยุตฺตญฺเจว โน จ คนฺโถ”ติปิ.\csromanspacer{} เทฺว สจฺจา น วตฺตพฺพา “คนฺถา เจว คนฺถสมฺปยุตฺตา จา”ติปิ, “คนฺถสมฺปยุตฺตา เจว โน จ คนฺโถ”ติปิ.\csromanspacer{} ทุกฺขสจฺจํ สิยา คนฺโถ เจว คนฺถสมฺปยุตฺตญฺจ, สิยา คนฺถสมฺปยุตฺตญฺเจว โน จ คนฺโถ, สิยา น วตฺตพฺพํ “คนฺโถ เจว คนฺถสมฺปยุตฺตญฺจา”ติปิ, “คนฺถสมฺปยุตฺตญฺเจว โน จ คนฺโถ”ติปิ.\csromanspacer{} เทฺว สจฺจา คนฺถวิปฺปยุตฺต-อคนฺถนิยา.\csromanspacer{} เทฺว สจฺจา สิยา คนฺถวิปฺปยุตฺตคนฺถนิยา, สิยา น วตฺตพฺพา “คนฺถวิปฺปยุตฺตคนฺถนิยา”ติปิ, “คนฺถวิปฺปยุตฺต-อคนฺถนิยา”ติปิ. (5)}

\prose{สมุทยสจฺจํ โอโฆ ฯเปฯ โยโค ฯเปฯ นีวรณํ.\csromanspacer{} เทฺว สจฺจา โน นีวรณา.\csromanspacer{} ทุกฺขสจฺจํ สิยา นีวรณํ, สิยา โน นีวรณํ.\csromanspacer{} เทฺว สจฺจา นีวรณิยา.\csromanspacer{} เทฺว สจฺจา อนีวรณิยา.\csromanspacer{} สมุทยสจฺจํ นีวรณสมฺปยุตฺตํ.\csromanspacer{} เทฺว สจฺจา นีวรณวิปฺปยุตฺตา.\csromanspacer{} ทุกฺขสจฺจํ สิยา นีวรณสมฺปยุตฺตํ สิยา นีวรณวิปฺปยุตฺตํ.\csromanspacer{} สมุทยสจฺจํ นีวรณญฺเจว นีวรณิยญฺจ.\csromanspacer{} เทฺว สจฺจา น วตฺตพฺพา “นีวรณา เจว นีวรณิยา จา”ติปิ, “นีวรณิยา เจว โน จ นีวรณา”ติปิ.\csromanspacer{} ทุกฺขสจฺจํ สิยา นีวรณญฺเจว นีวรณิยญฺจ, สิยา นีวรณิยญฺเจว โน จ นีวรณํ.\csromanspacer{} สมุทยสจฺจํ นีวรณญฺเจว นีวรณสมฺปยุตฺตญฺจ.\csromanspacer{} เทฺว สจฺจา น วตฺตพฺพา “นีวรณา เจว นีวรณสมฺปยุตฺตา จา”ติปิ, “นีวรณสมฺปยุตฺตา เจว โน จ นีวรณา”ติปิ.\csromanspacer{} ทุกฺขสจฺจํ สิยา นีวรณญฺเจว นีวรณสมฺปยุตฺตญฺจ, สิยา นีวรณสมฺปยุตฺตญฺเจว โน จ นีวรณํ, สิยา น วตฺตพฺพํ “นีวรณญฺเจว \csromanfolio{124}นีวรณสมฺปยุตฺตญฺจา”ติปิ, “นีวรณสมฺปยุตฺตญฺเจว โน จ นีวรณนฺ”ติปิ.\csromanspacer{} เทฺว สจฺจา นีวรณวิปฺปยุตฺต-อนีวรณิยา, สมุทยสจฺจํ น วตฺตพฺพํ “นีวรณวิปฺปยุตฺตนีวรณิยนฺ”ติปิ, “นีวรณวิปฺปยุตฺต- อนีวรณิยนฺ”ติปิ.\csromanspacer{} ทุกฺขสจฺจํ สิยา นีวรณวิปฺปยุตฺตนีวรณิยํ, สิยา น วตฺตพฺพํ “นีวรณวิปฺปยุตฺตนีวรณิยนฺ”ติปิ, “นีวรณวิปฺปยุตฺต- อนีวรณิยนฺ”ติปิ. (8)}

\prose{ตีณิ สจฺจานิ โน ปรามาสา.\csromanspacer{} ทุกฺขสจฺจํ สิยา ปรามาโส, สิยา โน ปรามาโส.\csromanspacer{} เทฺว สจฺจา ปรามฏฺฐา.\csromanspacer{} เทฺว สจฺจา อปรามฏฺฐา.\csromanspacer{} เทฺว สจฺจา ปรามาสวิปฺปยุตฺตา.\csromanspacer{} สมุทยสจฺจํ สิยา ปรามาสสมฺปยุตฺตํ, สิยา ปรามาสวิปฺปยุตฺตํ.\csromanspacer{} ทุกฺขสจฺจํ สิยา ปรามาสสมฺปยุตฺตํ, สิยา ปรามาสวิปฺปยุตฺตํ, สิยา น วตฺตพฺพํ “ปรามาสสมฺปยุตฺตนฺ”ติปิ, “ปรามาสวิปฺปยุตฺตนฺ”ติปิ.\csromanspacer{} สมุทยสจฺจํ น วตฺตพฺพํ “ปรามาโส เจว ปรามฏฺฐญฺจา”ติ, ปรามฏฺฐญฺเจว โน จ ปรามาโส.\csromanspacer{} เทฺว สจฺจา น วตฺตพฺพา “ปรามาสา เจว ปรามฏฺฐา จา”ติปิ, “ปรามฏฺฐา เจว โน จ ปรามาสา”ติปิ.\csromanspacer{} ทุกฺขสจฺจํ สิยา ปรามาโส เจว ปรามฏฺฐญฺจ, สิยา ปรามฏฺฐญฺเจว โน จ ปรามาโส.\csromanspacer{} เทฺว สจฺจา ปรามาสวิปฺปยุตฺต-อปรามฏฺฐา.\csromanspacer{} เทฺว สจฺจา สิยา ปรามาสวิปฺปยุตฺตปรามฏฺฐา, สิยา น วตฺตพฺพา “ปรามาสวิปฺปยุตฺตปรามฏฺฐา”ติปิ, “ปรามาสวิปฺปยุตฺต-อปรามฏฺฐา”ติปิ. (9)}

\prose{เทฺว สจฺจา สารมฺมณา.\csromanspacer{} นิโรธสจฺจํ อนารมฺมณํ.\csromanspacer{} ทุกฺขสจฺจํ สิยา สารมฺมณํ, สิยา อนารมฺมณํ.\csromanspacer{} ตีณิ สจฺจานิ โน จิตฺตา.\csromanspacer{} ทุกฺขสจฺจํ สิยา จิตฺตํ, สิยา โน จิตฺตํ.\csromanspacer{} เทฺว สจฺจา เจตสิกา.\csromanspacer{} นิโรธสจฺจํ อเจตสิกํ.\csromanspacer{} ทุกฺขสจฺจํ สิยา เจตสิกํ, สิยา อเจตสิกํ.\csromanspacer{} เทฺว สจฺจา จิตฺตสมฺปยุตฺตา.\csromanspacer{} นิโรธสจฺจํ จิตฺตวิปฺปยุตฺตํ.\csromanspacer{} ทุกฺขสจฺจํ สิยา จิตฺตสมฺปยุตฺตํ, สิยา จิตฺตวิปฺปยุตฺตํ, สิยา น วตฺตพฺพํ “จิตฺเตน สมฺปยุตฺตนฺ”ติปิ, “จิตฺเตน วิปฺปยุตฺตนฺ”ติปิ.\csromanspacer{} เทฺว สจฺจา จิตฺตสํสฏฺฐา.\csromanspacer{} นิโรธสจฺจํ จิตฺตวิสํสฏฺฐํ.\csromanspacer{} ทุกฺขสจฺจํ สิยา จิตฺตสํสฏฺฐํ, สิยา จิตฺตวิสํสฏฺฐํ, สิยา น วตฺตพฺพํ “จิตฺเตน สํสฏฺฐนฺ”ติปิ, “จิตฺเตน วิสํสฏฺฐนฺ”ติปิ.\csromanspacer{} เทฺว สจฺจา จิตฺตสมุฏฺฐานา.\csromanspacer{} นิโรธสจฺจํ โน จิตฺตสมุฏฺฐานํ.\csromanspacer{} ทุกฺขสจฺจํ สิยา จิตฺตสมุฏฺฐานํ, สิยา โน จิตฺตสมุฏฺฐานํ.\csromanspacer{} เทฺว สจฺจา จิตฺตสหภุโน.\csromanspacer{} นิโรธสจฺจํ โน จิตฺตสหภู.\csromanspacer{} ทุกฺขสจฺจํ สิยา จิตฺตสหภู, สิยา โน จิตฺตสหภู.\csromanspacer{} เทฺว สจฺจา จิตฺตานุปริวตฺติโน.\csromanspacer{} นิโรธสจฺจํ โน จิตฺตานุปริวตฺติ.\csromanspacer{} ทุกฺขสจฺจํ สิยา จิตฺตานุปริวตฺติ, สิยา โน จิตฺตานุปริวตฺติ.\csromanspacer{} เทฺว สจฺจา จิตฺตสํสฏฺฐสมุฏฺฐานา.\csromanspacer{} นิโรธสจฺจํ โน}

\vspace{\tipitakaparskip}
\csromancenter{สจฺจวิภงฺค จิตฺตสํสฏฺฐสมุฏฺฐานํ.\csromanspacer{} ทุกฺขสจฺจํ สิยา จิตฺตสํสฏฺฐสมุฏฺฐานํ, สิยา โน จิตฺตสํสฏฺฐสมุฏฺฐานํ.\csromanspacer{} เทฺว สจฺจา จิตฺตสํสฏฺฐสมุฏฺฐานสหภุโน.\csromanspacer{} นิโรธสจฺจํ โน จิตฺตสํสฏฺฐสมุฏฺฐานสหภู.\csromanspacer{} ทุกฺขสจฺจํ สิยา จิตฺตสํสฏฺฐสมุฏฺฐานสหภู, สิยา โน จิตฺตสํสฏฺฐสมุฏฺฐานสหภู.\csromanspacer{} เทฺว สจฺจา จิตฺตสํสฏฺฐสมุฏฺฐานานุปริวตฺติโน.\csromanspacer{} นิโรธสจฺจํ โน จิตฺตสํสฏฺฐสมุฏฺฐานานุปริวตฺติ.\csromanspacer{} ทุกฺขสจฺจํ สิยา จิตฺตสํสฏฺฐสมุฏฺฐานานุปริวตฺติ, สิยา โน จิตฺตสํสฏฺฐสมุฏฺฐานานุปริวตฺติ.\csromanspacer{} ตีณิ สจฺจานิ พาหิรา.\csromanspacer{} ทุกฺขสจฺจํ สิยา อชฺฌตฺตํ, สิยา พาหิรํ. (10)}
\prose{\csromanfolio{125}ตีณิ สจฺจานิ โน อุปาทา.\csromanspacer{} ทุกฺขสจฺจํ สิยา อุปาทา, สิยา โน อุปาทา.\csromanspacer{} ตีณิ สจฺจานิ อนุปาทินฺนา.\csromanspacer{} ทุกฺขสจฺจํ สิยา อุปาทินฺนํ, สิยา อนุปาทินฺนํ.\csromanspacer{} สมุทยสจฺจํ อุปาทานํ.\csromanspacer{} เทฺว สจฺจา โน อุปาทานา.\csromanspacer{} ทุกฺขสจฺจํ สิยา อุปาทานํ, สิยา โน อุปาทานํ.\csromanspacer{} เทฺว สจฺจา อุปาทานิยา.\csromanspacer{} เทฺว สจฺจา อนุปาทานิยา.\csromanspacer{} เทฺว สจฺจา อุปาทานวิปฺปยุตฺตา.\csromanspacer{} เทฺว สจฺจา สิยา อุปาทานสมฺปยุตฺตา, สิยา อุปาทานวิปฺปยุตฺตา.\csromanspacer{} สมุทยสจฺจํ อุปาทานญฺเจว อุปาทานิยญฺจ.\csromanspacer{} เทฺว สจฺจา น วตฺตพฺพา “อุปาทานา เจว อุปาทานิยา จา”ติปิ, “อุปาทานิยา เจว โน จ อุปาทานา”ติปิ.\csromanspacer{} ทุกฺขสจฺจํ สิยา อุปาทานญฺเจว อุปาทานิยญฺจ, สิยา อุปาทานิยญฺเจว โน จ อุปาทานํ.\csromanspacer{} สมุทยสจฺจํ สิยา อุปาทานญฺเจว อุปาทานสมฺปยุตฺตญฺจ.\csromanspacer{} สิยา น วตฺตพฺพํ “อุปาทานญฺเจว อุปาทานสมฺปยุตฺตญฺจา”ติปิ, “อุปาทานสมฺปยุตฺตญฺเจว โน จ อุปาทานนฺ”ติปิ.\csromanspacer{} เทฺว สจฺจา น วตฺตพฺพา “อุปาทานา เจว อุปาทานสมฺปยุตฺตา จา”ติปิ, “อุปาทานสมฺปยุตฺตา เจว โน จ อุปาทานา”ติปิ.\csromanspacer{} ทุกฺขสจฺจํ สิยา อุปาทานญฺเจว อุปาทานสมฺปยุตฺตญฺจ, สิยา อุปาทานสมฺปยุตฺตญฺเจว โน จ อุปาทานํ, สิยา น วตฺตพฺพํ “อุปาทานญฺเจว อุปาทานสมฺปยุตฺตญฺจา”ติปิ, “อุปาทานสมฺปยุตฺตญฺเจว โน จ อุปาทานนฺ”ติปิ.\csromanspacer{} เทฺว สจฺจา อุปาทานวิปฺปยุตฺต-อนุปาทานิยา.\csromanspacer{} เทฺว สจฺจา สิยา อุปาทานวิปฺปยุตฺต-อุปาทานิยา, สิยา น วตฺตพฺพา “อุปาทานวิปฺปยุตฺต- อุปาทานิยา”ติปิ, “อุปาทานวิปฺปยุตฺต-อนุปาทานิยา”ติปิ. (11)}

\prose{สมุทยสจฺจํ กิเลโส.\csromanspacer{} เทฺว สจฺจา โน กิเลสา.\csromanspacer{} ทุกฺขสจฺจํ สิยา กิเลโส, สิยา โน กิเลโส.\csromanspacer{} เทฺว สจฺจา สํกิเลสิกา.\csromanspacer{} เทฺว สจฺจา อสํกิเลสิกา.\csromanspacer{} สมุทยสจฺจํ สํกิลิฏฺฐํ.\csromanspacer{} เทฺว สจฺจา อสํกิลิฏฺฐา.\csromanspacer{} ทุกฺขสจฺจํ สิยา สํกิลิฏฺฐํ, สิยา อสํกิลิฏฺฐํ.\csromanspacer{} สมุทยสจฺจํ กิเลสสมฺปยุตฺตํ.\csromanspacer{} เทฺว สจฺจา กิเลสวิปฺปยุตฺตา.\csromanspacer{} ทุกฺขสจฺจํ สิยา กิเลสสมฺปยุตฺตํ, \csromanfolio{126}สิยา กิเลสวิปฺปยุตฺตํ.\csromanspacer{} สมุทยสจฺจํ กิเลโส เจว สํกิเลสิกญฺจ.\csromanspacer{} เทฺว สจฺจา น วตฺตพฺพา “กิเลสา เจว สํกิเลสิกา จา”ติปิ, “สํกิเลสิกา เจว โน จ กิเลสา”ติปิ.\csromanspacer{} ทุกฺขสจฺจํ สิยา กิเลโส เจว สํกิเลสิกญฺจ, สิยา สํกิเลสิกญฺเจว โน จ กิเลโส.\csromanspacer{} สมุทยสจฺจํ กิเลโส เจว สํกิลิฏฺฐญฺจ.\csromanspacer{} เทฺว สจฺจา น วตฺตพฺพา “กิเลสา เจว สํกิลิฏฺฐา จา”ติปิ, “สํกิลิฏฺฐา เจว โน จ กิเลสา”ติปิ.\csromanspacer{} ทุกฺขสจฺจํ สิยา กิเลโส เจว สํกิลิฏฺฐญฺจ, สิยา สํกิลิฏฺฐญฺเจว โน จ กิเลโส, สิยา น วตฺตพฺพํ “กิเลโส เจว สํกิลิฏฺฐญฺจา”ติปิ, “สํกิลิฏฺฐญฺเจว โน จ กิเลโส”ติปิ.\csromanspacer{} สมุทยสจฺจํ กิเลโส เจว กิเลสสมฺปยุตฺตญฺจ.\csromanspacer{} เทฺว สจฺจา น วตฺตพฺพา “กิเลสา เจว กิเลสสมฺปยุตฺตา จา”ติปิ, “กิเลสสมฺปยุตฺตา เจว โน จ กิเลสา”ติปิ.\csromanspacer{} ทุกฺขสจฺจํ สิยา กิเลโส เจว กิเลสสมฺปยุตฺตญฺจ, สิยา กิเลสสมฺปยุตฺตญฺเจว โน จ กิเลโส, สิยา น วตฺตพฺพํ “กิเลโส เจว กิเลสสมฺปยุตฺตญฺจา”ติปิ, “กิเลสสมฺปยุตฺตญฺเจว โน จ กิเลโส”ติปิ.\csromanspacer{} เทฺว สจฺจา กิเลสวิปฺปยุตฺต-อสํกิเลสิกา.\csromanspacer{} สมุทยสจฺจํ น วตฺตพฺพํ “กิเลสวิปฺปยุตฺตสํกิเลสิกนฺ”ติปิ, “กิเลสวิปฺปยุตฺต-อสํกิเลสิกนฺ”ติปิ.\csromanspacer{} ทุกฺขสจฺจํ สิยา กิเลสวิปฺปยุตฺตสํกิเลสิกํ, สิยา น วตฺตพฺพํ “กิเลสวิปฺปยุตฺตสํกิเลสิกนฺ”ติปิ, “กิเลสวิปฺปยุตฺต-อสํกิเลสิกนฺ”ติปิ. (12)}

\prose{เทฺว สจฺจา น ทสฺสเนน ปหาตพฺพา.\csromanspacer{} เทฺว สจฺจา สิยา ทสฺสเนน ปหาตพฺพา, สิยา น ทสฺสเนน ปหาตพฺพา.\csromanspacer{} เทฺว สจฺจา น ภาวนาย ปหาตพฺพา.\csromanspacer{} เทฺว สจฺจา สิยา ภาวนาย ปหาตพฺพา, สิยา น ภาวนาย ปหาตพฺพา.\csromanspacer{} เทฺว สจฺจา น ทสฺสเนน ปหาตพฺพเหตุกา.\csromanspacer{} เทฺว สจฺจา สิยา ทสฺสเนน ปหาตพฺพเหตุกา, สิยา น ทสฺสเนน ปหาตพฺพเหตุกา.\csromanspacer{} เทฺว สจฺจา น ภาวนาย ปหาตพฺพเหตุกา.\csromanspacer{} เทฺว สจฺจา สิยา ภาวนาย ปหาตพฺพเหตุกา, สิยา น ภาวนาย ปหาตพฺพเหตุกา.\csromanspacer{} สมุทยสจฺจํ สวิตกฺกํ.\csromanspacer{} นิโรธสจฺจํ อวิตกฺกํ.\csromanspacer{} เทฺว สจฺจา สิยา สวิตกฺกา, สิยา อวิตกฺกา.\csromanspacer{} สมุทยสจฺจํ สวิจารํ.\csromanspacer{} นิโรธสจฺจํ อวิจารํ.\csromanspacer{} เทฺว สจฺจา สิยา สวิจารา, สิยา อวิจารา.\csromanspacer{} นิโรธสจฺจํ อปฺปีติกํ.\csromanspacer{} ตีณิ สจฺจานิ สิยา สปฺปีติกา, สิยา อปฺปีติกา.\csromanspacer{} นิโรธสจฺจํ น ปีติสหคตํ.\csromanspacer{} ตีณิ สจฺจานิ สิยา ปีติสหคตา, สิยา น ปีติสหคตา.}

\vspace{\tipitakaparskip}
\csromancenter{สจฺจวิภงฺค นิโรธสจฺจํ น สุขสหคตํ.\csromanspacer{} ตีณิ สจฺจานิ สิยา สุขสหคตา, สิยา น สุขสหคตา.\csromanspacer{} นิโรธสจฺจํ น อุเปกฺขาสหคตํ.\csromanspacer{} ตีณิ สจฺจานิ สิยา อุเปกฺขาสหคตา, สิยา น อุเปกฺขาสหคตา.}
\prose{\csromanfolio{127}สมุทยสจฺจํ กามาวจรํ.\csromanspacer{} เทฺว สจฺจา น กามาวจรา.\csromanspacer{} ทุกฺขสจฺจํ สิยา กามาวจรํ, สิยา น กามาวจรํ.\csromanspacer{} ตีณิ สจฺจานิ น รูปาวจรา.\csromanspacer{} ทุกฺขสจฺจํ สิยา รูปาวจรํ, สิยา น รูปาวจรํ.\csromanspacer{} ตีณิ สจฺจานิ น อรูปาวจรา.\csromanspacer{} ทุกฺขสจฺจํ สิยา อรูปาวจรํ, สิยา น อรูปาวจรํ.\csromanspacer{} เทฺว สจฺจา ปริยาปนฺนา.\csromanspacer{} เทฺว สจฺจา อปริยาปนฺนา.\csromanspacer{} มคฺคสจฺจํ นิยฺยานิกํ.\csromanspacer{} ตีณิ สจฺจานิ อนิยฺยานิกา.\csromanspacer{} มคฺคสจฺจํ นิยตํ.\csromanspacer{} นิโรธสจฺจํ อนิยตํ.\csromanspacer{} เทฺว สจฺจา สิยา นิยตา, สิยา อนิยตา.\csromanspacer{} เทฺว สจฺจา ส- อุตฺตรา.\csromanspacer{} เทฺว สจฺจา อนุตฺตรา.\csromanspacer{} สมุทยสจฺจํ สรณํ.\csromanspacer{} เทฺว สจฺจา อรณา.\csromanspacer{} ทุกฺขสจฺจํ สิยา สรณํ, สิยา อรณนฺติ. (13)}

\vspace{\tipitakaparskip}
\csromancenter{ปญฺหาปุจฺฉกํ.}
\nitthitam{\textbf{สจฺจวิภงฺโค นิฏฺฐิโต.}}
\csromanmatikaanchor{mtk.39}
\csromantocmarknopage{paragraph}{5. อินฺทฺริยวิภงฺค}{mtk.39}
\bookmark[dest={mtk.39},level=4]{5. อินฺทฺริยวิภงฺค}
\csromanheader{4}{5. อินฺทฺริยวิภงฺค}
\csromanmatikaanchor{mtk.40}
\csromantocmarkref{subparagraph}{1. อภิธมฺมภาชนีย}{mtk.40}
\bookmark[dest={mtk.40},level=5]{1. อภิธมฺมภาชนีย}
\csromanheader{5}{1. อภิธมฺมภาชนีย}
\proseitem{219}{\csromanfolio{128}พาวีสตินฺทฺริยานิ จกฺขุนฺทฺริยํ โสตินฺทฺริยํ ฆานินฺทฺริยํ ชิวฺหินฺทฺริยํ กายินฺทฺริยํ มนินฺทฺริยํ อิตฺถินฺทฺริยํ ปุริสินฺทฺริยํ ชีวิตินฺทฺริยํ สุขินฺทฺริยํ ทุกฺขินฺทฺริยํ โสมนสฺสินฺทฺริยํ โทมนสฺสินฺทฺริยํ อุเปกฺขินฺทฺริยํ สทฺธินฺทฺริยํ วีริยินฺทฺริยํ\footnote{วิริยินฺทฺริยํ (สี, สฺยา)} สตินฺทฺริยํ สมาธินฺทฺริยํ ปญฺญินฺทฺริยํ อนญฺญาตญฺญสฺสามีตินฺทฺริยํ อญฺญินฺทฺริยํ อญฺญาตาวินฺทฺริยํ.}

\proseitem{220}{ตตฺถ กตมํ จกฺขุนฺทฺริยํ, ยํ จกฺขุ จตุนฺนํ มหาภูตานํ อุปาทาย ปสาโท ฯเปฯ สุญฺโญ คาโมเปโส.\csromanspacer{} อิทํ วุจฺจติ จกฺขุนฺทฺริยํ. (1)}

\prose{ตตฺถ กตมํ โสตินฺทฺริยํ ฯเปฯ ฆานินฺทฺริยํ ฯเปฯ ชิวฺหินฺทฺริยํ ฯเปฯ กายินฺทฺริยํ โย กาโย จตุนฺนํ มหาภูตานํ อุปาทาย ปสาโท ฯเปฯ สุญฺโญ คาโมเปโส.\csromanspacer{} อิทํ วุจฺจติ กายินฺทฺริยํ. (5)}

\prose{ตตฺถ กตมํ มนินฺทฺริยํ, เอกวิเธน มนินฺทฺริยํ\csromandash{} ผสฺสสมฺปยุตฺตํ ทุวิเธน มนินฺทฺริยํ\csromandash{}อตฺถิ สเหตุกํ, อตฺถิ อเหตุกํ.\csromanspacer{} ติวิเธน มนินฺทฺริยํ\csromandash{}อตฺถิ กุสลํ, อตฺถิ อกุสลํ, อตฺถิ อพฺยากตํ.\csromanspacer{} จตุพฺพิเธน มนินฺทฺริยํ\csromandash{}อตฺถิ กามาวจรํ, อตฺถิ รูปาวจรํ, อตฺถิ อรูปาวจรํ, อตฺถิ อปริยาปนฺนํ.\csromanspacer{} ปญฺจวิเธน มนินฺทฺริยํ\csromandash{}อตฺถิ สุขินฺทฺริยสมฺปยุตฺตํ, อตฺถิ ทุกฺขินฺทฺริยสมฺปยุตฺตํ, อตฺถิ โสมนสฺสินฺทฺริยสมฺปยุตฺตํ, อตฺถิ โทมนสฺสินฺทฺริยสมฺปยุตฺตํ, อตฺถิ อุเปกฺขินฺทฺริยสมฺปยุตฺตํ.\csromanspacer{} ฉพฺพิเธน มนินฺทฺริยํ\csromandash{}จกฺขุวิญฺญาณํ ฯเปฯ มโนวิญฺญาณํ.\csromanspacer{} เอวํ ฉพฺพิเธน มนินฺทฺริยํ.}

\prose{สตฺตวิเธน มนินฺทฺริยํ\csromandash{}จกฺขุวิญฺญาณํ ฯเปฯ กายวิญฺญาณํ, มโนธาตุ, มโนวิญฺญาณธาตุ.\csromanspacer{} เอวํ สตฺตวิเธน มนินฺทฺริยํ.}

\prose{อฏฺฐวิเธน มนินฺทฺริยํ\csromandash{}จกฺขุวิญฺญาณํ ฯเปฯ กายวิญฺญาณํ อตฺถิ สุขสหคตํ, อตฺถิ ทุกฺขสหคตํ, มโนธาตุ, มโนวิญฺญาณธาตุ.\csromanspacer{} เอวํ อฏฺฐวิเธน มนินฺทฺริยํ.}

\csromanheader{2}{5. อินฺทฺริยวิภงฺค}
\prose{\csromanfolio{129}นววิเธน มนินฺทฺริยํ\csromandash{}จกฺขุวิญฺญาณํ ฯเปฯ กายวิญฺญาณํ, มโนธาตุ, มโนวิญฺญาณธาตุ อตฺถิ กุสลา, อตฺถิ อกุสลา, อตฺถิ อพฺยากตา.\csromanspacer{} เอวํ นววิเธน มนินฺทฺริยํ.}

\prose{ทสวิเธน มนินฺทฺริยํ\csromandash{}จกฺขุวิญฺญาณํ ฯเปฯ กายวิญฺญาณํ อตฺถิ สุขสหคตํ, อตฺถิ ทุกฺขสหคตํ, มโนธาตุ, มโนวิญฺญาณธาตุ อตฺถิ กุสลา, อตฺถิ อกุสลา, อตฺถิ อพฺยากตา.\csromanspacer{} เอวํ ทสวิเธน มนินฺทฺริยํ ฯเปฯ.\csromanspacer{} เอวํ พหุวิเธน มนินฺทฺริยํ.\csromanspacer{} อิทํ วุจฺจติ มนินฺทฺริยํ. (6)}

\prose{ตตฺถ กตมํ อิตฺถินฺทฺริยํ, ยํ อิตฺถิยา อิตฺถิลิงฺคํ อิตฺถินิมิตฺตํ อิตฺถิกุตฺตํ อิตฺถากปฺโป อิตฺถตฺตํ อิตฺถิภาโว, อิทํ วุจฺจติ อิตฺถินฺทฺริยํ. (7)}

\prose{ตตฺถ กตมํ ปุริสินฺทฺริยํ, ยํ ปุริสสฺส ปุริสลิงฺคํ ปุริสนิมิตฺตํ ปุริสกุตฺตํ ปุริสากปฺโป ปุริสตฺตํ ปุริสภาโว, อิทํ วุจฺจติ ปุริสินฺทฺริยํ. (8)}

\prose{ตตฺถ กตมํ ชีวิตินฺทฺริยํ, ชีวิตินฺทฺริยํ ทุวิเธน อตฺถิ รูปชีวิตินฺทฺริยํ, อตฺถิ อรูปชีวิตินฺทฺริยํ.}

\prose{ตตฺถ กตมํ รูปชีวิตินฺทฺริยํ, โย เตสํ รูปีนํ ธมฺมานํ อายุ ฐิติ ยปนา ยาปนา อิริยนา วตฺตนา ปาลนา ชีวิตํ ชีวิตินฺทฺริยํ, อิทํ วุจฺจติ รูปชีวิตินฺทฺริยํ.}

\prose{ตตฺถ กตมํ อรูปชีวิตินฺทฺริยํ, โย เตสํ อรูปีนํ ธมฺมานํ อายุ ฐิติ ยปนา ยาปนา อิริยนา วตฺตนา ปาลนา ชีวิตํ ชีวิตินฺทฺริยํ, อิทํ วุจฺจติ อรูปชีวิตินฺทฺริยํ.\csromanspacer{} อิทํ วุจฺจติ ชีวิตินฺทฺริยํ. (9)}

\prose{ตตฺถ กตมํ สุขินฺทฺริยํ, ยํ กายิกํ สาตํ กายิกํ สุขํ กายสมฺผสฺสชํ สาตํ สุขํ เวทยิตํ กายสมฺผสฺสชา สาตา สุขา เวทนา, อิทํ วุจฺจติ สุขินฺทฺริยํ. (10)}

\prose{ตตฺถ กตมํ ทุกฺขินฺทฺริยํ, ยํ กายิกํ อสาตํ กายิยํ ทุกฺขํ กายสมฺผสฺสชํ อสาตํ ทุกฺขํ เวทยิตํ กายสมฺผสฺสชา อสาตา ทุกฺขา เวทนา, อิทํ วุจฺจติ ทุกฺขินฺทฺริยํ. (11)}

\prose{ตตฺถ กตมํ โสมนสฺสินฺทฺริยํ, ยํ เจตสิกํ สาตํ เจตสิกํ สุขํ เจโตสมฺผสฺสชํ สาตํ สุขํ เวทยิตํ เจโตสมฺผสฺสชา สาตา สุขา เวทนา, อิทํ วุจฺจติ โสมนสฺสินฺทฺริยํ. (12)}

\prose{\csromanfolio{130}ตตฺถ กตมํ โทมนสฺสินฺทฺริยํ, ยํ เจตสิกํ อสาตํ เจตสิกํ ทุกฺขํ เจโตสมฺผสฺสชํ อสาตํ ทุกฺขํ เวทยิตํ เจโตสมฺผสฺสชา อสาตา ทุกฺขา เวทนา, อิทํ วุจฺจติ โทมนสฺสินฺทฺริยํ. (13)}

\prose{ตตฺถ กตมํ อุเปกฺขินฺทฺริยํ, ยํ เจตสิกํ เนว สาตํ นาสาตํ เจโตสมฺผสฺสชํ อทุกฺขมสุขํ เวทยิตํ เจโตสมฺผสฺสชา อทุกฺขมสุขา เวทนา, อิทํ วุจฺจติ อุเปกฺขินฺทฺริยํ. (14)}

\prose{ตตฺถ กตมํ สทฺธินฺทฺริยํ, ยา สทฺธา สทฺทหนา โอกปฺปนา อภิปฺปสาโท สทฺธา สทฺธินฺทฺริยํ สทฺธาพลํ, อิทํ วุจฺจติ สทฺธินฺทฺริยํ. (15)}

\prose{ตตฺถ กตมํ วีริยินฺทฺริยํ, โย เจตสิโก วีริยารมฺโภ นิกฺกโม ปรกฺกโม อุยฺยาโม วายาโม อุสฺสาโห อุสฺโสฬฺหี ถาโม ฐิติ อสิถิลปรกฺกมตา อนิกฺขิตฺตฉนฺทตา อนิกฺขิตฺตธุรตา ธุรสมฺปคฺคาโห วีริยํ วีริยินฺทฺริยํ วีริยพลํ, อิทํ วุจฺจติ วีริยินฺทฺริยํ. (16)}

\prose{ตตฺถ กตมํ สตินฺทฺริยํ, ยา สติ อนุสฺสติ ปฏิสฺสติ สติ สรณตา ธารณตา อปิลาปนตา อสมฺมุสฺสนตา สติ สตินฺทฺริยํ สติพลํ สมฺมาสติ, อิทํ วุจฺจติ สตินฺทฺริยํ. (17)}

\prose{ตตฺถ กตมํ สมาธินฺทฺริยํ, ยา จิตฺตสฺส ฐิติ สณฺฐิติ อวฏฺฐิติ อวิสาหาโร อวิกฺเขโป อวิสาหฏมานสตา สมโถ สมาธินฺทฺริยํ สมาธิพลํ สมฺมาสมาธิ, อิทํ วุจฺจติ สมาธินฺทฺริยํ. (18)}

\prose{ตตฺถ กตมํ ปญฺญินฺทฺริยํ, ยา ปญฺญา ปชานนา ฯเปฯ อโมโห ธมฺมวิจโย สมฺมาทิฏฺฐิ, อิทํ วุจฺจติ ปญฺญินฺทฺริยํ. (19)}

\prose{ตตฺถ กตมํ อนญฺญาตญฺญสฺสามีตินฺทฺริยํ, ยา เตสํ ธมฺมานํ อนญฺญาตานํ อทิฏฺฐานํ อปฺปตฺตานํ อวิทิตานํ อสจฺฉิกตานํ สจฺฉิกิริยาย ปญฺญา ปชานนา ฯเปฯ อโมโห ธมฺมวิจโย สมฺมาทิฏฺฐิ ธมฺมวิจยสมฺโพชฺฌงฺโค มคฺคงฺคํ มคฺคปริยาปนฺนํ, อิทํ วุจฺจติ อนญฺญาตญฺญสฺสามีตินฺทฺริยํ. (20)}

\prose{ตตฺถ กตมํ อญฺญินฺทฺริยํ, ยา เตสํ ธมฺมานํ ญาตานํ ทิฏฺฐานํ ปตฺตานํ วิทิตานํ สจฺฉิกตานํ สจฺฉิกิริยาย ปญฺญา ปชานนา ฯเปฯ อโมโห ธมฺมวิจโย สมฺมาทิฏฺฐิ ธมฺมวิจยสมฺโพชฺฌงฺโค มคฺคงฺคํ มคฺคปริยาปนฺนํ, อิทํ วุจฺจติ อญฺญินฺทฺริยํ. (21)}

\csromanheader{2}{5. อินฺทฺริยวิภงฺค}
\prose{\csromanfolio{131}ตตฺถ กตมํ อญฺญาตาวินฺทฺริยํ, ยา เตสํ ธมฺมานํ อญฺญาตาวีนํ ทิฏฺฐานํ ปตฺตานํ วิทิตานํ สจฺฉิกตานํ สจฺฉิกิริยาย ปญฺญา ปชานนา ฯเปฯ อโมโห ธมฺมวิจโย สมฺมาทิฏฺฐิ ธมฺมวิจยสมฺโพชฺฌงฺโค มคฺคงฺคํ มคฺคปริยาปนฺนํ, อิทํ วุจฺจติ อญฺญาตาวินฺทฺริยํ. (22)}

\csromanheader{2}{อภิธมฺมภาชนียํ.}
\csromansectionrule
\csromanmatikaanchor{mtk.41}
\csromantocmarkref{subparagraph}{2. ปญฺหาปุจฺฉก}{mtk.41}
\bookmark[dest={mtk.41},level=5]{2. ปญฺหาปุจฺฉก}
\csromanheader{5}{2. ปญฺหาปุจฺฉก}
\proseitem{221}{พาวีสตินฺทฺริยานิ จกฺขุนฺทฺริยํ โสตินฺทฺริยํ ฆานินฺทฺริยํ ชิวฺหินฺทฺริยํ กายินฺทฺริยํ มนินฺทฺริยํ อิตฺถินฺทฺริยํ ปุริสินฺทฺริยํ ชีวิตินฺทฺริยํ สุขินฺทฺริยํ ทุกฺขินฺทฺริยํ โสมนสฺสินฺทฺริยํ โทมนสฺสินฺทฺริยํ อุเปกฺขินฺทฺริยํ สทฺธินฺทฺริยํ วีริยินฺทฺริยํ สตินฺทฺริยํ สมาธินฺทฺริยํ ปญฺญินฺทฺริยํ อนญฺญาตญฺญสฺสามีตินฺทฺริยํ อญฺญินฺทฺริยํ อญฺญาตาวินฺทฺริยํ.}

\proseitem{222}{พาวีสตีนํ อินฺทฺริยานํ กติ กุสลา, กติ อกุสลา, กติ อพฺยากตา ฯเปฯ กติ สรณา, กติ อรณา.}

\csromanmatikaanchor{mtk.42}
\csromantocmarkref{subparagraph}{1. ติก}{mtk.42}
\bookmark[dest={mtk.42},level=5]{1. ติก}
\csromanheader{6}{1. ติก}
\proseitem{223}{ทสินฺทฺริยา อพฺยากตา.\csromanspacer{} โทมนสฺสินฺทฺริยํ อกุสลํ.\csromanspacer{} อนญฺญาตญฺญสฺสามีตินฺทฺริยํ กุสลํ, จตฺตารินฺทฺริยา สิยา กุสลา, สิยา อพฺยากตา.\csromanspacer{} ฉ อินฺทฺริยา สิยา กุสลา, สิยา อกุสลา, สิยา อพฺยากตา.}

\prose{ทฺวาทสินฺทฺริยา น วตฺตพฺพา “สุขาย เวทนาย สมฺปยุตฺตา”ติปิ, “ทุกฺขาย เวทนาย สมฺปยุตฺตา”ติปิ, “อทุกฺขมสุขาย เวทนาย สมฺปยุตฺตา”ติปิ.\csromanspacer{} ฉ อินฺทฺริยา สิยา สุขาย เวทนาย สมฺปยุตฺตา, สิยา อทุกฺขมสุขาย เวทนาย สมฺปยุตฺตา.\csromanspacer{} ตีณินฺทฺริยา สิยา สุขาย เวทนาย สมฺปยุตฺตา, สิยา ทุกฺขาย เวทนาย สมฺปยุตฺตา, สิยา อทุกฺขมสุขาย เวทนาย สมฺปยุตฺตา.\csromanspacer{} ชีวิตินฺทฺริยํ สิยา สุขาย เวทนาย สมฺปยุตฺตํ, สิยา ทุกฺขาย เวทนาย สมฺปยุตฺตํ, สิยา อทุกฺขมสุขาย เวทนาย สมฺปยุตฺตํ, สิยา น วตฺตพฺพํ “สุขาย เวทนาย สมฺปยุตฺตนฺ”ติปิ, “ทุกฺขาย เวทนาย สมฺปยุตฺตนฺ”ติปิ, “อทุกฺขมสุขาย เวทนาย สมฺปยุตฺตนฺ”ติปิ.}

\prose{สตฺตินฺทฺริยา เนววิปากนวิปากธมฺมธมฺมา.\csromanspacer{} ตีณินฺทฺริยา วิปากา.\csromanspacer{} ทฺวินฺทฺริยา วิปากธมฺมธมฺมา.\csromanspacer{} อญฺญินฺทฺริยํ สิยา วิปากํ, สิยา วิปากธมฺมธมฺมํ.\csromanspacer{} นวินฺทฺริยา \csromanfolio{132}สิยา วิปากา, สิยา วิปากธมฺมธมฺมา, สิยา เนววิปากนวิปากธมฺมธมฺมา.\csromanspacer{} นวินฺทฺริยา อุปาทินฺนุปาทานิยา, โทมนสฺสินฺทฺริยํ อนุปาทินฺนุปาทานิยํ.\csromanspacer{} ตีณินฺทฺริยา อนุปาทินฺน- อนุปาทานิยา.\csromanspacer{} นวินฺทฺริยา สิยา อุปาทินฺนุปาทานิยา, สิยา อนุปาทินฺนุปาทานิยา, สิยา อนุปาทินฺน-อนุปาทานิยา.}

\prose{นวินฺทฺริยา อสํกิลิฏฺฐสํกิเลสิกา.\csromanspacer{} โทมนสฺสินฺทฺริยํ สํกิลิฏฺฐสํกิเลสิกํ.\csromanspacer{} ตีณินฺทฺริยา อสํกิลิฏฺฐ-อสํกิเลสิกา.\csromanspacer{} ตีณินฺทฺริยา สิยา อสํกิลิฏฺฐสํกิเลสิกา, สิยา อสํกิลิฏฺฐ-อสํกิเลสิกา.\csromanspacer{} ฉ อินฺทฺริยา สิยา สํกิลิฏฺฐสํกิเลสิกา, สิยา อสํกิลิฏฺฐสํกิเลสิกา, สิยา อสํกิลิฏฺฐ- อสํกิเลสิกา.\csromanspacer{} นวินฺทฺริยา อวิตกฺก-อวิจารา.\csromanspacer{} โทมนสฺสินฺทฺริยํ สวิตกฺกสวิจารํ.\csromanspacer{} อุเปกฺขินฺทฺริยํ สิยา สวิตกฺกสวิจารํ, สิยา อวิตกฺก-อวิจารํ.\csromanspacer{} เอกาทสินฺทฺริยา สิยา สวิตกฺกสวิจารา, สิยา อวิตกฺกวิจารมตฺตา, สิยา อวิตกฺก-อวิจารา.}

\prose{เอกาทสินฺทฺริยา น วตฺตพฺพา “ปีติสหคตา”ติปิ, “สุขสหคตา”ติปิ, “อุเปกฺขาสหคตา”ติปิ.\csromanspacer{} โสมนสฺสินฺทฺริยํ สิยา ปีติสหคตํ, น สุขสหคตํ, น อุเปกฺขาสหคตํ, สิยา น วตฺตพฺพํ “ปีติสหคตนฺ”ติ.\csromanspacer{} ฉ อินฺทฺริยา สิยา ปีติสหคตา, สิยา สุขสหคตา, สิยา อุเปกฺขาสหคตา.\csromanspacer{} จตฺตารินฺทฺริยา สิยา ปีติสหคตา, สิยา สุขสหคตา, สิยา อุเปกฺขาสหคตา, สิยา น วตฺตพฺพา “ปีติสหคตา”ติปิ, “สุขสหคตา”ติปิ, “อุเปกฺขาสหคตา”ติปิ.}

\prose{ปนฺนรสินฺทฺริยา เนว ทสฺสเนน น ภาวนาย ปหาตพฺพา.\csromanspacer{} โทมนสฺสินฺทฺริยํ สิยา ทสฺสเนน ปหาตพฺพํ, สิยา ภาวนาย ปหาตพฺพํ.\csromanspacer{} ฉ อินฺทฺริยา สิยา ทสฺสเนน ปหาตพฺพา, สิยา ภาวนาย ปหาตพฺพา, สิยา เนว ทสฺสเนน น ภาวนาย ปหาตพฺพา.\csromanspacer{} ปนฺนรสินฺทฺริยา เนว ทสฺสเนน น ภาวนาย ปหาตพฺพเหตุกา.\csromanspacer{} โทมนสฺสินฺทฺริยํ สิยา ทสฺสเนน ปหาตพฺพเหตุกํ, สิยา ภาวนาย ปหาตพฺพเหตุกํ.\csromanspacer{} ฉ อินฺทฺริยา สิยา ทสฺสเนน ปหาตพฺพเหตุกา, สิยา ภาวนาย ปหาตพฺพเหตุกา, สิยา เนว ทสฺสเนน น ภาวนาย ปหาตพฺพเหตุกา.}

\prose{ทสินฺทฺริยา เนวาจยคามินาปจยคามิโน.\csromanspacer{} โทมนสฺสินฺทฺริยํ อาจยคามิ.\csromanspacer{} อนญฺญาตญฺญสฺสามีตินฺทฺริยํ อปจยคามิ.\csromanspacer{} อญฺญินฺทฺริยํ สิยา อปจยคามิ, สิยา เนวาจยคามินาปจยคามิ.\csromanspacer{} นวินฺทฺริยา สิยา อาจยคามิโน,}

\vspace{\tipitakaparskip}
\csromancenter{อินฺทฺริยวิภงฺค สิยา อปจยคามิโน, สิยา เนวาจยคามินาปจยคามิโน.\csromanspacer{} ทสินฺทฺริยา เนวเสกฺขนาเสกฺขา.\csromanspacer{} ทฺวินฺทฺริยา เสกฺขา.\csromanspacer{} อญฺญาตาวินฺทฺริยํ อเสกฺขํ.\csromanspacer{} นวินฺทฺริยา สิยา เสกฺขา, สิยา อเสกฺขา, สิยา เนวเสกฺขนาเสกฺขา.}
\prose{\csromanfolio{133}ทสินฺทฺริยา ปริตฺตา.\csromanspacer{} ตีณินฺทฺริยา อปฺปมาณา.\csromanspacer{} นวินฺทฺริยา สิยา ปริตฺตา, สิยา มหคฺคตา, สิยา อปฺปมาณา.\csromanspacer{} สตฺตินฺทฺริยา อนารมฺมณา.\csromanspacer{} ทฺวินฺทฺริยา ปริตฺตารมฺมณา.\csromanspacer{} ตีณินฺทฺริยา อปฺปมาณารมฺมณา.\csromanspacer{} โทมนสฺสินฺทฺริยํ สิยา ปริตฺตารมฺมณํ, สิยา มหคฺคตารมฺมณํ, น อปฺปมาณารมฺมณํ, สิยา น วตฺตพฺพํ “ปริตฺตารมฺมณนฺ”ติปิ, “มหคฺคตารมฺมณนฺ”ติปิ.\csromanspacer{} นวินฺทฺริยา สิยา ปริตฺตารมฺมณา, สิยา มหคฺคตารมฺมณา, สิยา อปฺปมาณารมฺมณา, สิยา น วตฺตพฺพา “ปริตฺตารมฺมณา”ติปิ, “มหคฺคตารมฺมณา”ติปิ, “อปฺปมาณารมฺมณา”ติปิ.}

\prose{นวินฺทฺริยา มชฺฌิมา.\csromanspacer{} โทมนสฺสินฺทฺริยํ หีนํ.\csromanspacer{} ตีณินฺทฺริยา ปณีตา.\csromanspacer{} ตีณินฺทฺริยา สิยา มชฺฌิมา.\csromanspacer{} สิยา ปณีตา.\csromanspacer{} ฉ อินฺทฺริยา สิยา หีนา, สิยา มชฺฌิมา.\csromanspacer{} สิยา ปณีตา.\csromanspacer{} ทสินฺทฺริยา อนิยตา.\csromanspacer{} อนญฺญาตญฺญสฺสามีตินฺทฺริยํ สมฺมตฺตนิยตํ.\csromanspacer{} จตฺตารินฺทฺริยา สิยา สมฺมตฺตนิยตา, สิยา อนิยตา.\csromanspacer{} โทมนสฺสินฺทฺริยํ สิยา มิจฺฉตฺตนิยตํ, สิยา อนิยตํ.\csromanspacer{} ฉ อินฺทฺริยา สิยา มิจฺฉตฺตนิยตา, สิยา สมฺมตฺตนิยตา, สิยา อนิยตา.\csromanspacer{} สตฺตินฺทฺริยา อนารมฺมณา.\csromanspacer{} จตฺตารินฺทฺริยา น วตฺตพฺพา “มคฺคารมฺมณา”ติปิ, “มคฺคเหตุกา”ติปิ, “มคฺคาธิปติโน”ติปิ.\csromanspacer{} อนญฺญาตญฺญสฺสามีตินฺทฺริยํ น มคฺคารมฺมณํ, สิยา มคฺคเหตุกํ, สิยา มคฺคาธิปติ, สิยา น วตฺตพฺพํ “มคฺคเหตุกนฺ”ติปิ, “มคฺคาธิปตี”ติปิ.\csromanspacer{} อญฺญินฺทฺริยํ น มคฺคารมฺมณํ สิยา มคฺคเหตุกํ, สิยา มคฺคาธิปติ, สิยา น วตฺตพฺพํ “มคฺคเหตุกนฺ”ติปิ, “มคฺคาธิปตี”ติปิ.\csromanspacer{} นวินฺทฺริยา สิยา มคฺคารมฺมณา, สิยา มคฺคเหตุกา, สิยา มคฺคาธิปติโน, สิยา น วตฺตพฺพา “มคฺคารมฺมณา”ติปิ, “มคฺคเหตุกา”ติปิ, “มคฺคาธิปติโน”ติปิ.}

\prose{ทสินฺทฺริยา สิยา อุปฺปนฺนา, สิยา อุปฺปาทิโน, น วตฺตพฺพา “อนุปฺปนฺนา”ติ.\csromanspacer{} ทฺวินฺทฺริยา สิยา อุปฺปนฺนา, สิยา อนุปฺปนฺนา, น วตฺตพฺพา “อุปฺปาทิโน”ติ.\csromanspacer{} ทสินฺทฺริยา สิยา อุปฺปนฺนา, สิยา อนุปฺปนฺนา, สิยา อุปฺปาทิโน, สิยา อตีตา, สิยา อนาคตา, สิยา ปจฺจุปฺปนฺนา.\csromanspacer{} สตฺตินฺทฺริยา อนารมฺมณา.\csromanspacer{} ทฺวินฺทฺริยา ปจฺจุปฺปนฺนารมฺมณา.\csromanspacer{} ตีณินฺทฺริยา น วตฺตพฺพา “อตีตารมฺมณา”ติปิ, \csromanfolio{134}“อนาคตารมฺมณา”ติปิ, “ปจฺจุปฺปนฺนารมฺมณา”ติปิ.\csromanspacer{} ทสินฺทฺริยา สิยา อตีตารมฺมณา, สิยา อนาคตารมฺมณา, สิยา ปจฺจุปฺปนฺนารมฺมณา, สิยา น วตฺตพฺพา “อตีตารมฺมณา”ติปิ, “อนาคตารมฺมณา”ติปิ, “ปจฺจุปฺปนฺนารมฺมณา”ติปิ, สิยา อชฺฌตฺตา, สิยา พหิทฺธา, สิยา อชฺฌตฺตพหิทฺธา.\csromanspacer{} สตฺตินฺทฺริยา อนารมฺมณา.\csromanspacer{} ตีณินฺทฺริยา พหิทฺธารมฺมณา.\csromanspacer{} จตฺตารินฺทฺริยา สิยา อชฺฌตฺตารมฺมณา, สิยา พหิทฺธารมฺมณา, สิยา อชฺฌตฺตพหิทฺธารมฺมณา.\csromanspacer{} อฏฺฐินฺทฺริยา สิยา อชฺฌตฺตารมฺมณา, สิยา พหิทฺธารมฺมณา, สิยา อชฺฌตฺตพหิทฺธารมฺมณา, สิยา น วตฺตพฺพา “อชฺฌตฺตารมฺมณา”ติปิ, “พหิทฺธารมฺมณา”ติปิ, “อชฺฌตฺตพหิทฺธารมฺมณา”ติปิ.\csromanspacer{} ปญฺจินฺทฺริยา อนิทสฺสนสปฺปฏิฆา.\csromanspacer{} สตฺตรสินฺทฺริยา อนิทสฺสน-อปฺปฏิฆา.}

\csromanmatikaanchor{mtk.43}
\csromantocmarkref{subparagraph}{2. ทุก}{mtk.43}
\bookmark[dest={mtk.43},level=5]{2. ทุก}
\csromanheader{6}{2. ทุก}
\proseitem{224}{จตฺตารินฺทฺริยา เหตู.\csromanspacer{} อฏฺฐารสินฺทฺริยา น เหตู.\csromanspacer{} สตฺตินฺทฺริยา สเหตุกา.\csromanspacer{} นวินฺทฺริยา อเหตุกา.\csromanspacer{} ฉ อินฺทฺริยา สิยา สเหตุกา, สิยา อเหตุกา.\csromanspacer{} สตฺตินฺทฺริยา เหตุสมฺปยุตฺตา.\csromanspacer{} นวินฺทฺริยา เหตุวิปฺปยุตฺตา.\csromanspacer{} ฉ อินฺทฺริยา สิยา เหตุสมฺปยุตฺตา, สิยา เหตุวิปฺปยุตฺตา.\csromanspacer{} จตฺตารินฺทฺริยา เหตู เจว สเหตุกา จ.\csromanspacer{} นวินฺทฺริยา น วตฺตพฺพา “เหตู เจว สเหตุกา จา”ติปิ, “สเหตุกา เจว น จ เหตู”ติปิ.\csromanspacer{} ตีณินฺทฺริยา น วตฺตพฺพา “เหตู เจว สเหตุกา จา”ติ, สเหตุกา เจว น จ เหตู.\csromanspacer{} ฉ อินฺทฺริยา น วตฺตพฺพา “เหตู เจว สเหตุกา จา”ติ, สิยา สเหตุกา เจว น จ เหตู, สิยา น วตฺตพฺพา “สเหตุกา เจว น จ เหตู”ติ.}

\prose{จตฺตารินฺทฺริยา เหตู เจว เหตุสมฺปยุตฺตา จ.\csromanspacer{} นวินฺทฺริยา น วตฺตพฺพา “เหตู เจว เหตุสมฺปยุตฺตา จา”ติปิ, “เหตุสมฺปยุตฺตา เจว น จ เหตู”ติปิ.\csromanspacer{} ตีณินฺทฺริยา น วตฺตพฺพา “เหตู เจว เหตุสมฺปยุตฺตา จา”ติ, เหตุสมฺปยุตฺตา เจว น จ เหตู.\csromanspacer{} ฉ อินฺทฺริยา น วตฺตพฺพา “เหตู เจว เหตุสมฺปยุตฺตา จา”ติ, สิยา เหตุสมฺปยุตฺตา เจว น จ เหตู, สิยา น วตฺตพฺพา “เหตุสมฺปยุตฺตา เจว น จ เหตู”ติ.}

\prose{นวินฺทฺริยา น เหตู อเหตุกา.\csromanspacer{} ตีณินฺทฺริยา น เหตู สเหตุกา.\csromanspacer{} จตฺตารินฺทฺริยา น วตฺตพฺพา “น เหตู สเหตุกา”ติปิ, “น เหตู}

\vspace{\tipitakaparskip}
\csromancenter{อินฺทฺริยวิภงฺค อเหตุกา”ติปิ.\csromanspacer{} ฉ อินฺทฺริยา สิยา น เหตู สเหตุกา, สิยา น เหตู อเหตุกา. (1)}
\prose{\csromanfolio{135}สปฺปจฺจยา, สงฺขตา, อนิทสฺสนา.\csromanspacer{} ปญฺจินฺทฺริยา สปฺปฏิฆา.\csromanspacer{} สตฺตรสินฺทฺริยา อปฺปฏิฆา.\csromanspacer{} สตฺตินฺทฺริยา รูปา.\csromanspacer{} จุทฺทสินฺทฺริยา อรูปา.\csromanspacer{} ชีวิตินฺทฺริยํ สิยา รูปํ, สิยา อรูปํ.\csromanspacer{} ทสินฺทฺริยา โลกิยา.\csromanspacer{} ตีณินฺทฺริยา โลกุตฺตรา.\csromanspacer{} นวินฺทฺริยา สิยา โลกิยา, สิยา โลกุตฺตรา, เกนจิ วิญฺเญยฺยา, เกนจิ น วิญฺเญยฺยา. (2)}

\prose{โน อาสวา.\csromanspacer{} ทสินฺทฺริยา สาสวา.\csromanspacer{} ตีณินฺทฺริยา อนาสวา.\csromanspacer{} นวินฺทฺริยา สิยา สาสวา, สิยา อนาสวา.\csromanspacer{} ปนฺนรสินฺทฺริยา อาสววิปฺปยุตฺตา.\csromanspacer{} โทมนสฺสินฺทฺริยํ อาสวสมฺปยุตฺตํ.\csromanspacer{} ฉ อินฺทฺริยา สิยา อาสวสมฺปยุตฺตา, สิยา อาสววิปฺปยุตฺตา.\csromanspacer{} ทสินฺทฺริยา น วตฺตพฺพา “อาสวา เจว สาสวา จา”ติ, สาสวา เจว โน จ อาสวา.\csromanspacer{} ตีณินฺทฺริยา น วตฺตพฺพา “อาสวา เจว สาสวา จา”ติปิ, “สาสวา เจว โน จ อาสวา”ติปิ.\csromanspacer{} นวินฺทฺริยา น วตฺตพฺพา “อาสวา เจว สาสวา จา”ติ, สิยา สาสวา เจว โน จ อาสวา, สิยา น วตฺตพฺพา “สาสวา เจว โน จ อาสวา”ติ.}

\prose{ปนฺนรสินฺทฺริยา น วตฺตพฺพา “อาสวา เจว อาสวสมฺปยุตฺตา จา”ติปิ, “อาสวสมฺปยุตฺตา เจว โน จ อาสวา”ติปิ.\csromanspacer{} โทมนสฺสินฺทฺริยํ น วตฺตพฺพํ “อาสโว เจว อาสวสมฺปยุตฺตญฺจา”ติ, อาสวสมฺปยุตฺตญฺเจว โน จ อาสโว.\csromanspacer{} ฉ อินฺทฺริยา น วตฺตพฺพา “อาสวา เจว อาสวสมฺปยุตฺตา จา”ติ, สิยา อาสวสมฺปยุตฺตา เจว โน จ อาสวา, สิยา น วตฺตพฺพา “อาสวสมฺปยุตฺตา เจว โน จ อาสวา”ติ.\csromanspacer{} นวินฺทฺริยา อาสววิปฺปยุตฺตสาสวา.\csromanspacer{} ตีณินฺทฺริยา อาสววิปฺปยุตฺต-อนาสวา.\csromanspacer{} โทมนสฺสินฺทฺริยํ น วตฺตพฺพํ “อาสววิปฺปยุตฺตสาสวนฺ”ติปิ, “อาสววิปฺปยุตฺต-อนาสวนฺ”ติปิ.\csromanspacer{} ตีณินฺทฺริยา สิยา อาสววิปฺปยุตฺตสาสวา, สิยา อาสววิปฺปยุตฺต-อนาสวา.\csromanspacer{} ฉ อินฺทฺริยา สิยา อาสววิปฺปยุตฺตสาสวา, สิยา อาสววิปฺปยุตฺต-อนาสวา, สิยา น วตฺตพฺพา “อาสววิปฺปยุตฺตสาสวา”ติปิ, “อาสววิปฺปยุตฺต-อนาสวา”ติปิ. (3)}

\prose{โน สํโยชนา.\csromanspacer{} ทสินฺทฺริยา สํโยชนิยา.\csromanspacer{} ตีณินฺทฺริยา อสํโยชนิยา.\csromanspacer{} นวินฺทฺริยา สิยา สํโยชนิยา, สิยา อสํโยชนิยา.\csromanspacer{} ปนฺนรสินฺทฺริยา สํโยชนวิปฺปยุตฺตา.\csromanspacer{} โทมนสฺสินฺทฺริยํ สํโยชนสมฺปยุตฺตํ. \csromanfolio{136}ฉ อินฺทฺริยา สิยา สํโยชนสมฺปยุตฺตา, สิยา สํโยชนวิปฺปยุตฺตา.\csromanspacer{} ทสินฺทฺริยา น วตฺตพฺพา “สํโยชนา เจว สํโยชนิยา จา”ติ, สํโยชนิยา เจว โน จ สํโยชนา.\csromanspacer{} ตีณินฺทฺริยา น วตฺตพฺพา “สํโยชนา เจว สํโยชนิยา จา”ติปิ, “สํโยชนิยา เจว โน จ สํโยชนา”ติปิ.\csromanspacer{} นวินฺทฺริยา น วตฺตพฺพา “สํโยชนา เจว สํโยชนิยา จา”ติ, สิยา สํโยชนิยา เจว โน จ สํโยชนา, สิยา น วตฺตพฺพา “สํโยชนิยา เจว โน จ สํโยชนา”ติ.}

\prose{ปนฺนรสินฺทฺริยา น วตฺตพฺพา “สํโยชนา เจว สํโยชนสมฺปยุตฺตา จา”ติปิ, “สํโยชนสมฺปยุตฺตา เจว โน จ สํโยชนา”ติปิ.\csromanspacer{} โทมนสฺสินฺทฺริยํ น วตฺตพฺพํ “สํโยชนญฺเจว สํโยชนสมฺปยุตฺตญฺจา”ติ, สํโยชนสมฺปยุตฺตญฺเจว โน จ สํโยชนํ.\csromanspacer{} ฉ อินฺทฺริยา น วตฺตพฺพา “สํโยชนา เจว สํโยชนสมฺปยุตฺตา จา”ติ, สิยา สํโยชนสมฺปยุตฺตา เจว โน จ สํโยชนา, สิยา น วตฺตพฺพา “สํโยชนสมฺปยุตฺตา เจว โน จ สํโยชนา”ติ.}

\prose{นวินฺทฺริยา สํโยชนวิปฺปยุตฺตสํโยชนิยา.\csromanspacer{} ตีณินฺทฺริยา สํโยชนวิปฺปยุตฺต-อสํโยชนิยา.\csromanspacer{} โทมนสฺสินฺทฺริยํ น วตฺตพฺพํ “สํโยชนวิปฺปยุตฺตสํโยชนิยนฺ”ติปิ, “สํโยชนวิปฺปยุตฺต- อสํโยชนิยนฺ”ติปิ.\csromanspacer{} ตีณินฺทฺริยา สิยา สํโยชนวิปฺปยุตฺตสํโยชนิยา, สิยา สํโยชนวิปฺปยุตฺต-อสํโยชนิยา.\csromanspacer{} ฉ อินฺทฺริยา สิยา สํโยชนวิปฺปยุตฺตสํโยชนิยา, สิยา สํโยชนวิปฺปยุตฺต-อสํโยชนิยา, สิยา น วตฺตพฺพา “สํโยชนวิปฺปยุตฺตสํโยชนิยา”ติปิ, “สํโยชนวิปฺปยุตฺต-อสํโยชนิยา”ติปิ. (4)}

\prose{โน คนฺถา.\csromanspacer{} ทสินฺทฺริยา คนฺถนิยา.\csromanspacer{} ตีณินฺทฺริยา อคนฺถนิยา.\csromanspacer{} นวินฺทฺริยา สิยา คนฺถนิยา, สิยา อคนฺถนิยา.\csromanspacer{} ปนฺนรสินฺทฺริยา คนฺถวิปฺปยุตฺตา.\csromanspacer{} โทมนสฺสินฺทฺริยํ คนฺถสมฺปยุตฺตํ.\csromanspacer{} ฉ อินฺทฺริยา สิยา คนฺถสมฺปยุตฺตา, สิยา คนฺถวิปฺปยุตฺตา.\csromanspacer{} ทสินฺทฺริยา น วตฺตพฺพา “คนฺถา เจว คนฺถนิยา จา”ติ, คนฺถนิยา เจว โน จ คนฺถา.\csromanspacer{} ตีณินฺทฺริยา น วตฺตพฺพา “คนฺถา เจว คนฺถนิยา จา”ติปิ, “คนฺถนิยา เจว โน จ คนฺถา”ติปิ.\csromanspacer{} นวินฺทฺริยา น วตฺตพฺพา “คนฺถา เจว คนฺถนิยา จา”ติ, สิยา คนฺถนิยา เจว โน จ คนฺถา, สิยา น วตฺตพฺพา “คนฺถนิยา เจว โน จ คนฺถา”ติ.}

\csromanheader{2}{5. อินฺทฺริยวิภงฺค}
\prose{\csromanfolio{137}ปนฺนรสินฺทฺริยา น วตฺตพฺพา “คนฺถา เจว คนฺถสมฺปยุตฺตา จา”ติปิ, “คนฺถสมฺปยุตฺตา เจว โน จ คนฺถา”ติปิ.\csromanspacer{} โทมนสฺสินฺทฺริยํ น วตฺตพฺพํ “คนฺโถ เจว คนฺถสมฺปยุตฺตญฺจา”ติ, คนฺถสมฺปยุตฺตญฺเจว โน จ คนฺโถ.\csromanspacer{} ฉ อินฺทฺริยา น วตฺตพฺพา “คนฺถา เจว คนฺถสมฺปยุตฺตา จา”ติ, สิยา คนฺถสมฺปยุตฺตา เจว โน จ คนฺถา, สิยา น วตฺตพฺพา “คนฺถสมฺปยุตฺตา เจว โน จ คนฺถา”ติ.}

\prose{นวินฺทฺริยา คนฺถวิปฺปยุตฺตคนฺถนิยา, ตีณินฺทฺริยา คนฺถวิปฺปยุตฺต- อคนฺถนิยา.\csromanspacer{} โทมนสฺสินฺทฺริยํ น วตฺตพฺพํ “คนฺถวิปฺปยุตฺตคนฺถนิยนฺ”ติปิ, “คนฺถวิปฺปยุตฺต-อคนฺถนิยนฺ”ติปิ.\csromanspacer{} ตีณินฺทฺริยา สิยา คนฺถวิปฺปยุตฺตคนฺถนิยา, สิยา คนฺถวิปฺปยุตฺต- อคนฺถนิยา.\csromanspacer{} ฉ อินฺทฺริยา สิยา คนฺถวิปฺปยุตฺตคนฺถนิยา, สิยา คนฺถวิปฺปยุตฺต-อคนฺถนิยา, สิยา น วตฺตพฺพา “คนฺถวิปฺปยุตฺตคนฺถนิยา”ติปิ, “คนฺถวิปฺปยุตฺต-อคนฺถนิยา”ติปิ. (5)}

\prose{โน โอฆา ฯเปฯ.\csromanspacer{} โน โยคา ฯเปฯ.\csromanspacer{} โน นีวรณา.\csromanspacer{} ทสินฺทฺริยา นีวรณิยา.\csromanspacer{} ตีณินฺทฺริยา อนีวรณิยา.\csromanspacer{} นวินฺทฺริยา นีวรณิยา, สิยา อนีวรณิยา.\csromanspacer{} ปนฺนรสินฺทฺริยา นีวรณวิปฺปยุตฺตา.\csromanspacer{} โทมนสฺสินฺทฺริยํ นีวรณสมฺปยุตฺตํ.\csromanspacer{} ฉ อินฺทฺริยา สิยา นีวรณสมฺปยุตฺตา, สิยา นีวรณวิปฺปยุตฺตา.\csromanspacer{} ทสินฺทฺริยา น วตฺตพฺพา “นีวรณา เจว นีวรณิยา จา”ติ, นีวรณิยา เจว โน จ นีวรณา.\csromanspacer{} ตีณินฺทฺริยา น วตฺตพฺพา “นีวรณา เจว นีวรณิยา จา”ติปิ, “นีวรณิยา เจว โน จ นีวรณา”ติปิ.\csromanspacer{} นวินฺทฺริยา น วตฺตพฺพา “นีวรณา เจว นีวรณิยา จา”ติ, สิยา นีวรณิยา เจว โน จ นีวรณา, สิยา น วตฺตพฺพา “นีวรณิยา เจว โน จ นีวรณา”ติ.}

\prose{ปนฺนรสินฺทฺริยา น วตฺตพฺพา “นีวรณา เจว นีวรณสมฺปยุตฺตา จา”ติปิ, “นีวรณสมฺปยุตฺตา เจว โน จ นีวรณา”ติปิ.\csromanspacer{} โทมนสฺสินฺทฺริยํ น วตฺตพฺพํ “นีวรณญฺเจว นีวรณสมฺปยุตฺตญฺจา”ติ, นีวรณสมฺปยุตฺตญฺเจว โน จ นีวรณํ.\csromanspacer{} ฉ อินฺทฺริยา น วตฺตพฺพา “นีวรณา เจว นีวรณสมฺปยุตฺตา จา”ติ, สิยา นีวรณสมฺปยุตฺตา เจว โน จ นีวรณา, สิยา น วตฺตพฺพา “นีวรณสมฺปยุตฺตา เจว โน จ นีวรณา”ติ.}

\prose{นวินฺทฺริยา นีวรณวิปฺปยุตฺตนีวรณิยา.\csromanspacer{} ตีณินฺทฺริยา นีวรณวิปฺปยุตฺต-อนีวรณิยา.\csromanspacer{} โทมนสฺสินฺทฺริยํ น วตฺตพฺพํ “นีวรณวิปฺปยุตฺตนีวรณิยนฺ”ติปิ, “นีวรณวิปฺปยุตฺต-อนีวรณิยนฺ”ติปิ.\csromanspacer{} ตีณินฺทฺริยา สิยา นีวรณวิปฺปยุตฺตนีวรณิยา, สิยา นีวรณวิปฺปยุตฺต- อนีวรณิยา.\csromanspacer{} ฉ อินฺทฺริยา สิยา}

\prose{\csromanfolio{138}นีวรณวิปฺปยุตฺตนีวรณิยา, สิยา นีวรณวิปฺปยุตฺต-อนีวรณิยา, สิยา น วตฺตพฺพา “นีวรณวิปฺปยุตฺตนีวรณิยา”ติปิ, “นีวรณวิปฺปยุตฺต- อนีวรณิยา”ติปิ. (8)}

\prose{โน ปรามาสา.\csromanspacer{} ทสินฺทฺริยา ปรามฏฺฐา.\csromanspacer{} ตีณินฺทฺริยา อปรามฏฺฐา.\csromanspacer{} นวินฺทฺริยา สิยา ปรามฏฺฐา, สิยา อปรามฏฺฐา.\csromanspacer{} โสฬสินฺทฺริยา ปรามาสวิปฺปยุตฺตา.\csromanspacer{} ฉ อินฺทฺริยา สิยา ปรามาสสมฺปยุตฺตา, สิยา ปรามาสวิปฺปยุตฺตา.\csromanspacer{} ทสินฺทฺริยา น วตฺตพฺพา “ปรามาสา เจว ปรามฏฺฐา จา”ติ, ปรามฏฺฐา เจว โน จ ปรามาสา.\csromanspacer{} ตีณินฺทฺริยา น วตฺตพฺพา “ปรามาสา เจว ปรามฏฺฐา จา”ติปิ, “ปรามฏฺฐา เจว โน จ ปรามาสา”ติปิ.\csromanspacer{} นวินฺทฺริยา น วตฺตพฺพา “ปรามาสา เจว ปรามฏฺฐา จา”ติ, สิยา ปรามฏฺฐา เจว โน จ ปรามาสา, สิยา น วตฺตพฺพา “ปรามฏฺฐา เจว โน จ ปรามาสา”ติ.\csromanspacer{} ทสินฺทฺริยา ปรามาสวิปฺปยุตฺตปรามฏฺฐา.\csromanspacer{} ตีณินฺทฺริยา ปรามาสวิปฺปยุตฺต- อปรามฏฺฐา.\csromanspacer{} ตีณินฺทฺริยา สิยา ปรามาสวิปฺปยุตฺตปรามฏฺฐา, สิยา ปรามาสวิปฺปยุตฺต-อปรามฏฺฐา.\csromanspacer{} ฉ อินฺทฺริยา สิยา ปรามาสวิปฺปยุตฺตปรามฏฺฐา, สิยา ปรามาสวิปฺปยุตฺต-อปรามฏฺฐา, สิยา น วตฺตพฺพา “ปรามาสวิปฺปยุตฺตปรามฏฺฐา”ติปิ, “ปรามาสวิปฺปยุตฺต- อปรามฏฺฐา”ติปิ. (9)}

\prose{สตฺตินฺทฺริยา อนารมฺมณา.\csromanspacer{} จุทฺทสินฺทฺริยา สารมฺมณา.\csromanspacer{} ชีวิตินฺทฺริยา สิยา สารมฺมณํ, สิยา อนารมฺมณํ.\csromanspacer{} เอกวีสตินฺทฺริยา โน จิตฺตา.\csromanspacer{} มนินฺทฺริยํ จิตฺตํ.\csromanspacer{} เตรสินฺทฺริยา เจตสิกา.\csromanspacer{} อฏฺฐินฺทฺริยา อเจตสิกา.\csromanspacer{} ชีวิตินฺทฺริยํ สิยา เจตสิกํ, สิยา อเจตสิกํ.\csromanspacer{} เตรสินฺทฺริยา จิตฺตสมฺปยุตฺตา.\csromanspacer{} สตฺตินฺทฺริยา จิตฺตวิปฺปยุตฺตา.\csromanspacer{} ชีวิตินฺทฺริยํ สิยา จิตฺตสมฺปยุตฺตํ, สิยา จิตฺตวิปฺปยุตฺตํ.\csromanspacer{} มนินฺทฺริยํ น วตฺตพฺพํ “จิตฺเตน สมฺปยุตฺตนฺ”ติปิ, “จิตฺเตน วิปฺปยุตฺตนฺ”ติปิ.}

\prose{เตรสินฺทฺริยา จิตฺตสํสฏฺฐา.\csromanspacer{} สตฺตินฺทฺริยา จิตฺตวิสํสฏฺฐา.\csromanspacer{} ชีวิตินฺทฺริยํ สิยา จิตฺตสํสฏฺฐํ, สิยา จิตฺตวิสํสฏฺฐํ.\csromanspacer{} มนินฺทฺริยํ น วตฺตพฺพํ “จิตฺเตน สํสฏฺฐนฺ”ติปิ, “จิตฺเตน วิสํสฏฺฐนฺ”ติปิ.\csromanspacer{} เตรสินฺทฺริยา จิตฺตสมุฏฺฐานา.\csromanspacer{} อฏฺฐินฺทฺริยา โน จิตฺตสมุฏฺฐานา.\csromanspacer{} ชีวิตินฺทฺริยํ สิยา จิตฺตสมุฏฺฐานํ, สิยา โน จิตฺตสมุฏฺฐานํ.}

\prose{เตรสินฺทฺริยา จิตฺตสหภุโน, อฏฺฐินฺทฺริยา โน จิตฺตสหภุโน.\csromanspacer{} ชีวิตินฺทฺริยํ สิยา จิตฺตสหภู, สิยา โน จิตฺตสหภู.\csromanspacer{} เตรสินฺทฺริยา จิตฺตานุปริวตฺติโน.\csromanspacer{} อฏฺฐินฺทฺริยา โน จิตฺตานุปริวตฺติโน.\csromanspacer{} ชีวิตินฺทฺริยํ สิยา จิตฺตานุปริวตฺติ, สิยา โน จิตฺตานุปริวตฺติ.}

\csromanheader{2}{5. อินฺทฺริยวิภงฺค}
\prose{\csromanfolio{139}เตรสินฺทฺริยา จิตฺตสํสฏฺฐสมุฏฺฐานา.\csromanspacer{} อฏฺฐินฺทฺริยา โน จิตฺตสํสฏฺฐสมุฏฺฐานา.\csromanspacer{} ชีวิตินฺทฺริยํ สิยา จิตฺตสํสฏฺฐสมุฏฺฐานํ, สิยา โน จิตฺตสํสฏฺฐสมุฏฺฐานํ.\csromanspacer{} เตรสินฺทฺริยา จิตฺตสํสฏฺฐสมุฏฺฐานสหภุโน.\csromanspacer{} อฏฺฐินฺทฺริยา โน จิตฺตสํสฏฺฐสมุฏฺฐานสหภุโน.\csromanspacer{} ชีวิตินฺทฺริยํ สิยา จิตฺตสํสฏฺฐสมุฏฺฐานสหภู, สิยา โน จิตฺตสํสฏฺฐสมุฏฺฐานสหภู.\csromanspacer{} เตรสินฺทฺริยา จิตฺตสํสฏฺฐสมุฏฺฐานานุปริวตฺติโน.\csromanspacer{} อฏฺฐินฺทฺริยา โน จิตฺตสํสฏฺฐสมุฏฺฐานานุปริวตฺติโน.\csromanspacer{} ชีวิตินฺทฺริยํ สิยา จิตฺตสํสฏฺฐสมุฏฺฐานานุปริวตฺติ, สิยา โน จิตฺตสํสฏฺฐสมุฏฺฐานานุปริวตฺติ.\csromanspacer{} ฉ อินฺทฺริยา อชฺฌตฺติกา.\csromanspacer{} โสฬสินฺทฺริยา พาหิรา. (10)}

\prose{สตฺตินฺทฺริยา อุปาทา.\csromanspacer{} จุทฺทสินฺทฺริยา โน อุปาทา.\csromanspacer{} ชีวิตินฺทฺริยํ สิยา อุปาทา, สิยา โน อุปาทา.\csromanspacer{} นวินฺทฺริยา อุปาทินฺนา.\csromanspacer{} จตฺตารินฺทฺริยา อนุปาทินฺนา.\csromanspacer{} นวินฺทฺริยา สิยา อุปาทินฺนา, สิยา อนุปาทินฺนา, โน อุปาทานา.\csromanspacer{} ทสินฺทฺริยา อุปาทานิยา.\csromanspacer{} ตีณินฺทฺริยา อนุปาทานิยา.\csromanspacer{} นวินฺทฺริยา สิยา อุปาทานิยา, สิยา อนุปาทานิยา.\csromanspacer{} โสฬสินฺทฺริยา อุปาทานวิปฺปยุตฺตา.\csromanspacer{} ฉ อินฺทฺริยา สิยา อุปาทานสมฺปยุตฺตา, สิยา อุปาทานวิปฺปยุตฺตา.\csromanspacer{} ทสินฺทฺริยา น วตฺตพฺพา “อุปาทานา เจว อุปาทานิยา จา”ติ, อุปาทานิยา เจว โน จ อุปาทานา.\csromanspacer{} ตีณินฺทฺริยา น วตฺตพฺพา “อุปฺปาทานา เจว อุปาทานิยา จา”ติปิ, “อุปาทานิยา เจว โน จ อุปาทานา”ติปิ.\csromanspacer{} นวินฺทฺริยา น วตฺตพฺพา “อุปาทานา เจว อุปาทานิยา จา”ติ, สิยา อุปาทานิยา เจว โน จ อุปาทานา.\csromanspacer{} ทสินฺทฺริยา สิยา อุปาทานิยา เจว โน จ อุปาทานา, สิยา น วตฺตพฺพา “อุปาทานิยา เจว โน จ อุปาทานา”ติ.}

\prose{โสฬสินฺทฺริยา น วตฺตพฺพา “อุปาทานา เจว อุปาทานสมฺปยุตฺตา จา”ติปิ, “อุปาทานสมฺปยุตฺตา เจว โน จ อุปาทานา”ติปิ.\csromanspacer{} ฉ อินฺทฺริยา น วตฺตพฺพา “อุปาทานา เจว อุปาทานสมฺปยุตฺตา จา”ติ, สิยา อุปาทานสมฺปยุตฺตา เจว โน จ อุปาทานา, สิยา น วตฺตพฺพา “อุปาทานสมฺปยุตฺตา เจว โน จ อุปาทานา”ติ.\csromanspacer{} ทสินฺทฺริยา อุปาทานวิปฺปยุตฺต-อุปาทานิยา.\csromanspacer{} ตีณินฺทฺริยา อุปาทานวิปฺปยุตฺต- อนุปาทานิยา.\csromanspacer{} ตีณินฺทฺริยา สิยา อุปาทานวิปฺปยุตฺต-อุปาทานิยา, สิยา อุปาทานวิปฺปยุตฺต-อนุปาทานิยา.\csromanspacer{} ฉ อินฺทฺริยา สิยา อุปาทานวิปฺปยุตฺต- อุปาทานิยา, สิยา อุปาทานวิปฺปยุตฺต-อนุปาทานิยา, สิยา น วตฺตพฺพา “อุปาทานวิปฺปยุตฺต-อุปาทานิยา”ติปิ, “อุปาทานวิปฺปยุตฺต-อนุปาทานิยา”ติปิ. (11)}

\prose{\csromanfolio{140}โน กิเลสา.\csromanspacer{} ทสินฺทฺริยา สํกิเลสิกา, ตีณินฺทฺริยา อสํกิเลสิกา, นวินฺทฺริยา สิยา สํกิเลสิกา, สิยา อสํกิเลสิกา.\csromanspacer{} ปนฺนรสินฺทฺริยา อสํกิลิฏฺฐา.\csromanspacer{} โทมนสฺสินฺทฺริยํ สํกิลิฏฺฐํ.\csromanspacer{} ฉ อินฺทฺริยา สิยา สํกิลิฏฺฐา, สิยา อสํกิลิฏฺฐา.\csromanspacer{} ปนฺนรสินฺทฺริยา กิเลสวิปฺปยุตฺตา.\csromanspacer{} โทมนสฺสินฺทฺริยํ กิเลสสมฺปยุตฺตํ.\csromanspacer{} ฉ อินฺทฺริยา สิยา กิเลสสมฺปยุตฺตา, สิยา กิเลสวิปฺปยุตฺตา.\csromanspacer{} ทสินฺทฺริยา น วตฺตพฺพา “กิเลสา เจว สํกิเลสิกา จา”ติ.\csromanspacer{} สํกิเลสิกา เจว โน จ กิเลสา.\csromanspacer{} ตีณินฺทฺริยา น วตฺตพฺพา “กิเลสา เจว สํกิเลสิกา จา”ติปิ, “สํกิเลสิกา เจว โน จ กิเลสา”ติปิ.\csromanspacer{} นวินฺทฺริยา น วตฺตพฺพา “กิเลสา เจว สํกิเลสิกา จา”ติ, สิยา สํกิเลสิกา เจว โน จ กิเลสา, สิยา น วตฺตพฺพา “สํกิเลสิกา เจว โน จ กิเลสา”ติ.}

\prose{ปนฺนรสินฺทฺริยา น วตฺตพฺพา “กิเลสา เจว สํกิลิฏฺฐา จา”ติปิ, “สํกิลิฏฺฐา เจว โน จ กิเลสา”ติปิ.\csromanspacer{} โทมนสฺสินฺทฺริยํ น วตฺตพฺพํ “กิเลโส เจว สํกิลิฏฺฐญฺจา”ติ, สํกิลิฏฺฐญฺเจว โน จ กิเลโส.\csromanspacer{} ฉ อินฺทฺริยา น วตฺตพฺพา “กิเลสา เจว สํกิลิฏฺฐา จา”ติ, สิยา สํกิลิฏฺฐา เจว โน จ กิเลสา, สิยา น วตฺตพฺพา “สํกิลิฏฺฐา เจว โน จ กิเลสา”ติ.}

\prose{ปนฺนรสินฺทฺริยา น วตฺตพฺพา “กิเลสา เจว กิเลสสมฺปยุตฺตา จา”ติปิ, “กิเลสสมฺปยุตฺตา เจว โน จ กิเลสา”ติปิ.\csromanspacer{} โทมนสฺสินฺทฺริยํ น วตฺตพฺพํ “กิเลโส เจว กิเลสสมฺปยุตฺตญฺจา”ติ, กิเลสสมฺปยุตฺตญฺเจว โน จ กิเลโส.\csromanspacer{} ฉ อินฺทฺริยา น วตฺตพฺพา “กิเลสา เจว กิเลสสมฺปยุตฺตา จา”ติ.\csromanspacer{} สิยา กิเลสสมฺปยุตฺตา เจว โน จ กิเลสา, สิยา น วตฺตพฺพา “กิเลสสมฺปยุตฺตา เจว โน จ กิเลสา”ติ.\csromanspacer{} นวินฺทฺริยา กิเลสวิปฺปยุตฺตสํกิเลสิกา.\csromanspacer{} ตีณินฺทฺริยา กิเลสวิปฺปยุตฺต-อสํกิเลสิกา.\csromanspacer{} โทมนสฺสินฺทฺริยํ น วตฺตพฺพํ “กิเลสวิปฺปยุตฺตสํกิเลสิกนฺ”ติปิ, “กิเลสวิปฺปยุตฺต-อสํกิเลสิกนฺ”ติปิ.\csromanspacer{} ตีณินฺทฺริยา สิยา กิเลสวิปฺปยุตฺตสํกิเลสิกา, สิยา กิเลสวิปฺปยุตฺต- อสํกิเลสิกา.\csromanspacer{} ฉ อินฺทฺริยา สิยา กิเลสวิปฺปยุตฺตสํกิเลสิกา, สิยา กิเลสวิปฺปยุตฺต-อสํกิเลสิกา, สิยา น วตฺตพฺพา “กิเลสวิปฺปยุตฺตสํกิเลสิกา”ติปิ, “กิเลสวิปฺปยุตฺต-อสํกิเลสิกา”ติปิ. (12)}

\csromanheader{2}{5. อินฺทฺริยวิภงฺค}
\prose{\csromanfolio{141}ปนฺนรสินฺทฺริยา น ทสฺสเนน ปหาตพฺพา.\csromanspacer{} สตฺตินฺทฺริยา สิยา ทสฺสเนน ปหาตพฺพา, สิยา น ทสฺสเนน ปหาตพฺพา.\csromanspacer{} ปนฺนรสินฺทฺริยา น ภาวนาย ปหาตพฺพา.\csromanspacer{} สตฺตินฺทฺริยา สิยา ภาวนาย ปหาตพฺพา, สิยา น ภาวนาย ปหาตพฺพา.\csromanspacer{} ปนฺนรสินฺทฺริยา น ทสฺสเนน ปหาตพฺพเหตุกา.\csromanspacer{} สตฺตินฺทฺริยา สิยา ทสฺสเนน ปหาตพฺพเหตุกา, สิยา น ทสฺสเนน ปหาตพฺพเหตุกา.\csromanspacer{} ปนฺนรสินฺทฺริยา น ภาวนาย ปหาตพฺพเหตุกา.\csromanspacer{} สตฺตินฺทฺริยา สิยา ภาวนาย ปหาตพฺพเหตุกา, สิยา น ภาวนาย ปหาตพฺพเหตุกา.}

\prose{นวินฺทฺริยา อวิตกฺกา.\csromanspacer{} โทมนสฺสินฺทฺริยํ สวิตกฺกํ.\csromanspacer{} ทฺวาทสินฺทฺริยา สิยา สวิตกฺกา, สิยา อวิตกฺกา.\csromanspacer{} นวินฺทฺริยา อวิจารา.\csromanspacer{} โทมนสฺสินฺทฺริยํ สวิจารํ.\csromanspacer{} ทฺวาทสินฺทฺริยา สิยา สวิจารา, สิยา อวิจารา.\csromanspacer{} เอกาทสินฺทฺริยา อปฺปีติกา.\csromanspacer{} เอกาทสินฺทฺริยา สิยา สปฺปีติกา, สิยา อปฺปีติกา.\csromanspacer{} เอกาทสินฺทฺริยา น ปีติสหคตา.\csromanspacer{} เอกาทสินฺทฺริยา สิยา ปีติสหคตา, สิยา น ปีติสหคตา.\csromanspacer{} ทฺวาทสินฺทฺริยา น สุขสหคตา.\csromanspacer{} ทสินฺทฺริยา สิยา สุขสหคตา, สิยา น สุขสหคตา.\csromanspacer{} ทฺวาทสินฺทฺริยา น อุเปกฺขาสหคตา, ทสินฺทฺริยา สิยา อุเปกฺขาสหคตา, สิยา น อุเปกฺขาสหคตา.}

\prose{ทสินฺทฺริยา กามาวจรา.\csromanspacer{} ตีณินฺทฺริยา น กามาวจรา.\csromanspacer{} นวินฺทฺริยา สิยา กามาวจรา, สิยา น กามาวจรา.\csromanspacer{} เตรสินฺทฺริยา น รูปาวจรา.\csromanspacer{} นวินฺทฺริยา สิยา รูปาวจรา, สิยา น รูปาวจรา.\csromanspacer{} จุทฺทสินฺทฺริยา น อรูปาวจรา.\csromanspacer{} อฏฺฐินฺทฺริยา สิยา อรูปาวจรา, สิยา น อรูปาวจรา.\csromanspacer{} ทสินฺทฺริยา ปริยาปนฺนา.\csromanspacer{} ตีณินฺทฺริยา อปริยาปนฺนา.\csromanspacer{} นวินฺทฺริยา สิยา ปริยาปนฺนา, สิยา อปริยาปนฺนา.\csromanspacer{} เอกาทสินฺทฺริยา อนิยฺยานิกา.\csromanspacer{} อนญฺญาตญฺญสฺสามีตินฺทฺริยํ นิยฺยานิกํ.\csromanspacer{} ทสินฺทฺริยา สิยา นิยฺยานิกา, สิยา อนิยฺยานิกา.\csromanspacer{} ทสินฺทฺริยา อนิยตา.\csromanspacer{} อนญฺญาตญฺญสฺสามีตินฺทฺริยํ นิยตํ.\csromanspacer{} เอกาทสินฺทฺริยา สิยา นิยตา, สิยา อนิยตา.\csromanspacer{} ทสินฺทฺริยา ส-อุตฺตรา.\csromanspacer{} ตีณินฺทฺริยา อนุตฺตรา.\csromanspacer{} นวินฺทฺริยา สิยา ส-อุตฺตรา, สิยา อนุตฺตรา.\csromanspacer{} ปนฺนรสินฺทฺริยา อรณา.\csromanspacer{} โทมนสฺสินฺทฺริยํ สรณํ.\csromanspacer{} ฉ อินฺทฺริยา สิยา สรณา, สิยา อรณาติ. (13)}

\vspace{\tipitakaparskip}
\csromancenter{ปญฺหาปุจฺฉกํ.}
\nitthitam{\textbf{อินฺทฺริยวิภงฺโค นิฏฺฐิโต.}}
\csromanmatikaanchor{mtk.44}
\csromantocmarknopage{subparagraph}{6. ปฏิจฺจสมุปฺปาทวิภงฺค}{mtk.44}
\bookmark[dest={mtk.44},level=5]{6. ปฏิจฺจสมุปฺปาทวิภงฺค}
\csromanheader{6}{6. ปฏิจฺจสมุปฺปาทวิภงฺค}
\csromanmatikaanchor{mtk.45}
\csromantocmarkref{subparagraph}{1. สุตฺตนฺตภาชนีย}{mtk.45}
\bookmark[dest={mtk.45},level=5]{1. สุตฺตนฺตภาชนีย}
\csromanheader{6}{1. สุตฺตนฺตภาชนีย}
\proseitem{225}{\csromanfolio{142}อวิชฺชาปจฺจยา สงฺขารา, สงฺขารปจฺจยา วิญฺญาณํ, วิญฺญาณปจฺจยา นามรูปํ, นามรูปปจฺจยา สฬายตนํ, สฬายตนปจฺจยา ผสฺโส, ผสฺสปจฺจยา เวทนา, เวทนาปจฺจยา ตณฺหา, ตณฺหาปจฺจยา อุปาทานํ, อุปาทานปจฺจยา ภโว, ภวปจฺจยา ชาติ, ชาติปจฺจยา ชรามรณํ โสกปริเทวทุกฺขโทมนสฺสุปายาสา สมฺภวนฺติ, เอวเมตสฺส เกวลสฺส ทุกฺขกฺขนฺธสฺส สมุทโย โหติ.}

\proseitem{226}{ตตฺถ กตมา อวิชฺชา, ทุกฺเข อญฺญาณํ ทุกฺขสมุทเย อญฺญาณํ ทุกฺขนิโรเธ อญฺญาณํ ทุกฺขนิโรธคามินิยา ปฏิปทาย อญฺญาณํ, อยํ วุจฺจติ อวิชฺชา.}

\prose{ตตฺถ กตเม อวิชฺชาปจฺจยา สงฺขารา, ปุญฺญาภิสงฺขาโร อปุญฺญาภิสงฺขาโร อาเนญฺชาภิสงฺขาโร กายสงฺขาโร วจีสงฺขาโร จิตฺตสงฺขาโร.}

\prose{ตตฺถ กตโม ปุญฺญาภิสงฺขาโร, กุสลา เจตนา กามาวจรา รูปาวจรา ทานมยา สีลมยา ภาวนามยา, อยํ วุจฺจติ ปุญฺญาภิสงฺขาโร.}

\prose{ตตฺถ กตโม อปุญฺญาภิสงฺขาโร, อกุสลา เจตนา กามาวจรา, อยํ วุจฺจติ อปุญฺญาภิสงฺขาโร.}

\prose{ตตฺถ กตโม อาเนญฺชาภิสงฺขาโร, กุสลา เจตนา อรูปาวจรา, อยํ วุจฺจติ อาเนญฺชาภิสงฺขาโร.}

\prose{ตตฺถ กตโม กายสงฺขาโร, กายสญฺเจตนา กายสงฺขาโร, วจีสญฺเจตนา วจีสงฺขาโร, มโนสญฺเจตนา จิตฺตสงฺขาโร, อิเม วุจฺจนฺติ อวิชฺชาปจฺจยา สงฺขารา.}

\proseitem{227}{ตตฺถ กตมํ สงฺขารปจฺจยา วิญฺญาณํ, จกฺขุวิญฺญาณํ โสตวิญฺญาณํ ฆานวิญฺญาณํ ชิวฺหาวิญฺญาณํ กายวิญฺญาณํ มโนวิญฺญาณํ, อิทํ วุจฺจติ สงฺขารปจฺจยา วิญฺญาณํ.}

\proseitem{228}{ตตฺถ กตมํ วิญฺญาณปจฺจยา นามรูปํ, อตฺถิ นามํ, อตฺถิ รูปํ.\csromanspacer{} ตตฺถ กตมํ นามํ, เวทนากฺขนฺโธ สญฺญากฺขนฺโธ สงฺขารกฺขนฺโธ, อิทํ วุจฺจติ}

\vspace{\tipitakaparskip}
\csromancenter{ปฏิจฺจสมุปฺปาทวิภงฺค นามํ.\csromanspacer{} ตตฺถ กตมํ รูปํ, จตฺตาโร มหาภูตา จตุนฺนญฺจ มหาภูตานํ อุปาทายรูปํ, อิทํ วุจฺจติ รูปํ.\csromanspacer{} อิติ อิทญฺจ นามํ อิทญฺจ รูปํ, อิทํ วุจฺจติ วิญฺญาณปจฺจยา นามรูปํ.}
\proseitem{229}{\csromanfolio{143}ตตฺถ กตมํ นามรูปปจฺจยา สฬายตนํ, จกฺขายตนํ โสตายตนํ ฆานายตนํ ชิวฺหายตนํ กายายตนํ มนายตนํ, อิทํ วุจฺจติ นามรูปปจฺจยา สฬายตนํ.}

\proseitem{230}{ตตฺถ กตโม สฬายตนปจฺจยา ผสฺโส, จกฺขุสมฺผสฺโส โสตสมฺผสฺโส ฆานสมฺผสฺโส ชิวฺหาสมฺผสฺโส กายสมฺผสฺโส มโนสมฺผสฺโส, อยํ วุจฺจติ สฬายตนปจฺจยา ผสฺโส.}

\proseitem{231}{ตตฺถ กตมา ผสฺสปจฺจยา เวทนา, จกฺขุสมฺผสฺสชา เวทนา โสตสมฺผสฺสชา เวทนา ฆานสมฺผสฺสชา เวทนา ชิวฺหาสมฺผสฺสชา เวทนา กายสมฺผสฺสชา เวทนา มโนสมฺผสฺสชา เวทนา, อยํ วุจฺจติ ผสฺสปจฺจยา เวทนา.}

\proseitem{232}{ตตฺถ กตมา เวทนาปจฺจยา ตณฺหา, รูปตณฺหา สทฺทตณฺหา คนฺธตณฺหา รสตณฺหา โผฏฺฐพฺพตณฺหา ธมฺมตณฺหา, อยํ วุจฺจติ เวทนาปจฺจยา ตณฺหา.}

\proseitem{233}{ตตฺถ กตมํ ตณฺหาปจฺจยา อุปาทานํ, กามุปาทานํ ทิฏฺฐุปาทานํ สีลพฺพตุปาทานํ อตฺตวาทุปาทานํ, อิทํ วุจฺจติ ตณฺหาปจฺจยา อุปาทานํ.}

\proseitem{234}{ตตฺถ กตโม อุปาทานปจฺจยา ภโว, ภโว ทุวิเธน อตฺถิ กมฺมภโว, อตฺถิ อุปปตฺติภโว.\csromanspacer{} ตตฺถ กตโม กมฺมภโว, ปุญฺญาภิสงฺขาโร อปุญฺญาภิสงฺขาโร อาเนญฺชาภิสงฺขาโร, อยํ วุจฺจติ กมฺมภโว.\csromanspacer{} สพฺพมฺปิ ภวคามิกมฺมํ กมฺมภโว.}

\prose{ตตฺถ กตโม อุปปตฺติภโว, กามภโว รูปภโว อรูปภโว สญฺญาภโว อสญฺญาภโว เนวสญฺญานาสญฺญาภโว เอกโวการภโว จตุโวการภโว ปญฺจโวการภโว, อยํ วุจฺจติ อุปปตฺติภโว.\csromanspacer{} อยํ วุจฺจติ อุปาทานปจฺจยา ภโว.}

\proseitem{235}{\csromanfolio{144}ตตฺถ กตมา ภวปจฺจยา ชาติ, ยา เตสํ เตสํ สตฺตานํ ตมฺหิ ตมฺหิ สตฺตนิกาเย ชาติ สญฺชาติ โอกฺกนฺติ อภินิพฺพตฺติ ขนฺธานํ ปาตุภาโว อายตนานํ ปฏิลาโภ, อยํ วุจฺจติ ภวปจฺจยา ชาติ.}

\proseitem{236}{ตตฺถ กตมํ ชาติปจฺจยา ชรามรณํ, อตฺถิ ชรา, อตฺถิ มรณํ.\csromanspacer{} ตตฺถ กตมา ชรา, ยา เตสํ เตสํ สตฺตานํ ตมฺหิ ตมฺหิ สตฺตนิกาเย ชรา ชีรณตา ขณฺฑิจฺจํ ปาลิจฺจํ วลิตฺตจตา อายุโน สํหานิ อินฺทฺริยานํ ปริปาโก, อยํ วุจฺจติ ชรา.}

\prose{ตตฺถ กตมํ มรณํ, ยา เตสํ เตสํ สตฺตานํ ตมฺหา ตมฺหา สตฺตนิกายา จุติ จวนตา เภโท อนฺตรธานํ มจฺจุ มรณํ กาลกิริยา\footnote{กาลํกิริยา (ก)} ขนฺธานํ เภโท กเฬวรสฺส นิกฺเขโป ชีวิตินฺทฺริยสฺสุปจฺเฉโท, อิทํ วุจฺจติ มรณํ.\csromanspacer{} อิติ อยญฺจ ชรา อิทญฺจ มรณํ, อิทํ วุจฺจติ ชาติปจฺจยา ชรามรณํ.}

\proseitem{237}{ตตฺถ กตโม โสโก, ญาติพฺยสเนน วา ผุฏฺฐสฺส โภคพฺยสเนน วา ผุฏฺฐสฺส โรคพฺยสเนน วา ผุฏฺฐสฺส สีลพฺยสเนน วา ผุฏฺฐสฺส ทิฏฺฐิพฺยสเนน วา ผุฏฺฐสฺส อญฺญตรญฺญตเรน พฺยสเนน สมนฺนาคตสฺส อญฺญตรญฺญตเรน ทุกฺขธมฺเมน ผุฏฺฐสฺส โสโก โสจนา โสจิตตฺตํ อนฺโตโสโก อนฺโตปริโสโก เจตโส ปริชฺฌายนา โทมนสฺสํ โสกสลฺลํ, อยํ วุจฺจติ โสโก.}

\proseitem{238}{ตตฺถ กตโม ปริเทโว, ญาติพฺยสเนน วา ผุฏฺฐสฺส โภคพฺยสเนน วา ผุฏฺฐสฺส โรคพฺยสเนน วา ผุฏฺฐสฺส สีลพฺยสเนน วา ผุฏฺฐสฺส ทิฏฺฐิพฺยสเนน วา ผุฏฺฐสฺส อญฺญตรญฺญตเรน พฺยสเนน สมนฺนาคตสฺส อญฺญตรญฺญตเรน ทุกฺขธมฺเมน ผุฏฺฐสฺส อาเทโว ปริเทโว อาเทวนา ปริเทวนา อาเทวิตตฺตํ ปริเทวิตตฺตํ วาจา ปลาโป วิปฺปลาโป ลาลปฺโป ลาลปฺปนา ลาลปฺปิตตฺตํ, อยํ วุจฺจติ ปริเทโว.}

\proseitem{239}{ตตฺถ กตมํ ทุกฺขํ, ยํ กายิกํ อสาตํ กายิกํ ทุกฺขํ กายสมฺผสฺสชํ อสาตํ ทุกฺขํ เวทยิตํ กายสมฺผสฺสชา อสาตา ทุกฺขา เวทนา, อิทํ วุจฺจติ ทุกฺขํ.}

\csromanheader{2}{6. ปฏิจฺจสมุปฺปาทวิภงฺค}
\proseitem{240}{\csromanfolio{145}ตตฺถ กตมํ โทมนสฺสํ, ยํ เจตสิกํ อสาตํ เจตสิกํ ทุกฺขํ เจโตสมฺผสฺสชํ อสาตํ ทุกฺขํ เวทยิตํ เจโตสมฺผสฺสชา อสาตา ทุกฺขา เวทนา, อิทํ วุจฺจติ โทมนสฺสํ.}

\proseitem{241}{ตตฺถ กตโม อุปายาโส, ญาติพฺยสเนน วา ผุฏฺฐสฺส โภคพฺยสเนน วา ผุฏฺฐสฺส โรคพฺยสเนน วา ผุฏฺฐสฺส สีลพฺยสเนน วา ผุฏฺฐสฺส ทิฏฺฐิพฺยสเนน วา ผุฏฺฐสฺส อญฺญตรญฺญตเรน พฺยสเนน สมนฺนาคตสฺส อญฺญตรญฺญตเรน ทุกฺขธมฺเมน ผุฏฺฐสฺส อายาโส อุปายาโส อายาสิตตฺตํ อุปายาสิตตฺตํ, อยํ วุจฺจติ อุปายาโส.}

\proseitem{242}{เอวเมตสฺส เกวลสฺส ทุกฺขกฺขนฺธสฺส สมุทโย โหตีติ, เอวเมตสฺส เกวลสฺส ทุกฺขกฺขนฺธสฺส สงฺคติ โหติ, สมาคโม โหติ, สโมธานํ โหติ, ปาตุภาโว โหติ, เตน วุจฺจติ เอวเมตสฺส เกวลสฺส ทุกฺขกฺขนฺธสฺส สมุทโย โหตีติ.}

\csromanheader{2}{สุตฺตนฺตภาชนียํ.}
\csromansectionrule
\csromanmatikaanchor{mtk.46}
\csromantocmarkref{subparagraph}{2. อภิธมฺมภาชนีย}{mtk.46}
\bookmark[dest={mtk.46},level=5]{2. อภิธมฺมภาชนีย}
\csromanheader{6}{2. อภิธมฺมภาชนีย}
\csromanmatikaanchor{mtk.47}
\csromantocmarkref{subparagraph}{1. ปจฺจยจตุกฺก}{mtk.47}
\bookmark[dest={mtk.47},level=5]{1. ปจฺจยจตุกฺก}
\csromanheader{6}{1. ปจฺจยจตุกฺก}
\proseitem{243}{อวิชฺชาปจฺจยา สงฺขาโร, สงฺขารปจฺจยา วิญฺญาณํ, วิญฺญาณปจฺจยา นามํ, นามปจฺจยา ฉฏฺฐายตนํ, ฉฏฺฐายตนปจฺจยา ผสฺโส, ผสฺสปจฺจยา เวทนา, เวทนาปจฺจยา ตณฺหา, ตณฺหาปจฺจยา อุปาทานํ, อุปาทานปจฺจยา ภโว, ภวปจฺจยา ชาติ, ชาติปจฺจยา ชรามรณํ, เอวเมตสฺส เกวลสฺส ทุกฺขกฺขนฺธสฺส สมุทโย โหติ. (1)}

\prose{อวิชฺชาปจฺจยา สงฺขาโร, สงฺขารปจฺจยา วิญฺญาณํ, วิญฺญาณปจฺจยา นามํ, นามปจฺจยา ผสฺโส, ผสฺสปจฺจยา เวทนา, เวทนาปจฺจยา ตณฺหา, ตณฺหาปจฺจยา อุปาทานํ, อุปาทานปจฺจยา ภโว, ภวปจฺจยา ชาติ, ชาติปจฺจยา ชรามรณํ, เอวเมตสฺส เกวลสฺส ทุกฺขกฺขนฺธสฺส สมุทโย โหติ. (2)}

\prose{\csromanfolio{146}อวิชฺชาปจฺจยา สงฺขาโร, สงฺขารปจฺจยา วิญฺญาณํ, วิญฺญาณปจฺจยา นามรูปํ, นามรูปปจฺจยา ฉฏฺฐายตนํ, ฉฏฺฐายตนปจฺจยา ผสฺโส, ผสฺสปจฺจยา เวทนา, เวทนาปจฺจยา ตณฺหา, ตณฺหาปจฺจยา อุปาทานํ, อุปาทานปจฺจยา ภโว, ภวปจฺจยา ชาติ, ชาติปจฺจยา ชรามรณํ, เอวเมตสฺส เกวลสฺส ทุกฺขกฺขนฺธสฺส สมุทโย โหติ. (3)}

\prose{อวิชฺชาปจฺจยา สงฺขาโร, สงฺขารปจฺจยา วิญฺญาณํ, วิญฺญาณปจฺจยา นามรูปํ, นามรูปปจฺจยา สฬายตนํ, ฉฏฺฐายตนปจฺจยา ผสฺโส, ผสฺสปจฺจยา เวทนา, เวทนาปจฺจยา ตณฺหา, ตณฺหาปจฺจยา อุปาทานํ, อุปาทานปจฺจยา ภโว, ภวปจฺจยา ชาติ, ชาติปจฺจยา ชรามรณํ, เอวเมตสฺส เกวลสฺส ทุกฺขกฺขนฺธสฺส สมุทโย โหติ. (4)}

\csromanheader{2}{ปจฺจยจตุกฺกํ.}
\csromansectionrule
\csromanmatikaanchor{mtk.48}
\csromantocmarkref{subparagraph}{2. เหตุจตุกฺก}{mtk.48}
\bookmark[dest={mtk.48},level=5]{2. เหตุจตุกฺก}
\csromanheader{6}{2. เหตุจตุกฺก}
\proseitem{244}{อวิชฺชาปจฺจยา สงฺขาโร อวิชฺชาเหตุโก, สงฺขารปจฺจยา วิญฺญาณํ สงฺขารเหตุกํ, วิญฺญาณปจฺจยา นามํ วิญฺญาณเหตุกํ, นามปจฺจยา ฉฏฺฐายตนํ นามเหตุกํ, ฉฏฺฐายตนปจฺจยา ผสฺโส ฉฏฺฐายตนเหตุโก, ผสฺสปจฺจยา เวทนา ผสฺสเหตุกา, เวทนาปจฺจยา ตณฺหา เวทนาเหตุกา, ตณฺหาปจฺจยา อุปาทานํ ตณฺหาเหตุกํ, อุปาทานปจฺจยา ภโว, ภวปจฺจยา ชาติ, ชาติปจฺจยา ชรามรณํ, เอวเมตสฺส เกวลสฺส ทุกฺขกฺขนฺธสฺส สมุทโย โหติ. (1-5)}

\prose{อวิชฺชาปจฺจยา สงฺขาโร อวิชฺชาเหตุโก, สงฺขารปจฺจยา วิญฺญาณํ สงฺขารเหตุกํ, วิญฺญาณปจฺจยา นามํ วิญฺญาณเหตุกํ, นามปจฺจยา ผสฺโส นามเหตุโก, ผสฺสปจฺจยา เวทนา ผสฺสเหตุกา, เวทนาปจฺจยา ตณฺหา เวทนาเหตุกา, ตณฺหาปจฺจยา อุปาทานํ ตณฺหาเหตุกํ, อุปาทานปจฺจยา ภโว, ภวปจฺจยา ชาติ, ชาติปจฺจยา ชรามรณํ, เอวเมตสฺส เกวลสฺส ทุกฺขกฺขนฺธสฺส สมุทโย โหติ. (2-6)}

\prose{อวิชฺชาปจฺจยา สงฺขาโร อวิชฺชาเหตุโก, สงฺขารปจฺจยา วิญฺญาณํ สงฺขารเหตุกํ, วิญฺญาณปจฺจยา นามรูปํ วิญฺญาณเหตุกํ, นามรูปปจฺจยา}

\vspace{\tipitakaparskip}
\csromancenter{ปฏิจฺจสมุปฺปาทวิภงฺค ฉฏฺฐายตนํ นามรูปเหตุกํ, ฉฏฺฐายตนปจฺจยา ผสฺโส ฉฏฺฐายตนเหตุโก, ผสฺสปจฺจยา เวทนา ผสฺสเหตุกา, เวทนาปจฺจยา ตณฺหา เวทนาเหตุกา, ตณฺหาปจฺจยา อุปาทานํ ตณฺหาเหตุกํ, อุปาทานปจฺจยา ภโว, ภวปจฺจยา ชาติ, ชาติปจฺจยา ชรามรณํ, เอวเมตสฺส เกวลสฺส ทุกฺขกฺขนฺธสฺส สมุทโย โหติ. (3-7)}
\prose{\csromanfolio{147}อวิชฺชาปจฺจยา สงฺขาโร อวิชฺชาเหตุโก, สงฺขารปจฺจยา วิญฺญาณํ สงฺขารเหตุกํ, วิญฺญาณปจฺจยา นามรูปํ วิญฺญาณเหตุกํ, นามรูปปจฺจยา สฬายตนํ นามรูปเหตุกํ, ฉฏฺฐายตนปจฺจยา ผสฺโส ฉฏฺฐายนเหตุโก, ผสฺสปจฺจยา เวทนา ผสฺสเหตุกา, เวทนาปจฺจยา ตณฺหา เวทนาเหตุกา, ตณฺหาปจฺจยา อุปาทานํ ตณฺหาเหตุกํ, อุปาทานปจฺจยา ภโว, ภวปจฺจยา ชาติ, ชาติปจฺจยา ชรามรณํ.\csromanspacer{} เอวเมตสฺส เกวลสฺส ทุกฺขกฺขนฺธสฺส สมุทโย โหติ. (4-8)}

\csromanheader{2}{เหตุจตุกฺกํ.}
\csromansectionrule
\csromanmatikaanchor{mtk.49}
\csromantocmarkref{subparagraph}{3. สมฺปยุตฺตจตุกฺก}{mtk.49}
\bookmark[dest={mtk.49},level=5]{3. สมฺปยุตฺตจตุกฺก}
\csromanheader{6}{3. สมฺปยุตฺตจตุกฺก}
\proseitem{245}{อวิชฺชาปจฺจยา สงฺขาโร อวิชฺชาสมฺปยุตฺโต, สงฺขารปจฺจยา วิญฺญาณํ สงฺขารสมฺปยุตฺตํ, วิญฺญาณปจฺจยา นามํ วิญฺญาณสมฺปยุตฺตํ, นามปจฺจยา ฉฏฺฐายตนํ นามสมฺปยุตฺตํ, ฉฏฺฐายตนปจฺจยา ผสฺโส ฉฏฺฐายตนสมฺปยุตฺโต, ผสฺสปจฺจยา เวทนา ผสฺสสมฺปยุตฺตา, เวทนาปจฺจยา ตณฺหา เวทนาสมฺปยุตฺตา, ตณฺหาปจฺจยา อุปาทานํ ตณฺหาสมฺปยุตฺตํ, อุปาทานปจฺจยา ภโว, ภวปจฺจยา ชาติ, ชาติปจฺจยา ชรามรณํ, เอวเมตสฺส เกวลสฺส ทุกฺขกฺขนฺธสฺส สมุทโย โหติ. (1-9)}

\prose{อวิชฺชาปจฺจยา สงฺขาโร อวิชฺชาสมฺปยุตฺโต, สงฺขารปจฺจยา วิญฺญาณํ สงฺขารสมฺปยุตฺตํ, วิญฺญาณปจฺจยา นามํ วิญฺญาณสมฺปยุตฺตํ, นามปจฺจยา ผสฺโส นามสมฺปยุตฺโต, ผสฺสปจฺจยา เวทนา ผสฺสสมฺปยุตฺตา, เวทนาปจฺจยา ตณฺหา เวทนาสมฺปยุตฺตา, ตณฺหาปจฺจยา อุปาทานํ ตณฺหาสมฺปยุตฺตํ, อุปาทานปจฺจยา ภโว, ภวปจฺจยา ชาติ, ชาติปจฺจยา ชรามรณํ, เอวเมตสฺส เกวลสฺส ทุกฺขกฺขนฺธสฺส สมุทโย โหติ. (2-10)}

\prose{\csromanfolio{148}อวิชฺชาปจฺจยา สงฺขาโร อวิชฺชาสมฺปยุตฺโต, สงฺขารปจฺจยา วิญฺญาณํ สงฺขารสมฺปยุตฺตํ, วิญฺญาณปจฺจยา นามรูปํ วิญฺญาณสมฺปยุตฺตํ นามํ, นามรูปปจฺจยา ฉฏฺฐายตนํ นามรูปสมฺปยุตฺตํ, ฉฏฺฐายตนปจฺจยา ผสฺโส ฉฏฺฐายตนสมฺปยุตฺโต, ผสฺสปจฺจยา เวทนา ผสฺสสมฺปยุตฺตา, เวทนาปจฺจยา ตณฺหา เวทนาสมฺปยุตฺตา, ตณฺหาปจฺจยา อุปาทานํ ตณฺหาสมฺปยุตฺตํ, อุปาทานปจฺจยา ภโว, ภวปจฺจยา ชาติ, ชาติปจฺจยา ชรามรณํ, เอวเมตสฺส เกวลสฺส ทุกฺขกฺขนฺธสฺส สมุทโย โหติ. (3-11)}

\prose{อวิชฺชาปจฺจยา สงฺขาโร อวิชฺชาสมฺปยุตฺโต, สงฺขารปจฺจยา วิญฺญาณํ สงฺขารสมฺปยุตฺตํ, วิญฺญาณปจฺจยา นามรูปํ วิญฺญาณสมฺปยุตฺตํ นามํ, นามรูปปจฺจยา สฬายตนํ นามสมฺปยุตฺตํ ฉฏฺฐายตนํ, ฉฏฺฐายตนปจฺจยา ผสฺโส ฉฏฺฐายตนสมฺปยุตฺโต, ผสฺสปจฺจยา เวทนา ผสฺสสมฺปยุตฺตา, เวทนาปจฺจยา ตณฺหา เวทนาสมฺปยุตฺตา, ตณฺหาปจฺจยา อุปาทานํ ตณฺหาสมฺปยุตฺตํ, อุปาทานปจฺจยา ภโว, ภวปจฺจยา ชาติ, ชาติปจฺจยา ชรามรณํ, เอวเมตสฺส เกวลสฺส ทุกฺขกฺขนฺธสฺส สมุทโย โหติ. (4-12)}

\csromanheader{2}{สมฺปยุตฺตจตุกฺกํ.}
\csromansectionrule
\csromanmatikaanchor{mtk.50}
\csromantocmarkref{subparagraph}{4. อญฺญมญฺญจตุกฺก}{mtk.50}
\bookmark[dest={mtk.50},level=5]{4. อญฺญมญฺญจตุกฺก}
\csromanheader{6}{4. อญฺญมญฺญจตุกฺก}
\proseitem{246}{อวิชฺชาปจฺจยา สงฺขาโร, สงฺขารปจฺจยาปิ อวิชฺชา.\csromanspacer{} สงฺขารปจฺจยา วิญฺญาณํ, วิญฺญาณปจฺจยาปิ สงฺขาโร.\csromanspacer{} วิญฺญาณปจฺจยา นามํ, นามปจฺจยาปิ วิญฺญาณํ.\csromanspacer{} นามปจฺจยา ฉฏฺฐายตนํ, ฉฏฺฐายตนปจฺจยาปิ นามํ.\csromanspacer{} ฉฏฺฐายตนปจฺจยา ผสฺโส, ผสฺสปจฺจยาปิ ฉฏฺฐายตนํ.\csromanspacer{} ผสฺสปจฺจยา เวทนา, เวทนาปจฺจยาปิ ผสฺโส.\csromanspacer{} เวทนาปจฺจยา ตณฺหา, ตณฺหาปจฺจยาปิ เวทนา.\csromanspacer{} ตณฺหาปจฺจยา อุปาทานํ, อุปาทานปจฺจยาปิ ตณฺหา.\csromanspacer{} อุปาทานปจฺจยา ภโว, ภวปจฺจยา ชาติ, ชาติปจฺจยา ชรามรณํ, เอวเมตสฺส เกวลสฺส ทุกฺขกฺขนฺธสฺส สมุทโย โหติ. (1-13)}

\prose{อวิชฺชาปจฺจยา สงฺขาโร, สงฺขารปจฺจยาปิ อวิชฺชา.\csromanspacer{} สงฺขารปจฺจยา วิญฺญาณํ, วิญฺญาณปจฺจยาปิ สงฺขาโร.\csromanspacer{} วิญฺญาณปจฺจยา นามํ, นามปจฺจยาปิ วิญฺญาณํ.\csromanspacer{} นามปจฺจยา ผสฺโส, ผสฺสปจฺจยาปิ นามํ.\csromanspacer{} ผสฺสปจฺจยา เวทนา,}

\vspace{\tipitakaparskip}
\csromancenter{ปฏิจฺจสมุปฺปาทวิภงฺค เวทนาปจฺจยาปิ ผสฺโส.\csromanspacer{} เวทนาปจฺจยา ตณฺหา, ตณฺหาปจฺจยาปิ เวทนา.\csromanspacer{} ตณฺหาปจฺจยา อุปาทานํ, อุปาทานปจฺจยาปิ ตณฺหา.\csromanspacer{} อุปาทานปจฺจยา ภโว, ภวปจฺจยา ชาติ, ชาติปจฺจยา ชรามรณํ, เอวเมตสฺส เกวลสฺส ทุกฺขกฺขนฺธสฺส สมุทโย โหติ. (2-14)}
\prose{\csromanfolio{149}อวิชฺชาปจฺจยา สงฺขาโร, สงฺขารปจฺจยาปิ อวิชฺชา.\csromanspacer{} สงฺขารปจฺจยา วิญฺญาณํ, วิญฺญาณปจฺจยาปิ สงฺขาโร.\csromanspacer{} วิญฺญาณปจฺจยา นามรูปํ, นามรูปปจฺจยาปิ วิญฺญาณํ.\csromanspacer{} นามรูปปจฺจยา ฉฏฺฐายตนํ, ฉฏฺฐายตนปจฺจยาปิ นามรูปํ.\csromanspacer{} ฉฏฺฐายตนปจฺจยา ผสฺโส, ผสฺสปจฺจยาปิ ฉฏฺฐายตนํ.\csromanspacer{} ผสฺสปจฺจยา เวทนา, เวทนาปจฺจยาปิ ผสฺโส.\csromanspacer{} เวทนาปจฺจยา ตณฺหา, ตณฺหาปจฺจยาปิ เวทนา.\csromanspacer{} ตณฺหาปจฺจยา อุปาทานํ, อุปาทานปจฺจยาปิ ตณฺหา.\csromanspacer{} อุปาทานปจฺจยา ภโว, ภวปจฺจยา ชาติ, ชาติปจฺจยา ชรามรณํ, เอวเมตสฺส เกวลสฺส ทุกฺขกฺขนฺธสฺส สมุทโย โหติ. (3-15)}

\prose{อวิชฺชาปจฺจยา สงฺขาโร, สงฺขารปจฺจยาปิ อวิชฺชา.\csromanspacer{} สงฺขารปจฺจยา วิญฺญาณํ, วิญฺญาณปจฺจยาปิ สงฺขาโร.\csromanspacer{} วิญฺญาณปจฺจยา นามรูปํ, นามรูปปจฺจยาปิ วิญฺญาณํ.\csromanspacer{} นามรูปปจฺจยา สฬายตนํ, ฉฏฺฐายตนปจฺจยาปิ นามรูปํ.\csromanspacer{} ฉฏฺฐายตนปจฺจยา ผสฺโส, ผสฺสปจฺจยาปิ ฉฏฺฐายตนํ.\csromanspacer{} ผสฺสปจฺจยา เวทนา, เวทนาปจฺจยาปิ ผสฺโส.\csromanspacer{} เวทนาปจฺจยา ตณฺหา, ตณฺหาปจฺจยาปิ เวทนา.\csromanspacer{} ตณฺหาปจฺจยา อุปาทานํ, อุปาทานปจฺจยาปิ ตณฺหา.\csromanspacer{} อุปาทานปจฺจยา ภโว, ภวปจฺจยา ชาติ, ชาติปจฺจยา ชรามรณํ, เอวเมตสฺส เกวลสฺส ทุกฺขกฺขนฺธสฺส สมุทโย โหติ. (4-16)}

\csromanheader{2}{อญฺญมญฺญจตุกฺกํ.}
\csromansectionrule
\csromanheader{2}{\textbf{มาติกา}}
\proseitem{247}{สงฺขารปจฺจยา อวิชฺชา ฯเปฯ.\csromanspacer{} วิญฺญาณปจฺจยา อวิชฺชา ฯเปฯ.\csromanspacer{} นามปจฺจยา อวิชฺชา ฯเปฯ.\csromanspacer{} ฉฏฺฐายตนปจฺจยา อวิชฺชา ฯเปฯ.\csromanspacer{} ผสฺสปจฺจยา อวิชฺชา ฯเปฯ.\csromanspacer{} เวทนาปจฺจยา อวิชฺชา ฯเปฯ.\csromanspacer{} ตณฺหาปจฺจยา อวิชฺชา ฯเปฯ.\csromanspacer{} อุปาทานปจฺจยา อวิชฺชา ฯเปฯ.\csromanspacer{} อวิชฺชาปจฺจยา สงฺขาโร, สงฺขารปจฺจยา วิญฺญาณํ, วิญฺญาณปจฺจยา นามํ, นามปจฺจยา ฉฏฺฐายตนํ, ฉฏฺฐายตนปจฺจยา ผสฺโส, ผสฺสปจฺจยา เวทนา, เวทนาปจฺจยา ตณฺหา, \csromanfolio{150}ตณฺหาปจฺจยา อุปาทานํ, อุปาทานปจฺจยา ภโว, ภวปจฺจยา ชาติ, ชาติปจฺจยา ชรามรณํ, เอวเมตสฺส เกวลสฺส ทุกฺขกฺขนฺธสฺส สมุทโย โหติ.}

\csromanheader{2}{มาติกา.}
\csromansectionrule
\csromanmatikaanchor{mtk.51}
\csromantocmarkref{subparagraph}{5. ปจฺจยจตุกฺก (นิทฺเทส)}{mtk.51}
\bookmark[dest={mtk.51},level=5]{5. ปจฺจยจตุกฺก (นิทฺเทส)}
\csromanheader{6}{5. ปจฺจยจตุกฺก}
\proseitem{248}{กตเม ธมฺมา อกุสลา, ยสฺมิํ สมเย อกุสลํ จิตฺตํ อุปฺปนฺนํ โหติ โสมนสฺสสหคตํ ทิฏฺฐิคตสมฺปยุตฺตํ รูปารมฺมณํ วา สทฺทารมฺมณํ วา คนฺธารมฺมณํ วา รสารมฺมณํ วา โผฏฺฐพฺพารมฺมณํ วา ธมฺมารมฺมณํ วา ยํ ยํ วา ปนารพฺภ, ตสฺมิํ สมเย อวิชฺชาปจฺจยา สงฺขาโร, สงฺขารปจฺจยา วิญฺญาณํ, วิญฺญาณปจฺจยา นามํ, นามปจฺจยา ฉฏฺฐายตนํ, ฉฏฺฐายตนปจฺจยา ผสฺโส, ผสฺสปจฺจยา เวทนา, เวทนาปจฺจยา ตณฺหา, ตณฺหาปจฺจยา อุปาทานํ, อุปาทานปจฺจยา ภโว, ภวปจฺจยา ชาติ, ชาติปจฺจยา ชรามรณํ, เอวเมตสฺส เกวลสฺส ทุกฺขกฺขนฺธสฺส สมุทโย โหติ.}

\proseitem{249}{ตตฺถ กตมา อวิชฺชา, ยํ อญฺญาณํ อทสฺสนํ ฯเปฯ อวิชฺชาลงฺคี โมโห อกุสลมูลํ, อยํ วุจฺจติ อวิชฺชา.}

\prose{ตตฺถ กตโม อวิชฺชาปจฺจยา สงฺขาโร, ยา เจตนา สญฺเจตนา สญฺเจตยิตตฺตํ\footnote{เจตยิตตฺตํ (สี, ก)}, อยํ วุจฺจติ อวิชฺชาปจฺจยา สงฺขาโร.}

\prose{ตตฺถ กตมํ สงฺขารปจฺจยา วิญฺญาณํ, ยํ จิตฺตํ มโน มานสํ หทยํ ปณฺฑรํ มโน มนายตนํ มนินฺทฺริยํ วิญฺญาณํ วิญฺญาณกฺขนฺโธ ตชฺชามโนวิญฺญาณธาตุ, อิทํ วุจฺจติ สงฺขารปจฺจยา วิญฺญาณํ.}

\prose{ตตฺถ กตมํ วิญฺญาณปจฺจยา นามํ, เวทนากฺขนฺโธ สญฺญากฺขนฺโธ สงฺขารกฺขนฺโธ, อิทํ วุจฺจติ วิญฺญาณปจฺจยา นามํ.}

\prose{ตตฺถ กตมํ นามปจฺจยา ฉฏฺฐายตนํ, ยํ จิตฺตํ มโน มานสํ หทยํ ปณฺฑรํ มโน มนายตนํ มนินฺทฺริยํ วิญฺญาณํ วิญฺญาณกฺขนฺโธ ตชฺชามโนวิญฺญาณธาตุ, อิทํ วุจฺจติ นามปจฺจยา ฉฏฺฐายตนํ.}

\csromanheader{2}{6. ปฏิจฺจสมุปฺปาทวิภงฺค}
\prose{\csromanfolio{151}ตตฺถ กตโม ฉฏฺฐายตนปจฺจยา ผสฺโส, โย ผสฺโส ผุสนา สมฺผุสนา สมฺผุสิตตฺตํ, อยํ วุจฺจติ ฉฏฺฐายตนปจฺจยา ผสฺโส.}

\prose{ตตฺถ กตมา ผสฺสปจฺจยา เวทนา, ยํ เจตสิกํ สาตํ เจตสิกํ สุขํ เจโตสมฺผสฺสชํ สาตํ สุขํ เวทยิตํ เจโต สมฺผสฺสชา สาตา สุขา เวทนา, อยํ วุจฺจติ ผสฺสปจฺจยา เวทนา.}

\prose{ตตฺถ กตมา เวทนาปจฺจยา ตณฺหา, โย ราโค สาราโค อนุนโย อนุโรโธ นนฺที นนฺทิราโค จิตฺตสฺส สาราโค, อยํ วุจฺจติ เวทนาปจฺจยา ตณฺหา.}

\prose{ตตฺถ กตมํ ตณฺหาปจฺจยา อุปาทานํ, ยา ทิฏฺฐิ ทิฏฺฐิคตํ ทิฏฺฐิคหนํ ทิฏฺฐิกนฺตาโร ทิฏฺฐิวิสูกายิกํ ทิฏฺฐิวิปฺผนฺทิตํ ทิฏฺฐิสญฺโญชนํ คาโห ปติฏฺฐาโห อภินิเวโส ปรามาโส กุมฺมคฺโค มิจฺฉาปโถ มิจฺฉตฺตํ ติตฺถายตนํ วิปริยาสคฺคาโห\footnote{วิปริเยสคฺคาโห (พหูสุ)}, อิทํ วุจฺจติ ตณฺหาปจฺจยา อุปาทานํ.}

\prose{ตตฺถ กตโม อุปาทานปจฺจยา ภโว, ฐเปตฺวา อุปาทานํ เวทนากฺขนฺโธ สญฺญากฺขนฺโธ สงฺขารกฺขนฺโธ วิญฺญาณกฺขนฺโธ, อยํ วุจฺจติ อุปาทานปจฺจยา ภโว.}

\prose{ตตฺถ กตมา ภวปจฺจยา ชาติ, ยา เตสํ เตสํ ธมฺมานํ ชาติ สญฺชาติ นิพฺพตฺติ อภินิพฺพตฺติ ปาตุภาโว, อยํ วุจฺจติ ภวปจฺจยา ชาติ.}

\prose{ตตฺถ กตมํ ชาติปจฺจยา ชรามรณํ, อตฺถิ ชรา, อตฺถิ มรณํ.\csromanspacer{} ตตฺถ กตมา ชรา, ยา เตสํ เตสํ ธมฺมานํ ชรา ชีรณตา อายุโน สํหานิ, อยํ วุจฺจติ ชรา.\csromanspacer{} ตตฺถ กตมํ มรณํ, โย เตสํ เตสํ ธมฺมานํ ขโย วโย เภโท ปริเภโท อนิจฺจตา อนฺตรธานํ, อิทํ วุจฺจติ มรณํ.\csromanspacer{} อิติ อยญฺจ ชรา อิทญฺจ มรณํ, อิทํ วุจฺจติ ชาติปจฺจยา ชรามรณํ.}

\prose{เอวเมตสฺส เกวลสฺส ทุกฺขกฺขนฺธสฺส สมุทโย โหตีติ เอวเมตสฺส เกวลสฺส ทุกฺขกฺขนฺธสฺส สงฺคติ โหติ, สมาคโม โหติ, สโมธานํ โหติ, ปาตุภาโว โหติ, เตน วุจฺจติ “เอวเมตสฺส เกวลสฺส ทุกฺขกฺขนฺธสฺส สมุทโย โหตี”ติ. (1)}

\proseitem{250}{\csromanfolio{152}ตสฺมิํ สมเย อวิชฺชาปจฺจยา สงฺขาโร, สงฺขารปจฺจยา วิญฺญาณํ, วิญฺญาณปจฺจยา นามํ, นามปจฺจยา ผสฺโส, ผสฺสปจฺจยา เวทนา, เวทนาปจฺจยา ตณฺหา, ตณฺหาปจฺจยา อุปาทานํ, อุปาทานปจฺจยา ภโว, ภวปจฺจยา ชาติ, ชาติปจฺจยา ชรามรณํ, เอวเมตสฺส เกวลสฺส ทุกฺขกฺขนฺธสฺส สมุทโย โหติ.}

\proseitem{251}{ตตฺถ กตมา อวิชฺชา, ยํ อญฺญาณํ อทสฺสนํ ฯเปฯ อวิชฺชาลงฺคี โมโห อกุสลมูลํ, อยํ วุจฺจติ อวิชฺชา.}

\prose{ตตฺถ กตโม อวิชฺชาปจฺจยา สงฺขาโร, ยา เจตนา สญฺเจตนา สญฺเจตยิตตฺตํ, อยํ วุจฺจติ อวิชฺชาปจฺจยา สงฺขาโร.}

\prose{ตตฺถ กตมํ สงฺขารปจฺจยา วิญฺญาณํ, ยํ จิตฺตํ มโน มานสํ ฯเปฯ ตชฺชามโนวิญฺญาณธาตุ, อิทํ วุจฺจติ สงฺขารปจฺจยา วิญฺญาณํ.}

\prose{ตตฺถ กตมํ วิญฺญาณปจฺจยา นามํ, เวทนากฺขนฺโธ สญฺญากฺขนฺโธ สงฺขารกฺขนฺโธ, อิทํ วุจฺจติ วิญฺญาณปจฺจยา นามํ.}

\prose{นามปจฺจยา ผสฺโสติ ตตฺถ กตมํ นามํ, ฐเปตฺวา ผสฺสํ เวทนากฺขนฺโธ สญฺญากฺขนฺโธ สงฺขารกฺขนฺโธ วิญฺญาณกฺขนฺโธ, อิทํ วุจฺจติ นามํ.}

\prose{ตตฺถ กตโม นามปจฺจยา ผสฺโส, โย ผสฺโส ผุสนา สมฺผุสนา สมฺผุสิตตฺตํ, อยํ วุจฺจติ นามปจฺจยา ผสฺโส ฯเปฯ เตน วุจฺจติ “เอวเมตสฺส เกวลสฺส ทุกฺขกฺขนฺธสฺส สมุทโย โหตี”ติ. (2)}

\proseitem{252}{ตสฺมิํ สมเย อวิชฺชาปจฺจยา สงฺขาโร, สงฺขารปจฺจยา วิญฺญาณํ, วิญฺญาณปจฺจยา นามรูปํ, นามรูปปจฺจยา ฉฏฺฐายตนํ, ฉฏฺฐายตนปจฺจยา ผสฺโส, ผสฺสปจฺจยา เวทนา, เวทนาปจฺจยา ตณฺหา, ตณฺหาปจฺจยา อุปาทานํ, อุปาทานปจฺจยา ภโว, ภวปจฺจยา ชาติ, ชาติปจฺจยา ชรามรณํ, เอวเมตสฺส เกวลสฺเส ทุกฺขกฺขนฺธสฺส สมุทโย โหติ.}

\proseitem{253}{ตตฺถ กตมา อวิชฺชา, ยํ อญฺญาณํ อทสฺสนํ ฯเปฯ อวิชฺชาลงฺคี โมโห อกุสลมูลํ, อยํ วุจฺจติ อวิชฺชา.}

\prose{ตตฺถ กตโม อวิชฺชาปจฺจยา สงฺขาโร, ยา เจตนา สญฺเจตนา สญฺเจตยิตตฺตํ, อยํ วุจฺจติ อวิชฺชาปจฺจยา สงฺขาโร.}

\csromanheader{2}{6. ปฏิจฺจสมุปฺปาทวิภงฺค}
\prose{\csromanfolio{153}ตตฺถ กตมํ สงฺขารปจฺจยา วิญฺญาณํ, ยํ จิตฺตํ มโน มานสํ ฯเปฯ ตชฺชามโนวิญฺญาณธาตุ, อิทํ วุจฺจติ สงฺขารปจฺจยา วิญฺญาณํ.}

\prose{ตตฺถ กตมํ วิญฺญาณปจฺจยา นามรูปํ, อตฺถิ นามํ, อตฺถิ รูปํ.\csromanspacer{} ตตฺถ กตมํ นามํ, เวทนากฺขนฺโธ สญฺญากฺขนฺโธ สงฺขารกฺขนฺโธ, อิทํ วุจฺจติ นามํ.\csromanspacer{} ตตฺถ กตมํ รูปํ, จกฺขายตนสฺส อุปจโย, โสตายตนสฺส อุปจโย, ฆานายตนสฺส อุปจโย, ชิวฺหายตนสฺส อุปจโย, กายายตนสฺส อุปจโย, ยํ วา ปนญฺญมฺปิ อตฺถิ รูปํ จิตฺตชํ จิตฺตเหตุกํ จิตฺตสมุฏฺฐานํ, อิทํ วุจฺจติ รูปํ.\csromanspacer{} อิติ อิทญฺจ นามํ อิทญฺจ รูปํ, อิทํ วุจฺจติ วิญฺญาณปจฺจยา นามรูปํ.}

\prose{นามรูปปจฺจยา ฉฏฺฐายตนนฺติ อตฺถิ นามํ, อตฺถิ รูปํ.\csromanspacer{} ตตฺถ กตมํ นามํ, เวทนากฺขนฺโธ สญฺญากฺขนฺโธ สงฺขารกฺขนฺโธ, อิทํ วุจฺจติ นามํ.\csromanspacer{} ตตฺถ กตมํ รูปํ, ยํ รูปํ นิสฺสาย มโนวิญฺญาณธาตุ วตฺตติ, อิทํ วุจฺจติ รูปํ.\csromanspacer{} อิติ อิทญฺจ นามํ อิทญฺจ รูปํ, อิทํ วุจฺจติ นามรูปํ.}

\prose{ตตฺถ กตมํ นามรูปปจฺจยา ฉฏฺฐายตนํ, ยํ จิตฺตํ มโน มานสํ ฯเปฯ ตชฺชามโนวิญฺญาณธาตุ, อิทํ วุจฺจติ นามรูปปจฺจยา ฉฏฺฐายตนํ.}

\prose{ตตฺถ กตโม ฉฏฺฐายตนปจฺจยา ผสฺโส, โย ผสฺโส ผุสนา สมฺผุสนา สมฺผุสิตตฺตํ, อยํ วุจฺจติ ฉฏฺฐายตนปจฺจยา ผสฺโส ฯเปฯ เตน วุจฺจติ “เอวเมตสฺส เกวลสฺส ทุกฺขกฺขนฺธสฺส สมุทโย โหตี”ติ. (3)}

\proseitem{254}{ตสฺมิํ สมเย อวิชฺชาปจฺจยา สงฺขาโร, สงฺขารปจฺจยา วิญฺญาณํ, วิญฺญาณปจฺจยา นามรูปํ, นามรูปปจฺจยา สฬายตนํ, ฉฏฺฐายตนปจฺจยา ผสฺโส, ผสฺสปจฺจยา เวทนา, เวทนาปจฺจยา ตณฺหา, ตณฺหาปจฺจยา อุปาทานํ, อุปาทานปจฺจยา ภโว, ภวปจฺจยา ชาติ, ชาติปจฺจยา ชรามรณํ, เอวเมตสฺส เกวลสฺส ทุกฺขกฺขนฺธสฺส สมุทโย โหติ.}

\proseitem{255}{ตตฺถ กตมา อวิชฺชา, ยํ อญฺญาณํ อทสฺสนํ ฯเปฯ อวิชฺชาลงฺคี โมโห อกุสลมูลํ, อยํ วุจฺจติ อวิชฺชา.}

\prose{ตตฺถ กตโม อวิชฺชาปจฺจยา สงฺขาโร, ยา เจตนา สญฺเจตนา สญฺเจตยิตตฺตํ, อยํ วุจฺจติ อวิชฺชาปจฺจยา สงฺขาโร.}

\prose{\csromanfolio{154}ตตฺถ กตมํ สงฺขารปจฺจยา วิญฺญาณํ, ยํ จิตฺตํ มโน มานสํ ฯเปฯ ตชฺชามโนวิญฺญาณธาตุ.\csromanspacer{} อิทํ วุจฺจติ สงฺขารปจฺจยา วิญฺญาณํ.}

\prose{ตตฺถ กตมํ วิญฺญาณปจฺจยา นามรูปํ, อตฺถิ นามํ, อตฺถิ รูปํ.\csromanspacer{} ตตฺถ กตมํ นามํ, เวทนากฺขนฺโธ สญฺญากฺขนฺโธ สงฺขารกฺขนฺโธ, อิทํ วุจฺจติ นามํ.\csromanspacer{} ตตฺถ กตมํ รูปํ, จกฺขายตนสฺส อุปจโย, โสตายตนสฺส อุปจโย, ฆานายตนสฺส อุปจโย, ชิวฺหายตนสฺส อุปจโย, กายายตนสฺส อุปจโย, ยํ วา ปนญฺญมฺปิ อตฺถิ รูปํ จิตฺตชํ จิตฺตเหตุกํ จิตฺตสมุฏฺฐานํ, อิทํ วุจฺจติ รูปํ.\csromanspacer{} อิติ อิทญฺจ นามํ อิทญฺจ รูปํ, อิทํ วุจฺจติ วิญฺญาณปจฺจยา นามรูปํ.}

\prose{นามรูปปจฺจยา สฬายตนนฺติ อตฺถิ นามํ, อตฺถิ รูปํ.\csromanspacer{} ตตฺถ กตมํ นามํ, เวทนากฺขนฺโธ สญฺญากฺขนฺโธ สงฺขารกฺขนฺโธ, อิทํ วุจฺจติ นามํ.\csromanspacer{} ตตฺถ กตมํ รูปํ, จตฺตาโร จ มหาภูตา ยญฺจ รูปํ นิสฺสาย มโนวิญฺญาณธาตุ วตฺตติ, อิทํ วุจฺจติ รูปํ.\csromanspacer{} อิติ อิทญฺจ นามํ อิทญฺจ รูปํ, อิทํ วุจฺจติ นามรูปํ.}

\prose{ตตฺถ กตมํ นามรูปปจฺจยา สฬายตนํ, จกฺขายตนํ โสตายตนํ ฆานายตนํ ชิวฺหายตนํ กายายตนํ มนายตนํ, อิทํ วุจฺจติ นามรูปปจฺจยา สฬายตนํ.}

\prose{ตตฺถ กตโม ฉฏฺฐายตนปจฺจยา ผสฺโส, โย ผสฺโส ผุสนา สมฺผุสนา สมฺผุสิตตฺตํ, อยํ วุจฺจติ ฉฏฺฐายตนปจฺจยา ผสฺโส ฯเปฯ เตน วุจฺจติ “เอวเมตสฺส เกวลสฺส ทุกฺขกฺขนฺธสฺส สมุทโย โหตี”ติ. (4)}

\csromanheader{2}{ปจฺจยจตุกฺกํ.}
\csromansectionrule
\csromanmatikaanchor{mtk.52}
\csromantocmarkref{subparagraph}{6. เหตุจตุกฺก (นิทฺเทส)}{mtk.52}
\bookmark[dest={mtk.52},level=5]{6. เหตุจตุกฺก (นิทฺเทส)}
\csromanheader{6}{6. เหตุจตุกฺก}
\proseitem{256}{ตสฺมิํ สมเย อวิชฺชาปจฺจยา สงฺขาโร อวิชฺชาเหตุโก, สงฺขารปจฺจยา วิญฺญาณํ สงฺขารเหตุกํ, วิญฺญาณปจฺจยา นามํ วิญฺญาณเหตุกํ, นามปจฺจยา ฉฏฺฐายตนํ นามเหตุกํ, ฉฏฺฐายตนปจฺจยา ผสฺโส ฉฏฺฐายตนเหตุโก, ผสฺสปจฺจยา เวทนา ผสฺสเหตุกา, เวทนาปจฺจยา ตณฺหา เวทนาเหตุกา, ตณฺหาปจฺจยา อุปาทานํ ตณฺหาเหตุกํ, อุปาทานปจฺจยา ภโว, ภวปจฺจยา ชาติ, ชาติปจฺจยา ชรามรณํ, เอวเมตสฺส เกวลสฺส ทุกฺขกฺขนฺธสฺส สมุทโย โหติ.}

\csromanheader{2}{6. ปฏิจฺจสมุปฺปาทวิภงฺค}
\proseitem{257}{\csromanfolio{155}ตตฺถ กตมา อวิชฺชา, ยํ อญฺญาณํ อทสฺสนํ ฯเปฯ อวิชฺชาลงฺคี โมโห อกุสลมูลํ, อยํ วุจฺจติ อวิชฺชา.}

\prose{ตตฺถ กตโม อวิชฺชาปจฺจยา สงฺขาโร อวิชฺชาเหตุโก, ยา เจตนา สญฺเจตนา สญฺเจตยิตตฺตํ, อยํ วุจฺจติ อวิชฺชาปจฺจยา สงฺขาโร อวิชฺชาเหตุโก.}

\prose{ตตฺถ กตมํ สงฺขารปจฺจยา วิญฺญาณํ สงฺขารเหตุกํ, ยํ จิตฺตํ มโน มานสํ ฯเปฯ ตชฺชามโนวิญฺญาณธาตุ, อิทํ วุจฺจติ สงฺขารปจฺจยา วิญฺญาณํ สงฺขารเหตุกํ.}

\prose{ตตฺถ กตมํ วิญฺญาณปจฺจยา นามํ วิญฺญาณเหตุกํ, เวทนากฺขนฺโธ สญฺญากฺขนฺโธ สงฺขารกฺขนฺโธ, อิทํ วุจฺจติ วิญฺญาณปจฺจยา นามํ วิญฺญาณเหตุกํ.}

\prose{ตตฺถ กตมํ นามปจฺจยา ฉฏฺฐายตนํ นามเหตุกํ, ยํ จิตฺตํ มโน มานสํ ฯเปฯ ตชฺชามโนวิญฺญาณธาตุ, อิทํ วุจฺจติ นามปจฺจยา ฉฏฺฐายตนํ นามเหตุกํ.}

\prose{ตตฺถ กตโม ฉฏฺฐายตนปจฺจยา ผสฺโส ฉฏฺฐายตนเหตุโก, โย ผสฺโส ผุสนา สมฺผุสนา สมฺผุสิตตฺตํ, อยํ วุจฺจติ ฉฏฺฐายตนปจฺจยา ผสฺโส ฉฏฺฐายตนเหตุโก.}

\prose{ตตฺถ กตมา ผสฺสปจฺจยา เวทนา ผสฺสเหตุกา, ยํ เจตสิกํ สาตํ เจตสิกํ สุขํ เจโตสมฺผสฺสชํ สาตํ สุขํ เวทยิตํ เจโตสมฺผสฺสชา สาตา สุขา เวทนา, อยํ วุจฺจติ ผสฺสปจฺจยา เวทนา ผสฺสเหตุกา.}

\prose{ตตฺถ กตมา เวทนาปจฺจยา ตณฺหา เวทนาเหตุกา, โย ราโค สาราโค ฯเปฯ จิตฺตสฺส สาราโค, อยํ วุจฺจติ เวทนาปจฺจยา ตณฺหา เวทนาเหตุกา.}

\prose{ตตฺถ กตมํ ตณฺหาปจฺจยา อุปาทานํ ตณฺหาเหตุกํ, ยา ทิฏฺฐิ ทิฏฺฐิคตํ ฯเปฯ ติตฺถายตนํ วิปริยาสคฺคาโห, อิทํ วุจฺจติ ตณฺหาปจฺจยา อุปาทานํ ตณฺหาเหตุกํ ฯเปฯ เตน วุจฺจติ “เอวเมตสฺส เกวลสฺส ทุกฺขกฺขนฺธสฺส สมุทโย โหตี”ติ. (1-5)}

\proseitem{258}{\csromanfolio{156}ตสฺมิํ สมเย อวิชฺชาปจฺจยา สงฺขาโร อวิชฺชาเหตุโก, สงฺขารปจฺจยา วิญฺญาณํ สงฺขารเหตุกํ, วิญฺญาณปจฺจยา นามํ วิญฺญาณเหตุกํ, นามปจฺจยา ผสฺโส นามเหตุโก, ผสฺสปจฺจยา เวทนา ผสฺสเหตุกา, เวทนาปจฺจยา ตณฺหา เวทนาเหตุกา, ตณฺหาปจฺจยา อุปาทานํ ตณฺหาเหตุกํ, อุปาทานปจฺจยา ภโว, ภวปจฺจยา ชาติ, ชาติปจฺจยา ชรามรณํ, เอวเมตสฺส เกวลสฺส ทุกฺขกฺขนฺธสฺส สมุทโย โหติ.}

\proseitem{259}{ตตฺถ กตมา อวิชฺชา, ยํ อญฺญาณํ อทสฺสนํ ฯเปฯ อวิชฺชาลงฺคี โมโห อกุสลมูลํ, อยํ วุจฺจติ อวิชฺชา.}

\prose{ตตฺถ กตโม อวิชฺชาปจฺจยา สงฺขาโร อวิชฺชาเหตุโก, ยา เจตนา สญฺเจตนา สญฺเจตยิตตฺตํ, อยํ วุจฺจติ อวิชฺชาปจฺจยา สงฺขาโร อวิชฺชาเหตุโก.}

\prose{ตตฺถ กตมํ สงฺขารปจฺจยา วิญฺญาณํ สงฺขารเหตุกํ, ยํ จิตฺตํ มโน มานสํ ฯเปฯ ตชฺชามโนวิญฺญาณธาตุ, อิทํ วุจฺจติ สงฺขารปจฺจยา วิญฺญาณํ สงฺขารเหตุกํ.}

\prose{ตตฺถ กตมํ วิญฺญาณปจฺจยา นามํ วิญฺญาณเหตุกํ, เวทนากฺขนฺโธ สญฺญากฺขนฺโธ สงฺขารกฺขนฺโธ, อิทํ วุจฺจติ วิญฺญาณปจฺจยา นามํ วิญฺญาณเหตุกํ.}

\prose{นามปจฺจยา ผสฺโส นามเหตุโกติ ตตฺถ กตมํ นามํ, ฐเปตฺวา ผสฺสํ เวทนากฺขนฺโธ สญฺญากฺขนฺโธ สงฺขารกฺขนฺโธ วิญฺญาณกฺขนฺโธ, อิทํ วุจฺจติ นามํ.}

\prose{ตตฺถ กตโม นามปจฺจยา ผสฺโส นามเหตุโก, โย ผสฺโส ผุสนา สมฺผุสนา สมฺผุสิตตฺตํ, อยํ วุจฺจติ นามปจฺจยา ผสฺโส นามเหตุโก ฯเปฯ เตน วุจฺจติ “เอวเมตสฺส เกวลสฺส ทุกฺขกฺขนฺธสฺส สมุทโย โหตี”ติ. (2-6)}

\proseitem{260}{ตสฺมิํ สมเย อวิชฺชาปจฺจยา สงฺขาโร อวิชฺชาเหตุโก, สงฺขารปจฺจยา วิญฺญาณํ สงฺขารเหตุกํ, วิญฺญาณปจฺจยา นามรูปํ วิญฺญาณเหตุกํ, นามรูปปจฺจยา ฉฏฺฐายตนํ นามรูปเหตุกํ, ฉฏฺฐายตนปจฺจยา ผสฺโส ฉฏฺฐายตนเหตุโก, ผสฺสปจฺจยา เวทนา ผสฺสเหตุกา,}

\vspace{\tipitakaparskip}
\csromancenter{ปฏิจฺจสมุปฺปาทวิภงฺค เวทนาปจฺจยา ตณฺหา เวทนาเหตุกา, ตณฺหาปจฺจยา อุปาทานํ ตณฺหาเหตุกํ, อุปาทานปจฺจยา ภโว, ภวปจฺจยา ชาติ, ชาติปจฺจยา ชรามรณํ, เอวเมตสฺส เกวลสฺส ทุกฺขกฺขนฺธสฺส สมุทโย โหติ.}
\proseitem{261}{\csromanfolio{157}ตตฺถ กตมา อวิชฺชา, ยํ อญฺญาณํ อทสฺสนํ ฯเปฯ อวิชฺชาลงฺคี โมโห อกุสลมูลํ, อยํ วุจฺจติ อวิชฺชา.}

\prose{ตตฺถ กตโม อวิชฺชาปจฺจยา สงฺขาโร อวิชฺชาเหตุโก, ยา เจตนา สญฺเจตนา สญฺเจตยิตตฺตํ, อยํ วุจฺจติ อวิชฺชาปจฺจยา สงฺขาโร อวิชฺชาเหตุโก.}

\prose{ตตฺถ กตมํ สงฺขารปจฺจยา วิญฺญาณํ สงฺขารเหตุกํ, ยํ จิตฺตํ มโน มานสํ ฯเปฯ ตชฺชามโนวิญฺญาณธาตุ, อิทํ วุจฺจติ สงฺขารปจฺจยา วิญฺญาณํ สงฺขารเหตุกํ.}

\prose{ตตฺถ กตมํ วิญฺญาณปจฺจยา นามรูปํ วิญฺญาณเหตุกํ, อตฺถิ นามํ, อตฺถิ รูปํ.\csromanspacer{} ตตฺถ กตมํ นามํ, เวทนากฺขนฺโธ สญฺญากฺขนฺโธ สงฺขารกฺขนฺโธ, อิทํ วุจฺจติ นามํ.\csromanspacer{} ตตฺถ กตมํ รูปํ, จกฺขายตนสฺส อุปจโย, โสตายตนสฺส อุปจโย, ฆานายตนสฺส อุปจโย, ชิวฺหายตนสฺส อุปจโย, กายายตนสฺส อุปจโย, ยํ วา ปนญฺญมฺปิ อตฺถิ รูปํ จิตฺตชํ จิตฺตเหตุกํ จิตฺตสมุฏฺฐานํ, อิทํ วุจฺจติ รูปํ.\csromanspacer{} อิติ อิทญฺจ นามํ อิทญฺจ รูปํ, อิทํ วุจฺจติ วิญฺญาณปจฺจยา นามรูปํ วิญฺญาณเหตุกํ.}

\prose{นามรูปปจฺจยา ฉฏฺฐายตนํ นามรูปเหตุกนฺติ อตฺถิ นามํ, อตฺถิ รูปํ.\csromanspacer{} ตตฺถ กตมํ นามํ, เวทนากฺขนฺโธ สญฺญากฺขนฺโธ สงฺขารกฺขนฺโธ, อิทํ วุจฺจติ นามํ.\csromanspacer{} ตตฺถ กตมํ รูปํ, ยํ รูปํ นิสฺสาย มโนวิญฺญาณธาตุ วตฺตติ, อิทํ วุจฺจติ รูปํ.\csromanspacer{} อิติ อิทญฺจ นามํ อิทญฺจ รูปํ, อิทํ วุจฺจติ นามรูปํ.}

\prose{ตตฺถ กตมํ นามรูปปจฺจยา ฉฏฺฐายตนํ นามรูปเหตุกํ, ยํ จิตฺตํ มโน มานสํ ฯเปฯ ตชฺชามโนวิญฺญาณธาตุ, อิทํ วุจฺจติ นามรูปปจฺจยา ฉฏฺฐายตนํ นามรูปเหตุกํ.}

\prose{ตตฺถ กตโม ฉฏฺฐายตนปจฺจยา ผสฺโส ฉฏฺฐายตนเหตุโก, โย ผสฺโส ผุสนา สมฺผุสนา สมฺผุสิตตฺตํ, อยํ วุจฺจติ ฉฏฺฐายตนปจฺจยา ผสฺโส ฉฏฺฐายตนเหตุโก ฯเปฯ เตน วุจฺจติ “เอวเมตสฺส เกวลสฺส ทุกฺขกฺขนฺธสฺส สมุทโย โหตี”ติ. (3-7)}

\proseitem{262}{\csromanfolio{158}ตสฺมิํ สมเย อวิชฺชาปจฺจยา สงฺขาโร อวิชฺชาเหตุโก, สงฺขารปจฺจยา วิญฺญาณํ สงฺขารเหตุกํ, วิญฺญาณปจฺจยา นามรูปํ วิญฺญาณเหตุกํ, นามรูปปจฺจยา สฬายตนํ นามรูปเหตุกํ, ฉฏฺฐายตนปจฺจยา ผสฺโส ฉฏฺฐายตนเหตุโก, ผสฺสปจฺจยา เวทนา ผสฺสเหตุกา, เวทนาปจฺจยา ตณฺหา เวทนาเหตุกา, ตณฺหาปจฺจยา อุปาทานํ ตณฺหาเหตุกํ, อุปาทานปจฺจยา ภโว, ภวปจฺจยา ชาติ, ชาติปจฺจยา ชรามรณํ, เอวเมตสฺส เกวลสฺส ทุกฺขกฺขนฺธสฺส สมุทโย โหติ.}

\proseitem{263}{ตตฺถ กตมา อวิชฺชา, ยํ อญฺญาณํ อทสฺสนํ ฯเปฯ อวิชฺชาลงฺคี โมโห อกุสลมูลํ, อยํ วุจฺจติ อวิชฺชา.}

\prose{ตตฺถ กตโม อวิชฺชาปจฺจยา สงฺขาโร อวิชฺชาเหตุโก, ยา เจตนา สญฺเจตนา สญฺเจตยิตตฺตํ, อยํ วุจฺจติ อวิชฺชาปจฺจยา สงฺขาโร อวิชฺชาเหตุโก.}

\prose{ตตฺถ กตมํ สงฺขารปจฺจยา วิญฺญาณํ สงฺขารเหตุกํ, ยํ จิตฺตํ มโน มานสํ ฯเปฯ ตชฺชามโนวิญฺญาณธาตุ, อิทํ วุจฺจติ สงฺขารปจฺจยา วิญฺญาณํ สงฺขารเหตุกํ.}

\prose{ตตฺถ กตมํ วิญฺญาณปจฺจยา นามรูปํ วิญฺญาณเหตุกํ, อตฺถิ นามํ, อตฺถิ รูปํ.\csromanspacer{} ตตฺถ กตมํ นามํ, เวทนากฺขนฺโธ สญฺญากฺขนฺโธ สงฺขารกฺขนฺโธ อิทํ วุจฺจติ นามํ.\csromanspacer{} ตตฺถ กตมํ รูปํ, จกฺขายตนสฺส อุปจโย, โสตายตนสฺส อุปจโย, ฆานายตนสฺส อุปจโย, ชิวฺหายตนสฺส อุปจโย, กายายตนสฺส อุปจโย, ยํ วา ปนญฺญมฺปิ อตฺถิ รูปํ จิตฺตชํ จิตฺตเหตุกํ จิตฺตสมุฏฺฐานํ, อิทํ วุจฺจติ รูปํ.\csromanspacer{} อิติ อิทญฺจ นามํ อิทญฺจ รูปํ, อิทํ วุจฺจติ วิญฺญาณปจฺจยา นามรูปํ วิญฺญาณเหตุกํ.}

\prose{นามรูปปจฺจยา สฬายตนํ นามรูปเหตุกนฺติ อตฺถิ นามํ, อตฺถิ รูปํ.\csromanspacer{} ตตฺถ กตมํ นามํ, เวทนากฺขนฺโธ สญฺญากฺขนฺโธ สงฺขารกฺขนฺโธ, อิทํ วุจฺจติ นามํ.\csromanspacer{} ตตฺถ กตมํ รูปํ, จตฺตาโร จ มหาภูตา ยญฺจ รูปํ นิสฺสาย มโนวิญฺญาณธาตุ วตฺตติ, อิทํ วุจฺจติ รูปํ.\csromanspacer{} อิติ อิทญฺจ นามํ อิทญฺจ รูปํ, อิทํ วุจฺจติ นามรูปํ.}

\prose{ตตฺถ กตมํ นามรูปปจฺจยา สฬายตนํ นามรูปเหตุกํ, จกฺขายตนํ โสตายตนํ ฆานายตนํ ชิวฺหายตนํ กายายตนํ มนายตนํ, อิทํ วุจฺจติ นามรูปปจฺจยา สฬายตนํ นามรูปเหตุกํ.}

\csromanheader{2}{6. ปฏิจฺจสมุปฺปาทวิภงฺค}
\prose{\csromanfolio{159}ตตฺถ กตโม ฉฏฺฐายตนปจฺจยา ผสฺโส ฉฏฺฐายตนเหตุโก, โย ผสฺโส ผุสนา สมฺผุสนา สมฺผุสิตตฺตํ, อยํ วุจฺจติ ฉฏฺฐายตนปจฺจยา ผสฺโส ฉฏฺฐายตนเหตุโก.}

\prose{ตตฺถ กตมา ผสฺสปจฺจยา เวทนา ผสฺสเหตุกา, ยํ เจตสิกํ สาตํ เจตสิกํ สุขํ เจโตสมฺผสฺสชํ สาตํ สุขํ เวทยิตํ เจโตสมฺผสฺสชา สาตา สุขา เวทนา, อยํ วุจฺจติ ผสฺสปจฺจยา เวทนา ผสฺสเหตุกา.}

\prose{ตตฺถ กตมา เวทนาปจฺจยา ตณฺหา เวทนาเหตุกา, โย ราโค สาราโค ฯเปฯ จิตฺตสฺส สาราโค, อยํ วุจฺจติ เวทนาปจฺจยา ตณฺหา เวทนาเหตุกา.}

\prose{ตตฺถ กตมํ ตณฺหาปจฺจยา อุปาทานํ ตณฺหาเหตุกํ, ยา ทิฏฺฐิ ทิฏฺฐิคตํ ฯเปฯ ติตฺถายตนํ วิปริยาสคฺคาโห, อิทํ วุจฺจติ ตณฺหาปจฺจยา อุปาทานํ ตณฺหาเหตุกํ ฯเปฯ เตน วุจฺจติ “เอวเมตสฺส เกวลสฺส ทุกฺขกฺขนฺธสฺส สมุทโย โหตี”ติ. (4-8)}

\csromanheader{2}{เหตุจตุกฺกํ.}
\csromansectionrule
\csromanmatikaanchor{mtk.53}
\csromantocmarkref{subparagraph}{7. สมฺปยุตฺตจตุกฺก (นิทฺเทส)}{mtk.53}
\bookmark[dest={mtk.53},level=5]{7. สมฺปยุตฺตจตุกฺก (นิทฺเทส)}
\csromanheader{6}{7. สมฺปยุตฺตจตุกฺก}
\proseitem{264}{ตสฺมิํ สมเย อวิชฺชาปจฺจยา สงฺขาโร อวิชฺชาสมฺปยุตฺโต, สงฺขารปจฺจยา วิญฺญาณํ สงฺขารสมฺปยุตฺตํ, วิญฺญาณปจฺจยา นามํ วิญฺญาณสมฺปยุตฺตํ, นามปจฺจยา ฉฏฺฐายตนํ นามสมฺปยุตฺตํ, ฉฏฺฐายตนปจฺจยา ผสฺโส ฉฏฺฐายตนสมฺปยุตฺโต, ผสฺสปจฺจยา เวทนา ผสฺสสมฺปยุตฺตา, เวทนาปจฺจยา ตณฺหา เวทนาสมฺปยุตฺโต, ตณฺหาปจฺจยา อุปาทานํ ตณฺหาสมฺปยุตฺตํ, อุปาทานปจฺจยา ภโว, ภวปจฺจยา ชาติ, ชาติปจฺจยา ชรามรณํ, เอวเมตสฺส เกวลสฺส ทุกฺขกฺขนฺธสฺส สมุทโย โหติ.}

\proseitem{265}{ตตฺถ กตมา อวิชฺชา, ยํ อญฺญาณํ อทสฺสนํ ฯเปฯ อวิชฺชาลงฺคี โมโห อกุสลมูลํ, อยํ วุจฺจติ อวิชฺชา.}

\prose{\csromanfolio{160}ตตฺถ กตโม อวิชฺชาปจฺจยา สงฺขาโร อวิชฺชาสมฺปยุตฺโต, ยา เจตนา สญฺเจตนา สญฺเจตยิตตฺตํ, อยํ วุจฺจติ อวิชฺชาปจฺจยา สงฺขาโร อวิชฺชาสมฺปยุตฺโต.}

\prose{ตตฺถ กตมํ สงฺขารปจฺจยา วิญฺญาณํ สงฺขารสมฺปยุตฺตํ, ยํ จิตฺตํ มโน มานสํ ฯเปฯ ตชฺชามโนวิญฺญาณธาตุ, อิทํ วุจฺจติ สงฺขารปจฺจยา วิญฺญาณํ สงฺขารสมฺปยุตฺตํ.}

\prose{ตตฺถ กตมํ วิญฺญาณปจฺจยา นามํ วิญฺญาณสมฺปยุตฺตํ, เวทนากฺขนฺโธ สญฺญากฺขนฺโธ สงฺขารกฺขนฺโธ, อิทํ วุจฺจติ วิญฺญาณปจฺจยา นามํ วิญฺญาณสมฺปยุตฺตํ.}

\prose{ตตฺถ กตมํ นามปจฺจยา ฉฏฺฐายนํ ฯเปฯ นามสมฺปยุตฺตํ, ยํ จิตฺตํ มโน มานสํ ฯเปฯ ตชฺชามโนวิญฺญาณธาตุ, อิทํ วุจฺจติ นามปจฺจยา ฉฏฺฐายตนํ นามสมฺปยุตฺตํ.}

\prose{ตตฺถ กตโม ฉฏฺฐายตนปจฺจยา ผสฺโส ฉฏฺฐายตนสมฺปยุตฺโต, โย ผสฺโส ผุสนา สมฺผุสนา สมฺผุสิตตฺตํ, อยํ วุจฺจติ ฉฏฺฐายตนปจฺจยา ผสฺโส ฉฏฺฐายตนสมฺปยุตฺโต.}

\prose{ตตฺถ กตมา ผสฺสปจฺจยา เวทนา ผสฺสสมฺปยุตฺตา, ยํ เจตสิกํ สาตํ เจตสิกํ สุขํ เจโตสมฺผสฺสชํ สาตํ สุขํ เวทยิตํ เจโตสมฺผสฺสชา สาตา สุขา เวทนา, อยํ วุจฺจติ ผสฺสปจฺจยา เวทนา ผสฺสสมฺปยุตฺตา.}

\prose{ตตฺถ กตมา เวทนาปจฺจยา ตณฺหา เวทนาสมฺปยุตฺตา, โย ราโค สาราโค ฯเปฯ จิตฺตสฺส สาราโค, อยํ วุจฺจติ เวทนาปจฺจยา ตณฺหา เวทนาสมฺปยุตฺตา.}

\prose{ตตฺถ กตมํ ตณฺหาปจฺจยา อุปาทานํ ตณฺหาสมฺปยุตฺตํ, ยา ทิฏฺฐิ ทิฏฺฐิคตํ ฯเปฯ ติตฺถายตนํ วิปริยาสคฺคาโห, อิทํ วุจฺจติ ตณฺหาปจฺจยา อุปาทานํ ตณฺหาสมฺปยุตฺตํ ฯเปฯ เตน วุจฺจติ “เอวเมตสฺส เกวลสฺส ทุกฺขกฺขนฺธสฺส สมุทโย โหตี”ติ. (1-9)}

\proseitem{266}{ตสฺมิํ สมเย อวิชฺชาปจฺจยา สงฺขาโร อวิชฺชาสมฺปยุตฺโต, สงฺขารปจฺจยา วิญฺญาณํ สงฺขารสมฺปยุตฺตํ, วิญฺญาณปจฺจยา นามํ วิญฺญาณสมฺปยุตฺตํ, นามปจฺจยา ผสฺโส นามสมฺปยุตฺโต, ผสฺสปจฺจยา เวทนา}

\vspace{\tipitakaparskip}
\csromancenter{ปฏิจฺจสมุปฺปาทวิภงฺค ผสฺสสมฺปยุตฺตา, เวทนาปจฺจยา ตณฺหา เวทนาสมฺปยุตฺตา, ตณฺหาปจฺจยา อุปาทานํ ตณฺหาสมฺปยุตฺตํ.\csromanspacer{} อุปาทานปจฺจยา ภโว, ภวปจฺจยา ชาติ, ชาติปจฺจยา ชรามรณํ, เอวเมตสฺส เกลวสฺส ทุกฺขกฺขนฺธสฺส สมุทโย โหติ.}
\proseitem{267}{\csromanfolio{161}ตตฺถ กตมา อวิชฺชา, ยํ อญฺญาณํ อทสฺสนํ ฯเปฯ อวิชฺชาลงฺคี โมโห อกุสลมูลํ, อยํ วุจฺจติ อวิชฺชา.}

\prose{ตตฺถ กตโม อวิชฺชาปจฺจยา สงฺขาโร อวิชฺชาสมฺปยุตฺโต, ยา เจตนา สญฺเจตนา สญฺเจตยิตตฺตํ, อยํ วุจฺจติ อวิชฺชาปจฺจยา สงฺขาโร อวิชฺชาสมฺปยุตฺโต.}

\prose{ตตฺถ กตมํ สงฺขารปจฺจยา วิญฺญาณํ สงฺขารสมฺปยุตฺตํ, ยํ จิตฺตํ มโน มานสํ ฯเปฯ ตชฺชามโนวิญฺญาณธาตุ, อิทํ วุจฺจติ สงฺขารปจฺจยา วิญฺญาณํ สงฺขารสมฺปยุตฺตํ.}

\prose{ตตฺถ กตมํ วิญฺญาณปจฺจยา นามํ วิญฺญาณสมฺปยุตฺตํ, เวทนากฺขนฺโธ สญฺญากฺขนฺโธ สงฺขารกฺขนฺโธ, อิทํ วุจฺจติ วิญฺญาณปจฺจยา นามํ วิญฺญาณสมฺปยุตฺตํ.}

\prose{นามปจฺจยา ผสฺโส นามสมฺปยุตฺโตติ ตตฺถ กตมํ นามํ, ฐเปตฺวา ผสฺสํ เวทนากฺขนฺโธ สญฺญากฺขนฺโธ สงฺขารกฺขนฺโธ วิญฺญาณกฺขนฺโธ, อิทํ วุจฺจติ นามํ.}

\prose{ตตฺถ กตโม นามปจฺจยา ผสฺโส นามสมฺปยุตฺโต, โย ผสฺโส ผุสนา สมฺผุสนา สมฺผุสิตตฺตํ, อยํ วุจฺจติ นามปจฺจยา ผสฺโส นามสมฺปยุตฺโต ฯเปฯ เตน วุจฺจติ “เอวเมตสฺส เกวลสฺส ทุกฺขกฺขนฺธสฺส สมุทโย โหตี”ติ. (2-10)}

\proseitem{268}{ตสฺมิํ สมเย อวิชฺชาปจฺจยา สงฺขาโร อวิชฺชาสมฺปยุตฺโต, สงฺขารปจฺจยา วิญฺญาณํ สงฺขารสมฺปยุตฺตํ, วิญฺญาณปจฺจยา นามรูปํ วิญฺญาณสมฺปยุตฺตํ นามํ, นามรูปปจฺจยา ฉฏฺฐายตนํ นามสมฺปยุตฺตํ, ฉฏฺฐายตนปจฺจยา ผสฺโส ฉฏฺฐายตนสมฺปยุตฺโต, ผสฺสปจฺจยา เวทนา ผสฺสสมฺปยุตฺตา, เวทนาปจฺจยา ตณฺหา เวทนาสมฺปยุตฺตา, ตณฺหาปจฺจยา อุปาทานํ ตณฺหาสมฺปยุตฺตํ, อุปาทานปจฺจยา ภโว, ภวปจฺจยา ชาติ, \csromanfolio{162}ชาติปจฺจยา ชรามรณํ, เอวเมตสฺส เกวลสฺส ทุกฺขกฺขนฺธสฺส สมุทโย โหติ.}

\proseitem{269}{ตตฺถ กตมา อวิชฺชา, ยํ อญฺญาณํ อทสฺสนํ ฯเปฯ อวิชฺชาลงฺคี โมโห อกุสลมูลํ, อยํ วุจฺจติ อวิชฺชา.}

\prose{ตตฺถ กตโม อวิชฺชาปจฺจยา สงฺขาโร อวิชฺชาสมฺปยุตฺโต, ยา เจตนา สญฺเจตนา สญฺเจตยิตตฺตํ, อยํ วุจฺจติ อวิชฺชาปจฺจยา สงฺขาโร อวิชฺชาสมฺปยุตฺโต.}

\prose{ตตฺถ กตมํ สงฺขารปจฺจยา วิญฺญาณํ สงฺขารสมฺปยุตฺตํ, ยํ จิตฺตํ มโน มานสํ ฯเปฯ ตชฺชามโนวิญฺญาณธาตุ, อิทํ วุจฺจติ สงฺขารปจฺจยา วิญฺญาณํ สงฺขารสมฺปยุตฺตํ.}

\prose{ตตฺถ กตมํ วิญฺญาณปจฺจยา นามรูปํ วิญฺญาณสมฺปยุตฺตํ นามํ, อตฺถิ นามํ, อตฺถิ รูปํ.\csromanspacer{} ตตฺถ กตมํ นามํ, เวทนากฺขนฺโธ สญฺญากฺขนฺโธ สงฺขารกฺขนฺโธ, อิทํ วุจฺจติ นามํ.\csromanspacer{} ตตฺถ กตมํ รูปํ, จกฺขายตนสฺส อุปจโย, โสตายตนสฺส อุปจโย, ฆานายตนสฺส อุปจโย, ชิวฺหายตนสฺส อุปจโย, กายายตนสฺส อุปจโย, ยํ วา ปนญฺญมฺปิ อตฺถิ รูปํ จิตฺตชํ จิตฺตเหตุกํ จิตฺตสมุฏฺฐานํ, อิทํ วุจฺจติ รูปํ.\csromanspacer{} อิติ อิทญฺจ นามํ อิทญฺจ รูปํ, อิทํ วุจฺจติ วิญฺญาณปจฺจยา นามรูปํ วิญฺญาณสมฺปยุตฺตํ นามํ.}

\prose{นามรูปปจฺจยา ฉฏฺฐายตนํ นามสมฺปยุตฺตนฺติ อตฺถิ นามํ, อตฺถิ รูปํ.\csromanspacer{} ตตฺถ กตมํ นามํ, เวทนากฺขนฺโธ สญฺญากฺขนฺโธ สงฺขารกฺขนฺโธ, อิทํ วุจฺจติ นามํ.\csromanspacer{} ตตฺถ กตมํ รูปํ, ยํ รูปํ นิสฺสาย มโนวิญฺญาณธาตุ วตฺตติ, อิทํ วุจฺจติ รูปํ.\csromanspacer{} อิติ อิทญฺจ นามํ อิทญฺจ รูปํ, อิทํ วุจฺจติ นามรูปํ.}

\prose{ตตฺถ กตมํ นามรูปปจฺจยา ฉฏฺฐายตนํ นามสมฺปยุตฺตํ, ยํ จิตฺตํ มโน มานสํ ฯเปฯ ตชฺชามโนวิญฺญาณธาตุ, อิทํ วุจฺจติ นามรูปปจฺจยา ฉฏฺฐายตนํ นามสมฺปยุตฺตํ.}

\prose{ตตฺถ กตโม ฉฏฺฐายตนปจฺจยา ผสฺโส ฉฏฺฐายตนสมฺปยุตฺโต, โย ผสฺโส ผุสนา สมฺผสนา สมฺผุสิตตฺตํ, อยํ วุจฺจติ ฉฏฺฐายตนปจฺจยา ผสฺโส ฉฏฺฐายตนสมฺปยุตฺโต ฯเปฯ เตน วุจฺจติ “เอวเมตสฺส เกวลสฺส ทุกฺขกฺขนฺธสฺส สมุทโย โหตี”ติ. (3-11)}

\csromanheader{2}{6. ปฏิจฺจสมุปฺปาทวิภงฺค}
\proseitem{270}{\csromanfolio{163}ตสฺมิํ สมเย อวิชฺชาปจฺจยา สงฺขาโร อวิชฺชาสมฺปยุตฺโต, สงฺขารปจฺจยา วิญฺญาณํ สงฺขารสมฺปยุตฺตํ, วิญฺญาณปจฺจยา นามรูปํ วิญฺญาณสมฺปยุตฺตํ นามํ, นามรูปปจฺจยา สฬายตนํ นามสมฺปยุตฺตํ ฉฏฺฐายตนํ, ฉฏฺฐายตนปจฺจยา ผสฺโส ฉฏฺฐายตนสมฺปยุตฺโต, ผสฺสปจฺจยา เวทนา ผสฺสสมฺปยุตฺตา, เวทนาปจฺจยา ตณฺหา เวทนาสมฺปยุตฺตา, ตณฺหาปจฺจยา อุปาทานํ ตณฺหาสมฺปยุตฺตํ, อุปาทานปจฺจยา ภโว, ภวปจฺจยา ชาติ, ชาติปจฺจยา ชรามรณํ, เอวเมตสฺส เกวลสฺส ทุกฺขกฺขนฺธสฺส สมุทโย โหติ.}

\proseitem{271}{ตตฺถ กตมา อวิชฺชา, ยํ อญฺญาณํ อทสฺสนํ ฯเปฯ อวิชฺชาลงฺคี โมโห อกุสลมูลํ, อยํ วุจฺจติ อวิชฺชา.}

\prose{ตตฺถ กตโม อวิชฺชาปจฺจยา สงฺขาโร อวิชฺชาสมฺปยุตฺโต, ยา เจตนา สญฺเจตนา สญฺเจตยิตตฺตํ, อยํ วุจฺจติ อวิชฺชาปจฺจยา สงฺขาโร อวิชฺชาสมฺปยุตฺโต.}

\prose{ตตฺถ กตมํ สงฺขารปจฺจยา วิญฺญาณํ สงฺขารสมฺปยุตฺตํ, ยํ จิตฺตํ มโน มานสํ ฯเปฯ ตชฺชามโนวิญฺญาณธาตุ, อิทํ วุจฺจติ สงฺขารปจฺจยา วิญฺญาณํ สงฺขารสมฺปยุตฺตํ.}

\prose{ตตฺถ กตมํ วิญฺญาณปจฺจยา นามรูปํ วิญฺญาณสมฺปยุตฺตํ นามํ, อตฺถิ นามํ, อตฺถิ รูปํ.\csromanspacer{} ตตฺถ กตมํ นามํ, เวทนากฺขนฺโธ สญฺญากฺขนฺโธ สงฺขารกฺขนฺโธ, อิทํ วุจฺจติ นามํ.\csromanspacer{} ตตฺถ กตมํ รูปํ, จกฺขายตนสฺส อุปจโย, โสตายตนสฺส อุปจโย, ฆานายตนสฺส อุปจโย, ชิวฺหายตนสฺส อุปจโย, กายายตนสฺส อุปจโย, ยํ วา ปนญฺญมฺปิ อตฺถิ รูปํ จิตฺตชํ จิตฺตเหตุกํ จิตฺตสมุฏฺฐานํ, อิทํ วุจฺจติ รูปํ.\csromanspacer{} อิติ อิทญฺจ นามํ อิทญฺจ รูปํ, อิทํ วุจฺจติ วิญฺญาณปจฺจยา นามรูปํ วิญฺญาณสมฺปยุตฺตํ นามํ.}

\prose{นามรูปปจฺจยา สฬายตนํ นามสมฺปยุตฺตํ ฉฏฺฐายตนนฺติ อตฺถิ นามํ, อตฺถิ รูปํ.\csromanspacer{} ตตฺถ กตมํ นามํ, เวทนากฺขนฺโธ สญฺญากฺขนฺโธ สงฺขารกฺขนฺโธ, อิทํ วุจฺจติ นามํ.\csromanspacer{} ตตฺถ กตมํ รูปํ, จตฺตาโร จ มหาภูตา ยญฺจ รูปํ นิสฺสาย มโนวิญฺญาณธาตุ วตฺตติ, อิทํ วุจฺจติ รูปํ.\csromanspacer{} อิติ อิทญฺจ นามํ อิทญฺจ รูปํ, อิทํ วุจฺจติ นามรูปํ.}

\prose{ตตฺถ กตมํ นามรูปปจฺจยา สฬายตนํ นามสมฺปยุตฺตํ ฉฏฺฐายตนํ, จกฺขายตนํ โสตายตนํ ฆานายตนํ ชิวฺหายตนํ กายายตนํ \csromanfolio{164}มนายตนํ, อิทํ วุจฺจติ นามรูปปจฺจยา สฬายตนํ นามสมฺปยุตฺตํ ฉฏฺฐายตนํ.}

\prose{ตตฺถ กตโม ฉฏฺฐายตนปจฺจยา ผสฺโส ฉฏฺฐายตนสมฺปยุตฺโต, โย ผสฺโส ผุสนา สมฺผุสนา สมฺผุสิตตฺตํ, อยํ วุจฺจติ ฉฏฺฐายตนปจฺจยา ผสฺโส ฉฏฺฐายตนสมฺปยุตฺโต ฯเปฯ เตน วุจฺจติ “เอวเมตสฺส เกวลสฺส ทุกฺขกฺขนฺธสฺส สมุทโย โหตี”ติ. (4-12)}

\csromanheader{2}{สมฺปยุตฺตจตุกฺกํ.}
\csromansectionrule
\csromanmatikaanchor{mtk.54}
\csromantocmarkref{subparagraph}{8. อญฺญมญฺญจตุกฺก (นิทฺเทส)}{mtk.54}
\bookmark[dest={mtk.54},level=5]{8. อญฺญมญฺญจตุกฺก (นิทฺเทส)}
\csromanheader{6}{8. อญฺญมญฺญจตุกฺก}
\proseitem{272}{ตสฺมิํ สมเย อวิชฺชาปจฺจยา สงฺขาโร, สงฺขารปจฺจยาปิ อวิชฺชา, สงฺขารปจฺจยา วิญฺญาณํ, วิญฺญาณปจฺจยาปิ สงฺขาโร, วิญฺญาณปจฺจยา นามํ, นามปจฺจยาปิ วิญฺญาณํ, นามปจฺจยา ฉฏฺฐายตนํ, ฉฏฺฐายตนปจฺจยาปิ นามํ, ฉฏฺฐายตนปจฺจยา ผสฺโส, ผสฺสปจฺจยาปิ ฉฏฺฐายตนํ, ผสฺสปจฺจยา เวทนา, เวทนาปจฺจยาปิ ผสฺโส, เวทนาปจฺจยา ตณฺหา, ตณฺหาปจฺจยาปิ เวทนา, ตณฺหาปจฺจยา อุปาทานํ, อุปาทานปจฺจยาปิ ตณฺหา, อุปาทานปจฺจยา ภโว, ภวปจฺจยา ชาติ, ชาติปจฺจยา ชรามรณํ, เอวเมตสฺส เกวลสฺส ทุกฺขกฺขนฺธสฺส สมุทโย โหติ.}

\proseitem{273}{ตตฺถ กตมา อวิชฺชา, ยํ อญฺญาณํ อทสฺสนํ ฯเปฯ อวิชฺชาลงฺคี โมโห อกุสลมูลํ, อยํ วุจฺจติ อวิชฺชา.}

\prose{ตตฺถ กตโม อวิชฺชาปจฺจยา สงฺขาโร, ยา เจตนา สญฺเจตนา สญฺเจตยิตตฺตํ, อยํ วุจฺจติ อวิชฺชาปจฺจยา สงฺขาโร.}

\prose{ตตฺถ กตมา สงฺขารปจฺจยาปิ อวิชฺชา, ยํ อญฺญาณํ อทสฺสนํ ฯเปฯ อวิชฺชาลงฺคี โมโห อกุสลมูลํ, อยํ วุจฺจติ สงฺขารปจฺจยาปิ อวิชฺชา.}

\prose{ตตฺถ กตมํ สงฺขารปจฺจยา วิญฺญาณํ, ยํ จิตฺตํ มโน มานสํ ฯเปฯ ตชฺชามโนวิญฺญาณธาตุ, อิทํ วุจฺจติ สงฺขารปจฺจยา วิญฺญาณํ.}

\prose{ตตฺถ กตโม วิญฺญาณปจฺจยาปิ สงฺขาโร, ยา เจตนา สญฺเจตนา สญฺเจตยิตตฺตํ, อยํ วุจฺจติ วิญฺญาณปจฺจยาปิ สงฺขาโร.}

\csromanheader{2}{6. ปฏิจฺจสมุปฺปาทวิภงฺค}
\prose{\csromanfolio{165}ตตฺถ กตมํ วิญฺญาณปจฺจยา นามํ, เวทนากฺขนฺโธ สญฺญากฺขนฺโธ สงฺขารกฺขนฺโธ, อิทํ วุจฺจติ วิญฺญาณปจฺจยา นามํ.}

\prose{ตตฺถ กตมํ นามปจฺจยาปิ วิญฺญาณํ, ยํ จิตฺตํ มโน มานสํ ฯเปฯ ตชฺชามโนวิญฺญาณธาตุ, อิทํ วุจฺจติ นามปจฺจยาปิ วิญฺญาณํ.}

\prose{ตตฺถ กตมํ นามปจฺจยา ฉฏฺฐายตนํ, ยํ จิตฺตํ มโน มานสํ ฯเปฯ ตชฺชามโนวิญฺญาณธาตุ, อิทํ วุจฺจติ นามปจฺจยา ฉฏฺฐายตนํ.}

\prose{ตตฺถ กตมํ ฉฏฺฐายตนปจฺจยาปิ นามํ, เวทนากฺขนฺโธ สญฺญากฺขนฺโธ สงฺขารกฺขนฺโธ, อิทํ วุจฺจติ ฉฏฺฐายตนปจฺจยาปิ นามํ.}

\prose{ตตฺถ กตโม ฉฏฺฐายตนปจฺจยา ผสฺโส, โย ผสฺโส ผุสนา สมฺผุสนา สมฺผุสิตตฺตํ, อยํ วุจฺจติ ฉฏฺฐายตนปจฺจยา ผสฺโส.}

\prose{ตตฺถ กตมํ ผสฺสปจฺจยาปิ ฉฏฺฐายตนํ, ยํ จิตฺตํ มโน มานสํ ฯเปฯ ตชฺชามโนวิญฺญาณธาตุ, อิทํ วุจฺจติ ผสฺสปจฺจยาปิ ฉฏฺฐายตนํ.}

\prose{ตตฺถ กตมา ผสฺสปจฺจยา เวทนา, ยํ เจตสิกํ สาตํ เจตสิกํ สุขํ เจโตสมฺผสฺสชํ สาตํ สุขํ เวทยิตํ เจโตสมฺผสฺสชา สาตา สุขา เวทนา, อยํ วุจฺจติ ผสฺสปจฺจยา เวทนา.}

\prose{ตตฺถ กตโม เวทนาปจฺจยาปิ ผสฺโส, โย ผสฺโส ผุสนา สมฺผุสนา สมฺผุสิตตฺตํ, อยํ วุจฺจติ เวทนาปจฺจยาปิ ผสฺโส.}

\prose{ตตฺถ กตโม เวทนาปจฺจยา ตณฺหา, โย ราโค สาราโค ฯเปฯ จิตฺตสฺส สาราโค, อยํ วุจฺจติ เวทนาปจฺจยา ตณฺหา.}

\prose{ตตฺถ กตมา ตณฺหาปจฺจยาปิ เวทนา, ยํ เจตสิกํ สาตํ เจตสิกํ สุขํ เจโตสมฺผสฺสชํ สาตํ สุขํ เวทยิตํ เจโต สมฺผสฺสชา สาตา สุขา เวทนา, อยํ วุจฺจติ ตณฺหาปจฺจยาปิ เวทนา.}

\prose{ตตฺถ กตมํ ตณฺหาปจฺจยา อุปาทานํ, ยา ทิฏฺฐิ ทิฏฺฐิคตํ ฯเปฯ ติตฺถายตนํ วิปริยาสคฺคาโห, อิทํ วุจฺจติ ตณฺหาปจฺจยา อุปาทานํ.}

\prose{ตตฺถ กตมา อุปาทานปจฺจยาปิ ตณฺหา, โย ราโค ฯเปฯ จิตฺตสฺส สาราโค, อยํ วุจฺจติ อุปาทานปจฺจยาปิ ตณฺหา.}

\prose{\csromanfolio{166}ตตฺถ กตโม อุปาทานปจฺจยา ภโว, ฐเปตฺวา อุปาทานํ เวทนากฺขนฺโธ สญฺญากฺขนฺโธ สงฺขารกฺขนฺโธ วิญฺญาณกฺขนฺโธ, อยํ วุจฺจติ อุปาทานปจฺจยา ภโว.}

\prose{ตตฺถ กตมา ภวปจฺจยา ชาติ, ยา เตสํ เตสํ ธมฺมานํ ชาติ สญฺชาติ นิพฺพตฺติ อภินิพฺพตฺติ ปาตุภาโว, อยํ วุจฺจติ ภวปจฺจยา ชาติ.}

\prose{ตตฺถ กตมํ ชาติปจฺจยา ชรามรณํ, อตฺถิ ชรา, อตฺถิ มรณํ.\csromanspacer{} ตตฺถ กตมา ชรา, ยา เตสํ เตสํ ธมฺมานํ ชรา ชีรณตา อายุโน สํหานิ, อยํ วุจฺจติ ชรา.\csromanspacer{} ตตฺถ กตมํ มรณํ, โย เตสํ เตสํ ธมฺมานํ ขโย วโย เภโท ปริเภโท อนิจฺจตา อนฺตรธานํ, อิทํ วุจฺจติ มรณํ.\csromanspacer{} อิติ อยญฺจ ชรา อิทญฺจ มรณํ, อิทํ วุจฺจติ ชาติปจฺจยา ชรามรณํ.}

\prose{เอวเมตสฺส เกวลสฺส ทุกฺขกฺขนฺธสฺส สมุทโย โหตีติ เอวเมตสฺส เกวลสฺส ทุกฺขกฺขนฺธสฺส สงฺคติ โหติ, สมาคโม โหติ, สโมธานํ โหติ, ปาตุภาโว โหติ, เตน วุจฺจติ “เอวเมตสฺส เกวลสฺส ทุกฺขกฺขนฺธสฺส สมุทโย โหตี”ติ. (1-13)}

\proseitem{274}{ตสฺมิํ สมเย อวิชฺชาปจฺจยา สงฺขาโร, สงฺขารปจฺจยาปิ อวิชฺชา, สงฺขารปจฺจยา วิญฺญาณํ, วิญฺญาณปจฺจยาปิ สงฺขาโร, วิญฺญาณปจฺจยา นามํ, นามปจฺจยาปิ วิญฺญาณํ, นามปจฺจยา ผสฺโส, ผสฺสปจฺจยาปิ นามํ, ผสฺสปจฺจยา เวทนา, เวทนาปจฺจยาปิ ผสฺโส, เวทนาปจฺจยา ตณฺหา, ตณฺหาปจฺจยาปิ เวทนา, ตณฺหาปจฺจยา อุปาทานํ, อุปาทานปจฺจยาปิ ตณฺหา, อุปาทานปจฺจยา ภโว, ภวปจฺจยา ชาติ, ชาติปจฺจยา ชรามรณํ, เอวเมตสฺส เกวลสฺส ทุกฺขกฺขนฺธสฺส สมุทโย โหติ.}

\proseitem{275}{ตตฺถ กตมา อวิชฺชา, ยํ อญฺญาณํ อทสฺสนํ ฯเปฯ อวิชฺชาลงฺคี โมโห อกุสลมูลํ, อยํ วุจฺจติ อวิชฺชา.}

\prose{ตตฺถ กตโม อวิชฺชาปจฺจยา สงฺขาโร, ยา เจตนา สญฺเจตนา สญฺเจตยิตตฺตํ, อยํ วุจฺจติ อวิชฺชาปจฺจยา สงฺขาโร.}

\prose{ตตฺถ กตมา สงฺขารปจฺจยาปิ อวิชฺชา, ยํ อญฺญาณํ อทสฺสนํ ฯเปฯ อวิชฺชาลงฺคี โมโห อกุสลมูลํ, อยํ วุจฺจติ สงฺขารปจฺจยาปิ อวิชฺชา.}

\csromanheader{2}{6. ปฏิจฺจสมุปฺปาทวิภงฺค}
\prose{\csromanfolio{167}ตตฺถ กตมํ สงฺขารปจฺจยา วิญฺญาณํ, ยํ จิตฺตํ มโน มานสํ ฯเปฯ ตชฺชามโนวิญฺญาณธาตุ, อิทํ วุจฺจติ สงฺขารปจฺจยา วิญฺญาณํ.}

\prose{ตตฺถ กตโม วิญฺญาณปจฺจยาปิ สงฺขาโร, ยา เจตนา สญฺเจตนา สญฺเจตยิตตฺตํ, อยํ วุจฺจติ วิญฺญาณปจฺจยาปิ สงฺขาโร.}

\prose{ตตฺถ กตมํ วิญฺญาณปจฺจยา นามํ, เวทนากฺขนฺโธ สญฺญากฺขนฺโธ สงฺขารกฺขนฺโธ, อิทํ วุจฺจติ วิญฺญาณปจฺจยา นามํ.}

\prose{ตตฺถ กตมํ นามปจฺจยาปิ วิญฺญาณํ, ยํ จิตฺตํ มโน มานสํ ฯเปฯ ตชฺชามโนวิญฺญาณธาตุ, อิทํ วุจฺจติ นามปจฺจยาปิ วิญฺญาณํ.}

\prose{นามปจฺจยา ผสฺโสติ ตตฺถ กตมํ นามํ, ฐเปตฺวา ผสฺสํ เวทนากฺขนฺโธ สญฺญากฺขนฺโธ สงฺขารกฺขนฺโธ วิญฺญาณกฺขนฺโธ, อิทํ วุจฺจติ นามํ.}

\prose{ตตฺถ กตโม นามปจฺจยา ผสฺโส, โย ผสฺโส ผุสนา สมฺผุสนา สมฺผุสิตตฺตํ, อยํ วุจฺจติ นามปจฺจยา ผสฺโส.}

\prose{ตตฺถ กตมํ ผสฺสปจฺจยาปิ นามํ, เวทนากฺขนฺโธ สญฺญากฺขนฺโธ สงฺขารกฺขนฺโธ วิญฺญาณกฺขนฺโธ, อิทํ วุจฺจติ ผสฺสปจฺจยาปิ นามํ ฯเปฯ เตน วุจฺจติ “เอวเมตสฺส เกวลสฺส ทุกฺขกฺขนฺธสฺส สมุทโย โหตี”ติ. (2-14)}

\proseitem{276}{ตสฺมิํ สมเย อวิชฺชาปจฺจยา สงฺขาโร, สงฺขารปจฺจยาปิ อวิชฺชา, สงฺขารปจฺจยา วิญฺญาณํ, วิญฺญาณปจฺจยาปิ สงฺขาโร, วิญฺญาณปจฺจยา นามรูปํ, นามรูปปจฺจยาปิ วิญฺญาณํ, นามรูปปจฺจยา ฉฏฺฐายตนํ, ฉฏฺฐายตนปจฺจยาปิ นามรูปํ, ฉฏฺฐายตนปจฺจยา ผสฺโส, ผสฺสปจฺจยาปิ ฉฏฺฐายตนํ, ผสฺสปจฺจยา เวทนา, เวทนาปจฺจยาปิ ผสฺโส, เวทนาปจฺจยา ตณฺหา, ตณฺหาปจฺจยาปิ เวทนา, ตณฺหาปจฺจยา อุปาทานํ, อุปาทานปจฺจยาปิ ตณฺหา, อุปาทานปจฺจยา ภโว, ภวปจฺจยา ชาติ, ชาติปจฺจยา ชรามรณํ, เอวเมตสฺส เกวลสฺส ทุกฺขกฺขนฺธสฺส สมุทโย โหติ.}

\proseitem{277}{ตตฺถ กตมา อวิชฺชา, ยํ อญฺญาณํ อทสฺสนํ ฯเปฯ อวิชฺชาลงฺคี โมโห อกุสลมูลํ, อยํ วุจฺจติ อวิชฺชา.}

\prose{ตตฺถ กตโม อวิชฺชาปจฺจยา สงฺขาโร, ยา เจตนา สญฺเจตนา สญฺเจตยิตตฺตํ, อยํ วุจฺจติ อวิชฺชาปจฺจยา สงฺขาโร.}

\prose{\csromanfolio{168}ตตฺถ กตมา สงฺขารปจฺจยาปิ อวิชฺชา, ยํ อญฺญาณํ อทสฺสนํ ฯเปฯ อวิชฺชาลงฺคี โมโห อกุสลมูลํ, อยํ วุจฺจติ สงฺขารปจฺจยาปิ อวิชฺชา.}

\prose{ตตฺถ กตมํ สงฺขารปจฺจยา วิญฺญาณํ, ยํ จิตฺตํ มโน มานสํ ฯเปฯ ตชฺชามโนวิญฺญาณธาตุ, อิทํ วุจฺจติ สงฺขารปจฺจยา วิญฺญาณํ.}

\prose{ตตฺถ กตโม วิญฺญาณปจฺจยาปิ สงฺขาโร, ยา เจตนา สญฺเจตนา สญฺเจตยิตตฺตํ, อยํ วุจฺจติ วิญฺญาณปจฺจยาปิ สงฺขาโร.}

\prose{ตตฺถ กตมํ วิญฺญาณปจฺจยา นามรูปํ, อตฺถิ นามํ, อตฺถิ รูปํ.\csromanspacer{} ตตฺถ กตมํ นามํ, เวทนากฺขนฺโธ สญฺญากฺขนฺโธ สงฺขารกฺขนฺโธ, อิทํ วุจฺจติ นามํ.\csromanspacer{} ตตฺถ กตมํ รูปํ, จกฺขายตนสฺส อุปจโย, โสตายตนสฺส อุปจโย, ฆานายตนสฺส อุปจโย, ชิวฺหายตนสฺส อุปจโย, กายายตนสฺส อุปจโย, ยํ วา ปนญฺญมฺปิ อตฺถิ รูปํ จิตฺตชํ จิตฺตเหตุกํ จิตฺตสมุฏฺฐานํ, อิทํ วุจฺจติ รูปํ.\csromanspacer{} อิติ อิทญฺจ นามํ อิทญฺจ รูปํ, อิทํ วุจฺจติ วิญฺญาณปจฺจยา นามรูปํ.}

\prose{นามรูปปจฺจยาปิ วิญฺญาณนฺติ อตฺถิ นามํ, อตฺถิ รูปํ.\csromanspacer{} ตตฺถ กตมํ นามํ, เวทนากฺขนฺโธ สญฺญากฺขนฺโธ สงฺขารกฺขนฺโธ, อิทํ วุจฺจติ นามํ.\csromanspacer{} ตตฺถ กตมํ รูปํ, ยํ รูปํ นิสฺสาย มโนวิญฺญาณธาตุ วตฺตติ, อิทํ วุจฺจติ รูปํ.\csromanspacer{} อิติ อิทญฺจ นามํ อิทญฺจ รูปํ, อิทํ วุจฺจติ นามรูปํ.}

\prose{ตตฺถ กตมํ นามรูปปจฺจยาปิ วิญฺญาณํ, ยํ จิตฺตํ มโน มานสํ ฯเปฯ ตชฺชามโนวิญฺญาณธาตุ, อิทํ วุจฺจติ นามรูปปจฺจยาปิ วิญฺญาณํ.}

\prose{นามรูปปจฺจยา ฉฏฺฐายตนนฺติ อตฺถิ นามํ, อตฺถิ รูปํ ตตฺถ กตมํ นามํ, เวทนากฺขนฺโธ สญฺญากฺขนฺโธ สงฺขารกฺขนฺโธ, อิทํ วุจฺจติ นามํ.\csromanspacer{} ตตฺถ กตมํ รูปํ, ยํ รูปํ นิสฺสาย มโนวิญฺญาณธาตุ วตฺตติ, อิทํ วุจฺจติ รูปํ.\csromanspacer{} อิติ อิทญฺจ นามํ อิทญฺจ รูปํ, อิทํ วุจฺจติ นามรูปํ.}

\prose{ตตฺถ กตมํ นามรูปปจฺจยา ฉฏฺฐายตนํ, ยํ จิตฺตํ มโน มานสํ ฯเปฯ ตชฺชามโนวิญฺญาณธาตุ, อิทํ วุจฺจติ นามรูปปจฺจยา ฉฏฺฐายตนํ.}

\prose{ตตฺถ กตมํ ฉฏฺฐายตนปจฺจยาปิ นามรูปํ, อตฺถิ นามํ, อตฺถิ รูปํ.\csromanspacer{} ตตฺถ กตมํ นามํ, เวทนากฺขนฺโธ สญฺญากฺขนฺโธ สงฺขารกฺขนฺโธ, อิทํ}

\vspace{\tipitakaparskip}
\csromancenter{ปฏิจฺจสมุปฺปาทวิภงฺค วุจฺจติ นามํ.\csromanspacer{} ตตฺถ กตมํ รูปํ, จกฺขายตนสฺส อุปจโย, โสตายตนสฺส อุปจโย, ฆานายตนสฺส อุปจโย, ชิวฺหายตนสฺส อุปจโย, กายายตนสฺส อุปจโย, ยํ วา ปนญฺญมฺปิ อตฺถิ รูปํ จิตฺตชํ จิตฺตเหตุกํ จิตฺตสมุฏฺฐานํ, อิทํ วุจฺจติ รูปํ.\csromanspacer{} อิติ อิทญฺจ นามํ อิทญฺจ รูปํ, อิทํ วุจฺจติ ฉฏฺฐายตนปจฺจยาปิ นามรูปํ.}
\prose{\csromanfolio{169}ตตฺถ กตโม ฉฏฺฐายตนปจฺจยา ผสฺโส, โย ผสฺโส ผุสนา สมฺผุสนา สมฺผุสิตตฺตํ, อยํ วุจฺจติ ฉฏฺฐายตนปจฺจยา ผสฺโส.}

\prose{ตตฺถ กตมํ ผสฺสปจฺจยาปิ ฉฏฺฐายตนํ, ยํ จิตฺตํ มโน มานสํ ฯเปฯ ตชฺชามโนวิญฺญาณธาตุ, อิทํ วุจฺจติ ผสฺสปจฺจยาปิ ฉฏฺฐายตนํ ฯเปฯ เตน วุจฺจติ “เอวเมตสฺส เกวลสฺส ทุกฺขกฺขนฺธสฺส สมุทโย โหตี”ติ. (3-15)}

\proseitem{278}{ตสฺมิํ สมเย อวิชฺชาปจฺจยา สงฺขาโร, สงฺขารปจฺจยาปิ อวิชฺชา, สงฺขารปจฺจยา วิญฺญาณํ, วิญฺญาณปจฺจยาปิ สงฺขาโร, วิญฺญาณปจฺจยา นามรูปํ, นามรูปปจฺจยาปิ วิญฺญาณํ, นามรูปปจฺจยา สฬายตนํ, ฉฏฺฐายตนปจฺจยาปิ นามรูปํ, ฉฏฺฐายตนปจฺจยา ผสฺโส, ผสฺสปจฺจยาปิ ฉฏฺฐายตนํ, ผสฺสปจฺจยา เวทนา, เวทนาปจฺจยาปิ ผสฺโส, เวทนาปจฺจยา ตณฺหา, ตณฺหาปจฺจยาปิ เวทนา, ตณฺหาปจฺจยา อุปาทานํ, อุปาทานปจฺจยาปิ ตณฺหา, อุปาทานปจฺจยา ภโว, ภวปจฺจยา ชาติ, ชาติปจฺจยา ชรามรณํ, เอวเมตสฺส เกวลสฺส ทุกฺขกฺขนฺธสฺส สมุทโย โหติ.}

\proseitem{279}{ตตฺถ กตมา อวิชฺชา, ยํ อญฺญาณํ อทสฺสนํ ฯเปฯ อวิชฺชาลงฺคี โมโห อกุสลมูลํ, อยํ วุจฺจติ อวิชฺชา.}

\prose{ตตฺถ กตโม อวิชฺชาปจฺจยา สงฺขาโร, ยา เจตนา สญฺเจตนา สญฺเจตยิตตฺตํ, อยํ วุจฺจติ อวิชฺชาปจฺจยา สงฺขาโร.}

\prose{ตตฺถ กตมา สงฺขารปจฺจยาปิ อวิชฺชา, ยํ อญฺญาณํ อทสฺสนํ ฯเปฯ อวิชฺชาลงฺคี โมโห อกุสลมูลํ, อยํ วุจฺจติ สงฺขารปจฺจยาปิ อวิชฺชา.}

\prose{ตตฺถ กตมํ สงฺขารปจฺจยา วิญฺญาณํ, ยํ จิตฺตํ มโน มานสํ ฯเปฯ ตชฺชามโนวิญฺญาณธาตุ, อิทํ วุจฺจติ สงฺขารปจฺจยา วิญฺญาณํ.}

\prose{ตตฺถ กตโม วิญฺญาณปจฺจยาปิ สงฺขาโร, ยา เจตนา สญฺเจตนา สญฺเจตยิตตฺตํ, อยํ วุจฺจติ วิญฺญาณปจฺจยาปิ สงฺขาโร.}

\prose{\csromanfolio{170}ตตฺถ กตมํ วิญฺญาณปจฺจยา นามรูปํ, อตฺถิ นามํ, อตฺถิ รูปํ.\csromanspacer{} ตตฺถ กตมํ นามํ, เวทนากฺขนฺโธ สญฺญากฺขนฺโธ สงฺขารกฺขนฺโธ, อิทํ วุจฺจติ นามํ.\csromanspacer{} ตตฺถ กตมํ รูปํ, จกฺขายตนสฺส อุปจโย, โสตายตนสฺส อุปจโย, ฆานายตนสฺส อุปจโย, ชิวฺหายตนสฺส อุปจโย, กายายตนสฺส อุปจโย, ยํ วา ปนญฺญมฺปิ อตฺถิ รูปํ จิตฺตชํ จิตฺตเหตุกํ จิตฺตสมุฏฺฐานํ, อิทํ วุจฺจติ รูปํ.\csromanspacer{} อิติ อิทญฺจ นามํ อิทญฺจ รูปํ, อิทํ วุจฺจติ วิญฺญาณปจฺจยา นามรูปํ.}

\prose{นามรูปปจฺจยาปิ วิญฺญาณนฺติ อตฺถิ นามํ, อตฺถิ รูปํ.\csromanspacer{} ตตฺถ กตมํ นามํ, เวทนากฺขนฺโธ สญฺญากฺขนฺโธ สงฺขารกฺขนฺโธ, อิทํ วุจฺจติ นามํ.\csromanspacer{} ตตฺถ กตมํ รูปํ, ยํ รูปํ นิสฺสาย มโนวิญฺญาณธาตุ วตฺตติ, อิทํ วุจฺจติ รูปํ.\csromanspacer{} อิติ อิทญฺจ นามํ อิทญฺจ รูปํ, อิทํ วุจฺจติ นามรูปํ.}

\prose{ตตฺถ กตมํ นามรูปปจฺจยาปิ วิญฺญาณํ, ยํ จิตฺตํ มโน มานสํ ฯเปฯ ตชฺชามโนวิญฺญาณธาตุ, อิทํ วุจฺจติ นามรูปปจฺจยาปิ วิญฺญาณํ.}

\prose{นามรูปปจฺจยา สฬายตนนฺติ อตฺถิ นามํ, อตฺถิ รูปํ.\csromanspacer{} ตตฺถ กตมํ นามํ, เวทนากฺขนฺโธ สญฺญากฺขนฺโธ สงฺขารกฺขนฺโธ, อิทํ วุจฺจติ นามํ.\csromanspacer{} ตตฺถ กตมํ รูปํ, จตฺตาโร จ มหาภูตา ยญฺจ รูปํ นิสฺสาย มโนวิญฺญาณธาตุ วตฺตติ, อิทํ วุจฺจติ รูปํ.\csromanspacer{} อิติ อิทญฺจ นามํ อิทญฺจ รูปํ, อิทํ วุจฺจติ นามรูปํ.}

\prose{ตตฺถ กตมํ นามรูปปจฺจยา สฬายตนํ, จกฺขายตนํ โสตายตนํ ฆานายตนํ ชิวฺหายตนํ กายายตนํ มนายตนํ, อิทํ วุจฺจติ นามรูปปจฺจยา สฬายตนํ.}

\prose{ตตฺถ กตมํ ฉฏฺฐายตนปจฺจยาปิ นามรูปํ, อตฺถิ นามํ, อตฺถิ รูปํ.\csromanspacer{} ตตฺถ กตมํ นามํ, เวทนากฺขนฺโธ สญฺญากฺขนฺโธ สงฺขารกฺขนฺโธ, อิทํ วุจฺจติ นามํ.\csromanspacer{} ตตฺถ กตมํ รูปํ, จกฺขายตนสฺส อุปจโย, โสตายตนสฺส อุปจโย, ฆานายตนสฺส อุปจโย, ชิวฺหายตนสฺส อุปจโย, กายายตนสฺส อุปจโย, ยํ วา ปนญฺญมฺปิ อตฺถิ รูปํ จิตฺตชํ จิตฺตเหตุกํ จิตฺตสมุฏฺฐานํ, อิทํ วุจฺจติ รูปํ.\csromanspacer{} อิติ อิทญฺจ นามํ อิทญฺจ รูปํ, อิทํ วุจฺจติ ฉฏฺฐายตนปจฺจยาปิ นามรูปํ.}

\prose{ตตฺถ กตโม ฉฏฺฐายตนปจฺจยา ผสฺโส, โย ผสฺโส ผุสนา สมฺผุสนา สมฺผุสิตตฺตํ, อยํ วุจฺจติ ฉฏฺฐายตนปจฺจยา ผสฺโส.}

\csromanheader{2}{6. ปฏิจฺจสมุปฺปาทวิภงฺค}
\prose{\csromanfolio{171}ตตฺถ กตมํ ผสฺสปจฺจยาปิ ฉฏฺฐายตนํ, ยํ จิตฺตํ มโน มานสํ ฯเปฯ ตชฺชามโนวิญฺญาณธาตุ, อิทํ วุจฺจติ ผสฺสปจฺจยาปิ ฉฏฺฐายตนํ.}

\prose{ตตฺถ กตมา ผสฺสปจฺจยา เวทนา, ยํ เจตสิกํ สาตํ เจตสิกํ สุขํ เจโตสมฺผสฺสชํ สาตํ สุขํ เวทยิตํ เจโตสมฺผสฺสชา สาตา สุขา เวทนา, อยํ วุจฺจติ ผสฺสปจฺจยา เวทนา ฯเปฯ เตน วุจฺจติ “เอวเมตสฺส เกวลสฺส ทุกฺขกฺขนฺธสฺส สมุทโย โหตี”ติ. (4-16)}

\csromanheader{2}{อญฺญมญฺญจตุกฺกํ.}
\csromansectionrule
\csromanmatikaanchor{mtk.55}
\csromantocmarkref{subparagraph}{9. อกุสลนิทฺเทส}{mtk.55}
\bookmark[dest={mtk.55},level=5]{9. อกุสลนิทฺเทส}
\csromanheader{6}{9. อกุสลนิทฺเทส}
\proseitem{280}{กตเม ธมฺมา อกุสลา, ยสฺมิํ สมเย อกุสลํ จิตฺตํ อุปฺปนฺนํ โหติ โสมนสฺสสหคตํ ทิฏฺฐิคตสมฺปยุตฺตํ สสงฺขาเรน ฯเปฯ โสมนสฺสสหคตํ ทิฏฺฐิคตวิปฺปยุตฺตํ รูปารมฺมณํ วา ฯเปฯ โสมนสฺสสหคตํ ทิฏฺฐิคตวิปฺปยุตฺตํ สสงฺขาเรน รูปารมฺมณํ วา สทฺทารมฺมณํ วา คนฺธารมฺมณํ วา รสารมฺมณํ วา โผฏฺฐพฺพารมฺมณํ วา ธมฺมารมฺมณํ วา ยํ ยํ วา ปนารพฺภ, ตสฺมิํ สมเย อวิชฺชาปจฺจยา สงฺขาโร, สงฺขารปจฺจยา วิญฺญาณํ, วิญฺญาณปจฺจยา นามํ, นามปจฺจยา ฉฏฺฐายตนํ, ฉฏฺฐายตนปจฺจยา ผสฺโส, ผสฺสปจฺจยา เวทนา, เวทนาปจฺจยา ตณฺหา, ตณฺหาปจฺจยา อธิโมกฺโข, อธิโมกฺขปจฺจยา ภโว, ภวปจฺจยา ชาติ, ชาติปจฺจยา ชรามรณํ, เอวเมตสฺส เกวลสฺส ทุกฺขกฺขนฺธสฺส สมุทโย โหติ.}

\proseitem{281}{ตตฺถ กตมา อวิชฺชา, ยํ อญฺญาณํ อทสฺสนํ ฯเปฯ อวิชฺชาลงฺคี โมโห อกุสลมูลํ, อยํ วุจฺจติ อวิชฺชา.}

\prose{ตตฺถ กตโม อวิชฺชาปจฺจยา สงฺขาโร, ยา เจตนา สญฺเจตนา สญฺเจตยิตตฺตํ, อยํ วุจฺจติ อวิชฺชาปจฺจยา สงฺขาโร ฯเปฯ.}

\prose{ตตฺถ กตโม ตณฺหาปจฺจยา อธิโมกฺโข, โย จิตฺตสฺส อธิโมกฺโข อธิมุจฺจนา ตทธิมุตฺตตา, อยํ วุจฺจติ ตณฺหาปจฺจยา อธิโมกฺโข.}

\prose{ตตฺถ กตโม อธิโมกฺขปจฺจยา ภโว, ฐเปตฺวา อธิโมกฺขํ เวทนากฺขนฺโธ สญฺญากฺขนฺโธ สงฺขารกฺขนฺโธ วิญฺญาณกฺขนฺโธ, อยํ วุจฺจติ \csromanfolio{172}อธิโมกฺขปจฺจยา ภโว ฯเปฯ เตน วุจฺจติ “เอวเมตสฺส เกวลสฺส ทุกฺขกฺขนฺธสฺส สมุทโย โหตี”ติ.}

\proseitem{282}{กตเม ธมฺมา อกุสลา, ยสฺมิํ สมเย อกุสลํ จิตฺตํ อุปฺปนฺนํ โหติ อุเปกฺขาสหคตํ ทิฏฺฐิคตสมฺปยุตฺตํ รูปารมฺมณํ วา สทฺทารมฺมณํ วา คนฺธารมฺมณํ วา รสารมฺมณํ วา โผฏฺฐพฺพารมฺมณํ วา ธมฺมารมฺมณํ วา ยํ ยํ วา ปนารพฺภ, ตสฺมิํ สมเย อวิชฺชาปจฺจยา สงฺขาโร, สงฺขารปจฺจยา วิญฺญาณํ, วิญฺญาณปจฺจยา นามํ, นามปจฺจยา ฉฏฺฐายตนํ, ฉฏฺฐายตนปจฺจยา ผสฺโส, ผสฺสปจฺจยา เวทนา, เวทนาปจฺจยา ตณฺหา, ตณฺหาปจฺจยา อุปาทานํ, อุปาทานปจฺจยา ภโว, ภวปจฺจยา ชาติ, ชาติปจฺจยา ชรามรณํ, เอวเมตสฺส เกวลสฺส ทุกฺขกฺขนฺธสฺส สมุทโย โหติ.}

\proseitem{283}{ตตฺถ กตมา อวิชฺชา, ยํ อญฺญาณํ อทสฺสนํ ฯเปฯ อวิชฺชาลงฺคี โมโห อกุสลมูลํ, อยํ วุจฺจติ อวิชฺชา ฯเปฯ.}

\prose{ตตฺถ กตมา ผสฺสปจฺจยา เวทนา, ยํ เจตสิกํ เนว สาตํนาสาตํ เจโตสมฺผสฺสชํ อทุกฺขมสุขํ เวทยิตํ เจโตสมฺผสฺสชา อทุกฺขมสุขา เวทนา, อยํ วุจฺจติ ผสฺสปจฺจยา เวทนา ฯเปฯ เตน วุจฺจติ “เอวเมตสฺส เกวลสฺส ทุกฺขกฺขนฺธสฺส สมุทโย โหตี”ติ.}

\proseitem{284}{กตเม ธมฺมา อกุสลา, ยสฺมิํ สมเย อกุสลํ จิตฺตํ อุปฺปนฺนํ โหติ อุเปกฺขาสหคตํ ทิฏฺฐิคตสมฺปยุตฺตํ สสงฺขาเรน ฯเปฯ อุเปกฺขาสหคตํ ทิฏฺฐิคตวิปฺปยุตฺตํ รูปารมฺมณํ วา ฯเปฯ อุเปกฺขาสหคตํ ทิฏฺฐิคตวิปฺปยุตฺตํ สสงฺขาเรน รูปารมฺมณํ วา สทฺทารมฺมณํ วา คนฺธารมฺมณํ วา รสารมฺมณํ วา โผฏฺฐพฺพารมฺมณํ วา ธมฺมารมฺมณํ วา ยํ ยํ วา ปนารพฺภ, ตสฺมิํ สมเย อวิชฺชาปจฺจยา สงฺขาโร, สงฺขารปจฺจยา วิญฺญาณํ, วิญฺญาณปจฺจยา นามํ, นามปจฺจยา ฉฏฺฐายตนํ, ฉฏฺฐายตนปจฺจยา ผสฺโส, ผสฺสปจฺจยา เวทนา, เวทนาปจฺจยา ตณฺหา, ตณฺหาปจฺจยา อธิโมกฺโข, อธิโมกฺขปจฺจยา ภโว, ภวปจฺจยา ชาติ, ชาติปจฺจยา ชรามรณํ, เอวเมตสฺส เกวลสฺส ทุกฺขกฺขนฺธสฺส สมุทโย โหติ.}

\proseitem{285}{ตตฺถ กตมา อวิชฺชา ฯเปฯ เตน วุจฺจติ “เอวเมตสฺส เกวลสฺส ทุกฺขกฺขนฺธสฺส สมุทโย โหตี”ติ.}

\csromanheader{2}{6. ปฏิจฺจสมุปฺปาทวิภงฺค}
\proseitem{286}{\csromanfolio{173}กตเม ธมฺมา อกุสลา, ยสฺมิํ สมเย อกุสลํ จิตฺตํ อุปฺปนฺนํ โหติ โทมนสฺสสหคตํ ปฏิฆสมฺปยุตฺตํ รูปารมฺมณํ วา ฯเปฯ โทมนสฺสสหคตํ ปฏิฆสมฺปยุตฺตํ สสงฺขาเรน รูปารมฺมณํ วา สทฺทารมฺมณํ วา คนฺธารมฺมณํ วา รสารมฺมณํ วา โผฏฺฐพฺพารมฺมณํ วา ธมฺมารมฺมณํ วา ยํ ยํ วา ปนารพฺภ, ตสฺมิํ สมเย อวิชฺชาปจฺจยา สงฺขาโร, สงฺขารปจฺจยา วิญฺญาณํ, วิญฺญาณปจฺจยา นามํ, นามปจฺจยา ฉฏฺฐายตนํ, ฉฏฺฐายตนปจฺจยา ผสฺโส, ผสฺสปจฺจยา เวทนา, เวทนาปจฺจยา ปฏิฆํ, ปฏิฆปจฺจยา อธิโมกฺโข, อธิโมกฺขปจฺจยา ภโว, ภวปจฺจยา ชาติ, ชาติปจฺจยา ชรามรณํ, เอวเมตสฺส เกวลสฺส ทุกฺขกฺขนฺธสฺส สมุทโย โหติ.}

\proseitem{287}{ตตฺถ กตมา อวิชฺชา, ยํ อญฺญาณํ อทสฺสนํ ฯเปฯ อวิชฺชาลงฺคี โมโห อกุสลมูลํ, อยํ วุจฺจติ อวิชฺชา ฯเปฯ อยํ วุจฺจติ ฉฏฺฐายตนปจฺจยา ผสฺโส.}

\prose{ตตฺถ กตมา ผสฺสปจฺจยา เวทนา, ยํ เจตสิกํ อสาตํ เจตสิกํ ทุกฺขํ เจโตสมฺผสฺสชํ อสาตํ ทุกฺขํ เวทยิตํ เจโตสมฺผสฺสชา อสาตา ทุกฺขา เวทนา, อยํ วุจฺจติ ผสฺสปจฺจยา เวทนา.}

\prose{ตตฺถ กตมํ เวทนาปจฺจยา ปฏิฆํ, โย จิตฺตสฺส อาฆาโต ฯเปฯ จณฺฑิกฺกํ อสุโรโป\footnote{อสุโลโป (สฺยา)} อนตฺตมนตา จิตฺตสฺส, อิทํ วุจฺจติ เวทนาปจฺจยา ปฏิฆํ.}

\prose{ตตฺถ กตโม ปฏิฆปจฺจยา อธิโมกฺโข, โย จิตฺตสฺส อธิโมกฺโข อธิมุจฺจนา ตทธิมุตฺตตา, อยํ วุจฺจติ ปฏิฆปจฺจยา อธิโมกฺโข.}

\prose{ตตฺถ กตโม อธิโมกฺขปจฺจยา ภโว, ฐเปตฺวา อธิโมกฺขํ เวทนากฺขนฺโธ สญฺญากฺขนฺโธ สงฺขารกฺขนฺโธ วิญฺญาณกฺขนฺโธ, อยํ วุจฺจติ อธิโมกฺขปจฺจยา ภโว ฯเปฯ เตน วุจฺจติ “เอวเมตสฺส เกวลสฺส ทุกฺขกฺขนฺธสฺส สมุทโย โหตี”ติ.}

\proseitem{288}{กตเม ธมฺมา อกุสลา, ยสฺมิํ สมเย อกุสลํ จิตฺตํ อุปฺปนฺนํ โหติ อุเปกฺขาสหคตํ วิจิกิจฺฉาสมฺปยุตฺตํ รูปารมฺมณํ วา สทฺทารมฺมณํ วา คนฺธารมฺมณํ วา รสารมฺมณํ วา โผฏฺฐพฺพารมฺมณํ วา ธมฺมารมฺมณํ วา ยํ ยํ วา \csromanfolio{174}ปนารพฺภ, ตสฺมิํ สมเย อวิชฺชาปจฺจยา สงฺขาโร, สงฺขารปจฺจยา วิญฺญาณํ, วิญฺญาณปจฺจยา นามํ, นามปจฺจยา ฉฏฺฐายตนํ, ฉฏฺฐายตนปจฺจยา ผสฺโส, ผสฺสปจฺจยา เวทนา, เวทนาปจฺจยา วิจิกิจฺฉา, วิจิกิจฺฉาปจฺจยา ภโว, ภวปจฺจยา ชาติ, ชาติปจฺจยา ชรามรณํ, เอวเมตสฺส เกวลสฺส ทุกฺขกฺขนฺธสฺส สมุทโย โหติ.}

\proseitem{289}{ตตฺถ กตมา อวิชฺชา, ยํ อญฺญาณํ อทสฺสนํ ฯเปฯ อวิชฺชาลงฺคี โมโห อกุสลมูลํ, อยํ วุจฺจติ อวิชฺชา ฯเปฯ อยํ วุจฺจติ ฉฏฺฐายตนปจฺจยา ผสฺโส.}

\prose{ตตฺถ กตมา ผสฺสปจฺจยา เวทนา, ยํ เจตสิกํ เนว สาตํ นาสาตํ เจโตสมฺผสฺสชํ อทุกฺขมสุขํ เวทยิตํ เจโตสมฺผสฺสชา อทุกฺขมสุขา เวทนา, อยํ วุจฺจติ ผสฺสปจฺจยา เวทนา.}

\prose{ตตฺถ กตมา เวทนาปจฺจยา วิจิกิจฺฉา, ยา กงฺขา กงฺขายนา กงฺขายิตตฺตํ วิมติ วิจิกิจฺฉา เทฺวฬฺหกํ ทฺวิธาปโถ\footnote{เทฺวธาปโถ (สี, สฺยา)} สํสโย อเนกํสคฺคาโห อาสปฺปนา ปริสปฺปนา อปริโยคาหณา\footnote{อปริโยคาหนา (สี, สฺยา, ก)} ฉมฺภิตตฺตํ จิตฺตสฺส มโนวิเลโข, อยํ วุจฺจติ เวทนาปจฺจยา วิจิกิจฺฉา.}

\prose{ตตฺถ กตโม วิจิกิจฺฉาปจฺจยา ภโว, ฐเปตฺวา วิจิกิจฺฉํ เวทนากฺขนฺโธ สญฺญากฺขนฺโธ สงฺขารกฺขนฺโธ วิญฺญาณกฺขนฺโธ, อยํ วุจฺจติ วิจิกิจฺฉาปจฺจยา ภโว ฯเปฯ เตน วุจฺจติ “เอวเมตสฺส เกวลสฺส ทุกฺขกฺขนฺธสฺส สมุทโย โหตี”ติ.}

\proseitem{290}{กตเม ธมฺมา อกุสลา, ยสฺมิํ สมเย อกุสลํ จิตฺตํ อุปฺปนฺนํ โหติ อุเปกฺขาสหคตํ อุทฺธจฺจสมฺปยุตฺตํ รูปารมฺมณํ วา สทฺทารมฺมณํ วา คนฺธารมฺมณํ วา รสารมฺมณํ วา โผฏฺฐพฺพารมฺมณํ วา ธมฺมารมฺมณํ วา ยํ ยํ วา ปนารพฺภ, ตสฺมิํ สมเย อวิชฺชาปจฺจยา สงฺขาโร, สงฺขารปจฺจยา วิญฺญาณํ, วิญฺญาณปจฺจยา นามํ, นามปจฺจยา ฉฏฺฐายตนํ, ฉฏฺฐายตนปจฺจยา ผสฺโส, ผสฺสปจฺจยา เวทนา, เวทนาปจฺจยา อุทฺธจฺจํ, อุทฺธจฺจปจฺจยา อธิโมกฺโข, อธิโมกฺขปจฺจยา ภโว, ภวปจฺจยา ชาติ, ชาติปจฺจยา ชรามรณํ, เอวเมตสฺส เกวลสฺส ทุกฺขกฺขนฺธสฺส สมุทโย โหติ.}

\csromanheader{2}{6. ปฏิจฺจสมุปฺปาทวิภงฺค}
\proseitem{291}{\csromanfolio{175}ตตฺถ กตมา อวิชฺชา, ยํ อญฺญาณํ อทสฺสนํ ฯเปฯ อวิชฺชาลงฺคี โมโห อกุสลมูลํ, อยํ วุจฺจติ อวิชฺชา ฯเปฯ อยํ วุจฺจติ ฉฏฺฐายตนปจฺจยา ผสฺโส.}

\prose{ตตฺถ กตมา ผสฺสปจฺจยา เวทนา, ยํ เจตสิกํ เนว สาตํ นาสาตํ เจโตสมฺผสฺสชํ อทุกฺขมสุขํ เวทยิตํ เจโตสมฺผสฺสชา อทุกฺขมสุขา เวทนา, อยํ วุจฺจติ ผสฺสปจฺจยา เวทนา.}

\prose{ตตฺถ กตมํ เวทนาปจฺจยา อุทฺธจฺจํ, ยํ จิตฺตสฺส อุทฺธจฺจํ อวูปสโม เจตโส วิกฺเขโป ภนฺตตฺตํ จิตฺตสฺส, อิทํ วุจฺจติ เวทนาปจฺจยา อุทฺธจฺจํ.}

\prose{ตตฺถ กตโม อุทฺธจฺจปจฺจยา อธิโมกฺโข, โย จิตฺตสฺส อธิโมกฺโข อธิมุจฺจนา ตทธิมุตฺตตา, อยํ วุจฺจติ อุทฺธจฺจปจฺจยา อธิโมกฺโข.}

\prose{ตตฺถ กตโม อธิโมกฺขปจฺจยา ภโว, ฐเปตฺวา อธิโมกฺขํ เวทนากฺขนฺโธ สญฺญากฺขนฺโธ สงฺขารกฺขนฺโธ วิญฺญาณกฺขนฺโธ, อยํ วุจฺจติ อธิโมกฺขปจฺจยา ภโว ฯเปฯ เตน วุจฺจติ “เอวเมตสฺส เกวลสฺส ทุกฺขกฺขนฺธสฺส สมุทโย โหตี”ติ.}

\csromanheader{2}{อกุสลนิทฺเทโส.}
\csromansectionrule
\csromanmatikaanchor{mtk.56}
\csromantocmarkref{subparagraph}{10. กุสลนิทฺเทส}{mtk.56}
\bookmark[dest={mtk.56},level=5]{10. กุสลนิทฺเทส}
\csromanheader{6}{10. กุสลนิทฺเทส}
\proseitem{292}{กตเม ธมฺมา กุสลา, ยสฺมิํ สมเย กามาวจรํ กุสลํ จิตฺตํ อุปฺปนฺนํ โหติ โสมนสฺสสหคตํ ญาณสมฺปยุตฺตํ รูปารมฺมณํ วา สทฺทารมฺมณํ วา คนฺธารมฺมณํ วา รสารมฺมณํ วา โผฏฺฐพฺพารมฺมณํ วา ธมฺมารมฺมณํ วา ยํ ยํ วา ปนารพฺภ, ตสฺมิํ สมเย กุสลมูลปจฺจยา สงฺขาโร, สงฺขารปจฺจยา วิญฺญาณํ, วิญฺญาณปจฺจยา นามํ, นามปจฺจยา ฉฏฺฐายตนํ, ฉฏฺฐายตนปจฺจยา ผสฺโส, ผสฺสปจฺจยา เวทนา, เวทนาปจฺจยา ปสาโท, ปสาทปจฺจยา อธิโมกฺโข, อธิโมกฺขปจฺจยา ภโว, ภวปจฺจยา ชาติ, ชาติปจฺจยา ชรามรณํ, เอวเมตสฺส เกวลสฺส ทุกฺขกฺขนฺธสฺส สมุทโย โหติ.}

\proseitem{293}{ตตฺถ กตเม กุสลมูลา, อโลโภ อโทโส อโมโห.}

\prose{\csromanfolio{176}ตตฺถ กตโม อโลโภ, โย อโลโภ อลุพฺภนา อลุพฺภิตตฺตํ อสาราโค อสารชฺชนา อสารชฺชิตตฺตํ อนภิชฺฌา อโลโภ กุสลมูลํ, อยํ วุจฺจติ อโลโภ.}

\prose{ตตฺถ กตโม อโทโส, โย อโทโส อทุสฺสนา อทุสฺสิตตฺตํ อพฺยาปาโท อพฺยาปชฺโช อโทโส กุสลมูลํ, อยํ วุจฺจติ อโทโส.}

\prose{ตตฺถ กตโม อโมโห, ยา ปญฺญา ปชานนา ฯเปฯ อโมโห ธมฺมวิจโย สมฺมาทิฏฺฐิ, อยํ วุจฺจติ อโมโห.\csromanspacer{} อิเม วุจฺจนฺติ กุสลมูลา.}

\prose{ตตฺถ กตโม กุสลมูลปจฺจยา สงฺขาโร, ยา เจตนา สญฺเจตนา สญฺเจตยิตตฺตํ, อยํ วุจฺจติ กุสลมูลปจฺจยา สงฺขาโร.}

\prose{ตตฺถ กตมํ สงฺขารปจฺจยา วิญฺญาณํ ฯเปฯ วิญฺญาณปจฺจยา นามํ ฯเปฯ นามปจฺจยา ฉฏฺฐายตนํ ฯเปฯ ฉฏฺฐายตนปจฺจยา ผสฺโส ฯเปฯ ผสฺสปจฺจยา เวทนา ฯเปฯ อยํ วุจฺจติ ผสฺสปจฺจยา เวทนา.}

\prose{ตตฺถ กตโม เวทนาปจฺจยา ปสาโท, ยา สทฺธา สทฺทหนา โอกปฺปนา อภิปฺปสาโท, อยํ วุจฺจติ เวทนาปจฺจยา ปสาโท.}

\prose{ตตฺถ กตโม ปสาทปจฺจยา อธิโมกฺโข, โย จิตฺตสฺส อธิโมกฺโข อธิมุจฺจนา ตทธิมุตฺตตา, อยํ วุจฺจติ ปสาทปจฺจยา อธิโมกฺโข.}

\prose{ตตฺถ กตโม อธิโมกฺขปจฺจยา ภโว, ฐเปตฺวา อธิโมกฺขํ เวทนากฺขนฺโธ สญฺญากฺขนฺโธ สงฺขารกฺขนฺโธ วิญฺญาณกฺขนฺโธ, อยํ วุจฺจติ อธิโมกฺขปจฺจยา ภโว ฯเปฯ เตน วุจฺจติ “เอวเมตสฺส เกวลสฺส ทุกฺขกฺขนฺธสฺส สมุทโย โหตี”ติ.}

\proseitem{294}{กตเม ธมฺมา กุสลา, ยสฺมิํ สมเย กามาวจรํ กุสลํ จิตฺตํ อุปฺปนฺนํ โหติ โสมนสฺสสหคตํ ญาณสมฺปยุตฺตํ สสงฺขาเรน รูปารมฺมณํ วา ฯเปฯ โสมนสฺสสหคตํ ญาณวิปฺปยุตฺตํ รูปารมฺมณํ วา ฯเปฯ โสมนสฺสสหคตํ ญาณวิปฺปยุตฺตํ สสงฺขาเรน รูปารมฺมณํ วา สทฺทารมฺมณํ วา คนฺธารมฺมณํ วา รสารมฺมณํ วา โผฏฺฐพฺพารมฺมณํ วา ธมฺมารมฺมณํ วา ยํ ยํ วา ปนารพฺภ, ตสฺมิํ สมเย กุสลมูลปจฺจยา สงฺขาโร, สงฺขารปจฺจยา}

\vspace{\tipitakaparskip}
\csromancenter{ปฏิจฺจสมุปฺปาทวิภงฺค วิญฺญาณํ, วิญฺญาณปจฺจยา นามํ, นามปจฺจยา ฉฏฺฐายตนํ, ฉฏฺฐายตนปจฺจยา ผสฺโส, ผสฺสปจฺจยา เวทนา, เวทนาปจฺจยา ปสาโท, ปสาทปจฺจยา อธิโมกฺโข, อธิโมกฺขปจฺจยา ภโว, ภวปจฺจยา ชาติ, ชาติปจฺจยา ชรามรณํ, เอวเมตสฺส เกวลสฺส ทุกฺขกฺขนฺธสฺส สมุทโย โหติ.}
\proseitem{295}{\csromanfolio{177}ตตฺถ กตเม กุสลมูลา, อโลโภ อโทโส.}

\prose{ตตฺถ กตโม อโลโภ, โย อโลโภ อลุพฺภนา อลุพฺภิตตฺตํ อสาราโค อสารชฺชนา อสารชฺชิตตฺตํ อนภิชฺฌา อโลโภ กุสลมูลํ, อยํ วุจฺจติ อโลโภ.}

\prose{ตตฺถ กตโม อโทโส, โย อโทโส อทุสฺสนา อทุสฺสิตตฺตํ อพฺยาปาโท อพฺยาปชฺโช อโทโส กุสลมูลํ, อยํ วุจฺจติ อโทโส.\csromanspacer{} อิเม วุจฺจนฺติ กุสลมูลา.}

\prose{ตตฺถ กตโม กุสลมูลปจฺจยา สงฺขาโร, ยา เจตนา สญฺเจตนา สญฺเจตยิตตฺตํ, อยํ วุจฺจติ กุสลมูลปจฺจยา สงฺขาโร ฯเปฯ อยํ วุจฺจติ ฉฏฺฐายตนปจฺจยา ผสฺโส.}

\prose{ตตฺถ กตมา ผสฺสปจฺจยา เวทนา, ยํ เจตสิกํ สาตํ เจตสิกํ สุขํ เจโตสมฺผสฺสชํ สาตํ สุขํ เวทยิตํ เจโตสมฺผสฺสชา สาตา สุขา เวทนา, อยํ วุจฺจติ ผสฺสปจฺจยา เวทนา ฯเปฯ เตน วุจฺจติ “เอวเมตสฺส เกวลสฺส ทุกฺขกฺขนฺธสฺส สมุทโย โหตี”ติ.}

\proseitem{296}{กตเม ธมฺมา กุสลา, ยสฺมิํ สมเย กามาวจรํ กุสลํ จิตฺตํ อุปฺปนฺนํ โหติ อุเปกฺขาสหคตํ ญาณสมฺปยุตฺตํ รูปารมฺมณํ วา ฯเปฯ อุเปกฺขาสหคตํ ญาณสมฺปยุตฺตํ สสงฺขาเรน รูปารมฺมณํ วา สทฺทารมฺมณํ วา คนฺธารมฺมณํ วา รสารมฺมณํ วา โผฏฺฐพฺพารมฺมณํ วา ธมฺมารมฺมณํ วา ยํ ยํ วา ปนารพฺภ, ตสฺมิํ สมเย กุสลมูลปจฺจยา สงฺขาโร, สงฺขารปจฺจยา วิญฺญาณํ, วิญฺญาณปจฺจยา นามํ, นามปจฺจยา ฉฏฺฐายตนํ, ฉฏฺฐายตนปจฺจยา ผสฺโส, ผสฺสปจฺจยา เวทนา, เวทนาปจฺจยา ปสาโท, ปสาทปจฺจยา อธิโมกฺโข, อธิโมกฺขปจฺจยา ภโว, ภวปจฺจยา ชาติ, ชาติปจฺจยา ชรามรณํ, เอวเมตสฺส เกวลสฺส ทุกฺขกฺขนฺธสฺส สมุทโย โหติ.}

\proseitem{297}{\csromanfolio{178}ตตฺถ กตเม กุสลมูลา, อโลโภ อโทโส อโมโห, อิเม วุจฺจนฺติ กุสลมูลา.}

\prose{ตตฺถ กตโม กุสลมูลปจฺจยา สงฺขาโร, ยา เจตนา สญฺเจตนา สญฺเจตยิตตฺตํ, อยํ วุจฺจติ กุสลมูลปจฺจยา สงฺขาโร ฯเปฯ อยํ วุจฺจติ ฉฏฺฐายตนปจฺจยา ผสฺโส.}

\prose{ตตฺถ กตมา ผสฺสปจฺจยา เวทนา, ยํ เจตสิกํ เนว สาตํ นาสาตํ เจโตสมฺผสฺสชํ อทุกฺขมสุขํ เวทยิตํ เจโตสมฺผสฺสชา อทุกฺขมสุขา เวทนา, อยํ วุจฺจติ ผสฺสปจฺจยา เวทนา ฯเปฯ เตน วุจฺจติ “เอวเมตสฺส เกวลสฺส ทุกฺขกฺขนฺธสฺส สมุทโย โหตี”ติ.}

\proseitem{298}{กตเม ธมฺมา กุสลา, ยสฺมิํ สมเย กามาวจรํ กุสลํ จิตฺตํ อุปฺปนฺนํ โหติ อุเปกฺขาสหคตํ ญาณวิปฺปยุตฺตํ รูปารมฺมณํ วา ฯเปฯ อุเปกฺขาสหคตํ ญาณวิปฺปยุตฺตํ สสงฺขาเรน รูปารมฺมณํ วา สทฺทารมฺมณํ วา คนฺธารมฺมณํ วา รสารมฺมณํ วา โผฏฺฐพฺพารมฺมณํ วา ธมฺมารมฺมณํ วา ยํ ยํ วา ปนารพฺภ, ตสฺมิํ สมเย กุสลมูลปจฺจยา สงฺขาโร, สงฺขารปจฺจยา วิญฺญาณํ, วิญฺญาณปจฺจยา นามํ, นามปจฺจยา ฉฏฺฐายตนํ, ฉฏฺฐายตนปจฺจยา ผสฺโส, ผสฺสปจฺจยา เวทนา, เวทนาปจฺจยา ปสาโท, ปสาทปจฺจยา อธิโมกฺโข, อธิโมกฺขปจฺจยา ภโว, ภวปจฺจยา ชาติ, ชาติปจฺจยา ชรามรณํ, เอวเมตสฺส เกวลสฺส ทุกฺขกฺขนฺธสฺส สมุทโย โหติ.}

\proseitem{299}{ตตฺถ กตเม กุสลมูลา, อโลโภ อโทโส, อิเม วุจฺจนฺติ กุสลมูลา.}

\prose{ตตฺถ กตโม กุสลมูลปจฺจยา สงฺขาโร, ยา เจตนา สญฺเจตนา สญฺเจตยิตตฺตํ, อยํ วุจฺจติ กุสลมูลปจฺจยา สงฺขาโร ฯเปฯ เตน วุจฺจติ “เอวเมตสฺส เกวลสฺส ทุกฺขกฺขนฺธสฺส สมุทโย โหตี”ติ.}

\proseitem{300}{กตเม ธมฺมา กุสลา, ยสฺมิํ สมเย รูปูปปตฺติยา มคฺคํ ภาเวติ วิวิจฺเจว กาเมหิ ฯเปฯ ปฐมํ ฌานํ อุปสมฺปชฺช วิหรติ ปถวีกสิณํ\footnote{ปฐวีกสิณํ (สี, สฺยา)}, ตสฺมิํ สมเย กุสลมูลปจฺจยา สงฺขาโร, สงฺขารปจฺจยา}

\vspace{\tipitakaparskip}
\csromancenter{ปฏิจฺจสมุปฺปาทวิภงฺค วิญฺญาณํ, วิญฺญาณปจฺจยา นามํ, นามปจฺจยา ฉฏฺฐายตนํ, ฉฏฺฐายตนปจฺจยา ผสฺโส, ผสฺสปจฺจยา เวทนา, เวทนาปจฺจยา ปสาโท, ปสาทปจฺจยา อธิโมกฺโข, อธิโมกฺขปจฺจยา ภโว, ภวปจฺจยา ชาติ, ชาติปจฺจยา ชรามรณํ, เอวเมตสฺส เกวลสฺส ทุกฺขกฺขนฺธสฺส สมุทโย โหติ.}
\proseitem{301}{\csromanfolio{179}ตตฺถ กตเม กุสลมูลา, อโลโภ อโทโส อโมโห, อิเม วุจฺจนฺติ กุสลมูลา.}

\prose{ตตฺถ กตโม กุสลมูลปจฺจยา สงฺขาโร, ยา เจตนา สญฺเจตนา สญฺเจตยิตตฺตํ, อยํ วุจฺจติ กุสลมูลปจฺจยา สงฺขาโร ฯเปฯ เตน วุจฺจติ “เอวเมตสฺส เกวลสฺส ทุกฺขกฺขนฺธสฺส สมุทโย โหตี”ติ.}

\proseitem{302}{กตเม ธมฺมา กุสลา, ยสฺมิํ สมเย อรูปูปปตฺติยา มคฺคํ ภาเวติ สพฺพโส อากิญฺจญฺญายตนํ สมติกฺกมฺม เนวสญฺญานาสญฺญายตนสญฺญาสหคตํ สุขสฺส จ ปหานา ฯเปฯ จตุตฺถํ ฌานํ อุปสมฺปชฺช วิหรติ, ตสฺมิํ สมเย กุสลมูลปจฺจยา สงฺขาโร, สงฺขารปจฺจยา วิญฺญาณํ, วิญฺญาณปจฺจยา นามํ, นามปจฺจยา ฉฏฺฐายตนํ, ฉฏฺฐายตนปจฺจยา ผสฺโส, ผสฺสปจฺจยา เวทนา, เวทนาปจฺจยา ปสาโท, ปสาทปจฺจยา อธิโมกฺโข, อธิโมกฺขปจฺจยา ภโว, ภวปจฺจยา ชาติ, ชาติปจฺจยา ชรามรณํ, เอวเมตสฺส เกวลสฺส ทุกฺขกฺขนฺธสฺส สมุทโย โหติ.}

\proseitem{303}{ตตฺถ กตเม กุสลมูลา, อโลโภ อโทโส อโมโห, อิเม วุจฺจนฺติ กุสลมูลา.}

\prose{ตตฺถ กตโม กุสลมูลปจฺจยา สงฺขาโร, ยา เจตนา สญฺเจตนา สญฺเจตยิตตฺตํ, อยํ วุจฺจติ กุสลมูลปจฺจยา สงฺขาโร ฯเปฯ อยํ วุจฺจติ ฉฏฺฐายตนปจฺจยา ผสฺโส.}

\prose{ตตฺถ กตมา ผสฺสปจฺจยา เวทนา, ยํ เจตสิกํ สาตํ เจตสิกํ สุขํ เจโตสมฺผสฺสชํ สาตํ สุขํ เวทยิตํ เจโตสมฺผสฺสชา สาตา สุขา เวทนา, อยํ วุจฺจติ ผสฺสปจฺจยา เวทนา ฯเปฯ เตน วุจฺจติ “เอวเมตสฺส เกวลสฺส ทุกฺขกฺขนฺธสฺส สมุทโย โหตี”ติ.}

\proseitem{304}{\csromanfolio{180}กตเม ธมฺมา กุสลา, ยสฺมิํ สมเย โลกุตฺตรํ ฌานํ ภาเวติ นิยฺยานิกํ อปจยคามิํ ทิฏฺฐิคตานํ ปหานาย ปฐมาย ภูมิยา ปตฺติยา วิวิจฺเจว กาเมหิ ฯเปฯ ปฐมํ ฌานํ อุปสมฺปชฺช วิหรติ ทุกฺขปฏิปทํ ทนฺธาภิญฺญํ, ตสฺมิํ สมเย กุสลมูลปจฺจยา สงฺขาโร, สงฺขารปจฺจยา วิญฺญาณํ, วิญฺญาณปจฺจยา นามํ, นามปจฺจยา ฉฏฺฐายตนํ, ฉฏฺฐายตนปจฺจยา ผสฺโส, ผสฺสปจฺจยา เวทนา, เวทนาปจฺจยา ปสาโท, ปสาทปจฺจยา อธิโมกฺโข, อธิโมกฺขปจฺจยา ภโว, ภวปจฺจยา ชาติ, ชาติปจฺจยา ชรามรณํ, เอวเมเตสํ ธมฺมานํ สมุทโย โหติ.}

\proseitem{305}{ตตฺถ กตเม กุสลมูลา, อโลโภ อโทโส อโมโห.}

\prose{ตตฺถ กตโม อโลโภ ฯเปฯ.\csromanspacer{} อโทโส ฯเปฯ.\csromanspacer{} อโมโห, ยา ปญฺญา ปชานนา ฯเปฯ.\csromanspacer{} อโมโห ธมฺมวิจโย สมฺมาทิฏฺฐิ ธมฺมวิจยสมฺโพชฺฌงฺโค มคฺคงฺคํ มคฺคปริยาปนฺนํ, อยํ วุจฺจติ อโมโห.\csromanspacer{} อิเม วุจฺจนฺติ กุสลมูลา.}

\prose{ตตฺถ กตโม กุสลมูลปจฺจยา สงฺขาโร, ยา เจตนา สญฺเจตนา สญฺเจตยิตตฺตํ, อยํ วุจฺจติ กุสลมูลปจฺจยา สงฺขาโร ฯเปฯ อยํ วุจฺจติ ฉฏฺฐายตนปจฺจยา ผสฺโส.}

\prose{ตตฺถ กตมา ผสฺสปจฺจยา เวทนา, ยํ เจตสิกํ สาตํ เจตสิกํ สุขํ เจโตสมฺผสฺสชํ สาตํ สุขํ เวทยิตํ เจโตสมฺผสฺสชา สาตา สุขา เวทนา, อยํ วุจฺจติ ผสฺสปจฺจยา เวทนา.}

\prose{ตตฺถ กตโม เวทนาปจฺจยา ปสาโท, ยา สทฺธา สทฺทหนา โอกปฺปนา อภิปฺปสาโท, อยํ วุจฺจติ เวทนาปจฺจยา ปสาโท.}

\prose{ตตฺถ กตโม ปสาทปจฺจยา อธิโมกฺโข, โย จิตฺตสฺส อธิโมกฺโข อธิมุจฺจนา ตทธิมุตฺตตา, อยํ วุจฺจติ ปสาทปจฺจยา อธิโมกฺโข.}

\prose{ตตฺถ กตโม อธิโมกฺขปจฺจยา ภโว, ฐเปตฺวา อธิโมกฺขํ เวทนากฺขนฺโธ สญฺญากฺขนฺโธ สงฺขารกฺขนฺโธ วิญฺญาณกฺขนฺโธ, อยํ วุจฺจติ อธิโมกฺขปจฺจยา ภโว ฯเปฯ อยํ วุจฺจติ ชาติปจฺจยา ชรามรณํ.}

\csromanheader{2}{6. ปฏิจฺจสมุปฺปาทวิภงฺค}
\prose{\csromanfolio{181}เอวเมเตสํ ธมฺมานํ สมุทโย โหตีติ เอวเมเตสํ ธมฺมานํ สงฺคติ โหติ, สมาคโม โหติ, สโมธานํ โหติ, ปาตุภาโว โหติ, เตน วุจฺจติ “เอวเมเตสํ ธมฺมานํ สมุทโย โหตี”ติ.}

\csromanheader{2}{กุสลนิทฺเทโส.}
\csromansectionrule
\csromanmatikaanchor{mtk.57}
\csromantocmarkref{subparagraph}{11. อพฺยากตนิทฺเทส}{mtk.57}
\bookmark[dest={mtk.57},level=5]{11. อพฺยากตนิทฺเทส}
\csromanheader{6}{11. อพฺยากตนิทฺเทส}
\proseitem{306}{กตเม ธมฺมา อพฺยากตา, ยสฺมิํ สมเย กามาวจรสฺส กุสลสฺส กมฺมสฺส กตตฺตา อุปจิตตฺตา วิปากํ จกฺขุวิญฺญาณํ อุปฺปนฺนํ โหติ อุเปกฺขาสหคตํ รูปารมฺมณํ, ตสฺมิํ สมเย สงฺขารปจฺจยา วิญฺญาณํ, วิญฺญาณปจฺจยา นามํ, นามปจฺจยา ฉฏฺฐายตนํ, ฉฏฺฐายตนปจฺจยา ผสฺโส, ผสฺสปจฺจยา เวทนา, เวทนาปจฺจยา ภโว, ภวปจฺจยา ชาติ, ชาติปจฺจยา ชรามรณํ, เอวเมตสฺส เกวลสฺส ทุกฺขกฺขนฺธสฺส สมุทโย โหติ.}

\proseitem{307}{ตตฺถ กตโม สงฺขาโร, ยา เจตนา สญฺเจตนา สญฺเจตยิตตฺตํ, อยํ วุจฺจติ สงฺขาโร.}

\prose{ตตฺถ กตมํ สงฺขารปจฺจยา วิญฺญาณํ, ยํ จิตฺตํ มโน มานสํ ฯเปฯ ตชฺชาจกฺขุวิญฺญาณธาตุ, อิทํ วุจฺจติ สงฺขารปจฺจยา วิญฺญาณํ.}

\prose{ตตฺถ กตมํ วิญฺญาณปจฺจยา นามํ, เวทนากฺขนฺโธ สญฺญากฺขนฺโธ สงฺขารกฺขนฺโธ, อิทํ วุจฺจติ วิญฺญาณปจฺจยา นามํ.}

\prose{ตตฺถ กตมํ นามปจฺจยา ฉฏฺฐายตนํ, ยํ จิตฺตํ มโน มานสํ ฯเปฯ ตชฺชาจกฺขุวิญฺญาณธาตุ, อิทํ วุจฺจติ นามปจฺจยา ฉฏฺฐายตนํ.}

\prose{ตตฺถ กตโม ฉฏฺฐายตนปจฺจยา ผสฺโส, โย ผสฺโส ผุสนา สมฺผุสนา สมฺผุสิตตฺตํ, อยํ วุจฺจติ ฉฏฺฐายตนปจฺจยา ผสฺโส.}

\prose{ตตฺถ กตมา ผสฺสปจฺจยา เวทนา, ยํ เจตสิกํ เนว สาตํ นาสาตํ เจโตสมฺผสฺสชํ อทุกฺขมสุขํ เวทยิตํ เจโตสมฺผสฺสชา อทุกฺขมสุขา เวทนา, อยํ วุจฺจติ ผสฺสปจฺจยา เวทนา.}

\prose{\csromanfolio{182}ตตฺถ กตโม เวทนาปจฺจยา ภโว, ฐเปตฺวา เวทนํ สญฺญากฺขนฺโธ สงฺขารกฺขนฺโธ วิญฺญาณกฺขนฺโธ, อยํ วุจฺจติ เวทนาปจฺจยา ภโว ฯเปฯ เตน วุจฺจติ “เอวเมตสฺส เกวลสฺส ทุกฺขกฺขนฺธสฺส สมุทโย โหตี”ติ.}

\proseitem{308}{ตสฺมิํ สมเย สงฺขารปจฺจยา วิญฺญาณํ สงฺขารเหตุกํ, วิญฺญาณปจฺจยา นามํ วิญฺญาณเหตุกํ, นามปจฺจยา ฉฏฺฐายตนํ นามเหตุกํ, ฉฏฺฐายตนปจฺจยา ผสฺโส ฉฏฺฐายตนเหตุโก, ผสฺสปจฺจยา เวทนา ผสฺสเหตุกา, เวทนาปจฺจยา ภโว, ภวปจฺจยา ชาติ, ชาติปจฺจยา ชรามรณํ, เอวเมตสฺส เกวลสฺส ทุกฺขกฺขนฺธสฺส สมุทโย โหติ.}

\proseitem{309}{ตสฺมิํ สมเย สงฺขารปจฺจยา วิญฺญาณํ สงฺขารสมฺปยุตฺตํ, วิญฺญาณปจฺจยา นามํ วิญฺญาณสมฺปยุตฺตํ, นามปจฺจยา ฉฏฺฐายตนํ นามสมฺปยุตฺตํ, ฉฏฺฐายตนปจฺจยา ผสฺโส ฉฏฺฐายตนสมฺปยุตฺโต, ผสฺสปจฺจยา เวทนา ผสฺสสมฺปยุตฺตา, เวทนาปจฺจยา ภโว, ภวปจฺจยา ชาติ, ชาติปจฺจยา ชรามรณํ, เอวเมตสฺส เกวลสฺส ทุกฺขกฺขนฺธสฺส สมุทโย โหติ.}

\proseitem{310}{ตสฺมิํ สมเย สงฺขารปจฺจยา วิญฺญาณํ, วิญฺญาณปจฺจยาปิ สงฺขาโร, วิญฺญาณปจฺจยา นามํ, นามปจฺจยาปิ วิญฺญาณํ, นามปจฺจยา ฉฏฺฐายตนํ, ฉฏฺฐายตนปจฺจยาปิ นามํ, ฉฏฺฐายตนปจฺจยา ผสฺโส, ผสฺสปจฺจยาปิ ฉฏฺฐายตนํ, ผสฺสปจฺจยา เวทนา, เวทนาปจฺจยาปิ ผสฺโส, เวทนาปจฺจยา ภโว, ภวปจฺจยา ชาติ, ชาติปจฺจยา ชรามรณํ, เอวเมตสฺส เกวลสฺส ทุกฺขกฺขนฺธสฺส สมุทโย โหติ.}

\proseitem{311}{กตเม ธมฺมา อพฺยากตา, ยสฺมิํ สมเย กามาวจรสฺส กุสลสฺส กมฺมสฺส กตตฺตา อุปจิตตฺตา วิปากํ โสตวิญฺญาณํ อุปฺปนฺนํ โหติ อุเปกฺขาสหคตํ สทฺทารมฺมณํ ฯเปฯ ฆานวิญฺญาณํ อุปฺปนฺนํ โหติ อุเปกฺขาสหคตํ คนฺธารมฺมณํ ฯเปฯ ชิวฺหาวิญฺญาณํ อุปฺปนฺนํ โหติ อุเปกฺขาสหคตํ รสารมฺมณํ ฯเปฯ กายวิญฺญาณํ อุปฺปนฺนํ โหติ สุขสหคตํ โผฏฺฐพฺพารมฺมณํ, ตสฺมิํ สมเย สงฺขารปจฺจยา วิญฺญาณํ, วิญฺญาณปจฺจยา นามํ, นามปจฺจยา ฉฏฺฐายตนํ, ฉฏฺฐายตนปจฺจยา ผสฺโส, ผสฺสปจฺจยา เวทนา, เวทนาปจฺจยา ภโว, ภวปจฺจยา ชาติ, ชาติปจฺจยา ชรามรณํ, เอวเมตสฺส เกวลสฺส ทุกฺขกฺขนฺธสฺส สมุทโย โหติ.}

\csromanheader{2}{6. ปฏิจฺจสมุปฺปาทวิภงฺค}
\proseitem{312}{\csromanfolio{183}ตตฺถ กตโม สงฺขาโร, ยา เจตนา สญฺเจตนา สญฺเจตยิตตฺตํ, อยํ วุจฺจติ สงฺขาโร ฯเปฯ อยํ วุจฺจติ ฉฏฺฐายตนปจฺจยา ผสฺโส.}

\prose{ตตฺถ กตมา ผสฺสปจฺจยา เวทนา, ยํ กายิกํ สาตํ กายิกํ สุขํ กายสมฺผสฺสชํ สาตํ สุขํ เวทยิตํ กายสมฺผสฺสชา สาตา สุขา เวทนา, อยํ วุจฺจติ ผสฺสปจฺจยา เวทนา.}

\prose{ตตฺถ กตโม เวทนาปจฺจยา ภโว, ฐเปตฺวา เวทนํ สญฺญากฺขนฺโธ สงฺขารกฺขนฺโธ วิญฺญาณกฺขนฺโธ, อยํ วุจฺจติ เวทนาปจฺจยา ภโว ฯเปฯ เตน วุจฺจติ “เอวเมตสฺส เกวลสฺส ทุกฺขกฺขนฺธสฺส สมุทโย โหตี”ติ.}

\proseitem{313}{กตเม ธมฺมา อพฺยากตา, ยสฺมิํ สมเย กามาวจรสฺส กุสลสฺส กมฺมสฺส กตตฺตา อุปจิตตฺตา วิปากา มโนธาตุ อุปฺปนฺนา โหติ อุเปกฺขาสหคตา รูปารมฺมณา วา สทฺทารมฺมณา วา คนฺธารมฺมณา วา รสารมฺมณา วา โผฏฺฐพฺพารมฺมณา วา ยํ ยํ วา ปนารพฺภ, ตสฺมิํ สมเย สงฺขารปจฺจยา วิญฺญาณํ, วิญฺญาณปจฺจยา นามํ, นามปจฺจยา ฉฏฺฐายตนํ, ฉฏฺฐายตนปจฺจยา ผสฺโส, ผสฺสปจฺจยา เวทนา, เวทนาปจฺจยา อธิโมกฺโข, อธิโมกฺขปจฺจยา ภโว, ภวปจฺจยา ชาติ, ชาติปจฺจยา ชรามรณํ, เอวเมตสฺส เกวลสฺส ทุกฺขกฺขนฺธสฺส สมุทโย โหติ.}

\proseitem{314}{ตตฺถ กตโม สงฺขาโร, ยา เจตนา สญฺเจตนา สญฺเจตยิตตฺตํ, อยํ วุจฺจติ สงฺขาโร.}

\prose{ตตฺถ กตมํ สงฺขารปจฺจยา วิญฺญาณํ, ยํ จิตฺตํ มโน มานสํ ฯเปฯ ตชฺชามโนธาตุ, อิทํ วุจฺจติ สงฺขารปจฺจยา วิญฺญาณํ ฯเปฯ อยํ วุจฺจติ ฉฏฺฐายตนปจฺจยา ผสฺโส.}

\prose{ตตฺถ กตมา ผสฺสปจฺจยา เวทนา, ยํ เจตสิกํ เนว สาตํ นาสาตํ เจโตสมฺผสฺสชํ อทุกฺขมสุขํ เวทยิตํ เจโตสมฺผสฺสชา อทุกฺขมสุขา เวทนา, อยํ วุจฺจติ ผสฺสปจฺจยา เวทนา.}

\prose{ตตฺถ กตโม เวทนาปจฺจยา อธิโมกฺโข, โย จิตฺตสฺส อธิโมกฺโข อธิมุจฺจนา ตทธิมุตฺตตา, อยํ วุจฺจติ เวทนาปจฺจยา อธิโมกฺโข.}

\prose{\csromanfolio{184}ตตฺถ กตโม อธิโมกฺขปจฺจยา ภโว, ฐเปตฺวา อธิโมกฺขํ เวทนากฺขนฺโธ สญฺญากฺขนฺโธ สงฺขารกฺขนฺโธ วิญฺญาณกฺขนฺโธ, อยํ วุจฺจติ อธิโมกฺขปจฺจยา ภโว ฯเปฯ เตน วุจฺจติ “เอวเมตสฺส เกวลสฺส ทุกฺขกฺขนฺธสฺส สมุทโย โหตี”ติ.}

\proseitem{315}{กตเม ธมฺมา อพฺยากตา, ยสฺมิํ สมเย กามาวจรสฺส กุสลสฺส กมฺมสฺส กตตฺตา อุปจิตตฺตา วิปากา มโนวิญฺญาณธาตุ อุปฺปนฺนา โหติ โสมนสฺสสหคตา รูปารมฺมณา วา สทฺทารมฺมณา วา คนฺธารมฺมณา วา รสารมฺมณา วา โผฏฺฐพฺพารมฺมณา วา ธมฺมารมฺมณา วา ยํ ยํ วา ปนารพฺภ, ตสฺมิํ สมเย สงฺขารปจฺจยา วิญฺญาณํ, วิญฺญาณปจฺจยา นามํ, นามปจฺจยา ฉฏฺฐายตนํ, ฉฏฺฐายตนปจฺจยา ผสฺโส, ผสฺสปจฺจยา เวทนา, เวทนาปจฺจยา อธิโมกฺโข, อธิโมกฺขปจฺจยา ภโว, ภวปจฺจยา ชาติ, ชาติปจฺจยา ชรามรณํ, เอวเมตสฺส เกวลสฺส ทุกฺขกฺขนฺธสฺส สมุทโย โหติ.}

\proseitem{316}{ตตฺถ กตโม สงฺขาโร, ยา เจตนา สญฺเจตนา สญฺเจตยิตตฺตํ, อยํ วุจฺจติ สงฺขาโร.}

\prose{ตตฺถ กตมํ สงฺขารปจฺจยา วิญฺญาณํ, ยํ จิตฺตํ มโน มานสํ ฯเปฯ ตชฺชามโนวิญฺญาณธาตุ, อิทํ วุจฺจติ สงฺขารปจฺจยา วิญฺญาณํ ฯเปฯ อยํ วุจฺจติ ฉฏฺฐายตนปจฺจยา ผสฺโส.}

\prose{ตตฺถ กตมา ผสฺสปจฺจยา เวทนา, ยํ เจตสิกํ สาตํ เจตสิกํ สุขํ เจโตสมฺผสฺสชํ สาตํ สุขํ เวทยิตํ เจโตสมฺผสฺสชา สาตา สุขา เวทนา, อยํ วุจฺจติ ผสฺสปจฺจยา เวทนา.}

\prose{ตตฺถ กตโม เวทนาปจฺจยา อธิโมกฺโข, โย จิตฺตสฺส อธิโมกฺโข อธิมุจฺจนา ตทธิมุตฺตตา, อยํ วุจฺจติ เวทนาปจฺจยา อธิโมกฺโข.}

\prose{ตตฺถ กตโม อธิโมกฺขปจฺจยา ภโว, ฐเปตฺวา อธิโมกฺขํ เวทนากฺขนฺโธ สญฺญากฺขนฺโธ สงฺขารกฺขนฺโธ วิญฺญาณกฺขนฺโธ, อยํ วุจฺจติ อธิโมกฺขปจฺจยา ภโว ฯเปฯ เตน วุจฺจติ “เอวเมตสฺส เกวลสฺส ทุกฺขกฺขนฺธสฺส สมุทโย โหตี”ติ.}

\csromanheader{2}{6. ปฏิจฺจสมุปฺปาทวิภงฺค}
\proseitem{317}{\csromanfolio{185}กตเม ธมฺมา อพฺยากตา, ยสฺมิํ สมเย กามาวจรสฺส กุสลสฺส กมฺมสฺส กตตฺตา อุปจิตตฺตา วิปากา มโนวิญฺญาณธาตุ อุปฺปนฺนา โหติ อุเปกฺขาสหคตา รูปารมฺมณา วา สทฺทารมฺมณา วา คนฺธารมฺมณา วา รสารมฺมณา วา โผฏฺฐพฺพารมฺมณา วา ธมฺมารมฺมณา วา ยํ ยํ วา ปนารพฺภ, ตสฺมิํ สมเย สงฺขารปจฺจยา วิญฺญาณํ, วิญฺญาณปจฺจยา นามํ, นามปจฺจยา ฉฏฺฐายตนํ, ฉฏฺฐายตนปจฺจยา ผสฺโส, ผสฺสปจฺจยา เวทนา, เวทนาปจฺจยา อธิโมกฺโข, อธิโมกฺขปจฺจยา ภโว, ภวปจฺจยา ชาติ, ชาติปจฺจยา ชรามรณํ, เอวเมตสฺส เกวลสฺส ทุกฺขกฺขนฺธสฺส สมุทโย โหติ.}

\proseitem{318}{ตตฺถ กตโม สงฺขาโร, ยา เจตนา สญฺเจตนา สญฺเจตยิตตฺตํ, อยํ วุจฺจติ สงฺขาโร.}

\prose{ตตฺถ กตมํ สงฺขารปจฺจยา วิญฺญาณํ, ยํ จิตฺตํ มโน มานสํ ฯเปฯ ตชฺชามโนวิญฺญาณธาตุ, อิทํ วุจฺจติ สงฺขารปจฺจยา วิญฺญาณํ ฯเปฯ อยํ วุจฺจติ ฉฏฺฐายตนปจฺจยา ผสฺโส.}

\prose{ตตฺถ กตมา ผสฺสปจฺจยา เวทนา, ยํ เจตสิกํ เนว สาตํ นาสาตํ เจโตสมฺผสฺสชํ อทุกฺขมสุขํ เวทยิตํ เจโตสมฺผสฺสชา อทุกฺขมสุขา เวทนา, อยํ วุจฺจติ ผสฺสปจฺจยา เวทนา.}

\prose{ตตฺถ กตโม เวทนาปจฺจยา อธิโมกฺโข, โย จิตฺตสฺส อธิโมกฺโข อธิมุจฺจนา ตทธิมุตฺตตา, อยํ วุจฺจติ เวทนาปจฺจยา อธิโมกฺโข.}

\prose{ตตฺถ กตโม อธิโมกฺขปจฺจยา ภโว, ฐเปตฺวา อธิโมกฺขํ เวทนากฺขนฺโธ สญฺญากฺขนฺโธ สงฺขารกฺขนฺโธ วิญฺญาณกฺขนฺโธ, อยํ วุจฺจติ อธิโมกฺขปจฺจยา ภโว ฯเปฯ เตน วุจฺจติ “เอวเมตสฺส เกวลสฺส ทุกฺขกฺขนฺธสฺส สมุทโย โหตี”ติ.}

\proseitem{319}{กตเม ธมฺมา อพฺยากตา, ยสฺมิํ สมเย กามาวจรสฺส กุสลสฺส กมฺมสฺส กตตฺตา อุปจิตตฺตา วิปากา มโนวิญฺญาณธาตุ อุปฺปนฺนา โหติ โสมนสฺสสหคตา ญาณสมฺปยุตฺตา ฯเปฯ โสมนสฺสสหคตา ญาณสมฺปยุตฺตา สสงฺขาเรน ฯเปฯ.\csromanspacer{} โสมนสฺสสหคตา ญาณวิปฺปยุตฺตา ฯเปฯ โสมนสฺสสหคตา ญาณวิปฺปยุตฺตา สสงฺขาเรน ฯเปฯ. \csromanfolio{186}อุเปกฺขาสหคตา ญาณสมฺปยุตฺตา ฯเปฯ อุเปกฺขาสหคตา ญาณสมฺปยุตฺตา สสงฺขาเรน ฯเปฯ อุเปกฺขาสหคตา ญาณวิปฺปยุตฺตา ฯเปฯ อุเปกฺขาสหคตา ญาณวิปฺปยุตฺตา สสงฺขาเรน รูปารมฺมณา วา ฯเปฯ ธมฺมารมฺมณา วา ยํ ยํ วา ปนารพฺภ, ตสฺมิํ สมเย สงฺขารปจฺจยา วิญฺญาณํ, วิญฺญาณปจฺจยา นามํ, นามปจฺจยา ฉฏฺฐายตนํ, ฉฏฺฐายตนปจฺจยา ผสฺโส, ผสฺสปจฺจยา เวทนา, เวทนาปจฺจยา ปสาโท, ปสาทปจฺจยา อธิโมกฺโข, อธิโมกฺขปจฺจยา ภโว, ภวปจฺจยา ชาติ, ชาติปจฺจยา ชรามรณํ, เอวเมตสฺส เกวลสฺส ทุกฺขกฺขนฺธสฺส สมุทโย โหติ.}

\proseitem{320}{ตตฺถ กตโม สงฺขาโร, ยา เจตนา สญฺเจตนา สญฺเจตยิตตฺตํ, อยํ วุจฺจติ สงฺขาโร.}

\prose{ตตฺถ กตมํ สงฺขารปจฺจยา วิญฺญาณํ, ยํ จิตฺตํ มโน มานสํ ฯเปฯ ตชฺชามโนวิญฺญาณธาตุ, อิทํ วุจฺจติ สงฺขารปจฺจยา วิญฺญาณํ ฯเปฯ อยํ วุจฺจติ ผสฺสปจฺจยา เวทนา.}

\prose{ตตฺถ กตโม เวทนาปจฺจยา ปสาโท, ยา สทฺธา สทฺทหนา โอกปฺปนา อภิปฺปสาโท, อยํ วุจฺจติ เวทนาปจฺจยา ปสาโท.}

\prose{ตตฺถ กตโม ปสาทปจฺจยา อธิโมกฺโข, โย จิตฺตสฺส อธิโมกฺโข อธิมุจฺจนา ตทธิมุตฺตตา, อยํ วุจฺจติ ปสาทปจฺจยา อธิโมกฺโข.}

\prose{ตตฺถ กตโม อธิโมกฺขปจฺจยา ภโว, ฐเปตฺวา อธิโมกฺขํ เวทนากฺขนฺโธ สญฺญากฺขนฺโธ สงฺขารกฺขนฺโธ วิญฺญาณกฺขนฺโธ, อยํ วุจฺจติ อธิโมกฺขปจฺจยา ภโว ฯเปฯ เตน วุจฺจติ “เอวเมตสฺส เกวลสฺส ทุกฺขกฺขนฺธสฺส สมุทโย โหตี”ติ.}

\proseitem{321}{กตเม ธมฺมา อพฺยากตา, ยสฺมิํ สมเย รูปูปปตฺติยา มคฺคํ ภาเวติ วิวิจฺเจว กาเมหิ ฯเปฯ ปฐมํ ฌานํ อุปสมฺปชฺช วิหรติ ปถวีกสิณํ, ตสฺมิํ สมเย ผสฺโส โหติ ฯเปฯ อวิกฺเขโป โหติ, อิเม ธมฺมา กุสลา.}

\prose{ตสฺเสว รูปาวจรสฺส กุสลสฺส กมฺมสฺส กตตฺตา อุปจิตตฺตา วิปากํ วิวิจฺเจว กาเมหิ ฯเปฯ ปฐมํ ฌานํ อุปสมฺปชฺช วิหรติ ปถวีกสิณํ, ตสฺมิํ สมเย สงฺขารปจฺจยา วิญฺญาณํ, วิญฺญาณปจฺจยา นามํ, นามปจฺจยา}

\vspace{\tipitakaparskip}
\csromancenter{ปฏิจฺจสมุปฺปาทวิภงฺค ฉฏฺฐายตนํ, ฉฏฺฐายตนปจฺจยา ผสฺโส, ผสฺสปจฺจยา เวทนา, เวทนาปจฺจยา ปสาโท, ปสาทปจฺจยา อธิโมกฺโข, อธิโมกฺขปจฺจยา ภโว, ภวปจฺจยา ชาติ, ชาติปจฺจายา ชรามรณํ, เอวเมตสฺส เกวลสฺส ทุกฺขกฺขนฺธสฺส สมุทโย โหติ.}
\proseitem{322}{\csromanfolio{187}กตเม ธมฺมา อพฺยากตา, ยสฺมิํ สมเย อรูปูปปตฺติยา มคฺคํ ภาเวติ สพฺพโส อากิญฺจญฺญายตนํ สมติกฺกมฺม เนวสญฺญานาสญฺญายตนสญฺญาสหคตํ สุขสฺส จ ปหานา ฯเปฯ จตุตฺถํ ฌานํ อุปสมฺปชฺช วิหรติ, ตสฺมิํ สมเย ผสฺโส โหติ ฯเปฯ อวิกฺเขโป โหติ, อิเม ธมฺมา กุสลา.}

\prose{ตสฺเสว อรูปาวจรสฺส กุสลสฺส กมฺมสฺส กตตฺตา อุปจิตตฺตา วิปากํ สพฺพโส อากิญฺจญฺญายตนํ สมติกฺกมฺม เนวสญฺญานาสญฺญายตนสญฺญาสหคตํ สุขสฺส จ ปหานา ฯเปฯ จตุตฺถํ ฌานํ อุปสมฺปชฺช วิหรติ, ตสฺมิํ สมเย สงฺขารปจฺจยา วิญฺญาณํ, วิญฺญาณปจฺจยา นามํ, นามปจฺจยา ฉฏฺฐายตนํ, ฉฏฺฐายตนปจฺจยา ผสฺโส, ผสฺสปจฺจยา เวทนา, เวทนาปจฺจยา ปสาโท, ปสาทปจฺจยา อธิโมกฺโข, อธิโมกฺขปจฺจยา ภโว, ภวปจฺจยา ชาติ, ชาติปจฺจยา ชรามรณํ, เอวเมตสฺส เกวลสฺส ทุกฺขกฺขนฺธสฺส สมุทโย โหติ.}

\proseitem{323}{กตเม ธมฺมา อพฺยากตา, ยสฺมิํ สมเย โลกุตฺตรํ ฌานํ ภาเวติ นิยฺยานิกํ อปจยคามิํ ทิฏฺฐิคตานํ ปหานาย ปฐมาย ภูมิยา ปตฺติยา วิวิจฺเจว กาเมหิ ฯเปฯ ปฐมํ ฌานํ อุปสมฺปชฺช วิหรติ ทุกฺขปฏิปทํ ทนฺธาภิญฺญํ, ตสฺมิํ สมเย ผสฺโส โหติ ฯเปฯ อวิกฺเขโป โหติ, อิเม ธมฺมา กุสลา.}

\prose{ตสฺเสว โลกุตฺตรสฺส กุสลสฺส ฌานสฺส กตตฺตา ภาวิตตฺตา วิปากํ วิวิจฺเจว กาเมหิ ฯเปฯ ปฐมํ ฌานํ อุปสมฺปชฺช วิหรติ ทุกฺขปฏิปทํ ทนฺธาภิญฺญํ, ตสฺมิํ สมเย สงฺขารปจฺจยา วิญฺญาณํ, วิญฺญาณปจฺจยา นามํ, นามปจฺจยา ฉฏฺฐายตนํ, ฉฏฺฐายตนปจฺจยา ผสฺโส, ผสฺสปจฺจยา เวทนา, เวทนาปจฺจยา ปสาโท, ปสาทปจฺจยา อธิโมกฺโข, อธิโมกฺขปจฺจยา ภโว, ภวปจฺจยา ชาติ, ชาติปจฺจยา ชรามรณํ, เอวเมเตสํ ธมฺมานํ สมุทโย โหติ.}

\proseitem{324}{\csromanfolio{188}กตเม ธมฺมา อพฺยากตา, ยสฺมิํ สมเย อกุสลสฺส กมฺมสฺส กตตฺตา อุปจิตตฺตา วิปากํ จกฺขุวิญฺญาณํ อุปฺปนฺนํ โหติ อุเปกฺขาสหคตํ รูปารมฺมณํ ฯเปฯ โสตวิญฺญาณํ อุปฺปนฺนํ โหติ อุเปกฺขาสหคตํ สทฺทารมฺมณํ ฯเปฯ ฆานวิญฺญาณํ อุปฺปนฺนํ โหติ อุเปกฺขาสหคตํ คนฺธารมฺมณํ ฯเปฯ ชิวฺหาวิญฺญาณํ อุปฺปนฺนํ โหติ อุเปกฺขาสหคตํ รสารมฺมณํ ฯเปฯ กายวิญฺญาณํ อุปฺปนฺนํ โหติ ทุกฺขสหคตํ โผฏฺฐพฺพารมฺมณํ, ตสฺมิํ สมเย สงฺขารปจฺจยา วิญฺญาณํ, วิญฺญาณปจฺจยา นามํ, นามปจฺจยา ฉฏฺฐายตนํ, ฉฏฺฐายตนปจฺจยา ผสฺโส, ผสฺสปจฺจยา เวทนา, เวทนาปจฺจยา ภโว, ภวปจฺจยา ชาติ, ชาติปจฺจยา ชรามรณํ, เอวเมตสฺส เกวลสฺส ทุกฺขกฺขนฺธสฺส สมุทโย โหติ.}

\proseitem{325}{ตตฺถ กตโม สงฺขาโร, ยา เจตนา สญฺเจตนา สญฺเจตยิตตฺตํ, อยํ วุจฺจติ สงฺขาโร.}

\prose{ตตฺถ กตมํ สงฺขารปจฺจยา วิญฺญาณํ, ยํ จิตฺตํ มโน มานสํ ฯเปฯ ตชฺชากายวิญฺญาณธาตุ, อิทํ วุจฺจติ สงฺขารปจฺจยา วิญฺญาณํ ฯเปฯ อยํ วุจฺจติ ฉฏฺฐายตนปจฺจยา ผสฺโส.}

\prose{ตตฺถ กตมา ผสฺสปจฺจยา เวทนา, ยํ กายิกํ อสาตํ กายิกํ ทุกฺขํ กายสมฺผสฺสชํ อสาตํ ทุกฺขํ เวทยิตํ กายสมฺผสฺสชา อสาตา ทุกฺขา เวทนา, อยํ วุจฺจติ ผสฺสปจฺจยา เวทนา.}

\prose{ตตฺถ กตโม เวทนาปจยา ภโว, ฐเปตฺวา เวทนํ สญฺญากฺขนฺโธ สงฺขารกฺขนฺโธ วิญฺญาณกฺขนฺโธ, อยํ วุจฺจติ เวทนาปจฺจยา ภโว ฯเปฯ เตน วุจฺจติ “เอวเมตสฺส เกวลสฺส ทุกฺขกฺขนฺธสฺส สมุทโย โหตี”ติ.}

\proseitem{326}{กตเม ธมฺมา อพฺยากตา, ยสฺมิํ สมเย อกุสลสฺส กมฺมสฺส กตตฺตา อุปจิตตฺตา วิปากา มโนธาตุ อุปฺปนฺนา โหติ อุเปกฺขาสหคตา รูปารมฺมณา วา สทฺทารมฺมณา วา คนฺธารมฺมณา วา รสารมฺมณา วา โผฏฺฐพฺพารมฺมณา วา ยํ ยํ วา ปนารพฺภ, ตสฺมิํ สมเย สงฺขารปจฺจยา วิญฺญาณํ, วิญฺญาณปจฺจยา นามํ, นามปจฺจยา ฉฏฺฐายตนํ, ฉฏฺฐายตนปจฺจยา ผสฺโส, ผสฺสปจฺจยา เวทนา, เวทนาปจฺจยา อธิโมกฺโข, อธิโมกฺขปจฺจยา ภโว, ภวปจฺจยา ชาติ, ชาติปจฺจยา ชรามรณํ, เอวเมตสฺส เกวลสฺส ทุกฺขกฺขนฺธสฺส สมุทโย โหติ.}

\csromanheader{2}{6. ปฏิจฺจสมุปฺปาทวิภงฺค}
\proseitem{327}{\csromanfolio{189}ตตฺถ กตโม สงฺขาโร, ยา เจตนา สญฺเจตนา สญฺเจตยิตตฺตํ, อยํ วุจฺจติ สงฺขาโร.}

\prose{ตตฺถ กตมํ สงฺขารปจฺจยา วิญฺญาณํ, ยํ จิตฺตํ มโน มานสํ ฯเปฯ ตชฺชามโนธาตุ, อิทํ วุจฺจติ สงฺขารปจฺจยา วิญฺญาณํ ฯเปฯ อยํ วุจฺจติ ฉฏฺฐายตนปจฺจยา ผสฺโส.}

\prose{ตตฺถ กตมา ผสฺสปจฺจยา เวทนา, ยํ เจตสิกํ เนว สาตํ นาสาตํ เจโตสมฺผสฺสชํ อทุกฺขมสุขํ เวทยิตํ เจโตสมฺผสฺสชา อทุกฺขมสุขา เวทนา, อยํ วุจฺจติ ผสฺสปจฺจยา เวทนา.}

\prose{ตตฺถ กตโม เวทนาปจฺจยา อธิโมกฺโข, โย จิตฺตสฺส อธิโมกฺโข อธิมุจฺจนา ตทธิมุตฺตตา, อยํ วุจฺจติ เวทนาปจฺจยา อธิโมกฺโข.}

\prose{ตตฺถ กตโม อธิโมกฺขปจฺจยา ภโว, ฐเปตฺวา อธิโมกฺขํ เวทนากฺขนฺโธ สญฺญากฺขนฺโธ สงฺขารกฺขนฺโธ วิญฺญาณกฺขนฺโธ, อยํ วุจฺจติ อธิโมกฺขปจฺจยา ภโว ฯเปฯ เตน วุจฺจติ “เอวเมตสฺส เกวลสฺส ทุกฺขกฺขนฺธสฺส สมุทโย โหตี”ติ.}

\proseitem{328}{กตเม ธมฺมา อพฺยากตา, ยสฺมิํ สมเย อกุสลสฺส กมฺมสฺส กตตฺตา อุปจิตตฺตา วิปากา มโนวิญฺญาณธาตุ อุปฺปนฺนา โหติ อุเปกฺขาสหคตา รูปารมฺมณา วา ฯเปฯ ธมฺมารมฺมณา วา ยํ ยํ วา ปนารพฺภ, ตสฺมิํ สมเย สงฺขารปจฺจยา วิญฺญาณํ, วิญฺญาณปจฺจยา นามํ, นามปจฺจยา ฉฏฺฐายตนํ, ฉฏฺฐายตนปจฺจยา ผสฺโส, ผสฺสปจฺจยา เวทนา, เวทนาปจฺจยา อธิโมกฺโข, อธิโมกฺขปจฺจยา ภโว, ภวปจฺจยา ชาติ, ชาติปจฺจยา ชรามรณํ, เอวเมตสฺส เกวลสฺส ทุกฺขกฺขนฺธสฺส สมุทโย โหติ.}

\proseitem{329}{ตตฺถ กตโม สงฺขาโร, ยา เจตนา สญฺเจตนา สญฺเจตยิตตฺตํ, อยํ วุจฺจติ สงฺขาโร.}

\prose{ตตฺถ กตมํ สงฺขารปจฺจยา วิญฺญาณํ, ยํ จิตฺตํ มโน มานสํ ฯเปฯ ตชฺชามโนวิญฺญาณธาตุ, อิทํ วุจฺจติ สงฺขารปจฺจยา วิญฺญาณํ ฯเปฯ เตน วุจฺจติ “เอวเมตสฺส เกวลสฺส ทุกฺขกฺขนฺธสฺส สมุทโย โหตี”ติ.}

\proseitem{330}{\csromanfolio{190}กตเม ธมฺมา อพฺยากตา, ยสฺมิํ สมเย มโนธาตุ อุปฺปนฺนา โหติ กิริยา เนว กุสลา นากุสลา น จ กมฺมวิปากา อุเปกฺขาสหคตา รูปารมฺมณา วา ฯเปฯ โผฏฺฐพฺพารมฺมณา วา ฯเปฯ มโนวิญฺญาณธาตุ อุปฺปนฺนา โหติ กิริยา เนว กุสลา นากุสลา น จ กมฺมวิปากา โสมนสฺสสหคตา รูปารมฺมณา วา ฯเปฯ ธมฺมารมฺมณา วา ฯเปฯ มโนวิญฺญาณธาตุ อุปฺปนฺนา โหติ กิริยา เนว กุสลา นากุสลา น จ กมฺมวิปากา อุเปกฺขาสหคตา รูปารมฺมณา วา ฯเปฯ ธมฺมารมฺมณา วา ยํ ยํ วา ปนารพฺภ, ตสฺมิํ สมเย สงฺขารปจฺจยา วิญฺญาณํ, วิญฺญาณปจฺจยา นามํ, นามปจฺจยา ฉฏฺฐายตนํ, ฉฏฺฐายตนปจฺจยา ผสฺโส, ผสฺสปจฺจยา เวทนา, เวทนาปจฺจยา อธิโมกฺโข, อธิโมกฺขปจฺจยา ภโว, ภวปจฺจยา ชาติ, ชาติปจฺจยา ชรามรณํ, เอวเมตสฺส เกวลสฺส ทุกฺขกฺขนฺธสฺส สมุทโย โหติ.}

\proseitem{331}{กตเม ธมฺมา อพฺยากตา, ยสฺมิํ สมเย มโนวิญฺญาณธาตุ อุปฺปนฺนา โหติ กิริยา เนว กุสลา นากุสลา น จ กมฺมวิปากา โสมนสฺสสหคตา ญาณสมฺปยุตฺตา ฯเปฯ โสมนสฺสสหคตา ญาณสมฺปยุตฺตา สสงฺขาเรน ฯเปฯ โสมนสฺสสหคตา ญาณวิปฺปยุตฺตา ฯเปฯ โสมนสฺสสหคตา ญาณวิปฺปยุตฺตา สสงฺขาเรน ฯเปฯ อุเปกฺขาสหคตา ญาณสมฺปยุตฺตา ฯเปฯ อุเปกฺขาสหคตา ญาณสมฺปยุตฺตา สสงฺขาเรน ฯเปฯ อุเปกฺขาสหคตา ญาณวิปฺปยุตฺตา ฯเปฯ อุเปกฺขาสหคตา ญาณวิปฺปยุตฺตา สสงฺขาเรน รูปารมฺมณา วา ฯเปฯ ธมฺมารมฺมณา วา ยํ ยํ วา ปนารพฺภ, ตสฺมิํ สมเย สงฺขารปจฺจยา วิญฺญาณํ, วิญฺญาณปจฺจยา นามํ, นามปจฺจยา ฉฏฺฐายตนํ, ฉฏฺฐายตนปจฺจยา ผสฺโส, ผสฺสปจฺจยา เวทนา, เวทนาปจฺจยา ปสาโท, ปสาทปจฺจยา อธิโมกฺโข, อธิโมกฺขปจฺจยา ภโว, ภวปจฺจยา ชาติ, ชาติปจฺจยา ชรามรณํ, เอวเมตสฺส เกวลสฺส ทุกฺขกฺขนฺธสฺส สมุทโย โหติ.}

\proseitem{332}{กตเม ธมฺมา อพฺยากตา, ยสฺมิํ สมเย รูปาวจรํ ฌานํ ภาเวติ กิริยํ เนว กุสลํ นากุสลํ น จ กมฺมวิปากํ ทิฏฺฐธมฺมสุขวิหารํ วิวิจฺเจว กาเมหิ ฯเปฯ ปฐมํ ฌานํ อุปสมฺปชฺช วิหรติ ปถวีกสิณํ, ตสฺมิํ สมเย สงฺขารปจฺจยา วิญฺญาณํ, วิญฺญาณปจฺจยา นามํ, นามปจฺจยา ฉฏฺฐายตนํ, ฉฏฺฐายตนปจฺจยา ผสฺโส, ผสฺสปจฺจยา เวทนา, เวทนาปจฺจยา ปสาโท, ปสาทปจฺจยา อธิโมกฺโข, อธิโมกฺขปจฺจยา ภโว,}

\vspace{\tipitakaparskip}
\csromancenter{ปฏิจฺจสมุปฺปาทวิภงฺค ภวปจฺจยา ชาติ, ชาติปจฺจยา ชรามรณํ, เอวเมตสฺส เกวลสฺส ทุกฺขกฺขนฺธสฺส สมุทโย โหติ.}
\proseitem{333}{\csromanfolio{191}กตเม ธมฺมา อพฺยากตา, ยสฺมิํ สมเย อรูปาวจรํ ฌานํ ภาเวติ กิริยํ เนว กุสลํ นากุสลํ น จ กมฺมวิปากํ ทิฏฺฐธมฺมสุขวิหารํ สพฺพโส อากิญฺจญฺญายตนํ สมติกฺกมฺม เนวสญฺญานาสญฺญายตนสญฺญาสหคตํ สุขสฺส จ ปหานา ฯเปฯ จตุตฺถํ ฌานํ อุปสมฺปชฺช วิหรติ, ตสฺมิํ สมเย สงฺขารปจฺจยา วิญฺญาณํ, วิญฺญาณปจฺจยา นามํ, นามปจฺจยา ฉฏฺฐายตนํ, ฉฏฺฐายตนปจฺจยา ผสฺโส, ผสฺสปจฺจยา เวทนา, เวทนาปจฺจยา ปสาโท, ปสาทปจฺจยา อธิโมกฺโข, อธิโมกฺขปจฺจยา ภโว, ภวปจฺจยา ชาติ, ชาติปจฺจยา ชรามรณํ, เอวเมตสฺส เกวลสฺส ทุกฺขกฺขนฺธสฺส สมุทโย โหติ.}

\csromanheader{2}{อพฺยากตนิทฺเทโส.}
\csromansectionrule
\csromanmatikaanchor{mtk.58}
\csromantocmarkref{subparagraph}{12. อวิชฺชามูลกกุสลนิทฺเทส}{mtk.58}
\bookmark[dest={mtk.58},level=5]{12. อวิชฺชามูลกกุสลนิทฺเทส}
\csromanheader{6}{12. อวิชฺชามูลกกุสลนิทฺเทส}
\proseitem{334}{กตเม ธมฺมา กุสลา, ยสฺมิํ สมเย กามาวจรํ กุสลํ จิตฺตํ อุปฺปนฺนํ โหติ โสมนสฺสสหคตํ ญาณสมฺปยุตฺตํ รูปารมฺมณํ วา ฯเปฯ ธมฺมารมฺมณํ วา ยํ ยํ วา ปนารพฺภ, ตสฺมิํ สมเย อวิชฺชาปจฺจยา สงฺขาโร, สงฺขารปจฺจยา วิญฺญาณํ, วิญฺญาณปจฺจยา นามํ, นามปจฺจยา ฉฏฺฐายตนํ, ฉฏฺฐายตนปจฺจยา ผสฺโส, ผสฺสปจฺจยา เวทนา, เวทนาปจฺจยา ปสาโท, ปสาทปจฺจยา อธิโมกฺโข, อธิโมกฺขปจฺจยา ภโว, ภวปจฺจยา ชาติ, ชาติปจฺจยา ชรามรณํ, เอวเมตสฺส เกวลสฺส ทุกฺขกฺขนฺธสฺส สมุทโย โหติ.}

\proseitem{335}{ตตฺถ กตโม อวิชฺชาปจฺจยา สงฺขาโร, ยา เจตนา สญฺเจตนา สญฺเจตยิตตฺตํ, อยํ วุจฺจติ อวิชฺชาปจฺจยา สงฺขาโร ฯเปฯ อยํ วุจฺจติ ฉฏฺฐายตนปจฺจยา ผสฺโส.}

\prose{ตตฺถ กตมา ผสฺสปจฺจยา เวทนา, ยํ เจตสิกํ สาตํ เจตสิกํ สุขํ เจโตสมฺผสฺสชํ สาตํ สุขํ เวทยิตํ เจโตสมฺผสฺสชา สาตา สุขา เวทนา, อยํ วุจฺจติ ผสฺสปจฺจยา เวทนา.}

\prose{\csromanfolio{192}ตตฺถ กตโม เวทนาปจฺจยา ปสาโท, ยา สทฺธา สทฺทหนา โอกปฺปนา อภิปฺปสาโท, อยํ วุจฺจติ เวทนาปจฺจยา ปสาโท.}

\prose{ตตฺถ กตโม ปสาทปจฺจยา อธิโมกฺโข, โย จิตฺตสฺส อธิโมกฺโข อธิมุจฺจนา ตทธิมุตฺตตา, อยํ วุจฺจติ ปสาทปจฺจยา อธิโมกฺโข.}

\prose{ตตฺถ กตโม อธิโมกฺขปจฺจยา ภโว, ฐเปตฺวา อธิโมกฺขํ เวทนากฺขนฺโธ สญฺญากฺขนฺโธ สงฺขารกฺขนฺโธ วิญฺญาณกฺขนฺโธ, อยํ วุจฺจติ อธิโมกฺขปจฺจยา ภโว ฯเปฯ เตน วุจฺจติ “เอวเมตสฺส เกวลสฺส ทุกฺขกฺขนฺธสฺส สมุทโย โหตี”ติ.}

\proseitem{336}{ตสฺมิํ สมเย อวิชฺชาปจฺจยา สงฺขาโร, สงฺขารปจฺจยา วิญฺญาณํ, วิญฺญาณปจฺจยา นามํ, นามปจฺจยา ผสฺโส, ผสฺสปจฺจยา เวทนา, เวทนาปจฺจยา ปสาโท, ปสาทปจฺจยา อธิโมกฺโข, อธิโมกฺขปจฺจยา ภโว, ภวปจฺจยา ชาติ, ชาติปจฺจยา ชรามรณํ, เอวเมตสฺส เกวลสฺส ทุกฺขกฺขนฺธสฺส สมุทโย โหติ.}

\proseitem{337}{ตสฺมิํ สมเย อวิชฺชาปจฺจยา สงฺขาโร, สงฺขารปจฺจยา วิญฺญาณํ, วิญฺญาณปจฺจยา นามรูปํ, นามรูปปจฺจยา ฉฏฺฐายตนํ, ฉฏฺฐายตนปจฺจยา ผสฺโส, ผสฺสปจฺจยา เวทนา, เวทนาปจฺจยา ปสาโท, ปสาทปจฺจยา อธิโมกฺโข, อธิโมกฺขปจฺจยา ภโว, ภวปจฺจยา ชาติ, ชาติปจฺจยา ชรามรณํ, เอวเมตสฺส เกวลสฺส ทุกฺขกฺขนฺธสฺส สมุทโย โหติ.}

\proseitem{338}{ตสฺมิํ สมเย อวิชฺชาปจฺจยา สงฺขาโร, สงฺขารปจฺจยา วิญฺญาณํ, วิญฺญาณปจฺจยา นามรูปํ, นามรูปปจฺจยา ฉฏฺฐายตนํ, ฉฏฺฐายตนปจฺจยา ผสฺโส, ผสฺสปจฺจยา เวทนา, เวทนาปจฺจยา ปสาโท, ปสาทปจฺจยา อธิโมกฺโข, อธิโมกฺขปจฺจยา ภโว, ภวปจฺจยา ชาติ, ชาติปจฺจยา ชรามรณํ, เอวเมตสฺส เกวลสฺส ทุกฺขกฺขนฺธสฺส สมุทโย โหติ.}

\proseitem{339}{กตเม ธมฺมา กุสลา, ยสฺมิํ สมเย กามาวจรํ กุสลํ จิตฺตํ อุปฺปนฺนํ โหติ โสมนสฺสสหคตํ ญาณสมฺปยุตฺตํ สสงฺขาเรน ฯเปฯ โสมนสฺสสหคตํ ญาณวิปฺปยุตฺตํ ฯเปฯ โสมนสฺสสหคตํ ญาณวิปฺปยุตฺตํ สสงฺขาเรน ฯเปฯ อุเปกฺขาสหคตํ ญาณสมฺปยุตฺตํ ฯเปฯ อุเปกฺขาสหคตํ}

\vspace{\tipitakaparskip}
\csromancenter{ปฏิจฺจสมุปฺปาทวิภงฺค ญาณสมฺปยุตฺตํ สสงฺขาเรน ฯเปฯ อุเปกฺขาสหคตํ ญาณวิปฺปยุตฺตํ ฯเปฯ อุเปกฺขาสหคตํ ญาณวิปฺปยุตฺตํ สสงฺขาเรน รูปารมฺมณํ วา ฯเปฯ ธมฺมรมฺมณํ วา ยํ ยํ วา ปนารพฺภ, ตสฺมิํ สมเย อวิชฺชาปจฺจยา สงฺขาโร, สงฺขารปจฺจยา วิญฺญาณํ, วิญฺญาณปจฺจยา นามํ, นามปจฺจยา ฉฏฺฐายตนํ, ฉฏฺฐายตนปจฺจยา ผสฺโส, ผสฺสปจฺจยา เวทนา, เวทนาปจฺจยา ปสาโท, ปสาทปจฺจยา อธิโมกฺโข, อธิโมกฺขปจฺจยา ภโว, ภวปจฺจยา ชาติ, ชาติปจฺจยา ชรามรณํ, เอวเมตสฺส เกวลสฺส ทุกฺขกฺขนฺธสฺส สมุทโย โหติ.}
\proseitem{340}{\csromanfolio{193}กตเม ธมฺมา กุสลา, ยสฺมิํ สมเย รูปูปปตฺติยา มคฺคํ ภาเวติ วิวิจฺเจว กาเมหิ ฯเปฯ ปฐมํ ฌานํ อุปสมฺปชฺช วิหรติ ปถวีกสิณํ, ตสฺมิํ สมเย อวิชฺชาปจฺจยา สงฺขาโร, สงฺขารปจฺจยา วิญฺญาณํ, วิญฺญาณปจฺจยา นามํ, นามปจฺจยา ฉฏฺฐายตนํ, ฉฏฺฐายตนปจฺจยา ผสฺโส, ผสฺสปจฺจยา เวทนา, เวทนาปจฺจยา ปสาโท, ปสาทปจฺจยา อธิโมกฺโข, อธิโมกฺขปจฺจยา ภโว, ภวปจฺจยา ชาติ, ชาติปจฺจยา ชรามรณํ, เอวเมตสฺส เกวลสฺส ทุกฺขกฺขนฺธสฺส สมุทโย โหติ.}

\proseitem{341}{กตเม ธมฺมา กุสลา, ยสฺมิํ สมเย อรูปูปปตฺติยา มคฺคํ ภาเวติ สพฺพโส อากิญฺจญฺญายตนํ สมติกฺกมฺม เนวสญฺญานาสญฺญายตนสญฺญาสหคตํ สุขสฺส จ ปหานา ฯเปฯ จตุตฺถํ ฌานํ อุปสมฺปชฺช วิหรติ, ตสฺมิํ สมเย อวิชฺชาปจฺจยา สงฺขาโร, สงฺขารปจฺจยา วิญฺญาณํ, วิญฺญาณปจฺจยา นามํ, นามปจฺจยา ฉฏฺฐายตนํ, ฉฏฺฐายตนปจฺจยา ผสฺโส, ผสฺสปจฺจยา เวทนา, เวทนาปจฺจยา ปสาโท ปสาทปจฺจยา อธิโมกฺโข, อธิโมกฺขปจฺจยา ภโว, ภวปจฺจยา ชาติ, ชาติปจฺจยา ชรามรณํ, เอวเมตสฺส เกวลสฺส ทุกฺขกฺขนฺธสฺส สมุทโย โหติ.}

\proseitem{342}{กตเม ธมฺมา กุสลา, ยสฺมิํ สมเย โลกุตฺตรํ ฌานํ ภาเวติ นิยฺยานิกํ อปจยคามิํ ทิฏฺฐิคตานํ ปหานาย ปฐมาย ภูมิยา ปตฺติยา วิวิจฺเจว กาเมหิ ฯเปฯ ปฐมํ ฌานํ อุปสมฺปชฺช วิหรติ ทุกฺขปฏิปทํ ทนฺธาภิญฺญํ, ตสฺมิํ สมเย อวิชฺชาปจฺจยา สงฺขาโร, สงฺขารปจฺจยา วิญฺญาณํ, วิญฺญาณปจฺจยา นามํ, นามปจฺจยา ฉฏฺฐายตนํ, ฉฏฺฐายตนปจฺจยา ผสฺโส, ผสฺสปจฺจยา เวทนา, เวทนาปจฺจยา ปสาโท, ปสาทปจฺจยา \csromanfolio{194}อธิโมกฺโข, อธิโมกฺขปจฺจยา ภโว, ภวปจฺจยา ชาติ, ชาติปจฺจยา ชรามรณํ, เอวเมเตสํ ธมฺมานํ สมุทโย โหติ.}

\csromanheader{2}{อวิชฺชามูลกกุสลนิทฺเทโส.}
\csromansectionrule
\csromanmatikaanchor{mtk.59}
\csromantocmarkref{subparagraph}{13. กุสลมูลกวิปากนิทฺเทส}{mtk.59}
\bookmark[dest={mtk.59},level=5]{13. กุสลมูลกวิปากนิทฺเทส}
\csromanheader{6}{13. กุสลมูลกวิปากนิทฺเทส}
\proseitem{343}{กตเม ธมฺมา อพฺยากตา, ยสฺมิํ สมเย กามาวจรสฺส กุสลสฺส กมฺมสฺส กตตฺตา อุปจิตตฺตา วิปากํ จกฺขุวิญฺญาณํ อุปฺปนฺนํ โหติ อุเปกฺขาสหคตํ รูปารมฺมณํ, ตสฺมิํ สมเย กุสลมูลปจฺจยา สงฺขาโร, สงฺขารปจฺจยา วิญฺญาณํ, วิญฺญาณปจฺจยา นามํ, นามปจฺจยา ฉฏฺฐายตนํ, ฉฏฺฐายตนปจฺจยา ผสฺโส, ผสฺสปจฺจยา เวทนา, เวทนาปจฺจยา ภโว, ภวปจฺจยา ชาติ, ชาติปจฺจยา ชรามรณํ, เอวเมตสฺส เกวลสฺส ทุกฺขกฺขนฺธสฺส สมุทโย โหติ.}

\proseitem{344}{ตตฺถ กตมา กุสลมูลปจฺจยา สงฺขารา, ยา เจตนา สญฺเจตนา สญฺเจตยิตตฺตํ, อยํ วุจฺจติ กุสลมูลปจฺจยา สงฺขารา ฯเปฯ เตน วุจฺจติ “เอวเมตสฺส เกวลสฺส ทุกฺขกฺขนฺธสฺส สมุทโย โหตี”ติ.}

\proseitem{345}{กตเม ธมฺมา อพฺยากตา, ยสฺมิํ สมเย กามาวจรสฺส กุสลสฺส กมฺมสฺส กตตฺตา อุปจิตตฺตา วิปากํ โสตวิญฺญาณํ อุปฺปนฺนํ โหติ อุเปกฺขาสหคตํ สทฺทารมฺมณํ ฯเปฯ ฆานวิญฺญาณํ อุปฺปนฺนํ โหติ อุเปกฺขาสหคตํ คนฺธารมฺมณํ ฯเปฯ ชิวฺหาวิญฺญาณํ อุปฺปนฺนํ โหติ อุเปกฺขาสหคตํ รสารมฺมณํ ฯเปฯ กายวิญฺญาณํ อุปฺปนฺนํ โหติ สุขสหคตํ โผฏฺฐพฺพารมฺมณํ ฯเปฯ มโนธาตุ อุปฺปนฺนา โหติ อุเปกฺขาสหคตา รูปารมฺมณา วา ฯเปฯ โผฏฺฐพฺพารมฺมณา วา ฯเปฯ มโนวิญฺญาณธาตุ อุปฺปนฺนา โหติ โสมนสฺสสหคตา รูปารมฺมณา วา ฯเปฯ ธมฺมารมฺมณา วา ฯเปฯ มโนวิญฺญาณธาตุ อุปฺปนฺนา โหติ อุเปกฺขาสหคตา รูปารมฺมณา วา ฯเปฯ ธมฺมารมฺมณา วา ยํ ยํ วา ปนารพฺภ, ตสฺมิํ สมเย กุสลมูลปจฺจยา}

\vspace{\tipitakaparskip}
\csromancenter{ปฏิจฺจสมุปฺปาทวิภงฺค สงฺขาโร, สงฺขารปจฺจยา วิญฺญาณํ, วิญฺญาณปจฺจยา นามํ, นามปจฺจยา ฉฏฺฐายตนํ, ฉฏฺฐายตนปจฺจยา ผสฺโส, ผสฺสปจฺจยา เวทนา, เวทนาปจฺจยา อธิโมกฺโข, อธิโมกฺขปจฺจยา ภโว, ภวปจฺจยา ชาติ, ชาติปจฺจยา ชรามรณํ, เอวเมตสฺส เกวลสฺส ทุกฺขกฺขนฺธสฺส สมุทโย โหติ.}
\proseitem{346}{\csromanfolio{195}กตเม ธมฺมา อพฺยากตา, ยสฺมิํ สมเย กามาวจรสฺส กุสลสฺส กมฺมสฺส กตตฺตา อุปจิตตฺตา วิปากา มโนวิญฺญาณธาตุ อุปฺปนฺนา โหติ โสมนสฺสสหคตา ญาณสมฺปยุตฺตา ฯเปฯ โสมนสฺสสหคตา ญาณสมฺปยุตฺตา สสงฺขาเรน ฯเปฯ โสมนสฺสสหคตา ญาณวิปฺปยุตฺตา ฯเปฯ โสมนสฺสสหคตา ญาณวิปฺปยุตฺตา สสงฺขาเรน ฯเปฯ.\csromanspacer{} อุเปกฺขาสหคตา ญาณสมฺปยุตฺตา ฯเปฯ อุเปกฺขาสหคตา ญาณสมฺปยุตฺตา สสงฺขาเรน ฯเปฯ อุเปกฺขาสหคตา ญาณวิปฺปยุตฺตา ฯเปฯ อุเปกฺขาสหคตา ญาณวิปฺปยุตฺตา สสงฺขาเรน รูปารมฺมณา วา ฯเปฯ ธมฺมารมฺมณา วา ยํ ยํ วา ปนารพฺภ, ตสฺมิํ สมเย กุสลมูลปจฺจยา สงฺขาโร, สงฺขารปจฺจยา วิญฺญาณํ, วิญฺญาณปจฺจยา นามํ, นามปจฺจยา ฉฏฺฐายตนํ, ฉฏฺฐายตนปจฺจยา ผสฺโส, ผสฺสปจฺจยา เวทนา, เวทนาปจฺจยา ปสาโท ปสาทปจฺจยา อธิโมกฺโข, อธิโมกฺขปจฺจยา ภโว, ภวปจฺจยา ชาติ, ชาติปจฺจยา ชรามรณํ, เอวเมตสฺส เกวลสฺส ทุกฺขกฺขนฺธสฺส สมุทโย โหติ.}

\proseitem{347}{กตเม ธมฺมา อพฺยากตา, ยสฺมิํ สมเย รูปูปปตฺติยา มคฺคํ ภาเวติ วิวิจฺเจว กาเมหิ ฯเปฯ ปฐมํ ฌานํ อุปสมฺปชฺช วิหรติ ปถวีกสิณํ, ตสฺมิํ สมเย ผสฺโส โหติ ฯเปฯ อวิกฺเขโป โหติ, อิเม ธมฺมา กุสลา.}

\prose{ตสฺเสว รูปาวจรสฺส กุสลสฺส กมฺมสฺส กตตฺตา อุปจิตตฺตา วิปากํ วิวิจฺเจว กาเมหิ ฯเปฯ ปฐมํ ฌานํ อุปสมฺปชฺช วิหรติ ปถวีกสิณํ, ตสฺมิํ สมเย กุสลมูลปจฺจยา สงฺขาโร, สงฺขารปจฺจยา วิญฺญาณํ, วิญฺญาณปจฺจยา นามํ, นามปจฺจยา ฉฏฺฐายตนํ, ฉฏฺฐายตนปจฺจยา ผสฺโส, ผสฺสปจฺจยา เวทนา, เวทนาปจฺจยา ปสาโท, ปสาทปจฺจยา อธิโมกฺโข, อธิโมกฺขปจฺจยา ภโว, ภวปจฺจยา ชาติ, ชาติปจฺจยา ชรามรณํ, เอวเมตสฺส เกวลสฺส ทุกฺขกฺขนฺธสฺส สมุทโย โหติ.}

\proseitem{348}{\csromanfolio{196}กตเม ธมฺมา อพฺยากตา, ยสฺมิํ สมเย อรูปูปปตฺติยา มคฺคํ ภาเวติ สพฺพโส อากิญฺจญฺญายตนํ สมติกฺกมฺม เนวสญฺญานาสญฺญายตนสญฺญาสหคตํ สุขสฺส จ ปหานา ฯเปฯ จตุตฺถํ ฌานํ อุปสมฺปชฺช วิหรติ, ตสฺมิํ สมเย ผสฺโส โหติ ฯเปฯ อวิกฺเขโป โหติ, อิเม ธมฺมา กุสลา.}

\prose{ตสฺเสว อรูปาวจรสฺส กุสลสฺส กมฺมสฺส กตตฺตา อุปจิตตฺตา วิปากํ สพฺพโส อากิญฺจญฺญายตนํ สมติกฺกมฺม เนวสญฺญานาสญฺญายตนสญฺญาสญฺญาสหคตํ สุขสฺส จ ปหานา ฯเปฯ จตุตฺถํ ฌานํ อุปสมฺปชฺช วิหรติ, ตสฺมิํ สมเย กุสลมูลปจฺจยา สงฺขาโร, สงฺขารปจฺจยา วิญฺญาณํ, วิญฺญาณปจฺจยา นามํ, นามปจฺจยา ฉฏฺฐายตนํ, ฉฏฺฐายตนปจฺจยา ผสฺโส, ผสฺสปจฺจยา เวทนา, เวทนาปจฺจยา ปสาโท, ปสาทปจฺจยา อธิโมกฺโข, อธิโมกฺขปจฺจยา ภโว, ภวปจฺจยา ชาติ, ชาติปจฺจยา ชรามรณํ, เอวเมตสฺส เกวลสฺส ทุกฺขกฺขนฺธสฺส สมุทโย โหติ.}

\proseitem{349}{กตเม ธมฺมา อพฺยากตา, ยสฺมิํ สมเย โลกุตฺตรํ ฌานํ ภาเวติ นิยฺยานิกํ อปจยคามิํ ทิฏฺฐิคตานํ ปหานาย ปฐมาย ภูมิยา ปตฺติยา วิวิจฺเจว กาเมหิ ฯเปฯ ปฐมํ ฌานํ อุปสมฺปชฺช วิหรติ ทุกฺขปฏิปทํ ทนฺธาภิญฺญํ, ตสฺมิํ สมเย ผสฺโส โหติ ฯเปฯ อวิกฺเขโป โหติ, อิเม ธมฺมา กุสลา.}

\prose{ตสฺเสว โลกุตฺตรสฺส กุสลสฺส ฌานสฺส กตตฺตา ภาวิตตฺตา วิปากํ วิวิจฺเจว กาเมหิ ฯเปฯ ปฐมํ ฌานํ อุปสมฺปชฺช วิหรติ ทุกฺขปฏิปทํ ทนฺธาภิญฺญํ สุญฺญตํ, ตสฺมิํ สมเย กุสลมูลปจฺจยา สงฺขาโร, สงฺขารปจฺจยา วิญฺญาณํ, วิญฺญาณปจฺจยา นามํ, นามปจฺจยา ฉฏฺฐายตนํ, ฉฏฺฐายตนปจฺจยา ผสฺโส, ผสฺสปจฺจยา เวทนา, เวทนาปจฺจยา ปสาโท, ปสาทปจฺจยา อธิโมกฺโข, อธิโมกฺขปจฺจยา ภโว, ภวปจฺจยา ชาติ, ชาติปจฺจยา ชรามรณํ, เอวเมเตสํ ธมฺมานํ สมุทโย โหติ.}

\vspace{\tipitakaparskip}
\csromancenter{กุสลมูลกวิปากนิทฺเทโส.}
\csromanmatikaanchor{mtk.60}
\csromantocmarkref{subparagraph}{14. อกุสลมูลกวิปากนิทฺเทส}{mtk.60}
\bookmark[dest={mtk.60},level=5]{14. อกุสลมูลกวิปากนิทฺเทส}
\csromanheader{6}{6. ปฏิจฺจสมุปฺปาทวิภงฺค \textbf{14. อกุสลมูลกวิปากนิทฺเทส}}
\proseitem{350}{\csromanfolio{197}กตเม ธมฺมา อพฺยากตา ยสฺมิํ สมเย อกุสลสฺส กมฺมสฺส กตตฺตา อุปจิตตฺตา วิปากํ จกฺขุวิญฺญาณํ อุปฺปนฺนํ โหติ อุเปกฺขาสหคตํ รูปารมฺมณํ, ตสฺมิํ สมเย อกุสลมูลปจฺจยา สงฺขาโร, สงฺขารปจฺจยา วิญฺญาณํ, วิญฺญาณปจฺจยา นามํ, นามปจฺจยา ฉฏฺฐายตนํ, ฉฏฺฐายตนปจฺจยา ผสฺโส, ผสฺสปจฺจยา เวทนา, เวทนาปจฺจยา ภโว, ภวปจฺจยา ชาติ, ชาติปจฺจยา ชรามรณํ, เอวเมตสฺส เกวลสฺส ทุกฺขกฺขนฺธสฺส สมุทโย โหติ.}

\proseitem{351}{ตตฺถ กตโม อกุสลมูลปจฺจยา สงฺขาโร, ยา เจตนา สญฺเจตนา สญฺเจตยิตตฺตํ, อยํ วุจฺจติ อกุสลมูลปจฺจยา สงฺขาโร ฯเปฯ เตน วุจฺจติ “เอวเมตสฺส เกวลสฺส ทุกฺขกฺขนฺธสฺส สมุทโย โหตี”ติ.}

\proseitem{352}{กตเม ธมฺมา อพฺยากตา, ยสฺมิํ สมเย อกุสลสฺส กมฺมสฺส กตตฺตา อุปจิตตฺตา วิปากํ โสตวิญฺญาณํ อุปฺปนฺนํ โหติ อุเปกฺขาสหคตํ สทฺทารมฺมณํ ฯเปฯ ฆานวิญฺญาณํ อุปฺปนฺนํ โหติ อุเปกฺขาสหคตํ คนฺธารมฺมณํ ฯเปฯ ชิวฺหาวิญฺญาณํ อุปฺปนฺนํ โหติ อุเปกฺขาสหคตํ รสารมฺมณํ ฯเปฯ กายวิญฺญาณํ อุปฺปนฺนํ โหติ ทุกฺขสหคตํ โผฏฺฐพฺพารมฺมณํ ฯเปฯ มโนธาตุ อุปฺปนฺนา โหติ อุเปกฺขาสหคตา รูปารมฺมณา วา ฯเปฯ โผฏฺฐพฺพารมฺมณา วา ยํ ยํ วา ปนารพฺภ, ตสฺมิํ สมเย อกุสลมูลปจฺจยา สงฺขาโร, สงฺขารปจฺจยา วิญฺญาณํ, วิญฺญาณปจฺจยา นามํ, นามปจฺจยา ฉฏฺฐายตนํ, ฉฏฺฐายตนปจฺจยา ผสฺโส, ผสฺสปจฺจยา เวทนา, เวทนาปจฺจยา อธิโมกฺโข, อธิโมกฺขปจฺจยา ภโว, ภวปจฺจยา ชาติ, ชาติปจฺจยา ชรามรณํ, เอวเมตสฺส เกวลสฺส ทุกฺขกฺขนฺธสฺส สมุทโย โหติ.}

\proseitem{353}{กตเม ธมฺมา อพฺยากตา, ยสฺมิํ สมเย อกุสลสฺส กมฺมสฺส กตตฺตา อุปจิตตฺตา วิปากา มโนวิญฺญาณธาตุ อุปฺปนฺนา โหติ อุเปกฺขาสหคตา รูปารมฺมณา วา ฯเปฯ ธมฺมารมฺมณา วา ยํ ยํ วา ปนารพฺภ, ตสฺมิํ สมเย อกุสลมูลปจฺจยา สงฺขาโร, สงฺขารปจฺจยา วิญฺญาณํ, วิญฺญาณปจฺจยา นามํ, นามปจฺจยา ฉฏฺฐายตนํ, ฉฏฺฐายตนปจฺจยา \csromanfolio{198}ผสฺโส, ผสฺสปจฺจยา เวทนา, เวทนาปจฺจยา อธิโมกฺโข, อธิโมกฺขปจฺจยา ภโว, ภวปจฺจยา ชาติ, ชาติปจฺจยา ชรามรณํ, เอวเมตสฺส เกวลสฺส ทุกฺขกฺขนฺธสฺส สมุทโย โหติ.}

\proseitem{354}{ตตฺถ กตโม อกุสลมูลปจฺจยา สงฺขาโร, ยา เจตนา สญฺเจตนา สญฺเจตยิตตฺตํ, อยํ วุจฺจติ อกุสลมูลปจฺจยา สงฺขาโร.}

\prose{ตตฺถ กตมํ สงฺขารปจฺจยา วิญฺญาณํ, ยํ จิตฺตํ มโน มานสํ ฯเปฯ ตชฺชามโนวิญฺญาณธาตุ, อิทํ วุจฺจติ สงฺขารปจฺจยา วิญฺญาณํ.}

\prose{ตตฺถ กตมํ วิญฺญาณปจฺจยา นามํ, เวทนากฺขนฺโธ สญฺญากฺขนฺโธ สงฺขารกฺขนฺโธ, อิทํ วุจฺจติ วิญฺญาณปจฺจยา นามํ.}

\prose{ตตฺถ กตมํ นามปจฺจยา ฉฏฺฐายตนํ, ยํ จิตฺตํ มโน มานสํ ฯเปฯ ตชฺชามโนวิญฺญาณธาตุ, อิทํ วุจฺจติ นามปจฺจยา ฉฏฺฐายตนํ.}

\prose{ตตฺถ กตโม ฉฏฺฐายตนปจฺจยา ผสฺโส, โย ผสฺโส ผุสนา สมฺผุสนา สมฺผุสิตตฺตํ, อยํ วุจฺจติ ฉฏฺฐายตนปจฺจยา ผสฺโส.}

\prose{ตตฺถ กตมา ผสฺสปจฺจยา เวทนา, ยํ เจตสิกํ เนว สาตํ นาสาตํ เจโตสมฺผสฺสชํ อทุกฺขมสุขํ เวทยิตํ เจโตสมฺผสฺสชา อทุกฺขมสุขา เวทนา, อยํ วุจฺจติ ผสฺสปจฺจยา เวทนา.}

\prose{ตตฺถ กตโม เวทนาปจฺจยา อธิโมกฺโข, โย จิตฺตสฺส อธิโมกฺโข อธิมุจฺจนา ตทธิมุตฺตตา, อยํ วุจฺจติ เวทนาปจฺจยา อธิโมกฺโข.}

\prose{ตตฺถ กตโม อธิโมกฺขปจฺจยา ภโว, ฐเปตฺวา อธิโมกฺขํ เวทนากฺขนฺโธ สญฺญากฺขนฺโธ สงฺขารกฺขนฺโธ วิญฺญาณกฺขนฺโธ, อยํ วุจฺจติ อธิโมกฺขปจฺจยา ภโว.}

\prose{ตตฺถ กตมา ภวปจฺจยา ชาติ, ยา เตสํ เตสํ ธมฺมานํ ชาติ สญฺชาติ นิพฺพตฺติ อภินิพฺพตฺติ ปาตุภาโว, อยํ วุจฺจติ ภวปจฺจยา ชาติ.}

\prose{ตตฺถ กตมํ ชาติปจฺจยา ชรามรณํ, อตฺถิ ชรา, อตฺถิ มรณํ.\csromanspacer{} ตตฺถ กตมา ชรา, ยา เตสํ เตสํ ธมฺมานํ ชรา ชีรณตา อายุโน}

\vspace{\tipitakaparskip}
\csromancenter{ปฏิจฺจสมุปฺปาทวิภงฺค สํหานิ, อยํ วุจฺจติ ชรา.\csromanspacer{} ตตฺถ กตมํ มรณํ, โย เตสํ เตสํ ธมฺมานํ ขโย วโย เภโท ปริเภโท อนิจฺจตา อนฺตรธานํ, อิทํ วุจฺจติ มรณํ.\csromanspacer{} อิติ อยญฺจ ชรา อิทญฺจ มรณํ, อิทํ วุจฺจติ ชาติปจฺจยา ชรามรณํ.}
\prose{\csromanfolio{199}เอวเมตสฺส เกวลสฺส ทุกฺขกฺขนฺธสฺส สมุทโย โหตีติ เอวเมตสฺส เกวลสฺส ทุกฺขกฺขนฺธสฺส สงฺคติ โหติ, สมาคโม โหติ, สโมธานํ โหติ, ปาตุภาโว โหติ, เตน วุจฺจติ “เอวเมตสฺส เกวลสฺส ทุกฺขกฺขนฺธสฺส สมุทโย โหตี”ติ.}

\vspace{\tipitakaparskip}
\csromancenter{อกุสลมูลกวิปากนิทฺเทโส.}
\csromancenter{อภิธมฺมภาชนียํ.}
\nitthitam{\textbf{ปฏิจฺจสมุปฺปาทวิภงฺโค}1 \textbf{นิฏฺฐิโต.}}
\csromanmatikaanchor{mtk.61}
\csromantocmarknopage{subparagraph}{7. สติปฏฺฐานวิภงฺค}{mtk.61}
\bookmark[dest={mtk.61},level=5]{7. สติปฏฺฐานวิภงฺค}
\csromanheader{6}{7. สติปฏฺฐานวิภงฺค}
\csromanheader{2}{1. สุตฺตนฺตภาชนีย}
\proseitem{355}{\csromanfolio{200}จตฺตาโร สติปฏฺฐานา\csromandash{}อิธ ภิกฺขุ อชฺฌตฺตํ กาเย กายานุปสฺสี วิหรติ, พหิทฺธา กาเย กายานุปสฺสี วิหรติ, อชฺฌตฺตพหิทฺธา กาเย กายานุปสฺสี วิหรติ อาตาปี สมฺปชาโน สติมา วิเนยฺย โลเก อภิชฺฌาโทมนสฺสํ.\csromanspacer{} อชฺฌตฺตํ เวทนาสุ เวทนานุปสฺสี วิหรติ, พหิทฺธา เวทนาสุ เวทนานุปสฺสี วิหรติ, อชฺฌตฺตพหิทฺธา เวทนาสุ เวทนานุปสฺสี วิหรติ อาตาปี สมฺปชาโน สติมา วิเนยฺย โลเก อภิชฺฌาโทมนสฺสํ.\csromanspacer{} อชฺฌตฺตํ จิตฺเต จิตฺตานุปสฺสี วิหรติ, พหิทฺธา จิตฺเต จิตฺตานุปสฺสี วิหรติ, อชฺฌตฺตพหิทฺธา จิตฺเต จิตฺตานุปสฺสี วิหรติ อาตาปี สมฺปชาโน สติมา วิเนยฺย โลเก อภิชฺฌาโทมนสฺสํ.\csromanspacer{} อชฺฌตฺตํ ธมฺเมสุ ธมฺมานุปสฺสี วิหรติ, พหิทฺธา ธมฺเมสุ ธมฺมานุปสฺสี วิหรติ, อชฺฌตฺตพหิทฺธา ธมฺเมสุ ธมฺมานุปสฺสี วิหรติ อาตาปี สมฺปชาโน สติมา วิเนยฺย โลเก อภิชฺฌาโทมนสฺสํ.}

\csromanmatikaanchor{mtk.62}
\csromanmatikaanchor{mtk.63}
\csromantocmarknopage{subparagraph}{1. สุตฺตนฺตภาชนีย}{mtk.62}
\bookmark[dest={mtk.62},level=5]{1. สุตฺตนฺตภาชนีย}
\csromantocmarkref{subparagraph}{1. กายานุปสฺสนานิทฺเทส}{mtk.63}
\bookmark[dest={mtk.63},level=5]{1. กายานุปสฺสนานิทฺเทส}
\csromanheader{6}{1. กายานุปสฺสนานิทฺเทส}
\proseitem{356}{กถญฺจ ภิกฺขุ อชฺฌตฺตํ กาเย กายานุปสฺสี วิหรติ, อิธ ภิกฺขุ อชฺฌตฺตํ กายํ อุทฺธํ ปาทตลา อโธ เกสมตฺถกา ตจปริยนฺตํ ปูรํ นานปฺปการสฺส อสุจิโน ปจฺจเวกฺขติ “อตฺถิ อิมสฺมิํ กาเย เกสา โลมา นขา ทนฺตา ตโจ, มํสํ นฺหารุ\footnote{นหารุ (สี)} อฏฺฐิ อฏฺฐิมิญฺชํ\footnote{อฏฺฐิมิญฺชา (สี)} วกฺกํ, หทยํ ยกนํ กิโลมกํ ปิหกํ ปปฺผาสํ, อนฺตํ อนฺตคุณํ อุทริยํ กรีสํ, ปิตฺตํ เสมฺหํ ปุพฺโพ โลหิตํ เสโท เมโท, อสฺสุ วสา เขโฬ สิงฺฆาณิกา ลสิกา มุตฺตนฺ”ติ.\csromanspacer{} โส ตํ นิมิตฺตํ อาเสวติ ภาเวติ พหุลีกโรติ, สฺวาวตฺถิตํ ววตฺถเปติ, โส ตํ นิมิตฺตํ อาเสวิตฺวา ภาเวตฺวา พหุลีกริตฺวา สฺวาวตฺถิตํ ววตฺถเปตฺวา พหิทฺธา กาเย จิตฺตํ อุปสํหรติ.}

\prose{กถญฺจ ภิกฺขุ พหิทฺธา กาเย กายานุปสฺสี วิหรติ, อิธ ภิกฺขุ พหิทฺธา กายํ อุทฺธํ ปาทตลา อโธ เกสมตฺถกา ตจปริยนฺตํ ปูรํ นานปฺปการสฺส อสุจิโน ปจฺจเวกฺขติ “อตฺถิสฺส กาเย เกสา}

\vspace{\tipitakaparskip}
\csromancenter{สติปฏฺฐานวิภงฺค โลมา นขา ทนฺตา ตโจ, มํสํ นฺหารุ อฏฺฐิ อฏฺฐิมิญฺชํ วกฺกํ, หทยํ ยกนํ กิโลมกํ ปิหกํ ปปฺผาสํ, อนฺตํ อนฺตคุณํ อุทริยํ กรีสํ, ปิตฺตํ เสมฺหํ ปุพฺโพ โลหิตํ เสโท เมโท, อสฺสุ วสา เขโฬ สิงฺฆาณิกา ลสิกา มุตฺตนฺ”ติ.\csromanspacer{} โส ตํ นิมิตฺตํ อาเสวติ ภาเวติ พหุลีกโรติ, สฺวาวตฺถิตํ ววตฺถเปติ, โส ตํ นิมิตฺตํ อาเสวิตฺวา ภาเวตฺวา พหุลีกริตฺวา สฺวาวตฺถิตํ ววตฺถเปตฺวา อชฺฌตฺตพหิทฺธา กาเย จิตฺตํ อุปสํหรติ.}
\prose{\csromanfolio{201}กถญฺจ ภิกฺขุ อชฺฌตฺตพหิทฺธา กาเย กายานุปสฺสี วิหรติ, อิธ ภิกฺขุ อชฺฌตฺตพหิทฺธา กายํ อุทฺธํ ปาทตลา อโธ เกสมตฺถกา ตจปริยนฺตํ ปูรํ นานปฺปการสฺส อสุจิโน ปจฺจเวกฺขติ “อตฺถิ อิมสฺมิํ กาเย เกสา โลมา นขา ทนฺตา ตโจ, มํสํ นฺหารุ อฏฺฐิ อฏฺฐิมิญฺชํ วกฺกํ, หทยํ ยกนํ กิโลมกํ ปิหกํ ปปฺผาสํ, อนฺตํ อนฺตคุณํ อุทริยํ กรีสํ, ปิตฺตํ เสมฺหํ ปุพฺโพ โลหิตํ เสโท เมโท, อสฺสุ วสา เขโฬ สิงฺฆาณิกา ลสิกา มุตฺตนฺ”ติ.\csromanspacer{} เอวํ ภิกฺขุ อชฺฌตฺตพหิทฺธา กาเย กายานุปสฺสี วิหรติ อาตาปี สมฺปชาโน สติมา วิเนยฺย โลเก อภิชฺฌาโทมนสฺสํ.}

\proseitem{357}{“อนุปสฺสี”ติ ตตฺถ กตมา อนุปสฺสนา, ยา ปญฺญา ปชานนา ฯเปฯ อโมโห ธมฺมวิจโย สมฺมาทิฏฺฐิ, อยํ วุจฺจติ อนุปสฺสนา.\csromanspacer{} อิมาย อนุปสฺสนาย อุเปโต โหติ สมุเปโต อุปาคโต สมุปาคโต อุปปนฺโน สมฺปนฺโน สมนฺนาคโต, เตน วุจฺจติ “อนุปสฺสี”ติ.}

\proseitem{358}{“วีหรตี”ติ อิริยติ วตฺตติ ปาเลติ ยเปติ ยาเปติ จรติ วิหรติ, เตน วุจฺจติ “วิหรตี”ติ.}

\proseitem{359}{“อาตาปี”ติ ตตฺถ กตโม อาตาโป\footnote{กตมํ อาตาปํ (สพฺพตฺถ)}, โย เจตสิโก วีริยารมฺโภ ฯเปฯ สมฺมาวายาโม, อยํ วุจฺจติ อาตาโป.\csromanspacer{} อิมินา อาตาเปน อุเปโต โหติ สมุเปโต อุปาคโต สมุปาคโต อุปปนฺโน สมฺปนฺโน สมนฺนาคโต, เตน วุจฺจติ “อาตาปี”ติ.}

\proseitem{360}{\csromanfolio{202}“สมฺปชาโน”ติ ตตฺถ กตมํ สมฺปชญฺญํ, ยา ปญฺญา ปชานนา ฯเปฯ อโมโห ธมฺมวิจโย สมฺมาทิฏฺฐิ, อิทํ วุจฺจติ สมฺปชญฺญํ.\csromanspacer{} อิมินา สมฺปชญฺเญน อุเปโต โหติ สมุเปโต อุปาคโต สมุปาคโต อุปปนฺโน สมฺปนฺโน สมนฺนาคโต, เตน วุจฺจติ “สมฺปชาโน”ติ.}

\proseitem{361}{“สติมา”ติ ตตฺถ กตมา สติ, ยา สติ อนุสฺสติ ฯเปฯ สมฺมาสติ, อยํ วุจฺจติ สติ.\csromanspacer{} อิมาย สติยา อุเปโต โหติ สมุเปโต อุปาคโต สมุปาคโต อุปปนฺโน สมฺปนฺโน สมนฺนาคโต, เตน วุจฺจติ “สติมา”ติ.}

\proseitem{362}{“วิเนยฺย โลเก อภิชฺฌาโทมนสฺสนฺ”ติ ตตฺถ กตโม โลโก, เสฺวว กาโย โลโก, ปญฺจปิ อุปาทานกฺขนฺธา โลโก, อยํ วุจฺจติ โลโก.\csromanspacer{} ตตฺถ กตมา อภิชฺฌา, โย ราโค สาราโค ฯเปฯ จิตฺตสฺส สาราโค, อยํ วุจฺจติ อภิชฺฌา.\csromanspacer{} ตตฺถ กตมํ โทมนสฺสํ, ยํ เจตสิกํ อสาตํ เจตสิกํ ทุกฺขํ เจโตสมฺปสฺสชํ อสาตํ ทุกฺขํ เวทยิตํ เจโตสมฺผสฺสชา อสาตา ทุกฺขา เวทนา, อิทํ วุจฺจติ โทมนสฺสํ.\csromanspacer{} อิติ อยญฺจ อภิชฺฌา อิทญฺจ โทมนสฺสํ อิมมฺหิ โลเก วินีตา โหนฺติ ปฏิวินีตา สนฺตา สมิตา วูปสนฺตา อตฺถงฺคตา อพฺภตฺถงฺคตา อปฺปิตา พฺยปฺปิตา โสสิตา วิโสสิตา พฺยนฺตีกตา, เตน วุจฺจติ “เวเนยฺย โลเก อภิชฺฌาโทมนสฺสนฺ”ติ.}

\csromanheader{2}{กายานุปสฺสนานิทฺเทโส.}
\csromansectionrule
\csromanmatikaanchor{mtk.64}
\csromantocmarkref{subparagraph}{2. เวทนานุปสฺสนานิทฺเทส}{mtk.64}
\bookmark[dest={mtk.64},level=5]{2. เวทนานุปสฺสนานิทฺเทส}
\csromanheader{6}{2. เวทนานุปสฺสนานิทฺเทส}
\proseitem{363}{กถญฺจ ภิกฺขุ อชฺฌตฺตํ เวทนาสุ เวทนานุปสฺสี วิหรติ, อิธ ภิกฺขุ สุขํ เวทนํ เวทยมาโน “สุขํ เวทนํ เวทยามี”ติ ปชานาติ, ทุกฺขํ เวทนํ เวทยมาโน “ทุกฺขํ เวทนํ เวทยามี”ติ ปชานาติ, อทุกฺขมสุขํ เวทนํ เวทยมาโน “อทุกฺขมสุขํ เวทนํ เวทยามี”ติ ปชานาติ.\csromanspacer{} สามิสํ วา สุขํ เวทนํ เวทยมาโน “สามิสํ สุขํ เวทนํ เวทยามี”ติ ปชานาติ, นิรามิสํ วา สุขํ เวทนํ เวทยมาโน “นิรามิสํ สุขํ เวทนํ เวทยามี”ติ ปชานาติ.\csromanspacer{} สามิสํ วา ทุกฺขํ เวทนํ เวทยมาโน “สามิสํ}

\vspace{\tipitakaparskip}
\csromancenter{สติปฏฺฐานวิภงฺค ทุกฺขํ เวทนํ เวทยามี”ติ ปชานาติ, นิรามิสํ วา ทุกฺขํ เวทนํ เวทยมาโน “นิรามิสํ ทุกฺขํ เวทนํ เวทยามี”ติ ปชานาติ.\csromanspacer{} สามิสํ วา อทุกฺขมสุขํ เวทนํ เวทยมาโน “สามิสํ อทุกฺขมสุขํ เวทนํ เวทยามี”ติ ปชานาติ, นิรามิสํ วา อทุกฺขมสุขํ เวทนํ เวทยมาโน “นิรามิสํ อทุกฺขมสุขํ เวทนํ เวทยามี”ติ ปชานาติ.\csromanspacer{} โส ตํ นิมิตฺตํ อาเสวติ ภาเวติ พหุลีกโรติ, สฺวาวตฺถิตํ ววตฺถเปติ, โส ตํ นิมิตฺตํ อาเสวิตฺวา ภาเวตฺวา พหุลีกริตฺวา สฺวาวตฺถิตํ ววตฺถเปตฺวา พหิทฺธา เวทนาสุ จิตฺตํ อุปสํหรติ.}
\prose{\csromanfolio{203}กถญฺจ ภิกฺขุ พหิทฺธา เวทนาสุ เวทนานุปสฺสี วิหรติ, อิธ ภิกฺขุ สุขํ เวทนํ เวทยมานํ “สุขํ เวทนํ เวทยตี”ติ ปชานาติ, ทุกฺขํ เวทนํ เวทยมานํ “ทุกฺขํ เวทนํ เวทยตี”ติ ปชานาติ, อทุกฺขมสุขํ เวทนํ เวทยมานํ “อทุกฺขมสุขํ เวทนํ เวทยตี”ติ ปชานาติ.\csromanspacer{} สามิสํ วา สุขํ เวทนํ เวทยมานํ “สามิสํ สุขํ เวทนํ เวทยตี”ติ ปชานาติ, นิรามิสํ วา สุขํ เวทนํ เวทยมานํ “นิรามิสํ สุขํ เวทนํ เวทยตี”ติ ปชานาติ.\csromanspacer{} สามิสํ วา ทุกฺขํ เวทนํ เวทยมานํ “สามิสํ ทุกฺขํ เวทนํ เวทยตี”ติ ปชานาติ, นิรามิสํ วา ทุกฺขํ เวทนํ เวทยมานํ “นิรามิสํ ทุกฺขํ เวทนํ เวทยตี”ติ ปชานาติ.\csromanspacer{} สามิสํ วา อทุกฺขมสุขํ เวทนํ เวทยมานํ “สามิสํ อทุกฺขมสุขํ เวทนํ เวทยตี”ติ ปชานาติ, นิรามิสํ วา อทุกฺขมสุขํ เวทนํ เวทยมานํ “นิรามิสํ อทุกฺขมสุขํ เวทนํ เวทยตี”ติ ปชานาติ.\csromanspacer{} โส ตํ นิมิตฺตํ อาเสวติ ภาเวติ พหุลีกโรติ, สฺวาวตฺถิตํ ววตฺถเปติ, โส ตํ นิมิตฺตํ อาเสวิตฺวา ภาเวตฺวา พหุลีกริตฺวา สฺวาวตฺถิตํ ววตฺถเปตฺวา อชฺฌตฺตพหิทฺธา เวทนาสุ จิตฺตํ อุปสํหรติ.}

\prose{กถญฺจ ภิกฺขุ อชฺฌตฺตพหิทฺธา เวทนาสุ เวทนานุปสฺสี วิหรติ, อิธ ภิกฺขุ สุขํ เวทนํ “สุขา เวทนา”ติ ปชานาติ, ทุกฺขํ เวทนํ “ทุกฺขา เวทนา”ติ ปชานาติ, อทุกฺขมสุขํ เวทนํ “อทุกฺขมสุขา เวทนา”ติ ปชานาติ.\csromanspacer{} สามิสํ วา สุขํ เวทนํ “สามิสา สุขา เวทนา”ติ ปชานาติ, นิรามิสํ วา สุขํ เวทนํ “นิรามิสา สุขา เวทนา”ติ ปชานาติ, สามิสํ วา ทุกฺขํ เวทนํ “สามิสา ทุกฺขา เวทนา”ติ ปชานาติ, นิรามิสํ วา ทุกฺขํ เวทนํ “นิรามิสา ทุกฺขา เวทนา”ติ ปชานาติ, สามิสํ วา อทุกฺขมสุขํ เวทนํ “สามิสา อทุกฺขมสุขา เวทนา”ติ ปชานาติ, นิรามิสํ วา อทุกฺขมสุขํ \csromanfolio{204}เวทนํ “นิรามิสา อทุกฺขมสุขา เวทนา”ติ ปชานาติ.\csromanspacer{} เอวํ ภิกฺขุ อชฺฌตฺตพหิทฺธา เวทนาสุ เวทนานุปสฺสี วิหรติ อาตาปี สมฺปชาโน สติมา วิเนยฺย โลเก อภิชฺฌาโทมนสฺสํ.}

\proseitem{364}{“อนุปสฺสี”ติ ฯเปฯ. “วิหรตี”ติ ฯเปฯ. “อาตาปี”ติ ฯเปฯ. “สมฺปชาโน”ติ ฯเปฯ. “สติมา”ติ ฯเปฯ. “วิเนยฺย โลเก อภิชฺฌาโทมนสฺสนฺ”ติ ตตฺถ กตโม โลโก, สาเยว เวทนา โลโก, ปญฺจปิ อุปาทานกฺขนฺธา โลโก, อยํ วุจฺจติ โลโก.\csromanspacer{} ตตฺถ กตมา อภิชฺฌา, โย ราโค สาราโค ฯเปฯ จิตฺตสฺส สาราโค, อยํ วุจฺจติ อภิชฺฌา.\csromanspacer{} ตตฺถ กตมํ โทมนสฺสํ, ยํ เจตสิกํ อสาตํ เจตสิกํ ทุกฺขํ เจโตสมฺผสฺสชํ อสาตํ ทุกฺขํ เวทยิตํ เจโตสมฺผสฺสชา อสาตา ทุกฺขา เวทนา, อิทํ วุจฺจติ โทมนสฺสํ.\csromanspacer{} อิติ อยญฺจ อภิชฺฌา อิทญฺจ โทมนสฺสํ อิมมฺหิ โลเก วินีตา โหนฺติ ปฏิวินีตา สนฺตา สมิตา วูปสนฺตา อตฺถงฺคตา อพฺภตฺถงฺคตา อปฺปิตา พฺยปฺปิตา โสสิตา วิโสสิตา พฺยนฺตีกตา, เตน วุจฺจติ “วิเนยฺย โลเก อภิชฺฌาโทมนสฺสนฺ”ติ.}

\csromanheader{2}{เวทนานุปสฺสนานิทฺเทโส.}
\csromansectionrule
\csromanmatikaanchor{mtk.65}
\csromantocmarkref{subparagraph}{3. จิตฺตานุปสฺสนานิทฺเทส}{mtk.65}
\bookmark[dest={mtk.65},level=5]{3. จิตฺตานุปสฺสนานิทฺเทส}
\csromanheader{6}{3. จิตฺตานุปสฺสนานิทฺเทส}
\proseitem{365}{กถญฺจ ภิกฺขุ อชฺฌตฺตํ จิตฺเต จิตฺตานุปสฺสี วิหรติ, อิธ ภิกฺขุ สราคํ วา จิตฺตํ “สราคํ เม จิตฺตนฺ”ติ ปชานาติ, วีตราคํ วา จิตฺตํ “วีตราคํ เม จิตฺตนฺ”ติ ปชานาติ.\csromanspacer{} สโทสํ วา จิตฺตํ “สโทสํ เม จิตฺตนฺ”ติ ปชานาติ, วีตโทสํ วา จิตฺตํ “วีตโทสํ เม จิตฺตนฺ”ติ ปชานาติ.\csromanspacer{} สโมหํ วา จิตฺตํ “สโมหํ เม จิตฺตนฺ”ติ ปชานาติ, วีตโมหํ วา จิตฺตํ “วีตโมหํ เม จิตฺตนฺ”ติ ปชานาติ.\csromanspacer{} สํขิตฺตํ วา จิตฺตํ “สํขิตฺตํ เม จิตฺตนฺ”ติ ปชานาติ, วิกฺขิตฺตํ วา จิตฺตํ “วิกฺขิตฺตํ เม จิตฺตนฺ”ติ ปชานาติ.\csromanspacer{} มหคฺคตํ วา จิตฺตํ “มหคฺคตํ เม จิตฺตนฺ”ติ ปชานาติ, อมหคฺคตํ วา จิตฺตํ “อมหคฺคตํ เม จิตฺตนฺ”ติ ปชานาติ.\csromanspacer{} ส-อุตฺตรํ วา จิตฺตํ “ส- อุตฺตรํ เม จิตฺตนฺ”ติ ปชานาติ, อนุตฺตรํ วา จิตฺตํ “อนุตฺตรํ เม จิตฺตนฺ”ติ ปชานาติ.\csromanspacer{} สมาหิตํ วา จิตฺตํ “สมาหิตํ เม จิตฺตนฺ”ติ ปชานาติ,}

\vspace{\tipitakaparskip}
\csromancenter{สติปฏฺฐานวิภงฺค อสมาหิตํ วา จิตฺตํ “อสมาหิตํ เม จิตฺตนฺ”ติ ปชานาติ.\csromanspacer{} วิมุตฺตํ วา จิตฺตํ “วิมุตฺตํ เม จิตฺตนฺ”ติ ปชานาติ, อวิมุตฺตํ วา จิตฺตํ “อวิมุตฺตํ เม จิตฺตนฺ”ติ ปชานาติ.\csromanspacer{} โส ตํ นิมิตฺตํ อาเสวติ ภาเวติ พหุลีกโรติ, สฺวาวตฺถิตํ ววตฺถเปติ, โส ตํ นิมิตฺตํ อาเสวิตฺวา ภาเวตฺวา พหุลีกริตฺวา สฺวาวตฺถิตํ ววตฺถเปตฺวา พหิทฺธา จิตฺเต จิตฺตํ อุปสํหรติ.}
\prose{\csromanfolio{205}กถญฺจ ภิกฺขุ พหิทฺธา จิตฺเต จิตฺตานุปสฺสี วิหรติ, อิธ ภิกฺขุ สราคํ วาสฺส จิตฺตํ “สราคมสฺส จิตฺตนฺ”ติ ปชานาติ, วีตราคํ วาสฺส จิตฺตํ “วีตราคมสฺส จิตฺตนฺ”ติ ปชานาติ.\csromanspacer{} สโทสํ วาสฺส จิตฺตํ “สโทสมสฺส จิตฺตนฺ”ติ ปชานาติ, วีตโทสํ วาสฺส จิตฺตํ “วีตโทสมสฺส จิตฺตนฺ”ติ ปชานาติ.\csromanspacer{} สโมหํ วาสฺส จิตฺตํ “สโมหมสฺส จิตฺตนฺ”ติ ปชานาติ, วีตโมหํ วาสฺส จิตฺตํ “วีตโมหมสฺส จิตฺตนฺ”ติ ปชานาติ.\csromanspacer{} สํขิตฺตํ วาสฺส จิตฺตํ “สํขิตฺตมสฺส จิตฺตนฺ”ติ ปชานาติ, วิกฺขิตฺตํ วาสฺส จิตฺตํ “วิกฺขิตฺตมสฺส จิตฺตนฺ”ติ ปชานาติ.\csromanspacer{} มหคฺคตํ วาสฺส จิตฺตํ “มหคฺคตมสฺส จิตฺตนฺ”ติ ปชานาติ, อมหคฺคตํ วาสฺส จิตฺตํ “อมหคฺคตมสฺส จิตฺตนฺ”ติ ปชานาติ.\csromanspacer{} ส-อุตฺตรํ วาสฺส จิตฺตํ “ส-อุตฺตรมสฺส จิตฺตนฺ”ติ ปชานาติ, อนุตฺตรํ วาสฺส จิตฺตํ “อนุตฺตรมสฺส จิตฺตนฺ”ติ ปชานาติ.\csromanspacer{} สมาหิตํ วาสฺส จิตฺตํ “สมาหิตมสฺส จิตฺตนฺ”ติ ปชานาติ, อสมาหิตํ วาสฺส จิตฺตํ “อสมาหิตมสฺส จิตฺตนฺ”ติ ปชานาติ.\csromanspacer{} วิมุตฺตํ วาสฺส จิตฺตํ “วิมุตฺตมสฺส จิตฺตนฺ”ติ ปชานาติ, อวิมุตฺตํ วาสฺส จิตฺตํ “อวิมุตฺตมสฺส จิตฺตนฺ”ติ ปชานาติ.\csromanspacer{} โส ตํ นิมิตฺตํ อาเสวติ ภาเวติ พหุลีกโรติ, สฺวาวตฺถิตํ ววตฺถเปติ, โส ตํ นิมิตฺตํ อาเสวิตฺวา ภาเวตฺวา พหุลีกริตฺวา สฺวาวตฺถิตํ ววตฺถเปตฺวา อชฺฌตฺตพหิทฺธา จิตฺเต จิตฺตํ อุปสํหรติ.}

\prose{กถญฺจ ภิกฺขุ อชฺฌตฺตพหิทฺธา จิตฺเต จิตฺตานุปสฺสี วิหรติ, อิธ ภิกฺขุ สราคํ วา จิตฺตํ “สราคํ จิตฺตนฺ”ติ ปชานาติ, วีตราคํ วา จิตฺตํ “วีตราคํ จิตฺตนฺ”ติ ปชานาติ.\csromanspacer{} สโทสํ วา จิตฺตํ “สโทสํ จิตฺตนฺ”ติ ปชานาติ, วีตโทสํ วา จิตฺตํ “วีตโทสํ จิตฺตนฺ”ติ ปชานาติ.\csromanspacer{} สโมหํ วา จิตฺตํ “สโมหํ จิตฺตนฺ”ติ ปชานาติ, วีตโมหํ วา จิตฺตํ “วีตโมหํ จิตฺตนฺ”ติ ปชานาติ.\csromanspacer{} สํขิตฺตํ วา จิตฺตํ “สํขิตฺตํ จิตฺตนฺ”ติ ปชานาติ, วิกฺขิตฺตํ วา จิตฺตํ “วิกฺขิตฺตํ จิตฺตนฺ”ติ ปชานาติ.\csromanspacer{} มหคฺคตํ วา จิตฺตํ “มหคฺคตํ จิตฺตนฺ”ติ ปชานาติ, อมหคฺคตํ วา จิตฺตํ “อมหคฺคตํ จิตฺตนฺ”ติ ปชานาติ.\csromanspacer{} ส-อุตฺตรํ วา จิตฺตํ “ส-อุตฺตรํ จิตฺตนฺ”ติ ปชานาติ, อนุตฺตรํ วา จิตฺตํ “อนุตฺตรํ จิตฺตนฺ”ติ \csromanfolio{206}ปชานาติ.\csromanspacer{} สมาหิตํ วา จิตฺตํ “สมาหิตํ จิตฺตนฺ”ติ ปชานาติ, อสมาหิตํ วา จิตฺตํ “อสมาหิตํ จิตฺตนฺ”ติ ปชานาติ.\csromanspacer{} วิมุตฺตํ วา จิตฺตํ “วิมุตฺตํ จิตฺตนฺ”ติ ปชานาติ, อวิมุตฺตํ วา จิตฺตํ “อวิมุตฺตํ จิตฺตนฺ”ติ ปชานาติ.\csromanspacer{} เอวํ ภิกฺขุ อชฺฌตฺตพหิทฺธา จิตฺเต จิตฺตานุปสฺสี วิหรติ อาตาปี สมฺปชาโน สติมา วิเนยฺย โลเก อภิชฺฌาโทมนสฺสํ.}

\proseitem{366}{“อนุปสฺสี”ติ ฯเปฯ. “วิหรตี”ติ ฯเปฯ. “อาตาปี”ติ ฯเปฯ. “สมฺปชาโน”ติ ฯเปฯ. “สติมา”ติ ฯเปฯ. “วิเนยฺย โลเก อภิชฺฌาโทมนสฺสนฺ”ติ ตตฺถ กตโม โลโก, ตํเยว จิตฺตํ โลโก, ปญฺจปิ อุปาทานกฺขนฺธา โลโก.\csromanspacer{} ตตฺถ กตมา อภิชฺฌา, โย ราโค สาราโค ฯเปฯ จิตฺตสฺส สาราโค, อยํ วุจฺจติ อภิชฺฌา.\csromanspacer{} ตตฺถ กตมํ โทมนสฺสํ, ยํ เจตสิกํ อสาตํ เจตสิกํ ทุกฺขํ เจโตสมฺผสฺสชํ อสาตํ ทุกฺขํ เวทยิตํ เจโตสมฺผสฺสชา อสาตา ทุกฺขา เวทนา, อิทํ วุจฺจติ โทมนสฺสํ.\csromanspacer{} อิติ อยญฺจ อภิชฺฌา อิทญฺจ โทมนสฺสํ อิมมฺหิ โลเก วินีตา โหนฺติ ปฏิวินีตา สนฺตา สมิตา วูปสนฺตา อตฺถงฺคตา อพฺภตฺถงฺคตา อปฺปิตา พฺยปฺปิตา โสสิตา วิโสสิตา พฺยนฺตีกตา, เตน วุจฺจติ “วิเนยฺย โลเก อภิชฺฌาโทมนสฺสนฺ”ติ.}

\csromanheader{2}{จิตฺตานุปสฺสนานิทฺเทโส.}
\csromansectionrule
\csromanmatikaanchor{mtk.66}
\csromantocmarkref{subparagraph}{4. ธมฺมานุปสฺสนานิทฺเทส}{mtk.66}
\bookmark[dest={mtk.66},level=5]{4. ธมฺมานุปสฺสนานิทฺเทส}
\csromanheader{6}{4. ธมฺมานุปสฺสนานิทฺเทส}
\proseitem{367}{กถญฺจ ภิกฺขุ อชฺฌตฺตํ ธมฺเมสุ ธมฺมานุปสฺสี วิหรติ, อิธ ภิกฺขุ สนฺตํ วา อชฺฌตฺตํ กามจฺฉนฺทํ “อตฺถิ เม อชฺฌตฺตํ กามจฺฉนฺโท”ติ ปชานาติ, อสนฺตํ วา อชฺฌตฺตํ กามจฺฉนฺทํ “นตฺถิ เม อชฺฌตฺตํ กามจฺฉนฺโท”ติ ปชานาติ, ยถา จ อนุปฺปนฺนสฺส กามจฺฉนฺทสฺส อุปฺปาโท โหติ, ตญฺจ ปชานาติ, ยถา จ อุปฺปนฺนสฺส กามจฺฉนฺทสฺส ปหานํ โหติ, ตญฺจ ปชานาติ, ยถา จ ปหีนสฺส กามจฺฉนฺทสฺส อายติํ อนุปฺปาโท โหติ, ตญฺจ ปชานาติ.\csromanspacer{} สนฺตํ วา อชฺฌตฺตํ พฺยาปาทํ ฯเปฯ.\csromanspacer{} สนฺตํ วา อชฺฌตฺตํ ถินมิทฺธํ\footnote{ถีนมิทฺธํ (สี, สฺยา)} ฯเปฯ.\csromanspacer{} สนฺตํ วา อชฺฌตฺตํ อุทฺธจฺจกุกฺกุจฺจํ ฯเปฯ.\csromanspacer{} สนฺตํ วา อชฺฌตฺตํ วิจิกิจฺฉํ “อตฺถิ เม อชฺฌตฺตํ วิจิกิจฺฉา”ติ ปชานาติ, อสนฺตํ วา อชฺฌตฺตํ วิจิกิจฺฉํ “นตฺถิ เม อชฺฌตฺตํ}

\vspace{\tipitakaparskip}
\csromancenter{สติปฏฺฐานวิภงฺค วิจิกิจฺฉา”ติ ปชานาติ, ยถา จ อนุปฺปนฺนาย วิจิกิจฺฉาย อุปฺปาโท โหติ, ตญฺจ ปชานาติ, ยถา จ อุปฺปนฺนาย วิจิกิจฺฉาย ปหานํ โหติ, ตญฺจ ปชานาติ, ยถา จ ปหีนาย วิจิกิจฺฉาย อายติํ อนุปฺปาโท โหติ, ตญฺจ ปชานาติ.}
\prose{\csromanfolio{207}สนฺตํ วา อชฺฌตฺตํ สติสมฺโพชฺฌงฺคํ “อตฺถิ เม อชฺฌตฺตํ สติสมฺโพชฺฌงฺโค”ติ ปชานาติ, อสนฺตํ วา อชฺฌตฺตํ สติสมฺโพชฺฌงฺคํ “นตฺถิ เม อชฺฌตฺตํ สติสมฺโพชฺฌงฺโค”ติ ปชานาติ, ยถา จ อนุปฺปนฺนสฺส สติสมฺโพชฺฌงฺคสฺส อุปฺปาโท โหติ, ตญฺจ ปชานาติ, ยถา จ อุปฺปนฺนสฺส สติสมฺโพชฺฌงฺคสฺส ภาวนาย ปาริปูรี โหติ, ตญฺจ ปชานาติ.\csromanspacer{} สนฺตํ วา อชฺฌตฺตํ ธมฺมวิจยสมฺโพชฺฌงฺคํ ฯเปฯ.\csromanspacer{} สนฺตํ วา อชฺฌตฺตํ วีริยสมฺโพชฺฌงฺคํ\footnote{วิริยสมฺโพชฺฌงฺคํ (สี, สฺยา)} ฯเปฯ.\csromanspacer{} สนฺตํ วา อชฺฌตฺตํ ปีติสมฺโพชฺฌงฺคํ ฯเปฯ.\csromanspacer{} สนฺตํ วา อชฺฌตฺตํ ปสฺสทฺธิสมฺโพชฺฌงฺคํ ฯเปฯ.\csromanspacer{} สนฺตํ วา อชฺฌตฺตํ สมาธิสมฺโพชฺฌงฺคํ ฯเปฯ.\csromanspacer{} สนฺตํ วา อชฺฌตฺตํ อุเปกฺขาสมฺโพชฺฌงฺคํ “อตฺถิ เม อชฺฌตฺตํ อุเปกฺขาสมฺโพชฺฌงฺโค”ติ ปชานาติ, อสนฺตํ วา อชฺฌตฺตํ อุเปกฺขาสมฺโพชฺฌงฺคํ “นตฺถิ เม อชฺฌตฺตํ อุเปกฺขาสมฺโพชฺฌงฺโค”ติ ปชานาติ, ยถา จ อนุปฺปนฺนสฺส อุเปกฺขาสมฺโพชฺฌงฺคสฺส อุปฺปาโท โหติ, ตญฺจ ปชานาติ, ยถา จ อุปฺปนฺนสฺส อุเปกฺขาสมฺโพชฺฌงฺคสฺส ภาวนาย ปาริปูรี โหติ, ตญฺจ ปชานาติ.\csromanspacer{} โส ตํ นิมิตฺตํ อาเสวติ ภาเวติ พหุลีกโรติ, สฺวาวตฺถิตํ ววตฺถเปติ, โส ตํ นิมิตฺตํ อาเสวิตฺวา ภาเวตฺวา พหุลีกริตฺวา สฺวาวตฺถิตํ ววตฺถเปตฺวา พหิทฺธา ธมฺเมสุ จิตฺตํ อุปสํหรติ.}

\prose{กถญฺจ ภิกฺขุ พหิทฺธา ธมฺเมสุ ธมฺมานุปสฺสี วิหรติ, อิธ ภิกฺขุ สนฺตํ วาสฺส กามจฺฉนฺทํ “อตฺถิสฺส กามจฺฉนฺโท”ติ ปชานาติ, อสนฺตํ วาสฺส กามจฺฉนฺทํ “นตฺถิสฺส กามจฺฉนฺโท”ติ ปชานาติ, ยถา จ อนุปฺปนฺนสฺส กามจฺฉนฺทสฺส อุปฺปาโท โหติ, ตญฺจ ปชานาติ, ยถา จ อุปฺปนฺนสฺส กามจฺฉนฺทสฺส ปหานํ โหติ, ตญฺจ ปชานาติ, ยถา จ ปหีนสฺส กามจฺฉนฺทสฺส อายติํ อนุปฺปาโท โหติ, ตญฺจ ปชานาติ.\csromanspacer{} สนฺตํ วาสฺส พฺยาปาทํ ฯเปฯ.\csromanspacer{} สนฺตํ วาสฺส ถินมิทฺธํ ฯเปฯ.\csromanspacer{} สนฺตํ วาสฺส อุทฺธจฺจกุกฺกุจฺจํ ฯเปฯ.\csromanspacer{} สนฺตํ วาสฺส วิจิกิจฺฉํ “อตฺถิสฺส วิจิกิจฺฉา”ติ ปชานาติ, อสนฺตํ วาสฺส วิจิกิจฺฉํ “นตฺถิสฺส วิจิกิจฺฉา”ติ ปชานาติ, ยถา จ อนุปฺปนฺนาย \csromanfolio{208}วิจิกิจฺฉาย อุปฺปาโท โหติ, ตญฺจ ปชานาติ, ยถา จ อุปฺปนฺนาย วิจิกิจฺฉาย ปหานํ โหติ, ตญฺจ ปชานาติ, ยถา จ ปหีนาย วิจิกิจฺฉาย อายติํ อนุปฺปาโท โหติ, ตญฺจ ปชานาติ.}

\prose{สนฺตํ วาสฺส สติสมฺโพชฺฌงฺคํ “อตฺถิสฺส สติสมฺโพชฺฌงฺโค”ติ, ปชานาติ, อสนฺตํ วาสฺส สติสมฺโพชฺฌงฺคํ “นตฺถิสฺส สติสมฺโพชฺฌงฺโค”ติ ปชานาติ, ยถา จ อนุปฺปนฺนสฺส สติสมฺโพชฺฌงฺคสฺส อุปฺปาโท โหติ, ตญฺจ ปชานาติ, ยถา จ อุปฺปนฺนสฺส สติสมฺโพชฺฌงฺคสฺส ภาวนาย ปาริปูรี โหติ, ตญฺจ ปชานาติ.\csromanspacer{} สนฺตํ วาสฺส ธมฺมวิจยสมฺโพชฺฌงฺคํ ฯเปฯ.\csromanspacer{} สนฺตํ วาสฺส วีริยสมฺโพชฺฌงฺคํ ฯเปฯ.\csromanspacer{} สนฺตํ วาสฺส ปีติสมฺโพชฺฌงฺคํ ฯเปฯ.\csromanspacer{} สนฺตํ วาสฺส ปสฺสทฺธิสมฺโพชฺฌงฺคํ ฯเปฯ.\csromanspacer{} สนฺตํ วาสฺส สมาธิสมฺโพชฺฌงฺคํ ฯเปฯ.\csromanspacer{} สนฺตํ วาสฺส อุเปกฺขาสมฺโพชฺฌงฺคํ “อตฺถิสฺส อุเปกฺขาสมฺโพชฺฌงฺโค”ติ ปชานาติ, อสนฺตํ วาสฺส อุเปกฺขาสมฺโพชฺฌงฺคํ “นตฺถิสฺส อุเปกฺขาสมฺโพชฺฌงฺโค”ติ ปชานาติ, ยถา จ อนุปฺปนฺนสฺส อุเปกฺขาสมฺโพชฺฌงฺคสฺส อุปฺปาโท โหติ, ตญฺจ ปชานาติ, ยถา จ อุปฺปนฺนสฺส อุเปกฺขาสมฺโพชฺฌงฺคสฺส ภาวนาย ปาริปูรี โหติ, ตญฺจ ปชานาติ.\csromanspacer{} โส ตํ นิมิตฺตํ อาเสวติ ภาเวติ พหุลีกโรติ, สฺวาวตฺถิตํ ววตฺถเปติ, โส ตํ นิมิตฺตํ อาเสวิตฺวา ภาเวตฺวา พหุลีกริตฺวา สฺวาวตฺถิตํ ววตฺถเปตฺวา อชฺฌตฺตพหิทฺธา ธมฺเมสุ จิตฺตํ อุปสํหรติ.}

\prose{กถญฺจ ภิกฺขุ อชฺฌตฺตพหิทฺธา ธมฺเมสุ ธมฺมานุปสฺสี วิหรติ, อิธ ภิกฺขุ สนฺตํ วา กามจฺฉนฺทํ “อตฺถิ กามจฺฉนฺโท”ติ ปชานาติ, อสนฺตํ วา กามจฺฉนฺทํ “นตฺถิ กามจฺฉนฺโท”ติ ปชานาติ, ยถา จ อนุปฺปนฺนสฺส กามจฺฉนฺทสฺส อุปฺปาโท โหติ, ตญฺจ ปชานาติ, ยถา จ อุปฺปนฺนสฺส กามจฺฉนฺทสฺส ปหานํ โหติ, ตญฺจ ปชานาติ, ยถา จ ปหีนสฺส กามจฺฉนฺทสฺส อายติํ อนุปฺปาโท โหติ, ตญฺจ ปชานาติ.\csromanspacer{} สนฺตํ วา พฺยาปาทํ ฯเปฯ.\csromanspacer{} สนฺตํ วา ถินมิทฺธํ ฯเปฯ.\csromanspacer{} สนฺตํ วา อุทฺธจฺจกุกฺกุจฺจํ ฯเปฯ.\csromanspacer{} สนฺตํ วา วิจิกิจฺฉํ “อตฺถิ วิจิกิจฺฉา”ติ ปชานาติ, อสนฺตํ วา วิจิกิจฺฉํ “นตฺถิ วิจิกิจฺฉา”ติ ปชานาติ, ยถา จ อนุปฺปนฺนาย วิจิกิจฺฉาย อุปฺปาโท โหติ, ตญฺจ ปชานาติ, ยถา จ อุปฺปนฺนาย วิจิกิจฺฉาย ปหานํ โหติ, ตญฺจ ปชานาติ, ยถา จ ปหีนาย วิจิกิจฺฉาย อายติํ อนุปฺปาโท โหติ, ตญฺจ ปชานาติ.}

\csromanheader{2}{7. สติปฏฺฐานวิภงฺค}
\prose{\csromanfolio{209}สนฺตํ วา สติสมฺโพชฺฌงฺคํ “อตฺถิ สติสมฺโพชฺฌงฺโค”ติ ปชานาติ, อสนฺตํ วา สติสมฺโพชฺฌงฺคํ “นตฺถิ สติสมฺโพชฺฌงฺโค”ติ ปชานาติ, ยถา จ อนุปฺปนฺนสฺส สติสมฺโพชฺฌงฺคสฺส อุปฺปาโท โหติ, ตญฺจ ปชานาติ, ยถา จ อุปฺปนฺนสฺส สติสมฺโพชฺฌงฺคสฺส ภาวนาย ปาริปูรี โหติ, ตญฺจ ปชานาติ.\csromanspacer{} สนฺตํ วา ธมฺมวิจยสมฺโพชฺฌงฺคํ ฯเปฯ.\csromanspacer{} สนฺตํ วา วีริยสมฺโพชฺฌงฺคํ ฯเปฯ.\csromanspacer{} สนฺตํ วา ปีติสมฺโพชฺฌงฺคํ ฯเปฯ.\csromanspacer{} สนฺตํ วา ปสฺสทฺธิสมฺโพชฺฌงฺคํ ฯเปฯ.\csromanspacer{} สนฺตํ วา สมาธิสมฺโพชฺฌงฺคํ ฯเปฯ.\csromanspacer{} สนฺตํ วา อุเปกฺขาสมฺโพชฺฌงฺคํ “อตฺถิ อุเปกฺขาสมฺโพชฺฌงฺโค”ติ ปชานาติ, อสนฺตํ วา อุเปกฺขาสมฺโพชฺฌงฺคํ “นตฺถิ อุเปกฺขาสมฺโพชฺฌงฺโค”ติ ปชานาติ, ยถา จ อนุปฺปนฺนสฺส อุเปกฺขาสมฺโพชฺฌงฺคสฺส อุปฺปาโท โหติ, ตญฺจ ปชานาติ, ยถา จ อุปฺปนฺนสฺส อุเปกฺขาสมฺโพชฺฌงฺคสฺส ภาวนาย ปาริปูรี โหติ, ตญฺจ ปชานาติ.\csromanspacer{} เอวํ ภิกฺขุ อชฺฌตฺตพหิทฺธา ธมฺเมสุ ธมฺมานุปสฺสี วิหรติ อาตาปี สมฺปชาโน สติมา วิเนยฺย โลเก อภิชฺฌาโทมนสฺสํ.}

\proseitem{368}{“อนุปสฺสี”ติ ตตฺถ กตมา อนุปสฺสนา, ยา ปญฺญา ปชานนา ฯเปฯ อโมโห ธมฺมวิจโย สมฺมาทิฏฺฐิ, อยํ วุจฺจติ อนุปสฺสนา.\csromanspacer{} อิมาย อนุปสฺสนาย อุเปโต โหติ สมุเปโต อุปาคโต สมุปาคโต อุปปนฺโน สมฺปนฺโน สมนฺนาคโต, เตน วุจฺจติ “อนุปสฺสี”ติ.}

\proseitem{369}{“วิหรตี”ติ อิริยติ วตฺตติ ปาเลติ ยเปติ ยาเปติ จรติ วิหรติ, เตน วุจฺจติ “วิหรตี”ติ.}

\proseitem{370}{“อาตาปี”ติ ตตฺถ กตโม อาตาโป, โย เจตสิโก วีริยารมฺโภ ฯเปฯ สมฺมาวายาโม, อยํ วุจฺจติ อาตาโป.\csromanspacer{} อิมินา อาตาเปน อุเปโต โหติ ฯเปฯ สมนฺนาคโต, เตน วุจฺจติ “อาตาปี”ติ.}

\proseitem{371}{“สมฺปชาโน”ติ ตตฺถ กตมํ สมฺปชญฺญํ, ยา ปญฺญา ปชานนา ฯเปฯ อโมโห ธมฺมวิจโย สมฺมาทิฏฺฐิ, อิทํ วุจฺจติ สมฺปชญฺญํ.\csromanspacer{} อิมินา สมฺปชญฺเญน อุเปโต โหติ ฯเปฯ สมนฺนาคโต, เตน วุจฺจติ “สมฺปชาโน”ติ.}

\proseitem{372}{\csromanfolio{210}“สติมา”ติ ตตฺถ กตมา สติ, ยา สติ อนุสฺสติ ฯเปฯ สมฺมาสติ, อยํ วุจฺจติ สติ.\csromanspacer{} อิมาย สติยา อุเปโต โหติ ฯเปฯ สมนฺนาคโต, เตน วุจฺจติ “สติมา”ติ.}

\proseitem{373}{“วิเนยฺย โลเก อภิชฺฌาโทมนสฺสนฺ”ติ ตตฺถ กตโม โลโก, เตว ธมฺมา โลโก, ปญฺจปิ อุปาทานกฺขนฺธา โลโก, อยํ วุจฺจติ โลโก.\csromanspacer{} ตตฺถ กตมา อภิชฺฌา, โย ราโค สาราโค ฯเปฯ จิตฺตสฺส สาราโค, อยํ วุจฺจติ อภิชฺฌา.\csromanspacer{} ตตฺถ กตมํ โทมนสฺสํ, ยํ เจตสิกํ อสาตํ เจตสิกํ ทุกฺขํ เจโตสมฺผสฺสชํ อสาตํ ทุกฺขํ เวทยิตํ เจโตสมฺผสฺสชา อสาตา ทุกฺขา เวทนา, อิทํ วุจฺจติ โทมนสฺสํ.\csromanspacer{} อิติ อยญฺจ อภิชฺฌา อิทญฺจ โทมนสฺสํ อิมมฺหิ โลเก วินีตา โหนฺติ ปฏิวินีตา สนฺตา สมิตา วูปสนฺตา อตฺถงฺคตา อพฺภตฺถงฺคตา อปฺปิตา พฺยปฺปิตา โสสิตา วิโสสิตา พฺยนฺตีกตา, เตน วุจฺจติ “วิเนยฺย โลเก อภิชฺฌาโทมนสฺสนฺ”ติ.}

\vspace{\tipitakaparskip}
\csromancenter{ธมฺมานุปสฺสนานิทฺเทโส.}
\csromanheader{2}{สุตฺตนฺตภาชนียํ.}
\csromansectionrule
\csromanmatikaanchor{mtk.67}
\csromantocmarkref{subparagraph}{2. อภิธมฺมภาชนีย}{mtk.67}
\bookmark[dest={mtk.67},level=5]{2. อภิธมฺมภาชนีย}
\csromanheader{6}{2. อภิธมฺมภาชนีย}
\proseitem{374}{จตฺตาโร สติปฏฺฐานา\csromandash{}อิธ ภิกฺขุ กาเย กายานุปสฺสี วิหรติ, เวทนาสุ เวทนานุปสฺสี วิหรติ, จิตฺเต จิตฺตานุปสฺสี วิหรติ, ธมฺเมสุ ธมฺมานุปสฺสี วิหรติ.}

\proseitem{375}{กถญฺจ ภิกฺขุ กาเย กายานุปสฺสี วิหรติ, อิธ ภิกฺขุ ยสฺมิํ สมเย โลกุตฺตรํ ฌานํ ภาเวติ นิยฺยานิกํ อปจยคามิํ ทิฏฺฐิคตานํ ปหานาย ปฐมาย ภูมิยา ปตฺติยา วิวิจฺเจว กาเมหิ ฯเปฯ ปฐมํ ฌานํ อุปสมฺปชฺช วิหรติ ทุกฺขปฏิปทํ ทนฺธาภิญฺญํ กาเย กายานุปสฺสี, ยา ตสฺมิํ สมเย สติ อนุสฺสติ สมฺมาสติ สติสมฺโพชฺฌงฺโค มคฺคงฺคํ มคฺคปริยาปนฺนํ, อิทํ วุจฺจติ สติปฏฺฐานํ.\csromanspacer{} อวเสสา ธมฺมา สติปฏฺฐานสมฺปยุตฺตา.}

\csromanheader{2}{7. สติปฏฺฐานวิภงฺค}
\proseitem{376}{\csromanfolio{211}กถญฺจ ภิกฺขุ เวทนาสุ เวทนานุปสฺสี วิหรติ, อิธ ภิกฺขุ ยสฺมิํ สมเย โลกุตฺตรํ ฌานํ ภาเวติ นิยฺยานิกํ อปจยคามิํ ทิฏฺฐิคตานํ ปหานาย ปฐมาย ภูมิยา ปตฺติยา วิวิจฺเจว กาเมหิ ฯเปฯ ปฐมํ ฌานํ อุปสมฺปชฺช วิหรติ ทุกฺขปฏิปทํ ทนฺธาภิญฺญํ เวทนาสุ เวทนานุปสฺสี, ยา ตสฺมิํ สมเย สติ อนุสฺสติ สมฺมาสติ สติสมฺโพชฺฌงฺโค มคฺคงฺคํ มคฺคปริยาปนฺนํ, อิทํ วุจฺจติ สติปฏฺฐานํ.\csromanspacer{} อวเสสา ธมฺมา สติปฏฺฐานสมฺปยุตฺตา.}

\proseitem{377}{กถญฺจ ภิกฺขุ จิตฺเต จิตฺตานุปสฺสี วิหรติ, อิธ ภิกฺขุ ยสฺมิํ สมเย โลกุตฺตรํ ฌานํ ภาเวติ นิยฺยานิกํ อปจยคามิํ ทิฏฺฐิคตานํ ปหานาย ปฐมาย ภูมิยา ปตฺติยา วิวิจฺเจว กาเมหิ ฯเปฯ ปฐมํ ฌานํ อุปสมฺปชฺช วิหรติ ทุกฺขปฏิปทํ ทนฺธาภิญฺญํ จิตฺเต จิตฺตานุปสฺสี, ยา ตสฺมิํ สมเย สติ อนุสฺสติ สมฺมาสติ สติสมฺโพชฺฌงฺโค มคฺคงฺคํ มคฺคปริยาปนฺนํ, อิทํ วุจฺจติ สติปฏฺฐานํ.\csromanspacer{} อวเสสา ธมฺมา สติปฏฺฐานสมฺปยุตฺตา.}

\proseitem{378}{กถญฺจ ภิกฺขุ ธมฺเมสุ ธมฺมานุปสฺสี วิหรติ, อิธ ภิกฺขุ ยสฺมิํ สมเย โลกุตฺตรํ ฌานํ ภาเวติ นิยฺยานิกํ อปจยคามิํ ทิฏฺฐิคตานํ ปหานาย ปฐมาย ภูมิยา ปตฺติยา วิวิจฺเจว กาเมหิ ฯเปฯ ปฐมํ ฌานํ อุปสมฺปชฺช วิหรติ ทุกฺขปฏิปทํ ทนฺธาภิญฺญํ ธมฺเมสุ ธมฺมานุปสฺสี, ยา ตสฺมิํ สมเย สติ อนุสฺสติ สมฺมาสติ สติสมฺโพชฺฌงฺโค มคฺคงฺคํ มคฺคปริยาปนฺนํ, อิทํ วุจฺจติ สติปฏฺฐานํ.\csromanspacer{} อวเสสา ธมฺมา สติปฏฺฐานสมฺปยุตฺตา.}

\proseitem{379}{ตตฺถ กตมํ สติปฏฺฐานํ, อิธ ภิกฺขุ ยสฺมิํ สมเย โลกุตฺตรํ ฌานํ ภาเวติ นิยฺยานิกํ อปจยคามิํ ทิฏฺฐิคตานํ ปหานาย ปฐมาย ภูมิยา ปตฺติยา วิวิจฺเจว กาเมหิ ฯเปฯ ปฐมํ ฌานํ อุปสมฺปชฺช วิหรติ ทุกฺขปฏิปทํ ทนฺธาภิญฺญํ 1ธมฺเมสุ ธมฺมานุปสฺสี1, ยา ตสฺมิํ สมเย สติ อนุสฺสติ สมฺมาสติ สติสมฺโพชฺฌงฺโค มคฺคงฺคํ มคฺคปริยาปนฺนํ, อิทํ วุจฺจติ สติปฏฺฐานํ.\csromanspacer{} อวเสสา ธมฺมา สติปฏฺฐานสมฺปยุตฺตา.}

\proseitem{380}{จตฺตาโร สติปฏฺฐานา\csromandash{}อิธ ภิกฺขุ กาเย กายานุปสฺสี วิหรติ, เวทนาสุ เวทนานุปสฺสี วิหรติ, จิตฺเต จิตฺตานุปสฺสี วิหรติ ธมฺเมสุ ธมฺมานุปสฺสี วิหรติ.}

\proseitem{381}{\csromanfolio{212}กถญฺจ ภิกฺขุ กาเย กายานุปสฺสี วิหรติ, อิธ ภิกฺขุ ยสฺมิํ สมเย โลกุตฺตรํ ฌานํ ภาเวติ นิยฺยานิกํ อปจยคามิํ ทิฏฺฐิคตานํ ปหานาย ปฐมาย ภูมิยา ปตฺติยา วิวิจฺเจว กาเมหิ ฯเปฯ ปฐมํ ฌานํ อุปสมฺปชฺช วิหรติ ทุกฺขปฏิปทํ ทนฺธาภิญฺญํ, ตสฺมิํ สมเย ผสฺโส โหติ ฯเปฯ อวิกฺเขโป โหติ, อิเม ธมฺมา กุสลา.\csromanspacer{} ตสฺเสว โลกุตฺตรสฺส กุสลสฺส ฌานสฺส กตตฺตา ภาวิตตฺตา วิปากํ วิวิจฺเจว กาเมหิ ฯเปฯ ปฐมํ ฌานํ อุปสมฺปชฺช วิหรติ ทุกฺขปฏิปทํ ทนฺธาภิญฺญํ สุญฺญตํ กาเย กายานุปสฺสี, ยา ตสฺมิํ สมเย สติ อนุสฺสติ สมฺมาสติ สติสมฺโพชฺฌงฺโค มคฺคงฺคํ มคฺคปริยาปนฺนํ, อิทํ วุจฺจติ สติปฏฺฐานํ.\csromanspacer{} อวเสสา ธมฺมา สติปฏฺฐานสมฺปยุตฺตา.}

\proseitem{382}{กถญฺจ ภิกฺขุ เวทนาสุ เวทนานุปสฺสี วิหรติ, อิธ ภิกฺขุ ยสฺมิํ สมเย โลกุตฺตรํ ฌานํ ภาเวติ นิยฺยานิกํ อปจยคามิํ ทิฏฺฐิคตานํ ปหานาย ปฐมาย ภูมิยา ปตฺติยา วิวิจฺเจว กาเมหิ ฯเปฯ ปฐมํ ฌานํ อุปสมฺปชฺช วิหรติ ทุกฺขปฏิปทํ ทนฺธาภิญฺญํ, ตสฺมิํ สมเย ผสฺโส โหติ ฯเปฯ อวิกฺเขโป โหติ, อิเม ธมฺมา กุสลา.\csromanspacer{} ตสฺเสว โลกุตฺตรสฺส กุสลสฺส ฌานสฺส กตตฺตา ภาวิตตฺตา วิปากํ วิวิจฺเจว กาเมหิ ฯเปฯ ปฐมํ ฌานํ อุปสมฺปชฺช วิหรติ ทุกฺขปฏิปทํ ทนฺธาภิญฺญํ สุญฺญตํ เวทนาสุ เวทนานุปสฺสี, ยา ตสฺมิํ สมเย สติ อนุสฺสติ สมฺมาสติ สติสมฺโพชฺฌงฺโค มคฺคงฺคํ มคฺคปริยาปนฺนํ, อิทํ วุจฺจติ สติปฏฺฐานํ.\csromanspacer{} อวเสสา ธมฺมา สติปฏฺฐานสมฺปยุตฺตา.}

\proseitem{383}{กถญฺจ ภิกฺขุ จิตฺเต จิตฺตานุปสฺสี วิหรติ, อิธ ภิกฺขุ ยสฺมิํ สมเย โลกุตฺตรํ ฌานํ ภาเวติ นิยฺยานิกํ อปจยคามิํ ทิฏฺฐิคตานํ ปหานาย ปฐมาย ภูมิยา ปตฺติยา วิวิจฺเจว กาเมหิ ฯเปฯ ปฐมํ ฌานํ อุปสมฺปชฺช วิหรติ ทุกฺขปฏิปทํ ทนฺธาภิญฺญํ, ตสฺมิํ สมเย ผสฺโส โหติ ฯเปฯ อวิกฺเขโป โหติ, อิเม ธมฺมา กุสลา.\csromanspacer{} ตสฺเสว โลกุตฺตรสฺส กุสลสฺส ฌานสฺส กตตฺตา ภาวิตตฺตา วิปากํ วิวิจฺเจว กาเมหิ ฯเปฯ ปฐมํ ฌานํ อุปสมฺปชฺช วิหรติ ทุกฺขปฏิปทํ ทนฺธาภิญฺญํ สุญฺญตํ จิตฺเต จิตฺตานุปสฺสี, ยา ตสฺมิํ สมเย สติ อนุสฺสติ สมฺมาสติ สติสมฺโพชฺฌงฺโค มคฺคงฺคํ มคฺคปริยาปนฺนํ, อิทํ วุจฺจติ สติปฏฺฐานํ.\csromanspacer{} อวเสสา ธมฺมา สติปฏฺฐานสมฺปยุตฺตา.}

\csromanheader{2}{7. สติปฏฺฐานวิภงฺค}
\proseitem{384}{\csromanfolio{213}กถญฺจ ภิกฺขุ ธมฺเมสุ ธมฺมานุปสฺสี วิหรติ, อิธ ภิกฺขุ ยสฺมิํ สมเย โลกุตฺตรํ ฌานํ ภาเวติ นิยฺยานิกํ อปจยคามิํ ทิฏฺฐิคตานํ ปหานาย ปฐมาย ภูมิยา ปตฺติยา วิวิจฺเจว กาเมหิ ฯเปฯ ปฐมํ ฌานํ อุปสมฺปชฺช วิหรติ ทุกฺขปฏิปทํ ทนฺธาภิญฺญํ, ตสฺมิํ สมเย ผสฺโส โหติ ฯเปฯ อวิกฺเขโป โหติ, อิเม ธมฺมา กุสลา.\csromanspacer{} ตสฺเสว โลกุตฺตรสฺส กุสลสฺส ฌานสฺส กตตฺตา ภาวิตตฺตา วิปากํ วิวิจฺเจว กาเมหิ ฯเปฯ ปฐมํ ฌานํ อุปสมฺปชฺช วิหรติ ทุกฺขปฏิปทํ ทนฺธาภิญฺญํ สุญฺญตํ ธมฺเมสุ ธมฺมานุปสฺสี, ยา ตสฺมิํ สมเย สติ อนุสฺสติ สมฺมาสติ สติสมฺโพชฺฌงฺโค มคฺคงฺคํ มคฺคปริยาปนฺนํ, อิทํ วุจฺจติ สติปฏฺฐานํ.\csromanspacer{} อวเสสา ธมฺมา สติปฏฺฐานสมฺปยุตฺตา.}

\proseitem{385}{ตตฺถ กตมํ สติปฏฺฐานํ, อิธ ภิกฺขุ ยสฺมิํ สมเย โลกุตฺตรํ ฌานํ ภาเวติ นิยฺยานิกํ อปจยคามิํ ทิฏฺฐิคตานํ ปหานาย ปฐมาย ภูมิยา ปตฺติยา วิวิจฺเจว กาเมหิ ฯเปฯ ปฐมํ ฌานํ อุปสมฺปชฺช วิหรติ ทุกฺขปฏิปทํ ทนฺธาภิญฺญํ, ตสฺมิํ สมเย ผสฺโส โหติ ฯเปฯ อวิกฺเขโป โหติ, อิเม ธมฺมา กุสลา.\csromanspacer{} ตสฺเสว โลกุตฺตรสฺส กุสลสฺส ฌานสฺส กตตฺตา ภาวิตตฺตา วิปากํ วิวิจฺเจว กาเมหิ ฯเปฯ ปฐมํ ฌานํ อุปสมฺปชฺช วิหรติ ทุกฺขปฏิปทํ ทนฺธาภิญฺญํ สุญฺญตํ, ยา ตสฺมิํ สมเย สติ อนุสฺสติ สมฺมาสติ สติสมฺโพชฺฌงฺโค มคฺคงฺคํ มคฺคปริยาปนฺนํ, อิทํ วุจฺจติ สติปฏฺฐานํ.\csromanspacer{} อวเสสา ธมฺมา สติปฏฺฐานสมฺปยุตฺตา.}

\csromanheader{2}{อภิธมฺมภาชนียํ.}
\csromansectionrule
\csromanmatikaanchor{mtk.68}
\csromantocmarkref{subparagraph}{3. ปญฺหาปุจฺฉก}{mtk.68}
\bookmark[dest={mtk.68},level=5]{3. ปญฺหาปุจฺฉก}
\csromanheader{6}{3. ปญฺหาปุจฺฉก}
\proseitem{386}{จตฺตาโร สติปฏฺฐานา\csromandash{}อิธ ภิกฺขุ กาเย กายานุปสฺสี วิหรติ อาตาปี สมฺปชาโน สติมา วิเนยฺย โลเก อภิชฺฌาโทมนสฺสํ, เวทนาสุ เวทนานุปสฺสี วิหรติ อาตาปี สมฺปชาโน สติมา วิเนยฺย โลเก อภิชฺฌาโทมนสฺสํ, จิตฺเต จิตฺตานุปสฺสี วิหรติ อาตาปี สมฺปชาโน สติมา วิเนยฺย โลเก อภิชฺฌาโทมนสฺสํ, ธมฺเมสุ ธมฺมานุปสฺสี วิหรติ อาตาปี สมฺปชาโน สติมา วิเนยฺย โลเก อภิชฺฌาโทมนสฺสํ.}

\proseitem{387}{\csromanfolio{214}จตุนฺนํ สติปฏฺฐานานํ กติ กุสลา, กติ อกุสลา, กติ อพฺยากตา ฯเปฯ กติ สรณา, กติ อรณา.}

\csromanmatikaanchor{mtk.69}
\csromantocmarkref{subparagraph}{1. ติก}{mtk.69}
\bookmark[dest={mtk.69},level=5]{1. ติก}
\csromanheader{6}{1. ติก}
\proseitem{388}{สิยา กุสลา, สิยา อพฺยากตา ฯเปฯ สิยา สุขาย เวทนาย สมฺปยุตฺตา สิยา อทุกฺขมสุขาย เวทนาย สมฺปยุตฺตา.\csromanspacer{} สิยา วิปากา, สิยา วิปากธมฺมธมฺมา.\csromanspacer{} อนุปาทินฺน-อนุปาทานิยา.\csromanspacer{} อสํกิลิฏฺฐ- อสํกิเลสิกา.\csromanspacer{} สิยา สวิตกฺกสวิจารา, สิยา อวิตกฺกวิจารมตฺตา, สิยา อวิตกฺก-อวิจารา.\csromanspacer{} สิยา ปีติสหคตา, สิยา สุขสหคตา, สิยา อุเปกฺขาสหคตา.\csromanspacer{} เนว ทสฺสเนน น ภาวนาย ปหาตพฺพา, เนว ทสฺสเนน น ภาวนาย ปหาตพฺพเหตุกา.\csromanspacer{} สิยา อปจยคามิโน, สิยา เนวาจยคามินาปจยคามิโน.\csromanspacer{} สิยา เสกฺขา, สิยา อเสกฺขา.\csromanspacer{} อปฺปมาณา.\csromanspacer{} อปฺปมาณารมฺมณา.\csromanspacer{} ปณีตา.\csromanspacer{} สิยา สมฺมตฺตนิยตา, สิยา อนิยตา.\csromanspacer{} น มคฺคารมฺมณา.\csromanspacer{} สิยา มคฺคเหตุกา, สิยา มคฺคาธิปติโน, สิยา น วตฺตพฺพา “มคฺคเหตุกา”ติปิ, “มคฺคาธิปติโน”ติปิ.\csromanspacer{} สิยา อุปฺปนฺนา, สิยา อนุปฺปนฺนา, สิยา อุปฺปาทิโน.\csromanspacer{} สิยา อตีตา, สิยา อนาคตา, สิยา ปจฺจุปฺปนฺนา.\csromanspacer{} น วตฺตพฺพา “อตีตารมฺมณา”ติปิ, “อนาคตารมฺมณา”ติปิ, “ปจฺจุปฺปนฺนารมฺมณา”ติปิ.\csromanspacer{} สิยา อชฺฌตฺตา, สิยา พหิทฺธา, สิยา อชฺฌตฺตพหิทฺธา.\csromanspacer{} พหิทฺธารมฺมณา.\csromanspacer{} อนิทสฺสน-อปฺปฏิฆา.}

\csromanmatikaanchor{mtk.70}
\csromantocmarkref{subparagraph}{2. ทุก}{mtk.70}
\bookmark[dest={mtk.70},level=5]{2. ทุก}
\csromanheader{6}{2. ทุก}
\proseitem{389}{น เหตู, สเหตุกา, เหตุสมฺปยุตฺตา.\csromanspacer{} น วตฺตพฺพา “เหตู เจว สเหตุกา จา”ติ, สเหตุกา เจว น จ เหตู.\csromanspacer{} น วตฺตพฺพา “เหตู เจว เหตุสมฺปยุตฺตา จา”ติ, เหตุสมฺปยุตฺตา เจว น จ เหตู.\csromanspacer{} น เหตู สเหตุกา.}

\prose{สปฺปจฺจยา, สงฺขตา, อนิทสฺสนา, อปฺปฏิฆา, อรูปา, โลกุตฺตรา, เกนจิ วิญฺเญยฺยา, เกนจิ น วิญฺเญยฺยา.\csromanspacer{} โน อาสวา, อนาสวา.\csromanspacer{} อาสววิปฺปยุตฺตา.\csromanspacer{} น วตฺตพฺพา “อาสวา เจว สาสวา จา”ติปิ, “สาสวา เจว โน จ อาสวา”ติปิ.\csromanspacer{} น วตฺตพฺพา “อาสวา เจว อาสวสมฺปยุตฺตา จา”ติปิ, “อาสวสมฺปยุตฺตา เจว โน จ อาสวา”ติปิ.\csromanspacer{} อาสววิปฺปยุตฺตา.\csromanspacer{} อนาสวา.\csromanspacer{} โน สํโยชนา ฯเปฯ.\csromanspacer{} โน คนฺถา ฯเปฯ.\csromanspacer{} โน โอฆา ฯเปฯ.}

\vspace{\tipitakaparskip}
\csromancenter{สติปฏฺฐานวิภงฺค โน โยคา ฯเปฯ.\csromanspacer{} โน นีวรณา ฯเปฯ.\csromanspacer{} โน ปรามาสา ฯเปฯ สารมฺมณา, โน จิตฺตา, เจตสิกา, จิตฺตสมฺปยุตฺตา, จิตฺตสํสฏฺฐา, จิตฺตสมุฏฺฐานา, จิตฺตสหภุโน, จิตฺตานุปริวตฺติโน, จิตฺตสํสฏฺฐสมุฏฺฐานา, จิตฺตสํสฏฺฐสมุฏฺฐานสหภุโน, จิตฺตสํสฏฺฐสมุฏฺฐานานุปริวตฺติโน, พาหิรา.\csromanspacer{} โน อุปาทา, อนุปาทินฺนา, โน อุปาทานา ฯเปฯ โน กิเลสา ฯเปฯ น ทสฺสเนน ปหาตพฺพา, น ภาวนาย ปหาตพฺพา.\csromanspacer{} น ทสฺสเนน ปหาตพฺพเหตุกา, น ภาวนาย ปหาตพฺพเหตุกา.}
\prose{\csromanfolio{215}สิยา สวิตกฺกา, สิยา อวิตกฺกา.\csromanspacer{} สิยา สวิจารา, สิยา อวิจารา.\csromanspacer{} สิยา สปฺปีติกา, สิยา อปฺปีติกา.\csromanspacer{} สิยา ปีติสหคตา, สิยา น ปีติสหคตา.\csromanspacer{} สิยา สุขสหคตา, สิยา น สุขสหคตา.\csromanspacer{} สิยา อุเปกฺขาสหคตา, สิยา น อุเปกฺขาสหคตา.\csromanspacer{} น กามาวจรา.\csromanspacer{} น รูปาวจรา.\csromanspacer{} น อรูปาวจรา.\csromanspacer{} อปริยาปนฺนา.\csromanspacer{} สิยา นิยฺยานิกา, สิยา อนิยฺยานิกา.\csromanspacer{} สิยา นิยตา, สิยา อนิยตา.\csromanspacer{} อนุตฺตรา.\csromanspacer{} อรณาติ.}

\vspace{\tipitakaparskip}
\csromancenter{ปญฺหาปุจฺฉกํ.}
\nitthitam{\textbf{สติปฏฺฐานวิภงฺโค นิฏฺฐิโต.}}
\csromanmatikaanchor{mtk.71}
\csromantocmarknopage{subparagraph}{8. สมฺมปฺปธานวิภงฺค}{mtk.71}
\bookmark[dest={mtk.71},level=5]{8. สมฺมปฺปธานวิภงฺค}
\csromanheader{6}{8. สมฺมปฺปธานวิภงฺค}
\csromanmatikaanchor{mtk.72}
\csromantocmarkref{subparagraph}{1. สุตฺตนฺตภาชนีย}{mtk.72}
\bookmark[dest={mtk.72},level=5]{1. สุตฺตนฺตภาชนีย}
\csromanheader{6}{1. สุตฺตนฺตภาชนีย}
\proseitem{390}{\csromanfolio{216}จตฺตาโร สมฺมปฺปธานา\csromandash{}อิธ ภิกฺขุ อนุปฺปนฺนานํ ปาปกานํ อกุสลานํ ธมฺมานํ อนุปฺปาทาย ฉนฺทํ ชเนติ วายมติ วีริยํ อารภติ จิตฺตํ ปคฺคณฺหาติ ปทหติ, อุปฺปนฺนานํ ปาปกานํ อกุสลานํ ธมฺมานํ ปหานาย ฉนฺทํ ชเนติ วายมติ วีริยํ อารภติ จิตฺตํ ปคฺคณฺหาติ ปทหติ อนุปฺปนฺนานํ กุสลานํ ธมฺมานํ อุปฺปาทาย ฉนฺทํ ชเนติ วายมติ วีริยํ อารภติ จิตฺตํ ปคฺคณฺหาติ ปทหติ, อุปฺปนฺนานํ กุสลานํ ธมฺมานํ ฐิติยา อสมฺโมสาย ภิยฺโยภาวาย เวปุลฺลาย ภาวนาย ปาริปูริยา ฉนฺทํ ชเนติ วายมติ วีริยํ อารภติ จิตฺตํ ปคฺคณฺหาติ ปทหติ.}

\proseitem{391}{กถญฺจ ภิกฺขุ อนุปฺปนฺนานํ ปาปกานํ อกุสลานํ ธมฺมานํ อนุปฺปาทาย ฉนฺทํ ชเนติ วายมติ วีริยํ อารภติ จิตฺตํ ปคฺคณฺหาติ ปทหติ.\csromanspacer{} ตตฺถ กตเม อนุปฺปนฺนา ปาปกา อกุสลา ธมฺมา, ตีณิ อกุสลมูลานิ โลโภ โทโส โมโห ตเทกฏฺฐา จ กิเลสา, ตํสมฺปยุตฺโต เวทนากฺขนฺโธ สญฺญากฺขนฺโธ สงฺขารกฺขนฺโธ วิญฺญาณกฺขนฺโธ, ตํสมุฏฺฐานํ กายกมฺมํ วจีกมฺมํ มโนกมฺมํ, อิเม วุจฺจนฺติ อนุปฺปนฺนา ปาปกา อกุสลา ธมฺมา.\csromanspacer{} อิติ อิเมสํ อนุปฺปนฺนานํ ปาปกานํ อกุสลานํ ธมฺมานํ อนุปฺปาทาย ฉนฺทํ ชเนติ วายมติ วีริยํ อารภติ จิตฺตํ ปคฺคณฺหาติ ปทหติ.}

\proseitem{392}{“ฉนฺทํ ชเนตี”ติ ตตฺถ กตโม ฉนฺโท, โย ฉนฺโท ฉนฺทิกตา\footnote{ฉนฺทีกตา (สฺยา)} กตฺตุกมฺยตา กุสโล ธมฺมจฺฉนฺโท, อยํ วุจฺจติ ฉนฺโท.\csromanspacer{} อิมํ ฉนฺทํ ชเนติ สญฺชเนติ อุฏฺฐเปติ สมุฏฺฐเปติ\footnote{อุฏฺฐาเปติ สมุฏฺฐาเปติ (สฺยา) เอวมุปริปิ.} นิพฺพตฺเตติ อภินิพฺพตฺเตติ, เตน วุจฺจติ “ฉนฺทํ ชเนตี”ติ.}

\proseitem{393}{“วายมตี”ติ ตตฺถ กตโม วายาโม, โย เจตสิโก วีริยารมฺโภ ฯเปฯ สมฺมาวายาโม, อยํ วุจฺจติ วายาโม.\csromanspacer{} อิมินา วายาเมน อุเปโต โหติ สมุเปโต อุปาคโต สมุปาคโต อุปปนฺโน สมฺปนฺโน สมนฺนาคโต, เตน วุจฺจติ “วายมตี”ติ.}

\csromanheader{2}{8. สมฺมปฺปธานวิภงฺค}
\proseitem{394}{\csromanfolio{217}“วีริยํ อารภตี”ติ ตตฺถ กตมํ วีริยํ, โย เจตสิโก วีริยารมฺโภ ฯเปฯ สมฺมาวายาโม, อิทํ วุจฺจติ วีริยํ.\csromanspacer{} อิมํ วีริยํ อารภติ สมารภติ อาเสวติ ภาเวติ พหุลีกโรติ, เตน วุจฺจติ “วีริยํ อารภตี”ติ.}

\proseitem{395}{“จิตฺตํ ปคฺคณฺหาตี”ติ ตตฺถ กตมํ จิตฺตํ, ยํ จิตฺตํ มโน มานสํ ฯเปฯ ตชฺชามโนวิญฺญาณธาตุ, อิทํ วุจฺจติ จิตฺตํ.\csromanspacer{} อิมํ จิตฺตํ ปคฺคณฺหาติ สมฺปคฺคณฺหาติ อุปตฺถมฺเภติ ปจฺจุปตฺถมฺเภติ, เตน วุจฺจติ “จิตฺตํ ปคฺคณฺหาตี”ติ.}

\proseitem{396}{“ปทหตี”ติ ตตฺถ กตมํ ปธานํ, โย เจตสิโก วีริยารมฺโภ ฯเปฯ สมฺมาวายาโม.\csromanspacer{} อิทํ วุจฺจติ ปธานํ, อิมินา ปธาเนน อุเปโต โหติ ฯเปฯ สมนฺนาคโต, เตน วุจฺจติ “ปทหตี”ติ.}

\proseitem{397}{กถญฺจ ภิกฺขุ อุปฺปนฺนานํ ปาปกานํ อกุสลานํ ธมฺมานํ ปหานาย ฉนฺทํ ชเนติ วายมติ วีริยํ อารภติ จิตฺตํ ปคฺคณฺหาติ ปทหติ.\csromanspacer{} ตตฺถ กตเม อุปฺปนฺนา ปาปกา อกุสลา ธมฺมา, ตีณิ อกุสลมูลานิ โลโภ โทโส โมโห ตเทกฏฺฐา จ กิเลสา, ตํสมฺปยุตฺโต เวทนากฺขนฺโธ สญฺญากฺขนฺโธ สงฺขารกฺขนฺโธ วิญฺญาณกฺขนฺโธ, ตํสมุฏฺฐานํ กายกมฺมํ วจีกมฺมํ มโนกมฺมํ, อิเม วุจฺจนฺติ อุปฺปนฺนา ปาปกา อกุสลา ธมฺมา.\csromanspacer{} อิติ อิเมสํ อุปฺปนฺนานํ ปาปกานํ อกุสลานํ ธมฺมานํ ปหานาย ฉนฺทํ ชเนติ วายมติ วีริยํ อารภติ จิตฺตํ ปคฺคณฺหาติ ปทหติ.}

\proseitem{398}{“ฉนฺทํ ชเนตี”ติ ตตฺถ กตโม ฉนฺโท, โย ฉนฺโท ฉนฺทิกตา กตฺตุกมฺยตา กุสโล ธมฺมจฺฉนฺโท, อยํ วุจฺจติ ฉนฺโท.\csromanspacer{} อิมํ ฉนฺทํ ชเนติ สญฺชเนติ อุฏฺฐเปติ สมุฏฺฐเปติ นิพฺพตฺเตติ อภินิพฺพตฺเตติ, เตน วุจฺจติ “ฉนฺทํ ชเนตี”ติ.}

\proseitem{399}{“วายมตี”ติ ตตฺถ กตโม วายาโม, โย เจตสิโก วีริยารมฺโภ ฯเปฯ สมฺมาวายาโม, อยํ วุจฺจติ วายาโม.\csromanspacer{} อิมินา วายาเมน อุเปโต โหติ ฯเปฯ สมนฺนาคโต, เตน วุจฺจติ “วายมตี”ติ.}

\proseitem{400}{\csromanfolio{218}“วีริยํ อารภตี”ติ ตตฺถ กตมํ วีริยํ, โย เจตสิโก วีริยารมฺโภ ฯเปฯ สมฺมาวายาโม, อิทํ วุจฺจติ วีริยํ.\csromanspacer{} อิมํ วีริยํ อารภติ สมารภติ อาเสวติ ภาเวติ พหุลีกโรติ, เตน วุจฺจติ “วีริยํ อารภตี”ติ.}

\proseitem{401}{“จิตฺตํ ปคฺคณฺหาตี”ติ ตตฺถ กตมํ จิตฺตํ, ยํ จิตฺตํ มโน มานสํ ฯเปฯ ตชฺชามโนวิญฺญาณธาตุ, อิทํ วุจฺจติ จิตฺตํ.\csromanspacer{} อิมํ จิตฺตํ ปคฺคณฺหาติ สมฺปคฺคณฺหาติ อุปตฺถมฺเภติ ปจฺจุปตฺถมฺเภติ, เตน วุจฺจติ “จิตฺตํ ปคฺคณฺหาตี”ติ.}

\proseitem{402}{“ปทหตี”ติ ตตฺถ กตมํ ปธานํ, โย เจตสิโก วีริยารมฺโภ ฯเปฯ สมฺมาวายาโม, อิทํ วุจฺจติ ปธานํ.\csromanspacer{} อิมินา ปธาเนน อุเปโต โหติ ฯเปฯ สมนฺนาคโต, เตน วุจฺจติ “ปทหตี”ติ.}

\proseitem{403}{กถญฺจ ภิกฺขุ อนุปฺปนฺนานํ กุสลานํ ธมฺมานํ อุปฺปาทาย ฉนฺทํ ชเนติ วายมติ วีริยํ อารภติ จิตฺตํ ปคฺคณฺหาติ ปทหติ.\csromanspacer{} ตตฺถ กตเม อนุปฺปนฺนา กุสลา ธมฺมา, ตีณิ กุสลมูลานิ อโลโภ อโทโส อโมโห, ตํสมฺปยุตฺโต เวทนากฺขนฺโธ สญฺญากฺขนฺโธ สงฺขารกฺขนฺโธ วิญฺญาณกฺขนฺโธ, ตํสมุฏฺฐานํ กายกมฺมํ วจีกมฺมํ มโนกมฺมํ, อิเม วุจฺจนฺติ อนุปฺปนฺนา กุสลา ธมฺมา.\csromanspacer{} อิติ อิเมสํ อนุปฺปนฺนานํ กุสลานํ ธมฺมานํ อุปฺปาทาย ฉนฺทํ ชเนติ วายมติ วีริยํ อารภติ จิตฺตํ ปคฺคณฺหาติ ปทหติ.}

\proseitem{404}{“ฉนฺทํ ชเนตี”ติ ฯเปฯ. “วายมตี”ติ ฯเปฯ. “วีริยํ อารภตี”ติ ฯเปฯ. “จิตฺตํ ปคฺคณฺหาตี”ติ ฯเปฯ. “ปทหตี”ติ ตตฺถ กตมํ ปธานํ, โย เจตสิโก วีริยารมฺโภ ฯเปฯ สมฺมาวายาโม, อิทํ วุจฺจติ ปธานํ.\csromanspacer{} อิมินา ปธาเนน อุเปโต โหติ ฯเปฯ สมนฺนาคโต, เตน วุจฺจติ “ปทหตี”ติ.}

\proseitem{405}{กถญฺจ ภิกฺขุ อุปฺปนฺนานํ กุสลานํ ธมฺมานํ ฐิติยา อสมฺโมสาย ภิยฺโยภาวาย เวปุลฺลาย ภาวนาย ปาริปูริยา ฉนฺทํ ชเนติ วายมติ วีริยํ อารภติ จิตฺตํ ปคฺคณฺหาติ ปทหติ.\csromanspacer{} ตตฺถ กตเม อุปฺปนฺนา กุสลา ธมฺมา, ตีณิ กุสลมูลานิ อโลโภ อโทโส อโมโห, ตํสมฺปยุตฺโต เวทนากฺขนฺโธ สญฺญากฺขนฺโธ สงฺขารกฺขนฺโธ}

\vspace{\tipitakaparskip}
\csromancenter{สมฺมปฺปธานวิภงฺค วิญฺญาณกฺขนฺโธ, ตํสมุฏฺฐานํ กายกมฺมํ วจีกมฺมํ มโนกมฺมํ, อิเม วุจฺจนฺติ อุปฺปนฺนา กุสลา ธมฺมา.\csromanspacer{} อิติ อิเมสํ อุปฺปนฺนานํ กุสลานํ ธมฺมานํ ฐิติยา อสมฺโมสาย ภิยฺโยภาวาย เวปุลฺลาย ภาวนาย ปาริปูริยา ฉนฺทํ ชเนติ วายมติ วีริยํ อารภติ จิตฺตํ ปคฺคณฺหาติ ปทหติ.}
\proseitem{406}{\csromanfolio{219}“ฐิติยา”ติ ยา ฐิติ โส อสมฺโมโส, โย อสมฺโมโส โส ภิยฺโยภาโว, โย ภิยฺโยภาโว ตํ เวปุลฺลํ, ยํ เวปุลฺลํ สา ภาวนา, ยา ภาวนา สา ปาริปูรี.}

\proseitem{407}{“ฉนฺทํ ชเนตี”ติ ฯเปฯ. “วายมตี”ติ ฯเปฯ. “วีริยํ อารภตี”ติ ฯเปฯ. “จิตฺตํ ปคฺคณฺหาตี”ติ ฯเปฯ. “ปทหตี”ติ ตตฺถ กตมํ ปธานํ, โย เจตสิโก วีริยารมฺโภ ฯเปฯ สมฺมาวายาโม, อิทํ วุจฺจติ ปธานํ.\csromanspacer{} อิมินา ปธาเนน อุเปโต โหติ ฯเปฯ สมนฺนาคโต, เตน วุจฺจติ “ปทหตี”ติ.}

\csromanheader{2}{สุตฺตนฺตภาชนียํ.}
\csromansectionrule
\csromanmatikaanchor{mtk.73}
\csromantocmarkref{subparagraph}{2. อภิธมฺมภาชนีย}{mtk.73}
\bookmark[dest={mtk.73},level=5]{2. อภิธมฺมภาชนีย}
\csromanheader{6}{2. อภิธมฺมภาชนีย}
\proseitem{408}{จตฺตาโร สมฺมปฺปธานา\csromandash{}อิธ ภิกฺขุ อนุปฺปนฺนานํ ปาปกานํ อกุสลานํ ธมฺมานํ อนุปฺปาทาย ฉนฺทํ ชเนติ วายมติ วีริยํ อารภติ จิตฺตํ ปคฺคณฺหาติ ปทหติ, อุปฺปนฺนานํ ปาปกานํ อกุสลานํ ธมฺมานํ ปหานาย ฉนฺทํ ชเนติ วายมติ วีริยํ อารภติ จิตฺตํ ปคฺคณฺหาติ ปทหติ, อนุปฺปนฺนานํ กุสลานํ ธมฺมานํ อุปฺปาทาย ฉนฺทํ ชเนติ วายมติ วีริยํ อารภติ จิตฺตํ ปคฺคณฺหาติ ปทหติ, อุปฺปนฺนานํ กุสลานํ ธมฺมานํ ฐิติยา อสมฺโมสาย ภิยฺโยภาวาย เวปุลฺลาย ภาวนาย ปาริปูริยา ฉนฺทํ ชเนติ วายมติ วีริยํ อารภติ จิตฺตํ ปคฺคณฺหาติ ปทหติ.}

\proseitem{409}{กถญฺจ ภิกฺขุ อนุปฺปนฺนานํ ปาปกานํ อกุสลานํ ธมฺมานํ อนุปฺปาทาย ฉนฺทํ ชเนติ วายมติ วีริยํ อารภติ จิตฺตํ ปคฺคณฺหาติ ปทหติ.\csromanspacer{} อิธ ภิกฺขุ ยสฺมิํ สมเย โลกุตฺตรํ ฌานํ ภาเวติ นิยฺยานิกํ อปจยคามิํ ทิฏฺฐิคตานํ ปหานาย ปฐมาย ภูมิยา \csromanfolio{220}ปตฺติยา วิวิจฺเจว กาเมหิ ฯเปฯ ปฐมํ ฌานํ อุปสมฺปชฺช วิหรติ ทุกฺขปฏิปทํ ทนฺธาภิญฺญํ, ตสฺมิํ สมเย อนุปฺปนฺนานํ ปาปกานํ อกุสลานํ ธมฺมานํ อนุปฺปาทาย ฉนฺทํ ชเนติ วายมติ วีริยํ อารภติ จิตฺตํ ปคฺคณฺหาติ ปทหติ.}

\proseitem{410}{“ฉนฺทํ ชเนตี”ติ ตตฺถ กตโม ฉนฺโท, โย ฉนฺโท ฉนฺทิกตา กตฺตุกมฺยตา กุสโล ธมฺมจฺฉนฺโท, อยํ วุจฺจติ ฉนฺโท.\csromanspacer{} อิมํ ฉนฺทํ ชเนติ สญฺชเนติ อุฏฺฐเปติ สมุฏฺฐเปติ นิพฺพตฺเตติ อภินิพฺพตฺเตติ, เตน วุจฺจติ “ฉนฺทํ ชเนตี”ติ.}

\proseitem{411}{“วายมตี”ติ ตตฺถ กตโม วายาโม, โย เจตสิโก วีริยารมฺโภ ฯเปฯ สมฺมาวายาโม วีริยสมฺโพชฺฌงฺโค มคฺคงฺคํ มคฺคปริยาปนฺนํ, อยํ วุจฺจติ วายาโม.\csromanspacer{} อิมินา วายาเมน อุเปโต โหติ สมุเปโต อุปาคโต สมุปาคโต อุปปนฺโน สมฺปนฺโน สมนฺนาคโต, เตน วุจฺจติ “วายมตี”ติ.}

\proseitem{412}{“วีริยํ อารภตี”ติ ตตฺถ กตมํ วีริยํ, โย เจตสิโก วีริยารมฺโภ ฯเปฯ สมฺมาวายาโม วีริยสมฺโพชฺฌงฺโค มคฺคงฺคํ มคฺคปริยาปนฺนํ, อิทํ วุจฺจติ วีริยํ.\csromanspacer{} อิมํ วีริยํ อารภติ สมารภติ อาเสวติ ภาเวติ พหุลีกโรติ, เตน วุจฺจติ “วีริยํ อารภตี”ติ.}

\proseitem{413}{“จิตฺตํ ปคฺคณฺหาตี”ติ ตตฺถ กตมํ จิตฺตํ, ยํ จิตฺตํ มโน มานสํ ฯเปฯ ตชฺชามโนวิญฺญาณธาตุ, อิทํ วุจฺจติ จิตฺตํ.\csromanspacer{} อิมํ จิตฺตํ ปคฺคณฺหาติ สมฺปคฺคณฺหาติ อุปตฺถมฺเภติ ปจฺจุปตฺถมฺเภติ, เตน วุจฺจติ “จิตฺตํ ปคฺคณฺหาตี”ติ.}

\proseitem{414}{“ปทหตี”ติ ตตฺถ กตมํ สมฺมปฺปธานํ, โย เจตสิโก วีริยารมฺโภ ฯเปฯ สมฺมาวายาโม วีริยสมฺโพชฺฌงฺโค มคฺคงฺคํ มคฺคปริยาปนฺนํ, อิทํ วุจฺจติ สมฺมปฺปธานํ.\csromanspacer{} อวเสสา ธมฺมา สมฺมปฺปธานสมฺปยุตฺตา.}

\proseitem{415}{กถญฺจ ภิกฺขุ อุปฺปนฺนานํ ปาปกานํ อกุสลานํ ธมฺมานํ ปหานาย ฉนฺทํ ชเนติ วายมติ วีริยํ อารภติ จิตฺตํ ปคฺคณฺหาติ ปทหติ, อิธ ภิกฺขุ ยสฺมิํ สมเย โลกุตฺตรํ ฌานํ ภาเวติ นิยฺยานิกํ อปจยคามิํ ทิฏฺฐิคตานํ ปหานาย ปฐมาย ภูมิยา ปตฺติยา วิวิจฺเจว กาเมหิ ฯเปฯ ปฐมํ ฌานํ อุปสมฺปชฺช วิหรติ ทุกฺขปฏิปทํ ทนฺธาภิญฺญํ, ตสฺมิํ}

\vspace{\tipitakaparskip}
\csromancenter{สมฺมปฺปธานวิภงฺค สมเย อุปฺปานฺนานํ ปาปกานํ อกุสลานํ ธมฺมานํ ปหานาย ฉนฺทํ ชเนติ วายมติ วีริยํ อารภติ จิตฺตํ ปคฺคณฺหาติ ปทหติ.}
\proseitem{416}{\csromanfolio{221}“ฉนฺทํ ชเนตี”ติ ฯเปฯ “วายมตี”ติ ฯเปฯ “วีริยํ อารภตี”ติ ฯเปฯ “จิตฺตํ ปคฺคณฺหาตี”ติ ฯเปฯ “ปทหตี”ติ ตตฺถ กตมํ สมฺมปฺปธานํ, โย เจตสิโก วีริยารมฺโภ ฯเปฯ สมฺมาวายาโม วีริยสมฺโพชฺฌงฺโค มคฺคงฺคํ มคฺคปริยาปนฺนํ, อิทํ วุจฺจติ สมฺมปฺปธานํ.\csromanspacer{} อวเสสา ธมฺมา สมฺมปฺปธานสมฺปยุตฺตา.}

\proseitem{417}{กถญฺจ ภิกฺขุ อนุปฺปนฺนานํ กุสลานํ ธมฺมานํ อุปฺปาทาย ฉนฺทํ ชเนติ วายมติ วีริยํ อารภติ จิตฺตํ ปคฺคณฺหาติ ปทหติ, อิธ ภิกฺขุ ยสฺมิํ สมเย โลกุตฺตรํ ฌานํ ภาเวติ นิยฺยานิกํ อปจยคามิํ ทิฏฺฐิคตานํ ปหานาย ปฐมาย ภูมิยา ปตฺติยา วิวิจฺเจว กาเมหิ ฯเปฯ ปฐมํ ฌานํ อุปสมฺปชฺช วิหรติ ทุกฺขปฏิปทํ ทนฺธาภิญฺญํ, ตสฺมิํ สมเย อนุปฺปนฺนานํ กุสลานํ ธมฺมานํ อุปฺปาทาย ฉนฺทํ ชเนติ วายมติ วีริยํ อารภติ จิตฺตํ ปคฺคณฺหาติ ปทหติ.}

\proseitem{418}{“ฉนฺทํ ชเนตี”ติ ฯเปฯ “วายมตี”ติ ฯเปฯ “วีริยํ อารภตี”ติ ฯเปฯ “จิตฺตํ ปคฺคณฺหาตี”ติ ฯเปฯ “ปทหตี”ติ ตตฺถ กตมํ สมฺมปฺปธานํ, โย เจตสิโก วีริยารมฺโภ ฯเปฯ สมฺมาวายาโม วีริยสมฺโพชฺฌงฺโค มคฺคงฺคํ มคฺคปริยาปนฺนํ, อิทํ วุจฺจติ สมฺมปฺปธานํ.\csromanspacer{} อวเสสา ธมฺมา สมฺมปฺปธานสมฺปยุตฺตา.}

\proseitem{419}{กถญฺจ ภิกฺขุ อุปฺปนฺนานํ กุสลานํ ธมฺมานํ ฐิติยา อสมฺโมสาย ภิยฺโยภาวาย เวปุลฺลาย ภาวนาย ปาริปูริยา ฉนฺทํ ชเนติ วายมติ วีริยํ อารภติ จิตฺตํ ปคฺคณฺหาติ ปทหติ, อิธ ภิกฺขุ ยสฺมิํ สมเย โลกุตฺตรํ ฌานํ ภาเวติ นิยฺยานิกํ อปจยคามิํ ทิฏฺฐิคตานํ ปหานาย ปฐมาย ภูมิยา ปตฺติยา วิวิจฺเจว กาเมหิ ฯเปฯ ปฐมํ ฌานํ อุปสมฺปชฺช วิหรติ ทุกฺขปฏิปทํ ทนฺธาภิญฺญํ, ตสฺมิํ สมเย อุปฺปนฺนานํ กุสลานํ ธมฺมานํ ฐิติยา อสมฺโมสาย ภิยฺโยภาวาย เวปุลฺลาย ภาวนาย ปาริปูริยา ฉนฺทํ ชเนติ วายมติ วีริยํ อารภติ จิตฺตํ ปคฺคณฺหาติ ปทหติ.}

\proseitem{420}{\csromanfolio{222}“ฐิติยา”ติ ยา ฐิติ โส อสมฺโมโส, โย อสมฺโมโส โส ภิยฺโยภาโว, โย ภิยฺโยภาโว ตํ เวปุลฺลํ, ยํ เวปุลฺลํ สา ภาวนา, ยา ภาวนา สา ปาริปูรี.}

\proseitem{421}{“ฉนฺทํ ชเนตี”ติ ตตฺถ กตโม ฉนฺโท, โย ฉนฺโท ฉนฺทิกตา กตฺตุกมฺยตา กุสโล ธมฺมจฺฉนฺโท, อยํ วุจฺจติ ฉนฺโท.\csromanspacer{} อิมํ ฉนฺทํ ชเนติ สญฺชเนติ อุฏฺฐเปติ สมุฏฺฐเปติ นิพฺพตฺเตติ อภินิพฺพตฺเตติ, เตน วุจฺจติ “ฉนฺทํ ชเนตี”ติ.}

\proseitem{422}{“วายมตี”ติ ตตฺถ กตโม วายาโม, โย เจตสิโก วีริยารมฺโภ ฯเปฯ สมฺมาวายาโม วีริยสมฺโพชฺฌงฺโค มคฺคงฺคํ มคฺคปริยาปนฺนํ, อยํ วุจฺจติ วายาโม.\csromanspacer{} อิมินา วายาเมน อุเปโต โหติ ฯเปฯ สมนฺนาคโต, เตน วุจฺจติ “วายมตี”ติ.}

\proseitem{423}{“วีริยํ อารภตี”ติ ตตฺถ กตมํ วีริยํ, โย เจตสิโก วีริยารมฺโภ ฯเปฯ สมฺมาวายาโม วีริยสมฺโพชฺฌงฺโค มคฺคงฺคํ มคฺคปริยาปนฺนํ, อิทํ วุจฺจติ วีริยํ.\csromanspacer{} อิมํ วีริยํ อารภติ สมารภติ อาเสวติ ภาเวติ พหุลีกโรติ, เตน วุจฺจติ “วีริยํ อารภตี”ติ.}

\proseitem{424}{“จิตฺตํ ปคฺคณฺหาตี”ติ ตตฺถ กตมํ จิตฺตํ, ยํ จิตฺตํ มโน มานสํ ฯเปฯ ตชฺชามโนวิญฺญาณธาตุ, อิทํ วุจฺจติ จิตฺตํ.\csromanspacer{} อิมํ จิตฺตํ ปคฺคณฺหาติ สมฺปคฺคณฺหาติ อุปตฺถมฺเภติ ปจฺจุปตฺถมฺเภติ, เตน วุจฺจติ “จิตฺตํ ปคฺคณฺหาตี”ติ.}

\proseitem{425}{“ปทหตี”ติ ตตฺถ กตมํ สมฺมปฺปธานํ, โย เจตสิโก วีริยารมฺโภ ฯเปฯ สมฺมาวายาโม วีริยสมฺโพชฺฌงฺโค มคฺคงฺคํ มคฺคปริยาปนฺนํ, อิทํ วุจฺจติ สมฺมปฺปธานํ.\csromanspacer{} อวเสสา ธมฺมา สมฺมปฺปธานสมฺปยุตฺตา.}

\proseitem{426}{ตตฺถ กตมํ สมฺมปฺปธานํ, อิธ ภิกฺขุ ยสฺมิํ สมเย โลกุตฺตรํ ฌานํ ภาเวติ นิยฺยานิกํ อปจยคามิํ ทิฏฺฐิคตานํ ปหานาย ปฐมาย ภูมิยา ปตฺติยา วิวิจฺเจว กาเมหิ ฯเปฯ ปฐมํ ฌานํ อุปสมฺปชฺช วิหรติ ทุกฺขปฏิปทํ ทนฺธาภิญฺญํ, โย ตสฺมิํ สมเย เจตสิโก วีริยารมฺโภ ฯเปฯ สมฺมาวายาโม วีริยสมฺโพชฺฌงฺโค มคฺคงฺคํ มคฺคปริยาปนฺนํ, อิทํ วุจฺจติ สมฺมปฺปธานํ.\csromanspacer{} อวเสสา ธมฺมา สมฺมปฺปธานสมฺปยุตฺตา.}

\vspace{\tipitakaparskip}
\csromancenter{อภิธมฺมภาชนียํ.}
\csromanmatikaanchor{mtk.74}
\csromantocmarkref{subparagraph}{3. ปญฺหาปุจฺฉก}{mtk.74}
\bookmark[dest={mtk.74},level=5]{3. ปญฺหาปุจฺฉก}
\csromanheader{6}{8. สมฺมปฺปธานวิภงฺค \textbf{3. ปญฺหาปุจฺฉก}}
\proseitem{427}{\csromanfolio{223}จตฺตาโร สมฺมปฺปธานา\csromandash{}อิธ ภิกฺขุ อนุปฺปนฺนานํ ปาปกานํ อกุสลานํ ธมฺมานํ อนุปฺปาทาย ฉนฺทํ ชเนติ วายมติ วีริยํ อารภติ จิตฺตํ ปคฺคณฺหาติ ปทหติ, อุปฺปนฺนานํ ปาปกานํ อกุสลานํ ธมฺมานํ ปหานาย ฯเปฯ อนุปฺปนฺนานํ กุสลานํ ธมฺมานํ อุปฺปาทาย ฯเปฯ อุปฺปนฺนานํ กุสลานํ ธมฺมานํ ฐิติยา อสมฺโมสาย ภิยฺโยภาวาย เวปุลฺลาย ภาวนาย ปาริปูริยา ฉนฺทํ ชเนติ วายมติ วีริยํ อารภติ จิตฺตํ ปคฺคณฺหาติ ปทหติ.}

\proseitem{428}{จตุนฺนํ สมฺมปฺปธานานํ กติ กุสลา, กติ อกุสลา, กติ อพฺยากตา ฯเปฯ กติ สรณา, กติ อรณา.}

\csromanmatikaanchor{mtk.75}
\csromantocmarkref{subparagraph}{1. ติก}{mtk.75}
\bookmark[dest={mtk.75},level=5]{1. ติก}
\csromanheader{6}{1. ติก}
\proseitem{429}{กุสลาเยว.\csromanspacer{} สิยา สุขาย เวทนาย สมฺปยุตฺตา, สิยา อทุกฺขมสุขาย เวทนาย สมฺปยุตฺตา.\csromanspacer{} วิปากธมฺมธมฺมา.\csromanspacer{} อนุปาทินฺน-อนุปาทานิยา.\csromanspacer{} อสํกิลิฏฺฐ-อสํกิเลสิกา.\csromanspacer{} สิยา สวิตกฺกสวิจารา, สิยา อวิตกฺกวิจารมตฺตา, สิยา อวิตกฺก-อวิจารา.\csromanspacer{} สิยา ปีติสหคตา, สิยา สุขสหคตา, สิยา อุเปกฺขาสหคตา.\csromanspacer{} เนว ทสฺสเนน น ภาวนาย ปหาตพฺพา, เนว ทสฺสเนน น ภาวนาย ปหาตพฺพเหตุกา.\csromanspacer{} อปจยคามิโน.\csromanspacer{} เสกฺขา.\csromanspacer{} อปฺปมาณา.\csromanspacer{} อปฺปมาณารมฺมณา.\csromanspacer{} ปณีตา.\csromanspacer{} สมฺมตฺตนิยตา.\csromanspacer{} น มคฺคารมฺมณา.\csromanspacer{} มคฺคเหตุกา.\csromanspacer{} สิยา มคฺคาธิปติโน, สิยา น วตฺตพฺพา “มคฺคาธิปติโน”ติ.\csromanspacer{} สิยา อุปฺปนฺนา สิยา อนุปฺปนฺนา, น วตฺตพฺพา “อุปฺปาทิโน”ติ.\csromanspacer{} สิยา อตีตา, สิยา อนาคตา, สิยา ปจฺจุปฺปนฺนา, น วตฺตพฺพา “อตีตารมฺมณา”ติปิ, “อนาคตารมฺมณา”ติปิ, “ปจฺจุปฺปนฺนารมฺมณา”ติปิ.\csromanspacer{} สิยา อชฺฌตฺตา, สิยา พหิทฺธา, สิยา อชฺฌตฺตพหิทฺธา.\csromanspacer{} พหิทฺธารมฺมณา.\csromanspacer{} อนิทสฺสน-อปฺปฏิฆา.}

\csromanmatikaanchor{mtk.76}
\csromantocmarkref{subparagraph}{2. ทุก}{mtk.76}
\bookmark[dest={mtk.76},level=5]{2. ทุก}
\csromanheader{6}{2. ทุก}
\proseitem{430}{น เหตู, สเหตุกา.\csromanspacer{} เหตุสมฺปยุตฺตา.\csromanspacer{} น วตฺตพฺพา “เหตู เจว สเหตุกา จา”ติ.\csromanspacer{} สเหตุกา เจว น จ เหตู.\csromanspacer{} น วตฺตพฺพา “เหตู เจว เหตุสมฺปยุตฺตา จา”ติ.\csromanspacer{} เหตุสมฺปยุตฺตา เจว น จ เหตู. \csromanfolio{224}น เหตู, สเหตุกา.\csromanspacer{} สปฺปจฺจยา, สงฺขตา, อนิทสฺสนา, อปฺปฏิฆา, อรูปา, โลกุตฺตรา, เกนจิ วิญฺเญยฺยา, เกนจิ น วิญฺเญยฺยา.\csromanspacer{} โน อาสวา.\csromanspacer{} อนาสวา.\csromanspacer{} อาสววิปฺปยุตฺตา.\csromanspacer{} น วตฺตพฺพา “อาสวา เจว สาสวา จา”ติปิ, “สาสวา เจว โน จ อาสวา”ติปิ.\csromanspacer{} น วตฺตพฺพา “อาสวา เจว อาสวสมฺปยุตฺตา จา”ติปิ, “อาสวสมฺปยุตฺตา เจว โน จ อาสวา”ติปิ.\csromanspacer{} อาสววิปฺปยุตฺตา อนาสวา.\csromanspacer{} โน สํโยชนา ฯเปฯ.\csromanspacer{} โน คนฺถา ฯเปฯ.\csromanspacer{} โน โอฆา ฯเปฯ.\csromanspacer{} โน โยคา ฯเปฯ.\csromanspacer{} โน นีวรณา ฯเปฯ.\csromanspacer{} โน ปรามาสา ฯเปฯ สารมฺมณา, โน จิตฺตา, เจตสิกา, จิตฺตสมฺปยุตฺตา, จิตฺตสํสฏฺฐา, จิตฺตสมุฏฺฐานา, จิตฺตสหภุโน, จิตฺตานุปริวตฺติโน, จิตฺตสํสฏฺฐสมุฏฺฐานา, จิตฺตสํสฏฺฐสมุฏฺฐานสหภุโน, จิตฺตสํสฏฺฐสมุฏฺฐานานุปริวตฺติโน, พาหิรา.\csromanspacer{} โน อุปาทา, อนุปาทินฺนา, โน อุปาทานา ฯเปฯ.\csromanspacer{} โน กิเลสา ฯเปฯ.\csromanspacer{} น ทสฺสเนน ปหาตพฺพา, น ภาวนาย ปหาตพฺพา.\csromanspacer{} น ทสฺสเนน ปหาตพฺพเหตุกา, น ภาวนาย ปหาตพฺพเหตุกา.\csromanspacer{} สิยา สวิตกฺกา, สิยา อวิตกฺกา.\csromanspacer{} สิยา สวิจารา, สิยา อวิจารา.\csromanspacer{} สิยา สปฺปีติกา, สิยา อปฺปีติกา.\csromanspacer{} สิยา ปีติสหคตา, สิยา น ปีติสหคตา.\csromanspacer{} สิยา สุขสหคตา, สิยา น สุขสหคตา.\csromanspacer{} สิยา อุเปกฺขาสหคตา, สิยา น อุเปกฺขาสหคตา.\csromanspacer{} น กามาวจรา.\csromanspacer{} น รูปาวจรา.\csromanspacer{} น อรูปาวจรา.\csromanspacer{} อปริยาปนฺนา.\csromanspacer{} นิยฺยานิกา.\csromanspacer{} นิยตา.\csromanspacer{} อนุตฺตรา.\csromanspacer{} อรณาติ.}

\vspace{\tipitakaparskip}
\csromancenter{ปญฺหาปุจฺฉกํ.}
\nitthitam{\textbf{สมฺมปฺปธานวิภงฺโค นิฏฺฐิโต.}}
\csromanmatikaanchor{mtk.77}
\csromantocmarknopage{subparagraph}{9. อิทฺธิปาทวิภงฺค}{mtk.77}
\bookmark[dest={mtk.77},level=5]{9. อิทฺธิปาทวิภงฺค}
\csromanheader{6}{9. อิทฺธิปาทวิภงฺค}
\csromanheader{2}{1. สุตฺตนฺตภาชนีย}
\proseitem{431}{\csromanfolio{225}จตฺตาโร อิทฺธิปาทา\csromandash{}อิธ ภิกฺขุ ฉนฺทสมาธิปธานสงฺขารสมนฺนาคตํ อิทฺธิปาทํ ภาเวติ, วีริยสมาธิปธานสงฺขารสมนฺนาคตํ อิทฺธิปาทํ ภาเวติ, จิตฺตสมาธิปธานสงฺขารสมนฺนาคตํ อิทฺธิปาทํ ภาเวติ, วีมํสาสมาธิปธานสงฺขารสมนฺนาคตํ อิทฺธิปาทํ ภาเวติ.}

\csromanmatikaanchor{mtk.78}
\csromanmatikaanchor{mtk.79}
\csromantocmarknopage{subparagraph}{1. สุตฺตนฺตภาชนีย}{mtk.78}
\bookmark[dest={mtk.78},level=5]{1. สุตฺตนฺตภาชนีย}
\csromantocmarkref{subparagraph}{1. ฉนฺทิทฺธิปาท}{mtk.79}
\bookmark[dest={mtk.79},level=5]{1. ฉนฺทิทฺธิปาท}
\csromanheader{6}{1. ฉนฺทิทฺธิปาท}
\proseitem{432}{กถญฺจ ภิกฺขุ ฉนฺทสมาธิปธานสงฺขารสมนฺนาคตํ อิทฺธิปาทํ ภาเวติ, ฉนฺทํ เจ ภิกฺขุ อธิปติํ กริตฺวา ลภติ สมาธิํ, ลภติ จิตฺตสฺเสกคฺคตํ\footnote{จิตฺตสฺส เอกคฺคตํ (สี, สฺยา)}, อยํ วุจฺจติ ฉนฺทสมาธิ.\csromanspacer{} โส อนุปฺปนฺนานํ ปาปกานํ อกุสลานํ ธมฺมานํ อนุปฺปาทาย ฉนฺทํ ชเนติ วายมติ วีริยํ อารภติ จิตฺตํ ปคฺคณฺหาติ ปทหติ, อุปฺปนฺนานํ ปาปกานํ อกุสลานํ ธมฺมานํ ปหานาย ฉนฺทํ ชเนติ วายมติ วีริยํ อารภติ จิตฺตํ ปคฺคณฺหาติ ปทหติ, อนุปฺปนฺนานํ กุสลานํ ธมฺมานํ อุปฺปาทาย ฉนฺทํ ชเนติ วายมติ วีริยํ อารภติ จิตฺตํ ปคฺคณฺหาติ ปทหติ, อุปฺปนฺนานํ กุสลานํ ธมฺมานํ ฐิติยา อสมฺโมสาย ภิยฺโยภาวาย เวปุลฺลาย ภาวนาย ปาริปูริยา ฉนฺทํ ชเนติ วายมติ วีริยํ อารภติ จิตฺตํ ปคฺคณฺหาติ ปทหติ, อิเม วุจฺจนฺติ ปธานสงฺขารา.\csromanspacer{} อิติ อยญฺจ ฉนฺทสมาธิ อิเม จ ปธานสงฺขารา ตเทกชฺฌํ อภิสญฺญูหิตฺวา อภิสงฺขิปิตฺวา ฉนฺทสมาธิปธานสงฺขาโรเตฺวว สงฺขฺยํ\footnote{สงฺขํ (สี)} คจฺฉติ.}

\proseitem{433}{ตตฺถ กตโม ฉนฺโท, โย ฉนฺโท ฉนฺทิกตา กตฺตุกมฺยตา กุสโล ธมฺมจฺฉนฺโท, อยํ วุจฺจติ ฉนฺโท.}

\prose{ตตฺถ กตโม สมาธิ, ยา จิตฺตสฺส ฐิติ สณฺฐิติ อวฏฺฐิติ อวิสาหาโร อวิกฺเขโป อวิสาหฏมานสตา สมโถ สมาธินฺทฺริยํ สมาธิพลํ สมฺมาสมาธิ, อยํ วุจฺจติ สมาธิ.}

\prose{ตตฺถ กตโม ปธานสงฺขาโร, โย เจตสิโก วีริยารมฺโภ นิกฺกโม ปรกฺกโม อุยฺยาโม วายาโม อุสฺสาโห อุสฺโสฬฺหี ถาโม ฐิติ อสิถิลปรกฺกมตา อนิกฺขิตฺตฉนฺทตา อนิกฺขิตฺตธุรตา \csromanfolio{226}ธุรสมฺปคฺคาโห วีริยํ วีริยินฺทฺริยํ วีริยพลํ สมฺมาวายาโม, อยํ วุจฺจติ ปธานสงฺขาโร.\csromanspacer{} อิติ อิมินา จ ฉนฺเทน อิมินา จ สมาธินา อิมินา จ ปธานสงฺขาเรน อุเปโต โหติ สมุเปโต อุปาคโต สมุปาคโต อุปปนฺโน สมฺปนฺโน สมนฺนาคโต, เตน วุจฺจติ “ฉนฺทสมาธิปธานสงฺขารสมนฺนาคโต”ติ.}

\proseitem{434}{“อิทฺธี”ติ ยา เตสํ ธมฺมานํ อิทฺธิ สมิทฺธิ อิชฺฌนา สมิชฺฌนา ลาโภ ปฏิลาโภ ปตฺติ สมฺปตฺติ ผุสนา สจฺฉิกิริยา อุปสมฺปทา.}

\prose{“อิทฺธิปาโท”ติ ตถาภูตสฺส เวทนากฺขนฺโธ สญฺญากฺขนฺโธ สงฺขารกฺขนฺโธ วิญฺญาณกฺขนฺโธ.}

\prose{“อิทฺธิปาทํ ภาเวตี”ติ เต ธมฺเม อาเสวติ ภาเวติ พหุลีกโรติ, เตน วุจฺจติ “อิทฺธิปาทํ ภาเวตี”ติ.}

\csromanmatikaanchor{mtk.80}
\csromantocmarkref{subparagraph}{2. วีริยิทฺธิปาท}{mtk.80}
\bookmark[dest={mtk.80},level=5]{2. วีริยิทฺธิปาท}
\csromanheader{6}{2. วีริยิทฺธิปาท}
\proseitem{435}{กถญฺจ ภิกฺขุ วีริยสมาธิปธานสงฺขารสมนฺนาคตํ อิทฺธิปาทํ ภาเวติ, วีริยํ เจ ภิกฺขุ อธิปติํ กริตฺวา ลภติ สมาธิํ, ลภติ จิตฺตสฺเสกคฺคตํ, อยํ วุจฺจติ วีริยสมาธิ, โส อนุปฺปนฺนานํ ปาปกานํ อกุสลานํ ธมฺมานํ อนุปฺปาทาย ฉนฺทํ ชเนติ วายมติ วีริยํ อารภติ จิตฺตํ ปคฺคณฺหาติ ปทหติ, อุปฺปนฺนานํ ปาปกานํ อกุสลานํ ธมฺมานํ ปหานาย ฯเปฯ อนุปฺปนฺนานํ กุสลานํ ธมฺมานํ อุปฺปาทาย ฯเปฯ อุปฺปนฺนานํ กุสลานํ ธมฺมานํ ฐิติยา อสมฺโมสาย ภิยฺโยภาวาย เวปุลฺลาย ภาวนาย ปาริปูริยา ฉนฺทํ ชเนติ วายมติ วีริยํ อารภติ จิตฺตํ ปคฺคณฺหาติ ปทหติ, อิเม วุจฺจนฺติ ปธานสงฺขารา.\csromanspacer{} อิติ อยญฺจ วีริยสมาธิ อิเม จ ปธานสงฺขารา ตเทกชฺฌํ อภิสญฺญูหิตฺวา อภิสงฺขิปิตฺวา วีริยสมาธิปธานสงฺขาโรเตฺวว สงฺขฺยํ คจฺฉติ.}

\proseitem{436}{ตตฺถ กตมํ วีริยํ, โย เจตสิโก วีริยารมฺโภ ฯเปฯ สมฺมาวายาโม, อิทํ วุจฺจติ วีริยํ.}

\prose{ตตฺถ กตโม สมาธิ, ยา จิตฺตสฺส ฐิติ สณฺฐิติ อวฏฺฐิติ อวิสาหาโร อวิกฺเขโป อวิสาหฏมานสตา สมโถ สมาธินฺทฺริยํ สมาธิพลํ สมฺมาสมาธิ, อยํ วุจฺจติ สมาธิ.}

\csromanheader{2}{9. อิทฺธิปาทวิภงฺค}
\prose{\csromanfolio{227}ตตฺถ กตโม ปธานสงฺขาโร, โย เจตสิโก วีริยารมฺโภ ฯเปฯ สมฺมาวายาโม, อยํ วุจฺจติ ปธานสงฺขาโร.\csromanspacer{} อิติ อิมินา จ วีริเยน อิมินา จ สมาธินา อิมินา จ ปธานสงฺขาเรน อุเปโต โหติ ฯเปฯ สมนฺนาคโต, เตน วุจฺจติ “วีริยสมาธิปธานสงฺขารสมนฺนาคโต”ติ.}

\proseitem{437}{“อิทฺธี”ติ ยา เตสํ ธมฺมานํ อิทฺธิ สมิทฺธิ อิชฺฌนา สมิชฺฌนา ลาโภ ปฏิลาโภ ปตฺติ สมฺปตฺติ ผุสนา สจฺฉิกิริยา อุปสมฺปทา.}

\prose{“อิทฺธิปาโท”ติ ตถาภูตสฺส เวทนากฺขนฺโธ ฯเปฯ วิญฺญาณกฺขนฺโธ.}

\prose{“อิทฺธิปาทํ ภาเวตี”ติ เต ธมฺเม อาเสวติ ภาเวติ พหุลีกโรติ, เตน วุจฺจติ “อิทฺธิปาทํ ภาเวตี”ติ.}

\csromanmatikaanchor{mtk.81}
\csromantocmarkref{subparagraph}{3. จิตฺติทฺธิปาท}{mtk.81}
\bookmark[dest={mtk.81},level=5]{3. จิตฺติทฺธิปาท}
\csromanheader{6}{3. จิตฺติทฺธิปาท}
\proseitem{438}{กถญฺจ ภิกฺขุ จิตฺตสมาธิปธานสงฺขารสมนฺนาคตํ อิทฺธิปาทํ ภาเวติ, จิตฺตํ เจ ภิกฺขุ อธิปติํ กริตฺวา ลภติ สมาธิํ, ลภติ จิตฺตสฺเสกคฺคตํ, อยํ วุจฺจติ จิตฺตสมาธิ.\csromanspacer{} โส อนุปฺปนฺนานํ ปาปกานํ อกุสลานํ ธมฺมานํ อนุปฺปาทาย ฉนฺทํ ชเนติ วายมติ วีริยํ อารภติ จิตฺตํ ปคฺคณฺหาติ ปทหติ, อุปฺปนฺนานํ ปาปกานํ อกุสลานํ ธมฺมานํ ปหานาย ฯเปฯ อนุปฺปนฺนานํ กุสลานํ ธมฺมานํ อุปฺปาทาย ฯเปฯ อุปฺปนฺนานํ กุสลานํ ธมฺมานํ ฐิติยา อสมฺโมสาย ภิยฺโยภาวาย เวปุลฺลาย ภาวนาย ปาริปูริยา ฉนฺทํ ชเนติ วายมติ วีริยํ อารภติ จิตฺตํ ปคฺคณฺหาติ ปทหติ, อิเม วุจฺจนฺติ ปธานสงฺขารา.\csromanspacer{} อิติ อยญฺจ จิตฺตสมาธิ อิเม จ ปธานสงฺขารา ตเทกชฺฌํ อภิสญฺญูหิตฺวา อภิสงฺขิปิตฺวา จิตฺตสมาธิปธานสงฺขาโรเตฺวว สงฺขฺยํ คจฺฉติ.}

\proseitem{439}{ตตฺถ กตมํ จิตฺตํ, ยํ จิตฺตํ มโน มานสํ ฯเปฯ ตชฺชามโนวิญฺญาณธาตุ, อิทํ วุจฺจติ จิตฺตํ.}

\prose{ตตฺถ กตโม สมาธิ, ยา จิตฺตสฺส ฐิติ สณฺฐิติ ฯเปฯ สมฺมาสมาธิ, อยํ วุจฺจติ สมาธิ.}

\prose{ตตฺถ กตโม ปธานสงฺขาโร, โย เจตสิโก วีริยารมฺโภ ฯเปฯ สมฺมาวายาโม, อยํ วุจฺจติ ปธานสงฺขาโร.\csromanspacer{} อิติ อิมินา จ จิตฺเตน อิมินา จ สมาธินา อิมินา จ ปธานสงฺขาเรน อุเปโต โหติ ฯเปฯ สมนฺนาคโต, เตน วุจฺจติ “จิตฺตสมาธิปธานสงฺขารสมนฺนาคโต”ติ.}

\proseitem{440}{\csromanfolio{228}“อิทฺธี”ติ ยา เตสํ ธมฺมานํ อิทฺธิ สมิทฺธิ อิชฺฌนา สมิชฺฌนา ลาโภ ปฏิลาโภ ปตฺติ สมฺปตฺติ ผุสนา สจฺฉิกิริยา อุปสมฺปทา.}

\prose{“อิทฺธิปาโท”ติ ตถาภูตสฺส เวทนากฺขนฺโธ ฯเปฯ วิญฺญาณกฺขนฺโธ.}

\prose{“อิทฺธิปาทํ ภาเวตี”ติ เต ธมฺเม อาเสวติ ภาเวติ พหุลีกโรติ, เตน วุจฺจติ “อิทฺธิปาทํ ภาเวตี”ติ.}

\csromanmatikaanchor{mtk.82}
\csromantocmarkref{subparagraph}{4. วีมํสิทฺธิปาท}{mtk.82}
\bookmark[dest={mtk.82},level=5]{4. วีมํสิทฺธิปาท}
\csromanheader{6}{4. วีมํสิทฺธิปาท}
\proseitem{441}{กถญฺจ ภิกฺขุ วีมํสาสมาธิปธานสงฺขารสมนฺนาคตํ อิทฺธิปาทํ ภาเวติ, วีมํสํ เจ ภิกฺขุ อธิปติํ กริตฺวา ลภติ สมาธิํ, ลภติ จิตฺตสฺเสกคฺคตํ, อยํ วุจฺจติ วีมํสาสมาธิ.\csromanspacer{} โส อนุปฺปนฺนานํ ปาปกานํ อกุสลานํ ธมฺมานํ อนุปฺปาทาย ฉนฺทํ ชเนติ วายมติ วีริยํ อารภติ จิตฺตํ ปคฺคณฺหาติ ปทหติ, อุปฺปนฺนานํ ปาปกานํ อกุสลานํ ธมฺมานํ ปหานาย ฯเปฯ อนุปฺปนฺนานํ กุสลานํ ธมฺมานํ อุปฺปาทาย ฯเปฯ อุปฺปนฺนานํ กุสลานํ ธมฺมานํ ฐิติยา อสมฺโมสาย ภิยฺโยภาวาย เวปุลฺลาย ภาวนาย ปาริปูริยา ฉนฺทํ ชเนติ วายมติ วีริยํ อารภติ จิตฺตํ ปคฺคณฺหาติ ปทหติ, อิเม วุจฺจนฺติ ปธานสงฺขารา.\csromanspacer{} อิติ อยญฺจ วีมํสาสมาธิ อิเม จ ปธานสงฺขารา ตเทกชฺฌํ อภิสญฺญูหิตฺวา อภิสงฺขิปิตฺวา วีมํสาสมาธิปธานสงฺขาโรเตฺวว สงฺขฺยํ คจฺฉติ.}

\proseitem{442}{ตตฺถ กตมา วีมํสา, ยา ปญฺญา ปชานนา ฯเปฯ อโมโห ธมฺมวิจโย สมฺมาทิฏฺฐิ, อยํ วุจฺจติ วีมํสา.}

\prose{ตตฺถ กตโม สมาธิ, ยา จิตฺตสฺส ฐิติ สณฺฐิติ ฯเปฯ สมฺมาสมาธิ, อยํ วุจฺจติ สมาธิ.}

\prose{ตตฺถ กตโม ปธานสงฺขาโร, โย เจตสิโก วีริยารมฺโภ ฯเปฯ สมฺมาวายาโม, อยํ วุจฺจติ ปธานสงฺขาโร.\csromanspacer{} อิติ อิมาย จ วีมํสาย อิมินา จ สมาธินา อิมินา จ ปธานสงฺขาเรน อุเปโต โหติ ฯเปฯ สมนฺนาคโต, เตน วุจฺจติ “วีมํสาสมาธิปธานสงฺขารสมนฺนาคโต”ติ.}

\proseitem{443}{“อิทฺธี”ติ ยา เตสํ ธมฺมานํ อิทฺธิ สมิทฺธิ อิชฺฌนา สมิชฺฌนา ลาโภ ปฏิลาโภ ปตฺติ สมฺปตฺติ ผุสนา สจฺฉิกิริยา อุปสมฺปทา.}

\csromanheader{2}{9. อิทฺธิปาทวิภงฺค}
\prose{\csromanfolio{229}“อิทฺธิปาโท”ติ ตถาภูตสฺส เวทนากฺขนฺโธ สญฺญากฺขนฺโธ สงฺขารกฺขนฺโธ วิญฺญาณกฺขนฺโธ.}

\prose{“อิทฺธิปาทํ ภาเวตี”ติ เต ธมฺเม อาเสวติ ภาเวติ พหุลีกโรติ, เตน วุจฺจติ “อิทฺธิปาทํ ภาเวตี”ติ.}

\csromanheader{2}{สุตฺตนฺตภาชนียํ.}
\csromansectionrule
\csromanheader{2}{2. อภิธมฺมภาชนีย}
\proseitem{444}{จตฺตาโร อิทฺธิปาทา\csromandash{}อิธ ภิกฺขุ ฉนฺทสมาธิปธานสงฺขารสมนฺนาคตํ อิทฺธิปาทํ ภาเวติ, วีริยสมาธิปธานสงฺขารสมนฺนาคตํ อิทฺธิปาทํ ภาเวติ, จิตฺตสมาธิปธานสงฺขารสมนฺนาคตํ อิทฺธิปาทํ ภาเวติ, วีมํสาสมาธิปธานสงฺขารสมนฺนาคตํ อิทฺธิปาทํ ภาเวติ.}

\csromanmatikaanchor{mtk.83}
\csromanmatikaanchor{mtk.84}
\csromantocmarknopage{subparagraph}{2. อภิธมฺมภาชนีย}{mtk.83}
\bookmark[dest={mtk.83},level=5]{2. อภิธมฺมภาชนีย}
\csromantocmarkref{subparagraph}{1. ฉนฺทิทฺธิปาท}{mtk.84}
\bookmark[dest={mtk.84},level=5]{1. ฉนฺทิทฺธิปาท}
\csromanheader{6}{1. ฉนฺทิทฺธิปาท}
\proseitem{445}{กถญฺจ ภิกฺขุ ฉนฺทสมาธิปธานสงฺขารสมนฺนาคตํ อิทฺธิปาทํ ภาเวติ, อิธ ภิกฺขุ ยสฺมิํ สมเย โลกุตฺตรํ ฌานํ ภาเวติ นิยฺยานิกํ อปจยคามิํ ทิฏฺฐิคตานํ ปหานาย ปฐมาย ภูมิยา ปตฺติยา วิวิจฺเจว กาเมหิ ฯเปฯ ปฐมํ ฌานํ อุปสมฺปชฺช วิหรติ ทุกฺขปฏิปทํ ทนฺธาภิญฺญํ, ตสฺมิํ สมเย ฉนฺทสมาธิปธานสงฺขารสมนฺนาคตํ อิทฺธิปาทํ ภาเวติ.}

\proseitem{446}{ตตฺถ กตโม ฉนฺโท, โย ฉนฺโท ฉนฺทิกตา กตฺตุกมฺยตา กุสโล ธมฺมจฺฉนฺโท, อยํ วุจฺจติ ฉนฺโท.}

\prose{ตตฺถ กตโม สมาธิ, ยา จิตฺตสฺส ฐิติ สณฺฐิติ ฯเปฯ สมฺมาสมาธิ สมาธิสมฺโพชฺฌงฺโค มคฺคงฺคํ มคฺคปริยาปนฺนํ, อยํ วุจฺจติ สมาธิ.}

\prose{ตตฺถ กตโม ปธานสงฺขาโร, โย เจตสิโก วีริยารมฺโภ ฯเปฯ สมฺมาวายาโม วีริยสมฺโพชฺฌงฺโค มคฺคงฺคํ มคฺคปริยาปนฺนํ, อยํ วุจฺจติ ปธานสงฺขาโร.\csromanspacer{} อิติ อิมินา จ ฉนฺเทน อิมินา จ สมาธินา อิมินา จ ปธานสงฺขาเรน อุเปโต โหติ ฯเปฯ สมนฺนาคโต, เตน วุจฺจติ “ฉนฺทสมาธิปธานสงฺขารสมนฺนาคโต”ติ.}

\proseitem{447}{\csromanfolio{230}“อิทฺธี”ติ ยา เตสํ ธมฺมานํ อิทฺธิ สมิทฺธิ อิชฺฌนา สมิชฺฌนา ลาโภ ปฏิลาโภ ปตฺติ สมฺปตฺติ ผุสนา สจฺฉิกิริยา อุปสมฺปทา.}

\prose{“อิทฺธิปาโท”ติ ตถาภูตสฺส ผสฺโส ฯเปฯ ปคฺคาโห อวิกฺเขโป.}

\prose{“อิทฺธิปาทํ ภาเวตี”ติ เต ธมฺเม อาเสวติ ภาเวติ พหุลีกโรติ, เตน วุจฺจติ “อิทฺธิปาทํ ภาเวตี”ติ.}

\csromanmatikaanchor{mtk.85}
\csromantocmarkref{subparagraph}{2. วีริยิทฺธิปาท}{mtk.85}
\bookmark[dest={mtk.85},level=5]{2. วีริยิทฺธิปาท}
\csromanheader{6}{2. วีริยิทฺธิปาท}
\proseitem{448}{กถญฺจ ภิกฺขุ วีริยสมาธิปธานสงฺขารสมนฺนาคตํ อิทฺธิปาทํ ภาเวติ, อิธ ภิกฺขุ ยสฺมิํ สมเย โลกุตฺตรํ ฌานํ ภาเวติ นิยฺยานิกํ อปจยคามิํ ทิฏฺฐิคตานํ ปหานาย ปฐมาย ภูมิยา ปตฺติยา วิวิจฺเจว กาเมหิ ฯเปฯ อฐมํ ฌานํ อุปสมฺปชฺช วิหรติ ทุกฺขปฏิปทํ ทนฺธาภิญฺญํ, ตสฺมิํ สมเย วีริยสมาธิปธานสงฺขารสมนฺนาคตํ อิทฺธิปาทํ ภาเวติ.}

\proseitem{449}{ตตฺถ กตมํ วีริยํ, โย เจตสิโก วีริยารมฺโภ ฯเปฯ สมฺมาวายาโม วีริยสมฺโพชฺฌงฺโค มคฺคงฺคํ มคฺคปริยาปนฺนํ, อิทํ วุจฺจติ วีริยํ.}

\prose{ตตฺถ กตโม สมาธิ, ยา จิตฺตสฺส ฐิติ สณฺฐิติ ฯเปฯ สมฺมาสมาธิ สมาธิสมฺโพชฺฌงฺโค มคฺคงฺคํ มคฺคปริยาปนฺนํ, อยํ วุจฺจติ สมาธิ.}

\prose{ตตฺถ กตโม ปธานสงฺขาโร, โย เจตสิโก วีริยารมฺโภ ฯเปฯ สมฺมาวายาโม วีริยสมฺโพชฺฌงฺโค มคฺคงฺคํ มคฺคปริยาปนฺนํ, อยํ วุจฺจติ ปธานสงฺขาโร.\csromanspacer{} อิติ อิมินา จ วีริเยน อิมินา จ สมาธินา อิมินา จ ปธานสงฺขาเรน อุเปโต โหติ ฯเปฯ สมนฺนาคโต, เตน วุจฺจติ “วีริยสมาธิปธานสงฺขารสมนฺนาคโต”ติ.}

\proseitem{450}{“อิทฺธี”ติ ยา เตสํ ธมฺมานํ อิทฺธิ สมิทฺธิ อิชฺฌนา สมิชฺฌนา ลาโภ ปฏิลาโภ ปตฺติ สมฺปตฺติ ผุสนา สจฺฉิกิริยา อุปสมฺปทา.}

\prose{“อิทฺธิปาโท”ติ ตถาภูตสฺส ผสฺโส ฯเปฯ ปคฺคาโห อวิกฺเขโป.}

\prose{“อิทฺธิปาทํ ภาเวตี”ติ เต ธมฺเม อาเสวติ ภาเวติ พหุลีกโรติ, เตน วุจฺจติ “อิทฺธิปาทํ ภาเวตี”ติ.}

\csromanmatikaanchor{mtk.86}
\csromantocmarkref{subparagraph}{3. จิตฺติทฺธิปาท}{mtk.86}
\bookmark[dest={mtk.86},level=5]{3. จิตฺติทฺธิปาท}
\csromanheader{6}{9. อิทฺธิปาทวิภงฺค \textbf{3. จิตฺติทฺธิปาท}}
\proseitem{451}{\csromanfolio{231}กถญฺจ ภิกฺขุ จิตฺตสมาธิปธานสงฺขารสมนฺนาคตํ อิทฺธิปาทํ ภาเวติ, อิธ ภิกฺขุ ยสฺมิํ สมเย โลกุตฺตรํ ฌานํ ภาเวติ นิยฺยานิกํ อปจยคามิํ ทิฏฺฐิคตานํ ปหานาย ปฐมาย ภูมิยา ปตฺติยา วิวิจฺเจว กาเมหิ ฯเปฯ ปฐมํ ฌานํ อุปสมฺปชฺช วิหรติ ทุกฺขปฏิปทํ ทนฺธาภิญฺญํ, ตสฺมิํ สมเย จิตฺตสมาธิปธานสงฺขารสมนฺนาคตํ อิทฺธิปาทํ ภาเวติ.}

\proseitem{452}{ตตฺถ กตมํ จิตฺตํ, ยํ จิตฺตํ มโน มานสํ ฯเปฯ ตชฺชามโนวิญฺญาณธาตุ, อิทํ วุจฺจติ จิตฺตํ.}

\prose{ตตฺถ กตโม สมาธิ, ยา จิตฺตสฺส ฐิติ สณฺฐิติ ฯเปฯ สมฺมาสมาธิ สมาธิสมฺโพชฺฌงฺโค มคฺคงฺคํ มคฺคปริยาปนฺนํ, อยํ วุจฺจติ สมาธิ.}

\prose{ตตฺถ กตโม ปธานสงฺขาโร, โย เจตสิโก วีริยารมฺโภ ฯเปฯ สมฺมาวายาโม วีริยสมฺโพชฺฌงฺโค มคฺคงฺคํ มคฺคปริยาปนฺนํ, อยํ วุจฺจติ ปธานสงฺขาโร.\csromanspacer{} อิติ อิมินา จ จิตฺเตน อิมินา จ สมาธินา อิมินา จ ปธานสงฺขาเรน อุเปโต โหติ ฯเปฯ สมนฺนาคโต, เตน วุจฺจติ “จิตฺตสมาธิปธานสงฺขารสมนฺนาคโต”ติ.}

\proseitem{453}{“อิทฺธี”ติ ยา เตสํ ธมฺมานํ อิทฺธิ สมิทฺธิ อิชฺฌนา สมิชฺฌนา ลาโภ ปฏิลาโภ ปตฺติ สมฺปตฺติ ผุสนา สจฺฉิกิริยา อุปสมฺปทา.}

\prose{“อิทฺธิปาโท”ติ ตถาภูตสฺส ผสฺโส ฯเปฯ ปคฺคาโห อวิกฺเขโป.}

\prose{“อิทฺธิปาทํ ภาเวตี”ติ เต ธมฺเม อาเสวติ ภาเวติ พหุลีกโรติ, เตน วุจฺจติ “อิทฺธิปาทํ ภาเวตี”ติ.}

\csromanmatikaanchor{mtk.87}
\csromantocmarkref{subparagraph}{4. วีมํสิทฺธิปาท}{mtk.87}
\bookmark[dest={mtk.87},level=5]{4. วีมํสิทฺธิปาท}
\csromanheader{6}{4. วีมํสิทฺธิปาท}
\proseitem{454}{กถญฺจ ภิกฺขุ วีมํสาสมาธิปธานสงฺขารสมนฺนาคตํ อิทฺธิปาทํ ภาเวติ, อิธ ภิกฺขุ ยสฺมิํ สมเย โลกุตฺตรํ ฌานํ ภาเวติ นิยฺยานิกํ อปจยคามิํ ทิฏฺฐิคตานํ ปหานาย ปฐมาย ภูมิยา ปตฺติยา วิวิจฺเจว กาเมหิ ฯเปฯ ปฐมํ ฌานํ อุปสมฺปชฺช วิหรติ ทุกฺขปฏิปทํ ทนฺธาภิญฺญํ, ตสฺมิํ สมเย วีมํสาสมาธิปธานสงฺขารสมนฺนาคตํ อิทฺธิปาทํ ภาเวติ.}

\proseitem{455}{\csromanfolio{232}ตตฺถ กตมา วีมํสา, ยา ปญฺญา ปชานนา ฯเปฯ อโมโห ธมฺมวิจโย สมฺมาทิฏฺฐิ ธมฺมวิจยสมฺโพชฺฌงฺโค มคฺคงฺคํ มคฺคปริยาปนฺนํ, อยํ วุจฺจติ วีมํสา.}

\prose{ตตฺถ กตโม สมาธิ, ยา จิตฺตสฺส ฐิติ สณฺฐิติ ฯเปฯ สมฺมาสมาธิ สมาธิสมฺโพชฺฌงฺโค มคฺคงฺคํ มคฺคปริยาปนฺนํ, อยํ วุจฺจติ สมาธิ.}

\prose{ตตฺถ กตโม ปธานสงฺขาโร, โย เจตสิโก วีริยารมฺโภ ฯเปฯ สมฺมาวายาโม วีริยสมฺโพชฺฌงฺโค มคฺคงฺคํ มคฺคปริยาปนฺนํ, อยํ วุจฺจติ ปธานสงฺขาโร.\csromanspacer{} อิติ อิมาย จ วีมํสาย อิมินา จ สมาธินา อิมินา จ ปธานสงฺขาเรน อุเปโต โหติ สมุเปโต อุปาคโต สมุปาคโต อุปปนฺโน สมฺปนฺโน สมนฺนาคโต, เตน วุจฺจติ “วีมํสาสมาธิปธานสงฺขารสมนฺนาคโต”ติ.}

\proseitem{456}{“อิทฺธี”ติ ยา เตสํ ธมฺมานํ อิทฺธิ สมิทฺธิ อิชฺฌนา สมิชฺฌนา ลาโภ ปฏิลาโภ ปตฺติ สมฺปตฺติ ผุสนา สจฺฉิกิริยา อุปสมฺปทา.}

\prose{“อิทฺธิปาโท”ติ ตถาภูตสฺส ผสฺโส ฯเปฯ ปคฺคาโห อวิกฺเขโป.}

\prose{“อิทฺธิปาทํ ภาเวตี”ติ เต ธมฺเม อาเสวติ ภาเวติ พหุลีกโรติ, เตน วุจฺจติ “อิทฺธิปาทํ ภาเวตี”ติ.}

\proseitem{457}{จตฺตาโร อิทฺธิปาทา ฉนฺทิทฺธิปาโท วีริยิทฺธิปาโท จิตฺติทฺธิปาโท วีมํสิทฺธิปาโท.}

\proseitem{458}{ตตฺถ กตโม ฉนฺทิทฺธิปาโท, อิธ ภิกฺขุ ยสฺมิํ สมเย โลกุตฺตรํ ฌานํ ภาเวติ นิยฺยานิกํ อปจยคามิํ ทิฏฺฐิคตานํ ปหานาย ปฐมาย ภูมิยา ปตฺติยา วิวิจฺเจว กาเมหิ ฯเปฯ ปฐมํ ฌานํ อุปสมฺปชฺช วิหรติ ทุกฺขปฏิปทํ ทนฺธาภิญฺญํ, โย ตสฺมิํ สมเย ฉนฺโท ฉนฺทิกตา กตฺตุกมฺยตา กุสโล ธมฺมจฺฉนฺโท, อยํ วุจฺจติ ฉนฺทิทฺธิปาโท.\csromanspacer{} อวเสสา ธมฺมา ฉนฺทิทฺธิปาทสมฺปยุตฺตา.}

\proseitem{459}{ตตฺถ กตโม วีริยิทฺธิปาโท, อิธ ภิกฺขุ ยสฺมิํ สมเย โลกุตฺตรํ ฌานํ ภาเวติ นิยฺยานิกํ อปจยคามิํ ทิฏฺฐิคตานํ ปหานาย ปฐมาย ภูมิยา ปตฺติยา วิวิจฺเจว กาเมหิ ฯเปฯ ปฐมํ ฌานํ}

\vspace{\tipitakaparskip}
\csromancenter{อิทฺธิปาทวิภงฺค อุปสมฺปชฺช วิหรติ ทุกฺขปฏิปทํ ทนฺธาภิญฺญํ, โย ตสฺมิํ สมเย เจตสิโก วีริยารมฺโภ ฯเปฯ สมฺมาวายาโม วีริยสมฺโพชฺฌงฺโค มคฺคงฺคํ มคฺคปริยาปนฺนํ, อยํ วุจฺจติ วีริยิทฺธิปาโท.\csromanspacer{} อวเสสา ธมฺมา วีริยิทฺธิปาทสมฺปยุตฺตา.}
\proseitem{460}{\csromanfolio{233}ตตฺถ กตโม จิตฺติทฺธิปาโท, อิธ ภิกฺขุ ยสฺมิํ สมเย โลกุตฺตรํ ฌานํ ภาเวติ นิยฺยานิกํ อปจยคามิํ ทิฏฺฐิคตานํ ปหานาย ปฐมาย ภูมิยา ปตฺติยา วิวิจฺเจว กาเมหิ ฯเปฯ ปฐมํ ฌานํ อุปสมฺปชฺช วิหรติ ทุกฺขปฏิปทํ ทนฺธาภิญฺญํ, ยํ ตสฺมิํ สมเย จิตฺตํ มโน มานสํ ฯเปฯ ตชฺชามโนวิญฺญาณธาตุ, อยํ วุจฺจติ จิตฺติทฺธิปาโท.\csromanspacer{} อวเสสา ธมฺมา จิตฺติทฺธิปาทสมฺปยุตฺตา.}

\proseitem{461}{ตตฺถ กตโม วีมํสิทฺธิปาโท, อิธ ภิกฺขุ ยสฺมิํ สมเย โลกุตฺตรํ ฌานํ ภาเวติ นิยฺยานิกํ อปจยคามิํ ทิฏฺฐิคตานํ ปหานาย ปฐมาย ภูมิยา ปตฺติยา วิวิจฺเจว กาเมหิ ฯเปฯ ปฐมํ ฌานํ อุปสมฺปชฺช วิหรติ ทุกฺขปฏิปทํ ทนฺธาภิญฺญํ, ยา ตสฺมิํ สมเย ปญฺญา ปชานนา ฯเปฯ อโมโห ธมฺมวิจโย สมฺมาทิฏฺฐิ ธมฺมวิจยสมฺโพชฺฌงฺโค มคฺคงฺคํ มคฺคปริยาปนฺนํ, อยํ วุจฺจติ วีมํสิทฺธิปาโท.\csromanspacer{} อวเสสา ธมฺมา วีมํสิทฺธิปาทสมฺปยุตฺตา.}

\csromanheader{2}{อภิธมฺมภาชนียํ.}
\csromansectionrule
\csromanheader{2}{3. ปญฺหาปุจฺฉก}
\proseitem{462}{จตฺตาโร อิทฺธิปาทา\csromandash{}อิธ ภิกฺขุ ฉนฺทสมาธิปธานสงฺขารสมนฺนาคตํ อิทฺธิปาทํ ภาเวติ, วีริยสมาธิ ฯเปฯ จิตฺตสมาธิ ฯเปฯ วีมํสาสมาธิปธานสงฺขารสมนฺนาคตํ อิทฺธิปาทํ ภาเวติ.}

\proseitem{463}{จตุนฺนํ อิทฺธิปาทานํ กติ กุสลา, กติ อกุสลา, กติ อพฺยากตา ฯเปฯ กติ สรณา, กติ อรณา.}

\csromanmatikaanchor{mtk.88}
\csromanmatikaanchor{mtk.89}
\csromantocmarknopage{subparagraph}{3. ปญฺหาปุจฺฉก}{mtk.88}
\bookmark[dest={mtk.88},level=5]{3. ปญฺหาปุจฺฉก}
\csromantocmarkref{subparagraph}{1. ติก}{mtk.89}
\bookmark[dest={mtk.89},level=5]{1. ติก}
\csromanheader{6}{1. ติก}
\proseitem{464}{กุสลาเยว.\csromanspacer{} สิยา สุขาย เวทนาย สมฺปยุตฺตา, สิยา อทุกฺขมสุขาย เวทนาย สมฺปยุตฺตา.\csromanspacer{} วิปากธมฺมธมฺมา.\csromanspacer{} อนุปาทินฺน-อนุปาทานิยา.\csromanspacer{} อสํกิลิฏฺฐ-อสํกิเลสิกา.\csromanspacer{} สิยา สวิตกฺกสวิจารา, สิยา \csromanfolio{234}อวิตกฺกวิจารมตฺตา.\csromanspacer{} สิยา อวิตกฺก-อวิจารา.\csromanspacer{} สิยา ปีติสหคตา, สิยา สุขสหคตา, สิยา อุเปกฺขาสหคตา.\csromanspacer{} เนว ทสฺสเนน น ภาวนาย ปหาตพฺพา.\csromanspacer{} เนว ทสฺสเนน น ภาวนาย ปหาตพฺพเหตุกา.\csromanspacer{} อปจยคามิโน.\csromanspacer{} เสกฺขา.\csromanspacer{} อปฺปมาณา.\csromanspacer{} อปฺปมาณารมฺมณา.\csromanspacer{} ปณีตา.\csromanspacer{} สมฺมตฺตนิยตา.\csromanspacer{} น มคฺคารมฺมณา.\csromanspacer{} มคฺคเหตุกา.\csromanspacer{} น มคฺคาธิปติโน.\csromanspacer{} สิยา อุปฺปนฺนา, สิยา อนุปฺปนฺนา, น วตฺตพฺพา “อุปฺปาทิโน”ติ.\csromanspacer{} สิยา อตีตา, สิยา อนาคตา, สิยา ปจฺจุปฺปนฺนา.\csromanspacer{} น วตฺตพฺพา “อตีตารมฺมณา”ติปิ, “อนาคตารมฺมณา”ติปิ, “ปจฺจุปฺปนฺนารมฺมณา”ติปิ.\csromanspacer{} สิยา อชฺฌตฺตา, สิยา พหิทฺธา, สิยา อชฺฌตฺตพหิทฺธา.\csromanspacer{} พหิทฺธารมฺมณา.\csromanspacer{} อนิทสฺสน-อปฺปฏิฆา.}

\csromanmatikaanchor{mtk.90}
\csromantocmarkref{subparagraph}{2. ทุก}{mtk.90}
\bookmark[dest={mtk.90},level=5]{2. ทุก}
\csromanheader{6}{2. ทุก}
\proseitem{465}{วีมํสิทฺธิปาโท เหตุ.\csromanspacer{} ตโย อิทฺธิปาทา น เหตู, สเหตุกา.\csromanspacer{} เหตุ สมฺปยุตฺตา.\csromanspacer{} วีมํสิทฺธิปาโท เหตุ เจว สเหตุโก จ.\csromanspacer{} ตโย อิทฺธิปาทา น วตฺตพฺพา “เหตู เจว สเหตุกา จา”ติ, สเหตุกา เจว น จ เหตู.\csromanspacer{} วีมํสิทฺธิปาโท เหตุ เจว เหตุสมฺปยุตฺโต จ.\csromanspacer{} ตโย อิทฺธิปาทา น วตฺตพฺพา “เหตู เจว เหตุสมฺปยุตฺตา จา”ติ, เหตุสมฺปยุตฺตา เจว น จ เหตู.\csromanspacer{} ตโย อิทฺธิปาทา น เหตู, สเหตุกา.\csromanspacer{} วีมํสิทฺธิปาโท น วตฺตพฺโพ “น เหตุ สเหตุโก”ติปิ, “น เหตุ อเหตุโก”ติปิ.\csromanspacer{} สปฺปจฺจยา, สงฺขตา, อนิทสฺสนา, อปฺปฏิฆา, อรูปา, โลกุตฺตรา, เกนจิ วิญฺเญยฺยา, เกนจิ น วิญฺเญยฺยา.\csromanspacer{} โน อาสวา.\csromanspacer{} อนาสวา.\csromanspacer{} อาสว วิปฺปยุตฺตา.\csromanspacer{} น วตฺตพฺพา “อาสวา เจว สาสวา จา”ติปิ, “สาสวา เจว โน จ อาสวา”ติปิ.\csromanspacer{} น วตฺตพฺพา “อาสวา เจว อาสวสมฺปยุตฺตา จา”ติปิ, “อาสวสมฺปยุตฺตา เจว โน จ อาสวา”ติปิ.\csromanspacer{} อาสววิปฺปยุตฺตา อนาสวา.}

\prose{โน สํโยชนา ฯเปฯ.\csromanspacer{} โน คนฺถา ฯเปฯ.\csromanspacer{} โน โอฆา ฯเปฯ.\csromanspacer{} โน โยคา ฯเปฯ.\csromanspacer{} โน นีวรณา ฯเปฯ.\csromanspacer{} โน ปรามาสา ฯเปฯ สารมฺมณา.\csromanspacer{} ตโย อิทฺธิปาทา โน จิตฺตา.\csromanspacer{} จิตฺติทฺธิปาโท จิตฺตํ.\csromanspacer{} ตโย อิทฺธิปาทา เจตสิกา.\csromanspacer{} จิตฺติทฺธิปาโท อเจตสิโก.\csromanspacer{} ตโย อิทฺธิปาทา จิตฺตสมฺปยุตฺตา.\csromanspacer{} จิตฺติทฺธิปาโท น วตฺตพฺโพ “จิตฺเตน สมฺปยุตฺโต”ติปิ, “จิตฺเตน วิปฺปยุตฺโต”ติปิ.\csromanspacer{} ตโย อิทฺธิปาทา จิตฺตสํสฏฺฐา.\csromanspacer{} จิตฺติทฺธิปาโท น วตฺตพฺโพ “จิตฺเตน สํสฏฺโฐ”ติปิ, “จิตฺเตน วิสํสฏฺโฐ”ติปิ.\csromanspacer{} ตโย อิทฺธิปาทา จิตฺตสมุฏฺฐานา.\csromanspacer{} จิตฺติทฺธิปาโท โน}

\vspace{\tipitakaparskip}
\csromancenter{อิทฺธิปาทวิภงฺค จิตฺตสมุฏฺฐาโน.\csromanspacer{} ตโย อิทฺธิปาทา จิตฺตสหภุโน.\csromanspacer{} จิตฺติทฺธิปาโท โน จิตฺตสหภู.\csromanspacer{} ตโย อิทฺธิปาทา จิตฺตานุปริวตฺติโน.\csromanspacer{} จิตฺติทฺธิปาโท โน จิตฺตานุปริวตฺติ.\csromanspacer{} ตโย อิทฺธิปาทา จิตฺตสํสฏฺฐสมุฏฺฐานา.\csromanspacer{} จิตฺติทฺธิปาโท โน จิตฺตสํสฏฺฐสมุฏฺฐาโน.\csromanspacer{} ตโย อิทฺธิปาทา จิตฺตสํสฏฺฐสมุฏฺฐานสหภุโน.\csromanspacer{} จิตฺติทฺธิปาโท โน จิตฺตสํสฏฺฐสมุฏฺฐานสหภู.\csromanspacer{} ตโย อิทฺธิปาทา จิตฺตสํสฏฺฐสมุฏฺฐานานุปริวตฺติโน.\csromanspacer{} จิตฺติทฺธิปาโท โน จิตฺตสํสฏฺฐสมุฏฺฐานานุปริวตฺติ.}
\prosecont{\csromanfolio{235}ตโย อิทฺธิปาทา พาหิรา.\csromanspacer{} จิตฺติทฺธิปาโท อชฺฌตฺติโก.\csromanspacer{} โน อุปาทา.\csromanspacer{} อนุปาทินฺนา.\csromanspacer{} โน อุปาทานา ฯเปฯ.\csromanspacer{} โน กิเลสา ฯเปฯ.\csromanspacer{} น ทสฺสเนน ปหาตพฺพา.\csromanspacer{} น ภาวนาย ปหาตพฺพา.\csromanspacer{} น ทสฺสเนน ปหาตพฺพเหตุกา.\csromanspacer{} น ภาวนาย ปหาตพฺพเหตุกา.\csromanspacer{} สิยา สวิตกฺกา, สิยา อวิตกฺกา, สิยา สวิจารา, สิยา อวิจารา.\csromanspacer{} สิยา สปฺปีติกา, สิยา อปฺปีติกา, สิยา ปีติสหคตา, สิยา น ปีติสหคตา.\csromanspacer{} สิยา สุขสหคตา, สิยา น สุขสหคตา.\csromanspacer{} สิยา อุเปกฺขาสหคตา, สิยา น อุเปกฺขาสหคตา.\csromanspacer{} น กามาวจรา.\csromanspacer{} น รูปาวจรา.\csromanspacer{} น อรูปาวจรา.\csromanspacer{} อปริยาปนฺนา.\csromanspacer{} นิยฺยานิกา.\csromanspacer{} นิยตา.\csromanspacer{} อนุตฺตรา.\csromanspacer{} อรณาติ.}

\vspace{\tipitakaparskip}
\csromancenter{ปญฺหาปุจฺฉกํ.}
\nitthitam{\textbf{อิทฺธิปาทวิภงฺโค นิฏฺฐิโต.}}
\csromanmatikaanchor{mtk.91}
\csromantocmarknopage{subparagraph}{10. โพชฺฌงฺควิภงฺค}{mtk.91}
\bookmark[dest={mtk.91},level=5]{10. โพชฺฌงฺควิภงฺค}
\csromanheader{6}{10. โพชฺฌงฺควิภงฺค}
\csromanmatikaanchor{mtk.92}
\csromantocmarkref{subparagraph}{1. สุตฺตนฺตภาชนีย}{mtk.92}
\bookmark[dest={mtk.92},level=5]{1. สุตฺตนฺตภาชนีย}
\csromanheader{6}{1. สุตฺตนฺตภาชนีย}
\proseitem{466}{\csromanfolio{236}สตฺต โพชฺฌงฺคา สติสมฺโพชฺฌงฺโค ธมฺมวิจยสมฺโพชฺฌงฺโค วีริยสมฺโพชฺฌงฺโค ปีติสมฺโพชฺฌงฺโค ปสฺสทฺธิสมฺโพชฺฌงฺโค สมาธิสมฺโพชฺฌงฺโค อุเปกฺขาสมฺโพชฺฌงฺโค.}

\proseitem{467}{ตตฺถ กตโม สติสมฺโพชฺฌงฺโค, อิธ ภิกฺขุ สติมา โหติ ปรเมน สติเนปกฺเกน สมนฺนาคโต จิรกตมฺปิ จิรภาสิตมฺปิ สริตา โหติ อนุสฺสริตา, อยํ วุจฺจติ สติสมฺโพชฺฌงฺโค. (1)}

\prose{โส ตถา สโต วิหรนฺโต ตํ ธมฺมํ ปญฺญาย ปวิจินติ ปวิจรติ ปริวีมํสมาปชฺชติ, อยํ วุจฺจติ ธมฺมวิจยสมฺโพชฺฌงฺโค. (2)}

\prose{ตสฺส ตํ ธมฺมํ ปญฺญาย ปวิจินโต ปวิจรโต ปริวีมํสมาปชฺชโต อารทฺธํ โหติ วีริยํ อสลฺลีนํ, อยํ วุจฺจติ วีริยสมฺโพชฺฌงฺโค. (3)}

\prose{อารทฺธวีริยสฺส อุปฺปชฺชติ ปีติ นิรามิสา, อยํ วุจฺจติ ปีติสมฺโพชฺฌงฺโค. (4)}

\prose{ปีติมนสฺส กาโยปิ ปสฺสมฺภติ จิตฺตมฺปิ ปสฺสมฺภติ, อยํ วุจฺจติ ปสฺสทฺธิสมฺโพชฺฌงฺโค. (5)}

\prose{ปสฺสทฺธกายสฺส สุขิโน จิตฺตํ สมาธิยติ, อยํ วุจฺจติ สมาธิสมฺโพชฺฌงฺโค. (6)}

\prose{โส ตถา สมาหิตํ จิตฺตํ สาธุกํ อชฺฌุเปกฺขิตา โหติ, อยํ วุจฺจติ อุเปกฺขาสมฺโพชฺฌงฺโค. (7)}

\proseitem{468}{สตฺต โพชฺฌงฺคา สติสมฺโพชฺฌงฺโค ธมฺมวิจยสมฺโพชฺฌงฺโค วีริยสมฺโพชฺฌงฺโค ปีติสมฺโพชฺฌงฺโค ปสฺสทฺธิสมฺโพชฺฌงฺโค สมาธิสมฺโพชฺฌงฺโค อุเปกฺขาสมฺโพชฺฌงฺโค.}

\proseitem{469}{ตตฺถ กตโม สติสมฺโพชฺฌงฺโค, อตฺถิ อชฺฌตฺตํ ธมฺเมสุ สติ, อตฺถิ พหิทฺธา ธมฺเมสุ สติ.\csromanspacer{} ยทปิ อชฺฌตฺตํ ธมฺเมสุ สติ, ตทปิ สติสมฺโพชฺฌงฺโค อภิญฺญาย สมฺโพธาย นิพฺพานาย สํวตฺตติ.\csromanspacer{} ยทปิ}

\vspace{\tipitakaparskip}
\csromancenter{โพชฺฌงฺควิภงฺค พหิทฺธา ธมฺเมสุ สติ, ตทปิ สติสมฺโพชฺฌงฺโค อภิญฺญาย สมฺโพธาย นิพฺพานาย สํวตฺตติ. (1)}
\prose{\csromanfolio{237}ตตฺถ กตโม ธมฺมวิจยสมฺโพชฺฌงฺโค, อตฺถิ อชฺฌตฺตํ ธมฺเมสุ ปวิจโย, อตฺถิ พหิทฺธา ธมฺเมสุ ปวิจโย.\csromanspacer{} ยทปิ อชฺฌตฺตํ ธมฺเมสุ ปวิจโย, ตทปิ ธมฺมวิจยสมฺโพชฺฌงฺโค อภิญฺญาย สมฺโพธาย นิพฺพานาย สํวตฺตติ.\csromanspacer{} ยทปิ พหิทฺธา ธมฺเมสุ ปวิจโย, ตทปิ ธมฺมวิจยสมฺโพชฺฌงฺโค อภิญฺญาย สมฺโพธาย นิพฺพานาย สํวตฺตติ. (2)}

\prose{ตตฺถ กตโม วีริยสมฺโพชฺฌงฺโค, อตฺถิ กายิกํ วีริยํ, อตฺถิ เจตสิกํ วีริยํ.\csromanspacer{} ยทปิ กายิกํ วีริยํ, ตทปิ วีริยสมฺโพชฺฌงฺโค อภิญฺญาย สมฺโพธาย นิพฺพานาย สํวตฺตติ.\csromanspacer{} ยทปิ เจตสิกํ วีริยํ, ตทปิ วีริยสมฺโพชฺฌงฺโค อภิญฺญาย สมฺโพธาย นิพฺพานาย สํวตฺตติ. (3)}

\prose{ตตฺถ กตโม ปีติสมฺโพชฺฌงฺโค, อตฺถิ สวิตกฺกสวิจารา ปีติ, อตฺถิ อวิตกฺก-อวิจารา ปีติ.\csromanspacer{} ยทปิ สวิตกฺกสวิจารา ปีติ, ตทปิ ปีติสมฺโพชฺฌงฺโค อภิญฺญาย สมฺโพธาย นิพฺพานาย สํวตฺตติ.\csromanspacer{} ยทปิ อวิตกฺก-อวิจารา ปีติ, ตทปิ ปีติสมฺโพชฺฌงฺโค อภิญฺญาย สมฺโพธาย นิพฺพานาย สํวตฺตติ. (4)}

\prose{ตตฺถ กตโม ปสฺสทฺธิสมฺโพชฺฌงฺโค, อตฺถิ กายปสฺสทฺธิ\footnote{กายปฺปสฺสทฺธิ (สฺยา, ก)}, อตฺถิ จิตฺตปสฺสทฺธิ\footnote{จิตฺตปฺปสฺสทฺธิ (สฺยา, ก)}.\csromanspacer{} ยทปิ กายปสฺสทฺธิ, ตทปิ ปสฺสทฺธิสมฺโพชฺฌงฺโค อภิญฺญาย สมฺโพธาย นิพฺพานาย สํวตฺตติ.\csromanspacer{} ยทปิ จิตฺตปสฺสทฺธิ, ตทปิ ปสฺสทฺธิสมฺโพชฺฌงฺโค อภิญฺญาย สมฺโพธาย นิพฺพานาย สํวตฺตติ. (5)}

\prose{ตตฺถ กตโม สมาธิสมฺโพชฺฌงฺโค, อตฺถิ สวิตกฺโก สวิจาโร สมาธิ, อตฺถิ อวิตกฺโก อวิจาโร สมาธิ.\csromanspacer{} ยทปิ สวิตกฺโก สวิจาโร สมาธิ, ตทปิ สมาธิสมฺโพชฺฌงฺโค อภิญฺญาย สมฺโพธาย นิพฺพานาย สํวตฺตติ.\csromanspacer{} ยทปิ อวิตกฺโก อวิจาโร สมาธิ, ตทปิ สมาธิสมฺโพชฺฌงฺโค อภิญฺญาย สมฺโพธาย นิพฺพานาย สํวตฺตติ. (6)}

\prose{ตตฺถ กตโม อุเปกฺขาสมฺโพชฺฌงฺโค, อตฺถิ อชฺฌตฺตํ ธมฺเมสุ อุเปกฺขา, อตฺถิ พหิทฺธา ธมฺเมสุ อุเปกฺขา.\csromanspacer{} ยทปิ อชฺฌตฺตํ ธมฺเมสุ อุเปกฺขา, ตทปิ อุเปกฺขาสมฺโพชฺฌงฺโค อภิญฺญาย สมฺโพธาย นิพฺพานาย สํวตฺตติ. \csromanfolio{238}ยทปิ พหิทฺธา ธมฺเมสุ อุเปกฺขา, ตทปิ อุเปกฺขาสมฺโพชฺฌงฺโค อภิญฺญาย สมฺโพธาย นิพฺพานาย สํวตฺตติ. (7)}

\proseitem{470}{สตฺต โพชฺฌงฺคา สติสมฺโพชฺฌงฺโค ธมฺมวิจยสมฺโพชฺฌงฺโค วีริยสมฺโพชฺฌงฺโค ปีติสมฺโพชฺฌงฺโค ปสฺสทฺธิสมฺโพชฺฌงฺโค สมาธิสมฺโพชฺฌงฺโค อุเปกฺขาสมฺโพชฺฌงฺโค.}

\proseitem{471}{ตตฺถ กตโม สติสมฺโพชฺฌงฺโค, อิธ ภิกฺขุ สติสมฺโพชฺฌงฺคํ ภาเวติ วิเวกนิสฺสิตํ วิราคนิสฺสิตํ นิโรธนิสฺสิตํ โวสฺสคฺคปริณามิํ.\csromanspacer{} ธมฺมวิจยสมฺโพชฺฌงฺคํ ภาเวติ ฯเปฯ.\csromanspacer{} วีริยสมฺโพชฺฌงฺคํ ภาเวติ.\csromanspacer{} ปีติสมฺโพชฺฌงฺคํ ภาเวติ.\csromanspacer{} ปสฺสทฺธิสมฺโพชฺฌงฺคํ ภาเวติ.\csromanspacer{} สมาธิสมฺโพชฺฌงฺคํ ภาเวติ.\csromanspacer{} อุเปกฺขาสมฺโพชฺฌงฺคํ ภาเวติ วิเวกนิสฺสิตํ วิราคนิสฺสิตํ นิโรธนิสฺสิตํ โวสฺสคฺคปริณามิํ.}

\csromanheader{2}{สุตฺตนฺตภาชนียํ.}
\csromansectionrule
\csromanmatikaanchor{mtk.93}
\csromantocmarkref{subparagraph}{2. อภิธมฺมภาชนีย}{mtk.93}
\bookmark[dest={mtk.93},level=5]{2. อภิธมฺมภาชนีย}
\csromanheader{6}{2. อภิธมฺมภาชนีย}
\proseitem{472}{สตฺต โพชฺฌงฺคา สติสมฺโพชฺฌงฺโค ธมฺมวิจยสมฺโพชฺฌงฺโค วีริยสมฺโพชฺฌงฺโค ปีติสมฺโพชฺฌงฺโค ปสฺสทฺธิสมฺโพชฺฌงฺโค สมาธิสมฺโพชฺฌงฺโค อุเปกฺขาสมฺโพชฺฌงฺโค.}

\proseitem{473}{ตตฺถ กตเม สตฺต โพชฺฌงฺคา, อิธ ภิกฺขุ ยสฺมิํ สมเย โลกุตฺตรํ ฌานํ ภาเวติ นิยฺยานิกํ อปจยคามิํ ทิฏฺฐิคตานํ ปหานาย ปฐมาย ภูมิยา ปตฺติยา วิวิจฺเจว กาเมหิ ฯเปฯ ปฐมํ ฌานํ อุปสมฺปชฺช วิหรติ ทุกฺขปฏิปทํ ทนฺธาภิญฺญํ, ตสฺมิํ สมเย สตฺต โพชฺฌงฺคา โหนฺติ สติสมฺโพชฺฌงฺโค ฯเปฯ อุเปกฺขาสมฺโพชฺฌงฺโค.}

\proseitem{474}{ตตฺถ กตโม สติสมฺโพชฺฌงฺโค, ยา สติ อนุสฺสติ ฯเปฯ สมฺมาสติ สติสมฺโพชฺฌงฺโค มคฺคงฺคํ มคฺคปริยาปนฺนํ, อยํ วุจฺจติ สติสมฺโพชฺฌงฺโค. (1)}

\csromanheader{2}{10. โพชฺฌงฺควิภงฺค}
\prose{\csromanfolio{239}ตตฺถ กตโม ธมฺมวิจยสมฺโพชฺฌงฺโค, ยา ปญฺญา ปชานนา ฯเปฯ อโมโห ธมฺมวิจโย สมฺมาทิฏฺฐิ ธมฺมวิจยสมฺโพชฺฌงฺโค มคฺคงฺคํ มคฺคปริยาปนฺนํ, อยํ วุจฺจติ ธมฺมวิจยสมฺโพชฺฌงฺโค. (2)}

\prose{ตตฺถ กตโม วีริยสมฺโพชฺฌงฺโค, โย เจตสิโก วีริยารมฺโภ ฯเปฯ สมฺมาวายาโม วีริยสมฺโพชฺฌงฺโค มคฺคงฺคํ มคฺคปริยาปนฺนํ, อยํ วุจฺจติ วีริยสมฺโพชฺฌงฺโค. (3)}

\prose{ตตฺถ กตโม ปีติสมฺโพชฺฌงฺโค, ยา ปีติ ปาโมชฺชํ อาโมทนา ปโมทนา หาโส ปหาโส วิตฺติ โอทคฺยํ\footnote{โอทคฺคํ (สี) อภิ 1. 18 ปิฏฺเฐปิ.} อตฺตมนตา จิตฺตสฺส ปีติสมฺโพชฺฌงฺโค, อยํ วุจฺจติ ปีติสมฺโพชฺฌงฺโค. (4)}

\prose{ตตฺถ กตโม ปสฺสทฺธิสมฺโพชฺฌงฺโค, ยา เวทนากฺขนฺธสฺส สญฺญากฺขนฺธสฺส สงฺขารกฺขนฺธสฺส วิญฺญาณกฺขนฺธสฺส ปสฺสทฺธิ ปฏิปฺปสฺสทฺธิ ปสฺสมฺภนา ปฏิปฺปสฺสมฺภนา ปฏิปฺปสฺสมฺภิ ตตฺตํ ปสฺสทฺธิสมฺโพชฺฌงฺโค, อยํ วุจฺจติ ปสฺสทฺธิสมฺโพชฺฌงฺโค. (5)}

\prose{ตตฺถ กตโม สมาธิสมฺโพชฺฌงฺโค, ยา จิตฺตสฺส ฐิติ ฯเปฯ สมฺมาสมาธิ สมาธิสมฺโพชฺฌงฺโค มคฺคงฺคํ มคฺคปริยาปนฺนํ, อยํ วุจฺจติ สมาธิสมฺโพชฺฌงฺโค. (6)}

\prose{ตตฺถ กตโม อุเปกฺขาสมฺโพชฺฌงฺโค, ยา อุเปกฺขา อุเปกฺขนา อชฺฌุเปกฺขนา มชฺฌตฺตตา จิตฺตตา อุเปกฺขาสมฺโพชฺฌงฺโค, อยํ วุจฺจติ อุเปกฺขาสมฺโพชฺฌงฺโค.\csromanspacer{} อิเม วุจฺจนฺติ สตฺต โพชฺฌงฺคา.\csromanspacer{} อวเสสา ธมฺมา สตฺตหิ โพชฺฌงฺเคหิ สมฺปยุตฺตา. (7)}

\proseitem{475}{สตฺต โพชฺฌงฺคา สติสมฺโพชฺฌงฺโค ฯเปฯ อุเปกฺขาสมฺโพชฺฌงฺโค.}

\proseitem{476}{ตตฺถ กตโม สติสมฺโพชฺฌงฺโค, อิธ ภิกฺขุ ยสฺมิํ สมเย โลกุตฺตรํ ฌานํ ภาเวติ นิยฺยานิกํ อปจยคามิํ ทิฏฺฐิคตานํ ปหานาย ปฐมาย ภูมิยา ปตฺติยา วิวิจฺเจว กาเมหิ ฯเปฯ ปฐมํ ฌานํ อุปสมฺปชฺช วิหรติ ทุกฺขปฏิปทํ ทนฺธาภิญฺญํ, ยา ตสฺมิํ สมเย สติ อนุสฺสติ สมฺมาสติ สติสมฺโพชฺฌงฺโค มคฺคงฺคํ มคฺคปริยาปนฺนํ, อยํ วุจฺจติ สติสมฺโพชฺฌงฺโค.\csromanspacer{} อวเสสา ธมฺมา สติสมฺโพชฺฌงฺคสมฺปยุตฺตา ฯเปฯ.\csromanspacer{} อวเสสา \csromanfolio{240}ธมฺมา ธมฺมวิจยสมฺโพชฺฌงฺคสมฺปยุตฺตา ฯเปฯ.\csromanspacer{} อวเสสา ธมฺมา วีริยสมฺโพชฺฌงฺคสมฺปยุตฺตา ฯเปฯ.\csromanspacer{} อวเสสา ธมฺมา ปีติสมฺโพชฺฌงฺคสมฺปยุตฺตา ฯเปฯ.\csromanspacer{} อวเสสา ธมฺมา ปสฺสทฺธิสมฺโพชฺฌงฺคสมฺปยุตฺตา ฯเปฯ.\csromanspacer{} อวเสสา ธมฺมา สมาธิสมฺโพชฺฌงฺคสมฺปยุตฺตา.}

\prose{ตตฺถ กตโม อุเปกฺขาสมฺโพชฺฌงฺโค, อิธ ภิกฺขุ ยสฺมิํ สมเย โลกุตฺตรํ ฌานํ ภาเวติ นิยฺยานิกํ อปจยคามิํ ทิฏฺฐิคตานํ ปหานาย ปฐมาย ภูมิยา ปตฺติยา วิวิจฺเจว กาเมหิ ฯเปฯ ปฐมํ ฌานํ อุปสมฺปชฺช วิหรติ ทุกฺขปฏิปทํ ทนฺธาภิญฺญํ, ยา ตสฺมิํ สมเย อุเปกฺขา อุเปกฺขนา อชฺฌุเปกฺขนา มชฺฌตฺตตา จิตฺตสฺส อุเปกฺขาสมฺโพชฺฌงฺโค, อยํ วุจฺจติ อุเปกฺขาสมฺโพชฺฌงฺโค.\csromanspacer{} อวเสสา ธมฺมา อุเปกฺขาสมฺโพชฺฌงฺคสมฺปยุตฺตา.}

\proseitem{477}{สตฺต โพชฺฌงฺคา สติสมฺโพชฺฌงฺโค ฯเปฯ อุเปกฺขาสมฺโพชฺฌงฺโค.}

\proseitem{478}{ตตฺถ กตเม สตฺต โพชฺฌงฺคา, อิธ ภิกฺขุ ยสฺมิํ สมเย โลกุตฺตรํ ฌานํ ภาเวติ นิยฺยานิกํ อปจยคามิํ ทิฏฺฐิคตานํ ปหานาย ปฐมาย ภูมิยา ปตฺติยา วิวิจฺเจว กาเมหิ ฯเปฯ ปฐมํ ฌานํ อุปสมฺปชฺช วิหรติ ทุกฺขปฏิปทํ ทนฺธาภิญฺญํ, โย ตสฺมิํ สมเย ผสฺโส โหติ ฯเปฯ อวิกฺเขโป โหติ, อิเม ธมฺมา กุสลา.\csromanspacer{} ตสฺเสว โลกุตฺตรสฺส กุสลสฺส ฌานสฺส กตตฺตา ภาวิตตฺตา วิปากํ วิวิจฺเจว กาเมหิ ฯเปฯ ปฐมํ ฌานํ อุปสมฺปชฺช วิหรติ ทุกฺขปฏิปทํ ทนฺธาภิญฺญํ สุญฺญตํ, ตสฺมิํ สมเย สตฺต โพชฺฌงฺคา โหนฺติ สติสมฺโพชฺฌงฺโค ฯเปฯ อุเปกฺขาสมฺโพชฺฌงฺโค.}

\proseitem{479}{ตตฺถ กตโม สติสมฺโพชฺฌงฺโค, ยา สติ อนุสฺสติ ฯเปฯ สมฺมาสติ สติสมฺโพชฺฌงฺโค มคฺคงฺคํ มคฺคปริยาปนฺนํ, อยํ วุจฺจติ สติสมฺโพชฺฌงฺโค ฯเปฯ.}

\prose{ตตฺถ กตโม อุเปกฺขาสมฺโพชฺฌงฺโค, ยา อุเปกฺขา อุเปกฺขนา อชฺฌุเปกฺขนา มชฺฌตฺตตา จิตฺตสฺส อุเปกฺขาสมฺโพชฺฌงฺโค, อยํ วุจฺจติ อุเปกฺขาสมฺโพชฺฌงฺโค.\csromanspacer{} อิเม วุจฺจนฺติ สตฺต โพชฺฌงฺคา.\csromanspacer{} อวเสสา ธมฺมา สตฺตหิ โพชฺฌงฺเคหิ สมฺปยุตฺตา.}

\proseitem{480}{สตฺต โพชฺฌงฺคา สติสมฺโพชฺฌงฺโค ฯเปฯ อุเปกฺขาสมฺโพชฺฌงฺโค.}

\csromanheader{2}{10. โพชฺฌงฺควิภงฺค}
\proseitem{481}{\csromanfolio{241}ตตฺถ กตโม สติสมฺโพชฺฌงฺโค, อิธ ภิกฺขุ ยสฺมิํ สมเย โลกุตฺตรํ ฌานํ ภาเวติ นิยฺยานิกํ อปจยคามิํ ทิฏฺฐิคตานํ ปหานาย ปฐมาย ภูมิยา ปตฺติยา วิวิจฺเจว กาเมหิ ฯเปฯ ปฐมํ ฌานํ อุปสมฺปชฺช วิหรติ ทุกฺขปฏิปทํ ทนฺธาภิญฺญํ, ตสฺมิํ สมเย ผสฺโส โหติ ฯเปฯ อวิกฺเขโป โหติ, อิเม ธมฺมา กุสลา.\csromanspacer{} ตสฺเสว โลกุตฺตรสฺส กุสลสฺส ฌานสฺส กตตฺตา ภาวิตตฺตา วิปากํ วิวิจฺเจว กาเมหิ ฯเปฯ ปฐมํ ฌานํ อุปสมฺปชฺช วิหรติ ทุกฺขปฏิปทํ ทนฺธาภิญฺญํ, ยา ตสฺมิํ สมเย สติ อนุสฺสติ สมฺมาสติ สติสมฺโพชฺฌงฺโค มคฺคงฺคํ มคฺคปริยาปนฺนํ, อยํ วุจฺจติ สติสมฺโพชฺฌงฺโค.\csromanspacer{} อวเสสา ธมฺมา สติสมฺโพชฺฌงฺคสมฺปยุตฺตา ฯเปฯ.\csromanspacer{} อวเสสา ธมฺมา ธมฺมวิจยสมฺโพชฺฌงฺคสมฺปยุตฺตา ฯเปฯ.\csromanspacer{} อวเสสา ธมฺมา วีริยสมฺโพชฺฌงฺคสมฺปยุตฺตา ฯเปฯ.\csromanspacer{} อวเสสา ธมฺมา ปีติสมฺโพชฺฌงฺคสมฺปยุตฺตา ฯเปฯ.\csromanspacer{} อวเสสา ธมฺมา ปสฺสทฺธิสมฺโพชฺฌงฺคสมฺปยุตฺตา ฯเปฯ.\csromanspacer{} อวเสสา ธมฺมา สมาธิสมฺโพชฺฌงฺคสมฺปยุตฺตา.}

\prose{ตตฺถ กตโม อุเปกฺขาสมฺโพชฺฌงฺโค, อิธ ภิกฺขุ ยสฺมิํ สมเย โลกุตฺตรํ ฌานํ ภาเวติ นิยฺยานิกํ อปจยคามิํ ทิฏฺฐิคตานํ ปหานาย ปฐมาย ภูมิยา ปตฺติยา วิวิจฺเจว กาเมหิ ฯเปฯ ปฐมํ ฌานํ อุปสมฺปชฺช วิหรติ ทุกฺขปฏิปทํ ทนฺธาภิญฺญํ, ตสฺมิํ สมเย ผสฺโส โหติ ฯเปฯ อวิกฺเขโป โหติ, อิเม ธมฺมา กุสลา.\csromanspacer{} ตสฺเสว โลกุตฺตรสฺส กุสลสฺส ฌานสฺส กตตฺตา ภาวิตตฺตา วิปากํ วิวิจฺเจว กาเมหิ ฯเปฯ ปฐมํ ฌานํ อุปสมฺปชฺช วิหรติ ทุกฺขปฏิปทํ ทนฺธาภิญฺญํ สุญฺญตํ, ยา ตสฺมิํ สมเย อุเปกฺขา อุเปกฺขนา อชฺฌุเปกฺขนา มชฺฌตฺตตา จิตฺตสฺส อุเปกฺขาสมฺโพชฺฌงฺโค, อยํ วุจฺจติ อุเปกฺขาสมฺโพชฺฌงฺโค.\csromanspacer{} อวเสสา ธมฺมา อุเปกฺขาสมฺโพชฺฌงฺคสมฺปยุตฺตา.}

\csromanheader{2}{อภิธมฺมภาชนียํ.}
\csromansectionrule
\csromanmatikaanchor{mtk.94}
\csromantocmarkref{subparagraph}{3. ปญฺหาปุจฺฉก}{mtk.94}
\bookmark[dest={mtk.94},level=5]{3. ปญฺหาปุจฺฉก}
\csromanheader{6}{3. ปญฺหาปุจฺฉก}
\proseitem{482}{สตฺต โพชฺฌงฺคา สติสมฺโพชฺฌงฺโค ธมฺมวิจยสมฺโพชฺฌงฺโค วีริยสมฺโพชฺฌงฺโค ปีติสมฺโพชฺฌงฺโค ปสฺสทฺธิสมฺโพชฺฌงฺโค สมาธิสมฺโพชฺฌงฺโค อุเปกฺขาสมฺโพชฺฌงฺโค.}

\proseitem{483}{\csromanfolio{242}สตฺตนฺนํ โพชฺฌงฺคานํ กติ กุสลา, กติ อกุสลา, กติ อพฺยากตา ฯเปฯ กติ สรณา, กติ อรณา.}

\csromanmatikaanchor{mtk.95}
\csromantocmarkref{subparagraph}{1. ติก}{mtk.95}
\bookmark[dest={mtk.95},level=5]{1. ติก}
\csromanheader{6}{1. ติก}
\proseitem{484}{สิยา กุสลา, สิยา อพฺยากตา.\csromanspacer{} ปีติสมฺโพชฺฌงฺโค สุขาย เวทนาย สมฺปยุตฺโต.\csromanspacer{} ฉ โพชฺฌงฺคา สิยา สุขาย เวทนาย สมฺปยุตฺตา, สิยา อทุกฺขมสุขาย เวทนาย สมฺปยุตฺตา.\csromanspacer{} สิยา วิปากา, สิยา วิปากธมฺมธมฺมา.\csromanspacer{} อนุปาทินฺน-อนุปาทานิยา.\csromanspacer{} อสํกิลิฏฺฐ- อสํกิเลสิกา.\csromanspacer{} สิยา สวิตกฺกสวิจารา, สิยา อวิตกฺกวิจารมตฺตา, สิยา อวิตกฺก-อวิจารา.\csromanspacer{} ปีติสมฺโพชฺฌงฺโค น ปีติสหคโต, สุขสหคโต, น อุเปกฺขาสหคโต.\csromanspacer{} ฉ โพชฺฌงฺคา สิยา ปีติสหคตา, สิยา สุขสหคตา, สิยา อุเปกฺขาสหคตา.\csromanspacer{} เนว ทสฺสเนน น ภาวนาย ปหาตพฺพา, เนว ทสฺสเนน น ภาวนาย ปหาตพฺพเหตุกา.\csromanspacer{} สิยา อปจยคามิโน, สิยา เนวาจยคามินาปจยคามิโน.\csromanspacer{} สิยา เสกฺขา, สิยา อเสกฺขา.\csromanspacer{} อปฺปมาณา.\csromanspacer{} อปฺปมาณารมฺมณา.\csromanspacer{} ปณีตา.\csromanspacer{} สิยา สมฺมตฺตนิยตา, สิยา อนิยตา.\csromanspacer{} น มคฺคารมฺมณา, สิยา มคฺคเหตุกา.\csromanspacer{} สิยา มคฺคาธิปติโน, สิยา น วตฺตพฺพา “มคฺคเหตุกา”ติปิ, “มคฺคาธิปติโน”ติปิ.\csromanspacer{} สิยา อุปฺปนฺนา, สิยา อนุปฺปนฺนา, สิยา อุปฺปาทิโน.\csromanspacer{} สิยา อตีตา, สิยา อนาคตา, สิยา ปจฺจุปฺปนฺนา.\csromanspacer{} น วตฺตพฺพา “อตีตารมฺมณา”ติปิ, “อนาคตารมฺมณา”ติปิ, “ปจฺจุปฺปนฺนารมฺมณา”ติปิ.\csromanspacer{} สิยา อชฺฌตฺตา, สิยา พหิทฺธา, สิยา อชฺฌตฺตพหิทฺธา.\csromanspacer{} พหิทฺธารมฺมณา.\csromanspacer{} อนิทสฺสน-อปฺปฏิฆา.}

\csromanmatikaanchor{mtk.96}
\csromantocmarkref{subparagraph}{2. ทุก}{mtk.96}
\bookmark[dest={mtk.96},level=5]{2. ทุก}
\csromanheader{6}{2. ทุก}
\proseitem{485}{ธมฺมวิจยสมฺโพชฺฌงฺโค เหตุ.\csromanspacer{} ฉ โพชฺฌงฺคา เหตุสมฺปยุตฺตา.\csromanspacer{} ธมฺมวิจยสมฺโพชฺฌงฺโค เหตุ เจว สเหตุโก จ, ฉ โพชฺฌงฺคา น วตฺตพฺพา “เหตู เจว สเหตุกา จา”ติ, สเหตุกา เจว น จ เหตู.\csromanspacer{} ธมฺมวิจยสมฺโพชฺฌงฺโค เหตุ เจว เหตุสมฺปยุตฺโต จ.\csromanspacer{} ฉ โพชฺฌงฺคา น วตฺตพฺพา “เหตู เจว เหตุสมฺปยุตฺตา จา”ติ, เหตุสมฺปยุตฺตา เจว น จ เหตู.\csromanspacer{} ฉ โพชฺฌงฺคา น เหตู สเหตุกา.\csromanspacer{} ธมฺมวิจยสมฺโพชฺฌงฺโค น วตฺตพฺโพ “น เหตุสเหตุโก”ติปิ, “น เหตุ-อเหตุโก”ติปิ.\csromanspacer{} สปฺปจฺจยา, สงฺขตา, อนิทสฺสนา, อปฺปฏิฆา, อรูปา, โลกุตฺตรา, เกนจิ วิญฺเญยฺยา,}

\vspace{\tipitakaparskip}
\csromancenter{โพชฺฌงฺควิภงฺค เกนจิ น วิญฺเญยฺยา.\csromanspacer{} โน อาสวา, อนาสวา, อาสววิปฺปยุตฺตา, น วตฺตพฺพา “อาสวา เจว สาสวา จา”ติปิ, “สาสวา เจว โน จ อาสวา”ติปิ.\csromanspacer{} น วตฺตพฺพา “อาสวา เจว อาสวสมฺปยุตฺตา จา”ติปิ, “อาสวสมฺปยุตฺตา เจว โน จ อาสวา”ติปิ, อาสววิปฺปยุตฺตา อนาสวา.\csromanspacer{} โน.\csromanspacer{} สํโยชนา ฯเปฯ.\csromanspacer{} โน คนฺถา ฯเปฯ.\csromanspacer{} โน โอฆา ฯเปฯ.\csromanspacer{} โน โยคา ฯเปฯ.\csromanspacer{} โน นีวรณา ฯเปฯ.\csromanspacer{} โน ปรามาสา ฯเปฯ สารมฺมณา, โน จิตฺตา, เจตสิกา, จิตฺตสมฺปยุตฺตา, จิตฺตสํสฏฺฐา, จิตฺตสมุฏฺฐานา, จิตฺตสหภุโน, จิตฺตานุปริวตฺติโน, จิตฺตสํสฏฺฐสมุฏฺฐานา, จิตฺตสํสฏฺฐสมุฏฺฐานสหภุโน, จิตฺตสํสฏฺฐสมุฏฺฐานานุปริวตฺติโน, พาหิรา.\csromanspacer{} โน อุปาทา, อนุปาทินฺนา, โน อุปาทานา ฯเปฯ.\csromanspacer{} โน กิเลสา ฯเปฯ.\csromanspacer{} น ทสฺสเนน ปหาตพฺพา, น ภาวนาย ปหาตพฺพา.\csromanspacer{} น ทสฺสเนน ปหาตพฺพเหตุกา, น ภาวนาย ปหาตพฺพเหตุกา.\csromanspacer{} สิยา สวิตกฺกา, สิยา อวิตกฺกา.\csromanspacer{} สิยา สวิจารา, สิยา อวิจารา.}
\prose{\csromanfolio{243}ปีติสมฺโพชฺฌงฺโค อปฺปีติโก.\csromanspacer{} ฉ โพชฺฌงฺคา สิยา สปฺปีติกา, สิยา อปฺปีติกา.\csromanspacer{} ปีติสมฺโพชฺฌงฺโค น ปีติสหคโต.\csromanspacer{} ฉ โพชฺฌงฺคา สิยา ปีติสหคตา, สิยา น ปีติสหคตา.\csromanspacer{} ปีติสมฺโพชฺฌงฺโค สุขสหคโต.\csromanspacer{} ฉ โพชฺฌงฺคา สิยา สุขสหคตา, สิยา น สุขสหคตา.\csromanspacer{} ปีติสมฺโพชฺฌงฺโค น อุเปกฺขาสหคโต.\csromanspacer{} ฉ โพชฺฌงฺคา สิยา อุเปกฺขาสหคตา, สิยา น อุเปกฺขาสหคตา, น กามาวจรา, น รูปาวจรา, น อรูปาวจรา, อปริยาปนฺนา, สิยา นิยฺยานิกา, สิยา อนิยฺยานิกา, สิยา นิยตา, สิยา อนิยตา, อนุตฺตรา, อรณาติ.}

\vspace{\tipitakaparskip}
\csromancenter{ปญฺหาปุจฺฉกํ.}
\nitthitam{\textbf{โพชฺฌงฺควิภงฺโค นิฏฺฐิโต.}}
\csromanmatikaanchor{mtk.97}
\csromantocmarknopage{subparagraph}{11. มคฺคงฺควิภงฺค}{mtk.97}
\bookmark[dest={mtk.97},level=5]{11. มคฺคงฺควิภงฺค}
\csromanheader{6}{11. มคฺคงฺควิภงฺค}
\csromanmatikaanchor{mtk.98}
\csromantocmarkref{subparagraph}{1. สุตฺตนฺตภาชนีย}{mtk.98}
\bookmark[dest={mtk.98},level=5]{1. สุตฺตนฺตภาชนีย}
\csromanheader{6}{1. สุตฺตนฺตภาชนีย}
\proseitem{486}{\csromanfolio{244}อริโย อฏฺฐงฺคิโก มคฺโค.\csromanspacer{} เสยฺยถิทํ, สมฺมาทิฏฺฐิ สมฺมาสงฺกปฺโป สมฺมาวาจา สมฺมากมฺมนฺโต สมฺมา-อาชีโว สมฺมาวายาโม สมฺมาสติ สมฺมาสมาธิ.}

\proseitem{487}{ตตฺถ กตมา สมฺมาทิฏฺฐิ, ทุกฺเข ญาณํ ทุกฺขสมุทเย ญาณํ ทุกฺขนิโรเธ ญาณํ ทุกฺขนิโรธคามินิยา ปฏิปทาย ญาณํ, อยํ วุจฺจติ สมฺมาทิฏฺฐิ. (1)}

\prose{ตตฺถ กตโม สมฺมาสงฺกปฺโป, เนกฺขมฺมสงฺกปฺโป อพฺยาปาทสงฺกปฺโป อวิหิํสาสงฺกปฺโป, อยํ วุจฺจติ สมฺมาสงฺกปฺโป. (2)}

\prose{ตตฺถ กตมา สมฺมาวาจา, มุสาวาทา เวรมณี, ปิสุณาย วาจาย เวรมณี, ผรุสาย วาจาย เวรมณี, สมฺผปฺปลาปา เวรมณี, อยํ วุจฺจติ สมฺมาวาจา. (3)}

\prose{ตตฺถ กตโม สมฺมากมฺมนฺโต, ปาณาติปาตา เวรมณี, อทินฺนาทานา เวรมณี, กาเมสุมิจฺฉาจารา เวรมณี, อยํ วุจฺจติ สมฺมากมฺมนฺโต. (4)}

\prose{ตตฺถ กตโม สมฺมา-อาชีโว, อิธ อริยสาวโก มิจฺฉา-อาชีวํ ปหาย สมฺมา-อาชีเวน ชีวิกํ กปฺเปติ, อยํ วุจฺจติ สมฺมา-อาชีโว. (5)}

\prose{ตตฺถ กตโม สมฺมาวายาโม, อิธ ภิกฺขุ อนุปฺปนฺนานํ ปาปกานํ อกุสลานํ ธมฺมานํ อนุปฺปาทาย ฉนฺทํ ชเนติ วายมติ วีริยํ อารภติ จิตฺตํ ปคฺคณฺหาติ ปทหติ, อุปฺปนฺนานํ ปาปกานํ อกุสลานํ ธมฺมานํ ปหานาย ฉนฺทํ ชเนติ วายมติ วีริยํ อารภติ จิตฺตํ ปคฺคณฺหาติ ปทหติ, อนุปฺปนฺนานํ กุสลานํ ธมฺมานํ อุปฺปาทาย ฉนฺทํ ชเนติ วายมติ วีริยํ อารภติ จิตฺตํ ปคฺคณฺหาติ ปทหติ, อุปฺปนฺนานํ กุสลานํ ธมฺมานํ ฐิติยา อสมฺโมสาย ภิยฺโยภาวาย เวปุลฺลาย ภาวนาย ปาริปูริยา ฉนฺทํ ชเนติ วายมติ วีริยํ อารภติ จิตฺตํ ปคฺคณฺหาติ ปทหติ, อยํ วุจฺจติ สมฺมาวายาโม. (6)}

\prose{ตตฺถ กตมา สมฺมาสติ, อิธ ภิกฺขุ กาเย กายานุปสฺสี วิหรติ อาตาปี สมฺปชาโน สติมา วิเนยฺย โลเก อภิชฺฌาโทมนสฺสํ,}

\vspace{\tipitakaparskip}
\csromancenter{มคฺคงฺควิภงฺค เวทนาสุ เวทนานุปสฺสี วิหรติ อาตาปี สมฺปชาโน สติมา วิเนยฺย โลเก อภิชฺฌาโทมนสฺสํ, จิตฺเต จิตฺตานุปสฺสี วิหรติ อาตาปี สมฺปชาโน สติมา วิเนยฺย โลเก อภิชฺฌาโทมนสฺสํ, ธมฺเมสุ ธมฺมานุปสฺสี วิหรติ อาตาปี สมฺปชาโน สติมา วิเนยฺย โลเก อภิชฺฌาโทมนสฺสํ, อยํ วุจฺจติ สมฺมาสติ. (7)}
\prose{\csromanfolio{245}ตตฺถ กตโม สมฺมาสมาธิ, อิธ ภิกฺขุ วิวิจฺเจว กาเมหิ วิวิจฺจ อกุสเลหิ ธมฺเมหิ สวิตกฺกํ สวิจารํ วิเวกชํ ปีติสุขํ ปฐมํ ชานํ อุปสมฺปชฺช วิหรติ, วิตกฺกวิจารานํ วูปสมา อชฺฌตฺตํ สมฺปสาทนํ เจตโส เอโกทิภาวํ อวิตกฺกํ อวิจารํ สมาธิชํ ปีติสุขํ ทุติยํ ฌานํ อุปสมฺปชฺช วิหรติ, ปีติยา จ วิราคา อุเปกฺขโก จ วิหรติ, สโต จ สมฺปชาโน สุขญฺจ กาเยน ปฏิสํเวเทติ, ยํ ตํ อริยา อาจิกฺขนฺติ “อุเปกฺขโก สติมา สุขวิหารี”ติ, ตติยํ ฌานํ อุปสมฺปชฺช วิหรติ, สุขสฺส จ ปหานา ทุกฺขสฺส จ ปหานา ปุพฺเพว โสมนสฺสโทมนสฺสานํ อตฺถงฺคมา อทุกฺขมสุขํ อุเปกฺขาสติปาริสุทฺธิํ จตุตฺถํ ฌานํ อุปสมฺปชฺช วิหรติ, อยํ วุจฺจติ สมฺมาสมาธิ. (8)}

\proseitem{488}{อริโย อฏฺฐงฺคิโก มคฺโค.\csromanspacer{} เสยฺยถิทํ, สมฺมาทิฏฺฐิ สมฺมาสงฺกปฺโป สมฺมาวาจา สมฺมากมฺมนฺโต สมฺมา-อาชีโว สมฺมาวายาโม สมฺมาสติ สมฺมาสมาธิ.}

\proseitem{489}{ตตฺถ กตมา สมฺมาทิฏฺฐิ, อิธ ภิกฺขุ สมฺมาทิฏฺฐิํ ภาเวติ วิเวกนิสฺสิตํ วิราคนิสฺสิตํ นิโรธนิสฺสิตํ โวสฺสคฺคปริณามิํ.\csromanspacer{} สมฺมาสงฺกปฺปํ ภาเวติ ฯเปฯ สมฺมาวาจํ ภาเวติ ฯเปฯ สมฺมากมฺมนฺตํ ภาเวติ ฯเปฯ สมฺมา อาชีวํ ภาเวติ ฯเปฯ สมฺมาวายามํ ภาเวติ ฯเปฯ สมฺมาสติํ ภาเวติ ฯเปฯ สมฺมาสมาธิํ ภาเวติ วิเวกนิสฺสิตํ วิราคนิสฺสิตํ นิโรธนิสฺสิตํ โวสฺสคฺคปริณามิํ.}

\csromanheader{2}{สุตฺตนฺตภาชนียํ.}
\csromansectionrule
\csromanmatikaanchor{mtk.99}
\csromantocmarkref{subparagraph}{2. อภิธมฺมภาชนีย}{mtk.99}
\bookmark[dest={mtk.99},level=5]{2. อภิธมฺมภาชนีย}
\csromanheader{6}{2. อภิธมฺมภาชนีย}
\proseitem{490}{อฏฺฐงฺคิโก มคฺโค สมฺมาทิฏฺฐิ สมฺมาสงฺกปฺโป สมฺมาวาจา สมฺมากมฺมนฺโต สมฺมา-อาชีโว สมฺมาวายาโม สมฺมาสติ สมฺมาสมาธิ.}

\proseitem{491}{\csromanfolio{246}ตตฺถ กตโม อฏฺฐงฺคิโก มคฺโค, อิธ ภิกฺขุ ยสฺมิํ สมเย โลกุตฺตรํ ฌานํ ภาเวติ นิยฺยานิกํ อปจยคามิํ ทิฏฺฐิคตานํ ปหานาย ปฐมาย ภูมิยา ปตฺติยา วิวิจฺเจว กาเมหิ ฯเปฯ ปฐมํ ฌานํ อุปสมฺปชฺช วิหรติ ทุกฺขปฏิปทํ ทนฺธาภิญฺญํ, ตสฺมิํ สมเย อฏฺฐงฺคิโก มคฺโค โหติ สมฺมาทิฏฺฐิ ฯเปฯ สมฺมาสมาธิ.}

\proseitem{492}{ตตฺถ กตมา สมฺมาทิฏฺฐิ, ยา ปญฺญา ปชานนา ฯเปฯ อโมโห ธมฺมวิจโย สมฺมาทิฏฺฐิ ธมฺมวิจยสมฺโพชฺฌงฺโค มคฺคงฺคํ มคฺคปริยาปนฺนํ, อยํ วุจฺจติ สมฺมาทิฏฺฐิ. (1)}

\prose{ตตฺถ กตโม สมฺมาสงฺกปฺโป, โย ตกฺโก วิตกฺโก สงฺกปฺโป อปฺปนา พฺยปฺปนา เจตโส อภินิโรปนา สมฺมาสงฺกปฺโป มคฺคงฺคํ มคฺคปริยาปนฺนํ, อยํ วุจฺจติ สมฺมาสงฺกปฺโป. (2)}

\prose{ตตฺถ กตมา สมาวาจา, ยา จตูหิ วจีทุจฺจริเตหิ อารติ วิรติ ปฏิวิรติ เวรมณี อกิริยา อกรณํ อนชฺฌาปตฺติ เวลา-อนติกฺกโม เสตุฆาโต สมฺมาวาจา มคฺคงฺคํ มคฺคปริยาปนฺนํ, อยํ วุจฺจติ สมฺมาวาจา. (3)}

\prose{ตตฺถ กตโม สมฺมากมฺมนฺโต, ยา ตีหิ กายทุจฺจริเตหิ อารติ วิรติ ปฏิวิรติ เวรมณี อกิริยา อกรณํ อนชฺฌาปตฺติ เวลา-อนติกฺกโม เสตุฆาโต สมฺมากมฺมนฺโต มคฺคงฺคํ มคฺคปริยาปนฺนํ, อยํ วุจฺจติ สมฺมากมฺมนฺโต. (4)}

\prose{ตตฺถ กตโม สมฺมา-อาชีโว, ยา มิจฺฉา-อาชีวา อารติ วิรติ ปฏิวิรติ เวรมณี อกิริยา อกรณํ อนชฺฌาปตฺติ เวลา-อนติกฺกโม เสตุฆาโต สมฺมา- อาชีโว มคฺคงฺคํ มคฺคปริยาปนฺนํ, อยํ วุจฺจติ สมฺมา-อาชีโว. (5)}

\prose{ตตฺถ กตโม สมฺมาวายาโม, เย เจตสิโก วีริยารมฺโภ ฯเปฯ สมฺมาวายาโม วีริยสมฺโพชฺฌงฺโค มคฺคงฺคํ มคฺคปริยาปนฺนํ, อยํ วุจฺจติ สมฺมาวายาโม. (6)}

\prose{ตตฺถ กตมา สมฺมาสติ, ยา สติ อนุสฺสติ ฯเปฯ สมฺมาสติ สติสมฺโพชฺฌงฺโค มคฺคงฺคํ มคฺคปริยาปนฺนํ, อยํ วุจฺจติ สมฺมาสติ. (7)}

\csromanheader{2}{11. มคฺคงฺควิภงฺค}
\prose{\csromanfolio{247}ตตฺถ กตโม สมฺมาสมาธิ, ยา จิตฺตสฺส ฐิติ ฯเปฯ สมฺมาสมาธิ สมาธิสมฺโพชฺฌงฺโค มคฺคงฺคํ มคฺคปริยาปนฺนํ, อยํ วุจฺจติ สมฺมาสมาธิ.\csromanspacer{} อยํ วุจฺจติ อฏฺฐงฺคิโก มคฺโค.\csromanspacer{} อวเสสา ธมฺมา อฏฺฐงฺคิเกน มคฺเคน สมฺปยุตฺตา. (8)}

\proseitem{493}{ปญฺจงฺคิโก มคฺโค สมฺมาทิฏฺฐิ สมฺมาสงฺกปฺโป สมฺมาวายาโม สมฺมาสติ สมฺมาสมาธิ.}

\proseitem{494}{ตตฺถ กตโม ปญฺจงฺคิโก มคฺโค, อิธ ภิกฺขุ ยสฺมิํ สมเย โลกุตฺตรํ ฌานํ ภาเวติ นิยฺยานิกํ อปจยคามิํ ทิฏฺฐิคตานํ ปหานาย ปฐมาย ภูมิยา ปตฺติยา วิวิจฺเจว กเมหิ ฯเปฯ ปฐมํ ฌานํ อุปสมฺปชฺช วิหรติ ทุกฺขปฏิปทํ ทนฺธาภิญฺญํ, ตสฺมิํ สมเย ปญฺจงฺคิโก มคฺโค โหติ สมฺมาทิฏฺฐิ สมฺมาสงฺกปฺโป สมฺมาวายาโม สมฺมาสติ สมฺมาสมาธิ.}

\proseitem{495}{ตตฺถ กตมา สมฺมาทิฏฺฐิ, ยา ปญฺญา ปชานนา ฯเปฯ อโมโห ธมฺมวิจโย สมฺมาทิฏฺฐิ ธมฺมวิจยสมฺโพชฺฌงฺโค มคฺคงฺคํ มคฺคปริยาปนฺนํ, อยํ วุจฺจติ สมฺมาทิฏฺฐิ.}

\prose{ตตฺถ กตโม สมฺมาสงฺกปฺโป, โย ตกฺโก วิตกฺโก ฯเปฯ สมฺมาสงฺกปฺโป มคฺคงฺคํ มคฺคปริยาปนฺนํ, อยํ วุจฺจติ สมฺมาสงฺกปฺโป.}

\prose{ตตฺถ กตโม สมฺมาวายาโม, โย เจตสิโก วีริยารมฺโภ ฯเปฯ สมฺมาวายาโม วีริยสมฺโพชฺฌงฺโค มคฺคงฺคํ มคฺคปริยาปนฺนํ, อยํ วุจฺจติ สมฺมาวายาโม.}

\prose{ตตฺถ กตมา สมฺมาสติ, ยา สติ อนุสฺสติ ฯเปฯ สมฺมาสติ สติสมฺโพชฺฌงฺโค มคฺคงฺคํ มคฺคปริยาปนฺนํ, อยํ วุจฺจติ สมฺมาสติ.}

\prose{ตตฺถ กตโม สมฺมาสมาธิ, ยา จิตฺตสฺส ฐิติ ฯเปฯ สมฺมาสมาธิ สมาธิสมฺโพชฺฌงฺโค มคฺคงฺคํ มคฺคปริยาปนฺนํ, อยํ วุจฺจติ สมฺมาสมาธิ.\csromanspacer{} อยํ วุจฺจติ ปญฺจงฺคิโก มคฺโค.\csromanspacer{} อวเสสา ธมฺมา ปญฺจงฺคิเกน มคฺเคน สมฺปยุตฺตา.}

\proseitem{496}{ปญฺจงฺคิโก มคฺโค สมฺมาทิฏฺฐิ สมฺมาสงฺกปฺโป สมฺมาวายาโม สมฺมาสติ สมฺมาสมาธิ.}

\proseitem{497}{\csromanfolio{248}ตตฺถ กตมา สมฺมาทิฏฺฐิ, อิธ ภิกฺขุ ยสฺมิํ สมเย โลกุตฺตรํ ฌานํ ภาเวติ นิยฺยานิกํ อปจยคามิํ ทิฏฺฐิคตานํ ปหานาย ปฐมาย ภูมิยา ปตฺติยา วิวิจฺเจว กาเมหิ ฯเปฯ ปฐมํ ฌานํ อุปสมฺปชฺช วิหรติ ทุกฺขปฏิปทํ ทนฺธาภิญฺญํ, ยา ตสฺมิํ สมเย ปญฺญา ปชานนา ฯเปฯ อโมโห ธมฺมวิจโย สมฺมาทิฏฺฐิ ธมฺมวิจยสมฺโพชฺฌงฺโค มคฺคงฺคํ มคฺคปริยาปนฺนํ, อยํ วุจฺจติ สมฺมาทิฏฺฐิ.\csromanspacer{} อวเสสา ธมฺมา สมฺมาทิฏฺฐิยา สมฺปยุตฺตา ฯเปฯ.\csromanspacer{} อวเสสา ธมฺมา สมฺมาสงฺกปฺเปน สมฺปยุตฺตา ฯเปฯ.\csromanspacer{} อวเสสา ธมฺมา สมฺมาวายาเมน สมฺปยุตฺตา ฯเปฯ.\csromanspacer{} อวเสสา ธมฺมา สมฺมาสติยา สมฺปยุตฺตา.}

\prose{ตตฺถ กตโม สมฺมาสมาธิ, อิธ ภิกฺขุ ยสฺมิํ สมเย โลกุตฺตรํ ฌานํ ภาเวติ นิยฺยานิกํ อปจยคามิํ ทิฏฺฐิคตานํ ปหานาย ปฐมาย ภูมิยา ปตฺติยา วิวิจฺเจว กาเมหิ ฯเปฯ ปฐมํ ฌานํ อุปสมฺปชฺช วิหรติ ทุกฺขปฏิปทํ ทนฺธาภิญฺญํ, ยา ตสฺมิํ สมเย จิตฺตสฺส ฐิติ ฯเปฯ สมฺมาสมาธิ สมาธิสมฺโพชฺฌงฺโค มคฺคงฺคํ มคฺคปริยาปนฺนํ, อยํ วุจฺจติ สมฺมาสมาธิ.\csromanspacer{} อวเสสา ธมฺมา สมฺมาสมาธินา สมฺปยุตฺตา.}

\proseitem{498}{อฏฺฐงฺคิโก มคฺโค สมฺมาทิฏฺฐิ ฯเปฯ สมฺมาสมาธิ.}

\proseitem{499}{ตตฺถ กตโม อฏฺฐงฺคิโก มคฺโค, อิธ ภิกฺขุ ยสฺมิํ สมเย โลกุตฺตรํ ฌานํ ภาเวติ นิยฺยานิกํ อปจยคามิํ ทิฏฺฐิคตานํ ปหานาย ปฐมาย ภูมิยา ปตฺติยา วิวิจฺเจว กาเมหิ ฯเปฯ ปฐมํ ฌานํ อุปสมฺปชฺช วิหรติ ทุกฺขปฏิปทํ ทนฺธาภิญฺญํ, ตสฺมิํ สมเย ผสฺโส โหติ ฯเปฯ อวิกฺเขโป โหติ, อิเม ธมฺมา กุสลา.\csromanspacer{} ตสฺเสว โลกุตฺตรสฺส กุสลสฺส ฌานสฺส กตตฺตา ภาวิตตฺตา วิปากํ วิวิจฺเจว กาเมหิ ฯเปฯ ปฐมํ ฌานํ อุปสมฺปชฺช วิหรติ ทุกฺขปฏิปทํ ทนฺธาภิญฺญํ สุญฺญตํ, ตสฺมิํ สมเย อฏฺฐงฺคิโก มคฺโค โหติ สมฺมาทิฏฺฐิ ฯเปฯ สมฺมาสมาธิ, อยํ วุจฺจติ อฏฺฐงฺคิโก มคฺโค.\csromanspacer{} อวเสสา ธมฺมา อฏฺฐงฺคิเกน มคฺเคน สมฺปยุตฺตา.}

\proseitem{500}{ปญฺจงฺคิโก มคฺโค สมฺมาทิฏฺฐิ สมฺมาสงฺกปฺโป สมฺมาวายาโม สมฺมาสติ สมฺมาสมาธิ.}

\proseitem{501}{ตตฺถ กตโม ปญฺจงฺคิโก มคฺโค, อิธ ภิกฺขุ ยสฺมิํ สมเย โลกุตฺตรํ ฌานํ ภาเวติ นิยฺยานิกํ อปจยคามิํ ทิฏฺฐิคตานํ ปหานาย ปฐมาย ภูมิยา ปตฺติยา วิวิจฺเจว กาเมหิ ฯเปฯ ปฐมํ ฌานํ อุปสมฺปชฺช}

\vspace{\tipitakaparskip}
\csromancenter{มคฺคงฺควิภงฺค วิหรติ ทุกฺขปฏิปทํ ทนฺธาภิญฺญํ, ตสฺมิํ สมเย ผสฺโส โหติ ฯเปฯ อวิกฺเขโป โหติ, อิเม ธมฺมา กุสลา.\csromanspacer{} ตสฺเสว โลกุตฺตรสฺส กุสลสฺส ฌานสฺส กตตฺตา ภาวิตตฺตา วิปากํ วิวิจฺเจว กาเมหิ ฯเปฯ ปฐมํ ฌานํ อุปสมฺปชฺช วิหรติ ทุกฺขปฏิปทํ ทนฺธาภิญฺญํ สุญฺญตํ, ตสฺมิํ สมเย ปญฺจงฺคิโก มคฺโค โหติ สมฺมาทิฏฺฐิ สมฺมาสงฺกปฺโป สมฺมาวายาโม สมฺมาสติ สมฺมาสมาธิ, อยํ วุจฺจติ ปญฺจงฺคิโก มคฺโค.\csromanspacer{} อวเสสา ธมฺมา ปญฺจงฺคิเกน มคฺเคน สมฺปยุตฺตา.}
\proseitem{502}{\csromanfolio{249}ปญฺจงฺคิโก มคฺโค สมฺมาทิฏฺฐิ สมฺมาสงฺกปฺโป สมฺมาวายาโม สมฺมาสติ สมฺมาสมาธิ.}

\proseitem{503}{ตตฺถ กตมา สมฺมาทิฏฺฐิ, อิธ ภิกฺขุ ยสฺมิํ สมเย โลกุตฺตรํ ฌานํ ภาเวติ นิยฺยานิกํ อปจยคามิํ ทิฏฺฐิคตานํ ปหานาย ปฐมาย ภูมิยา ปตฺติยา วิวิจฺเจว กาเมหิ ฯเปฯ ปฐมํ ฌานํ อุปสมฺปชฺช วิหรติ ทุกฺขปฏิปทํ ทนฺธาภิญฺญํ, ตสฺมิํ สมเย ผสฺโส โหติ ฯเปฯ อวิกฺเขโป โหติ, อิเม ธมฺมา กุสลา.\csromanspacer{} ตสฺเสว โลกุตฺตรสฺส กุสลสฺส ฌานสฺส กตตฺตา ภาวิตตฺตา วิปากํ วิวิจฺเจว กาเมหิ ฯเปฯ ปฐมํ ฌานํ อุปสมฺปชฺช วิหรติ ทุกฺขปฏิปทํ ทนฺธาภิญฺญํ สุญฺญตํ, ยา ตสฺมิํ สมเย ปญฺญา ปชานนา ฯเปฯ อโมโห ธมฺมวิจโย สมฺมาทิฏฺฐิ ธมฺมวิจยสมฺโพชฺฌงฺโค มคฺคงฺคํ มคฺคปริยาปนฺนํ, อยํ วุจฺจติ สมฺมาทิฏฺฐิ.\csromanspacer{} อวเสสา ธมฺมา สมฺมาทิฏฺฐิยา สมฺปยุตฺตา ฯเปฯ.\csromanspacer{} อวเสสา ธมฺมา สมฺมาสงฺกปฺเปน สมฺปยุตฺตา ฯเปฯ.\csromanspacer{} อวเสสา ธมฺมา สมฺมาวายาเมน สมฺปยุตฺตา ฯเปฯ.\csromanspacer{} อวเสสา ธมฺมา สมฺมาสติยา สมฺปยุตฺตา.}

\prose{ตตฺถ กตโม สมฺมาสมาธิ, อิธ ภิกฺขุ ยสฺมิํ สมเย โลกุตฺตรํ ฌานํ ภาเวติ นิยฺยานิกํ อปจยคามิํ ทิฏฺฐิคตานํ ปหานาย ปฐมาย ภูมิยา ปตฺติยา วิวิจฺเจว กาเมหิ ฯเปฯ ปฐมํ ฌานํ อุปสมฺปชฺช วิหรติ ทุกฺขปฏิปทํ ทนฺธาภิญฺญํ, ตสฺมิํ สมเย ผสฺโส โหติ ฯเปฯ อวิกฺเขโป โหติ, อิเม ธมฺมา กุสลา.\csromanspacer{} ตสฺเสว โลกุตฺตรสฺส กุสลสฺส ฌานสฺส กตตฺตา ภาวิตตฺตา วิปากํ วิวิจฺเจว กาเมหิ ฯเปฯ ปฐมํ ฌานํ อุปสมฺปชฺช วิหรติ ทุกฺขปฏิปทํ ทนฺธาภิญฺญํ สุญฺญตํ, ยา ตสฺมิํ สมเย จิตฺตสฺส ฐิติ สณฺฐิติ อวฏฺฐิติ อวิสาหาโร อวิกฺเขโป \csromanfolio{250}อวิสาหฏมานสตา สมโถ สมาธินฺทฺริยํ สมาธิพลํ สมฺมาสมาธิ สมาธิสมฺโพชฺฌงฺโค มคฺคงฺคํ มคฺคปริยาปนฺนํ, อยํ วุจฺจติ สมฺมาสมาธิ.\csromanspacer{} อวเสสา ธมฺมา สมฺมาสมาธินา สมฺปยุตฺตา.}

\csromanheader{2}{อภิธมฺมภาชนียํ.}
\csromansectionrule
\csromanmatikaanchor{mtk.100}
\csromantocmarkref{subparagraph}{3. ปญฺหาปุจฺฉก}{mtk.100}
\bookmark[dest={mtk.100},level=5]{3. ปญฺหาปุจฺฉก}
\csromanheader{6}{3. ปญฺหาปุจฺฉก}
\proseitem{504}{อริโย อฏฺฐงฺคิโก มคฺโค.\csromanspacer{} เสยฺยถิทํ, สมฺมาทิฏฺฐิ สมฺมาสงฺกปฺโป สมฺมาวาจา สมฺมากมฺมนฺโต สมฺมา-อาชีโว สมฺมาวายาโม สมฺมาสติ สมฺมาสมาธิ.}

\proseitem{505}{อฏฺฐนฺนํ มคฺคงฺคานํ กติ กุสลา, กติ อกุสลา, กติ อพฺยากตา ฯเปฯ กติ สรณา, กติ อรณา.}

\csromanmatikaanchor{mtk.101}
\csromantocmarkref{subparagraph}{1. ติก}{mtk.101}
\bookmark[dest={mtk.101},level=5]{1. ติก}
\csromanheader{6}{1. ติก}
\proseitem{506}{สิยา กุสลา, สิยา อพฺยากตา.\csromanspacer{} สมฺมาสงฺกปฺโป สุขาย เวทนาย สมฺปยุตฺตา.\csromanspacer{} สตฺต มคฺคงฺคา สิยา สุขาย เวทนาย สมฺปยุตฺตา, สิยา อทุกฺขมสุขาย เวทนาย สมฺปยุตฺตา.\csromanspacer{} สิยา วิปากา, สิยา วิปากธมฺมธมฺมา.\csromanspacer{} อนุปาทินฺน-อนุปาทานิยา.\csromanspacer{} อสํกิลิฏฺฐ- อสํกิเลสิกา.\csromanspacer{} สมฺมาสงฺกปฺโป อวิตกฺกวิจารมตฺโต.\csromanspacer{} สตฺต มคฺคงฺคา สิยา สวิตกฺกสวิจารา, สิยา อวิตกฺกวิจารมตฺตา, สิยา อวิตกฺก-อวิจารา.\csromanspacer{} สมฺมาสงฺกปฺโป ปีติสหคโต, สุขสหคโต, น อุเปกฺขาสหคโต.\csromanspacer{} สตฺต มคฺคงฺคา สิยา ปีติสหคตา, สิยา สุขสหคตา, สิยา อุเปกฺขาสหคตา.\csromanspacer{} เนว ทสฺสเนน น ภาวนาย ปหาตพฺพา, เนว ทสฺสเนน น ภาวนาย ปหาตพฺพเหตุกา.\csromanspacer{} สิยา อปจยคามิโน, สิยา เนวาจยคามินาปจยคามิโน.\csromanspacer{} สิยา เสกฺขา, สิยา อเสกฺขา.\csromanspacer{} อปฺปมาณา, อปฺปมาณารมฺมณา, ปณีตา, สิยา สมฺมตฺตนิยตา, สิยา อนิยตา, น มคฺคารมฺมณา, สิยา มคฺคเหตุกา, สิยา มคฺคาธิปติโน, สิยา น วตฺตพฺพา “มคฺคเหตุกา”ติปิ. “มคฺคาธิปติโน”ติปิ.\csromanspacer{} สิยา อุปฺปนฺนา, สิยา อนุปฺปนฺนา, สิยา อุปฺปาทิโน.\csromanspacer{} สิยา อตีตา, สิยา อนาคตา,}

\vspace{\tipitakaparskip}
\csromancenter{มคฺคงฺควิภงฺค สิยา ปจฺจุปฺปนฺนา, น วตฺตพฺพา “อตีตารมฺมณา”ติปิ, “อนาคตารมฺมณา”ติปิ, “ปจฺจุปฺปนฺนารมฺมณา”ติปิ.\csromanspacer{} สิยา อชฺฌตฺตา, สิยา พหิทฺธา, สิยา อชฺฌตฺตพหิทฺธา.\csromanspacer{} พหิทฺธารมฺมณา, อนิทสฺสน-อปฺปฏิฆา.}
\csromanmatikaanchor{mtk.102}
\csromantocmarkref{subparagraph}{2. ทุก}{mtk.102}
\bookmark[dest={mtk.102},level=5]{2. ทุก}
\csromanheader{6}{2. ทุก}
\proseitem{507}{\csromanfolio{251}สมฺมาทิฏฺฐิ เหตุ.\csromanspacer{} สตฺต มคฺคงฺคา น เหตู, สเหตุกา, เหตุสมฺปยุตฺตา.\csromanspacer{} สมฺมาทิฏฺฐิ เหตุ เจว สเหตุกา จ.\csromanspacer{} สตฺต มคฺคงฺคา น วตฺตพฺพา “เหตู เจว สเหตุกา จา”ติ, สเหตุกา เจว น จ เหตู.\csromanspacer{} สมฺมาทิฏฺฐิ เหตุ เจว เหตุสมฺปยุตฺตา จ.\csromanspacer{} สตฺต มคฺคงฺคา น วตฺตพฺพา “เหตู เจว เหตุสมฺปยุตฺตา จา”ติ, เหตุสมฺปยุตฺตา เจว น จ เหตู.\csromanspacer{} สตฺต มคฺคงฺคา น เหตู, สเหตุกา.\csromanspacer{} สมฺมาทิฏฺฐิ น วตฺตพฺพา “น เหตุ สเหตุกา”ติปิ, “น เหตุ อเหตุกา”ติปิ.}

\prose{สปฺปจฺจยา, สงฺขตา, อนิทสฺสนา, อปฺปฏิฆา, อรูปา, โลกุตฺตรา, เกนจิ วิญฺเญยฺยา, เกนจิ น วิญฺเญยฺยา.}

\prose{โน อาสวา, อนาสวา, อาสววิปฺปยุตฺตา, น วตฺตพฺพา “อาสวา เจว สาสวา จา”ติปิ, “สาสวา เจว โน จ อาสวา”ติปิ, น วตฺตพฺพา “อาสวา เจว อาสวสมฺปยุตฺตา จา”ติปิ, อาสวสมฺปยุตฺตา เจว โน จ อาสวา”ติปิ, อาสววิปฺปยุตฺตา อนาสวา.}

\prose{โน สํโยชนา ฯเปฯ.\csromanspacer{} โน คนฺถา ฯเปฯ.\csromanspacer{} โน โอฆา ฯเปฯ.\csromanspacer{} โน โยคา ฯเปฯ.\csromanspacer{} โน นีวรณา ฯเปฯ.\csromanspacer{} โน ปรามาสา ฯเปฯ สารมฺมณา, โน จิตฺตา, เจตสิกา, จิตฺตสมฺปยุตฺตา, จิตฺตสํสฏฺฐา, จิตฺตสมุฏฺฐานา, จิตฺตสหภุโน, จิตฺตานุปริวตฺติโน, จิตฺตสํสฏฺฐสมุฏฺฐานา, จิตฺตสํสฏฺฐสมุฏฺฐานสหภุโน, จิตฺตสํสฏฺฐสมุฏฺฐานานุปริวตฺติโน, พาหิรา, โน อุปาทา, อนุปาทินฺนา.}

\prose{โน อุปาทานา ฯเปฯ.\csromanspacer{} โน กิเลสา ฯเปฯ.\csromanspacer{} น ทสฺสเนน ปหาตพฺพา, น ภาวนาย ปหาตพฺพา, น ทสฺสเนน ปหาตพฺพเหตุกา, น ภาวนาย ปหาตพฺพเหตุกา.\csromanspacer{} สมฺมาสงฺกปฺโป อวิตกฺโก.\csromanspacer{} สตฺต มคฺคงฺคา สิยา สวิตกฺกา, สิยา อวิตกฺกา.\csromanspacer{} สมฺมาสงฺกปฺโป สวิจาโร.\csromanspacer{} สตฺต มคฺคงฺคา สิยา สวิจารา, สิยา อวิจารา.\csromanspacer{} สมฺมาสงฺกปฺโป \csromanfolio{252}สปฺปีติโก.\csromanspacer{} สตฺต มคฺคงฺคา สิยา สปฺปีติกา, สิยา อปฺปีติกา.\csromanspacer{} สมฺมาสงฺกปฺโป ปีติสหคโต.\csromanspacer{} สตฺต มคฺคงฺคา สิยา ปีติสหคตา, สิยา น ปีติสหคตา.\csromanspacer{} สมฺมาสงฺกปฺโป สุขสหคโต.\csromanspacer{} สตฺต มคฺคงฺคา สิยา สุขสหคตา, สิยา น สุขสหคตา.\csromanspacer{} สมฺมาสงฺกปฺโป น อุเปกฺขาสหคโต.\csromanspacer{} สตฺต มคฺคงฺคา สิยา อุเปกฺขาสหคตา, สิยา น อุเปกฺขาสหคตา, น กามาวจรา, น รูปาวจรา, น อรูปาวจรา, อปริยาปนฺนา, สิยา นิยฺยานิกา, สิยา อนิยฺยานิกา, สิยา นิยตา, สิยา อนิยตา, อนุตฺตรา, อรณาติ.}

\vspace{\tipitakaparskip}
\csromancenter{ปญฺหาปุจฺฉกํ.}
\nitthitam{\textbf{มคฺคงฺควิภงฺโค นิฏฺฐิโต.}}
\csromanmatikaanchor{mtk.103}
\csromantocmarknopage{subparagraph}{12. ฌานวิภงฺค}{mtk.103}
\bookmark[dest={mtk.103},level=5]{12. ฌานวิภงฺค}
\csromanheader{6}{12. ฌานวิภงฺค}
\csromanmatikaanchor{mtk.104}
\csromantocmarkref{subparagraph}{1. สุตฺตนฺตภาชนีย}{mtk.104}
\bookmark[dest={mtk.104},level=5]{1. สุตฺตนฺตภาชนีย}
\csromanheader{6}{1. สุตฺตนฺตภาชนีย}
\proseitem{508}{\csromanfolio{253}อิธ ภิกฺขุ ปาติโมกฺขสํวรสํวุโต วิหรติ อาจารโคจรสมฺปนฺโน อณุมตฺเตสุ วชฺเชสุ ภยทสฺสาวี สมาทาย สิกฺขติ สิกฺขาปเทสุ อินฺทฺริเยสุ คุตฺตทฺวาโร โภชเน มตฺตญฺญู ปุพฺพรตฺตาปรรตฺตํ ชาคริยานุโยคมนุยุตฺโต สาตจฺจํ เนปกฺกํ โพธิปกฺขิกานํ ธมฺมานํ ภาวนานุโยคมนุยุตฺโต, โส อภิกฺกนฺเต ปฏิกฺกนฺเต สมฺปชานการี โหติ, อาโลกิเต วิโลกิเต สมฺปชานการี โหติ, สมิญฺชิเต\footnote{สมฺมิญฺชิเต (สี, สฺยา)} ปสาริเต สมฺปชานการี โหติ, สงฺฆาฏิปตฺตจีวรธารเณ สมฺปชานการี โหติ, อสิเต ปีเต ขายิเต สายิเต สมฺปชานการี โหติ, อุจฺจารปสฺสาวกมฺเม สมฺปชานการี โหติ, คเต ฐิเต นิสินฺเน สุตฺเต ชาคริเต ภาสิเต ตุณฺหีภาเว สมฺปชานการี โหติ, โส วิวิตฺตํ เสนาสนํ ภชติ อรญฺญํ รุกฺขมูลํ ปพฺพตํ กนฺทรํ คิริคุหํ สุสานํ วนปตฺถํ อพฺโภกาสํ ปลาลปุญฺชํ อปฺปสทฺทํ อปฺปนิคฺโฆสํ วิชนวาตํ มนุสฺสราหสฺเสยฺยกํ\footnote{มนุสฺสราหเสยฺยกํ (สี, สฺยา)} ปฏิสลฺลานสารุปฺปํ, โส อรญฺญคโต วา รุกฺขมูลคโต วา สุญฺญาคารคโต วา นิสีทติ ปลฺลงฺกํ อาภุชิตฺวา อุชุํ กายํ ปณิธาย ปริมุขํ สติํ อุปฏฺฐเปตฺวา, โส อภิชฺฌํ โลเก ปหาย วิคตาภิชฺเฌน เจตสา วิหรติ, อภิชฺฌาย จิตฺตํ ปริโสเธติ.\csromanspacer{} พฺยาปาทปโทสํ ปหาย อพฺยาปนฺนจิตฺโต วิหรติ สพฺพปาณภูตหิตานุกมฺปี, พฺยาปาทปโทสา จิตฺตํ ปริโสเธติ.\csromanspacer{} ถินมิทฺธํ ปหาย วิคตถินมิทฺโธ วิหรติ อาโลกสญฺญี สโต สมฺปชาโน, ถินมิทฺธา จิตฺตํ ปริโสเธติ.\csromanspacer{} อุทฺธจฺจกุกฺกุจฺจํ ปหาย อนุทฺธโต วิหรติ อชฺฌตฺตํ วูปสนฺตจิตฺโต, อุทฺธจฺจกุกฺกุจฺจา จิตฺตํ ปริโสเธติ.\csromanspacer{} วิจิกิจฺฉํ ปหาย ติณฺณวิจิกิจฺโฉ วิหรติ อกถํกถี กุสเลสุ ธมฺเมสุ, วิจิกิจฺฉาย จิตฺตํ ปริโสเธติ.\csromanspacer{} โส อิเม ปญฺจ นีวรเณ ปหาย เจตโส อุปกฺกิเลเส ปญฺญาย ทุพฺพลีกรเณ วิวิจฺเจว กาเมหิ วิวิจฺจ อกุสเลหิ ธมฺเมหิ สวิตกฺกํ สวิจารํ วิเวกชํ ปีติสุขํ ปฐมํ ฌานํ อุปสมฺปชฺช วิหรติ \csromanfolio{254}วิตกฺกวิจารานํ วูปสมา อชฺฌตฺตํ สมฺปสาทนํ เจตโส เอโกทิภาวํ อวิตกฺกํ อวิจารํ สมาธิชํ ปีติสุขํ ทุติยํ ฌานํ อุปสมฺปชฺช วิหรติ, ปีติยา จ วิราคา อุเปกฺขโก จ วิหรติ สโต จ สมฺปชาโน สุขญฺจ กาเยน ปฏิสํเวเทติ, ยํ ตํ อริยา อาจิกฺขนฺติ “อุเปกฺขโก สติมา สุขวิหารี”ติ, ตติยํ ฌานํ อุปสมฺปชฺช วิหรติ, สุขสฺส จ ปหานา ทุกฺขสฺส จ ปหานา ปุพฺเพว โสมนสฺสโทมนสฺสานํ อตฺถงฺคมา อทุกฺขมสุขํ อุเปกฺขาสติปาริสุทฺธิํ จตุตฺถํ ฌานํ อุปสมฺปชฺช วิหรติ, สพฺพโส รูปสญฺญานํ สมติกฺกมา ปฏิฆสญฺญานํ อตฺถงฺคมา นานตฺตสญฺญานํ อมนสิการา “อนนฺโต อากาโส”ติ อากาสานญฺจายตนํ อุปสมฺปชฺช วิหรติ, สพฺพโส อากาสานญฺจายตนํ สมติกฺกมฺม “อนนฺตํ วิญฺญาณนฺ”ติ วิญฺญาณญฺจายตนํ อุปสมฺปชฺช วิหรติ, สพฺพโส วิญฺญาณญฺจายตนํ สมติกฺกมฺม “นตฺถิ กิญฺจี”ติ อากิญฺจญฺญายตนํ อุปสมฺปชฺช วิหรติ, สพฺพโส อากิญฺจญฺญายตนํ สมติกฺกมฺม เนวสญฺญานาสญฺญายตนํ อุปสมฺปชฺช วิหรติ.\csromanspacer{} มาติกา.}

\proseitem{509}{“อิธา”ติ อิมิสฺสา ทิฏฺฐิยา อิมิสฺสา ขนฺติยา อิมิสฺสา รุจิยา อิมสฺมิํ อาทาเย อิมสฺมิํ ธมฺเม อิมสฺมิํ วินเย อิมสฺมิํ ธมฺมวินเย อิมสฺมิํ ปาวจเน อิมสฺมิํ พฺรหฺมจริเย อิมสฺมิํ สตฺถุสาสเน, เตน วุจฺจติ “อิธา”ติ.}

\proseitem{510}{“ภิกฺขู”ติ สมญฺญาย ภิกฺขุ, ปฏิญฺญาย ภิกฺขุ, ภิกฺขตีติ ภิกฺขุ, ภิกฺขโกติ ภิกฺขุ, ภิกฺขาจริยํ อชฺฌุปคโตติ ภิกฺขุ, ภินฺนปฏธโรติ ภิกฺขุ, ภินฺทติ ปาปเก อกุสเล ธมฺเมติ ภิกฺขุ, ภินฺนตฺตา ปาปกานํ อกุสลานํ ธมฺมานํ ภิกฺขุ, โอธิโส กิเลสานํ ปหานา ภิกฺขุ, อโนธิโส กิเลสานํ ปหานา ภิกฺขุ, เสกฺโข ภิกฺขุ, อเสกฺโข ภิกฺขุ, เนวเสกฺขนาเสกฺโข ภิกฺขุ, อคฺโค ภิกฺขุ, ภโทฺร ภิกฺขุ, มณฺโฑ ภิกฺขุ, สาโร ภิกฺขุ, สมคฺเคน สํเฆน ญตฺติจตุตฺเถน กมฺเมน อกุปฺเปน ฐานารเหน อุปสมฺปนฺโน ภิกฺขุ.}

\proseitem{511}{“ปาติโมกฺขนฺ”ติ สีลํ ปติฏฺฐา อาทิ จรณํ สํยโม สํวโร โมกฺขํ ปาโมกฺขํ กุสลานํ ธมฺมานํ สมาปตฺติยา. “สํวโร”ติ}

\vspace{\tipitakaparskip}
\csromancenter{ฌานวิภงฺค กายิโก อวีติกฺกโม วาจสิโก อวีติกฺกโม กายิกวาจสิโก อวีติกฺกโม. “สํวุโต”ติ อิมินา ปาติโมกฺขสํวเรน อุเปโต โหติ สมุเปโต อุปาคโต สมุปาคโต อุปปนฺโน สมฺปนฺโน สมนฺนาคโต, เตน วุจฺจติ “ปาติโมกฺขสํวรสํวุโต”ติ.}
\proseitem{512}{\csromanfolio{255}“วิหรตี”ติ อิริยติ วตฺตติ ปาเลติ ยเปติ ยาเปติ จรติ วิหรติ, เตน วุจฺจติ “วิหรตี”ติ.}

\proseitem{513}{“อาจารโคจรสมฺปนฺโน”ติ อตฺถิ อาจาโร, อตฺถิ อนาจาโร.}

\prose{ตตฺถ กตโม อนาจาโร, กายิโก วีติกฺกโม วาจสิโก วีติกฺกโม กายิกวาจสิโก วีติกฺกโม, อยํ วุจฺจติ อนาจาโร.\csromanspacer{} สพฺพมฺปิ ทุสฺสีลฺยํ อนาจาโร, อิเธกจฺโจ เวฬุทาเนน วา ปตฺตทาเนน วา ปุปฺผทาเนน วา ผลทาเนน วา สินานทาเนน วา ทนฺตกฏฺฐทาเนน วา จาฏุกมฺยตาย วา มุคฺคสูปฺยตาย\footnote{มุคฺคสุปฺปตาย (สี)} วา ปาริภฏฺยตาย วา ชงฺฆเปสนิเกน วา อญฺญตรญฺญตเรน วา พุทฺธปฏิกุฏฺเฐน มิจฺฉา-อาชีเวน ชีวิกํ กปฺเปติ, อยํ วุจฺจติ อนาจาโร.}

\prose{ตตฺถ กตโม อาจาโร, กายิโก อวีติกฺกโม วาจสิโก อวีติกฺกโม กายิกวาจสิโก อวีติกฺกโม, อยํ วุจฺจติ อาจาโร.\csromanspacer{} สพฺโพปิ สีลสํวโร อาจาโร, อิเธกจฺโจ น เวฬุทาเนน น ปตฺตทาเนน น ปุปฺผทาเนน น ผลทาเนน น สินานทาเนน น ทนฺตกฏฺฐทาเนน น จาฏุกมฺยตาย น มุคฺคสูปฺยตาย น ปาริภฏฺยตาย น ชงฺฆเปสนิเกน น อญฺญตรญฺญตเรน พุทฺธปฏิกุฏฺเฐน มิจฺฉา-อาชีเวน ชีวิกํ กปฺเปติ, อยํ วุจฺจติ อาจาโร.}

\proseitem{514}{“โคจโร”ติ อตฺถิ โคจโร, อตฺถิ อโคจโร.}

\prose{ตตฺถ กตโม อโคจโร, อิเธกจฺโจ เวสิยาโคจโร วา โหติ วิธวาโคจโร วา ถุลฺลกุมาริโคจโร วา ปณฺฑกโคจโร วา ภิกฺขุนิโคจโร วา ปานาคารโคจโร วา, สํสฏฺโฐ วิหรติ ราชูหิ ราชมหามตฺเตหิ ติตฺถิเยหิ ติตฺถิยสาวเกหิ อนนุโลมิเกน \csromanfolio{256}สํสคฺเคน, ยานิ วา ปน ตานิ กุลานิ อสฺสทฺธานิ อปฺปสนฺนานิ อโนปานภูตานิ อกฺโกสกปริภาสกานิ อนตฺถกามานิ อหิตกามานิ อผาสุกกามานิ อโยคกฺเขมกามานิ ภิกฺขูนํ ภิกฺขุนีนํ อุปาสกานํ อุปาสิกานํ, ตถารูปานิ กุลานิ เสวติ ภชติ ปยิรุปาสติ, อยํ วุจฺจติ อโคจโร.}

\prose{ตตฺถ กตโม โคจโร, อิเธกจฺโจ น เวสิยาโคจโร โหติ น วิธวาโคจโร น ถุลฺลกุมาริโคจโร น ปณฺฑกโคจโร น ภิกฺขุนิโคจโร น ปานาคารโคจโร, อสํสฏฺโฐ วิหรติ ราชูหิ ราชมหามตฺเตหิ ติตฺถิเยหิ ติตฺถิยสาวเกหิ อนนุโลมิเกน สํสคฺเคน, ยานิ วา ปน ตานิ กุลานิ สทฺธานิ ปสนฺนานิ โอปานภูตานิ กาสาวปชฺโชตานิ อิสิวาตปฏิวาตานิ อตฺถกามานิ หิตกามานิ ผาสุกกามานิ โยคกฺเขมกามานิ ภิกฺขูนํ ภิกฺขุนีนํ อุปาสกานํ อุปาสิกานํ, ตถารูปานิ กุลานิ เสวติ ภชติ ปยิรุปาสติ, อยํ วุจฺจติ โคจโร.\csromanspacer{} อิติ อิมินา จ อาจาเรน อิมินา จ โคจเรน อุเปโต โหติ ฯเปฯ สมนฺนาคโต, เตน วุจฺจติ “อาจารโคจรสมฺปนฺโน”ติ.}

\proseitem{515}{“อณุมตฺเตสุ วชฺเชสุ ภยทสฺสาวี”ติ ตตฺถ กตเม อณุมตฺตา วชฺชา, ยานิ ตานิ วชฺชานิ อปฺปมตฺตกานิ โอรมตฺตกานิ ลหุสานิ ลหุสมฺมตานิ สํยมกรณียานิ สํวรกรณียานิ จิตฺตุปฺปาทกรณียานิ มนสิการปฏิพทฺธานิ, อิเม วุจฺจนฺติ อณุมตฺตา วชฺชา.\csromanspacer{} อิติ อิเมสุ อณุมตฺเตสุ วชฺเชสุ วชฺชทสฺสาวี จ โหติ ภยทสฺสาวี จ อาทีนวทสฺสาวี จ นิสฺสรณทสฺสาวี จ, เตน วุจฺจติ “อณุมตฺเตสุ วชฺเชสุ ภยทสฺสาวี”ติ.}

\proseitem{516}{“สมาทาย สิกฺขติ สิกฺขาปเทสู”ติ ตตฺถ กตมา สิกฺขา, จตสฺโส สิกฺขา ภิกฺขูนํ ภิกฺขุสิกฺขา, ภิกฺขุนีนํ ภิกฺขุนิสิกฺขา, อุปาสกานํ อุปาสกสิกฺขา, อุปาสิกานํ อุปาสิกสิกฺขา, อิมา วุจฺจนฺติ สิกฺขาโย.\csromanspacer{} อิติ อิมาสุ สิกฺขาสุ สพฺเพน สพฺพํ สพฺพถา สพฺพํ อเสสํ นิสฺเสสํ สมาทาย วตฺตติ, เตน วุจฺจติ “สมาทาย สิกฺขติ สิกฺขาปเทสู”ติ.}

\proseitem{517}{“อินฺทฺริเยสุ คุตฺตทฺวาโร”ติ อตฺถิ อินฺทฺริเยสุ คุตฺตทฺวารตา, อตฺถิ อคุตฺตทฺวารตา.}

\csromanheader{2}{12. ฌานวิภงฺค}
\prose{\csromanfolio{257}ตตฺถ กตมา อินฺทฺริเยสุ อคุตฺตทฺวารตา, อิเธกจฺโจ จกฺขุนา รูปํ ทิสฺวา นิมิตฺตคฺคาหี โหติ อนุพฺยญฺชนคฺคาหี, ยตฺวาธิกรณเมนํ จกฺขุนฺทฺริยํ อสํวุตํ วิหรนฺตํ อภิชฺฌาโทมนสฺสา ปาปกา อกุสลา ธมฺมา อนฺวาสฺสเวยฺยุํ, ตสฺส สํวราย น ปฏิปชฺชติ, น รกฺขติ จกฺขุนฺทฺริยํ, จกฺขุนฺทฺริเย น สํวรํ อาปชฺชติ.\csromanspacer{} โสเตน สทฺทํ สุตฺวา ฯเปฯ.\csromanspacer{} ฆาเนน คนฺธํ ฆายิตฺวา ฯเปฯ.\csromanspacer{} ชิวฺหาย รสํ สายิตฺวา ฯเปฯ.\csromanspacer{} กาเยน โผฏฺฐพฺพํ ผุสิตฺวา ฯเปฯ.\csromanspacer{} มนสา ธมฺมํ วิญฺญาย นิมิตฺตคฺคาหี โหติ อนุพฺยญฺชนคฺคาหี, ยตฺวาธิกรณเมนํ มนินฺทฺริยํ อสํวุตํ วิหรนฺตํ อภิชฺฌาโทมนสฺสา ปาปกา อกุสลา ธมฺมา อนฺวาสฺสเวยฺยุํ, ตสฺส สํวราย น ปฏิปชฺชติ, น รกฺขติ มนินฺทฺริยํ, มนินฺทฺริเย น สํวรํ อาปชฺชติ.\csromanspacer{} ยา อิเมสํ ฉนฺนํ อินฺทฺริยานํ อคุตฺติ อโคปนา อนารกฺโข อสํวโร, อยํ วุจฺจติ อินฺทฺริเยสุ อคุตฺตทฺวารตา.}

\prose{ตตฺถ กตมา อินฺทฺริเยสุ คุตฺตทฺวารตา, อิเธกจฺโจ จกฺขุนา รูปํ ทิสฺวา น นิมิตฺตคฺคาหี โหติ นานุพฺยญฺชนคฺคาหี, ยตฺวาธิกรณเมนํ จกฺขุนฺทฺริยํ อสํวุตํ วิหรนฺตํ อภิชฺฌาโทมนสฺสา ปาปกา อกุสลา ธมฺมา อนฺวาสฺสเวยฺยุํ, ตสฺส สํวราย ปฏิปชฺชติ, รกฺขติ จกฺขุนฺทฺริยํ, จกฺขุนฺทฺริเย สํวรํ อาปชฺชติ.\csromanspacer{} โสเตน สทฺทํ สุตฺวา ฯเปฯ.\csromanspacer{} ฆาเนน คนฺธํ ฆายิตฺวา ฯเปฯ.\csromanspacer{} ชิวฺหาย รสํ สายิตฺวา ฯเปฯ.\csromanspacer{} กาเยน โผฏฺฐพฺพํ ผุสิตฺวา ฯเปฯ.\csromanspacer{} มนสา ธมฺมํ วิญฺญาย น นิมิตฺตคฺคาหี โหติ นานุพฺยญฺชนคฺคาหี, ยตฺวาธิกรณเมนํ มนินฺทฺริยํ อสํวุตํ วิหรนฺตํ อภิชฺฌาโทมนสฺสา ปาปกา อกุสลา ธมฺมา อนฺวาสฺสเวยฺยุํ, ตสฺส สํวราย ปฏิปชฺชติ, รกฺขติ มนินฺทฺริยํ, มนินฺทฺริเย สํวรํ อาปชฺชติ.\csromanspacer{} ยา อิเมสํ ฉนฺนํ อินฺทฺริยานํ คุตฺติ โคปนา อารกฺโข สํวโร, อยํ วุจฺจติ อินฺทฺริเยสุ คุตฺตทฺวารตา.\csromanspacer{} อิมาย อินฺทฺริเยสุ คุตฺตทฺวารตาย อุเปโต โหติ สมุเปโต ฯเปฯ สมนฺนาคโต, เตน วุจฺจติ “อินฺทฺริเยสุ คุตฺตทฺวาโร”ติ.}

\proseitem{518}{“โภชเน มตฺตญฺญู”ติ อตฺถิ โภชเน มตฺตญฺญุตา, อตฺถิ โภชเน อมตฺตญฺญุตา.}

\prose{ตตฺถ กตมา โภชเน อมตฺตญฺญุตา, อิเธกจฺโจ อปฺปฏิสงฺขา อโยนิโส อาหารํ อาหาเรติ ทวาย มทาย มณฺฑนาย วิภูสนาย.\csromanspacer{} ยา ตตฺถ อสนฺตุฏฺฐิตา อมตฺตญฺญุตา อปฺปฏิสงฺขา โภชเน, อยํ วุจฺจติ “โภชเน อมตฺตญฺญุตา”ติ.}

\prose{\csromanfolio{258}ตตฺถ กตมา โภชเน มตฺตญฺญุตา, อิเธกจฺโจ ปฏิสงฺขา โยนิโส อาหารํ อาหาเรติ เนว ทวาย น มทาย น มณฺฑนาย น วิภูสนาย ยาวเทว อิมสฺส กายสฺส ฐิติยา ยาปนาย วิหิํสูปรติยา พฺรหฺมจริยานุคฺคหาย, อิติ ปุราณญฺจ เวทนํ ปฏิหงฺขามิ, นวญฺจ เวทนํ น อุปฺปาเทสฺสามิ, ยาตฺรา จ เม ภวิสฺสติ อนวชฺชตา จ ผาสุ วิหาโร จาติ.\csromanspacer{} ยา ตตฺถ สนฺตุฏฺฐิตา มตฺตญฺญุตา ปฏิสงฺขา โภชเน, อยํ วุจฺจติ โภชเน มตฺตญฺญุตา.\csromanspacer{} อิมาย โภชเน มตฺตญฺญุตาย อุเปโต โหติ ฯเปฯ สมนฺนาคโต, เตน วุจฺจติ “โภชเน มตฺตญฺญู”ติ.}

\proseitem{519}{กถญฺจ ภิกฺขุ ปุพฺพรตฺตาปรรตฺตํ ชาคริยานุโยคมนุยุตฺโต โหติ, อิธ ภิกฺขุ ทิวสํ จงฺกเมน นิสชฺชาย อาวรณีเยหิ ธมฺเมหิ จิตฺตํ ปริโสเธติ, รตฺติยา ปฐมยามํ จงฺกเมน นิสชฺชาย อาวรณีเยหิ ธมฺเมหิ จิตฺตํ ปริโสเธติ, รตฺติยา มชฺฌิมยามํ ทกฺขิเณน ปสฺเสน สีหเสยฺยํ กปฺเปติ ปาเท ปาทํ\footnote{ปาเทน ปาทํ (สฺยา)} อจฺจาธาย สโต สมฺปชาโน อุฏฺฐานสญฺญํ มนสิกริตฺวา, รตฺติยา ปจฺฉิมยามํ ปจฺจุฏฺฐาย จงฺกเมน นิสชฺชาย อาวรณีเยหิ ธมฺเมหิ จิตฺตํ ปริโสเธติ, เอวํ ภิกฺขุ ปุพฺพรตฺตาปรรตฺตํ ชาคริยานุโยคมนุยุตฺโต.}

\proseitem{520}{“สาตจฺจนฺ”ติ โย เจตสิโก วีริยารมฺโภ ฯเปฯ สมฺมาวายาโม.}

\proseitem{521}{“เนปกฺกนฺ”ติ ยา ปญฺญา ปชานนา ฯเปฯ อโมโห ธมฺมวิจโย สมฺมาทิฏฺฐิ.}

\proseitem{522}{“โพธิปกฺขิกานํ ธมฺมานํ ภาวนานุโยคมนุยุตฺโต”ติ ตตฺถ กตเม โพธิปกฺขิกา ธมฺมา, สตฺต โพชฺฌงฺคา สติสมฺโพชฺฌงฺโค ธมฺมวิจยสมฺโพชฺฌงฺโค วีริยสมฺโพชฺฌงฺโค ปีติสมฺโพชฺฌงฺโค ปสฺสทฺธิสมฺโพชฺฌงฺโค สมาธิสมฺโพชฺฌงฺโค อุเปกฺขาสมฺโพชฺฌงฺโค, อิเม วุจฺจนฺติ โพธิปกฺขิกา ธมฺมา.\csromanspacer{} อิติ เต โพธิปกฺขิเก ธมฺเม อาเสวติ ภาเวติ พหุลีกโรติ, เตน วุจฺจติ “โพธิปกฺขิกานํ ธมฺมานํ ภาวนานุโยคมนุยุตฺโต”ติ.}

\csromanheader{2}{12. ฌานวิภงฺค}
\proseitem{523}{\csromanfolio{259}กถญฺจ ภิกฺขุ อภิกฺกนฺเต ปฏิกฺกนฺเต สมฺปชานการี โหติ, อาโลกิเต วิโลกิเต สมฺปชานการี โหติ, สมิญฺชิเต ปสาริเต สมฺปชานการี โหติ, สงฺฆาฏิปตฺตจีวรธารเณ สมฺปชานการี โหติ, อสิเต ปีเต ขายิเต สายิเต สมฺปชานการี โหติ, อุจฺจารปสฺสาวกมฺเม สมฺปชานการี โหติ, คเต ฐิเต นิสินฺเน สุตฺเต ชาคริเต ภาสิเต ตุณฺหีภาเว สมฺปชานการี โหติ, อิธ ภิกฺขุ สโต สมฺปชาโน อภิกฺกมติ, สโต สมฺปชาโน ปฏิกฺกมติ, สโต สมฺปชาโน อาโลเกติ, สโต สมฺปชาโน วิโลเกติ, สโต สมฺปชาโน สมิญฺเชติ, สโต สมฺปชาโน ปสาเรติ, สโต สมฺปชานการี โหติ, สงฺฆาฏิปตฺตจีวรธารเณ สโต สมฺปชานการี โหติ, อสิเต ปีเต ขายิเต สายิเต สโต สมฺปชานการี โหติ, อุจฺจารปสฺสาวกมฺเม สโต สมฺปชานการี โหติ, คเต ฐิเต นิสินฺเน สุตฺเต ชาคริเต ภาสิเต ตุณฺหีภาเว สโต สมฺปชานการี โหตีติ.}

\proseitem{524}{ตตฺถ กตโม สติ, ยา สติ อนุสฺสติ ปฏิสฺสติ สติ สรณตา ธารณตา อปิลาปนตา อสมฺมุสนตา สติ สตินฺทฺริยํ สติพลํ สมฺมาสติ, อยํ วุจฺจติ สติ.}

\proseitem{525}{“สมฺปชาโน”ติ ตตฺถ กตมํ สมฺปชญฺญํ, ยา ปญฺญา ปชานนา วิจโย ปวิจโย ธมฺมวิจโย สลฺลกฺขณา อุปลกฺขณา ปจฺจุปลกฺขณา ปณฺฑิจฺจํ โกสลฺลํ เนปุญฺญํ เวภพฺยา จินฺตา-อุปปริกฺขา ภูรี เมธา ปริณายิกา วิปสฺสนา สมฺปชญฺญํ ปโตโท ปญฺญา ปญฺญินฺทฺริยํ ปญฺญาพลํ ปญฺญาสตฺถํ ปญฺญาปาสาโท ปญฺญา-อาโลโก ปญฺญา-โอภาโส ปญฺญาปชฺโชโต ปญฺญารตนํ อโมโห ธมฺมวิจโย สมฺมาทิฏฺฐิ, อิทํ วุจฺจติ สมฺปชญฺญํ.\csromanspacer{} อิติ อิมาย จ สติยา อิมินา จ สมฺปชญฺเญน อุเปโต โหติ ฯเปฯ สมนฺนาคโต.\csromanspacer{} เอวํ ภิกฺขุ สโต สมฺปชาโน อภิกฺกมติ, สโต สมฺปชาโน ปฏิกฺกมติ, สโต สมฺปชาโน อาโลเกติ, สโต สมฺปชาโน วิโลเกติ, สโต สมฺปชาโน สมิญฺเชติ, สโต สมฺปชาโน ปสาเรติ, สโต สมฺปชานการี โหติ, สงฺฆาฏิปตฺตจีวรธารเณ สโต สมฺปชานการี โหติ, อสิเต ปีเต ขายิเต สายิเต สโต สมฺปชานการี โหติ, อุจฺจารปสฺสาวกมฺเม สโต \csromanfolio{260}สมฺปชานการี โหติ, คเต ฐิเต นิสินฺเน สุตฺเต ชาคริเต ภาสิเต ตุณฺหีภาเว สมฺปชานการี โหติ.}

\proseitem{526}{“วิวิตฺตนฺ”ติ สนฺติเก เจปิ เสนาสนํ โหติ, ตญฺจ อนากิณฺณํ คหฏฺเฐหิ ปพฺพชิเตหิ, เตน ตํ วิวิตฺตํ.\csromanspacer{} ทูเร เจปิ เสนาสนํ โหติ, ตญฺจ อนากิณฺณํ คหฏฺเฐหิ ปพฺพชิเตหิ, เตน ตํ วิวิตฺตํ.}

\proseitem{527}{“เสนาสนนฺ”ติ มญฺโจปิ เสนาสนํ, ปีฐมฺปิ เสนาสนํ, ภิสิปิ เสนาสนํ, พิพฺโพหนมฺปิ\footnote{พิมฺโพหนมฺปิ (สี, สฺยา)} เสนาสนํ, วิหาโรปิ เสนาสนํ, อฑฺฒโยโคปิ เสนาสนํ, ปาสาโทปิ เสนาสนํ, อฏฺโฏปิ เสนาสนํ, มาโฬปิ เสนาสนํ, เลณมฺปิ เสนาสนํ, คุหาปิ เสนาสนํ, รุกฺขมูลมฺปิ เสนาสนํ, เวฬุคุมฺโพปิ เสนาสนํ, ยตฺถ วา ปน ภิกฺขู ปฏิกฺกมนฺติ, สพฺพเมตํ เสนาสนํ.}

\proseitem{528}{“วิวิตฺตํ เสนาสนํ ภชตี”ติ อิมํ วิวิตฺตํ เสนาสนํ ภชติ สมฺภชติ เสวติ นิเสวติ สํเสวติ, เตน วุจฺจติ “วิวิตฺตํ เสนาสนํ ภชตี”ติ.}

\proseitem{529}{“อรญฺญนฺ”ติ นิกฺขมิตฺวา พหิ อินฺทขีลา, สพฺพเมตํ อรญฺญํ.}

\proseitem{530}{“รุกฺขมูลนฺ”ติ รุกฺขมูลํเยว รุกฺขมูลํ, ปพฺพโตเยว ปพฺพโต, กนฺทราเยว กนฺทรา, คิริคุหาเยว คิริคุหา, สุสานํเยว สุสานํ, อพฺโภกาโสเยว อพฺโภกาโส, ปลาลปุญฺโชเยว ปลาลปุญฺโช.}

\proseitem{531}{“วนปตฺถนฺ”ติ ทูรานเมตํ เสนาสนานํ อธิวจนํ, “วนปตฺถนฺ”ติ วนสณฺฑานเมตํ เสนาสนานํ อธิวจนํ, “วนปตฺถนฺ”ติ ภีสนกานเมตํ\footnote{ภิํสนกานเมตํ (สี, สฺยา)} เสนาสนานํ อธิวจนํ, “วนปตฺถนฺ”ติ สโลมหํสานเมตํ เสนาสนานํ อธิวจนํ, “วนปตฺถนฺ”ติ ปริยนฺตานเมตํ เสนาสนานํ อธิวจนํ, “วนปตฺถนฺ”ติ น มนุสฺสูปจารานเมตํ เสนาสนานํ อธิวจนํ, “วนปตฺถนฺ”ติ ทุรภิสมฺภวานเมตํ เสนาสนานํ อธิวจนํ.}

\csromanheader{2}{12. ฌานวิภงฺค}
\proseitem{532}{\csromanfolio{261}“อปฺปสทฺทนฺ”ติ สนฺติเก เจปิ เสนาสนํ โหติ, ตญฺจ อนากิณฺณํ คหฏฺเฐหิ ปพฺพชิเตหิ, เตน ตํ อปฺปสทฺทํ.\csromanspacer{} ทูเร เจปิ เสนาสนํ โหติ, ตญฺจ อนากิณฺณํ คหฏฺเฐหิ ปพฺพชิเตหิ, เตน ตํ อปฺปสทฺทํ.}

\proseitem{533}{“อปฺปนิคฺโฆสนฺ”ติ ยเทว ตํ อปฺปสทฺทํ, ตเทว ตํ อปฺปนิคฺโฆสํ.\csromanspacer{} ยเทว ตํ อปฺปนิคฺโฆสํ, ตเทว ตํ วิชนวาตํ.\csromanspacer{} ยเทว ตํ วิชนวาตํ, ตเทว ตํ มนุสฺสราหสฺเสยฺยกํ.\csromanspacer{} ยเทว ตํ มนุสฺสราหสฺเสยฺยกํ, ตเทว ตํ ปฏิสลฺลานสารุปฺปํ.}

\proseitem{534}{“อรญฺญคโต วา รุกฺขมูลคโต วา สุญฺญาคารคโต วา”ติ อรญฺญคโต วา โหติ รุกฺขมูลคโต วา สุญฺญาคารคโต วา.}

\proseitem{535}{“นิสีทติ ปลฺลงฺกํ อาภุชิตฺวา”ติ นิสินฺโน โหติ ปลฺลงฺกํ อาภุชิตฺวา.}

\proseitem{536}{“อุชุํ กายํ ปณิธายา”ติ อุชุโก โหติ กาโย ฐิโต ปณิหิโต.}

\proseitem{537}{“ปริมุขํ สติํ อุปฏฺฐเปตฺวา”ติ ตตฺถ กตมา สติ, ยา สติ อนุสฺสติ ปฏิสฺสติ ฯเปฯ สมฺมาสติ, อยํ วุจฺจติ สติ.\csromanspacer{} อยํ สติ อุปฏฺฐิตา โหติ สุปฏฺฐิตา นาสิกคฺเค วา มุขนิมิตฺเต วา, เตน วุจฺจติ “ปริมุขํ สติํ อุปฏฺฐเปตฺวา”ติ.}

\proseitem{538}{“อภิชฺฌํ โลเก ปหายา”ติ ตตฺถ กตมา อภิชฺฌา, โย ราโค สาราโค ฯเปฯ จิตฺตสฺส สาราโค, อยํ วุจฺจติ อภิชฺฌา.}

\prose{ตตฺถ กตโม โลโก, ปญฺจุปาทานกฺขนฺธา โลโก, อยํ วุจฺจติ โลโก.\csromanspacer{} อยํ อภิชฺฌา อิมมฺหิ โลเก สนฺตา โหติ สมิตา วูปสนฺตา อตฺถงฺคตา อพฺภตฺถงฺคตา อปฺปิตา พฺยปฺปิตา โสสิตา วิโสสิตา พฺยนฺตีกตา, เตน วุจฺจติ “อภิชฺฌํ โลเก ปหายา”ติ.}

\proseitem{539}{“วิคตาภิชฺเฌน เจตสา”ติ ตตฺถ กตมํ จิตฺตํ, ยํ จิตฺตํ มโน มานสํ ฯเปฯ ตชฺชามโนวิญฺญาณธาตุ, อิทํ วุจฺจติ จิตฺตํ.\csromanspacer{} อิทํ จิตฺตํ วิคตาภิชฺฌํ โหติ, เตน วุจฺจติ “วิคตาภิชฺเฌน เจตสา”ติ.}

\proseitem{540}{\csromanfolio{262}“วิหรตี”ติ อิริยติ วตฺตติ ปาเลติ ยเปติ ยาเปติ จรติ วิหรติ, เตน วุจฺจติ “วิหรตี”ติ.}

\proseitem{541}{“อภิชฺฌาย จิตฺตํ ปริโสเธตี”ติ ตตฺถ กตมา อภิชฺฌา, โย ราโค สาราโค ฯเปฯ จิตฺตสฺส สาราโค, อยํ วุจฺจติ อภิชฺฌา.}

\prose{ตตฺถ กตมํ จิตฺตํ, ยํ จิตฺตํ มโน มานสํ ฯเปฯ ตชฺชามโนวิญฺญาณธาตุ, อิทํ วุจฺจติ จิตฺตํ.\csromanspacer{} อิทํ จิตฺตํ อิมาย อภิชฺฌาย โสเธติ วิโสเธติ ปริโสเธติ โมเจติ วิโมเจติ ปริโมเจติ, เตน วุจฺจติ “อภิชฺฌาย จิตฺตํ ปริโสเธตี”ติ.}

\proseitem{542}{“พฺยาปาทปโทสํ ปหายา”ติ อตฺถิ พฺยาปาโท, อตฺถิ ปโทโส.}

\prose{ตตฺถ กตโม พฺยาปาโท, โย จิตฺตสฺส อาฆาโต ปฏิฆาโต ปฏิฆํ ปฏิวิโรโธ โกปิ ปโกโป สมฺปโกโป โทโส ปโทโส สมฺปโทโส จิตฺตสฺส พฺยาปตฺติ มโนปโทโส โกโธ กุชฺฌนา กุชฺฌิตตฺตํ โทโส ทุสฺสนา ทุสฺสิตตฺตํ\footnote{ทูสนา ทูสิตตฺตํ (สฺยา)} พฺยาปตฺติ พฺยาปชฺชนา พฺยาปชฺชิตตฺตํ วิโรโธ ปฏิวิโรโธ จณฺฑิกฺกํ อสุโรโป อนตฺตมนตา จิตฺตสฺส, อยํ วุจฺจติ พฺยาปาโท.}

\prose{ตตฺถ กตโม ปโทโส, โย พฺยาปาโท โส ปโทโส, โย ปโทโส โส พฺยาปาโท, อิติ อยญฺจ พฺยาปาโท อยญฺจ ปโทโส สนฺตา โหนฺติ สมิตา วูปสนฺตา อตฺถงฺคตา อพฺภตฺถงฺคตา อปฺปิตา พฺยปฺปิตา โสสิตา วิโสสิตา พฺยนฺตีกตา, เตน วุจฺจติ “พฺยาปาทปโทสํ ปหายา”ติ.}

\proseitem{543}{“อพฺยาปนฺนจิตฺโต”ติ ตตฺถ กตมํ จิตฺตํ, ยํ จิตฺตํ มโน มานสํ ฯเปฯ ตชฺชามโนวิญฺญาณธาตุ, อิทํ วุจฺจติ จิตฺตํ.\csromanspacer{} อิทํ จิตฺตํ อพฺยาปนฺนํ โหติ, เตน วุจฺจติ “อพฺยาปนฺนจิตฺโต”ติ.}

\proseitem{544}{“วิหรตี”ติ ฯเปฯ เตน วุจฺจติ “วิหรตี”ติ.}

\csromanheader{2}{12. ฌานวิภงฺค}
\proseitem{545}{\csromanfolio{263}“พฺยาปาทปโทสา จิตฺตํ ปริโสเธตี”ติ อตฺถิ พฺยาปาโท, อตฺถิ ปโทโส.}

\prose{ตตฺถ กตโม พฺยาปาโท, โย จิตฺตสฺส อาฆาโต ฯเปฯ จณฺฑิกฺกํ อสุโรโป อนตฺตมนตา จิตฺตสฺส, อยํ วุจฺจติ พฺยาปาโท.}

\prose{ตตฺถ กตโม ปโทโส, โย พฺยาปาโท โส ปโทโส, โย ปโทโส โส พฺยาปาโท.}

\prose{ตตฺถ กตมํ จิตฺตํ, ยํ จิตฺตํ มโน มานสํ ฯเปฯ ตชฺชามโนวิญฺญาณธาตุ, อิทํ วุจฺจติ จิตฺตํ.\csromanspacer{} อิทํ จิตฺตํ อิมมฺหา พฺยาปาทปโทสา โสเธติ วิโสเธติ ปริโสเธติ โมเจติ วิโมเจติ ปริโมเจติ, เตน วุจฺจติ “พฺยาปาทปโทสา จิตฺตํ ปริโสเธตี”ติ.}

\proseitem{546}{“ถินมิทฺธํ ปหายา”ติ อตฺถิ ถินํ\footnote{ถีนํ (สี, สฺยา)}, อตฺถิ มิทฺธํ.}

\prose{ตตฺถ กตมํ ถินํ, ยา จิตฺตสฺส อกลฺยตา อกมฺมญฺญตา โอลียนา สลฺลียนา ลีนํ ลียนา ลียิตตฺตํ ถินํ ถิยนา ถิยิตตฺตํ\footnote{ถียนา ถียิตตฺตํ (สี, สฺยา)} จิตฺตสฺส, อิทํ วุจฺจติ ถินํ.}

\prose{ตตฺถ กตมํ มิทฺธํ, ยา กายสฺส อกลฺยตา อกมฺมญฺญตา โอนาโห ปริโยนาโห อนฺโตสโมโรโธ มิทฺธํ สุปฺปํ ปจลายิกํ สุปฺปํ สุปฺปนา สุปฺปิตตฺตํ, อิทํ วุจฺจติ มิทฺธํ.\csromanspacer{} อิติ อิทญฺจ ถินํ อิทญฺจ มิทฺธํ สนฺตา โหนฺติ สมิตา วูปสนฺตา อตฺถงฺคตา อพฺภตฺถงฺคตา อปฺปิตา พฺยปฺปิตา โสสิตา วิโสสิตา พฺยนฺตีกตา, เตน วุจฺจติ “ถินมิทฺธํ ปหายา”ติ.}

\proseitem{547}{“วิคตถินมิทฺโธ”ติ ตสฺส ถินมิทฺธสฺส จตฺตตฺตา วนฺตตฺตา มุตฺตตฺตา ปหีนตฺตา ปฏินิสฺสฏฺฐตฺตา ปหีนปฏินิสฺสฏฺฐตฺตา, เตน วุจฺจติ “วิคตถินมิทฺโธ”ติ.}

\proseitem{548}{“วิหรตี”ติ ฯเปฯ เตน วุจฺจติ “วิหรตี”ติ.}

\proseitem{549}{“อาโลกสญฺญี”ติ ตตฺถ กตมา สญฺญา, ยา สญฺญา สญฺชานนา สญฺชานิตตฺตํ, อยํ วุจฺจติ สญฺญา.\csromanspacer{} อยํ สญฺญา อาโลกา \csromanfolio{264}โหติ วิวฏา ปริสุทฺธา ปริโยทาตา, เตน วุจฺจติ “อาโลกสญฺญี”ติ.}

\proseitem{550}{“สโต สมฺปชาโน”ติ ตตฺถ กตมา สติ, ยา สติ อนุสฺสติ ฯเปฯ สมฺมาสติ, อยํ วุจฺจติ สติ.}

\prose{ตตฺถ กตมํ สมฺปชญฺญํ, ยา ปญฺญา ปชานนา ฯเปฯ อโมโห ธมฺมวิจโย สมฺมาทิฏฺฐิ, อิทํ วุจฺจติ สมฺปชญฺญํ.\csromanspacer{} อิติ อิมาย จ สติยา อิมินา จ สมฺปชญฺเญน อุเปโต โหติ ฯเปฯ สมนฺนาคโต, เตน วุจฺจติ “สโต สมฺปชาโน”ติ.}

\proseitem{551}{“ถินมิทฺธา จิตฺตํ ปริโสเธตี”ติ อตฺถิ ถินํ, อตฺถิ มิทฺธํ.}

\prose{ตตฺถ กตมํ ถินํ ฯเปฯ อิทํ วุจฺจติ ถินํ.}

\prose{ตตฺถ กตมํ มิทฺธํ ฯเปฯ อิทํ วุจฺจติ มิทฺธํ.}

\prose{ตตฺถ กตมํ จิตฺตํ ฯเปฯ อิทํ วุจฺจติ จิตฺตํ.\csromanspacer{} อิทํ จิตฺตํ อิมมฺหา ถินมิทฺธา โสเธติ วิโสเธติ ปริโสเธติ โมเจติ วิโมเจติ ปริโมเจติ, เตน วุจฺจติ “ถินมิทฺธา จิตฺตํ ปริโสเธตี”ติ.}

\proseitem{552}{“อุทฺธจฺจกุกฺกุจฺจํ ปหายา”ติ อตฺถิ อุทฺธจฺจํ, อตฺถิ กุกฺกุจฺจํ.}

\prose{ตตฺถ กตมํ อุทฺธจฺจํ, ยํ จิตฺตสฺส อุทฺธจฺจํ อวูปสโม เจตโส วิกฺเขโป ภนฺตตฺตํ จิตฺตสฺส, อิทํ วุจฺจติ อุทฺธจฺจํ.}

\prose{ตตฺถ กตมํ กุกฺกุจฺจํ, อกปฺปิเย กปฺปิยสญฺญิตา กปฺปิเย อกปฺปิยสญฺญิตา อวชฺเช วชฺชสญฺญิตา วชฺเช อวชฺชสญฺญิตา, ยํ เอวรูปํ กุกฺกุจฺจํ กุกฺกุจฺจายนา กุกฺกุจฺจายิตตฺตํ เจตโส วิปฺปฏิสาโร มโนวิเลขา, อิทํ วุจฺจติ กุกฺกุจฺจํ.\csromanspacer{} อิติ อิทญฺจ อุทฺธจฺจํ อิทญฺจ กุกฺกุจฺจํ สนฺตา โหนฺติ สมิตา วูปสนฺตา อตฺถงฺคตา อพฺภตฺถงฺคตา อปฺปิตา พฺยปฺปิตา โสสิตา วิโสสิตา พฺยนฺตีกตา, เตน วุจฺจติ “อุทฺธจฺจกุกฺกุจฺจํ ปหายา”ติ.}

\proseitem{553}{“อนุทฺธโต”ติ ตสฺส อุทฺธจฺจกุกฺกุจฺจสฺส จตฺตตฺตา วนฺตตฺตา มุตฺตตฺตา ปหีนตฺตา ปฏินิสฺสฏฺฐตฺตา ปหีนปฏินิสฺสฏฺฐตฺตา, เตน วุจฺจติ “อนุทฺธโต”ติ.}

\proseitem{554}{“วิหรตี”ติ ฯเปฯ เตน วุจฺจติ “วิหรตี”ติ.}

\csromanheader{2}{12. ฌานวิภงฺค}
\proseitem{555}{\csromanfolio{265}“วูปสนฺตจิตฺโต”ติ ตตฺถ กตมํ จิตฺตํ, ยํ จิตฺตํ มโน มานสํ ฯเปฯ ตชฺชามโนวิญฺญาณธาตุ, อิทํ วุจฺจติ จิตฺตํ.\csromanspacer{} อิทํ จิตฺตํ อชฺฌตฺตํ สนฺตํ โหติ สมิตํ วูปสนฺตํ, เตน วุจฺจติ “อชฺฌตฺตํ วูปสนฺตจิตฺโต”ติ.}

\proseitem{556}{“อุทฺธจฺจกุกฺกุจฺจา จิตฺตํ ปริโสเธตี”ติ อตฺถิ อุทฺธจฺจํ, อตฺถิ กุกฺกุจฺจํ.}

\prose{ตตฺถ กตมํ อุทฺธจฺจํ, ยํ จิตฺตสฺส อุทฺธจฺจํ อวูปสโม เจตโส วิกฺเขโป ภนฺตตฺตํ จิตฺตสฺส, อิทํ วุจฺจติ อุทฺธจฺจํ.}

\prose{ตตฺถ กตมํ กุกฺกุจฺจํ ฯเปฯ อิทํ วุจฺจติ กุกฺกุจฺจํ.}

\prose{ตตฺถ กตมํ จิตฺตํ ฯเปฯ อิทํ วุจฺจติ จิตฺตํ.\csromanspacer{} อิทํ จิตฺตํ อิมมฺหา อุทฺธจฺจกุกฺกุจฺจา โสเธติ วิโสเธติ ปริโสเธติ โมเจติ วิโมเจติ ปริโมเจติ, เตน วุจฺจติ “อุทฺธจฺจกุกฺกุจฺจา จิตฺตํ ปริโสเธตี”ติ.}

\proseitem{557}{“วิจิกิจฺฉํ ปหายา”ติ ตตฺถ กตมา วิจิกิจฺฉา, ยา กงฺขา กงฺขายนา กงฺขายิตตฺตํ วิมติ วิจิกิจฺฉา เทฺวฬฺหกํ ทฺวิธาปโถ สํสโย อเนกํสคฺคาโห อาสปฺปนา ปริสปฺปนา อปริโยคาหณา ฉมฺภิตตฺตํ จิตฺตสฺส มโนวิเลโข, อยํ วุจฺจติ วิจิกิจฺฉา.\csromanspacer{} อยํ วิจิกิจฺฉา สนฺตา โหติ สมิตา วูปสนฺตา อตฺถงฺคตา อพฺภตฺถงฺคตา อปฺปิตา พฺยปฺปิตา โสสิตา วิโสสิตา พฺยนฺตีกตา, เตน วุจฺจติ “วิจิกิจฺฉํ ปหายา”ติ.}

\proseitem{558}{“ติณฺณวิจิกิจฺโฉ”ติ อิมํ วิจิกิจฺฉํ ติณฺโณ โหติ อุตฺติณฺโณ นิตฺติณฺโณ ปารงฺคโต ปารมนุปฺปตฺโต, เตน วุจฺจติ “ติณฺณวิจิกิจฺโฉ”ติ.}

\proseitem{559}{“อกถํกถี กุสเลสุ ธมฺเมสู”ติ อิมาย วิจิกิจฺฉาย กุสเลสุ ธมฺเมสุ น กงฺขติ น วิจิกิจฺฉติ อกถํกถี โหติ นิกฺกถํกถี วิกถํกโถ, เตน วุจฺจติ “อกถํกถี กุสเลสุ ธมฺเมสู”ติ.}

\proseitem{560}{“วิจิกิจฺฉาย จิตฺตํ ปริโสเธตี”ติ ตตฺถ กตมา วิจิกิจฺฉา, ยา กงฺขา กงฺขายนา กงฺขายิตตฺตํ ฉมฺภิตตฺตํ จิตฺตสฺส มโนวิเลโข, อยํ วุจฺจติ วิจิกิจฺฉา.}

\prose{\csromanfolio{266}ตตฺถ กตมํ จิตฺตํ, ยํ จิตฺตํ มโน มานสํ ฯเปฯ ตชฺชามโนวิญฺญาณธาตุ, อิทํ วุจฺจติ จิตฺตํ.\csromanspacer{} อิทํ จิตฺตํ อิมาย วิจิกิจฺฉาย โสเธติ วิโสเธติ ปริโสเธติ โมเจติ วิโมเจติ ปริโมเจติ, เตน วุจฺจติ “วิจิกิจฺฉาย จิตฺตํ ปริโสเธตี”ติ.}

\proseitem{561}{“อิเม ปญฺจ นีวรเณ ปหายา”ติ อิเม ปญฺจ นีวรณา สนฺตา โหนฺติ สมิตา วูปสนฺตา อตฺถงฺคตา อพฺภตฺถงฺคตา อปฺปิตา พฺยปฺปิตา โสสิตา วิโสสิตา พฺยนฺตีกตา, เตน วุจฺจติ “อิเม ปญฺจ นีวรเณ ปหายา”ติ.}

\proseitem{562}{“เจตโส อุปกฺกิเลเส”ติ อิเม ปญฺจ นีวรณา จิตฺตสฺส อุปกฺกิเลสา.}

\proseitem{563}{“ปญฺญาย ทุพฺพลีกรเณ”ติ อิเมหิ ปญฺจหิ นีวรเณหิ อนุปฺปนฺนา เจว ปญฺญา น อุปฺปชฺชติ, อุปฺปนฺนา จ ปญฺญา นิรุชฺฌติ, เตน วุจฺจติ “ปญฺญาย ทุพฺพลีกรเณ”ติ.}

\proseitem{564}{“วิวิจฺเจว กาเมหิ วิวิจฺจ อกุสเลหิ ธมฺเมหี”ติ ตตฺถ กตเม กามา, ฉนฺโท กาโม ราโค กาโม ฉนฺทราโค กาโม สงฺกปฺโป กาโม ราโค กาโม สงฺกปฺปราโค กาโม, อิเม วุจฺจนฺติ กามา.}

\prose{ตตฺถ กตเม อกุสลา ธมฺมา, กามจฺฉนฺโท พฺยาปาโท ถินํ มิทฺธํ อุทฺธจฺจํ กุกฺกุจฺจํ วิจิกิจฺฉา, อิเม วุจฺจนฺติ อกุสลา ธมฺมา.\csromanspacer{} อิติ อิเมหิ จ กาเมหิ อิเมหิ จ อกุสเลหิ ธมฺเมหิ วิวิตฺโต โหติ, เตน วุจฺจติ “วิวิจฺเจว กาเมหิ วิวิจฺจ อกุสเลหิ ธมฺเมหี”ติ.}

\proseitem{565}{“สวิตกฺกํ สวิจารนฺ”ติ อตฺถิ วิตกฺโก, อตฺถิ วิจาโร.}

\prose{ตตฺถ กตโม วิตกฺโก, โย ตกฺโก วิตกฺโก สงฺกปฺโป อปฺปนา พฺยปฺปนา เจตโส อภินิโรปนา สมฺมาสงฺกปฺโป, อยํ วุจฺจติ วิตกฺโก.}

\prose{ตตฺถ กตโม วิจาโร, โย จาโร วิจาโร อนุวิจาโร อุปวิจาโร จิตฺตสฺส อนุสนฺธนตา อนุเปกฺขนตา, อยํ วุจฺจติ วิจาโร.}

\vspace{\tipitakaparskip}
\csromancenter{ฌานวิภงฺค อิติ อิมินา จ วิตกฺเกน อิมินา จ วิจาเรน อุเปโต โหติ ฯเปฯ สมนฺนาคโต, เตน วุจฺจติ “สวิตกฺกํ สวิจารนฺ”ติ.}
\proseitem{566}{\csromanfolio{267}“วิเวกชนฺ”ติ วิตกฺโก วิจาโร ปีติ สุขํ จิตฺตสฺเสกคฺคตา, เต อิมสฺมิํ วิเวเก ชาตา โหนฺติ สญฺชาตา นิพฺพตฺตา อภินิพฺพตฺตา ปาตุภูตา, เตน วุจฺจติ “วิเวกชนฺ”ติ.}

\proseitem{567}{“ปีติสุขนฺ”ติ อตฺถิ ปีติ, อตฺถิ สุขํ.}

\prose{ตตฺถ กตมา ปีติ, ยา ปีติ ปาโมชฺชํ อาโมทนา ปโมทนา หาโส ปหาโส วิตฺติ โอทคฺยํ อตฺตมนตา จิตฺตสฺส, อยํ วุจฺจติ ปีติ.}

\prose{ตตฺถ กตมํ สุขํ, ยํ เจตสิกํ สาตํ เจตสิกํ สุขํ เจโตสมฺผสฺสชํ สาตํ สุขํ เวทยิตํ เจโตสมฺผสฺสชา สาตา สุขา เวทนา, อิทํ วุจฺจติ สุขํ.\csromanspacer{} อิทํ สุขํ อิมาย ปีติยา สหคตํ โหติ สหชาตํ สํสฏฺฐํ สมฺปยุตฺตํ, เตน วุจฺจติ “ปีติสุขนฺ”ติ.}

\proseitem{568}{“ปฐมนฺ”ติ คณนานุปุพฺพตา\footnote{คณนานุปุพฺพโต (สฺยา) เอวมุปริปิ.} ปฐมํ, อิทํ ปฐมํ สมาปชฺชตีติ ปฐมํ.}

\proseitem{569}{“ฌานนฺ”ติ วิตกฺโก วิจาโร ปีติ สุขํ จิตฺตสฺเสกคฺคตา.}

\proseitem{570}{“อุปสมฺปชฺชา”ติ โย ปฐมสฺส ฌานสฺส ลาโภ ปฏิลาโภ ปตฺติ สมฺปตฺติ ผุสนา สจฺฉิกิริยา อุปสมฺปทา.}

\proseitem{571}{“วิหรตี”ติ ฯเปฯ เตน วุจฺจติ “วิหรตี”ติ.}

\proseitem{572}{“วิตกฺกวิจารานํ วูปสมา”ติ อตฺถิ วิตกฺโก, อตฺถิ วิจาโร.}

\prose{ตตฺถ กตโม วิตกฺโก, โย ตกฺโก วิตกฺโก ฯเปฯ สมฺมาสงฺกปฺโป, อยํ วุจฺจติ วิตกฺโก.}

\prose{ตตฺถ กตโม วิจาโร, โย จาโร วิจาโร อนุวิจาโร อุปวิจาโร จิตฺตสฺส อนุสนฺธนตา อนุเปกฺขนตา, อยํ วุจฺจติ วิจาโร.\csromanspacer{} อิติ อยญฺจ วิตกฺโก อยญฺจ วิจาโร สนฺตา โหนฺติ สมิตา วูปสนฺตา \csromanfolio{268}อตฺถงฺคตา อพฺภตฺถงฺคตา อปฺปิตา พฺยปฺปิตา โสสิตา วิโสสิตา พฺยนฺตีกตา, เตน วุจฺจติ “วิตกฺกวิจารานํ วูปสมา”ติ.}

\proseitem{573}{“อชฺฌตฺตนฺ”ติ ยํ อชฺฌตฺตํ ปจฺจตฺตํ.}

\proseitem{574}{“สมฺปสาทนนฺ”ติ ยา สทฺธา สทฺทหนา โอกปฺปนา อภิปฺปสาโท.}

\proseitem{575}{“เจตโส เอโกทิภาวนฺ”ติ ยา จิตฺตสฺส ฐิติ ฯเปฯ สมฺมาสมาธิ.}

\proseitem{576}{“อวิตกฺกํ อวิจารนฺ”ติ อตฺถิ วิตกฺโก, อตฺถิ วิจาโร.}

\prose{ตตฺถ กตโม วิตกฺโก, โย ตกฺโก วิตกฺโก ฯเปฯ สมฺมาสงฺกปฺโป, อยํ วุจฺจติ วิตกฺโก.}

\prose{ตตฺถ กตโม วิจาโร, โย จาโร อนุจาโร วิจาโร อนุวิจาโร อุปวิจาโร จิตฺตสฺส อนุสนฺธนตา อนุเปกฺขนตา, อยํ วุจฺจติ วิจาโร.\csromanspacer{} อิติ อยญฺจ วิตกฺโก อยญฺจ วิจาโร สนฺตา โหนฺติ สมิตา วูปสนฺตา อตฺถงฺคตา อพฺภตฺถงฺคตา อปฺปิตา พฺยปฺปิตา โสสิตา วิโสสิตา พฺยนฺตีกตา, เตน วุจฺจติ “อวิตกฺกํ อวิจารนฺ”ติ.}

\proseitem{577}{“สมาธิชนฺ”ติ สมฺปสาโท ปีติสุขํ, เต อิมสฺมิํ สมาธิมฺหิ ชาตา โหนฺติ สญฺชาตา นิพฺพตฺตา อภินิพฺพตฺตา ปาตุภูตา, เตน วุจฺจติ “สมาธิชนฺ”ติ.}

\proseitem{578}{“ปีติสุขนฺ”ติ อตฺถิ ปีติ, อตฺถิ สุขํ.}

\prose{ตตฺถ กตมา ปีติ ฯเปฯ อยํ วุจฺจติ ปีติ.}

\prose{ตตฺถ กตมํ สุขํ ฯเปฯ อิทํ วุจฺจติ สุขํ.\csromanspacer{} อิทํ สุขํ อิมาย ปีติยา สหคตํ โหติ สหชาตํ สํสฏฺฐํ สมฺปยุตฺตํ, เตน วุจฺจติ “ปีติสุขนฺ”ติ.}

\proseitem{579}{“ทุติยนฺ”ติ คณนานุปุพฺพตา ทุติยํ, อิทํ ทุติยํ สมาปชฺชตีติ ทุติยํ.}

\proseitem{580}{“ฌานนฺ”ติ สมฺปสาโท ปีติสุขํ จิตฺตสฺเสกคฺคตา.}

\csromanheader{2}{12. ฌานวิภงฺค}
\proseitem{581}{\csromanfolio{269}“อุปสมฺปชฺชา”ติ โย ทุติยสฺส ฌานสฺส ลาโภ ปฏิลาโภ ปตฺติ สมฺปตฺติ ผุสนา สจฺฉิกิริยา อุปสมฺปทา.}

\proseitem{582}{“วิหรตี”ติ ฯเปฯ เตน วุจฺจติ “วิหรตี”ติ.}

\proseitem{583}{“ปีติยา จ วิราคา”ติ ตตฺถ กตมา ปีติ, ยา ปีติ ปาโมชฺชํ อาโมทนา ปโมทนา หาโส ปหาโส วิตฺติ โอทคฺยํ อตฺตมนตา จิตฺตสฺส, อยํ วุจฺจติ ปีติ.\csromanspacer{} อยํ ปีติ สนฺตา โหติ สมิตา วูปสนฺตา อตฺถงฺคตา อพฺภตฺถงฺคตา อปฺปิตา พฺยปฺปิตา โสสิตา วิโสสิตา พฺยนฺตีกตา, เตน วุจฺจติ “ปีติยา จ วิราคา”ติ.}

\proseitem{584}{“อุเปกฺขโก”ติ ตตฺถ กตมา อุเปกฺขา, ยา อุเปกฺขา อุเปกฺขนา อชฺฌุเปกฺขนา มชฺฌตฺตตา จิตฺตสฺส, อยํ วุจฺจติ อุเปกฺขา.\csromanspacer{} อิมาย อุเปกฺขาย อุเปโต โหติ ฯเปฯ สมนฺนาคโต, เตน วุจฺจติ “อุเปกฺขโก”ติ.}

\proseitem{585}{“วิหรตี”ติ ฯเปฯ เตน วุจฺจติ “วิหรตี”ติ.}

\proseitem{586}{“สโต จ สมฺปชาโน”ติ ตตฺถ กตมา สติ, ยา สติ อนุสฺสติ ฯเปฯ สมฺมาสติ, อยํ วุจฺจติ สติ.}

\prose{ตตฺถ กตมํ สมฺปชญฺญํ, ยา ปญฺญา ปชานนา ฯเปฯ อโมโห ธมฺมวิจโย สมฺมาทิฏฺฐิ, อิทํ วุจฺจติ สมฺปชญฺญํ.\csromanspacer{} อิติ อิมาย จ สติยา อิมินา จ สมฺปชญฺเญน อุเปโต โหติ ฯเปฯ สมนฺนาคโต, เตน วุจฺจติ “สโต จ สมฺปชาโน”ติ.}

\proseitem{587}{“สุขญฺจ กาเยน ปฏิสํเวเทตี”ติ ตตฺถ กตมํ สุขํ, ยํ เจตสิกํ สาตํ เจตสิกํ สุขํ เจโตสมฺผสฺสชํ สาตํ สุขํ เวทยิตํ เจโตสมฺผสฺสชา สาตา สุขา เวทนา, อิทํ วุจฺจติ สุขํ.}

\prose{ตตฺถ กตโม กาโย, สญฺญากฺขนฺโธ สงฺขารกฺขนฺโธ วิญฺญาณกฺขนฺโธ, อยํ วุจฺจติ กาโย.\csromanspacer{} อิทํ สุขํ อิมินา กาเยน ปฏิสํเวเทติ, เตน วุจฺจติ “สุขญฺจ กาเยน ปฏิสํเวเทตี”ติ.}

\proseitem{588}{\csromanfolio{270}“ยํ ตํ อริยา อาจิกฺขนฺตี”ติ ตตฺถ กตเม อริยา, อริยา วุจฺจนฺติ พุทฺธา จ พุทฺธสาวกา จ, เต อิมํ อาจิกฺขนฺติ เทเสนฺติ ปญฺญเปนฺติ\footnote{ปญฺญาเปนฺติ (สี, สฺยา)} ปฏฺฐเปนฺติ วิวรนฺติ วิภชนฺติ อุตฺตานิํ กโรนฺติ ปกาเสนฺติ, เตน วุจฺจติ “ยํ ตํ อริยา อาจิกฺขนฺตี”ติ.}

\proseitem{589}{“อุเปกฺขโก สติมา สุขวิหารี”ติ ตตฺถ กตมา อุเปกฺขา, ยา อุเปกฺขา อุเปกฺขนา อชฺฌุเปกฺขนา มชฺฌตฺตตา จิตฺตสฺส, อยํ วุจฺจติ อุเปกฺขา.}

\prose{ตตฺถ กตมา สติ, ยา สติ อนุสฺสติ ฯเปฯ สมฺมาสติ, อยํ วุจฺจติ สติ.}

\prose{ตตฺถ กตมํ สติ, ยํ เจตสิกํ สาตํ เจตสิกํ สุขํ เจโตสมฺผสฺสชํ สาตํ สุขํ เวทยิตํ เจโตสมฺผสฺสชา สาตา สุขา เวทนา, อิทํ วุจฺจติ สุขํ.\csromanspacer{} อิติ อิมาย จ อุเปกฺขาย อิมาย จ สติยา อิมินา จ สุเขน สมนฺนาคโต อิริยติ วตฺตติ ปาเลติ ยเปติ ยาเปติ จรติ วิหรติ, เตน วุจฺจติ “อุเปกฺขโก สติมา สุขวิหารี”ติ.}

\proseitem{590}{“ตติยนฺ”ติ คณนานุปุพฺพตา ตติยํ, อิทํ ตติยํ สมาปชฺชตีติ ตติยํ.}

\proseitem{591}{“ฌานนฺ”ติ อุเปกฺขา สติ สมฺปชญฺญํ สุขํ จิตฺตสฺเสกคฺคตา.}

\proseitem{592}{“อุปสมฺปชฺชา”ติ โย ตติยสฺส ฌานสฺส ลาโภ ปฏิลาโภ ปตฺติ สมฺปตฺติ ผุสนา สจฺฉิกิริยา อุปสมฺปทา.}

\proseitem{593}{“วิหรตี”ติ ฯเปฯ เตน วุจฺจติ “วิหรตี”ติ.}

\proseitem{594}{“สุขสฺส จ ปหานา ทุกฺขสฺส จ ปหานา”ติ อตฺถิ สุขํ, อตฺถิ ทุกฺขํ.}

\prose{ตตฺถ กตมํ สุขํ, ยํ กายิกํ สาตํ กายิกํ สุขํ กายสมฺผสฺสชํ สาตํ สุขํ เวทยิตํ กายสมฺผสฺสชา สาตา สุขา เวทนา, อิทํ วุจฺจติ สุขํ.}

\csromanheader{2}{12. ฌานวิภงฺค}
\prose{\csromanfolio{271}ตตฺถ กตมํ ทุกฺขํ, ยํ กายิกํ อสาตํ กายิกํ ทุกฺขํ กายสมฺผสฺสชํ อสาตํ ทุกฺขํ เวทยิตํ กายสมฺผสฺสชา อสาตา ทุกฺขา เวทนา, อิทํ วุจฺจติ ทุกฺขํ.\csromanspacer{} อิติ อิทญฺจ สุขํ อิทญฺจ ทุกฺขํ สนฺตา โหนฺติ สมิตา วูปสนฺตา อตฺถงฺคตา อพฺภตฺถงฺคตา อปฺปิตา พฺยปฺปิตา โสสิตา วิโสสิตา พฺยนฺตีกตา, เตน วุจฺจติ “สุขสฺส จ ปหานา ทุกฺขสฺส จ ปหานา”ติ.}

\proseitem{595}{“ปุพฺเพว โสมนสฺสโทมนสฺสานํ อตฺถงฺคมา”ติ อตฺถิ โสมนสฺสํ, อตฺถิ โทมนสฺสํ.}

\prose{ตตฺถ กตมํ โสมนสฺสํ, ยํ เจตสิกํ สาตํ เจตสิกํ สุขํ เจโตสมฺผสฺสชํ สาตํ สุขํ เวทยิตํ เจโตสมฺผสฺสชา สาตา สุขา เวทนา, อิทํ วุจฺจติ โสมนสฺสํ.}

\prose{ตตฺถ กตมํ โทมนสฺสํ, ยํ เจตสิกํ อสาตํ เจตสิกํ ทุกฺขํ เจโตสมฺผสฺสชํ อสาตํ ทุกฺขํ เวทยิตํ เจโตสมฺผสฺสชา อสาตา ทุกฺขา เวทนา, อิทํ วุจฺจติ โทมนสฺสํ.\csromanspacer{} อิติ อิทญฺจ โสมนสฺสํ อิทญฺจ โทมนสฺสํ ปุพฺเพว สนฺตา โหนฺติ สมิตา วูปสนฺตา อตฺถงฺคตา อพฺภตฺถงฺคตา อปฺปิตา พฺยปฺปิตา โสสิตา วิโสสิตา พฺยนฺตีกตา, เตน วุจฺจติ “ปุพฺเพว โสมนสฺสโทมนสฺสานํ อตฺถงฺคมา”ติ.}

\proseitem{596}{“อทุกฺขมสุขนฺ”ติ ยํ เจตสิกํ เนว สาตํ นาสาตํ เจโตสมฺผสฺสชํ อทุกฺขมสุขํ เวทยิตํ เจโตสมฺผสฺสชา อทุกฺขมสุขา เวทนา, เตน วุจฺจติ “อทุกฺขมสุขนฺ”ติ.}

\proseitem{597}{“อุเปกฺขาสติปาริสุทฺธินฺ”ติ ตตฺถ กตมา อุเปกฺขา, ยา อุเปกฺขา อุเปกฺขนา อชฺฌุเปกฺขนา มชฺฌตฺตตา จิตฺตสฺส, อยํ วุจฺจติ อุเปกฺขา.}

\prose{ตตฺถ กตมา สติ, ยา สติ อนุสฺสติ ฯเปฯ สมฺมาสติ, อยํ วุจฺจติ สติ.\csromanspacer{} อยํ สติ อิมาย อุเปกฺขาย วิวฏา โหติ ปริสุทฺธา ปริโยทาตา, เตน วุจฺจติ “อุเปกฺขาสติปาริสุทฺธินฺ”ติ.}

\proseitem{598}{“จตุตฺถนฺ”ติ คณนานุปุพฺพตา จตุตฺถํ, อิทํ จตุตฺถํ สมาปชฺชตีติ จตุตฺถํ.}

\proseitem{599}{“ฌานนฺ”ติ อุเปกฺขา สติ จิตฺตสฺเสกคฺคตา.}

\proseitem{600}{\csromanfolio{272}“อุปสมฺปชฺชา”ติ โย จตุตฺถสฺส ฌานสฺส ลาโภ ปฏิลาโภ ปตฺติ สมฺปตฺติ ผุสนา สจฺฉิกิริยา อุปสมฺปทา.}

\proseitem{601}{“วิหรตี”ติ ฯเปฯ เตน วุจฺจติ “วิหรตี”ติ.}

\proseitem{602}{“สพฺพโส รูปสญฺญานํ สมติกฺกมา”ติ ตตฺถ กตมา รูปสญฺญา, รูปาวจรสมาปตฺติํ สมาปนฺนสฺส วา อุปปนฺนสฺส วา ทิฏฺฐธมฺมสุขวิหาริสฺส วา สญฺญา สญฺชานนา สญฺชานิตตฺตํ, อิมา วุจฺจนฺติ รูปสญฺญาโย.\csromanspacer{} อิมา รูปสญฺญาโย อติกฺกนฺโต โหติ วีติกฺกนฺโต สมติกฺกนฺโต, เตน วุจฺจติ “สพฺพโส รูปสญฺญานํ สมติกฺกมา”ติ.}

\proseitem{603}{“ปฏิฆสญฺญานํ อตฺถงฺคมา”ติ ตตฺถ กตมา ปฏิฆสญฺญา, รูปสญฺญา สทฺทสญฺญา ฯเปฯ โผฏฺฐพฺพสญฺญา, อิมา วุจฺจนฺติ ปฏิฆสญฺญาโย.\csromanspacer{} อิมา ปฏิฆสญฺญาโย สนฺตา โหนฺติ สมิตา วูปสนฺตา อตฺถงฺคตา อพฺภตฺถงฺคตา อปฺปิตา พฺยปฺปิตา โสสิตา วิโสสิตา พฺยนฺตีกตา, เตน วุจฺจติ “ปฏิฆสญฺญานํ อตฺถงฺคมา”ติ.}

\proseitem{604}{“นานตฺตสญฺญานํ อมนสิการา”ติ ตตฺถ กตมา นานตฺตสญฺญา, อสมาปนฺนสฺส มโนธาตุสมงฺคิสฺส วา มโนวิญฺญาณธาตุสมงฺคิสฺส วา สญฺญา สญฺชานนา สญฺชานิตตฺตํ, อิมา วุจฺจนฺติ นานตฺตสญฺญาโย.\csromanspacer{} อิมา นานตฺตสญฺญาโย น มนสิ กโรติ, เตน วุจฺจติ “นานตฺตสญฺญานํ อมนสิการา”ติ.}

\proseitem{605}{“อนนฺโต อากาโส”ติ ตตฺถ กตโม อากาโส, โย อากาโส อากาสคตํ อฆํ อฆคตํ วิวโร วิวรคตํ อสมฺผุฏฺฐํ จตูหิ มหาภูเตหิ, อยํ วุจฺจติ อากาโส.\csromanspacer{} ตสฺมิํ อากาเส จิตฺตํ ฐเปติ สณฺฐเปติ อนนฺตํ ผรติ, เตน วุจฺจติ “อนนฺโต อากาโส”ติ.}

\proseitem{606}{“อากาสานญฺจายตนนฺ”ติ อากาสานญฺจายตนํ สมาปนฺนสฺส วา อุปปนฺนสฺส วา ทิฏฺฐธมฺมสุขวิหาริสฺส วา จิตฺตเจตสิกา ธมฺมา.}

\proseitem{607}{“อุปสมฺปชฺชา”ติ โย อากาสานญฺจายตนสฺส ลาโภ ปฏิลาโภ ปตฺติ สมฺปตฺติ ผุสนา สจฺฉิกิริยา อุปสมฺปทา.}

\csromanheader{2}{12. ฌานวิภงฺค}
\proseitem{608}{\csromanfolio{273}“วิหรตี”ติ ฯเปฯ เตน วุจฺจติ “วิหรตี”ติ.}

\proseitem{609}{“สพฺพโส อากาสานญฺจายตนํ สมติกฺกมฺมา”ติ อิมํ อากาสานญฺจายตนํ อติกฺกนฺโต โหติ วีติกฺกนฺโต สมติกฺกนฺโต, เตน วุจฺจติ “สพฺพโส อากาสานญฺจายตนํ สมติกฺกมฺมา”ติ.}

\proseitem{610}{“อนนฺตํ วิญฺญาณนฺ”ติ ตํเยว อากาสํ วิญฺญาเณน ผุฏฺฐํ มนสิ กโรติ อนนฺตํ ผรติ, เตน วุจฺจติ “อนนฺตํ วิญฺญาณนฺ”ติ.}

\proseitem{611}{“วิญฺญาณญฺจายตนนฺ”ติ วิญฺญาณญฺจายตนํ สมาปนฺนสฺส วา อุปปนฺนสฺส วา ทิฏฺฐธมฺมสุขวิหาริสฺส วา จิตฺตเจตสิกา ธมฺมา.}

\proseitem{612}{“อุปสมฺปชฺชา”ติ โย วิญฺญาณญฺจายตนสฺส ลาโภ ปฏิลาโภ ปตฺติ สมฺปตฺติ ผุสนา สจฺฉิกิริยา อุปสมฺปทา.}

\proseitem{613}{“วิหรตี”ติ ฯเปฯ เตน วุจฺจติ “วิหรตี”ติ.}

\proseitem{614}{“สพฺพโส วิญฺญาณญฺจายตนํ สมติกฺกมฺมา”ติ อิมํ วิญฺญาณญฺจายตนํ อติกฺกนฺโต โหติ วีติกฺกนฺโต สมติกฺกนฺโต, เตน วุจฺจติ “สพฺพโส วิญฺญาณญฺจายตนํ สมติกฺกมฺมา”ติ.}

\proseitem{615}{“นตฺถิ กิญฺจี”ติ ตํเยว วิญฺญาณํ ภาเวติ วิภาเวติ อนฺตรธาเปติ “นตฺถิ กิญฺจี”ติ ปสฺสติ, เตน วุจฺจติ “นตฺถิ กิญฺจี”ติ.}

\proseitem{616}{“อากิญฺจญฺญายตนนฺ”ติ อากิญฺจญฺญายตนํ สมาปนฺนสฺส วา อุปปนฺนสฺส วา ทิฏฺฐธมฺมสุขวิหาริสฺส วา จิตฺตเจตสิกา ธมฺมา.}

\proseitem{617}{“อุปสมฺปชฺชา”ติ โย อากิญฺจญฺญายตนสฺส ลาโภ ปฏิลาโภ ปตฺติ สมฺปตฺติ ผุสนา สจฺฉิกิริยา อุปสมฺปทา.}

\proseitem{618}{“วิหรตี”ติ ฯเปฯ เตน วุจฺจติ “วิหรตี”ติ.}

\proseitem{619}{“สพฺพโส อากิญฺจญฺญายตนํ สมติกฺกมฺมา”ติ อิมํ อากิญฺจญฺญายตนํ อติกฺกนฺโต โหติ วีติกฺกนฺโต สมติกฺกนฺโต, เตน วุจฺจติ “สพฺพโส อากิญฺจญฺญายตนํ สมติกฺกมฺมา”ติ.}

\prose{\csromanfolio{274}เนวสญฺญีนาสญฺญีติ ตํเยว อากิญฺจญฺญายตนํ สนฺตโต มนสิ กโรติ สงฺขาราวเสสสมาปตฺติํ ภาเวติ, เตน วุจฺจติ เนวสญฺญีนาสญฺญีติ\footnote{อยํ ปาโฐ มาติกายํ นตฺถิ, นิทฺเทเส ปน สพฺพโปตฺถเกสุ ทิสฺสติ.}.}

\proseitem{620}{“เนวสญฺญานาสญฺญายตนนฺ”ติ เนวสญฺญานาสญฺญายตนํ สมาปนฺนสฺส วา อุปปนฺนสฺส วา ทิฏฺฐธมฺมสุขวิหาริสฺส วา จิตฺตเจตสิกา ธมฺมา.}

\proseitem{621}{“อุปสมฺปชฺชา”ติ โย เนวสญฺญานาสญฺญายตนสฺส ลาโภ ปฏิลาโภ ปตฺติ สมฺปตฺติ ผุสนา สจฺฉิกิริยา อุปสมฺปทา.}

\proseitem{622}{“วิหรตี”ติ อิริยติ วตฺตติ ปาเลติ ยเปติ ยาเปติ จรติ วิหรติ, เตน วุจฺจติ “วิหรตี”ติ.}

\csromanheader{2}{สุตฺตนฺตภาชนียํ.}
\csromansectionrule
\csromanmatikaanchor{mtk.105}
\csromantocmarkref{subparagraph}{2. อภิธมฺมภาชนีย}{mtk.105}
\bookmark[dest={mtk.105},level=5]{2. อภิธมฺมภาชนีย}
\csromanheader{6}{2. อภิธมฺมภาชนีย}
\csromanmatikaanchor{mtk.106}
\csromantocmarkref{subparagraph}{1. รูปาวจรกุสล}{mtk.106}
\bookmark[dest={mtk.106},level=5]{1. รูปาวจรกุสล}
\csromanheader{6}{1. รูปาวจรกุสล}
\proseitem{623}{จตฺตาริ ฌานานิ, ปฐมํ ฌานํ ทุติยํ ฌานํ ตติยํ ฌานํ จตุตฺถํ ฌานํ.}

\proseitem{624}{ตตฺถ กตมํ ปฐมํ ฌานํ, อิธ ภิกฺขุ ยสฺมิํ สมเย รูปูปปตฺติยา มคฺคํ ภาเวติ วิวิจฺเจว กาเมหิ ฯเปฯ ปฐมํ ฌานํ อุปสมฺปชฺช วิหรติ ปถวีกสิณํ, ตสฺมิํ สมเย ปญฺจงฺคิกํ ฌานํ โหติ วิตกฺโก วิจาโร ปีติ สุขํ จิตฺตสฺเสกคฺคตา, อิทํ วุจฺจติ ปฐมํ ฌานํ.\csromanspacer{} อวเสสา ธมฺมา ฌานสมฺปยุตฺตา. (1)}

\prose{ตตฺถ กตมํ ทุติยํ ฌานํ, อิธ ภิกฺขุ ยสฺมิํ สมเย รูปูปปตฺติยา มคฺคํ ภาเวติ วิตกฺกวิจารานํ วูปสมา ฯเปฯ ทุติยํ ฌานํ อุปสมฺปชฺช วิหรติ}

\vspace{\tipitakaparskip}
\csromancenter{ฌานวิภงฺค ปถวีกสิณํ, ตสฺมิํ สมเย ติวงฺคิกํ ฌานํ โหติ ปีติ สุขํ จิตฺตสฺเสกคฺคตา, อิทํ วุจฺจติ ทุติยํ ฌานํ.\csromanspacer{} อวเสสา ธมฺมา ฌานสมฺปยุตฺตา. (2)}
\prose{\csromanfolio{275}ตตฺถ กตมํ ตติยํ ฌานํ, อิธ ภิกฺขุ ยสฺมิํ สมเย รูปูปปตฺติยา มคฺคํ ภาเวติ ปีติยา จ วิราคา ฯเปฯ ตติยํ ฌานํ อุปสมฺปชฺช วิหรติ ปถวีกสิณํ, ตสฺมิํ สมเย ทุวงฺคิกํ ฌานํ โหติ สุขํ จิตฺตสฺเสกคฺคตา, อิทํ วุจฺจติ ตติยํ ฌานํ.\csromanspacer{} อวเสสา ธมฺมา ฌานสมฺปยุตฺตา. (3)}

\prose{ตตฺถ กตมํ จตุตฺถํ ฌานํ, อิธ ภิกฺขุ ยสฺมิํ สมเย รูปูปปตฺติยา มคฺคํ ภาเวติ สุขสฺส จ ปหานา ฯเปฯ จตุตฺถํ ฌานํ อุปสมฺปชฺช วิหรติ ปถวีกสิณํ, ตสฺมิํ สมเย ทุวงฺคิกํ ฌานํ โหติ อุเปกฺขา จิตฺตสฺเสกคฺคตา, อิทํ วุจฺจติ จตุตฺถํ ฌานํ.\csromanspacer{} อวเสสา ธมฺมา ฌานสมฺปยุตฺตา. (4) จตุกฺกํ.}

\proseitem{625}{อิธ ภิกฺขุ ยสฺมิํ สมเย รูปูปปตฺติยา มคฺคํ ภาเวติ วิวิจฺเจว กาเมหิ ฯเปฯ ปฐมํ ฌานํ อุปสมฺปชฺช วิหรติ ปถวีกสิณํ, ตสฺมิํ สมเย ปญฺจงฺคิกํ ฌานํ โหติ วิตกฺโก วิจาโร ปีติ สุขํ จิตฺตสฺเสกคฺคตา, อิทํ วุจฺจติ ปฐมํ ฌานํ.\csromanspacer{} อวเสสา ธมฺมา ฌานสมฺปยุตฺตา. (1)}

\prose{อิธ ภิกฺขุ ยสฺมิํ สมเย รูปูปปตฺติยา มคฺคํ ภาเวติ วิวิจฺเจว กาเมหิ วิวิจฺจ อกุสเลหิ ธมฺเมหิ อวิตกฺกํ วิจารมตฺตํ วิเวกชํ ปีติสุขํ ทุติยํ ฌานํ อุปสมฺปชฺช วิหรติ ปถวีกสิณํ, ตสฺมิํ สมเย จตุรงฺคิกํ ฌานํ โหติ วิจาโร ปีติ สุขํ จิตฺตสฺเสกคฺคตา, อิทํ วุจฺจติ ทุติยํ ฌานํ.\csromanspacer{} อวเสสา ธมฺมา ฌานสมฺปยุตฺตา. (2)}

\prose{อิธ ภิกฺขุ ยสฺมิํ สมเย รูปูปปตฺติยา มคฺคํ ภาเวติ วิตกฺกวิจารานํ วูปสมา ฯเปฯ ตติยํ ฌานํ อุปสมฺปชฺช วิหรติ ปถวีกสิณํ, ตสฺมิํ สมเย ติวงฺคิกํ ฌานํ โหติ ปีติ สุขํ จิตฺตสฺเสกคฺคตา, อิทํ วุจฺจติ ตติยํ ฌานํ.\csromanspacer{} อวเสสา ธมฺมา ฌานสมฺปยุตฺตา. (3)}

\prose{อิธ ภิกฺขุ ยสฺมิํ สมเย รูปูปปตฺติยา มคฺคํ ภาเวติ ปีติยา จ วิราคา ฯเปฯ จตุตฺถํ ฌานํ อุปสมฺปชฺช วิหรติ ปถวีกสิณํ, ตสฺมิํ สมเย ทุวงฺคิกํ ฌานํ โหติ สุขํ จิตฺตสฺเสกคฺคตา, อิทํ วุจฺจติ จตุตฺถํ ฌานํ.\csromanspacer{} อวเสสา ธมฺมา ฌานสมฺปยุตฺตา. (4)}

\prose{\csromanfolio{276}อิธ ภิกฺขุ ยสฺมิํ สมเย รูปูปปตฺติยา มคฺคํ ภาเวติ สุขสฺส จ ปหานา ฯเปฯ ปญฺจมํ ฌานํ อุปสมฺปชฺช วิหรติ ปถวีกสิณํ, ตสฺมิํ สมเย ทุวงฺคิกํ ฌานํ โหติ อุเปกฺขา จิตฺตสฺเสกคฺคตา, อิทํ วุจฺจติ ปญฺจมํ ฌานํ.\csromanspacer{} อวเสสา ธมฺมา ฌานสมฺปยุตฺตา. (5)}

\csromanheader{2}{ปญฺจกํ.}
\csromansectionrule
\csromanmatikaanchor{mtk.107}
\csromantocmarkref{subparagraph}{2. อรูปาวจรกุสล}{mtk.107}
\bookmark[dest={mtk.107},level=5]{2. อรูปาวจรกุสล}
\csromanheader{6}{2. อรูปาวจรกุสล}
\proseitem{626}{อิธ ภิกฺขุ ยสฺมิํ สมเย อรูปูปปตฺติยา มคฺคํ ภาเวติ สพฺพโส อากิญฺจญฺญายตนํ สมติกฺกมฺม เนวสญฺญานาสญฺญายตนสญฺญาสหคตํ สุขสฺส จ ปหานา ฯเปฯ จตุตฺถํ ฌานํ อุปสมฺปชฺช วิหรติ, ตสฺมิํ สมเย ทุวงฺคิกํ ฌานํ โหติ อุเปกฺขา จิตฺตสฺเสกคฺคตา, อิทํ วุจฺจติ จตุตฺถํ ฌานํ.\csromanspacer{} อวเสสา ธมฺมา ฌานสมฺปยุตฺตา.}

\csromanmatikaanchor{mtk.108}
\csromantocmarkref{subparagraph}{3. โลกุตฺตรกุสล}{mtk.108}
\bookmark[dest={mtk.108},level=5]{3. โลกุตฺตรกุสล}
\csromanheader{6}{3. โลกุตฺตรกุสล}
\proseitem{627}{จตฺตาริ ฌานานิ, ปฐมํ ฌานํ ทุติยํ ฌานํ ตติยํ ฌานํ จตุตฺถํ ฌานํ.}

\proseitem{628}{ตตฺถ กตมํ ปฐมํ ฌานํ, อิธ ภิกฺขุ ยสฺมิํ สมเย โลกุตฺตรํ ฌานํ ภาเวติ นิยฺยานิกํ อปจยคามิํ ทิฏฺฐิคตานํ ปหานาย ปฐมาย ภูมิยา ปตฺติยา วิวิจฺเจว กาเมหิ ฯเปฯ ปฐมํ ฌานํ อุปสมฺปชฺช วิหรติ ทุกฺขปฏิปทํ ทนฺธาภิญฺญํ, ตสฺมิํ สมเย ปญฺจงฺคิกํ ฌานํ โหติ วิตกฺโก วิจาโร ปีติ สุขํ จิตฺตสฺเสกคฺคตา, อิทํ วุจฺจติ ปฐมํ ฌานํ.\csromanspacer{} อวเสสา ธมฺมา ฌานสมฺปยุตฺตา.}

\prose{ตตฺถ กตมํ ทุติยํ ฌานํ, อิธ ภิกฺขุ ยสฺมิํ สมเย โลกุตฺตรํ ฌานํ ภาเวติ นิยฺยานิกํ อปจยคามิํ ทิฏฺฐิคตานํ ปหานาย ปฐมาย ภูมิยา ปตฺติยา วิตกฺกวิจารานํ วูปสมา ฯเปฯ ทุติยํ ฌานํ อุปสมฺปชฺช วิหรติ ทุกฺขปฏิปทํ ทนฺธาภิญฺญํ, ตสฺมิํ สมเย ติวงฺคิกํ ฌานํ โหติ ปีติ สุขํ จิตฺตสฺเสกคฺคตา, อิทํ วุจฺจติ ทุติยํ ฌานํ.\csromanspacer{} อวเสสา ธมฺมา ฌานสมฺปยุตฺตา.}

\csromanheader{2}{12. ฌานวิภงฺค}
\prose{\csromanfolio{277}ตตฺถ กตมํ ตติยํ ฌานํ, อิธ ภิกฺขุ ยสฺมิํ สมเย โลกุตฺตรํ ฌานํ ภาเวติ นิยฺยานิกํ อปจยคามิํ ทิฏฺฐิคตานํ ปหานาย ปฐมาย ภูมิยา ปตฺติยา ปีติยา จ วิราคา ฯเปฯ ตติยํ ฌานํ อุปสมฺปชฺช วิหรติ ทุกฺขปฏิปทํ ทนฺธาภิญฺญํ, ตสฺมิํ สมเย ทุวงฺคิกํ ฌานํ โหติ สุขํ จิตฺตสฺเสกคฺคตา, อิทํ วุจฺจติ ตติยํ ฌานํ.\csromanspacer{} อวเสสา ธมฺมา ฌานสมฺปยุตฺตา.}

\prose{ตตฺถ กตมํ จตุตฺถํ ฌานํ, อิธ ภิกฺขุ ยสฺมิํ สมเย โลกุตฺตรํ ฌานํ ภาเวติ นิยฺยานิกํ อปจยคามิํ ทิฏฺฐิคตานํ ปหานาย ปฐมาย ภูมิยา ปตฺติยา สุขสฺส จ ปหานา ฯเปฯ จตุตฺถํ ฌานํ อุปสมฺปชฺช วิหรติ ทุกฺขปฏิปทํ ทนฺธาภิญฺญํ, ตสฺมิํ สมเย ทุวงฺคิกํ ฌานํ โหติ อุเปกฺขา จิตฺตสฺเสกคฺคตา, อิทํ วุจฺจติ จตุตฺถํ ฌานํ.\csromanspacer{} อวเสสา ธมฺมา ฌานสมฺปยุตฺตา.\csromanspacer{} จตุกฺกํ.}

\proseitem{629}{อิธ ภิกฺขุ ยสฺมิํ สมเย โลกุตฺตรํ ฌานํ ภาเวติ นิยฺยานิกํ อปจยคามิํ ทิฏฺฐิคตานํ ปหานาย ปฐมาย ภูมิยา ปตฺติยา วิวิจฺเจว กาเมหิ ฯเปฯ ปฐมํ ฌานํ อุปสมฺปชฺช วิหรติ ทุกฺขปฏิปทํ ทนฺธาภิญฺญํ, ตสฺมิํ สมเย ปญฺจงฺคิกํ ฌานํ โหติ วิตกฺโก วิจาโร ปีติ สุขํ จิตฺตสฺเสกคฺคตา, อิทํ วุจฺจติ ปฐมํ ฌานํ.\csromanspacer{} อวเสสา ธมฺมา ฌานสมฺปยุตฺตา.}

\prose{อิธ ภิกฺขุ ยสฺมิํ สมเย โลกุตฺตรํ ฌานํ ภาเวติ นิยฺยานิกํ อปจยคามิํ ทิฏฺฐิคตานํ ปหานาย ปฐมาย ภูมิยา ปตฺติยา วิวิจฺเจว กาเมหิ วิวิจฺจ อกุสเลหิ ธมฺเมหิ อวิตกฺกํ วิจารมตฺตํ วิเวกชํ ปีติสุขํ ทุติยํ ฌานํ อุปสมฺปชฺช วิหรติ ทุกฺขปฏิปทํ ทนฺธาภิญฺญํ, ตสฺมิํ สมเย จตุรงฺคิกํ ฌานํ โหติ วิจาโร ปีติ สุขํ จิตฺตสฺเสกคฺคตา, อิทํ วุจฺจติ ทุติยํ ฌานํ.\csromanspacer{} อวเสสา ธมฺมา ฌานสมฺปยุตฺตา.}

\prose{อิธ ภิกฺขุ ยสฺมิํ สมเย โลกุตฺตรํ ฌานํ ภาเวติ นิยฺยานิกํ อปจยคามิํ ทิฏฺฐิคตานํ ปหานาย ปฐมาย ภูมิยา ปตฺติยา วิตกฺกวิจารานํ วูปสมา ฯเปฯ ตติยํ ฌานํ อุปสมฺปชฺช วิหรติ ทุกฺขปฏิปทํ ทนฺธาภิญฺญํ, ตสฺมิํ สมเย ติวงฺคิกํ ฌานํ โหติ ปีติ สุขํ จิตฺตสฺเสกคฺคตา, อิทํ วุจฺจติ ตติยํ ฌานํ.\csromanspacer{} อวเสสา ธมฺมา ฌานสมฺปยุตฺตา.}

\prose{อิธ ภิกฺขุ ยสฺมิํ สมเย โลกุตฺตรํ ฌานํ ภาเวติ นิยฺยานิกํ อปจยคามิํ ทิฏฺฐิคตานํ ปหานาย ปฐมาย ภูมิยา ปตฺติยา ปีติยา จ \csromanfolio{278}วิราคา ฯเปฯ จตุตฺถํ ฌานํ อุปสมฺปชฺช วิหรติ ทุกฺขปฏิปทํ ทนฺธาภิญฺญํ, ตสฺมิํ สมเย ทุวงฺคิกํ ฌานํ โหติ สุขํ จิตฺตสฺเสกคฺคตา, อิทํ วุจฺจติ จตุตฺถํ ฌานํ.\csromanspacer{} อวเสสา ธมฺมา ฌานสมฺปยุตฺตา.}

\prose{อิธ ภิกฺขุ ยสฺมิํ สมเย โลกุตฺตรํ ฌานํ ภาเวติ นิยฺยานิกํ อปจยคามิํ ทิฏฺฐิคตานํ ปหานาย ปฐมาย ภูมิยา ปตฺติยา สุขสฺส จ ปหานา ฯเปฯ ปญฺจมํ ฌานํ อุปสมฺปชฺช วิหรติ ทุกฺขปฏิปทํ ทนฺธาภิญฺญํ, ตสฺมิํ สมเย ทุวงฺคิกํ ฌานํ โหติ อุเปกฺขา จิตฺตสฺเสกคฺคตา, อิทํ วุจฺจติ ปญฺจมํ ฌานํ.\csromanspacer{} อวเสสา ธมฺมา ฌานสมฺปยุตฺตา.}

\csromanheader{2}{ปญฺจกํ.}
\csromansectionrule
\csromanmatikaanchor{mtk.109}
\csromantocmarkref{subparagraph}{4. รูปาวจรวิปาก}{mtk.109}
\bookmark[dest={mtk.109},level=5]{4. รูปาวจรวิปาก}
\csromanheader{6}{4. รูปาวจรวิปาก}
\proseitem{630}{จตฺตาริ ฌานานิ, ปฐมํ ฌานํ ทุติยํ ฌานํ ตติยํ ฌานํ จตุตฺถํ ฌานํ.}

\proseitem{631}{ตตฺถ กตมํ ปฐมํ ฌานํ, อิธ ภิกฺขุ ยสฺมิํ สมเย รูปูปปตฺติยา มคฺคํ ภาเวติ วิวิจฺเจว กาเมหิ ฯเปฯ ปฐมํ ฌานํ อุปสมฺปชฺช วิหรติ ปถวีกสิณํ, ตสฺมิํ สมเย ผสฺโส โหติ ฯเปฯ อวิกฺเขโป โหติ, อิเม ธมฺมา กุสลา.\csromanspacer{} ตสฺเสว รูปาวจรสฺส กุสลสฺส กมฺมสฺส กตตฺตา อุปจิตตฺตา วิปากํ วิวิจฺเจว กาเมหิ ฯเปฯ ปฐมํ ฌานํ อุปสมฺปชฺช วิหรติ ปถวีกสิณํ, ตสฺมิํ สมเย ปญฺจงฺคิกํ ฌานํ โหติ วิตกฺโก วิจาโร ปีติ สุขํ จิตฺตสฺเสกคฺคตา, อิทํ วุจฺจติ ปฐมํ ฌานํ.\csromanspacer{} อวเสสา ธมฺมา ฌานสมฺปยุตฺตา.}

\prose{ตตฺถ กตมํ ทุติยํ ฌานํ, อิธ ภิกฺขุ ยสฺมิํ สมเย รูปูปปตฺติยา มคฺคํ ภาเวติ วิตกฺกวิจารานํ วูปสมา ฯเปฯ ทุติยํ ฌานํ อุปสมฺปชฺช วิหรติ ปถวีกสิณํ, ตสฺมิํ สมเย ผสฺโส โหติ ฯเปฯ อวิกฺเขโป โหติ, อิเม ธมฺมา กุสลา.\csromanspacer{} ตสฺเสว รูปาวจรสฺส กุสลสฺส กมฺมสฺส กตตฺตา อุปจิตตฺตา วิปากํ วิตกฺกวิจารานํ วูปสมา ฯเปฯ ทุติยํ ฌานํ ฯเปฯ ตติยํ ฌานํ ฯเปฯ จตุตฺถํ ฌานํ ฯเปฯ ปฐมํ ฌานํ ฯเปฯ ปญฺจมํ ฌานํ อุปสมฺปชฺช วิหรติ ปถวีกสิณํ, ตสฺมิํ สมเย ทุวงฺคิกํ ฌานํ โหติ อุเปกฺขา จิตฺตสฺเสกคฺคตา, อิทํ วุจฺจติ ปญฺจมํ ฌานํ.\csromanspacer{} อวเสสา ธมฺมา ฌานสมฺปยุตฺตา ฯเปฯ.}

\csromanmatikaanchor{mtk.110}
\csromantocmarkref{subparagraph}{5. อรูปาวจรวิปาก}{mtk.110}
\bookmark[dest={mtk.110},level=5]{5. อรูปาวจรวิปาก}
\csromanheader{6}{12. ฌานวิภงฺค \textbf{5. อรูปาวจรวิปาก}}
\proseitem{632}{\csromanfolio{279}อิธ ภิกฺขุ ยสฺมิํ สมเย อรูปูปปตฺติยา มคฺคํ ภาเวติ สพฺพโส อากิญฺจญฺญายตนํ สมติกฺกมฺม เนวสญฺญานาสญฺญายตนสญฺญาสหคตํ สุขสฺส จ ปหานา ฯเปฯ จตุตฺถํ ฌานํ อุปสมฺปชฺช วิหรติ, ตสฺมิํ สมเย ผสฺโส โหติ ฯเปฯ อวิกฺเขโป โหติ, อิเม ธมฺมา กุสลา.\csromanspacer{} ตสฺเสว อรูปาวจรสฺส กุสลสฺส กมฺมสฺส กตตฺตา อุปจิตตฺตา วิปากํ สพฺพโส อากิญฺจญฺญายตนํ สมติกฺกมฺม เนวสญฺญานาสญฺญายตนสญฺญาสหคตํ สุขสฺส จ ปหานา ฯเปฯ จตุตฺถํ ฌานํ อุปสมฺปชฺช วิหรติ, ตสฺมิํ สมเย ทุวงฺคิกํ ฌานํ โหติ อุเปกฺขา จิตฺตสฺเสกคฺคตา, อิทํ วุจฺจติ จตุตฺถํ ฌานํ.\csromanspacer{} อวเสสา ธมฺมา ฌานสมฺปยุตฺตา.}

\csromanmatikaanchor{mtk.111}
\csromantocmarkref{subparagraph}{6. โลกุตฺตรวิปาก}{mtk.111}
\bookmark[dest={mtk.111},level=5]{6. โลกุตฺตรวิปาก}
\csromanheader{6}{6. โลกุตฺตรวิปาก}
\proseitem{633}{จตฺตาริ ฌานานิ, ปฐมํ ฌานํ ทุติยํ ฌานํ ตติยํ ฌานํ จตุตฺถํ ฌานํ.}

\proseitem{634}{ตตฺถ กตมํ ปฐมํ ฌานํ, อิธ ภิกฺขุ ยสฺมิํ สมเย โลกุตฺตรํ ฌานํ ภาเวติ นิยฺยานิกํ อปจยคามิํ ทิฏฺฐิคตานํ ปหานาย ปฐมาย ภูมิยา ปตฺติยา วิวิจฺเจว กาเมหิ ฯเปฯ ปฐมํ ฌานํ อุปสมฺปชฺช วิหรติ ทุกฺขปฏิปทํ ทนฺธาภิญฺญํ, ตสฺมิํ สมเย ผสฺโส โหติ ฯเปฯ อวิกฺเขโป โหติ, อิเม ธมฺมา กุสลา.\csromanspacer{} ตสฺเสว โลกุตฺตรสฺส กุสลสฺส ฌานสฺส กตตฺตา ภาวิตตฺตา วิปากํ วิวิจฺเจว กาเมหิ ฯเปฯ ปฐมํ ฌานํ อุปสมฺปชฺช วิหรติ ทุกฺขปฏิปทํ ทนฺธาภิญฺญํ สุญฺญตํ, ตสฺมิํ สมเย ปญฺจงฺคิกํ ฌานํ โหติ วิตกฺโก วิจาโร ปีติ สุขํ จิตฺตสฺเสกคฺคตา, อิทํ วุจฺจติ ปฐมํ ฌานํ.\csromanspacer{} อวเสสา ธมฺมา ฌานสมฺปยุตฺตา.}

\prose{ตตฺถ กตมํ ทุติยํ ฌานํ, อิธ ภิกฺขุ ยสฺมิํ สมเย โลกุตฺตรํ ฌานํ ภาเวติ นิยฺยานิกํ อปจยคามิํ ทิฏฺฐิคตานํ ปหานาย ปฐมาย ภูมิยา ปตฺติยา วิตกฺกวิจารานํ วูปสมา ฯเปฯ ทุติยํ ฌานํ อุปสมฺปชฺช วิหรติ ทุกฺขปฏิปทํ ทนฺธาภิญฺญํ, ตสฺมิํ สมเย ผสฺโส โหติ ฯเปฯ อวิกฺเขโป โหติ, อิเม ธมฺมา กุสลา, ตสฺเสว โลกุตฺตรสฺส กุสลสฺส \csromanfolio{280}ฌานสฺส กตตฺตา ภาวิตตฺตา วิปากํ วิตกฺกวิจารานํ วูปสมา ฯเปฯ ทุติยํ ฌานํ ฯเปฯ ตติยํ ฌานํ ฯเปฯ จตุตฺถํ ฌานํ ฯเปฯ ปฐมํ ฌานํ ฯเปฯ ปญฺจมํ ฌานํ อุปสมฺปชฺช วิหรติ ทุกฺขปฏิปทํ ทนฺธาภิญฺญํ สุญฺญตํ, ตสฺมิํ สมเย ทุวงฺคิกํ ฌานํ โหติ อุเปกฺขา จิตฺตสฺเสกคฺคตา, อิทํ วุจฺจติ ปญฺจมํ ฌานํ.\csromanspacer{} อวเสสา ธมฺมา ฌานสมฺปยุตฺตา.}

\csromanmatikaanchor{mtk.112}
\csromantocmarkref{subparagraph}{7. รูปารูปาวจรกิริย}{mtk.112}
\bookmark[dest={mtk.112},level=5]{7. รูปารูปาวจรกิริย}
\csromanheader{6}{7. รูปารูปาวจรกิริย}
\proseitem{635}{จตฺตาริ ฌานานิ, ปฐมํ ฌานํ ทุติยํ ฌานํ ตติยํ ฌานํ จตุตฺถํ ฌานํ.}

\proseitem{636}{ตตฺถ กตมํ ปฐมํ ฌานํ, อิธ ภิกฺขุ ยสฺมิํ สมเย รูปาวจรํ ฌานํ ภาเวติ กิริยํ เนว กุสลํ นากุสลํ น จ กมฺมวิปากํ ทิฏฺฐธมฺมสุขวิหารํ วิวิจฺเจว กาเมหิ ฯเปฯ ปฐมํ ฌานํ อุปสมฺปชฺช วิหรติ ปถวีกสิณํ, ตสฺมิํ สมเย ปญฺจงฺคิกํ ฌานํ โหติ วิตกฺโก วิจาโร ปีติ สุขํ จิตฺตสฺเสกคฺคตา, อิทํ วุจฺจติ ปฐมํ ฌานํ.\csromanspacer{} อวเสสา ธมฺมา ฌานสมฺปยุตฺตา.}

\prose{ตตฺถ กตมํ ทุติยํ ฌานํ, อิธ ภิกฺขุ ยสฺมิํ สมเย รูปาวจรํ ฌานํ ภาเวติ กิริยํ เนว กุสลํ นากุสลํ น จ กมฺมวิปากํ ทิฏฺฐธมฺมสุขวิหารํ วิตกฺกวิจารานํ วูปสมา ฯเปฯ ทุติยํ ฌานํ ฯเปฯ ตติยํ ฌานํ ฯเปฯ จตุตฺถํ ฌานํ ฯเปฯ ปฐมํ ฌานํ ฯเปฯ ปญฺจมํ ฌานํ อุปสมฺปชฺช วิหรติ ปถวีกสิณํ, ตสฺมิํ สมเย ทุวงฺคิกํ ฌานํ โหติ อุเปกฺขา จิตฺตสฺเสกคฺคตา, อิทํ วุจฺจติ ปญฺจมํ ฌานํ, อวเสสา ธมฺมา ฌานสมฺปยุตฺตา.}

\proseitem{637}{อิธ ภิกฺขุ ยสฺมิํ สมเย อรูปาวจรํ ฌานํ ภาเวติ กิริยํ เนว กุสลํ นากุสลํ น จ กมฺมวิปากํ ทิฏฺฐธมฺมสุขวิหารํ สพฺพโส อากิญฺจญฺญายตนํ สมติกฺกมฺม เนวสญฺญานาสญฺญายตนสญฺญาสหคตํ สุขสฺส จ ปหานา ฯเปฯ จตุตฺถํ ฌานํ อุปสมฺปชฺช วิหรติ, ตสฺมิํ สมเย ทุวงฺคิกํ ฌานํ โหติ อุเปกฺขา จิตฺตสฺเสกคฺคตา, อิทํ วุจฺจติ จตุตฺถํ ฌานํ.\csromanspacer{} อวเสสา ธมฺมา ฌานสมฺปยุตฺตาติ.}

\vspace{\tipitakaparskip}
\csromancenter{อภิธมฺมภาชนียํ.}
\csromanheader{2}{12. ฌานวิภงฺค \textbf{3. ปญฺหาปุจฺฉก}}
\proseitem{638}{\csromanfolio{281}จตฺตาริ ฌานานิ, อิธ ภิกฺขุ วิวิจฺเจว กาเมหิ วิวิจฺจ อกุสเลหิ ธมฺเมหิ สวิตกฺกํ สวิจารํ วิเวกชํ ปีติสุขํ ปฐมํ ฌานํ อุปสมฺปชฺช วิหรติ ฯเปฯ วิตกฺกวิจารานํ วูปสมา อชฺฌตฺตํ สมฺปสาทนํ เจตโส เอโกทิภาวํ อวิตกฺกํ อวิจารํ สมาธิชํ ปีติสุขํ ทุติยํ ฌานํ อุปสมฺปชฺช วิหรติ ฯเปฯ ปีติยา จ วิราคา อุเปกฺขโก จ วิหรติ สโต จ สมฺปชาโน สุขญฺจ กาเยน ปฏิสํเวเทติ, ยํ ตํ อริยา อาจิกฺขนฺติ “อุเปกฺขโก สติมา สุขวิหารี”ติ, ตติยํ ฌานํ อุปสมฺปชฺช วิหรติ ฯเปฯ สุขสฺส จ ปหานา ทุกฺขสฺส จ ปหานา ปุพฺเพว โสมนสฺสโทมนสฺสานํ อตฺถงฺคมา อทุกฺขมสุขํ อุเปกฺขาสติปาริสุทฺธิํ จตุตฺถํ ฌานํ อุปสมฺปชฺช วิหรติ.}

\proseitem{639}{จตุนฺนํ ฌานานํ กติ กุสลา, กติ อกุสลา, กติ อพฺยากตา ฯเปฯ กติ สรณา, กติ อรณา.}

\csromanmatikaanchor{mtk.113}
\csromanmatikaanchor{mtk.114}
\csromantocmarknopage{subparagraph}{3. ปญฺหาปุจฺฉก}{mtk.113}
\bookmark[dest={mtk.113},level=5]{3. ปญฺหาปุจฺฉก}
\csromantocmarkref{subparagraph}{1. ติก}{mtk.114}
\bookmark[dest={mtk.114},level=5]{1. ติก}
\csromanheader{6}{1. ติก}
\proseitem{640}{สิยา กุสลา, สิยา อพฺยากตา.\csromanspacer{} ตีณิ ฌานานิ เอตฺถุปฺปนฺนํ สุขํ เวทนํ ฐเปตฺวา สุขาย เวทนาย สมฺปยุตฺตา.\csromanspacer{} จตุตฺถํ ฌานํ เอตฺถุปฺปนฺนํ อทุกฺขมสุขํ เวทนํ ฐเปตฺวา อทุกฺขมสุขาย เวทนาย สมฺปยุตฺตํ, สิยา วิปากา, สิยา วิปากธมฺมธมฺมา, สิยา เนว วิปาก น วิปากธมฺมธมฺมา, สิยา อุปาทินฺนุปาทานิยา, สิยา อนุปาทินฺนุปาทานิยา, สิยา อนุปาทินฺน-อนุปาทานิยา, สิยา อสํกิลิฏฺฐสํกิเลสิกา, สิยา อสํกิลิฏฺฐ-อสํกิเลสิกา.\csromanspacer{} ปฐมํ ฌานํ เอตฺถุปฺปนฺเน วิตกฺกวิจาเร ฐเปตฺวา สวิตกฺกํ สวิจารํ.\csromanspacer{} ตีณิ ฌานานิ อวิตกฺก-อวิจารา.\csromanspacer{} เทฺว ฌานานิ เอตฺถุปฺปนฺนํ ปีติํ ฐเปตฺวา ปีติสหคตา.\csromanspacer{} ตีณิ ฌานานิ เอตฺถุปฺปนฺนํ สุขํ ฐเปตฺวา สุขสหคตา.\csromanspacer{} จตุตฺถํ ฌานํ เอตฺถุปฺปนฺนํ อุเปกฺขํ ฐเปตฺวา อุเปกฺขาสหคตํ, เนว ทสฺสเนน น ภาวนาย ปหาตพฺพา, เนว ทสฺสเนน น ภาวนาย ปหาตพฺพา, เนว ทสฺสเนน น ภาวนาย ปหาตพฺพเหตุกา, สิยา อาจยคามิโน, สิยา อปจยคามิโน, สิยา เนวาจยคามินาปจยคามิโน, สิยา เสกฺขา, สิยา อเสกฺขา, สิยา เนวสกฺขนาเสกฺขา, สิยา มหคฺคตา, สิยา อปฺปมาณา.\csromanspacer{} ตีณิ ฌานานิ น วตฺตพฺพา “ปริตฺตารมฺมณา”ติปิ, “มหคฺคตารมฺมณา”ติปิ, สิยา \csromanfolio{282}อปฺปมาณารมฺมณา, สิยา น วตฺตพฺพา “อปฺปมาณารมฺมณา”ติ.\csromanspacer{} จตุตฺถํ ฌานํ สิยา ปริตฺตารมฺมณํ, สิยา มหคฺคตารมฺมณํ, สิยา อปฺปมาณารมฺมณํ, สิยา น วตฺตพฺพํ “ปริตฺตารมฺมณนฺ”ติปิ, “มหคฺคตารมฺมณนฺ”ติปิ, “อปฺปมาณารมฺมณนฺ”ติปิ, สิยา มชฺฌิมา, สิยา ปณีตา, สิยา สมฺมตฺตนิยตา, สิยา อนิยตา.\csromanspacer{} ตีณิ ฌานานิ น มคฺคารมฺมณา, สิยา มคฺคเหตุกา, สิยา มคฺคาธิปติโน, สิยา น วตฺตพฺพา “มคฺคเหตุกา”ติปิ, “มคฺคาธิปติโน”ติปิ.\csromanspacer{} จตุตฺถํ ฌานํ สิยา มคฺคารมฺมณํ, สิยา มคฺคเหตุกํ, สิยา มคฺคาธิปติ, สิยา น วตฺตพฺพํ “มคฺคารมฺมณนฺ”ติปิ, “มคฺคเหตุกนฺ”ติปิ, “มคฺคาธิปตี”ติปิ, สิยา อุปฺปนฺนา, สิยา อนุปฺปนฺนา, สิยา อุปฺปาทิโน, สิยา อตีตา, สิยา อนาคตา, สิยา ปจฺจุปฺปนฺนา, ตีณิ ฌานานิ น วตฺตพฺพา “อตีตารมฺมณา”ติปิ, “อนาคตารมฺมณา”ติปิ, “ปจฺจุปฺปนฺนารมฺมณา”ติปิ.\csromanspacer{} จตุตฺถํ ฌานํ สิยา อตีตารมฺมณํ, สิยา อนาคตารมฺมณํ, สิยา ปจฺจุปฺปนฺนารมฺมณํ, สิยา น วตฺตพฺพํ “อตีตารมฺมณนฺ”ติปิ, “อนาคตารมฺมณนฺ”ติปิ, “ปจฺจุปฺปนฺนารมฺมณนฺ”ติปิ, สิยา อชฺฌตฺตา, สิยา พหิทฺธา, สิยา อชฺฌตฺตพหิทฺธา.\csromanspacer{} ตีณิ ฌานานิ พหิทฺธารมฺมณา.\csromanspacer{} จตุตฺถํ ฌานํ สิยา อชฺฌตฺตารมฺมณํ, สิยา พหิทฺธารมฺมณํ, สิยา อชฺฌตฺตพหิทฺธารมฺมณํ, สิยา น วตฺตพฺพํ “อชฺฌตฺตารมฺมณนฺ”ติปิ, “พหิทฺธารมฺมณนฺ”ติปิ, “อชฺฌตฺตพหิทฺธารมฺมณนฺ”ติปิ, อนิทสฺสน- อปฺปฏิฆา.}

\csromanmatikaanchor{mtk.115}
\csromantocmarkref{subparagraph}{2. ทุก}{mtk.115}
\bookmark[dest={mtk.115},level=5]{2. ทุก}
\csromanheader{6}{2. ทุก}
\proseitem{641}{น เหตู, สเหตุกา, เหตุสมฺปยุตฺตา, น วตฺตพฺพา “เหตู เจว สเหตุกา จา”ติ, สเหตุกา เจว น จ เหตู, น วตฺตพฺพา “เหตู เจว เหตุสมฺปยุตฺตา จา”ติ, เหตุสมฺปยุตฺตา เจว น จ เหตู, น เหตู, สเหตุกา.}

\prose{สปฺปจฺจยา, สงฺขตา, อนิทสฺสนา, อปฺปฏิฆา, อรูปา, สิยา โลกิยา, สิยา โลกุตฺตรา, เกนจิ วิญฺเญยฺยา, เกนจิ น วิญฺเญยฺยา.}

\prose{โน อาสวา, สิยา สาสวา, สิยา อนาสวา, อาสววิปฺปยุตฺตา, น วตฺตพฺพา “อาสวา เจว สาสวา จา”ติ, สิยา สาสวา เจว โน จ อาสวา, สิยา น วตฺตพฺพา “สาสวา เจว โน จ อาสวา”ติ, น วตฺตพฺพา “อาสวา เจว อาสวสมฺปยุตฺตา จา”ติปิ, “อาสวสมฺปยุตฺตา}

\vspace{\tipitakaparskip}
\csromancenter{ฌานวิภงฺค เจว โน จ อาสวา”ติปิ, สิยา อาสววิปฺปยุตฺตา สาสวา, สิยา อาสววิปฺปยุตฺตา อนาสวา.}
\prose{\csromanfolio{283}โน สํโยชนา ฯเปฯ.\csromanspacer{} โน คนฺถา ฯเปฯ.\csromanspacer{} โน โอฆา ฯเปฯ.\csromanspacer{} โน โยคา ฯเปฯ.\csromanspacer{} โน นีวรณา ฯเปฯ.\csromanspacer{} โน ปรามาสา ฯเปฯ สารมฺมณา, โน จิตฺตา, เจตสิกา, จิตฺตสมฺปยุตฺตา, จิตฺตสํสฏฺฐา, จิตฺตสมุฏฺฐานา, จิตฺตสหภุโน, จิตฺตานุปริวตฺติโน, จิตฺตสํสฏฺฐสมุฏฺฐานา, จิตฺตสํสฏฺฐสมุฏฺฐานสหภุโน, จิตฺตสํสฏฺฐสมุฏฺฐานานุปริวตฺติโน, พาหิรา, โน อุปาทา, สิยา อุปาทินฺนา, สิยา อนุปาทินฺนา.}

\prose{โน อุปาทานา ฯเปฯ.\csromanspacer{} โน กิเลสา ฯเปฯ.\csromanspacer{} น ทสฺสเนน ปหาตพฺพา, น ภาวนาย ปหาตพฺพา, น ทสฺสเนน ปหาตพฺพเหตุกา, น ภาวนาย ปหาตพฺพเหตุกา.\csromanspacer{} ปฐมํ ฌานํ เอตฺถุปฺปนฺนํ วิตกฺกํ ฐเปตฺวา สวิตกฺกํ.\csromanspacer{} ตีณิ ฌานานิ อวิตกฺกา.\csromanspacer{} ปฐมํ ฌานํ เอตฺถุปฺปนฺนํ วิจารํ ฐเปตฺวา สวิจารํ.\csromanspacer{} ตีณิ ฌานานิ อวิจารา.\csromanspacer{} เทฺว ฌานานิ เอตฺถุปฺปนฺนํ ปีติํ ฐเปตฺวา สปฺปีติกา.\csromanspacer{} เทฺว ฌานานิ อปฺปีติกา.\csromanspacer{} เทฺว ฌานานิ เอตฺถุปฺปนฺนํ ปีติํ ฐเปตฺวา ปีติสหคตา.\csromanspacer{} เทฺว ฌานานิ น ปีติสหคตา.\csromanspacer{} ตีณิ ฌานานิ เอตฺถุปฺปนฺนํ สุขํ ฐเปตฺวา สุขสหคตา.\csromanspacer{} จตุตฺถํ ฌานํ น สุขสหคตํ.\csromanspacer{} จตุตฺถํ ฌานํ เอตฺถุปฺปนฺนํ อุเปกฺขํ ฐเปตฺวา อุเปกฺขาสหคตํ.\csromanspacer{} ตีณิ ฌานานิ อุเปกฺขาสหคตา, น กามาวจรา, สิยา รูปาวจรา, สิยา น รูปาวจรา.\csromanspacer{} ตีณิ ฌานานิ น อรูปาวจรา.\csromanspacer{} จตุตฺถํ ฌานํ สิยา อรูปาวจรํ, สิยา น อรูปาวจรํ, สิยา ปริยาปนฺนา, สิยา อปริยาปนฺนา, สิยา นิยฺยานิกา, สิยา อนิยฺยานิกา, สิยา นิยตา, สิยา อนิยตา, สิยา ส-อุตฺตรา, สิยา อนุตฺตรา, อรณาติ.}

\vspace{\tipitakaparskip}
\csromancenter{ปญฺหาปุจฺฉกํ.}
\nitthitam{\textbf{ฌานวิภงฺโค นิฏฺฐิโต.}}
\csromanmatikaanchor{mtk.116}
\csromantocmarknopage{subparagraph}{13. อปฺปมญฺญาวิภงฺค}{mtk.116}
\bookmark[dest={mtk.116},level=5]{13. อปฺปมญฺญาวิภงฺค}
\csromanheader{6}{13. อปฺปมญฺญาวิภงฺค}
\csromanheader{2}{1. สุตฺตนฺตภาชนีย}
\proseitem{642}{\csromanfolio{284}จตสฺโส อปฺปมญฺญาโย\csromandash{}อิธ ภิกฺขุ เมตฺตาสหคเตน เจตสา เอกํ ทิสํ ผริตฺวา วิหรติ, ตถา ทุติยํ.\csromanspacer{} ตถา ตติยํ.\csromanspacer{} ตถา จตุตฺถํ\footnote{จตุตฺถิํ (สี)}.\csromanspacer{} อิติ อุทฺธมโธ ติริยํ สพฺพธิ สพฺพตฺตตาย สพฺพาวนฺตํ โลกํ เมตฺตาสหคเตน เจตสา วิปุเลน มหคฺคเตน อปฺปมาเณน อเวเรน อพฺยาปชฺเชน\footnote{อพฺยาปชฺเฌน (สี, สฺยา)} ผริตฺวา วิหรติ.\csromanspacer{} กรุณาสหคเตน เจตสา เอกํ ทิสํ ผริตฺวา วิหรติ, ตถา ทุติยํ.\csromanspacer{} ตถา ตติยํ.\csromanspacer{} ตถา จตุตฺถํ.\csromanspacer{} อิติ อุทฺธมโธ ติริยํ สพฺพธิ สพฺพตฺตตาย สพฺพาวนฺตํ โลกํ กรุณาสหคเตน เจตสา วิปุเลน มหคฺคเตน อปฺปมาเณน อเวเรน อพฺยาปชฺเชน ผริตฺวา วิหรติ.\csromanspacer{} มุทิตาสหคเตน เจตสา เอกํ ทิสํ ผริตฺวา วิหรติ, ตถา ทุติยํ.\csromanspacer{} ตถา ตติยํ.\csromanspacer{} ตถา จตุตฺถํ.\csromanspacer{} อิติ อุทฺธมโธ ติริยํ สพฺพธิ สพฺพตฺตตาย สพฺพาวนฺตํ โลกํ มุทิตาสหคเตน เจตสา วิปุเลน มหคฺคเตน อปฺปมาเณน อเวเรน อพฺยาปชฺเชน ผริตฺวา วิหรติ.\csromanspacer{} อุเปกฺขาสหคเตน เจตสา เอกํ ทิสํ ผริตฺวา วิหรติ, ตถา ทุติยํ.\csromanspacer{} ตถา ตติยํ.\csromanspacer{} ตถา จตุตฺถํ.\csromanspacer{} อิติ อุทฺธมโธ ติริยํ สพฺพธิ สพฺพตฺตตาย สพฺพาวนฺตํ โลกํ อุเปกฺขาสหคเตน เจตสา วิปุเลน มหคฺคเตน อปฺปมาเณน อเวเรน อพฺยาปชฺเชน ผริตฺวา วิหรติ.}

\csromanmatikaanchor{mtk.117}
\csromanmatikaanchor{mtk.118}
\csromantocmarknopage{subparagraph}{1. สุตฺตนฺตภาชนีย}{mtk.117}
\bookmark[dest={mtk.117},level=5]{1. สุตฺตนฺตภาชนีย}
\csromantocmarkref{subparagraph}{1. เมตฺตา}{mtk.118}
\bookmark[dest={mtk.118},level=5]{1. เมตฺตา}
\csromanheader{6}{1. เมตฺตา}
\proseitem{643}{กถญฺจ ภิกฺขุ เมตฺตาสหคเตน เจตสา เอกํ ทิสํ ผริตฺวา วิหรติ, เสยฺยถาปิ นาม เอกํ ปุคฺคลํ ปิยํ มนาปํ ทิสฺวา เมตฺตาเยยฺย, เอวเมว สพฺเพ สตฺเต เมตฺตาย ผรติ.}

\prose{ตตฺถ กตมา เมตฺตา, ยา สตฺเตสุ เมตฺติ เมตฺตายนา เมตฺตายิตตฺตํ เมตฺตาเจโตวิมุตฺติ, อยํ วุจฺจติ เมตฺตา.}

\prose{ตตฺถ กตมํ จิตฺตํ, ยํ จิตฺตํ มโน มานสํ หทยํ ปณฺฑรํ มโน มนายตนํ มนินฺทฺริยํ วิญฺญาณํ วิญฺญาณกฺขนฺโธ ตชฺชามโนวิญฺญาณธาตุ,}

\vspace{\tipitakaparskip}
\csromancenter{อปฺปมญฺญาวิภงฺค อิทํ วุจฺจติ จิตฺตํ.\csromanspacer{} อิทํ จิตฺตํ อิมาย เมตฺตาย สหคตํ โหติ สหชาตํ สํสฏฺฐํ สมฺปยุตฺตํ, เตน วุจฺจติ “เมตฺตาสหคเตน เจตสา”ติ.}
\proseitem{644}{\csromanfolio{285}“เอกํ ทิสนฺ”ติ ปุรตฺถิมํ วา ทิสํ ปจฺฉิมํ วา ทิสํ อุตฺตรํ วา ทิสํ ทกฺขิณํ วา ทิสํ อุทฺธํ วา อโธ วา ติริยํ วา วิทิสํ วา.}

\proseitem{645}{“ผริตฺวา”ติ ผริตฺวา อธิมุจฺจิตฺวา.}

\proseitem{646}{“วิหรตี”ติ อิริยติ วตฺตติ ปาเลติ ยเปติ ยาเปติ จรติ วิหรติ, เตน วุจฺจติ “วิหรตี”ติ.}

\proseitem{647}{“ตถา ทุติยนฺ”ติ ยเถว เอกํ ทิสํ, ตถา ทุติยํ ทิสํ, ตถา ตติยํ ทิสํ, ตถา จตุตฺถํ ทิสํ, ตถา อุทฺธํ, ตถา อโธ, ตถา ติริยํ, ตถา วิทิสํ.}

\proseitem{648}{“สพฺพธิ สพฺพตฺตตาย สพฺพาวนฺตํ โลกนฺ”ติ สพฺเพน สพฺพํ สพฺพถา สพฺพํ อเสสํ นิสฺเสสํ ปริยาทายวจนเมตํ “สพฺพธิ สพฺพตฺตตาย สพฺพาวนฺตํ โลกนฺ”ติ.}

\proseitem{649}{“เมตฺตาสหคเตน เจตสา”ติ ตตฺถ กตมา เมตฺตา, ยา สตฺเตสุ เมตฺติ เมตฺตายนา เมตฺตายิตตฺตํ เมตฺตาเจโตวิมุตฺติ, อยํ วุจฺจติ เมตฺตา.}

\prose{ตตฺถ กตมํ จิตฺตํ, ยํ จิตฺตํ มโน มานสํ ฯเปฯ ตชฺชามโนวิญฺญาณธาตุ, อิทํ วุจฺจติ จิตฺตํ อิทํ จิตฺตํ อิมาย เมตฺตาย สหคตํ โหติ สหชาตํ สํสฏฺฐํ สมฺปยุตฺตํ, เตน วุจฺจติ “เมตฺตาสหคเตน เจตสา”ติ.}

\proseitem{650}{“วิปุเลนา”ติ ยํ วิปุลํ ตํ มหคฺคตํ, ยํ มหคฺคตํ ตํ อปฺปมาณํ, ยํ อปฺปมาณํ โส อเวโร, โย อเวโร, โส อพฺยาปชฺโช\footnote{อพฺยาปชฺโฌ (สี, สฺยา)}.}

\proseitem{651}{“ผริตฺวา”ติ ผริตฺวา อธิมุจฺจิตฺวา.}

\proseitem{652}{“วิหรตี”ติ ฯเปฯ เตน วุจฺจติ “วิหรตี”ติ.}

\csromanmatikaanchor{mtk.119}
\csromantocmarkref{subparagraph}{2. กรุณา}{mtk.119}
\bookmark[dest={mtk.119},level=5]{2. กรุณา}
\csromanheader{6}{2. กรุณา}
\proseitem{653}{\csromanfolio{286}กถญฺจ ภิกฺขุ กรุณาสหคเตน เจตสา เอกํ ทิสํ ผริตฺวา วิหรติ, เสยฺยถาปิ นาม เอกํ ปุคฺคลํ ทุคฺคตํ ทุรูเปตํ ทิสฺวา กรุณาเยยฺย, เอวเมว สพฺเพ สตฺเต กรุณาย ผรติ.}

\prose{ตตฺถ กตมา กรุณา, ยา สตฺเตสุ กรุณา กรุณายนา กรุณายิตตฺตํ กรุณาเจโตวิมุตฺติ, อยํ วุจฺจติ กรุณา.}

\prose{ตตฺถ กตมํ จิตฺตํ, ยํ จิตฺตํ มโน มานสํ ฯเปฯ ตชฺชามโนวิญฺญาณธาตุ, อิทํ วุจฺจติ จิตฺตํ.\csromanspacer{} อิทํ จิตฺตํ อิมายา กรุณาย สหคตํ โหติ สหชาตํ สํสฏฺฐํ สมฺปยุตฺตํ, เตน วุจฺจติ “กรุณาสหคเตน เจตสา”ติ.}

\proseitem{654}{“เอกํ ทิสนฺ”ติ ปุรตฺถิมํ วา ทิสํ ปจฺฉิมํ วา ทิสํ อุตฺตรํ วา ทิสํ ทกฺขิณํ วา ทิสํ อุทฺธํ วา อโธ วา ติริยํ วา วิทิสํ วา.}

\proseitem{655}{“ผริตฺวา”ติ ผริตฺวา อธิมุจฺจิตฺวา.}

\proseitem{656}{“วิหรตี”ติ อิริยติ วตฺตติ ปาเลติ ยเปติ ยาเปติ จรติ วิหรติ, เตน วุจฺจติ “วิหรตี”ติ.}

\proseitem{657}{“ตถา ทุติยนฺ”ติ ยเถว เอกํ ทิสํ, ตถา ทุติยํ ทิสํ, ตถา ตติยํ ทิสํ, ตถา จตุตฺถํ ทิสํ, ตถา อุทฺธํ, ตถา อโธ, ตถา ติริยํ, ตถา วิทิสํ.}

\proseitem{658}{“สพฺพธิ สพฺพตฺตตาย สพฺพาวนฺตํ โลกนฺ”ติ สพฺเพน สพฺพํ สพฺพถา สพฺพํ อเสสํ นิสฺเสสํ ปริยาทายวจนเมตํ “สพฺพธิ สพฺพตฺตตาย สพฺพาวนฺตํ โลกนฺ”ติ.}

\vspace{\tipitakaparskip}
\csromancenter{“กรุณาสหคเตน เจตสา”ติ ตตฺถ กตมา กรุณา, ยา สตฺเตสุ กรุณา กรุณายนา กรุณายิตตฺตํ กรุณาเจโตวิมุตฺติ, อยํ วุจฺจติ กรุณา.}
\prose{ตตฺถ กตมํ จิตฺตํ, ยํ จิตฺตํ มโน มานสํ ฯเปฯ ตชฺชามโนวิญฺญาณธาตุ, อิทํ วุจฺจติ จิตฺตํ.\csromanspacer{} อิทํ จิตฺตํ อิมาย กรุณาย สหคตํ โหติ}

\vspace{\tipitakaparskip}
\csromancenter{อปฺปมญฺญาวิภงฺค สหชาตํ สํสฏฺฐํ สมฺปยุตฺตํ, เตน วุจฺจติ “กรุณาสหคเตน เจตสา”ติ.}
\proseitem{660}{\csromanfolio{287}“วิปุเลนา”ติ ยํ วิปุลํ ตํ มหคฺคตํ, ยํ มหคฺคตํ ตํ อปฺปมาณํ, ยํ อปฺปมาณํ โส อเวโร, โย อเวโร โส อพฺยาปชฺโช.}

\proseitem{661}{“ผริตฺวา”ติ ผริตฺวา อธิมุจฺจิตฺวา.}

\proseitem{662}{“วิหรตี”ติ ฯเปฯ เตน วุจฺจติ “วิหรตี”ติ.}

\csromanmatikaanchor{mtk.120}
\csromantocmarkref{subparagraph}{3. มุทิตา}{mtk.120}
\bookmark[dest={mtk.120},level=5]{3. มุทิตา}
\csromanheader{6}{3. มุทิตา}
\proseitem{663}{กถญฺจ ภิกฺขุ มุทิตาสหคเตน เจตสา เอกํ ทิสํ ผริตฺวา วิหรติ, เสยฺยถาปิ นาม เอกํ ปุคฺคลํ ปิยํ มนาปํ ทิสฺวา มุทิโต อสฺส, เอวเมว สพฺเพ สตฺเต มุทิตาย ผรติ.}

\prose{ตตฺถ กตมา มุทิตา, ยา สตฺเตสุ มุทิตา มุทิตายนา มุทิตายิตตฺตํ มุทิตาเจโตวิมุตฺติ, อยํ วุจฺจติ มุทิตา.}

\prose{ตตฺถ กตมํ จิตฺตํ, ยํ จิตฺตํ มโน มานสํ ฯเปฯ ตชฺชามโนวิญฺญาณธาตุ, อิทํ วุจฺจติ จิตฺตํ.\csromanspacer{} อิทํ จิตฺตํ อิมาย มุทิตาย สหคตํ โหติ สหชาตํ สํสฏฺฐํ สมฺปยุตฺตํ, เตน วุจฺจติ “มุทิตาสหคเตน เจตสา”ติ.}

\proseitem{664}{“เอกํ ทิสนฺ”ติ ปุรตฺถิมํ วา ทิสํ ปจฺฉิมํ วา ทิสํ อุตฺตรํ วา ทิสํ ทกฺขิณํ วา ทิสํ อุทฺธํ วา อโธ วา ติริยํ วา วิทิสํ วา.}

\proseitem{665}{“ผริตฺวา”ติ ผริตฺวา อธิมุจฺจิตฺวา.}

\proseitem{666}{“วิหรตี”ติ ฯเปฯ เตน วุจฺจติ “วิหรตี”ติ.}

\proseitem{667}{“ตถา ทุติยนฺ”ติ ยเถว เอกํ ทิสํ, ตถา ทุติยํ ทิสํ, ตถา ตติยํ ทิสํ, ตถา จตุตฺถํ ทิสํ, ตถา อุทฺธํ, ตถา อโธ, ตถา ติริยํ, ตถา วิทิสํ.}

\proseitem{668}{\csromanfolio{288}“สพฺพธิ สพฺพตฺตตาย สพฺพาวนฺตํ โลกนฺ”ติ สพฺเพน สพฺพํ สพฺพถา สพฺพํ อเสสํ นิสฺเสสํ ปริยาทายวจนเมตํ “สพฺพธิ สพฺพตฺตตาย สพฺพาวนฺตํ โลกนฺ”ติ.}

\proseitem{669}{“มุทิตาสหคเตน เจตสา”ติ ตตฺถ กตมา มุทิตา, ยา สตฺเตสุ มุทิตา มุทิตายนา มุทิตายิตตฺตํ มุทิตาเจโตวิมุตฺติ, อยํ วุจฺจติ มุทิตา.}

\prose{ตตฺถ กตมํ จิตฺตํ, ยํ จิตฺตํ มโน มานสํ ฯเปฯ ตชฺชามโนวิญฺญาณธาตุ, อิทํ วุจฺจติ จิตฺตํ.\csromanspacer{} อิทํ จิตฺตํ อิมาย มุทิตาย สหคตํ โหติ สหชาตํ สํสฏฺฐํ สมฺปยุตฺตํ, เตน วุจฺจติ “มุทิตาสหคเตน เจตสา”ติ.}

\proseitem{670}{“วิปุเลนา”ติ ยํ วิปุลํ ตํ มหคฺคตํ, ยํ มหคฺคตํ ตํ อปฺปมาณํ, ยํ อปฺปมาณํ โส อเวโร, โย อเวโร โส อพฺยาปชฺโช.}

\proseitem{671}{“ผริตฺวา”ติ ผริตฺวา อธิมุจฺจิตฺวา.}

\proseitem{672}{“วิหรตี”ติ ฯเปฯ เตน วุจฺจติ “วิหรตี”ติ.}

\csromanmatikaanchor{mtk.121}
\csromantocmarkref{subparagraph}{4. อุเปกฺขา}{mtk.121}
\bookmark[dest={mtk.121},level=5]{4. อุเปกฺขา}
\csromanheader{6}{4. อุเปกฺขา}
\proseitem{673}{กถญฺจ ภิกฺขุ อุเปกฺขาสหคเตน เจตสา เอกํ ทิสํ ผริตฺวา วิหรติ, เสยฺยถาปิ นาม เอกํ ปุคฺคลํ เนว มนาปํ น อมนาปํ ทิสฺวา อุเปกฺขโก อสฺส, เอวเมว สพฺเพ สตฺเต อุเปกฺขาย ผรติ.}

\prose{ตตฺถ กตมา อุเปกฺขา, ยา สตฺเตสุ อุเปกฺขา อุเปกฺขายนา อุเปกฺขายิตตฺตํ อุเปกฺขาเจโตวิมุตฺติ, อยํ วุจฺจติ อุเปกฺขา.}

\prose{ตตฺถ กตมํ จิตฺตํ, ยํ จิตฺตํ มโน มานสํ ฯเปฯ ตชฺชามโนวิญฺญาณธาตุ, อิทํ วุจฺจติ จิตฺตํ.\csromanspacer{} อิทํ จิตฺตํ อิมาย อุเปกฺขาย สหคตํ โหติ สหชาตํ สํสฏฺฐํ สมฺปยุตฺตํ, เตน วุจฺจติ “อุเปกฺขาสหคเตน เจตสา”ติ.}

\proseitem{674}{“เอกํ ทิสนฺ”ติ ปุรตฺถิมํ วา ทิสํ ปจฺฉิมํ วา ทิสํ อุตฺตรํ วา ทิสํ ทกฺขิณํ วา ทิสํ อุทฺธํ วา อโธ วา ติริยํ วา วิทิสํ วา.}

\csromanheader{2}{13. อปฺปมญฺญาวิภงฺค}
\proseitem{675}{\csromanfolio{289}“ผริตฺวา”ติ ผริตฺวา อธิมุจฺจิตฺวา.}

\proseitem{676}{“วิหรตี”ติ ฯเปฯ เตน วุจฺจติ “วิหรตี”ติ.}

\proseitem{677}{“ตถา ทุติยนฺ”ติ ยเถว เอกํ ทิสํ, ตถา ทุติยํ ทิสํ, ตถา ตติยํ ทิสํ, ตถา จตุตฺถํ ทิสํ, ตถา อุทฺธํ, ตถา อโธ, ตถา ติริยํ, ตถา วิทิสํ.}

\proseitem{678}{“สพฺพธิ สพฺพตฺตตาย สพฺพาวนฺตํ โลกนฺ”ติ สพฺเพน สพฺพํ สพฺพถา สพฺพํ อเสสํ นิสฺเสสํ ปริยาทายวจนเมตํ “สพฺพธิ สพฺพตฺตตาย สพฺพาวนฺตํ โลกนฺ”ติ.}

\proseitem{679}{“อุเปกฺขาสหคเตน เจตสา”ติ ตตฺถ กตมา อุเปกฺขา, ยา สตฺเตสุ อุเปกฺขา อุเปกฺขายนา อุเปกฺขายิตตฺตํ อุเปกฺขาเจโตวิมุตฺติ, อยํ วุจฺจติ อุเปกฺขา.}

\prose{ตตฺถ กตมํ จิตฺตํ, ยํ จิตฺตํ มโน มานสํ ฯเปฯ ตชฺชามโนวิญฺญาณธาตุ, อิทํ วุจฺจติ จิตฺตํ.\csromanspacer{} อิทํ จิตฺตํ อิมาย อุเปกฺขาย สหคตํ โหติ สหชาตํ สํสฏฺฐํ สมฺปยุตฺตํ, เตน วุจฺจติ “อุเปกฺขาสหคเตน เจตสา”ติ.}

\proseitem{680}{“วิปุเลนา”ติ ยํ วิปุลํ ตํ มหคฺคตํ, ยํ มหคฺคตํ ตํ อปฺปมาณํ, ยํ อปฺปมาณํ โส อเวโร, โย อเวโร โส อพฺยาปชฺโช.}

\proseitem{681}{“ผริตฺวา”ติ ผริตฺวา อธิมุจฺจิตฺวา.}

\proseitem{682}{“วิหรตี”ติ ฯเปฯ เตน วุจฺจติ “วิหรตี”ติ.}

\csromanheader{2}{สุตฺตนฺตภาชนียํ.}
\csromansectionrule
\csromanmatikaanchor{mtk.122}
\csromantocmarkref{subparagraph}{2. อภิธมฺมภาชนีย}{mtk.122}
\bookmark[dest={mtk.122},level=5]{2. อภิธมฺมภาชนีย}
\csromanheader{6}{2. อภิธมฺมภาชนีย}
\proseitem{683}{จตสฺโส อปฺปมญฺญาโย เมตฺตา กรุณา มุทิตา อุเปกฺขา.}

\proseitem{684}{ตตฺถ กตมา เมตฺตา, อิธ ภิกฺขุ ยสฺมิํ สมเย รูปูปปตฺติยา มคฺคํ ภาเวติ วิวิจฺเจว กาเมหิ ฯเปฯ ปฐมํ ฌานํ อุปสมฺปชฺช วิหรติ \csromanfolio{290}เมตฺตาสหคตํ.\csromanspacer{} ยา ตสฺมิํ สมเย เมตฺติ เมตฺตายนา เมตฺตายิตตฺตํ เมตฺตาเจโตวิมุตฺติ, อยํ วุจฺจติ เมตฺตา.\csromanspacer{} อวเสสา ธมฺมา เมตฺตาย สมฺปยุตฺตา. (1)}

\prose{ตตฺถ กตมา เมตฺตา, อิธ ภิกฺขุ ยสฺมิํ สมเย รูปูปปตฺติยา มคฺคํ ภาเวติ วิตกฺกวิจารานํ วูปสมา ฯเปฯ ทุติยํ ฌานํ อุปสมฺปชฺช วิหรติ เมตฺตาสหคตํ.\csromanspacer{} ยา ตสฺมิํ สมเย เมตฺติ เมตฺตายนา เมตฺตายิตตฺตํ เมตฺตาเจโตวิมุตฺติ, อยํ วุจฺจติ เมตฺตา.\csromanspacer{} อวเสสา ธมฺมา เมตฺตาย สมฺปยุตฺตา. (2)}

\prose{ตตฺถ กตมา เมตฺตา, อิธ ภิกฺขุ ยสฺมิํ สมเย รูปูปปตฺติยา มคฺคํ ภาเวติ ปีติยา จ วิราคา ฯเปฯ ตติยํ ฌานํ อุปสมฺปชฺช วิหรติ เมตฺตาสหคตํ, ยา ตสฺมิํ สมเย เมตฺติ เมตฺตายนา เมตฺตายิตตฺตํ เมตฺตาเจโตวิมุตฺติ, อยํ วุจฺจติ เมตฺตา.\csromanspacer{} อวเสสา ธมฺมา เมตฺตาย สมฺปยุตฺตา. (3)}

\proseitem{685}{อิธ ภิกฺขุ ยสฺมิํ สมเย รูปูปปตฺติยา มคฺคํ ภาเวติ วิวิจฺเจว กาเมหิ ฯเปฯ ปฐมํ ฌานํ อุปสมฺปชฺช วิหรติ เมตฺตาสหคตํ, ยา ตสฺมิํ สมเย เมตฺติ เมตฺตายนา เมตฺตายิตตฺตํ เมตฺตาเจโตวิมุตฺติ, อยํ วุจฺจติ เมตฺตา.\csromanspacer{} อวเสสา ธมฺมา เมตฺตาย สมฺปยุตฺตา. (1-4)}

\prose{อิธ ภิกฺขุ ยสฺมิํ สมเย รูปูปปตฺติยา มคฺคํ ภาเวติ วิวิจฺเจว กาเมหิ วิวิจฺจ อกุสเลหิ ธมฺเมหิ อวิตกฺกํ วิจารมตฺตํ วิเวกชํ ปีติสุขํ ทุติยํ ฌานํ อุปสมฺปชฺช วิหรติ เมตฺตาสหคตํ, ยา ตสฺมิํ สมเย เมตฺติ เมตฺตายนา เมตฺตายิตตฺตํ เมตฺตาเจโตวิมุตฺติ, อยํ วุจฺจติ เมตฺตา.\csromanspacer{} อวเสสา ธมฺมา เมตฺตาย สมฺปยุตฺตา. (2-5)}

\prose{อิธ ภิกฺขุ ยสฺมิํ สมเย รูปูปปตฺติยา มคฺคํ ภาเวติ วิตกฺกวิจารานํ วูปสมา ฯเปฯ ตติยํ ฌานํ อุปสมฺปชฺช วิหรติ เมตฺตาสหคตํ, ยา ตสฺมิํ สมเย เมตฺติ เมตฺตายนา เมตฺตายิตตฺตํ เมตฺตาเจโตวิมุตฺติ, อยํ วุจฺจติ เมตฺตา.\csromanspacer{} อวเสสา ธมฺมา เมตฺตาย สมฺปยุตฺตา. (3-6)}

\prose{อิธ ภิกฺขุ ยสฺมิํ สมเย รูปูปปตฺติยา มคฺคํ ภาเวติ ปีติยา จ วิราคา ฯเปฯ จตุตฺถํ ฌานํ อุปสมฺปชฺช วิหรติ เมตฺตาสหคตํ, ยา ตสฺมิํ}

\vspace{\tipitakaparskip}
\csromancenter{อปฺปมญฺญาวิภงฺค สมเย เมตฺติ เมตฺตายนา เมตฺตายิตตฺตํ เมตฺตาเจโตวิมุตฺติ, อยํ วุจฺจติ เมตฺตา.\csromanspacer{} อวเสสา ธมฺมา เมตฺตาย สมฺปยุตฺตา. (4-7)}
\proseitem{686}{\csromanfolio{291}ตตฺถ กตมา กรุณา, อิธ ภิกฺขุ ยสฺมิํ สมเย รูปูปปตฺติยา มคฺคํ ภาเวติ วิวิจฺเจว กาเมหิ ฯเปฯ ปฐมํ ฌานํ อุปสมฺปชฺช วิหรติ กรุณาสหคตํ, ยา ตสฺมิํ สมเย กรุณา กรุณายนา กรุณายิตตฺตํ กรุณาเจโตวิมุตฺติ, อยํ วุจฺจติ กรุณา.\csromanspacer{} อวเสสา ธมฺมา กรุณาย สมฺปยุตฺตา. (1)}

\prose{ตตฺถ กตมา กรุณา, อิธ ภิกฺขุ ยสฺมิํ สมเย รูปูปปตฺติยา มคฺคํ ภาเวติ ปีติยา จ วิราคา ฯเปฯ ทุติยํ ฌานํ อุปสมฺปชฺช วิหรติ กรุณาสหคตํ, ยา ตสฺมิํ สมเย กรุณา กรุณายนา กรุณายิตตฺตํ กรุณาเจโตวิมุตฺติ, อยํ วุจฺจติ กรุณา.\csromanspacer{} อวเสสา ธมฺมา กรุณาย สมฺปยุตฺตา. (2)}

\prose{ตตฺถ กตมา กรุณา, อิธ ภิกฺขุ ยสฺมิํ สมเย รูปูปปตฺติยา มคฺคํ ภาเวติ ปีติยา จ วิราคา ฯเปฯ ตติยํ ฌานํ อุปสมฺปชฺช วิหรติ กรุณาสหคตํ, ยา ตสฺมิํ สมเย กรุณา กรุณายนา กรุณายิตตฺตํ กรุณาเจโตวิมุตฺติ, อยํ วุจฺจติ กรุณา.\csromanspacer{} อวเสสา ธมฺมา กรุณาย สมฺปยุตฺตา. (3)}

\proseitem{687}{อิธ ภิกฺขุ ยสฺมิํ สมเย รูปูปปตฺติยา มคฺคํ ภาเวติ วิวิจฺเจว กาเมหิ ฯเปฯ ปฐมํ ฌานํ อุปสมฺปชฺช วิหรติ กรุณาสหคตํ, ยา ตสฺมิํ สมเย กรุณา กรุณายนา กรุณายิตตฺตํ กรุณาเจโตวิมุตฺติ, อยํ วุจฺจติ กรุณา.\csromanspacer{} อวเสสา ธมฺมา กรุณาย สมฺปยุตฺตา. (1-4)}

\prose{อิธ ภิกฺขุ ยสฺมิํ สมเย รูปูปปตฺติยา มคฺคํ ภาเวติ วิวิจฺเจว กาเมหิ วิวิจฺจ อกุสเลหิ ธมฺเมหิ อวิตกฺกํ วิจารมตฺตํ วิเวกชํ ปีติสุขํ ทุติยํ ฌานํ อุปสมฺปชฺช วิหรติ กรุณาสหคตํ, ยา ตสฺมิํ สมเย กรุณา กรุณายนา กรุณายิตตฺตํ กรุณาเจโตวิมุตฺติ, อยํ วุจฺจติ กรุณา.\csromanspacer{} อวเสสา ธมฺมา กรุณาย สมฺปยุตฺตา. (2-5)}

\prose{อิธ ภิกฺขุ ยสฺมิํ สมเย รูปูปปตฺติยา มคฺคํ ภาเวติ วิตกฺกวิจารานํ วูปสมา ฯเปฯ ตติยํ ฌานํ อุปสมฺปชฺช วิหรติ กรุณาสหคตํ, ยา \csromanfolio{292}ตสฺมิํ สมเย กรุณา กรุณายนา กรุณายิตตฺตํ กรุณาเจโตวิมุตฺติ, อยํ วุจฺจติ กรุณา.\csromanspacer{} อวเสสา ธมฺมา กรุณาย สมฺปยุตฺตา. (3-6)}

\prose{อิธ ภิกฺขุ ยสฺมิํ สมเย รูปูปปตฺติยา มคฺคํ ภาเวติ ปีติยา จ วิราคา ฯเปฯ จตุตฺถํ ฌานํ อุปสมฺปชฺช วิหรติ กรุณาสหคตํ, ยา ตสฺมิํ สมเย กรุณา กรุณายนา กรุณายิตตฺตํ กรุณาเจโตวิมุตฺติ, อยํ วุจฺจติ กรุณา.\csromanspacer{} อวเสสา ธมฺมา กรุณาย สมฺปยุตฺตา. (4-7)}

\proseitem{688}{ตตฺถ กตมา มุทิตา, อิธ ภิกฺขุ ยสฺมิํ สมเย รูปูปปตฺติยา มคฺคํ ภาเวติ วิวิจฺเจว กาเมหิ ฯเปฯ ปฐมํ ฌานํ อุปสมฺปชฺช วิหรติ มุทิตาสหคตํ, ยา ตสฺมิํ สมเย มุทิตา มุทิตายนา มุทิตายิตตฺตํ มุทิตาเจโตวิมุตฺติ, อยํ วุจฺจติ มุทิตา.\csromanspacer{} อวเสสา ธมฺมา มุทิตาย สมฺปยุตฺตา. (1)}

\prose{ตตฺถ กตมา มุทิตา, อิธ ภิกฺขุ ยสฺมิํ สมเย รูปูปปตฺติยา มคฺคํ ภาเวติ วิตกฺกวิจารานํ วูปสมา ฯเปฯ ทุติยํ ฌานํ อุปสมฺปชฺช วิหรติ มุทิตาสหคตํ, ยา ตสฺมิํ สมเย มุทิตา มุทิตายนา มุทิตายิตตฺตํ มุทิตาเจโตวิมุตฺติ, อยํ วุจฺจติ มุทิตา.\csromanspacer{} อวเสสา ธมฺมา มุทิตาย สมฺปยุตฺตา. (2)}

\prose{ตตฺถ กตมา มุทิตา, อิธ ภิกฺขุ ยสฺมิํ สมเย รูปูปปตฺติยา มคฺคํ ภาเวติ ปีติยา จ วิราคา ฯเปฯ ตติยํ ฌานํ อุปสมฺปชฺช วิหรติ มุทิตาสหคตํ, ยา ตสฺมิํ สมเย มุทิตา มุทิตายนา มุทิตายิตตฺตํ มุทิตาเจโตวิมุตฺติ, อยํ วุจฺจติ มุทิตา.\csromanspacer{} อวเสสา ธมฺมา มุทิตาย สมฺปยุตฺตา. (3)}

\proseitem{689}{อิธ ภิกฺขุ ยสฺมิํ สมเย รูปูปปตฺติยา มคฺคํ ภาเวติ วิวิจฺเจว กาเมหิ ฯเปฯ ปฐมํ ฌานํ อุปสมฺปชฺช วิหรติ มุทิตาสหคตํ, ยา ตสฺมิํ สมเย มุทิตา มุทิตายนา มุทิตายิตตฺตํ มุทิตาเจโตวิมุตฺติ, อยํ วุจฺจติ มุทิตา.\csromanspacer{} อวเสสา ธมฺมา มุทิตาย สมฺปยุตฺตา. (1-4)}

\prose{อิธ ภิกฺขุ ยสฺมิํ สมเย รูปูปปตฺติยา มคฺคํ ภาเวติ วิวิจฺเจว กาเมหิ วิวิจฺจ อกุสเลหิ ธมฺเมหิ อวิตกฺกํ วิจารมตฺตํ วิเวกชํ}

\vspace{\tipitakaparskip}
\csromancenter{อปฺปมญฺญาวิภงฺค ปีติสุขํ ทุติยํ ฌานํ อุปสมฺปชฺช วิหรติ มุทิตาสหคตํ, ยา ตสฺมิํ สมเย มุทิตา มุทิตายนา มุทิตายิตตฺตํ มุทิตาเจโตวิมุตฺติ, อยํ วุจฺจติ มุทิตา.\csromanspacer{} อวเสสา ธมฺมา มุทิตาย สมฺปยุตฺตา. (2-5)}
\prose{\csromanfolio{293}อิธ ภิกฺขุ ยสฺมิํ สมเย รูปูปปตฺติยา มคฺคํ ภาเวติ วิตกฺกวิจารานํ วูปสมา ฯเปฯ ตติยํ ฌานํ อุปสมฺปชฺช วิหรติ มุทิตาสหคตํ, ยา ตสฺมิํ สมเย มุทิตา มุทิตายนา มุทิตายิตตฺตํ มุทิตาเจโตวิมุตฺติ, อยํ วุจฺจติ มุทิตา.\csromanspacer{} อวเสสา ธมฺมา มุทิตาย สมฺปยุตฺตา. (3-6)}

\prose{อิธ ภิกฺขุ ยสฺมิํ สมเย รูปูปปตฺติยา มคฺคํ ภาเวติ ปีติยา จ วิราคา ฯเปฯ จตุตฺถํ ฌานํ อุปสมฺปชฺช วิหรติ มุทิตาสหคตํ, ยา ตสฺมิํ สมเย มุทิตา มุทิตายนา มุทิตายิตตฺตํ มุทิตาเจโตมุตฺติ, อยํ วุจฺจติ มุทิตา.\csromanspacer{} อวเสสา ธมฺมา มุทิตาย สมฺปยุตฺตา. (4-7)}

\proseitem{690}{ตตฺถ กตมา อุเปกฺขา, อิธ ภิกฺขุ ยสฺมิํ สมเย รูปูปปตฺติยา มคฺคํ ภาเวติ สุขสฺส จ ปหานา ฯเปฯ จตุตฺถํ ฌานํ อุปสมฺปชฺช วิหรติ อุเปกฺขาสหคตํ, ยา ตสฺมิํ สมเย อุเปกฺขา อุเปกฺขายนา อุเปกฺขายิตตฺตํ อุเปกฺขาเจโตวิมุตฺติ, อยํ วุจฺจติ อุเปกฺขา.\csromanspacer{} อวเสสา ธมฺมา อุเปกฺขาย สมฺปยุตฺตา.}

\proseitem{691}{จตสฺโส อปฺปมญฺญาโย เมตฺตา กรุณา มุทิตา อุเปกฺขา.}

\proseitem{692}{ตตฺถ กตมา เมตฺตา, อิธ ภิกฺขุ ยสฺมิํ สมเย รูปูปปตฺติยา มคฺคํ ภาเวติ วิวิจฺเจว กาเมหิ ฯเปฯ ปฐมํ ฌานํ อุปสมฺปชฺช วิหรติ เมตฺตาสหคตํ, ตสฺมิํ สมเย ผสฺโส โหติ ฯเปฯ อวิกฺเขโป โหติ, อิเม ธมฺมา กุสลา.\csromanspacer{} ตสฺเสว รูปาวจรสฺส กุสลสฺส กมฺมสฺส กตตฺตา อุปจิตตฺตา วิปากํ วิวิจฺเจว กาเมหิ ฯเปฯ ปฐมํ ฌานํ อุปสมฺปชฺช วิหรติ เมตฺตาสหคตํ, ยา ตสฺมิํ สมเย เมตฺติ เมตฺตายนา เมตฺตายิตตฺตํ เมตฺตาเจโตวิมุตฺติ, อยํ วุจฺจติ เมตฺตา.\csromanspacer{} อวเสสา ธมฺมา เมตฺตาย สมฺปยุตฺตา.}

\prose{ตตฺถ กตมา เมตฺตา, อิธ ภิกฺขุ ยสฺมิํ สมเย รูปูปปตฺติยา มคฺคํ ภาเวติ วิตกฺกวิจารานํ วูปสมา ฯเปฯ ทุติยํ ฌานํ อุปสมฺปชฺช วิหรติ เมตฺตาสหคตํ, ตสฺมิํ สมเย ผสฺโส โหติ ฯเปฯ อวิกฺเขโป โหติ, อิเม ธมฺมา กุสลา.\csromanspacer{} ตสฺเสว รุปาวจรสฺส กุสลสฺส กมฺมสฺส \csromanfolio{294}กตตฺตา อุปจิตตฺตา วิปากํ วิตกฺกวิจารานํ วูปสมา ฯเปฯ ทุติยํ ฌานํ ฯเปฯ ตติยํ ฌานํ ฯเปฯ ปฐมํ ฌานํ ฯเปฯ ทุติยํ ฌานํ ฯเปฯ ตติยํ ฌานํ ฯเปฯ จตุตฺถํ ฌานํ อุปสมฺปชฺช วิหรติ เมตฺตาสหคตํ, ยา ตสฺมิํ สมเย เมตฺติ เมตฺตายนา เมตฺตายิตตฺตํ เมตฺตาเจโตวิมุตฺติ, อยํ วุจฺจติ เมตฺตา.\csromanspacer{} อวเสสา ธมฺมา เมตฺตาย สมฺปยุตฺตา.}

\proseitem{693}{ตตฺถ กตมา กรุณา, อิธ ภิกฺขุ ยสฺมิํ สมเย รูปูปปตฺติยา มคฺคํ ภาเวติ วิวิจฺเจว กาเมหิ ฯเปฯ ปฐมํ ฌานํ อุปสมฺปชฺช วิหรติ กรุณาสหคตํ, ตสฺมิํ สมเย ผสฺโส โหติ ฯเปฯ อวิกฺเขโป โหติ, อิเม ธมฺมา กุสลา.\csromanspacer{} ตสฺเสว รูปาวจรสฺส กุสลสฺส กมฺมสฺส กตตฺตา อุปจิตตฺตา วิปากํ วิวิจฺเจว กาเมหิ ฯเปฯ ปฐมํ ฌานํ อุปสมฺปชฺช วิหรติ กรุณาสหคตํ, ยา ตสฺมิํ สมเย กรุณา กรุณายนา กรุณายิตตฺตํ กรุณาเจโตวิมุตฺติ, อยํ วุจฺจติ กรุณา.\csromanspacer{} อวเสสา ธมฺมา กรุณาย สมฺปยุตฺตา.}

\prose{ตตฺถ กตมา กรุณา, อิธ ภิกฺขุ ยสฺมิํ สมเย รูปูปปตฺติยา มคฺคํ ภาเวติ วิตกฺกวิจารานํ วูปสมา ฯเปฯ ทุติยํ ฌานํ อุปสมฺปชฺช วิหรติ กรุณาสหคตํ, ตสฺมิํ สมเย ผสฺโส โหติ ฯเปฯ อวิกฺเขโป โหติ, อิเม ธมฺมา กุสลา.\csromanspacer{} ตสฺเสว รูปาวจรสฺส กุสลสฺส กมฺมสฺส กตตฺตา อุปจิตตฺตา วิปากํ วิตกฺกวิจารานํ วูปสมา ฯเปฯ ทุติยํ ฌานํ ฯเปฯ ตติยํ ฌานํ ฯเปฯ ปฐมํ ฌานํ ฯเปฯ ทุติยํ ฌานํ ฯเปฯ ตติยํ ฌานํ ฯเปฯ จตุตฺถํ ฌานํ อุปสมฺปชฺช วิหรติ กรุณาสหคตํ, ยา ตสฺมิํ สมเย กรุณา กรุณายนา กรุณายิตตฺตํ กรุณาเจโตวิมุตฺติ, อยํ วุจฺจติ กรุณา.\csromanspacer{} อวเสสา ธมฺมา กรุณาย สมฺปยุตฺตา.}

\proseitem{694}{ตตฺถ กตมา มุทิตา, อิธ ภิกฺขุ ยสฺมิํ สมเย รูปูปปตฺติยา มคฺคํ ภาเวติ วิวิจฺเจว กาเมหิ ฯเปฯ ปฐมํ ฌานํ อุปสมฺปชฺช วิหรติ มุทิตาสหคตํ, ตสฺมิํ สมเย ผสฺโส โหติ ฯเปฯ อวิกฺเขโป โหติ, อิเม ธมฺมา กุสลา.\csromanspacer{} ตสฺเสว รูปาวจรสฺส กุสลสฺส กมฺมสฺส กตตฺตา อุปจิตตฺตา วิปากํ วิวิจฺเจว กาเมหิ ฯเปฯ ปฐมํ ฌานํ อุปสมฺปชฺช วิหรติ มุทิตาสหคตํ, ยา ตสฺมิํ สมเย มุทิตา มุทิตายนา มุทิตายิตตฺตํ มุทิตาเจโตวิมุตฺติ, อยํ วุจฺจติ มุทิตา.\csromanspacer{} อวเสสา ธมฺมา มุทิตาย สมฺปยุตฺตา.}

\csromanheader{2}{13. อปฺปมญฺญาวิภงฺค}
\prose{\csromanfolio{295}ตตฺถ กตมา มุทิตา, อิธ ภิกฺขุ ยสฺมิํ สมเย รูปูปปตฺติยา มคฺคํ ภาเวติ วิตกฺกวิจารานํ วูปสมา ฯเปฯ ทุติยํ ฌานํ อุปสมฺปชฺช วิหรติ มุทิตาสหคตํ, ตสฺมิํ สมเย ผสฺโส โหติ ฯเปฯ อวิกฺเขโป โหติ, อิเม ธมฺมา กุสลา.\csromanspacer{} ตสฺเสว รูปาวจรสฺส กุสลสฺส กมฺมสฺส กตตฺตา อุปจิตตฺตา วิปากํ วิตกฺกวิจารานํ วูปสมา ฯเปฯ ทุติยํ ฌานํ ฯเปฯ ตติยํ ฌานํ ฯเปฯ ปฐมํ ฌานํ ฯเปฯ ทุติยํ ฌานํ ฯเปฯ ตติยํ ฌานํ ฯเปฯ จตุตฺถํ ฌานํ อุปสมฺปชฺช วิหรติ มุทิตาสหคตํ, ยา ตสฺมิํ สมเย มุทิตา มุทิตายนา มุทิตายิตตฺตํ มุทิตาเจโตวิมุตฺติ, อยํ วุจฺจติ มุทิตา.\csromanspacer{} อวเสสา ธมฺมา มุทิตาย สมฺปยุตฺตา.}

\proseitem{695}{ตตฺถ กตมา อุเปกฺขา, อิธ ภิกฺขุ ยสฺมิํ สมเย รูปูปปตฺติยา มคฺคํ ภาเวติ สุขสฺส จ ปหานา ฯเปฯ จตุตฺถํ ฌานํ อุปสมฺปชฺช วิหรติ อุเปกฺขาสหคตํ, ตสฺมิํ สมเย ผสฺโส โหติ ฯเปฯ อวิกฺเขโป โหติ, อิเม ธมฺมา กุสลา.\csromanspacer{} ตสฺเสว รูปาวจรสฺส กุสลสฺส กมฺมสฺส กตตฺตา อุปจิตตฺตา วิปากํ สุขสฺส จ ปหานา ทุกฺขสฺส จ ปหานา ฯเปฯ จตุตฺถํ ฌานํ อุปสมฺปชฺช วิหรติ อุเปกฺขาสหคตํ, ยา ตสฺมิํ สมเย อุเปกฺขา อุเปกฺขายนา อุเปกฺขายิตตฺตํ อุเปกฺขาเจโตวิมุตฺติ, อยํ วุจฺจติ อุเปกฺขา.\csromanspacer{} อวเสสา ธมฺมา อุเปกฺขาย สมฺปยุตฺตา.}

\proseitem{696}{จตสฺโส อปฺปมญฺญาโย เมตฺตา กรุณา มุทิตา อุเปกฺขา.}

\proseitem{697}{ตตฺถ กตมา เมตฺตา, อิธ ภิกฺขุ ยสฺมิํ สมเย รูปาวจรํ ฌานํ ภาเวติ กิริยํ เนว กุสลํ นากุสลํ น จ กมฺมวิปากํ ทิฏฺฐธมฺมสุขวิหารํ วิวิจฺเจว กาเมหิ ฯเปฯ ปฐมํ ฌานํ อุปสมฺปชฺช วิหรติ เมตฺตาสหคตํ, ยา ตสฺมิํ สมเย เมตฺติ เมตฺตายนา เมตฺตายิตตฺตํ เมตฺตาเจโตวิมุตฺติ, อยํ วุจฺจติ เมตฺตา.\csromanspacer{} อวเสสา ธมฺมา เมตฺตาย สมฺปยุตฺตา.}

\prose{ตตฺถ กตมา เมตฺตา, อิธ ภิกฺขุ ยสฺมิํ สมเย รูปาวจรํ ฌานํ ภาเวติ กิริยํ เนว กุสลํ นากุสลํ น จ กมฺมวิปากํ ทิฏฺฐธมฺมสุขวิหารํ วิตกฺกวิจารานํ วูปสมา ฯเปฯ ทุติยํ ฌานํ ฯเปฯ ตติยํ ฌานํ ฯเปฯ ปฐมํ ฌานํ ฯเปฯ ทุติยํ ฌานํ ฯเปฯ ตติยํ ฌานํ ฯเปฯ จตุตฺถํ ฌานํ อุปสมฺปชฺช วิหรติ เมตฺตาสหคตํ, ยา ตสฺมิํ สมเย เมตฺติ เมตฺตายนา เมตฺตายิตตฺตํ เมตฺตาเจโตวิมุตฺติ, อยํ วุจฺจติ เมตฺตา.\csromanspacer{} อวเสสา ธมฺมา เมตฺตาย สมฺปยุตฺตา.}

\proseitem{698}{\csromanfolio{296}ตตฺถ กตมา กรุณา ฯเปฯ.\csromanspacer{} ตตฺถ กตมา มุทิตา ฯเปฯ.\csromanspacer{} ตตฺถ กตมา อุเปกฺขา, อิธ ภิกฺขุ ยสฺมิํ สมเย รูปาวจรํ ฌานํ ภาเวติ กิริยํ เนว กุสลํ นากุสลํ น จ กมฺมวิปากํ ทิฏฺฐธมฺมสุขวิหารํ สุขสฺส จ ปหานา ทุกฺขสฺส จ ปหานา ฯเปฯ จตุตฺถํ ฌานํ อุปสมฺปชฺช วิหรติ อุเปกฺขาสหคตํ, ยา ตสฺมิํ สมเย อุเปกฺขา อุเปกฺขายนา อุเปกฺขายิตตฺตํ อุเปกฺขาเจโตวิมุตฺติ, อยํ วุจฺจติ อุเปกฺขา.\csromanspacer{} อวเสสา ธมฺมา อุเปกฺขาย สมฺปยุตฺตา.}

\csromanheader{2}{อภิธมฺมภาชนียํ.}
\csromansectionrule
\csromanmatikaanchor{mtk.123}
\csromantocmarkref{subparagraph}{3. ปญฺหาปุจฺฉก}{mtk.123}
\bookmark[dest={mtk.123},level=5]{3. ปญฺหาปุจฺฉก}
\csromanheader{6}{3. ปญฺหาปุจฺฉก}
\proseitem{699}{จตสฺโส อปฺปมญฺญาโย, อิธ ภิกฺขุ เมตฺตาสหคเตน เจตสา เอกํ ทิสํ ผริตฺวา วิหรติ, ตถา ทุติยํ.\csromanspacer{} ตถา ตติยํ.\csromanspacer{} ตถา จตุตฺถํ.\csromanspacer{} อิติ อุทฺธมโธ ติริยํ สพฺพธิ สพฺพตฺตตาย สพฺพาวนฺตํ โลกํ เมตฺตาสหคเตน เจตสา วิปุเลน มหคฺคเตน อปฺปมาเณน อเวเรน อพฺยาปชฺเชน ผริตฺวา วิหรติ.\csromanspacer{} กรุณาสหคเตน เจตสา เอกํ ทิสํ ผริตฺวา วิหรติ, ตถา ทุติยํ.\csromanspacer{} ตถา ตติยํ.\csromanspacer{} ตถา จตุตฺถํ.\csromanspacer{} อิติ อุทฺธมโธ ติริยํ สพฺพธิ สพฺพตฺตตาย สพฺพาวนฺตํ โลกํ กรุณาสหคเตน เจตสา วิปุเลน มหคฺคเตน อปฺปมาเณน อเวเรน อพฺยาปชฺเชน ผริตฺวา วิหรติ.\csromanspacer{} มุทิตาสหคเตน เจตสา เอกํ ทิสํ ผริตฺวา วิหรติ, ตถา ทุติยํ.\csromanspacer{} ตถา ตติยํ.\csromanspacer{} ตถา จตุตฺถํ.\csromanspacer{} อิติ อุทฺธมโธ ติริยํ สพฺพธิ สพฺพตฺตตาย สพฺพาวนฺตํ โลกํ มุทิตาสหคเตน เจตสา วิปุเลน มหคฺคเตน อปฺปมาเณน อเวเรน อพฺยาปชฺเชน ผริตฺวา วิหรติ.\csromanspacer{} อุเปกฺขาสหคเตน เจตสา เอกํ ทิสํ ผริตฺวา วิหรติ, ตถา ทุติยํ.\csromanspacer{} ตถา ตติยํ.\csromanspacer{} ตถา จตุตฺถํ.\csromanspacer{} อิติ อุทฺธมโธ ติริยํ สพฺพธิ สพฺพตฺตตาย สพฺพาวนฺตํ โลกํ อุเปกฺขาสหคเตน เจตสา วิปุเลน มหคฺคเตน อปฺปมาเณน อเวเรน อพฺยาปชฺเชน ผริตฺวา วิหรติ.}

\proseitem{700}{จตุนฺนํ อปฺปมญฺญานํ กติ กุสลา, กติ อกุสลา, กติ อพฺยากตา ฯเปฯ กติ สรณา, กติ อรณา.}

\csromanmatikaanchor{mtk.124}
\csromantocmarkref{subparagraph}{1. ติก}{mtk.124}
\bookmark[dest={mtk.124},level=5]{1. ติก}
\csromanheader{6}{13. อปฺปมญฺญาวิภงฺค \textbf{1. ติก}}
\proseitem{701}{\csromanfolio{297}สิยา กุสลา, สิยา อพฺยากตา.\csromanspacer{} ติสฺโส อปฺปมญฺญาโย สุขาย เวทนาย สมฺปยุตฺตา.\csromanspacer{} อุเปกฺขา อทุกฺขมสุขาย เวทนาย สมฺปยุตฺตา, สิยา วิปากา, สิยา วิปากธมฺมธมฺมา, สิยา เนว วิปากนวิปากธมฺมธมฺมา, สิยา อุปาทินฺนุปาทานิยา, สิยา อนุปาทินฺนุปาทานิยา, อสํกิลิฏฺฐสํกิเลสิกา.\csromanspacer{} ติสฺโส อปฺปมญฺญาโย สิยา สวิตกฺกสวิจารา, สิยา อวิตกฺกวิจารมตฺตา, สิยา อวิตกฺก-อวิจารา, อุเปกฺขา อวิตกฺก-อวิจารา.\csromanspacer{} ติสฺโส อปฺปมญฺญาโย สิยา ปีติสหคตา, สิยา สุขสหคตา, น อุเปกฺขาสหคตา.\csromanspacer{} สิยา น วตฺตพฺพา “ปีติสหคตา”ติ.\csromanspacer{} อุเปกฺขา อุเปกฺขาสหคตา.\csromanspacer{} เนว ทสฺสเนน น ภาวนาย ปหาตพฺพา, เนว ทสฺสเนน น ภาวนาย ปหาตพฺพเหตุกา, สิยา อาจยคามิโน, สิยา เนวาจยคามินาปจยคามิโน, เนวเสกฺขนาเสกฺขา, มหคฺคตา, น วตฺตพฺพา “ปริตฺตารมฺมณา”ติปิ, “มหคฺคตารมฺมณา”ติปิ, “อปฺปมาณารมฺมณา”ติปิ, มชฺฌิมา, อนิยตา, น วตฺตพฺพา “มคฺคารมฺมณา”ติปิ, “มคฺคเหตุกา”ติปิ, “มคฺคาธิปติโน”ติปิ, สิยา อุปฺปนฺนา, สิยา อนุปฺปนฺนา, สิยา อุปฺปาทิโน.\csromanspacer{} สิยา อตีตา, สิยา อนาคตา, สิยา ปจฺจุปฺปนฺนา, น วตฺตพฺพา “อตีตารมฺมณา”ติปิ, “อนาคตารมฺมณา”ติปิ, “ปจฺจุปฺปนฺนารมฺมณา”ติปิ, สิยา อชฺฌตฺตา, สิยา พหิทฺธา, สิยา อชฺฌตฺตพหิทฺธา, พหิทฺธารมฺมณา, อนิทสฺสน-อปฺปฏิฆา.}

\csromanmatikaanchor{mtk.125}
\csromantocmarkref{subparagraph}{2. ทุก}{mtk.125}
\bookmark[dest={mtk.125},level=5]{2. ทุก}
\csromanheader{6}{2. ทุก}
\proseitem{702}{เมตฺตา เหตู.\csromanspacer{} ติสฺโส อปฺปมญฺญาโย น เหตู, สเหตุกา, เหตุสมฺปยุตฺตา.\csromanspacer{} เมตฺตา เหตุ เจว สเหตุกา จ.\csromanspacer{} ติสฺโส อปฺปมญฺญาโย น วตฺตพฺพา “เหตุ เจว สเหตุกา จา”ติ, สเหตุกา เจว น จ เหตู.\csromanspacer{} เมตฺตา เหตุ เจว เหตุสมฺปยุตฺตา จ.\csromanspacer{} ติสฺโส อปฺปมญฺญาโย น วตฺตพฺพา “เหตู เจว เหตุสมฺปยุตฺตา จา”ติ, เหตุสมฺปยุตฺตา เจว น จ เหตู.\csromanspacer{} ติสฺโส อปฺปมญฺญาโย น เหตู, สเหตุกา.\csromanspacer{} เมตฺตา น วตฺตพฺพา “น เหตุ สเหตุกา”ติปิ, “น เหตุ อเหตุกา”ติปิ.}

\prose{สปฺปจฺจยา, สงฺขตา, อนิทสฺสนา, อปฺปฏิฆา, อรูปา, โลกิยา, เกนจิ วิญฺเญยฺยา, เกนจิ น วิญฺเญยฺยา, โน อาสวา, สาสวา, \csromanfolio{298}อาสววิปฺปยุตฺตา, น วตฺตพฺพา “อาสวา เจว สาสวา จา”ติ, สาสวา เจว โน จ อาสวา, น วตฺตพฺพา “อาสวา เจว อาสวสมฺปยุตฺตา จา”ติปิ, “อาสวสมฺปยุตฺตา เจว โน จ อาสวา”ติปิ, อาสววิปฺปยุตฺตา สาสวา.}

\prose{โน สํโยชนา ฯเปฯ.\csromanspacer{} โน คนฺถา ฯเปฯ.\csromanspacer{} โน โอฆา ฯเปฯ.\csromanspacer{} โน โยคา ฯเปฯ.\csromanspacer{} โน นีวรณา ฯเปฯ.\csromanspacer{} โน ปรามาสา ฯเปฯ สารมฺมณา, โน จิตฺตา, เจตสิกา, จิตฺตสมฺปยุตฺตา, จิตฺตสํสฏฺฐา, จิตฺตสมุฏฺฐานา, จิตฺตสหภุโน, จิตฺตานุปริวตฺติโน, จิตฺตสํสฏฺฐสมุฏฺฐานา, จิตฺตสํสฏฺฐสมุฏฺฐานสหภุโน.\csromanspacer{} จิตฺตสํสฏฺฐสมุฏฺฐานานุปริวตฺติโน, พาหิรา, โน อุปาทา, สิยา อุปาทินฺนา, สิยา อนุปาทินฺนา.}

\prose{โน อุปาทานา ฯเปฯ.\csromanspacer{} โน กิเลสา ฯเปฯ.\csromanspacer{} น ทสฺสเนน ปหาตพฺพา, น ภาวนาย ปหาตพฺพา, น ทสฺสเนน ปหาตพฺพเหตุกา, น ภาวนาย ปหาตพฺพเหตุกา.\csromanspacer{} ติสฺโส อปฺปมญฺญาโย สิยา สวิตกฺกา, สิยา อวิตกฺกา.\csromanspacer{} อุเปกฺขา อวิตกฺกา.\csromanspacer{} ติสฺโส อปฺปมญฺญาโย สิยา สวิจารา, สิยา อวิจารา.\csromanspacer{} อุเปกฺขา อวิจารา.\csromanspacer{} ติสฺโส อปฺปมญฺญาโย สิยา สปฺปีติกา, สิยา อปฺปีติกา, อุเปกฺขา อปฺปีติกา.\csromanspacer{} ติสฺโส อปฺปมญฺญาโย สิยา ปีติสหคตา, สิยา น ปีติสหคตา.\csromanspacer{} อุเปกฺขา น ปีติสหคตา.\csromanspacer{} ติสฺโส อปฺปมญฺญาโย สุขสหคตา, อุเปกฺขา น สุขสหคตา.\csromanspacer{} อุเปกฺขา อุเปกฺขาสหคตา.\csromanspacer{} ติสฺโส อปฺปมญฺญาโย น อุเปกฺขาสหคตา, น กามาวจรา, รูปาวจรา, น อรูปาวจรา, ปริยาปนฺนา, อนิยฺยานิกา, อนิยโต, ส-อุตฺตรา, อรณาติ.}

\vspace{\tipitakaparskip}
\csromancenter{ปญฺหาปุจฺฉกํ.}
\nitthitam{\textbf{อปฺปมญฺญาวิภงฺโค นิฏฺฐิโต.}}
\csromanmatikaanchor{mtk.126}
\csromantocmarknopage{subparagraph}{14. สิกฺขาปทวิภงฺค}{mtk.126}
\bookmark[dest={mtk.126},level=5]{14. สิกฺขาปทวิภงฺค}
\csromanheader{6}{14. สิกฺขาปทวิภงฺค}
\csromanmatikaanchor{mtk.127}
\csromantocmarkref{subparagraph}{1. อภิธมฺมภาชนีย}{mtk.127}
\bookmark[dest={mtk.127},level=5]{1. อภิธมฺมภาชนีย}
\csromanheader{6}{1. อภิธมฺมภาชนีย}
\proseitem{703}{\csromanfolio{299}ปญฺจ สิกฺขาปทานิ ปาณาติปาตา เวรมณี สิกฺขาปทํ, อทินฺนาทานา เวรมณี สิกฺขาปทํ, กาเมสุมิจฺฉาจารา เวรมณี สิกฺขาปทํ, มุสาวาทา เวรมณี สิกฺขาปทํ, สุราเมรยมชฺชปมาทฏฺฐานา เวรมณี สิกฺขาปทํ.}

\proseitem{704}{ตตฺถ กตมํ ปาณาติปาตา เวรมณี สิกฺขาปทํ, ยสฺมิํ สมเย กามาวจรํ กุสลํ จิตฺตํ อุปฺปนฺนํ โหติ โสมนสฺสสหคตํ ญาณสมฺปยุตฺตํ ปาณาติปาตา วิรมนฺตสฺส, ยา ตสฺมิํ สมเย ปาณาติปาตา อารติ วิรติ ปฏิวิรติ เวรมณี อกิริยา อกรณํ อนชฺฌาปตฺติ เวลา- อนติกฺกโม เสตุฆาโต, อิทํ วุจฺจติ ปาณาติปาตา เวรมณี สิกฺขาปทํ.\csromanspacer{} อวเสสา ธมฺมา เวรมณิยา สมฺปยุตฺตา.}

\prose{ตตฺถ กตมํ ปาณาติปาตา เวรมณี สิกฺขาปทํ, ยสฺมิํ สมเย กามาวจรํ กุสลํ จิตฺตํ อุปฺปนฺนํ โหติ โสมนสฺสสหคตํ ญาณสมฺปยุตฺตํ ปาณาติปาตา วิรมนฺตสฺส, ยา ตสฺมิํ สมเย เจตนา สญฺเจตนา สญฺเจตยิตตฺตํ, อิทํ วุจฺจติ ปาณาติปาตา เวรมณี สิกฺขาปทํ.\csromanspacer{} อวเสสา ธมฺมา เจตนาย สมฺปยุตฺตา.}

\prose{ตตฺถ กตมํ ปาณาติปาตา เวรมณี สิกฺขาปทํ, ยสฺมิํ สมเย กามาวจรํ กุสลํ จิตฺตํ อุปฺปนฺนํ โหติ โสมนสฺสสหคตํ ญาณสมฺปยุตฺตํ ปาณาติปาตา วิรมนฺตสฺส, โย ตสฺมิํ สมเย ผสฺโส ฯเปฯ ปคฺคาโห อวิกฺเขโป, อิทํ วุจฺจติ ปาณาติปาตา เวรมณี สิกฺขาปทํ.}

\prose{ตตฺถ กตมํ ปาณาติปาตา เวรมณี สิกฺขาปทํ, ยสฺมิํ สมเย กามาวจรํ กุสลํ จิตฺตํ อุปฺปนฺนํ โหติ โสมนสฺสสหคตํ ญาณสมฺปยุตฺตํ สสงฺขาเรน ฯเปฯ โสมนสฺสสหคตํ ญาณวิปฺปยุตฺตํ ฯเปฯ โสมนสฺสสหคตํ ญาณวิปฺปยุตฺตํ สสงฺขาเรน ฯเปฯ อุเปกฺขาสหคตํ ญาณสมยุตฺตํ ฯเปฯ อุเปกฺขาสหคตํ ญาณสมฺปยุตฺตํ สสงฺขาเรน ฯเปฯ อุเปกฺขาสหคตํ ญาณวิปฺปยุตฺตํ ฯเปฯ อุเปกฺขาสหคตํ ญาณวิปฺปยุตฺตํ สสงฺขาเรน ปาณาติปาตา วิรมนฺตสฺส, ยา ตสฺมิํ สมเย ปาณาติปาตา อารติ วิรติ ปฏิวิรติ เวรมณี อกิริยา อกรณํ อนชฺฌาปตฺติ เวลา-อนติกฺกโม \csromanfolio{300}เสตุฆาโต, อิทํ วุจฺจติ ปาณาติปาตา เวรมณี สิกฺขาปทํ, อวเสสา ธมฺมา เวรมณิยา สมฺปยุตฺตา.}

\proseitem{705}{ตตฺถ กตมํ ปาณาติปาตา เวรมณี สิกฺขาปทํ, ยสฺมิํ สมเย กามาวจรํ กุสลํ จิตฺตํ อุปฺปนฺนํ โหติ อุเปกฺขาสหคตํ ญาณวิปฺปยุตฺตํ สสงฺขาเรน ปาณาติปาตา วิรมนฺตสฺส, ยา ตสฺมิํ สมเย เจตนา สญฺเจตนา สญฺเจตยิตตฺตํ, อิทํ วุจฺจติ ปาณาติปาตา เวรมณี สิกฺขาปทํ, อวเสสา ธมฺมา เจตนาย สมฺปยุตฺตา.}

\prose{ตตฺถ กตมํ ปาณาติปาตา เวรมณี สิกฺขาปทํ, ยสฺมิํ สมเย กามาวจรํ กุสลํ จิตฺตํ อุปฺปนฺนํ โหติ อุเปกฺขาสหคตํ ญาณวิปฺปยุตฺตํ สสงฺขาเรน ปาณาติปาตา วิรมนฺตสฺส, ผสฺโส ฯเปฯ ปคฺคาโห อวิกฺเขโป, อิทํ วุจฺจติ ปาณาติปาตา เวรมณี สิกฺขาปทํ.}

\proseitem{706}{ตตฺถ กตมํ อทินฺนาทานา เวรมณี สิกฺขาปทํ ฯเปฯ กาเมสุมิจฺฉาจารา เวรมณี สิกฺขาปทํ ฯเปฯ มุสาวาทา เวรมณี สิกฺขาปทํ ฯเปฯ สุราเมรยมชฺชปมาทฏฺฐานา เวรมณี สิกฺขาปทํ ยสฺมิํ สมเย กามาวจรํ กุสลํ จิตฺตํ อุปฺปนฺนํ โหติ โสมนสฺสสหคตํ ญาณสมฺปยุตฺตํ สุราเมรยมชฺชปมาทฏฺฐานา วิรมนฺตสฺส, ยา ตสฺมิํ สมเย สุราเมรยมชฺชปมาทฏฺฐานา อารติ วิรติ ปฏิวิรติ เวรมณี อกิริยา อกรณํ อนชฺฌาปตฺติ เวลา-อนติกฺกโม เสตุฆาโต, อิทํ วุจฺจติ สุราเมรยมชฺชปมาทฏฺฐานา เวรมณี สิกฺขาปทํ.\csromanspacer{} อวเสสา ธมฺมา เวรมณิยา สมฺปยุตฺตา.}

\prose{ตตฺถ กตมํ สุราเมรยมชฺชปมาทฏฺฐานา เวรมณี สิกฺขาปทํ, ยสฺมิํ สมเย กามาวจรํ กุสลํ จิตฺตํ อุปฺปนฺนํ โหติ โสมนสฺสสหคตํ ญาณสมฺปยุตฺตํ สุราเมรยมชฺชปมาทฏฺฐานา วิรมนฺตสฺส, ยา ตสฺมิํ สมเย เจตนา สญฺเจตนา สญฺเจตยิตตฺตํ, อิทํ วุจฺจติ สุราเมรยมชฺชปมาทฏฺฐานา เวรมณี สิกฺขาปทํ.\csromanspacer{} อวเสสา ธมฺมา เจตนาย สมฺปยุตฺตา.}

\prose{ตตฺถ กตมํ สุราเมรยมชฺชปมาทฏฺฐานา เวรมณี สิกฺขาปทํ, ยสฺมิํ สมเย กามาวจรํ กุสลํ จิตฺตํ อุปฺปนฺนํ โหติ โสมนสฺสสหคตํ ญาณสมฺปยุตฺตํ สุราเมรยมชฺชปมาทฏฺฐานา วิรมนฺตสฺส, โย ตสฺมิํ สมเย ผสฺโส ฯเปฯ ปคฺคาโห อวิกฺเขโป, อิทํ วุจฺจติ สุราเมรยมชฺชปมาทฏฺฐานา เวรมณี สิกฺขาปทํ.}

\csromanheader{2}{14. สิกฺขาปทวิภงฺค}
\prose{\csromanfolio{301}ตตฺถ กตมํ สุราเมรยมชฺชปมาทฏฺฐานา เวรมณี สิกฺขาปทํ, ยสฺมิํ สมเย กามาวจรํ กุสลํ จิตฺตํ อุปฺปนฺนํ โหติ โสมนสฺสสหคตํ ญาณสมฺปยุตฺตํ สสงฺขาเรน ฯเปฯ โสมนสฺสสหคตํ ญาณวิปฺปยุตฺตํ ฯเปฯ โสมนสฺสสหคตํ ญาณวิปฺปยุตฺตํ สสงฺขาเรน ฯเปฯ อุเปกฺขาสหคตํ ญาณสมฺปยุตฺตํ ฯเปฯ อุเปกฺขาสหคตํ ญาณสมฺปยุตฺตํ สสงฺขาเรน ฯเปฯ อุเปกฺขาสหคตํ ญาณสมฺปยุตฺตํ ฯเปฯ อุเปกฺขาสหคตํ ญาณวิปฺปยุตฺตํ สสงฺขาเรน สุราเมรยมชฺชปมาทฏฺฐานา วิรมนฺตสฺส, ยา ตสฺมิํ สมเย สุราเมรยมชฺชปมาทฏฺฐานา อารติ วิรติ ปฏิวิรติ เวรมณี อกิริยา อกรณํ อนชฺฌาปตฺติ เวลา-อนติกฺกโม เสตุฆาโต, อิทํ วุจฺจติ สุราเมรยมชฺชปมาทฏฺฐานา เวรมณี สิกฺขาปทํ.\csromanspacer{} อวเสสา ธมฺมา เวรมณียา สมฺปยุตฺตา.}

\proseitem{707}{ตตฺถ กตมํ สุราเมรยมชฺชปมาทฏฺฐานา เวรมณี สิกฺขาปทํ, ยสฺมิํ สมเย กามาวจรํ กุสลํ จิตฺตํ อุปฺปนฺนํ โหติ อุเปกฺขาสหคตํ ญาณวิปฺปยุตฺตํ สสงฺขาเรน สุราเมรยมชฺชปมาทฏฺฐานา วิรมนฺตสฺส, ยา ตสฺมิํ สมเย เจตนา สญฺเจตนา สญฺเจตยิตตฺตํ, อิทํ วุจฺจติ สุราเมรยมชฺชปมาทฏฺฐานา เวรมณี สิกฺขาปทํ.\csromanspacer{} อวเสสา ธมฺมา เจตนาย สมฺปยุตฺตา.}

\prose{ตตฺถ กตมํ สุราเมรยมชฺชปมาทฏฺฐานา เวรมณี สิกฺขาปทํ, ยสฺมิํ สมเย กามาวจรํ กุสลํ จิตฺตํ อุปฺปนฺนํ โหติ อุเปกฺขาสหคตํ ญาณวิปฺปยุตฺตํ สสงฺขาเรน สุราเมรยมชฺชปมาทฏฺฐานา วิรมนฺตสฺส, โย ตสฺมิํ สมเย ผสฺโส ฯเปฯ ปคฺคาโห อวิกฺเขโป, อิทํ วุจฺจติ สุราเมรยมชฺชปมาทฏฺฐานา เวรมณี สิกฺขาปทํ.}

\proseitem{708}{ปญฺจ สิกฺขาปทานิ ปาณาติปาตา เวรมณี สิกฺขาปทํ, อทินฺนาทานา เวรมณี สิกฺขาปทํ, กาเมสุมิจฺฉาจารา เวรมณี สิกฺขาปทํ, มุสาวาทา เวรมณี สิกฺขาปทํ, สุราเมรยมชฺชปมาทฏฺฐานา เวรมณี สิกฺขาปทํ.}

\proseitem{709}{ตตฺถ กตมํ ปาณาติปาตา เวรมณี สิกฺขาปทํ, ยสฺมิํ สมเย กามาวจรํ กุสลํ จิตฺตํ อุปฺปนฺนํ โหติ โสมนสฺสสหคตํ ญาณสมฺปยุตฺตํ หีนํ มชฺฌิมํ ปณีตํ, ฉนฺทาธิปเตยฺยํ, วีริยาธิปเตยฺยํ, จิตฺตาธิปเตยฺยํ, วีมํสาธิปเตยฺยํ.\csromanspacer{} ฉนฺทาธิปเตยฺยํ หีนํ มชฺฌิมํ ปณีตํ, \csromanfolio{302}วีริยาธิปเตยฺยํ หีนํ มชฺฌิมํ ปณีตํ, จิตฺตาธิปเตยฺยํ หีนํ มชฺฌิมํ ปณีตํ, วีมํสาธิปเตยฺยํ หีนํ มชฺฌิมํ ปณีตํ ปาณาติปาตา วิรมนฺตสฺส, ยา ตสฺมิํ สมเย ปาณาติปาตา อารติ วิรติ ปฏิวิรติ เวรมณี อกิริยา อกรณํ อนชฺฌาปตฺติ เวลา-อนติกฺกโม เสตุฆาโต, อิทํ วุจฺจติ ปาณาติปาตา เวรมณี สิกฺขาปทํ.\csromanspacer{} อวเสสา ธมฺมา เวรมณิยา สมฺปยุตฺตา.}

\prose{ตตฺถ กตมํ ปาณาติปาตา เวรมณี สิกฺขาปทํ, ยสฺมิํ สมเย กามาวจรํ กุสลํ จิตฺตํ อุปฺปนฺนํ โหติ โสมนสฺสสหคตํ ญาณสมฺปยุตฺตํ หีนํ มชฺฌิมํ ปณีตํ, ฉนฺทาธิปเตยฺยํ, วีริยาธิปเตยฺยํ, จิตฺตาธิปเตยฺยํ, วีมํสาธิปเตยฺยํ.\csromanspacer{} ฉนฺทาธิปเตยฺยํ หีนํ มชฺฌิมํ ปณีตํ, วีริยาธิปเตยฺยํ หีนํ มชฺฌิมํ ปณีตํ, จิตฺตาธิปเตยฺยํ หีนํ มชฺฌิมํ ปณีตํ, วีมํสาธิปเตยฺยํ หีนํ มชฺฌิมํ ปณีตํ ปาณาติปาตา วิรมนฺตสฺส, ยา ตสฺมิํ สมเย เจตนา สญฺเจตนา สญฺเจตยิตตฺตํ, อิทํ วุจฺจติ ปาณาติปาตา เวรมณี สิกฺขาปทํ.\csromanspacer{} อวเสสา ธมฺมา เจตนาย สมฺปยุตฺตา.}

\prose{ตตฺถ กตมํ ปาณาติปาตา เวรมณี สิกฺขาปทํ, ยสฺมิํ สมเย กามาวจรํ กุสลํ จิตฺตํ อุปฺปนฺนํ โหติ โสมนสฺสสหคตํ ญาณสมฺปยุตฺตํ หีนํ มชฺฌิมํ ปณีตํ, ฉนฺทาธิปเตยฺยํ, วีริยาธิปเตยฺยํ, จิตฺตาธิปเตยฺยํ, วีมํสาธิปเตยฺยํ.\csromanspacer{} ฉนฺทาธิปเตยฺยํ หีนํ มชฺฌิมํ ปณีตํ, วีริยาธิปเตยฺยํ หีนํ มชฺฌิมํ ปณีตํ, จิตฺตาธิปเตยฺยํ หีนํ มชฺฌิมํ ปณีตํ, วีมํสาธิปเตยฺยํ หีนํ มชฺฌิมํ ปณีตํ ปาณาติปาตา วิรมนฺตสฺส, โย ตสฺมิํ สมเย ผสฺโส ฯเปฯ ปคฺคาโห อวิกฺเขโป, อิทํ วุจฺจติ ปาณาติปาตา เวรมณี สิกฺขาปทํ.}

\prose{ตตฺถ กตมํ ปาณาติปาตา เวรมณี สิกฺขาปทํ, ยสฺมิํ สมเย กามาวจรํ กุสลํ จิตฺตํ อุปฺปนฺนํ โหติ โสมนสฺสสหคตํ ญาณสมฺปยุตฺตํ สสงฺขาเรน ฯเปฯ โสมนสฺสสหคตํ ญาณวิปฺปยุตฺตํ ฯเปฯ โสมนสฺสสหคตํ ญาณวิปฺปยุตฺตํ สสงฺขาเรน ฯเปฯ อุเปกฺขาสหคตํ ญาณสมฺปยุตฺตํ ฯเปฯ อุเปกฺขาสหคตํ ญาณสมฺปยุตฺตํ สสงฺขาเรน ฯเปฯ อุเปกฺขาสหคตํ ญาณวิปฺปยุตฺตํ ฯเปฯ อุเปกฺขาสหคตํ ญาณวิปฺปยุตฺตํ สสงฺขาเรน หีนํ มชฺฌิมํ ปณีตํ, ฉนฺทาธิปเตยฺยํ, วีริยาธิปเตยฺยํ, จิตฺตาธิปเตยฺยํ.\csromanspacer{} ฉนฺทาธิปเตยฺยํ หีนํ มชฺฌิมํ ปณีตํ, วีริยาธิปเตยฺยํ หีนํ มชฺฌิมํ ปณีตํ, จิตฺตาธิปเตยฺยํ หีนํ มชฺฌิมํ ปณีตํ ปาณาติปาตา วิรมนฺตสฺส, ยา}

\csromanheader{2}{14. สิกฺขาปทวิภงฺค}
\prose{\csromanfolio{303}ตสฺมิํ สมเย ปาณาติปาตา อารติ วิรติ ปฏิวิรติ เวรมณี อกิริยา อกรณํ อนชฺฌาปตฺติ เวลา-อนติกฺกโม เสตุฆาโต, อิทํ วุจฺจติ ปาณาติปาตา เวรมณี สิกฺขาปทํ.\csromanspacer{} อวเสสา ธมฺมา เวรมณิยา สมฺปยุตฺตา ฯเปฯ.\csromanspacer{} อวเสสา ธมฺมา เจตนาย สมฺปยุตฺตา ฯเปฯ ผสฺโส ฯเปฯ ปคฺคาโห อวิกฺเขโป, อิทํ วุจฺจติ ปาณาติปาตา เวรมณี สิกฺขาปทํ.}

\proseitem{710}{ตตฺถ กตมํ อทินฺนาทานา เวรมณี สิกฺขาปทํ ฯเปฯ กาเมสุมิจฺฉาจารา เวรมณี สิกฺขาปทํ ฯเปฯ มุสาวาทา เวรมณี สิกฺขาปทํ ฯเปฯ สุราเมรยมชฺชปมาทฏฺฐานา เวรมณี สิกฺขาปทํ, ยสฺมิํ สมเย กามาวจรํ กุสลํ จิตฺตํ อุปฺปนฺนํ โหติ โสมนสฺสสหคตํ ญาณสมฺปยุตฺตํ หีนํ มชฺฌิมํ ปณีตํ, ฉนฺทาธิปเตยฺยํ, วีริยาธิปเตยฺยํ, จิตฺตาธิปเตยฺยํ, วีมํสาธิปเตยฺยํ.\csromanspacer{} ฉนฺทาธิปเตยฺยํ หีนํ มชฺฌิมํ ปณีตํ, วีริยาธิปเตยฺยํ หีนํ มชฺฌิมํ ปณีตํ, จิตฺตาธิปเตยฺยํ หีนํ มชฺฌิมํ ปณีตํ, วีมํสาธิปเตยฺยํ หีนํ มชฺฌิมํ ปณีตํ สุราเมรยมชฺชปมาทฏฺฐานา วิรมนฺตสฺส, ยา ตสฺมิํ สมเย สุราเมรยมชฺชปมาทฏฺฐานา อารติ วิรติ ปฏิวิรติ เวรมณี อกิริยา อกรณํ อนชฺฌาปตฺติ เวลา-อนติกฺกโม เสตุฆาโต, อิทํ วุจฺจติ สุราเมรยมชฺชปมาทฏฺฐานา เวรมณี สิกฺขาปทํ.\csromanspacer{} อวเสสา ธมฺมา เวรมณิยา สมฺปยุตฺตา ฯเปฯ.\csromanspacer{} อวเสสา ธมฺมา เจตนาย สมฺปยุตฺตา ฯเปฯ ผสฺโส ฯเปฯ ปคฺคาโห อวิกฺเขโป, อิทํ วุจฺจติ สุราเมรยมชฺชปมาทฏฺฐานา เวรมณี สิกฺขาปทํ.}

\proseitem{711}{ตตฺถ กตมํ สุราเมรยมชฺชปมาทฏฺฐานา เวรมณี สิกฺขาปทํ, ยสฺมิํ สมเย กามาวจรํ กุสลํ จิตฺตํ อุปฺปนฺนํ โหติ โสมนสฺสสหคตํ ญาณสมฺปยุตฺตํ สสงฺขาเรน ฯเปฯ โสมนสฺสสหคตํ ญาณวิปฺปยุตฺตํ ฯเปฯ โสมนสฺสสหคตํ ญาณวิปฺปยุตฺตํ สสงฺขาเรน ฯเปฯ อุเปกฺขาสหคตํ ญาณสมฺปยุตฺตํ ฯเปฯ อุเปกฺขาสหคตํ ญาณสมฺปยุตฺตํ สสงฺขาเรน ฯเปฯ อุเปกฺขาสหคตํ ญาณวิปฺปยุตฺตํ ฯเปฯ อุเปกฺขาสหคตํ ญาณวิปฺปยุตฺตํ สสงฺขาเรน หีนํ มชฺฌิมํ ปณีตํ, ฉนฺทาธิปเตยฺยํ, วีริยาธิปเตยฺยํ, จิตฺตาธิปเตยฺยํ.\csromanspacer{} ฉนฺทาธิปเตยฺยํ หีนํ มชฺฌิมํ ปณีตํ, วิริยาธิปเตยฺยํ หีนํ มชฺฌิมํ ปณีตํ, จิตฺตาธิปเตยฺยํ หีนํ มชฺฌิมํ ปณีตํ สุราเมรยมชฺชปมาทฏฺฐานา วิรมนฺตสฺส, ยา ตสฺมิํ สมเย สุราเมรยมชฺชปมาทฏฺฐานา อารติ วิรติ ปฏิวิรติ เวรมณี อกิริยา อกรณํ อนชฺฌาปตฺติ \csromanfolio{304}เวลา-อนติกฺกโม เสตุฆาโต, อิทํ วุจฺจติ สุราเมรยมชฺชปมาทฏฺฐานา เวรมณี สิกฺขาปทํ.\csromanspacer{} อวเสสา ธมฺมา เวรมณิยา สมฺปยุตฺตา ฯเปฯ.\csromanspacer{} อวเสสา ธมฺมา เจตนาย สมฺปยุตฺตา ฯเปฯ ผสฺโส ฯเปฯ ปคฺคาโห อวิกฺเขโป, อิทํ วุจฺจติ สุราเมรยมชฺชปมาทฏฺฐานา เวรมณี สิกฺขาปทํ.}

\proseitem{712}{กตเม ธมฺมา สิกฺขา, ยสฺมิํ สมเย กามาวจรํ กุสลํ จิตฺตํ อุปฺปนฺนํ โหติ โสมนสฺสสหคตํ ญาณสมฺปยุตฺตํ รูปารมฺมณํ วา ฯเปฯ ธมฺมารมฺมณํ วา ยํ ยํ วา ปนารพฺภ, ตสฺมิํ สมเย ผสฺโส โหติ ฯเปฯ อวิกฺเขโป โหติ, อิเม ธมฺมา สิกฺขา.}

\prose{กตเม ธมฺมา สิกฺขา, ยสฺมิํ สมเย กามาวจรํ กุสลํ จิตฺตํ อุปฺปนฺนํ โหติ โสมนสฺสสหคตํ ญาณสมฺปยุตฺตํ สสงฺขาเรน ฯเปฯ โสมนสฺสสหคตํ ญาณวิปฺปยุตฺตํ ฯเปฯ โสมนสฺสสหคตํ ญาณวิปฺปยุตฺตํ สสงฺขาเรน ฯเปฯ อุเปกฺขาสหคตํ ญาณสมฺปยุตฺตํ ฯเปฯ อุเปกฺขาสหคตํ ญาณสมฺปยุตฺตํ สสงฺขาเรน ฯเปฯ อุเปกฺขาสหคตํ ญาณวิปฺปยุตฺตํ ฯเปฯ อุเปกฺขาสหคตํ ญาณวิปฺปยุตฺตํ สสงฺขาเรน รูปารมฺมณํ วา ฯเปฯ ธมฺมารมฺมณํ วา ยํ ยํ วา ปนารพฺภ, ตสฺมิํ สมเย ผสฺโส โหติ ฯเปฯ อวิกฺเขโป โหติ, อิเม ธมฺมา สิกฺขา.}

\proseitem{713}{กตเม ธมฺมา สิกฺขา, ยสฺมิํ สมเย รูปูปปตฺติยา มคฺคํ ภาเวติ ฯเปฯ อรูปูปปตฺติยา มคฺคํ ภาเวติ ฯเปฯ โลกุตฺตรํ ฌานํ ภาเวติ นิยฺยานิกํ อปจยคามิํ ทิฏฺฐิคตานํ ปหานาย ปฐมาย ภูมิยา ปตฺติยา วิวิจฺเจว กาเมหิ ฯเปฯ ปฐมํ ฌานํ อุปสมฺปชฺช วิหรติ ทุกฺขปฏิปทํ ทนฺธาภิญฺญํ, ตสฺมิํ สมเย ผสฺโส โหติ ฯเปฯ อวิกฺเขโป โหติ, อิเม ธมฺมา สิกฺขา.}

\csromanheader{2}{อภิธมฺมภาชนียํ.}
\csromansectionrule
\csromanmatikaanchor{mtk.128}
\csromantocmarkref{subparagraph}{2. ปญฺหาปุจฺฉก}{mtk.128}
\bookmark[dest={mtk.128},level=5]{2. ปญฺหาปุจฺฉก}
\csromanheader{6}{2. ปญฺหาปุจฺฉก}
\proseitem{714}{ปญฺจ สิกฺขาปทานิ ปาณาติปาตา เวรมณี สิกฺขาปทํ, อทินฺนาทานา เวรมณี สิกฺขาปทํ, กาเมสุมิจฺฉาจารา เวรมณี สิกฺขาปทํ, มุสาวาทา เวรมณี สิกฺขาปทํ, สุราเมรยมชฺชปมาทฏฺฐานา เวรมณี สิกฺขาปทํ.}

\csromanheader{2}{14. สิกฺขาปทวิภงฺค}
\proseitem{715}{\csromanfolio{305}ปญฺจนฺนํ สิกฺขาปทานํ กติ กุสลา, กติ อกุสลา, กติ อพฺยากตา ฯเปฯ กติ สรณา, กติ อรณา.}

\csromanmatikaanchor{mtk.129}
\csromantocmarkref{subparagraph}{1. ติก}{mtk.129}
\bookmark[dest={mtk.129},level=5]{1. ติก}
\csromanheader{6}{1. ติก}
\proseitem{716}{กุสลาเยว.\csromanspacer{} สิยา สุขาย เวทนาย สมฺปยุตฺตา, สิยา อทุกฺขมสุขาย เวทนาย สมฺปยุตฺตา, วิปากธมฺมธมฺมา, อนุปาทินฺนุปาทานิยา, อสํกิลิฏฺฐสํกิเลสิกา, สวิตกฺกสวิจารา, สิยา ปีติสหคตา, สิยา สุขสหคตา, สิยา อุเปกฺขาสหคตา.}

\prose{เนว ทสฺสเนน นภาวนาย ปหาตพฺพา, เนว ทสฺสเนน นภาวนาย ปหาตพฺพเหตุกา, อาจยคามิโน, เนวเสกฺขนาเสกฺขา, ปริตฺตา, ปริตฺตารมฺมณา, มชฺฌิมา, อนิยตา, น วตฺตพฺพา “มคฺคารมฺมณา”ติปิ, “มคฺคเหตุกา”ติปิ, “มคฺคาธิปติโน”ติปิ, สิยา อุปฺปนฺนา, สิยา อนุปฺปนฺนา, น วตฺตพฺพา “อุปฺปาทิโน”ติ, สิยา อตีตา, สิยา อนาคตา, สิยา ปจฺจุปฺปนฺนา, ปจฺจุปฺปนฺนารมฺมณา, สิยา อชฺฌตฺตา, สิยา พหิทฺธา, สิยา อชฺฌตฺตพหิทฺธา, พหิทฺธารมฺมณา, อนิทสฺสน-อปฺปฏิฆา.}

\csromanmatikaanchor{mtk.130}
\csromantocmarkref{subparagraph}{2. ทุก}{mtk.130}
\bookmark[dest={mtk.130},level=5]{2. ทุก}
\csromanheader{6}{2. ทุก}
\proseitem{717}{น เหตู สเหตุกา, เหตุสมฺปยุตฺตา, น วตฺตพฺพา “เหตู เจว สเหตุกา จา”ติ, สเหตุกา เจว น จ เหตู, น วตฺตพฺพา “เหตู เจว เหตุสมฺปยุตฺตา จา”ติ, เหตุสมฺปยุตฺตา เจว น จ เหตู, น เหตุสเหตุกา, สปฺปจฺจยา, สงฺขตา, อนิทสฺสนา, อปฺปฏิฆา, อรูปา, โลกิยา, เกนจิ วิญฺเญยฺยา, เกนจิ น วิญฺเญยฺยา.}

\prose{โน อาสวา, สาสวา, อาสววิปฺปยุตฺตา, น วตฺตพฺพา “อาสวา เจว สาสวา จา”ติ, สาสวา เจว โน จ อาสวา, น วตฺตพฺพา “อาสวา เจว อาสวสมฺปยุตฺตา จา”ติปิ, “อาสวสมฺปยุตฺตา เจว โน จ อาสวา”ติปิ, อาสววิปฺปยุตฺตา สาสวา, โน สํโยชนา ฯเปฯ.\csromanspacer{} โน คนฺถา ฯเปฯ.\csromanspacer{} โน โอฆา ฯเปฯ.\csromanspacer{} โน โยคา ฯเปฯ.\csromanspacer{} โน นีวรณา ฯเปฯ.\csromanspacer{} โน ปรามาสา ฯเปฯ สารมฺมณา, โน จิตฺตา, เจตสิกา, จิตฺตสมฺปยุตฺตา, \csromanfolio{306}จิตฺตสํสฏฺฐา, จิตฺตสมุฏฺฐานา, จิตฺตสหภุโน, จิตฺตานุปริวตฺติโน, จิตฺตสํสฏฺฐสมุฏฺฐานา, จิตฺตสํสฏฺฐสมุฏฺฐานสหภุโน, จิตฺตสํสฏฺฐสมุฏฺฐานานุปริวตฺติโน, พาหิรา, โน อุปาทา, อนุปาทินฺนา, โน อุปาทานา ฯเปฯ.\csromanspacer{} โน กิเลสา ฯเปฯ.}

\prose{น ทสฺสเนน ปหาตพฺพา, น ภาวนายปหาตพฺพา, น ทสฺสเนน ปหาตพฺพเหตุกา, น ภาวนายปหาตพฺพเหตุกา, สวิตกฺกา, สวิจารา, สิยา สปฺปีติกา, สิยา อปฺปีติกา, สิยา ปีติสหคตา, สิยา น ปีติสหคตา, สิยา สุขสหคตา, สิยา น สุขสหคตา, สิยา อุเปกฺขาสหคตา, สิยา น อุเปกฺขาสหคตา, กามาวจรา, น รูปาวจรา, น อรูปาวจรา, ปริยาปนฺนา, อนิยฺยานิกา, อนิยตา, ส-อุตฺตรา, อรณาติ.}

\vspace{\tipitakaparskip}
\csromancenter{ปญฺหาปุจฺฉกํ.}
\nitthitam{\textbf{สิกฺขาปทวิภงฺโค นิฏฺฐิโต.}}
\csromanmatikaanchor{mtk.131}
\csromantocmarknopage{subparagraph}{15. ปฏิสมฺภิทาวิภงฺค}{mtk.131}
\bookmark[dest={mtk.131},level=5]{15. ปฏิสมฺภิทาวิภงฺค}
\csromanheader{6}{15. ปฏิสมฺภิทาวิภงฺค}
\csromanheader{2}{1. สุตฺตนฺตภาชนีย}
\csromanmatikaanchor{mtk.132}
\csromanmatikaanchor{mtk.133}
\csromantocmarknopage{subparagraph}{1. สุตฺตนฺตภาชนีย}{mtk.132}
\bookmark[dest={mtk.132},level=5]{1. สุตฺตนฺตภาชนีย}
\csromantocmarkref{subparagraph}{1. สงฺคหวาร}{mtk.133}
\bookmark[dest={mtk.133},level=5]{1. สงฺคหวาร}
\csromanheader{6}{1. สงฺคหวาร}
\proseitem{718}{\csromanfolio{307}จตสฺโส ปฏิสมฺภิทา, อตฺถปฏิสมฺภิทา ธมฺมปฏิสมฺภิทา นิรุตฺติปฏิสมฺภิทา ปฏิภานปฏิสมฺภิทา\footnote{ปฏิภาณปฏิสมฺภิทา (สฺยา) เอวมุปริปิ.}.\csromanspacer{} อตฺเถ ญาณํ อตฺถปฏิสมฺภิทา, ธมฺเม ญาณํ ธมฺมปฏิสมฺภิทา, ตตฺร ธมฺมนิรุตฺตาภิลาเป ญาณํ นิรุตฺติปฏิสมฺภิทา, ญาเณสุ ญาณํ ปฏิภานปฏิสมฺภิทา.\csromanspacer{}\csromanspacer{}อยํ สงฺคหวาโร.}

\csromanmatikaanchor{mtk.134}
\csromantocmarkref{subparagraph}{2. สจฺจวาร}{mtk.134}
\bookmark[dest={mtk.134},level=5]{2. สจฺจวาร}
\csromanheader{6}{2. สจฺจวาร}
\proseitem{719}{จตสฺโส ปฏิสมฺภิทา, อตฺถปฏิสมฺภิทา ธมฺมปฏิสมฺภิทา นิรุตฺติปฏิสมฺภิทา ปฏิภานปฏิสมฺภิทา.\csromanspacer{} ทุกฺเข ญาณํ อตฺถปฏิสมฺภิทา, ทุกฺขสมุทเย ญาณํ ธมฺมปฏิสมฺภิทา, ทุกฺขนิโรเธ ญาณํ อตฺถปฏิสมฺภิทา, ทุกฺขนิโรธคามินิยา ปฏิปทาย ญาณํ ธมฺมปฏิสมฺภิทา, ตตฺร ธมฺมนิรุตฺตาภิลาเป ญาณํ นิรุตฺติปฏิสมฺภิทา, ญาเณสุ ญาณํ ปฏิภานปฏิสมฺภิทา.\csromanspacer{}\csromanspacer{}อยํ สจฺจวาโร.}

\csromanmatikaanchor{mtk.135}
\csromantocmarkref{subparagraph}{3. เหตุวาร}{mtk.135}
\bookmark[dest={mtk.135},level=5]{3. เหตุวาร}
\csromanheader{6}{3. เหตุวาร}
\proseitem{720}{จตสฺโส ปฏิสมฺภิทา, อตฺถปฏิสมฺภิทา ธมฺมปฏิสมฺภิทา นิรุตฺติปฏิสมฺภิทา ปฏิภานปฏิสมฺภิทา.\csromanspacer{} เหตุมฺหิ ญาณํ ธมฺมปฏิสมฺภิทา, เหตุผเล ญาณํ อตฺถปฏิสมฺภิทา, ตตฺร ธมฺมนิรุตฺตาภิลาเป ญาณํ นิรุตฺติปฏิสมฺภิทา, ญาเณสุ ญาณํ ปฏิภานปฏิสมฺภิทา.\csromanspacer{}\csromanspacer{}อยํ เหตุวาโร.}

\csromanmatikaanchor{mtk.136}
\csromantocmarkref{subparagraph}{4. ธมฺมวาร}{mtk.136}
\bookmark[dest={mtk.136},level=5]{4. ธมฺมวาร}
\csromanheader{6}{4. ธมฺมวาร}
\proseitem{721}{จตสฺโส ปฏิสมฺภิทา, อตฺถปฏิสมฺภิทา ธมฺมปฏิสมฺภิทา นิรุตฺติปฏิสมฺภิทา ปฏิภานปฏิสมฺภิทา.\csromanspacer{} เย ธมฺมา ชาตา ภูตา สญฺชาตา นิพฺพตฺตา อภินิพฺพตฺตา ปาตุภูตา, อิเมสุ ธมฺเมสุ ญาณํ อตฺถปฏิสมฺภิทา, ยมฺหา ธมฺมา เต ธมฺมา ชาตา ภูตา สญฺชาตา นิพฺพตฺตา อภินิพฺพตฺตา ปาตุภูตา, \csromanfolio{308}เตสุ ธมฺเมสุ ญาณํ ธมฺมปฏิสมฺภิทา, ตตฺร ธมฺมนิรุตฺตาภิลาเป ญาณํ นิรุตฺติปฏิสมฺภิทา, ญาเณสุ ญาณํ ปฏิภานปฏิสมฺภิทา.\csromanspacer{}\csromanspacer{}อยํ ธมฺมวาโร.}

\csromanmatikaanchor{mtk.137}
\csromantocmarkref{subparagraph}{5. ปฏิจฺจสมุปฺปาทวาร}{mtk.137}
\bookmark[dest={mtk.137},level=5]{5. ปฏิจฺจสมุปฺปาทวาร}
\csromanheader{6}{5. ปฏิจฺจสมุปฺปาทวาร}
\proseitem{722}{จตสฺโส ปฏิสมฺภิทา, อตฺถปฏิสมฺภิทา ธมฺมปฏิสมฺภิทา นิรุตฺติปฏิสมฺภิทา ปฏิภานปฏิสมฺภิทา.\csromanspacer{} ชรามรเณ ญาณํ อตฺถปฏิสมฺภิทา, ชรามรณสมุทเย ญาณํ ธมฺมปฏิสมฺภิทา, ชรามรณนิโรเธ ญาณํ อตฺถปฏิสมฺภิทา, ชรามรณนิโรธคามินิยา ปฏิปทาย ญาณํ ธมฺมปฏิสมฺภิทา, ตตฺร ธมฺมนิรุตฺตาภิลาเป ญาณํ นิรุตฺติปฏิสมฺภิทา, ญาเณสุ ญาณํ ปฏิภานปฏิสมฺภิทา.}

\proseitem{723}{จตสฺโส ปฏิสมฺภิทา, อตฺถปฏิสมฺภิทา ธมฺมปฏิสมฺภิทา นิรุตฺติปฏิสมฺภิทา ปฏิภานปฏิสมฺภิทา.\csromanspacer{} ชาติยา ญาณํ ฯเปฯ ภเว ญาณํ ฯเปฯ อุปาทาเน ญาณํ ฯเปฯ ตณฺหาย ญาณํ ฯเปฯ เวทนาย ญาณํ ฯเปฯ ผสฺเส ญาณํ ฯเปฯ สฬายตเน ญาณํ ฯเปฯ นามรูเป ญาณํ ฯเปฯ วิญฺญาเณ ญาณํ ฯเปฯ สงฺขาเรสุ ญาณํ อตฺถปฏิสมฺภิทา, สงฺขารสมุทเย ญาณํ ธมฺมปฏิสมฺภิทา, สงฺขารนิโรเธ ญาณํ อตฺถปฏิสมฺภิทา, สงฺขารนิโรธคามินิยา ปฏิปทาย ญาณํ ธมฺมปฏิสมฺภิทา, ตตฺร ธมฺมนิรุตฺตาภิลาเป ญาณํ นิรุตฺติปฏิสมฺภิทา, ญาเณสุ ญาณํ ปฏิภานปฏิสมฺภิทา.\csromanspacer{}\csromanspacer{}อยํ ปฏิจฺจสมุปฺปาทวาโร.}

\csromanmatikaanchor{mtk.138}
\csromantocmarkref{subparagraph}{6. ปริยตฺติวาร}{mtk.138}
\bookmark[dest={mtk.138},level=5]{6. ปริยตฺติวาร}
\csromanheader{6}{6. ปริยตฺติวาร}
\proseitem{724}{จตสฺโส ปฏิสมฺภิทา, อตฺถปฏิสมฺภิทา ธมฺมปฏิสมฺภิทา นิรุตฺติปฏิสมฺภิทา ปฏิภานปฏิสมฺภิทา.}

\prose{ตตฺถ กตมา ธมฺมปฏิสมฺภิทา, อิธ ภิกฺขุ ธมฺมํ ชานาติ, สุตฺตํ, เคยฺยํ, เวยฺยากรณํ, คาถํ, อุทานํ, อิติวุตฺตกํ, ชาตกํ, อพฺภุตธมฺมํ, เวทลฺลํ, อยํ วุจฺจติ ธมฺมปฏิสมฺภิทา.\csromanspacer{} โส ตสฺส ตสฺเสว ภาสิตสฺส อตฺถํ ชานาติ “อยํ อิมสฺส ภาสิตสฺส อตฺโถ, อยํ อิมสฺส ภาสิตสฺส อตฺโถ”ติ, อยํ วุจฺจติ อตฺถปฏิสมฺภิทา, ตตฺร ธมฺมนิรุตฺตาภิลาเป ญาณํ นิรุตฺติปฏิสมฺภิทา, ญาเณสุ ญาณํ ปฏิภานปฏิสมฺภิทา.\csromanspacer{}\csromanspacer{}อยํ ปริยตฺติวาโร.}

\vspace{\tipitakaparskip}
\csromancenter{สุตฺตนฺตภาชนียํ.}
\csromanheader{2}{15. ปฏิสมฺภิทาวิภงฺค}
\csromanheader{2}{2. อภิธมฺมภาชนีย}
\csromanmatikaanchor{mtk.139}
\csromanmatikaanchor{mtk.140}
\csromantocmarknopage{subparagraph}{2. อภิธมฺมภาชนีย}{mtk.139}
\bookmark[dest={mtk.139},level=5]{2. อภิธมฺมภาชนีย}
\csromantocmarkref{subparagraph}{1. กุสลวาร}{mtk.140}
\bookmark[dest={mtk.140},level=5]{1. กุสลวาร}
\csromanheader{6}{1. กุสลวาร}
\proseitem{725}{\csromanfolio{309}จตสฺโส ปฏิสมฺภิทา, อตฺถปฏิสมฺภิทา ธมฺมปฏิสมฺภิทา นิรุตฺติปฏิสมฺภิทา ปฏิภานปฏิสมฺภิทา.}

\prose{กตเม ธมฺมา กุสลา, ยสฺมิํ สมเย กามาวจรํ กุสลํ จิตฺตํ อุปฺปนฺนํ โหติ โสมนสฺสสหคตํ ญาณสมฺปยุตฺตํ รูปารมฺมณํ วา ฯเปฯ ธมฺมารมฺมณํ วา ยํ ยํ วา ปนารพฺภ, ตสฺมิํ สมเย ผสฺโส โหติ ฯเปฯ อวิกฺเขโป โหติ, อิเม ธมฺมา กุสลา.\csromanspacer{} อิเมสุ ธมฺเมสุ ญาณํ ธมฺมปฏิสมฺภิทา, เตสํ วิปาเก ญาณํ อตฺถปฏิสมฺภิทา, ยาย นิรุตฺติยา เตสํ ธมฺมานํ ปญฺญตฺติ โหติ, ตตฺร ธมฺมนิรุตฺตาภิลาเป ญาณํ นิรุตฺติปฏิสมฺภิทา, เยน ญาเณน ตานิ ญาณานิ ชานาติ “อิมานิ ญาณานิ อิทมตฺถโชตกานี”ติ ญาเณสุ ญาณํ ปฏิภานปฏิสมฺภิทา.}

\proseitem{726}{จตสฺโส ปฏิสมฺภิทา, อตฺถปฏิสมฺภิทา ธมฺมปฏิสมฺภิทา นิรุตฺติปฏิสมฺภิทา ปฏิภานปฏิสมฺภิทา.}

\prose{กตเม ธมฺมา กุสลา, ยสฺมิํ สมเย กามาวจรํ กุสลํ จิตฺตํ อุปฺปนฺนํ โหติ โสมนสฺสสหคตํ ญาณสมฺปยุตฺตํ สสงฺขาเรน ฯเปฯ โสมนสฺสสหคตํ ญาณวิปฺปยุตฺตํ ฯเปฯ โสมนสฺสสหคตํ ญาณวิปฺปยุตฺตํ สสงฺขาเรน ฯเปฯ อุเปกฺขาสหคตํ ญาณสมฺปยุตฺตํ ฯเปฯ อุเปกฺขาสหคตํ ญาณสมฺปยุตฺตํ สสงฺขาเรน ฯเปฯ อุเปกฺขาสหคตํ ญาณวิปฺปยุตฺตํ ฯเปฯ อุเปกฺขาสหคตํ ญาณวิปฺปยุตฺตํ สสงฺขาเรน รูปารมฺมณํ วา ฯเปฯ ธมฺมารมฺมณํ วา ยํ ยํ วา ปนารพฺภ, ตสฺมิํ สมเย ผสฺโส โหติ ฯเปฯ อวิกฺเขโป โหติ.\csromanspacer{} อิเม ธมฺมา กุสลา.\csromanspacer{} อิเมสุ ธมฺเมสุ ญาณํ ธมฺมปฏิสมฺภิทา, เตสํ วิปาเก ญาณํ อตฺถปฏิสมฺภิทา, ยาย นิรุตฺติยา เตสํ ธมฺมานํ ปญฺญตฺติ โหติ, ตตฺร ธมฺมนิรุตฺตาภิลาเป ญาณํ นิรุตฺติปฏิสมฺภิทา, เยน ญาเณน ตานิ ญาณานิ ชานาติ “อิมานิ ญาณานิ อิทมตฺถโชตกานี”ติ ญาเณสุ ญาณํ ปฏิภานปฏิสมฺภิทา.}

\proseitem{727}{จตสฺโส ปฏิสมฺภิทา, อตฺถปฏิสมฺภิทา ธมฺมปฏิสมฺภิทา นิรุตฺติปฏิสมฺภิทา ปฏิภานปฏิสมฺภิทา.}

\prose{\csromanfolio{310}กตเม ธมฺมา กุสลา, ยสฺมิํ สมเย รูปูปปตฺติยา มคฺคํ ภาเวติ วิวิจฺเจว กาเมหิ ฯเปฯ ปฐมํ ฌานํ อุปสมฺปชฺช วิหรติ ปถวีกสิณํ, ตสฺมิํ สมเย ผสฺโส โหติ ฯเปฯ อวิกฺเขโป โหติ, อิเม ธมฺมา กุสลา.\csromanspacer{} อิเมสุ ธมฺเมสุ ญาณํ ธมฺมปฏิสมฺภิทา, เตสํ วิปาเก ญาณํ อตฺถปฏิสมฺภิทา, ยาย นิรุตฺติยา เตสํ ธมฺมานํ ปญฺญตฺติ โหติ, ตตฺร ธมฺมนิรุตฺตาภิลาเป ญาณํ นิรุตฺติปฏิสมฺภิทา, เยน ญาเณน ตานิ ญาณานิ ชานาติ “อิมานิ ญาณานิ อิทมตฺถโชตกานี”ติ ญาเณสุ ญาณํ ปฏิภานปฏิสมฺภิทา.}

\proseitem{728}{จตสฺโส ปฏิสมฺภิทา, อตฺถปฏิสมฺภิทา ธมฺมปฏิสมฺภิทา นิรุตฺติปฏิสมฺภิทา ปฏิภานปฏิสมฺภิทา.}

\prose{กตเม ธมฺมา กุสลา, ยสฺมิํ สมเย อรูปูปปตฺติยา มคฺคํ ภาเวติ สพฺพโส อากิญฺจญฺญายตนํ สมติกฺกมฺม เนวสญฺญานาสญฺญายตนสญฺญาสหคตํ สุขสฺส จ ปหานา ฯเปฯ จตุตฺถํ ฌานํ อุปสมฺปชฺช วิหรติ, ตสฺมิํ สมเย ผสฺโส โหติ ฯเปฯ อวิกฺเขโป โหติ.\csromanspacer{} อิเม ธมฺมา กุสลา.\csromanspacer{} อิเมสุ ธมฺเมสุ ญาณํ ธมฺมปฏิสมฺภิทา, เตสํ วิปาเก ญาณํ อตฺถปฏิสมฺภิทา, ยาย นิรุตฺติยา เตสํ ธมฺมานํ ปญฺญตฺติ โหติ ตตฺร ธมฺมนิรุตฺตาภิลาเป ญาณํ นิรุตฺติปฏิสมฺภิทา, เยน ญาเณน ตานิ ญาณานิ ชานาติ “อิมานิ ญาณานิ อิทมตฺถโชตกานี”ติ ญาเณสุ ญาณํ ปฏิภานปฏิสมฺภิทา.}

\proseitem{729}{จตสฺโส ปฏิสมฺภิทา, อตฺถปฏิสมฺภิทา ธมฺมปฏิสมฺภิทา นิรุตฺติ ปฏิสมฺภิทา ปฏิภานปฏิสมฺภิทา.}

\prose{กตเม ธมฺมา กุสลา, ยสฺมิํ สมเย โลกุตฺตรํ ฌานํ ภาเวติ นิยฺยานิกํ อปจยคามิํ ทิฏฺฐิคตานํ ปหานาย ปฐมาย ภูมิยา ปตฺติยา วิวิจฺเจว กาเมหิ ฯเปฯ ปฐมํ ฌานํ อุปสมฺปชฺช วิหรติ ทุกฺขปฏิปทํ ทนฺธาภิญฺญํ, ตสฺมิํ สมเย ผสฺโส โหติ ฯเปฯ อวิกฺเขโป โหติ, อิเม ธมฺมา กุสลา.\csromanspacer{} อิเมสุ ธมฺเมสุ ญาณํ ธมฺมปฏิสมฺภิทา, เตสํ วิปาเก ญาณํ อตฺถปฏิสมฺภิทา, ยาย นิรุตฺติยา เตสํ ธมฺมานํ ปญฺญตฺติ โหติ, ตตฺร ธมฺมนิรุตฺตาภิลาเป ญาณํ นิรุตฺติปฏิสมฺภิทา, เยน ญาเณน ตานิ ญาณานิ ชานาติ “อิมานิ ญาณานิ อิทมตฺถโชตกานี”ติ ญาเณสุ ญาณํ ปฏิภานปฏิสมฺภิทา.}

\csromanheader{2}{15. ปฏิสมฺภิทาวิภงฺค}
\csromanmatikaanchor{mtk.141}
\csromantocmarkref{subparagraph}{2. อกุสลวาร}{mtk.141}
\bookmark[dest={mtk.141},level=5]{2. อกุสลวาร}
\csromanheader{6}{2. อกุสลวาร}
\proseitem{730}{\csromanfolio{311}จตสฺโส ปฏิสมฺภิทา, อตฺถปฏิสมฺภิทา ธมฺมปฏิสมฺภิทา นิรุตฺติปฏิสมฺภิทา ปฏิภานปฏิสมฺภิทา.}

\prose{กตเม ธมฺมา อกุสลา, ยสฺมิํ สมเย อกุสลํ จิตฺตํ อุปฺปนฺนํ โหติ โสมนสฺสสหคตํ ทิฏฺฐิคตสมฺปยุตฺตํ รูปารมฺมณํ วา ฯเปฯ ธมฺมารมฺมณํ วา ยํ ยํ วา ปนารพฺภ, ตสฺมิํ สมเย ผสฺโส โหติ ฯเปฯ อวิกฺเขโป โหติ, อิเม ธมฺมา อกุสลา.\csromanspacer{} อิเมสุ ธมฺเมสุ ญาณํ ธมฺมปฏิสมฺภิทา, เตสํ วิปาเก ญาณํ อตฺถปฏิสมฺภิทา, ยาย นิรุตฺติยา เตสํ ธมฺมานํ ปญฺญตฺติ โหติ, ตตฺร ธมฺมนิรุตฺตาภิลาเป ญาณํ นิรุตฺติปฏิสมฺภิทา, เยน ญาเณน ตานิ ญาณานิ ชานาติ “อิมานิ ญาณานิ อิทมตฺถโชตกานี”ติ ญาเณสุ ญาณํ ปฏิภานปฏิสมฺภิทา.}

\proseitem{731}{จตสฺโส ปฏิสมฺภิทา, อตฺถปฏิสมฺภิทา ธมฺมปฏิสมฺภิทา นิรุตฺติปฏิสมฺภิทา ปฏิภานปฏิสมฺภิทา.}

\prose{กตเม ธมฺมา อกุสลา, ยสฺมิํ สมเย อกุสลํ จิตฺตํ อุปฺปนฺนํ โหติ โสมนสฺสสหคตํ ทิฏฺฐิคตสมฺปยุตฺตํ สสงฺขาเรน ฯเปฯ โสมนสฺสสหคตํ ทิฏฺฐิคตวิปฺปยุตฺตํ ฯเปฯ โสมนสฺสสหคตํ ทิฏฺฐิคตวิปฺปยุตฺตํ สสงฺขาเรน ฯเปฯ อุเปกฺขาสหคตํ ทิฏฺฐิคตสมฺปยุตฺตํ ฯเปฯ อุเปกฺขาสหคตํ ทิฏฺฐิคตสมฺปยุตฺตํ สสงฺขาเรน ฯเปฯ อุเปกฺขาสหคตํ ทิฏฺฐิคตวิปฺปยุตฺตํ ฯเปฯ อุเปกฺขาสหคตํ ทิฏฺฐิคตวิปฺปยุตฺตํ สสงฺขาเรน ฯเปฯ โทมนสฺสสหคตํ ปฏิฆสมฺปยุตฺตํ ฯเปฯ โทมนสฺสสหคตํ ปฏิฆสมฺปยุตฺตํ สสงฺขาเรน ฯเปฯ อุเปกฺขาสหคตํ วิจิกิจฺฉาสมฺปยุตฺตํ ฯเปฯ อุเปกฺขาสหคตํ อุทฺธจฺจสมฺปยุตฺตํ รูปารมฺมณํ วา ฯเปฯ ธมฺมารมฺมณํ วา ยํ ยํ วา ปนารพฺภ, ตสฺมิํ สมเย ผสฺโส โหติ ฯเปฯ อวิกฺเขโป โหติ, อิเม ธมฺมา อกุสลา.\csromanspacer{} อิเมสุ ธมฺเมสุ ญาณํ ธมฺมปฏิสมฺภิทา, เตสํ วิปาเก ญาณํ อตฺถปฏิสมฺภิทา, ยาย นิรุตฺติยา เตสํ ธมฺมานํ ปญฺญตฺติ โหติ, ตตฺร ธมฺมนิรุตฺตาภิลาเป ญาณํ นิรุตฺติปฏิสมฺภิทา, เยน ญาเณน ตานิ ญาณานิ ชานาติ “อิมานิ ญาณานิ อิทมตฺตโชตกานี”ติ ญาเณสุ ญาณํ ปฏิภานปฏิสมฺภิทา.}

\csromanmatikaanchor{mtk.142}
\csromantocmarkref{subparagraph}{3. วิปากวาร}{mtk.142}
\bookmark[dest={mtk.142},level=5]{3. วิปากวาร}
\csromanheader{6}{3. วิปากวาร}
\proseitem{732}{\csromanfolio{312}ติสฺโส ปฏิสมฺภิทา, อตฺถปฏิสมฺภิทา นิรุตฺติปฏิสมฺภิทา ปฏิภานปฏิสมฺภิทา.}

\prose{กตเม ธมฺมา อพฺยากตา, ยสฺมิํ สมเย กามาวจรสฺส กุสลสฺส กมฺมสฺส กตตฺตา อุปจิตตฺตา วิปากํ จกฺขุวิญฺญาณํ อุปฺปนฺนํ โหติ อุเปกฺขาสหคตํ รูปารมฺมณํ, ตสฺมิํ สมเย ผสฺโส โหติ, เวทนา โหติ, สญฺญา โหติ, เจตนา โหติ, จิตฺตํ โหติ, อุเปกฺขา โหติ, จิตฺตสฺเสกคฺคตา โหติ, มนินฺทฺริยํ โหติ, อุเปกฺขินฺทฺริยํ โหติ, ชีวิตินฺทฺริยํ โหติ, เย วา ปน ตสฺมิํ สมเย อญฺเญปิ อตฺถิ ปฏิจฺจสมุปฺปนฺนา อรูปิโน ธมฺมา, อิเม ธมฺมา อพฺยากตา.\csromanspacer{} อิเมสุ ธมฺเมสุ ญาณํ อตฺถปฏิสมฺภิทา, ยาย นิรุตฺติยา เตสํ ธมฺมานํ ปญฺญตฺติ โหติ, ตตฺร ธมฺมนิรุตฺตาภิลาเป ญาณํ นิรุตฺติปฏิสมฺภิทา, เยน ญาเณน ตานิ ญาณานิ ชานาติ “อิมานิ ญาณานิ อิทมตฺถโชตกานี”ติ ญาเณสุ ญาณํ ปฏิภานปฏิสมฺภิทา.}

\proseitem{733}{ติสฺโส ปฏิสมฺภิทา, อตฺถปฏิสมฺภิทา นิรุตฺติปฏิสมฺภิทา ปฏิภานปฏิสมฺภิทา.}

\prose{กตเม ธมฺมา อพฺยากตา, ยสฺมิํ สมเย กามาวจรสฺส กุสลสฺส กมฺมสฺส กตตฺตา อุปจิตตฺตา วิปากํ โสตวิญฺญาณํ อุปฺปนฺนํ โหติ อุเปกฺขาสหคตํ สทฺทารมฺมณํ ฯเปฯ ฆานวิญฺญาณํ อุปฺปนฺนํ โหติ อุเปกฺขาสหคตํ คนฺธารมฺมณํ ฯเปฯ ชิวฺหาวิญฺญาณํ อุปฺปนฺนํ โหติ อุเปกฺขาสหคตํ รสารมฺมณํ ฯเปฯ กายวิญฺญาณํ อุปฺปนฺนํ โหติ สุขสหคตํ โผฏฺฐพฺพารมฺมณํ, ตสฺมิํ สมเย ผสฺโส โหติ, เวทนา โหติ, สญฺญา โหติ, เจตนา โหติ, จิตฺตํ โหติ, สุขํ โหติ, จิตฺตสฺเสกคฺคตา โหติ, มนินฺทฺริยํ โหติ, สุขินฺทฺริยํ โหติ, ชีวิตินฺทฺริยํ โหติ, เย วา ปน ตสฺมิํ สมเย อญฺเญปิ อตฺถิ ปฏิจฺจสมุปฺปนฺนา อรูปิโน ธมฺมา, อิเม ธมฺมา อพฺยากตา.\csromanspacer{} อิเมสุ ธมฺเมสุ ญาณํ อตฺถปฏิสมฺภิทา, ยาย นิรุตฺติยา เตสํ ธมฺมานํ ปญฺญตฺติ โหติ ตตฺร ธมฺมนิรุตฺตาภิลาเป ญาณํ นิรุตฺติปฏิสมฺภิทา, เยน ญาเณน ตานิ ญาณานิ ชานาติ “อิมานิ ญาณานิ อิทมตฺถโชตกานี”ติ ญาเณสุ ญาณํ ปฏิภานปฏิสมฺภิทา.}

\csromanheader{2}{15. ปฏิสมฺภิทาวิภงฺค}
\proseitem{734}{\csromanfolio{313}ติสฺโส ปฏิสมฺภิทา ฯเปฯ ปฏิภานปฏิสมฺภิทา.\csromanspacer{} กตเม ธมฺมา อพฺยากตา, ยสฺมิํ สมเย กามาวจรสฺส กุสลสฺส กมฺมสฺส กตตฺตา อุปจิตตฺตา วิปากา มโนธาตุ อุปฺปนฺนา โหติ อุเปกฺขาสหคตา รูปารมฺมณา วา ฯเปฯ โผฏฺฐพฺพารมฺมณา วา ยํ ยํ วา ปนารพฺภ, ตสฺมิํ สมเย ผสฺโส โหติ, เวทนา โหติ, สญฺญา โหติ, เจตนา โหติ, จิตฺตํ โหติ, วิตกฺโก โหติ, วิจาโร โหติ, อุเปกฺขา โหติ, จิตฺตสฺเสกคฺคตา โหติ, มนินฺทฺริยํ โหติ, อุเปกฺขินฺทฺริยํ โหติ, ชีวิตินฺทฺริยํ โหติ, เย วา ปน ตสฺมิํ สมเย อญฺเญปิ อตฺถิ ปฏิจฺจสมุปฺปนฺนา อรูปิโน ธมฺมา, อิเม ธมฺมา อพฺยากตา.\csromanspacer{} อิเมสุ ธมฺเมสุ ญาณํ อตฺถปฏิสมฺภิทา, ยาย นิรุตฺติยา เตสํ ธมฺมานํ ปญฺญตฺติ โหติ, ตตฺร ธมฺมนิรุตฺตาภิลาเป ญาณํ นิรุตฺติปฏิสมฺภิทา, เยน ญาเณน ตานิ ญาณานิ ชานาติ “อิมานิ ญาณานิ อิทมตฺถโชตกานี”ติ ญาเณสุ ญาณํ ปฏิภานปฏิสมฺภิทา.}

\proseitem{735}{ติสฺโส ปฏิสมฺภิทา, อตฺถปฏิสมฺภิทา นิรุตฺติปฏิสมฺภิทา ปฏิภานปฏิสมฺภิทา.}

\prose{กตเม ธมฺมา อพฺยากตา, ยสฺมิํ สมเย กามาวจรสฺส กุสลสฺส กมฺมสฺส กตตฺตา อุปจิตตฺตา วิปากา มโนวิญฺญาณธาตุ อุปฺปนฺนา โหติ โสมนสฺสสหคตา รูปารมฺมณา วา ฯเปฯ ธมฺมารมฺมณา วา ยํ ยํ วา ปนารพฺภ, ตสฺมิํ สมเย ผสฺโส โหติ, เวทนา โหติ, สญฺญา โหติ, เจตนา โหติ, จิตฺตํ โหติ, วิตกฺโก โหติ, วิจาโร โหติ, ปีติ โหติ, สุขํ โหติ, จิตฺตสฺเสกคฺคตา โหติ, มนินฺทฺริยํ โหติ, โสมนสฺสินฺทฺริยํ โหติ, ชีวิตินฺทฺริยํ โหติ, เย วา ปน ตสฺมิํ สมเย อญฺเญปิ อตฺถิ ปฏิจฺจสมุปฺปนฺนา อรูปิโน ธมฺมา, อิเม ธมฺมา อพฺยากตา.\csromanspacer{} อิเมสุ ธมฺเมสุ ญาณํ อตฺถปฏิสมฺภิทา, ยาย นิรุตฺติยา เตสํ ธมฺมานํ ปญฺญตฺติ โหติ, ตตฺร ธมฺมนิรุตฺตาภิลาเป ญาณํ นิรุตฺติปฏิสมฺภิทา, เยน ญาเณน ตานิ ญาณานิ ชานาติ “อิมานิ ญาณานิ อิทมตฺถโชตกานี”ติ ญาเณสุ ญาณํ ปฏิภานปฏิสมฺภิทา.}

\proseitem{736}{ติสฺโส ปฏิสมฺภิทา, อตฺถปฏิสมฺภิทา นิรุตฺติปฏิสมฺภิทา ปฏิภานปฏิสมฺภิทา.}

\prose{กตเม ธมฺมา อพฺยากตา, ยสฺมิํ สมเย กามาวจรสฺส กุสลสฺส กมฺมสฺส กตตฺตา อุปจิตตฺตา วิปากา มโนวิญฺญาณธาตุ \csromanfolio{314}อุปฺปนฺนา โหติ อุเปกฺขาสหคตา รูปารมฺมณา วา ฯเปฯ ธมฺมารมฺมณา วา ยํ ยํ วา ปนารพฺภ, ตสฺมิํ สมเย ผสฺโส โหติ, เวทนา โหติ, สญฺญา โหติ, เจตนา โหติ, จิตฺตํ โหติ, วิตกฺโก โหติ, วิจาโร โหติ, อุเปกฺขา โหติ, จิตฺตสฺเสกคฺคตา โหติ, มนินฺทฺริยํ โหติ, อุเปกฺขินฺทฺริยํ โหติ, เย วา ปน ตสฺมิํ สมเย อญฺเญปิ อตฺถิ ปฏิจฺจสมุปฺปนฺนา อรูปิโน ธมฺมา, อิเม ธมฺมา อพฺยากตา.\csromanspacer{} อิเมสุ ธมฺเมสุ ญาณํ อตฺถปฏิสมฺภิทา, ยาย นิรุตฺติยา เตสํ ธมฺมานํ ปญฺญตฺติ โหติ, ตตฺร ธมฺมนิรุตฺตาภิลาเป ญาณํ นิรุตฺติปฏิสมฺภิทา, เยน ญาเณน ตานิ ญาณานิ ชานาติ “อิมานิ ญาณานิ อิทมตฺถโชตกานี”ติ ญาเณสุ ญาณํ ปฏิภานปฏิสมฺภิทา.}

\proseitem{737}{ติสฺโส ปฏิสมฺภิทา, อตฺถปฏิสมฺภิทา นิรุตฺติปฏิสมฺภิทา ปฏิภานปฏิสมฺภิทา.}

\prose{กตเม ธมฺมา อพฺยากตา, ยสฺมิํ สมเย กามาวจรสฺส กุสลสฺส กมฺมสฺส กตตฺตา อุปจิตตฺตา วิปากา มโนวิญฺญาณธาตุ อุปฺปนฺนา โหติ โสมนสฺสสหคตา ญาณสมฺปยุตฺตา ฯเปฯ โสมนสฺสสหคตา ญาณสมฺปยุตฺตา สสงฺขาเรน ฯเปฯ โสมนสฺสสหคตา ญาณวิปฺปยุตฺตา ฯเปฯ โสมนสฺสสหคตา ญาณวิปฺปยุตฺตา สสงฺขาเรน ฯเปฯ อุเปกฺขาสหคตา ญาณสมฺปยุตฺตา ฯเปฯ อุเปกฺขาสหคตา ญาณสมฺปยุตฺตา สสงฺขาเรน ฯเปฯ อุเปกฺขาสหคตา ญาณวิปฺปยุตฺตา ฯเปฯ อุเปกฺขาสหคตา ญาณวิปฺปยุตฺตา สสงฺขาเรน รูปารมฺมณา วา ฯเปฯ ธมฺมารมฺมณา วา ยํ ยํ วา ปนารพฺภ, ตสฺมิํ สมเย ผสฺโส โหติ ฯเปฯ อวิกฺเขโป โหติ, อิเม ธมฺมา อพฺยากตา.\csromanspacer{} อิเมสุ ธมฺเมสุ ญาณํ อตฺถปฏิสมฺภิทา, ยาย นิรุตฺติยา เตสํ ธมฺมานํ ปญฺญตฺติ โหติ, ตตฺร ธมฺมนิรุตฺตาภิลาเป ญาณํ นิรุตฺติปฏิสมฺภิทา, เยน ญาเณน ตานิ ญาณานิ ชานาติ “อิมานิ ญาณานิ อิทมตฺถโชตกานี”ติ ญาเณสุ ญาณํ ปฏิภานปฏิสมฺภิทา.}

\proseitem{738}{ติสฺโส ปฏิสมฺภิทา, อตฺถปฏิสมฺภิทา นิรุตฺติปฏิสมฺภิทา ปฏิภานปฏิสมฺภิทา.}

\prose{กตเม ธมฺมา อพฺยากตา, ยสฺมิํ สมเย รูปูปปตฺติยา มคฺคํ ภาเวติ วิวิจฺเจว กาเมหิ ฯเปฯ ปฐมํ ฌานํ อุปสมฺปชฺช วิหรติ ปถวีกสิณํ, ตสฺมิํ สมเย ผสฺโส โหติ ฯเปฯ อวิกฺเขโป โหติ, อิเม ธมฺมา}

\csromanheader{2}{15. ปฏิสมฺภิทาวิภงฺค}
\prose{\csromanfolio{315}กุสลา.\csromanspacer{} ตสฺเสว รูปาวจรสฺส กุสลสฺส กมฺมสฺส กตตฺตา อุปจิตตฺตา วิปากํ วิวิจฺเจว กาเมหิ ฯเปฯ ปฐมํ ฌานํ อุปสมฺปชฺช วิหรติ ปถวีกสิณํ, ตสฺมิํ สมเย ผสฺโส โหติ ฯเปฯ อวิกฺเขโป โหติ, อิเม ธมฺมา อพฺยากตา.\csromanspacer{} อิเมสุ ธมฺเมสุ ญาณํ อตฺถปฏิสมฺภิทา, ยาย นิรุตฺติยา เตสํ ธมฺมานํ ปญฺญตฺติ โหติ, ตตฺร ธมฺมนิรุตฺตาภิลาเป ญาณํ นิรุตฺติปฏิสมฺภิทา, เยน ญาเณน ตานิ ญาณานิ ชานาติ “อิมานิ ญาณานิ อิทมตฺถโชตกานี”ติ ญาเณสุ ญาณํ ปฏิภานปฏิสมฺภิทา.}

\proseitem{739}{ติสฺโส ปฏิสมฺภิทา, อตฺถปฏิสมฺภิทา นิรุตฺติปฏิสมฺภิทา ปฏิภานปฏิสมฺภิทา.}

\prose{กตเม ธมฺมา อพฺยากตา, ยสฺมิํ สมเย อรูปูปปตฺติยา มคฺคํ ภาเวติ สพฺพโส อากิญฺจญฺญายตนํ สมติกฺกมฺม เนวสญฺญานาสญฺญายตนสญฺญาสหคตํ สุขสฺส จ ปหานา ฯเปฯ จตุตฺถํ ฌานํ อุปสมฺปชฺช วิหรติ, ตสฺมิํ สมเย ผสฺโส โหติ ฯเปฯ อวิกฺเขโป โหติ.\csromanspacer{} อิเม ธมฺมา กุสลา.\csromanspacer{} ตสฺเสว อรูปาวจรสฺส กุสลสฺส กมฺมสฺส กตตฺตา อุปจิตตฺตา วิปากํ สพฺพโส อากิญฺจญฺญายตนํ สมติกฺกมฺม เนวสญฺญานาสญฺญายตนสญฺญาสหคตํ สุขสฺส จ ปหานา ฯเปฯ จตุตฺถํ ฌานํ อุปสมฺปชฺช วิหรติ, ตสฺมิํ สมเย ผสฺโส โหติ ฯเปฯ อวิกฺเขโป โหติ, อิเม ธมฺมา อพฺยากตา.\csromanspacer{} อิเมสุ ธมฺเมสุ ญาณํ อตฺถปฏิสมฺภิทา, ยาย นิรุตฺติยา เตสํ ธมฺมานํ ปญฺญตฺติ โหติ, ตตฺร ธมฺมนิรุตฺตาภิลาเป ญาณํ นิรุตฺติปฏิสมฺภิทา, เยน ญาเณน ตานิ ญาณานิ ชานาติ “อิมานิ ญาณานิ อิทมตฺถโชตกานี”ติ ญาเณสุ ญาณํ ปฏิภานปฏิสมฺภิทา.}

\proseitem{740}{ติสฺโส ปฏิสมฺภิทา, อตฺถปฏิสมฺภิทา นิรุตฺติปฏิสมฺภิทา ปฏิภานปฏิสมฺภิทา.}

\prose{กตเม ธมฺมา อพฺยากตา, ยสฺมิํ สมเย โลกุตฺตรํ ฌานํ ภาเวติ นิยฺยานิกํ อปจยคามิํ ทิฏฺฐิคตานํ ปหานาย ปฐมาย ภูมิยา ปตฺติยา วิวิจฺเจว กาเมหิ ฯเปฯ ปฐมํ ฌานํ อุปสมฺปชฺช วิหรติ ทุกฺขปฏิปทํ ทนฺธาภิญฺญํ, ตสฺมิํ สมเย ผสฺโส โหติ ฯเปฯ อวิกฺเขโป โหติ, อิเม ธมฺมา กุสลา.\csromanspacer{} ตสฺเสว โลกุตฺตรสฺส กุสลสฺส กมฺมสฺส กตตฺตา ภาวิตตฺตา วิปากํ วิวิจฺเจว กาเมหิ ฯเปฯ ปฐมํ ฌานํ อุปสมฺปชฺช \csromanfolio{316}วิหรติ ทุกฺขปฏิปทํ ทนฺธาภิญฺญํ สุญฺญตํ, ตสฺมิํ สมเย ผสฺโส โหติ ฯเปฯ อวิกฺเขโป โหติ, อิเม ธมฺมา อพฺยากตา.\csromanspacer{} อิเมสุ ธมฺเมสุ ญาณํ อตฺถปฏิสมฺภิทา, ยาย นิรุตฺติยา เตสํ ธมฺมานํ ปญฺญตฺติ โหติ, ตตฺร ธมฺมนิรุตฺตาภิลาเป ญาณํ นิรุตฺติปฏิสมฺภิทา, เยน ญาเณน ตานิ ญาณานิ ชานาติ “อิมานิ ญาณานิ อิทมตฺถโชตกานี”ติ ญาเณสุ ญาณํ ปฏิภานปฏิสมฺภิทา.}

\proseitem{741}{ติสฺโส ปฏิสมฺภิทา, อตฺถปฏิสมฺภิทา นิรุตฺติปฏิสมฺภิทา ปฏิภานปฏิสมฺภิทา.}

\prose{กตเม ธมฺมา อพฺยากตา, ยสฺมิํ สมเย อกุสลสฺส กมฺมสฺส กตตฺตา อุปจิตตฺตา วิปากํ จกฺขุวิญฺญาณํ อุปฺปนฺนํ โหติ อุเปกฺขาสหคตํ รูปารมฺมณํ ฯเปฯ โสตวิญฺญาณํ อุปฺปนฺนํ โหติ อุเปกฺขาสหคตํ สทฺทารมฺมณํ ฯเปฯ ฆานวิญฺญาณํ อุปฺปนฺนํ โหติ อุเปกฺขาสหคตํ คนฺธารมฺมณํ ฯเปฯ ชิวฺหาวิญฺญาณํ อุปฺปนฺนํ โหติ อุเปกฺขาสหคตํ รสารมฺมณํ ฯเปฯ กายวิญฺญาณํ อุปฺปนฺนํ โหติ ทุกฺขสหคตํ โผฏฺฐพฺพารมฺมณํ, ตสฺมิํ สมเย ผสฺโส โหติ, เวทนา โหติ, สญฺญา โหติ, เจตนา โหติ, จิตฺตํ โหติ, ทุกฺขํ โหติ, จิตฺตสฺเสกคฺคตา โหติ, มนินฺทฺริยํ โหติ, ทุกฺขินฺทฺริยํ โหติ, ชีวิตินฺทฺริยํ โหติ, เย วา ปน ตสฺมิํ สมเย อญฺเญปิ อตฺถิ ปฏิจฺจสมุปฺปนฺนา อรูปิโน ธมฺมา, อิเม ธมฺมา อพฺยากตา.\csromanspacer{} อิเมสุ ธมฺเมสุ ญาณํ อตฺถปฏิสมฺภิทา, ยาย นิรุตฺติยา เตสํ ธมฺมานํ ปญฺญตฺติ โหติ, ตตฺร ธมฺมนิรุตฺตาภิลาเป ญาณํ นิรุตฺติปฏิสมฺภิทา.\csromanspacer{} เยน ญาเณน ตานิ ญาณานิ ชานาติ “อิมานิ ญาณานิ อิทมตฺถโชตกานี”ติ ญาเณสุ ญาณํ ปฏิภานปฏิสมฺภิทา.}

\proseitem{742}{ติสฺโส ปฏิสมฺภิทา, อตฺถปฏิสมฺภิทา นิรุตฺติปฏิสมฺภิทา ปฏิภานปฏิสมฺภิทา.}

\prose{กตเม ธมฺมา อพฺยากตา, ยสฺมิํ สมเย อกุสลสฺส กมฺมสฺส กตตฺตา อุปจิตตฺตา วิปากา มโนธาตุ อุปฺปนฺนา โหติ อุเปกฺขาสหคตา รูปารมฺมณา วา ฯเปฯ โผฏฺฐพฺพารมฺมณา วา ฯเปฯ มโนวิญฺญาณธาตุ อุปฺปนฺนา โหติ อุเปกฺขาสหคตา รูปารมฺมณา วา ฯเปฯ ธมฺมารมฺมณา วา ยํ ยํ วา ปนารพฺภ, ตสฺมิํ สมเย ผสฺโส โหติ, เวทนา โหติ, สญฺญา โหติ, เจตนา โหติ, จิตฺตํ โหติ, วิตกฺโก โหติ,}

\csromanheader{2}{15. ปฏิสมฺภิทาวิภงฺค}
\prose{\csromanfolio{317}วิจาโร โหติ, อุเปกฺขา โหติ, จิตฺตสฺเสกคฺคตา โหติ, มนินฺทฺริยํ โหติ, อุเปกฺขินฺทฺริยํ โหติ, ชีวิตินฺทฺริยํ โหติ, เย วา ปน ตสฺมิํ สมเย อญฺเญปิ อตฺถิ ปฏิจฺจสมุปฺปนฺนา อรูปิโน ธมฺมา, อิเม ธมฺมา อพฺยากตา.\csromanspacer{} อิเมสุ ธมฺเมสุ ญาณํ อตฺถปฏิสมฺภิทา, ยาย นิรุตฺติยา เตสํ ธมฺมานํ ปญฺญตฺติ โหติ, ตตฺร ธมฺมนิรุตฺตาภิลาเป ญาณํ นิรุตฺติปฏิสมฺภิทา, เยน ญาเณน ตานิ ญาณานิ ชานาติ “อิมานิ ญาณานิ อิทมตฺถโชตกานี”ติ ญาเณสุ ญาณํ ปฏิภานปฏิสมฺภิทา.}

\csromanmatikaanchor{mtk.143}
\csromantocmarkref{subparagraph}{4. กิริยวาร}{mtk.143}
\bookmark[dest={mtk.143},level=5]{4. กิริยวาร}
\csromanheader{6}{4. กิริยวาร}
\proseitem{743}{ติสฺโส ปฏิสมฺภิทา, อตฺถปฏิสมฺภิทา นิรุตฺติปฏิสมฺภิทา ปฏิภานปฏิสมฺภิทา.}

\prose{กตเม ธมฺมา อพฺยากตา, ยสฺมิํ สมเย มโนธาตุ อุปฺปนฺนา โหติ กิริยา เนว กุสลา นากุสลา น จ กมฺมวิปากา อุเปกฺขาสหคตา รูปารมฺมณา วา ฯเปฯ โผฏฺฐพฺพารมฺมณา วา ยํ ยํ วา ปนารพฺภ, ตสฺมิํ สมเย ผสฺโส โหติ, เวทนา โหติ, สญฺญา โหติ, เจตนา โหติ, จิตฺตํ โหติ, วิตกฺโก โหติ, วิจาโร, โหติ, อุเปกฺขา โหติ, จิตฺตสฺเสกคฺคตา โหติ, มนินฺทฺริยํ โหติ, อุเปกฺขินฺทฺริยํ โหติ, ชีวิตินฺทฺริยํ โหติ, เย วา ปน ตสฺมิํ สมเย อญฺเญปิ อตฺถิ ปฏิจฺจสมุปฺปนฺนา อรูปิโน ธมฺมา, อิเม ธมฺมา อพฺยากตา.\csromanspacer{} อิเมสุ ธมฺเมสุ ญาณํ อตฺถปฏิสมฺภิทา.\csromanspacer{} ยาย นิรุตฺติยา เตสํ ธมฺมานํ ปญฺญตฺติ โหติ, ตตฺร ธมฺมนิรุตฺตาภิลาเป ญาณํ นิรุตฺติปฏิสมฺภิทา.\csromanspacer{} เยน ญาเณน ตานิ ญาณานิ ชานาติ “อิมานิ ญาณานิ อิทมตฺถโชตกานี”ติ ญาเณสุ ญาณํ ปฏิภานปฏิสมฺภิทา.}

\proseitem{744}{ติสฺโส ปฏิสมฺภิทา, อตฺถปฏิสมฺภิทา นิรุตฺติปฏิสมฺภิทา ปฏิภานปฏิสมฺภิทา.}

\prose{กตเม ธมฺมา อพฺยากตา, ยสฺมิํ สมเย มโนวิญฺญาณธาตุ อุปฺปนฺนา โหติ กิริยา เนว กุสลา นากุสลา น จ กมฺมวิปากา โสมนสฺสสหคตา รูปารมฺมณา วา ฯเปฯ ธมฺมารมฺมณา วา ฯเปฯ มโนวิญฺญาณธาตุ อุปฺปนฺนา โหติ กิริยา เนว กุสลา นากุสลา น จ กมฺมวิปากา อุเปกฺขาสหคตา รูปารมฺมณา วา ฯเปฯ ธมฺมารมฺมณา วา \csromanfolio{318}ยํ ยํ วา ปนารพฺภ, ตสฺมิํ สมเย ผสฺโส โหติ, เวทนา โหติ, สญฺญา โหติ, เจตนา โหติ, จิตฺตํ โหติ, วิตกฺโก โหติ, วิจาโร โหติ, อุเปกฺขา โหติ, จิตฺตสฺเสกคฺคตา โหติ, วีริยินฺทฺริยํ โหติ, สมาธินฺทฺริยํ โหติ, มนินฺทฺริยํ โหติ, อุเปกฺขินฺทฺริยํ โหติ, ชีวิตินฺทฺริยํ โหติ, เย วา ปน ตสฺมิํ สมเย อญฺเญปิ อตฺถิ ปฏิจฺจสมุปฺปนฺนา อรูปิโน ธมฺมา, อิเม ธมฺมา อพฺยากตา.\csromanspacer{} อิเมสุ ธมฺเมสุ ญาณํ อตฺถปฏิสมฺภิทา, ยาย นิรุตฺติยา เตสํ ธมฺมานํ ปญฺญตฺติ โหติ, ตตฺร ธมฺมนิรุตฺตาภิลาเป ญาณํ นิรุตฺติปฏิสมฺภิทา, เยน ญาเณน ตานิ ญาณานิ ชานาติ “อิมานิ ญาณานิ อิทมตฺถโชตกานี”ติ ญาเณสุ ญาณํ ปฏิภานปฏิสมฺภิทา.}

\proseitem{745}{ติสฺโส ปฏิสมฺภิทา, อตฺถปฏิสมฺภิทา นิรุตฺติปฏิสมฺภิทา ปฏิภานปฏิสมฺภิทา.}

\prose{กตเม ธมฺมา อพฺยากตา, ยสฺมิํ สมเย มโนวิญฺญาณาธาตุ อุปฺปนฺนา โหติ กิริยา เนว กุสลา นากุสลา น จ กมฺมวิปากา โสมนสฺสสหคตา ญาณสมฺปยุตฺตา ฯเปฯ โสมนสฺสสหคตา ญาณสมฺปยุตฺตา สสงฺขาเรน ฯเปฯ โสมนสฺสสหคตา ญาณวิปฺปยุตฺตา ฯเปฯ โสมนสฺสสหคตา ญาณวิปฺปยุตฺตา สสงฺขาเรน ฯเปฯ อุเปกฺขาสหคตา ญาณสมฺปยุตฺตา ฯเปฯ อุเปกฺขาสหคตา ญาณสมฺปยุตฺตา สสงฺขาเรน ฯเปฯ อุเปกฺขาสหคตา ญาณวิปฺปยุตฺตา ฯเปฯ อุเปกฺขาสหคตา ญาณวิปฺปยุตฺตา สสงฺขาเรน ฯเปฯ รูปาวจรํ ฌานํ ภาเวติ ฯเปฯ อรูปาวจรํ ฌานํ ภาเวติ กิริยํ เนว กุสลํ นากุสลํ น จ กมฺมวิปากํ ทิฏฺฐธมฺมสุขวิหารํ ฯเปฯ สพฺพโส อากิญฺจญฺญายตนํ สมติกฺกมฺม เนวสญฺญานาสญฺญายตนสญฺญาสหคตํ สุขสฺส จ ปหานา ฯเปฯ จตุตฺถํ ฌานํ อุปสมฺปชฺช วิหรติ, ตสฺมิํ สมเย ผสฺโส โหติ ฯเปฯ อวิกฺเขโป โหติ, อิเม ธมฺมา อพฺยากตา.\csromanspacer{} อิเมสุ ธมฺเมสุ ญาณํ อตฺถปฏิสมฺภิทา, ยาย นิรุตฺติยา เตสํ ธมฺมานํ ปญฺญตฺติ โหติ, ตตฺร ธมฺมนิรุตฺตาภิลาเป ญาณํ นิรุตฺติปฏิสมฺภิทา, เยน ญาเณน ตานิ ญาณานิ ชานาติ “อิมานิ ญาณานิ อิทมตฺถโชตกานี”ติ ญาเณสุ ญาณํ ปฏิภานปฏิสมฺภิทา.}

\proseitem{746}{จตสฺโส ปฏิสมฺภิทา, อตฺถปฏิสมฺภิทา ธมฺมปฏิสมฺภิทา นิรุตฺติปฏิสมฺภิทา ปฏิภานปฏิสมฺภิทา.\csromanspacer{} ติสฺโส ปฏิสมฺภิทา กามาวจรกุสลโต}

\csromanheader{2}{15. ปฏิสมฺภิทาวิภงฺค}
\prose{\csromanfolio{319}จตูสุ ญาณสมฺปยุตฺเตสุ จิตฺตุปฺปาเทสุ, กิริยโต จตูสุ ญาณสมฺปยุตฺเตสุ จิตฺตุปฺปาเทสุ อุปฺปชฺชนฺติ.\csromanspacer{} อตฺถปฏิสมฺภิทา เอเตสุ เจว อุปฺปชฺชติ, จตูสุ มคฺเคสุ จตูสุ ผเลสุ จ อุปฺปชฺชติ.}

\csromanheader{2}{อภิธมฺมภาชนียํ.}
\csromansectionrule
\csromanheader{2}{3. ปญฺหาปุจฺฉก}
\proseitem{747}{จตสฺโส ปฏิสมฺภิทา, อตฺถปฏิสมฺภิทา ธมฺมปฏิสมฺภิทา นิรุตฺติปฏิสมฺภิทา ปฏิภานปฏิสมฺภิทา.}

\proseitem{748}{จตุนฺนํ ปฏิสมฺภิทานํ กติ กุสลา, กติ อกุสลา, กติ อพฺยากตา ฯเปฯ กติ สรณา, กติ อรณา.}

\csromanmatikaanchor{mtk.144}
\csromanmatikaanchor{mtk.145}
\csromantocmarknopage{subparagraph}{3. ปญฺหาปุจฺฉก}{mtk.144}
\bookmark[dest={mtk.144},level=5]{3. ปญฺหาปุจฺฉก}
\csromantocmarkref{subparagraph}{1. ติก}{mtk.145}
\bookmark[dest={mtk.145},level=5]{1. ติก}
\csromanheader{6}{1. ติก}
\proseitem{749}{สิยา กุสลา, สิยา อพฺยากตา, สิยา สุขาย เวทนาย สมฺปยุตฺตา, สิยา อทุกฺขมสุขาย เวทนาย สมฺปยุตฺตา.\csromanspacer{} ติสฺโส ปฏิสมฺภิทา สิยา วิปากธมฺมธมฺมา, สิยา เนววิปากนวิปากธมฺมธมฺมา, อตฺถปฏิสมฺภิทา สิยา วิปากา, สิยา วิปากธมฺมธมฺมา, สิยา เนววิปากนวิปากธมฺมธมฺมา.\csromanspacer{} ติสฺโส ปฏิสมฺภิทา อนุปาทินฺนุปาทานิยา.\csromanspacer{} อตฺถปฏิสมฺภิทา สิยา อนุปาทินฺนุปาทานิยา, สิยา อนุปาทินฺน-อนุปาทานิยา.\csromanspacer{} ติสฺโส ปฏิสมฺภิทา อสํกิลิฏฺฐ-อสํกิเลสิกา.\csromanspacer{} อตฺถปฏิสมฺภิทา สิยา อสํกิลิฏฺฐสํกิเลสิกา, สิยา อสํกิลิฏฺฐ- อสํกิเลสิกา.}

\prose{ติสฺโส ปฏิสมฺภิทา สวิตกฺกสวิจารา.\csromanspacer{} อตฺถปฏิสมฺภิทา สิยา สวิตกฺกสวิจารา, สิยา อวิตกฺกวิจารมตฺตา, สิยา อวิตกฺก-อวิจารา, สิยา ปีติสหคตา, สิยา สุขสหคตา, สิยา อุเปกฺขาสหคตา, เนว ทสฺสเนน น ภาวนาย ปหาตพฺพา, เนว ทสฺสเนน น ภาวนาย ปหาตพฺพเหตุกา.}

\prose{ติสฺโส ปฏิสมฺภิทา สิยา อาจยคามิโน, สิยา เนวาจยคามินาปจยคามิโน.\csromanspacer{} อตฺถปฏิสมฺภิทา สิยา อาจยคามินี, สิยา อปจยคามินี, สิยา เนวาจยคามินาปจยคามินี.}

\prose{\csromanfolio{320}ติสฺโส ปฏิสมฺภิทา เนวเสกฺขนาเสกฺขา.\csromanspacer{} อตฺถปฏิสมฺภิทา สิยา เสกฺขา, สิยา อเสกฺขา, สิยา เนวเสกฺขนาเสกฺขา.\csromanspacer{} ติสฺโส ปฏิสมฺภิทา ปริตฺตา.\csromanspacer{} อตฺถปฏิสมฺภิทา สิยา ปริตฺตา, สิยา อปฺปมาณา.\csromanspacer{} นิรุตฺติปฏิสมฺภิทา ปริตฺตารมฺมณา.\csromanspacer{} ติสฺโส ปฏิสมฺภิทา สิยา ปริตฺตารมฺมณา, สิยา มหคฺคตารมฺมณา, สิยา อปฺปมาณารมฺมณา.}

\prose{ติสฺโส ปฏิสมฺภิทา มชฺฌิมา.\csromanspacer{} อตฺถปฏิสมฺภิทา สิยา มชฺฌิมา, สิยา ปณีตา.\csromanspacer{} ติสฺโส ปฏิสมฺภิทา อนิยตา.\csromanspacer{} อตฺถปฏิสมฺภิทา สิยา สมฺมตฺตนิยตา, สิยา อนิยตา.\csromanspacer{} นิรุตฺติปฏิสมฺภิทา น วตฺตพฺพา “มคฺคารมฺมณา”ติปิ, “มคฺคเหตุกา”ติปิ, “มคฺคาธิปตินี”ติปิ.\csromanspacer{} อตฺถปฏิสมฺภิทา น มคฺคารมฺมณา, สิยา มคฺคเหตุกา, สิยา มคฺคาธิปตินี, สิยา น วตฺตพฺพา “มคฺคเหตุกา”ติปิ, “มคฺคาธิปตินี”ติปิ.\csromanspacer{} เทฺว ปฏิสมฺภิทา สิยา มคฺคารมฺมณา, น มคฺคเหตุกา, สิยา มคฺคาธิปติโน, สิยา น วตฺตพฺพา “มคฺคารมฺมณา”ติปิ, “มคฺคาธิปติโน”ติปิ.}

\prose{ติสฺโส ปฏิสมฺภิทา สิยา อุปฺปนฺนา, สิยา อนุปฺปนฺนา, น วตฺตพฺพา “อุปฺปาทิโน”ติ.\csromanspacer{} อตฺถปฏิสมฺภิทา สิยา อุปฺปนฺนา, สิยา อนุปฺปนฺนา, สิยา อุปฺปาทินี, สิยา อตีตา, สิยา อนาคตา, สิยา ปจฺจุปฺปนฺนา.\csromanspacer{} นิรุตฺติปฏิสมฺภิทา ปจฺจุปฺปนฺนารมฺมณา.\csromanspacer{} เทฺว ปฏิสมฺภิทา สิยา อตีตารมฺมณา, สิยา อนาคตารมฺมณา, สิยา ปจฺจุปฺปนฺนารมฺมณา.\csromanspacer{} อตฺถปฏิสมฺภิทา สิยา อตีตารมฺมณา, สิยา อนาคตารมฺมณา, สิยา ปจฺจุปฺปนฺนารมฺมณา, สิยา น วตฺตพฺพา “อตีตารมฺมณา”ติปิ, “อนาคตารมฺมณา”ติปิ, “ปจฺจุปฺปนฺนารมฺมณา”ติปิ, สิยา อชฺฌตฺตา, สิยา พหิทฺธา, สิยา อชฺฌตฺตพหิทฺธา.\csromanspacer{} นิรุตฺติปฏิสมฺภิทา พหิทฺธารมฺมณา.\csromanspacer{} ติสฺโส ปฏิสมฺภิทา สิยา อชฺฌตฺตารมฺมณา, สิยา พหิทฺธารมฺมณา, สิยา อชฺฌตฺตพหิทฺธารมฺมณา, อนิทสฺสน-อปฺปฏิฆา.}

\csromanmatikaanchor{mtk.146}
\csromantocmarkref{subparagraph}{2. ทุก}{mtk.146}
\bookmark[dest={mtk.146},level=5]{2. ทุก}
\csromanheader{6}{2. ทุก}
\proseitem{750}{เหตู, สเหตุกา, เหตุสมฺปยุตฺตา, เหตู เจว สเหตุกา จ, เหตู เจว เหตุสมฺปยุตฺตา จ, น วตฺตพฺพา “นเหตุสเหตุกา”ติปิ, “นเหตุ-อเหตุกา”ติปิ.}

\csromanheader{2}{15. ปฏิสมฺภิทาวิภงฺค}
\prose{\csromanfolio{321}สปฺปจฺจยา, สงฺขตา, อนิทสฺสนา, อปฺปฏิฆา, อรูปา.\csromanspacer{} ติสฺโส ปฏิสมฺภิทา โลกิยา.\csromanspacer{} อตฺถปฏิสมฺภิทา สิยา โลกิยา, สิยา โลกุตฺตรา, เกนจิ วิญฺเญยฺยา, เกนจิ น วิญฺเญยฺยา.}

\prose{โน อาสวา.\csromanspacer{} ติสฺโส ปฏิสมฺภิทา สาสวา.\csromanspacer{} อตฺถปฏิสมฺภิทา สิยา สาสวา, สิยา อนาสวา, อาสววิปฺปยุตฺตา.\csromanspacer{} ติสฺโส ปฏิสมฺภิทา น วตฺตพฺพา “อาสวา เจว สาสวา จา”ติ, สาสวา เจว โน จ อาสวา.\csromanspacer{} อตฺถปฏิสมฺภิทา น วตฺตพฺพา “อาสโว\footnote{อาสวา (สี) ธาตุวิภงฺเค ปน ปาฐนานตฺตํ นตฺถิ.} เจว สาสวา จา”ติ, สิยา สาสวา เจว โน จ อาสโว1, สิยา น วตฺตพฺพา “สาสวา เจว โน จ อาสโว”ติ, น วตฺตพฺพา “อาสวา เจว อาสวสมฺปยุตฺตา จา”ติปิ, อาสวสมฺปยุตฺตา เจว โน จ อาสวา”ติปิ.\csromanspacer{} ติสฺโส ปฏิสมฺภิทา อาสววิปฺปยุตฺตา สาสวา.\csromanspacer{} อตฺถปฏิสมฺภิทา สิยา อาสววิปฺปยุตฺตา สาสวา, สิยา อาสววิปฺปยุตฺตา อนาสวา.}

\prose{โน สํโยชนา ฯเปฯ.\csromanspacer{} โน คนฺถา ฯเปฯ.\csromanspacer{} โน โอฆา ฯเปฯ.\csromanspacer{} โน โยคา ฯเปฯ.\csromanspacer{} โน นีวรณา ฯเปฯ.\csromanspacer{} โน ปรามาสา ฯเปฯ สารมฺมณา, โน จิตฺตา, เจตสิกา, จิตฺตสมฺปยุตฺตา, จิตฺตสํสฏฺฐา, จิตฺตสมุฏฺฐานา, จิตฺตสหภุโน, จิตฺตานุปริวตฺติโน, จิตฺตสํสฏฺฐสมุฏฺฐานา, จิตฺตสํสฏฺฐสมุฏฺฐานสหภุโน, จิตฺตสํสฏฺฐสมุฏฺฐานานุปริวตฺติโน, พาหิรา, โน อุปาทา, อนุปาทินฺนา.}

\prose{โน อุปาทานา ฯเปฯ.\csromanspacer{} โน กิเลสา ฯเปฯ.\csromanspacer{} น ทสฺสเนน ปหาตพฺพา, น ภาวนาย ปหาตพฺพา, น ทสฺสเนน ปหาตพฺพเหตุกา, น ภาวนาย ปหาตพฺพเหตุกา.\csromanspacer{} ติสฺโส ปฏิสมฺภิทา สวิตกฺกา.\csromanspacer{} อตฺถปฏิสมฺภิทา สิยา สวิตกฺกา, สิยา อวิตกฺกา.\csromanspacer{} ติสฺโส ปฏิสมฺภิทา สวิจารา.\csromanspacer{} อตฺถปฏิสมฺภิทา สิยา สวิจารา, สิยา อวิจารา, สิยา สปฺปีติกา, สิยา อปฺปีติกา, สิยา ปีติสหคตา, สิยา น ปีติสหคตา, สิยา สุขสหคตา, สิยา น สุขสหคตา, สิยา อุเปกฺขาสหคตา, สิยา น อุเปกฺขาสหคตา.\csromanspacer{} ติสฺโส ปฏิสมฺภิทา กามาวจรา.\csromanspacer{} อตฺถปฏิสมฺภิทา สิยา กามาวจรา, สิยา น กามาวจรา น รูปาวจรา น}

\prose{\csromanfolio{322}อรูปาวจรา.\csromanspacer{} ติสฺโส ปฏิสมฺภิทา ปริยาปนฺนา.\csromanspacer{} อตฺถปฏิสมฺภิทา สิยา ปริยาปนฺนา, สิยา อปริยาปนฺนา.\csromanspacer{} ติสฺโส ปฏิสมฺภิทา อนิยฺยานิกา, อตฺถปฏิสมฺภิทา สิยา นิยฺยานิกา, สิยา อนิยฺยานิกา.\csromanspacer{} ติสฺโส ปฏิสมฺภิทา อนิยตา.\csromanspacer{} อตฺถปฏิสมฺภิทา สิยา นิยตา, สิยา อนิยตา.\csromanspacer{} ติสฺโส ปฏิสมฺภิทา ส- อุตฺตรา.\csromanspacer{} อตฺถปฏิสมฺภิทา สิยา ส-อุตฺตรา, สิยา อนุตฺตรา, อรณาติ.}

\vspace{\tipitakaparskip}
\csromancenter{ปญฺหาปุจฺฉกํ.}
\nitthitam{\textbf{ปฏิสมฺภิทาวิภงฺโค นิฏฺฐิโต.}}
\csromanmatikaanchor{mtk.147}
\csromantocmarknopage{subparagraph}{16. ญาณวิภงฺค}{mtk.147}
\bookmark[dest={mtk.147},level=5]{16. ญาณวิภงฺค}
\csromanheader{6}{16. ญาณวิภงฺค}
\csromanmatikaanchor{mtk.148}
\csromantocmarkref{subparagraph}{1. เอกกมาติกา}{mtk.148}
\bookmark[dest={mtk.148},level=5]{1. เอกกมาติกา}
\csromanheader{6}{1. เอกกมาติกา}
\proseitem{751}{\csromanfolio{323}เอกวิเธน ญาณวตฺถุ\csromandash{}ปญฺจ วิญฺญาณา น เหตู, อเหตุกา, เหตุวิปฺปยุตฺตา, สปฺปจฺจยา, สงฺขตา, อรูปา, โลกิยา, สาสวา, สํโยชนิยา, คนฺถนิยา, โอฆนิยา, โยคนิยา, นีวรณิยา, ปรามฏฺฐา, อุปาทานิยา, สํกิเลสิกา, อพฺยากตา, สารมฺมณา, อเจตสิกา, วิปากา, อุปาทินฺนุปาทานิยา, อสํกิลิฏฺฐสํกิเลสิกา, น สวิตกฺกสวิจารา, น อวิตกฺกวิจารมตฺตา, อวิตกฺก-อวิจารา, น ปีติสหคตา, เนว ทสฺสเนน น ภาวนาย ปหาตพฺพา, เนว ทสฺสเนน น ภาวนาย ปหาตพฺพเหตุกา, เนวาจยคามินาปจยคามิโน, เนวเสกฺขนาเสกฺขา, ปริตฺตา, กามาวจรา, น รูปาวจรา, น อรูปาวจรา, ปริยาปนฺนา, โน อปริยาปนฺนา, อนิยตา, อนิยฺยานิกา.}

\prose{อุปฺปนฺนวตฺถุกา, อุปฺปนฺนารมฺมณา, ปุเรชาตวตฺถุกา, ปุเรชาตารมฺมณา, อชฺฌตฺติกวตฺถุกา, พาหิรารมฺมณา, อสมฺภินฺนวตฺถุกา, อสมฺภินฺนารมฺมณา, นานาวตฺถุกา, นานารมฺมณา, น อญฺญมญฺญสฺส โคจรวิสยํ ปจฺจนุโภนฺติ, น อสมนฺนาหารา อุปฺปชฺชนฺติ, น อมนสิการา อุปฺปชฺชนฺติ, น อพฺโพกิณฺณา อุปฺปชฺชนฺติ, น อปุพฺพํ อจริมํ อุปฺปชฺชนฺติ, น อญฺญมญฺญสฺส สมนนฺตรา อุปฺปชฺชนฺติ.}

\prose{ปญฺจ วิญฺญาณา อนาโภคา, ปญฺจหิ วิญฺญาเณหิ น กญฺจิ\footnote{กิญฺจิ (สี, ก)} ธมฺมํ ปฏิวิชานาติ อญฺญตฺร อภินิปาตมตฺตา, ปญฺจนฺนํ วิญฺญาณานํ สมนนฺตราปิ น กญฺจิ ธมฺมํ ปฏิวิชานาติ, ปญฺจหิ วิญฺญาเณหิ น กญฺจิ อิริยาปถํ กปฺเปติ, ปญฺจนฺนํ วิญฺญาณานํ สมนนฺตราปิ น กญฺจิ อิริยาปถํ กปฺเปติ, ปญฺจหิ วิญฺญาเณหิ น กายกมฺมํ น วจีกมฺมํ ปฏฺฐเปติ, ปญฺจนฺนํ วิญฺญาณานํ สมนนฺตราปิ น กายกมฺมํ น วจีกมฺมํ ปฏฺฐเปติ, ปญฺจหิ วิญฺญาเณหิ น กุสลากุสลํ ธมฺมํ สมาทิยติ, ปญฺจนฺนํ วิญฺญาณานํ สมนนฺตราปิ น กุสลากุสลํ ธมฺมํ สมาทิยติ, ปญฺจหิ วิญฺญาเณหิ น สมาปชฺชติ น วุฏฺฐาติ, ปญฺจนฺนํ วิญฺญาณานํ สมนนฺตราปิ น สมาปชฺชติ น วุฏฺฐาติ, ปญฺจหิ วิญฺญาเณหิ น วจติ น อุปฺปชฺชติ, ปญฺจนฺนํ วิญฺญาณานํ สมนนฺตราปิ น วจติ น อุปฺปชฺชติ, ปญฺจหิ วิญฺญาเณหิ น สุปติ \csromanfolio{324}น ปฏิพุชฺฌติ น สุปินํ ปสฺสติ, ปญฺจนฺนํ วิญฺญาณานํ สมนนฺตราปิ น สุปติ น ปฏิพุชฺฌติ น สุปินํ ปสฺสติ, ยาถาวกวตฺถุวิภาวนา ปญฺญา.\csromanspacer{} เอวํ เอกวิเธน ญาณวตฺถุ.}

\csromanmatikaanchor{mtk.149}
\csromantocmarkref{subparagraph}{2. ทุกมาติกา}{mtk.149}
\bookmark[dest={mtk.149},level=5]{2. ทุกมาติกา}
\csromanheader{6}{2. ทุกมาติกา}
\proseitem{752}{ทุวิเธน ญาณวตฺถุ\csromandash{}โลกิยา ปญฺญา, โลกุตฺตรา ปญฺญา.\csromanspacer{} เกนจิ วิญฺเญยฺยา ปญฺญา, เกนจิ น วิญฺเญยฺยา ปญฺญา.\csromanspacer{} สาสวา ปญฺญา, อนาสวา ปญฺญา.\csromanspacer{} อาสววิปฺปยุตฺตา สาสวา ปญฺญา, อาสววิปฺปยุตฺตา อนาสวา ปญฺญา.\csromanspacer{} สํโยชนิยา ปญฺญา, อสํโยชนิยา ปญฺญา.\csromanspacer{} สํโยชนวิปฺปยุตฺตา สํโยชนิยา ปญฺญา, สํโยชนวิปฺปยุตฺตา อสํโยชนิยา ปญฺญา.\csromanspacer{} คนฺถนิยา ปญฺญา, อคนฺถนิยา ปญฺญา.\csromanspacer{} คนฺถวิปฺปยุตฺตา คนฺถนิยา ปญฺญา, คนฺถวิปฺปยุตฺตา อคนฺถนิยา ปญฺญา.}

\prose{โอฆนิยา ปญฺญา, อโนฆนิยา ปญฺญา.\csromanspacer{} โอฆวิปฺปยุตฺตา โอฆนิยา ปญฺญา, โอฆวิปฺปยุตฺตา อโนฆนิยา ปญฺญา.\csromanspacer{} โยคนิยา ปญฺญา, อโยคนิยา ปญฺญา.\csromanspacer{} โยควิปฺปยุตฺตา โยคนิยา ปญฺญา, โยควิปฺปยุตฺตา อโยคนิยา ปญฺญา.\csromanspacer{} นีวรณิยา ปญฺญา, อนีวรณิยา ปญฺญา.\csromanspacer{} นีวรณวิปฺปยุตฺตา นีวรณิยา ปญฺญา, นีวรณวิปฺปยุตฺตา อนีวรณิยา ปญฺญา.\csromanspacer{} ปรามฏฺฐา ปญฺญา, อปรามฏฺฐา ปญฺญา.\csromanspacer{} ปรามาสวิปฺปยุตฺตา ปรามฏฺฐา ปญฺญา, ปรามาสวิปฺปยุตฺตา อปรามฏฺฐา ปญฺญา.\csromanspacer{} อุปาทินฺนา ปญฺญา, อนุปาทินฺนา ปญฺญา.\csromanspacer{} อุปาทานิยา ปญฺญา, อนุปาทานิยา ปญฺญา.\csromanspacer{} อุปาทานวิปฺปยุตฺตา อุปาทานิยา ปญฺญา, อุปาทานวิปฺปยุตฺตา อนุปาทานิยา ปญฺญา.}

\prose{สํกิเลสิกา ปญฺญา, อสํกิเลสิกา ปญฺญา.\csromanspacer{} กิเลสวิปฺปยุตฺตา สํกิเลสิกา ปญฺญา, กิเลสวิปฺปยุตฺตา อสํกิเลสิกา ปญฺญา.\csromanspacer{} สวิตกฺกา ปญฺญา, อวิตกฺกา ปญฺญา.\csromanspacer{} สวิจารา ปญฺญา, อวิจารา ปญฺญา.\csromanspacer{} สปฺปีติกา ปญฺญา, อปฺปีติกา ปญฺญา.\csromanspacer{} ปีติสหคตา ปญฺญา, น ปีติสหคตา ปญฺญา.\csromanspacer{} สุขสหคตา ปญฺญา, น สุขสหคตา ปญฺญา.\csromanspacer{} อุเปกฺขาสหคตา ปญฺญา, น อุเปกฺขาสหคตา ปญฺญา.\csromanspacer{} กามาวจรา ปญฺญา, น กามาวจรา ปญฺญา.\csromanspacer{} รูปาวจรา ปญฺญา, น รูปาวจรา ปญฺญา.\csromanspacer{} อรูปาวจรา ปญฺญา, น อรูปาวจรา ปญฺญา.\csromanspacer{} ปริยาปนฺนา ปญฺญา, อปริยาปนฺนา ปญฺญา.\csromanspacer{} นิยฺยานิกา ปญฺญา, อนิยฺยานิกา ปญฺญา.\csromanspacer{} นิยตา ปญฺญา, อนิยตา ปญฺญา.\csromanspacer{} ส-อุตฺตรา}

\csromanheader{2}{16. ญาณวิภงฺค}
\prose{\csromanfolio{325}ปญฺญา, อนุตฺตรา ปญฺญา.\csromanspacer{} อตฺถชาปิกา ปญฺญา, ชาปิตตฺตา ปญฺญา.\csromanspacer{} เอวํ ทุวิเธน ญาณวตฺถุ.}

\csromanmatikaanchor{mtk.150}
\csromantocmarkref{subparagraph}{3. ติกมาติกา}{mtk.150}
\bookmark[dest={mtk.150},level=5]{3. ติกมาติกา}
\csromanheader{6}{3. ติกมาติกา}
\proseitem{753}{ติวิเธน ญาณวตฺถุ\csromandash{}จินฺตามยา ปญฺญา, สุตมยา ปญฺญา, ภาวนามยา ปญฺญา.\csromanspacer{} ทานมยา ปญฺญา, สีลมยา ปญฺญา, ภาวนามยา ปญฺญา.\csromanspacer{} อธิสีเล ปญฺญา, อธิจิตฺเต ปญฺญา, อธิปญฺญาย ปญฺญา.\csromanspacer{} อายโกสลฺลํ, อปายโกสลฺลํ, อุปายโกสลฺลํ.\csromanspacer{} วิปากา ปญฺญา, วิปากธมฺมธมฺมา ปญฺญา, เนววิปากนวิปากธมฺมธมฺมา ปญฺญา.\csromanspacer{} อุปาทินฺนุปาทานิยา ปญฺญา, อนุปาทินฺนุปาทานิยา ปญฺญา, อนุปาทินฺน-อนุปาทานิยา ปญฺญา.\csromanspacer{} สวิตกฺกสวิจารา ปญฺญา, อวิตกฺกวิจารมตฺตา ปญฺญา, อวิตกฺก-อวิจารา ปญฺญา.}

\prose{ปีติสหคตา ปญฺญา, สุขสหคตา ปญฺญา, อุเปกฺขาสหคตา ปญฺญา.\csromanspacer{} อาจยคามินี ปญฺญา, อปจยคามินี ปญฺญา, เนวาจยคามินาปจยคามินี ปญฺญา.\csromanspacer{} เสกฺขา ปญฺญา, อเสกฺขา ปญฺญา, เนวเสกฺขนาเสกฺขา ปญฺญา.\csromanspacer{} ปริตฺตา ปญฺญา, มหคฺคตา ปญฺญา, อปฺปมาณา ปญฺญา.\csromanspacer{} ปริตฺตารมฺมณา ปญฺญา, มหคฺคตารมฺมณา ปญฺญา, อปฺปมาณารมฺมณา ปญฺญา.\csromanspacer{} มคฺคารมฺมณา ปญฺญา, มคฺคเหตุกา ปญฺญา, มคฺคาธิปตินี ปญฺญา.\csromanspacer{} อุปฺปนฺนา ปญฺญา, อนุปฺปนฺนา ปญฺญา, อุปฺปาทินี ปญฺญา.\csromanspacer{} อตีตา ปญฺญา, อนาคตา ปญฺญา, ปจฺจุปฺปนฺนา ปญฺญา.\csromanspacer{} อตีตารมฺมณา ปญฺญา, อนาคตารมฺมณา ปญฺญา, ปจฺจุปฺปนฺนารมฺมณา ปญฺญา.\csromanspacer{} อชฺฌตฺตา ปญฺญา, พหิทฺธา ปญฺญา, อชฺฌตฺตพหิทฺธา ปญฺญา.\csromanspacer{} อชฺฌตฺตารมฺมณา ปญฺญา, พหิทฺธารมฺมณา ปญฺญา, อชฺฌตฺตพหิทฺธารมฺมณา ปญฺญา.}

\prose{สวิตกฺกสวิจารา ปญฺญา อตฺถิ วิปากา, อตฺถิ วิปากธมฺมธมฺมา, อตฺถิ เนววิปากนวิปากธมฺมธมฺมา.\csromanspacer{} อตฺถิ อุปาทินฺนุปาทานิยา, อตฺถิ อนุปาทินฺนุปาทานิยา, อตฺถิ อนุปาทินฺน-อนุปาทานิยา.\csromanspacer{} อตฺถิ ปีติสหคตา, อตฺถิ สุขสหคตา, อตฺถิ อุเปกฺขาสหคตา.\csromanspacer{} อตฺถิ อาจยคามินี, อตฺถิ อปจยคามินี, อตฺถิ เนวาจยคามินาปจยคามินี.\csromanspacer{} อตฺถิ เสกฺขา, อตฺถิ อเสกฺขา, อตฺถิ เนวเสกฺขนาเสกฺขา.\csromanspacer{} อตฺถิ ปริตฺตา, อตฺถิ มหคฺคตา, อตฺถิ อปฺปมาณา.\csromanspacer{} อตฺถิ ปริตฺตารมฺมณา, อตฺถิ มหคฺคตารมฺมณา, อตฺถิ อปฺปมาณารมฺมณา.\csromanspacer{} อตฺถิ มคฺคารมฺมณา, อตฺถิ มคฺคเหตุกา, อตฺถิ มคฺคาธิปตินี.\csromanspacer{} อตฺถิ อุปฺปนฺนา, อตฺถิ อนุปฺปนฺนา, อตฺถิ อุปฺปาทินี.\csromanspacer{} อตฺถิ อตีตา, อตฺถิ \csromanfolio{326}อนาคตา, อตฺถิ ปจฺจุปฺปนฺนา.\csromanspacer{} อตฺถิ อตีตารมฺมณา, อตฺถิ อนาคตารมฺมณา, อตฺถิ ปจฺจุปฺปนฺนารมฺมณา.\csromanspacer{} อตฺถิ อชฺฌตฺตา, อตฺถิ พหิทฺธา, อตฺถิ อชฺฌตฺตพหิทฺธา.\csromanspacer{} อตฺถิ อชฺฌตฺตารมฺมณา, อตฺถิ พหิทฺธารมฺมณา, อตฺถิ อชฺฌตฺตพหิทฺธารมฺมณา.}

\prose{อวิตกฺกวิจารมตฺตา ปญฺญา อตฺถิ วิปากา, อตฺถิ วิปากธมฺมธมฺมา, อตฺถิ เนววิปากนวิปากธมฺมธมฺมา.\csromanspacer{} อตฺถิ อุปาทินฺนุปาทานิยา, อตฺถิ อนุปาทินฺนุปาทานิยา, อตฺถิ อนุปาทินฺน-อนุปาทานิยา.\csromanspacer{} อตฺถิ อาจยคามินี, อตฺถิ อปจยคามินี, อตฺถิ เนวาจยคามินาปจยคามินี.\csromanspacer{} อตฺถิ เสกฺขา, อตฺถิ อเสกฺขา, อตฺถิ เนวเสกฺขนาเสกฺขา.\csromanspacer{} อตฺถิ อุปฺปนฺนา, อตฺถิ อนุปฺปนฺนา, อตฺถิ อุปฺปาทินี.\csromanspacer{} อตฺถิ อตีตา, อตฺถิ อนาคตา, อตฺถิ ปจฺจุปฺปนฺนา.\csromanspacer{} อตฺถิ อชฺฌตฺตา, อตฺถิ พหิทฺธา, อตฺถิ อชฺฌตฺตพหิทฺธา.}

\prose{อวิตกฺก-อวิจารา ปญฺญา อตฺถิ วิปากา, อตฺถิ วิปากธมฺมธมฺมา, อตฺถิ เนววิปากนวิปากธมฺมธมฺมา.\csromanspacer{} อตฺถิ อุปาทินฺนุปาทานิยา, อตฺถิ อนุปาทินฺนุปาทานิยา, อตฺถิ อนุปาทินฺน-อนุปาทานิยา.\csromanspacer{} อตฺถิ ปีติสหคตา, อตฺถิ สุขสหคตา, อตฺถิ อุเปกฺขาสหคตา.\csromanspacer{} อตฺถิ อาจยคามินี, อตฺถิ อปจยคามินี, อตฺถิ เนวาจยคามินาปจยคามินี.\csromanspacer{} อตฺถิ เสกฺขา, อตฺถิ อเสกฺขา, อตฺถิ เนวเสกฺขนาเสกฺขา.\csromanspacer{} อตฺถิ ปริตฺตารมฺมณา, อตฺถิ มหคฺคตารมฺมณา, อตฺถิ อปฺปมาณารมฺมณา.\csromanspacer{} อตฺถิ มคฺคารมฺมณา, อตฺถิ มคฺคเหตุกา, อตฺถิ มคฺคาธิปตินี.\csromanspacer{} อตฺถิ อุปฺปนฺนา, อตฺถิ อนุปฺปนฺนา, อตฺถิ อุปฺปาทินี.\csromanspacer{} อตฺถิ อตีตา, อตฺถิ อนาคตา, อตฺถิ ปจฺจุปฺปนฺนา.\csromanspacer{} อตฺถิ อตีตารมฺมณา, อตฺถิ อนาคตารมฺมณา, อตฺถิ ปจฺจุปฺปนฺนารมฺมณา.\csromanspacer{} อตฺถิ อชฺฌตฺตา, อตฺถิ พหิทฺธา, อตฺถิ อชฺฌตฺตพหิทฺธา.\csromanspacer{} อตฺถิ อชฺฌตฺตารมฺมณา, อตฺถิ พหิทฺธารมฺมณา, อตฺถิ อชฺฌตฺตพหิทฺธารมฺมณา.}

\prose{ปีติสหคตา ปญฺญา สุขสหคตา ปญฺญา อตฺถิ วิปากา, อตฺถิ วิปากธมฺมธมฺมา, อตฺถิ เนววิปากนวิปากธมฺมธมฺมา.\csromanspacer{} อตฺถิ อุปาทินฺนุปาทานิยา, อตฺถิ อนุปาทินฺนุปาทานิยา, อตฺถิ อนุปาทินฺน-อนุปาทานิยา.\csromanspacer{} อตฺถิ สวิตกฺกสวิจารา, อตฺถิ อวิตกฺกวิจารมตฺตา, อตฺถิ อวิตกฺก-อวิจารา.\csromanspacer{} อตฺถิ อาจยคามินี, อตฺถิ อปจยคามินี, อตฺถิ เนวาจยคามินาปจยคามินี.\csromanspacer{} อตฺถิ เสกฺขา, อตฺถิ อเสกฺขา, อตฺถิ เนวเสกฺขนาเสกฺขา.\csromanspacer{} อตฺถิ ปริตฺตา, อตฺถิ มหคฺคตา, อตฺถิ อปฺปมาณา.\csromanspacer{} อตฺถิ ปริตฺตารมฺมณา, อตฺถิ มหคฺคตารมฺมณา, อตฺถิ อปฺปมาณารมฺมณา.\csromanspacer{} อตฺถิ}

\csromanheader{2}{16. ญาณวิภงฺค}
\prose{\csromanfolio{327}มคฺคารมฺมณา, อตฺถิ มคฺคเหตุกา, อตฺถิ มคฺคาธิปตินี.\csromanspacer{} อตฺถิ อุปฺปนฺนา, อตฺถิ อนุปฺปนฺนา, อตฺถิ อุปฺปาทินี.\csromanspacer{} อตฺถิ อตีตา, อตฺถิ อนาคตา, อตฺถิ ปจฺจุปฺปนฺนา.\csromanspacer{} อตฺถิ อตีตารมฺมณา, อตฺถิ อนาคตารมฺมณา, อตฺถิ ปจฺจุปฺปนฺนารมฺมณา.\csromanspacer{} อตฺถิ อชฺฌตฺตา, อตฺถิ พหิทฺธา, อตฺถิ อชฺฌตฺตพหิทฺธา.\csromanspacer{} อตฺถิ อชฺฌตฺตารมฺมณา, อตฺถิ พหิทฺธารมฺมณา, อตฺถิ อชฺฌตฺตพหิทฺธารมฺมณา.}

\prose{อุเปกฺขาสหคตา ปญฺญา อตฺถิ วิปากา, อตฺถิ วิปากธมฺมธมฺมา, อตฺถิ เนววิปากนวิปากธมฺมธมฺมา.\csromanspacer{} อตฺถิ อุปาทินฺนุปาทานิยา, อตฺถิ อนุปาทินฺนุปาทานิยา, อตฺถิ อนุปาทินฺน-อนุปาทานิยา.\csromanspacer{} อตฺถิ อาจยคามินี, อตฺถิ อปจยคามินี, อตฺถิ เนวาจยคามินาปจยคามินี.\csromanspacer{} อตฺถิ เสกฺขา, อตฺถิ อเสกฺขา, อตฺถิ เนวเสกฺขนาเสกฺขา.\csromanspacer{} อตฺถิ ปริตฺตา, อตฺถิ มหคฺคตา, อตฺถิ อปฺปมาณา.\csromanspacer{} อตฺถิ ปริตฺตารมฺมณา, อตฺถิ มหคฺคตารมฺมณา, อตฺถิ อปฺปมาณารมฺมณา.\csromanspacer{} อตฺถิ มคฺคารมฺมณา, อตฺถิ มคฺคเหตุกา, อตฺถิ มคฺคาธิปตินี.\csromanspacer{} อตฺถิ อุปฺปนฺนา, อตฺถิ อนุปฺปนฺนา, อตฺถิ อุปฺปาทินี.\csromanspacer{} อตฺถิ อตีตา, อตฺถิ อนาคตา, อตฺถิ ปจฺจุปฺปนฺนา.\csromanspacer{} อตฺถิ อตีตารมฺมณา, อตฺถิ อนาคตารมฺมณา, อตฺถิ ปจฺจุปฺปนฺนารมฺมณา.\csromanspacer{} อตฺถิ อชฺฌตฺตา, อตฺถิ พหิทฺธา, อตฺถิ อชฺฌตฺตพหิทฺธา.\csromanspacer{} อตฺถิ อชฺฌตฺตารมฺมณา, อตฺถิ พหิทฺธารมฺมณา, อตฺถิ อชฺฌตฺตพหิทฺธารมฺมณา.\csromanspacer{} เอวํ ติวิเธน ญาณวตฺถุ.}

\csromanmatikaanchor{mtk.151}
\csromantocmarkref{subparagraph}{4. จตุกฺกมาติกา}{mtk.151}
\bookmark[dest={mtk.151},level=5]{4. จตุกฺกมาติกา}
\csromanheader{6}{4. จตุกฺกมาติกา}
\proseitem{754}{จตุพฺพิเธน ญาณวตฺถุ\csromandash{}กมฺมสฺสกตญาณํ, สจฺจานุโลมิกํ ญาณํ, มคฺคสมงฺคิสฺส ญาณํ, ผลสมงฺคิสฺส ญาณํ.\csromanspacer{} ทุกฺเข ญาณํ, ทุกฺขสมุทเย ญาณํ, ทุกฺขนิโรเธ ญาณํ, ทุกฺขนิโรธคามินิยา ปฏิปทาย ญาณํ.\csromanspacer{} กามาวจรา ปญฺญา, รูปาวจรา ปญฺญา, อรูปาวจรา ปญฺญา, อปริยาปนฺนา ปญฺญา.\csromanspacer{} ธมฺเม ญาณํ, อนฺวเย ญาณํ, ปริเย\footnote{ปริจฺเจ (สพฺพตฺถ) ที 3. 188 ปิฏฺเฐปิ.} ญาณํ, สมฺมุติญาณํ\footnote{สมฺมติญาณํ (สฺยา) ที 3. 188 ปิฏฺเฐปิ.}.\csromanspacer{} อตฺถิ ปญฺญา อาจยาย โน อปจยาย, อตฺถิ ปญฺญา อปจยาย โน อาจยาย, อตฺถิ ปญฺญา อาจยาย เจว อปจยาย จ, อตฺถิ ปญฺญา เนวาจยาย โน อปจยาย.\csromanspacer{} อตฺถิ ปญฺญา นิพฺพิทาย โน ปฏิเวธาย, อตฺถิ ปญฺญา ปฏิเวธาย โน นิพฺพิทาย, อตฺถิ ปญฺญา นิพฺพิทาย เจว ปฏิเวธาย จ, อตฺถิ ปญฺญา เนว นิพฺพิทาย โน ปฏิเวธาย.\csromanspacer{} หานภาคินี ปญฺญา, ฐิติภาคินี ปญฺญา, วิเสสภาคินี ปญฺญา, นิพฺเพธภาคินี ปญฺญา.}

\prose{\csromanfolio{328}จตสฺโส ปฏิสมฺภิทา.\csromanspacer{} จตสฺโส ปฏิปทา.\csromanspacer{} จตฺตาริ อารมฺมณานิ.\csromanspacer{} ชรามรเณ ญาณํ, ชรามรณสมุทเย ญาณํ, ชรามรณนิโรเธ ญาณํ, ชรามรณนิโรธคามินิยา ปฏิปทาย ญาณํ.\csromanspacer{} ชาติยา ญาณํ ฯเปฯ ภเว ญาณํ ฯเปฯ อุปาทาเน ญาณํ ฯเปฯ ตณฺหาย ญาณํ ฯเปฯ เวทนาย ญาณํ ฯเปฯ ผสฺเส ญาณํ ฯเปฯ สฬายตเน ญาณํ ฯเปฯ นามรูเป ญาณํ ฯเปฯ วิญฺญาเณ ญาณํ ฯเปฯ สงฺขาเรสุ ญาณํ, สงฺขารสมุทเย ญาณํ, สงฺขารนิโรเธ ญาณํ, สงฺขารนิโรธคามินิยา ปฏิปทาย ญาณํ.\csromanspacer{} เอวํ จตุพฺพิเธน ญาณวตฺถุ.}

\csromanmatikaanchor{mtk.152}
\csromantocmarkref{subparagraph}{5. ปญฺจกมาติกา}{mtk.152}
\bookmark[dest={mtk.152},level=5]{5. ปญฺจกมาติกา}
\csromanheader{6}{5. ปญฺจกมาติกา}
\proseitem{755}{ปญฺจวิเธน ญาณวตฺถุ, ปญฺจงฺคิโก สมฺมาสมาธิ, ปญฺจญาณิโก สมฺมาสมาธิ.\csromanspacer{} เอวํ ปญฺจวิเธน ญาณวตฺถุ.}

\csromanmatikaanchor{mtk.153}
\csromantocmarkref{subparagraph}{6. ฉกฺกมาติกา}{mtk.153}
\bookmark[dest={mtk.153},level=5]{6. ฉกฺกมาติกา}
\csromanheader{6}{6. ฉกฺกมาติกา}
\vspace{\tipitakaparskip}
\csromancenter{ฉพฺพิเธน ญาณวตฺถุ, ฉสุ อภิญฺญาสุ ปญฺญา, เอวํ ฉพฺพิเธน ญาณวตฺถุ.}
\csromanmatikaanchor{mtk.154}
\csromantocmarkref{subparagraph}{7. สตฺตกมาติกา}{mtk.154}
\bookmark[dest={mtk.154},level=5]{7. สตฺตกมาติกา}
\csromanheader{6}{7. สตฺตกมาติกา}
\proseitem{757}{สตฺตวิเธน ญาณวตฺถุ, สตฺตสตฺตติ ญาณวตฺถูนิ, เอวํ สตฺตวิเธน ญาณวตฺถุ.}

\csromanmatikaanchor{mtk.155}
\csromantocmarkref{subparagraph}{8. อฏฺฐกมาติกา}{mtk.155}
\bookmark[dest={mtk.155},level=5]{8. อฏฺฐกมาติกา}
\csromanheader{6}{8. อฏฺฐกมาติกา}
\proseitem{758}{อฏฺฐวิเธน ญาณวตฺถุ, จตูสุ มคฺเคสุ จตูสุ ผเลสุ ปญฺญา, เอวํ อฏฺฐวิเธน ญาณวตฺถุ.}

\csromanmatikaanchor{mtk.156}
\csromantocmarkref{subparagraph}{9. นวกมาติกา}{mtk.156}
\bookmark[dest={mtk.156},level=5]{9. นวกมาติกา}
\csromanheader{6}{9. นวกมาติกา}
\proseitem{759}{นววิเธน ญาณวตฺถุ, นวสุ อนุปุพฺพวิหารสมาปตฺตีสุ ปญฺญา, เอวํ นววิเธน ญาณวตฺถุ.}

\csromanmatikaanchor{mtk.157}
\csromantocmarkref{subparagraph}{10. ทสกมาติกา}{mtk.157}
\bookmark[dest={mtk.157},level=5]{10. ทสกมาติกา}
\csromanheader{6}{10. ทสกมาติกา}
\proseitem{760}{ทสวิเธน ญาณวตฺถุ, ทส ตถาคตสฺส ตถาคตพลานิ, เยหิ พเลหิ สมนฺนาคโต ตถาคโต อาสภํ ฐานํ ปฏิชานาติ,}

\csromanheader{2}{16. ญาณวิภงฺค}
\prose{\csromanfolio{329}ปริสาสุ สีหนาทํ นทติ, พฺรหฺมจกฺกํ ปวตฺเตติ.\csromanspacer{} กตมานิ ทส, อิธ ตถาคโต ฐานญฺจ ฐานโต อฏฺฐานญฺจ อฏฺฐานโต ยถาภูตํ ปชานาติ, ยมฺปิ ตถาคโต ฐานญฺจ ฐานโต อฏฺฐานญฺจ อฏฺฐานโต ยถาภูตํ ปชานาติ, อิทมฺปิ ตถาคตสฺส ตถาคตพลํ โหติ.\csromanspacer{} ยํ พลํ อาคมฺม ตถาคโต อาสภํ ฐานํ ปฏิชานาติ, ปริสาสุ สีหนาทํ นทติ, พฺรหฺมจกฺกํ ปวตฺเตติ. (1)}

\prose{ปุน จปรํ ตถาคโต อตีตานาคตปจฺจุปฺปนฺนานํ กมฺมสมาทานานํ ฐานโส เหตุโส วิปากํ ยถาภูตํ ปชานาติ, ยมฺปิ ตถาคโต อตีตานาคตปจฺจุปฺปนฺนานํ กมฺมสมาทานานํ ฐานโส เหตุโส วิปากํ ยถาภูตํ ปชานาติ, อิทมฺปิ ตถาคตสฺส ตถาคตพลํ โหติ.\csromanspacer{} ยํ พลํ อาคมฺม ตถาคโต อาสภํ ฐานํ ปฏิชานาติ, ปริสาสุ สีหนาทํ นทติ, พฺรหฺมจกฺกํ ปวตฺเตติ. (2)}

\prose{ปุน จปรํ ตถาคโต สพฺพตฺถคามินิํ ปฏิปทํ ยถาภูตํ ปชานาติ, ยมฺปิ ตถาคโต สพฺพตฺถคามินิํ ปฏิปทํ ยถาภูตํ ปชานาติ, อิทมฺปิ ตถาคตสฺส ตถาคตพลํ โหติ.\csromanspacer{} ยํ พลํ อาคมฺม ตถาคโต อาสภํ ฐานํ ปฏิชานาติ, ปริสาสุ สีหนาทํ นทติ, พฺรหฺมจกฺกํ ปวตฺเตติ. (3)}

\prose{ปุน จปรํ ตถาคโต อเนกธาตุ นานาธาตุ โลกํ\footnote{อเนกธาตุํ นานาธาตุํ โลกํ (สฺยา) ม 1. 100 ปิฏฺเฐปิ.} ยถาภูตํ ปชานาติ, ยมฺปิ ตถาคโต อเนกธาตุ นานาธาตุ โลกํ ยถาภูตํ ปชานาติ, อิทมฺปิ ตถาคตสฺส ตถาคตพลํ โหติ.\csromanspacer{} ยํ พลํ อาคมฺม ตถาคโต อาสภํ ฐานํ ปฏิชานาติ, ปริสาสุ สีหนาทํ นทติ, พฺรหฺมจกฺกํ ปวตฺเตติ. (4)}

\prose{ปุน จปรํ ตถาคโต สตฺตานํ นานาธิมุตฺติกตํ ยถาภูตํ ปชานาติ, ยมฺปิ ตถาคโต สตฺตานํ นานาธิมุตฺติกตํ ยถาภูตํ ปชานาติ, อิทมฺปิ ตถาคตสฺส ตถาคตพลํ โหติ.\csromanspacer{} ยํ พลํ อาคมฺม ตถาคโต อาสภํ ฐานํ ปฏิชานาติ, ปริสาสุ สีหนาทํ นทติ, พฺรหฺมจกฺกํ ปวตฺเตติ. (5)}

\prose{ปุน จปรํ ตถาคโต ปรสตฺตานํ ปรปุคฺคลานํ อินฺทฺริยปโรปริยตฺตํ ยถาภูตํ ปชานาติ, ยมฺปิ ตถาคโต ปรสตฺตานํ ปรปุคฺคลานํ อินฺทฺริยปโรปริยตฺตํ ยถาภูตํ ปชานาติ, อิทมฺปิ ตถาคตสฺส}

\prose{\csromanfolio{330}ตถาคตพลํ โหติ.\csromanspacer{} ยํ พลํ อาคมฺม ตถาคโต อาสภํ ฐานํ ปฏิชานาติ, ปริสาสุ สีหนาทํ นทติ, พฺรหฺมจกฺกํ ปวตฺเตติ. (6)}

\prose{ปุน จปรํ ตถาคโต ฌานวิโมกฺขสมาธิสมาปตฺตีนํ สํกิเลสํ โวทานํ วุฏฺฐานํ ยถาภูตํ ปชานาติ, ยมฺปิ ตถาคโต ฌานวิโมกฺขสมาธิสมาปตฺตีนํ สํกิเลสํ โวทานํ วุฏฺฐานํ ยถาภูตํ ปชานาติ, อิทมฺปิ ตถาคตสฺส ตถาคตพลํ โหติ.\csromanspacer{} ยํ พลํ อาคมฺม ตถาคโต อาสภํ ฐานํ ปฏิชานาติ, ปริสาสุ สีหนาทํ นทติ, พฺรหฺมจกฺกํ ปวตฺเตติ. (7)}

\prose{ปุน จปรํ ตถาคโต ปุพฺเพนิวาสานุสฺสติํ ยถาภูตํ ปชานาติ, ยมฺปิ ตถาคโต ปุพฺเพนิวาสานุสฺสติํ ยถาภูตํ ปชานาติ, อิทมฺปิ ตถาคตสฺส ตถาคตพลํ โหติ.\csromanspacer{} ยํ พลํ อาคมฺม ตถาคโต อาสภํ ฐานํ ปฏิชานาติ, ปริสาสุ สีหนาทํ นทติ, พฺรหฺมจกฺกํ ปวตฺเตติ. (8)}

\prose{ปุน จปรํ ตถาคโต สตฺตานํ จุตูปปาตํ ยถาภูตํ ปชานาติ, ยมฺปิ ตถาคโต สตฺตานํ จุตูปปาตํ ยถาภูตํ ปชานาติ, อิทมฺปิ ตถาคตสฺส ตถาคตพลํ โหติ.\csromanspacer{} ยํ พลํ อาคมฺม ตถาคโต อาสภํ ฐานํ ปฏิชานาติ, ปริสาสุ สีหนาทํ นทติ, พฺรหฺมจกฺกํ ปวตฺเตติ. (9)}

\prose{ปุน จปรํ ตถาคโต อาสวานํ ขยํ ยถาภูตํ ปชานาติ, ยมฺปิ ตถาคโต อาสวานํ ขยํ ยถาภูตํ ปชานาติ, อิทมฺปิ ตถาคตสฺส ตถาคตพลํ โหติ.\csromanspacer{} ยํ พลํ อาคมฺม ตถาคโต อาสภํ ฐานํ ปฏิชานาติ, ปริสาสุ สีหนาทํ นทติ, พฺรหฺมจกฺกํ ปวตฺเตติ.\csromanspacer{} อิมานิ ทส ตถาคตสฺส ตถาคตพลานิ.\csromanspacer{} เยหิ พเลหิ สมนฺนาคโต ตถาคโต อาสภํ ฐานํ ปฏิชานาติ, ปริสาสุ สีหนาทํ นทติ, พฺรหฺมจกฺกํ ปวตฺเตติ. (10)}

\vspace{\tipitakaparskip}
\csromancenter{เอวํ ทสวิเธน ญาณวตฺถุ.}
\csromanheader{2}{มาติกา.}
\csromanheader{2}{16. ญาณวิภงฺค}
\csromanmatikaanchor{mtk.158}
\csromantocmarkref{subparagraph}{1. เอกกนิทฺเทส}{mtk.158}
\bookmark[dest={mtk.158},level=5]{1. เอกกนิทฺเทส}
\csromanheader{6}{1. เอกกนิทฺเทส}
\proseitem{761}{\csromanfolio{331}ปญฺจ วิญฺญาณา น เหตุเมว, อเหตุกเมว, เหตุวิปฺปยุตฺตเมว, สปฺปจฺจยเมว, สงฺขตเมว, อรูปเมว, โลกิยเมว, สาสวเมว, สํโยชนิยเมว, คนฺถนิยเมว, โอฆนิยเมว, โยคนิยเมว, นีวรณิยเมว, ปรามฏฺฐเมว, อุปาทานิยเมว, สํกิเลสิกเมว, อพฺยากตเมว, สารมฺมณเมว, อเจตสิกเมว, วิปากเมว, อุปาทินฺนุปาทานิยเมว, อสํกิลิฏฺฐสํกิเลสิกเมว, น สวิตกฺกสวิจารเมว, น อวิตกฺกวิจารมตฺตเมว, อวิตกฺก-อวิจารเมว, น ปีติสหคตเมว, เนว ทสฺสเนน น ภาวนาย ปหาตพฺพเมว, เนว ทสฺสเนน น ภาวนาย ปหาตพฺพเหตุกเมว, เนวาจยคามินาปจยคามิเมว, เนวเสกฺขนาเสกฺขเมว, ปริตฺตเมว, กามาวจรเมว, น รูปาวจรเมว, น อรูปาวจรเมว, ปริยาปนฺนเมว, โน อปริยาปนฺนเมว, อนิยตเมว, อนิยฺยานิกเมว, อุปฺปนฺนํ มโนวิญฺญาณวิญฺเญยฺยเมว, อนิจฺจเมว, ชราภิภูตเมว.}

\proseitem{762}{ปญฺจ วิญฺญาณา อุปฺปนฺนวตฺถุกา, อุปฺปนฺนารมฺมณาติ อุปฺปนฺนสฺมิํ วตฺถุสฺมิํ อุปฺปนฺเน อารมฺมเณ อุปฺปชฺชนฺติ.}

\prose{ปุเรชาตวตฺถุกา, ปุเรชาตารมฺมณาติ ปุเรชาตสฺมิํ วตฺถุสฺมิํ ปุเรชาเต อารมฺมเณ อุปฺปชฺชนฺติ.}

\prose{อชฺฌตฺติกวตฺถุกา, พาหิรารมฺมณาติ ปญฺจนฺนํ วิญฺญาณานํ วตฺถุ อชฺฌตฺติกา, อารมฺมณา พาหิรา.}

\prose{อสมฺภินฺนวตฺถุกา, อสมฺภินฺนารมฺมณาติ อสมฺภินฺนสฺมิํ วตฺถุสฺมิํ อสมฺภินฺเน อารมฺมเณ อุปฺปชฺชนฺติ.}

\prose{นานาวตฺถุกา, นานารมฺมณาติ อญฺญํ จกฺขุวิญฺญาณสฺส วตฺถุ จ อารมฺมณญฺจ, อญฺญํ โสตวิญฺญาณสฺส วตฺถุ จ อารมฺมณญฺจ, อญฺญํ ฆานวิญฺญาณสฺส วตฺถุ จ อารมฺมณญฺจ, อญฺญํ ชิวฺหาวิญฺญาณสฺส วตฺถุ จ อารมฺมณญฺจ, อญฺญํ กายวิญฺญาณสฺส วตฺถุ จ อารมฺมณญฺจ.}

\proseitem{763}{น อญฺญมญฺญสฺส โคจรวิสยํ ปจฺจนุโภนฺตีติ จกฺขุวิญฺญาณสฺส โคจรวิสยํ โสตวิญฺญาณํ น ปจฺจนุโภติ, โสตวิญฺญาณสฺส โคจรวิสยมฺปิ จกฺขุวิญฺญาณํ น ปจฺจนุโภติ.\csromanspacer{} จกฺขุวิญฺญาณสฺส โคจรวิสยํ ฆานวิญฺญาณํ น ปจฺจนุโภติ, ฆานวิญฺญาณสฺส โคจรวิสยมฺปิ \csromanfolio{332}จกฺขุวิญฺญาณํ น ปจฺจนุโภติ.\csromanspacer{} จกฺขุวิญฺญาณสฺส โคจรวิสยํ ชิวฺหาวิญฺญาณํ น ปจฺจนุโภติ, ชิวฺหาวิญฺญาณสฺส โคจรวิสยมฺปิ จกฺขุวิญฺญาณํ น ปจฺจนุโภติ.\csromanspacer{} จกฺขุวิญฺญาณสฺส โคจรวิสยํ กายวิญฺญาณํ น ปจฺจนุโภติ, กายวิญฺญาณสฺส โคจรวิสยมฺปิ จกฺขุวิญฺญาณํ น ปจฺจนุโภติ.\csromanspacer{} โสตวิญฺญาณสฺส ฯเปฯ.\csromanspacer{} ฆานวิญฺญาณสฺส ฯเปฯ.\csromanspacer{} ชิวฺหาวิญฺญาณสฺส ฯเปฯ.\csromanspacer{} กายวิญฺญาณสฺส โคจรวิสยํ จกฺขุวิญฺญาณํ น ปจฺจนุโภติ, จกฺขุวิญฺญาณสฺส โคจรวิสยมฺปิ กายวิญฺญาณํ น ปจฺจนุโภติ.\csromanspacer{} กายวิญฺญาณสฺส โคจรวิสยํ โสตวิญฺญาณํ น ปจฺจนุโภติ, โสตวิญฺญาณสฺส โคจรวิสยมฺปิ กายวิญฺญาณํ น ปจฺจนุโภติ.\csromanspacer{} กายวิญฺญาณสฺส โคจรวิสยํ ฆานวิญฺญาณํ น ปจฺจนุโภติ, ฆานวิญฺญาณสฺส โคจรวิสยมฺปิ กายวิญฺญาณํ น ปจฺจนุโภติ.\csromanspacer{} กายวิญฺญาณสฺส โคจรวิสยํ ชิวฺหาวิญฺญาณํ น ปจฺจนุโภติ, ชิวฺหาวิญฺญาณสฺส โคจรวิสยมฺปิ กายวิญฺญาณํ น ปจฺจนุโภติ.}

\proseitem{764}{น อสมนฺนาหารา อุปฺปชฺชนฺตีติ สมนฺนาหรนฺตสฺส อุปฺปชฺชนฺติ.}

\prose{น อมนสิการา อุปฺปชฺชนฺตีติ มนสิกโรนฺตสฺส อุปฺปชฺชนฺติ.}

\prose{น อพฺโพกิณฺณา อุปฺปชฺชนฺตีติ น ปฏิปาฏิยา อุปฺปชฺชนฺติ.}

\prose{น อปุพฺพํ อจริมํ อุปฺปชฺชนฺตีติ น เอกกฺขเณ อุปฺปชฺชนฺติ.}

\proseitem{765}{น อญฺญมญฺญสฺส สมนนฺตรา อุปฺปชฺชนฺตีติ จกฺขุวิญฺญาณสฺส อุปฺปนฺนสมนนฺตรา โสตวิญฺญาณํ น อุปฺปชฺชติ, โสตวิญฺญาณสฺส อุปฺปนฺนสมนนฺตราปิ จกฺขุวิญฺญาณํ น อุปฺปชฺชติ.\csromanspacer{} จกฺขุวิญฺญาณสฺส อุปฺปนฺนสมนนฺตรา ฆานวิญฺญาณํ น อุปฺปชฺชติ, ฆานวิญฺญาณสฺส อุปฺปนฺนสมนนฺตราปิ จกฺขุวิญฺญาณํ น อุปฺปชฺชติ.\csromanspacer{} จกฺขุวิญฺญาณสฺส อุปฺปนฺนสมนนฺตรา ชิวฺหาวิญฺญาณํ น อุปฺปชฺชติ, ชิวฺหาวิญฺญาณสฺส อุปฺปนฺนสมนนฺตราปิ จกฺขุวิญฺญาณํ น อุปฺปชฺชติ.\csromanspacer{} จกฺขุวิญฺญาณสฺส อุปฺปนฺนสมนนฺตรา กายวิญฺญาณํ น อุปฺปชฺชติ, กายวิญฺญาณสฺส อุปฺปนฺนสมนนฺตราปิ จกฺขุวิญฺญาณํ น อุปฺปชฺชติ.\csromanspacer{} โสตวิญฺญาณสฺส ฯเปฯ.\csromanspacer{} ฆานวิญฺญาณสฺส ฯเปฯ.\csromanspacer{} ชิวฺหาวิญฺญาณสฺส ฯเปฯ.\csromanspacer{} กายวิญฺญาณสฺส อุปฺปนฺนสมนนฺตรา จกฺขุวิญฺญาณํ น อุปฺปชฺชติ, จกฺขุวิญฺญาณสฺส อุปฺปนฺนสมนนฺตราปิ กายวิญฺญาณํ น อุปฺปชฺชติ.\csromanspacer{} กายวิญฺญาณสฺส อุปฺปนฺนสมนนฺตรา โสตวิญฺญาณํ น อุปฺปชฺชติ, โสตวิญฺญาณสฺส}

\csromanheader{2}{16. ญาณวิภงฺค}
\prose{\csromanfolio{333}อุปฺปนฺนสมนนฺตราปิ กายวิญฺญาณํ น อุปฺปชฺชติ.\csromanspacer{} กายวิญฺญาณสฺส อุปฺปนฺนสมนนฺตรา ฆานวิญฺญาณํ น อุปฺปชฺชติ, ฆานวิญฺญาณสฺส อุปฺปนฺนสมนนฺตราปิ กายวิญฺญาณํ น อุปฺปชฺชติ.\csromanspacer{} กายวิญฺญาณสฺส อุปฺปนฺนสมนนฺตรา ชิวฺหาวิญฺญาณํ น อุปฺปชฺชติ, ชิวฺหาวิญฺญาณสฺส อุปฺปนฺนสมนนฺตราปิ กายวิญฺญาณํ น อุปฺปชฺชติ.}

\proseitem{766}{ปญฺจ วิญฺญาณา อนาโภคาติ ปญฺจนฺนํ วิญฺญาณานํ นตฺถิ อาวฏฺฏนา วา อาโภโค วา สมนฺนาหาโร วา มนสิกาโร วา.}

\prose{ปญฺจหิ วิญฺญาเณหิ น กญฺจิ ธมฺมํ ปฏิวิชานาตีติ ปญฺจหิ วิญฺญาเณหิ น กญฺจิ ธมฺมํ ปฏิวิชานาติ.}

\prose{อญฺญตฺร อภินิปาตมตฺตาติ อญฺญตฺร อาปาตมตฺตา.}

\prose{ปญฺจนฺนํ วิญฺญาณานํ สมนนฺตราปิ น กญฺจิ ธมฺมํ ปฏิวิชานาตีติ ปญฺจนฺนํ วิญฺญาณานํ สมนนฺตรา มโนธาตุยาปิ น กญฺจิ ธมฺมํ ปฏิวิชานาติ.}

\prose{ปญฺจหิ วิญฺญาเณหิ น กญฺจิ อิริยาปถํ กปฺเปตีติ ปญฺจหิ วิญฺญาเณหิ น กญฺจิ อิริยาปถํ กปฺเปติ คมนํ วา ฐานํ วา นิสชฺชํ วา เสยฺยํ วา.}

\prose{ปญฺจนฺนํ วิญฺญาณานํ สมนนฺตราปิ น กญฺจิ อิริยาปถํ กปฺเปตีติ ปญฺจนฺนํ วิญฺญาณานํ สมนนฺตรา มโนธาตุยาปิ น กญฺจิ อิริยาปถํ กปฺเปติ คมนํ วา ฐานํ วา นิสชฺชํ วา เสยฺยํ วา.}

\prose{ปญฺจหิ วิญฺญาเณหิ น กายกมฺมํ น วจีกมฺมํ ปฏฺฐเปตีติ ปญฺจหิ วิญฺญาเณหิ น กายกมฺมํ น วจีกมฺมํ ปฏฺฐเปติ.}

\prose{ปญฺจนฺนํ วิญฺญาณานํ สมนนฺตราปิ น กายกมฺมํ น วจีกมฺมํ ปฏฺฐเปตีติ ปญฺจนฺนํ วิญฺญาณานํ สมนนฺตรา มโนธาตุยาปิ น กายกมฺมํ น วจีกมฺมํ ปฏฺฐเปติ.}

\prose{ปญฺจหิ วิญฺญาเณหิ น กุสลากุสลํ ธมฺมํ สมาทิยตีติ ปญฺจหิ วิญฺญาเณหิ น กุสลากุสลํ ธมฺมํ สมาทิยติ.}

\prose{ปญฺจนฺนํ วิญฺญาณานํ สมนนฺตราปิ น กุสลากุสลํ ธมฺมํ สมาทิยตีติ ปญฺจนฺนํ วิญฺญาณานํ สมนนฺตรา มโนธาตุยาปิ น กุสลากุสลํ ธมฺมํ สมาทิยติ.}

\prose{ปญฺจหิ วิญฺญาเณหิ น สมาปชฺชติ น วุฏฺฐาตีติ ปญฺจหิ วิญฺญาเณหิ น สมาปชฺชติ น วุฏฺฐาติ.}

\prose{\csromanfolio{334}ปญฺจนฺนํ วิญฺญาณานํ สมนนฺตราปิ น สมาปชฺชติ น วุฏฺฐาตีติ ปญฺจนฺนํ วิญฺญาณานํ สมนนฺตรา มโนธาตุยาปิ น สมาปชฺชติ น วุฏฺฐาติ.}

\prose{ปญฺจหิ วิญฺญาเณหิ น จวติ น อุปฺปชฺชตีติ ปญฺจหิ วิญฺญาเณหิ น จวติ น อุปฺปชฺชติ.}

\prose{ปญฺจนฺนํ วิญฺญาณานํ สมนนฺตราปิ น จวติ น อุปฺปชฺชตีติ ปญฺจนฺนํ วิญฺญาณานํ สมนนฺตรา มโนธาตุยาปิ น จวติ น อุปฺปชฺชติ.}

\prose{ปญฺจหิ วิญฺญาเณหิ น สุปติ น ปฏิพุชฺฌติ น สุปินํ ปสฺสตีติ ปญฺจหิ วิญฺญาเณหิ น สุปติ น ปฏิพุชฺฌติ น สุปินํ ปสฺสติ.}

\prose{ปญฺจนฺนํ วิญฺญาณานํ สมนนฺตราปิ น สุปติ น ปฏิพุชฺฌติ น สุปินํ ปสฺสตีติ ปญฺจนฺนํ วิญฺญาณานํ สมนนฺตรา มโนธาตุยาปิ น สุปติ น ปฏิพุชฺฌติ น สุปินํ ปสฺสติ, เอวํ ยาถาวกวตฺถุวิภาวนา ปญฺญา, เอวํ เอกวิเธน ญาณวตฺถุ.}

\csromanheader{2}{เอกกํ.}
\csromansectionrule
\csromanmatikaanchor{mtk.159}
\csromantocmarkref{subparagraph}{2. ทุกนิทฺเทส}{mtk.159}
\bookmark[dest={mtk.159},level=5]{2. ทุกนิทฺเทส}
\csromanheader{6}{2. ทุกนิทฺเทส}
\proseitem{767}{ตีสุ ภูมีสุ กุสลาพฺยากเต ปญฺญา โลกิยา ปญฺญา, จตูสุ มคฺเคสุ จตูสุ ผเลสุ ปญฺญา โลกุตฺตรา ปญฺญา.\csromanspacer{} สพฺพาว ปญฺญา เกนจิ วิญฺเญยฺยา, เกนจิ น วิญฺเญยฺยา.}

\prose{ตีสุ ภูมีสุ กุสลาพฺยากเต ปญฺญา สาสวา ปญฺญา, จตูสุ มคฺเคสุ จตูสุ ผเลสุ ปญฺญา อนาสวา ปญฺญา.\csromanspacer{} ตีสุ ภูมีสุ กุสลาพฺยากเต ปญฺญา อาสววิปฺปยุตฺตา สาสวา ปญฺญา, จตูสุ มคฺเคสุ จตูสุ ผเลสุ ปญฺญา อาสววิปฺปยุตฺตา อนาสวา ปญฺญา.}

\prose{ตีสุ ภูมีสุ กุสลาพฺยากเต ปญฺญา สํโยชนิยา ปญฺญา, จตูสุ มคฺเคสุ จตูสุ ผเลสุ ปญฺญา อสํโยชนิยา ปญฺญา.\csromanspacer{} ตีสุ ภูมีสุ กุสลาพฺยากเต ปญฺญา สํโยชนวิปฺปยุตฺตา สํโยชนิยา ปญฺญา, จตูสุ มคฺเคสุ จตูสุ ผเลสุ ปญฺญา สํโยชนวิปฺปยุตฺตา อสํโยชนิยา ปญฺญา.}

\csromanheader{2}{16. ญาณวิภงฺค}
\prose{\csromanfolio{335}ตีสุ ภูมีสุ กุสลาพฺยากเต ปญฺญา คนฺถนิยา ปญฺญา, จตูสุ มคฺเคสุ จตูสุ ผเลสุ ปญฺญา อคนฺถนิยา ปญฺญา.\csromanspacer{} ตีสุ ภูมีสุ กุสลาพฺยากเต ปญฺญา คนฺถวิปฺปยุตฺตา คนฺถนิยา ปญฺญา, จตูสุ มคฺเคสุ จตูสุ ผเลสุ ปญฺญา คนฺถวิปฺปยุตฺตา อคนฺถนิยา ปญฺญา.}

\prose{ตีสุ ภูมีสุ กุสลาพฺยากเต ปญฺญา โอฆนิยา ปญฺญา, จตูสุ มคฺเคสุ จตูสุ ผเลสุ ปญฺญา อโนฆนิยา ปญฺญา.\csromanspacer{} ตีสุ ภูมีสุ กุสลาพฺยากเต ปญฺญา โอฆวิปฺปยุตฺตา โอฆนิยา ปญฺญา, จตูสุ มคฺเคสุ จตูสุ ผเลสุ ปญฺญา โอฆวิปฺปยุตฺตา อโนฆนิยา ปญฺญา.}

\prose{ตีสุ ภูมีสุ กุสลาพฺยากเต ปญฺญา โยคนิยา ปญฺญา, จตูสุ มคฺเคสุ จตูสุ ผเลสุ ปญฺญา อโยคนิยา ปญฺญา.\csromanspacer{} ตีสุ ภูมีสุ กุสลาพฺยากเต ปญฺญา โยควิปฺปยุตฺตา โยคนิยา ปญฺญา, จตูสุ มคฺเคสุ จตูสุ ผเลสุ ปญฺญา โยควิปฺปยุตฺตา อโยคนิยา ปญฺญา.}

\prose{ตีสุ ภูมีสุ กุสลาพฺยากเต ปญฺญา นีวรณิยา ปญฺญา, จตูสุ มคฺเคสุ จตูสุ ผเลสุ ปญฺญา อนีวรณิยา ปญฺญา.\csromanspacer{} ตีสุ ภูมีสุ กุสลาพฺยากเต ปญฺญา นีวรณวิปฺปยุตฺตา นีวรณิยา ปญฺญา, จตูสุ มคฺเคสุ จตูสุ ผเลสุ ปญฺญา นีวรณวิปฺปยุตฺตา อนีวรณิยา ปญฺญา.}

\prose{ตีสุ ภูมีสุ กุสลาพฺยากเต ปญฺญา ปรามฏฺฐา ปญฺญา, จตูสุ มคฺเคสุ จตูสุ ผเลสุ ปญฺญา อปรามฏฺฐา ปญฺญา.\csromanspacer{} ตีสุ ภูมีสุ กุสลาพฺยากเต ปญฺญา ปรามาสวิปฺปยุตฺตา ปรามฏฺฐา ปญฺญา, จตูสุ มคฺเคสุ จตูสุ ผเลสุ ปญฺญา ปรามาสวิปฺปยุตฺตา อปรามฏฺฐา ปญฺญา.}

\prose{ตีสุ ภูมีสุ วิปาเก ปญฺญา อุปาทินฺนา ปญฺญา, ตีสุ ภูมีสุ กุสเล ตีสุ ภูมีสุ กิริยาพฺยากเต จตูสุ มคฺเคสุ จตูสุ ผเลสุ ปญฺญา อนุปาทินฺนา ปญฺญา.}

\prose{ตีสุ ภูมีสุ กุสลาพฺยากเต ปญฺญา อุปาทานิยา ปญฺญา, จตูสุ มคฺเคสุ จตูสุ ผเลสุ ปญฺญา อนุปาทานิยา ปญฺญา.\csromanspacer{} ตีสุ ภูมีสุ กุสลาพฺยากเต ปญฺญา อุปาทานวิปฺปยุตฺตา อุปาทานิยา ปญฺญา, จตูสุ มคฺเคสุ จตูสุ ผเลสุ ปญฺญา อุปาทานวิปฺปยุตฺตา อนุปาทานิยา ปญฺญา.}

\prose{ตีสุ ภูมีสุ กุสลาพฺยากเต ปญฺญา สํกิเลสิกา ปญฺญา, จตูสุ มคฺเคสุ จตูสุ ผเลสุ ปญฺญา อสํกิเลสิกา ปญฺญา.\csromanspacer{} ตีสุ ภูมีสุ}

\prose{\csromanfolio{336}กุสลาพฺยากเต ปญฺญา กิเลสวิปฺปยุตฺตา สํกิเลสิกา ปญฺญา, จตูสุ มคฺเคสุ จตูสุ ผเลสุ ปญฺญา กิเลสวิปฺปยุตฺตา อสํกิเลสิกา ปญฺญา.}

\prose{วิตกฺกสมฺปยุตฺตา ปญฺญา สวิตกฺกา ปญฺญา, วิตกฺกวิปฺปยุตฺตา ปญฺญา อวิตกฺกา ปญฺญา.\csromanspacer{} วิจารสมฺปยุตฺตา ปญฺญา สวิจารา ปญฺญา, วิจารวิปฺปยุตฺตา ปญฺญา อวิจารา ปญฺญา ปีติสมฺปยุตฺตา ปญฺญา สปฺปีติกา ปญฺญา, ปีติวิปฺปยุตฺตา ปญฺญา อปฺปีติกา ปญฺญา, ปีติสมฺปยุตฺตา ปญฺญา ปีติสหคตา ปญฺญา, ปีติวิปฺปยุตฺตา ปญฺญา น ปีติสหคตา ปญฺญา.\csromanspacer{} สุขสมฺปยุตฺตา ปญฺญา สุขสหคตา ปญฺญา, สุขวิปฺปยุตฺตา ปญฺญา น สุขสหคตา ปญฺญา.\csromanspacer{} อุเปกฺขาสมฺปยุตฺตา ปญฺญา อุเปกฺขาสหคตา ปญฺญา.\csromanspacer{} อุเปกฺขาวิปฺปยุตฺตา ปญฺญา น อุเปกฺขาสหคตา ปญฺญา.}

\prose{กามาวจรกุสลาพฺยากเต ปญฺญา กามาวจรา ปญฺญา, รูปาวจรา ปญฺญา อรูปาวจรา ปญฺญา อปริยาปนฺนา ปญฺญา น กามาวจรา ปญฺญา.\csromanspacer{} รูปาวจรกุสลาพฺยากเต ปญฺญา รูปาวจรา ปญฺญา, กามาวจรา ปญฺญา อรูปาวจรา ปญฺญา อปริยาปนฺนา ปญฺญา น รูปาวจรา ปญฺญา.\csromanspacer{} อรูปาวจรกุสลาพฺยากเต ปญฺญา อรูปาวจรา ปญฺญา, กามาวจรา ปญฺญา รูปาวจรา ปญฺญา อปริยาปนฺนา ปญฺญา น อรูปาวจรา ปญฺญา.}

\prose{ตีสุ ภูมีสุ กุสลาพฺยากเต ปญฺญา ปริยาปนฺนา ปญฺญา, จตูสุ มคฺเคสุ จตูสุ ผเลสุ ปญฺญา, อปริยาปนฺนา ปญฺญา.\csromanspacer{} จตูสุ มคฺเคสุ ปญฺญา นิยฺยานิกา ปญฺญา, ตีสุ ภูมีสุ กุสเล จตูสุ ภูมิสุ วิปาเก ตีสุ ภูมีสุ กิริยาพฺยากเต ปญฺญา อนิยฺยานิกา ปญฺญา.\csromanspacer{} จตูสุ มคฺเคสุ ปญฺญา นิยตา ปญฺญา, ตีสุ ภูมีสุ กุสเล จตูสุ ภูมีสุ วิปาเก ตีสุ ภูมีสุ กิริยาพฺยากเต ปญฺญา อนิยตา ปญฺญา.\csromanspacer{} ตีสุ ภูมีสุ กุสลาพฺยากเต ปญฺญา ส-อุตฺตรา ปญฺญา, จตูสุ มคฺเคสุ จตูสุ ผเลสุ ปญฺญา อนุตฺตรา ปญฺญา.}

\prose{ตตฺถ กตมา อตฺถชาปิกา ปญฺญา, จตูสุ ภูมีสุ กุสเล อรหโต อภิญฺญํ อุปฺปาเทนฺตสฺส สมาปตฺติํ อุปฺปาเทนฺตสฺส กิริยาพฺยากเต ปญฺญา อตฺถชาปิกา ปญฺญา, จตูสุ ภูมีสุ วิปาเก อรหโต อุปฺปนฺนาย อภิญฺญาย อุปฺปนฺนาย สมาปตฺติยา กิริยาพฺยากเต ปญฺญา ชาปิตตฺถา ปญฺญา.\csromanspacer{} เอวํ ทุวิเธน ญาณวตฺถุ.}

\vspace{\tipitakaparskip}
\csromancenter{ทุกํ.}
\csromanheader{2}{16. ญาณวิภงฺค}
\csromanmatikaanchor{mtk.160}
\csromantocmarkref{subparagraph}{3. ติกนิทฺเทส}{mtk.160}
\bookmark[dest={mtk.160},level=5]{3. ติกนิทฺเทส}
\csromanheader{6}{3. ติกนิทฺเทส}
\proseitem{768}{\csromanfolio{337}ตตฺถ กตมา จินฺตามยา ปญฺญา, โยควิหิเตสุ วา กมฺมายตเนสุ โยควิหิเตสุ วา สิปฺปายตเนสุ โยควิหิเตสุ วา วิชฺชาฏฺฐาเนสุ กมฺมสฺสกตํ วา สจฺจานุโลมิกํ วา รูปํ อนิจฺจนฺติ วา, เวทนา ฯเปฯ.\csromanspacer{} สญฺญา.\csromanspacer{} สงฺขารา.\csromanspacer{} วิญฺญาณํ อนิจฺจนฺติ วา, ยํ เอวรูปิํ อนุโลมิกํ ขนฺติํ ทิฏฺฐิํ รุจิํ มุทิํ เปกฺขํ ธมฺมนิชฺฌานกฺขนฺติํ ปรโต อสฺสุตฺวา ปฏิลภติ, อยํ วุจฺจติ จินฺตามยา ปญฺญา.}

\prose{ตตฺถ กตมา สุตมยา ปญฺญา, โยควิหิเตสุ วา กมฺมายตเนสุ โยควิหิเตสุ วา สิปฺปายตเนสุ โยควิหิเตสุ วา วิชฺชาฏฺฐาเนสุ กมฺมสฺสกตํ วา สจฺจานุโลมิกํ วา รูปํ อนิจฺจนฺติ วา, เวทนา ฯเปฯ.\csromanspacer{} สญฺญา.\csromanspacer{} สงฺขารา.\csromanspacer{} วิญฺญาณํ อนิจฺจนฺติ วา, ยํ เอวรูปิํ อนุโลมิกํ ขนฺติํ ทิฏฺฐิํ รุจิํ มุทิํ เปกฺขํ ธมฺมนิชฺฌานกฺขนฺติํ ปรโต สุตฺวา ปฏิลภติ, อยํ วุจฺจติ สุตมยา ปญฺญา.}

\prose{สพฺพาปิ สมาปนฺนสฺส ปญฺญา หาวนามยา ปญฺญา.}

\proseitem{769}{ตตฺถ กตมา ทานมยา ปญฺญา, ทานํ อารพฺภ ทานาธิคจฺฉ ยา อุปฺปชฺชติ ปญฺญา ปชานนา ฯเปฯ อโมโห ธมฺมวิจโย สมฺมาทิฏฺฐิ, อยํ วุจฺจติ ทานมยา ปญฺญา.}

\prose{ตตฺถ กตมา สีลมยา ปญฺญา, สีลํ อารพฺภ สีลาธิคจฺฉ ยา อุปฺปชฺชติ ปญฺญา ปชานนา ฯเปฯ อโมโห ธมฺมวิจโย สมฺมาทิฏฺฐิ, อยํ วุจฺจติ สีลมยา ปญฺญา.}

\prose{สพฺพาปิ สมาปนฺนสฺส ปญฺญา ภาวนามยา ปญฺญา.}

\proseitem{770}{ตตฺถ กตมา อธิสีเล ปญฺญา, ปาติโมกฺขสํวรํ สํวรนฺตสฺส ยา อุปฺปชฺชติ ปญฺญา ปชานนา ฯเปฯ อโมโห ธมฺมวิจโย สมฺมาทิฏฺฐิ, อยํ วุจฺจติ อธิสีเล ปญฺญา.}

\prose{ตตฺถ กตมา อธิจิตฺเต ปญฺญา, รูปาวจรารูปาวจรสมาปตฺติํ สมาปชฺชนฺตสฺส ยา อุปฺปชฺชติ ปญฺญา ปชานนา ฯเปฯ อโมโห ธมฺมวิจโย สมฺมาทิฏฺฐิ, อยํ วุจฺจติ อธิจิตฺเต ปญฺญา.}

\prose{\csromanfolio{338}ตตฺถ กตมา อธิปญฺญาย ปญฺญา, จตูสุ มคฺเคสุ จตูสุ ผเลสุ ปญฺญา, อยํ วุจฺจติ อธิปญฺญาย ปญฺญา.}

\proseitem{771}{ตตฺถ กตมํ อายโกสลฺลํ, อิเม ธมฺเม มนสิกโรโต อนุปฺปนฺนา เจว อกุสลา ธมฺมา น อุปฺปชฺชนฺติ, อุปฺปนฺนา จ อกุสลา ธมฺมา ปหียนฺติ, อิเม วา ปนิเม ธมฺเม มนสิกโรโต อนุปฺปนฺนา เจว กุสลา ธมฺมา อุปฺปชฺชนฺติ, อุปฺปนฺนา จ กุสลา ธมฺมา ภิยฺโยภาวาย เวปุลฺลาย ภาวนาย ปาริปูริยา สํวตฺตนฺตีติ, ยา ตตฺถ ปญฺญา ปชานนา ฯเปฯ อโมโห ธมฺมวิจโย สมฺมาทิฏฺฐิ, อิทํ วุจฺจติ อายโกสลฺลํ.}

\prose{ตตฺถ กตมํ อปายโกสลฺลํ, อิเม ธมฺเม มนสิกโรโต อนุปฺปนฺนา เจว กุสลา ธมฺมา น อุปฺปชฺชนฺติ, อุปฺปนฺนา จ กุสลา ธมฺมา นิรุชฺฌนฺติ, อิเม วา ปนิเม ธมฺเม มนสิกโรโต อนุปฺปนฺนา เจว อกุสลา ธมฺมา อุปฺปชฺชนฺติ, อุปฺปนฺนา จ อกุสลา ธมฺมา ภิยฺโยภาวาย เวปุลฺลาย สํวตฺตนฺตีติ, ยา ตตฺถ ปญฺญา ปชานนา ฯเปฯ อโมโห ธมฺมวิจโย สมฺมาทิฏฺฐิ, อิทํ วุจฺจติ อปายโกสลฺลํ.\csromanspacer{} สพฺพาปิ ตตฺรุปายา ปญฺญา อุปายโกสลฺลํ.}

\proseitem{772}{จตูสุ ภูมีสุ วิปาเก ปญฺญา วิปากา ปญฺญา, จตูสุ ภูมีสุ กุสเล ปญฺญา วิปากธมฺมธมฺมา ปญฺญา, ตีสุ ภูมีสุ กิริยาพฺยากเต ปญฺญา เนววิปากนวิปากธมฺมธมฺมา ปญฺญา.}

\proseitem{773}{ตีสุ ภูมีสุ วิปาเก ปญฺญา อุปาทินฺนุปาทานิยา ปญฺญา, ตีสุ ภูมีสุ กุสเล ตีสุ ภูมีสุ กิริยาพฺยากเต ปญฺญา อนุปาทินฺนุปาทานิยา ปญฺญา, จตูสุ มคฺเคสุ จตูสุ ผเลสุ ปญฺญา อนุปาทินฺน-อนุปาทานิยา ปญฺญา.}

\proseitem{774}{วิตกฺกวิจารสมฺปยุตฺตา ปญฺญา สวิตกฺกสวิจารา ปญฺญา, วิตกฺกวิปฺปยุตฺตา วิจารสมฺปยุตฺตา ปญฺญา อวิตกฺกวิจารมตฺตา ปญฺญา, วิตกฺกวิจารวิปฺปยุตฺตา ปญฺญา อวิตกฺก-อวิจารา ปญฺญา.}

\proseitem{775}{ปีติสมฺปยุตฺตา ปญฺญา ปีติสหคตา ปญฺญา, สุขสมฺปยุตฺตา ปญฺญา สุขสหคตา ปญฺญา, อุเปกฺขาสมฺปยุตฺตา ปญฺญา อุเปกฺขาสหคตา ปญฺญา.}

\csromanheader{2}{16. ญาณวิภงฺค}
\proseitem{776}{\csromanfolio{339}ตีสุ ภูมีสุ กุสเล ปญฺญา อาจยคามินี ปญฺญา, จตูสุ มคฺเคสุ ปญฺญา อปจยคามินี ปญฺญา, จตูสุ ภูมีสุ วิปาเก ตีสุ ภูมีสุ กิริยาพฺยากเต ปญฺญา เนวาจยคามินาปจยคามินี ปญฺญา.}

\proseitem{777}{จตูสุ มคฺเคสุ ตีสุ ผเลสุ ปญฺญา เสกฺขา ปญฺญา, อุปริฏฺฐิมา\footnote{อุปริฏฺฐิเม (สฺยา), อุปริฏฺฐิมํ (ก)} อรหตฺตผเล ปญฺญา อเสกฺขา ปญฺญา, ตีสุ ภูมีสุ กุสเล ตีสุ ภูมีสุ วิปาเก ตีสุ ภูมีสุ กิริยาพฺยากเต ปญฺญา เนวเสกฺขนาเสกฺขา ปญฺญา.}

\proseitem{778}{กามาวจรกุสลาพฺยากเต ปญฺญา ปริตฺตา ปญฺญา, รูปาวจรารูปาวจรกุสลาพฺยากเต ปญฺญา มหคฺคตา ปญฺญา, จตูสุ มคฺเคสุ จตูสุ ผเลสุ ปญฺญา อปฺปมาณา ปญฺญา.}

\proseitem{779}{ตตฺถ กตมา ปริตารมฺมณา ปญฺญา, ปริตฺเต ธมฺเม อารพฺภ ยา อุปฺปชฺชติ ปญฺญา ปชานนา ฯเปฯ อโมโห ธมฺมวิจโย สมฺมาทิฏฺฐิ, อยํ วุจฺจติ ปริตฺตารมฺมณา ปญฺญา.}

\proseitem{780}{ตตฺถ กตมา มหคฺคตารมฺมณา ปญฺญา, มหคฺคเต ธมฺเม อารพฺภ ยา อุปฺปชฺชติ ปญฺญา ปชานนา ฯเปฯ อโมโห ธมฺมวิจโย สมฺมาทิฏฺฐิ, อยํ วุจฺจติ มหคฺคตารมฺมณา ปญฺญา.}

\proseitem{781}{ตตฺถ กตมา อปฺปมาณารมฺมณา ปญฺญา, อปฺปมาเณ ธมฺเม อารพฺภ ยา อุปฺปชฺชติ ปญฺญา ปชานนา ฯเปฯ อโมโห ธมฺมวิจโย สมฺมาทิฏฺฐิ, อยํ วุจฺจติ อปฺปมาณารมฺมณา ปญฺญา.}

\proseitem{782}{ตตฺถ กตมา มคฺคารมฺมณา ปญฺญา, อริยมคฺคํ อารพฺภ ยา อุปฺปชฺชติ ปญฺญา ปชานนา ฯเปฯ อโมโห ธมฺมวิจโย สมฺมาทิฏฺฐิ, อยํ วุจฺจติ มคฺคารมฺมณา ปญฺญา.\csromanspacer{} จตูสุ มคฺเคสุ ปญฺญา มคฺคเหตุกา ปญฺญา.}

\proseitem{783}{ตตฺถ กตมา มคฺคาธิปตินี ปญฺญา, อริยมคฺคํ อธิปติํ กริตฺวา ยา อุปฺปชฺชติ ปญฺญา ปชานนา ฯเปฯ อโมโห ธมฺมวิจโย สมฺมาทิฏฺฐิ, อยํ วุจฺจติ มคฺคาธิปตินี ปญฺญา.}

\proseitem{784}{\csromanfolio{340}จตูสุ ภูมีสุ วิปาเก ปญฺญา, สิยา อุปฺปนฺนา, สิยา อุปฺปาทินี, น วตฺตพฺพา “อนุปฺปนฺนา”ติ.\csromanspacer{} จตูสุ ภูมีสุ กุสเล ตีสุ ภูมีสุ กิริยาพฺยากเต ปญฺญา, สิยา อุปฺปนฺนา, สิยา อนุปฺปนฺนา, น วตฺตพฺพา “อุปฺปาทินี”ติ.}

\proseitem{785}{สพฺพาว ปญฺญา สิยา อตีตา, สิยา อนาคตา, สิยา ปจฺจุปฺปนฺนา.}

\proseitem{786}{ตตฺถ กตมา อตีตารมฺมณา ปญฺญา, อตีเต ธมฺเม อารพฺภ ยา อุปฺปชฺชติ ปญฺญา ปชานนา ฯเปฯ อโมโห ธมฺมวิจโย สมฺมาทิฏฺฐิ, อยํ วุจฺจติ อตีตารมฺมณา ปญฺญา.}

\proseitem{787}{ตตฺถ กตมา อนาคตารมฺมณา ปญฺญา, อนาคเต ธมฺเม อารพฺภ ยา อุปฺปชฺชติ ปญฺญา ปชานนา ฯเปฯ อโมโห ธมฺมวิจโย สมฺมาทิฏฺฐิ, อยํ วุจฺจติ อนาคตารมฺมณา ปญฺญา.}

\proseitem{788}{ตตฺถ กตมา ปจฺจุปฺปนฺนารมฺมณา ปญฺญา, ปจฺจุปฺปนฺเน ธมฺเม อารพฺภ ยา อุปฺปชฺชติ ปญฺญา ปชานนา ฯเปฯ อโมโห ธมฺมวิจโย สมฺมาทิฏฺฐิ, อยํ วุจฺจติ ปจฺจุปฺปนฺนารมฺมณา ปญฺญา.}

\proseitem{789}{สพฺพาว ปญฺญา สิยา อชฺฌตฺตา, สิยา พหิทฺธา, สิยา อชฺฌตฺตพหิทฺธา.}

\proseitem{790}{ตตฺถ กตมา อชฺฌตฺตารมฺมณา ปญฺญา, อชฺฌตฺเต ธมฺเม อารพฺภ ยา อุปฺปชฺชติ ปญฺญา ปชานนา ฯเปฯ อโมโห ธมฺมวิจโย สมฺมาทิฏฺฐิ, อยํ วุจฺจติ อชฺฌตฺตารมฺมณา ปญฺญา.}

\proseitem{791}{ตตฺถ กตมา พหิทฺธารมฺมณา ปญฺญา, พหิทฺธาธมฺเม อารพฺภ ยา อุปฺปชฺชติ ปญฺญา ปชานนา ฯเปฯ อโมโห ธมฺมวิจโย สมฺมาทิฏฺฐิ, อยํ วุจฺจติ พหิทฺธารมฺมณา ปญฺญา.}

\proseitem{792}{ตตฺถ กตมา อชฺฌตฺตพหิทฺธารมฺมณา ปญฺญา, อชฺฌตฺตพหิทฺธา ธมฺเม อารพฺภ ยา อุปฺปชฺชติ ปญฺญา ปชานนา ฯเปฯ อโมโห ธมฺมวิจโย สมฺมาทิฏฺฐิ, อยํ วุจฺจติ อชฺฌตฺตพหิทฺธารมฺมณา ปญฺญา.\csromanspacer{} เอวํ ติวิเธน ญาณวตฺถุ.}

\vspace{\tipitakaparskip}
\csromancenter{ติกํ.}
\csromanheader{2}{16. ญาณวิภงฺค}
\csromanmatikaanchor{mtk.161}
\csromantocmarkref{subparagraph}{4. จตุกฺกนิทฺเทส}{mtk.161}
\bookmark[dest={mtk.161},level=5]{4. จตุกฺกนิทฺเทส}
\csromanheader{6}{4. จตุกฺกนิทฺเทส}
\proseitem{793}{\csromanfolio{341}ตตฺถ กตมํ กมฺมสฺสกตญาณํ, “อตฺถิ ทินฺนํ, อตฺถิ ยิฏฺฐํ, อตฺถิ หุตํ, อตฺถิ สุกตทุกฺกฏานํ กมฺมานํ ผลํ วิปาโก, อตฺถิ อยํ โลโก, อตฺถิ ปโร โลโก, อตฺถิ มาตา, อตฺถิ ปิตา, อตฺถิ สตฺตา โอปปาติกา, อตฺถิ โลเก สมณพฺราหฺมณา สมฺมคฺคตา สมฺมาปฏิปนฺนา, เย อิมญฺจ โลกํ ปรญฺจ โลกํ สยํ อภิญฺญา สจฺฉิกตฺวา ปเวเทนฺตี”ติ ยา เอวรูปา ปญฺญา ปชานนา ฯเปฯ อโมโห ธมฺมวิจโย สมฺมาทิฏฺฐิ, อิทํ วุจฺจติ กมฺมสฺสกตญาณํ.\csromanspacer{} ฐเปตฺวา สจฺจานุโลมิกํ ญาณํ สพฺพาปิ สาสวา กุสลา ปญฺญา กมฺมสฺสกตญาณํ.}

\prose{ตตฺถ กตมํ สจฺจานุโลมิกํ ญาณํ, “รูปํ อนิจฺจนฺ”ติ วา, เวทนา ฯเปฯ.\csromanspacer{} สญฺญา.\csromanspacer{} สงฺขารา. “วิญฺญาณํ อนิจฺจนฺ”ติ วา, ยา เอวรูปี อนุโลมิกา ขนฺติ ทิฏฺฐิ รุจิ มุทิ เปกฺขา ธมฺมนิชฺฌานกฺขนฺติ, อิทํ วุจฺจติ สจฺจานุโลมิกํ ญาณํ.\csromanspacer{} จตูสุ มคฺเคสุ ปญฺญา มคฺคสมงฺคิสฺส ญาณํ.\csromanspacer{} จตูสุ ผเลสุ ปญฺญา ผลสมงฺคิสฺส ญาณํ. (1)}

\proseitem{794}{มคฺคสมงฺคิสฺส ญาณํ ทุกฺเขเปตํ ญาณํ, ทุกฺขสมุทเยเปตํ ญาณํ, ทุกฺขนิโรเธเปตํ ญาณํ, ทุกฺขนิโรธคามินิยา ปฏิปทายเปตํ ญาณํ.}

\prose{ตตฺถ กตมํ ทุกฺเข ญาณํ, ทุกฺขํ อารพฺภ ยา อุปฺปชฺชติ ปญฺญา ปชานนา ฯเปฯ อโมโห ธมฺมวิจโย สมฺมาทิฏฺฐิ, อิทํ วุจฺจติ ทุกฺเข ญาณํ.\csromanspacer{} ทุกฺขสมุทยํ อารพฺภ ฯเปฯ ทุกฺขนิโรธํ อารพฺภ ฯเปฯ ทุกฺขนิโรธคามินิํ ปฏิปทํ อารพฺภ ยา อุปฺปชฺชติ ปญฺญา ปชานนา ฯเปฯ อโมโห ธมฺมวิจโย, สมฺมาทิฏฺฐิ, อิทํ วุจฺจติ ทุกฺขนิโรธคามินิยา ปฏิปทาย ญาณํ. (2)}

\proseitem{795}{กามาวจรกุสลาพฺยากเต ปญฺญา กามาวจรา ปญฺญา, รูปาวจรกุสลาพฺยากเต ปญฺญา รูปาวจรา ปญฺญา, อรูปาวจรกุสลาพฺยากเต ปญฺญา อรูปาวจรา ปญฺญา, จตูสุ มคฺเคสุ จตูสุ ผเลสุ ปญฺญา อปริยาปนฺนา ปญฺญา. (3)}

\proseitem{796}{ตตฺถ กตมํ ธมฺเม ญาณํ, จตูสุ มคฺเคสุ จตูสุ ผเลสุ ปญฺญา ธมฺเม ญาณํ, โส อิมินา ธมฺเมน ญาเตน ทิฏฺเฐน ปตฺเตน วิทิเตน ปริโยคาเฬฺหน อตีตานาคเตน นยํ เนติ, เยหิ เกจิ อตีตมทฺธานํ สมณา \csromanfolio{342}วา พฺราหฺมณา วา ทุกฺขํ อพฺภญฺญํสุ\footnote{อพฺภญฺญิํสุ (สฺยา) เอวมุปริปิ.}, ทุกฺขสมุทยํ อพฺภญฺญํสุ, ทุกฺขนิโรธํ อพฺภญฺญํสุ, ทุกฺขนิโรธคามินิํ ปฏิปทํ อพฺภญฺญํสุ, อิมญฺเญว เต ทุกฺขํ อพฺภญฺญํสุ, อิมญฺเญว เต ทุกฺขสมุทยํ อพฺภญฺญํสุ, อิมญฺเญว เต ทุกฺขนิโรธํ อพฺภญฺญํสุ, อิมญฺเญว เต ทุกฺขนิโรธคามินิํ ปฏิปทํ อพฺภญฺญํสุ.\csromanspacer{} เยหิ เกจิ อนาคตมทฺธานํ สมณา วา พฺราหฺมณา วา ทุกฺขํ อภิชานิสฺสนฺติ, ทุกฺขสมุทยํ อภิชานิสฺสนฺติ, ทุกฺขนิโรธํ อภิชานิสฺสนฺติ, ทุกฺขนิโรธคามินิํ ปฏิปทํ อภิชานิสฺสนฺติ, อิมญฺเญว เต ทุกฺขํ อภิชานิสฺสนฺติ, อิมญฺเญว เต ทุกฺขสมุทยํ อภิชานิสฺสนฺติ, อิมญฺเญว เต ทุกฺขนิโรธํ อภิชานิสฺสนฺติ, อิมญฺเญว เต ทุกฺขนิโรธคามินิํ ปฏิปทํ อภิชานิสฺสนฺตีติ ยา ตตฺถ ปญฺญา ปชานนา ฯเปฯ อโมโห ธมฺมวิจโย สมฺมาทิฏฺฐิ, อิทํ วุจฺจติ อนฺวเย ญาณํ.}

\prose{ตตฺถ กตมํ ปริเย ญาณํ, อิธ ภิกฺขุ ปรสตฺตานํ ปรปุคฺคลานํ เจตสา เจโต ปริจฺจ ปชานาติ, สราคํ วา จิตฺตํ “สราคํ จิตฺตนฺ”ติ ปชานาติ, วีตราคํ วา จิตฺตํ “วีตราคํ จิตฺตนฺ”ติ ปชานาติ, สโทสํ วา จิตฺตํ “สโทสํ จิตฺตนฺ”ติ ปชานาติ, วีตโทสํ วา จิตฺตํ “วีตโทสํ จิตฺตนฺ”ติ ปชานาติ, สโมหํ วา จิตฺตํ “สโมหํ จิตฺตนฺ”ติ ปชานาติ, วีตโมหํ วา จิตฺตํ “วีตโมหํ จิตฺตนฺ”ติ ปชานาติ, สํขิตฺตํ วา จิตฺตํ “สํขิตฺตํ จิตฺตนฺ”ติ ปชานาติ, วิกฺขิตฺตํ วา จิตฺตํ “วิกฺขิตฺตํ จิตฺตนฺ”ติ ปชานาติ, มหคฺคตํ วา จิตฺตํ “มหคฺคตํ จิตฺตนฺ”ติ ปชานาติ, อมหคฺคตํ วา จิตฺตํ “อมหคฺคตํ จิตฺตนฺ”ติ ปชานาติ, ส-อุตฺตรํ วา จิตฺตํ “ส-อุตฺตรํ จิตฺตนฺ”ติ ปชานาติ, อนุตฺตรํ วา จิตฺตํ “อนุตฺตรํ จิตฺตนฺ”ติ ปชานาติ, สมาหิตํ วา จิตฺตํ “สมาหิตํ จิตฺตนฺ”ติ ปชานาติ, อสมาหิตํ วา จิตฺตํ “อสมาหิตํ จิตฺตนฺ”ติ ปชานาติ, วิมุตฺตํ วา จิตฺตํ “วิมุตฺตํ จิตฺตนฺ”ติ ปชานาติ, อวิมุตฺตํ วา จิตฺตํ “อวิมุตฺตํ จิตฺตนฺ”ติ ปชานาตีติ ยา ตตฺถ ปญฺญา ปชานนา ฯเปฯ อโมโห ธมฺมวิจโย สมฺมาทิฏฺฐิ, อิทํ วุจฺจติ ปริเย ญาณํ.}

\prose{ฐเปตฺวา ธมฺเม ญาณํ อนฺวเย ญาณํ ปริเย ญาณํ อวเสสา ปญฺญา สมฺมุติญาณํ. (4)}

\proseitem{797}{ตตฺถ กตมา ปญฺญา อาจยาย โน อปจยาย, กามาวจรกุสเล ปญฺญา อาจยาย โน อปจยาย, จตูสุ มคฺเคสุ ปญฺญา}

\csromanheader{2}{16. ญาณวิภงฺค}
\prose{\csromanfolio{343}อปจยาย โน อาจยาย, รูปาวจรารูปาวจรกุสเล ปญฺญา อาจยาย เจว อปจยาย จ, อวเสสา ปญฺญา เนว อาจยาย โน อปจยาย. (5)}

\proseitem{798}{ตตฺถ กตมา ปญฺญา นิพฺพิทาย โน ปฏิเวธาย, ยาย ปญฺญาย กาเมสุ วีตราโค โหติ น จ อภิญฺญาโย ปฏิวิชฺฌติ น จ สจฺจานิ, อยํ วุจฺจติ ปญฺญา นิพฺพิทาย โน ปฏิเวธาย.\csromanspacer{} เสฺวว ปญฺญาย กาเมสุ วีตราโค สมาโน อภิญฺญาโย ปฏิวิชฺฌติ น จ สจฺจานิ, อยํ วุจฺจติ ปญฺญา ปฏิเวธาย โน นิพฺพิทาย.\csromanspacer{} จตูสุ มคฺเคสุ ปญฺญา นิพฺพิทาย เจว ปฏิเวธาย จ, อวเสสา ปญฺญา เนว นิพฺพิทาย โน ปฏิเวธาย. (6)}

\proseitem{799}{ตตฺถ กตมา หานภาคินี ปญฺญา, ปฐมสฺส ฌานสฺส ลาภิํ กามสหคตา สญฺญามนสิการา สมุทาจรนฺติ หานภาคินี ปญฺญา, ตทนุธมฺมตา สติ สนฺติฏฺฐติ ฐิติภาคินี ปญฺญา.\csromanspacer{} อวิตกฺกสหคตา สญฺญามนสิการา สมุทาจรนฺติ วิเสสภาคินี ปญฺญา, นิพฺพิทาสหคตา สญฺญามนสิการา สมุทาจรนฺติ วิราคูปสญฺหิตา นิพฺเพธภาคินี ปญฺญา.\csromanspacer{} ทุติยสฺส ฌานสฺส ลาภิํวิตกฺกสหคตา สญฺญามนสิการา สมุทาจรนฺติ หานภาคินี ปญฺญา, ตทนุธมฺมตา สติ สนฺติฏฺฐติ ฐิติภาคินี ปญฺญา, อุเปกฺขาสหคตา สญฺญามนสิการา สมุทาจรนฺติ วิเสสภาคินี ปญฺญา, นิพฺพิทาสหคตา สญฺญามนสิการา สมุทาจรนฺติ วิราคูปสญฺหิตา นิพฺเพธภาคินี ปญฺญา.\csromanspacer{} ตติยสฺส ฌานสฺส ลาภิํ ปีติสุขสหคตา สญฺญามนสิการา สมุทาจรนฺติ หานภาคินี ปญฺญา, ตทนุธมฺมตา สติ สนฺติฏฺฐติ ฐิติภาคินี ปญฺญา.\csromanspacer{} อทุกฺขมสุขสหคตา สญฺญามนสิการา สมุทาจรนฺติ วิเสสภาคินี ปญฺญา, นิพฺพิทาสหคตา สญฺญามนสิการา สมุทาจรนฺติ วิราคูปสญฺหิตา นิพฺเพธภาคินี ปญฺญา.\csromanspacer{} จตุตฺถสฺส ฌานสฺส ลาภิํ อุเปกฺขาสหคตา สญฺญามนสิการา สมุทาจรนฺติ หานภาคินี ปญฺญา, ตทนุธมฺมตา สติ สนฺติฏฺฐติ ฐิติภาคินี ปญฺญา.\csromanspacer{} อากาสานญฺจายตนสหคตา สญฺญามนสิการา สมุทาจรนฺติ วิเสสภาคินี ปญฺญา, นิพฺพิทาสหคตา สญฺญามนสิการา สมุทาจรนฺติ วิราคูปสญฺหิตา นิพฺเพธภาคินี ปญฺญา.\csromanspacer{} อากาสานญฺจายตนสฺส ลาภิํ รูปสหคตา สญฺญามนสิการา สมุทาจรนฺติ หานภาคินี ปญฺญา, ตทนุธมฺมตา สติ สนฺติฏฺฐติ \csromanfolio{344}ฐิติภาคินี ปญฺญา.\csromanspacer{} วิญฺญาณญฺจายตนสหคตา สญฺญามนสิการา สมุทาจรนฺติ วิเสสภาคินี ปญฺญา, นิพฺพิทาสหคตา สญฺญามนสิการา สมุทาจรนฺติ วิราคูปสญฺหิตา นิพฺเพธภาคินี ปญฺญา.\csromanspacer{} วิญฺญาณญฺจายตนสฺส ลาภิํ อากาสานญฺจายตนสหคตา สญฺญามนสิการา สมุทาจรนฺติ หานภาคินี ปญฺญา, ตทนุธมฺมตา สติ สนฺติฏฺฐติ ฐิติภาคินี ปญฺญา.\csromanspacer{} อากิญฺจญฺญายตนสหคตา สญฺญามนสิการา สมุทาจรนฺติ วิเสสภาคินี ปญฺญา, นิพฺพิทาสหคตา สญฺญามนสิการา สมุทาจรนฺติ วิราคูปสญฺหิตา นิพฺเพธภาคินี ปญฺญา.\csromanspacer{} อากิญฺจญฺญายตนสฺส ลาภิํ วิญฺญาณญฺจายตนสหคตา สญฺญามนสิการา สมุทาจรนฺติ หานภาคินี ปญฺญา, ตทนุธมฺมตา สติ สนฺติฏฺฐติ ฐิติภาคินี ปญฺญา, เนวสญฺญานาสญฺญายตนสหคตา สญฺญามนสิการา สมุทาจรนฺติ วิเสสภาคินี ปญฺญา, นิพฺพิทาสหคตา สญฺญามนสิการา สมุทาจรนฺติ วิราคูปสญฺหิตา นิพฺเพธภาคินี ปญฺญา. (7)}

\proseitem{800}{ตตฺถ กตมา จตสฺโส ปฏิสมฺภิทา, อตฺถปฏิสมฺภิทา ธมฺมปฏิสมฺภิทา นิรุตฺติปฏิสมฺภิทา ปฏิภานปฏิสมฺภิทา.\csromanspacer{} อตฺเถ ญาณํ อตฺถปฏิสมฺภิทา, ธมฺเม ญาณํ ธมฺมปฏิสมฺภิทา, ตตฺร ธมฺมนิรุตฺตาภิลาเป ญาณํ นิรุตฺติปฏิสมฺภิทา, ญาเณสุ ญาณํ ปฏิภานปฏิสมฺภิทา, อิมา จตสฺโส ปฏิสมฺภิทา. (8)}

\proseitem{801}{ตตฺถ กตมา จตสฺโส ปฏิปทา, ทุกฺขปฏิปทา ทนฺธาภิญฺญา ปญฺญา, ทุกฺขปฏิปทา ขิปฺปาภิญฺญา ปญฺญา, สุขปฏิปทา ทนฺธาภิญฺญา ปญฺญา, สุขปฏิปทา ขิปฺปาภิญฺญา ปญฺญา.}

\prose{ตตฺถ กตมา ทุกฺขปฏิปทา ทนฺธาภิญฺญา ปญฺญา, กิจฺเฉน กสิเรน สมาธิํ อุปฺปาเทนฺตสฺส ทนฺธํ ตณฺฐานํ อภิชานนฺตสฺส ยา อุปฺปชฺชติ ปญฺญา ปชานนา ฯเปฯ อโมโห ธมฺมวิจโย สมฺมาทิฏฺฐิ, อยํ วุจฺจติ ทุกฺขปฏิปทา ทนฺธาภิญฺญา ปญฺญา.}

\prose{ตตฺถ กตมา ทุกฺขปฏิปทา ขิปฺปาภิญฺญา ปญฺญา, กิจฺเฉน กสิเรน สมาธิํ อุปฺปาเทนฺตสฺส ขิปฺปํ ตณฺฐานํ อภิชานนฺตสฺส ยา อุปฺปชฺชติ ปญฺญา ปชานนา ฯเปฯ อโมโห ธมฺมวิจโย สมฺมาทิฏฺฐิ, อยํ วุจฺจติ ทุกฺขปฏิปทา ขิปฺปาภิญฺญา ปญฺญา.}

\csromanheader{2}{16. ญาณวิภงฺค}
\prose{\csromanfolio{345}ตตฺถ กตมา สุขปฏิปทา ทนฺธาภิญฺญา ปญฺญา, อกิจฺเฉน อกสิเรน สมาธิํ อุปฺปาเทนฺตสฺส ทนฺธํ ตณฺฐานํ อภิชานนฺตสฺส ยา อุปฺปชฺชติ ปญฺญา ปชานนา ฯเปฯ อโมโห ธมฺมวิจโย สมฺมาทิฏฺฐิ, อยํ วุจฺจติ สุขปฏิปทา ทนฺธาภิญฺญา ปญฺญา.}

\prose{ตตฺถ กตมา สุขปฏิปทา ขิปฺปาภิญฺญา ปญฺญา, อกิจฺเฉน อกสิเรน สมาธิํ อุปฺปาเทนฺตสฺส ขิปฺปํ ตณฺฐานํ อภิชานนฺตสฺส ยา อุปฺปชฺชติ ปญฺญา ปชานนา ฯเปฯ อโมโห ธมฺมวิจโย สมฺมาทิฏฺฐิ, อยํ วุจฺจติ สุขปฏิปทา ขิปฺปาภิญฺญา ปญฺญา.\csromanspacer{} อิมา จตสฺโส ปฏิปทา. (9)}

\proseitem{802}{ตตฺถ กตมานิ จตฺตาริ อารมฺมณานิ, ปริตฺตา ปริตฺตารมฺมณา ปญฺญา, ปริตฺตา อปฺปมาณารมฺมณา ปญฺญา, อปฺปมาณา ปริตฺตารมฺมณา ปญฺญา, อปฺปมาณา อปฺปมาณารมฺมณา ปญฺญา.}

\prose{ตตฺถ กตมา ปริตฺตา ปริตฺตารมฺมณา ปญฺญา, สมาธิสฺส น นิกามลาภิสฺส อารมฺมณํ โถกํ ผรนฺตสฺส ยา อุปฺปชฺชติ ปญฺญา ปชานนา ฯเปฯ อโมโห ธมฺมวิจโย สมฺมาทิฏฺฐิ, อยํ วุจฺจติ ปริตฺตา ปริตฺตารมฺมณา ปญฺญา.}

\prose{ตตฺถ กตมา ปริตฺตา อปฺปมาณารมฺมณา ปญฺญา, สมาธิสฺส น นิกามลาภิสฺส อารมฺมณํ วิปุลํ ผรนฺตสฺส ยา อุปฺปชฺชติ ปญฺญา ปชานนา ฯเปฯ อโมโห ธมฺมวิจโย สมฺมาทิฏฺฐิ, อยํ วุจฺจติ ปริตฺตา อปฺปมาณารมฺมณา ปญฺญา.}

\prose{ตตฺถ กตมา อปฺปมาณา ปริตฺตารมฺมณา ปญฺญา, สมาธิสฺส นิกามลาภิสฺส อารมฺมณํ โถกํ ผรนฺตสฺส ยา อุปฺปชฺชติ ปญฺญา ปชานนา ฯเปฯ อโมโห ธมฺมวิจโย สมฺมาทิฏฺฐิ, อยํ วุจฺจติ อปฺปมาณา ปริตฺตารมฺมณา ปญฺญา.}

\prose{ตตฺถ กตมา อปฺปมาณา อปฺปมาณารมฺมณา ปญฺญา, สมาธิสฺส นิกามลาภิสฺส อารมฺมณํ วิปุลํ ผรนฺตสฺส ยา อุปฺปชฺชติ ปญฺญา ปชานนา ฯเปฯ อโมโห ธมฺมวิจโย สมฺมาทิฏฺฐิ, อยํ วุจฺจติ อปฺปมาณา อปฺปมาณารมฺมณา ปญฺญา.\csromanspacer{} อิมานิ จตฺตาริ อารมฺมณานิ. (10)}

\prose{มคฺคสมงฺคิสฺส ญาณํ ชรามรเณเปตํ ญาณํ, ชรามรณสมุทเยเปตํ ญาณํ, ชรามรณนิโรเธเปตํ ญาณํ, ชรามรณนิโรธคามินิยา ปฏิปทายเปตํ ญาณํ.}

\prose{\csromanfolio{346}ตตฺถ กตมํ ชรามรเณ ญาณํ, ชรามรณํ อารพฺภ ยา อุปฺปชฺชติ ปญฺญา ปชานนา ฯเปฯ อโมโห ธมฺมวิจโย สมฺมาทิฏฺฐิ, อิทํ วุจฺจติ ชรามรเณ ญาณํ.}

\prose{ชรามรณสมุทยํ อารพฺภ ฯเปฯ.\csromanspacer{} ชรามรณนิโรธํ อารพฺภ ฯเปฯ.\csromanspacer{} ชรามรณนิโรธคามินิํ ปฏิปทํ อารพฺภ ยา อุปฺปชฺชติ ปญฺญา ปชานนา ฯเปฯ อโมโห ธมฺมวิจโย สมฺมาทิฏฺฐิ, อิทํ วุจฺจติ ชรามรณนิโรธคามินิยา ปฏิปทาย ญาณํ. (11)}

\proseitem{803}{ธมฺมสมงฺคิสฺส ญาณํ ชาติยาเปตํ ญาณํ ฯเปฯ ภเวเปตํ ญาณํ ฯเปฯ อุปาทาเนเปตํ ญาณํ ฯเปฯ ตณฺหายเปตํ ญาณํ ฯเปฯ เวทนายเปตํ ญาณํ ฯเปฯ ผสฺเสเปตํ ญาณํ ฯเปฯ สฬายตเนเปตํ ญาณํ ฯเปฯ นามรูเปเปตํ ญาณํ ฯเปฯ วิญฺญาเณเปตํ ญาณํ ฯเปฯ สงฺขาเรสุเปตํ ญาณํ, สงฺขารสมุทเยเปตํ ญาณํ, สงฺขารนิโรเธเปตํ ญาณํ, สงฺขารนิโรธคามินิยา ปฏิปทายเปตํ ญาณํ. (10-21)}

\prose{ตตฺถ กตมํ สงฺขาเรสุ ญาณํ, สงฺขาเร อารพฺภ ยา อุปฺปชฺชติ ปญฺญา ปชานนา ฯเปฯ อโมโห ธมฺมวิจโย สมฺมาทิฏฺฐิ, อิทํ วุจฺจติ สงฺขาเรสุ ญาณํ.}

\prose{สงฺขารสมุทยํ อารพฺภ ฯเปฯ สงฺขารนิโรธํ อารพฺภ ฯเปฯ สงฺขารนิโรธคามินิํ ปฏิปทํ อารพฺภ ยา อุปฺปชฺชติ ปญฺญา ปชานนา ฯเปฯ อโมโห ธมฺมวิจโย สมฺมาทิฏฺฐิ, อิทํ วุจฺจติ สงฺขารนิโรธคามินิยา ปฏิปทาย ญาณํ.\csromanspacer{} เอวํ จตุพฺพิเธน ญาณวตฺถุ.}

\csromanheader{2}{จตุกฺกํ.}
\csromansectionrule
\csromanmatikaanchor{mtk.162}
\csromantocmarkref{subparagraph}{5. ปญฺจกนิทฺเทส}{mtk.162}
\bookmark[dest={mtk.162},level=5]{5. ปญฺจกนิทฺเทส}
\csromanheader{6}{5. ปญฺจกนิทฺเทส}
\proseitem{804}{ตตฺถ กตโม ปญฺจงฺคิโก สมฺมาสมาธิ, ปีติผรณตา, สุขผรณตา, เจโตผรณตา, อาโลกผรณตา, ปจฺจเวกฺขณา นิมิตฺตํ.\csromanspacer{} ทฺวีสุ ฌาเนสุ ปญฺญา ปีติผรณตา, ตีสุ ฌาเนสุ ปญฺญา สุขผรณตา, ปรจิตฺเต ญาณํ เจโตผรณตา, ทิพฺพจกฺขุ อาโลกผรณตา, ตมฺหา ตมฺหา สมาธิมฺหา วุฏฺฐิตสฺส ปจฺจเวกฺขณาญาณํ ปจฺจเวกฺขณานิมิตฺตํ, อยํ วุจฺจติ ปญฺจงฺคิโก สมฺมาสมาธิ.}

\csromanheader{2}{16. ญาณวิภงฺค}
\prose{\csromanfolio{347}ตตฺถ กตโม ปญฺจญาณิโก สมฺมาสมาธิ, “อยํ สมาธิ ปจฺจุปฺปนฺนสุโข เจว อายติญฺจ สุขวิปาโก”ติ ปจฺจตฺตญฺเญว ญาณํ อุปฺปชฺชติ, “อยํ สมาธิ อริโย นิรามิโส”ติ ปจฺจตฺตญฺเญว ญาณํ อุปฺปชฺชติ, “อยํ สมาธิ อกาปุริสเสวิโต”ติ ปจฺจตฺตญฺเญว ญาณํ อุปฺปชฺชติ, “อยํ สมาธิ สนฺโต ปณีโต ปฏิปฺปสฺสทฺธลทฺโธ เอโกทิภาวาธิคโต น สสงฺขารนิคฺคยฺหวาริตคโต”ติ ปจฺจตฺตญฺเญว ญาณํ อุปฺปชฺชติ, “โส โข ปนาหํ อิมํ สมาธิํ สโต สมาปชฺชามิ สโต วุฏฺฐหามี”ติ ปจฺจตฺตญฺเญว ญาณํ อุปฺปชฺชติ, อยํ ปญฺจญาณิโก สมฺมาสมาธิ.\csromanspacer{} เอวํ ปญฺจวิเธน ญาณวตฺถุ.}

\csromanheader{2}{ปญฺจกํ.}
\csromansectionrule
\csromanmatikaanchor{mtk.163}
\csromantocmarkref{subparagraph}{6. ฉกฺกนิทฺเทส}{mtk.163}
\bookmark[dest={mtk.163},level=5]{6. ฉกฺกนิทฺเทส}
\csromanheader{6}{6. ฉกฺกนิทฺเทส}
\proseitem{805}{ตตฺถ กตมา ฉสุ อภิญฺญาสุ ปญฺญา, อิทฺธิวิเธ ญาณํ, โสตธาตุวิสุทฺธิยา ญาณํ, ปรจิตฺเต ญาณํ, ปุพฺเพนิวาสานุสฺสติยา ญาณํ, สตฺตานํ จุตูปปาเต ญาณํ, อาสวานํ ขเย ญาณํ, อิมา ฉสุ อภิญฺญาสุ ปญฺญา.\csromanspacer{} เอวํ ฉพฺพิเธน ญาณวตฺถุ.}

\csromanmatikaanchor{mtk.164}
\csromantocmarkref{subparagraph}{7. สตฺตกนิทฺเทส}{mtk.164}
\bookmark[dest={mtk.164},level=5]{7. สตฺตกนิทฺเทส}
\csromanheader{6}{7. สตฺตกนิทฺเทส}
\proseitem{806}{ตตฺถ กตมานิ สตฺตสตฺตติ ญาณวตฺถูนิ, “ชาติปจฺจยา ชรามรณนฺ”ติ ญาณํ, อสติ ชาติยา “นตฺถิ ชรามรณนฺ”ติ ญาณํ, อตีตมฺปิ อทฺธานํ “ชาติปจฺจยา ชรามรณนฺ”ติ ญาณํ, อสติ ชาติยา “นตฺถิ ชรา มรณนฺ”ติ ญาณํ, อนาคตมฺปิ “อทฺธานํ ชาติปจฺจยา ชรามรณนฺ”ติ ญาณํ, อสติ ชาติยา “นตฺถิ ชรามรณนฺ”ติ ญาณํ, ยมฺปิสฺส ตํ ธมฺมฏฺฐิติญาณํ, ตมฺปิ “ขยธมฺมํ วยธมฺมํ วิราคธมฺมํ นิโรธธมฺมนฺ”ติ ญาณํ, “ภวปจฺจยา ชาตี”ติ ญาณํ ฯเปฯ “อุปาทานปจฺจยา ภโว”ติ ญาณํ ฯเปฯ “ตณฺหาปจฺจยา อุปาทานนฺ”ติ ญาณํ ฯเปฯ “เวทนาปจฺจยา ตณฺหา”ติ ญาณํ ฯเปฯ “ผสฺสปจฺจยา เวทนา”ติ ญาณํ ฯเปฯ “สฬายตนปจฺจยา ผสฺโส”ติ ญาณํ ฯเปฯ “นามรูปปจฺจยา สฬายตนนฺ”ติ ญาณํ ฯเปฯ “วิญฺญาณปจฺจยา นามรูปนฺ”ติ ญาณํ ฯเปฯ “สงฺขารปจฺจยา วิญฺญาณนฺ”ติ ญาณํ ฯเปฯ “อวิชฺชาปจฺจยา สงฺขารา”ติ ญาณํ.\csromanspacer{} อสติ อวิชฺชาย “นตฺถิ สงฺขารา”ติ ญาณํ, “อตีตมฺปิ อทฺธานํ อวิชฺชาปจฺจยา สงฺขารา”ติ ญาณํ, อสติ อวิชฺชาย “นตฺถิ สงฺขารา”ติ ญาณํ, อนาคตมฺปิ \csromanfolio{348}อทฺธานํ “อวิชฺชาปจฺจยา สงฺขารา”ติ ญาณํ, อสติ อวิชฺชาย “นตฺถิ สงฺขารา”ติ ญาณํ, ยมฺปิสฺส ตํ ธมฺมฏฺฐิติญาณํ, ตมฺปิ “ขยธมฺมํ วยธมฺมํ วิราคธมฺมํ นิโรธธมฺมนฺ”ติ ญาณํ, อิมานิ สตฺตสตฺตติ ญาณวตฺถูนิ.\csromanspacer{} เอวํ สตฺตวิเธน ญาณวตฺถุ.}

\csromanmatikaanchor{mtk.165}
\csromantocmarkref{subparagraph}{8. อฏฺฐกนิทฺเทส}{mtk.165}
\bookmark[dest={mtk.165},level=5]{8. อฏฺฐกนิทฺเทส}
\csromanheader{6}{8. อฏฺฐกนิทฺเทส}
\proseitem{807}{ตตฺถ กตมา จตูสุ มคฺเคสุ จตูสุ ผเลสุ ปญฺญา, โสตาปตฺติมคฺเค ปญฺญา, โสตาปตฺติผเล ปญฺญา, สกทาคามิมคฺเค ปญฺญา, สกทาคามิผเล ปญฺญา, อนาคามิมคฺเค ปญฺญา, อนาคามิผเล ปญฺญา, อรหตฺตมคฺเค ปญฺญา, อรหตฺตผเล ปญฺญา, อิมา จตูสุ มคฺเคสุ จตูสุ ผเลสุ ปญฺญา.\csromanspacer{} เอวํ อฏฺฐวิเธน ญาณวตฺถุ.}

\csromanmatikaanchor{mtk.166}
\csromantocmarkref{subparagraph}{9. นวกนิทฺเทส}{mtk.166}
\bookmark[dest={mtk.166},level=5]{9. นวกนิทฺเทส}
\csromanheader{6}{9. นวกนิทฺเทส}
\proseitem{808}{ตตฺถ กตมา นวสุ อนุปุพฺพวิหารสมาปตฺตีสุ ปญฺญา, ปฐมชฺฌานสมาปตฺติยา ปญฺญา, ทุติยชฺฌานสมาปตฺติยา ปญฺญา, ตติยชฺฌานสมาปตฺติยา ปญฺญา, จตุตฺถชฺฌานสมาปตฺติยา ปญฺญา, อากาสานญฺจายตนสมาปตฺติยา ปญฺญา, วิญฺญาณญฺจายตนสมาปตฺติยา ปญฺญา, อากิญฺจญฺญายตนสมาปตฺติยา ปญฺญา, เนวสญฺญานาสญฺญายตนสมาปตฺติยา ปญฺญา, สญฺญาเวทยิตนิโรธสมาปตฺติยา วุฏฺฐิตสฺส ปจฺจเวกฺขณา ญาณํ, อิมา นวสุ อนุปุพฺพวิหารสมาปตฺตีสุ ปญฺญา.\csromanspacer{} เอวํ นววิเธน ญาณวตฺถุ.}

\csromanmatikaanchor{mtk.167}
\csromantocmarkref{subparagraph}{10. ทสกนิทฺเทส}{mtk.167}
\bookmark[dest={mtk.167},level=5]{10. ทสกนิทฺเทส}
\csromanheader{6}{10. ทสกนิทฺเทส}
\proseitem{809}{ตตฺถ กตมํ ตถาคตสฺส ฐานญฺจ ฐานโต อฏฺฐานญฺจ อฏฺฐานโต ยถาภูตํ ญาณํ.\csromanspacer{} อิธ ตถาคโต “อฏฺฐานเมตํ อนวกาโส ยํ ทิฏฺฐิสมฺปนฺโน ปุคฺคโล กญฺจิ สงฺขารํ นิจฺจโต อุปคจฺเฉยฺย เนตํ ฐานํ วิชฺชตี”ติ ปชานาติ, “ฐานญฺจ โข เอตํ วิชฺชติ ยํ ปุถุชฺชโน กญฺจิ สงฺขารํ นิจฺจโต อุปคจฺเฉยฺย ฐานเมตํ วิชฺชตี”ติ ปชานาติ, “อฏฺฐานเมตํ อนวกาโส ยํ ทิฏฺฐิสมฺปนฺโน ปุคฺคโล กญฺจิ สงฺขารํ สุขโต อุปคจฺเฉยฺย เนตํ ฐานํ วิชฺชตี”ติ ปชานาติ, “ฐานญฺจ โข เอตํ วิชฺชติ ยํ ปุถุชฺชโน กญฺจิ สงฺขารํ สุขโต อุปคจฺเฉยฺย “ฐานเมตํ วิชฺชตี”ติ ปชานาติ,}

\csromanheader{2}{16. ญาณวิภงฺค}
\prose{\csromanfolio{349}“อฏฺฐานเมตํ อนวกาโส ยํ ทิฏฺฐิสมฺปนฺโน ปุคฺคโล กญฺจิ ธมฺมํ อตฺถโต อุปคจฺเฉยฺย เนตํ ฐานํ วิชฺชตี”ติ ปชานาติ, “ฐานญฺจ โข เอตํ วิชฺชติ, ยํ ปุถุชฺชโน กญฺจิ ธมฺมํ อตฺถโต อุปคจฺเฉยฺย ฐานเมตํ วิชฺชตี”ติ ปชานาติ, “อฏฺฐานเมตํ อนวกาโส ยํ ทิฏฺฐิสมฺปนฺโน ปุคฺคโล มาตรํ ชีวิตา โวโรเปยฺย เนตํ ฐานํ วิชฺชตี”ติ ปชานาติ, “ฐานญฺจ โข เอตํ วิชฺชติ, ยํ ปุถุชฺชโน มาตรํ ชีวิตา โวโรเปยฺย ฐานเมตํ วิชฺชตี”ติ ปชานาติ, “อฏฺฐานเมตํ อนวกาโส ยํ ทิฏฺฐิสมฺปนฺโน ปุคฺคโล ปิตรํ ชีวิตา โวโรเปยฺย ฯเปฯ อรหนฺตํ ชีวิตา โวโรเปยฺย ฯเปฯ ปทุฏฺเฐน จิตฺเตน ตถาคตสฺส โลหิตํ อุปฺปาเทยฺย ฯเปฯ สํฆํ ภินฺเทยฺย ฯเปฯ อญฺญํ สตฺถารํ อุทฺทิเสยฺย ฯเปฯ อฏฺฐมํ ภวํ นิพฺพตฺเตยฺย เนตํ ฐานํ วิชฺชตี”ติ ปชานาติ, “ฐานญฺจ โข เอตํ วิชฺชติ ยํ ปุถุชฺชโน อฏฺฐมํ ภวํ นิพฺพตฺเตยฺย ฐานเมตํ วิชฺชตี”ติ ปชานาติ.}

\prose{“อฏฺฐานเมตํ อนวกาโส ยํ เอกิสฺสา โลกธาตุยา เทฺว อรหนฺโต สมฺมาสมฺพุทฺธา อปุพฺพํ อจริมํ อุปฺปชฺเชยฺยุํ เนตํ ฐานํ วิชฺชตี”ติ ปชานาติ, “ฐานญฺจ โข เอตํ วิชฺชติ ยํ เอกิสฺสา โลกธาตุยา เอโก อรหํ สมฺมาสมฺพุทฺโธ อุปฺปชฺเชยฺย ฐานเมตํ วิชฺชตี”ติ ปชานาติ, “อฏฺฐานเมตํ อนวกาโส ยํ เอกิสฺสา โลกธาตุยา เทฺว ราชาโน จกฺกวตฺตี\footnote{จกฺกวตฺติ (สี, สฺยา)} อปุพฺพํ อจริมํ อุปฺปชฺเชยฺยุํ เนตํ ฐานํ วิชฺชตี”ติ ปชานาติ, “ฐานญฺจ โข เอตํ วิชฺชติ ยํ เอกิสฺสา โลกธาตุยา เอโก ราชา จกฺกวตฺตี อุปฺปชฺเชยฺย ฐานเมตํ วิชฺชตี”ติ ปชานาติ.}

\prose{“อฏฺฐานเมตํ อนวกาโส ยํ อิตฺถี อรหํ อสฺส สมฺมาสมฺพุทฺโธ เนตํ ฐานํ วิชฺชตี”ติ ปชานาติ, “ฐานญฺจ โข เอตํ วิชฺชติ ยํ ปุริโส อรหํ สมฺมาสมฺพุทฺโธ ฐานเมตํ วิชฺชตี”ติ ปชานาติ, “อฏฺฐานเมตํ อนวกาโส ยํ อิตฺถี ราชา อสฺส จกฺกวตฺตี เนตํ ฐานํ วิชฺชตี”ติ ปชานาติ, “ฐานญฺจ โข เอตํ วิชฺชติ ยํ ปุริโส ราชา อสฺส จกฺกวตฺตี ฐานเมตํ วิชฺชตี”ติ ปชานาติ, “อฏฺฐานเมตํ อนวกาโส ยํ อิตฺถี สกฺกตฺตํ กเรยฺย, มารตฺตํ กเรยฺย, พฺรหฺมตฺตํ กเรยฺย, เนตํ ฐานํ วิชฺชตี”ติ ปชานาติ, “ฐานญฺจ โข เอตํ วิชฺชติ ยํ ปุริโส สกฺกตฺตํ กเรยฺย, มารตฺตํ กเรยฺย, พฺรหฺมตฺตํ กเรยฺย, ฐานเมตํ วิชฺชตี”ติ ปชานาติ.}

\prose{\csromanfolio{350}“อฏฺฐานเมตํ อนวกาโส ยํ กายทุจฺจริตสฺส อิฏฺโฐ กนฺโต มนาโป วิปาโก นิพฺพตฺเตยฺย เนตํ ฐานํ วิชฺชตี”ติ ปชานาติ, “ฐานญฺจ โข เอตํ วิชฺชติ ยํ กายทุจฺจริตสฺส อนิฏฺโฐ อกนฺโต อมนาโป วิปาโก นิพฺพตฺเตยฺย ฐานเมตํ วิชฺชตี”ติ ปชานาติ, “อฏฺฐานเมตํ อนวกาโส ยํ วจีทุจฺจริตสฺส ฯเปฯ ยํ มโนทุจฺจริตสฺส อิฏฺโฐ กนฺโต มนาโป วิปาโก นิพฺพตฺเตยฺย เนตํ ฐานํ วิชฺชตี”ติ ปชานาติ, “ฐานญฺจ โข เอตํ วิชฺชติ ยํ วจีทุจฺจริตสฺส ฯเปฯ ยํ มโนทุจฺจริตสฺส อนิฏฺโฐ อกนฺโต อมนาโป วิปาโก นิพฺพตฺเตยฺย ฐานเมตํ วิชฺชตี”ติ ปชานาติ.}

\prose{“อฏฺฐานเมตํ อนวกาโส ยํ กายสุจริตสฺส อนิฏฺโฐ อกนฺโต อมนาโป วิปาโก นิพฺพตฺเตยฺย เนตํ ฐานํ วิชฺชตี”ติ ปชานาติ, “ฐานญฺจ โข เอตํ วิชฺชติ ยํ กายสุจริตสฺส อิฏฺโฐ กนฺโต มนาโป วิปาโก นิพฺพตฺเตยฺย ฐานเมตํ วิชฺชตี”ติ ปชานาติ, “อฏฺฐานเมตํ อนวกาโส ยํ วจีสุจริตสฺส ฯเปฯ ยํ มโนสุจริตสฺส อนิฏฺโฐ อกนฺโต อมนาโป วิปาโก นิพฺพตฺเตยฺย เนตํ ฐานํ วิชฺชตี”ติ ปชานาติ, “ฐานญฺจ โข เอตํ วิชฺชติ ยํ วจีสุจริตสฺส ฯเปฯ ยํ มโนสุจริตสฺส อิฏฺโฐ กนฺโต มนาโป วิปาโก นิพฺพตฺเตยฺย ฐานเมตํ วิชฺชตี”ติ ปชานาติ.}

\prose{“อฏฺฐานเมตํ อนวกาโส ยํ กายทุจฺจริตสมงฺคี ตนฺนิทานํ ตปฺปจฺจยา กายสฺส เภทา ปรํ มรณา สุคติํ สคฺคํ โลกํ อุปปชฺเชยฺย เนตํ ฐานํ วิชฺชตี”ติ ปชานาติ, “ฐานญฺจ โข เอตํ วิชฺชติ ยํ กายทุจฺจริตสมงฺคี ตนฺนิทานํ ตปฺปจฺจยา กายสฺส เภทา ปรํ มรณา อปายํ ทุคฺคติํ วินิปาตํ นิรยํ อุปปชฺเชยฺย ฐานเมตํ วิชฺชตี”ติ ปชานาติ, “อฏฺฐานเมตํ อนวกาโส ยํ วจีทุจฺจริตสมงฺคี ฯเปฯ ยํ มโนทุจฺจริตสมงฺคี ตนฺนิทานํ ตปฺปจฺจยา กายสฺส เภทา ปรํ มรณา สุคติํ สคฺคํ โลกํ อุปปชฺเชยฺย เนตํ ฐานํ วิชฺชตี”ติ ปชานาติ, “ฐานญฺจ โข เอตํ วิชฺชติ ยํ วจีทุจฺจริตสมงฺคี ฯเปฯ ยํ มโนทุจฺจริตสมงฺคี ตนฺนิทานํ ตปฺปจฺจยา กายสฺส เภทา ปรํ มรณา อปายํ ทุคฺคติํ วินิปาตํ นิรยํ อุปปชฺเชยฺย ฐานเมตํ วิชฺชตี”ติ ปชานาติ.}

\prose{“อฏฺฐานเมตํ อนวกาโส ยํ กายสุจริตสมงฺคี ตนฺนิทานํ ตปฺปจฺจยา กายสฺส เภทา ปรํ มรณา อปายํ ทุคฺคติํ วินิปาตํ นิรยํ อุปปชฺเชยฺย เนตํ ฐานํ วิชฺชตี”ติ ปชานาติ, “ฐานญฺจ โข เอตํ วิชฺชติ ยํ กายสุจริตสมงฺคี ตนฺนิทานํ ตปฺปจฺจยา กายสฺส เภทา ปรํ มรณา}

\csromanheader{2}{16. ญาณวิภงฺค}
\prose{\csromanfolio{351}สุคติํ สคฺคํ โลกํ อุปปชฺเชยฺย ฐานเมตํ วิชฺชตี”ติ ปชานาติ, “อฏฺฐานเมตํ อนวกาโส ยํ วจีสุจริตสมงฺคี ฯเปฯ ยํ มโนสุจริตสมงฺคี ตนฺนิทานํ ตปฺปจฺจยา กายสฺส เภทา ปรํ มรณา อปายํ ทุคฺคติํ วินิปาตํ นิรยํ อุปปชฺเชยฺย เนตํ ฐานํ วิชฺชตี”ติ ปชานาติ, “ฐานญฺจ โข เอตํ วิชฺชติ ยํ วจีสุจริตสมงฺคี ฯเปฯ ยํ มโนสุจริตสมงฺคี ตนฺนิทานํ ตปฺปจฺจยา กายสฺส เภทา ปรํ มรณา สุคติํ สคฺคํ โลกํ อุปปชฺเชยฺย ฐานเมตํ วิชฺชตี”ติ ปชานาติ, เย เย ธมฺมา เยสํ เยสํ ธมฺมานํ เหตู ปจฺจยา อุปาทาย ตํ ตํ ฐานํ เย เย ธมฺมา เยสํ เยสํ ธมฺมานํ น เหตู อปฺปจฺจยา อุปาทาย ตํ ตํ อฏฺฐานนฺติ ยา ตตฺถ ปญฺญา ปชานนา ฯเปฯ อโมโห ธมฺมวิจโย สมฺมาทิฏฺฐิ, อิทํ ตถาคตสฺส ฐานญฺจ ฐานโต อฏฺฐานญฺจ อฏฺฐานโต ยถาภูตํ ญาณํ. (1)}

\proseitem{810}{ตตฺถ กตมํ ตถาคตสฺส อตีตานาคตปจฺจุปฺปนฺนานํ กมฺมสมาทานานํ ฐานโส เหตุโส วิปากํ ยถาภูตํ ญาณํ, อิธ ตถาคโต ปชานาติ, อตฺเถกจฺจานิ ปาปกานิ กมฺมสมาทานานิ คติสมฺปตฺติปฏิพาฬฺหานิ น วิปจฺจนฺติ, อตฺเถกจฺจานิ ปาปกานิ กมฺมสมาทานานิ อุปธิสมฺปตฺติปฏิพาฬฺหานิ น วิปจฺจนฺติ, อตฺเถกจฺจานิ ปาปกานิ กมฺมสมาทานานิ กาลสมฺปตฺติปฏิพาฬฺหานิ น วิปจฺจนฺติ, อตฺเถกจฺจานิ ปาปกานิ กมฺมสมาทานานิ ปโยคสมฺปตฺติปฏิพาฬฺหานิ น วิปจฺจนฺติ.}

\prose{อตฺเถกจฺจานิ ปาปกานิ กมฺมสมาทานานิ คติวิปตฺติํ อาคมฺม วิปจฺจนฺติ, อตฺเถกจฺจานิ ปาปกานิ กมฺมสมาทานานิ อุปธิวิปตฺติํ อาคมฺม วิปจฺจนฺติ, อตฺเถกจฺจานิ ปาปกานิ กมฺมสมาทานานิ กาลวิปตฺติํ อาคมฺม วิปจฺจนฺติ, อตฺเถกจฺจานิ ปาปกานิ กมฺมสมาทานานิ ปโยควิปตฺติํ อาคมฺม วิปจฺจนฺติ.}

\prose{อตฺเถกจฺจานิ กลฺยาณานิ กมฺมสมาทานานิ คติวิปตฺติปฏิพาฬฺหานิ น วิปจฺจนฺติ, อตฺเถกจฺจานิ กลฺยาณานิ กมฺมสมาทานานิ อุปธิวิปตฺติปฏิพาฬฺหานิ น วิปจฺจนฺติ, อตฺเถกจฺจานิ กลฺยาณานิ กมฺมสมาทานานิ กาลวิปตฺติปฏิพาฬฺหานิ น วิปจฺจนฺติ, อตฺเถกจฺจานิ กลฺยาณานิ กมฺมสมาทานานิ ปโยควิปตฺติปฏิพาฬฺหานิ น วิปจฺจนฺติ.}

\prose{อตฺเถกจฺจานิ กลฺยาณานิ กมฺมสมาทานานิ คติสมฺปตฺติํ อาคมฺม วิปจฺจนฺติ, อตฺเถกจฺจานิ กลฺยาณานิ กมฺมสมาทานานิ อุปธิสมฺปตฺติํ อาคมฺม วิปจฺจนฺติ, อตฺเถกจฺจานิ กลฺยาณานิ กมฺมสมาทานานิ กาลสมฺปตฺติํ อาคมฺม วิปจฺจนฺติ, \csromanfolio{352}อตฺเถกจฺจานิ กลฺยาณานิ กมฺมสมาทานานิ ปโยคสมฺปตฺติํ อาคมฺม วิปจฺจนฺตีติ ยา ตตฺถ ปญฺญา ปชานนา ฯเปฯ อโมโห ธมฺมวิจโย สมฺมาทิฏฺฐิ, อิทํ ตถาคตสฺส อตีตานาคตปจฺจุปฺปนฺนานํ กมฺมสมาทานานํ ฐานโส เหตุโส วิปากํ ยถาภูตํ ญาณํ. (2)}

\proseitem{811}{ตตฺถ กตมํ ตถาคตสฺส สพฺพตฺถคามินิํ ปฏิปทํ ยถาภูตํ ญาณํ, อิธ ตถาคโต “อยํ มคฺโค อยํ ปฏิปทา นิรยคามี”ติ\footnote{นิรยคามินีติ (สฺยา)} ปชานาติ, “อยํ มคฺโค อยํ ปฏิปทา ติรจฺฉานโยนิคามี”ติ\footnote{ติรจฺฉานคามินีติ (สฺยา) เอวมุปริปิ. อฏฺฐกถา โอโลเกตพฺพา.} ปชานาติ, “อยํ มคฺโค อยํ ปฏิปทา เปตฺติวิสยคามี”ติ ปชานาติ, “อยํ มคฺโค อยํ ปฏิปทา มนุสฺสโลกคามี”ติ ปชานาติ, “อยํ มคฺโค อยํ ปฏิปทา เทวโลกคามี”ติ ปชานาติ, “อยํ มคฺโค อยํ ปฏิปทา นิพฺพานคามี”ติ ปชานาตีติ ยา ตตฺถ ปญฺญา ปชานนา ฯเปฯ อโมโห ธมฺมวิจโย สมฺมาทิฏฺฐิ, อิทํ ตถาคตสฺส สพฺพตฺถคามินิํ ปฏิปทํ ยถาภูตํ ญาณํ. (3)}

\proseitem{812}{ตตฺถ กตมํ ตถาคตสฺส อเนกธาตุนานาธาตุโลกํ ยถาภูตํ ญาณํ, อิธ ตถาคโต ขนฺธนานตฺตํ ปชานาติ, อายตนนานตฺตํ ปชานาติ, ธาตุนานตฺตํ ปชานาติ, อเนกธาตุนานาธาตุโลกนานตฺตํ ปชานาตีติ ยา ตตฺถ ปญฺญา ปชานนา ฯเปฯ อโมโห ธมฺมวิจโย สมฺมาทิฏฺฐิ, อิทํ ตถาคตสฺส อเนกธาตุนานาธาตุโลกํ ยถาภูตํ ญาณํ. (4)}

\proseitem{813}{ตตฺถ กตมํ ตถาคตสฺส สตฺตานํ นานาธิมุตฺติกตํ ยถาภูตํ ญาณํ, อิธ ตถาคโต ปชานาติ “สนฺติ สตฺตา หีนาธิมุตฺติกา, สนฺติ สตฺตา ปณีตาธิมุตฺติกา.\csromanspacer{} หีนาธิมุตฺติกา สตฺตา หีนาธิมุตฺติเก สตฺเต เสวนฺติ ภชนฺติ ปยิรุปาสนฺติ.\csromanspacer{} ปณีตาธิมุตฺติกา สตฺตา ปณีตาธิมุตฺติเก สตฺเต เสวนฺติ ภชนฺติ ปยิรุปาสนฺติ.}

\prose{อตีตมฺปิ อทฺธานํ หีนาธิมุตฺติกา สตฺตา หีนาธิมุตฺติเก สตฺเต เสวิํสุ ภชิํสุ ปยิรุปาสิํสุ.\csromanspacer{} ปณีตาธิมุตฺติกา สตฺตา ปณีตาธิมุตฺติเก สตฺเต เสวิํสุ ภชิํสุ ปยิรุปาสิํสุ.}

\prose{อนาคตมฺปิ อทฺธานํ หีนาธิมุตฺติกา สตฺตา หีนาธิมุตฺติเก สตฺเต เสวิสฺสนฺติ ภชิสฺสนฺติ ปยิรุปาสิสฺสนฺติ.\csromanspacer{} ปณีตาธิมุตฺติกา สตฺตา}

\csromanheader{2}{16. ญาณวิภงฺค}
\prose{\csromanfolio{353}ปณีตาธิมุตฺติเก สตฺเต เสวิสฺสนฺติ ภชิสฺสนฺติ ปยิรุปาสิสฺสนฺตี”ติ ยา ตตฺถ ปญฺญา ปชานนา ฯเปฯ อโมโห ธมฺมวิจโย สมฺมาทิฏฺฐิ, อิทํ ตถาคตสฺส สตฺตานํ นานาธิมุตฺติกตํ ยถาภูตํ ญาณํ. (5)}

\proseitem{814}{ตตฺถ กตมํ ตถาคตสฺส ปรสตฺตานํ ปรปุคฺคลานํ อินฺทฺริยปโรปริยตฺตํ ยถาภูตํ ญาณํ, อิธ ตถาคโต สตฺตานํ อาสยํ ปชานาติ, อนุสยํ ปชานาติ, จริตํ ปชานาติ, อธิมุตฺติํ ปชานาติ, อปฺปรชกฺเข มหารชกฺเข ติกฺขินฺทฺริเย มุทินฺทฺริเย สฺวากาเร ทฺวากาเร สุวิญฺญาปเย ทุวิญฺญาปเย ภพฺพาภพฺเพ สตฺเต ปชานาติ.}

\proseitem{815}{กตโม จ สตฺตานํ อาสโย, “สสฺสโต โลโก”ติ วา “อสสฺสโต โลโก”ติ วา “อนฺตวา โลโก”ติ วา “อนนฺตวา โลโก”ติ วา “ตํ ชีวํ ตํ สรีรนฺ”ติ วา “อญฺญํ ชีวํ อญฺญํ สรีรนฺ”ติ วา “โหติ ตถาคโต ปรํ มรณา”ติ วา “น โหติ ตถาคโต ปรํ มรณา”ติ วา “โหติ จ น จ โหติ ตถาคโต ปรํ มรณา”ติ วา “เนว โหติ น น โหติ ตถาคโต ปรํ มรณา”ติ วา อิติ ภวทิฏฺฐิสนฺนิสฺสิตา วา สตฺตา โหนฺติ วิภวทิฏฺฐิสนฺนิสฺสิตา วา, เอเต วา ปน อุโภ อนฺเต อนุปคมฺม อิทปฺปจฺจยตา ปฏิจฺจสมุปฺปนฺเนสุ ธมฺเมสุ อนุโลมิกา ขนฺติ ปฏิลทฺธา โหติ ยถาภูตํ ญาณํ, อยํ สตฺตานํ อาสโย.}

\proseitem{816}{กตโม จ สตฺตานํ อนุสโย, สตฺตานุสยา กามราคานุสโย ปฏิฆานุสโย มานานุสโย ทิฏฺฐานุสโย วิจิกิจฺฉานุสโย ภวราคานุสโย อวิชฺชานุสโย, ยํ โลเก ปิยรูปํ สาตรูปํ, เอตฺถ สตฺตานํ ราคานุสโย อนุเสติ, ยํ โลเก อปฺปิยรูปํ อสาตรูปํ, เอตฺถ สตฺตานํ ปฏิฆานุสโย อนุเสติ, อิติ อิเมสุ ทฺวีสุ ธมฺเมสุ อวิชฺชานุปติตา ตเทกฏฺโฐ มาโน จ ทิฏฺฐิ จ วิจิกิจฺฉา จ ทฏฺฐพฺพา, อยํ สตฺตานํ อนุสโย.}

\proseitem{817}{กตมญฺจ สตฺตานํ จริตํ, ปุญฺญาภิสงฺขาโร อปุญฺญาภิสงฺขาโร อาเนญฺชาภิสงฺขาโร ปริตฺตภูมโก วา มหาภูมโก วา, อิทํ สตฺตานํ จริตํ.}

\proseitem{818}{\csromanfolio{354}กตมา จ สตฺตานํ อธิมุตฺติ, สนฺติ สตฺตา หีนาธิมุตฺติกา, สนฺติ สตฺตา ปณีตาธิมุตฺติกา, หีนาธิมุตฺติกา สตฺตา หีนาธิมุตฺติเก สตฺเต เสวนฺติ ภชนฺติ ปยิรุปาสนฺติ.\csromanspacer{} ปณีตาธิมุตฺติกา สตฺตา ปณีตาธิมุตฺติเก สตฺเต เสวนฺติ ภชนฺติ ปยิรุปาสนฺติ.}

\prose{อตีตมฺปิ อทฺธานํ หีนาธิมุตฺติกา สตฺตา หีนาธิมุตฺติเก สตฺเต เสวิํสุ ภชิํสุ ปยิรุปาสิํสุ.\csromanspacer{} ปณีตาธิมุตฺติกา สตฺตา ปณีตาธิมุตฺติเก สตฺเต เสวิํสุ ภชิํสุ ปยิรุปาสิํสุ.}

\prose{อนาคตมฺปิ อทฺธานํ หีนาธิมุตฺติกา สตฺตา หีนาธิมุตฺติเก สตฺเต เสวิสฺสนฺติ ภชิสฺสนฺติ ปยิรุปาสิสฺสนฺติ.\csromanspacer{} ปณีตาธิมุตฺติกา สตฺตา ปณีตาธิมุตฺติเก สตฺเต เสวิสฺสนฺติ ภชิสฺสนฺติ ปยิรุปาสิสฺสนฺติ, อยํ สตฺตานํ อธิมุตฺติ.}

\proseitem{819}{กตเม เต สตฺตา มหารชกฺขา, ทส กิเลสวตฺถูนิ โลโภ โทโส โมโห มาโน ทิฏฺฐิ วิจิกิจฺฉา ถินํ อุทฺธจฺจํ อหิริกํ อโนตฺตปฺปํ, เยสํ สตฺตานํ อิมานิ ทส กิเลสวตฺถูนิ อาเสวิตานิ ภาวิตานิ พหุลีกตานิ อุสฺสทคตานิ, อิเม เต สตฺตา มหารชกฺขา.}

\proseitem{820}{กตเม เต สตฺตา อปฺปรชกฺขา, เยสํ สตฺตานํ อิมานิ ทส กิเลสวตฺถูนิ อนาเสวิตานิ อภาวิตานิ อพหุลีกตานิ อนุสฺสทคตานิ, อิเม เต สตฺตา อปฺปรชกฺขา.}

\proseitem{821}{กตเม เต สตฺตา มุทินฺทฺริยา, ปญฺจินฺทฺริยานิ สทฺธินฺทฺริยํ วีริยินฺทฺริยํ สตินฺทฺริยํ สมาธินฺทฺริยํ ปญฺญินฺทฺริยํ, เยสํ สตฺตานํ อิมานิ ปญฺจินฺทฺริยานิ อนาเสวิตานิ อภาวิตานิ อพหุลีกตานิ อนุสฺสทคตานิ, อิเม เต สตฺตา มุทินฺทฺริยา.}

\proseitem{822}{กตเม เต สตฺตา ติกฺขินฺทฺริยา, เยสํ สตฺตานํ อิมานิ ปญฺจินฺทฺริยานิ อาเสวิตานิ ภาวิตานิ พหุลีกตานิ อุสฺสทคตานิ, อิเม เต สตฺตา ติกฺขินฺทฺริยา.}

\csromanheader{2}{16. ญาณวิภงฺค}
\proseitem{823}{\csromanfolio{355}กตเม เต สตฺตา ทฺวาการา, เย เต สตฺตา ปาปาสยา ปาปานุสยา ปาปจริตา ปาปาธิมุตฺติกา มหารชกฺขา มุทินฺทฺริยา, อิเม เต สตฺตา ทฺวาการา.}

\proseitem{824}{กตเม เต สตฺตา สฺวาการา, เย เต สตฺตา กลฺยาณาสยา กลฺยาณจริตา กลฺยาณาธิมุตฺติกา อปฺปรชกฺขา ติกฺขินฺทฺริยา, อิเม เต สตฺตา สฺวาการา.}

\proseitem{825}{กตเม เต สตฺตา ทุวิญฺญาปยา, เย จ เต สตฺตา ทฺวาการา, เตว เต สตฺตา ทุวิญฺญาปยา.\csromanspacer{} เย จ เต สตฺตา สฺวาการา, เตว เต สตฺตา สุวิญฺญาปยา.}

\proseitem{826}{กตเม เต สตฺตา อภพฺพา, เย เต สตฺตา กมฺมาวรเณน สมนฺนาคตา กิเลสาวรเณน สมนฺนาคตา วิปากาวรเณน สมนฺนาคตา อสฺสทฺธา อจฺฉนฺทิกา ทุปฺปญฺญา อภพฺพา นิยามํ โอกฺกมิตุํ กุสเลสุ ธมฺเมสุ สมฺมตฺตํ, อิเม เต สตฺตา อภพฺพา.}

\proseitem{827}{กตเม เต สตฺตา ภพฺพา, เย เต สตฺตา น กมฺมาวรเณน สมนฺนาคตา น กิเลสาวรเณน สมนฺนาคตา น วิปากาวรเณน สมนฺนาคตา สทฺธา ฉนฺทิกา ปญฺญวนฺโต ภพฺพา นิยามํ โอกฺกมิตุํ กุสเลสุ ธมฺเมสุ สมฺมตฺตํ, อิเม เต สตฺตา ภพฺพาติ ยา ตตฺถ ปญฺญา ปชานนา ฯเปฯ อโมโห ธมฺมวิจโย สมฺมาทิฏฺฐิ, อิทํ ตถาคตสฺส ปรสตฺตานํ ปรปุคฺคลานํ อินฺทฺริยปโรปริยตฺตํ ยถาภูตํ ญาณํ. (6)}

\proseitem{828}{ตตฺถ กตมํ ตถาคตสฺส ฌานวิโมกฺขสมาธิสมาปตฺตีนํ สํกิเลสํ โวทานํ วุฏฺฐานํ ยถาภูตํ ญาณํ, ฌายีติ จตฺตาโร ฌายี, อตฺเถกจฺโจ ฌายี สมฺปตฺติํเยว\footnote{สมฺปตฺติเยว (ก)} สมานํ “วิปตฺตี”ติ ปจฺเจติ, อตฺเถกจฺโจ ฌายี วิปตฺติํเยว\footnote{วิปตฺติเยว (ก)} สมานํ “สมฺปตฺตี”ติ ปจฺเจติ, อตฺเถกจฺโจ ฌายี สมฺปตฺติํเยว สมานํ “สมฺปตฺตี”ติ ปจฺเจติ, อตฺเถกจฺโจ ฌายี วิปตฺติํเยว สมานํ “วิปตฺตี”ติ ปจฺเจติ, อิเม จตฺตาโร ฌายี.}

\prose{อปเรปิ จตฺตาโร ฌายี, อตฺเถกจฺโจ ฌายี ทนฺธํ สมาปชฺชติ ขิปฺปํ วุฏฺฐาติ, อตฺเถกจฺโจ ฌายี ขิปฺปํ สมาปชฺชติ ทนฺธํ วุฏฺฐาติ, อตฺเถกจฺโจ \csromanfolio{356}ฌายี ทนฺธํ สมาปชฺชติ ทนฺธํ วุฏฺฐาติ, อตฺเถกจฺโจ ฌายี ขิปฺปํ สมาปชฺชติ ขิปฺปํ วุฏฺฐาติ, อิเม จตฺตาโร ฌายี.}

\prose{อปเรปิ จตฺตาโร ฌายี, อตฺเถกจฺโจ ฌายี สมาธิสฺมิํ สมาธิกุสโล โหติ น สมาธิสฺมิํ สมาปตฺติกุสโล, อตฺเถกจฺโจ ฌายี สมาธิสฺมิํ สมาปตฺติกุสโล โหติ น สมาธิสฺมิํ สมาธิกุสโล, อตฺเถกจฺโจ ฌายี สมาธิสฺมิํ สมาธิกุสโล จ โหติ สมาธิสฺมิํ สมาปตฺติกุสโล จ, อตฺเถกจฺโจ ฌายี เนว สมาธิสฺมิํ สมาธิกุสโล โหติ น สมาธิสฺมิํ สมาปตฺติกุสโล, อิเม จตฺตาโร ฌายี.}

\prose{ฌานนฺติ จตฺตาริ ฌานานิ, ปฐมํ ฌานํ ทุติยํ ฌานํ ตติยํ ฌานํ จตุตฺถํ ฌานํ.}

\prose{วิโมกฺโขติ อฏฺฐ วิโมกฺขา, รูปี รูปานิ ปสฺสติ, อยํ ปฐโม วิโมกฺโข.}

\prose{อชฺฌตฺตํ อรูปสญฺญี พหิทฺธา รูปานิ ปสฺสติ, อยํ ทุติโย วิโมกฺโข.}

\prose{สุภนฺเตว อธิมุตฺโต โหติ, อยํ ตติโย วิโมกฺโข.}

\prose{สพฺพโส รูปสญฺญานํ สมติกฺกมา ปฏิฆสญฺญานํ อตฺถงฺคมา นานตฺตสญฺญานํ อมนสิการา “อนนฺโต อากาโส”ติ อากาสานญฺจายตนํ อุปสมฺปชฺช วิหรติ, อยํ จตุตฺโถ วิโมกฺโข.}

\prose{สพฺพโส อากาสานญฺจายตนํ สมติกฺกมฺม “อนนฺตํ วิญฺญาณนฺ”ติ วิญฺญาณญฺจายตนํ อุปสมฺปชฺช วิหรติ, อยํ ปญฺจโม วิโมกฺโข.}

\prose{สพฺพโส วิญฺญาณญฺจายตนํ สมติกฺกมฺม “นตฺถิ กิญฺจิ”ติ อากิญฺจญฺญายตนํ อุปสมฺปชฺช วิหรติ, อยํ ฉฏฺโฐ วิโมกฺโข.}

\prose{สพฺพโส อากิญฺจญฺญายตนํ สมติกฺกมฺม เนวสญฺญานาสญฺญายตนํ อุปสมฺปชฺช วิหรติ, อยํ สตฺตโม วิโมกฺโข.}

\prose{สพฺพโส เนวสญฺญานาสญฺญายตนํ สมติกฺกมฺม สญฺญาเวทยิตนิโรธํ อุปสมฺปชฺช วิหรติ, อยํ อฏฺฐโม วิโมกฺโข.}

\prose{สมาธีติ ตโย สมาธี, สวิตกฺกสวิจาโร สมาธิ, อวิตกฺกวิจารมตฺโต สมาธิ, อวิตกฺก-อวิจาโร สมาธิ.}

\csromanheader{2}{16. ญาณวิภงฺค}
\prose{\csromanfolio{357}สมาปตฺตีติ นว อนุปุพฺพวิหารสมาปตฺติโย, ปฐมชฺฌานสมาปตฺติ ทุติยชฺฌานสมาปตฺติ ตติยชฺฌานสมาปตฺติ จตุตฺถชฺฌานสมาปตฺติ อากาสานญฺจายตนสมาปตฺติ วิญฺญาณญฺจายตนสมาปตฺติ อากิญฺจญฺญายตนสมาปตฺติ เนวสญฺญานาสญฺญายตนสมาปตฺติ สญฺญาเวทยิตนิโรธสมาปตฺติ.}

\prose{สํกิเลสนฺติ หานภาคิโย ธมฺโม.\csromanspacer{} โวทานนฺติ วิเสสภาคิโย ธมฺโม, วุฏฺฐานนฺติ โวทานมฺปิ วุฏฺฐานํ, ตมฺหา ตมฺหา สมาธิมฺหา วุฏฺฐานมฺปิ วุฏฺฐานนฺติ ยา ตตฺถ ปญฺญา ปชานนา ฯเปฯ อโมโห ธมฺมวิจโย สมฺมาทิฏฺฐิ, อิทํ ตถาคตสฺส ฌานวิโมกฺขสมาธิสมาปตฺตีนํ สํกิเลสํ โวทานํ วุฏฺฐานํ ยถาภูตํ ญาณํ. (7)}

\proseitem{829}{ตตฺถ กตมํ ตถาคตสฺส ปุพฺเพนิวาสานุสฺสติ ยถาภูตํ ญาณํ, อิธ ตถาคโต อเนกวิหิตํ ปุพฺเพนิวาสํ อนุสฺสรติ.\csromanspacer{} เสยฺยถิทํ, เอกมฺปิ ชาติํ เทฺวปิ ชาติโย ติสฺโสปิ ชาติโย จตสฺโสปิ ชาติโย ปญฺจปิ ชาติโย ทสปิ ชาติโย วีสมฺปิ ชาติโย ติํสมฺปิ ชาติโย จตฺตาลีสมฺปิ ชาติโย ปญฺญาสมฺปิ ชาติโย ชาติสตมฺปิ ชาติสหสฺสมฺปิ ชาติสตสหสฺสมฺปิ อเนเกปิ สํวฏฺฏกปฺเป อเนเกปิ วิวฏฺฏกปฺเป อเนเกปิ สํวฏฺฏวิวฏฺฏกปฺเป “อมุตฺราสิํ เอวํนาโม เอวํโคตฺโต เอวํวณฺโณ เอวมาหาโร เอวํสุขทุกฺขปฺปฏิสํเวที เอวมายุปริยนฺโต, โส ตโต จุโต อมุตฺร อุทปาทิํ, ตตฺราปาสิํ เอวํนาโม เอวํโคตฺโต เอวํวณฺโณ เอวมาหาโร เอวํสุขทุกฺขปฺปฏิสํเวที เอวมายุปริยนฺโต, โส ตโต จุโต อิธูปปนฺโน”ติ อิติ สาการํ ส-อุทฺเทสํ อเนกวิหิตํ ปุพฺเพนิวาสํ อนุสฺสรตีติ ยา ตตฺถ ปญฺญา ปชานนา ฯเปฯ อโมโห ธมฺมวิจโย สมฺมาทิฏฺฐิ, อิทํ ตถาคตสฺส ปุพฺเพนิวาสานุสฺสติ ยถาภูตํ ญาณํ. (8)}

\proseitem{830}{ตตฺถ กตมํ ตถาคตสฺส สตฺตานํ จุตูปปาตํ ยถาภูตํ ญาณํ, อิธ ตถาคโต ทิพฺเพน จกฺขุนา วิสุทฺเธน อติกฺกนฺตมานุสเกน สตฺเต ปสฺสติ จวมาเน อุปปชฺชมาเน หีเน ปณีเต สุวณฺเณ ทุพฺพณฺเณ สุคเต ทุคฺคเต, ยถากมฺมูปเค สตฺเต ปชานาติ “อิเม วต โภนฺโต สตฺตา กายทุจฺจริเตน สมนฺนาคตา วจีทุจฺจริเตน สมนฺนาคตา \csromanfolio{358}มโนทุจฺจริเตน สมนฺนาคตา อริยานํ อุปวาทกา มิจฺฉาทิฏฺฐิกา มิจฺฉาทิฏฺฐิกมฺมสมาทานา, เต กายสฺส เภทา ปรํ มรณา อปายํ ทุคฺคติํ วินิปาตํ นิรยํ อุปปนฺนา, อิเม วา ปน โภนฺโต สตฺตา กายสุจริเตน สมนฺนาคตา วจีสุจริเตน สมนฺนาคตา มโนสุจริเตน สมนฺนาคตา อริยานํ อนุปวาทกา สมฺมาทิฏฺฐิกา สมฺมาทิฏฺฐิกมฺมสมาทานา, เต กายสฺส เภทา ปรํ มรณา สุคติํ สคฺคํ โลกํ อุปปนฺนา”ติ, อิติ ทิพฺเพน จกฺขุนา วิสุทฺเธน อติกฺกนฺตมานุสเกน สตฺเต ปสฺสติ จวมาเน อุปปชฺชมาเน หีเน ปณีเต สุวณฺเณ ทุพฺพณฺเณ สุคเต ทุคฺคเต, ยถากมฺมูปเค สตฺเต ปชานาตีติ ยา ตตฺถ ปญฺญา ปชานนา ฯเปฯ อโมโห ธมฺมวิจโย สมฺมาทิฏฺฐิ, อิทํ ตถาคตสฺส สตฺตานํ จุตูปปาตํ ยถาภูตํ ญาณํ. (9)}

\proseitem{831}{ตตฺถ กตมํ ตถาคตสฺส อาสวานํ ขยํ ยถาภูตํ ญาณํ, อิธ ตถาคโต อาสวานํ ขยา อนาสวํ เจโตวิมุตฺติํ ปญฺญาวิมุตฺติํ ทิฏฺเฐว ธมฺเม สยํ อภิญฺญา สจฺฉิกตฺวา อุปสมฺปชฺช วิหรตีติ ยา ตตฺถ ปญฺญา ปชานนา ฯเปฯ อโมโห ธมฺมวิจโย สมฺมาทิฏฺฐิ, อิทํ ตถาคตสฺส อาสวานํ ขยํ ยถาภูตํ ญาณนฺติ. (10)}

\nitthitam{\textbf{ญาณวิภงฺโค นิฏฺฐิโต.}}
\csromanmatikaanchor{mtk.168}
\csromantocmarknopage{subparagraph}{17. ขุทฺทกวตฺถุวิภงฺค}{mtk.168}
\bookmark[dest={mtk.168},level=5]{17. ขุทฺทกวตฺถุวิภงฺค}
\csromanheader{6}{17. ขุทฺทกวตฺถุวิภงฺค}
\csromanmatikaanchor{mtk.169}
\csromantocmarkref{subparagraph}{1. เอกกมาติกา}{mtk.169}
\bookmark[dest={mtk.169},level=5]{1. เอกกมาติกา}
\csromanheader{6}{1. เอกกมาติกา}
\proseitem{832}{\csromanfolio{359}ชาติมโท, โคตฺตมโท, อาโรคฺยมโท, โยพฺพนมโท, ชีวิตมโท, ลาภมโท, สกฺการมโท, ครุการมโท, ปุเรกฺขารมโท, ปริวารมโท, โภคมโท, วณฺณมโท, สุตมโท, ปฏิภานมโท, รตฺตญฺญุมโท, ปิณฺฑปาติกมโท, อนวญฺญาตมโท, อิริยาปถมโท, อิทฺธิมโท, ยสมโท, สีลมโท, ฌานมโท, สิปฺปมโท, อาโรหมโท, ปริณาหมโท, สณฺฐานมโท, ปาริปูริมโท, มโท, ปมาโท, ถมฺโภ, สารมฺโภ, อตฺริจฺฉตา, มหิจฺฉตา, ปาปิจฺฉตา, สิงฺคํ, ตินฺติณํ, จาปลฺยํ, อสภาควุตฺติ, อรติ, ตนฺที, วิชมฺภิตา, ภตฺตสมฺมโท, เจตโส จ ลีนตฺตํ\footnote{เจตโส ลีนตฺตํ (สฺยา)}, กุหนา, ลปนา, เนมิตฺติกตา, นิปฺเปสิกตา, ลาเภน ลาภํ นิชิคีสนตา\footnote{นิชิคิํสนตา (สี), ชิคิํสนตา (สฺยา)} “เสยฺโยหมสฺมี”ติ มาโน, “สทิโสหมสฺมี”ติ มาโน, “หีโน หมสฺมี”ติ มาโน, เสยฺยสฺส “เสยฺโยหมสฺมี”ติ มาโน, เสยฺยสฺส “สทิโสหมสฺมี”ติ มาโน, เสยฺยสฺส “หีโนหมสฺมี”ติ มาโน, สทิสสฺส “เสยฺโยหมสฺมี”ติ มาโน, สทิสสฺส “สทิโสหมสฺมี”ติ มาโน, สทิสสฺส “หีโนหมสฺมี”ติ มาโน, หีนสฺส “เสยฺโยหมสฺมี”ติ มาโน, หีนสฺส “สทิโสหมสฺมี”ติ มาโน, “หีนสฺส หีโนหมสฺมี”ติ มาโน, มาโน, อติมาโน, มานาติมาโน, โอมาโน, อธิมาโน, อสฺมิมาโน, มิจฺฉามาโน, ญาติวิตกฺโก, ชนปทวิตกฺโก, อมรวิตกฺโก, ปรานุทฺทยตาปฏิสํยุตฺโต วิตกฺโก, ลาภสกฺการสิโลกปฏิสํยุตฺโต วิตกฺโก, อนวญฺญตฺติปฏิสํยุตฺโต วิตกฺโก.}

\csromanheader{2}{เอกกํ.}
\csromansectionrule
\csromanmatikaanchor{mtk.170}
\csromantocmarkref{subparagraph}{2. ทุกมาติกา}{mtk.170}
\bookmark[dest={mtk.170},level=5]{2. ทุกมาติกา}
\csromanheader{6}{2. ทุกมาติกา}
\proseitem{833}{โกโธ จ อุปนาโห จ.\csromanspacer{} มกฺโข จ ปฬาโส\footnote{ปลาโส (สี, สฺยา)} จ.\csromanspacer{} อิสฺสา จ มจฺฉริยญฺจ.\csromanspacer{} มายา จ สาเฐยฺยญฺจ.\csromanspacer{} อวิชฺชา จ ภวตณฺหา \csromanfolio{360}จ.\csromanspacer{} ภวทิฏฺฐิ จ วิภวทิฏฺฐิ จ.\csromanspacer{} สสฺสตทิฏฺฐิ จ อุจฺเฉททิฏฺฐิ จ.\csromanspacer{} อนฺตวาทิฏฺฐิ จ อนนฺตวาทิฏฺฐิ จ.\csromanspacer{} ปุพฺพนฺตานุทิฏฺฐิ จ อปรนฺตานุทิฏฺฐิ จ.\csromanspacer{} อหิริกญฺจ อโนตฺตปฺปญฺจ.\csromanspacer{} โทวจสฺสตา จ ปาปมิตฺตตา จ.\csromanspacer{} อนชฺชโว จ อมทฺทโว จ.\csromanspacer{} อกฺขนฺติ จ อโสรจฺจญฺจ.\csromanspacer{} อสาขลฺยญฺจ อปฺปฏิสนฺถาโร จ.\csromanspacer{} อินฺทฺริเยสุ อคุตฺตทฺวารตา จ โภชเน อมตฺตญฺญุตา จ.\csromanspacer{} มุฏฺฐสฺสจฺจญฺจ อสมฺปชญฺญญฺจ.\csromanspacer{} สีลวิปตฺติ จ ทิฏฺฐิวิปตฺติ จ.\csromanspacer{} อชฺฌตฺตสํโยชนญฺจ พหิทฺธาสํโยชนญฺจ.}

\csromanheader{2}{ทุกํ.}
\csromansectionrule
\csromanmatikaanchor{mtk.171}
\csromantocmarkref{subparagraph}{3. ติกมาติกา}{mtk.171}
\bookmark[dest={mtk.171},level=5]{3. ติกมาติกา}
\csromanheader{6}{3. ติกมาติกา}
\proseitem{834}{ตีณิ อกุสลมูลานิ, ตโย อกุสลวิตกฺกา, ติสฺโส อกุสลสญฺญา, ติสฺโส อกุสลธาตุโย, ตีณิ ทุจฺจริตานิ, ตโย อาสวา, ตีณิ สํโยชนานิ, ติสฺโส ตณฺหา, อปราปิ ติสฺโส ตณฺหา, ติสฺโส เอสนา, ติสฺโส วิธา, ตีณิ ภยานิ, ตีณิ ตมานิ, ตีณิ ติตฺถายตนานิ, ตโย กิญฺจนา, ตีณิ องฺคณานิ, ตีณิ มลานิ, ตีณิ วิสมานิ, อปรานิปิ ตีณิ วิสมานิ, ตโย อคฺคี, ตโย กสาวา, อปเรปิ ตโย กสาวา.}

\prose{อสฺสาททิฏฺฐิ อตฺตานุทิฏฺฐิ มิจฺฉาทิฏฺฐิ.\csromanspacer{} อรติ วิเหสา อธมฺมจริยา.\csromanspacer{} โทวจสฺสตา ปาปมิตฺตตา นานตฺตสญฺญา.\csromanspacer{} อุทฺธจฺจํ โกสชฺชํ ปมาโท.\csromanspacer{} อสนฺตุฏฺฐิตา อสมฺปชญฺญตา มหิจฺฉตา.\csromanspacer{} อหิริกํ อโนตฺตปฺปํ ปมาโท.\csromanspacer{} อนาทริยํ โทวจสฺสตา ปาปมิตฺตตา.\csromanspacer{} อสฺสทฺธิยํ อวทญฺญุตา โกสชฺชํ.\csromanspacer{} อุทฺธจฺจํ อสํวโร ทุสฺสีลฺยํ.\csromanspacer{} อริยานํ อทสฺสนกมฺยตา สทฺธมฺมํ อโสตุกมฺยตา อุปารมฺภจิตฺตตา.\csromanspacer{} มุฏฺฐสฺสจฺจํ อสมฺปชญฺญํ เจตโส วิกฺเขโป.\csromanspacer{} อโยนิโส มนสิกาโร กุมฺมคฺคเสวนา เจตโส จ ลีนตฺตํ.}

\csromanheader{2}{ติกํ.}
\csromansectionrule
\csromanmatikaanchor{mtk.172}
\csromantocmarkref{subparagraph}{4. จตุกฺกมาติกา}{mtk.172}
\bookmark[dest={mtk.172},level=5]{4. จตุกฺกมาติกา}
\csromanheader{6}{4. จตุกฺกมาติกา}
\proseitem{835}{จตฺตาโร อาสวา, จตฺตาโร คนฺถา, จตฺตาโร โอฆา, จตฺตาโร โยคา, จตฺตาริ อุปาทานานิ, จตฺตาโร ตณฺหุปฺปาทา, จตฺตาริ}

\vspace{\tipitakaparskip}
\csromancenter{ขุทฺทกวตฺถุวิภงฺค อคติคมนานิ, จตฺตาโร วิปริยาสา\footnote{วิปริเยสา (สี, สฺยา, ก)}, จตฺตาโร อนริยโวหารา, อปเรปิ จตฺตาโร อนริยโวหารา, จตฺตาริ ทุจฺจริตานิ, อปรานิปิ จตฺตาริ ทุจฺจริตานิ, จตฺตาริ ภยานิ, อปรานิปิ จตฺตาริ ภยานิ, จตสฺโส ทิฏฺฐิโย.}
\csromanheader{2}{จตุกฺกํ.}
\csromansectionrule
\csromanmatikaanchor{mtk.173}
\csromantocmarkref{subparagraph}{5. ปญฺจกมาติกา}{mtk.173}
\bookmark[dest={mtk.173},level=5]{5. ปญฺจกมาติกา}
\csromanheader{6}{5. ปญฺจกมาติกา}
\proseitem{836}{\csromanfolio{361}ปญฺโจรมฺภาคิยานิ สํโยชนานิ, ปญฺจุทฺธมฺภาคิยานิ สํโยชนานิ, ปญฺจ มจฺฉริยานิ, ปญฺจ สงฺคา, ปญฺจ สลฺลา, ปญฺจ เจโตขิลา, ปญฺจ เจตโสวินิพนฺธา\footnote{เจตโสวินิพทฺธา (ก)}, ปญฺจ นีวรณานิ, ปญฺจ กมฺมานิ อานนฺตริกานิ, ปญฺจ ทิฏฺฐิโย, ปญฺจ เวรา, ปญฺจ พฺยสนา, ปญฺจ อกฺขนฺติยา อาทีนวา, ปญฺจ ภยานิ, ปญฺจ ทิฏฺฐธมฺมนิพฺพานวาทา.}

\csromanheader{2}{ปญฺจกํ.}
\csromansectionrule
\csromanmatikaanchor{mtk.174}
\csromantocmarkref{subparagraph}{6. ฉกฺกมาติกา}{mtk.174}
\bookmark[dest={mtk.174},level=5]{6. ฉกฺกมาติกา}
\csromanheader{6}{6. ฉกฺกมาติกา}
\proseitem{837}{ฉ วิวาทมูลานิ, ฉ ฉนฺทราคา, ฉ วิโรธวตฺถูนิ, ฉ ตณฺหากายา, ฉ อคารวา, ฉ ปริหานิยา ธมฺมา, อปเรปิ ฉ ปริหานิยา ธมฺมา, ฉ โสมนสฺสุปวิจารา, ฉ โทมนสฺสุปวิจารา, ฉ อุเปกฺขุปวิจารา, ฉ เคหสิตานิ โสมนสฺสานิ, ฉ เคหสิตานิ โทมนสฺสานิ, ฉ เคหสิตา อุเปกฺขา, ฉ ทิฏฺฐิโย.}

\csromanheader{2}{ฉกฺกํ.}
\csromansectionrule
\csromanmatikaanchor{mtk.175}
\csromantocmarkref{subparagraph}{7. สตฺตกมาติกา}{mtk.175}
\bookmark[dest={mtk.175},level=5]{7. สตฺตกมาติกา}
\csromanheader{6}{7. สตฺตกมาติกา}
\proseitem{838}{สตฺต อนุสยา, สตฺต สํโยชนานิ, สตฺต ปริยุฏฺฐานานิ, สตฺต อสทฺธมฺมา, สตฺต ทุจฺจริตานิ, สตฺต มานา, สตฺต ทิฏฺฐิโย.\csromanspacer{} สตฺตกํ.}

\csromanmatikaanchor{mtk.176}
\csromantocmarkref{subparagraph}{8. อฏฺฐกมาติกา}{mtk.176}
\bookmark[dest={mtk.176},level=5]{8. อฏฺฐกมาติกา}
\csromanheader{6}{8. อฏฺฐกมาติกา}
\proseitem{839}{\csromanfolio{362}อฏฺฐ กิเลสวตฺถูนิ, อฏฺฐ กุสีตวตฺถูนิ, อฏฺฐสุ โลกธมฺเมสุ จิตฺตสฺส ปฏิฆาโต, อฏฺฐ อนริยโวหารา, อฏฺฐ มิจฺฉตฺตา, อฏฺฐ ปุริสโทสา, อฏฺฐ อสญฺญิวาทา, อฏฺฐ เนวสญฺญินาสญฺญิวาทา.}

\csromanheader{2}{อฏฺฐกํ.}
\csromansectionrule
\csromanmatikaanchor{mtk.177}
\csromantocmarkref{subparagraph}{9. นวกมาติกา}{mtk.177}
\bookmark[dest={mtk.177},level=5]{9. นวกมาติกา}
\csromanheader{6}{9. นวกมาติกา}
\proseitem{840}{นว อาฆาตวตฺถูนิ, นว ปุริสมลานิ, นววิธา มานา, นว ตณฺหามูลกา ธมฺมา, นว อิญฺชิตานิ, นว มญฺญิตานิ, นว ผนฺทิตานิ, นว ปปญฺจิตานิ, นว สงฺขตานิ.}

\csromanheader{2}{นวกํ.}
\csromansectionrule
\csromanmatikaanchor{mtk.178}
\csromantocmarkref{subparagraph}{10. ทสกมาติกา}{mtk.178}
\bookmark[dest={mtk.178},level=5]{10. ทสกมาติกา}
\csromanheader{6}{10. ทสกมาติกา}
\proseitem{841}{ทส กิเลสวตฺถูนิ, ทส อาฆาตวตฺถูนิ, ทส อกุสลกมฺมปถา, ทส สํโยชนานิ, ทส มิจฺฉตฺตา, ทสวตฺถุกา มิจฺฉาทิฏฺฐิ, ทสวตฺถุกา อนฺตคฺคาหิกา ทิฏฺฐิ.\csromanspacer{} ทสกํ.}

\proseitem{842}{อฏฺฐารส ตณฺหาวิจริตานิ อชฺฌตฺติกสฺส อุปาทาย, อฏฺฐรส ตณฺหาวิจริตานิ พาหิรสฺส อุปาทาย, ตเทกชฺฌํ อภิสญฺญุหิตฺวา อภิสงฺขิปิตฺวา ฉตฺติํส ตณฺหาวิจริตานิ โหนฺติ, อิติ อตีตานิ ฉตฺติํส ตณฺหาวิจริตานิ อนาคตานิ ฉตฺติํส ตณฺหาวิจริตานิ ปจฺจุปฺปนฺนานิ ฉตฺติํส ตณฺหาวิจริตานิ, ตเทกชฺฌํ อภิสญฺญุหิตฺวา อภิสงฺขิปิตฺวา อฏฺฐ ตณฺหาวิจริตสตํ โหติ, ยานิ จ ทฺวาสฏฺฐิ ทิฏฺฐิคตานิ พฺรหฺมชาเล เวยฺยากรเณ วุตฺตานิ ภควตา.}

\csromanheader{2}{มาติกา.}
\csromanmatikaanchor{mtk.179}
\csromantocmarkref{subparagraph}{1. เอกกนิทฺเทส}{mtk.179}
\bookmark[dest={mtk.179},level=5]{1. เอกกนิทฺเทส}
\csromanheader{6}{17. ขุทฺทกวตฺถุวิภงฺค \textbf{1. เอกกนิทฺเทส}}
\proseitem{843}{\csromanfolio{363}ตตฺถ กตโม ชาติมโท, ชาติํ ปฏิจฺจ มโท มชฺชนา มชฺชิตตฺตํ มาโน มญฺญนา มญฺญิตตฺตํ อุนฺนติ อุนฺนาโม\footnote{อุณฺณติ อุณฺณาโม (สฺยา, ก), อภิ 1. 225 ปิฏฺเฐปิ.} ธโช สมฺปคฺคาโห เกตุกมฺยตา จิตฺตสฺส, อยํ วุจฺจติ ชาติมโท.}

\proseitem{844}{ตตฺถ กตโม โคตฺตมโท, โคตฺตํ ปฏิจฺจ ฯเปฯ อาโรคฺยํ ปฏิจฺจ ฯเปฯ โยพฺพนํ ปฏิจฺจ ฯเปฯ ชีวิตํ ปฏิจฺจ ฯเปฯ ลาภํ ปฏิจฺจ ฯเปฯ สกฺการํ ปฏิจฺจ ฯเปฯ ครุการํ ปฏิจฺจ ฯเปฯ ปุเรกฺขารํ ปฏิจฺจ ฯเปฯ ปริวารํ ปฏิจฺจ ฯเปฯ โภคํ ปฏิจฺจ ฯเปฯ วณฺณํ ปฏิจฺจ ฯเปฯ สุตํ ปฏิจฺจ ฯเปฯ ปฏิภานํ ปฏิจฺจ ฯเปฯ รตฺตญฺญุตํ ปฏิจฺจ ฯเปฯ ปิณฺฑปาติกตฺตํ ปฏิจฺจ ฯเปฯ อนวญฺญาตํ ปฏิจฺจ ฯเปฯ อิริยาปถํ ปฏิจฺจ ฯเปฯ อิทฺธิํ ปฏิจฺจ ฯเปฯ ยสํ ปฏิจฺจ ฯเปฯ สีลํ ปฏิจฺจ ฯเปฯ ฌานํ ปฏิจฺจ ฯเปฯ สิปฺปํ ปฏิจฺจ ฯเปฯ อาโรหํ ปฏิจฺจ ฯเปฯ ปริณาหํ ปฏิจฺจ ฯเปฯ สณฺฐานํ ปฏิจฺจ ฯเปฯ ปาริปูริํ ปฏิจฺจ มโท มชฺชนา มชฺชิตตฺตํ มาโน จ มญฺญนา มญฺญิตตฺตํ อุนฺนติ อุนฺนาโม ธโช สมฺปคฺคาโห เกตุกมฺยตา จิตฺตสฺส, อยํ วุจฺจติ ปาริปูริมโท.}

\proseitem{845}{ตตฺถ กตโม มโท, โย มโท มชฺชนา มชฺชิตตฺตํ มาโน มญฺญนา มญฺญิตตฺตํ อุนฺนติ อุนฺนาโม ธโช สมฺปคฺคาโห เกตุกมฺยตา จิตฺตสฺส, อยํ วุจฺจติ มโท.}

\proseitem{846}{ตตฺถ กตโม ปมาโท, กายทุจฺจริเต วา วจีทุจฺจริเต วา มโนทุจฺจริเต วา ปญฺจสุ วา กามคุเณสุ จิตฺตสฺส โวสฺสคฺโค โวสฺสคฺคานุปฺปทานํ กุสลานํ วา ธมฺมานํ ภาวนาย อสกฺกจฺจกิริยตา อสาตจฺจกิริยตา อนฏฺฐิตกิริยตา โอลีนวุตฺติตา นิกฺขิตฺตฉนฺทตา นิกฺขิตฺตธุรตา อนาเสวนา อภาวนา อพหุลีกมฺมํ อนธิฏฺฐานํ อนนุโยโค ปมาโท, โย เอวรูโป ปมาโท ปมชฺชนา ปมชฺชิตตฺตํ, อยํ วุจฺจติ ปมาโท.}

\proseitem{847}{ตตฺถ กตโม ถมฺโภ, โย ถมฺโภ ถมฺภนา ถมฺภิตตฺตํ กกฺขฬิยํ ผารุสิยํ อุชุจิตฺตตา อมุทุตา, อยํ วุจฺจติ ถมฺโภ.}

\proseitem{848}{\csromanfolio{364}ตตฺถ กตโม สารมฺโภ, โย สารมฺโภ ปฏิสารมฺโภ สารมฺภนา ปฏิสารมฺภนา ปฏิสารมฺภิตตฺตํ, อยํ วุจฺจติ สารมฺโภ.}

\proseitem{849}{ตตฺถ กตมา อตฺริจฺฉตา, อิตรีตรจีวรปิณฺฑปาตเสนาสนคิลานปจฺจยเภสชฺชปริกฺขาเรหิ ปญฺจหิ วา กามคุเณหิ อสนฺตุฏฺฐสฺส ภิยฺโยกมฺยตา, ยา เอวรูปา อิจฺฉา อิจฺฉาคตา อตฺริจฺฉตา ราโค สาราโค จิตฺตสฺส สาราโค, อยํ วุจฺจติ อตฺริจฺฉตา.}

\proseitem{850}{ตตฺถ กตมา มหิจฺฉตา, อิตรีตรจีวรปิณฺฑปาตเสนาสนคิลานปจฺจยเภสชฺชปริกฺขาเรหิ ปญฺจหิ วา กามคุเณหิ อสนฺตุฏฺฐสฺส ภิยฺโยกมฺยตา, ยา เอวรูปา อิจฺฉา อิจฺฉาคตา มหิจฺฉตา ราโค สาราโค จิตฺตสฺส สาราโค, อยํ วุจฺจติ มหิจฺฉตา.}

\proseitem{851}{ตตฺถ กตมา ปาปิจฺฉตา, อิเธกจฺโจ อสฺสทฺโธ สมาโน “สทฺโธ”ติ มํ ชโน ชานาตูติ อิจฺฉติ, ทุสฺสีโล สมาโน “สีลวา”ติ มํ ชโน ชานาตูติ อิจฺฉติ, อปฺปสฺสุโต สมาโน “พหุสฺสุโต”ติ มํ ชโน ชานาตูติ อิจฺฉติ, สงฺคณิการาโม สมาโน “ปวิวิตฺโต”ติ มํ ชโน ชานาตูติ อิจฺฉติ, กุสีโต สมาโน “อารทฺธวีริโย”ติ มํ ชโน ชานาตูติ อิจฺฉติ, มุฏฺฐสฺสตี สมาโน “อุปฏฺฐิตสฺสตี”ติ มํ ชโน ชานาตูติ อิจฺฉติ, อสมาหิโต สมาโน “สมาหิโต”ติ มํ ชโน ชานาตูติ อิจฺฉติ, ทุปฺปญฺโญ สมาโน “ปญฺญวา”ติ มํ ชโน ชานาตูติ อิจฺฉติ, อขีณาสโว สมาโน “ขีณาสโว”ติ มํ ชโน ชานาตูติ อิจฺฉติ, ยา เอวรูปา อิจฺฉา อิจฺฉาคตา ปาปิจฺฉตา ราโค สาราโค จิตฺตสฺส สาราโค, อยํ วุจฺจติ ปาปิจฺฉตา.}

\proseitem{852}{ตตฺถ กตมํ สิงฺคํ, ยํ สิงฺคํ สิงฺคารตา จาตุรตา จาตุริยํ ปริกฺขตฺตตา ปาริกฺขตฺติยํ, อิทํ วุจฺจติ สิงฺคํ.}

\proseitem{853}{ตตฺถ กตมํ ตินฺติณํ, ยํ ตินฺติณํ ตินฺติณายนา ตินฺติณายิตตฺตํ โลลุปฺปํ โลลุปฺปายนา โลลุปฺปายิตตฺตํ ปุจฺฉญฺชิกตา สาธุกมฺยตา, อิทํ วุจฺจติ ตินฺติณํ.}

\csromanheader{2}{17. ขุทฺทกวตฺถุวิภงฺค}
\proseitem{854}{\csromanfolio{365}ตตฺถ กตมํ จาปลฺยํ, จีวรมณฺฑนา ปตฺตมณฺฑนา เสนาสนมณฺฑนา อิมสฺส วา ปูติกายสฺส พาหิรานํ วา ปริกฺขารานํ มณฺฑนา วิภูสนา เกฬนา ปริเกฬนา คิทฺธิกตา คิทฺธิกตฺตํ จปลตา จาปลฺยํ, อิทํ วุจฺจติ จาปลฺยํ.}

\proseitem{855}{ตตฺถ กตมํ อสภาควุตฺติ, มาตริ วา ปิตริ วา เชฏฺเฐ วา ภาตริ วา อาจริเยสุ วา อุปชฺฌาเย วา พุทฺเธ วา สาวเกสุ วา อญฺญตรญฺญตเรสุ ครุฏฺฐานิเยสุ วิปฺปฏิกูลคฺคาหิตา วิปจฺจนีกสาตตา อนาทริยํ อนาทริยตา อคารวตา อปฺปติสฺสวตา, อยํ วุจฺจติ อสภาควุตฺติ.}

\proseitem{856}{ตตฺถ กตมา อรติ, ปนฺเตสุ วา เสนาสเนสุ อญฺญตรญฺญตเรสุ วา อธิกุสเลสุ ธมฺเมสุ อรติ อรติตา อนภิรติ อนภิรมณา อุกฺกณฺฐิตา ปริตสฺสิตา, อยํ วุจฺจติ อรติ.}

\proseitem{857}{ตตฺถ กตมา ตนฺที, ยา ตนฺที ตนฺทิยนา ตนฺทิมนกตา, อาลสฺยํ อาลสฺยายนา อาลสฺยายิตตฺตํ, อยํ วุจฺจติ ตนฺที.}

\proseitem{858}{ตตฺถ กตมา วิชมฺภิตา, ยา กายสฺส ชมฺภนา วิชมฺภนา อานมนา วินมนา สนฺนมนา ปณมนา พฺยาธิยกํ, อยํ วุจฺจติ วิชมฺภิตา.}

\proseitem{859}{ตตฺถ กตโม ภตฺตสมฺมโท, ยา ภุตฺตาวิสฺส ภตฺตมุจฺฉา ภตฺตกิลมโถ ภตฺตปริฬาโห กายทุฏฺฐุลฺลํ, อยํ วุจฺจติ ภตฺตสมฺมโท.}

\proseitem{860}{ตตฺถ กตมํ เจตโส จ ลีนตฺตํ, ยา จิตฺตสฺส อกลฺยตา อกมฺมญฺญตา โอลียนา สลฺลียนา ลีนํ ลียนา ลียิตตฺตํ ถินํ ถิยนา ถิยิตตฺตํ จิตฺตสฺส, อิทํ วุจฺจติ เจตโส จ ลีนตฺตํ.}

\proseitem{861}{ตตฺถ กตมา กุหนา, ลาภสกฺการสิโลกสนฺนิสฺสิตสฺส ปาปิจฺฉสฺส อิจฺฉาปกตสฺส ปจฺจยปฏิเสวนสงฺขาเตน\footnote{ปจฺจยปฏิเสธนสงฺขาเตน (สี)} วา สามนฺตชปฺปิเตน วา อิริยาปถสฺส วา อฐปนา\footnote{อาฐปนา (ก)} ฐปนา สณฺฐปนา ภากุฏิตา ภากุฏิยํ กุหนา กุหายนา กุหิตตฺตํ, อยํ วุจฺจติ กุหนา.}

\proseitem{862}{\csromanfolio{366}ตตฺถ กตมา ลปนา, ลาภสกฺการสิโลกสนฺนิสฺสิตสฺส ปาปิจฺฉสฺส อิจฺฉาปกตสฺส ยา ปเรสํ อาลปนา ลปนา สลฺลปนา อุลฺลปนา สมุลฺลปนา อุนฺนหนา สมุนฺนหนา อุกฺกาจนา สมุกฺกาจนา อนุปฺปิยภาณิตา จาฏุกมฺยตา มุคฺคสูปฺยตา ปาริภฏฺยตา, อยํ วุจฺจติ ลปนา.}

\proseitem{863}{ตตฺถ กตมา เนมิตฺติกตา, ลาภสกฺการสิโลกสนฺนิสฺสิตสฺส ปาปิจฺฉสฺส อิจฺฉาปกตสฺส ยํ ปเรสํ นิมิตฺตํ นิมิตฺตกมฺมํ โอภาโส โอภาสกมฺมํ สามนฺตชปฺปา ปริกถา, อยํ วุจฺจติ เนมิตฺติกตา.}

\proseitem{864}{ตตฺถ กตมา นิปฺเปสิกตา, ลาภสกฺการสิโลกสนฺนิสฺสิตสฺส ปาปิจฺฉสฺส อิจฺฉาปกตสฺส ยา ปเรสํ อกฺโกสนา วมฺภนา ครหณา อุกฺเขปนา สมุกฺเขปนา ขิปนา สํขิปนา ปาปนา สมฺปาปนา อวณฺณหาริกา ปรปิฏฺฐิมํสิกตา, อยํ วุจฺจติ นิปฺเปสิกตา.}

\proseitem{865}{ตตฺถ กตมา ลาเภน ลาภํ นิชิคีสนตา, ลาภสกฺการสิโลกสนฺนิสฺสิโต ปาปิจฺโฉ อิจฺฉาปกโต อิโต ลทฺธํ อามิสํ อมุตฺร หรติ, อมุตฺร วา ลทฺธํ อามิสํ อิธ อาหรติ, ยา เอวรูปา อามิสสฺส เอฏฺฐิ คเวฏฺฐิ ปริเยฏฺฐิ เอสนา คเวสนา ปริเยสนา, อยํ วุจฺจติ ลาเภน ลาภํ นิชิคีสนตา.}

\proseitem{866}{ตตฺถ กตโม “เสยฺโยหมสฺมี”ติ มาโน, อิเธกจฺโจ ชาติยา วา โคตฺเตน วา โกลปุตฺติเยน วา วณฺณโปกฺขรตาย วา ธเนน วา อชฺเฌเนน วา กมฺมายตเนน วา สิปฺปายตเนน วา วิชฺชาฏฺฐาเนน\footnote{วิชฺชฏฺฐาเนน (สฺยา)} วา สุเตน วา ปฏิภาเนน วา อญฺญตรญฺญตเรน วตฺถุนา มานํ ชปฺเปติ, โย เอวรูโป มาโน มญฺญนา มญฺญิตตฺตํ อุนฺนติ อุนฺนาโม ธโช สมฺปคฺคาโห เกตุกมฺยตา จิตฺตสฺส, อยํ วุจฺจติ “เสยฺโยหมสฺมี”ติ มาโน.}

\proseitem{867}{ตตฺถ กตโม “สทิโสหมสฺมี”ติ มาโน, อิเธกจฺโจ ชาติยา วา โคตฺเตน วา โกลปุตฺติเยน วา วณฺณโปกฺขรตาย วา ธเนน วา อชฺเฌเนน วา กมฺมายตเนน วา สิปฺปายตเนน วา วิชฺชาฏฺฐาเนน}

\vspace{\tipitakaparskip}
\csromancenter{ขุทฺทกวตฺถุวิภงฺค วา สุเตน วา ปฏิภาเนน วา อญฺญตรญฺญตเรน วตฺถุนา มานํ ชปฺเปติ, โย เอวรูโป มาโน มญฺญนา มญฺญิตตฺตํ อุนฺนติ อุนฺนาโม ธโช สมฺปคฺคาโห เกตุกมฺยตา จิตฺตสฺส, อยํ วุจฺจติ “สทิโสหมสฺมี”ติ มาโน.}
\proseitem{868}{\csromanfolio{367}ตตฺถ กตโม “หีโนหมสฺมี”ติ มาโน, อิเธกจฺโจ ชาติยา วา โคตฺเตน วา โกลปุตฺติเยน วา วณฺณโปกฺขรตาย วา ธเนน วา อชฺเฌเนน วา กมฺมายตเนน วา สิปฺปายตเนน วา วิชฺชาฏฺฐาเนน วา สุเตน วา ปฏิภาเนน วา อญฺญตรญฺญตเรน วตฺถุนา โอมานํ ชปฺเปติ, โย เอวรูโป โอมาโน โอมญฺญนา โอมญฺญิตตฺตํ หีฬนา โอหีฬนา โอหีฬิตตฺตํ อตฺตุญฺญา อตฺตวญฺญา อตฺตปริภโว, อยํ วุจฺจติ “หีโนหมสฺมี”ติ มาโน.}

\proseitem{869}{ตตฺถ กตโม เสยฺยสฺส “เสยฺโยหมสฺมี”ติ มาโน, อิเธกจฺโจ เสยฺโย โหติ ชาติยา วา โคตฺเตน วา โกลปุตฺติเยน วา วณฺณโปกฺขรตาย วา ธเนน วา อชฺเฌเนน วา กมฺมายตเนน วา สิปฺปายตเนน วา วิชฺชาฏฺฐาเนน วา สุเตน วา ปฏิภาเนน วา อญฺญตรญฺญตเรน วตฺถุนา ปเรหิ เสยฺยํ อตฺตานํ ทหติ, โส ตํ นิสฺสาย มานํ ชปฺเปติ, โย เอวรูโป มาโน มญฺญนา มญฺญิตตฺตํ อุนฺนติ อุนฺนาโม ธโช สมฺปคฺคาโห เกตุกมฺยตา จิตฺตสฺส, อยํ วุจฺจติ เสยฺยสฺส “เสยฺโยหมสฺมี”ติ มาโน.}

\proseitem{870}{ตตฺถ กตโม เสยฺยสฺส “สทิโสหมสฺมี”ติ มาโน, อิเธกจฺโจ เสยฺโย โหติ ชาติยา วา โคตฺเตน วา โกลปุตฺติเยน วา วณฺณโปกฺขรตาย วา ธเนน วา อชฺเฌเนน วา กมฺมายตเนน วา สิปฺปายตเนน วา วิชฺชาฏฺฐาเนน วา สุเตน วา ปฏิภาเนน วา อญฺญตรญฺญตเรน วตฺถุนา ปเรหิ สทิสํ อตฺตานํ ทหติ, โส ตํ นิสฺสาย มานํ ชปฺเปติ, โย เอวรูโป มาโน มญฺญนา มญฺญิตตฺตํ อุนฺนติ อุนฺนาโม ธโช สมฺปคฺคาโห เกตุกมฺยตา จิตฺตสฺส, อยํ วุจฺจติ เสยฺยสฺส “สทิโสหมสฺมี”ติ มาโน.}

\proseitem{871}{ตตฺถ กตโม เสยฺยสฺส “หีโนหมสฺมี”ติ มาโน, อิเธกจฺโจ เสยฺโย โหติ ชาติยา วา โคตฺเตน วา โกลปุตฺติเยน \csromanfolio{368}วา วณฺณโปกฺขรตาย วา ธเนน วา อชฺเฌเนน วา กมฺมายตเนน วา สิปฺปายตเนน วา วิชฺชาฏฺฐาเนน วา สุเตน วา ปฏิภาเนน วา อญฺญตรญฺญตเรน วตฺถุนา ปเรหิ หีนํ อตฺตานํ ทหติ, โส ตํ นิสฺสาย โอมานํ ชปฺเปติ, โย เอวรูโป โอมาโน โอมญฺญนา โอมญฺญิตตฺตํ หีฬนา โอหีฬนา โอหีฬิตตฺตํ อตฺตุญฺญา อตฺตวญฺญา อตฺตปริภโว, อยํ วุจฺจติ เสยฺยสฺส “หีโนหมสฺมี”ติ มาโน.}

\proseitem{872}{ตตฺถ กตโม สทิสสฺส “เสยฺโยหมสฺมี”ติ มาโน, อิเธกจฺโจ สทิโส โหติ ชาติยา วา โคตฺเตน วา โกลปุตฺติเยน วา ฯเปฯ อญฺญตรญฺญตเรน วตฺถุนา ปเรหิ เสยฺยํ อตฺตานํ ทหติ, โส ตํ นิสฺสาย มานํ ชปฺเปติ, โย เอวรูโป มาโน มญฺญนา มญฺญิตตฺตํ ฯเปฯ เกตุกมฺยตา จิตฺตสฺส, อยํ วุจฺจติ สทิสสฺส “เสยฺโยหมสฺมี”ติ มาโน.}

\proseitem{873}{ตตฺถ กตโม สทิสสฺส “สทิโสหมสฺมี”ติ มาโน, อิเธกจฺโจ สทิโส โหติ ชาติยา วา โคตฺเตน วา โกลปุตฺติเยน วา ฯเปฯ อญฺญตรญฺญตเรน วตฺถุนา ปเรหิ สทิสํ อตฺตานํ ทหติ, โส ตํ นิสฺสาย มานํ ชปฺเปติ, โย เอวรูโป มาโน มญฺญนา มญฺญิตตฺตํ อุนฺนติ อุนฺนาโม ธโช สมฺปคฺคาโห เกตุกมฺยตา จิตฺตสฺส, อยํ วุจฺจติ สทิสสฺส “สทิโสหมสฺมี”ติ มาโน.}

\proseitem{874}{ตตฺถ กตโม สทิสสฺส “หีโนหมสฺมี”ติ มาโน.\csromanspacer{} อิเธกจฺโจ สทิโส โหติ ชาติยา วา โคตฺเตน วา โกลปุตฺติเยน วา ฯเปฯ อญฺญตรญฺญตเรน วตฺถุนา ปเรหิ หีนํ อตฺตานํ ทหติ, โส ตํ นิสฺสาย โอมานํ ชปฺเปติ, โย เอวรูโป โอมาโน โอมญฺญนา โอมญฺญิตตฺตํ หีฬนา โอหีฬนา โอหีฬิตตฺตํ อตฺตุญฺญา อตฺตวญฺญา อตฺตปริภโว, อยํ วุจฺจติ สทิสสฺส “หีโนหมสฺมี”ติ มาโน.}

\proseitem{875}{ตตฺถ กตโม หีนสฺส “เสยฺโยหมสฺมี”ติ มาโน, อิเธกจฺโจ หีโน โหติ ชาติยา วา โคตฺเตน วา โกลปุตฺติเยน วา ฯเปฯ อญฺญตรญฺญตเรน วตฺถุนา ปเรหิ เสยฺยํ อตฺตานํ ทหติ, โส ตํ นิสฺสาย มานํ ชปฺเปติ, โย เอวรูโป มาโน มญฺญนา มญฺญิตตฺตํ อุนฺนติ, อุนฺนาโม ธโช สมฺปคฺคาโห เกตุกมฺยตา จิตฺตสฺส, อยํ วุจฺจติ หีนสฺส “เสยฺโยหมสฺมี”ติ มาโน.}

\csromanheader{2}{17. ขุทฺทกวตฺถุวิภงฺค}
\proseitem{876}{\csromanfolio{369}ตตฺถ กตโม หีนสฺส “สทิโสหมสฺมี”ติ มาโน, อิเธกจฺโจ หีโน โหติ ชาติยา วา โคตฺเตน วา โกลปุตฺติเยน วา ฯเปฯ อญฺญตรญฺญตเรน วตฺถุนา ปเรหิ สทิสํ อตฺตานํ ทหติ, โส ตํ นิสฺสาย มานํ ชปฺเปติ, โย เอวรูโป มาโน มญฺญนา มญฺญิตตฺตํ อุนฺนติ อุนฺนาโม ธโช สมฺปคฺคาโห เกตุกมฺยตา จิตฺตสฺส, อยํ วุจฺจติ หีนสฺส “สทิโสหมสฺมี”ติ มาโน.}

\proseitem{877}{ตตฺถ กตโม หีนสฺส “หีโนหมสฺมี”ติ มาโน, อิเธกจฺโจ หีโน โหติ ชาติยา วา โคตฺเตน วา โกลปุตฺติเยน วา ฯเปฯ อญฺญตรญฺญตเรน วตฺถุนา ปเรหิ สทิสํ อตฺตานํ ทหติ, โส ตํ นิสฺสาย โอมานํ ชปฺเปติ, โย เอวรูโป โอมาโน โอมญฺญนา โอมญฺญิตตฺตํ หีฬนา โอหีฬนา โอหีฬิตตฺตํ อตฺตุญฺญา อตฺตวญฺญา อตฺตปริภโว, อยํ วุจฺจติ หีนสฺส “หีโนหมสฺมี”ติ มาโน.}

\proseitem{878}{ตตฺถ กตโม มาโน, โย มาโน มญฺญนา มญฺญิตตฺตํ อุนฺนติ อุนฺนาโม ธโช สมฺปคฺคาโห เกตุกมฺยตา จิตฺตสฺส, อยํ วุจฺจติ มาโน.}

\proseitem{879}{ตตฺถ กตโม อติมาโน, อิเธกจฺโจ ชาติยา วา โคตฺเตน วา โกลปุตฺติเยน วา ฯเปฯ อญฺญตรญฺญตเรน วตฺถุนา ปเรหิ อตฺตานํ อติมญฺญติ, โย เอวรูโป มาโน มญฺญนา มญฺญิตตฺตํ อุนฺนติ อุนฺนาโม ธโช สมฺปคฺคาโห เกตุกมฺยตา จิตฺตสฺส, อยํ วุจฺจติ อติมาโน.}

\proseitem{880}{ตตฺถ กตโม มานาติมาโน, อิเธกจฺโจ ชาติยา วา โคตฺเตน วา โกลปุตฺติเยน วา ฯเปฯ อญฺญตรญฺญตเรน วตฺถุนา ปุพฺพกาลํ ปเรหิ สทิสํ อตฺตานํ ทหติ, อปรกาลํ อตฺตานํ เสยฺยํ ทหติ, โย เอวรูโป มาโน มญฺญนา มญฺญิตตฺตํ อุนฺนติ อุนฺนาโม ธโช สมฺปคฺคาโห เกตุกมฺยตา จิตฺตสฺส, อยํ วุจฺจติ มานาติมาโน.}

\proseitem{881}{ตตฺถ กตโม โอมาโน, อิเธกจฺโจ ชาติยา วา โคตฺเตน วา โกลปุตฺติเยน วา วณฺณโปกฺขรตาย วา ธเนน วา อชฺเฌเนน วา \csromanfolio{370}กมฺมายตเนน วา สิปฺปายตเนน วา วิชฺชาฏฺฐาเนน วา สุเตน วา ปฏิภาเนน วา อญฺญตรญฺญตเรน วตฺถุนา โอมานํ ชปฺเปติ, โย เอวรูโป โอมาโน โอมญฺญนา โอมญฺญิตตฺตํ หีฬนา โอหีฬนา โอหีฬิตตฺตํ อตฺตุญฺญา อตฺตวญฺญา อตฺตปริภโว, อยํ วุจฺจติ โอมาโน.}

\proseitem{882}{ตตฺถ กตโม อธิมาโน, อปฺปตฺเต ปตฺตสญฺญิตา, อกเต กตสญฺญิตา, อนธิคเต อธิคตสญฺญิตา, อสจฺฉิกเต สจฺฉิกตสญฺญิตา, โย เอวรูโป มาโน มญฺญนา มญฺญิตตฺตํ อุนฺนติ อุนฺนาโม ธโช สมฺปคฺคาโห เกตุกมฺยตา จิตฺตสฺส, อยํ วุจฺจติ อธิมาโน.}

\proseitem{883}{ตตฺถ กตโม อสฺมิมาโน, รูปํ “อสฺมี”ติ มาโน “อสฺมี”ติ ฉนฺโท “อสฺมี”ติ อนุสโย เวทนา ฯเปฯ สญฺญา ฯเปฯ สงฺขารา ฯเปฯ วิญฺญาณํ “อสฺมี”ติ มาโน “อสฺมี”ติ ฉนฺโท “อสฺมี”ติ อนุสโย, โย เอวรูโป มาโน มญฺญนา มญฺญิตตฺตํ อุนฺนติ อุนฺนาโม ธโช สมฺปคฺคาโห เกตุกมฺยตา จิตฺตสฺส, อยํ วุจฺจติ อสฺมิมาโน.}

\proseitem{884}{ตตฺถ กตโม มิจฺฉามาโน, อิเธกจฺโจ ปาปเกน วา กมฺมายตเนน ปาปเกน วา สิปฺปายตเนน ปาปเกน วา วิชฺชาฏฺฐาเนน ปาปเกน วา สุเตน ปาปเกน วา ปฏิภาเนน ปาปเกน วา สีเลน ปาปเกน วา วเตน ปาปเกน วา สีลพฺพเตน ปาปิกาย วา ทิฏฺฐิยา อญฺญตรญฺญตเรน วตฺถุนา มานํ ชปฺเปติ, โย เอวรูโป มาโน มญฺญนา มญฺญิตตฺตํ อุนฺนติ อุนฺนาโม ธโช สมฺปคฺคาโห เกตุกมฺยตา จิตฺตสฺส, อยํ วุจฺจติ มิจฺฉามาโน.}

\proseitem{885}{ตตฺถ กตโม ญาติวิตกฺโก, ญาตเก อารพฺภ เคหสิโต ตกฺโก วิตกฺโก มิจฺฉาสงฺกปฺโป, อยํ วุจฺจติ ญาติวิตกฺโก.}

\proseitem{886}{ตตฺถ กตโม ชนปทวิตกฺโก, ชนปทํ อารพฺภ เคหสิโต ตกฺโก วิตกฺโก มิจฺฉาสงฺกปฺโป, อยํ วุจฺจติ ชนปทวิตกฺโก.}

\proseitem{887}{ตตฺถ กตโม อมรวิตกฺโก, ทุกฺกรการิกาปฏิสํยุตฺโต วา ทิฏฺฐิคตปฏิสํยุตฺโต วา เคหสิโต ตกฺโก วิตกฺโก มิจฺฉาสงฺกปฺโป, อยํ วุจฺจติ อมรวิตกฺโก.}

\csromanheader{2}{17. ขุทฺทกวตฺถุวิภงฺค}
\proseitem{888}{\csromanfolio{371}ตตฺถ กตโม ปรานุทฺทยตาปฏิสํยุตฺโต วิตกฺโก, อิเธกจฺโจ คิหีหิ สํสฏฺโฐ วิหรติ สหนนฺที สหโสกี สุขิเตสุ สุขิโต ทุกฺขิเตสุ ทุกฺขิโต อุปฺปนฺเนสุ กิจฺจกรณีเยสุ อตฺตนา วา โยคํ อาปชฺชติ, โย ตตฺต เคหสิโต ตกฺโก วิตกฺโก มิจฺฉาสงฺกปฺโป, อยํ วุจฺจติ ปรานุทฺทยตาปฏิสํยุตฺโต วิตกฺโก.}

\proseitem{889}{ตตฺถ กตโม ลาภสกฺการสิโลกปฏิสํยุตฺโต วิตกฺโก, ลาภสกฺการสิโลกํ อารพฺภ เคหสิโต ตกฺโก วิตกฺโก มิจฺฉาสงฺกปฺโป, อยํ วุจฺจติ ลาภสกฺการสิโลกปฏิสํยุตฺโต วิตกฺโก.}

\proseitem{890}{ตตฺถ กตโม อนวญฺญตฺติปฏิสํยุตฺโต วิตกฺโก, อิเธกจฺโจ ชาติยา วา โคตฺเตน วา โกลปุตฺติเยน วา วณฺณโปกฺขรตาย วา ธเนน วา อชฺเฌเนน วา กมฺมายตเนน วา สิปฺปายตเนน วา วิชฺชาฏฺฐาเนน วา สุเตน วา ปฏิภาเนน วา อญฺญตรญฺญตเรน วตฺถุนา “มา มํ ปเร อวชานิํสู”ติ โย ตตฺถ เคหสิโต ตกฺโก วิตกฺโก มิจฺฉาสงฺกปฺโป, อยํ วุจฺจติ อนวญฺญตฺติปฏิสํยุตฺโต วิตกฺโก.}

\csromanheader{2}{เอกกํ.}
\csromansectionrule
\csromanmatikaanchor{mtk.180}
\csromantocmarkref{subparagraph}{2. ทุกนิทฺเทส}{mtk.180}
\bookmark[dest={mtk.180},level=5]{2. ทุกนิทฺเทส}
\csromanheader{6}{2. ทุกนิทฺเทส}
\proseitem{891}{ตตฺถ กตโม โกโธ, โย โกโธ กุชฺฌนา กุชฺฌิตตฺตํ โทโส ทุสฺสนา ทุสฺสิตตฺตํ พฺยาปตฺติ พฺยาปชฺชนา พฺยาปชฺชิตตฺตํ วิโรโธ ปฏิวิโรโธ จณฺฑิกฺกํ อสุโรโป อนตฺตมนตา จิตฺตสฺส, อยํ วุจฺจติ โกโธ.}

\prose{ตตฺถ กตโม อุปนาโห, ปุพฺพกาลํ โกโธ อปรกาลํ อุปนาโห, โย เอวรูโป อุปนาโห อุปนยฺหนา อุปนยฺหิตตฺตํ อฏฺฐปนา ฐปนา สณฺฐปนา อนุสํสนฺทนา อนุปฺปพนฺธนา ทฬฺหีกมฺมํ โกธสฺส, อยํ วุจฺจติ อุปนาโห. (1)}

\proseitem{892}{ตตฺถ กตโม มกฺโข, โย มกฺโข มกฺขายนา มกฺขายิตตฺตํ\footnote{มกฺขิยนา มกฺขิยิตตฺตํ (สี, ก)} นิฏฺฐุริยํ นิฏฺฐุริยกมฺมํ, อยํ วุจฺจติ มกฺโข.}

\prose{\csromanfolio{372}ตตฺถ กตโม ปฬาโส, โย ปฬาโส ปฬาสายนา\footnote{ปลาสายนา ปลาสายิตตฺตํ (สฺยา) 2. ปคฺคหิตตฺตํ (สี, ก) อภิ 1. 226 ปิฏฺเฐปิ.} ปฬาสาหาโร วิวาทฏฺฐานํ ยุคคฺคาโห อปฺปฏินิสฺสคฺโค, อยํ วุจฺจติ ปฬาโส. (2)}

\proseitem{893}{ตตฺถ กตมา อิสฺสา, ยา ปรลาภสกฺการครุการมานนวนฺทนปูชนาสุ อิสฺสา อิสฺสายนา อิสฺสายิตตฺตํ อุสูยา อุสูยนา อุสูยิตตฺตํ, อยํ วุจฺจติ อิสฺสา.}

\prose{ตตฺถ กตมํ มจฺฉริยํ, ปญฺจ มจฺฉริยานิ อาวาสมจฺฉริยํ กุลมจฺฉริยํ ลาภมจฺฉริยํ วณฺณมจฺฉริยํ ธมฺมมจฺฉริยํ.\csromanspacer{} ยํ เอวรูปํ มจฺเฉรํ มจฺฉรายนา มจฺฉรายิตตฺตํ เววิจฺฉํ กทริยํ กฏุกญฺจุกตา อคฺคหิตตฺตํ2 จิตฺตสฺส, อิทํ วุจฺจติ มจฺฉริยํ. (3)}

\proseitem{894}{ตตฺถ กตมา มายา, อิเธกจฺโจ กาเยน ทุจฺจริตํ จริตฺวา วาจาย ทุจฺจริตํ จริตฺวา มนสา ทุจฺจริตํ จริตฺวา ตสฺส ปฏิจฺฉาทนเหตุํ ปาปิกํ อิจฺฉํ ปณิทหติ, “มา มํ ชญฺญา”ติ อิจฺฉติ, “มา มํ ชญฺญา”ติ สงฺกปฺเปติ, “มา มํ ชญฺญา”ติ วาจํ ภาสติ, “มา มํ ชญฺญา”ติ กาเยน ปรกฺกมติ, ยา เอวรูปา มายา มายาวิตา อจฺจาสรา วญฺจนา นิกติ วิกิรณา ปริหรณา คูหนา ปริคูหนา ฉาทนา ปฏิจฺฉาทนา อนุตฺตานีกมฺมํ อนาวิกมฺมํ โวจฺฉาทนา ปาปกิริยา, อยํ วุจฺจติ มายา.}

\prose{ตตฺถ กตมํ สาเฐยฺยํ, อิเธกจฺโจ สโฐ โหติ ปริสโฐ, ยํ ตตฺถ สฐํ สฐตา สาเฐยฺยํ กกฺกรตา กกฺกริยํ\footnote{กกฺขฬตา กกฺขฬิยํ (สฺยา)} ปริกฺขตฺตตา ปาริกฺขตฺติยํ, อิทํ วุจฺจติ สาเฐยฺยํ. (4)}

\proseitem{895}{ตตฺถ กตมา อวิชฺชา, ยํ อญฺญาณํ อทสฺสนํ ฯเปฯ อวิชฺชาลงฺคี โมโห อกุสลมูลํ, อยํ วุจฺจติ อวิชฺชา.}

\prose{ตตฺถ กตมา ภวตณฺหา, โย ภเวสุ ภวจฺฉนฺโท ภวราโค ภวนนฺที ภวตณฺหา ภวสิเนโห ภวปริฬาโห ภวมุจฺฉา ภวชฺโฌสานํ, อยํ วุจฺจติ ภวตณฺหา. (5)}

\csromanheader{2}{17. ขุทฺทกวตฺถุวิภงฺค}
\proseitem{896}{\csromanfolio{373}ตตฺถ กตมา ภวทิฏฺฐิ, “ภวิสฺสติ อตฺตา จ โลโก จา”ติ ยา เอวรูปา ทิฏฺฐิ ทิฏฺฐิคตํ ฯเปฯ วิปริยาสคฺคาโห, อยํ วุจฺจติ ภวทิฏฺฐิ.}

\prose{ตตฺถ กตมา วิภวทิฏฺฐิ, “น ภวิสฺสติ อตฺตา จ โลโก จา”ติ ยา เอวรูปา ทิฏฺฐิ ทิฏฺฐิคตํ ฯเปฯ วิปริยาสคฺคาโห, อยํ วุจฺจติ วิภวทิฏฺฐิ. (6)}

\proseitem{897}{ตตฺถ กตมา สสฺสตทิฏฺฐิ, “สสฺสโต อตฺตา จ โลโก จา”ติ ยา เอวรูปา ทิฏฺฐิ ทิฏฺฐิคตํ ฯเปฯ วิปริยาสคฺคาโห, อยํ วุจฺจติ สสฺสตทิฏฺฐิ.}

\prose{ตตฺถ กตมา อุจฺเฉททิฏฺฐิ, “อุจฺฉิชฺชิสฺสติ อตฺตา จ โลโก จา”ติ ยา เอวรูปา ทิฏฺฐิ ทิฏฺฐิคตํ ฯเปฯ วิปริยาสคฺคาโห, อยํ วุจฺจติ อุจฺเฉททิฏฺฐิ. (7)}

\proseitem{898}{ตตฺถ กตมา อนฺตวาทิฏฺฐิ, “อนฺตวา อตฺตา จ โลโก จา”ติ ยา เอวรูปา ทิฏฺฐิ ทิฏฺฐิคตํ ฯเปฯ วิปริยาสคฺคาโห, อยํ วุจฺจติ อนฺตวาทิฏฺฐิ.}

\prose{ตตฺถ กตมา อนนฺตวาทิฏฺฐิ, “อนนฺตวา อตฺตา จ โลโก จา”ติ ยา เอวรูปา ทิฏฺฐิ ทิฏฺฐิคตํ ฯเปฯ วิปริยาสคฺคาโห, อยํ วุจฺจติ อนนฺตวาทิฏฺฐิ. (8)}

\proseitem{899}{ตตฺถ กตมา ปุพฺพนฺตานุทิฏฺฐิ, ปุพฺพนฺตํ อารพฺภ ยา อุปฺปชฺชติ ทิฏฺฐิ ทิฏฺฐิคตํ ฯเปฯ วิปริยาสคฺคาโห, อยํ วุจฺจติ ปุพฺพนฺตานุทิฏฺฐิ.}

\prose{ตตฺถ กตมา อปรนฺตานุทิฏฺฐิ, อปรนฺตํ อารพฺภ ยา อุปฺปชฺชติ ทิฏฺฐิ ทิฏฺฐิคตํ ฯเปฯ วิปริยาสคฺคาโห, อยํ วุจฺจติ อปรนฺตานุทิฏฺฐิ. (9)}

\proseitem{900}{ตตฺถ กตมํ อหิริกํ, ยํ น หิรียติ หิริยิตพฺเพน น หิรียติ ปาปกานํ อกุสลานํ ธมฺมานํ สมาปตฺติยา, อิทํ วุจฺจติ อหิริกํ.}

\prose{ตตฺถ กตมํ อโนตฺตปฺปํ, ยํ น โอตฺตปฺปติ โอตฺตปฺปิตพฺเพน น โอตฺตปฺปติ ปาปกานํ อกุสลานํ ธมฺมานํ สมาปตฺติยา, อิทํ วุจฺจติ อโนตฺตปฺปํ. (10)}

\proseitem{901}{ตตฺถ กตมา โทวจสฺสตา, สหธมฺมิเก วุจฺจมาเน โทวจสฺสายํ โทวจสฺสิยํ โทวจสฺสตา วิปฺปฏิกุลคฺคาหิตา วิปจฺจนีกสาตตา อนาทริยํ อนาทรตา อคารวตา อปฺปติสฺสวตา, อยํ วุจฺจติ โทวจสฺสตา.}

\prose{\csromanfolio{374}ตตฺถ กตมา ปาปมิตฺตตา, เย เต ปุคฺคลา อสฺสทฺธา ทุสฺสีลา อปฺปสฺสุตา มจฺฉริโน ทุปฺปญฺญา.\csromanspacer{} ยา เตสํ เสวนา นิเสวนา สํเสวนา ภชนา สมฺภชนา ภตฺติ สมฺภตฺติ สมฺปวงฺกตา, อยํ วุจฺจติ ปาปมิตฺตตา. (11)}

\proseitem{902}{ตตฺถ กตโม อนชฺชโว, โย อนชฺชโว อนชฺชวตา ชิมฺหตา วงฺกตา กุฏิลตา, อยํ วุจฺจติ อนชฺชโว.}

\prose{ตตฺถ กตโม อมทฺทโว, ยา อมุทุตา อมทฺทวตา กกฺขฬิยํ ผารุสิยํ กกฺขฬตา กฐินตา\footnote{กถินตา (สฺยา, ก) อภิ 1. 206 ปิฏฺเฐปิ.} อุชุจิตฺตตา อมุทุตา, อยํ วุจฺจติ อมทฺทโว. (12)}

\proseitem{903}{ตตฺถ กตมา อกฺขนฺติ, ยา อกฺขนฺติ อกฺขมนตา อนธิวาสนตา จณฺฑิกฺกํ อสุโรโป อนตฺตมนตา จิตฺตสฺส, อยํ วุจฺจติ อกฺขนฺติ.}

\prose{ตตฺถ กตมํ อโสรจฺจํ, กายิโก วีติกฺกโม วาจสิโก วีติกฺกโม กายิกวาจสิโก วีติกฺกโม, อิทํ วุจฺจติ อโสรจฺจํ.\csromanspacer{} สพฺพมฺปิ ทุสฺสีลฺยํ อโสรจฺจํ. (13)}

\proseitem{904}{ตตฺถ กตมํ อสาขลฺยํ, ยา สา วาจา กณฺฑกา กกฺกสา ปรกฏุกา ปราภิสชฺชนี โกธสามนฺตา อสมาธิสํวตฺตนิกา, ตถารูปิํ วาจํ ภาสิตา โหติ.\csromanspacer{} ยา ตตฺถ อสณฺหวาจตา อสขิลวาจตา ผรุสวาจตา, อิทํ วุจฺจติ อสาขลฺยํ.}

\prose{ตตฺถ กตโม อปฺปฏิสนฺถาโร\footnote{อปฺปฏิสนฺธาโร (ก)}, เทฺว ปฏิสนฺถารา อามิสปฏิสนฺถาโร จ ธมฺมปฏิสนฺถาโร จา.\csromanspacer{} อิเธกจฺโจ อปฺปฏิสนฺถารโก โหติ อามิสปฏิสนฺถาเรน วา ธมฺมปฏิสนฺถาเรน วา, อยํ วุจฺจติ อปฺปฏิสนฺถาโร. (14)}

\proseitem{905}{ตตฺถ กตมา อินฺทฺริเยสุ อคุตฺตทฺวารตา, อิเธกจฺโจ จกฺขุนา รูปํ ทิสฺวา นิมิตฺตคฺคาหี โหติ อนุพฺยญฺชนคฺคาหี, ยตฺวาธิกรณเมนํ จกฺขุนฺทฺริยํ อสํวุตํ วิหรนฺตํ อภิชฺฌาโทมนสฺสา ปาปกา อกุสลา ธมฺมา อนฺวาสฺสเวยฺยุํ, ตสฺส สํวราย น ปฏิปชฺชติ, น รกฺขติ จกฺขุนฺทฺริยํ, จกฺขุนฺทฺริเย น สํวรํ อาปชฺชติ, โสเตน สทฺทํ สุตฺวา ฯเปฯ ฆาเนน คนฺธํ}

\vspace{\tipitakaparskip}
\csromancenter{ขุทฺทกวตฺถุวิภงฺค ฆายิตฺวา ฯเปฯ ชิวฺหาย รสํ สายิตฺวา ฯเปฯ กาเยน โผฏฺฐพฺพํ ผุสิตฺวา ฯเปฯ มนสา ธมฺมํ วิญฺญาย นิมิตฺตคฺคาหี โหติ อนุพฺยญฺชนคฺคาหี, ยตฺวาธิกรณเมนํ มนินฺทฺริยํ อสํวุตํ วิหรนฺตํ อภิชฺฌาโทมนสฺสา ปาปกา อกุสลา ธมฺมา อนฺวาสฺสเวยฺยุํ, ตสฺส สํวราย น ปฏิปชฺชติ, น รกฺขติ มนินฺทฺริยํ, มนินฺทฺริเย น สํวรํ อาปชฺชติ.\csromanspacer{} ยา อิเมสํ ฉนฺนํ อินฺทฺริยานํ อคุตฺติ อโคปนา อนารกฺโข อสํวโร, อยํ วุจฺจติ อินฺทฺริเยสุ อคุตฺตทฺวารตา.}
\prose{\csromanfolio{375}ตตฺถ กตมา โภชเน อมตฺตญฺญุตา, อิเธกจฺโจ อปฺปฏิสงฺขา อโยนิโส อาหารํ อาหาเรติ ทวาย มทาย มณฺฑนาย วิภูสนาย, ยา ตตฺถ อสนฺตุฏฺฐิตา อมตฺตญฺญุตา อปฺปฏิสงฺขา โภชเน, อยํ วุจฺจติ โภชเน อมตฺตญฺญุตา. (15)}

\proseitem{906}{ตตฺถ กตมํ มุฏฺฐสฺสจฺจํ, ยา อสฺสติ อนนุสฺสติ อปฺปฏิสฺสติ อสฺสติ อสฺสรณตา อธารณตา ปิลาปนตา สมฺมุสนตา, อิทํ วุจฺจติ มุฏฺฐสฺสจฺจํ.}

\prose{ตตฺถ กตมํ อสมฺปชญฺญํ, ยํ อญฺญาณํ อทสฺสนํ ฯเปฯ อวิชฺชาลงฺคี โมโห อกุสลมูลํ, อิทํ วุจฺจติ อสมฺปชญฺญํ. (16)}

\proseitem{907}{ตตฺถ กตมา สีลวิปตฺติ, โย กายิโก วีติกฺกโม วาจสิโก วีติกฺกโม กายิกวาจสิโก วีติกฺกโม, อยํ วุจฺจติ สีลวิปตฺติ, สพฺพมฺปิ ทุสฺสีลฺยํ สีลวิปตฺติ.}

\prose{ตตฺถ กตมา ทิฏฺฐิวิปตฺติ, “นตฺถิ ทินฺนํ, นตฺถิ ยิฏฺฐํ ฯเปฯ เย อิมญฺจ โลกํ ปรญฺจ โลกํ สยํ อภิญฺญา สจฺฉิกตฺวา ปเวเทนฺตี”ติ ยา เอวรูปา ทิฏฺฐิ ทิฏฺฐิคตํ ฯเปฯ วิปริยาสคฺคาโห, อยํ วุจฺจติ ทิฏฺฐิวิปตฺติ.\csromanspacer{} สพฺพาปิ มิจฺฉาทิฏฺฐิ ทิฏฺฐิวิปตฺติ. (17)}

\proseitem{908}{ตตฺถ กตมํ อชฺฌตฺตสํโยชนํ, ปญฺโจรมฺภาคิยานิ สํโยชนานิ อชฺฌตฺตสํโยชนํ, ปญฺจุทฺธมฺภาคิยานิ สํโยชนานิ พหิทฺธาสํโยชนํ. (18)}

\vspace{\tipitakaparskip}
\csromancenter{ทุกํ.}
\csromanmatikaanchor{mtk.181}
\csromantocmarkref{subparagraph}{3. ติกนิทฺเทส}{mtk.181}
\bookmark[dest={mtk.181},level=5]{3. ติกนิทฺเทส}
\csromanheader{6}{3. ติกนิทฺเทส}
\proseitem{909}{\csromanfolio{376}ตตฺถ กตมานิ ตีณิ อกุสลมูลานิ, โลโภ โทโส โมโห.}

\prose{ตตฺถ กตโม โลโภ, โย ราโค สาราโค อนุนโย อนุโรโธ นนฺที นนฺทิราโค จิตฺตสฺส สาราโค อิจฺฉา มุจฺฉา อชฺโฌสานํ เคโธ ปริเคโธ สงฺโค ปงฺโก เอชา มายา ชนิกา สญฺชนนี สิพฺพินี ชาลินี สริตา วิสตฺติกา โสตํ วิสฏา\footnote{วิสทา (สี, ก) อภิ 1. 229 ปิฏฺเฐปิ.} อายูหนี\footnote{อายูหินี (ก)} ทุติยา ปณิธิ ภวเนตฺติ วนํ วนโถ สนฺถโว\footnote{สนฺธโว (ก)} สิเนโห อเปกฺขา ปฏิพนฺธุ อาสา อาสีสนา อาสีสิตตฺตํ\footnote{อาสิํสนา อาสิํสิตตฺตํ (สี, สฺยา)} รูปาสา สทฺทาสา คนฺธาสา รสาสา โผฏฺฐพฺพาสา ลาภาสา ธนาสา ปุตฺตาสา ชีวิตาสา ชปฺปา อภิชปฺปา ชปฺปนา ชปฺปิตตฺตํ โลลุปฺปํ โลลุปฺปายนา โลลุปฺปายิตตฺตํ ปุจฺฉญฺชิกตา สาธุกมฺยตา อธมฺมราโค วิสมโลโภ นิกนฺติ นิกามนา ปตฺถนา ปิหนา สมฺปตฺถนา กามตณฺหา ภวตณฺหา วิภวตณฺหา รูปตณฺหา อรูปตณฺหา นิโรธตณฺหา สทฺทตณฺหา คนฺธตณฺหา รสตณฺหา โผฏฺฐพฺพตณฺหา ธมฺมตณฺหา โอโฆ โยโค คนฺโถ อุปาทานํ อาวรณํ นีวรณํ ฉทนํ พนฺธนํ อุปกฺกิเลโส อนุสโย ปริยุฏฺฐานํ ลตา เววิจฺฉํ ทุกฺขมูลํ ทุกฺขนิทานํ ทุกฺขปฺปภโว มารปาโส มารพฬิสํ มารวิสโย ตณฺหานที ตณฺหาชาลํ ตณฺหาคทฺทุลํ ตณฺหาสมุทฺโท อภิชฺฌา โลโภ อกุสลมูลํ, อยํ วุจฺจติ โลโภ.}

\prose{ตตฺถ กตโม โทโส, “อนตฺถํ เม อจรี”ติ อาฆาโต ชายติ, “อนตฺถํ เม จรตี”ติ อาฆาโต ชายติ, “อนตฺถํ เม จริสฺสตี”ติ อาฆาโต ชายติ, “ปิยสฺส เม มนาปสฺส อนตฺถํ อจริ ฯเปฯ อนตฺถํ จรติ ฯเปฯ อนตฺถํ จริสฺสตี”ติ อาฆาโต ชายติ, “อปฺปิยสฺส เม อมนาปสฺส อตฺถํ อจริ ฯเปฯ อตฺถํ จรติ ฯเปฯ อตฺถํ จริสฺสตี”ติ อาฆาโต ชายติ, อฏฺฐาเน วา ปน อาฆาโต ชายติ.\csromanspacer{} โย เอวรูโป จิตฺตสฺส อาฆาโต ปฏิฆาโต ปฏิฆํ ปฏิวิโรโธ โกโป ปโกโป สมฺปโกโป โทโส ปโทโส สมฺปโทโส จิตฺตสฺส พฺยาปตฺติ มโนปโทโส โกโธ กุชฺฌนา กุชฺฌิตตฺตํ โทโส ทุสฺสนา}

\vspace{\tipitakaparskip}
\csromancenter{ขุทฺทกวตฺถุวิภงฺค ทุสฺสิตตฺตํ พฺยาปตฺติ พฺยาปชฺชนา พฺยาปชฺชิตตฺตํ วิโรโธ ปฏิวิโรโธ จณฺฑิกฺกํ อสุโรโป อนตฺตมนตา จิตฺตสฺส, อยํ วุจฺจติ โทโส.}
\prose{\csromanfolio{377}ตตฺถ กตโม โมโห, ทุกฺเข อญฺญาณํ, ทุกฺขสมุทเย อญฺญาณํ, ทุกฺขนิโรเธ อญฺญาณํ, ทุกฺขนิโรธคามินิยา ปฏิปทาย อญฺญาณํ, ปุพฺพนฺเต อญฺญาณํ, อปรนฺเต อญฺญาณํ, ปุพฺพนฺตาปรนฺเต อญฺญาณํ, อิทปฺปจฺจยตา ปฏิจฺจสมุปฺปนฺเนสุ ธมฺเมสุ อญฺญาณํ, ยํ เอวรูปํ อญฺญาณํ อทสฺสนํ ฯเปฯ อวิชฺชาลงฺคี โมโห อกุสลมูลํ, อยํ วุจฺจติ โมโห.\csromanspacer{} อิมานิ ตีณิ อกุสลมูลานิ. (1)}

\proseitem{910}{ตตฺถ กตเม ตโย อกุสลวิตกฺกา, กามวิตกฺโก พฺยาปาทวิตกฺโก วิหิํสาวิตกฺโก.}

\prose{ตตฺถ กตโม กามวิตกฺโก, กามปฏิสํยุตฺโต ตกฺโก วิตกฺโก มิจฺฉาสงฺกปฺโป, อยํ วุจฺจติ กามวิตกฺโก.}

\prose{ตตฺถ กตโม พฺยาปาทวิตกฺโก, พฺยาปาทปฏิสํยุตฺโต ตกฺโก วิตกฺโก มิจฺฉาสงฺกปฺโป, อยํ วุจฺจติ พฺยาปาทวิตกฺโก.}

\prose{ตตฺถ กตโม วิหิํสาวิตกฺโก, วิหิํสาปฏิสํยุตฺโต ตกฺโก วิตกฺโก มิจฺฉาสงฺกปฺโป, อยํ วุจฺจติ วิหิํสาวิตกฺโก.\csromanspacer{} อิเม ตโย อกุสลวิตกฺกา. (2)}

\proseitem{911}{ตตฺถ กตมา ติสฺโส อกุสลสญฺญา, กามสญฺญา พฺยาปาทสญฺญา, วิหิํสาสญฺญา.}

\prose{ตตฺถ กตมา กามสญฺญา, กามปฏิสํยุตฺตา สญฺญา สญฺชานนา สญฺชานิตตฺตา, อยํ วุจฺจติ กามสญฺญา.}

\prose{ตตฺถ กตมา พฺยาปาทสญฺญา, พฺยาปาทปฏิสํยุตฺตา สญฺญา สญฺชานนา สญฺชานิตตฺตํ, อยํ วุจฺจติ พฺยาปาทสญฺญา.}

\prose{ตตฺถ กตมา วิหิํสาสญฺญา, วิหิํสาปฏิสํยุตฺตา สญฺญา สญฺชานนา สญฺชานิตตฺตํ, อยํ วุจฺจติ วิหิํสาสญฺญา.\csromanspacer{} อิมา ติสฺโส อกุสลสญฺญา. (3)}

\proseitem{912}{ตตฺถ กตมา ติสฺโส อกุสลธาตุโย, กามธาตุ พฺยาปาทธาตุ วิหิํสาธาตุ.}

\prose{\csromanfolio{378}ตตฺถ กตมา กามธาตุ, กามวิตกฺโก กามธาตุ, พฺยาปาทวิตกฺโก พฺยาปาทธาตุ, วิหิํสาวิตกฺโก วิหิํสาธาตุ.}

\prose{ตตฺถ กตโม กามวิตกฺโก, กามปฏิสํยุตฺโต ตกฺโก วิตกฺโก มิจฺฉาสงฺกปฺโป, อยํ วุจฺจติ กามวิตกฺโก.}

\prose{ตตฺถ กตโม พฺยาปาทวิตกฺโก, พฺยาปาทปฏิสํยุตฺโต ตกฺโก วิตกฺโก มิจฺฉาสงฺกปฺโป, อยํ วุจฺจติ พฺยาปาทวิตกฺโก.}

\prose{ตตฺถ กตโม วิหิํสาวิตกฺโก, วิหิํสาปฏิสํยุตฺโต ตกฺโก วิตกฺโก มิจฺฉาสงฺกปฺโป, อยํ วุจฺจติ วิหิํสาวิตกฺโก.\csromanspacer{} อิมา ติสฺโส อกุสลธาตุโย. (4)}

\proseitem{913}{ตตฺถ กตมานิ ตีณิ ทุจฺจริตานิ, กายทุจฺจริตํ วจีทุจฺจริตํ มโนทุจฺจริตํ.}

\prose{ตตฺถ กตมํ กายทุจฺจริตํ, ปาณาติปาโต อทินฺนาทานํ กาเมสุมิจฺฉาจาโร, อิทํ วุจฺจติ กายทุจฺจริตํ.}

\prose{ตตฺถ กตมํ วจีทุจฺจริตํ, มุสาวาโท ปิสุณา วาจา ผรุสา วาจา สมฺผปฺปลาโป, อิทํ วุจฺจติ วจีทุจฺจริตํ.}

\prose{ตตฺถ กตมํ มโนทุจฺจริตํ, อภิชฺฌา พฺยาปาโท มิจฺฉาทิฏฺฐิ, อิทํ วุจฺจติ มโนทุจฺจริตํ.}

\prose{ตตฺถ กตมํ กายทุจฺจริตํ, อกุสลํ กายกมฺมํ กายทุจฺจริตํ, อกุสลํ วจีกมฺมํ วจีทุจฺจริตํ, อกุสลํ มโนกมฺมํ มโนทุจฺจริตํ.}

\prose{ตตฺถ กตมํ อกุสลํ กายกมฺมํ, อกุสลา กายสญฺเจตนา อกุสลํ กายกมฺมํ, อกุสลา วจีสญฺเจตนา อกุสลํ วจีกมฺมํ, อกุสลา มโนสญฺเจตนา อกุสลํ มโนกมฺมํ, อิมานิ ตีณิ ทุจฺจริตานิ. (5)}

\proseitem{914}{ตตฺถ กตเม ตโย อาสวา, กามาสโว ภวาสโว อวิชฺชาสโว.}

\prose{ตตฺถ กตโม กามาสโว, โย กาเมสุ กามจฺฉนฺโท กามราโค กามนนฺที กามตณฺหา กามสิเนโห กามปริฬาโห กามมุจฺฉา กามชฺโฌสานํ, อยํ วุจฺจติ กามาสโว.}

\csromanheader{2}{17. ขุทฺทกวตฺถุวิภงฺค}
\prose{\csromanfolio{379}ตตฺถ กตโม ภวาสโว, โย ภเวสุ ภวจฺฉนฺโท ฯเปฯ ภวชฺโฌสานํ, อยํ วุจฺจติ ภวาสโว.}

\prose{ตตฺถ กตโม อวิชฺชาสโว, ทุกฺเข อญฺญาณํ ฯเปฯ อวิชฺชาลงฺคี โมโห อกุสลมูลํ, อยํ วุจฺจติ อวิชฺชาสโว.\csromanspacer{} อิเม ตโย อาสวา. (6)}

\proseitem{915}{ตตฺถ กตมานิ ตีณิ สํโยชนานิ, สกฺกายทิฏฺฐิ วิจิกิจฺฉา สีลพฺพตปรามาโส.}

\prose{ตตฺถ กตมา สกฺกายทิฏฺฐิ, อิธ อสฺสุตวา ปุถุชฺชโน อริยานํ อทสฺสาวี อริยธมฺมสฺส อโกวิโท อริยธมฺเม อวินีโต, สปฺปุริสานํ อทสฺสาวี สปฺปุริสธมฺมสฺส อโกวิโท สปฺปุริสธมฺเม อวินีโต รูปํ อตฺตโต สมนุปสฺสติ, รูปวนฺตํ วา อตฺตานํ, อตฺตนิ วา รูปํ, รูปสฺมิํ วา อตฺตานํ, เวทนํ ฯเปฯ สญฺญํ ฯเปฯ สงฺขาเร ฯเปฯ วิญฺญาณํ อตฺตโต สมนุปสฺสติ, วิญฺญาณวนฺตํ วา อตฺตานํ, อตฺตนิ วา วิญฺญาณํ, วิญฺญาณสฺมิํ วา อตฺตานํ, ยา เอวรูปา ทิฏฺฐิ ทิฏฺฐิคตํ ฯเปฯ วิปริยาสคฺคาโห, อยํ วุจฺจติ สกฺกายทิฏฺฐิ.}

\prose{ตตฺถ กตมา วิจิกิจฺฉา, สตฺถริ กงฺขติ วิจิกิจฺฉติ, ธมฺเม กงฺขติ วิจิกิจฺฉติ, สํเฆ กงฺขติ วิจิกิจฺฉติ, สิกฺขาย กงฺขติ วิจิกิจฺฉติ, ปุพฺพนฺเต กงฺขติ วิจิกิจฺฉติ, อปรนฺเต กงฺขติ วิจิกิจฺฉติ, ปุพฺพนฺตาปรนฺเต กงฺขติ วิจิกิจฺฉติ, อิทปฺปจฺจยตาปฏิจฺจสมุปฺปนฺเนสุ ธมฺเมสุ กงฺขติ วิจิกิจฺฉติ, ยา เอวรูปา กงฺขา กงฺขายนา กงฺขายิตตฺตํ ถมฺภิตตฺตํ จิตฺตสฺส มโนวิเลโข, อยํ วุจฺจติ วิจิกิจฺฉา.}

\prose{ตตฺถ กตโม สีลพฺพตปรามาโส, อิโต พหิทฺธา สมณพฺราหฺมณานํ สีเลน สุทฺธิ วเตน สุทฺธิ สีลพฺพเตน สุทฺธีติ ยา เอวรูปา ทิฏฺฐิ ทิฏฺฐิคตํ ฯเปฯ วิปริยาสคฺคาโห, อยํ วุจฺจติ สีลพฺพตปรามาโส.\csromanspacer{} อิมานิ ตีณิ สํโยชนานิ. (7)}

\proseitem{916}{ตตฺถ กตมา ติสฺโส ตณฺหา, กามตณฺหา ภวตณฺหา วิภวตณฺหา.}

\prose{ตตฺถ กตมา ภวตณฺหา, ภวทิฏฺฐิสหคโต ราโค สาราโค จิตฺตสฺส สาราโค, อยํ วุจฺจติ ภวตณฺหา.}

\prose{\csromanfolio{380}ตตฺถ กตมา วิภวตณฺหา, อุจฺเฉททิฏฺฐิสหคโต ราโค สาราโค จิตฺตสฺส สาราโค, อยํ วุจฺจติ วิภวตณฺหา.\csromanspacer{} อวเสสา ตณฺหา กามตณฺหา.}

\prose{ตตฺถ กตมา กามตณฺหา, กามธาตุปฏิสํยุตฺโต ราโค สาราโค จิตฺตสฺส สาราโค, อยํ วุจฺจติ กามตณฺหา.}

\prose{( )\footnote{(ตตฺถ กตมา ภวตณฺหา,) (?)} รูปธาตุ-อรูปธาตุปฏิสํยุตฺโต ราโค สาราโค จิตฺตสฺส สาราโค, อยํ วุจฺจติ ภวตณฺหา.}

\prose{( )\footnote{(ตตฺถ กตมา วิภวตณฺหา,) (?)} อุจฺเฉททิฏฺฐิสหคโต ราโค สาราโค จิตฺตสฺส สาราโค, อยํ วุจฺจติ วิภวตณฺหา อิมา ติสฺโส ตณฺหา. (8)}

\proseitem{917}{ตตฺถ กตมา อปราปิ ติสฺโส ตณฺหา, กามตณฺหา รูปตณฺหา อรูปตณฺหา.}

\prose{ตตฺถ กตมา กามตณฺหา, กามธาตุปฏิสํยุตฺโต ราโค สาราโค จิตฺตสฺส สาราโค, อยํ วุจฺจติ กามตณฺหา.}

\prose{ตตฺถ กตมา รูปตณฺหา, รูปธาตุปฏิสํยุตฺโต ราโค สาราโค จิตฺตสฺส สาราโค, อยํ วุจฺจติ รูปตณฺหา.}

\prose{ตตฺถ กตมา อรูปตณฺหา, อรูปธาตุปฏิสํยุตฺโต ราโค สาราโค จิตฺตสฺส สาราโค, อยํ วุจฺจติ อรูปตณฺหา.\csromanspacer{} อิมา ติสฺโส ตณฺหา. (9)}

\proseitem{918}{ตตฺถ กตมา อปราปิ ติสฺโส ตณฺหา, รูปตณฺหา อรูปตณฺหา นิโรธตณฺหา.}

\prose{ตตฺถ กตมา รูปตณฺหา, รูปธาตุปฏิสํยุตฺโต ราโค สาราโค จิตฺตสฺส สาราโค, อยํ วุจฺจติ รูปตณฺหา.}

\prose{ตตฺถ กตมา อรูปตณฺหา, อรูปธาตุปฏิสํยุตฺโต ราโค สาราโค จิตฺตสฺส สาราโค, อยํ วุจฺจติ อรูปตณฺหา.}

\prose{ตตฺถ กตมา นิโรธตณฺหา, อุจฺเฉททิฏฺฐิสหคโต ราโค สาราโค จิตฺตสฺส สาราโค, อยํ วุจฺจติ นิโรธตณฺหา.\csromanspacer{} อิมา ติสฺโส ตณฺหา. (10)}

\csromanheader{2}{17. ขุทฺทกวตฺถุวิภงฺค}
\proseitem{919}{\csromanfolio{381}ตตฺถ กตมา ติสฺโส เอสนา, กาเมสนา ภเวสนา พฺรหฺมจริเยสนา.}

\prose{ตตฺถ กตมา กาเมสนา, โย กาเมสุ กามจฺฉนฺโท ฯเปฯ กามชฺโฌสานํ, อยํ วุจฺจติ กาเมสนา.}

\prose{ตตฺถ กตมา ภเวสนา, โย ภเวสุ ภวจฺฉนฺโท ฯเปฯ ภวชฺโฌสานํ, อยํ วุจฺจติ ภเวสนา.}

\prose{ตตฺถ กตมา พฺรหฺมจริเยสนา, “สสฺสโต โลโก”ติ วา “อสสฺสโต โลโก”ติ วา ฯเปฯ “เนว โหติ น น โหติ ตถาคโต ปรํ มรณา”ติ วา, ยา เอวรูปา ทิฏฺฐิ ทิฏฺฐิคตํ ฯเปฯ วิปริยาสคฺคาโห, อยํ วุจฺจติ พฺรหฺมจริเยสนา.}

\prose{ตตฺถ กตมา กาเมสนา, กามราโค ตเทกฏฺฐํ อกุสลํ กายกมฺมํ วจีกมฺมํ มโนกมฺมํ, อยํ วุจฺจติ กาเมสนา.}

\prose{ตตฺถ กตมา ภเวสนา, ภวราโค ตเทกฏฺฐํ อกุสลํ กายกมฺมํ วจีกมฺมํ มโนกมฺมํ, อยํ วุจฺจติ ภเวสนา.}

\prose{ตตฺถ กตมา พฺรหฺมจริเยสนา, อนฺตคฺคาหิกา ทิฏฺฐิ ตเทกฏฺฐํ อกุสลํ กายกมฺมํ วจีกมฺมํ มโนกมฺมํ, อยํ วุจฺจติ พฺรหฺมจริเยสนา.\csromanspacer{} อิมา ติสฺโส เอสนา. (11)}

\proseitem{920}{ตตฺถ กตมา ติสฺโส วิธา, “เสยฺโยหมสฺมี”ติ วิธา, “สทิโสหมสฺมี”ติ วิธา, “หีโนหมสฺมี”ติ วิธา, อิมา ติสฺโส วิธา. (12)}

\proseitem{921}{ตตฺถ กตมานิ ตีณิ ภยานิ, ชาติภยํ ชราภยํ มรณภยํ.}

\prose{ตตฺถ กตมํ ชาติภยํ, ชาติํ ปฏิจฺจ ภยํ ภยานกํ ฉมฺภิตตฺตํ โลมหํโส เจตโส อุตฺราโส, อิทํ วุจฺจติ ชาติภยํ.}

\prose{ตตฺถ กตมํ ชราภยํ, ชรํ ปฏิจฺจ ภยํ ภยานกํ ฉมฺภิตตฺตํ โลมหํโส เจตโส อุตฺราโส, อิทํ วุจฺจติ ชราภยํ.}

\prose{\csromanfolio{382}ตตฺถ กตมํ มรณภยํ, มรณํ ปฏิจฺจ ภยํ ภยานกํ ฉมฺภิตตฺตํ โลมหํโส เจตโส อุตฺราโส, อิทํ วุจฺจติ มรณภยํ.\csromanspacer{} อิมานิ ตีณิ ภยานิ. (13)}

\proseitem{922}{ตตฺถ กตมานิ ตีณิ ตมานิ, อตีตํ วา อทฺธานํ อารพฺภ กงฺขติ วิจิกิจฺฉติ นาธิมุจฺจติ น สมฺปสีทติ, อนาคตํ วา อทฺธานํ อารพฺภ กงฺขติ วิจิกิจฺฉติ นาธิมุจฺจติ น สมฺปสีทติ, เอตรหิ วา ปจฺจุปฺปนฺนํ อทฺธานํ อารพฺภ กงฺขติ วิจิกิจฺฉติ นาธิมุจฺจติ น สมฺปสีทติ, อิมานิ ตีณิ ตมานิ. (14)}

\proseitem{923}{ตตฺถ กตมานิ ตีณิ ติตฺถายตนานิ, อิเธกจฺโจ สมโณ วา พฺราหฺมโณ วา เอวํวาที โหติ เอวํทิฏฺฐี “ยํ กิญฺจายํ ปุริสปุคฺคโล ปฏิสํเวเทติ สุขํ วา ทุกฺขํ วา อทุกฺขมสุขํ วา, สพฺพํ ตํ ปุพฺเพ กตเหตู”ติ, อิธ ปเนกจฺโจ สมโณ วา พฺราหฺมโณ วา เอวํวาที โหติ เอวํทิฏฺฐี “ยํ กิญฺจายํ ปุริสปุคฺคโล ปฏิสํเวเทติ สุขํ วา ทุกฺขํ วา อทุกฺขมสุขํ วา, สพฺพํ ตํ อิสฺสรนิมฺมานเหตู”ติ, อิธ ปเนกจฺโจ สมโณ วา พฺราหฺมโณ วา เอวํวาที โหติ เอวํทิฏฺฐี “ยํ กิญฺจายํ ปุริสปุคฺคโล ปฏิสํเวเทติ สุขํ วา ทุกฺขํ วา อทุกฺขมสุขํ วา, สพฺพํ ตํ อเหตุ อปฺปจฺจยาติ, อิมานิ ตีณิ ติตฺถายตนานิ. (15)}

\proseitem{924}{ตตฺถ กตเม ตโย กิญฺจนา\footnote{กตมานิ ตีณิ กิญฺจนานิ (?) ที 3. 182 ปิฏฺเฐปิ.}, ราโค กิญฺจนํ, โทโส กิญฺจนํ, โมโห กิญฺจนํ, อิเม ตโย กิญฺจนา. (16)}

\prose{ตตฺถ กตมานิ ตีณิ องฺคณานิ, ราโค องฺคณํ, โทโส องฺคณํ, โมโห องฺคณํ, อิมานิ ตีณิ องฺคณานิ. (17)}

\prose{ตตฺถ กตมานิ ตีณิ มลานิ, ราโค มลํ, โทโส มลํ, โมโห มลํ, อิมานิ ตีณิ มลานิ. (18)}

\prose{ตตฺถ กตมานิ ตีณิ วิสมานิ, ราโค วิสมํ, โทโส วิสมํ, โมโห วิสมํ, อิมานิ ตีณิ วิสมานิ. (19)}

\prose{ตตฺถ กตมานิ อปรานิปิ ตีณิ วิสมานิ, กายวิสมํ วจีวิสมํ มโนวิสมํ, อิมานิ ตีณิ วิสมานิ. (20)}

\csromanheader{2}{17. ขุทฺทกวตฺถุวิภงฺค}
\prose{\csromanfolio{383}ตตฺถ กตเม ตโย อคฺคี, ราคคฺคิ โทสคฺคิ โมหคฺคิ, อิเม ตโย อคฺคี. (21)}

\prose{ตตฺถ กตเม ตโย กสาวา, ราคกสาโว โทสกสาโว โมหกสาโว, อิเม ตโย กสาวา. (22)}

\prose{ตตฺถ กตเม อปเรปิ ตโย กสาวา, กายกสาโว วจีกสาโว มโนกสาโว, อิเม ตโย กสาวา. (23)}

\proseitem{925}{ตตฺถ กตมา อสฺสาททิฏฺฐิ, อิเธกจฺโจ สมโณ วา พฺราหฺมโณ วา เอวํวาที โหติ เอวํทิฏฺฐี “นตฺถิ กาเมสุ โทโส”ติ, โส กาเมสุ ปาตพฺยตํ อาปชฺชติ, อยํ วุจฺจติ อสฺสาททิฏฺฐิ.}

\prose{ตตฺถ กตมา อตฺตานุทิฏฺฐิ, อิธ อสฺสุตวา ปุถุชฺชโน อริยานํ อทสฺสาวี อริยธมฺมสฺส อโกวิโท อริยธมฺเม อวินีโต สปฺปุริสานํ อทสฺสาวี สปฺปุริสธมฺมสฺส อโกวิโท สปฺปุริสธมฺเม อวินีโต รูปํ อตฺตโต สมนุปสฺสติ, รูปวนฺตํ วา อตฺตานํ, อตฺตนิ วา รูปํ, รูปสฺมิํ วา อตฺตานํ.\csromanspacer{} เวทนํ ฯเปฯ สญฺญํ ฯเปฯ สงฺขาเร ฯเปฯ วิญฺญาณํ อตฺตโต สมนุปสฺสติ, วิญฺญาณวนฺตํ วา อตฺตานํ, อตฺตนิ วา วิญฺญาณํ, วิญฺญาณสฺมิํ วา อตฺตานํ, ยา เอวรูปา ทิฏฺฐิ ทิฏฺฐิคตํ ฯเปฯ วิปริยาสคฺคาโห, อยํ วุจฺจติ อตฺตานุทิฏฺฐิ.}

\prose{ตตฺถ กตมา มิจฺฉาทิฏฺฐิ, “นตฺถิ ทินฺนํ, นตฺถิ ยิฏฺฐํ ฯเปฯ เย อิมญฺจ โลกํ ปรญฺจ โลกํ สยํ อภิญฺญา สจฺฉิกตฺวา ปเวเทนฺตี”ติ ยา เอวรูปา ทิฏฺฐิ ทิฏฺฐิคตํ ฯเปฯ วิปริยาสคฺคาโห, อยํ วุจฺจติ มิจฺฉาทิฏฺฐิ.\csromanspacer{} สสฺสตทิฏฺฐิ.\csromanspacer{} อสฺสาททิฏฺฐิ.\csromanspacer{} สกฺกายทิฏฺฐิ.\csromanspacer{} อตฺตานุทิฏฺฐิ.\csromanspacer{} อุจฺเฉททิฏฺฐิ.\csromanspacer{} มิจฺฉาทิฏฺฐิ. (24)}

\proseitem{926}{ตตฺถ กตมา อรติ, ปนฺเตสุ วา เสนาสเนสุ อญฺญตรญฺญตเรสุ วา อธิกุสเลสุ ธมฺเมสุ อรติ อรติตา อนภิรติ อนภิรมณา อุกฺกณฺฐิตา ปริตสฺสิตา, อยํ วุจฺจติ อรติ.}

\prose{ตตฺถ กตมา วิเหสา, อิเธกจฺโจ ปาณินา วา เลฑฺฑุนา วา ทณฺเฑน วา สตฺเถน วา รชฺชุยา วา อญฺญตรญฺญตเรน สตฺเต วิเหเฐติ, ยา เอวรูปา เหฐนา วิเหฐนา หิํสนา วิหิํสนา โรสนา วิโรสนา ปรูปฆาโต, อยํ วุจฺจติ วิเหสา.}

\prose{\csromanfolio{384}ตตฺถ กตมา อธมฺมจริยา, กาเยน อธมฺมจริยาวิสมจริยา, วาจาย อธมฺมจริยาวิสมจริยา, มนสา อธมฺมจริยาวิสมจริยา, อยํ วุจฺจติ อธมฺมจริยา. (25)}

\proseitem{927}{ตตฺถ กตมา โทวจสฺสตา, สหธมฺมิเก วุจฺจมาเน โทวจสฺสายํ โทวจสฺสิยํ โทวจสฺสตา วิปฺปฏิกุลคฺคาหิตา วิปจฺจนีกสาตตา อนาทริยํ อนาทรตา อคารวตา อปฺปติสฺสวตา, อยํ วุจฺจติ โทวจสฺสตา.}

\prose{ตตฺถ กตมา ปาปมิตฺตตา, เย เต ปุคฺคลา อสฺสทฺธา ทุสฺสีลา อปฺปสฺสุตา มจฺฉริโน ทุปฺปญฺญา, ยา เตสํ เสวนา นิเสวนา สํเสวนา ภชนา สมฺภชนา ภตฺติ สมฺภตฺติ ตํสมฺปวงฺกตา, อยํ วุจฺจติ ปาปมิตฺตตา.}

\prose{ตตฺถ กตมา นานตฺตสญฺญา,กามสญฺญา พฺยาปาทสญฺญา วิหิํสาสญฺญา, อยํ วุจฺจติ นานตฺตสญฺญา.\csromanspacer{} สพฺพาปิ อกุสลา สญฺญา นานตฺตสญฺญา. (26)}

\proseitem{928}{ตตฺถ กตมํ อุทฺธจฺจํ, ยํ จิตฺตสฺส อุทฺธจฺจํ อวูปสโม เจตโส วิกฺเขโป ภนฺตตฺตํ จิตฺตสฺส, อิทํ วุจฺจติ อุทฺธจฺจํ.}

\prose{ตตฺถ กตมํ โกสชฺชํ, กายทุจฺจริเต วา วจีทุจฺจริเต วา มโนทุจฺจริเต วา ปญฺจสุ วา กามคุเณสุ จิตฺตสฺส โวสฺสคฺโค โวสฺสคฺคานุปฺปทานํ กุสลานํ ธมฺมานํ ภาวนาย อสกฺกจฺจกิริยตา อสาตจฺจกิริยตา อนฏฺฐิตกิริยตา โอลีนวุตฺติตา นิกฺขิตฺตฉนฺทตา นิกฺขิตฺตธุรตา อนาเสวนา อภาวนา อพหุลีกมฺมํ อนธิฏฺฐานํ อนนุโยโค ปมาโท, อิทํ วุจฺจติ โกสชฺชํ.}

\prose{ตตฺถ กตโม ปมาโท, กายทุจฺจริเต วา วจีทุจฺจริเต วา มโนทุจฺจริเต วา ปญฺจสุ วา กามคุเณสุ จิตฺตสฺส โวสฺสคฺโค โวสฺสคฺคานุปฺปทานํ กุสลานํ ธมฺมานํ ภาวนาย อสกฺกจฺจกิริยตา อสาตจฺจกิริยตา อนฏฺฐิตกิริยตา โอลีนวุตฺติตา นิกฺขิตฺตฉนฺทตา นิกฺขิตฺตธุรตา อนาเสวนา อภาวนา อพหุลีกมฺมํ อนธิฏฺฐานํ อนนุโยโค ปมาโท, โย เอวรูโป ปมาโท ปมชฺชนา ปมชฺชิตตฺตํ, อยํ วุจฺจติ ปมาโท. (27)}

\csromanheader{2}{17. ขุทฺทกวตฺถุวิภงฺค}
\proseitem{929}{\csromanfolio{385}ตตฺถ กตมา อสนฺตุฏฺฐิตา, อิตรีตรจีวรปิณฺฑปาตเสนาสนคิลานปจฺจยเภสชฺชปริกฺขาเรหิ ปญฺจหิ วา กามคุเณหิ อสนฺตุฏฺฐสฺส ภิยฺโยกมฺยตา, ยา เอวรูปา อิจฺฉา อิจฺฉาคตา อสนฺตุฏฺฐิตา ราโค สาราโค จิตฺตสฺส สาราโค, อยํ วุจฺจติ อสนฺตุฏฺฐิตา.}

\prose{ตตฺถ กตมา อสมฺปชญฺญตา, ยํ อญฺญาณํ อทสฺสนํ ฯเปฯ อวิชฺชาลงฺคี โมโห อกุสลมูลํ, อยํ วุจฺจติ อสมฺปชญฺญตา.}

\prose{ตตฺถ กตมา มหิจฺฉตา, อิตรีตรจีวรปิณฺฑปาตเสนาสนคิลานปจฺจยเภสชฺชปริกฺขาเรหิ ปญฺจหิ วา กามคุเณหิ อสนฺตุฏฺฐสฺส ภิยฺโยกมฺยตา, ยา เอวรูปา อิจฺฉา อิจฺฉาคตา มหิจฺฉตา ราโค สาราโค จิตฺตสฺส สาราโค, อยํ วุจฺจติ มหิจฺฉตา. (28)}

\proseitem{930}{ตตฺถ กตมํ อหิริกํ, ยํ น หิรียติ หิรียิตพฺเพน น หิรียติ ปาปกานํ อกุสลานํ ธมฺมานํ สมาปตฺติยา, อิทํ วุจฺจติ อหิริกํ.}

\prose{ตตฺถ กตมํ อโนตฺตปฺปํ, ยํ น โอตฺตปฺปติ โอตฺตปฺปิตพฺเพน น โอตฺตปฺปติ ปาปกานํ อกุสลานํ ธมฺมานํ สมาปตฺติยา, อิทํ วุจฺจติ อโนตฺตปฺปํ.}

\prose{ตตฺถ กตโม ปมาโท, กายทุจฺจริเต วา วจีทุจฺจริเต วา มโนทุจฺจริเต วา ปญฺจสุ วา กามคุเณสุ จิตฺตสฺส โวสฺสคฺโค โวสฺสคฺคานุปฺปทานํ กุสลานํ ธมฺมานํ ภาวนาย อสกฺกจฺจกิริยตา อสาตจฺจกิริยตา อนฏฺฐิตกิริยตา โอลีนวุตฺติตา นิกฺขิตฺตฉนฺทตา นิกฺขิตฺตธุรตา อนาเสวนา อภาวนา อพหุลีกมฺมํ อนธิฏฺฐานํ อนนุโยโค ปมาโท, โย เอวรูโป ปมาโท, ปมชฺชนา ปมชฺชิตตฺตํ, อยํ วุจฺจติ ปมาโท. (29)}

\proseitem{931}{ตตฺถ กตมํ อนาทริยํ, ยํ อนาทริยํ อนาทรตา อคารวตา อปฺปติสฺสวตา อนทฺทา อนทฺทายนา อนทฺทายิตตฺตํ\footnote{อนาทา อนาทายนา อนาทายิตตฺตํ (สฺยา)} อสีลฺยํ อจิตฺตีกาโร, อิทํ วุจฺจติ อนาทริยํ.}

\prose{ตตฺถ กตมา โทวจสฺสตา, สหธมฺมิเก วุจฺจมาเน โทวจสฺสายํ โทวจสฺสิยํ โทวจสฺสตา วิปฺปฏิกุลคฺคาหิตา วิปจฺจนีกสาตตา \csromanfolio{386}อนาทริยํ อนาทรตา อคารวตา อปฺปติสฺสวตา, อยํ วุจฺจติ โทวจสฺสตา.}

\prose{ตตฺถ กตมา ปาปมิตฺตตา, เย เต ปุคฺคลา อสฺสทฺธา ทุสฺสีลา อปฺปสฺสุตา มจฺฉริโน ทุปฺปญฺญา, ยา เตสํ เสวนา นิเสวนา สํเสวนา ปฏิเสวนา ภชนา สมฺภชนา ภตฺติ สมฺภตฺติ ตํสมฺปวงฺกตา, อยํ วุจฺจติ ปาปมิตฺตตา. (30)}

\proseitem{932}{ตตฺถ กตมํ อสฺสทฺธิยํ, อิเธกจฺโจ อสฺสทฺโธ โหติ น สทฺทหติ พุทฺธํ วา ธมฺมํ วา สํฆํ วา, ยํ เอวรูปํ อสฺสทฺธิยํ อสฺสทฺทหนา อโนกปฺปนา อนภิปฺปสาโท, อิทํ วุจฺจติ อสฺสทฺธิยํ.}

\prose{ตตฺถ กตมา อวทญฺญุตา, ปญฺจ มจฺฉริยานิ อาวาสมจฺฉริยํ กุลมจฺฉริยํ ลาภมจฺฉริยํ วณฺณมจฺฉริยํ ธมฺมมจฺฉริยํ, ยํ เอวรูปํ มจฺเฉรํ มจฺฉรายนา มจฺฉรายิตตฺตํ เววิจฺฉํ กทริยํ กฏุกญฺจุกตา อคฺคหิตตฺตํ จิตฺตสฺส, อยํ วุจฺจติ อวทญฺญุตา.}

\prose{ตตฺถ กตมํ โกสชฺชํ, กายทุจฺจริเต วา วจีทุจฺจริเต วา มโนทุจฺจริเต วา ปญฺจสุ วา กามคุเณสุ จิตฺตสฺส โวสฺสคฺโค โวสฺสคฺคานุปฺปทานํ กุสลานํ ธมฺมานํ ภาวนาย อสกฺกจฺจกิริยตา อสาตจฺจกิริยตา อนฏฺฐิตกิริยตา โอลีนวุตฺติตา นิกฺขิตฺตฉนฺทตา นิกฺขิตฺตธุรตา อนาเสวนา อภาวนา อพหุลีกมฺมํ อนธิฏฺฐานํ อนนุโยโค ปมาโท, อิทํ วุจฺจติ โกสชฺชํ. (31)}

\proseitem{933}{ตตฺถ กตมํ อุทฺธจฺจํ, ยํ จิตฺตสฺส อุทฺธจฺจํ อวูปสโม เจตโส วิกฺเขโป ภนฺตตฺตํ จิตฺตสฺส, อิทํ วุจฺจติ อุทฺธจฺจํ.}

\prose{ตตฺถ กตโม อสํวโร, อิเธกจฺโจ จกฺขุนา รูปํ ทิสฺวา นิมิตฺตคฺคาหี โหติ อนุพฺยญฺชนคฺคาหี, ยตฺวาธิกรณเมนํ จกฺขุนฺทฺริยํ อสํวุตํ วิหรนฺตํ อภิชฺฌาโทมนสฺสา ปาปกา อกุสลา ธมฺมา อนฺวาสฺสเวยฺยุํ, ตสฺส สํวราย น ปฏิปชฺชติ, น รกฺขติ จกฺขุนฺทฺริยํ, จกฺขุนฺทฺริเย น สํวรํ อาปชฺชติ.\csromanspacer{} โสเตน สทฺทํ สุตฺวา ฯเปฯ ฆาเนน คนฺธํ ฆายิตฺวา ฯเปฯ ชิวฺหาย รสํ สายิตฺวา ฯเปฯ กาเยน โผฏฺฐพฺพํ ผุสิตฺวา ฯเปฯ มนสา ธมฺมํ วิญฺญาย นิมิตฺตคฺคาหี โหติ อนุพฺยญฺชนคฺคาหี, ยตฺวาธิกรณเมนํ มนินฺทฺริยํ อสํวุตํ วิหรนฺตํ อภิชฺฌาโทมนสฺสา ปาปกา อกุสลา ธมฺมา อนฺวาสฺสเวยฺยุํ, ตสฺส}

\vspace{\tipitakaparskip}
\csromancenter{ขุทฺทกวตฺถุวิภงฺค สํวราย น ปฏิปชฺชติ, น รกฺขติ มนินฺทฺริยํ, มนินฺทฺริเย น สํวรํ อาปชฺชติ, อยํ วุจฺจติ อสํวโร.}
\prose{\csromanfolio{387}ตตฺถ กตมํ ทุสฺสีลฺยํ, กายิโก วีติกฺกโม วาจสิโก วีติกฺกโม กายิกวาจสิโก วีติกฺกโม, อิทํ วุจฺจติ ทุสฺสีลฺยํ. (32)}

\proseitem{934}{ตตฺถ กตมา อริยานํ อทสฺสนกมฺยตา.\csromanspacer{} ตตฺถ กตเม อริยา, อริยา วุจฺจนฺติ พุทฺธา จ พุทฺธสาวกา จ, ยา อิเมสํ อริยานํ อทสฺสนกมฺยตา อทฏฺฐุกมฺยตา อสเมตุกมฺยตา อสมาคนฺตุกมฺยตา.\csromanspacer{} อยํ วุจฺจติ อริยานํ อทสฺสนกมฺยตา.}

\prose{ตตฺถ กตมา สทฺธมฺมํ อโสตุกมฺยตา.\csromanspacer{} ตตฺถ กตโม สทฺธมฺโม, จตฺตาโร สติปฏฺฐานา จตฺตาโร สมฺมปฺปธานา จตฺตาโร อิทฺธิปาทา ปญฺจินฺทฺริยานิ ปญฺจ พลานิ สตฺต โพชฺฌงฺคา อริโย อฏฺฐงฺคิโก มคฺโค, อยํ วุจฺจติ สทฺธมฺโม.\csromanspacer{} ยา อิมสฺส สทฺธมฺมสฺส อโสตุกมฺยตา อสวนกมฺยตา อนุคฺคเหตุกมฺยตา อธาเรตุกมฺยตา, อยํ วุจฺจติ สทฺธมฺมํ อโสตุกมฺยตา.}

\prose{ตตฺถ กตมา อุปารมฺภจิตฺตตา.\csromanspacer{} ตตฺถ กตโม อุปารมฺโภ, โย อุปารมฺโภ อนุปารมฺโภ อุปารมฺภนา อนุปารมฺภนา อนุปารมฺภิ ตตฺตํ อุญฺญา อวญฺญา ปริภโว รนฺธคเวสิตา, อยํ วุจฺจติ อุปารมฺภจิตฺตตา. (33)}

\proseitem{935}{ตตฺถ กตมํ มุฏฺฐสฺสจฺจํ, ยา อสฺสติ อนนุสฺสติ อปฺปฏิสฺสติ อสฺสติ อสฺสรณตา อธารณตา ปิลาปนตา สมฺมุสนตา, อิทํ วุจฺจติ มุฏฺฐสฺสจฺจํ.}

\prose{ตตฺถ กตมํ อสมฺปชญฺญํ, ยํ อญฺญาณํ อทสฺสนํ ฯเปฯ อวิชฺชาลงฺคี โมโห อกุสลมูลํ, อิทํ วุจฺจติ อสมฺปชญฺญํ.}

\prose{ตตฺถ กตโม เจตโส วิกฺเขโป, ยํ จิตฺตสฺส อุทฺธจฺจํ อวูปสโม เจตโส วิกฺเขโป ภนฺตตฺตํ จิตฺตสฺส, อยํ วุจฺจติ เจตโส วิกฺเขโป. (34)}

\proseitem{936}{ตตฺถ กตโม อโยนิโส มนสิกาโร, “อนิจฺเจ นิจฺจนฺ”ติ อโยนิโส มนสิกาโร, “ทุกฺเข สุขนฺ”ติ อโยนิโส มนสิกาโร, “อนตฺตนิ อตฺตา”ติ อโยนิโส มนสิกาโร, \csromanfolio{388}“อสุเภ สุภนฺ”ติ อโยนิโส มนสิกาโร, สจฺจวิปฺปฏิกุเลน วา จิตฺตสฺส อาวฏฺฏนา อนาวฏฺฏนา อาโภโค สมนฺนาหาโร มนสิกาโร, อยํ วุจฺจติ อโยนิโส มนสิกาโร.}

\prose{ตตฺถ กตมา กุมฺมคฺคเสวนา.\csromanspacer{} ตตฺถ กตโม กุมฺมคฺโค, มิจฺฉาทิฏฺฐิ มิจฺฉาสงฺกปฺโป มิจฺฉาวาจา มิจฺฉากมฺมนฺโต มิจฺฉา-อาชีโว มิจฺฉาวายาโม มิจฺฉาสติ มิจฺฉาสมาธิ, อยํ วุจฺจติ กุมฺมคฺโค.\csromanspacer{} ยา อิมสฺส กุมฺมคฺคสฺส เสวนา นิเสวนา สํเสวนา ภชนา สมฺภชนา ภตฺติ สมฺภตฺติ ตํสมฺปวงฺกตา, อยํ วุจฺจติ กุมฺมคฺคเสวนา.}

\prose{ตตฺถ กตมํ เจตโส จ ลีนตฺตํ, ยา จิตฺตสฺส อกลฺยตา อกมฺมญฺญตา โอลียนา สลฺลียนา ลีนํ ลียนา ลียิตตฺตํ ถินํ ถิยนา ถิยิตตฺตํ จิตฺตสฺส, อิทํ วุจฺจติ เจตโส จ ลีนตฺตํ. (35)}

\csromanheader{2}{ติกํ.}
\csromansectionrule
\csromanmatikaanchor{mtk.182}
\csromantocmarkref{subparagraph}{4. จตุกฺกนิทฺเทส}{mtk.182}
\bookmark[dest={mtk.182},level=5]{4. จตุกฺกนิทฺเทส}
\csromanheader{6}{4. จตุกฺกนิทฺเทส}
\proseitem{937}{ตตฺถ กตเม จตฺตาโร อาสวา, กามาสโว ภวาสโว ทิฏฺฐาสโว อวิชฺชาสโว.}

\prose{ตตฺถ กตโม กามาสโว, โย กาเมสุ กามจฺฉนฺโท กามราโค กามนนฺที กามตณฺหา กามสิเนโห กามปริฬาโห กามมุจฺฉา กามชฺโฌสานํ, อยํ วุจฺจติ กามาสโว.}

\prose{ตตฺถ กตโม ภวาสโว, โย ภเวสุ ภวจฺฉนฺโท ฯเปฯ ภวชฺโฌสานํ, อยํ วุจฺจติ ภวาสโว.}

\prose{ตตฺถ กตโม ทิฏฺฐาสโว, “สสฺสโต โลโก”ติ วา “อสสฺสโต โลโก”ติ วา “อนฺตวา โลโก”ติ วา “อนนฺตวา โลโก”ติ วา “ตํ ชีวํ ตํ สรีรนฺ”ติ วา “อญฺญํ ชีวํ อญฺญํ สรีรนฺ”ติ วา “โหติ ตถาคโต ปรํ มรณา”ติ วา “น โหติ ตถาคโต ปรํ มรณา”ติ วา “โหติ จ น จ โหติ ตถาคโต ปรํ มรณา”ติ วา “เนว โหติ น น โหติ ตถาคโต ปรํ มรณา”ติ วา, ยา เอวรูปา}

\vspace{\tipitakaparskip}
\csromancenter{ขุทฺทกวตฺถุวิภงฺค ทิฏฺฐิ ทิฏฺฐิคตํ ฯเปฯ วิปริยาสคฺคาโห, อยํ วุจฺจติ ทิฏฺฐาสโว.\csromanspacer{} สพฺพาปิ มิจฺฉาทิฏฺฐิ ทิฏฺฐาสโว.}
\prose{\csromanfolio{389}ตตฺถ กตโม อวิชฺชาสโว, ทุกฺเข อญฺญาณํ, ทุกฺขสมุทเย อญฺญาณํ, ทุกฺขนิโรเธ อญฺญาณํ, ทุกฺขนิโรธคามินิยา ปฏิปทาย อญฺญาณํ, ปุพฺพนฺเต อญฺญาณํ, อปรนฺเต อญฺญาณํ, ปุพฺพนฺตาปรนฺเต อญฺญาณํ, อิทปฺปจฺจยตา ปฏิจฺจสมุปฺปนฺเนสุ ธมฺเมสุ อญฺญาณํ, ยํ เอวรูปํ อญฺญาณํ อทสฺสนํ ฯเปฯ อวิชฺชาลงฺคี โมโห อกุสลมูลํ, อยํ วุจฺจติ อวิชฺชาสโว.\csromanspacer{} อิเม จตฺตาโร อาสวา. (1)}

\proseitem{938}{ตตฺถ กตเม จตฺตาโร คนฺถา ฯเปฯ จตฺตาโร โอฆา ฯเปฯ จตฺตาโร โยคา ฯเปฯ จตฺตาริ อุปาทานานิ, กามุปาทานํ ทิฏฺฐุปาทานํ สีลพฺพตุปาทานํ อตฺตวาทุปาทานํ.}

\prose{ตตฺถ กตมํ กามุปาทานํ, โย กาเมสุ กามจฺฉนฺโท ฯเปฯ กามชฺโฌสานํ, อิทํ วุจฺจติ กามุปาทานํ.}

\prose{ตตฺถ กตมํ ทิฏฺฐุปาทานํ, “นตฺถิ ทินฺนํ, นตฺถิ ยิฏฺฐํ ฯเปฯ เย อิมญฺจ โลกํ ปรญฺจ โลกํ สยํ อภิญฺญา สจฺฉิกตฺวา ปเวเทนฺตี”ติ ยา เอวรูปา ทิฏฺฐิ ทิฏฺฐิคตํ ฯเปฯ วิปริยาสคฺคาโห, อิทํ วุจฺจติ ทิฏฺฐุปาทานํ.\csromanspacer{} ฐเปตฺวา สีลพฺพตุปาทานญฺจ อตฺตวาทุปาทานญฺจ สพฺพาปิ มิจฺฉาทิฏฺฐิ ทิฏฺฐุปาทานํ.}

\prose{ตตฺถ กตมํ สีลพฺพตุปาทานํ, อิโต พหิทฺธา สมณพฺราหฺมณานํ สีเลน สุทฺธิ วเตน สุทฺธิ สีลพฺพเตน สุทฺธีติ, ยา เอวรูปา ทิฏฺฐิ ทิฏฺฐิคตํ ฯเปฯ วิปริยาสคฺคาโห, อิทํ วุจฺจติ สีลพฺพตุปาทานํ.}

\prose{ตตฺถ กตมํ อตฺตวาทุปาทานํ, อิธ อสฺสุตวา ปุถุชฺชโน อริยานํ อทสฺสาวี อริยธมฺมสฺส อโกวิโท อริยธมฺเม อวินีโต, สปฺปุริสานํ อทสฺสาวี สปฺปุริสธมฺมสฺส อโกวิโท สปฺปุริสธมฺเม อวินีโต รูปํ อตฺตโต สมนุปสฺสติ, รูปวนฺตํ วา อตฺตานํ, อตฺตนิ วา รูปํ, รูปสฺมิํ วา อตฺตานํ.\csromanspacer{} เวทนํ ฯเปฯ.\csromanspacer{} สญฺญํ ฯเปฯ.\csromanspacer{} สงฺขาเร ฯเปฯ.\csromanspacer{} วิญฺญาณํ อตฺตโต สมนุปสฺสติ, วิญฺญาณวนฺตํ วา อตฺตานํ, อตฺตนิ วา วิญฺญาณํ, วิญฺญาณสฺมิํ วา อตฺตานํ, ยา เอวรูปา ทิฏฺฐิ ทิฏฺฐิคตํ ฯเปฯ วิปริยาสคฺคาโห, อิทํ วุจฺจติ อตฺตวาทุปาทานํ.\csromanspacer{} อิมานิ จตฺตาริ อุปาทานานิ. (5)}

\proseitem{939}{\csromanfolio{390}ตตฺถ กตเม จตฺตาโร ตณฺหุปฺปาทา, จีวรเหตุ วา ภิกฺขุโน ตณฺหา อุปฺปชฺชมานา อุปฺปชฺชติ, ปิณฺฑปาตเหตุ วา ภิกฺขุโน ตณฺหา อุปฺปชฺชมานา อุปฺปชฺชติ, เสนาสนเหตุ วา ภิกฺขุโน ตณฺหา อุปฺปชฺชมานา อุปฺปชฺชติ, อิติภวาภวเหตุ วา ภิกฺขุโน ตณฺหา อุปฺปชฺชมานา อุปฺปชฺชติ.\csromanspacer{} อิเม จตฺตาโร ตณฺหุปฺปาทา. (6)}

\prose{ตตฺถ กตมานิ จตฺตาริ อคติคมนานิ, ฉนฺทาคติํ คจฺฉติ, โทสาคติํ คจฺฉติ, โมหาคติํ คจฺฉติ, ภยาคติํ คจฺฉติ, ยา เอวรูปา อคติ อคติคมนํ ฉนฺทคมนํ วคฺคคมนํ วาริคมนํ, อิมานิ จตฺตาริ อคติคมนานิ. (7)}

\prose{ตตฺถ กตเม จตฺตาโร วิปริยาสา, “อนิจฺเจ นิจฺจนฺ”ติ สญฺญาวิปริยาโส จิตฺตวิปริยาโส ทิฏฺฐิวิปริยาโส, “ทุกฺเข สุขนฺ”ติ สญฺญาวิปริยาโส จิตฺตวิปริยาโส ทิฏฺฐิวิปริยาโส, “อนตฺตนิ อตฺตา”ติ สญฺญาวิปริยาโส จิตฺตวิปริยาโส ทิฏฺฐิวิปริยาโส, “อสุเภ สุภนฺ”ติ สญฺญาวิปริยาโส จิตฺตวิปริยาโส ทิฏฺฐิวิปริยาโส.\csromanspacer{} อิเม จตฺตาโร วิปริยาสา. (8)}

\prose{ตตฺถ กตเม จตฺตาโร อนริยโวหารา, อทิฏฺเฐ ทิฏฺฐวาทิตา, อสฺสุเต สุตวาทิตา, อมุเต มุตวาทิตา, อวิญฺญาเต วิญฺญาตวาทิตา, อิเม จตฺตาโร อนริยโวหารา. (9)}

\prose{ตตฺถ กตเม อปเรปิ จตฺตาโร อนริยโวหารา, ทิฏฺเฐ อทิฏฺฐวาทิตา, สุเต อสฺสุตวาทิตา, มุเต อมุตวาทิตา, วิญฺญาเต อวิญฺญาตวาทิตา, อิเม จตฺตาโร อนริยโวหารา. (10)}

\prose{ตตฺถ กตมานิ จตฺตาริ ทุจฺจริตานิ, ปาณาติปาโต อทินฺนาทานํ กาเมสุมิจฺฉาจาโร มุสาวาโท, อิมานิ จตฺตาริ ทุจฺจริตานิ. (11)}

\prose{ตตฺถ กตมานิ อปรานิปิ จตฺตาริ ทุจฺจริตานิ, มุสาวาโท ปิสุณา วาจา ผรุสา วาจา สมฺผปฺปลาโป, อิมานิ จตฺตาริ ทุจฺจริตานิ. (12)}

\prose{ตตฺถ กตมานิ จตฺตาริ ภยานิ, ชาติภยํ ชราภยํ พฺยาธิภยํ มรณภยํ, อิมานิ จตฺตาริ ภยานิ. (13)}

\prose{ตตฺถ กตมานิ อปรานิปิ จตฺตาริ ภยานิ, ราชภยํ โจรภยํ อคฺคิภยํ อุทกภยํ, อิมานิ จตฺตาริ ภยานิ. (14)}

\csromanheader{2}{17. ขุทฺทกวตฺถุวิภงฺค}
\prose{\csromanfolio{391}ตตฺถ กตมานิ อปรานิปิ จตฺตาริ ภยานิ, อูมิภยํ กุมฺภีลภยํ\footnote{กุมฺภีรภยํ, กุมฺภีฬภยํ, (?) (กุมฺภ + อีร + อ = ภย)} อาวฏฺฏภยํ สุสุกาภยํ, อิมานิ จตฺตาริ ภยานิ. (15)}

\prose{ตตฺถ กตมานิ อปรานิปิ จตฺตาริ ภยานิ, อตฺตานุวาทภยํ ปรานุวาทภยํ ทณฺฑภยํ ทุคฺคติภยํ, อิมานิ จตฺตาริ ภยานิ. (16)}

\prose{ตตฺถ กตมา จตสฺโส ทิฏฺฐิโย, “สยํกตํ สุขทุกฺขนฺ”ติ สจฺจโต เถตโต ทิฏฺฐิ อุปฺปชฺชติ, “ปรํกตํ สุขทุกฺขนฺ”ติ สจฺจโต เถตโต ทิฏฺฐิ อุปฺปชฺชติ, “สยํกตญฺจ ปรํกตญฺจ สุขทุกฺขนฺ”ติ สจฺจโต เถตโต ทิฏฺฐิ อุปฺปชฺชติ, “อสยํการํ อปรํการํ อธิจฺจสมุปฺปนฺนํ สุขทุกฺขนฺ”ติ สจฺจโต เถตโต ทิฏฺฐิ อุปฺปชฺชติ, อิมา จตสฺโส ทิฏฺฐิโย. (17)}

\csromanheader{2}{จตุกฺกํ.}
\csromansectionrule
\csromanmatikaanchor{mtk.183}
\csromantocmarkref{subparagraph}{5. ปญฺจกนิทฺเทส}{mtk.183}
\bookmark[dest={mtk.183},level=5]{5. ปญฺจกนิทฺเทส}
\csromanheader{6}{5. ปญฺจกนิทฺเทส}
\proseitem{940}{ตตฺถ กตมานิ ปญฺโจรมฺภาคิยานิ สํโยชนานิ, สกฺกายทิฏฺฐิ วิจิกิจฺฉา สีลพฺพตปรามาโส กามจฺฉนฺโท พฺยาปาโท, อิมานิ ปญฺโจรมฺภาคิยานิ สํโยชนานิ. (1)}

\prose{ตตฺถ กตมานิ ปญฺจุทฺธมฺภาคิยานิ สํโยชนานิ, รูปราโค อรูปราโค มาโน อุทฺธจฺจํ อวิชฺชา, อิมานิ ปญฺจุทฺธมฺภาคิยานิ สํโยชนานิ. (2)}

\prose{ตตฺถ กตมานิ ปญฺจ มจฺฉริยานิ, อาวาสมจฺฉริยํ กุลมจฺฉริยํ ลาภมจฺฉริยํ วณฺณมจฺฉริยํ ธมฺมมจฺฉริยํ, อิมานิ ปญฺจ มจฺฉริยานิ. (3)}

\prose{ตตฺถ กตเม ปญฺจ สงฺคา, ราคสงฺโค โทสสงฺโค โมหสงฺโค มานสงฺโค ทิฏฺฐิสงฺโค, อิเม ปญฺจ สงฺคา. (4)}

\prose{ตตฺถ กตเม ปญฺจ สลฺลา, ราคสลฺลํ โทสสลฺลํ โมหสลฺลํ มานสลฺลํ ทิฏฺฐิสลฺลํ, อิเม ปญฺจ สลฺลา. (5)}

\proseitem{941}{ตตฺถ กตเม ปญฺจ เจโตขิลา, สตฺถริ กงฺขติ วิจิกิจฺฉติ นาธิมุจฺจติ น สมฺปสีทติ, ธมฺเม กงฺขติ วิจิกิจฺฉติ นาธิมุจฺจติ น สมฺปสีทติ, \csromanfolio{392}สํเฆ กงฺขติ วิจิกิจฺฉติ นาธิมุจฺจติ น สมฺปสีทติ, สิกฺขาย กงฺขติ วิจิกิจฺฉติ นาธิมุจฺจติ น สมฺปสีทติ, สพฺรหฺมจารีสุ กุปิโต โหติ อนตฺตมโน อาหตจิตฺโต ขิลชาโต, อิเม ปญฺจ เจโตขิลา. (6)}

\prose{ตตฺถ กตเม ปญฺจ เจตโสวินิพนฺธา, กาเม อวีตราโค โหติ อวิคตจฺฉนฺโท อวิคตเปโม อวิคตปิปาโส อวิคตปริฬาโห อวิคตตโณฺห, กาเย อวีตราโค โหติ.\csromanspacer{} รูเป อวีตราโค โหติ.\csromanspacer{} ยาวทตฺถํ อุทราวเทหกํ ภุญฺชิตฺวา เสยฺยสุขํ ปสฺสสุขํ มิทฺธสุขํ อนุยุตฺโต วิหรติ, อญฺญตรํ เทวนิกายํ ปณิธาย พฺรหฺมจริยํ จรติ “อิมินาหํ สีเลน วา วเตน วา ตเปน วา พฺรหฺมจริเยน วา เทโว วา ภวิสฺสามิ เทวญฺญตโร วา”ติ, อิเม ปญฺจ เจตโสวินิพนฺธา. (7)}

\prose{ตตฺถ กตมานิ ปญฺจ นีวรณานิ, กามจฺฉนฺทนีวรณํ พฺยาปาทนีวรณํ ถินมิทฺธนีวรณํ อุทฺธจฺจกุกฺกุจฺจนีวรณํ วิจิกิจฺฉานีวรณํ, อิมานิ ปญฺจ นีวรณานิ. (8)}

\prose{ตตฺถ กตมานิ ปญฺจ กมฺมานิ อานนฺตริกานิ, มาตา ชีวิตา โวโรปิตา โหติ, ปิตา ชีวิตา โวโรปิโต โหติ, อรหนฺโต ชีวิตา โวโรปิโต โหติ, ทุฏฺเฐน จิตฺเตน ตถาคตสฺส โลหิตํ อุปฺปาทิตํ โหติ, สํโฆ ภินฺโน โหติ, อิมานิ ปญฺจ กมฺมานิ อานนฺตริกานิ. (9)}

\prose{ตตฺถ กตมา ปญฺจ ทิฏฺฐิโย, “สญฺญี อตฺตา โหติ อโรโค ปรํ มรณา”ติ อิตฺเถเก อภิวทนฺติ, “อสญฺญี อตฺตา โหติ อโรโค ปรํ มรณา”ติ อิตฺเถเก อภิวทนฺติ, “เนวสญฺญีนาสญฺญี อตฺตา โหติ อโรโค ปรํ มรณา”ติ อิตฺเถเก อภิวทนฺติ, สโต วา ปน สตฺตสฺส อุจฺเฉทํ วินาสํ วิภวํ ปญฺญเปนฺติ ทิฏฺฐธมฺมนิพฺพานํ วา ปเนเก อภิวทนฺติ, อิมา ปญฺจ ทิฏฺฐิโย. (10)}

\proseitem{942}{ตตฺถ กตเม ปญฺจ เวรา, ปาณาติปาโต อทินฺนาทานํ กาเมสุมิจฺฉาจาโร มุสาวาโท สุราเมรยมชฺชปมาทฏฺฐานา, อิเม ปญฺจ เวรา. (11)}

\csromanheader{2}{17. ขุทฺทกวตฺถุวิภงฺค}
\prose{\csromanfolio{393}ตตฺถ กตเม ปญฺจ พฺยสนา, ญาติพฺยสนํ โภคพฺยสนํ โรคพฺยสนํ สีลพฺยสนํ ทิฏฺฐิพฺยสนํ, อิเม ปญฺจ พฺยสนา. (12)}

\prose{ตตฺถ กตเม ปญฺจ อกฺขนฺติยา อาทีนวา, พหุโน ชนสฺส อปฺปิโย โหติ อมนาโป, เวรพหุโล จ โหติ วชฺชพหุโล จ, สมฺมูโฬฺห กาลํ กโรติ, กายสฺส เภทา ปรํ มรณา อปายํ ทุคฺคติํ วินิปาตํ นิรยํ อุปปชฺชติ, อิเม ปญฺจ อกฺขนฺติยา อาทีนวา. (13)}

\prose{ตตฺถ กตมานิ ปญฺจ ภยานิ, อาชีวกภยํ อสิโลกภยํ ปริสสารชฺชภยํ มรณภยํ ทุคฺคติภยํ, อิมานิ ปญฺจ ภยานิ. (14)}

\proseitem{943}{ตตฺถ กตเม ปญฺจ ทิฏฺฐธมฺมนิพฺพานวาทา, อิเธกจฺโจ สมโณ วา พฺราหฺมโณ วา เอวํวาที โหติ เอวํทิฏฺฐี\footnote{เอวํทิฏฺฐิ (สฺยา)} “ยโต โข โภ อยํ อตฺตา ปญฺจหิ กามคุเณหิ สมปฺปิโต สมงฺคีภูโต ปริจาเรติ, เอตฺตาวตา โข โภ อยํ อตฺตา ปรมทิฏฺฐธมฺมนิพฺพานปฺปตฺโต โหตี”ติ, อิตฺเถเก สโต สตฺตสฺส ปรมทิฏฺฐธมฺมนิพฺพานํ ปญฺญเปนฺติ.}

\prose{ตมญฺโญ เอวมาห “อตฺถิ โข โภ เอโส อตฺตา, ยํ ตฺวํ วเทสิ, ‘เนโส นตฺถี’ติ วทามิ.\csromanspacer{} โน จ โข โภ อยํ อตฺตา เอตฺตาวตา ปรมทิฏฺฐธมฺมนิพฺพานปฺปตฺโต โหติ.\csromanspacer{} ตํ กิสฺส เหตุ, กามา หิ โภ อนิจฺจา ทุกฺขา วิปริณามธมฺมา, เตสํ วิปริณามญฺญถาภาวา อุปฺปชฺชนฺติ, โสกปริเทวทุกฺขโทมนสฺสุปายาสา, ยโต โข โภ อยํ อตฺตา วิวิจฺเจว กาเมหิ ฯเปฯ ปฐมํ ฌานํ อุปสมฺปชฺช วิหรติ.\csromanspacer{} เอตฺตาวตา โข โภ อยํ อตฺตา ปรมทิฏฺฐธมฺมนิพฺพานปฺปตฺโต โหตี”ติ, อิตฺเถเก สโต สตฺตสฺส ปรมทิฏฺฐธมฺมนิพฺพานํ ปญฺญเปนฺติ.}

\prose{ตมญฺโญ เอวมาห “อตฺถิ โข โภ เอโส อตฺตา, ยํ ตฺวํ วเทสิ, ‘เนโส นตฺถี’ติ วทามิ.\csromanspacer{} โน จ โข โภ อยํ อตฺตา เอตฺตาวตา ปรมทิฏฺฐธมฺมนิพฺพานปฺปตฺโต โหติ.\csromanspacer{} ตํ กิสฺส เหตุ, ยเทว ตตฺถ วิตกฺกิตํ วิจาริตํ, เอเตน เอตํ โอฬาริกํ อกฺขายติ, ยโต โข โภ อยํ อตฺตา วิตกฺกวิจารานํ วูปสมา ฯเปฯ ทุติยํ ฌานํ อุปสมฺปชฺช วิหรติ.\csromanspacer{} เอตฺตาวตา โข โภ อยํ อตฺตา ปรมทิฏฺฐธมฺมนิพฺพานปฺปตฺโต โหตี”ติ, อิตฺเถเก สโต สตฺตสฺส ปรมทิฏฺฐธมฺมนิพฺพานํ ปญฺญเปนฺติ}

\prose{\csromanfolio{394}ตมญฺโญ เอวมาห “อตฺถิ โข โภ เอโส อตฺตา, ยํ ตฺวํ วเทสิ, ‘เนโส นตฺถี’ติ วทามิ.\csromanspacer{} โน จ โข โภ อยํ อตฺตา เอตฺตาวตา ปรมทิฏฺฐธมฺมนิพฺพานปฺปตฺโต โหติ.\csromanspacer{} ตํ กิสฺส เหตุ, ยเทว ตตฺถ ปีติคตํ เจตโส อุปฺปิลาวิตํ, เอเตน เอตํ โอฬาริกํ อกฺขายติ, ยโต โข โภ อยํ อตฺตา ปีติยา จ วิราคา ฯเปฯ ตติยํ ฌานํ อุปสมฺปชฺช วิหรติ.\csromanspacer{} เอตฺตาวตา โข โภ อยํ อตฺตา ปรมทิฏฺฐธมฺมนิพฺพานปฺปตฺโต โหตี”ติ, อิตฺเถเก สโต สตฺตสฺส ปรมทิฏฺฐธมฺมนิพฺพานํ ปญฺญเปนฺติ.}

\prose{ตมญฺโญ เอวมาห “อตฺถิ โข โภ เอโส อตฺตา, ยํ ตฺวํ วเทสิ, ‘เนโส นตฺถี’ติ วทามิ.\csromanspacer{} โน จ โข โภ อยํ อตฺตา เอตฺตาวตา ปรมทิฏฺฐธมฺมนิพฺพานปฺปตฺโต โหติ.\csromanspacer{} ตํ กิสฺส เหตุ, ยเทว ตตฺถ สุขปีติ เจตโส อาโภโค, เอเตน เอตํ โอฬาริกํ อกฺขายติ, ยโต โข โภ อยํ อตฺตา สุขสฺส จ ปหานา ฯเปฯ จตุตฺถํ ฌานํ อุปสมฺปชฺช วิหรติ.\csromanspacer{} เอตฺตาวตา โข โภ อยํ อตฺตา ปรมทิฏฺฐธมฺมนิพฺพานปฺปตฺโต โหตี”ติ, อิตฺเถเก สโต สตฺตสฺส ปรมทิฏฺฐธมฺมนิพฺพานํ ปญฺญเปนฺติ.\csromanspacer{} อิเม ปญฺจ ทิฏฺฐธมฺมนิพฺพานวาทา. (15)}

\csromanheader{2}{ปญฺจกํ.}
\csromansectionrule
\csromanmatikaanchor{mtk.184}
\csromantocmarkref{subparagraph}{6. ฉกฺกนิทฺเทส}{mtk.184}
\bookmark[dest={mtk.184},level=5]{6. ฉกฺกนิทฺเทส}
\csromanheader{6}{6. ฉกฺกนิทฺเทส}
\proseitem{944}{ตตฺถ กตมานิ ฉ วิวาทมูลานิ, โกโธ มกฺโข อิสฺสา สาเฐยฺยํ ปาปิจฺฉตา สนฺทิฏฺฐิปรามาสิตา, อิมานิ ฉ วิวาทมูลานิ. (1)}

\prose{ตตฺถ กตเม ฉ ฉนฺทราคา, ฉนฺทราคา เคหสิตา ธมฺมา มนาปิเยสุ รูเปสุ เคหสิโต ราโค สาราโค จิตฺตสฺส สาราโค, มนาปิเยสุ สทฺเทสุ ฯเปฯ มนาปิเยสุ คนฺเธสุ ฯเปฯ มนาปิเยสุ รเสสุ ฯเปฯ มนาปิเยสุ โผฏฺฐพฺเพสุ ฯเปฯ มนาปิเยสุ ธมฺเมสุ เคหสิโต ราโค สาราโค จิตฺตสฺส สาราโค, อิเม ฉ ฉนฺทราคา. (2)}

\prose{ตตฺถ กตมานิ ฉ วิโรธวตฺถูนิ, อมนาปิเยสุ รูเปสุ จิตฺตสฺส อาฆาโต ปฏิฆาโต จณฺฑิกฺกํ อสุโรโป อนตฺตมนตา จิตฺตสฺส, อมนาปิเยสุ สทฺเทสุ ฯเปฯ อมนาปิเยสุ คนฺเธสุ ฯเปฯ อมนาปิเยสุ รเสสุ ฯเปฯ อมนาปิเยสุ โผฏฺฐพฺเพสุ ฯเปฯ อมนาปิเยสุ ธมฺเมสุ}

\vspace{\tipitakaparskip}
\csromancenter{ขุทฺทกวตฺถุวิภงฺค จิตฺตสฺส อาฆาโต ปฏิฆาโต จณฺฑิกฺกํ อสุโรโป อนตฺตมนตา จิตฺตสฺส, อิมานิ ฉ วิโรธวตฺถูนิ. (3)}
\prose{\csromanfolio{395}ตตฺถ กตเม ฉ ตณฺหากายา, รูปตณฺหา สทฺทตณฺหา คนฺธตณฺหา รสตณฺหา โผฏฺฐพฺพตณฺหา ธมฺมตณฺหา, อิเม ฉ ตณฺหากายา. (4)}

\proseitem{945}{ตตฺถ กตเม ฉ อคารวา, สตฺถริ อคารโว วิหรติ อปฺปติสฺโส, ธมฺเม อคารโว วิหรติ อปฺปติสฺโส, สํเฆ อคารโว วิหรติ อปฺปติสฺโส, สิกฺขาย อคารโว วิหรติ อปฺปติสฺโส, อปฺปมาเท อคารโว วิหรติ อปฺปติสฺโส, ปฏิสนฺถาเร อคารโว วิหรติ อปฺปติสฺโส, อิเม ฉ อคารวา. (5)}

\prose{ตตฺถ กตเม ฉ ปริหานิยา ธมฺมา, กมฺมารามตา ภสฺสารามตา นิทฺทารามตา สงฺคณิการามตา สํสคฺคารามตา ปปญฺจารามตา, อิเม ฉ ปริหานิยา ธมฺมา. (6)}

\proseitem{946}{ตตฺถ กตเม อปเรปิ ฉ ปริหานิยา ธมฺมา, กมฺมารามตา ภสฺสารามตา นิทฺทารามตา สงฺคณิการามตา โทวจสฺสตา ปาปมิตฺตตา, อิเม ฉ ปริหานิยา ธมฺมา. (7)}

\prose{ตตฺถ กตเม ฉ โสมนสฺสุปวิจารา, จกฺขุนา รูปํ ทิสฺวา โสมนสฺสฏฺฐานิยํ รูปํ อุปวิจรติ.\csromanspacer{} โสเตน สทฺทํ สุตฺวา ฯเปฯ ฆาเนน คนฺธํ ฆายิตฺวา ฯเปฯ ชิวฺหาย รสํ สายิตฺวา ฯเปฯ กาเยน โผฏฺฐพฺพํ ผุสิตฺวา ฯเปฯ มนสา ธมฺมํ วิญฺญาย โสมนสฺสฏฺฐานิยํ ธมฺมํ อุปวิจรติ, อิเม ฉ โสมนสฺสุปวิจารา. (8)}

\prose{ตตฺถ กตเม ฉ โทมนสฺสุปวิจารา, จกฺขุนา รูปํ ทิสฺวา โทมนสฺสฏฺฐานิยํ รูปํ อุปวิจรติ.\csromanspacer{} โสเตน สทฺทํ สุตฺวา ฯเปฯ ฆาเนน คนฺธํ ฆายิตฺวา ฯเปฯ ชิวฺหาย รสํ สายิตฺวา ฯเปฯ กาเยน โผฏฺฐพฺพํ ผุสิตฺวา ฯเปฯ มนสา ธมฺมํ วิญฺญาย โทมนสฺสฏฺฐานิยํ ธมฺมํ อุปวิจรติ, อิเม ฉ โทมนสฺสุปวิจารา. (9)}

\prose{ตตฺถ กตเม ฉ อุเปกฺขุปวิจารา, จกฺขุนา รูปํ ทิสฺวา อุเปกฺขาฏฺฐานิยํ รูปํ อุปวิจรติ.\csromanspacer{} โสเตน สทฺทํ สุตฺวา ฯเปฯ ฆาเนน คนฺธํ ฆายิตฺวา ฯเปฯ ชิวฺหาย รสํ สายิตฺวา ฯเปฯ กาเยน โผฏฺฐพฺพํ ผุสิตฺวา ฯเปฯ มนสา ธมฺมํ วิญฺญาย อุเปกฺขาฏฺฐานิยํ ธมฺมํ อุปวิจรติ, อิเม ฉ อุเปกฺขุปวิจารา. (10)}

\proseitem{947}{\csromanfolio{396}ตตฺถ กตมานิ ฉ เคหสิตานิ โสมนสฺสานิ, มนาปิเยสุ รูเปสุ เคหสิตํ เจตสิกํ สาตํ เจตสิกํ สุขํ เจโตสมฺผสฺสชํ สาตํ สุขํ เวทยิตํ เจโตสมฺผสฺสชา สาตา สุขา เวทนา, มนาปิเยสุ สทฺเทสุ ฯเปฯ มนาปิเยสุ คนฺเธสุ ฯเปฯ มนาปิเยสุ รเสสุ ฯเปฯ มนาปิเยสุ โผฏฺฐพฺเพสุ ฯเปฯ มนาปิเยสุ ธมฺเมสุ เคหสิตํ เจตสิกํ สาตํ เจตสิกํ สุขํ เจโตสมฺผสฺสชํ สาตํ สุขํ เวทยิตํ เจโตสมฺผสฺสชา สาตา สุขา เวทนา, อิมานิ ฉ เคหสิตานิ โสมนสฺสานิ. (11)}

\prose{ตตฺถ กตมานิ ฉ เคหสิตานิ โทมนสฺสานิ, อมนาปิเยสุ รูเปสุ เคหสิตํ เจตสิกํ อสาตํ เจตสิกํ ทุกฺขํ เจโตสมฺผสฺสชํ อสาตํ ทุกฺขํ เวทยิตํ เจโตสมฺผสฺสชา อสาตา ทุกฺขา เวทนา, อมนาปิเยสุ สทฺเทสุ ฯเปฯ อมนาปิเยสุ คนฺเธสุ ฯเปฯ อมนาปิเยสุ รเสสุ ฯเปฯ อมนาปิเยสุ โผฏฺฐพฺเพสุ ฯเปฯ อมนาปิเยสุ ธมฺเมสุ เคหสิตํ เจตสิกํ อสาตํ เจตสิกํ ทุกฺขํ เจโตสมฺผสฺสชํ อสาตํ ทุกฺขํ เวทยิตํ เจโตสมฺผสฺสชา อสาตา ทุกฺขา เวทนา, อิมานิ ฉ เคหสิตานิ โทมนสฺสานิ. (12)}

\prose{ตตฺถ กตมา ฉ เคหสิตา อุเปกฺขา, อุเปกฺขาฏฺฐานิเยสุ รูเปสุ เคหสิตํ เจตสิกํ เนว สาตํ นาสาตํ เจโตสมฺผสฺสชํ อทุกฺขมสุขํ เวทยิตํ เจโตสมฺผสฺสชา อทุกฺขมสุขา เวทนา, อุเปกฺขาฏฺฐานิเยสุ สทฺเทสุ ฯเปฯ อุเปกฺขาฏฺฐานิเยสุ คนฺเธสุ ฯเปฯ อุเปกฺขาฏฺฐานิเยสุ รเสสุ ฯเปฯ อุเปกฺขาฏฺฐานิเยสุ โผฏฺฐพฺเพสุ ฯเปฯ อุเปกฺขาฏฺฐานิเยสุ ธมฺเมสุ เคหสิตํ เจตสิกํ เนว สาตํ นาสาตํ เจโตสมฺผสฺสชํ อทุกฺขมสุขํ เวทยิตํ เจโตสมฺผสฺสชา อทุกฺขมสุขา เวทนา, อิมา ฉ เคหสิตา อุเปกฺขา. (13)}

\proseitem{948}{ตตฺถ กตมา ฉ ทิฏฺฐิโย, “อตฺถิ เม อตฺตา”ติ วา อสฺส สจฺจโต เถตโต ทิฏฺฐิ อุปฺปชฺชติ, “นตฺถิ เม อตฺตา”ติ วา อสฺส สจฺจโต เถตโต ทิฏฺฐิ อุปฺปชฺชติ, “อตฺตนา วา อตฺตานํ สญฺชานามี”ติ วา อสฺส สจฺจโต เถตโต ทิฏฺฐิ อุปฺปชฺชติ, “อตฺตนา วา อนตฺตานํ สญฺชานามี”ติ วา อสฺส สจฺจโต เถตโต ทิฏฺฐิ อุปฺปชฺชติ, “อนตฺตนา วา อตฺตานํ สญฺชานามี”ติ วา อสฺส สจฺจโต เถตโต ทิฏฺฐิ อุปฺปชฺชติ, อถ วา}

\vspace{\tipitakaparskip}
\csromancenter{ขุทฺทกวตฺถุวิภงฺค ปนสฺส เอวํทิฏฺฐี โหติ “โส เม อยํ อตฺตา วโท เวเทยฺโย, ตตฺร ตตฺร ทีฆรตฺตํ กลฺยาณปาปกานํ กมฺมานํ วิปากํ ปจฺจนุโภติ, น โส ชาโต นาโหสิ, น โส ชาโต น ภวิสฺสติ, “นิจฺโจ ธุโว สสฺสโต อวิปริณามธมฺโม”ติ วา ปนสฺส สจฺจโต เถตโต ทิฏฺฐิ อุปฺปชฺชติ, อิมา ฉ ทิฏฺฐิโย. (14)}
\csromanheader{2}{ฉกฺกํ.}
\csromansectionrule
\csromanmatikaanchor{mtk.185}
\csromantocmarkref{subparagraph}{7. สตฺตกนิทฺเทส}{mtk.185}
\bookmark[dest={mtk.185},level=5]{7. สตฺตกนิทฺเทส}
\csromanheader{6}{7. สตฺตกนิทฺเทส}
\proseitem{949}{\csromanfolio{397}ตตฺถ กตเม สตฺตานุสยา, กามราคานุสโย ปฏิฆานุสโย มานานุสโย ทิฏฺฐานุสโย วิจิกิจฺฉานุสโย ภวราคานุสโย อวิชฺชานุสโย, อิเม สตฺต อนุสยา. (1)}

\prose{ตตฺถ กตมานิ สตฺต สํโยชนานิ, กามราคสํโยชนํ ปฏิฆสํโยชนํ มานสํโยชนํ ทิฏฺฐิสํโยชนํ วิจิกิจฺฉาสํโยชนํ ภวราคสํโยชนํ อวิชฺชาสํโยชนํ, อิมานิ สตฺต สํโยชนานิ. (2)}

\prose{ตตฺถ กตมานิ สตฺต ปริยุฏฺฐานานิ, กามราคปริยุฏฺฐานํ ปฏิฆปริยุฏฺฐานํ มานปริยุฏฺฐานํ ทิฏฺฐิปริยุฏฺฐานํ วิจิกิจฺฉาปริยุฏฺฐานํ ภวราคปริยุฏฺฐานํ อวิชฺชาปริยุฏฺฐานํ, อิมานิ สตฺต ปริยุฏฺฐานานิ. (3)}

\proseitem{950}{ตตฺถ กตเม สตฺต อสทฺธมฺมา, อสฺสทฺโธ โหติ, อหิริโก โหติ, อโนตฺตปฺปี โหติ, อปฺปสฺสุโต โหติ, กุสีโต โหติ, มุฏฺฐสฺสตี โหติ, ทุปฺปญฺโญ โหติ, อิเม สตฺต อสทฺธมฺมา. (4)}

\prose{ตตฺถ กตมานิ สตฺต ทุจฺจริตานิ, ปาณาติปาโต อทินฺนาทานํ กาเมสุมิจฺฉาจาโร มุสาวาโท ปิสุณา วาจา ผรุสา วาจา สมฺผปฺปลาโป, อิมานิ สตฺต ทุจฺจริตานิ. (5)}

\prose{ตตฺถ กตเม สตฺต มานา, มาโน อติมาโน มานาติมาโน โอมาโน อธิมาโน อสฺมิมาโน มิจฺฉามาโน, อิเม สตฺต มานา. (6)}

\proseitem{951}{ตตฺถ กตมา สตฺต ทิฏฺฐิโย, อิเธกจฺโจ สมโณ วา พฺราหฺมโณ วา เอวํวาที โหติ เอวํทิฏฺฐี “ยโต โข โภ อยํ อตฺตา \csromanfolio{398}รูปี จาตุมหาภูติโก\footnote{จาตุมฺมหาภูติโก (สี, สฺยา)} มาตาเปตฺติกสมฺภโว กายสฺส เภทา อุจฺฉิชฺชติ วินสฺสติ, น โหติ ปรํ มรณา, เอตฺตาวตา โข โภ อยํ อตฺตา สมฺมา สมุจฺฉินฺโน โหตี”ติ, อิตฺเถเก สโต สตฺตสฺส อุจฺเฉทํ วินาสํ วิภวํ ปญฺญเปนฺติ. (6-1)}

\prose{ตมญฺโญ เอวมาห “อตฺถิ โข โภ เอโส อตฺตา, ยํ ตฺวํ วเทสิ, ‘เนโส นตฺถี’ติ วทามิ.\csromanspacer{} โน จ โข โภ อยํ อตฺตา เอตฺตาวตา สมฺมา สมุจฺฉินฺโน โหติ.\csromanspacer{} อตฺถิ โข โภ อญฺโญ อตฺตา ทิพฺโพ รูปี กามาวจโร กพฬีการภกฺโข, ตํ ตฺวํ น ชานาสิ น ปสฺสสิ, ตมหํ ชานามิ ปสฺสามิ.\csromanspacer{} โส โข โภ อตฺตา ยโต กายสฺส เภทา อุจฺฉิชฺชติ วินสฺสติ, น โหติ ปรํ มรณา, เอตฺตาวตา โข โภ อยํ อตฺตา สมฺมา สมุจฺฉินฺโน โหตี”ติ, อิตฺเถเก โภ สตฺตสฺส อุจฺเฉทํ วินาสํ วิภวํ ปญฺญเปนฺติ. (6-2)}

\prose{ตมญฺโญ เอวมาห “อตฺถิ โข โภ เอโส อตฺตา, ยํ ตฺวํ วเทสิ, ‘เนโส นตฺถี’ติ วทามิ.\csromanspacer{} โน จ โข โภ อยํ เอตฺตาวตา สมฺมา สมุจฺฉินฺโน โหติ.\csromanspacer{} อตฺถิ โข โภ อญฺโญ อตฺตา ทิพฺโพ รูปี มโนมโย สพฺพงฺคปจฺจงฺคี อหีนินฺทฺริโย, ตํ ตฺวํ น ชานาสิ น ปสฺสสิ, ตมหํ ชานามิ ปสฺสามิ.\csromanspacer{} โส โข โภ อตฺตา ยโต กายสฺส เภทา อุจฺฉิชฺชติ วินสฺสติ, น โหติ ปรํ มรณา, เอตฺตาวตา โข โภ อยํ อตฺตา สมฺมา สมุจฺฉินฺโน โหตี”ติ, อิตฺเถเก สโต สตฺตสฺส อุจฺเฉทํ วินาสํ วิภวํ ปญฺญเปนฺติ. (6-3)}

\prose{ตมญฺโญ เอวมาห “อตฺถิ โข โภ เอโส อตฺตา, ยํ ตฺวํ วเทสิ, ‘เนโส นตฺถี’ติ วทามิ.\csromanspacer{} โน จ โข โภ อยํ อตฺตา เอตฺตาวตา สมฺมา สมุจฺฉินฺโน โหติ.\csromanspacer{} อตฺถิ โข โภ อญฺโญ อตฺตา สพฺพโส รูปสญฺญานํ สมติกฺกมา ปฏิฆสญฺญานํ อตฺถงฺคมา นานตฺตสญฺญานํ อมนสิการา ‘อนนฺโต อากาโส’ติ อากาสานญฺจายตนูปโค, ตํ ตฺวํ น ชานาสิ น ปสฺสสิ, ตมหํ ชานามิ ปสฺสามิ.\csromanspacer{} โส โข โภ อตฺตา ยโต กายสฺส เภทา อุจฺฉิชฺชติ วินสฺสติ, น โหติ ปรํ มรณา, เอตฺตาวตา โข โภ อยํ อตฺตา สมฺมา สมุจฺฉินฺโน โหตี”ติ, อิตฺเถเก สโต สตฺตสฺส อุจฺเฉทํ วินาสํ วิภวํ ปญฺญเปนฺติ. (6-4)}

\csromanheader{2}{17. ขุทฺทกวตฺถุวิภงฺค}
\prose{\csromanfolio{399}ตมญฺโญ เอวมาห “อตฺถิ โข โภ เอโส อตฺตา, ยํ ตฺวํ วเทสิ, ‘เนโส นตฺถี’ติ วทามิ.\csromanspacer{} โน จ โข โภ อยํ อตฺตา เอตฺตาวตา สมฺมา สมุจฺฉินฺโน โหติ.\csromanspacer{} อตฺถิ โข โภ อญฺโญ อตฺตา สพฺพโส อากาสานญฺจายตนํ สมติกฺกมฺม ‘อนนฺตํ วิญฺญาณนฺ’ติ วิญฺญาณญฺจายตนูปโค, ตํ ตฺวํ น ชานาสิ น ปสฺสสิ, ตมหํ ชานามิ ปสฺสามิ.\csromanspacer{} โส โข โภ อตฺตา ยโต กายสฺส เภทา อุจฺฉิชฺชติ วินสฺสติ, น โหติ ปรํ มรณา, เอตฺตาวตา โข โภ อยํ อตฺตา สมฺมา สมุจฺฉินฺโน โหตี”ติ, อิตฺเถเก สโต สตฺตสฺส อุจฺเฉทํ วินาสํ วิภวํ ปญฺญเปนฺติ. (6-5)}

\prose{ตมญฺโญ เอวมาห “อตฺถิ โข โภ เอโส อตฺตา, ยํ ตฺวํ วเทสิ, ‘เนโส นตฺถี’ติ วทามิ.\csromanspacer{} โน จ โข โภ อยํ อตฺตา เอตฺตาวตา สมฺมา สมุจฺฉินฺโน โหติ.\csromanspacer{} อตฺถิ โข โภ อญฺโญ อตฺตา สพฺพโส วิญฺญาณญฺจายตนํ สมติกฺกมฺม ‘นตฺถิ กิญฺจี’ติ อากิญฺจญฺญายตนูปโค, ตํ ตฺวํ น ชานาสิ น ปสฺสสิ, ตมหํ ชานามิ ปสฺสามิ.\csromanspacer{} โส โข โภ อตฺตา ยโต กายสฺส เภทา อุจฺฉิชฺชติ วินสฺสติ, น โหติ ปรํ มรณา, เอตฺตาวตา โข โภ อยํ อตฺตา สมฺมา สมุจฺฉินฺโน โหตี”ติ, อิตฺเถเก สโต สตฺตสฺส อุจฺเฉทํ วินาสํ วิภวํ ปญฺญเปนฺติ. (6-6)}

\prose{ตมญฺโญ เอวมาห “อตฺถิ โข ปน เอโส อตฺตา, ยํ ตฺวํ วเทสิ, ‘เนโส นตฺถี’ติ วทามิ.\csromanspacer{} โน จ โข โภ อยํ อตฺตา เอตฺตาวตา สมฺมา สมุจฺฉินฺโน โหติ.\csromanspacer{} อตฺถิ โข โภ อญฺโญ อตฺตา สพฺพโส อากิญฺจญฺญายตนํ สมติกฺกมฺม เนวสญฺญานาสญฺญายตนูปโค.\csromanspacer{} ตํ ตฺวํ น ชานาสิ น ปสฺสสิ, ตมหํ ชานามิ ปสฺสามิ.\csromanspacer{} โส โข โภ อตฺตา ยโต กายสฺส เภทา อุจฺฉิชฺชติ วินสฺสติ, น โหติ ปรํ มรณา, เอตฺตาวตา โข โภ อยํ อตฺตา สมฺมา สมุจฺฉินฺโน โหตี”ติ, อิตฺเถเก สโต สตฺตสฺส อุจฺเฉทํ วินาสํ วิภวํ ปญฺญเปนฺติ.\csromanspacer{} อิมา สตฺต ทิฏฺฐิโย. (7-7)}

\csromanheader{2}{สตฺตกํ.}
\csromansectionrule
\csromanmatikaanchor{mtk.186}
\csromantocmarkref{subparagraph}{8. อฏฺฐกนิทฺเทส}{mtk.186}
\bookmark[dest={mtk.186},level=5]{8. อฏฺฐกนิทฺเทส}
\csromanheader{6}{8. อฏฺฐกนิทฺเทส}
\proseitem{952}{ตตฺถ กตมานิ อฏฺฐ กิเลสวตฺถูนิ, โลโภ โทโส โมโห มาโน ทิฏฺฐิ วิจิกิจฺฉา ถินํ อุทฺธจฺจํ, อิมานิ อฏฺฐ กิเลสวตฺถูนิ.}

\proseitem{953}{\csromanfolio{400}ตตฺถ กตมานิ อฏฺฐ กุสีตวตฺถูนิ, อิธ ภิกฺขุนา กมฺมํ กาตพฺพํ โหติ, ตสฺส เอวํ โหติ “กมฺมํ โข เม กาตพฺพํ ภวิสฺสติ, กมฺมํ โข ปน เม กโรนฺตสฺส กาโย กิลมิสฺสติ, หนฺทาหํ นิปชฺชามี”ติ.\csromanspacer{} โส นิปชฺชติ น วีริยํ อารภติ อปฺปตฺตสฺส ปตฺติยา อนธิคตสฺส อธิคมาย อสจฺฉิกตสฺส สจฺฉิกิริยาย, อิทํ ปฐมํ กุสีตวตฺถุ.}

\prose{ปุน จปรํ ภิกฺขุนา กมฺมํ กตํ โหติ, ตสฺส เอวํ โหติ “อหํ โข กมฺมํ อกาสิํ, กมฺมํ โข ปน เม กโรนฺตสฺส กาโย กิลนฺโต, หนฺทาหํ นิปชฺชามี”ติ.\csromanspacer{} โส นิปชฺชติ น วีริยํ อารภติ อปฺปตฺตสฺส ปตฺติยา อนธิคตสฺส อธิคมาย อสจฺฉิกตสฺส สจฺฉิกิริยาย, อิทํ ทุติยํ กุสีตวตฺถุ.}

\prose{ปุน จปรํ ภิกฺขุนา มคฺโค คนฺตพฺโพ โหติ, ตสฺส เอวํ โหติ “มคฺโค โข เม คนฺตพฺโพ ภวิสฺสติ, มคฺคํ โข ปน เม คจฺฉนฺตสฺส กาโย กิลมิสฺสติ, หนฺทาหํ นิปชฺชามี”ติ.\csromanspacer{} โส นิปชฺชติ น วีริยํ อารภติ อปฺปตฺตสฺส ปตฺติยา อนธิคตสฺส อธิคมาย อสจฺฉิกตสฺส สจฺฉิกิริยาย, อิทํ ตติยํ กุสีตวตฺถุ.}

\prose{ปุน จปรํ ภิกฺขุนา มคฺโค คโต โหติ, ตสฺส เอวํ โหติ “อหํ โข มคฺคํ อคมาสิํ, มคฺคํ โข ปน เม คจฺฉนฺตสฺส กาโย กิลนฺโต, หนฺทาหํ นิปชฺชามี”ติ.\csromanspacer{} โส นิปชฺชติ น วีริยํ อารภติ อปฺปตฺตสฺส ปตฺติยา อนธิคตสฺส อธิคมาย อสจฺฉิกตสฺส สจฺฉิกิริยาย, อิทํ จตุตฺถํ กุสีตวตฺถุ.}

\prose{ปุน จปรํ ภิกฺขุ คามํ วา นิคมํ วา ปิณฺฑาย จรนฺโต น ลภติ ลูขสฺส วา ปณีตสฺส วา โภชนสฺส ยาวทตฺถํ ปาริปูริํ, ตสฺส เอวํ โหติ “อหํ โข คามํ วา นิคมํ วา ปิณฺฑาย จรนฺโต นาลตฺถํ ลูขสฺส วา ปณีตสฺส วา โภชนสฺส ยาวทตฺถํ ปาริปูริํ, ตสฺส เม กาโย กิลนฺโต อกมฺมญฺโญ, หนฺทาหํ นิปชฺชามี”ติ.\csromanspacer{} โส นิปชฺชติ น วีริยํ อารภติ อปฺปตฺตสฺส ปตฺติยา อนธิคตสฺส อธิคมาย อสจฺฉิกตสฺส สจฺฉิกิริยาย, อิทํ ปญฺจมํ กุสีตวตฺถุ.}

\prose{ปุน จปรํ ภิกฺขุ คามํ วา นิคมํ วา ปิณฺฑาย จรนฺโต ลภติ ลูขสฺส วา ปณีตสฺส วา โภชนสฺส ยาวทตฺถํ ปาริปูริํ, ตสฺส เอวํ โหติ “อหํ โข คามํ วา นิคมํ วา ปิณฺฑาย จรนฺโต อลตฺถํ ลูขสฺส วา ปณีตสฺส}

\vspace{\tipitakaparskip}
\csromancenter{ขุทฺทกวตฺถุวิภงฺค วา โภชนสฺส ยาวทตฺถํ ปาริปูริํ, ตสฺส เม กาโย กิลนฺโต อกมฺมญฺโญ, มาสาจิตํ มญฺเญ หนฺทาหํ นิปชฺชามี”ติ.\csromanspacer{} โส นิปชฺชติ น วีริยํ อารภติ อปฺปตฺตสฺส ปตฺติยา อนธิคตสฺส อธิคมาย อสจฺฉิกตสฺส สจฺฉิกิริยาย, อิทํ ฉฏฺฐํ กุสีตวตฺถุ.}
\prose{\csromanfolio{401}ปุน จปรํ ภิกฺขุโน อุปฺปนฺโน โหติ อปฺปมตฺตโก อาพาโธ, ตสฺส เอวํ โหติ “อุปฺปนฺโน โข เม อยํ อปฺปมตฺตโก อาพาโธ, อตฺถิ กปฺโป นิปชฺชิตุํ, หนฺทาหํ นิปชฺชามี”ติ.\csromanspacer{} โส นิปชฺชติ น วีริยํ อารภติ อปฺปตฺตสฺส ปตฺติยา อนธิคตสฺส อธิคมาย อสจฺฉิกตสฺส สจฺฉิกิริยาย, อิทํ สตฺตมํ กุสีตวตฺถุ.}

\prose{ปุน จปรํ ภิกฺขุ คิลานา วุฏฺฐิโต\footnote{คิลานวุฏฺฐิโต (สทฺทนีภิ). อํ 2. 263 ปิฏฺเฐ ปาฬิยา ฏีกา ปสฺสิตพฺพา.} โหติ อจิรวุฏฺฐิโต เคลญฺญา, ตสฺส เอวํ โหติ “อหํ โข คิลานา วุฏฺฐิโต อจิรวุฏฺฐิโต เคลญฺญา, ตสฺส เม กาโย ทุพฺพโล อกมฺมญฺโญ, หนฺทาหํ นิปชฺชามี”ติ.\csromanspacer{} โส นิปชฺชติ น วีริยํ อารภติ อปฺปตฺตสฺส ปตฺติยา อนธิคตสฺส อธิคมาย อสจฺฉิกตสฺส สจฺฉิกิริยาย, อิทํ อฏฺฐมํ กุสีตวตฺถุ.\csromanspacer{} อิมานิ อฏฺฐ กุสีตวตฺถูนิ.}

\proseitem{954}{ตตฺถ กตเมสุ อฏฺฐสุ โลกธมฺเมสุ จิตฺตสฺส ปฏิฆาโต, ลาเภ สาราโค, อลาเภ ปฏิวิโรโธ, ยเส สาราโค, อยเส ปฏิวิโรโธ, ปสํสาย สาราโค, นินฺทาย ปฏิวิโรโธ, สุเข สาราโค, ทุกฺเข ปฏิวิโรโธ.\csromanspacer{} อิเมสุ อฏฺฐสุ โลกธมฺเมสุ จิตฺตสฺส ปฏิฆาโต.}

\proseitem{955}{ตตฺถ กตเม อฏฺฐ อนริยโวหารา, อทิฏฺเฐ ทิฏฺฐวาทิตา, อสฺสุเต สุตวาทิตา, อมุเต มุตวาทิตา, อวิญฺญาเต วิญฺญาตวาทิตา, ทิฏฺเฐ อทิฏฺฐวาทิตา, สุเต อสฺสุตวาทิตา, มุเต อมุตวาทิตา, วิญฺญาเต อวิญฺญาตวาทิตา, อิเม อฏฺฐ อนริยโวหารา.}

\proseitem{956}{ตตฺถ กตเม อฏฺฐ มิจฺฉตฺตา, มิจฺฉาทิฏฺฐิ มิจฺฉาสงฺกปฺโป มิจฺฉาวาจา มิจฺฉากมฺมนฺโต มิจฺฉา-อาชีโว มิจฺฉาวายาโม มิจฺฉาสติ มิจฺฉาสมาธิ, อิเม อฏฺฐ มิจฺฉตฺตา.}

\proseitem{957}{\csromanfolio{402}ตตฺถ กตเม อฏฺฐ ปุริสโทสา, อิธ ภิกฺขู ภิกฺขุํ อาปตฺติยา โจเทนฺติ, โส ภิกฺขุ ภิกฺขูหิ อาปตฺติยา โจทียมาโน\footnote{โจทิยมาโน (สี, สฺยา) อํ 3. 34 ปิฏฺเฐปิ.} “น สรามิ น สรามี”ติ อสฺสติยาว นิพฺเพเฐติ, อยํ ปฐโม ปุริสโทโส.}

\prose{ปุน จปรํ ภิกฺขู ภิกฺขุํ อาปตฺติยา โจเทนฺติ, โส ภิกฺขุ ภิกฺขูหิ อาปตฺติยา โจทียมาโน โจทกํเยว ปฏิปฺผรติ “กิํ นุ โข ตุยฺหํ พาลสฺส อพฺยตฺตสฺส ภณิเตน, ตฺวมฺปิ นาม มํ ภณิตพฺพํ มญฺญสี”ติ, อยํ ทุติโย ปุริสโทโส.}

\prose{ปุน จปรํ ภิกฺขู ภิกฺขุํ อาปตฺติยา โจเทนฺติ, โส ภิกฺขุ ภิกฺขูหิ อาปตฺติยา โจทียมาโน โจทกํเยว\footnote{โจทกสฺเสว (สฺยา) อํ 3. 35 ปิฏฺเฐปิ.} ปจฺจาโรเปติ “ตฺวมฺปิ โขสิ อิตฺถนฺนามํ อาปตฺติํ อาปนฺโน ตฺวํ ตาว ปฐมํ ปฏิกโรหี”ติ, อยํ ตติโย ปุริสโทโส.}

\prose{ปุน จปรํ ภิกฺขู ภิกฺขุํ อาปตฺติยา โจเทนฺติ, โส ภิกฺขุ ภิกฺขูหิ อาปตฺติยา โจทียมาโน อญฺเญนาญฺญํ ปฏิจรติ, พหิทฺธา กถํ อปนาเมติ, โกปญฺจ โทสญฺจ อปฺปจฺจยญฺจ ปาตุกโรติ, อยํ จตุตฺโถ ปุริสโทโส.}

\prose{ปุน จปรํ ภิกฺขู ภิกฺขุํ อาปตฺติยา โจเทนฺติ, โส ภิกฺขุ ภิกฺขูหิ อาปตฺติยา โจทียมาโน สํฆมชฺเฌ พาหาวิกฺเขปกํ ภณติ, อยํ ปญฺจโม ปุริสโทโส.}

\prose{ปุน จปรํ ภิกฺขู ภิกฺขุํ อาปตฺติยา โจเทนฺติ, โส ภิกฺขุ ภิกฺขูหิ อาปตฺติยา โจทียมาโน “เนวาหํ อาปนฺโนมฺหิ, น ปนาหํ อนาปนฺโนมฺหี”ติ ตุณฺหีภูโต สํฆํ วิเหเสติ, อยํ ฉฏฺโฐ ปุริสโทโส.}

\prose{ปุน จปรํ ภิกฺขู ภิกฺขุํ อาปตฺติยา โจเทนฺติ, โส ภิกฺขุ ภิกฺขูหิ อาปตฺติยา โจทียมาโน อนาทิยิตฺวา สํฆํ อนาทิยิตฺวา โจทกํ สาปตฺติโกว\footnote{อาปตฺติโกว (ก) อํ 3. 36 ปิฏฺเฐปิ.} เยนกามํ ปกฺกมติ, อยํ สตฺตโม ปุริสโทโส.}

\prose{ปุน จปรํ ภิกฺขู ภิกฺขุํ อาปตฺติยา โจเทนฺติ, โส ภิกฺขุ ภิกฺขูหิ อาปตฺติยา โจทียมาโน เอวมาห “กิํ นุ โข ตุเมฺห อายสฺมนฺโต}

\vspace{\tipitakaparskip}
\csromancenter{ขุทฺทกวตฺถุวิภงฺค อติพาฬฺหํ มยิ พฺยาวฏา, อิทานาหํ สิกฺขํ ปจฺจกฺขาย หีนายาวตฺติสฺสามี”ติ, โส สิกฺขํ ปจฺจกฺขาย หีนายาวตฺติตฺวา เอวมาห “อิทานิ โข ตุเมฺห อายสฺมนฺโต อตฺตมนา โหถา”ติ, อยํ อฏฺฐโม ปุริสโทโส.\csromanspacer{} อิเม อฏฺฐ ปุริสโทสา.}
\proseitem{958}{\csromanfolio{403}ตตฺถ กตเม อฏฺฐ อสญฺญีวาทา, รูปี อตฺตา โหติ อโรโค ปรํ มรณาติ “อสญฺญี”ติ นํ ปญฺญเปนฺติ.\csromanspacer{} อรูปี อตฺตา ฯเปฯ.\csromanspacer{} รูปี จ อรูปี จ ฯเปฯ.\csromanspacer{} เนวรูปีนารูปี ฯเปฯ.\csromanspacer{} อนฺตวา อตฺตา โหติ อโรโค ปรํ มรณาติ “อสญฺญี”ติ นํ ปญฺญเปนฺติ.\csromanspacer{} อนนฺตวา อตฺตา โหติ อโรโค ปรํ มรณาติ “อสญฺญี”ติ นํ ปญฺญเปนฺติ.\csromanspacer{} อนฺตวา จ อนนฺตวา จ อตฺตา โหติ อโรโค ปรํ มรณาติ “อสญฺญี”ติ นํ ปญฺญเปนฺติ.\csromanspacer{} เนวนฺตวา นานนฺตวา อตฺตา โหติ อโรโค ปรํ มรณาติ, “อสญฺญี”ติ นํ ปญฺญเปนฺติ.\csromanspacer{} อิเม อฏฺฐ อสญฺญีวาทา.}

\proseitem{959}{ตตฺถ กตเม อฏฺฐ เนวสญฺญีนาสญฺญีวาทา, รูปี อตฺตา โหติ อโรโค ปรํ มรณาติ “เนวสญฺญีนาสญฺญี”ติ นํ ปญฺญเปนฺติ.\csromanspacer{} อรูปี อตฺตา โหติ อโรโค ปรํ มรณาติ “เนวสญฺญีนาสญฺญี”ติ นํ ปญฺญเปนฺติ.\csromanspacer{} รูปี จ อรูปี จ อตฺตา โหติ อโรโค ปรํ มรณาติ “เนวสญฺญีนาสญฺญี”ติ นํ ปญฺญเปนฺติ.\csromanspacer{} เนวรูปีนารูปี อตฺตา โหติ อโรโค ปรํ มรณาติ “เนวสญฺญีนาสญฺญี”ติ นํ ปญฺญเปนฺติ.\csromanspacer{} อนฺตวา อตฺตา โหติ อโรโค ปรํ มรณาติ “เนวสญฺญีนาสญฺญี”ติ นํ ปญฺญเปนฺติ.\csromanspacer{} อนฺตวา อตฺตา โหติ อโรโค ปรํ มรณาติ “เนวสญฺญีนาสญฺญี”ติ นํ ปญฺญเปนฺติ.\csromanspacer{} อนฺตวา จ อนนฺตวา จ อตฺตา โหติ อโรโค ปรํ มรณาติ “เนวสญฺญีนาสญฺญี”ติ นํ ปญฺญเปนฺติ.\csromanspacer{} เนวนฺตวา นานนฺตวา อตฺตา โหติ อโรโค ปรํ มรณาติ “เนวสญฺญีนาสญฺญี”ติ นํ ปญฺญเปนฺติ, อิเม อฏฺฐ เนวสญฺญีนาสญฺญีวาทา.}

\csromanheader{2}{อฏฺฐกํ.}
\csromansectionrule
\csromanmatikaanchor{mtk.187}
\csromantocmarkref{subparagraph}{9. นวกนิทฺเทส}{mtk.187}
\bookmark[dest={mtk.187},level=5]{9. นวกนิทฺเทส}
\csromanheader{6}{9. นวกนิทฺเทส}
\proseitem{960}{ตตฺถ กตมานิ นว อาฆาตวตฺถูนิ, “อนตฺถํ เม อจรี”ติ อาฆาโต ชายติ, “อนตฺถํ เม จรตี”ติ อาฆาโต ชายติ, \csromanfolio{404}“อนตฺถํ เม จริสฺสตี”ติ อาฆาโต ชายติ, “ปิยสฺส เม มนาปสฺส อนตฺถํ อจริ ฯเปฯ อนตฺถํ จรติ ฯเปฯ อนตฺถํ จริสฺสตี”ติ อาฆาโต ชายติ, “อปฺปิยสฺส เม อมนาปสฺส อตฺถํ อจริ ฯเปฯ อตฺถํ จรติ ฯเปฯ อตฺถํ จริสฺสตี”ติ อาฆาโต ชายติ.\csromanspacer{} อิมานิ นว อาฆาตวตฺถูนิ.}

\proseitem{961}{ตตฺถ กตมานิ นว ปุริสมลานิ, โกโธ มกฺโข อิสฺสา มจฺฉริยํ มายา สาเฐยฺยํ มุสาวาโท ปาปิจฺฉตา มิจฺฉาทิฏฺฐิ, อิมานิ นว ปุริสมลานิ.}

\proseitem{962}{ตตฺถ กตเม นววิธา มานา, เสยฺยสฺส “เสยฺโยหมสฺมี”ติ มาโน, เสยฺยสฺส “สทิโสหมสฺมี”ติ มาโน, เสยฺยสฺส “หีโนหมสฺมี”ติ มาโน, สทิสสฺส “เสยฺโยหมสฺมี”ติ มาโน, สทิสสฺส “สทิโสหมสฺมี”ติ มาโน, สทิสสฺส “หีโนหมสฺมี”ติ มาโน, หีนสฺส “เสยฺโยหมสฺมี”ติ มาโน, หีนสฺส “สทิโสหมสฺมี”ติ มาโน, หีนสฺส “หีโนหมสฺมี”ติ มาโน.\csromanspacer{} อิเม นววิธา มานา.}

\proseitem{963}{ตตฺถ กตเม นว ตณฺหามูลกา ธมฺมา, ตณฺหํ ปฏิจฺจ ปริเยสนา, ปริเยสนํ ปฏิจฺจ ลาโภ, ลาภํ ปฏิจฺจ วินิจฺฉโย, วินิจฺฉยํ ปฏิจฺจ ฉนฺทราโค, ฉนฺทราคํ ปฏิจฺจ อชฺโฌสานํ, อชฺโฌสานํ ปฏิจฺจ ปริคฺคโห, ปริคฺคหํ ปฏิจฺจ มจฺฉริยํ, มจฺฉริยํ ปฏิจฺจ อารกฺโข, อารกฺขาธิกรณํ ทณฺฑาทานสตฺถาทานกลหวิคฺคหวิวาทตุวํตุวํเปสุญฺญมุสาวาทา อเนเก ปาปกา อกุสลา ธมฺมา สมฺภวนฺติ, อิเม นว ตณฺหามูลกา ธมฺมา.}

\proseitem{964}{ตตฺถ กตมานิ นว อิญฺชิตานิ, “อสฺมี”ติ อิญฺชิตเมตํ, “อหมสฺมี”ติ อิญฺชิตเมตํ, “อยมหมสฺมี”ติ อิญฺชิตเมตํ, “ภวิสฺสนฺ”ติ อิญฺชิตเมตํ, “รูปี ภวิสฺสนฺ”ติ อิญฺชิตเมตํ, “อรูปี ภวิสฺสนฺ”ติ อิญฺชิตเมตํ, “สญฺญี ภวิสฺสนฺ”ติ อิญฺชิตเมตํ, “อสญฺญี ภวิสฺสนฺ”ติ อิญฺชิตเมตํ, “เนวสญฺญีนาสญฺญี ภวิสฺสนฺ”ติ อิญฺชิตเมตํ.\csromanspacer{} อิมานิ นว อิญฺชิตานิ.}

\proseitem{965}{ตตฺถ กตมานิ นว มญฺญิตานิ.\csromanspacer{} นว ผนฺทิตานิ.\csromanspacer{} นว ปปญฺจิตานิ.\csromanspacer{} นว สงฺขตานิ, “อสฺมี”ติ สงฺขตเมตํ, “อหมสฺมี”ติ สงฺขตเมตํ, “อยมหมสฺมี”ติ สงฺขตเมตํ, “ภวิสฺสนฺ”ติ สงฺขตเมตํ, “รูปี ภวิสฺสนฺ”ติ สงฺขตเมตํ, “อรูปี ภวิสฺสนฺ”ติ สงฺขตเมตํ, “สญฺญี ภวิสฺสนฺ”ติ}

\vspace{\tipitakaparskip}
\csromancenter{ขุทฺทกวตฺถุวิภงฺค สงฺขตเมตํ, “อสญฺญี ภวิสฺสนฺ”ติ สงฺขตเมตํ, “เนวสญฺญีนาสญฺญี ภวิสฺสนฺ”ติ สงฺขตเมตํ, อิมานิ นว สงฺขตานิ.}
\csromanheader{2}{นวกํ.}
\csromansectionrule
\csromanmatikaanchor{mtk.188}
\csromantocmarkref{subparagraph}{10. ทสกนิทฺเทส}{mtk.188}
\bookmark[dest={mtk.188},level=5]{10. ทสกนิทฺเทส}
\csromanheader{6}{10. ทสกนิทฺเทส}
\proseitem{966}{\csromanfolio{405}ตตฺถ กตมานิ ทส กิเลสวตฺถูนิ, โลโภ โทโส โมโห มาโน ทิฏฺฐิ วิจิกิจฺฉา ถินํ อุทฺธจฺจํ อหิริกํ อโนตฺตปฺปํ, อิมานิ ทส กิเลสวตฺถูนิ.}

\proseitem{967}{ตตฺถ กตมานิ ทส อาฆาตวตฺถูนิ, “อนตฺถํ เม อจรี”ติ อาฆาโต ชายติ, “อนตฺถํ เม จรตี”ติ อาฆาโต ชายติ, “อนตฺถํ เม จริสฺสตี”ติ อาฆาโต ชายติ, “ปิยสฺส เม มนาปสฺส อนตฺถํ อจริ ฯเปฯ อนตฺถํ จรติ ฯเปฯ อนตฺถํ จริสฺสตี”ติ อาฆาโต ชายติ, “อปฺปิยสฺส เม อมนาปสฺส อตฺถํ อจริ ฯเปฯ อตฺถํ จรติ ฯเปฯ อตฺถํ จริสฺสตี”ติ อาฆาโต ชายติ, อฏฺฐาเน วา ปน อาฆาโต ชายติ, อิมานิ ทส อาฆาตวตฺถูนิ.}

\proseitem{968}{ตตฺถ กตเม ทส อกุสลกมฺมปถา, ปาณาติปาโต อทินฺนาทานํ กาเมสุมิจฺฉาจาโร มุสาวาโท ปิสุณา วาจา ผรุสา วาจา สมฺผปฺปลาโป อภิชฺฌา พฺยาปาโท มิจฺฉาทิฏฺฐิ, อิเม ทส อกุสลกมฺมปถา.}

\proseitem{969}{ตตฺถ กตมานิ ทส สํโยชนานิ, กามราคสํโยชนํ ปฏิฆสํโยชนํ มานสํโยชนํ ทิฏฺฐิสํโยชนํ วิจิกิจฺฉาสํโยชนํ สีลพฺพตปรามาสสํโยชนํ ภวราคสํโยชนํ อิสฺสาสํโยชนํ มจฺฉริยสํโยชนํ อวิชฺชาสํโยชนํ, อิมานิ ทส สํโยชนานิ.}

\proseitem{970}{ตตฺถ กตเม ทส มิจฺฉตฺตา, มิจฺฉาทิฏฺฐิ มิจฺฉาสงฺกปฺโป มิจฺฉาวาจา มิจฺฉากมฺมนฺโต มิจฺฉา-อาชีโว มิจฺฉาวายาโม มิจฺฉาสติ มิจฺฉาสมาธิ มิจฺฉาญาณํ มิจฺฉาวิมุตฺติ, อิเม ทส มิจฺฉตฺตา.}

\proseitem{971}{ตตฺถ กตมา ทสวตฺถุกา มิจฺฉาทิฏฺฐิ, “นตฺถิ ทินฺนํ, นตฺถิ ยิฏฺฐํ, นตฺถิ หุตํ, นตฺถิ สุกตทุกฺกฏานํ กมฺมานํ ผลํ วิปาโก, นตฺถิ อยํ โลโก, \csromanfolio{406}นตฺถิ ปโร โลโก, นตฺถิ มาตา, นตฺถิ ปิตา, นตฺถิ สตฺตา โอปปาติกา, นตฺถิ โลเก สมณพฺราหฺมณา สมฺมคฺคตา\footnote{สมคฺคตา (ก)} สมฺมา ปฏิปนฺนา เย อิมญฺจ โลกํ ปรญฺจ โลกํ สยํ อภิญฺญา สจฺฉิกตฺวา ปเวเทนฺตี”ติ, อยํ ทสวตฺถุกา มิจฺฉาทิฏฺฐิ.}

\proseitem{972}{ตตฺถ กตมา ทสวตฺถุกา อนฺตคฺคาหิกา ทิฏฺฐิ, “สสฺสโต โลโก”ติ วา, “อสสฺสโต โลโก”ติ วา, “อนฺตวา โลโก”ติ วา, “อนนฺตวา โลโก”ติ วา, “ตํ ชีวํ ตํ สรีรนฺ”ติ วา, “อญฺญํ ชีวํ อญฺญํ สรีรนฺ”ติ วา, “โหติ ตถาคโต ปรํ มรณา”ติ วา, “น โหติ ตถาคโต ปรํ มรณา”ติ วา, “โหติ จ น จ โหติ ตถาคโต ปรํ มรณา”ติ วา, “เนว โหติ น น โหติ ตถาคโต ปรํ มรณา”ติ วา.\csromanspacer{} อยํ ทสวตฺถุกา อนฺตคฺคาหิกา ทิฏฺฐิ.}

\csromanheader{2}{ทสกํ.}
\csromansectionrule
\csromanmatikaanchor{mtk.189}
\csromantocmarkref{subparagraph}{11. ตณฺหาวิจริตนิทฺเทส}{mtk.189}
\bookmark[dest={mtk.189},level=5]{11. ตณฺหาวิจริตนิทฺเทส}
\csromanheader{6}{11. ตณฺหาวิจริตนิทฺเทส}
\proseitem{973}{ตตฺถ กตมานิ อฏฺฐารส ตณฺหาวิจริตานิ, อชฺฌตฺติกสฺส อุปาทาย “อสฺมี”ติ โหติ, “อิตฺถสฺมี”ติ โหติ, “เอวสฺมี”ติ โหติ, “อญฺญถาสฺมี”ติ โหติ, “ภวิสฺสนฺ”ติ โหติ, “อิตฺถํ ภวิสฺสนฺ”ติ โหติ, “เอวํ ภวิสฺสนฺ”ติ โหติ, “อญฺญถา ภวิสฺสนฺ”ติ โหติ, “อสสฺมี”ติ โหติ, “สาตสฺมี”ติ โหติ, “สิยนฺ”ติ โหติ, “อิตฺถํ สิยนฺ”ติ โหติ, “เอวํ สิยนฺ”ติ โหติ, “อญฺญถา สิยนฺ”ติ โหติ, “อปาหํ สิยนฺ”ติ โหติ, “อปาหํ อิตฺถํ สิยนฺ”ติ โหติ, “อปาหํ เอวํ สิยนฺ”ติ โหติ, “อปาหํ อญฺญถา สิยนฺ”ติ โหติ.}

\proseitem{974}{กถญฺจ “อสฺมี”ติ โหติ, กญฺจิ ธมฺมํ อนวการิํ\footnote{อนวการี (สี, ก)} กริตฺวา รูปํ ฯเปฯ.\csromanspacer{} เวทนํ.\csromanspacer{} สญฺญํ.\csromanspacer{} สงฺขาเร.\csromanspacer{} วิญฺญาณํ “อสฺมี”ติ ฉนฺทํ ปฏิลภติ, “อสฺมี”ติ มานํ ปฏิลภติ, “อสฺมี”ติ ทิฏฺฐิํ ปฏิลภติ, ตสฺมิํ สติ อิมานิ ปปญฺจิตานิ โหนฺติ “อิตฺถสฺมี”ติ วา, “เอวสฺมี”ติ วา, “อญฺญถาสฺมี”ติ วา. (1)}

\csromanheader{2}{17. ขุทฺทกวตฺถุวิภงฺค}
\prose{\csromanfolio{407}“กถญฺจ “อิตฺถสฺมี”ติ โหติ, “ขตฺติโยสฺมี”ติ วา, “พฺราหฺมโณสฺมี”ติ วา, “เวสฺโสสฺมี”ติ วา, “สุทฺโทสฺมี”ติ วา, “คหฏฺโฐสฺมี”ติ วา, “ปพฺพชิโตสฺมี”ติ วา, “เทโวสฺมี”ติ วา, “มนุสฺโสสฺมี”ติ วา, “รูปีสฺมี”ติ วา, “อรูปีสฺมี”ติ วา, “สญฺญีสฺมี”ติ วา, “อสญฺญีสฺมี”ติ วา, “เนวสญฺญีนาสญฺญีสฺมี”ติ วา, เอวํ “อิตฺถสฺมี”ติ โหติ. (2)}

\prose{กถญฺจ “เอวสฺมี”ติ โหติ, ปรปุคฺคลํ\footnote{ปรํ ปุคฺคลํ (สฺยา)} อุปนิธาย “ยถา โส ขตฺติโย ตถาหํ ขตฺติโยสฺมี”ติ วา, “ยถา โส พฺราหฺมโณ ตถาหํ พฺราหฺมโณสฺมี”ติ วา, “ยถา โส เวสฺโส ตถาหํ เวสฺโสสฺมี”ติ วา “ยถา โส สุทฺโท ตถาหํ สุทฺโทสฺมี”ติ วา, “ยถา โส คหฏฺโฐ ตถาหํ คหฏฺโฐสฺมี”ติ วา, “ยถา โส ปพฺพชิโต ตถาหํ ปพฺพชิโตสฺมี”ติ วา, “ยถา โส เทโว ตถาหํ เทโวสฺมี”ติ วา, “ยถา โส มนุสฺโส ตถาหํ มนุสฺโสสฺมี”ติ วา, “ยถา โส รูปี ตถาหํ รูปีสฺมี”ติ วา, “ยถา โส อรูปี ตถาหํ อรูปีสฺมี”ติ วา “ยถา โส สญฺญี ตถาหํ สญฺญีสฺมี”ติ วา, “ยถา โส อสญฺญี ตถาหํ อสญฺญีสฺมี”ติ วา, “ยถา โส เนวสญฺญีนาสญฺญี ตถาหํ เนวสญฺญีนาสญฺญีสฺมี”ติ วา, เอวํ “เอวสฺมี”ติ โหติ. (3)}

\prose{กถญฺจ “อญฺญถาสฺมี”ติ โหติ, ปรปุคฺคลํ อุปนิธาย “ยถา โส ขตฺติโย นาหํ ตถา ขตฺติโยสฺมี”ติ วา, “ยถา โส พฺราหฺมโณ นาหํ ตถา พฺราหฺมโณสฺมี”ติ วา, “ยถา โส เวสฺโส นาหํ ตถา เวสฺโสสฺมี”ติ วา, “ยถา โส สุทฺโท นาหํ ตถา สุทฺโทสฺมี”ติ วา, “ยถา โส คหฏฺโฐ นาหํ ตถา คหฏฺโฐสฺมี”ติ วา, “ยถา โส ปพฺพชิโต นาหํ ตถา ปพฺพชิโตสฺมี”ติ วา, “ยถา โส เทโว นาหํ ตถา เทโวสฺมี”ติ วา, “ยถา โส มนุสฺโส นาหํ ตถา มนุสฺโสสฺมี”ติ วา, “ยถา โส รูปี นาหํ ตถา รูปีสฺมี”ติ วา, “ยถา โส อรูปี นาหํ ตถา อรูปีสฺมี”ติ วา, “ยถา โส สญฺญี นาหํ ตถา สญฺญีสฺมี”ติ วา, “ยถา โส อสญฺญี นาหํ ตถา อสญฺญีสฺมี”ติ วา, “ยถา โส เนวสญฺญีนาสญฺญี นาหํ ตถา เนวสญฺญีนาสญฺญีสฺมี”ติ วา, เอวํ “อญฺญถาสฺมี”ติ โหติ. (4)}

\prose{\csromanfolio{408}กถญฺจ “ภวิสฺสนฺ”ติ โหติ, กญฺจิ ธมฺมํ อนวการิํ กริตฺวา รูปํ ฯเปฯ.\csromanspacer{} เวทนํ.\csromanspacer{} สญฺญํ.\csromanspacer{} สงฺขาเร.\csromanspacer{} วิญฺญาณํ “ภวิสฺสนฺ”ติ ฉนฺทํ ปฏิลภติ, “ภวิสฺสนฺ”ติ มานํ ปฏิลภติ, “ภวิสฺสนฺ”ติ ทิฏฺฐิํ ปฏิลภติ, ตสฺมิํ สติ อิมานิ ปปญฺจิตานิ โหนฺติ “อิตฺถํ ภวิสฺสนฺ”ติ วา, “เอวํ ภวิสฺสนฺ”ติ วา, “อญฺญถา ภวิสฺสนฺ”ติ วา. (5)}

\prose{กถญฺจ “อิตฺถํ ภวิสฺสนฺ”ติ โหติ, “ขตฺติโย ภวิสฺสนฺ”ติ วา, “พฺราหฺมโณ ภวิสฺสนฺ”ติ วา, “เวสฺโส ภวิสฺสนฺ”ติ วา, “สุทฺโท ภวิสฺสนฺ”ติ วา, “คหฏฺโฐ ภวิสฺสนฺ”ติ วา, “ปพฺพชิโต ภวิสฺสนฺ”ติ วา, “เทโว ภวิสฺสนฺ”ติ วา, “มนุสฺโส ภวิสฺสนฺ”ติ วา, “รูปี ภวิสฺสนฺ”ติ วา, “อรูปี ภวิสฺสนฺ”ติ วา, “สญฺญี ภวิสฺสนฺ”ติ วา, “อสญฺญี ภวิสฺสนฺ”ติ วา, “เนวสญฺญีนาสญฺญี ภวิสฺสนฺ”ติ วา.\csromanspacer{} เอวํ “อิตฺถํ ภวิสฺสนฺ”ติ โหติ. (6)}

\prose{กถญฺจ “เอวํ ภวิสฺสนฺ”ติ โหติ, ปรปุคฺคลํ อุปนิธาย “ยถา โส ขตฺติโย ตถาหํ ขตฺติโย ภวิสฺสนฺ”ติ วา, “ยถา โส พฺราหฺมโณ ตถาหํ พฺราหฺมโณ ภวิสฺสนฺ”ติ วา ฯเปฯ “ยถา โส “เนวสญฺญีนาสญฺญี ตถาหํ เนวสญฺญีนาสญฺญี ภวิสฺสนฺ”ติ วา. “เอวํ ภวิสฺสนฺ”ติ โหติ. (7)}

\prose{กถญฺจ “อญฺญถา ภวิสฺสนฺ”ติ โหติ, ปรปุคฺคลํ อุปนิธาย “ยถา โส ขตฺติโย นาหํ ตถา ขตฺติโย ภวิสฺสนฺ”ติ วา, “ยถา โส พฺราหฺมโณ นาหํ ตถา พฺราหฺมโณ ภวิสฺสนฺ”ติ วา ฯเปฯ “ยถา โส เนวสญฺญีนาสญฺญี นาหํ ตถา เนวสญฺญีนาสญฺญี ภวิสฺสนฺ”ติ วา.\csromanspacer{} เอวํ “อญฺญถา ภวิสฺสนฺ”ติ โหติ. (8)}

\prose{กถญฺจ “อสสฺมี”ติ โหติ, กญฺจิ ธมฺมํ อนวการิํ กริตฺวา รูปํ ฯเปฯ.\csromanspacer{} เวทนํ.\csromanspacer{} สญฺญํ.\csromanspacer{} สงฺขาเร.\csromanspacer{} วิญฺญาณํ “นิจฺโจสฺมิ, ธุโวสฺมิ, สสฺสโตสฺมิ, อวิปริณามธมฺโมสฺมี”ติ. “เอวํ อสสฺมี”ติ โหติ. (9)}

\prose{กถญฺจ “สาตสฺมี”ติ โหติ, กญฺจิ ธมฺมํ อนวการิํ กริตฺวา รูปํ ฯเปฯ.\csromanspacer{} เวทนํ.\csromanspacer{} สญฺญํ.\csromanspacer{} สงฺขาเร.\csromanspacer{} วิญฺญาณํ “อุจฺฉิชฺชิสฺสามิ วินสฺสิสฺสามิ น ภวิสฺสามี”ติ, เอวํ “สาตสฺมี”ติ โหติ. (10)}

\prose{กถญฺจ “สิยนฺ”ติ โหติ, กญฺจิ ธมฺมํ อนวการิํ กริตฺวา รูปํ ฯเปฯ.\csromanspacer{} เวทนํ.\csromanspacer{} สญฺญํ.\csromanspacer{} สงฺขาเร.\csromanspacer{} วิญฺญาณํ “สิยนฺ”ติ ฉนฺทํ ปฏิลภติ, “สิยนฺ”ติ มานํ}

\vspace{\tipitakaparskip}
\csromancenter{ขุทฺทกวตฺถุวิภงฺค ปฏิลภติ, “สิยนฺ”ติ ทิฏฺฐิํ ปฏิลภติ, ตสฺมิํ สติ อิมานิ ปปญฺจิตานิ โหนฺติ “อิตฺถํ สิยนฺ”ติ วา, “เอวํ สิยนฺ”ติ วา, “อญฺญถา สิยนฺ”ติ วา. (11)}
\prose{\csromanfolio{409}กถญฺจ “อิตฺถํ สิยนฺ”ติ โหติ, “ขตฺติโย สิยนฺ”ติ วา, “พฺราหฺมโณ สิยนฺ”ติ วา, “เวสฺโส สิยนฺ”ติ วา, “สุทฺโท สิยนฺ”ติ วา, “คหฏฺโฐ สิยนฺ”ติ วา, “ปพฺพชิโต สิยนฺ”ติ วา, “เทโว สิยนฺ”ติ วา, “มนุสฺโส สิยนฺ”ติ วา, “รูปี สิยนฺ”ติ วา, “อรูปี สิยนฺ”ติ วา, “สญฺญี สิยนฺ”ติ วา, “อสญฺญี สิยนฺ”ติ วา, “เนวสญฺญีนาสญฺญี สิยนฺ”ติ วา, เอวํ “อิตฺถํ สิยนฺ”ติ โหติ. (12)}

\prose{กถญฺจ “เอวํ สิยนฺ”ติ โหติ, ปรปุคฺคลํ อุปนิธาย “ยถา โส ขตฺติโย ตถาหํ ขตฺติโย สิยนฺ”ติ วา, “ยถา โส พฺราหฺมโณ ตถาหํ พฺรหฺมโณ สิยนฺ”ติ วา ฯเปฯ “ยถา โส เนวสญฺญีนาสญฺญี ตถาหํ เนวสญฺญีนาสญฺญี สิยนฺ”ติ วา, เอวํ “เอวํ สิยนฺ”ติ โหติ. (13)}

\prose{กถญฺจ “อญฺญถา สิยนฺ”ติ โหติ, ปรปุคฺคลํ อุปนิธาย “ยถา โส ขตฺติโย นาหํ ตถา ขตฺติโย สิยนฺ”ติ วา, “ยถา โส พฺราหฺมโณ นาหํ ตถา พฺราหฺมโณ สิยนฺ”ติ วา ฯเปฯ “ยถา โส เนวสญฺญีนาสญฺญี นาหํ ตถา เนวสญฺญีนาสญฺญี สิยนฺ”ติ วา, เอวํ “อญฺญถา สิยานฺ”ติ โหติ. (14)}

\prose{กถญฺจ “อปาหํ สิยนฺ”ติ โหติ, กญฺจิ ธมฺมํ อนวการิํ กริตฺวา รูปํ.\csromanspacer{} เวทนํ.\csromanspacer{} สญฺญํ.\csromanspacer{} สงฺขาเร.\csromanspacer{} วิญฺญาณํ “อปาหํ สิยนฺ”ติ ฉนฺทํ ปฏิลภติ, “อปาหํ สิยนฺ”ติ มานํ ปฏิลภติ, “อปาหํ สิยนฺ”ติ ทิฏฺฐิํ ปฏิลภติ, ตสฺมิํ สติ อิมานิ ปปญฺจิตานิ โหนฺติ “อปาหํ อิตฺถํ สิยนฺ”ติ วา, “อปาหํ เอวํ สิยนฺ”ติ วา, “อปาหํ อญฺญถา สิยนฺ”ติ วา. (15)}

\prose{กถญฺจ “อปาหํ อิตฺถํ สิยนฺ”ติ โหติ, “อปาหํ ขตฺติโย สิยนฺ”ติ วา, “อปาหํ พฺราหฺมโณ สิยนฺ”ติ วา, “อปาหํ เวสฺโส สิยนฺ”ติ วา, “อปาหํ สุทฺโท สิยนฺ”ติ วา, “อปาหํ คหฏฺโฐ สิยนฺ”ติ วา, “อปาหํ ปพฺพชิโต สิยนฺ”ติ วา, “อปาหํ เทโว สิยนฺ”ติ วา, “อปาหํ มนุสฺโส สิยนฺ”ติ วา, “อปาหํ รูปี สิยนฺ”ติ วา, “อปาหํ อรูปี สิยนฺ”ติ วา, “อปาหํ สญฺญี สิยนฺ”ติ วา, “อปาหํ อสญฺญี \csromanfolio{410}สิยนฺ”ติ วา, “อปาหํ เนวสญฺญีนาสญฺญี สิยนฺ”ติ วา, เอวํ “อปาหํ อิตฺถํ สิยนฺ”ติ โหติ. (16)}

\prose{กถญฺจ “อปาหํ เอวํ สิยนฺ”ติ โหติ, ปรปุคฺคลํ อุปนิธาย “ยถา โส ขตฺติโย อปาหํ ตถา ขตฺติโย สิยนฺ”ติ วา, “ยถา โส พฺราหฺมโณ อปาหํ ตถา พฺราหฺมโณ สิยนฺ”ติ วา ฯเปฯ “ยถา โส เนวสญฺญีนาสญฺญี อปาหํ ตถา เนวสญฺญีนาสญฺญี สิยนฺ”ติ วา, เอวํ “อปาหํ เอวํ สิยนฺ”ติ โหติ. (17)}

\prose{กถญฺจ “อปาหํ อญฺญถา สิยนฺ”ติ โหติ, ปรปุคฺคลํ อุปนิธาย “ยถา โส ขตฺติโย อปาหํ น ตถา ขตฺติโย สิยนฺ”ติ วา, “ยถา โส พฺราหฺมโณ อปาหํ น ตถา พฺราหฺมโณ สิยนฺ”ติ วา ฯเปฯ “ยถา โส เนวสญฺญีนาสญฺญี อปาหํ น ตถา เนวสญฺญีนาสญฺญี สิยนฺ”ติ วา, เอวํ “อปาหํ อญฺญถา สิยนฺ”ติ โหติ. (18)}

\prose{อิมานิ อฏฺฐารส ตณฺหาวิจริตานิ อชฺฌตฺติกสฺส อุปาทาย.}

\proseitem{975}{ตตฺถ กตมานิ อฏฺฐารส ตณฺหาวิจริตานิ พาหิรสฺส อุปาทาย, “อิมินา อสฺมี”ติ โหติ, “อิมินา อิตฺถสฺมี”ติ โหติ, “อิมินา เอวสฺมี”ติ โหติ, “อิมินา อญฺญถาสฺมี”ติ โหติ, “อิมินา ภวิสฺสนฺ”ติ โหติ, “อิมินา อิตฺถํ ภวิสฺสนฺ”ติ โหติ, “อิมินา เอวํ ภวิสฺสนฺ”ติ โหติ, “อิมินา อญฺญถา ภวิสฺสนฺ”ติ โหติ, “อิมินา อสสฺมี”ติ โหติ, “อิมินา สาตสฺมี”ติ โหติ, “อิมินา สิยนฺ”ติ โหติ, “อิมินา อิตฺถํ สิยนฺ”ติ โหติ, “อิมินา เอวํ สิยนฺ”ติ โหติ, “อิมินา อญฺญถา สิยนฺ”ติ โหติ, “อิมินา อปาหํ สิยนฺ”ติ โหติ, “อิมินา อปาหํ อิตฺถํ สิยนฺ”ติ โหติ, “อิมินา อปาหํ เอวํ สิยนฺ”ติ โหติ, “อิมินา อปาหํ อญฺญถา สิยนฺ”ติ โหติ.}

\proseitem{976}{กถญฺจ “อิมินา อสฺมี”ติ โหติ, กญฺจิ ธมฺมํ อวการิํ กริตฺวา รูปํ.\csromanspacer{} เวทนํ.\csromanspacer{} สญฺญํ.\csromanspacer{} สงฺขาเร.\csromanspacer{} วิญฺญาณํ “อิมินา อสฺมี”ติ ฉนฺทํ ปฏิลภติ, “อิมินา อสฺมี”ติ มานํ ปฏิลภติ, “อิมินา อสฺมี”ติ ทิฏฺฐิํ ปฏิลภติ, ตสฺมิํ สติ อิมานิ ปปญฺจิตานิ โหนฺติ “อิมินา อิตฺถสฺมี”ติ วา, “อิมินา เอวสฺมี”ติ วา, “อิมินา อญฺญถาสฺมี”ติ วา. (1)}

\csromanheader{2}{17. ขุทฺทกวตฺถุวิภงฺค}
\prose{\csromanfolio{411}กถญฺจ “อิมินา อิตฺถสฺมี”ติ โหติ, “อิมินา ขตฺติโยสฺมี”ติ วา, “อิมินา พฺราหฺมโณสฺมี”ติ วา, “อิมินา เวสฺโสสฺมี”ติ วา, “อิมินา สุทฺโทสฺมี”ติ วา, “อิมินา คหฏฺโฐสฺมี”ติ วา, “อิมินา ปพฺพชิโตสฺมี”ติ วา, “อิมินา เทโวสฺมี”ติ วา, “อิมินา มนุสฺโสสฺมี”ติ วา, “อิมินา รูปีสฺมี”ติ วา, “อิมินา อรูปีสฺมี”ติ วา, “อิมินา สญฺญีสฺมี”ติ วา, “อิมินา อสญฺญีสฺมี”ติ วา, “อิมินา เนวสญฺญีนาสญฺญีสฺมี”ติ วา, เอวํ “อิมินา อิตฺถสฺมี”ติ โหติ. (2)}

\prose{กถญฺจ “อิมินา เอวสฺมี”ติ โหติ, ปรปุคฺคลํ อุปนิธาย “ยถา โส ขตฺติโย อิมินา ตถาหํ ขตฺติโยสฺมี”ติ วา, “ยถา โส พฺราหฺมโณ อิมินา ตถาหํ พฺราหฺมโณสฺมี”ติ วา ฯเปฯ “ยถา โส เนวสญฺญีนาสญฺญี อิมินา ตถาหํ เนวสญฺญีนาสญฺญีสฺมี”ติ วา, เอวํ “อิมินา เอวสฺมี”ติ โหติ. (3)}

\prose{กถญฺจ “อิมินา อญฺญถาสฺมี”ติ โหติ, ปรปุคฺคลํ อุปนิธาย “ยถา โส ขตฺติโย อิมินา นาหํ ตถา ขตฺติโยสฺมี”ติ วา, “ยถา โส พฺราหฺมโณ อิมินา นาหํ ตถา พฺราหฺมโณสฺมี”ติ วา ฯเปฯ “ยถา โส เนวสญฺญีนาสญฺญี อิมินา นาหํ ตถา เนวสญฺญีนาสญฺญีสฺมี”ติ วา, เอวํ “อิมินา อญฺญถาสฺมี”ติ โหติ. (4)}

\prose{กถญฺจ “อิมินา ภวิสฺสนฺ”ติ โหติ, กญฺจิ ธมฺมํ อวการิํ กริตฺวา รูปํ.\csromanspacer{} เวทนํ.\csromanspacer{} สญฺญํ.\csromanspacer{} สงฺขาเร.\csromanspacer{} วิญฺญาณํ “อิมินา ภวิสฺสนฺ”ติ ฉนฺทํ ปฏิลภติ, “อิมินา ภวิสฺสนฺ”ติ มานํ ปฏิลภติ “อิมินา ภวิสฺสนฺ”ติ ทิฏฺฐิํ ปฏิลภติ, ตสฺมิํ สติ อิมานิ ปปญฺจิตานิ โหนฺติ “อิมินา อิตฺถํ ภวิสฺสนฺ”ติ วา, “อิมินา เอวํ ภวิสฺสนฺ”ติ วา, “อิมินา อญฺญถา ภวิสฺสนฺ”ติ วา. (5)}

\prose{กถญฺจ “อิมินา อิตฺถํ ภวิสฺสนฺ”ติ โหติ, “อิมินา ขตฺติโย ภวิสฺสนฺ”ติ วา ฯเปฯ “อิมินา อรูปี ภวิสฺสนฺ”ติ วา, “อิมินา สญฺญี ภวิสฺสนฺ”ติ วา, “อิมินา อสญฺญี ภวิสฺสนฺ”ติ วา, “อิมินา เนวสญฺญีนาสญฺญี ภวิสฺสนฺ”ติ วา, เอวํ “อิมินา อิตฺถํ ภวิสฺสนฺ”ติ โหติ. (6)}

\prose{กถญฺจ “อิมินา เอวํ ภวิสฺสนฺ”ติ โหติ, ปรปุคฺคลํ อุปนิธาย “ยถา โส ขตฺติโย อิมินา ตถาหํ ขตฺติโย ภวิสฺสนฺ”ติ วา “ยถา}

\prose{\csromanfolio{412}โส พฺราหฺมโณ อิมินา ตถาหํ พฺราหฺมโณ ภวิสฺสนฺ”ติ วา ฯเปฯ “ยถา โส เนวสญฺญีนาสญฺญี อิมินา ตถาหํ เนวสญฺญีนาสญฺญี ภวิสฺสนฺ”ติ วา, เอวํ “อิมินา เอวํ ภวิสฺสนฺ”ติ โหติ. (7)}

\prose{กถญฺจ “อิมินา อญฺญถา ภวิสฺสนฺ”ติ โหติ, ปรปุคฺคลํ อุปนิธาย “ยถา โส ขตฺติโย อิมินา นาหํ ตถา ขตฺติโย ภวิสฺสนฺ”ติ วา, “ยถา โส พฺราหฺมโณ อิมินา นาหํ ตถา พฺราหฺมโณ ภวิสฺสนฺ”ติ วา ฯเปฯ “ยถา โส เนวสญฺญีนาสญฺญี อิมินา นาหํ ตถา เนวสญฺญีนาสญฺญี ภวิสฺสนฺ”ติ วา, เอวํ “อิมินา อญฺญถา ภวิสฺสนฺ”ติ โหติ. (8)}

\prose{กถญฺจ “อิมินา อสสฺมี”ติ โหติ, กญฺจิ ธมฺมํ อวการิํ\footnote{อวการี (สี)} กริตฺวา รูปํ ฯเปฯ.\csromanspacer{} เวทนํ.\csromanspacer{} สญฺญํ.\csromanspacer{} สงฺขาเร.\csromanspacer{} วิญฺญาณํ “อิมินา นิจฺโจสฺมิ, ธุโวสฺมิ, สสฺสโตสฺมิ อวิปริณามธมฺโมสฺมี”ติ.\csromanspacer{} เอวํ “อิมินา อสสฺมี”ติ โหติ. (9)}

\prose{กถญฺจ “อิมินา สาตสฺมี”ติ โหติ.\csromanspacer{} กญฺจิ ธมฺมํ อวการิํ กริตฺวา รูปํ ฯเปฯ.\csromanspacer{} เวทนํ.\csromanspacer{} สญฺญํ.\csromanspacer{} สงฺขาเร.\csromanspacer{} วิญฺญาณํ “อิมินา อุจฺฉิชฺชิสฺสามิ วินสฺสิสฺสามิ, น ภวิสฺสามี”ติ, เอวํ “อิมินา สาตสฺมี”ติ โหติ. (10)}

\prose{กถญฺจ “อิมินา สิยนฺ”ติ โหติ, กญฺจิ ธมฺมํ อวการิํ กริตฺวา รูปํ ฯเปฯ.\csromanspacer{} เวทนํ.\csromanspacer{} สญฺญํ.\csromanspacer{} สงฺขาเร.\csromanspacer{} วิญฺญาณํ “อิมินา สิยนฺ”ติ ฉนฺทํ ปฏิลภติ, “อิมินา สิยนฺ”ติ มานํ ปฏิลภติ, “อิมินา สิยนฺ”ติ ทิฏฺฐิํ ปฏิลภติ, ตสฺมิํ สติ อิมินา ปปญฺจิตานิ โหนฺติ “อิมินา อิตฺถํ สิยนฺ”ติ วา, “อิมินา เอวํ สิยนฺ”ติ วา, “อิมินา อญฺญถา สิยนฺ”ติ วา. (11)}

\prose{กถญฺจ “อิมินา อิตฺถํ สิยนฺ”ติ โหติ, “อิมินา ขตฺติโย สิยนฺ”ติ วา, “อิมินา พฺราหฺมโณ สิยนฺ”ติ วา, “อิมินา เวสฺโส สิยนฺ”ติ วา, “อิมินา สุทฺโท สิยนฺ”ติ วา, “อิมินา คหฏฺโฐ สิยนฺ”ติ วา, “อิมินา ปพฺพชิโต สิยนฺ”ติ วา, “อิมินา เทโว สิยนฺ”ติ วา, “อิมินา มนุสฺโส สิยนฺ”ติ วา, “อิมินา รูปี สิยนฺ”ติ วา, “อิมินา อรูปี สิยนฺ”ติ วา, “อิมินา สญฺญี สิยนฺ”ติ วา, “อิมินา อสญฺญี สิยนฺ”ติ วา, “อิมินา เนวสญฺญีนาสญฺญี สิยนฺ”ติ วา, เอวํ “อิมินา อิตฺถํ สิยนฺ”ติ โหติ. (12)}

\csromanheader{2}{17. ขุทฺทกวตฺถุวิภงฺค}
\prose{\csromanfolio{413}กถญฺจ “อิมินา เอวํ สิยนฺ”ติ โหติ, ปรปุคฺคลํ อุปนิธาย “ยถา โส ขตฺติโย อิมินา ตถาหํ ขตฺติโย สิยนฺ”ติ วา, ‘ยถา โส พฺราหฺมโณ อิมินา ตถาหํ พฺราหฺมโณ สิยนฺ”ติ วา ฯเปฯ “ยถา โส เนวสญฺญีนาสญฺญี อิมินา ตถาหํ เนวสญฺญีนาสญฺญี สิยนฺ”ติ วา, เอวํ “อิมินา เอวํ สิยนฺ”ติ โหติ. (13)}

\prose{กถญฺจ “อิมินา อญฺญถา สิยนฺ”ติ โหติ, ปรปุคฺคลํ อุปนิธาย “ยถา โส ขตฺติโย อิมินา นาหํ ตถา ขตฺติโย สิยนฺ”ติ วา, “ยถา โส พฺราหฺมโณ อิมินา นาหํ ตถา พฺราหฺมโณ สิยนฺ”ติ วา “ยถา โส เนวสญฺญีนาสญฺญี อิมินา นาหํ ตถา เนวสญฺญีนาสญฺญี สิยนฺ”ติ วา, เอวํ “อิมินา อญฺญถา สิยนฺ”ติ โหติ. (14)}

\prose{กถญฺจ “อิมินา อปาหํ สิยนฺ”ติ โหติ, กญฺจิ ธมฺมํ อวการิํ กริตฺวา รูปํ ฯเปฯ.\csromanspacer{} เวทนํ.\csromanspacer{} สญฺญํ.\csromanspacer{} สงฺขาเร.\csromanspacer{} วิญฺญาณํ “อิมินา อปาหํ สิยนฺ”ติ ฉนฺทํ ปฏิลภติ, “อิมินา อปาหํ สิยนฺ”ติ มานํ ปฏิลภติ, “อิมินา อปาหํ สิยนฺ”ติ ทิฏฺฐิํ ปฏิลภติ, ตสฺมิํ สติ อิมานิ ปปญฺจิตานิ โหนฺติ “อิมินา อปาหํ อิตฺถํ สิยนฺ”ติ วา, “อิมินา อปาหํ เอวํ สิยนฺ”ติ วา, “อิมินา อปาหํ อญฺญถา สิยนฺ”ติ วา. (15)}

\prose{กถญฺจ “อิมินา อปาหํ อิตฺถํ สิยนฺ”ติ โหติ, “อิมินา อปาหํ ขตฺติโย สิยนฺ”ติ วา, “อิมินา อปาหํ พฺราหฺมโณ สิยนฺ”ติ วา, “อิมินา อปาหํ เวสฺโส สิยนฺ”ติ วา, “อิมินา อปาหํ สุทฺโท สิยนฺ”ติ วา, “อิมินา อปาหํ คหฏฺโฐ สิยนฺ”ติ วา, “อิมินา อปาหํ ปพฺพชิโต สิยนฺ”ติ วา, “อิมินา อปาหํ เทโว สิยนฺ”ติ วา, “อิมินา อปาหํ มนุสฺโส สิยนฺ”ติ วา, “อิมินา อปาหํ รูปี สิยนฺ”ติ วา, “อิมินา อปาหํ อรูปี สิยนฺ”ติ วา, “อิมินา อปาหํ สญฺญี สิยนฺ”ติ วา, “อิมินา อปาหํ อสญฺญี สิยนฺ”ติ วา, “อิมินา อปาหํ เนวสญฺญีนาสญฺญี สิยนฺ”ติ วา, เอวํ “อิมินา อปาหํ อิตฺถํ สิยนฺ”ติ โหติ. (16)}

\prose{กถญฺจ “อิมินา อปาหํ เอวํ สิยนฺ”ติ โหติ, ปรปุคฺคลํ อุปนิธาย “ยถา โส ขตฺติโย อิมินา อปาหํ ตถา ขตฺติโย สิยนฺ”ติ วา, “ยถา โส พฺราหฺมโณ อิมินา อปาหํ ตถา พฺราหฺมโณ สิยนฺ”ติ วา ฯเปฯ “ยถา โส เนวสญฺญีนาสญฺญี อิมินา อปาหํ ตถา เนวสญฺญีนาสญฺญี สิยนฺ”ติ วา.\csromanspacer{} เอวํ “อิมินา อปาหํ เอวํ สิยนฺ”ติ โหติ. (17)}

\prose{\csromanfolio{414}กถญฺจ “อิมินา อปาหํ อญฺญถา สิยนฺ”ติ โหติ, ปรปุคฺคลํ อุปนิธาย “ยถา โส ขตฺติโย อิมินา อปาหํ น ตถา ขตฺติโย สิยนฺ”ติ วา, “ยถา โส พฺราหฺมโณ อิมินา อปาหํ น ตถา พฺราหฺมโณ สิยนฺ”ติ วา ฯเปฯ “ยถา โส เนวสญฺญีนาสญฺญี อิมินา อปาหํ น ตถา เนวสญฺญีนาสญฺญี สิยนฺ”ติ วา. “เอวํ อิมินา อปาหํ อญฺญถา สิยนฺ”ติ โหติ. (18)}

\prose{อิมานิ อฏฺฐารส ตณฺหาวิจริตานิ พาหิรสฺส อุปาทาย.}

\prose{อิติ อิมานิ อฏฺฐารส ตณฺหาวิจริตานิ อชฺฌตฺติกสฺส อุปาทาย, อิมานิ อฏฺฐารส ตณฺหาวิจริตานิ พาหิรสฺส อุปาทาย ตเทกชฺฌํ อภิสญฺญูหิตฺวา อภิสงฺขิปิตฺวา ฉตฺติํส ตณฺหาวิจริตานิ โหนฺติ.\csromanspacer{} อิติ เอวรูปานิ อตีตานิ ฉตฺติํส ตณฺหาวิจริตานิ, อนาคตานิ ฉตฺติํส ตณฺหาวิจริตานิ, ปจฺจุปฺปนฺนานิ ฉตฺติํส ตณฺหาวิจริตานิ ตเทกชฺฌํ อภิสญฺญูหิตฺวา อภิสงฺขิปิตฺวา อฏฺฐตณฺหาวิจริตสตํ โหติ.}

\proseitem{977}{ตตฺถ กตมานิ ทฺวาสฏฺฐิ ทิฏฺฐิคตานิ พฺรหฺมชาเล เวยฺยากรเณ วุตฺตานิ ภควตา, จตฺตาโร สสฺสตวาทา, จตฺตาโร เอกจฺจสสฺสติกา, จตฺตาโร อนฺตานนฺติกา, จตฺตาโร อมราวิกฺเขปิกา, เทฺว อธิจฺจสมุปฺปนฺนิกา, โสฬส สญฺญีวาทา, อฏฺฐ อสญฺญีวาทา, อฏฺฐ เนวสญฺญีนาสญฺญีวาทา, สตฺต อุจฺเฉทวาทา, ปญฺจ ทิฏฺฐธมฺมนิพฺพานวาทา.\csromanspacer{} อิมานิ ทฺวาสฏฺฐิ ทิฏฺฐิคตานิ พฺรหฺมชาเล เวยฺยากรเณ วุตฺตานิ ภควตาติ.}

\nitthitam{\textbf{ขุทฺทกวตฺถุวิภงฺโค นิฏฺฐิโต.}}
\csromanmatikaanchor{mtk.190}
\csromantocmarknopage{subparagraph}{18. ธมฺมหทยวิภงฺค}{mtk.190}
\bookmark[dest={mtk.190},level=5]{18. ธมฺมหทยวิภงฺค}
\csromanheader{6}{18. ธมฺมหทยวิภงฺค}
\csromanmatikaanchor{mtk.191}
\csromantocmarkref{subparagraph}{1. สพฺพสงฺคาหิกวาร}{mtk.191}
\bookmark[dest={mtk.191},level=5]{1. สพฺพสงฺคาหิกวาร}
\csromanheader{6}{1. สพฺพสงฺคาหิกวาร}
\proseitem{978}{\csromanfolio{415}กติ ขนฺธา, กติ อายตนานิ, กติ ธาตุโย, กติ สจฺจานิ, กติ อินฺทฺริยานิ, กติ เหตู, กติ อาหารา, กติ ผสฺสา, กติ เวทนา, กติ สญฺญา, กติ เจตนา, กติ จิตฺตานิ.}

\prose{ปญฺจกฺขนฺธา, ทฺวาทสายตนานิ, อฏฺฐารส ธาตุโย, จตฺตาริ สจฺจานิ, พาวีสตินฺทฺริยานิ, นว เหตู, จตฺตาโร อาหารา, สตฺต ผสฺสา, สตฺต เวทนา, สตฺต สญฺญา, สตฺต เจตนา, สตฺต จิตฺตานิ.}

\proseitem{979}{ตตฺถ กตเม ปญฺจกฺขนฺธา, รูปกฺขนฺโธ เวทนากฺขนฺโธ สญฺญากฺขนฺโธ สงฺขารกฺขนฺโธ วิญฺญาณกฺขนฺโธ, อิเม วุจฺจนฺติ ปญฺจกฺขนฺธา. (1)}

\proseitem{980}{ตตฺถ กตมานิ ทฺวาทสายตนานิ, จกฺขายตนํ รูปายตนํ โสตายตนํ สทฺทายตนํ ฆานายตนํ คนฺธายตนํ ชิวฺหายตนํ รสายตนํ กายายตนํ โผฏฺฐพฺพายตนํ มนายตนํ ธมฺมายตนํ, อิมานิ วุจฺจนฺติ ทฺวาทสายตนานิ. (2)}

\proseitem{981}{ตตฺถ กตมา อฏฺฐารส ธาตุโย, จกฺขุธาตุ รูปธาตุ จกฺขุวิญฺญาณธาตุ โสตธาตุ สทฺทธาตุ โสตวิญฺญาณธาตุ ฆานธาตุ คนฺธธาตุ ฆานวิญฺญาณธาตุ ชิวฺหาธาตุ รสธาตุ ชิวฺหาวิญฺญาณธาตุ กายธาตุ โผฏฺฐพฺพธาตุ กายวิญฺญาณธาตุ มโนธาตุ ธมฺมธาตุ มโนวิญฺญาณธาตุ, อิมา วุจฺจนฺติ อฏฺฐารส ธาตุโย. (3)}

\proseitem{982}{ตตฺถ กตมานิ จตฺตาริ สจฺจานิ, ทุกฺขสจฺจํ สมุทยสจฺจํ นิโรธสจฺจํ มคฺคสจฺจํ, อิมานิ วุจฺจนฺติ จตฺตาริ สจฺจานิ. (4)}

\proseitem{983}{ตตฺถ กตมานิ พาวีสตินฺทฺริยานิ, จกฺขุนฺทฺริยํ โสตินฺทฺริยํ ฆานินฺทฺริยํ ชิวฺหินฺทฺริยํ กายินฺทฺริยํ มนินฺทฺริยํ อิตฺถินฺทฺริยํ ปุริสินฺทฺริยํ ชีวิตินฺทฺริยํ สุขินฺทฺริยํ ทุกฺขินฺทฺริยํ โสมนสฺสินฺทฺริยํ โทมนสฺสินฺทฺริยํ อุเปกฺขินฺทฺริยํ สทฺธินฺทฺริยํ วีริยินฺทฺริยํ สตินฺทฺริยํ สมาธินฺทฺริยํ ปญฺญินฺทฺริยํ อนญฺญาตญฺญสฺสามีตินฺทฺริยํ อญฺญินฺทฺริยํ อญฺญาตาวินฺทฺริยํ, อิมานิ วุจฺจนฺติ พาวีสตินฺทฺริยานิ. (5)}

\proseitem{984}{\csromanfolio{416}ตตฺถ กตเม นว เหตู, ตโย กุสลเหตู, ตโย อกุสลเหตู, ตโย อพฺยากตเหตู.}

\prose{ตตฺถ กตเม ตโย กุสลเหตู, อโลโภ กุสลเหตุ อโทโส กุสลเหตุ อโมโห กุสลเหตุ, อิเม ตโย กุสลเหตู.}

\prose{ตตฺถ กตเม ตโย อกุสลเหตู, โลโภ อกุสลเหตุ โทโส อกุสลเหตุ โมโห อกุสลเหตุ, อิเม ตโย อกุสลเหตู.}

\prose{ตตฺถ กตเม ตโย อพฺยากตเหตู, กุสลานํ วา ธมฺมานํ วิปากโต กิริยาพฺยากเตสุ วา ธมฺเมสุ อโลโภ อโทโส อโมโห, อิเม ตโย อพฺยากตเหตู.\csromanspacer{} อิเม วุจฺจนฺติ นว เหตู. (6)}

\proseitem{985}{ตตฺถ กตเม จตฺตาโร อาหารา, กพฬีการาหาโร\footnote{กพฬีกาโร อาหาโร (สี, สฺยา)} ผสฺสาหาโร มโนสญฺเจตนาหาโร วิญฺญาณาหาโร, อิเม วุจฺจนฺติ จตฺตาโร อาหารา. (7)}

\proseitem{986}{ตตฺถ กตเม สตฺต ผสฺสา, จกฺขุสมฺผสฺโส โสตสมฺผสฺโส ฆานสมฺผสฺโส ชิวฺหาสมฺผสฺโส กายสมฺผสฺโส มโนธาตุสมฺผสฺโส มโนวิญฺญาณธาตุสมฺผสฺโส, อิเม วุจฺจนฺติ สตฺต ผสฺสา. (8)}

\proseitem{987}{ตตฺถ กตมา สตฺต เวทนา, จกฺขุสมฺผสฺสชา เวทนา โสตสมฺผสฺสชา เวทนา ฆานสมฺผสฺสชา เวทนา ชิวฺหาสมฺผสฺสชา เวทนา กายสมฺผสฺสชา เวทนา มโนธาตุสมฺผสฺสชา เวทนา มโนวิญฺญาณธาตุสมฺผสฺสชา เวทนา, อิมา วุจฺจนฺติ สตฺต เวทนา. (9)}

\proseitem{988}{ตตฺถ กตมา สตฺต สญฺญา, จกฺขุสมฺผสฺสชา สญฺญา โสตสมฺผสฺสชา สญฺญา ฆานสมฺผสฺสชา สญฺญา ชิวฺหาสมฺผสฺสชา สญฺญา กายสมฺผสฺสชา สญฺญา มโนธาตุสมฺผสฺสชา สญฺญา มโนวิญฺญาณธาตุสมฺผสฺสชา สญฺญา, อิมา วุจฺจนฺติ สตฺต สญฺญา. (10)}

\csromanheader{2}{18. ธมฺมหทยวิภงฺค}
\proseitem{989}{\csromanfolio{417}ตตฺถ กตมา สตฺต เจตนา, จกฺขุสมฺผสฺสชา เจตนา โสตสมฺผสฺสชา เจตนา ฆานสมฺผสฺสชา เจตนา ชิวฺหาสมฺผสฺสชา เจตนา กายสมฺผสฺสชา เจตนา มโนธาตุสมฺผสฺสชา เจตนา มโนวิญฺญาณธาตุสมฺผสฺสชา เจตนา, อิมา วุจฺจนฺติ สตฺต เจตนา. (11)}

\proseitem{990}{ตตฺถ กตมานิ สตฺต จิตฺตานิ, จกฺขุวิญฺญาณํ โสตวิญฺญาณํ ฆานวิญฺญาณํ ชิวฺหาวิญฺญาณํ กายวิญฺญาณํ มโนธาตุ มโนวิญฺญาณธาตุ, อิมานิ วุจฺจนฺติ สตฺต จิตฺตานิ. (12)}
\csromansectionrule

\csromanmatikaanchor{mtk.192}
\csromantocmarkref{subparagraph}{2. อุปฺปตฺตานุปฺปตฺติวาร}{mtk.192}
\bookmark[dest={mtk.192},level=5]{2. อุปฺปตฺตานุปฺปตฺติวาร}
\csromanheader{6}{2. อุปฺปตฺตานุปฺปตฺติวาร}
\csromanmatikaanchor{mtk.193}
\csromantocmarkref{subparagraph}{1. กามธาตุ}{mtk.193}
\bookmark[dest={mtk.193},level=5]{1. กามธาตุ}
\csromanheader{6}{1. กามธาตุ}
\vspace{\tipitakaparskip}
\csromancenter{กามธาตุยา กติ ขนฺธา, กติ อายตนานิ, กติ ธาตุโย, กติ สจฺจานิ, กติ อินฺทฺริยานิ, กติ เหตู, กติ อาหารา, กติ ผสฺสา, กติ เวทนา, กติ สญฺญา, กติ เจตนา, กติ จิตฺตานิ.}
\csromancenter{กามธาตุยา ปญฺจกฺขนฺธา, ทฺวาทสายตนานิ, อฏฺฐารส ธาตุโย, ติณิ สจฺจานิ, พาวีสตินฺทฺริยานิ, นว เหตู, จตฺตาโร อาหารา, สตฺต ผสฺสา, สตฺต เวทนา, สตฺต สญฺญา, สตฺต เจตนา, สตฺต จิตฺตานิ.}
\proseitem{992}{ตตฺถ กตเม กามธาตุยา ปญฺจกฺขนฺธา, รูปกฺขนฺโธ เวทนากฺขนฺโธ สญฺญากฺขนฺโธ สงฺขารกฺขนฺโธ วิญฺญาณกฺขนฺโธ, อิเม วุจฺจนฺติ กามธาตุยา ปญฺจกฺขนฺธา. (1)}

\prose{ตตฺถ กตมานิ กามธาตุยา ทฺวาทสายตนานิ, จกฺขายตนํ รูปายตนํ โสตายตนํ สทฺทายตนํ ฆานายตนํ คนฺธายตนํ ชิวฺหายตนํ รสายตนํ กายายตนํ โผฏฺฐพฺพายตนํ มนายตนํ ธมฺมายตนํ, อิมานิ วุจฺจนฺติ กามธาตุยา ทฺวาทสายตนานิ. (2)}

\prose{ตตฺถ กตมา กามธาตุยา อฏฺฐารส ธาตุโย, จกฺขุธาตุ รูปธาตุ จกฺขุวิญฺญาณธาตุ โสตธาตุ สทฺทธาตุ โสตวิญฺญาณธาตุ ฆานธาตุ คนฺธธาตุ ฆานวิญฺญาณธาตุ ชิวฺหาธาตุ รสธาตุ ชิวฺหาวิญฺญาณธาตุ กายธาตุ โผฏฺฐพฺพธาตุ กายวิญฺญาณธาตุ มโนธาตุ \csromanfolio{418}ธมฺมธาตุ มโนวิญฺญาณธาตุ, อิมา วุจฺจนฺติ กามธาตุยา อฏฺฐารส ธาตุโย. (3)}

\prose{ตตฺถ กตมานิ กามธาตุยา ตีณิ สจฺจานิ, ทุกฺขสจฺจํ สมุทยสจฺจํ มคฺคสจฺจํ, อิมานิ วุจฺจนฺติ กามธาตุยา ตีณิ สจฺจานิ. (4)}

\prose{ตตฺถ กตมานิ กามธาตุยา พาวีสตินฺทฺริยานิ, จกฺขุนฺทฺริยํ โสตินฺทฺริยํ ฯเปฯ อญฺญาตาวินฺทฺริยํ, อิมานิ วุจฺจนฺติ กามธาตุยา พาวีสตินฺทฺริยานิ. (5)}

\prose{ตตฺถ กตเม กามธาตุยา นว เหตู, ตโย กุสลเหตู, ตโย อกุสลเหตู, ตโย อพฺยากตเหตู ฯเปฯ อิเม วุจฺจนฺติ กามธาตุยา นว เหตู. (6)}

\prose{ตตฺถ กตเม กามธาตุยา จตฺตาโร อาหารา, กพฬีการาหาโร ผสฺสาหาโร มโนสญฺเจตนาหาโร วิญฺญาณาหาโร, อิเม วุจฺจนฺติ กามธาตุยา จตฺตาโร อาหารา. (7)}

\prose{ตตฺถ กตเม กามธาตุยา สตฺต ผสฺสา, จกฺขุสมฺผสฺโส โสตสมฺผสฺโส ฆานสมฺผสฺโส ชิวฺหาสมฺผสฺโส กายสมฺผสฺโส มโนธาตุสมฺผสฺโส มโนวิญฺญาณธาตุสมฺผสฺโส, อิเม วุจฺจนฺติ กามธาตุยา สตฺต ผสฺสา. (8)}

\prose{ตตฺถ กตมา กามธาตุยา สตฺต เวทนา, จกฺขุสมฺผสฺสชา เวทนา โสตสมฺผสฺสชา เวทนา ฆานสมฺผสฺสชา เวทนา ชิวฺหาสมฺผสฺสชา เวทนา กายสมฺผสฺสชา เวทนา มโนธาตุสมฺผสฺสชา เวทนา มโนวิญฺญาณธาตุสมฺผสฺสชา เวทนา, อิมา วุจฺจนฺติ กามธาตุยา สตฺต เวทนา. (9)}

\prose{ตตฺถ กตมา กามธาตุยา สตฺต สญฺญา, จกฺขุสมฺผสฺสชา สญฺญา โสตสมฺผสฺสชา สญฺญา ฆานสมฺผสฺสชา สญฺญา ชิวฺหาสมฺผสฺสชา สญฺญา กายสมฺผสฺสชา สญฺญา มโนธาตุสมฺผสฺสชา สญฺญา มโนวิญฺญาณธาตุสมฺผสฺสชา สญฺญา, อิมา วุจฺจนฺติ กามธาตุยา สตฺต สญฺญา. (10)}

\csromanheader{2}{18. ธมฺมหทยวิภงฺค}
\prose{\csromanfolio{419}ตตฺถ กตมา กามธาตุยา สตฺต เจตนา, จกฺขุสมฺผสฺสชา เจตนา โสตสมฺผสฺสชา เจตนา ฆานสมฺผสฺสชา เจตนา ชิวฺหาสมฺผสฺสชา เจตนา กายสมฺผสฺสชา เจตนา มโนธาตุสมฺผสฺสชา เจตนา มโนวิญฺญาณธาตุสมฺผสฺสชา เจตนา, อิมา วุจฺจนฺติ กามธาตุยา สตฺต เจตนา. (11)}

\prose{ตตฺถ กตมานิ กามธาตุยา สตฺต จิตฺตานิ, จกฺขุวิญฺญาณํ โสตวิญฺญาณํ ฆานวิญฺญาณํ ชิวฺหาวิญฺญาณํ กายวิญฺญาณํ มโนธาตุ มโนวิญฺญาณธาตุ, อิมานิ วุจฺจนฺติ กามธาตุยา สตฺต จิตฺตานิ. (12)}

\csromanmatikaanchor{mtk.194}
\csromantocmarkref{subparagraph}{2. รูปธาตุ}{mtk.194}
\bookmark[dest={mtk.194},level=5]{2. รูปธาตุ}
\csromanheader{6}{2. รูปธาตุ}
\vspace{\tipitakaparskip}
\csromancenter{รูปธาตุยา กติ ขนฺธา, กติ อายตนา, กติ ธาตุโย, กติ สจฺจานิ, กติ อินฺทฺริยานิ ฯเปฯ กติ จิตฺตานิ.}
\csromancenter{รูปธาตุยา ปญฺจกฺขนฺธา, ฉ อายตนานิ, นว ธาตุโย, ตีณิ สจฺจานิ, จุทฺทสินฺทฺริยานิ, อฏฺฐ เหตู, ตโย อาหารา, จตฺตาโร ผสฺสา, จตสฺโส เวทนา, จตสฺโส สญฺญา, จตสฺโส เจตนา, จตฺตาริ จิตฺตานิ.}
\proseitem{994}{ตตฺถ กตเม รูปธาตุยา ปญฺจกฺขนฺธา, รูปกฺขนฺโธ เวทนากฺขนฺโธ สญฺญากฺขนฺโธ สงฺขารกฺขนฺโธ วิญฺญาณกฺขนฺโธ, อิเม วุจฺจนฺติ รูปธาตุยา ปญฺจกฺขนฺธา. (1)}

\prose{ตตฺถ กตมานิ รูปธาตุยา ฉ อายตนานิ, จกฺขายตนํ รูปายตนํ โสตายตนํ สทฺทายตนํ มนายตนํ ธมฺมายตนํ, อิมานิ วุจฺจนฺติ รูปธาตุยา ฉ อายตนานิ. (2)}

\prose{ตตฺถ กตมา รูปธาตุยา นว ธาตุโย, จกฺขุธาตุ รูปธาตุ จกฺขุวิญฺญาณธาตุ โสตธาตุ สทฺทธาตุ โสตวิญฺญาณธาตุ มโนธาตุ ธมฺมธาตุ มโนวิญฺญาณธาตุ, อิมา วุจฺจนฺติ รูปธาตุยา นว ธาตุโย. (3)}

\prose{ตตฺถ กตมานิ รูปธาตุยา ตีณิ สจฺจานิ, ทุกฺขสจฺจํ สมุทยสจฺจํ มคฺคสจฺจํ, อิมานิ วุจฺจนฺติ รูปธาตุยา ตีณิ สจฺจานิ. (4)}

\prose{\csromanfolio{420}ตตฺถ กตมานิ รูปธาตุยา จุทฺทสินฺทฺริยานิ, จกฺขุนฺทฺริยํ โสตินฺทฺริยํ มนินฺทฺริยํ ชีวิตินฺทฺริยํ โสมนสฺสินฺทฺริยํ อุเปกฺขินฺทฺริยํ สทฺธินฺทฺริยํ วีริยินฺทฺริยํ สตินฺทฺริยํ สมาธินฺทฺริยํ ปญฺญินฺทฺริยํ อนญฺญาตญฺญสฺสามีตินฺทฺริยํ อญฺญินฺทฺริยํ อญฺญาตาวินฺทฺริยํ, อิมานิ วุจฺจนฺติ รูปธาตุยา จุทฺทสินฺทฺริยานิ. (5)}

\prose{ตตฺถ กตเม รูปธาตุยา อฏฺฐ เหตู, ตโย กุสลเหตู, เทฺว อกุสลเหตู, ตโย อพฺยากตเหตู.}

\prose{ตตฺถ กตเม ตโย กุสลเหตู, อโลโภ กุสลเหตุ อโทโส กุสลเหตุ อโมโห กุสลเหตุ, อิเม ตโย กุสลเหตู.}

\prose{ตตฺถ กตเม เทฺว อกุสลเหตู, โลโภ อกุสลเหตุ โมโห อกุสลเหตุ, อิเม เทฺว อกุสลเหตู.}

\prose{ตตฺถ กตเม ตโย อพฺยากตเหตู, กุสลานํ วา ธมฺมานํ วิปากโต กิริยาพฺยากเตสุ วา ธมฺเมสุ อโลโภ อโทโส อโมโห, อิเม ตโย อพฺยากตเหตู.\csromanspacer{} อิเม วุจฺจนฺติ รูปธาตุยา อฏฺฐ เหตู. (6)}

\prose{ตตฺถ กตเม รูปธาตุยา ตโย อาหารา, ผสฺสาหาโร มโนสญฺเจตนาหาโร วิญฺญาณาหาโร, อิเม วุจฺจนฺติ รูปธาตุยา ตโย อาหารา. (7)}

\prose{ตตฺถ กตเม รูปธาตุยา จตฺตาโร ผสฺสา, จกฺขุสมฺผสฺโส โสตสมฺผสฺโส มโนธาตุสมฺผสฺโส มโนวิญฺญาณธาตุสมฺผสฺโส, อิเม วุจฺจนฺติ รูปธาตุยา จตฺตาโร ผสฺสา. (8)}

\prose{ตตฺถ กตมา รูปธาตุยา จตสฺโส เวทนา ฯเปฯ จตสฺโส สญฺญา ฯเปฯ จตสฺโส เจตนา ฯเปฯ จตฺตาริ จิตฺตานิ, จกฺขุวิญฺญาณํ โสตวิญฺญาณํ มโนธาตุ มโนวิญฺญาณธาตุ, อิมานิ วุจฺจนฺติ รูปธาตุยา จตฺตาริ จิตฺตานิ. (12)}

\csromanheader{2}{18. ธมฺมหทยวิภงฺค}
\csromanmatikaanchor{mtk.195}
\csromantocmarkref{subparagraph}{3. อรูปธาตุ}{mtk.195}
\bookmark[dest={mtk.195},level=5]{3. อรูปธาตุ}
\csromanheader{6}{3. อรูปธาตุ}
\vspace{\tipitakaparskip}
\csromancenter{อรูปธาตุยา กติ ขนฺธา ฯเปฯ กติ จิตฺตานิ.}
\csromancenter{อรูปธาตุยา จตฺตาโร ขนฺธา, เทฺว อายตนานิ, เทฺว ธาตุโย, ตีณิ สจฺจานิ, เอกาทสินฺทฺริยานิ, อฏฺฐ เหตู, ตโย อาหารา, เอโก ผสฺโส, เอกา เวทนา, เอกา สญฺญา, เอกา เจตนา, เอกํ จิตฺตํ.}
\proseitem{996}{\csromanfolio{421}ตตฺถ กตเม อรูปธาตุยา จตฺตาโร ขนฺธา, เวทนากฺขนฺโธ สญฺญากฺขนฺโธ สงฺขารกฺขนฺโธ วิญฺญาณกฺขนฺโธ, อิเม วุจฺจนฺติ อรูปธาตุยา จตฺตาโร ขนฺธา. (1)}

\prose{ตตฺถ กตมานิ อรูปธาตุยา เทฺว อายตนานิ, มนายตนํ ธมฺมายตนํ, อิมานิ วุจฺจนฺติ อรูปธาตุยา เทฺว อายตนานิ. (2)}

\prose{ตตฺถ กตมา อรูปธาตุยา เทฺว ธาตุโย, มโนวิญฺญาณธาตุ ธมฺมธาตุ, อิมา วุจฺจนฺติ อรูปธาตุยา เทฺว ธาตุโย. (3)}

\prose{ตตฺถ กตมานิ อรูปธาตุยา ตีณิ สจฺจานิ, ทุกฺขสจฺจํ สมุทยสจฺจํ มคฺคสจฺจํ, อิมานิ วุจฺจนฺติ อรูปธาตุยา ตีณิ สจฺจานิ. (4)}

\prose{ตตฺถ กตมานิ อรูปธาตุยา เอกาทสินฺทฺริยานิ, มนินฺทฺริยํ ชีวิตินฺทฺริยํ โสมนสฺสินฺทฺริยํ อุเปกฺขินฺทฺริยํ สทฺธินฺทฺริยํ วีริยินฺทฺริยํ สตินฺทฺริยํ สมาธินฺทฺริยํ ปญฺญินฺทฺริยํ อญฺญินฺทฺริยํ อญฺญาตาวินฺทฺริยํ, อิมานิ วุจฺจนฺติ อรูปธาตุยา เอกาทสินฺทฺริยานิ. (5)}

\prose{ตตฺถ กตเม อรูปธาตุยา อฏฺฐ เหตู, ตโย กุสลเหตู, เทฺว อกุสลเหตู, ตโย อพฺยากตเหตู ฯเปฯ อิเม วุจฺจนฺติ อรูปธาตุยา อฏฺฐ เหตู. (6)}

\prose{ตตฺถ กตเม อรูปธาตุยา ตโย อาหารา, ผสฺสาหาโร มโนสญฺเจตนาหาโร วิญฺญาณาหาโร, อิเม วุจฺจนฺติ อรูปธาตุยา ตโย อาหารา. (7)}

\prose{ตตฺถ กตโม อรูปธาตุยา เอโก ผสฺโส, มโนวิญฺญาณธาตุสมฺผสฺโส, อยํ วุจฺจติ อรูปธาตุยา เอโก ผสฺโส. (8)}

\prose{\csromanfolio{422}ตตฺถ กตมา อรูปธาตุยา เอกา เวทนา ฯเปฯ เอกา สญฺญา ฯเปฯ เอกา เจตนา ฯเปฯ เอกํ จิตฺตํ, มโนวิญฺญาณธาตุ, อิทํ วุจฺจติ อรูปธาตุยา เอกํ จิตฺตํ. (12)}

\csromanmatikaanchor{mtk.196}
\csromantocmarkref{subparagraph}{4. อปริยาปนฺน}{mtk.196}
\bookmark[dest={mtk.196},level=5]{4. อปริยาปนฺน}
\csromanheader{6}{4. อปริยาปนฺน}
\proseitem{997}{อปริยาปนฺเน กติ ขนฺธา ฯเปฯ กติ จิตฺตานิ.}

\prose{อปริยาปนฺเน จตฺตาโร ขนฺธา, เทฺว อายตนานิ, เทฺว ธาตุโย, เทฺว สจฺจานิ, ทฺวาทสินฺทฺริยานิ, ฉ เหตู, ตโย อาหารา, เอโก ผสฺโส, เอกา เวทนา, เอกา สญฺญา, เอกา เจตนา, เอกํ จิตฺตํ.}

\proseitem{998}{ตตฺถ กตเม อปริยาปนฺเน จตฺตาโร ขนฺธา, เวทนากฺขนฺโธ สญฺญากฺขนฺโธ สงฺขารกฺขนฺโธ วิญฺญาณกฺขนฺโธ, อิเม วุจฺจนฺติ อปริยาปนฺเน จตฺตาโร ขนฺธา. (1)}

\prose{ตตฺถ กตมานิ อปริยาปนฺเน เทฺว อายตนานิ, มนายตนํ ธมฺมายตนํ, อิมานิ วุจฺจนฺติ อปริยาปนฺเน เทฺว อายตนานิ. (2)}

\prose{ตตฺถ กตมา อปริยาปนฺเน เทฺว ธาตุโย, มโนวิญฺญาณธาตุ ธมฺมธาตุ, อิมา วุจฺจนฺติ อปริยาปนฺเน เทฺว ธาตุโย. (3)}

\prose{ตตฺถ กตมานิ อปริยาปนฺเน เทฺว สจฺจานิ, มคฺคสจฺจํ นิโรธสจฺจํ, อิมานิ วุจฺจนฺติ อปริยาปนฺเน เทฺว สจฺจานิ. (4)}

\prose{ตตฺถ กตมานิ อปริยาปนฺเน ทฺวาทสินฺทฺริยานิ, มนินฺทฺริยํ ชีวิตินฺทฺริยํ โสมนสฺสินฺทฺริยํ อุเปกฺขินฺทฺริยํ สทฺธินฺทฺริยํ วีริยินฺทฺริยํ สตินฺทฺริยํ สมาธินฺทฺริยํ ปญฺญินฺทฺริยํ อนญฺญาตญฺญสฺสามีตินฺทฺริยํ อญฺญินฺทฺริยํ อญฺญาตาวินฺทฺริยํ, อิมานิ วุจฺจนฺติ อปริยาปนฺเน ทฺวาทสินฺทฺริยานิ. (5)}

\prose{ตตฺถ กตเม อปริยาปนฺเน ฉ เหตู, ตโย กุสลเหตู ตโย อพฺยากตเหตู.}

\prose{ตตฺถ กตเม ตโย กุสลเหตู, อโลโภ กุสลเหตุ อโทโส กุสลเหตุ อโมโห กุสลเหตุ, อิเม ตโย กุสลเหตู.}

\csromanheader{2}{18. ธมฺมหทยวิภงฺค}
\prose{\csromanfolio{423}ตตฺถ กตเม ตโย อพฺยากตเหตู, กุสลานํ ธมฺมานํ วิปากโต อโลโภ อโทโส อโมโห, อิเม ตโย อพฺยากตเหตู, อิเม วุจฺจนฺติ อปริยาปนฺเน ฉ เหตู. (6)}

\prose{ตตฺถ กตเม อปริยาปนฺเน ตโย อาหารา, ผสฺสาหาโร มโนสญฺเจตนาหาโร วิญฺญาณาหาโร, อิเม วุจฺจนฺติ อปริยาปนฺเน ตโย อาหารา. (7)}

\prose{ตตฺถ กตโม อปริยาปนฺเน เอโก ผสฺโส, มโนวิญฺญาณธาตุสมฺผสฺโส, อยํ วุจฺจติ อปริยาปนฺเน เอโก ผสฺโส. (8)}

\prose{ตตฺถ กตมา อปริยาปนฺเน เอกา เวทนา ฯเปฯ เอกา สญฺญา ฯเปฯ เอกา เจตนา ฯเปฯ เอกํ จิตฺตํ, มโนวิญฺญาณธาตุ, อิทํ วุจฺจติ อปริยาปนฺเน เอกํ จิตฺตํ. (12)}
\csromansectionrule

\csromanmatikaanchor{mtk.197}
\csromantocmarknopage{subparagraph}{3. ปริยาปนฺนาปริยาปนฺนวาร}{mtk.197}
\bookmark[dest={mtk.197},level=5]{3. ปริยาปนฺนาปริยาปนฺนวาร}
\csromanheader{6}{3. ปริยาปนฺนาปริยาปนฺนวาร}
\csromanmatikaanchor{mtk.198}
\csromantocmarkref{subparagraph}{1. กามธาตุ}{mtk.198}
\bookmark[dest={mtk.198},level=5]{1. กามธาตุ}
\csromanheader{6}{1. กามธาตุ}
\proseitem{999}{ปญฺจนฺนํ ขนฺธานํ กติ กามธาตุปริยาปนฺนา, กติ น กามธาตุปริยาปนฺนา ฯเปฯ สตฺตนฺนํ จิตฺตานํ กติ กามธาตุปริยาปนฺนา, กติ น กามธาตุปริยาปนฺนา.}

\proseitem{1000}{รูปกฺขนฺโธ กามธาตุปริยาปนฺโน.\csromanspacer{} จตฺตาโร ขนฺธา สิยา กามธาตุปริยาปนฺนา, สิยา น กามธาตุปริยาปนฺนา.}

\prose{ทสายตนา กามธาตุปริยาปนฺนา.\csromanspacer{} เทฺว อายตนา สิยา กามธาตุปริยาปนฺนา, สิยา น กามธาตุปริยาปนฺนา.}

\prose{โสฬส ธาตุโย กามธาตุปริยาปนฺนา.\csromanspacer{} เทฺว ธาตุโย สิยา กามธาตุปริยาปนฺนา, สิยา น กามธาตุปริยาปนฺนา.}

\prose{สมุทยสจฺจํ กามธาตุปริยาปนฺนํ.\csromanspacer{} เทฺว สจฺจา น กามธาตุปริยาปนฺนา.\csromanspacer{} ทุกฺขสจฺจํ สิยา กามธาตุปริยาปนฺนํ, สิยา น กามธาตุปริยาปนฺนํ.}

\prose{ทสินฺทฺริยา กามธาตุปริยาปนฺนา.\csromanspacer{} ตีณินฺทฺริยา น กามธาตุปริยาปนฺนา.\csromanspacer{} นวินฺทฺริยา สิยา กามธาตุปริยาปนฺนา, สิยา น กามธาตุปริยาปนฺนา.}

\prose{ตโย อกุสลเหตู กามธาตุปริยาปนฺนา.\csromanspacer{} ฉ เหตู สิยา กามธาตุปริยาปนฺนา, สิยา น กามธาตุปริยาปนฺนา.}

\prose{\csromanfolio{424}กพฬีกาโร กาหาโร กามธาตุปริยาปนฺโน.\csromanspacer{} ตโย อาหารา สิยา กามธาตุปริยาปนฺนา, สิยา น กามธาตุปริยาปนฺนา.}

\prose{ฉ ผสฺสา กามธาตุปริยาปนฺนา.\csromanspacer{} มโนวิญฺญาณธาตุสมฺผสฺโส สิยา กามธาตุ ปริยาปนฺโน, สิยา น กามธาตุปริยาปนฺโน.}

\prose{ฉ เวทนา ฯเปฯ.\csromanspacer{} ฉ สญฺญา.\csromanspacer{} ฉ เจตนา.\csromanspacer{} ฉ จิตฺตา กามธาตุปริยาปนฺนา.\csromanspacer{} มโนวิญฺญาณธาตุ สิยา กามธาตุปริยาปนฺนา, สิยา น กามธาตุปริยาปนฺนา.}

\csromanmatikaanchor{mtk.199}
\csromantocmarkref{subparagraph}{2. รูปธาตุ}{mtk.199}
\bookmark[dest={mtk.199},level=5]{2. รูปธาตุ}
\csromanheader{6}{2. รูปธาตุ}
\proseitem{1001}{ปญฺจนฺนํ ขนฺธานํ กติ รูปธาตุปริยาปนฺนา, กติ น รูปธาตุปริยาปนฺนา ฯเปฯ สตฺตนฺนํ จิตฺตานํ กติ รูปธาตุปริยาปนฺนา, กติ น รูปธาตุปริยาปนฺนา.}

\proseitem{1002}{รูปกฺขนฺโธ น รูปธาตุปริยาปนฺโน, จตฺตาโร ขนฺธา สิยา รูปธาตุปริยาปนฺนา, สิยา น รูปธาตุปริยาปนฺนา.}

\prose{ทสายตนา น รูปธาตุปริยาปนฺนา.\csromanspacer{} เทฺว อายตนา สิยา รูปธาตุปริยาปนฺนา, สิยา น รูปธาตุปริยาปนฺนา.}

\prose{โสฬส ธาตุโย น รูปธาตุปริยาปนฺนา.\csromanspacer{} เทฺว ธาตุโย สิยา รูปธาตุปริยาปนฺนา, สิยา น รูปธาตุปริยาปนฺนา.}

\prose{ตีณิ สจฺจานิ น รูปธาตุปริยาปนฺนา.\csromanspacer{} ทุกฺขสจฺจํ สิยา รูปธาตุปริยาปนฺนํ, สิยา น รูปธาตุปริยาปนฺนํ.}

\prose{เตรสินฺทฺริยา น รูปธาตุปริยาปนฺนา.\csromanspacer{} นวินฺทฺริยา สิยา รูปธาตุปริยาปนฺนา, สิยา น รูปธาตุปริยาปนฺนา.}

\prose{ตโย อกุสลเหตู น รูปธาตุปริยาปนฺนา.\csromanspacer{} ฉ เหตู สิยา รูปธาตุปริยาปนฺนา, สิยา น รูปธาตุปริยาปนฺนา.}

\prose{กพฬีกาโร อาหาโร น รูปธาตุปริยาปนฺโน.\csromanspacer{} ตโย อาหารา สิยา รูปธาตุปริยาปนฺนา, สิยา น รูปธาตุปริยาปนฺนา.}

\prose{ฉ ผสฺสา น รูปธาตุปริยาปนฺนา.\csromanspacer{} มโนวิญฺญาณธาตุสมฺผสฺโส สิยา รูปธาตุปริยาปนฺโน, สิยา น รูปธาตุปริยาปนฺโน.}

\csromanheader{2}{18. ธมฺมหทยวิภงฺค}
\prose{\csromanfolio{425}ฉ เวทนา ฯเปฯ.\csromanspacer{} ฉ สญฺญา.\csromanspacer{} ฉ เจตนา.\csromanspacer{} ฉ จิตฺตา น รูปธาตุปริยาปนฺนา.\csromanspacer{} มโนวิญฺญาณธาตุ สิยา รูปธาตุปริยาปนฺนา, สิยา น รูปธาตุปริยาปนฺนา.}

\csromanmatikaanchor{mtk.200}
\csromantocmarkref{subparagraph}{3. อรูปธาตุ}{mtk.200}
\bookmark[dest={mtk.200},level=5]{3. อรูปธาตุ}
\csromanheader{6}{3. อรูปธาตุ}
\proseitem{1003}{ปญฺจนฺนํ ขนฺธานํ กติ อรูปธาตุปริยาปนฺนา, กติ น รูปธาตุปริยาปนฺนา ฯเปฯ สตฺตนฺนํ จิตฺตานํ กติ อรูปธาตุปริยาปนฺนา, กติ น อรูปธาตุปริยาปนฺนา.}

\proseitem{1004}{รูปกฺขนฺโธ น อรูปธาตุปริยาปนฺนา.\csromanspacer{} จตฺตาโร ขนฺธา สิยา อรูปธาตุปริยาปนฺนา, สิยา น อรูปธาตุปริยาปนฺนา.}

\prose{ทสายตนา น อรูปธาตุปริยาปนฺนา.\csromanspacer{} เทฺว อายตนา สิยา อรูปธาตุปริยาปนฺนา, สิยา น อรูปธาตุปริยาปนฺนา.}

\prose{โสฬสธาตุโย น อรูปธาตุปริยาปนฺนา.\csromanspacer{} เทฺว ธาตุโย สิยา อรูปธาตุปริยาปนฺนา, สิยา น อรูปธาตุปริยาปนฺนา.}

\prose{ตีณิ สจฺจานิ น อรูปธาตุปริยาปนฺนานิ.}

\prose{ทุกฺขสจฺจํ สิยา อรูปธาตุปริยาปนฺนํ, สิยา น อรูปธาตุปริยาปนฺนํ.}

\prose{จุทฺทสินฺทฺริยา น อรูปธาตุปริยาปนฺนา.\csromanspacer{} อฏฺฐินฺทฺริยา สิยา อรูปธาตุปริยาปนฺนา, สิยา น อรูปธาตุปริยาปนฺนา.}

\prose{ตโย อกุสลเหตู น อรูปธาตุปริยาปนฺนา.\csromanspacer{} ฉ เหตู สิยา อรูปธาตุปริยาปนฺนา.\csromanspacer{} สิยา น อรูปธาตุปริยาปนฺนา.}

\prose{กพฬีกาโร อาหาโร น อรูปธาตุปริยาปนฺโน.\csromanspacer{} ตโย อาหารา สิยา อรูปธาตุปริยาปนฺนา, สิยา น อรูปธาตุปริยาปนฺนา.}

\prose{ฉ ผสฺสา น อรูปธาตุปริยาปนฺนา.\csromanspacer{} มโนวิญฺญาณธาตุสมฺผสฺโส สิยา อรูปธาตุปริยาปนฺโน, สิยา น อรูปธาตุปริยาปนฺโน.}

\prose{ฉ เวทนา ฯเปฯ.\csromanspacer{} ฉ สญฺญา.\csromanspacer{} ฉ เจตนา.\csromanspacer{} ฉ จิตฺตา น อรูปธาตุปริยาปนฺนา.\csromanspacer{} มโนวิญฺญาณธาตุ สิยา อรูปธาตุปริยาปนฺนา, สิยา น อรูปธาตุ ปริยาปนฺนา.}

\csromanmatikaanchor{mtk.201}
\csromantocmarkref{subparagraph}{4. ปริยาปนฺนาปริยาปนฺน}{mtk.201}
\bookmark[dest={mtk.201},level=5]{4. ปริยาปนฺนาปริยาปนฺน}
\csromanheader{6}{4. ปริยาปนฺนาปริยาปนฺน}
\proseitem{1005}{\csromanfolio{426}ปญฺจนฺนํ ขนฺธานํ กติ ปริยาปนฺนา, กติ อปริยาปนฺนา ฯเปฯ สตฺตนฺนํ จิตฺตานํ กติ ปริยาปนฺนา, กติ อปริยาปนฺนา.}

\proseitem{1006}{รูปกฺขนฺโธ ปริยาปนฺโน.\csromanspacer{} จตฺตาโร ขนฺธา สิยา ปริยาปนฺนา, สิยา อปริยาปนฺนา.}

\prose{ทสายตนา ปริยาปนฺนา.\csromanspacer{} เทฺว อายตนา สิยา ปริยาปนฺนา, สิยา อปริยาปนฺนา.}

\prose{โสฬส ธาตุโย ปริยาปนฺนา.\csromanspacer{} เทฺว ธาตุโย สิยา ปริยาปนฺนา, สิยา อปริยาปนฺนา.}

\prose{เทฺว สจฺจา ปริยาปนฺนา.\csromanspacer{} เทฺว สจฺจา อปริยาปนฺนา.}

\prose{ทสินฺทฺริยา ปริยาปนฺนา.\csromanspacer{} ตีณินฺทฺริยา อปริยาปนฺนา.\csromanspacer{} นวินฺทฺริยา สิยา ปริยาปนฺนา, สิยา อปริยาปนฺนา.}

\prose{ตโย อกุสลเหตู ปริยาปนฺนา.\csromanspacer{} ฉ เหตู สิยา ปริยาปนฺนา, สิยา อปริยาปนฺนา.}

\prose{กพฬีกาโร อาหาโร ปริยาปนฺโน.\csromanspacer{} ตโย อาหารา สิยา ปริยาปนฺนา, สิยา อปริยาปนฺนา.}

\prose{ฉ ผสฺสา ปริยาปนฺนา.\csromanspacer{} มโนวิญฺญาณธาตุสมฺผสฺโส สิยา ปริยาปนฺโน, สิยา อปริยาปนฺโน.}

\prose{ฉ เวทนา ฯเปฯ.\csromanspacer{} ฉ สญฺญา.\csromanspacer{} ฉ เจตนา.\csromanspacer{} ฉ จิตฺตา ปริยาปนฺนา.\csromanspacer{} มโนวิญฺญาณธาตุ สิยา ปริยาปนฺนา, สิยา อปริยาปนฺนา.}
\csromansectionrule

\csromanmatikaanchor{mtk.202}
\csromantocmarknopage{subparagraph}{4. ธมฺมทสฺสนวาร}{mtk.202}
\bookmark[dest={mtk.202},level=5]{4. ธมฺมทสฺสนวาร}
\csromanheader{6}{4. ธมฺมทสฺสนวาร}
\csromanmatikaanchor{mtk.203}
\csromantocmarkref{subparagraph}{1. กามธาตุ}{mtk.203}
\bookmark[dest={mtk.203},level=5]{1. กามธาตุ}
\csromanheader{6}{1. กามธาตุ}
\vspace{\tipitakaparskip}
\csromancenter{กามธาตุยา อุปปตฺติกฺขเณ กติ ขนฺธา ปาตุภวนฺติ ฯเปฯ กติ จิตฺตานิ ปาตุภวนฺติ.}
\csromancenter{กามธาตุยา อุปปตฺติกฺขเณ สพฺเพสํ ปญฺจกฺขนฺธา ปาตุภวนฺติ.\csromanspacer{} กสฺสจิ เอกาทสายตนานิ ปาตุภวนฺติ, กสฺสจิ ทสายตนานิ ปาตุภวนฺติ,}
\csromanheader{2}{18. ธมฺมหทยวิภงฺค}
\prose{\csromanfolio{427}กสฺสจิ อปรานิ ทสายตนานิ ปาตุภวนฺติ, กสฺสจิ นวายตนานิ ปาตุภวนฺติ, กสฺสจิ สตฺตายตนานิ ปาตุภวนฺติ, กสฺสจิ เอกาทส ธาตุโย ปาตุภวนฺติ, กสฺสจิ ทส ธาตุโย ปาตุภวนฺติ, กสฺสจิ อปรา ทส ธาตุโย ปาตุภวนฺติ, กสฺสจิ นว ธาตุโย ปาตุภวนฺติ, กสฺสจิ สตฺต ธาตุโย ปาตุภวนฺติ.\csromanspacer{} สพฺเพสํ เอกํ สจฺจํ ปาตุภวติ.\csromanspacer{} กสฺสจิ จุทฺทสินฺทฺริยานิ ปาตุภวนฺติ, กสฺสจิ เตรสินฺทฺริยานิ ปาตุภวนฺติ, กสฺสจิ อปรานิ เตรสินฺทฺริยานิ ปาตุภวนฺติ, กสฺสจิ ทฺวาทสินฺทฺริยานิ ปาตุภวนฺติ, กสฺสจิ ทสินฺทฺริยานิ ปาตุภวนฺติ, กสฺสจิ นวินฺทฺริยานิ ปาตุภวนฺติ, กสฺสจิ อปรานิ นวินฺทฺริยานิ ปาตุภวนฺติ, กสฺสจิ อฏฺฐินฺทฺริยานิ ปาตุภวนฺติ, กสฺสจิ อปรานิ อฏฺฐินฺทฺริยานิ ปาตุภวนฺติ, กสฺสจิ สตฺตินฺทฺริยานิ ปาตุภวนฺติ, กสฺสจิ ปญฺจินฺทฺริยานิ ปาตุภวนฺติ, กสฺสจิ จตฺตารินฺทฺริยานิ ปาตุภวนฺติ.\csromanspacer{} กสฺสจิ ตโย เหตู ปาตุภวนฺติ, กสฺสจิ เทฺว เหตู ปาตุภวนฺติ, กสฺสจิ\footnote{เกจิ (สฺยา)} อเหตุกา ปาตุภวนฺติ.\csromanspacer{} สพฺเพสํ จตฺตาโร อาหารา ปาตุภวนฺติ.\csromanspacer{} สพฺเพสํ เอโก ผสฺโส ปาตุภวติ.\csromanspacer{} สพฺเพสํ เอกา เวทนา.\csromanspacer{} เอกา สญฺญา.\csromanspacer{} เอกา เจตนา.\csromanspacer{} เอกํ จิตฺตํ ปาตุภวติ.}

\proseitem{1008}{กามธาตุยา อุปปตฺติกฺขเณ สพฺเพสํ กตเม ปญฺจกฺขนฺธา ปาตุภวนฺติ.\csromanspacer{} รูปกฺขนฺโธ ฯเปฯ จิญฺญาณกฺขนฺโธ.\csromanspacer{} กามธาตุยา อุปปตฺติกฺขเณ สพฺเพสํ อิเม ปญฺจกฺขนฺธา ปาตุภวนฺติ. (1)}

\proseitem{1009}{กามธาตุยา อุปปตฺติกฺขเณ กสฺส เอกาทสายตนานิ ปาตุภวนฺติ.\csromanspacer{} กามาวจรานํ เทวานํ, ปฐมกปฺปิกานํ มนุสฺสานํ, โอปปาติกานํ เปตานํ, โอปปาติกานํ อสุรานํ, โอปปาติกานํ ติรจฺฉานคตานํ, เนรยิกานํ, ปริปุณฺณายตนานํ อุปปตฺติกฺขเณ เอกาทสายตนานิ ปาตุภวนฺติ จกฺขายตนํ รูปายตนํ โสตายตนํ ฆานายตนํ คนฺธายตนํ ชิวฺหายตนํ รสายตนํ กายายตนํ โผฏฺฐพฺพายตนํ มนายตนํ ธมฺมายตนํ.\csromanspacer{} กามธาตุยา อุปปตฺติกฺขเณ เอเตสํ อิมานิ เอกาทสายตนานิ ปาตุภวนฺติ.}

\prose{กามธาตุยา อุปปตฺติกฺขเณ กสฺส ทสายตนานิ ปาตุภวนฺติ.\csromanspacer{} โอปปาติกานํ เปตานํ, โอปปาติกานํ อสุรานํ, โอปปาติกานํ ติรจฺฉานคตานํ, เนรยิกานํ, ชจฺจนฺธานํ อุปปตฺติกฺขเณ ทสายตนานิ \csromanfolio{428}ปาตุภวนฺติ รูปายตนํ โสตายตนํ ฆานายตนํ คนฺธายตนํ ชิวฺหายตนํ รสายตนํ กายายตนํ โผฏฺฐพฺพายตนํ มนายตนํ ธมฺมายตนํ.\csromanspacer{} กามธาตุยา อุปปตฺติกฺขเณ เอเตสํ อิมานิ ทสายตนานิ ปาตุภวนฺติ.}

\prose{กามธาตุยา อุปปตฺติกฺขเณ กสฺส อปรานิ ทสายตนานิ ปาตุภวนฺติ.\csromanspacer{} โอปปาติกานํ เปตานํ, โอปปาติกานํ อสุรานํ, โอปปาติกานํ ติรจฺฉานคตานํ, เนรยิกานํ, ชจฺจพธิรานํ อุปปตฺติกฺขเณ ทสายตนานิ ปาตุภวนฺติ จกฺขายตนํ รูปายตนํ ฆานายตนํ คนฺธายตนํ ชิวฺหายตนํ รสายตนํ กายายตนํ โผฏฺฐพฺพายตนํ มนายตนํ ธมฺมายตนํ.\csromanspacer{} กามธาตุยา อุปปตฺติกฺขเณ เอเตสํ อิมานิ ทสายตนานิ ปาตุภวนฺติ.}

\prose{กามธาตุยา อุปปตฺติกฺขเณ กสฺส นวายตนานิ ปาตุภวนฺติ.\csromanspacer{} โอปปาติกานํ เปตานํ, โอปปาติกานํ อสุรานํ, โอปปาติกานํ ติรจฺฉานคตานํ, เนรยิกานํ, ชจฺจนฺธพธิรานํ อุปปตฺติกฺขเณ นวายตนานิ ปาตุภวนฺติ รูปายตนํ ฆานายตนํ คนฺธายตนํ ชิวฺหายตนํ รสายตนํ กายายตนํ โผฏฺฐพฺพายตนํ มนายตนํ ธมฺมายตนํ.\csromanspacer{} กามธาตุยา อุปปตฺติกฺขเณ เอเตสํ อิมานิ นวายตนานิ ปาตุภวนฺติ.}

\prose{กามธาตุยา อุปปตฺติกฺขเณ กสฺส สตฺตายตนานิ ปาตุภวนฺติ.\csromanspacer{} คพฺภเสยฺยกานํ สตฺตานํ อุปปตฺติกฺขเณ สตฺตายตนานิ ปาตุภวนฺติ รูปายตนํ คนฺธายตนํ รสายตนํ กายายตนํ โผฏฺฐพฺพายตนํ มนายตนํ ธมฺมายตนํ.\csromanspacer{} กามธาตุยา อุปปตฺติกฺขเณ เอเตสํ อิมานิ สตฺตายตนานิ ปาตุภวนฺติ. (2)}

\proseitem{1010}{กามธาตุยา อุปปตฺติกฺขเณ กสฺส เอกาทส ธาตุโย ปาตุภวนฺติ.\csromanspacer{} กามาวจรานํ เทวานํ, ปฐมกปฺปิกานํ มนุสฺสานํ, โอปปาติกานํ เปตานํ, โอปปาติกานํ อสุรานํ, โอปปาติกานํ ติรจฺฉานคตานํ, เนรยิกานํ, ปริปุณฺณายตนานํ อุปปตฺติกฺขเณ เอกาทส ธาตุโย ปาตุภวนฺติ จกฺขุธาตุ รูปธาตุ โสตธาตุ ฆานธาตุ คนฺธธาตุ ชิวฺหาธาตุ รสธาตุ กายธาตุ โผฏฺฐพฺพธาตุ มโนวิญฺญาณธาตุ ธมฺมธาตุ.\csromanspacer{} กามธาตุยา อุปปตฺติกฺขเณ เอเตสํ อิมา เอกาทส ธาตุโย ปาตุภวนฺติ.}

\csromanheader{2}{18. ธมฺมหทยวิภงฺค}
\prose{\csromanfolio{429}กามธาตุยา อุปปตฺติกฺขเณ กสฺส ทส ธาตุโย ปาตุภวนฺติ.\csromanspacer{} โอปปาติกานํ เปตานํ, โอปปาติกานํ อสุรานํ, โอปปาติกานํ ติรจฺฉานคตานํ, เนรยิกานํ, ชจฺจนฺธานํ อุปปตฺติกฺขเณ ทส ธาตุโย ปาตุภวนฺติ รูปธาตุ โสตธาตุ ฆานธาตุ คนฺธธาตุ ชิวฺหาธาตุ รสธาตุ กายธาตุ โผฏฺฐพฺพธาตุ มโนวิญฺญาณธาตุ ธมฺมธาตุ.\csromanspacer{} กามธาตุยา อุปปตฺติกฺขเณ เอเตสํ อิมา ทส ธาตุโย ปาตุภวนฺติ.}

\prose{กามธาตุยา อุปปตฺติกฺขเณ กสฺส อปรา ทส ธาตุโย ปาตุภวนฺติ.\csromanspacer{} โอปปาติกานํ เปตานํ, โอปปาติกานํ อสุรานํ, โอปปาติกานํ ติรจฺฉานคตานํ, เนรยิกานํ, ชจฺจพธิรานํ อุปปตฺติกฺขเณ ทส ธาตุโย ปาตุภวนฺติ จกฺขุธาตุ รูปธาตุ ฆานธาตุ คนฺธธาตุ ชิวฺหาธาตุ รสธาตุ กายธาตุ โผฏฺฐพฺพธาตุ มโนวิญฺญาณธาตุ ธมฺมธาตุ.\csromanspacer{} กามธาตุยา อุปปตฺติกฺขเณ เอเตสํ อิมา ทส ธาตุโย ปาตุภวนฺติ.}

\prose{กามธาตุยา อุปปตฺติกฺขเณ กสฺส นว ธาตุโย ปาตุภวนฺติ.\csromanspacer{} โอปปาติกานํ เปตานํ, โอปปาติกานํ อสุรานํ, โอปปาติกานํ ติรจฺฉานคตานํ, เนรยิกานํ, ชจฺจนฺธพธิรานํ อุปปตฺติกฺขเณ นว ธาตุโย ปาตุภวนฺติ รูปธาตุ ฆานธาตุ คนฺธธาตุ ชิวฺหาธาตุ รสธาตุ กายธาตุ โผฏฺฐพฺพธาตุ มโนวิญฺญาณธาตุ ธมฺมธาตุ.\csromanspacer{} กามธาตุยา อุปปตฺติกฺขเณ เอเตสํ อิมา นว ธาตุโย ปาตุภวนฺติ.}

\prose{กามธาตุยา อุปปตฺติกฺขเณ กสฺส สตฺต ธาตุโย ปาตุภวนฺติ.\csromanspacer{} คพฺภเสยฺยกานํ สตฺตานํ อุปปตฺติกฺขเณ สตฺต ธาตุโย ปาตุภวนฺติ รูปธาตุ คนฺธธาตุ รสธาตุ กายธาตุ โผฏฺฐพฺพธาตุ มโนวิญฺญาณธาตุ ธมฺมธาตุ.\csromanspacer{} กามธาตุยา อุปปตฺติกฺขเณ เอเตสํ อิมา สตฺต ธาตุโย ปาตุภวนฺติ. (3)}

\proseitem{1011}{กามธาตุยา อุปปตฺติกฺขเณ สพฺเพสํ กตมํ เอกํ สจฺจํ ปาตุภวติ ทุกฺขสจฺจํ.\csromanspacer{} กามธาตุยา อุปปตฺติกฺขเณ สพฺเพสํ อิทํ เอกํ สจฺจํ ปาตุภวติ. (4)}

\proseitem{1012}{กามธาตุยา อุปปตฺติกฺขเณ กสฺส จุทฺทสินฺทฺริยานิ ปาตุภวนฺติ.\csromanspacer{} กามาวจรานํ เทวานํ, สเหตุกานํ ญาณสมฺปยุตฺตานํ อุปปตฺติกฺขเณ จุทฺทสินฺทฺริยานิ ปาตุภวนฺติ จกฺขุนฺทฺริยํ โสตินฺทฺริยํ ฆานินฺทฺริยํ ชิวฺหินฺทฺริยํ \csromanfolio{430}กายินฺทฺริยํ มนินฺทฺริยํ อิตฺถินฺทฺริยํ วา ปุริสินฺทฺริยํ วา ชีวิตินฺทฺริยํ โสมนสฺสินฺทฺริยํ วา อุเปกฺขินฺทฺริยํ วา สทฺธินฺทฺริยํ วีริยินฺทฺริยํ สตินฺทฺริยํ สมาธินฺทฺริยํ ปญฺญินฺทฺริยํ.\csromanspacer{} กามธาตุยา อุปปตฺติกฺขเณ เอเตสํ อิมานิ จุทฺทสินฺทฺริยานิ ปาตุภวนฺติ.}

\prose{กามธาตุยา อุปปตฺติกฺขเณ กสฺส เตรสินฺทฺริยานิ ปาตุภวนฺติ.\csromanspacer{} กามาวจรานํ เทวานํ, สเหตุกานํ ญาณวิปฺปยุตฺตานํ อุปปตฺติกฺขเณ เตรสินฺทฺริยานิ ปาตุภวนฺติ จกฺขุนฺทฺริยํ โสตินฺทฺริยํ ฆานินฺทฺริยํ ชิวฺหินฺทฺริยํ กายินฺทฺริยํ มนินฺทฺริยํ อิตฺถินฺทฺริยํ วา ปุริสินฺทฺริยํ วา ชีวิตินฺทฺริยํ โสมนสฺสินฺทฺริยํ วา อุเปกฺขินฺทฺริยํ วา สทฺธินฺทฺริยํ วีริยินฺทฺริยํ สตินฺทฺริยํ สมาธินฺทฺริยํ.\csromanspacer{} กามธาตุยา อุปปตฺติกฺขเณ เอเตสํ อิมานิ เตรสินฺทฺริยานิ ปาตุภวนฺติ.}

\prose{กามธาตุยา อุปปตฺติกฺเขเณ กสฺส อปรานิ เตรสินฺทฺริยานิ ปาตุภวนฺติ.\csromanspacer{} ปฐมกปฺปิกานํ มนุสฺสานํ, สเหตุกานํ ญาณสมฺปยุตฺตานํ อุปปตฺติกฺขเณ เตรสินฺทฺริยานิ ปาตุภวนฺติ จกฺขุนฺทฺริยํ โสตินฺทฺริยํ ฆานินฺทฺริยํ ชิวฺหินฺทฺริยํ กายินฺทฺริยํ มนินฺทฺริยํ ชีวิตินฺทฺริยํ โสมนสฺสินฺทฺริยํ วา อุเปกฺขินฺทฺริยํ วา สทฺธินฺทฺริยํ วีริยินฺทฺริยํ สตินฺทฺริยํ สมาธินฺทฺริยํ ปญฺญินฺทฺริยํ.\csromanspacer{} กามธาตุยา อุปปตฺติกฺขเณ เอเตสํ อิมานิ เตรสินฺทฺริยานิ ปาตุภวนฺติ.}

\prose{กามธาตุยา อุปปตฺติกฺขเณ กสฺส ทฺวาทสินฺทฺริยานิ ปาตุภวนฺติ.\csromanspacer{} ปฐมกปฺปิกานํ มนุสฺสานํ, สเหตุกานํ ญาณวิปฺปยุตฺตานํ อุปปตฺติกฺขเณ ทฺวาทสินฺทฺริยานิ ปาตุภวนฺติ จกฺขุนฺทฺริยํ โสตินฺทฺริยํ ฆานินฺทฺริยํ ชิวฺหินฺทฺริยํ กายินฺทฺริยํ มนินฺทฺริยํ ชีวิตินฺทฺริยํ โสมนสฺสินฺทฺริยํ วา อุเปกฺขินฺทฺริยํ วา สทฺธินฺทฺริยํ วีริยินฺทฺริยํ สตินฺทฺริยํ สมาธินฺทฺริยํ.\csromanspacer{} กามธาตุยา อุปปตฺติกฺขเณ เอเตสํ อิมานิ ทฺวาทสินฺทฺริยานิ ปาตุภวนฺติ.}

\prose{กมาธาตุยา อุปปตฺติกฺขเณ กสฺส ทสินฺทฺริยานิ ปาตุภวนฺติ.\csromanspacer{} คพฺภเสยฺยกานํ สตฺตานํ, สเหตุกานํ ญาณสมฺปยุตฺตานํ อุปปตฺติกฺขเณ ทสินฺทฺริยานิ ปาตุภวนฺติ กายินฺทฺริยํ มนินฺทฺริยํ อิตฺถินฺทฺริยํ วา ปุริสินฺทฺริยํ วา ชีวิตินฺทฺริยํ โสมนสฺสินฺทฺริยํ วา อุเปกฺขินฺทฺริยํ วา สทฺธินฺทฺริยํ วีริยินฺทฺริยํ สตินฺทฺริยํ สมาธินฺทฺริยํ ปญฺญินฺทฺริยํ.\csromanspacer{} กามธาตุยา อุปปตฺติกฺขเณ เอเตสํ อิมานิ ทสินฺทฺริยานิ ปาตุภวนฺติ.}

\prose{กามธาตุยา อุปปตฺติกฺขเณ กสฺส นวินฺทฺริยานิ ปาตุภวนฺติ.\csromanspacer{} คพฺภเสยฺยกานํ สตฺตานํ, สเหตุกานํ ญาณวิปฺปยุตฺตานํ อุปปตฺติกฺขเณ}

\csromanheader{2}{18. ธมฺมหทยวิภงฺค}
\prose{\csromanfolio{431}นวินฺทฺริยานิ ปาตุภวนฺติ กายินฺทฺริยํ มนินฺทฺริยํ อิตฺถินฺทฺริยํ วา ปุริสินฺทฺริยํ วา ชีวิตินฺทฺริยํ โสมนสฺสินฺทฺริยํ วา อุเปกฺขินฺทฺริยํ วา สทฺธินฺทฺริยํ วีริยินฺทฺริยํ สตินฺทฺริยํ สมาธินฺทฺริยํ.\csromanspacer{} กามธาตุยา อุปปตฺติกฺขเณ เอเตสํ อิมานิ นวินฺทฺริยานิ ปาตุภวนฺติ.}

\prose{กามธาตุยา อุปปตฺติกฺขเณ กสฺส อปรานิ นวินฺทฺริยานิ ปาตุภวนฺติ.\csromanspacer{} โอปปาติกานํ เปตานํ, โอปปาติกานํ อสุรานํ, โอปปาติกานํ ติรจฺฉานคตานํ, เนรยิกานํ, ปริปุณฺณายตนานํ อุปปตฺติกฺขเณ นวินฺทฺริยานิ ปาตุภวนฺติ จกฺขุนฺทฺริยํ โสตินฺทฺริยํ ฆานินฺทฺริยํ ชิวฺหินฺทฺริยํ กายินฺทฺริยํ มนินฺทฺริยํ อิตฺถินฺทฺริยํ วา ปุริสินฺทฺริยํ วา ชีวิตินฺทฺริยํ อุเปกฺขินฺทฺริยํ.\csromanspacer{} กามธาตุยา อุปปตฺติกฺขเณ เอเตสํ อิมานิ นวินฺทฺริยานิ ปาตุภวนฺติ.}

\prose{กามธาตุยา อุปปตฺติกฺขเณ กสฺส อฏฺฐินฺทฺริยานิ ปาตุภวนฺติ.\csromanspacer{} โอปปาติกานํ เปตานํ, โอปปาติกานํ อสุรานํ, โอปปาติกานํ ติรจฺฉานคตานํ, เนรยิกานํ, ชจฺจนฺธานํ อุปปตฺติกฺขเณ อฏฺฐินฺทฺริยานิ ปาตุภวนฺติ โสตินฺทฺริยํ ฆานินฺทฺริยํ ชิวฺหินฺทฺริยํ กายินฺทฺริยํ มนินฺทฺริยํ อิตฺถินฺทฺริยํ วา ปุริสินฺทฺริยํ วา ชีวิตินฺทฺริยํ อุเปกฺขินฺทฺริยํ.\csromanspacer{} กามธาตุยา อุปปตฺติกฺขเณ เอเตสํ อิมานิ อฏฺฐินฺทฺริยานิ ปาตุภวนฺติ.}

\prose{กามธาตุยา อุปปตฺติกฺขเณ กสฺส อปรานิ อฏฺฐินฺทฺริยานิ ปาตุภวนฺติ.\csromanspacer{} โอปปาติกานํ เปตานํ, โอปปาติกานํ อสุรานํ, โอปปาติกานํ ติรจฺฉานคตานํ, เนรยิกานํ, ชจฺจพธิรานํ อุปปตฺติกฺขเณ อฏฺฐินฺทฺริยานิ ปาตุภวนฺติ จกฺขุนฺทฺริยํ ฆานินฺทฺริยํ ชิวฺหินฺทฺริยํ กายินฺทฺริยํ มนินฺทฺริยํ อิตฺถินฺทฺริยํ วา ปุริสินฺทฺริยํ วา ชีวิตินฺทฺริยํ อุเปกฺขินฺทฺริยํ.\csromanspacer{} กามธาตุยา อุปปตฺติกฺขเณ เอเตสํ อิมานิ อฏฺฐินฺทฺริยานิ ปาตุภวนฺติ.}

\prose{กามธาตุยา อุปปตฺติกฺขเณ กสฺส สตฺตินฺทฺริยานิ ปาตุภวนฺติ.\csromanspacer{} โอปปาติกานํ เปตานํ, โอปปาติกานํ อสุรานํ, โอปปาติกานํ ติรจฺฉานคตานํ, เนรยิกานํ, ชจฺจนฺธพธิรานํ อุปปตฺติกฺขเณ สตฺตินฺทฺริยานิ ปาตุภวนฺติ ฆานินฺทฺริยํ ชิวฺหินฺทฺริยํ กายินฺทฺริยํ มนินฺทฺริยํ อิตฺถินฺทฺริยํ วา ปุริสินฺทฺริยํ วา ชีวิตินฺทฺริยํ อุเปกฺขินฺทฺริยํ.\csromanspacer{} กามธาตุยา อุปปตฺติกฺขเณ เอเตสํ อิมานิ สตฺตินฺทฺริยานิ ปาตุภวนฺติ.}

\prose{กามธาตุยา อุปปตฺติกฺขเณ กสฺส ปญฺจินฺทฺริยานิ ปาตุภวนฺติ.\csromanspacer{} คพฺภเสยฺยกานํ สตฺตานํ, อเหตุกานํ ฐเปตฺวา นปุํสกานํ อุปปตฺติกฺขเณ}

\prose{\csromanfolio{432}ปญฺจินฺทฺริยานิ ปาตุภวนฺติ กายินฺทฺริยํ มนินฺทฺริยํ อิตฺถินฺทฺริยํ วา ปุริสินฺทฺริยํ วา ชีวิตินฺทฺริยํ อุเปกฺขินฺทฺริยํ.\csromanspacer{} กามธาตุยา อุปปตฺติกฺขเณ เอเตสํ อิมานิ ปญฺจินฺทฺริยานิ ปาตุภวนฺติ.}

\prose{กามธาตุยา อุปปตฺติกฺขเณ กสฺส จตฺตารินฺทฺริยานิ ปาตุภวนฺติ.\csromanspacer{} คพฺภเสยฺยกานํ สตฺตานํ, อเหตุกานํ, นปุํสกานํ อุปปตฺติกฺขเณ จตฺตารินฺทฺริยานิ ปาตุภวนฺติ กายินฺทฺริยํ มนินฺทฺริยํ ชีวิตินฺทฺริยํ อุเปกฺขินฺทฺริยํ.\csromanspacer{} กามธาตุยา อุปปตฺติกฺขเณ เอเตสํ อิมานิ จตฺตารินฺทฺริยานิ ปาตุภวนฺติ. (5)}

\proseitem{1013}{กามธาตุยา อุปปตฺติกฺขเณ กสฺส ตโย เหตู ปาตุภวนฺติ.\csromanspacer{} กามาวจรานํ เทวานํ, ปฐมกปฺปิกานํ มนุสฺสานํ, คพฺภเสยฺยกานํ สตฺตานํ, สเหตุกานํ ญาณสมฺปยุตฺตานํ อุปปตฺติกฺขเณ ตโย เหตู ปาตุภวนฺติ อโลโภ วิปากเหตุ อโทโส วิปากเหตุ อโมโห วิปากเหตุ.\csromanspacer{} กามธาตุยา อุปปตฺติกฺขเณ เอเตสํ อิเม ตโย เหตู ปาตุภวนฺติ.}

\prose{กามธาตุยา อุปปตฺติกฺขเณ กสฺส เทฺว เหตู ปาตุภวนฺติ.\csromanspacer{} กามาวจรานํ เทวานํ, ปฐมกปฺปิกานํ มนุสฺสานํ, คพฺภเสยฺยกานํ สตฺตานํ, สเหตุกานํ ญาณวิปฺปยุตฺตานํ อุปปตฺติกฺขเณ เทฺว เหตู ปาตุภวนฺติ อโลโภ วิปากเหตุ อโทโส วิปากเหตุ.\csromanspacer{} กามธาตุยา อุปปตฺติกฺขเณ เอเตสํ อิเม เทฺว เหตู ปาตุภวนฺติ.\csromanspacer{} อวเสสานํ สตฺตานํ\footnote{อวเสสา สตฺตา (?)} อเหตุกา ปาตุภวนฺติ. (6)}

\proseitem{1014}{กามธาตุยา อุปปตฺติกฺขเณ สพฺเพสํ กตเม จตฺตาโร อาหารา ปาตุภวนฺติ.\csromanspacer{} กพฬีกาโร อาหาโร ผสฺสาหาโร มโนสญฺเจตนาหาโร วิญฺญาณาหาโร.\csromanspacer{} กามธาตุยา อุปปตฺติกฺขเณ สพฺเพสํ อิเม จตฺตาโร อาหารา ปาตุภวนฺติ. (7)}

\prose{กามธาตุยา อุปปตฺติกฺขเณ สพฺเพสํ กตโม เอโก ผสฺโส ปาตุภวติ.\csromanspacer{} มโนวิญฺญาณธาตุสมฺผสฺโส.\csromanspacer{} กามธาตุยา อุปปตฺติกฺขเณ สพฺเพสํ อยํ เอโก ผสฺโส ปาตุภวติ. (8)}

\csromanheader{2}{18. ธมฺมหทยวิภงฺค}
\prose{\csromanfolio{433}กามธาตุยา อุปปตฺติกฺขเณ สพฺเพสํ กตมา เอกา เวทนา.\csromanspacer{} เอกา สญฺญา.\csromanspacer{} เอกา เจตนา.\csromanspacer{} เอกํ จิตฺตํ ปาตุภวติ.\csromanspacer{} มโนวิญฺญาณธาตุ.\csromanspacer{} กามธาตุยา อุปปตฺติกฺขเณ สพฺเพสํ อิทํ เอกํ จิตฺตํ ปาตุภวติ. (12)}

\csromanmatikaanchor{mtk.204}
\csromantocmarkref{subparagraph}{2. รูปธาตุ}{mtk.204}
\bookmark[dest={mtk.204},level=5]{2. รูปธาตุ}
\csromanheader{6}{2. รูปธาตุ}
\vspace{\tipitakaparskip}
\csromancenter{รูปธาตุยา อุปปตฺติกฺขเณ กติ ขนฺธา ปาตุภวนฺติ ฯเปฯ กติ จิตฺตานิ ปาตุภวนฺติ.}
\csromancenter{รูปธาตุยา อุปปตฺติกฺขเณ ฐเปตฺวา อสญฺญสตฺตานํ เทวานํ ปญฺจกฺขนฺธา ปาตุภวนฺติ.\csromanspacer{} ปญฺจายตนานิ ปาตุภวนฺติ.\csromanspacer{} ปญฺจ ธาตุโย ปาตุภวนฺติ.\csromanspacer{} เอกํ สจฺจํ ปาตุภวติ.\csromanspacer{} ทสินฺทฺริยานิ ปาตุภวนฺติ.\csromanspacer{} ตโย เหตู ปาตุภวนฺติ.\csromanspacer{} ตโย อาหารา ปาตุภวนฺติ.\csromanspacer{} เอโก ผสฺโส ปาตุภวติ.\csromanspacer{} เอกา เวทนา.\csromanspacer{} เอกา สญฺญา.\csromanspacer{} เอกา เจตนา.\csromanspacer{} เอกํ จิตฺตํ ปาตุภวติ.}
\csromancenter{รูปธาตุยา อุปปตฺติกฺขเณ กตเม ปญฺจกฺขนฺธา ปาตุภวนฺติ.\csromanspacer{} รูปกฺขนฺโธ เวทนากฺขนฺโธ สญฺญากฺขนฺโธ สงฺขารกฺขนฺโธ วิญฺญาณกฺขนฺโธ.\csromanspacer{} รูปธาตุยา อุปปตฺติกฺขเณ อิเม ปญฺจกฺขนฺธา ปาตุภวนฺติ. (1)}
\csromancenter{รูปธาตุยา อุปปตฺติกฺขเณ กตมานิ ปญฺจายตนานิ ปาตุภวนฺติ.\csromanspacer{} จกฺขายตนํ รูปายตนํ โสตายตนํ มนายตนํ ธมฺมายตนํ.\csromanspacer{} รูปธาตุยา อุปปตฺติกฺขเณ อิมานิ ปญฺจายตนานิ ปาตุภวนฺติ. (2)}
\csromancenter{รูปธาตุยา อุปปตฺติกฺขเณ กตมา ปญฺจ ธาตุโย ปาตุภวนฺติ.\csromanspacer{} จกฺขุธาตุ รูปธาตุ โสตธาตุ มโนวิญฺญาณธาตุ ธมฺมธาตุ.\csromanspacer{} รูปธาตุยา อุปปตฺติกฺขเณ อิมา ปญฺจ ธาตุโย ปาตุภวนฺติ. (3)}
\csromancenter{รูปธาตุยา อุปปตฺติกฺขเณ กตมํ เอกํ สจฺจํ ปาตุภวติ.\csromanspacer{} ทุกฺขสจฺจํ.\csromanspacer{} รูปธาตุยา อุปปตฺติกฺขเณ อิทํ เอกํ สจฺจํ ปาตุภวติ. (4)}
\csromancenter{รูปธาตุยา อุปปตฺติกฺขเณ กตมานิ ทสินฺทฺริยานิ ปาตุภวนฺติ.\csromanspacer{} จกฺขุนฺทฺริยํ โสตินฺทฺริยํ มนินฺทฺริยํ ชีวิตินฺทฺริยํ โสมนสฺสินฺทฺริยํ วา อุเปกฺขินฺทฺริยํ วา สทฺธินฺทฺริยํ วีริยินฺทฺริยํ สตินฺทฺริยํ สมาธินฺทฺริยํ ปญฺญินฺทฺริยํ.\csromanspacer{} รูปธาตุยา อุปปตฺติกฺขเณ อิมานิ ทสินฺทฺริยานิ ปาตุภวนฺติ. (5)}
\csromancenter{รูปธาตุยา อุปปตฺติกฺขเณ กตเม ตโย เหตู ปาตุภวนฺติ.\csromanspacer{} อโลโภ วิปากเหตุ อโทโส วิปากเหตุ อโมโห}
\prosecont{\csromanfolio{434}วิปากเหตุ.\csromanspacer{} รูปธาตุยา อุปปตฺติกฺขเณ อิเม ตโย เหตู ปาตุภวนฺติ. (6)}

\prose{รูปธาตุยา อุปปตฺติกฺขเณ กตเม ตโย อาหารา ปาตุภวนฺติ.\csromanspacer{} ผสฺสาหาโร มโนสญฺเจตนาหาโร วิญฺญาณาหาโร.\csromanspacer{} รูปธาตุยา อุปปตฺติกฺขเณ อิเม ตโย อาหารา ปาตุภวนฺติ. (7)}

\prose{รูปธาตุยา อุปปตฺติกฺขเณ กตโม เอโก ผสฺโส ปาตุภวติ.\csromanspacer{} มโนวิญฺญาณธาตุสมฺผสฺโส.\csromanspacer{} รูปธาตุยา อุปปตฺติกฺขเณ อยํ เอโก ผสฺโส ปาตุภวติ. (8)}

\prose{รูปธาตุยา อุปปตฺติกฺขเณ กตมา เอกา เวทนา.\csromanspacer{} เอกา สญฺญา.\csromanspacer{} เอกา เจตนา.\csromanspacer{} เอกํ จิตฺตํ ปาตุภวติ.\csromanspacer{} มโนวิญฺญาณธาตุ.\csromanspacer{} รูปธาตุยา อุปปตฺติกฺขเณ อิทํ เอกํ จิตฺตํ ปาตุภวติ. (12)}

\csromanmatikaanchor{mtk.205}
\csromantocmarkref{subparagraph}{3. อสญฺญสตฺต}{mtk.205}
\bookmark[dest={mtk.205},level=5]{3. อสญฺญสตฺต}
\csromanheader{6}{3. อสญฺญสตฺต}
\proseitem{1017}{อสญฺญสตฺตานํ เทวานํ อุปปตฺติกฺขเณ กติ ขนฺธา ปาตุภวนฺติ ฯเปฯ กติ จิตฺตานิ ปาตุภวนฺติ.}

\prose{อสญฺญสตฺตานํ เทวานํ อุปปตฺติกฺขเณ เอโก ขนฺโธ ปาตุภวติ รูปกฺขนฺโธ.\csromanspacer{} เทฺว อายตนานิ ปาตุภวนฺติ รูปายตนํ ธมฺมายตนํ.\csromanspacer{} เทฺว ธาตุโย ปาตุภวนฺติ รูปธาตุ ธมฺมธาตุ.\csromanspacer{} เอกํ สจฺจํ ปาตุภวติ ทุกฺขสจฺจํ.\csromanspacer{} เอกินฺทฺริยํ ปาตุภวติ รูปชีวิตินฺทฺริยํ.\csromanspacer{} อสญฺญสตฺตา เทวา อเหตุกา อนาหารา อผสฺสกา อเวทนกา อสญฺญกา อเจตนกา อจิตฺตกา ปาตุภวนฺติ.}

\csromanmatikaanchor{mtk.206}
\csromantocmarkref{subparagraph}{4. อรูปธาตุ}{mtk.206}
\bookmark[dest={mtk.206},level=5]{4. อรูปธาตุ}
\csromanheader{6}{4. อรูปธาตุ}
\vspace{\tipitakaparskip}
\csromancenter{อรูปธาตุยา อุปปตฺติกฺขเณ กติ ขนฺธา ปาตุภวนฺติ ฯเปฯ กติ จิตฺตานิ ปาตุภวนฺติ.}
\csromancenter{อรูปธาตุยา อุปปตฺติกฺขเณ จตฺตาโร ขนฺธา ปาตุภวนฺติ, เทฺว อายตนานิ ปาตุภวนฺติ, เทฺว ธาตุโย ปาตุภวนฺติ, เอกํ สจฺจํ ปาตุภวติ, อฏฺฐินฺทฺริยานิ ปาตุภวนฺติ ตโย เหตู ปาตุภวนฺติ, ตโย อาหารา ปาตุภวนฺติ, เอโก ผสฺโส ปาตุภวติ, เอกา เวทนา.\csromanspacer{} เอกา สญฺญา.\csromanspacer{} เอกา เจตนา.\csromanspacer{} เอกํ จิตฺตํ ปาตุภวติ.}
\csromanheader{2}{18. ธมฺมหทยวิภงฺค}
\proseitem{1019}{\csromanfolio{435}อรูปธาตุยา อุปปตฺติกฺขเณ กตเม จตฺตาโร ขนฺธา ปาตุภวนฺติ.\csromanspacer{} เวทนากฺขนฺโธ สญฺญากฺขนฺโธ สงฺขารกฺขนฺโธ วิญฺญาณกฺขนฺโธ.\csromanspacer{} อรูปธาตุยา อุปปตฺติกฺขเณ อิเม จตฺตาโร ขนฺธา ปาตุภวนฺติ. (1)}

\prose{อรูปธาตุยา อุปปตฺติกฺขเณ กตมานิ เทฺว อายตนานิ ปาตุภวนฺติ.\csromanspacer{} มนายตนํ ธมฺมายตนํ.\csromanspacer{} อรูปธาตุยา อุปปตฺติกฺขเณ อิมานิ เทฺว อายตนานิ ปาตุภวนฺติ. (2)}

\prose{อรูปธาตุยา อุปปตฺติกฺขเณ กตมา เทฺว ธาตุโย ปาตุภวนฺติ.\csromanspacer{} มโนวิญฺญาณธาตุ ธมฺมธาตุ.\csromanspacer{} อรูปธาตุยา อุปปตฺติกฺขเณ อิมา เทฺว ธาตุโย ปาตุภวนฺติ. (3)}

\prose{อรูปธาตุยา อุปปตฺติกฺขเณ กตมํ เอกํ สจฺจํ ปาตุภวติ.\csromanspacer{} ทุกฺขสจฺจํ.\csromanspacer{} อรูปธาตุยา อุปปตฺติกฺขเณ อิทํ เอกํ สจฺจํ ปาตุภวติ. (4)}

\prose{อรูปธาตุยา อุปปตฺติกฺขเณ กตมานิ อฏฺฐินฺทฺริยานิ ปาตุภวนฺติ.\csromanspacer{} มนินฺทฺริยํ ชีวิตินฺทฺริยํ อุเปกฺขินฺทฺริยํ สทฺธินฺทฺริยํ วีริยินฺทฺริยํ สตินฺทฺริยํ สมาธินฺทฺริยํ ปญฺญินฺทฺริยํ.\csromanspacer{} อรูปธาตุยา อุปปตฺติกฺขเณ อิมานิ อฏฺฐินฺทฺริยานิ ปาตุภวนฺติ. (5)}

\prose{อรูปธาตุยา อุปปตฺติกฺขเณ กตเม ตโย เหตู ปาตุภวนฺติ.\csromanspacer{} อโลโภ วิปากเหตุ อโทโส วิปากเหตุ อโมโห วิปากเหตุ.\csromanspacer{} อรูปธาตุยา อุปปตฺติกฺขเณ อิเม ตโย เหตู ปาตุภวนฺติ. (6)}

\prose{อรูปธาตุยา อุปปตฺติกฺขเณ กตเม ตโย อาหารา ปาตุภวนฺติ.\csromanspacer{} ผสฺสาหาโร มโนสญฺเจตนาหาโร วิญฺญาณาหาโร.\csromanspacer{} อรูปธาตุยา อุปปตฺติกฺขเณ อิเม ตโย อาหารา ปาตุภวนฺติ. (7)}

\prose{อรูปธาตุยา อุปปตฺติกฺขเณ กตโม เอโก ผสฺโส ปาตุภวติ.\csromanspacer{} มโนวิญฺญาณธาตุสมฺผสฺโส.\csromanspacer{} อรูปธาตุยา อุปปตฺติกฺขเณ อยํ เอโก ผสฺโส ปาตุภวติ. (8)}

\prose{อรูปธาตุยา อุปปตฺติกฺขเณ กตมา เอกา เวทนา ฯเปฯ เอกา สญฺญา.\csromanspacer{} เอกา เจตนา.\csromanspacer{} เอกํ จิตฺตํ ปาตุภวติ.\csromanspacer{} มโนวิญฺญาณธาตุ.\csromanspacer{} อรูปธาตุยา อุปปตฺติกฺขเณ อิทํ เอกํ จิตฺตํ ปาตุภวติ. (12)}

\csromanmatikaanchor{mtk.207}
\csromantocmarkref{subparagraph}{5. ภูมนฺตรทสฺสนวาร}{mtk.207}
\bookmark[dest={mtk.207},level=5]{5. ภูมนฺตรทสฺสนวาร}
\csromanheader{6}{5. ภูมนฺตรทสฺสนวาร}
\proseitem{1020}{\csromanfolio{436}กามาวจรา ธมฺมา, น กามาวจรา ธมฺมา.\csromanspacer{} รูปาวจรา ธมฺมา, น รูปาวจรา ธมฺมา.\csromanspacer{} อรูปาวจรา ธมฺมา, น อรูปาวจรา ธมฺมา.\csromanspacer{} ปริยาปนฺนา ธมฺมา, อปริยาปนฺนา ธมฺมา.}

\prose{กตเม ธมฺมา กามาวจรา, เหฏฺฐโต อวีจินิรยํ ปริยนฺตํ กริตฺวา อุปริโต ปรนิมฺมิตวสวตฺตี เทเว อนฺโตกริตฺวา ยํ เอตสฺมิํ อนฺตเร เอตฺถาวจรา เอตฺถ ปริยาปนฺนา ขนฺธธาตุ-อายตนา รูปํ เวทนา สญฺญา สงฺขารา วิญฺญาณํ.\csromanspacer{} อิเม ธมฺมา กามาวจรา.}

\prose{กตเม ธมฺมา น กามาวจรา, รูปาวจรา อรูปาวจรา อปริยาปนฺนา, อิเม ธมฺมา น กามาวจรา.}

\prose{กตเม ธมฺมา รูปาวจรา, เหฏฺฐโต พฺรหฺมโลกํ ปริยนฺตํ กริตฺวา อุปริโต อกนิฏฺเฐ เทเว อนฺโตกริตฺวา ยํ เอตสฺมิํ อนฺตเร เอตฺถาวจรา เอตฺถ ปริยาปนฺนา สมาปนฺนสฺส วา อุปปนฺนสฺส วา ทิฏฺฐธมฺมสุขวิหาริสฺส วา จิตฺตเจตสิกา ธมฺมา, อิเม ธมฺมา รูปาวจรา.}

\prose{กตเม ธมฺมา น รูปาวจรา, กามาวจรา อรูปาวจรา อปริยาปนฺนา, อิเม ธมฺมา น รูปาวจรา.}

\prose{กตเม ธมฺมา อรูปาวจรา, เหฏฺฐโต อากาสานญฺจายตนูปเค เทเว ปริยนฺตํ กริตฺวา อุปริโต เนวสญฺญานาสญฺญายตนูปเค เทเว อนฺโตกริตฺวา ยํ เอตสฺมิํ อนฺตเร เอตฺถาวจรา เอตฺถ ปริยาปนฺนา สมาปนฺนสฺส วา อุปปนฺนสฺส วา ทิฏฺฐธมฺมสุขวิหาริสฺส วา จิตฺตเจตสิกา ธมฺมา, อิเม ธมฺมา อรูปาวจรา.}

\prose{กตเม ธมฺมา น อรูปาวจรา, กามาวจรา รูปาวจรา อปริยาปนฺนา, อิเม ธมฺมา น อรูปาวจรา.}

\prose{กตเม ธมฺมา ปริยาปนฺนา, สาสวา กุสลากุสลาพฺยากตา ธมฺมา กามาวจรา รูปาวจรา อรูปาวจรา รูปกฺขนฺโธ เวทนากฺขนฺโธ สญฺญากฺขนฺโธ สงฺขารกฺขนฺโธ วิญฺญาณกฺขนฺโธ, อิเม ธมฺมา ปริยาปนฺนา.}

\prose{กตเม ธมฺมา อปริยาปนฺนา, มคฺคา จ มคฺคผลานิ จ อสงฺขตา จ ธาตุ, อิเม ธมฺมา อปริยาปนฺนา.}

\csromanheader{2}{18. ธมฺมหทยวิภงฺค}
\csromanmatikaanchor{mtk.208}
\csromantocmarkref{subparagraph}{6. อุปฺปาทกกมฺม-อายุปฺปมาณวาร}{mtk.208}
\bookmark[dest={mtk.208},level=5]{6. อุปฺปาทกกมฺม-อายุปฺปมาณวาร}
\csromanheader{6}{6. อุปฺปาทกกมฺม-อายุปฺปมาณวาร}
\csromanmatikaanchor{mtk.209}
\csromantocmarkref{subparagraph}{1. อุปฺปาทกกมฺม}{mtk.209}
\bookmark[dest={mtk.209},level=5]{1. อุปฺปาทกกมฺม}
\csromanheader{6}{1. อุปฺปาทกกมฺม}
\proseitem{1021}{\csromanfolio{437}เทวาติ ตโย เทวา สมฺมุติเทวา\footnote{สมฺมติเทวา (สฺยา)} อุปปตฺติเทวา วิสุทฺธิเทวา.}

\prose{สมฺมุติเทวา นาม ราชาโน เทวิโย กุมารา.}

\prose{อุปปตฺติเทวา นาม จาตุมหาราชิเก\footnote{จาตุมฺมหาราชิเก (สี, สฺยา) 3. อปฺปํ วา ภิยฺโย วา (สฺยา, ก) ที 2. 3 ปิฏฺเฐปิ. 4. รตฺติทิโว (ก) อํ 1. 213 ปิฏฺเฐปิ.} เทเว อุปาทาย ตทุปริ เทวา.}

\prose{วิสุทฺธิเทวา นาม อรหนฺโต วุจฺจนฺติ.}

\prose{ทานํ ทตฺวา สีลํ สมาทิยิตฺวา อุโปสถกมฺมํ กตฺวา กตฺถ อุปปชฺชนฺติ.\csromanspacer{} ทานํ ทตฺวา สีลํ สมาทิยิตฺวา อุโปสถกมฺมํ กตฺวา อปฺเปกจฺเจ ขตฺติยมหาสาลานํ สหพฺยตํ อุปปชฺชนฺติ, อปฺเปกจฺเจ พฺราหฺมณมหาสาลานํ สหพฺยตํ อุปปชฺชนฺติ, อปฺเปกจฺเจ คหปติมหาสาลานํ สหพฺยตํ อุปปชฺชนฺติ, อปฺเปกจฺเจ จาตุมหาราชิกานํ เทวานํ สหพฺยตํ อุปปชฺชนฺติ, อปฺเปกจฺเจ ตาวติํสานํ เทวานํ สหพฺยตํ อุปปชฺชนฺติ, อปฺเปกจฺเจ ยามานํ เทวานํ สหพฺยตํ อุปปชฺชนฺติ, อปฺเปกจฺเจ ตุสิตานํ เทวานํ สหพฺยตํ อุปปชฺชนฺติ, อปฺเปกจฺเจ นิมฺมานรตีนํ เทวานํ สหพฺยตํ อุปปชฺชนฺติ, อปฺเปกจฺเจ ปรนิมฺมิตวสวตฺตีนํ เทวานํ สหพฺยตํ อุปปชฺชนฺติ.}

\csromanmatikaanchor{mtk.210}
\csromantocmarkref{subparagraph}{2. อายุปฺปมาณ}{mtk.210}
\bookmark[dest={mtk.210},level=5]{2. อายุปฺปมาณ}
\csromanheader{6}{2. อายุปฺปมาณ}
\proseitem{1022}{มนุสฺสานํ กิตฺตกํ อายุปฺปมาณํ, วสฺสสตํ อปฺปํ วา ภิยฺโย3. (1)}

\proseitem{1023}{จาตุมหาราชิกานํ เทวานํ กิตฺตกํ อายุปฺปมาณํ, ยานิ มานุสกานิ ปญฺญาส วสฺสานิ จาตุมหาราชิกานํ เทวานํ เอโส เอโก รตฺตินฺทิโว4, ตาย รตฺติยา ติํส รตฺติโย มาโส, เตน มาเสน ทฺวาทสมาสิโย สํวจฺฉโร, เตน สํวจฺฉเรน ทิพฺพานิ ปญฺจ วสฺสสตานิ จาตุมหาราชิกานํ เทวานํ อายุปฺปมาณํ, มนุสฺสานํ คณนาย กิตฺตกํ โหติ, นวุติ วสฺสสตสหสฺสานิ. (2)}

\prose{\csromanfolio{438}ตาวติํสานํ เทวานํ กิตฺตกํ อายุปฺปมาณํ, ยํ มานุสกํ วสฺสสตํ ตาวติํสานํ เทวานํ เอโส เอโก รตฺตินฺทิโว, ตาย รตฺติยา ติํสรตฺติโย มาโส, เตน มาเสน ทฺวาทสมาสิโย สํวจฺฉโร, เตน สํวจฺฉเรน ทิพฺพํ วสฺสสหสฺสํ ตาวติํสานํ เทวานํ อายุปฺปมาณํ, มนุสฺสานํ คณนาย กิตฺตกํ โหติ, ติสฺโส จ วสฺสโกฏิโย สฏฺฐิ จ วสฺสสตสหสฺสานิ. (3)}

\prose{ยามานํ เทวานํ กิตฺตกํ อายุปฺปมาณํ, ยานิ มานุสกานิ เทฺว วสฺสสตานิ ยามานํ เทวานํ เอโส เอโก รตฺตินฺทิโว, ตาย รตฺติยา ติํส รตฺติโย มาโส, เตน มาเสน ทฺวาทสมาสิโย สํวจฺฉโร, เตน สํวจฺฉเรน ทิพฺพานิ เทฺว วสฺสสหสฺสานิ ยามานํ เทวานํ อายุปฺปมาณํ, มนุสฺสานํ คณนาย กิตฺตกํ โหติ, จุทฺทสญฺจ วสฺสโกฏิโย จตฺตารีสญฺจ วสฺสสตสหสฺสานิ. (4)}

\prose{ตุสิตานํ เทวานํ กิตฺตกํ อายุปฺปมาณํ, ยานิ มานุสกานิ จตฺตาริ วสฺสสตานิ ตุสิตานํ เทวานํ เอโส เอโก รตฺตินฺทิโว, ตาย รตฺติยา ติํส รตฺติโย มาโส, เตน มาเสน ทฺวาทสมาสิโย สํวจฺฉโร, เตน สํวจฺฉเรน ทิพฺพานิ จตฺตาริ วสฺสสหสฺสานิ ตุสิตานํ เทวานํ อายุปฺปมาณํ, มนุสฺสานํ คณนาย กิตฺตกํ โหติ, สตฺตปญฺญาส วสฺสโกฏิโย สฏฺฐิ จ วสฺสสตสหสฺสานิ. (5)}

\prose{นิมฺมานรตีนํ เทวานํ กิตฺตกํ อายุปฺปมาณํ, ยานิ มานุสกานิ อฏฺฐ วสฺสสตานิ นิมฺมานรตีนํ เทวานํ เอโส เอโก รตฺตินฺทิโว, ตาย รตฺติยา ติํส รตฺติโย มาโส, เตน มาเสน ทฺวาทสมาสิโย สํวจฺฉโร, เตน สํวจฺฉเรน ทิพฺพานิ อฏฺฐ วสฺสสหสฺสานิ นิมฺมานรตีนํ เทวานํ อายุปฺปมาณํ, มนุสฺสานํ คณนาย กิตฺตกํ โหติ, เทฺว วสฺสโกฏิสตานิ ติํสญฺจ วสฺสโกฏิโย จตฺตารีสญฺจ วสฺสสตสหสฺสานิ. (6)}

\prose{ปรนิมฺมิตวสวตฺตีนํ เทวานํ กิตฺตกํ อายุปฺปมาณํ, ยานิ มานุสกานิ โสฬส วสฺสสตานิ ปรนิมฺมิตวสวตฺตีนํ เทวานํ เอโส เอโก รตฺตินฺทิโว, ตาย รตฺติยา ติํส รตฺติโย มาโส, เตน มาเสน ทฺวาทสมาสิโย สํวจฺฉโร, เตน สํวจฺฉเรน ทิพฺพานิ โสฬส วสฺสสหสฺสานิ ปรนิมฺมิตวสวตฺตีนํ เทวานํ อายุปฺปมาณํ, มนุสฺสานํ คณนาย}

\csromanheader{2}{18. ธมฺมหทยวิภงฺค}
\prose{\csromanfolio{439}กิตฺตกํ โหติ, นว จ วสฺสโกฏิสตานิ เอกวีสญฺจ วสฺสโกฏิโย สฏฺฐิ จ วสฺสสตสหสฺสานีติ. (7)}

\csromangathameasure{ฉ เอเต กามาวจรา, สพฺพกามสมิทฺธิโน. \\ สพฺเพสํ เอกสงฺขาโต, อายุ ภวติ กิตฺตโก. \\ ทฺวาทส โกฏิสตํ เตสํ, อฏฺฐวีสญฺจ โกฏิโย. \\ ปญฺญาส สตสหสฺสานิ, วสฺสคฺเคน ปกาสิตาติ.}
\csromangathagroup{\csromangathastanza{ฉ เอเต\footnote{ฉปิ เต (สฺยา)} กามาวจรา, สพฺพกามสมิทฺธิโน. \\ สพฺเพสํ เอกสงฺขาโต, อายุ ภวติ กิตฺตโก.}\csromangathastanza{ทฺวาทส โกฏิสตํ เตสํ, อฏฺฐวีสญฺจ โกฏิโย. \\ ปญฺญาส สตสหสฺสานิ, วสฺสคฺเคน ปกาสิตาติ.}}

\proseitem{1024}{ปฐมํ ฌานํ ปริตฺตํ ภาเวตฺวา กตฺถ อุปปชฺชนฺติ, ปฐมํ ฌานํ ปริตฺตํ ภาเวตฺวา พฺรหฺมปาริสชฺชานํ เทวานํ สหพฺยตํ อุปปชฺชนฺติ.\csromanspacer{} เตสํ กิตฺตกํ อายุปฺปมาณํ, กปฺปสฺส ตติโย ภาโค. (8)}

\prose{ปฐมํ ฌานํ มชฺฌิมํ ภาเวตฺวา กตฺถ อุปปชฺชนฺติ, ปฐมํ ฌานํ มชฺฌิมํ ภาเวตฺวา พฺรหฺมปุโรหิตานํ เทวานํ สหพฺยตํ อุปปชฺชนฺติ.\csromanspacer{} เตสํ กิตฺตกํ อายุปฺปมาณํ, อุปฑฺฒกปฺโป. (9)}

\prose{ปฐมํ ฌานํ ปณีตํ ภาเวตฺวา กตฺถ อุปปชฺชนฺติ, ปฐมํ ฌานํ ปณีตํ ภาเวตฺวา มหาพฺรหฺมานํ เทวานํ สหพฺยตํ อุปปชฺชนฺติ.\csromanspacer{} เตสํ กิตฺตกํ อายุปฺปมาณํ, กปฺโป\footnote{เอโก กปฺโป (สฺยา)}. (10)}

\proseitem{1025}{ทุติยํ ฌานํ ปริตฺตํ ภาเวตฺวา กตฺถ อุปปชฺชนฺติ, ทุติยํ ฌานํ ปริตฺตํ ภาเวตฺวา ปริตฺตาภานํ เทวานํ สหพฺยตํ อุปปชฺชนฺติ.\csromanspacer{} เตสํ กิตฺตกํ อายุปฺปมาณํ, เทฺว กปฺปา. (11)}

\prose{ทุติยํ ฌานํ มชฺฌิมํ ภาเวตฺวา กตฺถ อุปปชฺชนฺติ, ทุติยํ ฌานํ มชฺฌิมํ ภาเวตฺวา อปฺปมาณาภานํ เทวานํ สหพฺยตํ อุปปชฺชนฺติ.\csromanspacer{} เตสํ กิตฺตกํ อายุปฺปมาณํ, จตฺตาโร กปฺปา. (12)}

\prose{ทุติยํ ฌานํ ปณีตํ ภาเวตฺวา กตฺถ อุปปชฺชนฺติ, ทุติยํ ฌานํ ปณีตํ ภาเวตฺวา อาภสฺสรานํ เทวานํ สหพฺยตํ อุปปชฺชนฺติ.\csromanspacer{} เตสํ กิตฺตกํ อายุปฺปมาณํ, อฏฺฐ กปฺปา. (13)}

\proseitem{1026}{ตติยํ ฌานํ ปริตฺตํ ภาเวตฺวา กตฺถ อุปปชฺชนฺติ, ตติยํ ฌานํ ปริตฺตํ ภาเวตฺวา ปริตฺตสุภานํ เทวานํ สหพฺยตํ อุปปชฺชนฺติ.\csromanspacer{} เตสํ กิตฺตกํ อายุปฺปมาณํ, โสฬส กปฺปา. (14)}

\prose{\csromanfolio{440}ตติยํ ฌานํ มชฺฌิมํ ภาเวตฺวา กตฺถ อุปปชฺชนฺติ, ตติยํ ฌานํ มชฺฌิมํ ภาเวตฺวา อปฺปมาณสุภานํ เทวานํ สหพฺยตํ อุปปชฺชนฺติ.\csromanspacer{} เตสํ กิตฺตกํ อายุปฺปมาณํ, พาตฺติํส กปฺปา. (15)}

\prose{ตติยํ ฌานํ ปณีตํ ภาเวตฺวา กตฺถ อุปปชฺชนฺติ, ตติยํ ฌานํ ปณีตํ ภาเวตฺวา สุภกิณฺหานํ เทวานํ สหพฺยตํ อุปปชฺชนฺติ.\csromanspacer{} เตสํ กิตฺตกํ อายุปฺปมาณํ, จตุสฏฺฐิ กปฺปา. (16)}

\proseitem{1027}{จตุตฺถํ ฌานํ ภาเวตฺวา อารมฺมณนานตฺตตา, มนสิการนานตฺตตา, ฉนฺทนานตฺตตา, ปณิธินานตฺตตา, อธิโมกฺขนานตฺตตา, อภินีหารนานตฺตตา, ปญฺญานานตฺตตา อปฺเปกจฺเจ อสญฺญสตฺตานํ เทวานํ สหพฺยตํ อุปปชฺชนฺติ, อปฺเปกจฺเจ เวหปฺผลานํ เทวานํ สหพฺยตํ อุปปชฺชนฺติ, อปฺเปกจฺเจ อวิหานํ เทวานํ สหพฺยตํ อุปปชฺชนฺติ, อปฺเปกจฺเจ อตปฺปานํ เทวานํ สหพฺยตํ อุปปชฺชนฺติ, อปฺเปกจฺเจ สุทสฺสานํ เทวานํ สหพฺยตํ อุปปชฺชนฺติ, อปฺเปกจฺเจ สุทสฺสีนํ เทวานํ สหพฺยตํ อุปปชฺชนฺติ, อปฺเปกจฺเจ อกนิฏฺฐานํ เทวานํ สหพฺยตํ อุปปชฺชนฺติ.\csromanspacer{} อปฺเปกจฺเจ อากาสานญฺจายตนูปคานํ เทวานํ สหพฺยตํ อุปปชฺชนฺติ, อปฺเปกจฺเจ วิญฺญาณญฺจายตนูปคานํ เทวานํ สหพฺยตํ อุปปชฺชนฺติ, อปฺเปกจฺเจ อากิญฺจญฺญายตนูปคานํ เทวานํ สหพฺยตํ อุปปชฺชนฺติ, อปฺเปกจฺเจ เนวสญฺญานาสญฺญายตนูปคานํ เทวานํ สหพฺยตํ อุปปชฺชนฺติ.}

\prose{อสญฺญสตฺตานญฺจ เวหปฺผลานญฺจ เทวานํ กิตฺตกํ อายุปฺปมาณํ, ปญฺจกปฺปสตานิ. (18)}

\prose{อวิหานํ เทวานํ กิตฺตกํ อายุปฺปมาณํ, กปฺปสหสฺสํ. (19)}

\prose{อตปฺปานํ เทวานํ กิตฺตกํ อายุปฺปมาณํ, เทฺว กปฺปสหสฺสานิ. (20)}

\prose{สุทสฺสานํ เทวานํ กิตฺตกํ อายุปฺปมาณํ, จตฺตาริ กปฺปสหสฺสานิ. (21)}

\prose{สุทสฺสีนํ เทวานํ กิตฺตกํ อายุปฺปมาณํ, อฏฺฐ กปฺปสหสฺสานิ. (22)}

\prose{อกนิฏฺฐานํ เทวานํ กิตฺตกํ อายุปฺปมาณํ, โสฬส กปฺปสหสฺสานิ. (23)}

\proseitem{1028}{อากาสานญฺจายตนูปคานํ เทวานํ กิตฺตกํ อายุปฺปมาณํ, วีสติ กปฺปสหสฺสานิ. (24)}

\csromanheader{2}{18. ธมฺมหทยวิภงฺค}
\prose{\csromanfolio{441}วิญฺญาณญฺจายตนูปคานํ เทวานํ กิตฺตกํ อายุปฺปมาณํ, จตฺตารีส กปฺปสหสฺสานิ. (25)}

\prose{อากิญฺจญฺญายตนูปคานํ เทวานํ กิตฺตกํ อายุปฺปมาณํ, สฏฺฐิ กปฺปสหสฺสานิ. (26)}

\prose{เนวสญฺญานาสญฺญายตนูปคานํ เทวานํ กิตฺตกํ อายุปฺปมาณํ, จตุราสีติ กปฺปสหสฺสานีติ. (27)}

\proseitem{1029}{อุกฺขิตฺตา ปุญฺญเตเชน, กามรูปคติํคตา.}

\csromangathameasure{ภวคฺคตมฺปิ สมฺปตฺตา, ปุนาคจฺฉนฺติ ทุคฺคติํ. \\ ตาว ทีฆายุกา สตฺตา, จวนฺติ อายุสงฺขยา. \\ นตฺถิ โกจิ ภโว นิจฺโจ, อิติ วุตฺตํ มเหสินา. \\ ตสฺมา หิ ธีรา นิปกา, นิปุณา อตฺถจินฺตกา. \\ ชรามรณโมกฺขาย, ภาเวนฺติ มคฺคมุตฺตมํ. \\ ภาวยิตฺวา สุจิํ มคฺคํ, นิพฺพาโนคธคามินํ. \\ สพฺพาสเว ปริญฺญาย, ปรินิพฺพนฺติ อนาสวาติ.}
\csromangathagroup{\csromangathastanza{ภวคฺคตมฺปิ\footnote{ภวคฺคํ วาปี (สฺยา)} สมฺปตฺตา, ปุนาคจฺฉนฺติ\footnote{ปุน คจฺฉนฺติ (สฺยา)} ทุคฺคติํ. \\ ตาว ทีฆายุกา สตฺตา, จวนฺติ อายุสงฺขยา.}\csromangathastanza{นตฺถิ โกจิ ภโว นิจฺโจ, อิติ วุตฺตํ มเหสินา. \\ ตสฺมา หิ ธีรา นิปกา, นิปุณา อตฺถจินฺตกา.}\csromangathastanza{ชรามรณโมกฺขาย, ภาเวนฺติ มคฺคมุตฺตมํ. \\ ภาวยิตฺวา สุจิํ มคฺคํ, นิพฺพาโนคธคามินํ. \\ สพฺพาสเว ปริญฺญาย, ปรินิพฺพนฺติ อนาสวาติ.}}

\csromanmatikaanchor{mtk.211}
\csromantocmarkref{subparagraph}{7. อภิญฺเญยฺยาทิวาร}{mtk.211}
\bookmark[dest={mtk.211},level=5]{7. อภิญฺเญยฺยาทิวาร}
\csromanheader{6}{7. อภิญฺเญยฺยาทิวาร}
\proseitem{1030}{ปญฺจนฺนํ ขนฺธานํ กติ อภิญฺเญยฺยา, กติ ปริญฺเญยฺยา, กติ ปหาตพฺพา, กติ ภาเวตพฺพา, กติ สจฺฉิกาตพฺพา, กติ น ปหาตพฺพา น ภาเวตพฺพา น สจฺฉิกาตพฺพา ฯเปฯ สตฺตนฺนํ จิตฺตานํ กติ อภิญฺเญยฺยา, กติ ปริญฺเญยฺยา, กติ ปหาตพฺพา, กติ ภาเวตพฺพา, กติ สจฺฉิกาตพฺพา, กติ น ปหาตพฺพา น ภาเวตพฺพา น สจฺฉิกาตพฺพา.}

\proseitem{1031}{รูปกฺขนฺโธ อภิญฺเญยฺโย ปริญฺเญยฺโย, น ปหาตพฺโพ น ภาเวตพฺโพ น สจฺฉิกาตพฺโพ.\csromanspacer{} จตฺตาโร ขนฺธา อภิญฺเญยฺยา ปริญฺเญยฺยา, สิยา ปหาตพฺพา, สิยา ภาเวตพฺพา, สิยา สจฺฉิกาตพฺพา, สิยา น ปหาตพฺพา น ภาเวตพฺพา น สจฺฉิกาตพฺพา. (1)}

\prose{\csromanfolio{442}ทสายตนา อภิญฺเญยฺยา ปริญฺเญยฺยา, น ปหาตพฺพา น ภาเวตพฺพา น สจฺฉิกาตพฺพา.\csromanspacer{} เทฺว อายตนา อภิญฺเญยฺยา ปริญฺเญยฺยา, สิยา ปหาตพฺพา, สิยา ภาเวตพฺพา, สิยา สจฺฉิกาตพฺพา, สิยา น ปหาตพฺพา น ภาเวตพฺพา น สจฺฉิกาตพฺพา. (2)}

\prose{โสฬส ธาตุโย อภิญฺเญยฺยา ปริญฺเญยฺยา, น ปหาตพฺพา น ภาเวตพฺพา น สจฺฉิกาตพฺพา.\csromanspacer{} เทฺว ธาตุโย อภิญฺเญยฺยา ปริญฺเญยฺยา, สิยา ปหาตพฺพา, สิยา ภาเวตพฺพา, สิยา สจฺฉิกาตพฺพา, สิยา น ปหาตพฺพา น ภาเวตพฺพา น สจฺฉิกาตพฺพา. (3)}

\prose{สมุทยสจฺจํ อภิญฺเญยฺยํ ปริญฺเญยฺยํ ปหาตพฺพํ, น ภาเวตพฺพํ น สจฺฉิกาตพฺพํ.\csromanspacer{} มคฺคสจฺจํ อภิญฺเญยฺยํ ปริญฺเญยฺยํ, น ปหาตพฺพํ, ภาเวตพฺพํ, น สจฺฉิกาตพฺพํ.\csromanspacer{} นิโรธสจฺจํ อภิญฺเญยฺยํ ปริญฺเญยฺยํ, น ปหาตพฺพํ น ภาเวตพฺพํ, สจฺฉิกาตพฺพํ.\csromanspacer{} ทุกฺขสจฺจํ อภิญฺเญยฺยํ ปริญฺเญยฺยํ, สิยา ปหาตพฺพํ, น ภาเวตพฺพํ น สจฺฉิกาตพฺพํ, สิยา น ปหาตพฺพํ. (4)}

\prose{นวินฺทฺริยา อภิญฺเญยฺยา ปริญฺเญยฺยา, น ปหาตพฺพา น ภาเวตพฺพา น สจฺฉิกาตพฺพา.\csromanspacer{} โทมนสฺสินฺทฺริยํ อภิญฺเญยฺยํ ปริญฺเญยฺยํ ปหาตพฺพํ, น ภาเวตพฺพํ น สจฺฉิกาตพฺพํ.\csromanspacer{} อนญฺญาตญฺญสฺสามีตินฺทฺริยํ อภิญฺเญยฺยํ ปริญฺเญยฺยํ, น ปหาตพฺพํ, ภาเวตพฺพํ, น สจฺฉิกาตพฺพํ.\csromanspacer{} อญฺญินฺทฺริยํ อภิญฺเญยฺยํ ปริญฺเญยฺยํ, น ปหาตพฺพํ, สิยา ภาเวตพฺพํ, สิยา สจฺฉิกาตพฺพํ.\csromanspacer{} อญฺญาตาวินฺทฺริยํ อภิญฺเญยฺยํ ปริญฺเญยฺยํ, น ปหาตพฺพํ น ภาเวตพฺพํ, สจฺฉิกาตพฺพํ.\csromanspacer{} ตีณินฺทฺริยา อภิญฺเญยฺยา ปริญฺเญยฺยา, น ปหาตพฺพา, สิยา ภาเวตพฺพา, สิยา สจฺฉิกาตพฺพา, สิยา น ภาเวตพฺพา, สจฺฉิกาตพฺพา.\csromanspacer{} ฉ อินฺทฺริยา อภิญฺเญยฺยา ปริญฺเญยฺยา, สิยา ปหาตพฺพา, สิยา ภาเวตพฺพา, สิยา สจฺฉิกาตพฺพา, สิยา น ปหาตพฺพา น ภาเวตพฺพา น สจฺฉิกาตพฺพา. (5)}

\prose{ตโย อกุสลเหตู อภิญฺเญยฺยา ปริญฺเญยฺยา ปหาตพฺพา, น ภาเวตพฺพา น สจฺฉิกาตพฺพา.\csromanspacer{} ตโย กุสลเหตู อภิญฺเญยฺยา ปริญฺเญยฺยา, น ปหาตพฺพา, สิยา ภาเวตพฺพา, น สจฺฉิกาตพฺพา, สิยา น ภาเวตพฺพา.\csromanspacer{} ตโย อพฺยากตเหตู อภิญฺเญยฺยา ปริญฺเญยฺยา, น ปหาตพฺพา น ภาเวตพฺพา, สิยา สจฺฉิกาตพฺพา, สิยา น สจฺฉิกาตพฺพา. (6)}

\csromanheader{2}{18. ธมฺมหทยวิภงฺค}
\prose{\csromanfolio{443}กพฬีกาโร อาหาโร อภิญฺเญยฺโย ปริญฺเญยฺโย, น ปหาตพฺโพ น ภาเวตพฺโพ น สจฺฉิกาตพฺโพ.\csromanspacer{} ตโย อาหารา อภิญฺเญยฺยา ปริญฺเญยฺยา, สิยา ปหาตพฺพา, สิยา ภาเวตพฺพา, สิยา สจฺฉิกาตพฺพา, สิยา น ปหาตพฺพา น ภาเวตพฺพา น สจฺฉิกาตพฺพา. (7)}

\prose{ฉ ผสฺสา อภิญฺเญยฺยา ปริญฺเญยฺยา, น ปหาตพฺพา น ภาเวตพฺพา น สจฺฉิกาตพฺพา.\csromanspacer{} มโนวิญฺญาณธาตุสมฺผสฺโส อภิญฺเญยฺโย ปริญฺเญยฺโย, สิยา ปหาตพฺโพ, สิยา ภาเวตพฺโพ, สิยา สจฺฉิกาตพฺโพ, สิยา น ปหาตพฺโพ น ภาเวตพฺโพ น สจฺฉิกาตพฺโพ. (8)}

\prose{ฉ เวทนา ฯเปฯ.\csromanspacer{} ฉ สญฺญา.\csromanspacer{} ฉ เจตนา.\csromanspacer{} ฉ จิตฺตา อภิญฺเญยฺยา ปริญฺเญยฺยา, น ปหาตพฺพา น ภาเวตพฺพา น สจฺฉิกาตพฺพา.\csromanspacer{} มโนวิญฺญาณธาตุ อภิญฺเญยฺยา ปริญฺเญยฺยา, สิยา ปหาตพฺพา, สิยา ภาเวตพฺพา, สิยา สจฺฉิกาตพฺพา, สิยา น ปหาตพฺพา น ภาเวตพฺพา น สจฺฉิกาตพฺพา. (12)}
\csromansectionrule

\csromanmatikaanchor{mtk.212}
\csromantocmarkref{subparagraph}{8. สารมฺมณานารมฺมณวาร ...}{mtk.212}
\bookmark[dest={mtk.212},level=5]{8. สารมฺมณานารมฺมณวาร ...}
\csromanheader{6}{8. สารมฺมณานารมฺมณวาร}
\proseitem{1032}{ปญฺจนฺนํ ขนฺธานํ กติ สารมฺมณา, กติ อนารมฺมณา ฯเปฯ สตฺตนฺนํ จิตฺตานํ กติ สารมฺมณา, กติ อนารมฺมณา.}

\csromanheader{2}{1033. รูปกฺขนฺโธ อนารมฺมโณ. จตฺตาโร ขนฺธา สารมฺมณา. (1)}
\prose{ทสายตนา อนารมฺมณา.\csromanspacer{} มนายตนํ สารมฺมณํ.\csromanspacer{} ธมฺมายตนํ สิยา สารมฺมณํ, สิยา อนารมฺมณํ. (2)}

\prose{ทส ธาตุโย อนารมฺมณา.\csromanspacer{} สตฺต ธาตุโย สารมฺมณา.\csromanspacer{} ธมฺมธาตุ สิยา สารมฺมณา, สิยา อนารมฺมณา. (3)}

\prose{เทฺว สจฺจา สารมฺมณา.\csromanspacer{} นิโรธสจฺจํ อนารมฺมณํ.\csromanspacer{} ทุกฺขสจฺจํ สิยา สารมฺมณํ, สิยา อนารมฺมณํ. (4)}

\prose{สตฺตินฺทฺริยา อนารมฺมณา.\csromanspacer{} จุทฺทสินฺทฺริยา สารมฺมณา.\csromanspacer{} ชีวิตินฺทฺริยํ สิยา สารมฺมณํ, สิยา อนารมฺมณํ.\csromanspacer{} นว เหตู สารมฺมณา.\csromanspacer{} กพฬีกาโร อาหาโร อนารมฺมโณ.\csromanspacer{} ตโย อาหารา สารมฺมณา.\csromanspacer{} สตฺต ผสฺสา, สตฺต เวทนา, สตฺต สญฺญา, สตฺต เจตนา, สตฺต จิตฺตา สารมฺมณา. (12)}

\proseitem{1034}{\csromanfolio{444}ปญฺจนฺนํ ขนฺธานํ กติ สารมฺมณารมฺมณา.\csromanspacer{} กติ อนารมฺมณารมฺมณา ฯเปฯ สตฺตนฺนํ จิตฺตานํ กติ สารมฺมณารมฺมณา, กติ อนารมฺมณารมฺมณา.}

\proseitem{1035}{รูปกฺขนฺโธ อนารมฺมโณ.\csromanspacer{} จตฺตาโร ขนฺธา สิยา สารมฺมณารมฺมณา, สิยา อนารมฺมณารมฺมณา. (1)}

\prose{ทสายตนา อนารมฺมณา.\csromanspacer{} มนายตนํ สิยา สารมฺมณารมฺมณํ, สิยา อนารมฺมณารมฺมณํ.\csromanspacer{} ธมฺมายตนํ สิยา สารมฺมณารมฺมณํ, สิยา อนารมฺมณารมฺมณํ, สิยา อนารมฺมณํ. (2)}

\prose{ทส ธาตุโย อนารมฺมณา.\csromanspacer{} ฉ ธาตุโย อนารมฺมณารมฺมณา.\csromanspacer{} มโนวิญฺญาณธาตุ สิยา สารมฺมณารมฺมณา สิยา อนารมฺมณารมฺมณา.\csromanspacer{} ธมฺมธาตุ สิยา สารมฺมณารมฺมณา, สิยา อนารมฺมณารมฺมณา, สิยา อนารมฺมณา. (3)}

\prose{นิโรธสจฺจํ อนารมฺมณํ.\csromanspacer{} มคฺคสจฺจํ อนารมฺมณารมฺมณํ.\csromanspacer{} สมุทยสจฺจํ สิยา สารมฺมณารมฺมณํ, สิยา อนารมฺมณารมฺมณํ.\csromanspacer{} ทุกฺขสจฺจํ สิยา สารมฺมณารมฺมณํ, สิยา อนารมฺมณารมฺมณํ, สิยา อนารมฺมณํ. (4)}

\prose{สตฺตินฺทฺริยา อนารมฺมณา.\csromanspacer{} ปญฺจินฺทฺริยา อนารมฺมณารมฺมณา.\csromanspacer{} นวินฺทฺริยา สิยา สารมฺมณารมฺมณา, สิยา อนารมฺมณารมฺมณา.\csromanspacer{} ชีวิตินฺทฺริยํ สิยา สารมฺมณารมฺมณํ, สิยา อนารมฺมณารมฺมณํ, สิยา อนารมฺมณํ. (5)}

\prose{นว เหตู สิยา สารมฺมณารมฺมณา, สิยา อนารมฺมณารมฺมณา.\csromanspacer{} กพฬีกาโร อาหาโร อนารมฺมโณ.\csromanspacer{} ตโย อาหารา สิยา สารมฺมณารมฺมณา, สิยา อนารมฺมณารมฺมณา.\csromanspacer{} ฉ ผสฺสา อนารมฺมณารมฺมณา.\csromanspacer{} มโนวิญฺญาณธาตุสมฺผสฺโส สิยา สารมฺมณารมฺมโณ, สิยา อนารมฺมณารมฺมโณ.\csromanspacer{} ฉ เวทนา, ฉ สญฺญา, ฉ เจตนา, ฉ จิตฺตา อนารมฺมณารมฺมณา.\csromanspacer{} มโนวิญฺญาณธาตุ สิยา สารมฺมณารมฺมณา, สิยา อนารมฺมณารมฺมณา. (12)}
\csromansectionrule

\csromanmatikaanchor{mtk.213}
\csromantocmarkref{subparagraph}{9. ทิฏฺฐสุตาทิทสฺสนวาร}{mtk.213}
\bookmark[dest={mtk.213},level=5]{9. ทิฏฺฐสุตาทิทสฺสนวาร}
\csromanheader{6}{9. ทิฏฺฐสุตาทิทสฺสนวาร}
\proseitem{1036}{ปญฺจนฺนํ ขนฺธานํ กติ ทิฏฺฐา, กติ สุตา, กติ มุตา, กติ วิญฺญาตา.\csromanspacer{} กติ น ทิฏฺฐา น สุตา น มุตา น วิญฺญาตา ฯเปฯ สตฺตนฺนํ จิตฺตานํ}

\csromanheader{2}{18. ธมฺมหทยวิภงฺค}
\prose{\csromanfolio{445}กติ ทิฏฺฐา, กติ สุตา, กติ มุตา, กติ วิญฺญาตา.\csromanspacer{} กติ น ทิฏฺฐา น สุตา น มุตา น วิญฺญาตา.}

\proseitem{1037}{รูปกฺขนฺโธ สิยา ทิฏฺโฐ, สิยา สุโต, สิยา มุโต, สิยา วิญฺญาโต.\csromanspacer{} สิยา น ทิฏฺโฐ น สุโต น มุโต วิญฺญาโต.\csromanspacer{} จตฺตาโร ขนฺธา น ทิฏฺฐา น สุตา น มุตา วิญฺญาตา.}

\prose{รูปายตนํ ทิฏฺฐํ น สุตํ น มุตํ วิญฺญาตํ.\csromanspacer{} สทฺทายตนํ น ทิฏฺฐํ สุตํ น มุตํ วิญฺญาตํ.\csromanspacer{} คนฺธายตนํ, รสายตนํ, โผฏฺฐพฺพายตนํ น ทิฏฺฐํ น สุตํ มุตํ วิญฺญาตํ.\csromanspacer{} สตฺตายตนา น ทิฏฺฐา น สุตา น มุตา วิญฺญาตา.}

\prose{รูปธาตุ ทิฏฺฐา น สุตา น มุตา วิญฺญาตา.\csromanspacer{} สทฺทธาตุ น ทิฏฺฐา สุตา น มุตา วิญฺญาตา.\csromanspacer{} คนฺธธาตุ, รสธาตุ, โผฏฺฐพฺพธาตุ น ทิฏฺฐา น สุตา มุตา วิญฺญาตา.\csromanspacer{} เตรส ธาตุโย น ทิฏฺฐา น สุตา น มุตา วิญฺญาตา.}

\prose{ตีณิ สจฺจานิ น ทิฏฺฐา น สุตา น มุตา วิญฺญาตา.\csromanspacer{} ทุกฺขสจฺจํ สิยา ทิฏฺฐํ, สิยา สุตํ, สิยา มุตํ, สิยา วิญฺญาตํ, สิยา น ทิฏฺฐํ น สุตํ น มุตํ วิญฺญาตํ.}

\prose{พาวีสตินฺทฺริยา น ทิฏฺฐา น สุตา น มุตา วิญฺญาตา.\csromanspacer{} นว เหตู น ทิฏฺฐา น สุตา น มุตา วิญฺญาตา.\csromanspacer{} จตฺตาโร อาหารา น ทิฏฺฐา น สุตา น มุตา วิญฺญาตา.\csromanspacer{} สตฺต ผสฺสา น ทิฏฺฐา น สุตา น มุตา วิญฺญาตา.\csromanspacer{} สตฺต เวทนา, สตฺต สญฺญา, สตฺต เจตนา, สตฺต จิตฺตา น ทิฏฺฐา น สุตา น มุตา วิญฺญาตา.}
\csromansectionrule

\csromanmatikaanchor{mtk.214}
\csromantocmarkref{subparagraph}{10. ติกาทิทสฺสนวาร}{mtk.214}
\bookmark[dest={mtk.214},level=5]{10. ติกาทิทสฺสนวาร}
\csromanheader{6}{10. ติกาทิทสฺสนวาร}
\csromanmatikaanchor{mtk.215}
\csromantocmarkref{subparagraph}{1. กุสลตฺติก}{mtk.215}
\bookmark[dest={mtk.215},level=5]{1. กุสลตฺติก}
\csromanheader{6}{1. กุสลตฺติก}
\proseitem{1038}{ปญฺจนฺนํ ขนฺธานํ กติ กุสลา, กติ อกุสลา, กติ อพฺยากตา ฯเปฯ สตฺตนฺนํ จิตฺตานํ กติ กุสลา, กติ อกุสลา, กติ อพฺยากตา.}

\prose{รูปกฺขนฺโธ อพฺยากโต.\csromanspacer{} จตฺตาโร ขนฺธา สิยา กุสลา, สิยา อกุสลา, สิยา อพฺยากตา.\csromanspacer{} ทสายตนา อพฺยากตา.\csromanspacer{} ทฺวายตนา สิยา กุสลา, สิยา อกุสลา, สิยา อพฺยากตา.\csromanspacer{} โสฬส \csromanfolio{446}ธาตุโย อพฺยากตา.\csromanspacer{} เทฺว ธาตุโย สิยา กุสลา, สิยา อกุสลา, สิยา อพฺยากตา.\csromanspacer{} สมุทยสจฺจํ อกุสลํ.\csromanspacer{} มคฺคสจฺจํ กุสลํ.\csromanspacer{} นิโรธสจฺจํ อพฺยากตํ.\csromanspacer{} ทุกฺขสจฺจํ สิยา กุสลํ, สิยา อกุสลํ, สิยา อพฺยากตํ.}

\prose{ทสินฺทฺริยา อพฺยากตา.\csromanspacer{} โทมนสฺสินฺทฺริยํ อกุสลํ.\csromanspacer{} อนญฺญาตญฺญสฺสามีตินฺทฺริยํ กุสลํ.\csromanspacer{} จตฺตารินฺทฺริยา สิยา กุสลา, สิยา อพฺยากตา.\csromanspacer{} ฉ อินฺทฺริยา สิยา กุสลา, สิยา อกุสลา, สิยา อพฺยากตา.}

\prose{ตโย กุสลเหตู กุสลา.\csromanspacer{} ตโย อกุสลเหตู อกุสลา.\csromanspacer{} ตโย อพฺยากตเหตู อพฺยากตา.\csromanspacer{} กพฬีกาโร อาหาโร อพฺยากโต.\csromanspacer{} ตโย อาหารา สิยา กุสลา, สิยา อกุสลา, สิยา อพฺยากตา.\csromanspacer{} ฉ ผสฺสา อพฺยากตา.\csromanspacer{} มโนวิญฺญาณธาตุสมฺผสฺโส สิยา กุสโล, สิยา อกุสโล, สิยา อพฺยากโต.\csromanspacer{} ฉ เวทนา, ฉ สญฺญา, ฉ เจตนา, ฉ จิตฺตา อพฺยากตา.\csromanspacer{} มโนวิญฺญาณธาตุ สิยา กุสลา, สิยา อกุสลา, สิยา อพฺยากตา.}

\csromanmatikaanchor{mtk.216}
\csromantocmarkref{subparagraph}{2. เวทนาติก}{mtk.216}
\bookmark[dest={mtk.216},level=5]{2. เวทนาติก}
\csromanheader{6}{2. เวทนาติก}
\proseitem{1039}{ปญฺจนฺนํ ขนฺธานํ กติ สุขาย เวทนาย สมฺปยุตฺตา, กติ ทุกฺขาย เวทนาย สมฺปยุตฺตา, กติ อทุกฺขมสุขาย เวทนาย สมฺปยุตฺตา ฯเปฯ สตฺตนฺนํ จิตฺตานํ กติ สุขาย เวทนาย สมฺปยุตฺตา, กติ ทุกฺขาย เวทนาย สมฺปยุตฺตา, กติ อทุกฺขมสุขาย เวทนาย สมฺปยุตฺตา.}

\prose{เทฺว ขนฺธา น วตฺตพฺพา “สุขาย เวทนาย สมฺปยุตฺตา”ติปิ, “ทุกฺขาย เวทนาย สมฺปยุตฺตา”ติปิ, “อทุกฺขมสุขาย เวทนาย สมฺปยุตฺตา”ติปิ.\csromanspacer{} ตโย ขนฺธา สิยา สุขาย เวทนาย สมฺปยุตฺตา, สิยา ทุกฺขาย เวทนาย สมฺปยุตฺตา, สิยา อทุกฺขมสุขาย เวทนาย สมฺปยุตฺตา. (1)}

\prose{ทสายตนา น วตฺตพฺพา “สุขาย เวทนาย สมฺปยุตฺตา”ติปิ, “ทุกฺขาย เวทนาย สมฺปยุตฺตา”ติปิ, “อทุกฺขมสุขาย เวทนาย สมฺปยุตฺตา”ติปิ.\csromanspacer{} มนายตนํ สิยา สุขาย เวทนาย สมฺปยุตฺตํ, สิยา ทุกฺขาย เวทนาย สมฺปยุตฺตํ, สิยา อทุกฺขมสุขาย เวทนาย สมฺปยุตฺตํ.\csromanspacer{} ธมฺมายตนํ สิยา สุขาย เวทนาย สมฺปยุตฺตํ, สิยา ทุกฺขาย เวทนาย สมฺปยุตฺตํ, สิยา อทุกฺขมสุขาย เวทนาย สมฺปยุตฺตํ, สิยา น วตฺตพฺพํ “สุขาย}

\csromanheader{2}{18. ธมฺมหทยวิภงฺค}
\prose{\csromanfolio{447}เวทนาย สมฺปยุตฺตนฺ”ติปิ, “ทุกฺขาย เวทนาย สมฺปยุตฺตนฺ”ติปิ, “อทุกฺขมสุขาย เวทนาย สมฺปยุตฺตนฺ”ติปิ. (2)}

\prose{ทส ธาตุโย น วตฺตพฺพา “สุขาย เวทนาย สมฺปยุตฺตา”ติปิ, “ทุกฺขาย เวทนาย สมฺปยุตฺตา”ติปิ, อทุกฺขมสุขาย เวทนาย สมฺปยุตฺตา”ติปิ.\csromanspacer{} ปญฺจ ธาตุโย อทุกฺขมสุขาย เวทนาย สมฺปยุตฺตา.\csromanspacer{} กายวิญฺญาณธาตุ สิยา สุขาย เวทนาย สมฺปยุตฺตา, สิยา ทุกฺขาย เวทนาย สมฺปยุตฺตา.\csromanspacer{} มโนวิญฺญาณธาตุ สิยา สุขาย เวทนาย สมฺปยุตฺตา, สิยา ทุกฺขาย เวทนาย สมฺปยุตฺตา, สิยา อทุกฺขมสุขาย เวทนาย สมฺปยุตฺตา.\csromanspacer{} ธมฺมธาตุ สิยา สุขาย เวทนาย สมฺปยุตฺตา, สิยา ทุกฺขาย เวทนาย สมฺปยุตฺตา, สิยา อทุกฺขมสุขาย เวทนาย สมฺปยุตฺตา, สิยา น วตฺตพฺพา “สุขาย เวทนาย สมฺปยุตฺตา”ติปิ, “ทุกฺขาย เวทนาย สมฺปยุตฺตา”ติปิ, “อทุกฺขมสุขาย เวทนาย สมฺปยุตฺตา”ติปิ. (3)}

\prose{เทฺว สจฺจา สิยา สุขาย เวทนาย สมฺปยุตฺตา, สิยา อทุกฺขมสุขาย เวทนาย สมฺปยุตฺตา.\csromanspacer{} นิโรธสจฺจํ น วตฺตพฺพํ “สุขาย เวทนาย สมฺปยุตฺตนฺ”ติปิ, “ทุกฺขาย เวทนาย สมฺปยุตฺตนฺ”ติปิ, “อทุกฺขมสุขาย เวทนาย สมฺปยุตฺตนฺ”ติปิ.\csromanspacer{} ทุกฺขสจฺจํ สิยา สุขาย เวทนาย สมฺปยุตฺตํ, สิยา ทุกฺขาย เวทนาย สมฺปยุตฺตํ, สิยา อทุกฺขมสุขาย เวทนาย สมฺปยุตฺตํ, สิยา น วตฺตพฺพํ “สุขาย เวทนาย สมฺปยุตฺตนฺ”ติปิ, “ทุกฺขาย เวทนาย สมฺปยุตฺตนฺ”ติปิ, “อทุกฺขมสุขาย เวทนาย สมฺปยุตฺตนฺ”ติปิ. (4)}

\prose{ทฺวาทสินฺทฺริยา น วตฺตพฺพา “สุขาย เวทนาย สมฺปยุตฺตา”ติปิ, “ทุกฺขาย เวทนาย สมฺปยุตฺตา”ติปิ, “อทุกฺขมสุขาย เวทนาย สมฺปยุตฺตา”ติปิ.\csromanspacer{} ฉ อินฺทฺริยา สิยา สุขาย เวทนาย สมฺปยุตฺตา, สิยา อทุกฺขมสุขาย เวทนาย สมฺปยุตฺตา, ตีณินฺทฺริยา สิยา สุขาย เวทนาย สมฺปยุตฺตา, สิยา ทุกฺขาย เวทนาย สมฺปยุตฺตา, สิยา อทุกฺขมสุขาย เวทนาย สมฺปยุตฺตา.\csromanspacer{} ชีวิตินฺทฺริยํ สิยา สุขาย เวทนาย สมฺปยุตฺตํ, สิยา ทุกฺขาย เวทนาย สมฺปยุตฺตํ, สิยา อทุกฺขมสุขาย เวทนาย สมฺปยุตฺตํ, สิยา น วตฺตพฺพํ “สุขาย เวทนาย สมฺปยุตฺตนฺ”ติปิ, “ทุกฺขาย เวทนาย สมฺปยุตฺตนฺ”ติปิ, “อทุกฺขมสุขาย เวทนาย สมฺปยุตฺตนฺ”ติปิ. (5)}

\prose{\csromanfolio{448}โทโส อกุสลเหตุ ทุกฺขาย เวทนาย สมฺปยุตฺโต.\csromanspacer{} สตฺต เหตู สิยา สุขาย เวทนาย สมฺปยุตฺตา, สิยา อทุกฺขมสุขาย เวทนาย สมฺปยุตฺตา, โมโห อกุสลเหตุ สิยา สุขาย เวทนาย สมฺปยุตฺโต, สิยา ทุกฺขาย เวทนาย สมฺปยุตฺโต, สิยา อทุกฺขมสุขาย เวทนาย สมฺปยุตฺโต. (6)}

\prose{กพฬีกาโร อาหาโร น วตฺตพฺโพ “สุขาย เวทนาย สมฺปยุตฺโต”ติปิ, “ทุกฺขาย เวทนาย สมฺปยุตฺโต”ติปิ, “อทุกฺขมสุขาย เวทนาย สมฺปยุตฺโต”ติปิ.\csromanspacer{} ตโย อาหารา สิยา สุขาย เวทนาย สมฺปยุตฺตา, สิยา ทุกฺขาย เวทนาย สมฺปยุตฺตา, สิยา อทุกฺขมสุขาย เวทนาย สมฺปยุตฺตา. (7)}

\prose{ปญฺจ ผสฺสา อทุกฺขมสุขาย เวทนาย สมฺปยุตฺตา.\csromanspacer{} กายวิญฺญาณธาตุสมฺผสฺโส สิยา สุขาย เวทนาย สมฺปยุตฺโต, สิยา ทุกฺขาย เวทนาย สมฺปยุตฺโต.\csromanspacer{} มโนวิญฺญาณธาตุสมฺผสฺโส สิยา สุขาย เวทนาย สมฺปยุตฺโต, สิยา ทุกฺขาย เวทนาย สมฺปยุตฺโต, สิยา อทุกฺขมสุขาย เวทนาย สมฺปยุตฺโต. (8)}

\prose{สตฺต เวทนา น วตฺตพฺพา “สุขาย เวทนาย สมฺปยุตฺตา”ติปิ, “ทุกฺขาย เวทนาย สมฺปยุตฺตา”ติปิ, อทุกฺขมสุขาย เวทนาย สมฺปยุตฺตา”ติปิ.\csromanspacer{} ปญฺจ สญฺญา, ปญฺจ เจตนา, ปญฺจ จิตฺตา อทุกฺขมสุขาย เวทนาย สมฺปยุตฺตา.\csromanspacer{} กายวิญฺญาณธาตุ สิยา สุขาย เวทนาย สมฺปยุตฺตา, สิยา ทุกฺขาย เวทนาย สมฺปยุตฺตา.\csromanspacer{} มโนวิญฺญาณธาตุ สิยา สุขาย เวทนาย สมฺปยุตฺตา, สิยา ทุกฺขาย เวทนาย สมฺปยุตฺตา, สิยา อทุกฺขมสุขาย เวทนาย สมฺปยุตฺตา. (12)}

\csromanmatikaanchor{mtk.217}
\csromantocmarkref{subparagraph}{3. วิปากตฺติก}{mtk.217}
\bookmark[dest={mtk.217},level=5]{3. วิปากตฺติก}
\csromanheader{6}{3. วิปากตฺติก}
\proseitem{1040}{ปญฺจนฺนํ ขนฺธานํ กติ วิปากา, กติ วิปากธมฺมธมฺมา, กติ เนว วิปาก น วิปากธมฺมธมฺมา ฯเปฯ สตฺตนฺนํ จิตฺตานํ กติ วิปากา กติ วิปากธมฺมธมฺมา, กติ เนว วิปาก น วิปากธมฺมธมฺมา.}

\prose{รูปกฺขนฺโธ เนว วิปาก น วิปากธมฺมธมฺโม.\csromanspacer{} จตฺตาโร ขนฺธา สิยา วิปากา, สิยา วิปากธมฺมธมฺมา, สิยา เนว วิปาก น วิปากธมฺมธมฺมา. (1)}

\csromanheader{2}{18. ธมฺมหทยวิภงฺค}
\prose{\csromanfolio{449}ทสายตนา เนววิปากนวิปากธมฺมธมฺมา, ทฺวายตนา สิยา วิปากา, สิยา วิปากธมฺมธมฺมา, สิยา เนว วิปากนวิปากธมฺมธมฺมา. (2)}

\prose{ทส ธาตุโย เนววิปากนวิปากธมฺมธมฺมา.\csromanspacer{} ปญฺจ ธาตุโย วิปากา.\csromanspacer{} มโนธาตุ สิยา วิปากา, สิยา เนววิปากนวิปากธมฺมธมฺมา.\csromanspacer{} เทฺว ธาตุโย สิยา วิปากา, สิยา วิปากธมฺมธมฺมา, สิยา เนว วิปากนวิปากธมฺมธมฺมา. (3)}

\prose{เทฺว สจฺจา วิปากธมฺมธมฺมา.\csromanspacer{} นิโรธสจฺจํ เนววิปากนวิปากธมฺมธมฺมํ.\csromanspacer{} ทุกฺขสจฺจํ สิยา วิปากํ, สิยา วิปากธมฺมธมฺมํ, สิยา เนววิปากนวิปากธมฺมธมฺมํ. (4)}

\prose{สตฺตินฺทฺริยา เนววิปากนวิปากธมฺมธมฺมา.\csromanspacer{} ตีณินฺทฺริยา วิปากา.\csromanspacer{} ทฺวินฺทฺริยา วิปากธมฺมธมฺมา.\csromanspacer{} อญฺญินฺทฺริยํ สิยา วิปากํ, สิยา วิปากธมฺมธมฺมํ.\csromanspacer{} นวินฺทฺริยา สิยา วิปากา, สิยา วิปากธมฺมธมฺมา, สิยา เนววิปากนวิปากธมฺมธมฺมา. (5)}

\prose{ฉ เหตู วิปากธมฺมธมฺมา.\csromanspacer{} ตโย อพฺยากตเหตู สิยา วิปากา, สิยา เนววิปากนวิปากธมฺมธมฺมา. (6)}

\prose{กพฬีกาโร อาหาโร เนววิปากนวิปากธมฺมธมฺโม.\csromanspacer{} ตโย อาหารา สิยา วิปากา, สิยา วิปากธมฺมธมฺมา, สิยา เนววิปากนวิปากธมฺมธมฺมา.\csromanspacer{} ปญฺจ ผสฺสา วิปากา.\csromanspacer{} มโนธาตุสมฺผสฺโส สิยา วิปาโก, สิยา เนววิปากนวิปากธมฺมธมฺโม.\csromanspacer{} มโนวิญฺญาณธาตุสมฺผสฺโส สิยา วิปาโก, สิยา วิปากธมฺมธมฺโม, สิยา เนววิปากนวิปากธมฺมธมฺโม.\csromanspacer{} ปญฺจ เวทนา, ปญฺจ สญฺญา, ปญฺจ เจตนา, ปญฺจ จิตฺตา วิปากา.\csromanspacer{} มโนธาตุ สิยา วิปากา, สิยา เนววิปากนวิปากธมฺมธมฺมา.\csromanspacer{} มโนวิญฺญาณธาตุ สิยา วิปากา, สิยา วิปากธมฺมธมฺมา, สิยา เนววิปากนวิปากธมฺมธมฺมา. (12)}

\csromanmatikaanchor{mtk.218}
\csromantocmarkref{subparagraph}{4. อุปาทินฺนตฺติก}{mtk.218}
\bookmark[dest={mtk.218},level=5]{4. อุปาทินฺนตฺติก}
\csromanheader{6}{4. อุปาทินฺนตฺติก}
\proseitem{1041}{ปญฺจนฺนํ ขนฺธานํ กติ อุปาทินฺนุปาทานิยา, กติ อนุปาทินฺนุปาทานิยา, กติ อนุปาทินฺน-อนุปาทานิยา ฯเปฯ สตฺตนฺนํ จิตฺตานํ กติ อุปาทินฺนุปาทานิยา, กติ อนุปาทินฺนุปาทานิยา, กติ อนุปาทินฺน- อนุปาทานิยา.}

\prose{\csromanfolio{450}รูปกฺขนฺโธ สิยา อุปาทินฺนุปาทานิโย, สิยา อนุปาทินฺนุปาทานิโย.\csromanspacer{} จตฺตาโร ขนฺธา สิยา อุปาทินฺนุปาทานิยา, สิยา อนุปาทินฺนุปาทานิยา, สิยา อนุปาทินฺน-อนุปาทานิยา. (1)}

\prose{ปญฺจายตนา อุปาทินฺนุปาทานิยา.\csromanspacer{} สทฺทายตนํ อนุปาทินฺนุปาทานิยํ.\csromanspacer{} จตฺตาโร อายตนา สิยา อุปาทินฺนุปาทานิยา สิยา อนุปาทินฺนุปาทานิยา.\csromanspacer{} ทฺวายตนา สิยา อุปาทินฺนุปาทานิยา, สิยา อนุปาทินฺนุปาทานิยา, สิยา อนุปาทินฺน-อนุปาทานิยา. (2)}

\prose{ทส ธาตุโย อุปาทินฺนุปาทานิยา.\csromanspacer{} สทฺทธาตุ อนุปาทินฺนุปาทานิยา.\csromanspacer{} ปญฺจ ธาตุโย สิยา อุปาทินฺนุปาทานิยา, สิยา อนุปาทินฺนุปาทานิยา.\csromanspacer{} เทฺว ธาตุโย สิยา อุปาทินฺนุปาทานิยา, สิยา อนุปาทินฺนุปาทานิยา, สิยา อนุปาทินฺน- อนุปาทานิยา. (3)}

\prose{สมุทยสจฺจํ อนุปาทินฺนุปาทานิยํ.\csromanspacer{} เทฺว สจฺจา อนุปาทินฺน- อนุปาทานิยา.\csromanspacer{} ทุกฺขสจฺจํ สิยา อุปาทินฺนุปาทานิยํ, สิยา อนุปาทินฺนุปาทานิยํ. (4)}

\prose{นวินฺทฺริยา อุปาทินฺนุปาทานิยา.\csromanspacer{} โทมนสฺสินฺทฺริยํ อนุปาทินฺนุปาทานิยํ.\csromanspacer{} ตีณินฺทฺริยา อนุปาทินฺน-อนุปาทานิยา.\csromanspacer{} นวินฺทฺริยา สิยา อุปาทินฺนุปาทานิยา, สิยา อนุปาทินฺนุปาทานิยา, สิยา อนุปาทินฺน-อนุปาทานิยา.\csromanspacer{} ตโย อกุสลเหตู อนุปาทินฺนุปาทานิยา.\csromanspacer{} ตโย กุสลเหตู สิยา อนุปาทินฺนุปาทานิยา, สิยา อนุปาทินฺน-อนุปาทานิยา.\csromanspacer{} ตโย อพฺยากตเหตู สิยา อุปาทินฺนุปาทานิยา, สิยา อนุปาทินฺนุปาทานิยา, สิยา อนุปาทินฺน-อนุปาทานิยา. (5)}

\prose{กพฬีกาโร อาหาโร สิยา อุปาทินฺนุปาทานิโย, สิยา อนุปาทินฺนุปาทานิโย.\csromanspacer{} ตโย อาหารา สิยา อุปาทินฺนุปาทานิยา, สิยา อนุปาทินฺนุปาทานิยา, สิยา อนุปาทินฺน-อนุปาทานิยา. (6)}

\prose{ปญฺจ ผสฺสา อุปาทินฺนุปาทานิยา.\csromanspacer{} มโนธาตุสมฺผสฺโส สิยา อุปาทินฺนุปาทานิโย, สิยา อนุปาทินฺนุปาทานิโย.\csromanspacer{} มโนวิญฺญาณธาตุสมฺผสฺโส สิยา อุปาทินฺนุปาทานิโย, สิยา อนุปาทินฺนุปาทานิโย, สิยา อนุปาทินฺน- อนุปาทานิโย.\csromanspacer{} ปญฺจ เวทนา, ปญฺจ สญฺญา, ปญฺจ เจตนา, ปญฺจ จิตฺตา อุปาทินฺนุปาทานิยา.\csromanspacer{} มโนธาตุ สิยา อุปาทินฺนุปาทานิยา, สิยา อนุปาทินฺนุปาทานิยา.\csromanspacer{} มโนวิญฺญาณธาตุ สิยา อุปาทินฺนุปาทานิยา, สิยา อนุปาทินฺนุปาทานิยา, สิยา อนุปาทินฺน-อนุปาทานิยา. (12)}

\csromanheader{2}{18. ธมฺมหทยวิภงฺค}
\csromanmatikaanchor{mtk.219}
\csromantocmarkref{subparagraph}{5. วิตกฺกตฺติก}{mtk.219}
\bookmark[dest={mtk.219},level=5]{5. วิตกฺกตฺติก}
\csromanheader{6}{5. วิตกฺกตฺติก}
\proseitem{1042}{\csromanfolio{451}ปญฺจนฺนํ ขนฺธานํ กติ สวิตกฺกสวิจารา, กติ อวิตกฺกวิจารมตฺตา, กติ อวิตกฺก-อวิจารา ฯเปฯ สตฺตนฺนํ จิตฺตานํ กติ สวิตกฺกสวิจารา, กติ อวิตกฺกวิจารมตฺตา, กติ อวิตกฺก-อวิจารา.}

\prose{รูปกฺขนฺโธ อวิตกฺก-อวิจาโร.\csromanspacer{} ตโย ขนฺธา สิยา สวิตกฺกสวิจารา, สิยา อวิตกฺกวิจารมตฺตา, สิยา อวิตกฺก-อวิจารา.\csromanspacer{} สงฺขารกฺขนฺโธ อวิตกฺก-อวิจาโร, สิยา น วตฺตพฺโพ “สวิตกฺกสวิจาโร”ติปิ, “อวิตกฺกวิจารมตฺโต”ติปิ, “อวิตกฺก-อวิจาโร”ติปิ. (1)}

\prose{ทสายตนา อวิตกฺก-อวิจารา.\csromanspacer{} มนายตนํ สิยา สวิตกฺกสวิจารํ, สิยา อวิตกฺกวิจารมตฺตํ, สิยา อวิตกฺก-อวิจารํ.\csromanspacer{} ธมฺมายตนํ สิยา สวิตกฺกสวิจารํ, สิยา อวิตกฺกวิจารมตฺตํ, สิยา อวิตกฺก-อวิจารํ, สิยา น วตฺตพฺพํ “สวิตกฺกสวิจารนฺ”ติปิ, “อวิตกฺกวิจารมตฺตนฺ”ติปิ, “อวิตกฺก-อวิจารนฺ”ติปิ. (2)}

\prose{ปนฺนรส ธาตุโย อวิตกฺก-อวิจารา.\csromanspacer{} มโนธาตุ สวิตกฺกสวิจารา.\csromanspacer{} มโนวิญฺญาณธาตุ สิยา สวิตกฺกสวิจารา, สิยา อวิตกฺกวิจารมตฺตา, สิยา อวิตกฺก-อวิจารา.\csromanspacer{} ธมฺมธาตุ สิยา สวิตกฺกสวิจารา, สิยา อวิตกฺกวิจารมตฺตา, สิยา อวิตกฺก-อวิจารา, สิยา น วตฺตพฺพา “สวิตกฺกสวิจารา”ติปิ, “อวิตกฺกวิจารมตฺตา”ติปิ, “อวิตกฺข-อวิจารา”ติปิ. (3)}

\prose{สมุทยสจฺจํ สวิตกฺกสวิจารํ.\csromanspacer{} นิโรธสจฺจํ อวิตกฺก- อวิจารํ.\csromanspacer{} มคฺคสจฺจํ สิยา สวิตกฺกสวิจารํ, สิยา อวิตกฺกวิจารมตฺตํ, สิยา อวิตกฺก-อวิจารํ.\csromanspacer{} ทุกฺขสจฺจํ สิยา สวิตกฺกสวิจารํ, สิยา อวิตกฺกวิจารมตฺตํ, สิยา อวิตกฺก-อวิจารํ, สิยา น วตฺตพฺพํ “สวิตกฺกสวิจารนฺ”ติปิ, “อวิตกฺกวิจารมตฺตนฺ”ติ, “อวิตกฺก-อวิจารนฺ”ติปิ. (4)}

\prose{นวินฺทฺริยา อวิตกฺก-อวิจารา.\csromanspacer{} โทมนสฺสินฺทฺริยํ สวิตกฺกสวิจารํ.\csromanspacer{} อุเปกฺขินฺทฺริยํ สิยา สวิตกฺกสวิจารํ, สิยา อวิตกฺก-อวิจารํ.\csromanspacer{} เอกาทสินฺทฺริยา สิยา สวิตกฺกสวิจารา, สิยา อวิตกฺกวิจารมตฺตา, สิยา อวิตกฺก-อวิจารา. (5)}

\prose{\csromanfolio{452}ตโย อกุสลเหตู สวิตกฺกสวิจารา.\csromanspacer{} ฉ เหตู สิยา สวิตกฺกสวิจารา, สิยา อวิตกฺกวิจารมตฺตา, สิยา อวิตกฺก-อวิจารา.\csromanspacer{} กพฬีกาโร อาหาโร อวิตกฺก- อวิจาโร.\csromanspacer{} ตโย อาหารา สิยา สวิตกฺกสวิจารา, สิยา อวิตกฺกวิจารมตฺตา, สิยา อวิตกฺก-อวิจารา.\csromanspacer{} ปญฺจ ผสฺสา อวิตกฺก-อวิจารา.\csromanspacer{} มโนธาตุสมฺผสฺโส สวิตกฺกสวิจาโร.\csromanspacer{} มโนวิญฺญาณธาตุสมฺปสฺโส สิยา สวิตกฺกสวิจาโร, สิยา อวิตกฺกวิจารมตฺโต, สิยา อวิตกฺก-อวิจาโร.\csromanspacer{} ปญฺจ เวทนา, ปญฺจ สญฺญา, ปญฺจ เจตนา, ปญฺจ จิตฺตา อวิตกฺก-อวิจารา.\csromanspacer{} มโนธาตุ สจิตกฺกสวิจารา.\csromanspacer{} มโนวิญฺญาณธาตุ สิยา สวิตกฺกสวิจารา, สิยา อวิตกฺกวิจารมตฺตา, สิยา อวิตกฺก-อวิจารา. (12)}

\csromanmatikaanchor{mtk.220}
\csromantocmarkref{subparagraph}{1. รูปทุก}{mtk.220}
\bookmark[dest={mtk.220},level=5]{1. รูปทุก}
\csromanheader{6}{1. รูปทุก}
\proseitem{1043}{ปญฺจนฺนํ ขนฺธานํ กติ รูปา, กติ อรูปา ฯเปฯ สตฺตนฺนํ จิตฺตานํ กติ รูปา, กติ อรูปา.}

\prose{รูปกฺขนฺโธ รูปํ, จตฺตาโร ขนฺธา อรูปา.\csromanspacer{} ทสายตนา รูปา.\csromanspacer{} มนายตนํ อรูปํ, ธมฺมายตนํ สิยา รูปํ, สิยา อรูปํ.\csromanspacer{} ทส ธาตุโย รูปา.\csromanspacer{} สตฺต ธาตุโย อรูปา.\csromanspacer{} ธมฺมธาตุ สิยา รูปา, สิยา อรูปา.\csromanspacer{} ตีณิ สจฺจานิ อรูปา.\csromanspacer{} ทุกฺขสจฺจํ สิยา รูปํ, สิยา อรูปํ.\csromanspacer{} สตฺตินฺทฺริยา รูปา.\csromanspacer{} จุทฺทสินฺทฺริยา อรูปา.\csromanspacer{} ชีวิตินฺทฺริยํ สิยา รูปํ, สิยา อรูปํ.\csromanspacer{} นว เหตู อรูปา.\csromanspacer{} กพฬีกาโร อาหาโร รูปํ.\csromanspacer{} ตโย อาหารา อรูปา.\csromanspacer{} สตฺต ผสฺสา อรูปา.\csromanspacer{} สตฺต เวทนา, สตฺต สญฺญา, สตฺต เจตนา, สตฺต จิตฺตา อรูปา.}

\csromanmatikaanchor{mtk.221}
\csromantocmarkref{subparagraph}{2. โลกิยทุก}{mtk.221}
\bookmark[dest={mtk.221},level=5]{2. โลกิยทุก}
\csromanheader{6}{2. โลกิยทุก}
\proseitem{1044}{ปญฺจนฺนํ ขนฺธานํ กติ โลกิยา, กติ โลกุตฺตรา.\csromanspacer{} ทฺวาทสนฺนํ อายตนานํ กติ โลกิยา, กติ โลกุตฺตรา.\csromanspacer{} อฏฺฐารสนฺนํ ธาตูนํ กติ โลกิยา, กติ โลกุตฺตรา.\csromanspacer{} จตุนฺนํ สจฺจานํ กติ โลกิยา, กติ โลกุตฺตรา ฯเปฯ สตฺตนฺนํ จิตฺตานํ กติ โลกิยา, กติ โลกุตฺตรา.}

\prose{รูปกฺขนฺโธ โลกิโย.\csromanspacer{} จตฺตาโร ขนฺธา สิยา โลกิยา, สิยา โลกุตฺตรา.\csromanspacer{} ทสายตนา โลกิยา.\csromanspacer{} เทฺว อายตนา สิยา}

\csromanheader{2}{18. ธมฺมหทยวิภงฺค}
\prose{\csromanfolio{453}โลกิยา, สิยา โลกุตฺตรา.\csromanspacer{} โสฬส ธาตุโย โลกิยา.\csromanspacer{} เทฺว ธาตุโย สิยา โลกิยา, สิยา โลกุตฺตรา.\csromanspacer{} เทฺว สจฺจา โลกิยา, เทฺว สจฺจา โลกุตฺตรา.}

\prose{ทสินฺทฺริยา โลกิยา, ตีณินฺทฺริยา โลกุตฺตรา.\csromanspacer{} นวินฺทฺริยา สิยา โลกิยา, สิยา โลกุตฺตรา.\csromanspacer{} ตโย อกุสลเหตู โลกิยา.\csromanspacer{} ฉ เหตู สิยา โลกิยา, สิยา โลกุตฺตรา.\csromanspacer{} กพฬีกาโร อาหาโร โลกิโย.\csromanspacer{} ตโย อาหารา สิยา โลกิยา, สิยา โลกุตฺตรา.\csromanspacer{} ฉ ผสฺสา โลกิยา.\csromanspacer{} มโนวิญฺญาณธาตุสมฺผสฺโส สิยา โลกิโย, สิยา โลกุตฺตโร.\csromanspacer{} ฉ เวทนา โลกิยา.\csromanspacer{} มโนวิญฺญาณธาตุสมฺผสฺสชา เวทนา สิยา โลกิยา, สิยา โลกุตฺตรา.\csromanspacer{} ฉ สญฺญา โลกิยา.\csromanspacer{} มโนวิญฺญาณธาตุสมฺผสฺสชา สญฺญา สิยา โลกิยา, สิยา โลกุตฺตรา.\csromanspacer{} ฉ เจตนา โลกิยา.\csromanspacer{} มโนวิญฺญาณธาตุสมฺผสฺสชา เจตนา สิยา โลกิยา, สิยา โลกุตฺตรา.\csromanspacer{} ฉ จิตฺตา โลกิยา, มโนวิญฺญาณธาตุ สิยา โลกิยา, สิยา โลกุตฺตราติ.}

\csromangathameasure{อภิญฺญา เทฺว สารมฺมณา, ทิฏฺฐา กุสลเวทนา. \\ วิปากา จ อุปาทินฺนา, วิตกฺกํ รูปโลกิยาติ.}
\csromangathagroup{\csromangathastanza{อภิญฺญา เทฺว สารมฺมณา, ทิฏฺฐา กุสลเวทนา. \\ วิปากา จ อุปาทินฺนา, วิตกฺกํ รูปโลกิยาติ.}}

\vspace{\tipitakaparskip}
\csromancenter{ธมฺมหทยวิภงฺโค นิฏฺโฐโต.}
\nitthitam{\textbf{วิภงฺคปกรณํ นิฏฺฐิตํ.}}
\orphannote{ถิสฺ ผฺรเส โสฺหอุลฺทฺ โนตฺ เพ เอxอิสฺเตทฺ เหเร.}

